% ============================================================
%  The Essential Kierkegaard — danske originaltekster
%  Udvalg efter Howard V. & Edna H. Hong (Princeton, 2000).
%  Kilde: Søren Kierkegaards Skrifter (SKS), tekster.kb.dk/sks (CC0).
%  Formatering efter danish-texts-projektets konventioner.
% ============================================================
\documentclass[12pt,a4paper]{article}
\usepackage[utf8]{inputenc}
\usepackage[T1]{fontenc}
\usepackage{libertinus}
\usepackage{libertinust1math}
\usepackage[danish]{babel}
\usepackage[protrusion=true,expansion=true]{microtype}
\usepackage{setspace}
\usepackage[margin=1.2in]{geometry}
\usepackage{fancyhdr}
\usepackage{graphicx}
\usepackage{newunicodechar}
% --- occasional Greek + Kierkegaard's "ɔ:" under pdflatex ---------------
\newunicodechar{α}{\ensuremath{\alpha}}\newunicodechar{β}{\ensuremath{\beta}}
\newunicodechar{γ}{\ensuremath{\gamma}}\newunicodechar{δ}{\ensuremath{\delta}}
\newunicodechar{ε}{\ensuremath{\epsilon}}\newunicodechar{η}{\ensuremath{\eta}}
\newunicodechar{ζ}{\ensuremath{\zeta}}
\newunicodechar{ϑ}{\ensuremath{\vartheta}}\newunicodechar{θ}{\ensuremath{\theta}}
\newunicodechar{ι}{\ensuremath{\iota}}\newunicodechar{κ}{\ensuremath{\kappa}}
\newunicodechar{ϰ}{\ensuremath{\kappa}}\newunicodechar{λ}{\ensuremath{\lambda}}
\newunicodechar{μ}{\ensuremath{\mu}}\newunicodechar{ν}{\ensuremath{\nu}}
\newunicodechar{ξ}{\ensuremath{\xi}}\newunicodechar{ο}{o}
\newunicodechar{π}{\ensuremath{\pi}}\newunicodechar{ρ}{\ensuremath{\rho}}
\newunicodechar{σ}{\ensuremath{\sigma}}\newunicodechar{ς}{\ensuremath{\varsigma}}
\newunicodechar{τ}{\ensuremath{\tau}}\newunicodechar{υ}{\ensuremath{\upsilon}}
\newunicodechar{φ}{\ensuremath{\varphi}}\newunicodechar{χ}{\ensuremath{\chi}}
\newunicodechar{ψ}{\ensuremath{\psi}}\newunicodechar{ω}{\ensuremath{\omega}}
\newunicodechar{ɔ}{\reflectbox{c}}
\newunicodechar{Γ}{\ensuremath{\Gamma}}\newunicodechar{Δ}{\ensuremath{\Delta}}
\newunicodechar{Θ}{\ensuremath{\Theta}}\newunicodechar{Λ}{\ensuremath{\Lambda}}
\newunicodechar{Ξ}{\ensuremath{\Xi}}\newunicodechar{Π}{\ensuremath{\Pi}}
\newunicodechar{Σ}{\ensuremath{\Sigma}}\newunicodechar{Φ}{\ensuremath{\Phi}}
\newunicodechar{Ψ}{\ensuremath{\Psi}}\newunicodechar{Ω}{\ensuremath{\Omega}}
\newunicodechar{Η}{\ensuremath{\mathrm{H}}}\newunicodechar{Α}{\ensuremath{\mathrm{A}}}
\newunicodechar{Ε}{\ensuremath{\mathrm{E}}}\newunicodechar{Ι}{\ensuremath{\mathrm{I}}}
\newunicodechar{Κ}{\ensuremath{\mathrm{K}}}\newunicodechar{Ο}{\ensuremath{\mathrm{O}}}
\newunicodechar{Τ}{\ensuremath{\mathrm{T}}}\newunicodechar{Ν}{\ensuremath{\mathrm{N}}}
\usepackage{hyperref}
\hypersetup{
  pdftitle={The Essential Kierkegaard — danske originaltekster},
  pdfauthor={Søren Kierkegaard},
  hidelinks
}

\pagestyle{fancy}
\fancyhf{}
\rhead{\emph{The Essential Kierkegaard}}
\cfoot{\thepage}
\setlength{\headheight}{14pt}

\setstretch{1.3}

% --- project helpers ---------------------------------------------------
\newcommand{\el}{[\dots]}

% --- headings ----------------------------------------------------------
% Each anthology work opens a new page; sub-selections use \subsection*.
\newcommand{\workheading}[2]{%
  \clearpage
  \section*{#1\ {\normalfont\normalsize (#2)}}%
  \addcontentsline{toc}{section}{#1}\markright{#1}}
\newcommand{\worknote}[1]{{\small\itshape #1\par}\medskip\hrule\bigskip}
\newcommand{\sectionheading}[1]{\subsection*{\normalfont\itshape #1}}

\title{\textbf{The Essential Kierkegaard}\\[0.5ex]
  {\large danske originaltekster}}
\author{Udvalg efter Howard V.\ \& Edna H.\ Hong (red.),\\
  \emph{The Essential Kierkegaard} (Princeton, 2000)}
\date{Kildetekst: \emph{Søren Kierkegaards Skrifter} (SKS),\\
  Det Kgl.\ Bibliotek — tekster.kb.dk/sks (CC0)}

\begin{document}
\maketitle
\thispagestyle{fancy}

\section*{Om denne udgave}
For hvert uddrag i \emph{The Essential Kierkegaard} (red.\ Howard V.\ og Edna
H.\ Hong, Princeton University Press, 2000) gengiver denne udgave den
tilsvarende originaltekst fra \emph{Søren Kierkegaards Skrifter} (SKS), i
antologiens rækkefølge. Grænserne er afstemt efter Hongs engelske tekst og
deres margental (SV\textsuperscript{1}, Samlede Værker, 1.\ udgave). Oprindelig
ortografi er bevaret; tekstkritisk apparat, varianter og sidetal er udeladt.
Hvor antologien samler et uddrag fra flere adskilte steder i et værk, er
udeladelser markeret med \el.

\tableofcontents
\workheading{Af Søren Kierkegaards Journaler og Papirer}{1835–43}

\worknote{Antologiens ''Early Journal Entries''. Journaloptegnelse AA:12 — Kierkegaards brevudkast til P. W. Lund (juni–august 1835), hvorfra Hong-udvalget tager ''Copenhagen, June 1'', ''Nonnulla desunt'' og ''Gilleleie, August 1''; samt JJ:167 (1843).}

\sectionheading{Kjøbenhavn, den 1ste Juni 1835  (AA:12)}

De veed, med hvor megen Begeistring jeg i sin Tid hørte Dem tale; hvor enthusiastisk jeg var for Deres Beskrivelse af Deres Ophold i Brasilien, og i den Henseende igjen ikke saa meget med Hensyn til den Masse af enkelte Iagttagelser, hvormed De havde beriget Dem selv og Deres Videnskab, som med Hensyn til det Indtryk, Deres første Udtræden i den vidunderlige Natur giorde paa Dem; Deres paradisiske Lykke og Glæde. Noget Sligt maa altid virke sympathetisk paa ethvert Menneske, som dog ellers har nogen Følelse og Varme, om han endog mener at finde sin Tilfredshed, sin Virksomhed i en ganske anden Sphære, men især paa den Yngre, der endnu kun drømmer om sin Bestemmelse. Vor første Ungdom staaer som Blomsten i Morgendæmring med en deilig Dugdraabe i sit Bæger, hvori alle Omgivelser harmonisk-melancholsk speile sig. Men snart hæver Solen sig paa Horizonten, og Dugdraaben udtørres; med den forsvinder Livets Drømmerier, og nu gjælder det, om Mennesket er istand til, for atter at tage et Billede fra Blomsterne, at udvikle – ved sin egen Kraft ligesom en Nereum, – en Draabe, der kan staae som Frugten af hans Liv. Dertil hører først og fremmest, at man kommer til at staae i den Jordbund, hvor man egenlig hører hjemme; men den er ikke altid saa let at finde. Der gives i den Henseende lykkelige Naturer, som have en saa afgiort Tilbøielighed til en vis Retning, at de troligen arbeide frem ad den engang saaledes anviste Vei, uden at nogensinde den Tanke faaer nogen Magt over dem, at det maaskee egenlig var en ganske anden Vei, de skulde betræde. Der gives Andre, som saa aldeles lade sig styre af deres Omgivelser, at det aldrig gaaer ret op for dem, hvor de egenlig stræbe hen. Ligesom den foregaaende Classe har sit indvortes categoriske Imperativ, saaledes anerkjender denne sidste et udvortes categorisk Imperativ. Men hvor Faa ere ikke de første, og til de sidstes Classe ønsker jeg ikke at høre. Større er deres Tal, som i Livet faae at prøve, hvad egenlig denne Hegelske Dialectik vil sige. Det er nu forresten ganske i sin Orden, at Vinen gjærer, inden den bliver klar; men ubehagelig er desuagtet ofte denne Tilstand i dens enkelte Momenter, ihvorvel den naturligviis betragtet i sin Helhed har sin Behagelighed, forsaavidt den indenfor den almindelige Tvivlen dog har sine relative Resultater. Især har den stor Betydning for den, der igjennem den er kommet paa det Rene med sin Bestemmelse, ikke blot for den paafølgende Ro i Modsætning til den foregaaende Storm: men fordi man da har Livet i en ganske anden Betydning end før. Det er dette Faustiske Element, som for en Deel mere eller mindre gjør sig gjældende i enhver intellectuel Udvikling, hvorfor det altid har forekommet mig, at man burde indrømme Ideen til * Faust* en Verdens-Betydning. Saaledes som vore Forfædre havde en Gudinde for Længselen, saaledes, mener jeg, staaer * Faust* som den personificerede Tvivl. Og mere skal han ikke være, og det er vistnok en Synden mod Ideen, naar Goethe har ladet Faust omvende sig, ligesaavel som naar Merimée har ladet Don Juan omvende sig. Man indvende mig ikke, at Faust dog i det Øieblik, han henvendte sig til Djævelen, gjorde et positivt Skridt, thi heri synes mig netop at ligge Noget af det Dybeste i Sagnet om Faust. Han netop hengav sig til ham for at faae Oplysning og havde den altsaa ikke iforveien, og netop fordi han hengav sig til Djævelen, forøgedes Tvivlen (ligesom en Syg, der falder i en Qvaksalvers Hænder, gjerne bliver endnu værre); thi vel lod Mephistopheles ham see igjennem sine Briller ind i Menneskene og ind i Jordens forborgne Gjemmer, men Faust maatte bestandig nære Tvivl til ham, fordi han aldrig kunde oplyse ham om det allerdybeste i intellectuel Henseende. Til Gud kunde han ifølge sin Idee aldrig komme til at henvende sig, forsaavidt han i det Øieblik, han gjorde det, maatte sige sig selv, at her var i Sandhed Oplysning at finde; men i samme Øieblik havde han jo fornægtet sin Charakteer som Tvivler.

Men en saadan Tvivlen kan ogsaa vise sig i andre Sphærer. Om Mennesket ogsaa er kommet paa det Rene med sig selv om enkelte af saadanne Hovedpuncter, saa gives der i Livet andre ogsaa betydelige Spørgsmaal. Ethvert Menneske ønsker naturligviis at virke efter sine Evner i Verden, men deraf følger da igjen, at han ønsker at uddanne sine Evner i en bestemt Retning, i den nemlig, som er meest passende for hans Individualitet. Men hvilken er denne? Her staaer jeg ved et stort Spørgsmaalstegn. Her staaer jeg som Hercules men ikke paa Skilleveien – nei, her viser sig en langt større Mangfoldighed af Veie, og desto vanskeligere er det altsaa at gribe den rette. Det er maaskee netop en Ulykke ved min Existents, at jeg interesserer mig for altfor Meget og ikke afgjort for Noget; mine Interesser staae ikke alle subordinerede een, men alle staae coordinerede.

Jeg vil forsøge at vise, hvorledes Sagerne staae for mig.

1. Naturvidenskaberne. Seer jeg for det første hen til hele denne Retning (idet jeg derunder indbefatter alle dem, der søge at forstaaeliggjøre og tyde Naturens Runer, ligefra den, der beregner Stjernernes Fart og saa at sige standser dem deri for nærmere at undersøge dem, til den, der beskriver et enkelt Dyrs Physiologie; fra den, der fra Bjergenes Høider overskuer Jordfladen, til den, der stiger ned i Afgrundens Dybder; fra den, der forfølger Menneskelegemets Dannelse igjennem de utallige Nuancer, til den, der iagttager Indvoldsormene), da har jeg naturligviis paa denne Vei som paa enhver anden (dog fornemlig paa denne) seet Exempler paa Folk, der have skabt sig et Navn i Literaturen ved en uhyre Samlerflid. En stor Mængde af Enkeltheder kjende de, og de have opdaget mange nye; men heller ikke mere. De have blot leveret et Substrat for Andres Tænkning og Bearbeidelse. Og disse Mennesker staae nu der tilfredse med deres Enkeltheder, og dog forekomme de mig at staae ligesom den rige Bonde i * Evangeliet*; en stor Mængde have de samlet i Laden, men Videnskaben kan sige til dem: »Imorgen vil jeg kræve Dit Liv«, forsaavidt det er den, der afgjør, hvad Betydning hvert enkelt Resultat skal have i det Hele. Forsaavidt nu der var et Slags ubevidst Liv i en saadan Mands Viden, forsaavidt kan Videnskaberne siges at kræve hans Liv; forsaavidt det ikke var, er hans Virksomhed ligesom Menneskets, der ved sit døde Legemes Hensmulren bidrager til Jordens Vedligeholdelse. Anderledes har det naturligviis været med andre Phænomener, med saadanne Naturforskere, som ved deres Speculation have fundet eller stræbt at finde hiint archimediske Punct, som er intetsteds i Verden, og nu derfra have betragtet det Hele og seet Enkelthederne i det rette Lys. Og med Hensyn til saadanne skal jeg ikke nægte, at de have gjort et høist velgjørende Indtryk paa mig. Den Ro, den Harmoni, den Glæde, som man finder hos dem, finder man ellers sjelden. Vi have her i Byen 3 værdige Repræsentanter: en Ørsted, hvis Ansigt bestandig har forekommet mig som en Klangfigur, som Naturen netop har anstrøget paa den rette Maade; en Schouw, der afgiver et Studium for en Maler, der vilde male Adam, hvor han giver alle Dyrene Navne; og endelig en Horneman, der fortrolig med enhver Plante staaer som en Patriarch i Naturen. I den Henseende erindrer jeg ogsaa med Glæde det Indtryk, som De gjorde hos mig, der stod som Repræsentant for en stor Natur, der ogsaa skulde have sin Stemme med paa Rigsdagen. Begeistret har jeg været for Naturvidenskaberne og er det endnu; men dog synes mig, at jeg ikke vil gjøre mig dem til mit Hovedstudium. Mig har altid Livet i Kraft af Fornuft og Frihed interesseret meest, at klare og løse Livets Gaade har bestandig været mit Ønske. De 40 Aar i Ørkenen, inden jeg naaede Videnskabernes forjættede Land, forekomme mig for kostbare og det saa meget mere, som jeg troer, at Naturen ogsaa har en Side at betragtes fra, hvortil ikke hører Indsigt i Videnskabens Hemmeligheder. Hvad enten saa jeg i det enkelte Blomster skuer hele Verden, eller jeg lytter til de mange Vink, som Naturen frembyder med Hensyn til Menneskelivet; eller jeg beundrer hine dristige Frihaandstegninger paa Firmamentet; eller jeg f. Ex. ved hine Naturlyd paa Ceylon erindres om hine Toner i den aandelige Verden, eller ved Trækfuglenes Bortgang erindres om de dybere Længsler i Menneskets Bryst.

2. Theologien. Det synes at være det, jeg nærmest har grebet; men ogsaa her møder store Vanskeligheder. Her staaer i Christendommen selv saa store Modsætninger, at det idetmindste meget forhindrer det frie Blik. Orthodoxien er jeg saa at sige voxet op i, som De nok veed; men saasnart jeg selv begyndte at tænke, begyndte efterhaanden den uhyre Colos at vakle. Jeg kalder den med Flid en uhyre Colos, thi den har virkelig i det Hele taget megen Consequents, og de enkelte Dele ere i de mange henrundne Aarhundreder smeltede saa tæt sammen, at det er vanskeligt at komme den tillivs. Jeg kunde nu vel være enig med den i enkelte Puncter, men disse bleve da at betragte som de Spirer, man ofte kan finde i Klipperifter. Paa den anden Side kunde jeg vel ogsaa see det Skjæve, der laae i mange enkelte Puncter, men Hovedbasis maatte jeg for en Tid lade staae in dubio. I det Øieblik den forandredes, antog naturligviis det Hele en ganske anden Skikkelse, og saaledes henledes min Opmærksomhed paa et andet Phænomen: Rationalismen, som i det Hele taget gjør en temmelig maadelig Figur. Forsaavidt nemlig Fornuften consequent forfølger sig selv og – ved at gjøre sig Rede for Forholdet mellem Gud og Verden – nu igjen kommer til at see Mennesket i dets dybeste og inderligste Forhold til Gud og i den Henseende ogsaa fra sit Standpunct af betragter Christendommen som den, der i saa mange Aarhundreder har tilfredsstillet Menneskets dybeste Trang, forsaavidt er der Intet at indvende imod den, men det bliver da heller ikke Rationalisme, thi Rationalisme faaer sin egenlige Farve af Christendommen og staaer altsaa paa et ganske andet Gebet og danner ikke et System men en Noahs Ark (for at benytte et af Prof. Heiberg ved en anden Leilighed anvendt Udtryk), hvori de rene og urene Dyr ligge ved Siden af hinanden. Den gjør omtrent samme Indtryk som vor Borgervæbning i ældre Tid vilde gjøre ved Siden af Potsdammer-Garden. Den søger derfor væsentlig at knytte sig til Christendommen, baserer sine Udviklinger paa Skriften og skikker en Legion af Skriftsteder forud for hvert enkelt Punct; men selve Udviklingen er ikke gjennemtrængt deraf. De bære dem ad som Cambyses, der paa hans Tog mod Ægypten skikkede de hellige Høns og Katte iforveien, men de ere derfor ogsaa parate til ligesom den romerske Consul at kaste de hellige Høns overbord, naar de ikke vil æde. Feilen ligger altsaa deri, at de, naar de finde Eenhed med Skriften, basere sig derpaa, men ellers ikke, og saaledes staae paa to fremmedartede Standpuncter.

\sectionheading{Nonnulla desunt}

Nonnulla desunt.

– Hvad smaa Ubehageligheder angaaer, vil jeg kun bemærke, at jeg er ifærd med at læse til theologisk Attestats, en Beskjæftigelse, som slet ikke interesserer og som derfor heller ikke gaaer synderlig rask fra Haanden. Jeg har altid holdt mere af det frie, maaskee derfor ogsaa lidt ubestemte Studium end af den Beværtning ved sluttet Bord, hvor man veed forud, hvilke Gjæster man træffer, og hvilke Retter Mad man faaer hver Dag i Ugen. Da det imidlertid engang er en Nødvendighed, og man knap faaer Lov at komme ind paa de videnskabelige Fælleder uden at være indbrændt, og jeg saavel anseer det med Hensyn til min nærværende Sindstilstand for noget Gavnligt for mig, som ogsaa veed, at jeg dermed kan gjøre Fader en stor Glæde (han mener nemlig, at det egenlige Canaan ligger paa den anden Side af theologisk Attestats, men bestiger derhos ligesom Moses fordum Thabor og beretter, at jeg aldrig kommer derind – dog vil jeg haabe, at det ikke ogsaa denne Gang gaaer i Opfyldelse), saa faaer jeg vel at gribe mig an. Hvor lykkelig er ikke De, som i Brasilien har fundet en uhyre Mark for Deres Iagttagelser, hvor der ved hvert Skridt frembyder sig nye Mærkeligheder, hvor den øvrige lærde Republiques Skrig ikke forstyrrer Deres Ro. Thi mig forekommer den lærde theologiske Verden som Strandveien en Søndag Eftermiddag i Dyrehaugstiden – de storme hverandre forbi, huje og skrige, lee og gjøre Nar af hverandre, kjøre Hestene ihjel, vælte og blive overkjørte, og naar de saa endelig komme overstøvede og forpustede til Bakken – ja saa see de paa hverandre – og vende hjem.

Hvad Deres Tilbagekomst angaaer, saa vilde det være barnagtigt af mig at paaskynde den, ligesaa barnagtigt som naar Achilles Moder søgte at skjule ham for at undgaae den hurtige, hæderlige Død. – – Lev vel!

\sectionheading{Gilleleie, den 1. August 1835}

Gilleleied. 1. Aug. 1835.

Saaledes som jeg har søgt at vise det paa de foregaaende Sider, stode virkelig Sagerne for mig. Naar jeg nu derimod søger at bringe mig selv paa det Rene om mit Liv, forekommer det mig anderledes. Ligesom det varer længe inden Barnet lærer at skjelne sig selv fra Gjenstandene, og det saaledes i lang Tid saa lidet udsondrer sig fra sine Omgivelser, at det med Fremhævelse af den lidende Side siger f. Ex.: »mig slaaer Hesten«, saaledes gjentager det samme Phænomen sig i en høiere aandelig Sphære. Jeg troede derfor, at jeg muligen vilde komme mere til Ro ved at gribe et andet Fag, ved at rette mine Kræfter paa et andet Maal. For en Tid vilde det vel og have lykkedes mig at fordrive en vis Uro, men forstærket vilde den vistnok være vendt tilbage som Feberanfaldet efter Nydelsen af koldt Vand. Det, der egenlig mangler mig, er at komme paa det Rene med mig selv om, hvad jeg skal gjøre, ikke om hvad jeg skal erkjende, uden forsaavidt en Erkjenden maa gaae forud for enhver Handlen. Det kommer an paa at forstaae min Bestemmelse, at see, hvad Guddommen egenlig vil, at jeg skal gjøre; det gjælder om at finde en Sandhed, som er Sandhed for mig, at finde den Idee, for hvilken jeg vil leve og døe. Og hvad nyttede det mig dertil, om jeg udfandt en saakaldet objectiv Sandhed; om jeg arbeidede mig igjennem Philosophernes Systemer og kunde, naar forlangtes, holde Revue over dem; at jeg kunde paavise Inconsequentser indenfor hver enkelt Kreds; – hvad nyttede det mig dertil, at jeg kunde udvikle en Statstheorie og af de mangestedsfra hentede Enkeltheder combinere en Totalitet, construere en Verden, hvori jeg saa igjen ikke levede, men som jeg blot holdt frem til Skue for Andre; – hvad nyttede det mig at kunne udvikle Christendommens Betydning, at kunne forklare mange enkelte Phænomener, naar den for mig selv og mit Liv ikke havde nogen dybere Betydning? Og jo mere jeg kunde det, jo mere jeg saae Andre tilegne sig mine Tankefostre, desto sørgeligere blev min Stilling, ei ulig de Forældres, der af Armod maa skikke deres Børn ud i Verden og overlade dem til Andres Pleie. Hvad nyttede det mig, at Sandheden stod der for mig kold og nøgen, ligegyldig ved, om jeg anerkjendte den eller ikke, snarere bevirkende en ængstelig Gysen end en tillidsfuld Hengivelse? Vel vil jeg ikke nægte, at jeg endnu antager et Erkjendelsens Imperativ; og at der igjennem det ogsaa lader sig virke paa Menneskene, men da maa den levende optages i mig, og det er det, jeg nu anerkjender for Hovedsagen. Det er derefter min Sjæl tørster, som Africas Ørkener efter Vand. Det er det, der mangler, og derfor staaer jeg som en Mand, der samlede Indbo og havde leiet Værelser, men havde endnu ikke fundet den Elskede, der skulde dele Livets Med- og Modgang med ham. Men for saaledes at finde hiin Idee eller rettere sagt finde mig selv, nytter det mig ikke at styrte mig endnu mere ind i Verden. Og det var netop det, jeg før gjorde. Derfor troede jeg, at det vilde være godt at kaste mig ind i Jurisprudentsen, for at jeg kunde udvikle min Skarpsindighed i de mange Livets Forviklinger. Her tilbød sig netop en stor Masse af Enkeltheder, hvori jeg kunde fortabe mig; her kunde jeg maaskee fra de givne Facta udconstruere en Totalitet, en Organisme af Tyve-Liv, forfølge det i hele dets dunkle Side (ogsaa her er en vis Associations-Aand i høi Grad mærkelig). Derfor kunde jeg ønske at blive Acteur, for at jeg ved at sætte mig ind i en Andens Rolle kunde faae, saa at sige, et Surrogat for mit eget Liv, og ved den udvortes Afvexling finde en vis Adspredelse. Det var det, der manglede mig, at føre et fuldkommen menneskeligt Liv og ikke blot et Erkjendelsens, saa jeg derved kommer til ikke at basere mine Tanke-Udviklinger paa – ja paa Noget, man kalder Objectivt, – Noget, som dog i ethvert Tilfælde ikke er mit eget, men paa Noget, som hænger sammen med min Existents's dybeste Rod, hvorigjennem jeg saa at sige er indvoxet i det Guddommelige, hænger fast dermed, om saa hele Verden styrter sammen. See det er det, jeg mangler, og derhen stræber jeg. Med Glæde og inderlig Styrkelse betragter jeg derfor de store Mænd, som saaledes have fundet hiin Ædelsteen, for hvilken de sælge Alt, endog deres Liv, hvad enten jeg seer dem med Kraft gribe ind i Livet, med sikkre Skridt, uden at vakle, at gaae frem paa deres bestemte Bane; eller jeg opdager dem afsides fra den alfare Vei, fordybede i sig selv og i Arbeiden for deres ophøiede Maal. Med Ærbødighed betragter jeg endog de Afveie, der ogsaa her ligge saa nær. Det er denne Menneskets indvortes Handlen, denne Menneskets Guds-Side det kommer an paa, ikke en Masse af Erkjendelser; thi de ville vel følge og ville da ikke vise sig som tilfældige Agregater eller som en Række af Enkeltheder, den ene ved Siden af den anden, uden et System, uden et Brændpunct, hvori alle Radier samles. Et saadant Brændpunct har jeg vel ogsaa søgt. Saavel paa Forlystelsernes bundløse Hav har jeg forgjeves søgt en Ankerplads som i Erkjendelsens Dybder. Jeg har følt den næsten uimodstaaelige Magt, hvormed den ene Forlystelse rækker den anden Haanden; jeg har følt den Art af uægte Begeistring, som den er istand til at fremkalde; jeg har ogsaa følt den Kjedsomhed, den Sønderrevethed, der følger derpaa. Jeg har smagt Frugterne af Kundskabens Træ og tidt og ofte glædet mig ved deres Velsmag. Men denne Glæde var kun i Erkjendelsens Øieblik og efterlod ikke noget dybere Mærke i mig selv. Det forekommer mig, at jeg ikke har drukket af Viisdommens Bæger men er faldet ned deri. Jeg har søgt at finde dette Princip for mit Liv ved Resignation, ved at forestille mig selv, at da Alt gik efter urandsagelige Love, det ikke kunde være anderledes, ved at sløve min Ærgjerrighed og Forfængeligheds Følehorn. Fordi jeg ikke kunde faae Alt til at passe efter mit Hoved, trak jeg mig ud deraf med Bevidsthed om min egen Dygtighed, omtrent som naar en aflægs Præst resignerer med Pension. Hvad fandt jeg? Ikke mit Jeg; thi det søgte jeg netop at finde paa hine Veie (jeg tænkte mig, om jeg saa maa sige, min Sjæl som indesluttet i en Kasse med Springlaas for, og som nu de udvortes Omgivelser ved et Tryk paa Fjedrene skulde faae til at springe op). – Det var altsaa det første, der skulde afgjøres, det var hiin Søgen og Finden af Himmerigets Rige. Ligesaalidet som et Verdenslegeme, naar det tænktes at skulle danne sig, først vilde fatte Bestemmelse om, hvorledes dets Overflade skulde være, til hvilke Legemer det skulde vende den lyse, til hvilke den mørke Side, men først vilde lade Centrifugal- og Centripetal-Kræfternes Harmoni realisere dets Existents og lade Resten komme af sig selv – ligesaalidet nytter det Mennesket først at ville bestemme det Udvortes og siden det Grundconstituerende. Man maa først lære at kjende sig selv, inden man kjender noget Andet (γνωϑι σεαυτον). Først naar Mennesket saaledes inderlig har forstaaet sig selv og nu seer sin Gang hen af sin Bane, først da faaer hans Liv Ro og Betydning, først da bliver han fri for hiin besværlige, uheldsvangre Reisekammerat – hiin Livs-Ironi, der viser sig i Erkjendelsens Sphære og byder den sande Erkjenden at begynde med en Ikke-Erkjenden (Socrates), ligesom Gud skabte Verden af Intet. Men fornemlig hører den hjemme paa Sædelighedens Farvande for dem, som endnu ikke ere komne ind under Dydens Passat. Her tumler den Mennesket paa det rædsomste omkring; snart lader den Mennesket føle sig lykkelig og tilfreds ved Forsættet om at vandre frem ad den rigtige Vei, snart styrter den ham ned i Fortvivlelsens Afgrunde. Ofte dysser den Mennesket i Søvn ved den Tanke: »det kan nu eengang ikke være anderledes«, for pludselig at vække ham til et strængt Forhør. Ofte lader den ligesom Glemselens Slør falde over det Forbigangne, for da atter igjen at lade hver enkelt Ubetydelighed træde frem i et levende Lys. Naar Mennesket her kjæmper sig frem paa den rette Vei, glæder sig over at have beseiret Fristelsernes Magt, kommer maaskee næsten i samme Moment ovenpaa den fuldkomneste Seir en tilsyneladende ringe udvortes Omstændighed, der støder ham lig Sisyphus ned fra Klippens Spidse. Ofte naar Mennesket har concentreret sin Kraft paa Noget, møder en lille udvortes Omstændighed, der tilintetgjør det Hele. (Jeg vilde sige: som en Mand, der kjed af Livet vilde styrte sig i Themsen, og som nu et Myggestik standser netop i det afgjørende Øieblik.) Ofte lader den Mennesket ligesom den Brystsvage befinde sig allerbedst, naar det er allerværst. Forgjeves søger han at arbeide imod; han har ikke Kraft nok dertil, intet hjælper det ham, at han saa ofte gjennemgaaer det Samme; den Art af Øvelse, som derved opnaaes, kommer det ei an paa. Ligesaa lidet som den, der har nok saa megen Øvelse i at svømme, vil kunne holde sig oppe i et Uveir, men kun den, der inderlig er overbeviist om og har erfaret, at han virkelig er lettere end Vandet, ligesaalidet kan det Menneske, der mangler det indre Holdningspunct, holde sig oppe i Livets Storme. – Kun naar Mennesket saaledes har forstaaet sig selv, er han istand til at hævde sig en selvstændig Existents og saaledes undgaae at opgive sit eget Jeg. Hvor ofte see vi ikke – (i en Tid, hvor vi ved vore Lovtaler ophøie hiin græske Historieskriver, fordi han vidste at adoptere en fremmed Stiil til den meest skuffende Lighed med dens Ophavsmand, istedetfor at han burde dadles, da den første Priis for en Skribent altid er at have en egen Stiil ɔ: en ved sin Individualitet modificeret Udtryks- og Fremstillings-Maade) – hvor ofte see vi ikke Folk, der enten af aandelig Ladhed leve af de Smuler, der falde fra Andres Bord, eller af mere egoistiske Grunde søge at leve sig ind i Andre, for da tilsidst ligesom en Løgner ved den hyppige Gjentagelse af sine Historier selv at troe dem. Uagtet jeg endnu langt fra ikke saaledes inderlig er kommen til at forstaae mig selv, har jeg med dyb Agtelse for dens Betydning søgt at frede om min Individualitet – dyrket den ubekjendte Guddom. Med en utidig Ængstelighed har jeg søgt at undgaae at komme i for nær Berøring med de Phænomener, hvis Tiltrækningskraft maaskee vilde udøve for stor Magt over mig. Jeg har søgt at tilegne mig meget af dem, studeret deres Individualitet og Betydning i Menneskelivet, men dog tillige vogtet mig for ikke som Myggen at komme Lyset for nær. Ved Omgang med de sædvanlige Mennesker har jeg kun havt Lidet at vinde eller tabe. Deels har deres hele Syslen – det saakaldte practiske Liv – ikkun lidet interesseret mig; deels har den Kulde og den Mangel paa al Sympathie, hvormed de betragte de aandelige og dybere Bevægelser i Mennesket, fjernet mig endmere fra dem. Mine Omgangsvenner have med faa Undtagelser ikke udøvet nogen synderlig Magt over mig. Et Liv, der ikke er kommen paa det Rene med sig selv, maa nødvendigviis vise en ujævn Side-Overflade frem; her ere de nu kun blevne staaende ved enkelte Facta og deres tilsyneladende Disharmoni, thi til at søge at opløse den i en høiere Harmoni eller til at indsee Nødvendigheden deraf, dertil havde de ikke Interesse nok for mig. Deres Dom over mig var derfor altid eensidig, og jeg har afvexlende enten lagt for megen eller for liden Vægt paa deres Udsagn. Ogsaa deres Indflydelse og de deraf mulige Misvisninger paa mit Livs Compas har jeg nu unddraget mig. Og saaledes staaer jeg da atter paa det Punct, at jeg maa begynde paa en anden Maade. Jeg vil nu forsøge rolig at fæste Blikket paa mig selv og begynde inderlig at handle; thi kun derved bliver jeg istand til, ligesom Barnet ved sin første med Bevidsthed foretagne Handling kalder sig: »jeg«, saaledes istand til i dybere Betydning at kalde mig: »jeg«.

Dog dertil udfordres Udholdenhed, og man kan ikke strax høste, hvor man har saaet. Jeg vil erindre hiin Philosophs Methode: at lade sine Disciple tie i tre Aar, saa kommer det vel. Ligesom man ikke begynder en Fest med Solens Opgang men med dens Nedgang, saaledes maa man ogsaa i den aandelige Verden først arbeide i nogen Tid frem, førend Solen ret kan skinne for os og gaae op i al sin Herlighed; thi vel hedder det: at Gud lader sin Sol opgaae over Gode og Onde og lader regne over Retfærdige og Uretfærdige, men ikke saa i den aandelige Verden. Saa være da Loddet kastet – jeg gaaer over Rubicon! Vel fører denne Vei mig til Kamp; men jeg vil ikke forsage. Jeg vil ikke sørge over den forbigangne Tid – thi hvortil Sorg? Med Kraft vil jeg arbeide mig frem og ikke spilde Tiden med at sørge ligesom den, der i en Hængesæk først vilde beregne, hvor dybt han var sunken, uden at erindre, at han i den Tid, han anvender hertil, synker endnu dybere. Jeg vil ile frem paa den fundne Vei og tilraabe Enhver, som jeg møder: ikke som Loths Hustru at see sig tilbage, men huske paa, at det er en Bakke, vi stræbe op af.

\sectionheading{[Livet maa leves forlænds]  (JJ:167, 1843)}

Det er ganske sandt, hvad Philosophien siger, at Livet maa forstaaes baglænds. Men derover glemmer man den anden Sætning, at det maa leves forlænds. Hvilken Sætning, jo meer den gjennemtænkes, netop ender med, at Livet i Timeligheden aldrig ret bliver forstaaeligt, netop fordi jeg intet Øieblik kan faae fuldelig Ro til at indtage Stillingen: baglænds.

\workheading{Af en endnu Levendes Papirer}{1838}

\worknote{Antologiens ''From the Papers of One Still Living'' — Kierkegaards anmeldelse af H.C. Andersens roman ''Kun en Spillemand''. Hong-udvalget tager kritikken af Andersens manglende livs-anskuelse og den afsluttende personlige bemærkning; den detaljerede gennemgang af romanen (SV¹ XIII, 76–90) er udeladt.}

\sectionheading{Andersen mangler en Livs-Anskuelse  ·  SV¹ XIII, 68–75}

Naar vi nu sige, at Andersen aldeles mangler Livs-Anskuelse, saa begrundes denne Yttring ligesaa meget i det Foregaaende, som den selv, efterviist i sin Sandhed, begrunder dette. En Livs-Anskuelse er nemlig mere end et Indbegreb eller en Sum af Sætninger, fastholdt i sin abstrakte Hverkenhed; den er mere end Erfaringen, der som saadan altid er atomistisk, den er nemlig Erfaringens Transsubstantiation, den er en tilkæmpet af al Empirie urokkelig Sikkerhed i sig selv, hvad enten den saa blot har orienteret sig i alle verdenslige Forhold (et blot humant Standpunkt, Stoicisme for Exempel), der derved holder sig uden for Berøring med en dybere Empirie – eller den har i sin Retning mod Himlen (det Religieuse) deri fundet det Centrale, saavel for sin himmelske som sin jordiske Existents, har vundet den sande christelige Forvisning: »at hverken Død, ei heller Liv, ei heller Engle, ei heller Fyrstendømmer, ei heller Magter, ei heller det Nærværende, ei heller det Tilkommende, ei heller det Høie, ei heller det Dybe, ei heller nogen anden Skabning skal kunne skille os fra Guds Kjærlighed i Christo Jesu vor Herre.« See vi nu hen til, hvorledes det i denne Henseende staaer sig med Andersen, da finde vi Forholdet ganske som vi havde ventet det: paa den ene Side stikke enkelte Sætninger frem som Hieroglypher, der stundom ere Gjenstand for en from Veneration; paa den anden Side dvæler han ved sin egen Erfarings enkelte Phænomener, som nu snart igjen potenseres til Sætninger, og ere da at subsumere under den foregaaende Classe, snart mere afvindes den Side at være oplevede, uden at man derfor med Rette kan, saalænge disse forblive i deres Pebersvende-Stand, deraf drage nogen videre Følge. Vil man nu maaskee sige, at den af os skildrede Livs-Anskuelse er et Standpunkt, man kun efterhaanden kan nærme sig til, og at det er ubilligt at gjøre saa store Fordringer til en saa ung Mand som Andersen, saa ville vi, hvad det Sidste angaaer, medens vi gjerne indrømme, at Andersen er en ung Mand, dog erindre om, at vi kun have med Andersen som Roman-Digter at gjøre, og anticipando tilføie, at en saadan Livs-Anskuelse er for en Roman-Digter af den Classe, til hvilken Andersen hører, conditio sine qua non; og, hvad det Første angaaer, gjerne indrømme en vis Aproximation i sin fulde Betydning, men dog ogsaa i Tide sige Stop, inden vi faae den hele vor Betragtning annullerende Consequents paa Halsen, at den egentlige Livs-Anskuelse indtræder først (demum) i Ens Dødsstund, eller vel endog paa en af Planeterne. Spørge vi nu, hvorledes en saadan Livs-Anskuelse kommer istand, da svare vi, at for den, der ikke tillader sit Liv altfor meget at futte ud, men saavidt som muligt søger at føre dets enkelte Yttringer tilbage i sig selv igjen, maa der nødvendigviis indtræde et Øieblik, i hvilket der udbreder sig et besynderligt Lys over Livet, uden at man derfor i fjerneste Maade behøver at have forstaaet alle mulige Enkeltheder, til hvis successive Forstaaelse man imidlertid nu har Nøglen, maa, siger jeg, indtræde det Øieblik, da, som Daub bemærker, Livet forstaaes baglænds gjennem Ideen. Har man nu ikke bragt det hertil, ja mangler man endog aldeles Forstand paa, hvad dette vil sige, saa kommer man til parodisk at stille sig en Livs-Opgave, enten derved, at denne allerede er løst, om man saa vil, skjøndt den i en anden Forstand aldrig har været knyttet, eller derved, at den aldrig bliver at løse. Begge Dele finde vi nu til yderligere Bestyrkelse hos Andersen, idet begge Anskuelser saavel foredrages som Laanesætninger, som ogsaa anskueliggjøres til en vis Grad i enkelte poetiske Personligheder. Paa den ene Side læres nemlig, at der er skrevet paa ethvert Menneske et mene mene etcetera; i Analogie hermed optræde Individer, hvis egentlige Opgave ligger bagved dem, uden at dette derfor hjælper dem til at komme i den for Livets Betragtning rigtige Stilling »baglænds«, da denne Opgave snarere som en Pukkel er anbragt paa deres egen Ryg, og som de derfor ret egentlig aldrig faae at see, eller i aandelig Forstand umulig nogensinde kunne blive sig bevidste, med mindre Andersen til Afvexling lægger en den hele Opfattelse forstyrrende Bevidsthed ind i dem, Individer, der som andre Himmellegemer gaae deres engang anviste Vei med en ufravigelig Nøiagtighed, eller han taber sig i ikke saa meget høittravende, som langttravende Betragtninger, hvis Helt er en storartet Peripatetiker, der, da han ingen nødvendig Grund har til nogetsteds at standse, men Tilværelsen derimod altid er en Cirkel, kommer til at gaae i Cirkel, om end Andersen og Andre, der i mange Aar have boet paa Bjerget, troe at han gaaer lige ud, da Jorden er saa flad som en Pandekage. Imellem, det vil sige i Eenheden af disse Standpunkter, ligger den sande Midte; men deraf følger ingenlunde, at der jo ved en ny Inconsequents, der, vel at mærke, ikke hæver den foregaaende (thi det var det heldigste) kan opstaae nye Phænomener, idet for Exempel Andersen pludselig afbryder deres ufortrødne Vandring, idømmer dem en arbitrair Straf, skjærer Næse og Øren af dem, sender dem til Siberien, og saa maa vor Herre, eller hvo der ellers vil, tage sig af dem. –

Men er da en saadan Livs-Anskuelse absolut nødvendig for en Roman-Digter, eller gives der ikke en vis poetisk Stemning, der som saadan kan i Forening med en livlig Skildring udrette det Samme? Vort Svar herpaa ligger for største Delen i det af os tidligere med Hensyn til Blicher Udviklede, hvor vi jo netop have søgt at paavise Betydningen af en saadan af Stemning udgaaet og af Læserne gjennem en Succession af Modificationer i et heelt Billede opfattet Eenhed, og forøvrigt ville vi, forsaavidt man vil gjøre en lignende Betragtning gjældende ved en Utallighed af i Reflexionen givne Standpunkter, hvorved man da maa erindre, at alle disse Standpunkter som saadanne med Hensyn til Produktivitæt sætte væsentlig af og meer og mere lader den oprindelige Stemning fordampe, blot tilføie, at der paa alle disse vel lader sig producere, men at man dog, naar man er en Smule kræsen med sine Titler, hellere maatte kalde dem Studier til Noveller etcetera end Noveller, som det jo ogsaa paa dette Trin vil mislykkes i samme Grad, som man egentlig sætter sig en novellistisk eller Roman-Opgave. Maaskee vil man gaae videre, og idet man beraaber sig paa, at der jo dog er een Idee, der idelig forekommer i Andersens Romaner, derved vindicere Andersen en Livs-Anskuelse, og forekaste mig, der jo selv indrømmer dette, min Inconseqvents. Hertil maa jeg svare, at jeg aldrig har paastaaet, at en Idee som saadan (allermindst en fix Idee) er at ansee for en Livs-Anskuelse, og dernæst, at jeg, for at indlade mig paa denne Undersøgelse, maa vide lidt nærmere om denne Idees Indhold. Gaaer den nu ud paa, at Livet ikke er en Udviklings-Proces, men en Undergangs-Proces af det Store og Udmærkede, der vilde spire frem, saa troer jeg dog vel med Rette at turde protestere mod Prædikatet Livs-Anskuelses Anvendelse herpaa, forsaavidt man ellers vil give mig Ret i, at Skepsis som saadan ikke er en Erkjendelsens-Theorie, eller for at blive i mit Forehavende, at en saadan Mistillid til Livet vel indeholder en Sandhed, forsaavidt den leder til at finde en Tillid (for Exempel naar Salomo siger, at Alt er forfængeligt), men derimod i samme Moment, som den determinerer sig til en endelig Afgjørelse af Livets Spørgsmaal, indeholder en Usandhed. Men vi gaae videre. Vi ville for et Øieblik antage, at man kunde have Ret i at kalde en saadan i Reflexion vilkaarligt standset og nu til endelig Sandhed potenseret Betragtning en Livs-Anskuelse; vi ville tænke os et Individ, der, stærkt omtumlet af en hæftig bevæget Samtid, endelig afslutter sig paa et saadant Standpunkt; vi ville lade ham producere Noveller. Et Modermærke ville de alle faae; men i samme Grad, han havde oplevet Meget, i samme Grad han ret havde deeltaget i Livets Svingninger, i samme Grad vilde han ogsaa være istand til i sine Noveller at udfolde en stor Suite af rædsomme Conseqventser, der alle sigtede til hans Helts endelige Undergang; i samme Grad vilde man en lang Tid føle sig fristet til at troe paa hans Livs-Opfattelses Sandhed. Men er dette Tilfældet hos Andersen? det vil vist Ingen paastaae. Tvertimod springer Andersen den egentlige Udvikling over, lægger et behørigt Tidsrum imellem, lader først see, saa godt han kan, de store Kræfter og Anlæg, og derpaa disses Fortabelse. Her vil man dog vel give os Ret i, at det ingen Livs-Anskuelse er. Til yderligere Opklarelse af vor Mening ville vi blot tilføie, at en saadan Fortabelsens-Theorie kan udgaae deels fra et alvorligt foretaget, men dog mislykket Forsøg paa at forstaae Verden, derved, at Individet, deprimeret af Verden, skjøndt længe arbeidende imod, tilsidst ligger under; eller den kan fremkaldes derved, at man ved den allerførste Vaagnen af Reflexionen nu ikke kaster Blikket ud efter, men øieblikkelig ind i sig selv og i sin saa kaldte Betragtning af Verden blot nøiagtigt gjennemfører sin egen Liden. Det første er en mislykket Activitæt, det andet en oprindelig Passivitæt, det første er knækket Mandighed, det andet gjennemført Qvindelighed. Vi vende tilbage til vort Forehavende: gjennem en kort Antydning af en Livs-Anskuelses Nødvendighed for Roman- og Novelle-Digteren, at udvikle, hvorledes det i denne Henseende staaer sig med Andersen. En Livs-Anskuelse er egentlig Forsynet i Romanen, den er dens dybere Eenhed, der giver denne at have Tyngdepunktet i sig; den befrier denne fra at blive vilkaarlig eller hensigtsløs, idet Hensigten overalt er immanent tilstede i Konstværket. Naar derimod en saadan Livs-Anskuelse mangler, da søger Romanen enten paa Poesiens Bekostning at insinuere en eller anden Theorie (dogmatiske, doctrinaire Noveller), eller den træder i endeligt og tilfældigt Forhold til Forfatterens Kjød og Blod. Dette Sidste kan imidlertid gaae for sig igjennem en stor Mangfoldighed af Modificationer, lige fra en uvilkaarlig Overstrømmen af det Dygtigere i Personligheden, lige indtil det Punkt, hvor Forfatteren maler sig selv med, ligesom Landskabsmalerne undertiden ynde det, ja endog glemme, at dette kun havde Betydning, forsaavidt man opfattede det som Situation, og derfor gaae over til, aldeles forglemmende Landskabet, at udpensle sig selv i sin egen forfængelige salomoniske Pragt-Herlighed, der klæder Blomsterne, men ikke Menneskene. (Naar jeg i denne Henseende danner en Modsætning mellem de doctrinaire og de subjektive Romaner, saa seer jeg meget godt, at det dog kun er gjennem en Subdivision, at disse blive hinanden coordinerede; thi ogsaa de doctrinaire Romaner staae i et tilfældigt Forhold til Personligheden, da deres Forfattere gjennem en tilfældig Villiesbestemmelse acquiescere i Sætninger, som de endnu ikke tilstrækkelig have gjennemlevet.) Medens imidlertid begge Classer af Romaner staae i et endeligt og forkeert Forhold til Personligheden, saa mener jeg ingenlunde, at Romanen, i en vis prosaisk Forstand skal abstrahere fra Personligheden, eller at man fra et andet Standpunkt kunde med Rette gjøre Fordringer til den, som til stræng Speculation; men at Digteren først og fremmest maa selv tilkæmpe sig en dygtig Personlighed, og at det kun er denne saaledes døde og forklarede Personlighed, der bør og kan producere, ikke den mangekantede, jordiske, palpable. Hvor vanskeligt det er at tilkæmpe sig en saadan Personlighed, kan man ogsaa see deraf, at der i mange ellers dygtige Romaner, findes ligesom et Residuum af Forfatterens endelige Charakteer, der som en næsviis tredie Person, som et uopdragent Barn ofte snakker med paa upassende Steder. Anvende vi nu dette paa Andersen, saa troe vi med Rette, ikke saa meget argumenterende af som apostropherende til den temmelig gængse Conversations-Dom, »det er altid Andersen, han er altid sig selv etcetera«, at turde regne ham til den Classe af Roman-Digtere, der give en upoetisk Tilgift af deres egen blot phænomenologiske Personlighed, uden at vi derfor ville tillægge ham en saa stærk Villies-Determination i denne Henseende, at den forhindrede ham i ogsaa at streife over paa de doctrinaire Novellers Gebeet, dog ikke saaledes, at han skulde have en eller anden større Theorie, han vilde gjøre gjældende, men snarere, som ovenfor viist, som en Forkjærlighed for og Overvurderen af visse enkelte Sætninger, der i et saa strængt Coelibat, som Forfatteren holder dem i, ikke have stort at sige. I hvor stærk Forstand Andersens Romaner saaledes staae i et forkeert Forhold til hans Person, vil man bedst overbevise sig om ved at reproducere det Total-Indtryk, hans Romaner efterlade sig. Vi mene ingenlunde, at det er forkeert, at det enkelte Individ gaaer under i Romanen; men dette maa blot være en poetisk Sandhed, ikke, som hos nogle Digtere, en Opdragelsens pia fraus, eller som hos Andersen en sidste Villie; vi fordre ingenlunde i nogen større Betydning Forstandighed og Klarhed over Livet i ethvert enkelt af hans poetisk skabte Individer, tvertimod ville vi indrømme ham fuld Magt til at lade dem, om galt skulde være, blive afsindige; kun maa det ikke skee saaledes, at en Galskab i 3die Person afløses af en i første, at Forfatteren selv overtager den Gales Partes. Der maa i Romanen være en udødelig Aand, der overlever det Hele. Hos Andersen derimod er der aldeles intet Hold: naar Helten døer, døer Andersen med og afnøder i det Høieste Læseren som sidste Indtryk et Suk over dem begge. – Idet vi saaledes her flere Gange har omtalt Andersens Person og Personlighed, vil jeg – til Svar paa en dog kun ved Misforstaaelse og Mistydning mulig Indvending, som om jeg her overtraadte Grændsen for min æsthetiske Jurisdiction og den indenfor denne indrømmede Competence ved at omtale Andersen som Menneske – vil jeg, uden at beraabe mig paa den Omstændighed, at jeg saa godt som ikke kjender Andersen personlig, alene bemærke, at da den egentlige poetiske Frembringelse, især paa Novellens og Romanens Gebeet, ikke er andet end en fyldig, i en friere Verden sig indskabende og i denne sig rørende, reproducerende 2den Potens af det i første Potens allerede paa mangfoldig Maade poetisk Oplevede, og da man desuden i Andersens Romaner deels savner den constituerende Total-Overskuelse (en Livs-Anskuelse), deels støder atter og atter paa Situationer, Bemærkninger etcetera, hvilke man vel ikke kan nægte at være poetiske, men som dog hos Andersen staae som ufordøiede og som i poetisk Forstand (ikke saaledes i mercantilsk) ubenyttede, utilegnede, uaffiltrerede, saa kunde man med Rette gjøre den Slutning tilbage, at Andersen selv ikke har i første Potens levet poetisk-klart, da det Poetiske i anden Potens ikke har vundet større Consolidation i det Hele, og ikke er mere udvadsket i det Enkelte; man kunde gjøre denne Slutning, siger jeg, uden at jeg derfor har gjort den eller vil gjøre den, men jeg vil blot gjøre Andersen opmærksom paa, at dersom han paa en eller anden Maade skulde føle sig personligt (som bosiddende Mand i Kjøbenhavn) afficeret, Grunden da ikke ligger hos mig, men hos Andersen, hvis Romaner staae i et saa physisk Forhold til ham selv, at deres Tilblivelse ikke er saa meget at ansee for en Produktion som for en Amputation af ham selv; og det veed man jo nok, at om det saaledes Amputerede end er langt borte, føler man stundom uvilkaarligt en reen physisk Smerte deri.

\sectionheading{Slutning — med sympathetisk Blæk til Andersen  ·  SV¹ XIII, 91–92}

Hvad jeg til Slutning har at sige, foranlediget ved det vistnok i det Hele taget som factisk indrømmede Misforhold mellem en læsende og en critiserende Verdens Dom over Andersen, forsaavidt dette ogsaa har gjentaget sig i min Bevidsthed, da vilde jeg ønske, at det maa lykkes mig herom at tale ligesaa personligt, som jeg har søgt at holde det Foregaaende frit for ethvert obliquet Forhold til min Personlighed. Idet jeg nemlig reproducerer det første Stadium, levendegjøres Erindringen om en Mangfoldighed af poetiske Stemninger, som ethvert digterisk Liv, endogsaa det mest uklare, og i en vis Forstand dette maaskee mest, maa være gjennemflettet med, og idet jeg endnu en Gang søger at fastholde hver enkelt, fortrænger den ene den anden saa stærkt, at Totalitæten af dem samler sig ligesom til Afsked i en eneste Concentration, i et Nærværende, der dog i samme Øieblik føler i sig Nødvendigheden af at blive et Forbigangent, og derved fremkalder et vist veemodigt Smiil hos mig, idet jeg betragter dem, en Taknemlighedsfølelse, idet jeg erindrer mig den Mand, hvem jeg skylder det Hele, en Følelse, som jeg hellere ønskede at hvidske Andersen i Øret end betroe til Papiret, ikke som om det noget Øieblik har været Andet for mig end en Glæde at kunne yde ham, hvad han har Ret til, og som fornemmelig i vor Tid Enhver, der endnu føler lidt for Poesien i den ecclesia pressa, i hvilken vi leve, næsten maa fristes til at yde varmere, end Sandheden maaskee kunde fordre det; ei heller som om en saadan Yttring ikke lod sig bringe i Samklang med min hele tidligere udtalte Anskuelse af Andersen, da jeg altid for mit Vedkommende har gjort mere af Poesie trods al dens Liggen under end af den mest seierrige Prosa, eller for at blive ved Andersen, da jeg, trods al hans Omtumlen, al hans Bøien for enhver poetisk Luftning, dog altid glæder mig over, at han endnu ikke er kommen ind under Politikens altomfattende fandenivoldske Passat – men fordi en saadan Yttring i det Hele taget er meget udsat for Misforstaaelse, Noget jeg imidlertid haaber, at jeg skal kunne finde mig i, naar blot Andersen for at undgaae den, vil holde, hvad jeg med sympathetisk Blæk har skrevet, for det Lys, der ene gjør Skriften læselig og Meningen tydelig.

\workheading{Om Begrebet Ironi, med stadigt Hensyn til Socrates}{1841}

\worknote{Antologiens ''The Concept of Irony''. Hong-udvalget samler tre adskilte stykker: fra Første Deel afsnittet ''Xenophon og Plato'' (afbrudt midt i Aristophanes-afsnittet, markeret med […]), og fra Anden Deel dennes ''Indledning'' samt afsnittet om ironiens verdenshistoriske gyldighed; de mellemliggende afsnit er udeladt.}

\sectionheading{Første Deel: Xenophon og Plato (Aristophanes)  ·  SV¹ XIII, 212–216}

Skulde man med faa Ord udtrykke Platos Opfattelse af Socrates, saa kunde man sige, at han giver ham Ideen. Hvor Empirien ender, der begynder Socrates, hans Virksomhed er at føre Speculationen ud af Endelighedens Bestemmelser, at tabe Endeligheden af Sigte og styre ud paa det Okeanos, i hvilket den ideale Stræben og den ideale Uendelighed intet fremmed Hensyn kjender, men er sig selv sit uendelige Maal. Som derfor den lavere Sandsning blegner ved Siden af denne høiere Erkjendelse, ja bliver et Bedrag, en Skuffelse i Sammenligning med denne, saaledes bliver ethvert Hensyn til et endeligt Maal en Forkleinelse, en Profanation af det Hellige. Kort sagt, Socrates har vundet Idealiteten, erobret disse uhyre Regioner, der hidtil var et terra incognita. Han foragter derfor det Nyttige, er indifferent mod det Bestaaende, en afgjort Fjende af den Middelmaadighed, der hos Empirien er, som det Høieste, Gjenstand for en from Tilbedelse, men for Speculationen en Skifting af Troldeæt. Erindrer man nu derimod det Resultat, vi kom til gjennem Xenophon, at vi her fandt Socrates i travl Virksomhed som en Endelighedens Apostel, som en geskjeftig Probenreuter fra en Maadelighedens Propaganda, ufortrødent anbefalende sit ene saliggjørende jordiske Evangelium, at vi der fandt for det Gode det Nyttige, for det Skjønne det Brugbare, for det Sande det Bestaaende, for det Sympathetiske det Lucrative, for harmonisk Eenhed Nøkternhed, saa vil man vist indrømme, at disse to Opfattelser ikke godt lade sig forene. Enten maatte man sigte Xenophon for en total Vilkaarlighed, for et ubegribeligt Had til Socrates, der søgte en Tilfredsstillelse i en saadan Bagvadskelse; eller man maatte tilskrive Plato en ligesaa gaadefuld Idiosynkrasie for sin Modsætning, der ligesaa uforklarligt realiserede sig i at forvandle denne i Lighed med sig selv. Ville vi nu et Øieblik lade Socrates' Virkelighed staae som en ubekjendt Størrelse, saa kan man med Hensyn til disse to Opfattelser sige, Xenophon har som en Kræmmer pruttet sig sin Socrates til, Plato har som en Kunstner skabt sig sin Socrates i oaturlig Størrelse. Men hvorledes har Socrates virkelig været, hvad har været Udgangspunktet for hans Virksomhed? Svaret herpaa maa naturligviis tillige hjælpe os ud af den Nød, hvori vi hidtil have været bestedte. Svaret herpaa er: Socrates' Existents er Ironi. Ligesom nu dette Svar, saa troer jeg, hæver Vanskeligheden, saaledes gjør den Omstændighed, at det hæver Vanskeligheden, det tillige til det rigtige Svar, saa at det i samme Øieblik viser sig som Hypothese og som det Sande. Det Punkt nemlig, den Streeg, der gjør Ironien til Ironi, er saare vanskelig at faae fat paa. Med Xenophon kan man derfor gjerne antage, at Socrates har yndet at gaae omkring og tale med allehaande Mennesker, fordi enhver udvortes Begivenhed eller Ting for den altid slagfærdige Ironiker er en Anledning; med Plato kan man gjerne lade Socrates berøre Ideen, kun at Ideen ikke aabner sig for ham, men tvertimod er Grændse. Hver af disse to Fremstillere har nu naturligviis søgt at fuldstændiggjøre Socrates, Xenophon ved at drage ham ned i det Nyttiges lave Egne, Plato ved at løfte ham op i Ideens overjordiske Regioner. Men det Punkt, som ligger herimellem, usynligt og saare vanskeligt at fastholde, er Ironien. Paa den ene Side er netop Virkelighedens Mangfoldighed Ironikerens Element, men paa den anden Side er hans Gang henover Virkeligheden svævende og ætherisk, han berører bestandig blot Jorden; men da Idealitetens egentlige Rige endnu staaer ham fremmed, er han endnu ikke udvandret dertil, men ligesom i ethvert Øieblik færdig til Afmarsch. Ironien oscillerer mellem det ideelle Jeg og det empiriske Jeg; det ene vilde gjøre Socrates til Philosoph, det andet til Sophist; men det, der gjør ham til Mere end Sophist, er, at hans empiriske Jeg har universel Gyldighed.

Aristophanes' Opfattelse vil netop afgive den nødvendige Modsætning til Platos, og just ved denne Modsætning tilveiebringe Muligheden af en ny Vei for vor Beregning. Ja det vilde været et stort Savn, om vi manglede den aristophaniske Vurdering af Socrates; thi ligesom enhver Udvikling i Almindelighed ender med at parodiere sig selv, og en saadan Parodi er en Sikkerhed for, at denne Udvikling har overlevet sig selv, saaledes er den comiske Opfattelse et Moment, et i mange Maader uendeligt berigtigende Moment med i den totale Anskueliggjørelse af en Personlighed eller en Stræben. Mangler man derfor end det umiddelbare Vidnesbyrd om Socrates, mangler man end en ganske paalidelig Opfattelse af ham, saa har man til Vederlag alle Misforstaaelsens forskjellige Nuancer, og ved en saadan Personlighed som Socrates troer jeg, at vi ere bedst tjente dermed. Plato og Aristophanes have nu det tilfælleds, at deres Fremstilling er ideal, men i Forhold til hinanden omvendt, Plato har den tragiske Idealitet, Aristophanes den comiske. Hvad der har bevæget Aristophanes til at opfatte Socrates saaledes, om han var bestukken dertil af Socrates' Anklagere, om han var forbittret over Socrates' venskabelige Forhold til Euripides, om han har i ham bekjæmpet Anaxagoras' Naturspeculationer, om han har identificeret ham med Sophisterne, kort sagt, om der har været nogen endelig og jordisk Grund, der har bestemt ham dertil, det er denne Undersøgelse aldeles uvedkommende, og forsaavidt som den skulde give et Svar i denne Henseende, maatte det naturligviis blive et benægtende, da den vedkjender sig den Overbeviisning, at Aristophanes' Opfattelse er ideal, hvorved den allerede er frigjort fra ethvert saadant Hensyn, ikke ludende lister sig langs Jorden, men fri og let overflyver den. Blot at opfatte Socrates' empiriske Virkelighed, at bringe ham paa Scenen, saaledes som han gik og stod i Livet, vilde have været under Aristophanes' Værdighed og forvandlet hans Comoedie til et satirisk Digt; paa den anden Side, at idealisere ham efter en Maalestok, hvorved han blev aldeles ukjendelig, vilde ligge aldeles udenfor den græske Comoedies Interesse. At dette sidste heller ikke var Tilfældet, derom vidner jo Oldtiden, der beretter, at Opførelsen af Skyerne beæredes med Nærværelsen af den i denne Henseende strængeste Kritiker, Socrates selv, der til Publikums Fornøielse stod op under Forestillingen, for at den i Theatret forsamlede Mængde kunde overbevise sig om den behørige Lighed. At en saadan blot excentrisk ideel Opfattelse heller ei vilde ligge i den græske Comoedies Interesse, deri vil man vist ogsaa give den skarpsindige Rötscher Ret, der saa fortræffeligt har udviklet, hvorledes dennes Væsen netop laae i at opfatte Virkeligheden idealt, at bringe en virkelig Personlighed paa Scenen, dog saaledes at denne netop saaes som Ideens Repræsentant, hvorfor man jo og hos Aristophanes finder de tre store comiske Paradigmer, Kleon, Euripides, Socrates, hvis Person comisk repræsenterer Tidens Stræben i dens tredobbelte Retning. Som derfor den indtil Detail nøiagtige Opfattelse af Virkeligheden udfyldte Afstanden mellem de Skuespilbesøgende og Theatret, saaledes fjernede dog igjen den ideale Opfattelse disse to Magter fra hinanden saavidt, som Kunst altid maa gjøre det. At Socrates nu virkelig i sit Liv har frembudt mange comiske Sider, at han, for engang for alle at sige Ordet ud, indtil en vis Grad har været en Sonderling, det lader sig ikke nægte; at deri allerede laae en Berettigelse for en comisk Digter, lader sig heller ei nægte; men at det vilde have været for lidet for en Aristophanes, lader sig heller ei fragaae. […]

\sectionheading{Anden Deel: Indledning  ·  SV¹ XIII, 317–321}

Hvad der egentlig skal udgjøre Gjenstanden for denne Deel af Undersøgelsen, er til en vis Grad allerede givet i det Foregaaende, forsaavidt som dette under Contemplationens Form lod en Side af Begrebet komme tilsyne. Jeg har derfor ikke saameget forudsat Begrebet i den tidligere Deel af Afhandlingen, som jeg har ladet det blive til, idet jeg søgte at orientere mig i Phænomenet. Jeg har derved fundet en ubekjendt Størrelse, et Standpunkt, der viste sig som det, der maa have været eiendommeligt for Socrates. Jeg har kaldt dette Standpunkt Ironi; men Navnet paa det er i den første Deel af Afhandlingen det mindre Vigtige: Hovedsagen er, at der intet Moment, intet Træk er blevet overseet, samt at alle Momenter og alle Træk have ordnet sig sammen i en Totalitet. Om dette Standpunkt er Ironi, det skal først nu afgjøres, idet jeg i Udviklingen af Begrebet ogsaa vil komme til det Moment, hvor Socrates maa passe, saasandt hans Standpunkt ellers har været Ironi. Men som jeg i den første Deel af Afhandlingen alene havde med Socrates at gjøre, saaledes vil det i Udviklingen af Begrebet vise sig, i hvilken Forstand han er et Moment i Begrebets Udvikling, med andre Ord, det vil vise sig, om Begrebet Ironi er absolut udtømt i ham, eller om der ikke er andre Fremtrædelsesformer deraf, til hvilke der ogsaa maa tages et Hensyn, førend man kan sige, at Begrebet tilstrækkeligt er opfattet. Som derfor i den første Deel af Afhandlingen Begrebet stedse svævede i Baggrunden, med en bestandig Trang til at tage Skikkelse i Phænomenet, saaledes vil i denne Deel af Afhandlingen Begrebets phænomenale Fremtrædelse, som en bestandig Mulighed til at tage Bolig iblandt os, følge Udviklingen. Disse to Momenter ere uadskillelige; thi hvis Begrebet ikke var i Phænomenet, eller rettere, Phænomenet først var forstaaeligt, først virkeligt i og med Begrebet, og hvis Phænomenet ikke var i Begrebet, eller rettere, Begrebet først forstaaeligt, først virkeligt i og med Phænomenet, saa var al Erkjendelse umulig, idet jeg i det ene Tilfælde vilde mangle Sandheden, i det andet Virkeligheden. Er nu Ironi en Subjectivitetens Bestemmelse, saa vil man strax see Nødvendigheden af to Fremtrædelser af dette Begreb; og Virkeligheden har ogsaa knyttet Navnet til dem. Den første er naturligviis den, hvor Subjectiviteten første Gang gjør sin Ret gjeldende i Verdenshistorien. Her have vi Socrates, det vil sige, herved ere vi anviste, hvor vi skal søge Begrebet i dets historiske Fremtrædelse. Men da Subjectiviteten havde viist sig i Verden, forsvandt den ikke atter sporløs, Verden sank ikke atter tilbage i den tidligere Udviklingsform, tvertimod, det Gamle forsvandt, og Alt blev nyt. Skal der nu kunne vise sig en ny Fremtrædelsesform af Ironien, da maa det være, idet Subjectiviteten i en endnu høiere Form gjør sig gjeldende. Det maa da være en Subjectivitetens anden Potents, en Subjectivitetens Subjectivitet, der svarer til Reflexionens Reflexion. Herved ere vi atter verdenshistorisk orienterede, vi ere nemlig henviste til den Udvikling, som det nyere Philosophi erholdt i Kant, og som fuldendtes i Fichte, og nærmere igjen til de Standpunkter, der efter Fichte søgte at gjøre Subjectiviteten i anden Potents gjeldende. At dette nu forholder sig rigtigt, det viser Virkeligheden ogsaa, thi her møde vi atter Ironien. Men da dette Standpunkt er en potenseret subjectiv Bevidsthed, saa følger deraf, at den bliver sig Ironien tydelig og bestemt bevidst, at den udtaler Ironien som sit Standpunkt. Dette skete nu ogsaa i Fr. Schlegel, der søgte at gjøre den gjeldende i Forhold til Virkeligheden, i Tieck, der søgte at gjøre den gjeldende i Poesien, i Solger, der blev sig dette æsthetisk og philosophisk bevidst. Endelig traf ogsaa Ironien her sin Mester i Hegel. Medens den første Form af Ironi ikke blev bekjæmpet, men beroliget derved at Subjectiviteten skete sin Ret, saa blev den anden Form af Ironi bekjæmpet og tilintetgjort; thi da den var uberettiget, kunde den kun skee sin Ret derved, at den blev ophævet.

Er man nu ved disse Betragtninger tilstrækkelig orienteret med Hensyn til dette Begrebs Historie, saa skal dermed dog ingenlunde være sagt, at en Opfattelse af dette Begreb, forsaavidt den søger Tilhold og Støtte i de tidligere Udviklinger, ikke er forbunden med Vanskelighed. Forsaavidt som man nemlig søger en fuldstændig og sammenhængende Udvikling af dette Begreb, vil man snart overbevise sig om, at det har en besynderlig eller rettere ingen Historie. I den Periode efter Fichte, hvor det især blev gjort gjeldende, finder man det atter og atter nævnet, atter og atter antydet, atter og atter forudsat. Søger man derimod en klar Udvikling, saa søger man forgjæves. Solger klager over at A. W. v. Schlegel i hans Vorlesungen über dramatische Kunst und Litteratur, hvor man snarest maa vente at finde den tilstrækkelige Oplysning, kun flygtigt antyder den et eneste Sted. Hegel klager over, at det er gaaet Solger paa samme Maade, og Tieck ikke bedre. Og da Alle klage, hvorfor skulde jeg ikke ogsaa klage? Jeg klager over, at det er gaaet Hegel omvendt. Paa det Punkt i alle hans Systemer, hvor man kunde vente at finde Ironien udviklet, finder man den omtalt, men uagtet, hvis man vilde lade det Alt aftrykke, man maatte tilstaae, at det var ikke Lidet, der var sagt om Ironi, saa er det dog i en anden Forstand ikke Meget; thi han siger omtrent det Samme paa alle Punkter. Hertil kommer, at han retter sit Angreb mod de enkelte ofte forskjellige Tanker, man har knyttet til dette Ord, hvoraf bliver en Følge, at da Sprogbrugen ikke er constant, hans Polemik ikke altid bliver aldeles klar. Imidlertid er det dog langtfra, at jeg med Rette kan klage over Hegel i samme Forstand, som Hegel klager over sine Forgjængere. Isærdeleshed findes der fortrinlige Bemærkninger i den Anmeldelse af Solgers efterladte Skrifter, der findes i det 16de Bind af hans samtlige Værker. Og er end ikke altid Fremstillingen og Skildringen af negative Standpunkter (thi ved disse gjelder det især om Skildring: loquere, ut videam te) saa udtømmende, saa indholdsriig, som man kunde ønske det, saa veed Hegel desto bedre at magte dem, og forsaavidt vil man i den Positivitet, han gjør gjeldende, have et indirecte Bidrag til Skildringen. Medens Schlegelerne og Tieck havde deres største Betydning ved den Polemik, hvorved de tilintetgjorde en foregaaende Udvikling, og medens netop af denne Grund deres Standpunkt blev noget adspredt, fordi det ikke var et Hovedslag, de vandt, men en Mangfoldighed af Træfninger, saa har Hegel derimod absolut Betydning ved i sin positive Total-Anskuelse at beseire den polemiske Knibskhed, der, som Dronning Brynhilds Jomfruelighed, behøvede en mere end almindelig Ægtemand, behøvede en Sigurd Svend for at betvinges. Hos J. Paul er der ogsaa ofte Tale om Ironi, og i hans Æsthetik findes der Et og Andet, men uden philosophisk eller ægte æsthetisk Myndighed. Han taler overhovedet, som Æsthetiker, mere ud af en riig æsthetisk Erfaring, end han egentlig begrunder sit æsthetiske Standpunkt. For ham er Ironi, Humor, Lune ligesom forskjellige Sprog, og hans Skildring indskrænker sig til, at han udtrykker den samme Tanke ironisk, humoristisk, i Lunets Sprog, omtrent som naar Fr. Baader stundom, efterat have givet en Fremstilling af enkelte mystiske Sætninger, nu oversætter denne paa mystisk.

Men da det Begreb Ironi saaledes ofte har faaet en forskjellig Betydning, gjelder det om, at man ikke med eller mod sit Vidende kommer til at bruge det aldeles vilkaarligt, gjelder det om, at man, idet man knytter sig til den almindelige Sprogbrug, faaer at see, at de forskjellige Betydninger, det i Tidens Løb har faaet, dog alle passe ind herunder.

\sectionheading{Anden Deel: Ironiens verdenshistoriske Gyldighed, Socrates' Ironi  ·  SV¹ XIII, 333–344}

Men vende vi tilbage til den i det Foregaaende givne almindelige Betegnelse af Ironi, som den uendelige absolute Negativitet, saa er heri tilstrækkeligen antydet, at Ironien nu ikke længer vender sig mod dette eller hiint enkelte Phænomen, mod et enkelt Tilværende, men at hele Tilværelsen er bleven fremmed for det ironiske Subject, og dette igjen fremmed for Tilværelsen, at idet Virkeligheden har tabt sin Gyldighed for det, det selv til en vis Grad er blevet uvirkeligt. Det Ord »Virkelighed« maa her imidlertid for det Første tages i Betydning af den historiske Virkelighed, det vil sige, den til en vis Tid, under visse Forhold givne Virkelighed. Dette Ord kan nemlig deels tages i metaphysisk Forstand, saaledes som, naar man behandler det metaphysiske Problem om Ideens Forhold til Virkelighed, hvor der da ikke bliver Spørgsmaal om denne eller hiin Virkelighed, men om Ideens Concretion, det er, dens Virkelighed; deels kan Ordet Virkelighed bruges om den historisk virkeliggjorte Idee. Denne sidstnævnte Virkelighed er til forskjellige Tider en anden. Hermed skal ingenlunde være meent, at den historiske Virkelighed ikke skulde have i sin hele Sum af Existents en evig Sammenhæng i sig selv, men for de ved Tid og Rum adskilte Slægter er den givne Virkelighed en forskjellig. Om nu end Verdens-Aanden i enhver Udvikling bestandig dog er i sig selv, saa er det ikke saaledes Tilfældet med Slægten til en vis Tid og de i samme til en vis Tid givne Individer. For disse frembyder sig en given Virkelighed, som det ikke staaer i deres Magt at vrage; thi Verdensudviklingen leder den, der selv vil med, river den med, der ikke vil. Men forsaavidt som Ideen er concret i sig, er det den nødvendigt, bestandigt at blive hvad den er, det er – blive concret. Men dette kan den dog kun blive igjennem Slægt og Individer.

Herved viser sig en Modsigelse, igjennem hvilken Verdens-Udviklingen foregaaer. Den til en vis Tid givne Virkelighed er den for Slægten og Individerne i Slægten gyldige, og dog maa, saafremt man ikke vil sige, at Udviklingen er forbi, denne Virkelighed fortrænges af en anden Virkelighed, og dette maa skee gjennem og ved Individerne og Slægten. For den med Reformationen samtidige Slægt var saaledes Katholicismen den givne Virkelighed, og dog var den tillige den Virkelighed, der ikke mere havde Gyldighed som saadan. Her colliderer altsaa en Virkelighed med en anden Virkelighed. Heri ligger det dybe Tragiske i Verdenshistorien. Et Individ kan paa eengang være verdenshistorisk berettiget og dog ubeføiet. Forsaavidt som han er det Sidste, maa han blive et Offer, forsaavidt som han er det Første, maa han seire, det vil sige, han maa seire derved, at han bliver et Offer. Man seer her, hvor conseqvent Verdens-Udviklingen er i sig; thi idet den sandere Virkelighed skal frem, agter den dog selv den forbigangne; det er ingen Revolution, men en Evolution; den forbigangne Virkelighed viser sig endnu berettiget derved, at den kræver et Offer, den nye Virkelighed derved, at den bringer et Offer. Men et Offer maa der til, fordi der virkelig skal et nyt Moment frem, fordi den nye Virkelighed ikke er en blot Conclusion af det Forbigangne, men indeholder et Mere i sig, ikke er et blot Correctiv til det Forbigangne, men tillige er en ny Begyndelse.

I ethvert saadant Vendepunkt i Historien er der to Bevægelser, man maa bemærke. Paa den ene Side skal det Nye frem, paa den anden Side skal det Gamle fortrænges. Forsaavidt det Nye skal frem, da møde vi her det prophetiske Individ, der øiner det Nye i det Fjerne, i dunkle og ubestemte Omrids. Det prophetiske Individ eier ikke det Tilkommende, han aner det blot. Han kan ikke gjøre det gjeldende, men ogsaa han er tabt for den Virkelighed, han tilhører. Hans Forhold til denne er imidlertid et fredeligt Forhold; thi den givne Virkelighed føler ingen Modsætning. Derpaa følger den egentlig tragiske Helt. Han kjæmper for det Nye, han stræber at tilintetgjøre det, der for ham er et Forsvindende, men hans Opgave er dog ikke saa meget at tilintetgjøre, som at gjøre det Nye gjeldende og derved indirecte tilintetgjøre det Forbigangne. Men paa den anden Side skal det Gamle fortrænges, det Gamle maa sees i sin hele Ufuldkommenhed. Her møde vi det ironiske Subject. For det ironiske Subject har den givne Virkelighed aldeles tabt sin Gyldighed, den er bleven ham en ufuldkommen Form, som overalt generer. Men paa den anden Side, det Nye eier han ikke. Han veed blot dette, at det Nærværende ikke svarer til Ideen. Han er den, der skal holde Dom. Ironikeren er vel i en vis Forstand prophetisk, thi han peger bestandig hen paa noget Tilkommende, men hvad det er, veed han ei. Han er prophetisk; men hans Stilling og Situation er det Modsatte af Prophetens. Propheten gaaer Haand i Haand med sin Samtid, og fra dette Standpunkt øiner han det som skal komme. Propheten er, som ovenfor bemærket, tabt for sin Samtid, men det er han dog egentlig kun, fordi han hensynker i sine Syner. Ironikeren derimod er traadt ud af Samtidens Rækker, har gjort Front imod denne. Det, der skal komme, er skjult for ham, ligger bag ved hans Ryg, men den Virkelighed, han har stillet sig fjendtligt ligeoverfor, er det, han skal tilintetgjøre, mod den vender hans fortærende Blik sig, paa hans Forhold til sin Samtid kan man anvende Skriftens Ord: see deres Fødder ere udenfor, der skal bære Dig bort. Ironikeren er ogsaa et Offer, som Verdens-Udviklingen kræver, ikke som om Ironikeren altid behøver i strængere Forstand at falde som et Offer, men Nidkjærhed i Verdensaandens Tjeneste fortærer ham.

Her have vi altsaa Ironien som den uendelige absolute Negativitet. Den er Negativitet, thi den negerer blot; den er uendelig, thi den negerer ikke dette eller hiint Phænomen; den er absolut, thi det, i Kraft af hvilket den negerer, er et Høiere, der dog ikke er. Ironien etablerer Intet; thi det, der skal etableres, ligger bag ved den. Den er et guddommeligt Vanvid, der raser som en Tamerlan og ikke lader Steen paa Steen tilbage. Her have vi altsaa Ironien. Til en vis Grad maa nu ethvert verdenshistorisk Vendepunkt ogsaa have denne Formation, og det kunde vist ikke være uden historisk Interesse at forfølge denne Formation gjennem Verdenshistorien. Uden imidlertid at indlade mig herpaa, skal jeg blot, som Exempler tagne fra den nærmest Reformationen liggende Tid, anføre Cardanus, Campanella, Bruno. Ogsaa Erasmus Rotterodamus var til en vis Grad Ironi. Betydningen af denne Formation troer jeg, at man hidtil ikke tilstrækkelig har paaagtet; og det er saameget desto besynderligere, som Hegel med en saa afgjort Forkjærlighed har behandlet det Negative. Men til det Negative i Systemet svarer Ironien i den historiske Virkelighed. I den historiske Virkelighed er det Negative til, hvilket det aldrig er i Systemet.

Ironien er en Subjectivitetens Bestemmelse. I Ironien er Subjectet negativ frit; thi den Virkelighed, der skal give det Indhold, er der ikke, det er frit fra den Bundethed, i hvilken den givne Virkelighed holder Subjectet, men det er negativ frit og som saadant svævende, fordi der Intet er, der holder det. Men netop denne Frihed, denne Svæven, giver Ironikeren en vis Enthusiasme, idet han ligesom beruser sig i Mulighedernes Uendelighed, idet han, forsaavidt han behøver en Trøst over alt det, der gaaer under, kan tage sin Tilflugt til Mulighedens uhyre Reservefond. Denne Enthusiasme hengiver han sig imidlertid ikke til, den begeistrer og nærer blot Tilintetgjørelsens Enthusiasme i ham. – Men da Ironikeren ikke har det Nye i sin Magt, saa kunde man spørge, hvorved han da tilintetgjør det Gamle, og dertil maa man svare: han tilintetgjør den givne Virkelighed ved selve den givne Virkelighed, men dog tillige erindre, at det nye Princip i ham er tilstede ϰατα δυναμιν, som Mulighed. Men idet han tilintetgjør Virkeligheden ved selve Virkeligheden, træder han i Verdens-Ironiens Tjeneste. Hegel bemærker i sin Gesch. der Philos. 2 Bind Pagina 62: Alle Dialectik läßt das gelten, was gelten soll, als ob es gelte, läßt die innere Zerstörung selbst sich daran entwickeln – allgemeine Ironie der Welt, og heri er Verdens-Ironien meget rigtigt opfattet. Netop fordi enhver enkelt historisk Virkelighed bestandig dog kun er Moment i Ideens Virkeliggjørelse, bærer den i sig selv Spiren til sin Undergang. Dette viser sig isærdeleshed ganske tydeligt ved Jødedommen, hvis Betydning som Gjennemgangsmoment især er paafaldende. Det var saaledes allerede en dyb Ironi over Verden, naar Loven, efterat have forkyndt Buddene, tilføiede Forjættelsen: dersom Du gjør disse, da skal Du blive salig, da det jo netop viste sig, at Menneskene ikke kunde fuldkomme Loven, og altsaa en Salighed, der blev knyttet til dette Vilkaar, blev meer end hypothetisk. Men at Jødedommen tilintetgjorde sig ved sig selv, det viser sig netop ved dens historiske Forhold til Christendommen. Naar vi, uden videre at gaae ind i en Undersøgelse om Betydningen af Christi Fremtræden, blot vil fastholde denne som et Vendepunkt i Verdenshistorien, saa kan man jo heller ikke der savne den ironiske Formation. Den er nu og givet ved Johannes den Døber. Han var ikke den, som skulde komme, han var ikke vidende om, hvad der skulde komme, og dog tilintetgjorde han Jødedommen. Han tilintetgjorde den altsaa ikke ved det Nye, men han tilintetgjorde den ved den selv. Han fordrede af Jødedommen, det Jødedommen vilde give – Retfærdighed; men dette var den ikke istand til at give, og derved gik den under. Han lod altsaa Jødedommen bestaae og udviklede paa samme Tid Spiren til dens Undergang i den. I Johannes den Døber træder imidlertid hans Personlighed aldeles i Skygge, i ham seer man ligesom Verdens-Ironien i dens objective Skikkelse, saa at han blot bliver et Redskab i dens Haand. Men for at den ironiske Formation skal være fuldkommen udviklet, fordres der, at Subjectet tillige bliver sig sin Ironi bevidst, føler sig, idet han fordømmer den givne Virkelighed, negativ fri, og nyder denne negative Frihed. For at dette skal kunne skee, maa Subjectiviteten være udviklet, eller rettere, idet Subjectiviteten gjør sig gjeldende, fremtræder Ironien. Subjectiviteten føler sig selv ligeoverfor den givne Virkelighed, føler sin Kraft, sin Gyldighed og Betydning. Men idet den føler dette, frelser den sig ligesom ud af den Relativitet, hvori den givne Virkelighed vil holde den. Forsaavidt som denne Ironi nu er verdenshistorisk berettiget, saa er Subjectivitetens Frigjørelse foretaget i Ideens Tjeneste, om end det ironiske Subject ikke er sig dette tydelig bevidst. Dette er det Geniale i den berettigede Ironi. Om den uberettigede Ironi gjelder det, at den, som vil frelse sin Sjæl, skal miste den. Men om Ironien er berettiget eller ikke, derom kan kun Historien dømme.

Men fordi Subjectet seer Virkeligheden ironisk, deraf følger ingenlunde, at han i at gjøre denne sin Opfattelse af Virkeligheden gjeldende forholder sig ironisk. Saaledes har der i den nyere Tid været Tale nok om Ironi og om den ironiske Opfattelse af Virkeligheden; men denne Opfattelse har sjeldnere gestaltet sig ironisk. Men jo mere dette skeer, desto sikkrere og uundgaaeligere er ogsaa Virkelighedens Undergang, desto mere Overvægt har det ironiske Subject over den Virkelighed, han vil tilintetgjøre, og desto mere fri er han ogsaa. Her foretager han nu i al Stilhed den samme Operation som Verdens-Ironien. Han lader det Bestaaende bestaae, men for ham har det ingen Gyldighed; imidlertid lader han dog, som om det havde Gyldighed for ham, og under denne Maske fører han det sin visse Undergang imøde. Forsaavidt som det ironiske Subject er verdenshistorisk berettiget, er her en Eenhed af det Geniale og den kunstneriske Besindighed.

Men er Ironien en Subjectivitetens Bestemmelse, saa maa den ogsaa vise sig, første Gang Subjectiviteten fremtræder i Verdenshistorien. Ironien er nemlig Subjectivitetens første og abstracteste Bestemmelse. Dette peger hen paa det historiske Vendepunkt, hvor Subjectiviteten første Gang traadte frem, og herved ere vi komne til Socrates.

Hvorledes det nu stod sig med Socrates' Ironi, det er tilstrækkeligen bleven oplyst i den foregaaende Deel af denne Undersøgelse. Hele den givne Virkelighed havde for ham tabt sin Gyldighed, hele Substantialitetens Virkelighed var han bleven fremmed for. Det er den ene Side af Ironien; men paa den anden Side, idet han tilintetgjorde Græciteten, brugte han Ironien; hans Adfærd mod den var bestandig ironisk; han var uvidende og vidste Intet, men søgte bestandig Oplysning hos Andre, men idet han saaledes lod det Bestaaende bestaae, gik det under. Denne Taktik vedligeholdt han til det Yderste, hvilket især viste sig, da han var bleven anklaget. Men Nidkjærhed i denne Tjeneste fortærede ham, og tilsidst greb Ironien ham, det svimler for ham, Alt taber sin Realitet. Denne Anskuelse af Socrates, og af Betydningen af hans Standpunkt i Verdenshistorien, forekommer mig saa naturligt at afrunde sig i sig selv, at den vil, saa haaber jeg, skaffe sig Indgang hos en og anden Læser. Da imidlertid Hegel erklærer sig imod at opfatte Socrates' Standpunkt som Ironi, bliver det nødvendigt at tage et Hensyn til de Indvendinger herimod, der findes hist og her i hans Skrifter.

Forinden vil jeg imidlertid, saa godt det staaer til mig, stræbe at belyse en Svaghed, hvoraf hele Hegels Opfattelse af det Begreb Ironi synes at lide. Hegel omtaler altid Ironien paa en meget afvisende Maade, Ironien er en Vederstyggelighed i hans Øine. Samtidig med Hegels Fremtræden falder Schlegels meest glimrende Periode. Men som Schlegelernes Ironi havde i Æsthetiken holdt Dom over en om sig gribende Sentimentalitet, saaledes var Hegel den, der skulde berigtige den Misviisning, der laa i Ironien. Det er overhovedet en af Hegels store Fortjenester, at han har standset, eller idetmindste har villet standse Speculationens forlorne Sønner paa deres Fortabelses Vei. Men hertil brugte han ikke altid de lempeligste Midler, og hans Stemme var, naar han kaldte ad dem, ikke altid en mild og faderlig Røst, men havde ofte noget Barskt og Skolemesteragtigt. Ironiens Tilhængere vare de, der voldte ham meest Uleilighed, og han opgav derfor snart Haabet om deres Frelse, og behandler dem nu som uforbederlige og forstokkede Syndere. Ved enhver Leilighed kommer Hegel til at tale om disse Ironikere, og altid blive de omtalte paa den meest afvisende Maade, ja Hegel seer med en uhyre Haan og Fornemhed ned paa disse, som han ofte nok kalder »fornemme Folk.« Men den Omstændighed, at Hegel saaledes forsaae sig paa den Form af Ironi, som laae ham nærmest, har naturligviis skadet hans Opfattelse af Begrebet. Udvikling faaer man ofte ikke – derimod faaer Schlegel altid Skjend. Hermed skal nu ingenlunde være sagt, at Hegel ikke har Ret mod Schlegelerne, og at den Schlegel-Schlegelske Ironi ikke var en saare betænkelig Afvei; heller ikke skal hermed være sagt Andet, end at Hegel vist har stiftet stort Gavn ved den Alvor, hvormed han træder op mod enhver Isolation, en Alvor, der gjør, at man kan læse mangen Udvikling med megen Opbyggelse og Styrkelse; men derimod skal dermed være sagt, at Hegel, ved eensidig at vende sig mod den efterfichtiske Ironi, har overseet Ironiens Sandhed, og idet han identificerede al Ironi med denne, har gjort Ironien Uret. Saasnart Hegel nævner det Ord Ironi, kommer han strax til at tænke paa Schlegel og Tieck, og hans Stiil bærer da øieblikkeligen Præg af en vis Forbittrelse. Hvori det Forkerte og Ubeføiede i Schlegels Ironi laae, skal paa sit Sted blive oplyst, saavelsom Hegels Fortjeneste i denne Henseende. Vi vende tilbage til hans Betragtningsmaade af Socrates' Ironi.

I det Foregaaende have vi gjort opmærksom paa, at Hegel i at fremstille Socrates' Methode især udhæver to Former, hans Ironi og hans Gjordemoderkunst. Hans Fremstilling heraf findes i Geschichte der Philos. 2 Bind Pagina 59-67. Udviklingen af den sokratiske Ironi er imidlertid meget kort, hvorimod Hegel benytter Leiligheden til at ivre mod Ironi, som almindeligt Princip, og tilføier Pagina 62: Friedrich von Schlegel ist es, der diese Gedanken zuerst aufgebracht, Ast hat es nachgesprochen; og derpaa følge de Alvorsord, som Hegel i den Anledning pleier at fremsige. Socrates lader, som han var uvidende, og under Skin af at lade sig belære, belærer han Andre. Pagina 60: Dieses ist dann die Seite der berühmten socratischen Ironie. Sie hat bei ihm die subjective Gestalt der Dialectik, sie ist Benehmungsweise im Umgang; die Dialectik ist Gründe der Sache, die Ironie ist besondere Benehmungsweise von Person zu Person. Men da der lidt iforveien bemærkes, at Socrates bruger den samme Ironi, wenn er die Manier der Sophisten zu Schanden machen will, saa møder her strax en Vanskelighed; thi i det ene Tilfælde vil han jo belære, i det andet blot beskjæmme. Hegel gjør nu opmærksom paa, at denne socratiske Ironi synes at indeholde noget Usandt, men viser derpaa Rigtigheden af hans Adfærd. Endelig viser han den egentlige Betydning af den socratiske Ironi, det Store ved den. Dette er, at den leder til at gjøre de abstracte Forestillinger concrete og udviklede. Han tilføier derpaa Pagina 62: Wenn ich sage, ich weiß, was Vernunft, was Glaube ist, so sind dieß nur ganz abstracte Vorstellungen; daß sie nun concret werden, dazu gehört, daß sie explicirt werden, daß vorausgesetzt werde, es sey nicht bekannt, was es eigentlich sey. Diese Explication solcher Vorstellungen bewirkt nun Socrates; und dieß ist das Wahrhafte der socratischen Ironie. Men hermed er nu Alt forvirret, Fremstillingen af den socratiske Ironi taber ganske sin historiske Tyngde, og den her citerede Passus er saa moderne, at den neppe erindrer om Socrates. Det var nemlig ingenlunde Socrates' Sag, at gjøre det Abstracte concret, og de anførte Exempler ere vistnok temmelig uheldig valgte; thi jeg troer ikke, at Hegel skal kunne anføre Analogier hertil, medmindre han vil tage hele Plato og skyde sig ind under, at Socrates' Navn endnu bestandig bruges i Plato, hvorved han da vil komme i Strid baade med sig selv og med Alle. Socrates' Sag var det ikke at gjøre det Abstracte concret, men igjennem det umiddelbar Concrete at lade det Abstracte komme tilsyne. Det vil derfor mod disse hegelske Betragtninger være tilstrækkeligt at erindre, deels om den dobbelte Art Ironi, vi fandt i Plato (thi det er aabenbart den Ironi, vi have kaldt den platoniske, Hegel har meent, hvormed han og Pagina 64 identificerer den socratiske Ironi), deels om Bevægelsesloven i hele Socrates' Liv, at den ikke var fra det Abstracte at komme til det Concrete, men fra det Concrete at komme til det Abstracte, og bestandig at komme dertil. Naar derfor den hele Undersøgelse om den socratiske Ironi hos Hegel ende med, at den socratiske Ironi bliver identisk med den platoniske, og saavel den socratiske som den platoniske bliver mehr Manier der Conversation, die gesellige Heiterkeit, als daß jene reine Negation, jenes negative Verhalten darunter verstanden wäre (Pagina 64), saa har ogsaa denne Bemærkning fundet sin Besvarelse i det Foregaaende. – Stort bedre gaaer det heller ikke Hegel i at fremstille Socrates' Gjordemoderkunst. Her udvikler han Betydningen af, at Socrates spurgte, og denne Udvikling er baade skjøn og sand; men den Forskjel, vi i vort Foregaaende har gjort gjeldende mellem at spørge for at faae Svar og at spørge for at beskjæmme, oversees her. Det Exempel, han til Slutning vælger fra det Begreb at vorde, er atter aldeles usocratisk, medmindre han vil i Parmenides finde en socratisk Udvikling. – Forsaavidt han endelig taler om Socrates' tragiske Ironi, saa maa man erindre, at det ikke er Socrates' Ironi, men det er Verdens Ironi med Socrates. Det kan derfor Intet oplyse med Hensyn til Spørgsmaalet om den socratiske Ironi.

I Anmeldelsen af Solgers Skrifter gjør Hegel Pagina 488 atter opmærksom paa Forskjellen mellem den schlegelske Ironi og den socratiske. At der er Forskjel, det har vi alt indrømmet og skal yderligere paavise paa sit Sted, men deraf følger ingenlunde, at Socrates' Standpunkt ikke var Ironi. Han bebreider Fr. Schlegel, at han, med Uforstand paa det Speculative og Tilsidesættelse af dette, har revet den fichtiske Sætning om Jegets constitutive Gyldighed ud af dens metaphysiske Sammenhæng, revet den ud af Tænkningens Gebeet og directe anvendt den paa Virkeligheden zum Verneinen der Lebendigkeit der Vernunft und Wahrheit, und zur Herabsetzung derselben zum Schein im Subject und zum Scheinen für Andere. Han gjør nu opmærksom paa, at man, for at betegne denne Sandhedens Forvanskning til Skin, har tilladt sig at forvanske Navnet paa den uskyldige socratiske Ironi. Naar man nemlig fortrinsviis har sat Ligheden deri, at Socrates altid indlod sig paa sin Undersøgelse med den Forsikkring, at han Intet vidste, for at beskjæmme Sophisterne, saa er Resultatet af denne Adfærd altid noget Negativt, der bliver uden videnskabeligt Resultat. Forsaavidt bliver altsaa den Forsikkring af Socrates, at han Intet vidste, given for ramme Alvor, og er altsaa ikke ironisk. Jeg skal her ikke videre indlade mig paa den Vanskelighed, der fremgaaer af, at Hegel her viser, at Socrates' Underviisning endte uden Resultat, naar man hermed sammenholder, hvad tidligere efter Hegel blev udviklet, at Socrates ved sin ironiske Underviisning gjorde det Abstracte concret, men derimod skal jeg lidt udførligere undersøge, hvorvidt det var Socrates Alvor med sin Uvidenhed.

At Socrates, idet han sagde, han var uvidende, dog var vidende, da han var vidende om sin Uvidenhed, og at paa den anden Side denne Viden dog ikke var en Viden om Noget, det vil sige, ikke havde noget positivt Indhold, og at forsaavidt hans Uvidenhed var ironisk, det har viist sig i det Foregaaende; og da Hegel, som det synes mig, forgjeves har søgt at vindicere ham et positivt Indhold, saa troer jeg, at Læseren deri maa give mig Ret. Havde hans Viden været en Viden om Noget, da havde hans Uvidenhed blot været en Conversationsform. Nu er derimod hans Ironi fuldendt i sig selv. Forsaavidt er det altsaa paa eengang Alvor med hans Uvidenhed og dog atter ikke Alvor, og paa denne Spidse maa man fastholde Socrates. Det at man veed, at man er uvidende, er Begyndelsen til at blive vidende, men naar man ikke veed Mere, er det en blot Begyndelse. Denne Viden er det, der holder Socrates ironisk oppe. Naar dernæst Hegel, ved at gjøre opmærksom paa, at det var Socrates' Alvor med hans Uvidenhed, troer at vise, at hans Uvidenhed ikke var Ironi, saa synes det atter her, at Hegel ikke er constant. Naar nemlig Ironien skal opstille en høieste Sætning, saa gaaer det den som ethvert negativt Standpunkt, saa udsiger den noget Positivt, det er den Alvor med det, den siger. For Ironien er Intet et Bestaaende, den skalter og valter ad libitum med Alt; men naar den vil udsige dette, siger den noget Positivt, hvormed dens Souverainitet forsaavidt er til Ende. Naar derfor Schlegel eller Solger siger: Virkeligheden er blot Skin, blot Tilsyneladelse, blot Forfængelighed, et Intet, saa mener han det aabenbart alvorligt hermed, og dog antager jo Hegel, at dette er Ironi. Vanskeligheden, som her møder, er egentlig den, at Ironien i strængere Forstand aldrig kan opstille en Sætning, fordi Ironien er en Bestemmelse af det forsigværende Subject, der i idelig Agilitet Intet lader bestaae, og paa Grund af denne Agilitet ikke kan samle sig i den Total-Anskuelse, at den Intet lader bestaae. Schlegels og Solgers Bevidsthed om, at Endeligheden er et Intet, er aabenbare ligesaa alvorligt meent som Socrates' Uvidenhed. I sidste Instants maa Ironikeren altid ponere Noget, men det, han saaledes ponerer, er Intet. Men nu er det umuligt, at man kan tage Intet alvorligt, uden enten at komme til Noget (dette skeer, naar man tager det speculativ alvorligt), eller at fortvivle (naar man tager det personligt alvorligt). Ingen af Delene gjør imidlertid Ironikeren, og forsaavidt kan man ogsaa sige, at det er ham ikke Alvor dermed. Ironien er det uendelig lette Spil med Intet, der ikke forfærdes ved det, men endnu engang stikker Hovedet i Veiret. Naar man nu ikke tager Intet speculativt eller personligt alvorligt, saa tager man det aabenbart letsindigt, og forsaavidt tager man det ikke alvorligt. Dersom Hegel mener, at det ikke var Schlegels Alvor med, at Tilværelsen var et Intet uden Realitet, saa maa der jo have været Noget, der har havt Gyldighed for ham, men saa var hans Ironi jo blot Form. Man kan derfor sige om Ironien, at det er den Alvor med Intet, forsaavidtsom det ikke er den Alvor med Noget. Den opfatter bestandig Intet i Modsætning til Noget, og, for at frigjøre sig for Alvor med Noget, griber den Intet. Men Intet bliver det den heller ikke Alvor med, uden forsaavidt som det ikke er den Alvor med Noget. Saaledes forholder det sig ogsaa med Socrates' Uvidenhed, hans Uvidenhed er det Intet, hvorved han tilintetgjør enhver Viden. Det kan man bedst see af hans Opfattelse af Døden. Han er uvidende om, hvad Døden er, og om hvad der er efter Døden, om der er Noget eller slet Intet, han er altsaa uvidende; men denne Uvidenhed tager han sig ikke videre nær, tvertimod, han føler sig ret egentlig fri i denne Uvidenhed, det er ham altsaa ikke Alvor med denne Uvidenhed, og dog er det jo hans ramme Alvor, at han er uvidende. – Jeg troer derfor, at man vil give mig Ret i, at der ved disse hegelske Betragtninger ikke er opstillet Noget, der forbyder at antage, at Socrates' Standpunkt var Ironi.

Sammenfatte vi nu, hvad der i den første Deel af denne Afhandling blev udhævet som det Eiendommelige ved Socrates' Standpunkt, at hele det substantielle Liv i Græciteten havde tabt sin Gyldighed for ham, det vil altsaa sige, at den bestaaende Virkelighed var for ham uvirkelig, og det ikke i en eller anden enkelt Retning, men i sin hele Totalitet som saadan; at han i Forhold til denne ugyldige Virkelighed lod det Bestaaende paa Skrømt bestaae og derved bragte det til at gaae under; at han under Alt dette stedse blev lettere og lettere, stedse mere og mere negativ fri; saa see vi jo, at dette Socrates' Standpunkt, ifølge den nu givne Udvikling, var, som den uendelige absolute Negativitet, Ironi. Imidlertid var det ikke Virkeligheden overhovedet, han negerede, men det var den til en vis Tid givne Virkelighed, Substantialitetens, saaledes som den var i Grækenland, og det, hans Ironi fordrede, var Subjectivitetens, Idealitetens. Herom har nu Historien dømt, at Socrates var verdenshistorisk berettiget. Han blev et Offer. Dette er nu vistnok en tragisk Skjebne, men dog er Socrates' Død egentlig ikke tragisk; og den græske Stat kommer igrunden bagefter med sin Dødsdom, og har paa den anden Side heller ei stor Opbyggelse af Dødsstraffens Execution, thi Døden har ingen Realitet for Socrates. For den tragiske Helt har Døden Gyldighed, for ham er Døden i Sandhed den sidste Strid og den sidste Lidelse. Herved kan derfor den Samtid, han vilde tilintetgjøre, faae Hevnens Vrede tilfredsstillet. Men denne Tilfredsstillelse kunde aabenbart den græske Stat ikke faae ved Socrates' Død; thi Socrates havde ved sin Uvidenhed forhindret al betydningsfuldere Communication med Tanken om Døden. Vel frygter den tragiske Helt ikke Døden, men han veed den dog som en Smerte, som en tung og besværlig Gang, og forsaavidt har det Gyldighed, naar han dømmes fra Livet, men Socrates veed slet Intet, og forsaavidt er det Ironi over Staten, at den dømmer ham fra Livet, og derved troer at have taget Straf over ham.

\workheading{Enten – Eller. Et Livs-Fragment, Første Deel}{1843}

\worknote{Antologiens ''Either/Or … I'' (A.'s Papirer). Hong-udvalget bringer et udvalg af Diapsalmata (38 af aforismerne), den intetsigende indledning til ''Det Musikalsk-Erotiske'', hele ''Vexel-Driften'' samt seks stykker af ''Forførerens Dagbog''. Udeladelser er markeret med […].}

\sectionheading{Diapsalmata (udvalg)  ·  SV¹ I, 3–27}

Hvad er en Digter? Et ulykkeligt Menneske, der gjemmer dybe Qvaler i sit Hjerte, men hvis Læber ere dannede saaledes, at idet Sukket og Skriget strømme ud over dem, lyde de som en skjøn Musik. Det gaaer ham som de Ulykkelige, der i Phalaris's Oxe langsomt pintes ved en sagte Ild, deres Skrig kunde ikke naae hen til Tyrannens Øre for at forfærde ham, for ham løde de som en sød Musik. Og Menneskene flokkes om Digteren og sige til ham: syng snart igjen; det vil sige, gid nye Lidelser maae martre Din Sjæl, og gid Læberne maae vedblive at være dannede som forhen; thi Skriget vilde blot ængste os, men Musikken den er liflig. Og Recensenterne træde til, de sige: det er rigtigt, saaledes skal det være efter Æsthetikens Regler. Nu det forstaaer sig, en Recensent ligner ogsaa en Digter paa et Haar, kun har han ikke Qvalerne i Hjertet, ikke Musikken paa Læberne. See derfor vil jeg hellere være Svinehyrde paa Amagerbro og være forstaaet af Svinene, end være Digter og være misforstaaet af Menneskene.

Menneskene ere dog urimelige. De bruge aldrig de Friheder, de har, men fordre dem, de ikke har; de har Tænkefrihed, de fordre Yttringsfrihed.

Jeg gider slet ikke. Jeg gider ikke ride, det er for stærk en Bevægelse; jeg gider ikke gaae, det er for anstrængende; jeg gider ikke lægge mig ned; thi enten skulde jeg blive liggende, og det gider jeg ikke, eller jeg skulde reise mig op igjen, og det gider jeg heller ikke. Summa Summarum jeg gider slet ikke.

Probat Raad for Forfattere

Man nedskriver sine egne Betragtninger skjødesløst, man lader dem trykke, i de forskjellige Correcturer vil man da efterhaanden faae en Mængde gode Indfald. Fatter derfor Mod I, som endnu ikke har dristet Eder til at lade Noget trykke, ogsaa Trykfeil ere ikke at foragte, og at blive vittig ved Hjælp af Trykfeil maa ansees for en lovlig Maade at blive det paa.

Alderdommen realiserer Ungdommens Drømme: det seer man paa Swift, han byggede i sin Ungdom en Daarekiste, i sin Alderdom gik han selv i den.

Jeg siger om min Sorg hvad Engelskmanden siger om sit Huus: min Sorg is my castle. Mange Mennesker ansee det at have Sorg for at være en af Livets Bequemmeligheder.

Jeg er tilmode som en Brik i Schakspillet maa være det, naar Modspilleren siger om den: den Brik kan ikke røres.

Deri yttrer sig blandt Andet Folke-Litteraturens uhyre poetiske Kraft, at den har Styrke til at begjære. I Sammenligning med denne er vor Tids Begjær paa engang syndig og kjedsommelig, fordi den begjærer hvad Næstens er. Hiin er sig meget godt bevidst, at Næsten eier ligesaa lidet det, den søger, som den selv gjør det. Og skal den begjære syndigt, saa er den saa himmelskrigende, at den maa ryste Mennesket. Den lader sig Intet afprutte af en nøktern Forstands kolde Sandsynlighedsberegninger. Endnu skrider Don Juan over Scenen med sine 1003 Elskerinder. Ingen vover at smile af Ærbødighed for Traditionens Ærværdighed. Havde en Digter vovet det i vor Tid, da var han bleven udleet.

Mod har jeg til at tvivle, jeg troer om Alt; jeg har Mod til at kæmpe, jeg troer med Alt; men jeg har ikke Mod til at erkjende Noget; ikke Mod til at besidde, til at eie Noget. De Fleste klage over, at Verden er saa prosaisk, at det gaaer ikke til i Livet som i Romanen, hvor Leiligheden altid er saa gunstig; jeg klager over, at det ikke er i Livet, som i Romanen, hvor man har haardhjertede Fædre, og Nisser og Trolde at bekæmpe, fortryllede Prindsesser at befrie. Hvad er alle saadanne Fiender tilsammentagne imod de blege, blodløse, seiglivede, natlige Skikkelser, hvilke jeg kæmper med og som jeg selv giver Liv og Tilvær.

Hvor er min Sjæl og min Tanke saa ufrugtbar, og dog idelig piint af indholdsløse vellystige og qvalfulde Veer! Skal da Aandens Tungebaand aldrig løsnes paa mig, skal jeg altid lalle? Det, jeg behøver, er en Stemme gjennemtrængende som et Lynceus-Blik, forfærdende som Giganternes Suk, udholdende som en Naturlyd, spottende som et riimfrostigt Vindpust, ondskabsfuldt som Echos hjerteløse Haanen, af Omfang fra den dybeste Bas til de mest smeltende Brysttoner, modeleret fra en hellig-sagte Hvidsken til Raseriets Energi. Det er det jeg behøver for at faae Luft, for at faae udtalt, hvad der ligger mig paa Sinde, for at faae baade Vredens og Sympathiens Indvolde rystede. – Men min Stemme er kun hæs som et Maageskrig, eller hendøende som Velsignelsen paa den Stummes Læber.

Hvad skal der komme? Hvad skal Fremtiden bringe? Jeg veed det ikke, jeg ahner Intet. Naar en Edderkop fra et fast Punkt styrter sig ned i sine Conseqventser, da seer den bestandig et tomt Rum foran sig, hvori den ikke kan finde Fodfæste, hvormeget den end sprætter. Saaledes gaaer det mig; foran bestandig et tomt Rum, hvad der driver mig frem er en Consequents, der ligger bag mig. Dette Liv er bagvendt og rædsomt, ikke til at udholde.

Af alle latterlige Ting forekommer det mig at være det allerlatterligste at have travlt i Verden, at være en Mand, der er rask til sin Mad og rask til sin Gjerning. Naar jeg derfor seer en Flue i det afgjørende Øieblik sætte sig paa en saadan Forretningsmands Næse, eller han bliver overstænket af en Vogn, der i endnu større Hast kjører ham forbi, eller Knippelsbro gaaer op, eller der falder en Tagsteen ned og slaaer ham ihjel, da leer jeg af Hjertens Grund. Og hvo kunde vel bare sig for at lee? Hvad udrette de vel, disse travle Hastværkere? Gaaer det dem ikke som det gik hiin Kone, der i Befippelse over, at der var Ildløs i Huset, reddede Ildtangen? Hvad Mere redde de vel ud af Livets store Ildebrand?

Ingen vender tilbage fra de Døde, Ingen er gaaet uden grædende ind i Verden; Ingen spørger Een, naar man vil ind, Ingen naar man vil ud.

Troldmanden Virgilius lod sig selv hugge i Stykker og putte i en Gryde for at koges i 8 Dage, og ved denne Proces blive forynget. Han fik en Anden til at passe paa, at ingen Uvedkommende kigede i Gryden. Paapasseren kunde imidlertid ikke modstaae Fristelsen; det var for tidligt, Virgilius forsvandt som et lille Barn med et Skrig. Jeg har nok ogsaa kiget for tidlig i Gryden, i Livets og den historiske Udviklings Gryde, og driver det nok aldrig til mere end at blive et Barn.

Lad Andre klage over, at Tiden er ond; jeg klager over, at den er ussel; thi den er uden Lidenskab. Menneskenes Tanker ere tynde og skrøbelige som Kniplinger, de selv ynkværdige som Kniplings-Piger. Deres Hjertes Tanker ere for usle til at være syndige. For en Orm vilde det maaskee kunne ansees for Synd at nære saadanne Tanker, ikke for et Menneske, der er skabt i Guds Billede. Deres Lyster ere adstadige og dorske, deres Lidenskaber søvnige; de gjøre deres Pligter, disse Kræmmersjæle; men tillade sig dog ligesom Jøderne at beklippe Mønten en lille Smule; de mene, at om vor Herre holder nok saa ordentlig Bog, kan man nok slippe godt fra at narre ham lidt. Pfui over dem! Derfor vender min Sjæl altid tilbage til det gamle Testamente og til Shakespeare. Der føler man dog, at det er Mennesker, der taler; der hader man, der elsker man, myrder sin Fjende, forbander hans Afkom gjennem alle Slægter, der synder man.

Mit Livs-Resultat bliver slet Intet, en Stemning, en enkelt Farve. Mit Resultat faaer en Lighed med hiin Kunstners Maleri, der skulde male Jødernes Overgang over det røde Hav, og til den Ende malede hele Væggen rød, idet han forklarede, at Jøderne vare gaaede over, Ægypterne vare druknede.

Den menneskelige Værdighed anerkjendes dog endnu i Naturen; thi naar man vil holde Fuglene borte fra Træerne, saa sætter man Noget op, der skal ligne et Menneske, og endogsaa den fjerne Lighed med et Menneske, som en Fugleskræmsel har, er nok til at indgyde Respekt.

De fleste Mennesker haste saa stærkt efter Nydelsen, at de haste den forbi. Det gaaer dem, som det gik hiin Dverg, der bevogtede en bortført Prindsesse i sit Slot. En Dag tog han sig en Middagssøvn. Da han en Time efter vaagnede op, var hun borte. Hurtigt trækker han sine Syvmiils-Støvler paa; med eet Skridt er han hende langt forbi.

Hvor er Livet tomt og betydningsløst. – Man begraver et Menneske; man følger ham til Jorden, man kaster tre Spader Jord paa ham; man kjører ud i Kareth, man kjører hjem i Kareth; man trøster sig med, at der ligger et langt Liv for Een. Hvor langt er vel 7 × 10 Aar? Hvorfor gjør man det ikke af paa eengang, hvorfor bliver man ikke derude, og gaaer ned med i Graven, og trækker Lod om, hvem den Ulykke skal times, at være den sidst Levende, der kaster de sidste 3 Spader Jord paa den sidste Døde?

Disse to bekjendte Violinstrøg! Disse to bekjendte Violinstrøg her i dette Øieblik, midt paa Gaden. Har jeg tabt Forstanden, er det mit Øre, der af Kjærlighed til Mozarts Musik har ophørt at høre, er det en Belønning af Guderne at skjænke mig Ulykkelige, der sidder som en Betler ved Templets Dør, et Øre, der selv foredrager, hvad det selv hører. Kun disse to Violinstrøg; thi nu hører jeg Intet mere. Som de i hiin udødelige Ouverture bryde ud af de dybe Choral-Toner, saaledes vikle de sig her ud af Støien og Larmen paa Gaden, med en Aabenbarelses hele Overraskelse. – Det maa dog være her i Nærheden; thi nu hører jeg de lette Dandsetoner. – Det er altsaa Eder, ulykkelige Kunstnerpar, jeg skylder denne Glæde. – Den ene af dem var vel en 17 Aar gammel, iført en grøn Kalmuks-Frakke, med store Beenknapper. Frakken var meget for stor til ham. Han holdt Violinen tæt op under Hagen; Kaskjetten var trykket ned i Øinene; hans Haand var skjult af en Handske uden Fingre, Fingrene røde og blaae af Kulde. Den Anden var ældre, havde en Chenille paa. Begge vare de blinde. En lille Pige, som formodentlig ledede dem, stod foran, puttede sine Hænder ind under Halstørklædet. Vi samledes efterhaanden nogle Beundrere af disse Toner, et Postbud med sin Brev-Pakke, en lille Dreng, en Tjenestepige, et Par Sjouere. De herskabelige Vogne rullede larmende forbi, Arbeidsvognene overdøvede disse Toner, som glimtviis dukkede frem af dem. Ulykkelige Kunstner-Par veed I, at disse Toner gjemme alle Verdens Herligheder i sig. – Var det ikke som et Stevnemøde. –

Det hendte paa et Theater, at der gik Ild i Coulisserne. Bajads kom for at underrette Publicum derom. Man troede, det var en Vittighed og applauderede; han gjentog det; man jublede endnu mere. Saaledes tænker jeg, at Verden vil gaae til Grunde under almindelig Jubel af vittige Hoveder, der troer, at det er en Witz.

Hvad er overhovedet Betydningen af dette Liv? Deler man Menneskene i to store Classer, saa kan man sige, den ene arbeider for at leve, den anden har ikke dette behov. Men det at arbeide for at leve kan jo ikke være Livets Betydning, da det jo er en Modsigelse, at det, bestandig at tilveiebringe Betingelserne, skal være Svaret paa Spørgsmaalet om dets Betydning, der ved Hjælp heraf skal betinges. De Øvriges Liv har i Almindelighed heller ingen Betydning uden den at fortære Betingelserne. Vil man sige, at Livets Betydning er at døe, saa synes dette atter en Modsigelse.

Den egentlige Nydelse ligger ikke i hvad man nyder; men i Forestillingen. Hvis jeg i min Tjeneste havde en underdanig Aand, der, naar jeg forlangte et Glas Vand, vilde bringe mig al Verdens kostbareste Vine lifligt blandede i en Pokal, da vilde jeg give ham hans Afsked, indtil han lærte, at Nydelsen ikke ligger i hvad jeg nyder, men i at faae min Villie.

Det, Philosopherne tale om Virkeligheden, er ofte lige saa skuffende, som naar man hos en Marchandiser læser paa et Skildt: her rulles. Vilde man komme med sit Tøi for at faae det rullet, saa var man narret; thi Skildtet er blot tilsalgs.

For mig er intet farligere end at erindre. Har jeg først erindret et Livsforhold, saa er Forholdet selv ophørt. Man siger, at Adskillelse hjælper til at opfriske Kjærlighed. Det er ganske sandt, men den opfrisker den paa en reen poetisk Maade. At leve i Erindringen er det fuldendteste Liv, der lader sig tænke, Erindringen mætter rigeligere end al Virkelighed, og den har en Tryghed, som ingen Virkelighed eier. Et erindret Livsforhold er allerede gaaet ind i Evigheden og har ingen timelig Interesse mere.

Der hører dog en stor Naivetet til at troe, at det skal hjælpe, at raabe og skrige i Verden, som om derved Ens Skjæbne blev forandret. Man tage den som den bydes, og afholde sig fra alle Vidtløftigheder. Naar jeg i min Ungdom kom paa en Restauration, da sagde jeg ogsaa til Opvarteren: et godt Stykke, et meget godt Stykke, af Ryggen, ikke for fedt. Opvarteren hørte maaskee neppe mit Raab, end mindre at han skulde agte derpaa, end mindre, at min Røst skulde kunne trænge ud i Kjøkkenet, bevæge Forskæreren, og selv om alt dette skete, saa var der maaskee ikke et godt Stykke paa den hele Steg. Nu raaber jeg aldrig mere.

Som det efter Sagnet gik Parmeniskus, der i den trophoniske Hule tabte Evnen til at lee, men fik den igjen paa Delos ved Synet af en uformelig Klods, der blev fremstillet som Gudinden Letos Billede, saaledes er det gaaet mig. Da jeg var meget ung, da glemte jeg i den trophoniske Hule at lee, da jeg blev ældre, da jeg slog Øiet op og betragtede Virkeligheden, da kom jeg til at lee og har siden den Tid ikke ophørt dermed. Jeg saae, at det var Livets Betydning at faae et Levebrød, dets Maal at blive Justitsraad, at det var Elskovens rige Lyst at faae en velhavende Pige, at det var Venskabets Salighed at hjælpe hinanden i Pengeforlegenheder; at det var Viisdommen, hvad de Fleste antoge derfor; at det var Begeistring at holde en Tale; at det var Mod at vove at blive mulkteret paa 10 Rigsbankdaler; at det var Hjertelighed at sige Velbekomme efter et Middagsmaaltid; at det var Gudsfrygt een Gang om Aaret at gaae til Alters. Det saae jeg og jeg loe.

Hvad er det, der binder mig? Hvoraf var den Lænke dannet, hvormed Fenris-Ulven blev bunden? Den var forfærdiget af den Larm, Kattens Been gjør, naar den gaaer paa Jorden, af Qvinders Skæg, af Klippernes Rødder, af Bjørnens Græs, af Fiskenes Aande og Fuglenes Spyt. Saaledes er ogsaa jeg bunden i en Lænke, der er dannet af mørke Indbildninger, af ængstende Drømme, af urolige Tanker, af bange Anelser, af uforklarede Angester. Denne Lænke er »saare smidig, blød som Silke, giver efter for den stærkeste Anspændelse, og kan ikke slides itu.«

Mit Liv er aldeles meningsløst. Naar jeg betragter dets forskjellige Epocher, da gaaer det med mit Liv som med det Ord: Schnur i Lexikonnet, der for det første betyder en Snor, for det andet en Sønnekone. Der manglede blot, at det Ord: Schnur for det tredie skulde betyde en Kameel, for det fjerde en Støvekost.

Hvor er dog Kjedsommelighed rædsom – rædsom kjedsommelig; jeg veed intet stærkere Udtryk, intet sandere; thi kun det Lige erkjendes af det Lige. Gid der var et høiere Udtryk, et stærkere, saa var der dog endnu een Bevægelse. Jeg ligger henstrakt, uvirksom; det Eneste, jeg seer, er Tomhed, det Eneste, jeg lever af, er Tomhed, det Eneste, jeg bevæger mig i, er Tomhed. End ikke Smerte lider jeg. Gribben hakkede dog bestandig i Prometheus's Lever; Loke dryppede der dog bestandig Gift ned paa; det var dog en Afbrydelse om end eensformig. Smerten selv har tabt sin Vederqvægelse for mig. Om man bød mig al Verdens Herligheder eller al Verdens Qvaler, de røre mig lige meget, jeg vilde ikke vende mig om paa den anden Side hverken for at naae eller for at flye. Jeg døer Døden. Og hvad skulde kunne adsprede mig? Ja hvis jeg fik en Trofasthed at see, der bestod enhver Prøvelse, en Begeistring, der bar Alt; en Tro, der flyttede Bjerge; hvis jeg fornam en Tanke, der forbandt det Endelige og det Uendelige. Men min Sjæls giftige Tvivl fortærer Alt. Min Sjæl er som det døde Hav, over hvilket ingen Fugl kan flyve; naar den er kommen midtveis, synker den mat ned i Død og Undergang.

Enten – Eller

Et exstatisk Foredrag

Gift Dig, Du vil fortryde det; gift Dig ikke, Du vil ogsaa fortryde det; gift Dig eller gift Dig ikke, Du vil fortryde begge Dele; enten Du gifter Dig, eller Du ikke gifter Dig, Du fortryder begge Dele. Lee ad Verdens Daarskaber, Du vil fortryde det; græd over dem, Du vil ogsaa fortryde det; lee ad Verdens Daarskaber eller græd over dem, Du vil fortryde begge Dele; enten Du leer ad Verdens Daarskaber, eller Du græder over dem, Du fortryder begge Dele. Troe en Pige, Du vil fortryde det; troe hende ikke, Du vil ogsaa fortryde det; troe en Pige eller troe hende ikke, Du vil fortryde begge Dele; enten Du troer en Pige eller Du ikke troer hende, Du vil fortryde begge Dele. Hæng Dig, Du vil fortryde det; hæng Dig ikke, Du vil ogsaa fortryde det; hæng Dig eller hæng Dig ikke, Du vil fortryde begge Dele; enten Du hænger Dig, eller Du ikke hænger Dig, Du vil fortryde begge Dele. Dette, mine Herrer, er Indbegrebet af al Leve-Viisdom. Det er ikke blot i enkelte Øieblikke, at jeg, som Spinoza siger, betragter Alt æterno modo, men jeg er bestandig æterno modo. Dette troer Mange, at de ogsaa ere, naar de efter at have gjort det Ene eller det Andet forene eller mediere disse Modsætninger. Dog dette er en Misforstand; thi den sande Evighed ligger ikke bag ved enten – eller, men foran. Deres Evighed vil derfor ogsaa være en smertelig Tids-Succession, da de vil have den dobbelte Fortrydelse at tære paa. Min Viisdom er da let at fatte; thi jeg har kun een Grundsætning, som jeg end ikke gaaer ud fra. Man maa skjelne mellem den efterfølgende Dialektik i enten – eller og den her antydede evige. Naar jeg saaledes her siger, at jeg ikke gaaer ud fra min Grundsætning, saa har dette ikke Modsætningen i en Gaaenudderfra, men er blot det negative Udtryk for min Grundsætning, det hvorved den fatter sig selv i Modsætning til en Gaaenudderfra eller en Ikke-Gaaenudderfra. Jeg gaaer ikke ud fra min Grundsætning; thi gik jeg ud fra den, vilde jeg fortryde det, gik jeg ikke ud fra den, vilde jeg ogsaa fortryde det. Skulde det derfor forekomme en eller anden af mine høistærede Tilhørere, at der dog var Noget i det, jeg sagde, saa beviser han blot derved, at hans Hoved ikke er skikket for Philosophi; skulde det synes ham, at der var Bevægelse i det Sagte, saa beviser dette det Samme. For de Tilhørere derimod, der ere istand til at følge mig, uagtet jeg ingen Bevægelse gjør, vil jeg nu udvikle den evige Sandhed, hvorved denne Philosophi bliver i sig selv, og ikke indrømmer nogen høiere. Dersom jeg nemlig gik ud fra min Grundsætning, saa vilde jeg ikke kunne holde op igjen; thi holdt jeg ikke op, saa vilde jeg fortryde det; holdt jeg op, saa vilde jeg ogsaa fortryde det og saa videre Nu derimod da jeg aldrig gaaer ud, saa kan jeg altid høre op; thi min evige Udgang er mit evige Ophør. Erfaring har viist, at det ingenlunde er Philosophien saa vanskeligt at begynde. Langtfra; den begynder jo med Intet, og kan altsaa altid begynde. Det, der derimod falder Philosophien og Philosopherne vanskeligt, er at høre op. Ogsaa denne Vanskelighed har jeg undgaaet; thi dersom Nogen skulde troe, at jeg idet jeg nu holder op, virkelig holder op, saa viser han, at han ikke har spekulativt Begreb. Jeg holder nemlig ikke nu op; men jeg holdt op, den Gang jeg begyndte. Min Philosophi har derfor den fortrinlige Egenskab, at den er kort, og at den er uimodsigelig; thi dersom Nogen siger mig imod, saa turde jeg vel have Ret i at erklære ham for gal. Philosophen er da bestandig æterno modo og har ikke som salig Sintenis blot enkelte Timer, der ere levede for Evigheden.

Vinen fryder ikke mere mit Hjerte; lidt af den gjør mig veemodig; meget – tungsindig. Min Sjæl er mat og kraftesløs, forgjeves hugger jeg Lystens Spore i dens Side, den kan ikke meer, den reiser sig ikke mere i sit kongelige Spring. Jeg har tabt al min Illusion. Forgjeves søger jeg at hengive mig i Glædens Uendelighed, den kan ikke løfte mig, eller rettere jeg kan ikke løfte mig selv. Fordum naar den blot vinkede, da steeg jeg let og sund og freidig. Naar jeg reed langsomt gjennem Skoven, da var det som om jeg fløi; naar nu Hesten skummer nærved at styrte, da synes jeg, at jeg ikke kommer af Stedet. Ene er jeg, det har jeg altid været; forladt, ikke af Menneskene, det vilde ikke smerte mig, men af Glædens lykkelige Genier, der i talrig Skare omringede mig, der overalt traf Bekjendte, overalt viste mig en Leilighed. Som en beruset Mand samler Ungdommens kaade Vrimmel om sig, saaledes flokkedes de om mig, Glædens Alfer, og mit Smiil gjaldt dem. Min Sjæl har tabt Muligheden. Skulde jeg ønske mig Noget, da vilde jeg ikke ønske mig Rigdom eller Magt, men Mulighedens Lidenskab, det Øie, der overalt evigt ungt, evigt brændende seer Muligheden. Nydelsen skuffer, Muligheden ikke. Og hvilken Viin er saa skummende, hvilken saa duftende, hvilken saa berusende!

Hvor Solens Straaler ikke naae hen, der naae dog Tonerne. Mit Værelse er mørkt og skummelt, en høi Muur holder næsten Dagens Lys borte. Det maa være i Nabogaarden, formodentlig en omvandrende Musikant. Hvad er det for et Instrument? En Rørfløite?.... hvad hører jeg – Menuetten af Don Juan. Saa bærer mig da atter bort I rige og stærke Toner til Pigernes Kreds, til Dandsens Lyst. – Apothekeren støder i sin Morter, Pigen skurer sin Gryde, Staldkarlen strigler sin Hest og banker Striglen af paa Brostenene; kun mig gjelder disse Toner, kun mig vinker de. O! hav Tak hvo Du end er, hav Tak! min Sjæl er saa riig, saa sund, saa glædedrukken.

Solen skinner saa skjønt og livligt ind i mit Værelse, Vinduet staaer aabent i det næste; paa Gaden er Alt stille, det er Søndag-Eftermiddag: jeg hører tydelig en Lærke, der uden for et Vindue i en af Naboegaardene slaaer sine Triller, uden for det Vindue, hvor den smukke Pige boer; langt borte fra en fjern Gade hører jeg en Mand raabe med Reier; Luften er saa varm, og dog er hele Byen som uddød. – Da mindes jeg min Ungdom og min første Kjærlighed – da længtes jeg, nu længes jeg kun efter min første Længsel. Hvad er Ungdom? En Drøm. Hvad er Kjærligheden? Drømmens Indhold.

Noget Vidunderligt er der hendt mig. Jeg blev henrykket i den syvende Himmel. Der sad alle Guderne forsamlede. Af særlig Naade tilstodes den Gunst mig at gjøre et Ønske. »Vil Du«, sagde Mercur, »vil Du have Ungdom, eller Skjønhed, eller Magt, eller et langt Liv, eller den skjønneste Pige, eller en anden Herlighed af de mange, vi har i Kramkisten, saa vælg, men kun een Ting.« Jeg var et Øieblik raadvild, derpaa henvendte jeg mig til Guderne saaledes: Høistærede Samtidige, jeg vælger een Ting, at jeg altid maa have Latteren paa min Side. Der var ikke en Gud, der svarede et Ord, derimod gave de sig alle til at lee. Deraf sluttede jeg, at min Bøn var opfyldt, og fandt, at Guderne vidste at udtrykke sig med Smag; thi det havde jo dog været upassende, alvorligt at svare: det er Dig indrømmet.

\sectionheading{De umiddelbare erotiske Stadier eller det Musikalsk-Erotiske: Intetsigende Indledning  ·  SV¹ I, 48–56}

Dersom jeg tænkte mig to Riger, der grændsede op til hinanden, af hvilke jeg var temmelig nøie bekjendt med det ene, og aldeles ubekjendt med det andet, og det ikke tillodes mig at komme ind i hiint ubekjendte Rige, hvormeget jeg end ønskede det, saa vilde jeg dog være istand til at gjøre mig en Forestilling derom. Jeg vilde da vandre ud til det mig bekjendte Riges Grændse, bestandig følge den, og idet jeg gjorde dette, vilde jeg ved denne Bevægelse beskrive hiint ubekjendte Lands Omrids og saaledes have en almindelig Forestilling derom, skjøndt jeg aldrig havde sat min Fod deri. Og naar nu det var et Arbeide, der beskæftigede mig meget, naar jeg var utrættelig i min Nøiagtighed, saa vilde det vel ogsaa stundom hændes, at idet jeg stod vemodig i mit Riges Grændse, og skuede længselsfuldt ind i hiint ubekjendte Land, der var mig saa nært og dog saa fjernt, at en enkelt lille Aabenbarelse blev mig til Deel. Og føler jeg end, at Musik er en Kunst, der i høi Grad fordrer Erfaring, for at man ret skal kunne have en Mening om den, saa trøster jeg mig igjen som saa ofte med Paradoxet, at man ogsaa i Ahnelse og i Uvidenhed kan gjøre en Art Erfaring, jeg trøster mig med, at Diana, der ikke selv havde født, kom de Fødende tilhjælp, ja at hun havde det som en medfødt Gave fra Barn af, saa at hun selv kom Latone tilhjælp i hendes Fødselssmerte, da hun selv blev født.

Det mig bekjendte Rige, til hvis yderste Grændse jeg vil gaae for at opdage Musikken, er Sproget. Vil man ordne de forskjellige Medier i en bestemt Udviklingsproces, saa nødsages man ogsaa til at sætte Sproget og Musikken nærmest ved hinanden, hvorfor man jo ogsaa har sagt, at Musikken er et Sprog. Dette er nemlig mere end en aandrig Bemærkning. Naar man nemlig havde Lyst til at behage sig selv i Aandrighed, saa kunde man sige, at ogsaa Sculptur og Maleri ere en Art Sprog, forsaavidt som ethvert Udtryk for Ideen altid er et Sprog, da Ideens Væsen er Sproget. Aandrige Folk tale derfor om Naturens Sprog, og blødsødne Præster slaae af og til Naturens Bog op for os og læse Noget, som da hverken de selv eller deres Tilhørere forstaae. Naar det ikke var bedre bevendt med hiin Bemærkning, at Musikken er et Sprog, saa skulde jeg ikke uleilige den, men upaatalt lade den gaae og gjælde for hvad den var. Saaledes er det imidlertid ikke. Først idet Aanden er sat, er Sproget indsat i sine Rettigheder, men idet Aanden er sat, er Alt det, som ikke er Aand, udelukket. Men denne Udelukkelse er Aandens Bestemmelse, og forsaavidt altsaa det Udelukkede skal gjøre sig gjældende, fordrer det et Medium, der er aandelig bestemmet, og dette er Musikken. Men et Medium, der er aandelig bestemmet, er væsenlig Sprog, da nu Musikken er aandelig bestemmet, saa er den med Rette bleven kaldet et Sprog.

Sproget er som Medium betragtet det absolut aandelig bestemte Medium, det er derfor Ideens egentlige Medium. Dybere at udvikle dette, ligger hverken i min Competence eller i denne lille Undersøgelses Interesse. Kun en enkelt Bemærkning, der atter fører mig ind paa Musikken, finde her sin Plads. I Sproget er saaledes det Sandselige som Medium nedsat til blot Redskab, og bestandig negeret. Saaledes er det ikke med de andre Medier. Hverken i Sculptur eller Maleri er det Sandselige et blot Redskab, men det er et Medhenhørende, ikke heller skal det bestandig negeres, thi det skal bestandig sees med. Det vilde være en besynderlig bagvendt Betragtning af et Billedhuggerarbeide eller af et Maleri, om jeg vilde betragte det saaledes, at jeg gjorde mig Umage for at see det Sandselige bort, hvorved jeg vilde aldeles hæve dets Skjønhed. I Sculptur, Architektur, Maleri, er Ideen bunden i Mediet, men dette, at Ideen ikke nedsætter Mediet til blot Redskab, bestandig negerer det, det er ligesom et Udtryk for, at dette Medium ikke kan tale. Saaledes ogsaa med Naturen. Derfor siger man med Rette, at Naturen er stum, og Architektur og Sculptur og Maleri; det siger man med Rette tiltrods for alle de fine, følsomme Øren, der kan høre dem tale. Derfor er det en Taabelighed at sige, at Naturen er et Sprog, saa sandt som det er inept at sige, at det Stumme er talende, da det end ikke i den Forstand er et Sprog, som Fingersproget er det. Saaledes er det derimod ikke med Sproget. Det Sandselige er nedsat til blot Redskab og saaledes ophævet. Hvis et Menneske talte saaledes, at man hørte Tungens Slag og saa videre, saa talte han slet; hvis han hørte saaledes, at han hørte Luftsvingningerne istedetfor Ordet, saa hørte han slet; hvis En læste en Bog saaledes, at han bestandig saae hvert enkelt Bogstav, saa læste han slet. Netop da er Sproget det fuldkomme Medium, naar alt det Sandselige er negeret deri. Saaledes er det ogsaa med Musikken: det, der egentlig skal høres, frigjør sig bestandig fra det Sandselige. At Musikken som Medium ikke staaer saa høit som Sproget, er allerede før erindret, og derfor sagde jeg ogsaa, at i en vis Forstand var Musikken et Sprog.

Sproget henvender sig til Øret. Dette gjør intet andet Medium. Øret er igjen den mest aandelig bestemte Sands. Dette troer jeg de Fleste vil indrømme mig; ønsker Nogen nærmere Oplysning herom, saa vil jeg henvise til Fortalen til »Karrikaturen des Heiligsten« af Steffens. Musikken er foruden Sproget det eneste Medium, der henvender sig til Øret. Heri ligger atter en Analogi og et Vidne om, i hvilken Forstand Musikken er et Sprog. Der er Meget i Naturen, der henvender sig til Øret, men det, der her berører Øret, er det reent Sandselige og derfor er Naturen stum, og det er en latterlig Indbildning, at man hører Noget, fordi man hører en Ko brøle eller, hvad maaskee gjør større Prætension, en Nattergal slaae; det er Indbildning, at man hører Noget, en Indbildning, at det Ene er mere værd end det Andet, da det Alt er Hips om Haps.

Sproget har i Tiden sit Element, alle øvrige Medier har Rummet til Element. Kun Musikken foregaaer ogsaa i Tiden. Men det, at den foregaaer i Tiden, det er igjen en Negation af det Sandselige. Hvad de øvrige Kunster frembringe, antyder netop derved deres Sandselighed, at det har sin Bestaaen i Rummet. Nu er der igjen meget i Naturen, der foregaaer i Tiden. Naar saaledes en Bæk risler og bliver ved at risle, saa synes deri at ligge en Bestemmelse af Tid. Imidlertid er det ikke saa, og forsaavidt man endelig vil have, at Tidens Bestemmelse her skal være tilstede, saa maa man sige, at den vel er det, men den er rummelig bestemmet. Musikken existerer ikke uden i det Øieblik, den foredrages, thi om man endog nok saa godt kunde læse Noder og har en nok saa levende Indbildningskraft, saa kan man dog ikke negte, at det er kun i uegentlig Forstand at den er til, idet den læses. Egentlig existerer den kun idet den foredrages. Dette kunde synes en Ufuldkommenhed ved denne Kunst i Sammenligning med de andre Kunster, hvis Frembringelser bestandig bestaae, fordi de have deres Bestaaen i det Sandselige. Dog er det ikke saa. Det er netop et Beviis paa, at det er en høiere, en aandeligere Kunst.

Gaaer jeg nu ud fra Sproget, for ved en Bevægelse igjennem dette ligesom at lytte mig Musikken ud, saa viser Sagen sig omtrent saaledes. Antager jeg, at Prosa er den Sprogform, der er mest fjernet fra Musikken, saa bemærker jeg allerede i det oratoriske Foredrag, i den sonore Bygning af Perioder en Anklang af det Musikalske, der træder stærkere og stærkere frem igjennem forskjellige Trin i det poetiske Foredrag, i Versets Bygning, i Rimet, indtil endelig det Musikalske har udviklet sig saa stærkt, at Sproget hører op og Alt bliver Musik. Dette er jo et Yndlingsudtryk, som Digterne bruge for at betegne, at de ligesom give Afkald paa Ideen, den forsvinder for dem, Alt ender med Musik. Heri kunde nu synes at ligge, at Musik er et endnu fuldkomnere Medium end Sproget. Imidlertid er dette en af disse smægtende Misforstaaelser, der kun opkomme i tomme Hoveder. At det er en Misforstaaelse, skal senere igjen blive paaviist, her vil jeg blot gjøre opmærksom paa den mærkelige Omstændighed, at jeg nemlig ved en Bevægelse i modsat Retning igjen støder paa Musikken, naar jeg nemlig fra den af Begrebet gjennemtrængte Prosa gaaer ned efter, til jeg lander ved Interjectioner, der igjen ere musikalske, saaledes som jo Barnets første Lallen igjen er musikalsk. Her kan der dog vel ikke være Tale om, at Musikken er et fuldkomnere Medium end Sproget, eller at Musikken er et rigere Medium end Sproget, med mindre man antager, at det at sige: uh, er mere værd end en heel Tanke. Men hvad følger nu heraf: overalt hvor Sproget hører op møder jeg det Musikalske? Dette er dog vel det fuldkomneste Udtryk for, at Musikken overalt begrændser Sproget. Heraf vil man tillige see, hvorledes det hænger sammen med hiin Misforstaaelse, at Musikken skulde være et rigere Medium end Sproget. Idet nemlig Sproget hører op, Musikken begynder, idet, som man siger, Alt er musikalsk, saa gaaer man ikke frem, men man gaaer tilbage. Deraf kommer det, at jeg, og deri vil maaskee ogsaa Kyndige give mig Ret, aldrig har havt Sympathi for den sublimerede Musik, der mener ikke at behøve Ordet. Den mener nemlig i Reglen at være høiere end Ordet, uagtet den er ringere. Nu kunde man vel gjøre mig følgende Indvending. Hvis det er sandt, at Sproget er et rigere Medium end Musikken, saa er det ubegribeligt, at det er forbundet med saa stor Vanskelighed æsthetisk at gjøre Rede for det Musikalske, ubegribeligt, at Sproget her bestandig viser sig som et fattigere Medium end Musikken. Dette er imidlertid hverken ubegribeligt eller uforklarligt. Musikken udtrykker nemlig bestandig det Umiddelbare i dets Umiddelbarhed; deraf kommer det ogsaa, at Musikken i Forhold til Sproget viser sig først og sidst, men deraf indsees ogsaa, at det er en Misforstaaelse at sige, at Musikken er et fuldkomnere Medium. I Sproget ligger Reflexionen, og derfor kan Sproget ikke udsige det Umiddelbare. Reflexionen dræber det Umiddelbare, og derfor er det umuligt i Sproget at udsige det Musikalske, men denne Sprogets tilsyneladende Fattigdom er netop dets Rigdom. Det Umiddelbare er nemlig det Ubestemmelige, og derfor kan Sproget ikke opfatte det; men det, at det er det Ubestemmelige, er ikke dets Fuldkommenhed men en Mangel ved det. Dette anerkjender man indirecte paa mange Maader. Saaledes, for blot at anføre et Exempel, siger man: jeg veed egentlig ikke at gjøre Rede for, hvorfor jeg gjør dette eller hiint saa eller saa, jeg gjør det efter Gehøret. Man bruger her ofte om Ting, der staae i intet Forhold til det Musikalske, et Ord, hentet fra Musikken, men betegner tillige herved det Dunkle, Uforklarede, Umiddelbare.

Er nu det Umiddelbare aandelig bestemmet, det, der egentlig kommer til Udtryk i det Musikalske, saa kan man atter nærmere spørge, hvad det er for en Art af det Umiddelbare, der væsentlig er Musikkens Gjenstand. Det Umiddelbare, aandelig bestemmet, kan enten være bestemmet saaledes, at det falder ind under Aand, da kan det vel finde sit Udtryk i det Musikalske, men dette Umiddelbare kan dog ikke være Musikkens absolute Gjenstand, thi da det er bestemmet saaledes, at det skal falde ind under Aand, saa er derved antydet, at Musikken er paa et fremmed Gebeet, den danner et Forspil, som bestandig ophæves. Er derimod det Umiddelbare aandelig bestemmet, bestemmet saaledes, at det falder udenfor Aand, da har Musikken heri sin absolute Gjenstand. For det første Umiddelbare er det et Uvæsentligt, at det udtrykkes i Musik, medens det er væsentligt for det, at det bliver Aand og altsaa udtrykkes i Sproget; for dette derimod er det et Væsentligt, at det udtrykkes i Musik, det kan alene udtrykkes deri, kan ikke udtrykkes i Sproget, da det er aandeligt bestemmet saaledes, at det falder udenfor Aanden og altsaa udenfor Sproget. Men det Umiddelbare, der saaledes udelukkes af Aanden, er den sandselige Umiddelbarhed. Denne tilhører Christendommen. Den har i Musikken sit absolute Medium, og heraf lader det sig ogsaa forklare, at Musikken i den antike Verden ikke egentlig blev udviklet, men tilhører den christelige. Den er da Mediet for det Umiddelbare, der, aandelig bestemt, er bestemt saaledes, at det ligger udenfor Aanden. Naturligviis kan Musikken udtrykke meget Andet, men dette er dens absolute Gjenstand. Man mærker ogsaa let, at Musikken er et sandseligere Medium end Sproget, da jo her lægges meget mere Vægt paa den sandselige Lyden end ved Sproget.

Sandselig Genialitet er altsaa Musikkens absolute Gjenstand. Sandselig Genialitet er absolut lyrisk, og i Musikken kommer den til Udbrud i hele sin lyriske Utaalmodighed; den er nemlig aandelig bestemmet og er derfor Kraft, Liv, Bevægelse, bestandig Uro, bestandig Succession, men denne Uro, denne Succession beriger den ikke, den bliver bestandig den samme, den udfolder sig ikke, men stormer uafbrudt frem ligesom i eet Aandedrag. Naar jeg nu med et eneste Prædikat skulde betegne denne Lyriskhed, saa maatte jeg sige: den toner; og hermed er jeg da atter kommen tilbage paa den sandselige Genialitet som den, der viser sig umiddelbar musikalsk.

[…]

De Vanskeligheder, som altid møde, naar man vil gjøre Musik til Gjenstand for æsthetisk Betragtning, udeblive naturligviis heller ikke her. Vanskeligheden i det Foregaaende laae fornemmelig deri, at medens jeg ad Tankens Vei vilde bevise, at sandselig Genialitet er Musikkens væsentlige Gjenstand, dette dog egentlig kun kan bevises ved Musik, ligesom jeg jo ogsaa selv gjennem Musik er kommen til Erkjendelse deraf. Den Vanskelighed, det Følgende har at kæmpe med, er nærmest den, at da det, den Musik udtrykker, der her bliver Gjenstand for Omtale, væsentlig er Musikkens egentlige Gjenstand, saa udtrykker den det langt fuldkomnere end Sproget formaaer det, der tager sig meget fattigt ud ved Siden af den. Ja, havde jeg med forskjellige Bevidsthedstrin at gjøre, da blev Fordelen naturligviis paa min og paa Sprogets Side, men det er her ikke Tilfælde. Det, der altsaa her bliver at udvikle, kan kun have sin Betydning for den, der har hørt og bestandig vedbliver at høre. For ham kan det maaskee indeholde et enkelt Vink, der kan bevæge ham til atter at høre.

\sectionheading{Vexel-Driften. Forsøg til en social Klogskabslære  ·  SV¹ I, 257–272}

At gaae ud fra en Grundsætning paastaae erfarne Folk skal være meget forstandigt; jeg føier dem og gaaer ud fra den Grundsætning, at alle Mennesker ere kjedsommelige. Eller skulde der være Nogen, der vilde være kjedsommelig nok til at sige mig imod heri? Denne Grundsætning har nu i allerhøieste Grad den frastødende Kraft, man altid fordrer hos det Negative, der egentlig er Bevægelses-Principet; den er ikke blot frastødende men uendeligt afskrækkende, og den, der har denne Grundsætning bag ved sig, maa nødvendigviis have en uendelig Fart til at gjøre Opdagelser med. Naar nemlig min Sætning er sand, saa behøver man blot i samme Grad som man vil dæmpe eller fremskynde sin impetus, mere eller mindre tempereret at overveie med sig selv, hvor fordærvelig Kjedsommelighed er for Mennesket, og vil man næsten med Fare for Locomotivet drive Bevægelsens Hurtighed til det Høieste, saa behøver man blot at sige til sig selv: Kjedsommelighed er en Rod til alt Ondt. Det er besynderligt nok, at Kjedsommelighed, der selv er et saa roligt og adstadigt Væsen, kan have en saadan Kraft til at sætte i Bevægelse. Det er en aldeles magisk Virkning, Kjedsommeligheden udøver, kun at denne Virkning ikke er tiltrækkende men frastødende.

Hvor fordærvelig Kjedsommelighed er, det anerkjende nu ogsaa alle Mennesker i Forhold til Børn. Saalænge Børn more sig, saalænge ere de altid artige, det kan man sige i allerstrengeste Forstand, thi blive de endog i Legen stundom ustyrlige, saa er det egentlig, fordi de begynde at kjede sig; Kjedsommeligheden er allerede i Anmarsch, kun paa en anden Maade. Naar man derfor vælger en Barnepige, saa seer man altid væsentlig paa, ikke blot at hun er ædruelig, trofast og skikkelig, men man tager tillige et æsthetisk Hensyn til, om hun veed at more Børnene, og man vilde ikke tage i Betænkning at give en Barnepige Afsked, naar hun befandtes ikke at have denne Egenskab, om hun end havde alle andre fortrinlige Dyder. Her anerkjendes jo tydeligt nok Principet, men saa besynderligt gaaer det til i Verden, i den Grad har Vane og Kjedsommelighed taget Overhaand, at Barnepige er det eneste Forhold, i hvilket Æsthetiken skeer sin Ret. Naar Een vilde fordre Skilsmisse, fordi hans Kone var kjedsommelig, eller man vilde fordre en Konge afsat, fordi han var kjedsommelig at see paa, eller en Præst jaget i Landflygtighed, fordi han var kjedsommelig at høre paa, eller en Minister afskediget, eller en Journalist straffet paa Livet, fordi han var rasende kjedsommelig, saa vilde man ikke være istand til at trænge igjennem. Hvad Under da, at Verden gaaer tilbage, at det Onde mere og mere griber om sig, da Kjedsommeligheden tiltager og Kjedsommelighed er en Rod til alt Ondt. Dette kan man forfølge lige fra Verdens Begyndelse. Guderne kjedede sig, derfor skabte de Menneskene. Adam kjedede sig, fordi han var alene, derfor skabtes Eva. Fra det Øieblik af kom Kjedsommeligheden ind i Verden, og voxede i Størrelse paa det Nøiagtigste alt eftersom Folkemængden voxede. Adam kjedede sig alene, derpaa kjedede Adam og Eva sig i Forening, derpaa kjedede Adam og Eva og Kain og Abel sig en famille, derpaa tiltog Folkemængden i Verden, og Folkene kjedede sig en masse. For at adsprede sig fattede de den Tanke at bygge et Taarn, der var saa langt, at det ragede op i Skyen. Denne Tanke er ligesaa kjedsommelig som Taarnet var langt, og et forfærdeligt Beviis paa, hvorledes Kjedsommeligheden havde taget Overhaand. Derpaa bleve de adspredte over Verden, ligesom man nu reiser udenlands, men de vedbleve at kjede sig. Og hvilke Følger havde ikke denne Kjedsommelighed. Mennesket stod høit og faldt dybt, først ved Eva, saa fra det babyloniske Taarn. Hvad var det paa den anden Side, der opholdt Roms Undergang, det var panis og circenses. Hvad gjør man i vor Tid? Er man betænkt paa noget Adspredelses-Middel? Tvertimod man fremskynder Undergangen. Man tænker paa at sammenkalde en Stænderforsamling. Kan der tænkes noget Kjedsommeligere saa vel for de Herrer Deeltagere som for den, der skal læse og høre om dem? Man vil forbedre Statens Finanser ved Besparelser. Kan der tænkes noget Kjedsommeligere. Istedenfor at forøge Gjælden, vil man afbetale den. Efter hvad jeg kjender til de politiske Forhold, vil det være Danmark let at aabne et Laan paa 15 Millioner. Hvorfor tænker Ingen derpaa. At et Menneske er Geni og ikke betaler sin Gjæld, det hører man dog af og til, hvorfor skulde en Stat ikke kunne gjøre det Samme, naar der blot er Enighed? Man aabne da et Laan paa 15 Millioner, man anvende det ikke til Afbetaling, men til offentlig Forlystelse. Lader os feire det tusindaarige Rige med Fryd og Gammen. Som der nu overalt staaer Bøsser, hvori man kan lægge Penge, saa skulde der da overalt staae Skaale, hvori der laae Penge. Alt vilde blive gratis; man gik gratis i Theatret, gratis til de offentlige Fruentimmer, man kjørte gratis til Dyrehaugen, man blev begraven gratis, gratis blev der holdt Tale over Een; jeg siger gratis; thi naar man altid har Penge ved Haanden, saa er paa en Maade Alt gratis. Ingen maatte eie fast Eiendom. Kun med mig skulde der gjøres en Undtagelse. Jeg forbeholder mig 100 Rigsbankdaler. om Dagen fast i Londoner Bank, deels fordi jeg ikke kan have Mindre, deels fordi det er mig, der har givet Ideen, endelig fordi man ikke kan vide, om jeg ikke kunde hitte paa en ny Idee, naar de 15 Millioner vare forbrugte. Hvad vilde Følgen af denne Velstand blive? alt Stort vilde strømme til Kjøbenhavn, de største Kunstnere, Skuespillere og Dandserinder. Kjøbenhavn vilde blive et andet Athen. Hvad vilde Følgen blive? Alle Rigmænd vilde nedsætte sig i denne Stad. Blandt Andre vilde vel Keiseren af Persien og Kongen af Engelland ogsaa komme hertil. See her er min anden Idee. Man bemægtiger sig Keiserens Person. Maaskee vil Een sige: der bliver Oprør i Persien, man sætter en ny Keiser paa Thronen, det er skeet saa ofte før, og den gamle Keiser falder i Prisen. I saa Fald er det min Idee, at man sælger ham til Tyrken, han vil nok vide at gjøre ham i Penge. Dertil kommer endnu en Omstændighed, som vore Politikere ganske synes at oversee. Danmark er Ligevægten i Europa. Der kan ikke tænkes nogen lykkeligere Existens. Jeg veed det af egen Erfaring. Jeg var engang Ligevægten i en Familie, jeg kunde gjøre lige hvad jeg vilde, det gik aldrig ud over mig, men altid over de Andre. O! maatte min Tale trænge hen til Eders Øren, I som ere satte paa de høie Steder, til at raade og raade, I Kongens og Folkets Mænd, vise og forstandige Statsborgere af alle Classer! Seer Eder dog for! Gamle Danmark gaaer under, det er fatalt, det gaaer under af Kjedsommelighed, det er det Allerfataleste. I Oldtiden blev den Konge, der skjønnest besang den afdøde; i vor Tid burde den være Konge, der sagde den bedste Vittighed, den Kronprinds, der gav Anledning til, at den bedste Vittighed blev sagt.

Dog hvor river Du mig hen, skjønne følsomme Sværmeri! Skulde jeg oplade min Mund saaledes for at tiltale Samtiden, for at indvie den i min Viisdom? Ingenlunde; thi min Viisdom er egentlig ikke zum Gebrauch für Jedermann, og Klogskabsregler er det altid det Klogeste at tie med. Disciple ønsker jeg derfor ikke, men dersom Een stod ved mit Dødsleie, da vilde jeg maaskee, naar jeg var sikker paa, at det var forbi med mig, i et Anfald af philanthropisk Delirium, hvidske ham min Lære i Øret, uvis, om jeg havde gjort ham en Tjeneste eller ikke. Man taler saa meget om, at Mennesket er et selskabeligt Dyr, i Grunden er det et Rovdyr, Noget man ikke blot forvisser sig om ved Betragtning af dets Tænder. Den hele Snak om Selskabelighed og Samfund er derfor deels et nedarvet Hykleri, deels en udspeculeret Lumskhed.

Alle Mennesker ere da kjedsommelige. Ordet selv viser Muligheden af en Inddeling. Det Ord: kjedsommelig kan ligesaa godt betegne et Menneske, der kjeder Andre, som et, der kjeder sig selv. De, der kjede Andre, ere Plebs, Hoben, Menneskets uendelige Slæng i Almindelighed; de, der kjede sig selv, ere de Udvalgte, Adelen; og saa besynderligt er det, de, som ikke kjede sig selv, kjede i Almindelighed Andre, de derimod, der kjede sig selv, underholde Andre. De, der ikke kjede sig, ere i Almindelighed de, der paa en eller anden Maade have travlt i Verden, men disse ere netop derfor de allerkjedsommeligste, de utaaleligste. Denne Dyreclasse er vistnok ikke Frugt af Mandens Begjær og Qvindens Lyst. Den udmærker sig som alle lavere Dyreclasser ved en høi Grad af Frugtbarhed, og formerer sig i det Utrolige. Ubegribeligt var det ogsaa, om Naturen behøvede 9 Maaneder for at frembringe saadanne Væsener, der vel snarere lod sig frembringe i Sneseviis. Den anden Classe Mennesker, de fornemme, er de, som kjede sig selv. Som ovenfor bemærket more de i Almindelighed Andre, paa en vis udvortes Maade stundom Pøbelen, i dybere Forstand de Medindviede. Jo grundigere de kjede sig selv, desto kraftigere et Adspredelses-Middel frembyde de for disse, ogsaa naar Kjedsommeligheden naaer sit Høieste, idet de enten (den passive Bestemmelse) døe af Kjedsommelighed, eller (den aktive Bestemmelse) skyde sig af Nysgjerrighed.

Lediggang, pleier man at sige, er en Rod til alt Ondt. For at forhindre det Onde anbefaler man Arbeide. Man seer imidlertid let saa vel af den befrygtede Anledning som af det anbefalede Middel, at den hele Betragtning er af meget plebeiisk Extraktion. Lediggang som saadan er ingenlunde en Rod til Ondt, tvertimod den er et sandt guddommeligt Liv, naar man ikke kjeder sig. Ja Lediggang kan give Anledning til at man taber sin Formue og saa videre; men den adelige Natur frygter ikke Sligt, men vel at kjede sig. De olympiske Guder kjedede sig ikke, de levede lykkelige i lykkelig Lediggang. En qvindelig Skjønhed, der hverken syer eller spinder eller stryger eller læser, eller musicerer, hun er lykkelig i Lediggang; thi hun kjeder sig ikke. Lediggang er da saa langt fra at være Roden til det Onde, at den snarere er det sande Gode. Kjedsommeligheden er Roden til det Onde, den er det, der maa holdes borte. Lediggang er ikke det Onde, ja man maa sige, at ethvert Menneske, der ikke har Sands derfor, derved viser, at han ikke har hævet sig til det Humane. Der gives en utrættelig Virksomhed, som udelukker et Menneske fra Aandens Verden og sætter ham i Classe med Dyrene, der instinktagtigt altid maae være i Bevægelse. Der gives Mennesker, som have en overordentlig Gave til at forvandle Alt til Forretning, hvis hele Liv er Forretning, som forelske og gifte sig, høre en Vittighed og beundre et Kunststykke med samme Forretnings-Iver som den, med hvilken de arbeide i Contoiret. Det latinske Ordsprog otium est pulvinar diaboli er ganske rigtigt; men Djævelen faaer ikke Tid at lægge sit Hoved paa denne Pude, naar man ikke kjeder sig. Da imidlertid Folk troe, at det er Menneskets Bestemmelse at arbeide, saa er Modsætningen rigtig: Lediggang – Arbeide. Jeg antager, at det er Menneskets Bestemmelse at more sig, min Modsætning er derfor ikke mindre rigtig.

Kjedsommelighed er den dæmoniske Pantheisme. Bliver man staaende ved den som saadan, da bliver den det Onde, saasnart den derimod ophæves, saa er den sand; men den ophæves kun ved at more sig – ergo bør man more sig. At sige, at den ophæves ved at arbeide røber Uklarhed; thi Lediggang kan visselig hæves ved Arbeide, da dette er dens Modsætning, men Kjedsommelighed ikke, som man jo ogsaa seer, at de allertravleste Arbeidere, de i deres emsige Brummen mest snurrende Insekter, ere de allerkjedsommeligste; og naar de ikke kjede dem, saa kommer det deraf, at de ingen Forestilling have om, hvad Kjedsommelighed er; men saa er Kjedsommeligheden ikke ophævet.

Kjedsommelighed er deels en umiddelbar Genialitet, deels en erhvervet Umiddelbarhed. Den engelske Nation er i det Hele den paradigmatiske Nation. Den sande geniale Indolens træffer man sjeldnere, i Naturen findes den ikke, den tilhører Aandens Verden. Man træffer undertiden en reisende Englænder, der er som en Incarnation af denne Genialitet, et tungt ubevægeligt Murmeldyr, hvis hele Sprog-Riigdom gaaer op i et eneste Eenstavelsesord, en Interjektion, hvormed han betegner sin høieste Beundring og sin dybeste Ligegyldighed, fordi Beundring og Ligegyldighed har indifferentieret sig i Kjedsommelighedens Eenhed. Ingen anden Nation end den engelske frembringer saadanne Naturmærkværdigheder; ethvert Menneske, der tilhører en anden Nation, vil altid være lidt mere livlig, ikke saa absolut dødfødt. Den eneste Analogi, jeg kjender, er den tomme Begeistrings Apostle, der ligeledes reise Livet igjennem paa en Interjektion, Mennesker, der overalt gjøre Profession af at være begeistrede, overalt ere tilstede, og hvad enten der skeer noget Betydeligt eller noget Ubetydeligt, raabe: Ei! eller Oi!, fordi det Betydeliges og det Ubetydeliges Differens for dem har indifferentieret sig i den blindt larmende Begeistrings Tomhed. Den senere Kjedsommelighed er gjerne en Frugt af en misforstaaet Adspredelse. At saaledes det, der er Midlet mod Kjedsommelighed, kan fremkalde den, synes betænkeligt; men det kan ogsaa kun forsaavidt fremkalde den, som det bliver urigtigt anvendt. En forkeert, i Almindelighed excentrisk Adspredelse har ogsaa Kjedsommeligheden i sig, og det er paa den Maade den arbeider sig frem, og viser sig som det Umiddelbare. Som man med Hensyn til Heste skjelner mellem Død-Kuller og Flyve-Kuller, men dog kalder begge Dele Kuller, saaledes kan man ogsaa gjøre Forskjel mellem to Arter af Kjedsommelighed, som dog begge enes i Bestemmelsen af Kjedsommelighed.

I Pantheismen ligger i Almindelighed Bestemmelsen af Fylde, med Kjedsommelighed er det omvendt, den er bygget paa Tomhed, men er netop derfor en pantheistisk Bestemmelse. Kjedsommelighed hviler paa det Intet, der slynger sig gjennem Tilværelsen, dens Svimmelhed er som den, der fremkommer ved at skue ned i en uendelig Afgrund, uendelig. At derfor hiin excentriske Adspredelse er bygget paa Kjedsommelighed, kan man ogsaa see deraf, at Adspredelsen gjenlyder uden Efterklang, netop fordi der i Intet end ikke er saa meget, at en Gjenlyd bliver mulig.

Naar nu som ovenfor udviklet Kjedsommelighed er en Rod til alt Ondt, hvad er da naturligere end at man søger at overvinde den. Det gjelder imidlertid her som allevegne fornemlig om roligt Overlæg, at man ikke dæmonisk besat af Kjedsommelighed, idet man vil undflye den, arbeider sig ind i den. Forandring raaber Alle, der kjede sig, paa. Heri er jeg ganske enig med dem, kun gjelder det om at handle efter Princip.

Min Afvigelse fra den almindelige Anskuelse er tilstrækkeligt udtrykt ved det Ord: Vexeldrift. I dette Ord kunde synes at ligge en Tvetydighed, og naar jeg i dette Ord vilde finde Rum for en Betegnelse af den almindelige Methode, saa maatte jeg sige, at Vexeldriften bestod i bestandig at vexle Jordbund. Saaledes bruger Landmanden imidlertid ikke dette Udtryk. Jeg vil dog et Øieblik benytte det saaledes, for at tale om den Vexeldrift, der beroer paa Forandringens grændseløse Uendelighed, dens extensive Dimension.

Denne Vexeldrift er den vulgaire, den ukunstneriske, og ligger i en Illusion. Man er kjed af at leve paa Landet, man reiser til Hovedstaden; man er kjed af sit Fødeland, man reiser udenlands; man er »europamüde« man reiser til America og saa videre, man hengiver sig til et sværmerisk Haab om en uendelig Reisen fra Stjerne til Stjerne. Eller Bevægelsen er en anden, men dog extensiv. Man er kjed af at spise paa Porcellain, man spiser paa Sølv; man er kjed heraf, man spiser paa Guld, man brænder det halve Rom af for at see Troias Ildebrand. Denne Methode hæver sig selv og er den slette Uendelighed. Hvad opnaaede ogsaa Nero? Nei da var Keiser Antonin klogere, han siger: αναβιωναι σοι εξεστιν· ιδε παλιν τα πραγματα, ως εωρας· εν τουτω γαρ το αναβιωναι. (Bιβλιον Z., β.)

Den Methode, jeg foreslaaer, ligger ikke i at forandre Jordbunden, men som den sande Vexeldrift i at skifte Drifts-Methoden og Sædearter. Her ligger strax det Begrændsnings Princip, som er det ene frelsende i Verden. Jo mere man begrændser sig selv, desto meer opfindsom bliver man. En eensom Livsfange er saare opfindsom, en Ædderkop kan være ham til stor Moro. Man tænke paa sin Skoletid, da man er indtraadt i den Alder, hvor der intet æsthetisk Hensyn tages i Valget af dem, der skal belære Een, og disse desaarsag ofte ere meget kjedsommelige; hvor opfindsom er man da ikke? Hvor kan det ikke more Een, at have fanget en Flue og holde den fangen under en Nøddeskal, og see, hvorledes den kan løbe om med denne; hvor kan det glæde at have skaaret et Hul i Bordet, at spærre en Flue ind deri, og gjennem et Stykke Papir at kigge ned paa den? Hvor underholdende kan det ikke være at høre det eensformige Tagdryp? Hvilken grundig Iagttager bliver man ikke, ikke den mindste Larm eller Bevægelse undgaaer Een. Her er Yderspidsen af det Princip, der ikke ved Extensitet men ved Intensitet søger efter Beroligelse.

Jo opfindsommere et Menneske kan være i at forandre Drifts-Methoden, desto bedre; men enhver enkelt Forandring ligger dog indenfor den almindelige Regel af Forholdet mellem at erindre og at glemme. Det er i disse to Strømninger at hele Livet bevæger sig, og det gjelder derfor om ret at have dem i sin Magt. Først naar man har kastet Haabet over Bord, først da begynder man at leve kunstnerisk, saa længe man haaber, kan man ikke begrændse sig. Det er ret smukt at see et Menneske stikke i Søen med Haabets Medbør, man kan benytte Leiligheden til at lade sig tage med paa Slæbetouget, men selv bør man aldrig have det ombord paa sin Skude, allermindst som Lods; thi det er en troløs Skibsfører. Haabet var derfor ogsaa den ene af Prometheus's betænkelige Gaver; istedenfor de Udødeliges Forudvidenhed gav han Menneskene Haabet.

Glemme – det ville alle Mennesker; og naar der møder dem noget Ubehageligt, saa sige de altid, hvo der dog kunde glemme. Men det at glemme er en Kunst, som maatte være indøvet iforveien. At kunne glemme beroer altid paa, hvorledes man husker; men hvorledes man husker, beroer atter paa, hvorledes man oplever Virkeligheden. Den, der med Haabets Fart løber sig fast, han vil erindre saaledes, at han ikke formaaer at glemme. Nil admirari er derfor den egentlige Leveviisdom. Ethvert Livs-Moment maa ikke have mere Betydning for Een end at man kan, hvad Øieblik man vil, glemme det; ethvert enkelt Livs-Moment maa fra den anden Side have saa megen Betydning for Een, at man hvert Øieblik kan huske det. Den Alder, der husker bedst, er tillige den glemsomste, det er Barnealderen. Jo mere poetisk man husker, jo lettere glemmer man, thi det at huske poetisk er egentlig blot et Udtryk for at glemme. Naar jeg husker poetisk, saa er der allerede foregaaet den Forandring med det Oplevede, hvorved det har tabt alt det Piinlige. For saaledes at kunne erindre, maa man være opmærksom paa, hvorledes man lever, især hvorledes man nyder. Nyder man frisk væk indtil det Sidste, tager man bestandig det Høieste med, som Nydelsen kan give, saa vil man hverken være istand til at erindre eller til at glemme. Man har da nemlig intet Andet at erindre end en Overmæthed, som man kun ønsker at glemme, men som nu plager med en ufrivillig Erindring. Naar man derfor mærker, at Nydelsen eller et Livs-Moment for stærkt river Een hen, da holder man et Øieblik inde og erindrer. Der er intet Middel, der bedre giver Afsmag for at vedblive for længe. Man holder fra Begyndelsen af Styr paa Nydelsen; sætter ikke alle Seil til for enhver Beslutning; man hengiver sig med en vis Mistro, først da er man istand til at gjøre det Ordsprog til Løgner, der siger, at man ikke paa eengang kan have baade i Sæk og i Pose. Vel forbyder Politiet at bære hemmelige Vaaben, og dog er intet Vaaben saa farligt som den Kunst at kunne erindre. Det er en egen Følelse, naar man midt i Nydelsen seer paa den for at erindre.

Naar man saaledes har perfektioneret sig i den Kunst at glemme og den Kunst at erindre, saa er man istand til at spille Fjæderbold med hele Tilværelsen.

Paa Kraften til at glemme kan man egentlig maale et Menneskes Elasticitet. Den, der ikke kan glemme, ham bliver der ikke Meget af. Om der nogensteds skulde sprudle en Lethe, veed jeg ikke; men det veed jeg, at denne Kunst kan udvikles. Den bestaaer dog ingenlunde i, at det enkelte Indtryk skulde sporløst forsvinde; thi Glemsomhed er ikke identisk med den Kunst at kunne glemme. Man seer ogsaa let, hvor liden Forstand Folk i Almindelighed have paa denne Kunst; thi de ville som oftest kun glemme det Ubehagelige, ikke det Behagelige. Dette røber en fuldkommen Eensidighed. Glemsel er nemlig det rette Udtryk for den egentlige Assimilation, der nedsætter det Oplevede til Resonansbund. Derfor er Naturen saa stor, fordi den har glemt, at den var Chaos, men denne Tanke kan fremkomme naar det skal være. Da man som oftest kun tænker sig Glemsel i Forhold til det Ubehagelige, saa forestiller man sig den oftest som en vild Magt, der overdøver. Men Glemsel er tvertimod en stille Syslen, og den bør være ligesaa vel i Forhold til det Behagelige som til det Ubehagelige. Ogsaa det Behagelige har som forbigangent, netop som forbigangent, en Ubehagelighed i sig, hvorved det kan vække Savnet; denne Ubehagelighed hæves ved Glemsel. Det Ubehagelige har en Braad, det tilstaae Alle. Ogsaa den tages bort ved Glemsel. Bærer man sig imidlertid ad som Mange af dem, der fuske i den Kunst at glemme, slaaer det Ubehagelige aldeles bort, saa vil man snart see, hvad det hjælper til. I et ubevogtet Øieblik overrasker det En ofte med det Pludseliges hele Magt. Dette er aldeles stridende mod den velordnede Indretning i et forstandigt Hoved. Ingen Ulykke, ingen Gjenvordighed er saa lidet affabel, saa døv, at den ikke lader sig smigre lidt; selv Cerberus tog mod Honningkager, og det er ikke blot Pigebørn, man bedaarer. Man besnakker den og berøver den derved sin Skarphed, og ønsker ingenlunde at glemme den, men glemmer den for at erindre den. Ja selv med Minder af den Natur, at man skulde synes evig Glemsel var det eneste Middel mod dem, tillader man sig slig Underfundighed, og Falskneriet lykkes den Behændige. Glemsel er den Sax, med hvilken man klipper bort, hvad man ikke kan bruge, vel at mærke under Erindringens allerhøieste Opsyn. Glemsel og Erindring ere saaledes identiske, og den kunstneriske tilveiebragte Identitet det archimediske Punkt, med hvilket man løfter hele Verden. Naar man siger at skrive Noget i Glemmebogen, saa antyder man jo paa eengang at det glemmes og at det dog opbevares.

Den Kunst at erindre og at glemme vil da ogsaa forebygge, at man ikke løber sig fast i noget enkelt Livsforhold og sikkre Een den fuldkomne Svæven.

Man vogte sig da for Venskab. Hvorledes defineres en Ven? En Ven er ikke hvad Philosophien kalder det nødvendige Andet, men det overflødige Tredie. Hvilke ere Venskabs Ceremonier? Man drikker Duus, man aabner en Aare, man blander sit Blod med Vennens. Naar dette Øieblik kommer er vanskeligt at bestemme; men det forkynder sig selv paa en gaadefuld Maade, man føler det, man kan ikke mere sige De til hinanden. Naar denne Følelse har været tilstede, saa kan det aldrig vise sig, at man har taget feil som Geert Vestphaler, da han drak Duus med Skarpretteren. – Hvilket er Venskabets sikkre Kjendetegn? Oldtiden svarer: idem velle, idem nolle ea demum firma amicitia, og tillige yderst kjedsommeligt. Hvilken er Venskabets Betydning? Gjensidig Assistance med Raad og Daad. Derfor slutte to Venner sig nøie sammen, for at være Alt for hinanden; og det uagtet det ene Menneske slet ikke kan være Andet for det andet Menneske end være i Veien for ham. Ja man kan hjælpe hinanden med Penge, hjælpe hinanden i og af Frakken, være hinandens ydmyge Tjenere, møde til en oprigtig Nytaars-Gratulation, item til Bryllup, Barsel og Begravelse.

Men fordi man afholder sig fra Venskab, derfor skal man ikke leve uden Berøring med Menneskene. Tvertimod disse Forhold kunne undertiden ogsaa antage et dybere Sving, kun at man altid, skjøndt man en tidlang deler Bevægelsens Fart, dog har saa meget mere Hurtighed, at man kan løbe fra det. Man mener vel, at en saadan Adfærd efterlader ubehagelige Erindringer, at det Ubehagelige ligger i, at et Forhold fra at have været Noget for Een svinder hen til at være Intet. Dette er dog en Misforstaaelse. Det Ubehagelige er nemlig en pikant Ingrediens i Livets Tværhed. Desuden kan det samme Forhold atter faae Betydning paa en anden Maade. Det, man skal passe paa, er, aldrig at løbe fast, og til den Ende altid at have Glemselen bag Øret. Den erfarne Landmand brakker af og til, den sociale Klogskabslære anbefaler det Samme. Alt kommer nok igjen men paa en anden Maade; hvad der eengang har været optaget i Rotationen bliver deri, men varieres ved Drifts-Maaden. Man haaber derfor ganske consequent at træffe sine gamle Venner og Bekjendtere i en bedre Verden, men man deler ikke Mængdens Frygt for, at de skulde have forandret sig saa meget, at man ikke kunde kjende dem igjen; man frygter snarere for, at de skulde være aldeles uforandrede. Det er utroligt, hvad selv det ubetydeligste Menneske kan vinde ved en saadan fornuftig Drift.

Man indlade sig aldrig i Ægteskab. Ægtefolkene de love hinanden Kjærlighed for evig. Det er nu nemt nok, men har heller ikke stort at betyde; thi bliver man færdig med Tiden, Evigheden bliver man nok færdig med. Naar derfor Vedkommende istedenfor at sige for evig, sagde: til Paaske, eller til førstkommende Maidag, saa var der dog Mening i deres Tale; thi baade sagde man Noget, og Noget, man maaskee kunde holde. Hvorledes gaaer det ogsaa i Ægteskabet. Efter kort Tids Forløb mærker først den ene Part, at det er galt; nu klager den anden Part og raaber høit i Sky: Troløshed, Troløshed. Efter nogen Tids Forløb kommer den anden Part paa det samme Punkt, og der tilveiebringes en Neutralitet, idet den gjensidige Troløshed qvitter til fælleds Contentement og Fornøielse. Imidlertid er det dog bag efter; thi en Skilsmisse er forbunden med store Vanskeligheder.

Naar det forholder sig saaledes med Ægteskabet, saa er det ikke forunderligt, at man paa saa mange Maader maa stive det af med moralske Støtter. Naar Een vil skilles fra sin Kone, saa raaber man: han er et nedrigt Menneske, en Skurk og saa videre Hvor taabeligt, og hvilket indirecte Angreb paa Ægteskabet. Enten har Ægteskabet Realitet i sig selv, og saa er han jo straffet tilstrækkelig ved at gaae Glip af denne; eller det har ingen Realitet, og saa er det jo urimeligt at udskjelde ham, fordi han er visere end Andre. Naar Een blev kjed af sine Penge, og kastede dem ud af Vinduet, saa vilde Ingen sige, han var et nedrigt Menneske; thi enten har Penge Realitet og saa er han jo tilstrækkelig straffet ved at berøve sig dem, eller de have ingen Realitet, og da er han jo viis.

Man maa altid vogte sig for at contrahere et Livsforhold, hvorved man kan blive til Flere. Derfor er Venskabet allerede farligt, end mere Ægteskabet. Vel siger man, at Ægtefolkene blive til een; men dette er en saare dunkel og mystisk Tale. Naar man er Flere, saa har man tabt sin Frihed og kan ikke bestille Reisestøvler, naar man vil, kan ikke flakke ustadigt om. Har man Kone er det vanskeligt, har man Kone og maaskee Børn, er det besværligt, har man Kone og Børn, er det umuligt. Vel har man havt Exempel paa, at en Zigeunerinde har baaret sin Mand paa Ryggen gjennem Livet, men deels er dette en Sjeldenhed, deels ogsaa i Længden trættende – for Manden. Ved Ægteskabet geraader man desuden i en høist fatal Continuitet med Skik og Brug, og Skik og Brug er ligesom Veir og Vind noget aldeles Ubestemmeligt. I Japan er det, saavidt jeg veed, Skik og Brug, at ogsaa Mændene ligge i Barselseng. Hvorfor kunde den Tid ikke komme, da Europa vilde indføre fremmede Landes Skikke.

Venskabet er allerede farligt, Ægteskabet er det endnu meer; thi Qvinden er og bliver Mandens Ruin, saa snart man contraherer et vedvarende Forhold med hende. Tag et ungt Menneske, fyrigt som en arabisk Hest, lad ham gifte sig, han er tabt. Først er Qvinden stolt, saa er hun svag, saa besvimer hun, saa besvimer han, saa besvimer hele Familien. En Qvindes Kjærlighed er kun Forstillelse og Svaghed.

Fordi man ikke indlader sig i Ægteskab, derfor behøver Eens Liv ikke at være uden Erotik. Det Erotiske bør ogsaa have Uendelighed, men poetisk Uendelighed, der ligesaa godt kan begrændses i een Time som i en Maaned. Naar to Mennesker forelske sig i hinanden, og ahne, at de ere bestemte for hinanden, da gjelder det om at have Mod til at bryde af; thi ved at vedblive er kun Alt at tabe, Intet at vinde. Det synes et Paradox og er det ogsaa for Følelsen, ikke for Forstanden. Paa dette Gebeet gjelder det især om at kunne bruge Stemninger, kan man det, da kan man tilveiebringe en uudtømmelig Afvexling af Combinationer.

Man paatage sig aldrig nogen Kaldsforretning. Gjør man det, saa bliver man en slet og ret Peter Flere, en lille bitte Tap med i Statslegemets Maskine; man ophører selv at være Drifts-Herren, og da kan Theorien kun lidet hjælpe. Man faaer en Titel, og heri ligger Syndens og det Ondes hele Consequens. Den Lov, man da træller under, er lige kjedsommelig, enten Avancementet er hurtigt eller langsomt. En Titel bliver man aldrig af med igjen, det skulde da være ved en Forbrydelse, der paadrog Een Kagstrygning, og selv da er man ikke sikker for, at man jo kan blive benaadet ved en kongelig Resolution og faae sin Titel igjen.

Om man end afholder sig fra Kaldsforretninger bør man dog ikke være uvirksom, men lægge Vægt paa al den Beskæftigelse, der er identisk med Ørkesløshed, man maa drive allehaande brødløse Kunster. Dog bør man i denne Henseende ikke udvikle sig saa meget extensivt som intensivt, og skjøndt ældre i Aar, vise Rigtigheden af det gamle Ord, at det er Lidet, der kan fornøie Børn.

Som man nu, ifølge den sociale Klogskabslære, til en vis Grad varierer Jordbunden; thi dersom man kun vilde leve i Forhold til eet Menneske, saa maatte Vexeldriften falde slet ud, ligesom hvis en Landmand kun havde een Tønde Land, hvoraf Følgen vilde blive, at han aldrig kunde brakke, noget, der er saa yderlig vigtigt; saaledes maa man ogsaa bestandigt variere sig selv; og dette er egentlig Hemmeligheden. Til den Ende maa man nødvendig have Stemningerne i sin Magt. I den Forstand at have dem i sin Magt, at man skulde kunne frembringe dem, naar man vilde, er en Umulighed, men Klogskaben lærer at benytte Øieblikket. Som en erfaren Sømand altid forskende seer hen over Vandet, og seer en Kastevind længe forud, saaledes bør man altid see Stemningen lidt forud. Man maa vide, hvorledes Stemningen virker paa Een selv og efter Sandsynlighed paa Andre, førend man ifører sig den. Man stryger først for at fremkalde rene Toner og see, hvad der stikker i et Menneske, senere følge Mellemtonerne. Jo mere Praxis man har, jo lettere vil man overbevise sig om, at der ofte er meget i et Menneske, som man aldrig tænker paa. Naar følsomme Mennesker, der som saadanne ere yderst kjedsommelige, blive arrige, saa ere de ofte morsomme. Drilleri er især et ypperligt Explorations-Middel.

I Vilkaarligheden ligger hele Hemmeligheden. Man troer det er ingen Kunst at være vilkaarlig, og dog hører der et dybt Studium til at være vilkaarlig saaledes, at man ikke selv løber vild deri, at man selv har Fornøielse deraf. Man nyder ikke umiddelbart, men noget ganske Andet, som man selv vilkaarligt lægger ind. Man seer Midten af et Theater-Stykke, læser tredie Deel af en Bog. Derved faaer man en ganske anden Nydelse end den, Forfatteren har været saa god at tiltænke Een. Man nyder noget aldeles Tilfældigt, man betragter hele Tilværelsen fra dette Standpunkt, lader dens Realitet strande derpaa. Jeg vil anføre et Exempel. Der var et Menneske, hvis Passiar et Livsforhold gjorde det nødvendigt for mig at høre paa. Han var ved enhver Leilighed til Rede med et lille philosophisk Foredrag, der var yderst kjedsommeligt. Nær ved at fortvivle opdager jeg pludselig, at han svedte ualmindelig stærkt, naar han talte. Denne Sved tildrog sig nu min Opmærksomhed. Jeg saae, hvorledes Sved-Perlerne samlede sig paa hans Pande, derpaa forenede sig i Bække, gled ned ad hans Næse, og endte i et draabeformigt Legeme, der blev hængende paa Næsens yderste Spidse. Fra dette Øieblik af var Alt forandret, jeg kunde endogsaa have min Glæde af at anspore ham til at begynde sin philosophiske Belæring, alene for at iagttage Sveden paa hans Pande og paa hans Næse. Baggesen siger etsteds om en Mand, at han vist er et meget skikkeligt Menneske, men at han har den Ting at indvende imod ham, at Intet rimer paa hans Navn. Det er yderst velgjørende saaledes at lade Livets Realiteter indifferentiere sig paa en saadan vilkaarlig Interesse. Man gjør noget Tilfældigt til det Absolute og som saadant til Gjenstand for absolut Beundring. Dette virker især ypperligt, naar Gemytterne ere i Bevægelse. I Forhold til mange Mennesker er denne Methode et ypperligt Pirrings-Middel. Man betragter Alt i Livet som et Veddemaal og saa videre Jo mere consequent man veed at fastholde sin Vilkaarlighed, desto morsommere blive Combinationerne. Consequensens Grad viser altid, om man er Kunstner eller Fusker; thi til en vis Grad gjøre alle Mennesker det Samme. Det Øie, hvormed man seer Virkeligheden, maa bestandig forandres. Neo-Platonikerne antoge, at de Mennesker, der havde været mindre fuldkomne i Verden, bleve efter Døden til meer eller mindre fuldkomne Dyr, alt efter deres Fortjenester; de, der for Exempel havde øvet borgerlige Dyder efter en mindre Grad (Detaillisterne), bleve til borgerlige Dyr f. E. til Bier. En saadan Livs-Betragtning, der her i Verden seer alle Mennesker forvandlede til Dyr eller til Planter (ogsaa dette meente Plotin, at nogle bleve forvandlede til Planter) frembyder en riig Mangfoldighed af Afvexling. Maleren Tischbein har forsøgt at idealisere ethvert Menneske til et Dyr. Hans Methode har den Feil, at den er for alvorlig, og stræber at udfinde en virkelig Lighed.

Med Vilkaarligheden i Een selv corresponderer Tilfældigheden udenfor Een. Man bør derfor altid have et aabent Øie for det Tilfældige, altid være expeditus, hvis Nogen skulde tilbyde sig. De saa kaldte selskabelige Glæder, til hvilke man forbereder sig otte eller fjorten Dage forud, have ikke stort at betyde, derimod kan selv den ubetydeligste Ting ved et Tilfælde blive et riigt Stof til Morskab. At gaae her i Detail lader sig ikke gjøre, saavidt kan ingen Theori naae. Selv den udførligste Theori er dog kun Fattigdom imod hvad Geniet i sin Ubiquitet let opdager.

\sectionheading{Forførerens Dagbog (seks uddrag)  ·  SV¹ I, 276–412}

Titlen paa Bogen frapperede mig ikke i og for sig; jeg tænkte det var en Samling af Excerpter, hvilket forekom mig ganske naturligt, da jeg vidste, at han med Iver altid havde omfattet sine Studier. Den indeholdt imidlertid ganske andre Ting. Det var hverken meer eller mindre end en Dagbog, omhyggeligt ført; og som jeg, efter hvad jeg tilforn havde kjendt til ham, ikke fandt, at hans Liv saa høilig trængte til Commentar, saa negter jeg ikke, efter det Indblik, jeg nu gjorde, at Titlen er valgt med megen Smag og megen Forstand, med sand æsthetisk, objektiv Overlegenhed over ham selv og over Situationen. Denne Titel staaer i fuldkommen Harmoni med hele Indholdet. Hans Liv har været et Forsøg paa at realisere den Opgave at leve poetisk. Med et skarpt udviklet Organ for at udfinde det Interessante i Livet, har han vidst at finde det, og efterat have fundet det bestandigt halvt digterisk reproduceret det Oplevede. Hans Dagbog er derfor ikke historisk nøiagtig eller simpelt fortællende, ikke indikativisk, men conjunktivisk. Uagtet det Oplevede naturligviis er optegnet efter at det er oplevet, stundom maaskee endog længere Tid derefter, saa er det dog ofte fremstillet saaledes som om det foregik i samme Øieblik, saa dramatisk levende, at det stundom er som foregik Alt for Ens Øine. At han nu skulde have gjort det, fordi han havde nogensomhelst anden Hensigt med denne Dagbog, er høist usandsynligt; at den i strængeste Forstand blot har havt personlig Betydning for ham, er iøinefaldende; og at ville antage, at jeg har et Digterværk for mig, maaskee endog bestemt til at trykkes, forbyder saavel det Hele som det Enkelte. Vel behøvede han ikke at befrygte Noget for sin Person ved at udgive den; thi de fleste Navne ere saa besynderlige, at der aldeles ingen Sandsynlighed er for, at de ere historiske; kun har jeg fattet en Mistanke om, at Fornavnet er historisk rigtigt, saaledes, at han altid selv har været sikker paa at gjenkjende den virkelige Person, medens enhver Uvedkommende maatte vildledes af Familie-Navnet. Saaledes er det idetmindste Tilfælde med den Pige, jeg kjendte, om hvem Hovedinteressen dreier sig, Cordelia, hun heed meget rigtigt Cordelia, men derimod ikke Wahl.

Hvorledes lader det sig nu forklare, at Dagbogen desuagtet har faaet et saadant digterisk Anstrøg? Svaret herpaa er ikke vanskeligt, det lader sig forklare af den digteriske Natur, der er i ham, der, om man saa vil, ikke er riig nok, om man saa vil, ikke fattig nok til at adskille Poesi og Virkelighed fra hinanden. Det Poetiske var det Mere, han selv bragte med sig. Dette Mere var det Poetiske han nød i Virkelighedens poetiske Situation; dette tog han atter tilbage i Form af digterisk Reflexion. Det var den anden Nydelse, og Nydelse var hele hans Liv beregnet paa. I første Tilfælde nød han personlig det Æsthetiske, i andet Tilfælde nød han æsthetisk sin Personlighed. I første Tilfælde var Pointet det, at han egoistisk personlig nød det, som deels Virkeligheden gav ham, deels han selv havde besvangret Virkeligheden med; i andet Tilfælde blev hans Personlighed forflygtiget, og han nød da Situationen og sig selv i Situationen. I første Tilfælde behøvede han bestandig Virkeligheden som Anledning, som Moment; i andet Tilfælde var Virkeligheden druknet i det Poetiske. Det første Stadiums Frugt er saaledes den Stemning, af hvilken Dagbogen er fremgaaet som det andet Stadiums Frugt, dette Ord taget i sidste Tilfælde i en noget anden Betydning end i første. Det Poetiske har han saaledes bestandigt havt ved den Tvetydighed, i hvilken hans Liv gik hen.

[…]

Min Cordelia!

Forelsket er jeg i mig selv, siger man om mig. Det forundrer mig ikke; thi hvorledes skulde man kunne mærke, at jeg kan elske, da jeg kun elsker Dig, hvorledes skulde nogen Anden ahne det, da jeg kun elsker Dig. Forelsket er jeg i mig selv, hvorfor? fordi jeg er forelsket i Dig; thi Dig elsker jeg, Dig alene og Alt, hvad der tilhører Dig i Sandhed, og saaledes elsker jeg mig selv, fordi dette mit Jeg tilhører Dig, saa dersom jeg ophørte at elske Dig, ophørte jeg at elske mig selv. Hvad der da i Verdens profane Øine er Udtryk for den største Egoisme, det er for Dit indviede Blik Udtryk for den reneste Sympathi, hvad der i Verdens profane Øine er Udtryk for den mest prosaiske Selvopholdelse, det er for Dit helligede Syn Udtryk for den mest begeistrede Tilintetgjørelse af sig selv.

Din Johannes.

[…]

En gammel Philosoph har sagt, naar man nøiagtigt nedskriver Alt, hvad man oplever, saa er man, inden man veed et Ord af det, Philosoph. Jeg har nu i længere Tid levet i Forbindelse med de Forlovedes Menighed. Nogen Frugt maa jo dog et saadant Forhold yde. Jeg har da tænkt paa at samle Materialier til et Skrift, betitlet: Bidrag til Kyssets Theori, alle ømt Elskende helliget. Det er iøvrigt mærkeligt, at der intet Skrift existerer om denne Sag. Hvis det da lykkes mig at blive færdig, vil jeg tillige afhjælpe et længe følt Savn. Skulde denne Mangel i Litteraturen have sin Grund i, at Philosopherne ikke tænke over Sligt, eller i, at de ikke forstaae sig paa Sligt? – Enkelte Vink vil jeg allerede være istand til at meddele. Til et fuldstændigt Kys fordres, at det er en Pige og en Mand, der ere de Handlende. Et Mandfolke-Kys er smagløst, eller har, hvad værre er, Afsmag. – Dernæst troer jeg, at et Kys kommer Ideen nærmere, naar en Mand kysser en Pige end naar en Pige kysser en Mand. Hvor der i Aarenes Løb er tilveiebragt Indifferens i dette Forhold, der har Kysset tabt sin Betydning. Dette gjelder om det ægteskabelige Huskys, hvormed Ægtefolkene, i Mangel af Serviet, tørre hinanden om Munden, idet der siges: velbekom's. – Er Afstanden i Alder meget stor, saa ligger Kysset udenfor Ideen. Jeg erindrer i en Pigeskole i en af Provindserne havde den ældste Classe en egen Terminus: at kysse Justitsraaden, et Udtryk, hvormed de forbandt en intet mindre end behagelig Forestilling. Oprindelsen til denne Terminus var følgende. Lærerinden havde en Svoger, der levede i Huset hos hende, han havde været Justitsraad, var en ældre Mand, og tog sig nu i Kraft heraf den Frihed at ville kysse de unge Piger. – Kysset maa være Udtryk for en bestemt Lidenskab. Naar en Broder og en Søster, der ere Tvillinger, kysse hinanden, er det Kys intet rigtigt Kys. Om et Kys, der falder af i en Juleleeg, gjelder det Samme, item om et stjaalent Kys. Et Kys er en symbolsk Handling, der Intet har at betyde, naar den Følelse, den skal betegne, ikke er tilstede, og denne Følelse kan kun være tilstede under bestemte Forhold. – Vil man gjøre et Forsøg paa at inddele Kysset, saa kan man tænke sig flere Inddelings-Principer. Man kan inddele det med Hensyn til Lyden. Destoværre strækker Sproget her ikke til i Forhold til mine Iagttagelser. Jeg troer ikke, at al Verdens Sprog havde et fornødent Forraad af Onomatopoietika, for at betegne de Forskjelligheder, jeg blot i min Onkels Huus har lært at kjende. Snart er det smækkende, snart hvislende, snart kladskende, snart knaldende, snart drønende, snart fuldt, snart huult, snart som Kattun og saa videre, og saa videre – Man kan inddele Kysset med Hensyn til Berøringen, i det tangerende Kys, eller Kysset en passant, og det cohærerende. – Man kan inddele det med Hensyn til Tiden i det korte og i det lange. Med Hensyn til Tiden gives der ogsaa en anden Inddeling, og denne er egentlig den eneste, der har behaget mig. Man gjør da en Forskjel mellem det første Kys og alle andre. Det, hvorpaa der her reflekteres, er incommensurabelt for det, der kommer tilsyne ved de øvrige Inddelinger, det er indifferent mod Lyden, Berøringen, Tiden i Almindelighed. Det første Kys er imidlertid qualitativt forskjelligt fra alle andre. Det er der kun faa Mennesker, der tænke paa, det var da Synd Andet end at der var Een, der tænker derover.

[…]

Min Cordelia!

Et godt Svar er som et sødt Kys siger Salomo. Du veed jeg er slem til at spørge; jeg maa næsten høre ilde derfor. Det kommer deraf, at man ikke forstaaer, hvorom jeg spørger; thi Du og Du alene forstaaer, hvorom jeg spørger, og Du og Du alene forstaaer at svare, og Du og Du alene forstaaer at give et godt Svar; thi et godt Svar er som et sødt Kys, siger Salomo.

Din Johannes.

[…]

Har jeg nu i mit Forhold til Cordelia bestandig været min Pagt tro? Det vil sige, min Pagt med det Æsthetiske; thi det er det, der gjør mig stærk, at jeg bestandig har Ideen paa min Side. Det er en Hemmelighed ligesom Samsons Haar, som ingen Dalila skal fravriste mig. Slet og ret at bedrage en Pige, det vil jeg vist ikke have Udholdenhed til; men det, at Ideen er med i Bevægelse, at det er i dens Tjeneste jeg handler, til dens Tjeneste jeg helliger mig, det giver mig Strenghed mod mig selv, Afholdenhed fra enhver forbuden Nydelse. Er det Interessante altid blevet bevaret? Ja, det tør jeg sige frit og aabent i denne hemmelige Samtale. Forlovelsen selv var det Interessante netop derved, at den ikke gav det, man i Almindelighed forstaaer ved det Interessante. Den bevarede det Interessante netop derved, at den ydre Tilsyneladelse var i Modsigelse med det indre Liv. Havde jeg været hemmelig forbunden med hende, saa havde det kun været interessant i første Potens. Dette er derimod det Interessante i anden Potens og derfor for hende først det Interessante. Forlovelsen brister, men derved, at hun selv hæver den, for at svinge sig op i en høiere Sphære. Saaledes skal det være; dette er nemlig den Form af det Interessante, der mest vil beskæftige hende.

[…]

den 25. September

Hvorfor kan en saadan Nat ikke vare længere? Kunde Alektryon forglemme sig, hvorfor kan da Solen ikke være medlidende nok dertil? Dog nu er det forbi, og jeg ønsker aldrig mere at see hende. Naar en Pige har bortgivet Alt, da er hun svag, da har hun tabt Alt; thi Uskyld er hos Manden et negativt Moment, hos Qvinden er det hendes Væsens Gehalt. Nu er al Modstand umulig, og kun saa længe den er til, er det skjønt at elske, naar den er ophørt, er det Svaghed og Vane. Jeg ønsker ikke at mindes om mit Forhold til hende; hun har tabt Duften, og de Tider ere forbi, da en Pige af Smerte over sin troløse Elsker forvandles til en Heliotrop. Afskeed vil jeg ikke tage med hende; Intet er mig modbydeligere end Qvindegraad og Qvindebønner, der forandre Alt og dog egentlig ikke have Noget at betyde. Jeg har elsket hende; men fra nu af kan hun ikke mere beskæftige min Sjæl. Hvis jeg var en Gud, da vilde jeg gjøre for hende, hvad Neptun gjorde for en Nymphe, forvandle hende til en Mand.

Det var dog virkeligt værd at vide, om man ikke skulde være istand til saaledes at digte sig ud af en Pige, at man kunde gjøre hende saa stolt, at hun bildte sig ind, at det var hende, der var kjed af Forholdet. Det kunde blive et ret interessant Efterspil, der i og for sig kunde have psychologisk Interesse og ved Siden deraf berige Een med mange erotiske Iagttagelser.

\workheading{Enten – Eller. Et Livs-Fragment, Anden Deel}{1843}

\worknote{Antologiens ''Either/Or … II'' (B.'s Papirer — Assessor Wilhelms breve). Hong-udvalget bringer uddrag af de to breve: ''Ægteskabets æsthetiske Gyldighed'' og ''Ligevægten mellem det Æsthetiske og Ethiske''. Begge er stærkt sammensatte af adskilte steder; udeladelser markeret med […]. Slutstykket ''Ultimatum'' indgår ikke i antologien.}

\sectionheading{Ægteskabets æsthetiske Gyldighed (uddrag)  ·  SV¹ II, 5–127}

Min Ven!

De Linier, Dit Øie her først falder paa, ere skrevne sidst. Hensigten med dem er endnu engang at gjøre et Forsøg paa at tvinge den udførligere Undersøgelse, der herved tilstilles Dig, ind i et Brevs Form. Disse Linier svare da til de sidste Linier og danne tilsammen en Convolut og antyde saaledes paa en udvortes Maade, hvad indvortes Beviser paa mange Maader ville overtyde Dig om, at det er et Brev, Du læser. Denne Tanke, at det var et Brev, jeg skrev til Dig, har jeg ikke villet opgive, deels fordi min Tid ikke har tilladt den omhyggeligere Udarbeidelse, en Afhandling kræver, deels fordi jeg nødig vilde gaae Glip af den Leilighed til mere formanende og indtrængende at tiltale Dig, som Brevformen medfører. Du er altfor dreven i den Kunst, ganske i Almindelighed at kunne tale om Alt uden at lade Dig personligt berøre deraf, til at jeg skulde friste Dig ved at sætte Din dialektiske Kraft i Bevægelse. Du veed nok, hvorledes Propheten Nathan bar sig ad med Kong David, da han vel vilde forstaae Lignelsen, Propheten havde fremsat, men ikke vilde forstaae, at den gjaldt ham. Nathan tilføiede da for en Forsigtigheds Skyld: Du est Manden, Herre Konge. Saaledes har jeg ogsaa bestandig søgt at minde Dig om, at det er Dig, der tales om, og Dig der tales til.

[…]

Med Hensyn til det individuelle Liv er der to Arter af Historie, udvortes og indvortes. Det er to Arter af Strømninger, hvis Bevægelse er modsat. Den første har igjen to Sider i sig. Individet har ikke det, han stræber efter, og Historien er den Strid, i hvilken han erhverver det. Eller Individet har det, men kan dog ikke komme i Besiddelse deraf, fordi der bestandig er noget Udvortes, der vil forhindre ham heri. Historien er da den Strid, i hvilken han beseirer disse Hindringer. Den anden Art af Historie begynder med Besiddelsen, og Historien er den Udvikling, gjennem hvilken han erhverver Besiddelsen. Da nu i første Tilfælde Historien er udvortes, og det, hvortil der stræbes, ligger udenfor, saa har Historien ikke sand Realitet, og den digteriske og kunstneriske Fremstilling handler ganske rigtigt i at forkorte denne og haste til det intensive Moment. Lad os, for at blive ved, hvad vi nærmest have at gjøre med, tænke os en romantisk Kjærlighed. Tænk Dig altsaa en Ridder, der har nedlagt 5 Vildsviin, 4 Dværge, befriet 3 fortryllede Prindser, Brødre til den Prindsesse, han tilbeder. Dette har i den romantiske Tankegang sin fuldkomne Realitet. Kunstneren og Digteren er det imidlertid slet ikke af Vigtighed, om der er 5 eller kun 4. Kunstneren er i det Hele mere indskrænket end Digteren, men selv denne vil det ikke interessere at fortælle omstændeligt, hvorledes det gik til med hvert enkelt Vildsviins Nedlæggelse. Han haster til Momentet. Han indskrænker maaskee Antallet, concentrerer Besværligheder og Farer i digterisk Intensitet og iler til Momentet, til Besiddelses-Momentet. Den hele historiske Succession er ham af mindre Vigtighed. Hvor der derimod er Tale om indvortes Historie, der er hvert enkelt lille Moment af yderste Vigtighed. Den indvortes Historie er først den sande Historie, men den sande Historie kæmper med det, der er Livs-Principet i Historien – med Tiden, men naar man kæmper med Tiden, saa har derved netop det Timelige og ethvert lille Moment sin store Realitet. Overalt hvor Individualitetens indre Blomstring ikke er begyndt, hvor Individualiteten endnu er lukket, der bliver der Tale om udvortes Historie. Saasnart derimod denne saa at sige springer ud, saa begynder den indre Historie. Tænk nu paa, hvad vi gik ud fra, Forskjellen mellem den erobrende og den besiddende Natur. Den erobrende Natur er bestandig udenfor sig, den besiddende i sig selv, derfor faaer den første udvortes Historie, den anden indvortes. Men da den udvortes Historie netop uden Skade lader sig concentrere, saa er det naturligt, at Kunst og Poesi nærmest vælge den, og altsaa igjen vælge den uaabnede Individualitet og Alt hvad der tilhører den til Fremstilling. Nu siger man vel, at Kjærlighed aabner Individualiteten, men ikke naar Kjærligheden opfattes saaledes som det skeer i Romantiken, da bringes den blot til det Punkt, at den skal aabne sig, og der ender man, eller den er ifærd med at aabne sig, men bliver afbrudt. Men som den udvortes Historie og den lukkede Individualitet nærmest vil blive Gjenstand for den kunstneriske og digteriske Fremstilling, saa vil Alt, hvad der udgjør en saadan Individualitets Indhold ogsaa blive Gjenstand derfor. Men dette er i Grunden Alt, hvad der tilhører det naturlige Menneske. Et Par Exempler. Stolthed lader sig ypperligt fremstille; thi det Essentielle i Stolthed er ikke Succession, men Intensitet i Momentet. Ydmyghed lader sig vanskeligt fremstille, fordi den netop er Succession, og medens Betragteren ikke behøver mere end at see Stoltheden i sin Culmination, saa fordrer han egentlig i det andet Tilfælde, det som Poesi og Kunst ikke kunne give, at see den i sin stadige Tilbliven; thi dette hører Ydmygheden væsentlig til, at den stadig bliver; og naar man viser ham den i dens ideale Moment, saa savner han Noget, fordi han føler, at dens sande Idealitet bestaaer ikke i at den er ideal i Momentet, men at den er bestandig. Den romantiske Kjærlighed lader sig ypperlig fremstille i Momentet, den ægteskabelige ikke; thi en idealiseret Ægtemand er ikke en, der engang i sit Liv er det, men som hver Dag er det. Naar jeg vil fremstille en Helt, der erobrer Riger og Lande, da lader det sig ypperligt fremstille i Momentet, men en Korsdrager, der hver Dag tager sit Kors op, lader sig aldrig fremstille hverken i Poesi eller Kunst, fordi Pointet er, at han hver Dag gjør det. Naar jeg vil tænke mig en Helt, der sætter sit Liv til, saa lader det sig ypperlig concentrere i Momentet, derimod ikke det at døe hver Dag, fordi Hovedsagen er, at det skeer hver Dag. Mod lader sig ypperligt concentrere i Momentet, Taalmod ikke, netop fordi Taalmod strider mod Tiden. Du vil sige, at Kunsten dog har fremstillet Christus som Billede paa Taalmod, som bærende al Verdens Synd, at religiøse Digtninger have concentreret al Livets Bitterhed i een Kalk og ladet et Individ i eet Moment udtømme den. Det er sandt; men det er fordi man næsten rummeligt har concentreret den. Den, der derimod veed lidt Beskeed om Taalmod, veed meget godt, at dens egentlige Modsætning ikke er Lidelsens Intensitet (thi da nærmer det sig mere til Mod) men Tiden, og at den sande Taalmod er den, der viser sig stridende mod Tid, eller egentlig er Langmod, men Langmod lader sig ikke kunstnerisk fremstille; thi dets Point er incommensurabelt for Kunst, det lader sig heller ikke digte; thi det fordrer Tidens Langvarighed.

Hvad jeg nu her videre vil fremføre, det kan Du ansee for en stakkels Ægtemands ringe Offer paa Æsthetikens Alter, og om Du og alle Æsthetikens Præster vil forsmaae det, saa skal jeg vel vide at trøste mig, og det saa meget mere som det, jeg bringer, ikke er et Skuebrød, ene spiseligt for Præsterne, men et Huusbagerbrød, der som al Huuskost er simpel og ukrydret, men sund og nærende.

Naar man ligesaa meget dialektisk som historisk forfølger det Æsthetisk-Skjønnes Udvikling, vil man finde, at Retningen i denne Bevægelse er fra Rummets til Tidens Bestemmelser, og at Kunstens Fuldkommengjørelse afhænger af den successive Mulighed af mere og mere at løsrive sig fra Rummet og bestemme sig hen til Tiden. Heri ligger Overgangen og Betydningen af Overgangen fra Skulptur til Maleri, saaledes som Schelling tidligt har peget hen herpaa. Musikken har Tiden til sit Element, men vinder ingen Bestaaen i den, dens Betydning er den bestandige Forsvinden i Tiden, den toner i Tiden, men fortoner sig tillige og har ingen Bestaaen. Poesien endelig er den fuldkomneste af alle Kunster og derfor ogsaa den, der mest veed at gjøre Tidens Betydning gjeldende. Den behøver ikke i den Forstand at indskrænke sig til Momentet, som Maleriet gjør det, den forsvinder heller ikke i den Forstand sporløst som Musikken. Men uagtet dette nødsages ogsaa denne, som vi have seet, til at concentrere sig i Momentet. Den har derfor sin Grændse, og kan, som ovenfor viist, ikke fremstille det, hvis Sandhed netop er den timelige Succession. Og dog er det, at Tiden bliver gjort gjeldende, ikke nogen Forkleinelse af det Æsthetiske, tvertimod bliver det æsthetiske Ideal rigere og fyldigere, jo mere dette skeer. Det Æsthetiske altsaa, der bliver incommensurabelt endog for Poesiens Fremstilling, hvorledes lader det sig fremstille? Svar: derved at det leves. Det faaer derved en Lighed med Musik, der blot er, fordi den bestandig gjentages, blot er i Udførelsens Øieblik. Derfor gjorde jeg i det Foregaaende opmærksom paa den fordærvelige Forvexling af det Æsthetiske og det, som æsthetisk lader sig fremstille i digterisk Reproduktion. Alt, hvad jeg nemlig her taler om, lader sig visselig æsthetisk fremstille, men ikke i digterisk Reproduktion, men derved, at man lever det, i Virkelighedens Liv realiserer det. Saaledes hæver Æsthetiken sig selv og forsoner sig med Livet; thi som i een Forstand Poesi og Kunst netop er en Forsoning med Livet, saa ere de i en anden Forstand Fjendskab til Livet, fordi de kun forsone een Side af Sjælen. Her er jeg ved det Høieste i det Æsthetiske. Og i Sandhed, den, der har Ydmyghed og Mod nok til her at lade sig æsthetisk forklare, den, der føler sig med som en Person i det Skuespil, Guddommen digter, hvor Digteren og Souffleuren ikke ere forskjellige Personer, hvor Individet, som den øvede Skuespiller, der har levet sig ind i sin Charakteer og sin Replik, ikke forstyrres af Souffleuren, men føler, at det, der tilhvidskes ham, er det, han selv vil sige, saa det næsten bliver tvivlsomt, om han lægger Souffleuren Ordene i Munden, eller Souffleuren ham, den, der i dybeste Forstand føler sig paa eengang digtende og digtet, der i det Øieblik, han føler sig digtende, har Replikkens oprindelige Pathos, i det Øieblik, han føler sig digtet, det erotiske Øre, der opfanger enhver Lyd, han og først han har realiseret det Høieste i Æsthetiken. Men denne Historie, der viser sig som incommensurabel selv for Poesien, det er den indvortes Historie. Denne har Ideen i sig og netop derfor er den æsthetisk. Den begynder derfor, som jeg udtrykte det, med Besiddelsen, og dens Fortgang er Erhvervelsen af denne Besiddelse. Den er en Evighed, i hvilken det Timelige ikke er forsvunden som et idealt Moment, men i hvilket det som realt Moment bestandig er nærværende. Naar saaledes Taalmod i Taalmod erhverver sig selv, saa er det indvortes Historie.

Lad os nu see hen paa Forholdet mellem den romantiske og den ægteskabelige Kjærlighed; thi Forholdet mellem den erobrende og den besiddende Natur vil slet ingen Vanskeligheder formaae at frembyde. Den romantiske Kjærlighed bliver bestandig abstrakt i sig selv, og naar den ingen udvortes Historie kan faae, saa lurer Døden allerede paa den, fordi dens Evighed er illusorisk. Den ægteskabelige Kjærlighed begynder med Besiddelsen, og faaer indre Historie. Den er trofast, det er den romantiske Kjærlighed ogsaa, men see nu Forskjellen. Den trofaste romantiske Elsker han venter for Exempel i 15 Aar, nu kommer da Øieblikket, der belønner ham. Her seer Poesien meget rigtigt, at de 15 Aar ypperligen lade sig concentrere, den haster nu til Momentet. En Ægtemand er trofast i 15 Aar, og dog har han i de 15 Aar havt Besiddelsen, han har altsaa i denne lange Succession bestandig erhvervet den Trofasthed, han besad, da jo den ægteskabelige Kjærlighed i sig har den første Kjærlighed og derved dennes Trofasthed. Men en saadan ideal Ægtemand lader sig ikke fremstille; thi Pointet er Tiden i sin Extension. Ved Enden af de 15 Aar er han tilsyneladende slet ikke kommen videre end han var i Begyndelsen, og dog har han levet i høi Grad æsthetisk. Hans Besiddelse har ikke været ham en død Eiendom, men han har bestandig erhvervet sin Besiddelse. Han har ikke kæmpet med Løver og Trolde, men med den farligste Fjende, som er Tiden. Men nu kommer ikke Evigheden bag efter som for Ridderen; men han har havt Evigheden i Tiden, bevaret Evigheden i Tiden. Han først har derfor seiret over Tiden; thi Ridderen kan man sige om, han har dræbt Tiden, som man jo stedse ønsker at slaae Tiden ihjel, naar den ingen Realitet har for Een; men dette er aldrig den rigtige Seir. Ægtemanden har som en sand Seierherre ikke dræbt Tiden, men frelst og bevaret den i Evigheden. Den Ægtemand, der gjør det, han lever i Sandhed poetisk, han løser den store Gaade, at leve i Evigheden og dog høre Stueuhret slaae, saaledes, at dets Slag ikke forkorter, men forlænger hans Evighed, en Modsigelse, der er ligesaa dyb, men langt herligere end den, der indeholdes i den bekjendte Situation, som skyldes Middelalderen, der fortæller om en Ulykkelig, at han vaagnede op i Helvede og raabte, hvad er Klokken, hvorpaa Djævelen svarede en Evighed. Og lader Sligt sig nu end ikke kunstnerisk fremstille, saa være det Din Trøst som det er min, at det Høieste og Skjønneste i Livet, det skal man ikke læse om, ikke høre om, ikke see, men om man saa vil, leve. Naar jeg derfor gjerne indrømmer, at den romantiske Kjærlighed meget bedre egner sig til kunstnerisk Fremstilling end den ægteskabelige, saa skal dermed ingenlunde være sagt, at denne er mindre æsthetisk end hiin, tvertimod, den er mere æsthetisk. Der forekommer i en af den romantiske Skoles genialeste Fortællinger en Person, der ikke har Lyst til, som de Andre, med hvilke han lever, at digte, fordi det er Tidsspilde og berøver ham den sande Nydelse; han vil derimod leve. Havde han nu havt en rigtigere Forestilling om, hvad det er at leve, saa havde han været min Mand.

Den ægteskabelige Kjærlighed har da sin Fjende i Tiden, sin Seir i Tiden, sin Evighed i Tiden, saa den stedse vilde have sin Opgave, om jeg end tænkte mig alle saakaldte udvortes og indvortes Anfægtelser borte. I Almindelighed har den ogsaa disse, men naar man ret vil opfatte dem, maa man passe paa to Ting, at de bestandig ere Bestemmelser ind efter, og at de bestandig have Tidens Bestemmelse i sig. Ogsaa af denne Grund indsees let, at denne Kjærlighed ikke lader sig fremstille. Den drager sig bestandig ind efter, og hendrager (i god Forstand) sig i Tiden; men hvad der skal fremstilles ved Reproduktion, det maa lade sig fremlokke, og dets Tid maa lade sig forkorte. Det vil Du yderligere overbevise Dig om, ved at overveie de Prædikater, man maa bruge om den ægteskabelige Kjærlighed. Den er trofast, bestandig, ydmyg, taalmodig, langmodig, overbærende, oprigtig, nøisom, aarvaagen, varagtig, villig, glad. Alle disse Dyder have den Egenskab, at de i Individet ere Bestemmelser ind efter. Individet kæmper ikke mod udvortes Fjender, men udkæmper sig selv, sin Kjærlighed af sig selv; og de have Bestemmelse af Tid; thi deres Sandhed bestaaer ikke i, at de ere engang for alle, men at de bestandig ere. Og ved disse Dyder erhverves ikke noget Andet, kun de selv erhverves. Den ægteskabelige Kjærlighed er derfor paa eengang, hvad Du spotviis ofte har kaldt den, den dagligdagske, og tillige den guddommelige (i græsk Forstand), og er den guddommelige derved, at den er den dagligdags. Den ægteskabelige Kjærlighed kommer ikke med udvortes Tegn, ikke som den rige Fugl med Susen og Brusen, men den er en stille Aands uforkrænkelige Væsen.

\sectionheading{Ligevægten mellem det Æsthetiske og Ethiske i Personlighedens Udarbeidelse (uddrag)  ·  SV¹ II, 141–297}

Min Ven!

Hvad jeg saa ofte har sagt Dig, det siger jeg endnu engang, eller rettere det tilraaber jeg Dig: enten – eller; aut – aut; thi et enkelt aut, der berigtigende træder til, gjør ikke Sagen klar, da det, hvorom her er Talen, er for betydningsfuldt til, at man kan nøies med en Deel deraf, i sig selv for sammenhængende til, at kunne besiddes partielt. Der gives Livsforhold, paa hvilke det vilde være en Latterlighed eller en Art Sindssvaghed at anvende et: enten – eller; men der gives ogsaa Mennesker, hvis Sjæl er for dissolut til at fatte, hvad der ligger i et saadant Dilemma, hvis Personlighed mangler Energi til med Pathos at kunne sige: enten – eller. Paa mig har disse Ord altid gjort et stærkt Indtryk, og gjøre det endnu, især naar jeg saaledes nævner dem blankt og bart, hvori jo ligger en Mulighed af at sætte de forfærdeligste Modsætninger i Bevægelse. De virke paa mig som en Besværgelses-Formular, og min Sjæl bliver i høi Grad alvorlig, stundom næsten rystet. Jeg tænker paa en tidlig Ungdom, hvor jeg uden ret at fatte, hvad det er at vælge i Livet, med barnlig Tillid hørte paa Ældres Tale, og Valgets Øieblik blev mig høitideligt og ærværdigt, uagtet jeg i at vælge kun fulgte en Andens Anviisning. Jeg tænker paa de Øieblikke i et senere Liv, hvor jeg stod paa Skilleveien, hvor min Sjæl modnedes i Afgjørelsens Stund. Jeg tænker paa de mange mindre vigtige men for mig ikke ligegyldige Tilfælde i Livet, hvor det gjaldt om at vælge; thi om der end kun er eet Forhold, hvor dette Ord har sin absolute Betydning, hver Gang der nemlig paa den ene Side viser sig Sandhed, Retfærdighed og Hellighed; paa den anden Side Lyst og Tilbøieligheder, og dunkle Lidenskaber og Fortabelse, saa er det dog altid af Vigtighed ogsaa i Ting, hvor det i og for sig er uskyldigt, hvilket man vælger, at vælge rigtigt, at prøve sig selv, at man ikke engang med Smerte skal begynde et Tilbagetog til det Punkt, man gik ud fra, og takke Gud, om man ikke har mere at bebreide sig end at have spildt sin Tid. I daglig Tale bruger jeg disse Ord som Andre bruge dem, og det vilde jo ogsaa være et taabeligt Pedanteri at aflade dermed; men dog kan det stundom hænde mig, at det falder mig paa, at jeg har brugt dem om aldeles ligegyldige Ting. De afføre sig da den ringe Dragt, jeg glemmer de ubetydelige Tanker, de adskilte, de træde op for mig i al deres Værdighed, i deres Ornat. Som en Øvrigheds-Person til daglig Brug viser sig i civil Dragt og blander sig mellem Mængden, uden videre Forskjel, saaledes gaaer det hine Ord i den daglige Tale; naar han derimod træder op i sin Myndighed da adskiller han sig fra Alle. Som en saadan Øvrigheds-Person, jeg kun er vant til at see ved høitidelige Leiligheder, vise da disse Ord sig, og min Sjæl bliver altid alvorlig. Og skjøndt nu mit Liv til en vis Grad har sit: enten – eller bag ved sig, saa veed jeg dog meget vel, at der endnu kan møde mangt et Tilfælde, hvor det vil have sin fulde Betydning. Jeg haaber imidlertid, at disse Ord skulle finde mig idetmindste værdigen stemt, idet de standse mig paa min Vei, og jeg haaber, at det skal lykkes mig at vælge det Rette; men i ethvert Tilfælde skal jeg bestræbe mig for at vælge med uskrømtet Alvor; jeg tør da saa idetmindste trøste mig til, at jeg snarere vil komme bort fra min Afvei.

Og nu Du, Du bruger jo ofte nok dette Ord, ja det er næsten blevet Dig et Mundheld, hvad Betydning har det for Dig? Slet ingen. For Dig er det, at jeg skal erindre om Dine egne Udtryk: et Blink, en Haandevending, et coup des mains, et Abracadabra. Ved enhver Leilighed veed Du at anbringe det, og det bliver heller ikke uden Virkning; paa Dig virker det nemlig som en stærk Drik paa en Nervesvag; Du bliver fuldkommen beruset i, hvad Du selv kalder, høiere Galenskab. »Deri indeholdes al Leveviisdom, men aldrig har endnu noget Menneske foredraget den med saa megen Fynd, som var det en Gud i en Popans Skikkelse, der talte til den lidende Menneskehed, som hiin store Tænker og sande Livs-Philosoph, der sagde til en Mand, der havde revet hans Hat paa Gulvet: tager Du den op, faaer Du Prygl, tager Du den ikke op, faaer Du ogsaa Prygl, nu kan Du vælge.« Du har Din store Glæde af at »trøste« Folk, naar de henvende sig til Dig i kritiske Tilfælde; Du hører paa deres Udvikling og nu siger Du: ja jeg indseer det nu fuldelig, der er to Tilfælde mulige, man kan enten gjøre det eller det, min oprigtige Mening og mit venskabelige Raad er som følger, gjør det eller gjør det ikke, Du vil fortryde begge Dele. Dog den, der spotter Andre, han spotter sig selv, og det er ikke for Intet, men en dyb Spot over Dig, et sørgeligt Beviis paa, hvor ledeløs Din Sjæl er, at Din Livsanskuelse concentrerer sig i een eneste Sætning »jeg siger bare enten – eller.« Dersom det nu virkelig var Din Alvor, saa var der Intet med Dig at udrette, man maatte lade Dig gaae for hvad Du var, og beklage, at Tungsind eller Letsind havde svækket Din Aand. Nu derimod, da man vel veed, at det ikke forholder sig saaledes, fristes man ikke til at beklage Dig, men vel til at ønske, at Dit Livs Forhold engang maae skrue Dig ind i sine Bøiler, og tvinge Dig til at lade Det komme frem, der boer i Dig, maa begynde den skarpere Examination, der ikke lader sig nøie med Snak eller med Vittigheder. Livet er en Masquerade, oplyser Du, og dette er Dig et uudtømmeligt Stof til Moro, og endnu er det ikke lykkedes Nogen at kjende Dig; thi enhver Aabenbarelse er altid et Bedrag, kun paa den Maade kan Du aande, og forhindre, at Folk ikke trænge ind paa Dig og besvære Respirationen. Deri har Du Din Virksomhed, at bevare Dit Skjul, og dette lykkes Dig; thi Din Maske er den gaadefuldeste af alle; Du er nemlig Ingenting, og er bestandig blot i Forhold til Andre, og hvad Du er, er Du ved dette Forhold. En kjelen Hyrdinde rækker Du smægtende Haanden, og er i samme Nu masqueret i al mulig Schäfer-Sentimentalitet; en ærværdig geistlig Fader bedrager Du med et Broderkys og saa videre Du selv er Ingenting, en gaadefuld Skikkelse, paa hvis Pande staaer enten eller; »thi dette er mit Valgsprog, og disse Ord ere ikke, som Grammatikerne troe, disjunktive Conjunktioner, nei de høre uadskilleligt sammen og bør derfor skrives i eet Ord, da de i Forening danne en Interjection, jeg tilraaber Menneskeheden ligesom man raaber Hep efter en Jøde.« Uagtet nu enhver Yttring fra Dig af den Art paa mig bliver uden Virkning, eller forsaavidt den virker Noget, i det Høieste virker til at fremkalde en retfærdig Indignation, saa vil jeg dog for Din egen Skyld svare Dig: veed Du da ikke, at der kommer en Midnatstime, hvor Enhver skal demasquere sig, troer Du, at Livet altid lader sig spøge med, troer Du, at man kan liste bort lidt før Midnat for at undgaae det? Eller forfærdes Du ikke herved? Jeg har seet Mennesker i Livet, der saa længe bedrog Andre, at deres sande Væsen derved tilsidst ikke kunde aabenbare sig; jeg har seet Mennesker, der saa længe legede Skjul, at tilsidst Vanvid gjennem dem ligesaa modbydeligt paanødte Andre deres lønlige Tanker, som de hidtil stolt havde skjult dem. Eller kan Du tænke Dig noget Forfærdeligere end at det endte med, at Dit Væsen opløste sig i en Mangfoldighed, at Du virkelig blev Flere, blev ligesom hine ulykkelige Dæmoniske en Legio, og Du saaledes havde tabt det Inderste, Helligste i et Menneske, Personlighedens bindende Magt? Du skulde i Sandhed ikke spøge med, hvad der ikke blot er alvorligt, men rædsomt. I ethvert Menneske er der Noget, der til en vis Grad forhindrer ham i at blive sig selv fuldelig gjennemsigtig; men dette kan i saa høi Grad være Tilfælde, han kan saa uforklarligt være indflettet i Livs-Forhold, der ligge ud over ham selv, at han næsten ikke kan aabenbare sig; men den, der ikke kan aabenbare sig, kan ikke elske, og den, der ikke kan elske, han er den Ulykkeligste af Alle. Og Du gjør i Kaadhed det Samme, Du øver Dig i den Kunst at blive gaadefuld for Alle. Min unge Ven! sæt der var Ingen, der brød sig om at gjætte Din Gaade, hvad Glæde havde Du da deraf? Men fremfor Alt for Din egen Skyld, for Din Frelses Skyld; thi jeg veed ingen Sjæls-Tilstand, man bedre kan betegne som Fortabelse; stands denne vilde Flugt, denne Tilintetgjørelsens Lidenskab, der raser i Dig; thi det er det, Du vil, Du vil tilintetgjøre Alt; Du vil mætte Tvivlens Hunger i Dig paa Tilværelsen. Dertil danner Du Dig, dertil forhærder Du Dit Sind; thi det tilstaaer Du gjerne, Du duer til Intet, kun dette forlyster Dig, at gaae 7 Gange omkring Tilværelsen, og blæse paa Basunen, og derpaa lade det Hele gaae under, at Din Sjæl kan blive beroliget ja veemodig, at Du kan faae Ekcho kaldet frem; thi Ekcho lyder kun i Tomhed.

Dog ad denne Vei kommer jeg vel ikke videre med Dig, desuden er mit Hoved om Du saa vil for svagt, til at kunne udholde, eller som jeg vil, for stærkt til at finde Behag i, at det bestandig svimler for mit Øie. Jeg vil derfor begynde Sagen fra en anden Side. Tænk Dig et ungt Menneske, i den Alder da Livet ret begynder at faae Betydning for ham; han er sund, reen, glad, aandelig begavet, selv riig paa Haab, Enhvers Haab, der kjender ham, tænk Dig, ja det er haardt at jeg skal sige det, at han tog Feil af Dig, at han troede, Du var et alvorligt, prøvet, erfarent Menneske, hos hvem man sikkert kunde søge Oplysning om Livets Gaader; tænk Dig, at han henvendte sig til Dig med den elskelige Tillid, der er Ungdommens Prydelse, med den uafviselige Fordring, der er Ungdommens Ret – hvad vilde Du svare ham? Vilde Du svare: ja jeg siger bare enten – eller; det vilde Du vel neppe? Vilde Du, som Du pleier at udtrykke Dig, naar Du vil betegne Din Modbydelighed for, at Andre besvære Dig med deres Hjertes-Anliggender, stikke Hovedet ud af Vinduet og sige: Huus forbi; eller vilde Du behandle ham som Andre, der ville beraadføre sig med Dig, eller søge Oplysning hos Dig, hvilke Du afviser ligesom de, der kræve Præstepenge med de Ord, at Du kun er Logerende i Livet, ikke en bosiddende Mand og Familiefader? Dette vilde Du vel heller ikke. Et ungt Menneske, aandelig begavet, er Noget, Du kun sætter altfor megen Priis paa. Men Dit Forhold til ham var ikke ganske som Du ellers ønskede det, det var ikke et tilfældigt Sammenstød, der havde bragt Dig i Berøring med ham, Din Ironi var ikke fristet. Skjøndt han var den Yngre, Du den Ældre, saa havde han dog ved sin ædle Ungdommelighed gjort Øieblikket alvorligt. Ikke sandt, Du vilde selv blive ung, Du vilde føle, at der ligger noget Skjønt i at være ung, men ogsaa noget saare Alvorligt, at det ingenlunde er en ligegyldig Sag, hvorledes man bruger sin Ungdom, at der ligger et Valg for Een, et virkeligt enten – eller? Du vilde føle, at det, det dog kommer an paa, ikke saa meget er at danne sin Aand, som at modne sin Personlighed. Din Godmodighed, Din Sympathi var sat i Bevægelse, ud af den vilde Du tale til ham; Du vilde styrke hans Sjæl, bestyrke ham i den Tillid, han havde til Verden, Du vilde forsikkre ham, at der er en Magt i et Menneske, der kan trodse hele Verden, Du vilde lægge ham ret stærkt paa Hjerte at bruge Tiden. Alt dette kan Du gjøre, og naar Du vil kan Du gjøre det smukt. Men agt nu vel paa, hvad jeg vil sige Dig, unge Menneske; thi skjøndt Du ikke er ung, nødsages man dog altid til at kalde Dig saa. Hvad gjorde Du nu her? Du anerkjendte, hvad Du ellers ikke vil anerkjende, Betydningen af et enten – eller, og hvorfor, fordi Din Sjæl var bevæget af Kjærlighed til det unge Menneske; og dog bedrog Du jo ham paa en Maade, thi han vil maaskee træffe sammen med Dig til andre Tider, hvor det ingenlunde er Din Leilighed at anerkjende det. Der seer Du en sørgelig Følge af, at et Menneskes Væsen ikke kan harmonisk aabenbare sig. Du troede at gjøre det Bedste, og dog har Du maaskee skadet ham; maaskee han snarere vilde have kunnet holde sig ligeoverfor Din Mistillid til Livet, end finde Hvile i den subjektiv svigagtige Tillid, Du bibragte ham. Tænk Dig, at Du efter nogle Aars Forløb atter traf sammen med hiint unge Menneske; han var livlig, vittig, aandrig, driftig i sin Tanke, kjæk i sit Udtryk; men Dit fine Øre opdagede let Tvivlen i hans Sjæl, Du fattede en Mistanke om, at ogsaa han var kommen til den tvetydige Viisdom: jeg siger bare enten – eller; ikke sandt det vilde gjøre Dig ondt for ham, Du vilde føle, at han havde tabt Noget, og noget saare Væsentligt. Men over Dig selv vil Du ikke sørge, Du er tilfreds, ja stolt af Din tvetydige Viisdom, ja saa stolt af den, at Du ikke kan tillade en Anden at dele den, da Du vil være ene om den. Og dog finder Du i en anden Henseende det beklageligt, og det er Din oprigtige Mening, at det er beklageligt, at hiint unge Menneske var kommet til samme Viisdom. Hvilken uhyre Modsigelse! Dit hele Væsen modsiger sig selv. Men ud af denne Modsigelse kan Du ikkun komme ved et enten – eller; og jeg, som elsker Dig oprigtigere end Du elskede hiint unge Menneske, jeg, der i mit Liv har erfaret Valgets Betydning, jeg ønsker Dig til Lykke med, at Du endnu er saa ung, at om Du end altid vil gaae Glip af Noget, Du dog, hvis Du har eller rettere Du vil have Energi dertil, kan vinde, hvad der er Hovedsagen i Livet, vinde Dig selv, erhverve Dig selv.

Dersom nu et Menneske bestandig kunde holde sig paa Valg-Øieblikkets Spidse, dersom han kunde ophøre at være et Menneske, dersom han i sit inderste Væsen kun var en luftig Tanke, dersom Personligheden ikke havde videre at betyde end at være en Nisse, der vel deeltog i Bevægelserne, men dog blev uforandret, dersom det forholdt sig saaledes, saa vilde det være en Taabelighed, at det kunde være for sildigt for et Menneske at vælge, da der i dybere Forstand slet ikke kunne være Tale om et Valg. Valget selv er afgjørende for Personlighedens Indhold; ved Valget synker den ned i det Valgte, og naar den ikke vælger, henvisner den i Tæring. Et Øieblik er det saa, et Øieblik kan det synes saa, at det, der skal vælges mellem, ligger udenfor den Vælgende; han staaer i intet Forhold dertil, han kan bevare sig i Indifferens ligeover derfor. Dette er Overveielsens Øieblik, men det er ligesom det platoniske egentlig slet ikke, og allermindst i den abstracte Forstand, hvori Du vil fastholde det; og jo længere man stirrer derpaa, jo mindre er det. Det, der skal vælges, staaer i det dybeste Forhold til den Vælgende, og naar der er Tale om Valg, der angaaer et Livs-Spørgsmaal, saa skal Individet jo paa samme Tid leve, og kommer derved, jo længere han skyder Valget ud, let til at alterere det, uagtet han bestandig overveier og overveier, og derved troer ret at holde Valgets Modsætninger ude fra hinanden. Naar man betragter Livets Enten-Eller saaledes, da fristes man ikke let til at spøge dermed. Man seer da, at Personlighedens indre Drift ikke har Tid til Tanke-Experimenter, at den bestandig iler frem og paa en eller anden Maade ponerer enten det Ene eller det Andet, hvorved da Valget i næste Øieblik bliver vanskeligere; thi det, der er poneret, skal tages tilbage. Dersom Du tænker Dig en Styrmand i hans Skib i det Øieblik, der skal gjøres et Slag, saa vil han maaskee kunne sige, jeg kan enten gjøre det eller det; men dersom han ikke er en maadelig Styrmand, saa vil han tillige blive sig bevidst, at Skibet under alt dette skyder sin sædvanlige Fart, og at det saaledes kun er et Øieblik, hvor det er ligegyldigt, om han gjør dette eller dette. Saaledes med et Menneske, glemmer han at beregne denne Fart, saa kommer der tilsidst et Øieblik, hvor der ikke mere er Tale om et Enten-Eller, ikke fordi han har valgt, men fordi han har ladet det være, hvilket ogsaa kan udtrykkes saaledes, fordi Andre have valgt for ham, fordi han har tabt sig selv.

Af det her Udviklede vil Du ogsaa see, hvori min Betragtning af et Valg er væsentlig forskjellig fra Din, ifald jeg ellers kan tale om denne; thi Din er netop deri forskjellig, at den forhindrer et Valg. Valgets Øieblik er for mig saare alvorligt, ikke saa meget paa Grund af den strenge Gjennemtænkning af det, der i Valget viser sig adskilt, ikke paa Grund af den Mangfoldighed af Tanker, der knytte sig til det enkelte Led, som fordi der er Fare paa Færde, at det i næste Øieblik ikke saaledes staaer til min Raadighed at vælge, at der da allerede er levet Noget, som maa leves om igjen; thi naar man troer, at man noget Øieblik kan holde sin Personlighed blank og bar, eller at man i strengere Forstand kan standse og afbryde det personlige Liv, saa er man i en Vildfarelse. Personligheden er allerede, før man vælger, interesseret i Valget, og naar man udsætter Valget, saa vælger Personligheden ubevidst eller de dunkle Magter i den. Naar man da endelig faaer valgt, hvis man ikke, hvad jeg før bemærkede, bliver aldeles forflygtiget, saa opdager man, at der er Noget, der maa gjøres om igjen, der maa tages tilbage, og dette er ofte saare vanskeligt. Der tales i Eventyr om Mennesker, som Havfruer eller Havmænd ved deres dæmoniske Musik rev ind i deres Magt. For at løse Trylleriet, lærer Eventyret, var det nødvendigt, at den Fortryllede spillede det samme Stykke baglænds, uden at tage Feil en eneste Gang. Dette er saare dybsindigt tænkt, men saare vanskeligt udført, og dog er det saaledes; det Forfeilede, man har faaet ind i sig, maa man paa den Maade udrydde, og hver Gang man tager Feil, maa man begynde forfra. See derfor er det af Vigtighed at vælge og at vælge itide. Du derimod har en anden Methode; thi det veed jeg vel, at den polemiske Side, Du vender ud imod Verden, ikke er Dit sande Væsen. Ja hvis det at overveie var Opgaven for et Menneskeliv, saa var Du Fuldkommenheden nær. Jeg vil tage et Exempel. Det maa naturligviis være dristige Modsætninger for at det kan passe paa Dig: enten Præst – eller Skuespiller. Her er Dilemmaet. Nu vaagner hele Din lidenskabelige Energi; Reflexionen griber med sine hundrede Arme den Tanke at være Præst. Du finder ingen Hvile, Dag og Nat tænker Du derover; Du læser alle de Skrifter, Du kan faae fat paa, gaaer 3 Gange hver Søndag i Kirke, stifter Bekjendtskab med Præster, skriver selv Prædikener, holder dem for Dig selv, i et halvt Aar er du Død for hele Verden. Nu er Du færdig; Du kan nu tale indsigtsfuldere og tilsyneladende erfarnere om det at være Præst end Mangen, der har været Præst i 20 Aar. Det vækker Din Forbittrelse, naar Du træffer sammen med Saadanne, at de ikke veed med ganske anderledes Veltalenhed at expectorere sig; er det Begeistring, siger Du, jeg som ikke er Præst, som ikke har helliget mig dertil, jeg taler med Englerøst i Sammenligning med dem. Det er nu maaskee ogsaa sandt nok, imidlertid er Du dog ikke bleven Præst. Nu bærer Du Dig ligeledes ad med det andet Problem, og Din Kunst-Begeistring overgaaer næsten Din geistlige Veltalenhed. Du er da færdig til at vælge. Imidlertid kan man være sikker paa, at der i den uhyre Tankebedrift, i hvilken Du har levet, er faldet Meget af, mange Smaabemærkninger og Iagttagelser. I det Øieblik derfor Du skal til at vælge, kommer der Liv og Røre i dette Affald, der viser sig et nyt enten – eller: Jurist; Advocat maaskee, det har Noget tilfælles med begge Dele. Nu er Du tabt. I samme Moment er Du nemlig strax Advokat nok til at kunne bevise det Rigtige i at tage det Tredie med. Saaledes gaaer Dit Liv hen. Efter at have spildt halvandet Aar paa disse Overveielser, efterat have anstrengt al Din Sjæls Kraft med en beundringsværdig Energi, er Du ikke kommen et Skridt videre. Da brister Tankens Traad, Du bliver utaalmodig, lidenskabelig, skjender og brænder, og nu vedbliver Du: eller Haarskærer, eller Tæller i Banken, jeg siger bare enten–eller. Hvad Under da, at dette Ord er bleven Dig en Forargelse og en Daarskab, »at det forekommer Dig at være ligesom Armene paa den Jomfru, hvis Favnetag var Dødsstraf.« Du overseer Menneskene, Du driver Din Spot med dem, og det, Du er bleven, er det, Du af Alt mest afskyer – Kritiker, en Universal-Kritiker ved alle Faculteter. Stundom kan jeg ikke lade være at smile ad Dig, og dog er det sørgeligt, at Dine i Sandhed udmærkede Aandsevner saaledes adsplittes. Dog her er atter den samme Modsigelse i Dit Væsen; thi det Latterlige seer Du meget godt, og Gud trøste den, der falder i Din Haand, naar det hænger saaledes sammen med ham, og dog er hele Forskjellen den, at han maaskee bliver nedbøiet og knækket, Du derimod bliver let og rank og lystigere end nogensinde, og lyksaliggjør Dig selv og Andre med det Evangelium: vanitas vanitatum vanitas juchhe! Men dette er intet Valg, det er hvad man siger paa Dansk: lad gaae, eller en Mediation som den at lade fem være lige. Nu føler Du Dig fri, siger Verden Farvel.

So zieh' ich hin in alle Ferne

Ueber meiner Mütze nur die Sterne.

See dermed har Du valgt, vistnok ikke, det vil Du vel selv tilstaae, den bedre Deel; men Du har egentlig slet ikke valgt, eller Du har valgt i uegentlig Forstand. Dit Valg er et æsthetisk Valg; men et æsthetisk Valg er intet Valg. Overhovedet er det at vælge et egentligt og stringent Udtryk for det Ethiske. Overalt hvor der i strengere Forstand er Tale om et Enten–Eller, der kan man altid være sikker paa, at det Ethiske er med i Spillet. Det eneste absolute Enten–Eller, der gives, er Valget mellem Godt og Ondt, men det er ogsaa absolut ethisk. Det æsthetiske Valg er enten aldeles umiddelbart og forsaavidt intet Valg, eller det fortaber sig i Mangfoldigheden. Naar saaledes en ung Pige følger sit Hjertes Valg, saa er dette Valg, hvor skjønt det end forresten er, i strengere Forstand intet Valg, da det er aldeles umiddelbart. Naar et Menneske æsthetisk overveier en Mængde Livs-Opgaver, saaledes som Du i det Foregaaende, saa faaer han ikke let et Enten–Eller, men en heel Mangfoldighed, fordi det Selvbestemmende i Valget ikke her bliver ethisk accentueret, og fordi, naar man ikke vælger absolut, da vælger man kun for Momentet, og kan desaarsag i næste Øieblik vælge noget Andet. Det ethiske Valg er derfor i en vis Forstand langt lettere, langt simplere, men i en anden Forstand er det uendeligt langt sværere. Den, der ethisk vil bestemme sig sin Livs-Opgave, har i Almindelighed ikke et saa betydeligt Udvalg; derimod har Valgets Akt langt mere for ham at betyde. Naar Du da vil forstaae mig rigtigt, saa kan jeg gjerne sige, at det i at vælge ikke saa meget kommer an paa at vælge det Rigtige, som paa den Energi, den Alvor og Pathos, hvormed man vælger. Deri forkynder Personligheden sig i sin indre Uendelighed, og derved consolideres igjen Personligheden. Selv derfor om et Menneske valgte det Urette, saa vil han dog, netop paa Grund af den Energi, hvormed han valgte, opdage, at han valgte det Urette. Idet nemlig Valget er foretaget med hele Personlighedens Inderlighed, er hans Væsen luttret, og han selv bragt i et umiddelbart Forhold til den evige Magt, der allestedsnærværende gjennemtrænger hele Tilværelsen. Denne Forklarelse, denne høiere Indvielse opnaaer aldrig den, der blot vælger æsthetisk. Rhythmen i hans Sjæl er trods al dens Lidenskab dog kun en spiritus lenis.

Som en Cato tilraaber jeg Dig da mit Enten–Eller, og dog ikke som en Cato; thi min Sjæl har endnu ikke erhvervet den resignerede Kulde, han var i Besiddelse af. Men jeg veed det, kun denne Besværgelse, hvis jeg har den tilstrækkelige Kraft, vil være istand til at vække Dig, ikke til en Tankens Virksomhed, thi den mangler Du ikke, men til Aandens Alvor. Maaskee vil det lykkes Dig ogsaa uden den at udrette Meget, maaskee endog at forbause Verden (thi jeg er ikke knap), og dog vil Du gaae Glip af det Høieste, af det Eneste, der i Sandhed giver Livet Betydning, maaskee vil Du vinde den ganske Verden og tabe Dig selv.

Hvad er det da jeg adskiller i mit Enten–Eller? Er det Godt og Ondt? Nei, jeg vil kun bringe Dig paa det Punkt, at dette Valg i Sandhed faaer Betydning for Dig. Derom er det at Alt dreier sig. Naar man først kan faae et Menneske til at staae paa Skilleveien saaledes, at der ingen Udvei er for ham uden ved at vælge, saa vælger han det Rette. Skulde det derfor skee, at Du, inden Du fik læst denne lidt udførligere Undersøgelse, der atter tilsendes Dig i et Brevs Form, følte, at Valgets Øieblik var der, saa kast det Øvrige bort, bryd Dig aldrig om det, Du har Intet tabt; men vælg, og Du skal see, hvilken Gyldighed der ligger deri, ja ingen ung Pige kan være saa lykkelig ved sit Hjertes Valg, som en Mand, der har vidst at vælge. Enten skal man altsaa leve æsthetisk, eller man skal leve ethisk. Her er, som sagt, endnu ikke i strengere Forstand Tale om et Valg; thi den, som lever æsthetisk, vælger ikke, og den som, efterat det Ethiske har viist sig for ham, vælger det Æsthetiske, han lever ikke æsthetisk, thi han synder, og ligger under ethiske Bestemmelser, om hans Liv end maa betegnes som uethisk. See dette er ligesom en character indelebilis ved det Ethiske, at det, uagtet det beskedent stiller sig hen paa lige Trin med det Æsthetiske, dog egentlig er det, der gjør Valget til et Valg. Og dette er det Sørgelige, naar man betragter Menneskenes Liv, at saa Mange henleve deres Liv i stille Fortabelse; de udleve dem selv, ikke i den Betydning, at Livets Indhold successivt udfolder sig, og nu eies i denne Udfoldelse, men de leve dem ligesom ud af dem selv, forsvinde som Skygger; deres udødelige Sjæl henveires, og de ængstes ikke ved Spørgsmaal om dens Udødelighed; thi de ere jo allerede opløste før de døe. De leve ikke æsthetisk, men det Ethiske har heller ei viist sig for dem i sin Heelhed; de have da heller ei egentlig forkastet det, de synde derfor heller ikke, uden forsaavidt det er en Synd, at de hverken er det ene eller det andet; de tvivle da heller ei om deres Udødelighed; thi den, som dybt og inderligt paa egne Vegne tvivler derom, han finder nok det Rette. Paa egne Vegne, siger jeg, og det er vel paa den høie Tid, at man advarer mod den storhjertede, heltemodige Objektivitet, hvormed mange Tænkere tænke paa Alles Vegne, ikke paa deres egne. Vil man kalde hvad jeg her fordrer Selvkjærlighed, saa vil jeg svare: det kommer deraf, at man ingen Forestilling har om, hvad dette »Selv« er, og at det hjalp et Menneske saare lidet, at han vandt den ganske Verden, men tabte sig selv, samt at det nødvendigt maa være et slet Beviis, der ikke først og fremmest overbeviser den, der fremsætter det.

Mit Enten–Eller betegner ikke nærmest Valget mellem Godt og Ondt, det betegner det Valg, hvorved man vælger Godt og Ondt eller udelukker dem. Spørgsmaalet er her, under hvilke Bestemmelser man vil betragte hele Tilværelsen og selv leve. At den, der vælger Godt og Ondt, vælger det Gode, er vel sandt, men dette viser sig først bag efter; thi det Æsthetiske er ikke det Onde, men Indifferensen, og derfor var det jeg sagde, at det Ethiske constituerer Valget. Der er derfor ikke saa meget Tale om at vælge mellem at ville det Gode eller det Onde, som om at vælge det at ville, men hermed er igjen det Gode og det Onde sat. Den, der vælger det Ethiske, vælger det Gode, men det Gode er her aldeles abstrakt, dets Væren er derved blot sat, og deraf følger ingenlunde, at den Vælgende ikke igjen kan vælge det Onde, uagtet han valgte det Gode. Her seer Du igjen, hvor vigtigt det er, at der vælges, og at det, det kommer an paa, ikke saa meget er Overveielsen som den Villiens Daab, der indoptager denne i det Ethiske. Jo længere Tiden gaaer hen, jo vanskeligere bliver det at vælge; thi Sjælen er bestandig i den ene af Dilemmaets Dele, og det bliver derfor vanskeligere og vanskeligere at rive sig løs. Og dog er dette nødvendigt, hvis der skal vælges, og altsaa af yderste Vigtighed, hvis et Valg har Noget at betyde, og at dette er Tilfælde, skal jeg senere vise.

Du veed, jeg har aldrig givet mig ud for Philosoph, allermindst naar jeg underholder mig med Dig. Deels for at drille Dig lidt, deels fordi det virkelig er min kjæreste og dyrebareste, og i en vis Forstand betydningsfuldeste Stilling i Livet, pleier jeg at træde op som Ægtemand. Jeg har ikke opoffret mit Liv til Kunst og Videnskab, hvad jeg har offret mig for er i Sammenligning hermed Ubetydeligheder; jeg offrer mig for min Gjerning, min Kone, mine Børn, eller rettere sagt, jeg offrer mig ikke derfor, men jeg finder min Tilfredshed og Glæde deri. Det er Ubetydeligheder i Sammenligning med hvad Du lever for, og dog, min unge Ven, pas vel paa, at det Store, Du virkelig offrer Dig for, ikke bedrager Dig. Skjøndt jeg nu ikke er Philosoph, saa nødsages jeg dog her til at vove mig ind paa en lille philosophisk Overveielse, som jeg vil bede Dig ikke saa meget at kritisere som at tage Dig ad notam. Det polemiske Resultat nemlig, hvoraf alle Dine Seirshymner over Tilværelsen gjenlyde, har en underlig Lighed med den nyere Philosophies Yndlings-Theori, at Modsigelsens Grundsætning er hævet. Vel veed jeg, at det Standpunkt, Du indtager, er Philosophien en Vederstyggelighed, og dog forekommer det mig, at den selv gjør sig skyldig i den samme Feil, ja at Grunden, hvorfor man ikke strax mærker det, er den, at den end ikke er stillet saa rigtig som Du. Du ligger paa Gjerningens Gebet, den paa Contemplationens. Saasnart man derfor vil føre den over paa det Praktiske, maa den komme til samme Resultat som Du, om den end ikke udtrykker sig saaledes. Du medierer Modsætningerne i en høiere Galskab, Philosophien i en høiere Eenhed. Du vender Dig mod den tilkommende Tid, thi Handling er væsentlig futurisk; Du siger, jeg kan enten gjøre det eller gjøre det, men hvilken af Delene jeg gjør er det lige galt, ergo gjør jeg slet Intet. Philosophien vender sig mod den forbigangne Tid, mod hele den oplevede Verdenshistorie, den viser, hvorledes de discursive Momenter gaae sammen i en høiere Eenhed, den medierer og medierer. Derimod synes den mig slet ikke at svare paa det, hvorom jeg spørger; thi jeg spørger om den tilkommende Tid. Du svarer dog paa en Maade, om Du end svarer Nonsens. Jeg antager nu, at Philosophien har Ret, at Modsigelsens Grundsætning virkelig er hævet, eller at Philosopherne i ethvert Øieblik hæve den i den høiere Eenhed, der er for Tanken. Dette kan jo dog ikke gjelde om den tilkommende Tid; thi Modsætningerne maae jo dog først have været der, førend jeg kan mediere dem. Men er Modsætningen tilstede, saa er der et Enten–Eller. Philosophen siger: saaledes er det gaaet hidtil; jeg spørger: hvad har jeg at gjøre, hvis jeg ikke vil være Philosoph, thi dersom jeg vil det, saa seer jeg jo nok, at jeg, som de andre Philosopher, kommer til at mediere den forbigangne Tid. Deels er dette intet Svar paa mit Spørgsmaal, hvad jeg skal gjøre; thi om jeg saa end var det mest begavede philosophiske Hoved, der nogensinde har levet i Verden, saa maa der være en Ting mere jeg har at gjøre, end at sidde og betragte det Forbigangne; deels er jeg Ægtemand, og ingenlunde et philosophisk Hoved, men henvender mig i al Ærbødighed til denne Videnskabs Dyrkere for at faae at vide, hvad jeg skal gjøre. Jeg faaer imidlertid intet Svar; thi Philosophen medierer det Forbigangne og er i dette, Philosophen haster i den Grad ind i Fortiden, at, som en Digter siger om en Antiqvar, kun hans Kjoleskøder ere blevne tilbage i Nutiden. See her enes Du med Philosopherne. Det, I enes i, er, at Livet gaaer istaa. For Philosophen er Verdenshistorien afsluttet, og han medierer. Derfor hører det ulystelige Syn til Dagens Orden i vor Tid, at see unge Mennesker, der kunne mediere Christendom og Hedenskab, der kunne lege med Historiens titaniske Kræfter, og som ikke kunne sige et eenfoldigt Menneske, hvad han har at gjøre her i Livet, og som heller ikke veed, hvad de selv have at gjøre. Du er saare mangfoldig i Udtryk for Dit Yndlings-Resultat, jeg vil her fremhæve eet, fordi Du deri har en paafaldende Lighed med Philosophen, om end hans virkelige eller paatagne Alvor vil forbyde ham at deeltage i det obligate Opsving, hvori Du forlyster Dig. Spørger man Dig, om Du vil underskrive en Adresse til Kongen, eller om Du ønsker en Constitution, eller Skattebevilgningsret, eller om Du vil forene Dig til dette eller hiint goddædige Øiemed, saa svarer Du: »Høistærede Samtidige! I misforstaae mig, jeg er slet ikke med, jeg er ude, jeg er ude som et lille bitte spansk s.« Saaledes gaaer det ogsaa Philosophen, han er ude, han er ikke med, han sidder og ældes ved at høre paa Fortidens Sange, han lytter til Mediationens Harmonier. Jeg hædrer Videnskaben, jeg ærer dens Dyrkere, men Livet har ogsaa sine Fordringer, og om jeg end, hvis jeg saae et enkelt ualmindeligt begavet Hoved eensidigen fortabe sig i det Forbigangne, vilde være raadvild om, hvad jeg skulde dømme, hvilken Mening jeg skulde have ved Siden af den Ærbødighed, jeg vilde nære for hans Aandsdygtighed; i vor Tid bliver jeg ikke raadvild, da jeg seer en Skare af unge Mennesker, der dog umulig alle kunne være philosophiske Hoveder, fortabte i Tidens Yndlings-Philosophi, eller som jeg snart fristes til at kalde den, Tidens Ynglings-Philosophi. Jeg har en gyldig Fordring ligeoverfor Philosophien, saaledes som Enhver har den, hvem den ikke tør afvise paa Grund af total Mangel paa Evne. Jeg er Ægtemand, jeg har Børn. Hvad om jeg nu i deres Navn spurgte den om, hvad et Menneske har at gjøre i Livet. Du vil maaskee smile, i ethvert Tilfælde vil den philosophiske Ungdom smile ad en Familiefader, og dog mener jeg, at det i Sandhed er et forfærdeligt Argument mod den, hvis den Intet har at svare. Er Livets Gang standset, kan maaskee den Generation, der nu er, leve af at betragte, hvad skal den følgende leve af? Af at betragte det Samme? Den sidste Generation har jo Intet udrettet, Intet efterladt, der skal medieres. See her kan jeg atter slaae Dig sammen med Philosopherne, og sige til Eder: I gaae dog Glip af det Høieste. Min Stilling som Ægtemand kommer mig her til Hjælp, til bedre at kunne oplyse, hvad jeg mener. Hvis en Ægtemand vilde sige, det fuldkomne Ægteskab er det, hvor der ingen Børn er, saa vilde han gjøre sig skyldig i den samme Misforstaaelse, som Philosopherne. Han gjør sig selv til det Absolute, og dog vil enhver Ægtemand føle, at det er usandt og uskjønt, og at det, at han selv bliver til Moment, som han bliver ved et Barn, er langt sandere.

Dog jeg er maaskee allerede gaaet for vidt, jeg har indladt mig paa Undersøgelser, jeg maaskee slet ikke burde, deels fordi jeg ikke er Philosoph, deels fordi det ingenlunde er min Hensigt at underholde mig med Dig om et eller andet Phænomen i Tiden, men vel at tiltale Dig, at lade Dig paa enhver Maade føle, at Du er den Tiltalte. Da jeg imidlertid engang er kommen saa vidt, saa vil jeg dog lidt nøiere overveie, hvorledes det hænger sammen med den philosophiske Mediation af Modsætningerne. Skulde det, jeg her har at sige, mangle Stringens, saa har det maaskee lidt mere Alvor, og det er ogsaa alene af den Grund, det her bliver fremsat; thi jeg agter ikke at concurrere til nogen philosophisk Værdighed, men vel, da jeg engang har faaet Pennen i Haanden, ogsaa med den at forsvare, hvad jeg ellers forsvarer paa andre og bedre Maader.

Saasandt der altsaa er en tilkommende Tid, saa sandt er der et enten – eller. Den Tid, i hvilken Philosophen lever, er da ikke den absolute Tid, den er selv et Moment, og det er altid en betænkelig Omstændighed, naar en Philosophi er ufrugtbar, ja det maa ansees for en Vanære for den, ligesom man i Orienten anseer Ufrugtbarhed for en Skjendsel. Tiden bliver altsaa selv Moment, og Philosophen bliver selv Moment i Tiden. Vor Tid vil da igjen for en senere Tid vise sig som et discursivt Moment, og en senere Tids Philosoph vil igjen mediere vor Tid, og saaledes videre. Forsaavidt er da Philosophien i sin Ret, og det blev at betragte som en tilfældig Feil ved vor Tids Philosophi, at den forvexlede vor Tid med den absolute Tid. Imidlertid er det dog let at indsee, at Mediationens Kategori derved har lidt et betydeligt Knæk, og at den absolute Mediation først bliver mulig, naar Historien bliver færdig, med andre Ord, at Systemet er i en bestandig Vorden. Det, som derimod Philosophien har beholdt, er Anerkjendelsen af, at der er en absolut Mediation. Dette er den naturligviis af yderste Vigtighed; thi opgiver man Mediationen, saa opgiver man Speculationen. Paa den anden Side er det en betænkelig Sag at indrømme det; thi indrømmer man Mediationen, saa er der intet absolut Valg, og er der ikke et saadant, saa er der intet absolut enten – eller. Dette er Vanskeligheden; dog troer jeg, at den for en Deel ligger deri, at man forvexler to Sphærer med hinanden, Tænkningens og Frihedens. For Tanken bestaaer Modsætningen ikke, den gaaer over i det Andet og derpaa sammen i en høiere Eenhed. For Friheden bestaaer Modsætningen; thi den udelukker den. Jeg forvexler ingenlunde liberum arbitrium med den sande positive Frihed; thi selv denne har i al Evighed det Onde udenfor sig, om end som en afmægtig Mulighed, og den bliver ikke fuldkommen derved, at den mere og mere optager det Onde, men derved, at den mere og mere udelukker det, men Udelukkelse er netop Modsætning til Mediation. At jeg herved ikke kommer til at antage et radikalt Onde, skal jeg senere vise.

De Sphærer, som Philosophien egentlig har med at gjøre, de Sphærer, der egentlig ere for Tanken, er det Logiske, Naturen, Historien. Her hersker Nødvendigheden, og derfor har Mediationen sin Gyldighed. At dette er Tilfældet med det Logiske og Naturen vil vel Ingen negte; med Historien har det derimod sin Vanskelighed; thi, siger man, her hersker Friheden. Imidlertid troer jeg, at man betragter Historien urigtigt, og at Vanskeligheden kommer heraf. Historien er nemlig mere end et Produkt af de frie Individers frie Handlinger. Individet handler, men denne Handling gaaer ind i den Tingenes Orden, der bærer hele Tilværelsen. Hvad der skal komme ud deraf, veed den Handlende egentlig ikke. Men denne høiere Tingenes Orden, der saa at sige fordøier de frie Handlinger, og sammenarbeider dem i sine evige Love, er Nødvendigheden, og denne Nødvendighed er Bevægelsen i Verdenshistorien, og det er derfor ganske rigtigt, at Philosophien anvender Mediationen, det vil da sige den relative Mediation. Betragter jeg en verdenshistorisk Individualitet, saa kan jeg skjelne mellem de Gjerninger, om hvilke Skriften siger, at de følge ham, og de Gjerninger, ved hvilke han tilhører Historien. Hvad man kunde kalde den indvortes Gjerning, har Philosophien slet ikke med at gjøre; men den indvortes Gjerning er Frihedens sande Liv. Philosophien betragter den udvortes Gjerning, og denne seer den igjen ikke isoleret, men seer den indoptagen og forvandlet i den verdenshistoriske Proces. Denne Proces er egentlig Philosophiens Gjenstand, og denne betragter den under Nødvendighedens Bestemmelse. Den holder derfor den Reflexion borte, der vil gjøre opmærksom paa, at Alt kunde være anderledes, den betragter Verdenshistorien saaledes, at der intet Spørgsmaal er om et enten – eller. At der i denne dens Betragtning indblander sig megen taabelig og inept Tale, det forekommer idetmindste mig; at især de unge Hexemestere, der ville besværge Historiens Aander, synes mig latterlige, negter jeg ikke, men jeg bøier mig ogsaa dybt for de storartede Præstationer, som vor Tid har at opvise. Som sagt, Philosophien seer Historien under Nødvendighedens Bestemmelse, ikke under Frihedens; thi om man end kalder den verdenshistoriske Proces fri, saa er det dog i samme Forstand, som man taler om den organiserende Proces i Naturen. For den historiske Proces er der intet Spørgsmaal om et enten – eller; men dog kan det vel ikke falde nogen Philosoph ind at negte, at der for det handlende Individ var det. Deraf igjen den Sorgløshed, den Forsonlighed, hvormed Philosophien betragter Historien og dens Helte; thi den seer dem under Nødvendighedens Bestemmelse. Deraf igjen dens Uformuenhed til at faae et Menneske til at handle; dens Tilbøielighed til at lade Alt gaae istaa; thi den fordrer egentlig, at man skal handle nødvendigt, hvilket er en Modsigelse.

Selv det ringeste Individ har saaledes en Dobbelt-Existens. Ogsaa han har en Historie, og denne er ikke blot et Produkt af hans egne frie Handlinger. Den indvortes Gjerning derimod tilhører ham selv og skal i al Evighed tilhøre ham; den kan Historien eller Verdenshistorien ikke tage fra ham, den følger ham enten til Glæde eller til Bedrøvelse. I denne Verden hersker der et absolut enten – eller; men denne Verden har Philosophien ikke med at gjøre. Forestiller jeg mig en ældre Mand, der skuer tilbage over et bevæget Liv, saa faaer han for Tanken ogsaa en Mediation deraf; thi hans Historie var indflettet i Tidens; men inderst inde faaer han ingen Mediation. Der adskiller endnu bestandigt et enten – eller det, der var adskilt, da han valgte. Skal her være Tale om en Mediation, saa kunde man sige, det er Angeren; men Angeren er ingen Mediation, den seer ikke attraaende paa det, der skal medieres, dens Vrede fortærer det; men dette er ligesom Udelukkelse, Modsætningen til Mediation. Her viser det sig tillige, at jeg ikke antager et radikalt Onde, da jeg statuerer Angerens Realitet; men Anger er vel et Udtryk for Forsoning; men tillige et absolut uforsonligt Udtryk.

Dog alt dette indrømmer Du mig maaskee; Du, der dog paa saa mange Maader gjør fælleds Sag med Philosopherne, undtagen forsaavidt Du for privat Regning paatager Dig at gjøre Nar af dem; Du mener maaskee, at jeg som Ægtemand kan lade mig nøie dermed, og gjøre Brug deraf i min Husholdning. Oprigtigt talt, jeg forlanger ikke mere; men jeg gad dog vide, hvilket Liv der var det høieste, enten Philosophens eller den frie Mands. Dersom Philosophen blot er Philosoph, fortabt heri uden at kjende Frihedens salige Liv, da mangler han et saare vigtigt Punkt, han vinder den ganske Verden, og han taber sig selv; dette kan aldrig hænde den, der lever for Friheden, om han end tabte nok saa meget.

For Friheden kæmper jeg derfor (deels her i dette Brev; deels og fornemlig i mig selv), for den tilkommende Tid, for enten – eller. Det er den Skat, jeg agter at efterlade dem, jeg elsker i Verden. Ja dersom min lille Søn i dette Øieblik var i den Alder, at han ret kunde forstaae mig, og min sidste Time var kommen, da vilde jeg sige til ham: jeg efterlader Dig ikke Formue, ikke Titler og Værdigheder; men jeg veed, hvor der ligger en Skat begraven, der kan gjøre Dig rigere end hele Verden, og denne Skat tilhører Dig, og Du skal end ikke takke mig for den, at Du ikke skal tage Skade paa Din Sjæl ved at skylde et Menneske Alt; denne Skat er nedlagt i Dit eget Indre: der er et enten – eller, der gjør et Menneske større end Englene.

Her vil jeg afbryde denne Betragtning. Den tilfredsstiller Dig maaskee ikke, Dit graadige Øie sluger den uden at Du mættes, men det ligger deri, at Øiet er det, der sidst mættes, især naar man som Du ikke hungrer, men blot lider af en Øiets Lyst, der ikke kan tilfredsstilles.

Det, der da træder frem ved mit enten – eller er det Ethiske. Der er derfor endnu ikke Tale om Valget af Noget, ikke Tale om Realiteten af det Valgte, men om Realiteten af det at vælge. Dette er imidlertid det Afgjørende, og det er dertil jeg vil stræbe at vække Dig. Til dette Punkt kan det ene Menneske hjælpe det andet; naar han er kommen hertil, da bliver den Betydning, det ene Menneske kan have for det andet, mere underordnet. Jeg har i et foregaaende Brev bemærket, at det at have elsket giver et Menneskes Væsen en Harmoni, som aldrig tabes ganske; nu vil jeg sige, det at vælge giver et Menneskes Væsen en Høitidelighed, en stille Værdighed, som aldrig tabes ganske. Der gives Mange, der sætte en overordentlig Priis paa at have skuet en eller anden mærkelig verdenshistorisk Individualitet Ansigt til Ansigt. Dette Indtryk glemme de aldrig, det har givet deres Sjæl et idealt Billede, som adler deres Væsen; og dog er selv dette Øieblik, hvor betydningsfuldt det end kan være, Intet mod Valgets Øieblik. Naar da Alt er blevet stille omkring Een, høitideligt som en stjerneklar Nat, naar Sjælen bliver ene i den hele Verden, da viser der sig ligeoverfor den ikke et udmærket Menneske, men den evige Magt selv, da skiller Himlen sig ligesom ad, og Jeget vælger sig selv, eller rettere, det modtager sig selv. Da har Sjælen seet det Høieste, hvad intet dødeligt Øie kan see, og som aldrig kan glemmes, da modtager Personligheden det Ridderslag, der adler den for en Evighed. Han bliver ikke en anden end han var før, men han bliver sig selv; Bevidstheden slutter sig sammen, og han er sig selv. Som en Arving, om han end var Arving til Alverdens Skatte, dog ikke eier dem, før han er bleven myndig, saaledes er selv den rigeste Personlighed Intet, før han har valgt sig selv, og paa den anden Side er, selv hvad man maatte kalde den fattigste Personlighed Alt, naar han har valgt sig selv; thi det Store er ikke at være Dette eller Hiint; men at være sig selv, og dette kan ethvert Menneske, naar han vil det.

At der i en vis Forstand ikke er Tale om et Valg af Noget, vil Du see deraf, at det, der viser sig paa den anden Side, er det Æsthetiske, der er Indifferensen. Og dog er her Tale om et Valg, ja et absolut Valg; thi kun ved at vælge absolut, kan man vælge det Ethiske. Ved det absolute Valg er altsaa det Ethiske sat; men deraf følger ingenlunde, at det Æsthetiske er udelukket. I det Ethiske er Personligheden centraliseret i sig selv, absolut er altsaa det Æsthetiske udelukket eller det er udelukket som det Absolute, men relativt bliver det bestandigt tilbage. Idet Personligheden vælger sig selv, vælger den sig selv ethisk og udelukker absolut det Æsthetiske; men da han dog vælger sig selv, og ikke ved at vælge sig selv bliver et andet Væsen, men bliver sig selv, saa vender hele det Æsthetiske tilbage i sin Relativitet.

Det enten – eller, jeg har opstillet, er altsaa i en vis Forstand absolut; thi det er det mellem at vælge og ikke at vælge. Men da Valget er et absolut Valg, saa er enten – eller absolut, i en anden Forstand indtræder dog først det absolute enten – eller med Valget; thi nu viser Valget mellem Godt og Ondt sig. Dette i og med det første Valg satte Valg skal ikke beskæftige mig her, jeg vil blot tvinge Dig hen paa det Punkt, hvor Nødvendigheden af Valget viser sig, og derpaa betragte Tilværelsen under ethiske Bestemmelser. Jeg er ingen ethisk Rigorist, begeistret for en formel abstrakt Frihed; naar blot Valget er sat, saa vender alt det Æsthetiske tilbage, og Du skal see, at først herved bliver Tilværelsen skjøn, og at det først ad denne Vei kan lykkes et Menneske at frelse sin Sjæl og vinde den hele Verden, at bruge Verden uden at misbruge den.

Men hvad er det at leve æsthetisk, og hvad er det at leve ethisk? Hvad er det Æsthetiske i et Menneske og hvad er det Ethiske? Herpaa vilde jeg svare: det Æsthetiske i et Menneske er det, hvorved han umiddelbar er det, han er; det Ethiske er det, hvorved han bliver det, han bliver. Den, der lever i og ved og af og for det Æsthetiske i ham, han lever æsthetisk.

[…]

Du har adskillige gode Ideer, mange snurrige Indfald, en Mængde Taabeligheder, behold det Altsammen, jeg forlanger det ikke, men Du har een Idee, jeg vil bede Dig holde fast paa, en Idee, der forvisser mig om, at min Aand er beslægtet med Din. Du har ofte sagt, at Du vilde hellere være alt Andet i Verden end være Digter, da i Reglen en Digter-Existens er en Menneske-Offring. For mit Vedkommende skal ingenlunde være negtet, at der har levet Digtere, der havde vundet dem selv, før de begyndte at digte, eller som vandt sig selv ved at digte, men paa den anden Side er det ogsaa vist, at en Digter-Existens som saadan ligger i den Dunkelhed, der er Følge af, at en Fortvivlelse ikke bliver gjennemført, at Sjælen bestandig zittrer i Fortvivlelse og Aanden ikke kan vinde sin sande Forklarelse. Det digteriske Ideal er altid et usandt Ideal, thi det sande Ideal er altid det virkelige. Naar Aanden da ikke faaer Lov til at svinge sig op i Aandens evige Verden, saa bliver den underveis og fryder sig ved de Billeder, der speile sig i Skyerne og græde over deres Forgængelighed. En Digter-Existens er derfor som saadan en ulykkelig Existens, den er høiere end Endeligheden og dog ikke Uendeligheden. Digteren seer Idealerne, men han maa flygte bort fra Verden for at glæde sig ved dem, han kan ikke bære disse Gudebilleder i sig midt i Livets Forvirring, ikke gaae rolig sin Gang uanfægtet af Karrikaturen, der viser sig rundtom, end sige at han skulde have Styrke til at iføre sig dem. Digterens Liv er derfor ofte Gjenstand for en ussel Medlidenhed af de Mennesker, der mene at have Deres paa det Tørre, fordi de ere forblevne i Endeligheden. Du yttrede engang i et mismodigt Øieblik, at der vistnok allerede var dem, der i deres stille Sind havde gjort Regningen op med Dig, og som vare villige til at qvittere paa følgende Vilkaar: Du blev erkjendt for et godt Hoved, til Gjengjæld gik Du under og blev intet nævenyttigt Medlem af Samfundet. Ja det er unegteligt, der gives en saadan Usselhed i Verden, der paa den Maade vil seire over Alt, hvad der rager blot en Tomme frem. Lad det imidlertid ikke forstyrre Dig, trods dem ikke, foragt dem ikke, her vil jeg sige, som Du pleier at sige: det er ikke Umagen værd. Men vil Du ikke være Digter, saa er der ingen anden Vei for Dig end den, jeg har viist Dig: fortvivl!

Saa vælg da Fortvivlelsen, thi Fortvivlelsen selv er et Valg, thi man kan tvivle uden at vælge det, men fortvivle kan man ikke uden at vælge det. Og idet man fortvivler, vælger man atter, og hvad vælger man da, man vælger sig selv, ikke i sin Umiddelbarhed, ikke som dette tilfældige Individ, men man vælger sig selv i sin evige Gyldighed.

Dette Punkt skal jeg lidt nærmere stræbe at belyse med Hensyn til Dig. Der har i den nyere Philosophi været meer end tilstrækkelig Tale om, at al Speculation begynder med Tvivl; derimod har jeg, forsaavidt jeg leilighedsviis kunde beskæftige mig med saadanne Overveielser, forgjæves søgt en Oplysning om, hvori Tvivlen er forskjellig fra Fortvivlelse. Jeg vil her forsøge at oplyse denne Forskjel, i det Haab, at det vil bidrage til at orientere Dig og stille Dig rigtigt. Jeg er langt fra at tiltroe mig nogen egentlig philosophisk Dygtighed, jeg har ikke Din Virtuositet i at spille med Kategorier, men hvad der i dybeste Forstand er Livets Betydning, det maa dog vel ogsaa kunne fattes af et mere eenfoldigt Menneske. Tvivl er Tankens Fortvivlelse, Fortvivlelse er Personlighedens Tvivl, derfor er det, jeg holder saa fast paa den Bestemmelse at vælge, der er mit Løsen, Nerven i min Livs-Anskuelse, og en saadan har jeg, om jeg end ingenlunde anmasser mig at have et System. Tvivl er den indre Bevægelse i Tanken selv, og i min Tvivl forholder jeg mig saa upersonligt som muligt. Jeg antager nu, at Tanken, idet Tvivlen gjennemføres, finder det Absolute og hviler deri, saa hviler den deri ikke ifølge et Valg men ifølge samme Nødvendighed, ifølge hvilken den tvivlede; thi Tvivlen selv er en Bestemmelse af Nødvendighed, og Hvilen ligeledes. Det er det Ophøiede i Tvivlen, hvorfor den saa ofte er bleven anpriist og udskreget af Folk, der neppe forstode, hvad de sagde. Men dette, at det er en Nødvendighedens Bestemmelse, viser, at den hele Personlighed ikke er med i Bevægelse. Der er derfor noget meget sandt i, at et Menneske siger: jeg vilde gjerne troe, jeg kan ikke, jeg maa tvivle. Man seer derfor ogsaa ofte, at en Tvivler dog kan eie et positivt Gehalt i sig selv, som lever udenfor al Communication med Tanken, at han kan være et høist samvittighedsfuldt Menneske, der ingenlunde tvivler om Pligtens Gyldighed, og om Reglen for sin Handling, ingenlunde tvivler om en Mængde sympathetiske Følelser og Stemninger. Man seer paa den anden Side, især i vor Tid, Mennesker, der have Fortvivlelsen i Hjertet, og som dog have beseiret Tvivlen. Det har især været mig paafaldende ved at betragte enkelte af Tydsklands Philosopher. Deres Tanke er beroliget, den objective logiske Tanke er bragt til Hvile i sin tilsvarende Objectivitet, og dog ere de Fortvivlede, om de end adsprede sig ved den objective Tænkning, thi et Menneske kan adsprede sig paa mange Maader, og der er neppe noget saa bedøvende Middel som abstrakt Tænkning, fordi det der gjelder om at forholde sig saa upersonligt som muligt. Tvivl og Fortvivlelse høre altsaa hjemme i aldeles forskjellige Sphærer, det er forskjellige Sider af Sjælen, der blive satte i Bevægelse. Dog hermed er jeg ingenlunde tilfredsstillet, thi saa bleve Tvivl og Fortvivlelse hinanden sideordnede, og dette er ikke Tilfældet. Fortvivlelse er et langt dybere og fuldstændigere Udtryk, dens Bevægelse langt mere omfattende end Tvivlens. Fortvivlelse er netop et Udtryk for den hele Personlighed, Tvivl kun for Tanken. Den formeentlige Objectivitet, Tvivlen har, hvoraf den er saa fornem, er netop et Udtryk for dens Ufuldkommenhed. Tvivl ligger derfor i Differensen, Fortvivlelse i det Absolute. Der hører Talent til at tvivle, men der hører slet ikke Talent til at fortvivle; men Talent er som saadant en Differens, og hvad der, for at gjøre sig gjeldende, fordrer en Differens, kan aldrig være det Absolute; thi det Absolute kan kun være som det Absolute for det Absolute. Det ringeste, mindst begavede Menneske kan fortvivle, en ung Pige, der er Intet mindre end Tænker, kan fortvivle, hvorimod Enhver let føler det Taabelige i at sige om disse, at de ere Tvivlere. Grunden, hvorfor et Menneskes Tvivl kan være beroliget, og han dog kan være fortvivlet, og dette saaledes kan gaae hen, er, at han ikke i dybere Forstand vil Fortvivlelsen. Man kan overhovedet slet ikke fortvivle, uden man vil det, men for i Sandhed at fortvivle maa man i Sandhed ville det, men naar man i Sandhed vil det, saa er man i Sandhed ude over Fortvivlelsen; naar man i Sandhed har valgt Fortvivlelsen, saa har man i Sandhed valgt det, Fortvivlelsen vælger: sig selv i sin evige Gyldighed. Først i Fortvivlelsen er Personligheden beroliget, ikke med Nødvendighed, thi jeg fortvivler aldrig nødvendigt, men med Frihed, og først deri er det Absolute vundet. I denne Henseende mener jeg, at vor Tid vil gjøre et Fremskridt, forsaavidt jeg ellers kan have nogen Mening om vor Tid, da jeg kun kjender den af min Avis-Læsning og at et enkelt Skrift, eller af Samtale med Dig. Den Tid vil ikke være langt borte, da man, dyrt nok maaskee, vil erfare, at det sande Udgangspunkt for at finde det Absolute ikke er Tvivl, men Fortvivlelse.

Dog, jeg vender tilbage til min Kategori, jeg er ikke Logiker, jeg har kun een saadan, men jeg forsikkrer Dig til, den er baade mit Hjertes og min Tankes Valg, min Sjæls Lyst og min Salighed – jeg vender tilbage til Betydningen af det at vælge. Idet jeg da vælger absolut, vælger jeg Fortvivlelsen, og i Fortvivlelsen vælger jeg det Absolute, thi jeg er selv det Absolute, jeg sætter det Absolute og jeg er selv det Absolute; men som aldeles identisk hermed maa jeg sige: jeg vælger det Absolute, der vælger mig, jeg sætter det Absolute, der sætter mig; thi erindrer jeg ikke, at dette andet Udtryk er ligesaa absolut, saa er min Kategori at vælge usand; thi den er netop Identiteten af begge. Hvad jeg vælger, det sætter jeg ikke, thi dersom det ikke var sat, saa kunde jeg ikke vælge det, og dog, dersom jeg ikke satte det derved, at jeg valgte det, saa valgte jeg det ikke. Det er, thi dersom det ikke var, kunde jeg ikke vælge det; det er ikke, thi det bliver først derved, at jeg vælger det, og ellers var mit Valg en Illusion.

Men hvad er det da, jeg vælger, er det Dette eller Hiint? Nei, thi jeg vælger absolut, og absolut vælger jeg jo netop derved, at jeg har valgt ikke at vælge Dette eller Hiint. Jeg vælger det Absolute, og hvad er det Absolute? Det er mig selv i min evige Gyldighed. Noget Andet end mig selv kan jeg aldrig vælge som det Absolute, thi vælger jeg noget Andet, da vælger jeg det som en Endelighed, og vælger det altsaa ikke absolut. Selv Jøden, der valgte Gud, valgte ikke absolut, thi vel valgte han det Absolute, men han valgte det ikke absolut, og derved ophørte det at være det Absolute, og blev en Endelighed.

Men hvad er da dette mit Selv? Dersom jeg vilde tale om et første Øieblik, et første Udtryk derfor, saa er mit Svar: det er det Abstrakteste af Alt, der dog tillige i sig er det Concreteste af Alt – det er Friheden. Lad mig her anstille en lille psychologisk Iagttagelse. Man hører ofte nok Folk give deres Misfornøielse Luft i Klage over Livet, man hører dem ofte nok ønske. Tænk Dig nu en saadan Stymper; lad os springe de Ønsker over, der her Intet oplyse, fordi de ligge i det aldeles Tilfældige. Han ønsker: gid jeg havde hiint Menneskes Aand, eller hiin Mands Talent og saa videre, ja, for at gaae til den yderste Spidse –: gid jeg havde hiint Menneskes Fasthed. Slige Ønsker hører man hyppigt nok, men har Du nogensinde hørt et Menneske for Alvor ønske, at han kunde blive en Anden, det er saa langt fra, at det netop er charakteristisk for hvad man kalder ulykkelige Individualiteter, at de klamre sig mest fast om sig selv, at de trods alle deres Lidelser dog for al Verdens Pris ikke vilde ønske at være en Anden, Noget, der har sin Grund i, at saadanne Individualiteter ere Sandheden meget nær, og de føle Personlighedens evige Gyldighed, ikke i dens Velsignelse men i dens Qval, om de end have beholdt dette aldeles abstrakte Udtryk for Glæden deri tilbage, at de dog helst ville blive sig selv. Men nu han med de mange Ønsker, han mener dog bestandig at blive sig selv, uagtet Alt blev forandret. Altsaa er der i ham selv Noget, der er absolut i Forhold til alt Andet, Noget, hvorved han er den, han er, om saa end Forandringen, han ved sit Ønske opnaaede, var den størst mulige. At han er i en Misforstaaelse, skal jeg senere vise, men her vil jeg blot finde det abstrakteste Udtryk for dette »Selv«, der gjør ham til den, han er. Og dette er ikke Andet end Friheden. Der lod sig virkelig ad denne Vei føre et høist plausibelt Beviis for Personlighedens evige Gyldighed; ja selv en Selvmorder vil dog egentlig ikke af med sit Selv, ogsaa han ønsker, han ønsker en anden Form for sit Selv, og man kan derfor vel finde en Selvmorder, der i høieste Grad var overbeviist om Sjælens Udødelighed, men hvis hele Væsen var saa hildet, at han ved dette Skridt troede at finde den absolute Form for sin Aand.

Dog, Grunden, hvorfor det kan forekomme et Individ saaledes, som om han bestandig kunde forandres og dog blev den samme, som om hans inderste Væsen var en algebraisk Størrelse, der kunde betegne, hvad det skulde være, ligger i, at han er urigtigt stillet, at han ikke har valgt sig selv, ikke har Forestilling derom, og dog ligger der selv i hans Uforstand en Anerkjendelse af Personlighedens evige Gyldighed. Den, der derimod er rigtigt stillet, ham gaaer det anderledes. Han vælger sig selv, ikke i endelig Forstand, thi saa blev jo dette »Selv« en Endelighed, der gik med mellem andre Endeligheder, men i absolut Forstand, og dog vælger han jo sig selv og ikke en Anden. Dette Selv, han saaledes vælger, er uendelig concret, thi det er ham selv, og dog er det absolut forskjelligt fra hans tidligere Selv, thi han har valgt det absolut. Dette »Selv« har ikke været til før, thi det blev til ved Valget, og dog har det været til, thi det var jo »ham selv.«

Valget gjør her paa eengang de to dialektiske Bevægelser, det, der vælges, er ikke til og bliver til ved Valget, det, der vælges, er til, ellers var det ikke et Valg. Dersom nemlig det, jeg valgte, ikke var til, men blev absolut til ved Valget, saa valgte jeg ikke, saa skabte jeg; men jeg skaber ikke mig selv, jeg vælger mig selv. Medens derfor Naturen er skabt af Intet, medens jeg selv som umiddelbar Personlighed er skabt af Intet, saa er jeg som fri Aand født af Modsigelsens Grundsætning, eller født derved, at jeg valgte mig selv.

Han opdager nu, at det »Selv«, han vælger, har en uendelig Mangfoldighed i sig, forsaavidt det har en Historie, en Historie, i hvilken han vedkjender sig Identiteten med sig selv. Denne Historie er af forskjellig Art, thi i denne Historie staaer han i Forhold til andre Individer i Slægten og til hele Slægten, og denne Historie indeholder noget Smerteligt, og dog er han kun den, han er, ved denne Historie. Derfor hører der Mod til at vælge sig selv; thi paa samme Tid, som det synes, at han isolerer sig allermest, paa samme Tid fordyber han sig allermest i den Rod, ved hvilken han hænger sammen med det Hele. Det ængster ham, og dog maa det være saa; thi naar Frihedens Lidenskab er vaagnet i ham, – og den er vaagnet i Valget, ligesom den forudsætter sig selv i Valget – saa vælger han sig selv og kæmper for denne Besiddelse som for sin Salighed, og det er hans Salighed. Han kan Intet opgive af det Altsammen, ikke det Smerteligste, ikke det Tungeste, og dog er Udtrykket for denne Kamp, for denne Erhverven – Anger. Han angrer sig tilbage i sig selv, tilbage i Familien, tilbage i Slægten, indtil han finder sig selv i Gud. Kun paa dette Vilkaar kan han vælge sig selv, og dette er det eneste Vilkaar, han vil, thi kun saaledes kan han absolut vælge sig selv. Hvad er dog et Menneske uden Kjærlighed? Men der er mange Arter af Kjærlighed; jeg elsker en Fader anderledes end en Moder, min Hustru atter anderledes, og enhver forskjellig Kjærlighed har sit forskjellige Udtryk, men der gives ogsaa en Kjærlighed, hvormed jeg elsker Gud, og denne har kun eet Udtryk i Sproget, det er: Anger. Naar jeg ikke elsker ham saaledes, da elsker jeg ham ikke absolut, ikke af mit inderste Væsen, enhver anden Kjærlighed til det Absolute er en Misforstaaelse, thi for at tage, hvad man ellers saa høit anpriser, og hvad jeg selv ærer, naar Tanken af al sin Kjærlighed hænger fast ved det Absolute, saa er det ikke det Absolute, jeg elsker, jeg elsker ikke absolut, thi jeg elsker nødvendigt; saasnart jeg elsker frit og elsker Gud, da angrer jeg. Og skulde der ingen anden Grund være til, at Udtrykket for min Kjærlighed til Gud var Anger, saa er der den, at han har elsket mig først. Og dog er dette en ufuldkommen Betegnelse; thi kun naar jeg vælger mig selv som skyldig, vælger jeg absolut mig selv, hvis jeg overhovedet skal vælge mig selv absolut saaledes, at det ikke er identisk med at skabe sig selv; og om det var Faderens Brøde, der var gaaet i Arv paa Sønnen, han angrer denne med; thi kun saaledes kan han vælge sig selv, vælge sig absolut, og om Taarerne næsten vilde udslette ham Alt, han holder ved at angre; thi kun saaledes vælger han sig selv. Hans Selv er ligesom udenfor ham, og det skal erhverves, og Angeren er hans Kjærlighed dertil, fordi han vælger det absolut, af den evige Guds Haand.

Hvad jeg her har fremsat er ikke Katheder-Viisdom, det er Noget, ethvert Menneske kan fremsætte, der vil det, og som ethvert Menneske kan ville, naar han vil. Jeg har ikke lært det i Høresalene, jeg har lært det i Dagligstuen, eller om Du vil, i Børnestuen, thi naar jeg seer min lille Søn løbe over Gulvet, saa glad, saa lykkelig, da tænker jeg, hvo veed, om jeg dog ikke har havt megen skadelig Indflydelse paa ham. Gud veed det, jeg bærer al mulig Omsorg for ham, men denne Tanke beroliger mig ikke. Da siger jeg til mig selv, der vil komme et Øieblik i hans Liv, da ogsaa hans Aand vil modnes i Valgets Øieblik, da vil han vælge sig selv, da vil han angre ogsaa hvad Skyld der fra mig kan hvile paa ham. Og det er smukt, at en Søn angrer Faderens Brøde, og dog vil han ikke gjøre det for min Skyld, men fordi han kun saaledes kan vælge sig selv. Lad der saa skee, hvad der skal skee, ofte kan det, man anseer for det Bedste, have de skadeligste Følger for et Menneske, men Alt dette er dog Intet. Jeg kan gavne ham meget, det skal jeg bestræbe mig for, men til det Høieste kan han kun gjøre sig selv. See, derfor er det, at Mennesker have saa ondt ved at vælge sig selv, fordi den absolute Isolation er her identisk med den dybeste Continuitet, fordi der, saalænge man ikke har valgt sig selv, er ligesom en Mulighed af at blive noget anderledes, enten paa den ene eller paa den anden Maade.

See, her har Du min ringe Mening om hvad det er at vælge og at angre. Det er usømmeligt at elske en ung Pige som var det Eens Moder, eller Eens Moder, som var det en ung Pige, hver Kjærlighed har sin Eiendommelighed, Kjærlighed til Gud har sin absolute Eiendommelighed, dens Udtryk er Anger. Og hvad er dog al anden Kjærlighed i Sammenligning med denne, den er kun Barnelallen derimod. Jeg er ikke et ophidset ungt Menneske, der søger at anbefale sine Theorier, jeg er Ægtemand, jeg tør vel lade min Kone høre paa, at dog al Kjærlighed i Sammenligning med Angeren kun er Barnelallen, og dog veed jeg, jeg er en brav Ægtemand, »jeg, der endnu som Ægtemand kæmper under den første Kjærligheds seirrige Banner«; jeg veed, hun deler min Anskuelse, og derfor elsker jeg hende endnu høiere, og derfor vilde jeg ikke elskes af hiin unge Pige, fordi hun ikke delte denne Anskuelse.

At der her atter viser sig nye og forfærdelige Afveie, at den ikke saa let er udsat for at falde, der kryber langs Jorden, som den, der bestiger Bjergenes Toppe, den ikke saa let udsat for at fare vild, der bliver i Kakkelovnskrogen, som den, der vover sig ud i Verden, det veed jeg, men derfor bliver jeg lige trøstig i mit Valg.

Fra dette Punkt vil nu en Theolog finde et Udgangspunkt for en Mangfoldighed af Betragtninger; jeg vil ikke gaae videre ind paa disse, da jeg kun er Lægmand. Kun det Foregaaende vil jeg søge at belyse med den Bemærkning, at først i Christendommen har Angeren fundet sit sande Udtryk. Den fromme Jøde følte Fædrenes Brøde hvile paa sig, og dog følte han det ikke nær saa dybt som den Christne; thi den fromme Jøde kunde ikke angre det, fordi han ikke absolut kunde vælge sig selv. Forfædrenes Skyld tyngede paa ham, rugede over ham, han segnede under denne Byrde, han sukkede, men han kunde ikke løfte den op, det kan kun den, der absolut vælger sig selv, ved Hjelp af Angeren. Jo større Frihed, jo større Skyld, og det er Salighedens Hemmelighed, og det er om ikke Feighed, saa dog en Sjælens Kleinmodighed, ikke at ville angre Forfædrenes Skyld, det er, om ikke Usselhed, saa dog Smaalighed og Mangel paa Høimod.

Fortvivlelsens Valg er altsaa »mig selv«; thi vel er det sandt, at idet jeg fortvivler, fortvivler jeg, som over alt Andet, saaledes ogsaa over mig selv; men det Selv, jeg fortvivler over, er en Endelighed, ligesom enhver anden Endelighed, det »Selv«, jeg vælger, er det absolute »Selv«, eller mit Selv efter sin absolute Gyldighed. Naar dette forholder sig saaledes, saa vil Du her igjen see, hvorfor jeg i det Foregaaende og endnu bestandig siger, at det Enten–Eller, jeg opstillede mellem at leve æsthetisk og ethisk, ikke er et fuldstændigt Dilemma, fordi der egentlig kun er Tale om eet Valg. Ved dette Valg vælger jeg egentlig ikke mellem Godt og Ondt, men jeg vælger det Gode, men idet jeg vælger det Gode, vælger jeg eo ipso Valget mellem Godt og Ondt. Det oprindelige Valg er stedse tilstede i ethvert følgende Valg.

Saa fortvivl da, og Dit Letsind skal aldrig mere bringe Dig til at vandre som en ustadig Aand, som et Gjenfærd mellem Ruinerne af en Verden, der dog er tabt for Dig; fortvivl, og Din Aand skal aldrig mere sukke i Tungsind, thi Verden skal atter blive Dig skjøn og glædelig, om Du end seer paa den med andre Øine end før, og Din Aand skal frigjort svinge sig op i Frihedens Verden.

[…]

Deri ligger nemlig et Menneskes evige Værdighed, at han kan faae Historie, deri ligger det Guddommelige i ham, at han selv, hvis han vil, kan give denne Historie Continuitet; thi det faaer den først, naar den ikke er Indbegrebet af hvad der er skeet eller hændt mig, men min egen Gjerning, saaledes, at selv det, der er hændt mig, ved mig er forvandlet og overført fra Nødvendighed til Frihed. Dette er det Misundelsesværdige ved et Menneskeliv, at man kan komme Guddommen til Hjælp, kan forstaae ham, og det er igjen den eneste et Menneske værdige Maade at forstaae ham paa, at man i Frihed tilegner sig Alt, hvad der tilkommer Een, baade det Glade og det Sørgelige. Eller synes det Dig ikke saa? Mig forekommer det saaledes, ja, jeg synes, at man blot behøvede at sige det høit til et Menneske, for at gjøre ham misundelig paa sig selv.

[…]

Her vil jeg nu erindre om den Bestemmelse, jeg i det Foregaaende gav af det Ethiske, at det er det, hvorved et Menneske bliver det, han bliver. Det vil da ikke gjøre Individet til et andet, men til det selv; det vil ikke tilintetgjøre det Æsthetiske, men forklare det. For at et Menneske skal leve ethisk, er det nødvendigt, at han bliver sig bevidst, saa gjennemgribende, at ingen Tilfældighed undgaaer ham. Denne Concretion vil det Ethiske ikke udslette, men seer i denne sin Opgave, seer det, hvoraf det skal danne, og det, det skal danne. I Almindelighed betragter man det Ethiske aldeles abstrakt og har derfor en hemmelig Gru for det. Det Ethiske betragtes da som noget Personligheden Fremmed, og man krymper sig ved at hengive sig til det, da man dog ikke ret kan være sikker paa, hvad det i Tidens Længde kan føre til. Saaledes frygte ogsaa mange Mennesker for Døden, fordi de nære dunkle og uklare Forestillinger om, at Sjælen i Døden skulde gaae over i en anden Tingenes Orden, hvor der herskede Love og Vedtægter, der vare aldeles forskjellige fra dem, de i denne Verden havde lært at kjende. Grunden til en saadan Frygt for Døden er da Individets Utilbøielighed til at blive sig selv gjennemsigtigt, thi saafremt man vil det, saa indseer man let det Urimelige i denne Frygt. Saaledes ogsaa med det Ethiske; naar et Menneske frygter Gjennemsigtighed, saa flyer han altid det Ethiske, thi Andet vil dette egentlig ikke.

I Modsætning til en æsthetisk Livs-Anskuelse, der vil nyde Livet, hører man ofte omtale en anden Livs-Anskuelse, der sætter Livets Betydning i, at leve for Opfyldelsen af sine Pligter. Hermed vil man da betegne en ethisk Livs-Anskuelse. Udtrykket er imidlertid meget ufuldkomment, og man skulde næsten troe, at det var opfundet for at bringe det Ethiske i Miscredit; saa meget er vist, at man i vor Tid ofte seer det brugt saaledes, at man næsten kommer til at smile, saaledes naar Scribe lader denne Sætning foredrage med en vis lavcomisk Alvor, der danner en meget misrecommanderende Modsætning til Nydelsens Glæde og Munterhed. Feilen er den, at Individet sættes i et udvortes Forhold til Pligten. Det Ethiske bestemmes som Pligt, og Pligt igjen som en Mangfoldighed af enkelte Sætninger, men Individet og Pligten staae udenfor hinanden. Et saadant Pligtliv er naturligviis saare uskjønt og kjedsommeligt, og dersom det Ethiske ikke havde et langt dybere Sammenhæng med Personligheden, saa vilde det altid blive saare vanskeligt at forfægte det ligeoverfor det Æsthetiske. At der gives mange Mennesker, der ikke komme videre, vil jeg ikke negte; men det ligger ikke i Pligten, men i Menneskene.

Besynderligt nok er det, at man ved det Ord: Pligt, kan komme til at tænke paa et udvortes Forhold, da dette Ords Derivation selv antyder, at det er et indvortes Forhold; thi det, der paaligger mig, ikke som dette tilfældige Individ, men efter mit sande Væsen, det staaer jo dog vel i det inderligste Forhold til mig selv. Pligten er nemlig ikke et Paalæg, men Noget, der paaligger. Naar Pligten sees saaledes, da er dette et Tegn paa, at Individet er orienteret i sig selv. Pligten vil da ikke adsplitte sig for ham i en Mangfoldighed af enkelte Bestemmelser; thi dette tyder altid paa, at han kun staaer i et udvortes Forhold til den. Han har iført sig Pligten, den er ham Udtrykket for hans inderste Væsen. Naar han saaledes har orienteret sig i sig selv, da har han fordybet sig i det Ethiske, og han vil ikke jage sig selv stakaandet paa at opfylde sine Pligter. Det sande ethiske Individ har derfor en Ro og Sikkerhed i sig selv, fordi han ikke har Pligten udenfor sig men i sig. Jo dybere et Menneske ethisk har anlagt sit Liv, desto mindre vil han føle Trang til, hvert Øieblik at tale om Pligten, hvert Øieblik at ængstes for, om han opfylder den, hvert Øieblik at raadføre sig med Andre om, hvad der dog er hans Pligt. Naar det Ethiske sees rigtigt, gjør det Individet uendelig sikkert i sig selv, naar det ikke sees rigtigt, gjør det Individet aldeles usikkert, og jeg kan ikke tænke mig en ulykkeligere eller qualfuldere Existens, end naar et Menneske har faaet Pligten udenfor sig, og dog bestandig vil realisere den.

Seer man det Ethiske udenfor Personligheden og i et udvortes Forhold til denne, saa har man opgivet Alt, saa har man fortvivlet. Det Æsthetiske som saadant er Fortvivlelse, det Ethiske er det Abstrakte og som saadant uformuende til at frembringe det Mindste. Naar man derfor stundom seer Mennesker med en vis redelig Iver slide og slæbe for at realisere det Ethiske, der som en Skygge bestandig flygter, saasnart de gribe efter den, saa er det baade comisk og tragisk.

Det Ethiske er det Almene og saaledes det Abstrakte. I sin fuldkomne Abstraktion er derfor det Ethiske altid forbydende. Saaledes viser det Ethiske sig som Lov. Saasnart det Ethiske er befalende, saa har det allerede Noget af det Æsthetiske i sig. Jøderne vare Lovens Folk. De forstode derfor ypperlig de fleste Bud i Moseloven; men det Bud, de ikke synes at have forstaaet, var det Bud, som Christendommen nærmest knyttede sig til: Du skal elske Gud af Dit ganske Hjerte. Dette Bud er heller ikke negativt, er heller ikke abstrakt, det er i høieste Grad positivt og i høieste Grad concret. Naar det Ethiske bliver concretere, da gaaer det over i Bestemmelsen af Sæder. Men Realiteten af det i denne Henseende Ethiske ligger i Realiteten af en folkelig Individualitet, og her har da allerede det Ethiske optaget et æsthetisk Moment i sig. Dog er det Ethiske endnu abstrakt og lader sig ikke realisere tilfulde, fordi det ligger udenfor Individet. Først naar Individet selv er det Almene, først da lader det Ethiske sig realisere. Det er denne Hemmelighed, der ligger i Samvittigheden, det er denne Hemmelighed, det individuelle Liv har med sig selv, at det paa eengang er et individuelt Liv og tillige det Almene, om ikke umiddelbart som saadant, saa dog efter sin Mulighed. Den, der betragter Livet ethisk, han seer det Almene, og den, der lever ethisk, han udtrykker i sit Liv det Almene, han gjør sig til det almene Menneske, ikke derved, at han affører sig sin Concretion, thi saa bliver han til slet Intet; men derved, at han ifører sig den og gjennemtrænger den med det Almene. Det almene Menneske er nemlig ikke et Phantom, men ethvert Menneske er det almene Menneske, det vil sige, ethvert Menneske er anviist den Vei, ad hvilken han bliver det almene Menneske. Den, der lever æsthetisk, han er det tilfældige Menneske, han troer at være det fuldkomne Menneske derved, at han er det eneste Menneske; den, der lever ethisk, arbeider hen til at blive det almene Menneske. Naar saaledes et Menneske er æsthetisk forelsket, saa spiller det Tilfældige en uhyre Rolle, og det er ham af Vigtighed, at der er Ingen, der har elsket saaledes, med de Nuancer, som han; naar den, der lever ethisk, gifter sig, saa realiserer han det Almene. Derfor bliver han ingen Hader af det Concrete, men han har et Udtryk mere, dybere end ethvert æsthetisk Udtryk, idet han i Kjærligheden seer en Aabenbarelse af det Almeen-Menneskelige. Den, der lever ethisk, har da sig selv som sin Opgave. Hans »Selv« er som umiddelbart tilfældigt bestemt, og Opgaven er den, at sammenarbeide det Tilfældige og det Almene.

Det ethiske Individ har da ikke Pligten udenfor sig, men i sig; i Fortvivlelsens Øieblik kommer dette tilsyne og arbeider sig nu frem gjennem det Æsthetiske i og med dette. Om det ethiske Individ kan man sige, at han er som det stille Vand, der har den dybe Grund, hvorimod den, der lever æsthetisk, blot er overfladisk bevæget. Naar derfor det ethiske Individ har fuldendt sin Opgave, har stridt den gode Strid, saa er han kommen dertil, at han er bleven det eneste Menneske, det vil sige, at der er intet Menneske saaledes som han, og tillige, at han er bleven det almene Menneske. At være det eneste Menneske, er i og for sig ikke noget saa Stort, thi det har ethvert Menneske tilfælleds med enhver Naturfrembringelse; men at være det saaledes, at han deri tillige er det Almene, det er den sande Livs-Kunst.

Personligheden har da det Ethiske ikke udenfor sig, men i sig, og det bryder frem af dette Dyb. Det gjælder da, som sagt, om, at den ikke i en abstrakt og indholdsløs Stormen tilintetgjør det Concrete, men assimilerer sig det. Da det Ethiske saaledes ligger dybest i Sjælen, saa falder det ikke altid i Øinene, og et Menneske, der lever ethisk, kan gjøre aldeles det Samme som den, der lever æsthetisk, saaledes, at det længe kan skuffe, men tilsidst kommer der et Øieblik, hvor det viser sig, at den, der lever ethisk, har en Grændse, som den Anden ikke kjender. I denne Forvisning om, at hans Liv er anlagt ethisk, hviler Individet med tryg Sikkerhed og plager derfor ikke sig selv og Andre med spidsfindig Ængstelighed over dette eller hiint. At nemlig den, der lever ethisk, har et heelt Spatium for det Indifferente, finder jeg ganske i sin Orden, og det er netop en Veneration for det Ethiske, at man ikke vil tvinge det ind i enhver Ubetydelighed. En saadan Stræben, der altid mislykkes, finder man ogsaa kun hos dem, der ikke have Mod til at troe paa det Ethiske, og som i dybere Forstand mangle indre Sikkerhed. Der gives Mennesker, hvis Pussillanimitet man netop kjender derpaa, at de aldrig kunne blive færdige med det Totale, fordi dette netop for dem er det Mangfoldige, men disse ligge ogsaa udenfor det Ethiske, naturligviis ikke af nogen anden Grund end formedelst Villiens Svaghed, der som al anden Aandssvaghed kan betragtes som en Art Afsindighed. Saadanne Menneskers Liv gaae op i at afsie Myggen. De have hverken Forestilling om det Ethiskes skjønne og rene Alvor, eller om det Indifferentes sorgløse Glæde. Dog er naturligviis det Indifferente for det ethiske Individ dethroniseret, og han kan hvert Øieblik sætte det Grændse. Saaledes troer man ogsaa, at der er et Forsyn til, og Sjælen hviler trygt i denne Forvisning, og dog vilde man ikke falde paa at forsøge at gjennemtrænge enhver Tilfældighed med denne Tanke, eller hvert Minut at blive sig denne Tro bevidst. At ville det Ethiske uden at forstyrres af det Indifferente, at troe paa et Forsyn uden at forstyrres af Tilfældighed er en Sundhed, der kan erhverves og bevares, naar et Menneske selv vil det. Ogsaa i denne Henseende gjælder det om at see Opgaven, at denne, forsaavidt et Menneske har Tilbøielighed til saaledes at adspredes, er den, at gjøre Modstand, at holde det Uendelige fast, og ikke at rende med Liimstangen.

Den, der ethisk vælger sig selv, han har sig selv som sin Opgave, ikke som en Mulighed, ikke som et Legetøi for sin Vilkaarligheds Spil. Ethisk kan han kun vælge sig, naar han vælger sig i Continuitet, og han har saaledes sig selv som en mangfoldig bestemmet Opgave. Denne Mangfoldighed søger han ikke at udslette eller forflygtige, tvertimod angrer han sig fast i den, fordi denne Mangfoldighed er ham selv, og kun ved angrende at fordybe sig i den, kan han komme til sig selv, da han ikke antager, at Verden begynder med ham, eller at han skaber sig selv; det Sidste har Sproget selv stemplet med Foragt, og man siger altid foragteligt om et Menneske: han skaber sig. Men idet han angrende vælger sig selv, er han handlende, ikke i Retning af Isolation, men i Retning af Continuitet.

Lad os nu engang sammenstille et ethisk og et æsthetisk Individ. Hovedforskjellen, hvorom Alt dreier sig, er den, at det ethiske Individ er sig selv gjennemsigtigt og ikke lever ins Blaue hinein, som det æsthetiske Individ gjør det. Med denne Forskjel er Alt givet. Den, der lever ethisk, har seet sig selv, kjender sig selv, gjennemtrænger med sin Bevidsthed sin hele Concretion, tillader ikke ubestemte Tanker at pusle om i ham, ikke fristende Muligheder at adsprede ham med deres Gjøgleværk, han er ikke sig selv som et Hexebrev, der snart kan komme Eet ud af, snart et Andet, alt eftersom man vender og dreier det. Han kjender sig selv. Det Udtryk γνωϑι σεαυτον, er ofte nok blevet gjentaget, og man har deri seet Maalet for hele Menneskets Stræben. Det er ogsaa ganske rigtigt, men dog er det ligesaa vist, at det ikke kan være Maalet, dersom det ikke tillige er Begyndelsen. Det ethiske Individ kjender sig selv, men dette Kjendskab er ikke en blot Contemplation, thi saa bliver Individet bestemmet efter dets Nødvendighed, det er en Besindelse paa sig selv, der selv er en Handling, og derfor har jeg med Flid brugt det Udtryk: at vælge sig selv, istedetfor at kjende sig selv. Idet da Individet kjender sig selv, er han ikke færdig, tvertimod er dette Kjendskab i høi Grad frugtbart, og af dette Kjendskab fremgaaer det sande Individ. Dersom jeg vilde være aandrig, kunde jeg her sige, at Individet paa en lignende Maade kjendte sig selv, som naar det i det Gamle Testamente hedder, at Adam kjendte Eva. Ved Individets Omgang med sig selv besvangres Individet med sig selv og føder sig selv. Det Selv, Individet kjender, er paa eengang det virkelige Selv og det ideale Selv, som Individet har udenfor sig som det Billede, i hvis Lighed det skal danne sig, og som det dog paa den anden Side har i sig, da det er det selv. Kun i sig selv har Individet det Formaal, hvorefter han skal stræbe, og dog har han dette Formaal udenfor sig, idet han stræber derefter. Troer nemlig Individet, at det almene Menneske ligger udenfor ham, at det udenfra skal komme ham imøde, saa er han desorienteret, saa har han en abstrakt Forestilling, og hans Methode bliver altid en abstrakt Tilintetgjørelse af det oprindelige Selv. Kun i sig selv kan Individet faae Oplysning om sig selv. Derfor har det ethiske Liv denne Dobbelthed, at Individet har sig selv udenfor sig selv i sig selv. Det typiske Selv er imidlertid det ufuldkomne Selv, thi det er kun en Propheti, og derfor ikke det virkelige. Imidlertid ledsager det ham bestandigt; men jo mere han realiserer det, jo mere svinder det ind i ham, indtil det tilsidst, istedetfor at vise sig foran ham, ligger bag ved ham som en afbleget Mulighed. Det gaaer med dette Billede som med Menneskets Skygge. Om Morgenen kaster Mennesket sin Skygge foran sig, om Middagen gaaer den næsten ubemærkelig ved Siden af ham, om Aftenen falder den bag ved ham. Naar Individet har kjendt sig selv og har valgt sig selv, saa er han ifærd med at realisere sig selv, men da han frit skal realisere sig selv, saa maa han vide, hvad det er, han vil realisere. Det han vil realisere, er da vel sig selv, men det er hans ideale Selv, som han dog intet andetsteds faaer end i sig selv. Holder man ikke fast paa, at Individet har det ideale Selv i sig selv, saa bliver hans Digten og Tragten abstrakt. Den, der vil copiere et andet Menneske, og den, der vil copiere det normale Menneske, blive, om end paa forskjellig Maade, begge lige affecteerte.

Det æsthetiske Individ betragter sig selv i sin Concretion og distinguerer nu inter et inter. Han seer Noget som tilfældigt ham tilhørende, Andet som væsentligt. Denne Distinction er imidlertid yderst relativ; thi saalænge et Menneske blot lever æsthetisk, tilhører Alt ham egentlig lige tilfældigt, og det er blot Mangel paa Energi, naar et æsthetisk Individ fastholder denne Distinction. Det ethiske Individ har i Fortvivlelsen lært dette, han har derfor en anden Distinction; thi ogsaa han adskiller mellem det Væsentlige og det Tilfældige. Alt hvad der er sat ved hans Frihed tilhører ham væsentligt, hvor tilfældigt det end kan synes at være, Alt hvad der ikke er det, er ham tilfældigt, hvor væsentligt det end kan synes at være. Denne Distinction er imidlertid for det ethiske Individ ikke en Frugt af hans Vilkaarlighed, saa det kunde synes, at han havde Magtfuldkommenhed til at gjøre sig selv til hvad han vilde. Vel tør nemlig det ethiske Individ bruge det Udtryk, at han er sin egen Redacteur, men han er sig tillige fuldelig bevidst, at han er ansvarhavende, ansvarhavende for sig selv i personlig Betydning, forsaavidt det vil have afgjørende Indflydelse paa ham selv hvad han vælger; ansvarhavende ligeover for den Tingenes Orden, i hvilken han lever, ansvarhavende ligeover for Gud. Naar man seer det saaledes, saa troer jeg, Distinctionen er rigtig; thi væsentlig tilhører dog kun det mig, hvad jeg ethisk overtager som en Opgave. Negter jeg at overtage det, da tilhører det mig væsentligt, at jeg har negtet det. Naar et Menneske æsthetisk betragter sig selv, saa distinguerer han maaskee saaledes. Han siger: jeg har Talent til at male, det anseer jeg for en Tilfældighed; men jeg har Vittighed og Skarpsindighed, det anseer jeg som det Væsentlige, som ikke kan tages fra mig, uden at jeg bliver en Anden. Hertil vilde jeg svare: denne hele Distinction er en Illusion; thi naar Du ikke overtager denne Vittighed og Skarpsindighed ethisk, som en Opgave, som Noget, Du er ansvarlig for, saa tilhører den Dig ikke væsentlig, og det fornemmelig af den Grund, at saalænge Du blot lever æsthetisk, saalænge er Dit Liv totalt uvæsenligt. Den, der lever ethisk, hæver nu til en vis Grad Distinctionen mellem det Tilfældige og det Væsentlige; thi han overtager sig heel og holden som lige væsentlig; men den kommer igjen; thi efterat han har gjort det, adskiller han, dog saaledes, at det, han udelukker som det Tilfældige, overtager han et væsentligt Ansvar for med Hensyn til, at han udelukkede det.

Forsaavidt det æsthetiske Individ med »æsthetisk Alvor« sætter sig en Opgave for sit Liv, saa er denne egentlig den, at fordybe sig i sin egen Tilfældighed, at blive et Individ, hvis Paradoxi og Uregelmæssighed man aldrig har seet Mage til, en Grimace af et Menneske. Grunden, hvorfor man sjeldnere træffer saadanne Skikkelser i Livet, er, at man saa sjeldent træffer Folk, der have en Forestilling om det at leve. Da derimod Mange have en afgjort Forkjærlighed for at snakke, saa træffer man paa Gaden, i Selskaber og i Bøger megen Snak, der umiskjendeligt bærer Præg af den Originalitets-Wuth, der, overført paa Livet, vilde berige Verden med en Mængde Kunstproducter, hvoraf det ene var latterligere end det andet. Den Opgave, det ethiske Individ sætter sig, er at forvandle sig selv til det almene Individ. Kun det ethiske Individ gjør sig for Alvor Rede for sig selv, og har derfor Redelighed mod sig selv, kun han har den paradigmatiske Anstand og Sømmelighed, der er skjønnere end alt Andet. Men at forvandle sig selv til det almene Menneske er kun muligt, hvis jeg allerede ϰατα δυναμιν har det i mig. Det Almene kan nemlig godt bestaae med og i det Særegne, uden at fortære dette; det er som hiin Ild, der brændte, uden at fortære Tornebusken. Ligger det almene Menneske udenfor mig, saa er kun een Methode mulig, og den er at afføre mig selv hele min Concretion. Denne Stræben ud i Abstraktionens Tøilesløshed finder man ofte. Der var en Sect blandt Hussiterne, der meente, at det, det egentlig kom an paa, for at blive det normale Menneske, var at gaae nøgen ligesom Adam og Eva i Paradiis. Man træffer i vor Tid ikke sjelden Folk, der i aandelig Henseende lære det Samme, at man bliver det normale Menneske ved at blive aldeles splitternøgen, hvad man da kan opnaae, naar man affører sig hele sin Concretion. Men saaledes forholder det sig ikke. I Fortvivlelsens Act kom det almene Menneske frem, og er nu bag ved Concretionen og bryder frem derigjennem. I et Sprog er der mange flere paradigmatiske Verber end det ene, der i en Grammatik opstilles som Paradigma, det er tilfældigt, at det stilles frem, alle andre regelmæssige Verber kunne ligesaagodt; saaledes ogsaa med Menneskene. Ethvert Menneske kan, naar han vil, blive et paradigmatisk Menneske, ikke derved, at han afstryger sig sin Tilfældighed; men derved, at han bliver i den, og forædler den. Men han forædler den derved, at han vælger den.

Du vil nu let have indseet, at det ethiske Individ i sit Liv gjennemløber de Stadier, vi før have udviist som særskilte Stadier; han vil i sit Liv udvikle de personlige, de borgerlige, de religiøse Dyder, og hans Liv gaaer for sig derved, at han bestandig oversætter sig selv fra det ene Stadium til det andet. Saasnart man mener, at eet af disse Stadier er tilstrækkeligt, og at man eensidigen tør samle sig derpaa, saa har man ikke valgt sig selv ethisk, men enten overseet Isolationens eller Continuitetens Betydning, og fremfor Alt ikke fattet, at Sandheden ligger i disses Identitet.

Den, der ethisk har valgt og fundet sig selv, han har sig selv bestemmet i hele sin Concretion. Han har sig da som et Individ, der har disse Evner, disse Lidenskaber, disse Tilbøieligheder, disse Vaner, der staaer under disse ydre Indflydelser, der paavirkes i een Retning saaledes, i en anden saaledes. Her har han da sig selv som Opgave saaledes, at denne nærmest er den, at ordne, danne, temperere, opflamme, tilbagetrykke, kort at tilveiebringe en Ligelighed i Sjælen, en Harmoni, der er de personlige Dyders Frugt. Formaalet for hans Virksomhed er her ham selv, men dog ikke vilkaarligt bestemmet, thi han har sig selv som en Opgave, der er sat ham, om den end er bleven hans derved, at han har valgt den. Men skjøndt han selv er sit Formaal, saa er dog dette Formaal tillige et andet; thi det »Selv«, der er Formaalet, er ikke et abstrakt »Selv«, der passer allevegne, og derfor intetsteds, men et concret Selv, der staaer i levende Vexelvirkning med disse bestemte Omgivelser, disse Livsforhold, denne Tingenes Orden. Det Selv, der er Formaalet, er ikke blot et personligt Selv, men et socialt, et borgerligt Selv. Han har da sig selv som Opgave for en Virksomhed, hvorved han som denne bestemte Personlighed griber ind i Livets Forholde. Hans Opgave er ikke her at danne sig selv, men at virke, og dog danner han paa samme Tid sig selv; thi, som jeg ovenfor bemærkede, det ethiske Individ lever saaledes, at han bestandig sætter sig selv over fra et Stadium til et andet. Har Individet ikke oprindelig fattet sig selv som en concret Personlighed i Continuitet, saa vil han heller ei vinde denne senere Continuitet. Mener han, at Kunsten er, at begynde som en Robinson, saa bliver han en Eventyrer hele sit Liv. Indseer han derimod, at naar han ikke begynder concret, saa kommer han aldrig til at begynde, og naar han ikke kommer til at begynde, kommer han aldrig til at ende, saa vil han paa eengang være i Continuitet med det Forbigangne og det Tilkommende. Fra det personlige Liv sætter han sig over i det borgerlige, fra dette i det personlige. Det personlige Liv som saadant var en Isolation og derfor ufuldkomment, men idet han gjennem det borgerlige Liv vender tilbage i sin Personlighed, saa viser det personlige Liv sig i en høiere Skikkelse. Personligheden viser sig selv som det Absolute, der har sin Teleologi i sig selv. Man har, idet man gjorde det til Opgave for et Menneskes Liv, at leve for Pligtens Opfyldelse, ofte erindret om den Skepsis, at Pligten selv er vaklende, at Lovene kunne forandres. Du seer let, at hvad det sidste Udtryk angaaer, er der nærmest tænkt paa de Fluctuationer, som de borgerlige Dyder altid ere udsatte for. Dog denne Skepsis rammer ikke det Negativ-Moralske; thi det vedbliver uforandret. Derimod er der en anden Skepsis, der rammer enhver Pligt, det er den, at jeg overhovedet slet ikke kan gjøre Pligten. Pligten er det Almene; det, der fordres af mig, er det Almene; det, jeg kan gjøre, er det Enkelte. Denne Skepsis har imidlertid sin store Betydning, forsaavidt den viser, at Personligheden selv er det Absolute. Dette maa dog lidt nærmere bestemmes. Det er mærkeligt nok, at Sproget selv fremhæver denne Skepsis. Jeg siger aldrig om et Menneske: han gjør Pligten eller Pligterne, men jeg siger: han gjør sin Pligt; jeg siger: jeg gjør min Pligt, gjør Du Din. Dette viser, at Individet paa eengang er det Almene og det Enkelte. Pligten er det Almene, den fordres af mig; er jeg altsaa ikke det Almene, saa kan jeg heller ikke gjøre Pligten. Paa den anden Side er min Pligt det Enkelte, noget for mig alene, og dog er det Pligten og altsaa det Almene. Her viser Personligheden sig i sin høieste Gyldighed. Den er ikke lovløs, giver sig heller ikke selv sin Lov; thi Bestemmelsen af Pligt vedbliver, men Personligheden viser sig som Eenheden af det Almene og det Enkelte. At det forholder sig saaledes, det er klart, det kan man gjøre et Barn begribeligt; thi jeg kan gjøre Pligten, og dog ikke gjøre min Pligt, og jeg kan gjøre min Pligt og dog ikke gjøre Pligten. At Verden derfor skulde synke hen i Skepsis, indseer jeg slet ikke; thi Forskjellen mellem Godt og Ondt bliver altid, Ansvaret og Pligten ligesaa, om det end bliver umuligt for et andet Menneske at sige, hvad der er min Pligt, hvorimod det altid vil være muligt for ham at sige, hvad der er hans, hvilket ikke var Tilfældet, hvis Eenheden af det Almene og Enkelte ikke var sat. Man synes maaskee at have fjernet al Skepsis, naar man faaer gjort Pligten til noget Udvortes, Fast og Bestemt, om hvilket der kan siges: det er Pligten. Dette er imidlertid en Misforstaaelse; thi Tvivlen ligger ikke i det Udvortes, men i det Indvortes, i mit Forhold til det Almene. Som enkelt Individ er jeg ikke det Almene, og fordrer man det af mig, er det en Urimelighed; skal jeg altsaa kunne gjøre det Almene, maa jeg paa samme Tid som jeg er det Enkelte, være det Almene, men saa ligger Pligtens Dialektik i mig selv. Som sagt, denne Lære fører ingen Fare med sig for det Ethiske, tvertimod, den hævder det. Naar man ikke antager dette, saa bliver Personligheden abstrakt, dens Forhold til Pligten abstrakt, dens Udødelighed abstrakt. Forskjellen mellem Godt og Ondt hæves heller ikke; thi jeg tvivler paa, at der nogensinde har været et Menneske, der har paastaaet, at det er Pligt at gjøre Ondt. At han gjorde Ondt, er noget Andet, men han søgte tillige at indbilde sig selv og Andre, at det var godt. At han skulde kunne forblive i denne Indbildning, er utænkeligt, da han selv er det Almene; han har da ikke Fjenden udenfor sig, men i sig. Antager jeg derimod, at Pligten er noget Udvortes, saa er Forskjellen mellem Godt og Ondt hævet, thi naar jeg ikke selv er det Almene, kan jeg kun træde i et abstrakt Forhold til det, men Forskjellen mellem Godt og Ondt er incommensurabel for et abstrakt Forhold.

Netop naar man indseer, at Personligheden er det absolute, er sit eget Formaal, er Eenheden af det Almene og det Enkelte, netop da vil enhver Skepsis, der gjør det Historiske til sit Udgangspunkt, være overvunden. Fritænkerne have ofte nok søgt at forvirre Begreberne ved at gjøre opmærksom paa, hvorledes stundom Folk erklærede det for Helligt og Lovligt, hvad der i andre Folks Øine var en Afskyelighed og en Misgjerning. Man har her selv ladet sig blænde af det Udvortes; men ved det Ethiske er der aldrig Spørgsmaal om det Udvortes, men om det Indvortes. Men om det Udvortes end nok saa meget forandres, saa kan dog Handlingens sædelige Gehalt blive den samme. Der har saaledes vist aldrig været noget Folk til, der har meent, at Børn skulle hade deres Forældre. Imidlertid har man, for at nære Tvivlen, gjort opmærksom paa, at medens alle dannede Nationer gjorde det til Børnenes Pligt, at sørge for deres Forældre, saa havde de Vilde den Skik at slaae deres gamle Forældre ihjel. Det er vel muligt, at det forholder sig saaledes; men dermed er man dog ikke kommen videre, thi Spørgsmaalet bliver, om de Vilde mene at gjøre noget Ondt dermed. Det Ethiske ligger altid i denne Bevidsthed, hvorimod det er et andet Spørgsmaal, om ikke en mangelfuld Erkjendelse er tilregnelig. Fritænkeren indseer meget godt, at det er den Maade, paa hvilken det Ethiske allerlettest forflygtiges, naar man aabner Døren for den historiske Uendelighed. Og dog ligger der noget Sandt i hans Adfærd, thi naar ikke Individet til syvende og sidst selv er det Absolute, saa er Empirien den eneste Vei, der er ham anviist, og den Vei har med Hensyn til Udløb den samme Egenskab som den Flod Niger med Hensyn til Udspring, at Ingen veed, hvor det er. Er jeg anviist Endeligheden, saa er det en Vilkaarlighed at blive staaende paa noget enkelt Punkt. Ad den Vei kommer man derfor aldrig til at begynde, thi for at skulle begynde maatte man være kommen til Enden, men dette er en Umulighed. Naar Personligheden er det Absolute, da er den selv det archimediske Punkt, fra hvilket man kan løfte Verden. At denne Bevidsthed ikke kan forlede Individet til at ville kaste Virkeligheden af sig, er let at indsee, thi vil han saaledes være det Absolute, saa er han slet Intet, en Abstraktion. Kun som den Enkelte er han den Absolute, og denne Bevidsthed vil frelse ham fra al revolutionair Radikalisme.

[…]

Den, der derimod med Smerte forvissede sig om, at han var et ualmindeligt Menneske, ved sin Sorg derover atter forsonede sig med det Almene, han vil maaskee engang opleve den Glæde, at hvad der forvoldte ham Smerte og gjorde ham ringe i hans egne Øine, viser sig som Anledning til, at han atter opløftes og i ædlere Forstand bliver et ualmindeligt Menneske. Hvad han nemlig tabte i Omfang, det vandt han maaskee i intensiv Inderlighed. Ikke ethvert Menneske nemlig, hvis Liv middelmaadigt udtrykker det Almene, er derfor et ualmindeligt Menneske, thi det vilde jo være en Forgudelse af Trivialiteten; for at han i Sandhed skal kaldes saa, maa der ogsaa spørges om den intensive Kraft, med hvilken han gjør det. Denne Kraft vil nu hiin Anden være i Besiddelse af paa de Punkter, hvor han kan realisere det Almene. Hans Sorg vil da atter forsvinde, den vil opløse sig i Harmoni; thi han vil indsee, at han var kommen til Grændsen for sin Individualitet. Han veed vel, at ethvert Menneske udvikler sig med Frihed, men han veed ogsaa, at et Menneske ikke skaber sig selv af Intet, at han har sig selv i sin Concretion som sin Opgave; han vil atter forsones med Tilværelsen, idet han vil indsee, at i en vis Forstand ethvert Menneske er en Undtagelse, og at det er lige sandt, at ethvert Menneske er det Almeen-Menneskelige og tillige en Undtagelse.

Her har Du min Mening om, hvad det er at være et ualmindeligt Menneske. Jeg elsker Tilværelsen og det, at være et Menneske, altfor meget til at troe, at Veien til at blive et ualmindeligt Menneske er let eller uden Anfægtelser. Men selv om et Menneske saaledes i ædlere Forstand er et ualmindeligt Menneske, saa vil han dog bestandigt tilstaae, at det var endnu Fuldkomnere at optage hele det Almene i sig.

Saa modtag min Hilsen, tag mit Venskab; thi skjøndt jeg ikke i strengeste Forstand tør betegne vort Forhold saaledes, saa haaber jeg dog, at min unge Ven engang vil blive saa meget ældre, at jeg i Sandhed tør bruge dette Ord; vær forsikkret om min Deeltagelse. Modtag en Hilsen fra hende, hvem jeg elsker, hvis Tanker ere skjulte i mine Tanker, modtag en Hilsen, der er uadskillelig fra min, men modtag ogsaa en særlig Hilsen fra hende, venlig og oprigtig som altid.

Da Du for nogle Dage siden var hos os, tænkte Du maaskee ikke paa, at jeg atter var færdig med en saa stor Skrivelse. Jeg veed, Du er ingen Ven af, at man taler til Dig om Din indre Historie, derfor har jeg valgt at skrive, og skal aldrig tale til Dig derom. At Du modtager et saadant Brev, bliver en Hemmelighed, og jeg ønskede ikke gjerne, at den skulde have nogen Indflydelse til at forandre Dit Forhold til mig og min Familie. At Du har Virtuositet nok til at gjøre det, naar Du vil, det veed jeg, og derfor beder jeg Dig derom for Din og for min Skyld. Jeg har aldrig villet trænge mig ind i Dig, og kan godt elske Dig paa Afstand, uagtet vi ofte sees. Dit Væsen er for indesluttet til at jeg troer, det vilde gavne at tale til Dig, derimod haaber jeg, at mine Breve ikke skulle være uden Betydning. Naar Du da forarbeider Dig selv i Din Personligheds lukkede Maskineri, da stikker jeg mine Indlæg ind og er forvisset om, at de komme med i Bevægelsen.

Da saaledes vort skriftlige Forhold bliver en Hemmelighed, saa iagttager jeg alle Formaliteter, ønsker Dig et Levvel, som om vi vare langt fjernede fra hinanden, uagtet jeg haaber at see Dig ligesaa ofte hos mig som før.

\workheading{Fire opbyggelige Taler (1844)}{1844}

\worknote{Antologiens ''Four Upbuilding Discourses''. Hong-udvalget bringer bogens forord samt begyndelsen af den første tale, ''At trænge til Gud er Menneskets høieste Fuldkommenhed'' (til og med ''… da har man erkjendt sig selv''); talens fortsættelse og de øvrige tre taler er ikke medtaget.}

\sectionheading{Forord  ·  SV¹ V, 79}

Skjøndt denne lille Bog (som derfor blev kalden »Taler« ikke Prædikener, fordi dens Forfatter ikke har Myndighed til at prædike, »opbyggelige« Taler ikke Taler til Opbyggelse, fordi den Talende ingenlunde fordrer at være Lærer) nu atter gaaer ud i Verden, frygter den dog end mindre, end da den første Gang tiltraadte Vandringen, at skulle drage nogen forsinkende Opmærksomhed paa sig; den haaber tvertimod, at de Forbigaaende formedelst Gjentagelsen neppe bemærke den eller dog kun for at lade den skjøtte sig selv. Saaledes gaaer stundom et Bud til bestemte Tider sin vante Vei; snart er han kjendt, kjendt, saa den Forbigaaende neppe seer ham, end sige seer efter ham – og saaledes gaaer denne lille Bog ud som et Bud, men ikke saaledes som et Bud, at den atter vender tilbage. Den søger hiin Enkelte, hvem jeg med Glæde og Taknemlighed kalder min Læser, for at besøge ham, ja for at blive hos ham, thi Den, hvem man elsker, til ham kommer man og gjør Bolig hos ham og bliver hos ham, hvis det forundes. Saasnart nemlig han har taget den imod, da haver den ophørt at være: den er Intet for sig selv og ved sig selv, men Alt, hvad den er, er den kun for ham og ved ham. Og uagtet saaledes Sporet bestandig fører hen til min Læser, ikke tilbage, og uagtet det foregaaende Bud aldrig vender hjem, og uagtet Den, som udsender det, aldrig erfarer Noget om dets Skjebne; saa gaaer dog det næste Bud freidig gjennem Døden til Livet, trøstig sin Gang for at forsvinde, glad over aldrig at skulle vende hjem igjen – og dette er netop Dens Glæde, som udsender det, som bestandigt kun kom til sin Læser for at tage Afskeed, og nu tager den for sidste Gang.

Kjøbenhavn den 9de August 1844.

S. K.

\sectionheading{At trænge til Gud er Menneskets høieste Fuldkommenhed (uddrag)  ·  SV¹ V, 81–100}

»Et Menneske behøver kun Lidet for at leve, og dette Lidet behøver han kun en liden Tid« – dette er et høimodigt Ord, der er værd at annammes og at forstaaes som det vil forstaaes; thi det er for alvorligt til at ville beundres som et skjønt Udtryk eller en siirlig Vending. Saaledes henkastes det vel stundom; man tilraaber den Nødlidende det, maaskee for i Forbigaaende at trøste ham, maaskee ogsaa kun for at sige Noget; man siger det til sig selv endog paa Lykkens Dag, thi det menneskelige Hjerte er meget svigefuldt og tager kun altfor gjerne Høimodet forfængeligt, og er stolt af kun at behøve Lidet – medens man bruger meget; man siger det til sig selv paa Nødens Dag, og iler forud for beundrende at modtage sig selv ved Maalet – naar man har fuldkommet det Herlige; men hermed er man ligesaa lidt selv tjent, som Ordet er det. »Man behøver det kun en liden Tid«; men som det stundom hænder, at det er, naar Dagen længes, at ogsaa Vinteren strenges, saa er det altid saaledes, at Savnets og Gjenvordighedens Vinter gjør Dagene lange, om end Tiden og Livet er kort. Hvor Meget er da det Lidet, som et Menneske behøver? I Almindelighed kan man egentligen slet ikke besvare dette Spørgsmaal; og selv Den, der har oplevet, om ikke at behøve Lidet, saa dog at maatte hjælpe sig med Lidet, selv han vil ikke i Almindelighed kunne bestemme, hvad det Lidet er. Som nemlig Tiden ofte bringer den Sørgende ny Trøst, den Nedbøiede ny Opreisning, Den, der tabte meget, nyt Vederlag, saa er den, selv naar den vedbliver at tage, dog omhyggelig mod den Lidende; tager sjeldent Alt paa eengang, men lidt efter lidt, og vænner ham saaledes lidt efter lidt til at undvære, indtil han selv med Forundring seer, at han endog behøver mindre, end hvad han engang ansaae for det Mindste, ja for saa Lidet, at han med Forfærdelse tænkte paa at maatte behøve saa Lidet (om han end ikke udtrykte sig ganske klart, thi Forfærdelsen ligger sandeligen ikke i kun at behøve Lidet) og næsten ophidsedes ved den modsigende Forestilling, at skulle behøve det, for at kunne vedblive at behøve det; om han end ikke ganske forstod sig selv, da Fuldkommenheden just ikke bestod i at komme til at behøve Mere. Men hvor meget er da det Lidet, som et Menneske behøver? Lad Livet svare, og Talen gjøre hvad Livets Nød og Gjenvordighed stundom gjør, afklæde et Menneske for at see, hvor Lidet det er, han behøver. Og Du, min Tilhører, deeltage deri, som Du nu i Forhold til Din særlige Tilstand maa eller vil deeltage; thi Forfærdelse skal Talen ikke vække, dersom den ellers er betænkt paa at finde en Trøst, og ikke har til Hensigt at ville bedrage Dig, som naar Fortvivlelsens Iis og Snee danner den skuffende Bjergstrøm, efter hvilken Karavanen bøier Veien, begiver sig ind i Ørkenen og omkommer (Job, VI, 15-18). Din Forfærdelse kan den ikke vække, dersom Du selv haver oplevet det, og fundet Trøsten; og hvis Du ikke haver oplevet det, da kan dog Talen derom kun forfærde Dig, hvis Din Styrke er i den Betragtning, at Sligt kun hænder sjeldent. Men hvo er da den Elendigste: Den, der haver oplevet det, eller den feige og blødagtige Daare, der ikke mærker, at hans Trøst er Svig, at det, naar Ulykken rammer, kun lidet baader, at det er den sjeldne Gang. Saa tag det da fra ham: Rigdom og Magt og Vælde, og falske Venners forræderske Tjenstagtighed, og Lysternes Underdanighed under Ønskets Lune, og Forfængelighedens Triumpher over den tilbedende Beundring, og Menneskevrimlens smigrende Opmærksomhed og hele hans Fremtrædens misundte Herlighed; han har tabt det og nøies med mindre. Som Verden ikke kan kjende ham igjen paa Grund af den uhyre Forandring, saa kan han neppe kjende sig selv – saa forandret er han: at han, der behøvede saa meget, nu behøver saa langt mindre. Sandeligen, det er langt lettere og langt glædeligere at forstaae, hvorledes denne Forandring kan gjøre ham ukjendelig for sig selv, end hvorledes hiin anden kan gjøre ham ukjendelig for Menneskene; thi er det ikke Daarlighed, at det er Klæderne, der gjør en Mand ukjendelig, saa man ikke kan kjende ham, naar han er afklædt; og er det ikke sørgeligt, at det er Klæderne man beundrer, ikke Manden! Men en gudeligere Betragtning seer let, at han er ifærd med at klædes om og at iføres Høitidsdragten, thi den jordiske Bryllupsdragt og den himmelske ere yderst forskjellige. Dog lidet Gods med Nøisomhed er jo allerede en stor Vinding; saa tag det da fra ham – ikke Nøisomheden – men det Sidste, han eiede. Han lider ikke Nød, han gaaer ikke hungrig til Sengs; men hvor han skal faae det Fornødne, det veed han ikke, ikke om Aftenen naar han sover ind fra Bekymringen, ikke om Morgenen naar han vaagner til den; men dog faaer han det – det Lidet han behøver for at leve. Saa er han da fattig, og det Ord, der er saa tungt at høre, det maa han høre sagt om sig selv. Dobbelt tungt føler han det; thi han valgte ikke selv denne Tilstand som Den, der henkastede sit Gods for at forsøge sig, og som maaskee lettere finder sig i selvvalgt Armod, men derfor ikke altid bedre, dersom han kun forsagede Forfængeligheden ved Hjælp af en endnu større Forfængelighed. »Kun at behøve Lidet«, saa talte Ordet; men at vide, at man kun behøver Lidet, uden dog i noget Øieblik med Vished at vide, at man kan faae det Lidet, man behøver – Den, der bærer dette, han behøver kun Lidet, han behøver end ikke, hvad der dog er Noget, at vide sig dette Lidet sikkret. Mere behøver da et Menneske ikke, dersom det er sandt, at han kun behøver Lidet – for at leve; thi en Grav finder han vel, og i Graven behøver ethvert Menneske lige lidet. Hvad enten den døde Mand eier den Grav, i hvilken han ligger, maaskee paa hundrede Aar (ak, hvilken besynderlig Modsigelse!), eller han har maattet trænge sig ind mellem Andre, selv i Døden har maattet kæmpe sig frem for at faae en lille Plads: de eie lige Meget og behøve lige Lidet, og behøve det kun for en liden Tid. Men den liden Tid forinden, om hvilken Ordet taler, den kan maaskee blive lang, thi om Veien til Graven end ikke var lang, om Du maaskee ikke sjeldent saae ham gaae mødig sin Gang derud, for med Øiet at erobre det lidet Land, han som Død agtede at indtage; kunde Veien derfor ikke i en anden Forstand blive meget lang? Hvis han da stundom blev mistrøstig, hvis han ikke altid forstod, at et Menneske kun behøver Lidet, havde Du da slet ikke Andet at sige ham end at gjentage hiint Ord? Eller Du sagde vel til ham, hvad der tilbød sig ganske naturligt, saa naturligt, at Du maaskee end ikke selv i Hjertet ret havde Tillid til den Trøst, du bød en Anden: lad Dig saa nøies med Guds Naade.

Stands nu et Øieblik, at ikke Alt forvirres: Tanken og Talen og Sproget; at ikke Alt forvirres, idet Forholdet vel vedbliver, men kun derved, at det er ombyttet, saa hiin Mand er Den, der eier Trøsten, Du Den, der behøver den; hiin Mand den Rige, Du den Nødlidende, uagtet det var aldeles omvendt, lige indtil Du hørte det lille Trylle-Ord, der forvandlede Alt. Maaskee mærkede Du det end ikke selv; thi saaledes omgaaes man mangt et Ord, som Barnet det Betydningsfulde, uden at opdage den Tankens Braad deri, som saarer indtil Døden for at frelse Livet. At nøies med Guds Naade! Guds Naade er jo det Herligste af Alt, derom skulle vi dog vel ikke stride; thi dette er i Grunden ethvert Menneskes inderste og saligste Forvisning. Men Tanken om den kalder han saa sjeldent frem; og tilsidst, naar han ret vil være ærlig, da anvender han vel i al Stilhed uden dog at gjøre sig det tydeligt, hvad han gjør, paa denne Forestilling hiint gamle Ord: at for Lidet og for Meget fordærver Alt. Skal han tænke Tanken i sin evige Gyldighed, da sigter den øieblikkeligt dødbringende paa al hans verdslige Tænken, Digten og Tragten, vender ham Alt om, og dette kan han ikke udholde længe. Da vender han tilbage til det Verdsliges lave Egne, til sin sædvanlige Tale og Tankegang. Og jo ældre et Menneske bliver, desto vanskeligere vorder det ham at lære et nyt Sprog, især et saa høist forskjelligt. Dette mærker han nu nok af og til, at det ikke hænger rigtigt sammen med den Maade, paa hvilken han bruger dette Ord, blandet ind mellem al verdslig Tale; han faaer en ond Samvittighed mod Ordet, og har ingen Velsignelse af det. Dog er jo Guds Naade det Herligste af Alt. Men dersom nu et Menneske eiede, hvad der vel ikke er saa herligt, eiede alle Jordens Skatte, og du da vilde sige til ham: lad Dig nu nøies dermed; da vilde han vel smile af Dig. Dersom han selv sagde til Dig: nu vil jeg lade mig nøies dermed; da vilde det vel oprøre Dig; thi hvad mere kunde han forlange, og hvilken Frækhed at ville nøies med det Meste. Hvad man nøies med, det maa jo være det Lidet; men at nøies med det Herligste af Alt, det synes en besynderlig Tale; og at dette Trøstens Forslag gjøres af et Menneske, der ikke selv fatter den, er igjen besynderligt, som hvis Nogen ikke uden deeltagende Bekymring gav den Nødlidende en liden Skilling, formanende ham, at nøies med den, medens dog denne Skilling gjorde Modtageren til Eier af hele Verden. Thi var det ikke besynderligt, at Giveren selv kunde tænke saa ringe om Gaven, han bød, at han lod Formaningen om Nøisomhed følge med! Eller var det ikke som, hvis en Mand, der, selv indbuden af den Mægtige, gik til Gjestebudshuus, mødte en Ringe, til hvem han sagde, for dog at byde ham en Slags Trøst: lad dig nøies med at sidde tilbords i Himmeriges Rige! Eller dersom den Ringe selv var talende og sagde: ak! jeg bliver ikke indbuden af den Mægtige og kan da heller ei modtage hans Indbydelse, da jeg er andetsteds indbuden og maa nøies med at sidde tilbords i Himmerigets Rige – var dette ikke en underlig Tale? Jo mere man tænker derover, desto besynderligere bliver Jordlivet og det menneskelige Sprog; thi midt i al den jordiske og verdslige Forskjellighed, der er nidkjær nok paa sig selv, blander Guds-Forskjelligheden sig tankeløst ind, ja vel endog saaledes, at den i Grunden er udelukket. At komme i Kongens Navn, det aabner alle Døre for En, men at komme i Guds Navn, det er det Sidste, hvori et Menneske skal forsøge sig; og Den, der maa nøies dermed, han maa nøies med Lidet. Hvis han kom til den Mægtiges Dør, hvis Tjeneren end ikke forstod, fra hvem det var han hilsede, hvis den Mægtige utaalmodig selv gik ud og saae den ringe Mand, der skulde hilse fra Gud i Himlene – maaskee blev Døren lukket for ham.

Dog vil Talen ikke overrumple Dig, min Tilhører, eller frembringe en pludselig Virkning. Naar Ordet siger: at nøies med Guds Naade, da er vel Grunden, at Guds Naade ikke udtrykker sig saaledes som et Menneske saa gjerne vil forstaae den, men taler vanskeligere. Saasnart nemlig Guds Naade giver Mennesket hvad han ønsker og begjerer, da nøies han ikke blot med Naaden, men er glad ved hvad han modtager, og forstaaer, efter sin Formening, let, at Gud er ham naadig. At dette nu er en Misforstaaelse, som Ingen skal forhaste sig med at ivre mod, det er vist nok, men derfor bør man dog ikke glemme at øve sig i den beleilige Tid til at forstaae det Vanskeligere og det Sande. Dersom nemlig et Menneske kan være forvisset om Guds Naade, uden at behøve det timelige Vidnesbyrd som Mellemmand, eller den efter hans Begreb ham gavnlige Tilskikkelse som Tolk, da er det jo ham vist, at Guds Naade er det Herligste af Alt, og da vil han bestræbe sig for, saaledes at glædes ved den, at han ikke blot nøies med den, saaledes at takke for den, at han ikke nøies med Naaden: ikke sørger over Det, som blev nægtet, ikke over den Sprogforskjellighed der var mellem Guds evige Tilforladelighed og hans barnlige Lidettroenhed, men som nu ikke mere er, nu da »hans Hjerte styrkes ved Naaden og ikke ved Mad« (Hebr. 13, 9). Dersom en Nødlidende turde glæde sig ved en Mægtigs Venskab, men denne mægtige Mand kunde Intet gjøre for ham (dette er jo det Tilsvarende til at Guds Naade lader det jordiske Vidnesbyrd udeblive): det, at han havde et saadant Venskab, var dog allerede saare meget. Men her ligger maaskee Vanskeligheden; thi den Nødlidende kunde jo overbevises om, at hiin mægtige Mand virkeligen ikke formaaede at gjøre Noget for ham; men hvorledes skulde han paa endelig Maade overbevises om, at Gud ikke kan, han er jo almægtig! Deraf kommer det vel, at Utaalmodighedens Tanke bestandigen ligesom pukker paa, at Gud kan nok; og derfor er det, fordi Mennesket er saa utaalmodigt, derfor er det, at Sproget siger: at nøies med Guds Naade. I Begyndelsen, naar Utaalmodigheden er meest høirøstet, da kan den neppe forstaae, at det er en priselig Nøisomhed; alt eftersom den sagtnes og dysses i det indvortes Menneskes stille Uforkrænkelighed, da fatter den det bedre og bedre, indtil Hjertet bevæges og stundom idetmindste seer den guddommelige Herlighed, der havde paataget sig en ringe Skikkelse. Og forsvinder denne Herlighed atter for Mennesket, saa han igjen er, hvad han dog ogsaa var medens han saae Herligheden, nødlidende, synes det ham atter, at der dog hører Nøisomhed til for at nøies med Guds Naade, da tilstaaer han dog vel stundom sig selv i Beskæmmelse, at Guds Naade nok er værd at nøies med, ja alene værd at eftertragte, ja ene saligt at besidde.

Saa vil da lidt efter lidt, thi Guds Naade tages aldrig med Vold, det menneskelige Hjerte i skjøn Forstand blive mere og mere unøisomt, det vil sige mere og mere higende, mere og mere længselsfuldt efter at vide sig Naaden sikkret. Og see, nu er Alt blevet nyt, Alt blevet forandret! I Forhold til det Jordiske gjelder det, at behøve Lidet, og i samme Grad man behøver mindre, i samme Grad er man fuldkomnere, som en Hedning, der jo kun vidste at tale om det Jordiske, har sagt det, at Guddommen var salig, fordi han Intet behøvede, næst ham den Vise, fordi han behøvede Lidet. I Menneskets Forhold til Gud er det omvendt: jo mere han behøver Gud, jo dybere han fatter at han trænger til Gud, og nu i sin Trang trænger sig frem til Gud, desto fuldkomnere er han. Ordet at nøies med Guds Naade vil derfor ikke blot trøste et Menneske, og da atter trøste, hvergang jordisk Savn og Nød gjør ham, jordisk talt, Trøsten fornøden; men naar han ret er bleven opmærksom paa Ordet, da kalder det ham afsides, hvor han ikke mere hører det verdslige Sinds jordiske Modersmaal, ikke Menneskenes Taler, ikke Larmen af de Handlende, men hvor Ordet forklarer sig for ham, betroer ham Fuldkommenhedens Hemmelighed: at det at behøve Gud, ikke er til at skamme sig ved, men netop er Fuldkommenheden, og at dette er det Sørgeligste, om et Menneske gik gjennem Livet uden at opdage, at han behøvede Gud.

Saa ville vi da tydeliggjøre os denne opbyggelige Tanke:

At trænge til Gud er Menneskets høieste Fuldkommenhed.

At det er saa, at det at trænge til Gud er en Fuldkommenhed, derom synes en Omstændighed, der vel er Alle bekjendt, at minde Enhver idetmindste med en forbigaaende Paamindelse. I de forskjellige Landes Kirker bedes jo efter Prædikenen for Kongen og det kongelige Huus. At der bedes for Dem, som syge og sorrigfulde ere, kan ikke bevise, at det at trænge til Gud er en Fuldkommenhed; thi disse ere jo de Lidende. Men Kongen, han er den mægtige Mand, ja den mægtigste; og dog bedes der ganske særligen for ham, kun i Almindelighed for de Syge og Sorrigfulde, om Kirken end haaber og fortrøster sig til, at Gud i Himlene vil forstaae det ganske særligen, at han, medens Kirken end ikke tænker paa nogen Bestemt, tænker paa enhver Bestemt i Besynderlighed. Og dersom det er anderledes med Guds Forstaaen, dersom hans Regjeringssorger ogsaa kun tillod ham i Almindelighed at bekymre sig om den Enkelte – ja, da Gud hjælpe os! Ak, dette er dog det Sidste et Menneske siger i sin Elendighed; selv naar han ikke kan udholde denne sidste Tanke, at Gud kun i Almindelighed skulde bære Omsorg for den Enkelte, selv da siger han: Gud hjælpe mig til at udholde denne Tanke, og faaer dog saaledes alligevel Gud til at bekymre sig særligen om ham. Men hvorfor bedes der nu for Kongen i Særdeleshed? Er det fordi han har den jordiske Magt og holder Manges Skjebne i sin Haand; er det fordi hans Velvære betinger Utalliges; er det fordi enhver »Modgangsskygge,« der gaaer over Kongens Huus, ogsaa gaaer over hele Folket; er det fordi hans Sygdom standser Statens Virksomhed, hans Død forstyrrer dens Liv? En saadan reen verdslig Bekymring kan vistnok, og ei heller uskjønt, beskæftige Mange; men dog vel neppe bringe Nogen til at bede paa nogen anden Maade end med den Tilbageholdenhed, der er fornøden, naar vi bede om jordiske Goder; thi et saadant jordisk Gode bliver jo i saa Fald en Konge. I denne Forstand vilde da og Forbønnen blive inderligere og inderligere, alt eftersom den Bedendes eget Liv var knyttet nærmere og nærmere til ham, indtil Bønnen tilsidst ikke mere var en Forbøn, saa lidet som Hustruens Bøn for Manden er en Forbøn. Men saaledes kan jo Kirken ikke være bedende, netop fordi den gjør Forbøn; men dette gjør den vel, fordi den er overbeviist om, at jo høiere et Menneske staaer, desto mere behøver han Gud.

Dog om der end bedes i alle Kirker for en Konge, deraf følger jo ikke, at den Konge, for hvem der bedes, selv fatter, at det at trænge til Gud er Menneskets høieste Fuldkommenhed. Og om den Enkelte i Kirken stiltiende giver sit Minde til Forbønnen, ak, og om de Mange, der ikke besøge Kirken, dog Intet have at indvende mod Forbønnen, deraf følger jo ikke, at disse forstaae det gudeligt, at jo høiere man stiger i jordisk Magt og Vælde, desto nærmere kommer man kun til Forbønnen. At tage den forfængeligt, er for den Mægtige kun altfor let; at sige den forfængeligt, er for den Bedende kun altfor let; hvorimod det, ret for Alvor at ville forstaae, hvad man allerede ikke forstaaer alvorligt nok, naar man selv vil forsøge sig til denne Bevidsthed, og ikke lade Gud om det, der bedst veed at ængste al Selvtillid ud af et Menneske og forhindre ham i, naar han skal til at synke i sin egen Intethed, at vedligeholde Dykkerens Forbindelse med det Jordiske selv; det, ret for Alvor at ville forstaae det, gjør Livet vanskeligt, vi nægte det ikke. Lader os kun tilstaae det, uden derfor at blive mismodige eller feige, saa vi ville sove os til, hvad Andre have maattet arbeide for; lader os ikke tage det forfængeligt, naar den Troende begeistret siger, at al hans Lidelse kun er stakket og kort, at Selvfornægtelsens Aag er saa let at bære. Men lader os heller ei tvivle om, at Selvfornægtelsens Aag er gavnligt, at Lidelsernes Kors adler et Menneske trods Noget, og haabe til Gud, engang at skulle komme saavidt, at vi ogsaa kunne tale begeistrede. Men lader os ikke forlange det for tidligen, at ikke den Troendes begeistrede Tale skal gjøre os mismodige, fordi det ikke strax vil lykkes. Ofte gaaer det et Menneske saaledes, at han indprenter i sin Hukommelse et enkelt stærkt Ord. Naar nu Lidelsen kommer til hans Huus, da mindes han Ordet, og mener strax at skulle seire i Ordets Glæde. Men selv en Apostel taler ikke altid stærke Ord, ogsaa han er stundom svag, ogsaa han ængstes, og giver derved at forstaae, at det stærke Ord er dyrekjøbt, og aldrig besiddes saaledes, at man jo igjen kan faae Leilighed til at forvisse sig om, hvor dyrt det er kjøbt.

Men om nu end denne Forstaaelse gjør Livet vanskeligere, ikke blot for den Lykkeliges Letsindighed og for de Mange, hvis Higen er at blive som han er det; men ogsaa for de Ulykkelige, da dog Forstaaelsen ikke virker som en Trolddom, ikke virker paa nogen udvortes afgjørende Maade; skulle vi derfor tvivlsomt anprise den, eller tvesindede attraae den? Og dog er det jo en betænkelig Sag, at det, der skal bydes som Trøsten i Livet, begynder med at gjøre det vanskeligere, for da – ja for da i Sandhed at gjøre det lettere; thi saaledes er det med enhver Sandhedens Undergjerning, som det var med hiint Under ved Brylluppet i Cana: først iskjænker Sandheden den slette Viin og gjemmer den bedste til Sidst; hvorimod den bedragerske Verden skjænker den bedste Viin først. Fordi nemlig et Menneske blev ulykkelig, »grændseløs ulykkelig«, som han jo selv siger, deraf følger slet ikke, at den Forstaaelse, der betinger Trøsten, den Forstaaelse, at han selv slet Intet formaaer, endnu er modnet i ham. Dersom han saaledes troer, at han blot mangler Midlet, da troer han endnu paa sig selv, dersom han troer, at hvis Magten gaves ham, eller Menneskenes Beundring, eller Besiddelsen af det Ønskede, dersom han mener, at Klagen indeholder en berettiget Fordring paa noget Timeligt, som bliver mere berettiget, jo heftigere Klagen bliver – da har han endnu en, menneskelig talt, bitter Kalk at tømme, før Trøsten kommer. Derfor er det altid en vanskelig Sag for det ene Menneske at tilbyde det andet en saadan Trøst; thi naar den Bekymrede tyer til ham, og han da vil sige: »jeg veed vel, hvor der er en Trøst at finde, en ubeskrivelig Trøst, ja hvad mere er, den forvandler sig lidt efter lidt i Din Sjel til den høieste Glæde,« da vil vel den Bekymrede lytte opmærksomt; men naar der da tilføies: »førend dog denne Trøst kan komme, da maa man forstaae, at man selv er slet Intet; da maa man afhugge den Sandsynlighedens Bro, der vil sætte Ønsket og Utaalmodigheden og Begjeringen og Forventningen i Forbindelse med det Ønskede, det Begjerede, det Forventede; da maa man forsage det verdslige Sinds Omgængelse med det Tilkommende; da maa man trække sig tilbage i sig selv, ikke som i en Fæstning, der endnu byder Verden Trods, medens dog den Indesluttede har hos sig i Fæstningen sin farligste Fjende, ja medens det maaskee var Fjendens Raad, han fulgte, da han saaledes sluttede sig inde, men trække sig tilbage i sig selv, synkende i sin egen Intethed, overgivende sig selv paa Naade og Unaade«, da var der vel Den, som gik bedrøvet bort liig hiin rige Yngling, der havde meget Gods, og det uagtet han ikke havde meget Gods, men dog lignede ham, saa man ikke kunde skjelne dem. Eller dersom den Bekymrede var løben vild og løben fast i Overveielse, saa han ikke mægtede at handle, fordi der fra begge Sider lod sig sige ligemeget, og et andet Menneske da vilde sige til ham: »jeg veed en Udvei, og Du skal være Seiren vis; opgiv Ønsket, handel, handel i den Overbeviisning, at selv naar netop det Modsatte af Ønsket skeer, har Du dog seiret«, da var der vel Den, der vendte sig utaalmodigen bort, fordi en saadan Seier syntes ham et Nederlag, og fordi en saadan Udvei tyktes ham endnu tungere end den tvivlende Sjels mangfoldige Uro.

Hvad er dog et Menneske? Er han kun en Prydelse mere i Skabningens Række; eller har han ingen Magt, formaaer han selv Intet? Og hvilken er da denne Magt; hvad er det Høieste han formaaer at ville? Hvorledes lyder Svaret paa dette Spørgsmaal, naar Ungdommens Dristighed forener sig med Manddommens Kraft for at spørge, naar denne herlige Forening er villig til at offre Alt for at fuldkomme det Store, naar den brændende i Iver siger: »om end ingen Anden før naaede det i Verden, saa vil dog jeg naae det; om Millioner vanartede og glemte Opgaven, saa vil dog jeg stride – men hvad er det Høieste«? Nu vel, vi ville ikke besvige det Høieste for sin Priis, vi dølge ikke, at det sjeldent blev naaet i Verden; thi det Høieste er: at et Menneske fuldelig overbevises om, at han selv slet Intet formaaer, slet Intet. O, sjeldne Vælde, ikke saaledes sjelden, at kun en Enkelt fødtes kongebaaren, thi Enhver fødtes til den! O, sjeldne Viisdom, ikke saaledes sjelden, at den kun tilbødes Enkelte, da den dog tilbydes Alle! O, vidunderlige Sjeldenhed, der ikke nedværdiges ved at tilbydes Alle, og ved at kunne besiddes af Alle! Ja naar Mennesket vil vende sig ud efter, da synes det vel som formaaede han hvad der var mere forbausende, hvad der vilde tilfredsstille ham selv ganske anderledes, hvad Beundringen jublende vilde flokkes om; thi hiin sjeldne Ophøiethed er ikke god at beundre, frister ikke det sandselige Menneske, da den tvertimod dømmer den Beundrende, at han er en Daare, der ikke selv veed hvad han beundrer, og byder ham at gaae hjem; eller at han er en svigagtig Sjæl, og byder ham at gaae i sig selv. For den udvortes Betragtning er vel Mennesket den herligste Skabning, men al hans Herlighed er dog kun i det Udvortes og til det Udvortes: thi sigter ikke Øiet med sin Piil ud efter, hver Gang Lidenskaben og Lysten spænder Strengen, griber ikke Haanden ud efter, er ikke hans Arm udrakt, er ikke hans Snilde erobrende! Men dersom han dog ikke vil være som et Krigsredskab i uforklarede Drifters Tjeneste, ja i Verdens Tjeneste, fordi Verden, til hvilken hans Attraa er, selv vækker Driften; dersom han dog ikke vil være som en Strengeleg i uforklarede Stemningers Haand, eller rettere i Verdens Haand, fordi hans Sjæls Bevægelse er i Forhold til som Verden griber i Strengene; dersom han ikke vil være som et Speil, hvori han opfanger Verden, eller rettere, hvori Verden speiler sig selv; dersom han ikke vil dette, dersom han, førend Øiet sigter efter Noget for at erobre det, først selv vil fange Øiet, at det maa tilhøre ham, ikke han Øiet, dersom han griber Haanden, førend den griber til efter det Udvortes, at den maa tilhøre ham, ikke han Haanden, dersom han vil det saa alvorligt, at han ikke frygter for at udrive Øiet, at afhugge Haanden, at lukke Sandsens Vindve, hvis det skulde gjøres fornødent – ja, da er Alt forandret, Magten taget fra ham og Herligheden. Han strider ikke med Verden, men med sig selv. Betragt ham nu, hans kraftige Skikkelse holdes omsluttet af en anden Skikkelse, og saa fast holde de hinanden omslyngede, lige smidigt og lige stærkt sammenknuge de hinanden, saa Brydningen end ikke kan begynde, fordi hiin anden Skikkelse i samme Øieblik vilde overvælde ham; men denne anden Skikkelse er ham selv. Saaledes formaaer han slet Intet; selv det svageste Menneske, der ikke forsøges i denne Strid, formaaer langt mere end han. Og denne Strid er ikke blot anstrængende, men tillige meget forfærdelig (dersom det dog ikke er ham selv, der følgende sit eget Paafund har vovet sig ud i den, og hvis saa er, da forsøges han ikke i den Strid, om hvilken vi tale) naar Livet ved Guds Styrelse kaster et Menneske ud for at styrkes i denne Tilintetgjørelse, der ingen Skuffelse kjender, ingen Udflugt tillader, intet Selvbedrag foranlediger, som formaaede han under andre Omstændigheder noget Mere; thi da han strider med sig selv, kan Omstændigheder ikke gjøre Udslaget. Denne er et Menneskes Tilintetgjørelse, og Tilintetgjørelsen hans Sandhed. Han skal ikke slippe ud af denne Erkjendelse; thi han er jo sig selv sit eget Vidne, sin Anklager, sin Dommer, han er sig selv den Eneste der formaaede at trøste, da han forstaaer Tilintetgjørelsens Nød, og den Eneste, der ikke kan trøste, da han jo selv er Tilintetgjørelsens Redskab. At fatte denne Tilintetgjørelse, er det Høieste, hvad et Menneske formaaer, at ruge over denne Forstaaelse, fordi den er ham et betroet Gode, af Gud i Himlene nemlig betroet ham som Sandhedens Hemmelighed, er det Høieste og det Vanskeligste hvad et Menneske formaaer; thi Bedraget og Falskneriet er let gjort, saa han selv paa Sandhedens Bekostning bliver til Noget. Dette er det Høieste og Vanskeligste et Menneske formaaer, dog hvad siger jeg, end ikke dette formaaer han, han formaaer i det Høieste at ville forstaae, at denne tørre Brand kun fortærer indtil Gudskjærlighedens Ild tænder Luen i hvad den tørre Brand ikke kunde fortære. – Saaledes er Mennesket en hjælpeløs Skabning; thi al anden Forstaaelse, hvorved han forstaaer, at han kan hjælpe sig selv, er kun en Misforstaaelse, om han end i Verdens Øine blev anseet for modig – ved at have Mod til at blive i en Misforstaaelse, det er, ved ikke at have Mod til at forstaae Sandheden.

Men i Himlene, min Tilhører, der boer den Gud, som formaaer Alt, eller rettere, han boer jo allevegne, om Menneskene end ikke mærke det; »ja var Du, o Herre, et afmægtigt, livløst Legeme som en Blomst, der visner; var Du som en Bæk, der rinder forbi; var Du som en Bygning, der synker hen i Tiden; saa vilde Menneskene agte paa Dig, da var Du en passende Gjenstand for vore lave og dyriske Tanker«; men nu er det ikke saaledes, og Din Storhed gjør Dig netop usynlig; thi i Din Viisdom er Du for langt borte fra Menneskets Tanke, til at han kan see Dig, og i Din Allestedsnærværelse er Du ham for nær til at han kan see Dig; i Din Godhed skjuler Du Dig for ham, og Din Almagt gjør, at han ikke kan see Dig, thi saa blev han selv til Intet! Men Gud i Himlene formaaer Alt, og Mennesket slet Intet.

Min Tilhører, er det nu ikke saaledes, at disse Tvende passe for hinanden: Gud og Mennesket? Men passe de for hinanden, da er jo kun Spørgsmaalet, om Du vil være glad over denne vidunderlige Lykke, at I Tvende passe for hinanden; eller om Du hellere vil være en Saadan, der slet ikke passede for Gud, en Saadan, der selv formaaede Noget – og altsaa ikke ganske passede for Gud: thi Gud kan Du jo ikke og kan Du jo ikke ville forandre, at han ikke skulde formaae Alt. At blive til Intet, det synes haardt, o, men selv i det Menneskelige tale vi anderledes. Thi hvis Ulykken lærte to Mennesker, at de passede for hinanden i Venskab eller i Elskov, hvor ringe syntes dem dog ikke den Nød, som Ulykken bragte, i Sammenligning med den Glæde, som Ulykken ogsaa bragte: at disse Tvende passede for hinanden! Og hvis Tvende først i Døden forstode, at de passede for hinanden i al Evighed, o, hvad var da hiint korte om end bittre Adskillelsens Øieblik, der er Dødens, i Sammenligning med en evig Forstaaelse!

Saaledes er Mennesket stort og paa sit Høieste, naar han passer for Gud ved selv at være slet Intet; men lader os ikke letsindigen beundre eller tage Beundringen forfængeligt. Gik ikke Moses som Herrens Udsending til et vanartet Folk for at udfrie dette fra det selv, dets Trællesind, og fra dets Trællestand under en Tyrans Aag? Hvad er selv den største Helts Daad i Sammenligning med hvad man saaledes kalder Moses's Gjerninger; thi hvad er det at sløife Bjerge og fylde Floder mod at lade et Mørke falde over hele Ægypten! Men dette var jo ogsaa kun Moses's saakaldte Gjerninger, thi han formaaede slet Intet, og Gjerningen var Herrens. Seer her Forskjellen. Moses, han fatter ikke Beslutninger, udkaster ikke Planer, medens de Forstandiges Raad opmærksomt lytter til, fordi Anføreren er den Viseste; Moses formaaer slet Intet. Hvis Folket vilde sige til ham: gaa Du til Pharao, thi Dit Ord er mægtigt, Din Røst seierrig, Din Veltalenhed uimodstaaelig, da havde han vel svaret: »O, I Daarer! jeg formaaer slet Intet, end ikke at lade mit Liv for Eder, hvis Herren ikke saa vil det, jeg formaaer kun at henstille Alt til Herren«. Da træder han frem for Pharao, og hvilket er hans Vaaben? Den Afmægtiges – Bønner, og selv naar Bønnens sidste Ord allerede har naaet Himlen, da veed han endnu ikke, hvad der skal skee, om han end troer, at hvad der saa end skeer, skeer dog det Bedste. Da vender han hjem til Folket, men hvis det vilde prise og takke ham, da vilde han vel svare: »jeg formaaer slet Intet.« Eller naar Folket tørster i Ørkenen, og maaskee tyer til Moses og siger: tag Din Stav og byd Klippen at give Vand, da skulde vel Moses svare: »hvad er min Stav andet end en Kjæp?« Og vedblev Folket: men i Din Haand er Staven mægtig; da maatte vel Moses sige: »jeg formaaer slet Intet; men da dog Folket begjerer det, da jeg selv ikke kan udholde Synet af de Vansmægtendes Elendighed; saa slaaer jeg Klippen, om jeg selv end ikke troer, at der vil springe Vand af den« – og Klippen gav ikke Vand. Altsaa: om den Stav, han holder i sin Haand, skal være Almagtens Finger, eller den skal være Moses's Kjæp, veed han ikke, end ikke i det Øieblik, da Staven allerede berører Klippen; han veed det først bag efter, som han bestandig kun seer Herrens Ryg. O! menneskeligt talt formaaer den Svageste i Israel mere end Moses, thi denne mener dog, at der er Noget, han formaaer, men Moses formaaer slet Intet. I det ene Øieblik ligesom at være stærkere end den Stærkeste, end alle Mennesker, end hele Verden, forsaavidt Underet skeer ved hans Haand; i det næste Øieblik, ja i samme Øieblik at være svagere end den Svageste, forsaavidt denne dog endnu bestandigen mener, at der er Noget, han formaaer – en saadan Storhed skal ikke friste Forfængelighedens Higen, forsaavidt den da giver sig Tid til at forstaae, hvori Storheden ligger; thi ellers skulde den vel flux være rede til med sin væmmelige Feighed at ønske at være i Moses's Sted.

Dog gjør denne Betragtning, at det at trænge til Gud er Menneskets høieste Fuldkommenhed, end Livet vanskeligere, da gjør den det jo kun, fordi den vil betragte Mennesket efter hans Fuldkommenhed, og bringe ham selv til at betragte sig saaledes; thi i og ved denne Betragtning lærer Mennesket at kjende sig selv. Og det Menneske, der ikke kjender sig selv, hans Liv er jo i dybere Forstand en Skuffelse. Dog en saadan Skuffelse forskyldes vel sjeldent, at et Menneske ikke opdager, hvilke Evner der ere ham betroede, ikke søger at udvikle disse saa meget som muligt i Overeensstemmelse med de Livsforhold, der ere ham opgivne, saa han ret slaaer dybe Rødder i Tilværelsen, ikke omgaaes sig selv letsindigen liig det lykkeligt begavede Barn, der ikke forstaaer, hvor meget der dog er det betroet, liig den letsindige rige Yngling, der ikke veed hvad Guld har at betyde – og saaledes er jo ogsaa det vi kalde et Menneskes Selv ligesom Penges Værd; og Den, der kjender sig selv, han veed indtil det mindste, hvor meget han lyder paa, og veed at sætte sig selv om, saa han faaer den hele Værdi ud. Gjør han ikke dette, da kjender han ikke sig selv, og han er skuffet, hvad den Forstandige nok vil sige ham, og sige ham Skridt for Skridt, eftersom Livet gaaer frem: at han ikke fryder sig ved Livet i Dagenes Vaar; at han ikke gjør sig gjeldende for hvad han virkelig er; at han ikke veed at Menneskene tage En for hvad man selv giver sig for, at han ikke har vidst at gjøre sig selv betydningsfuld og derved at give Livet Betydning for sig. Ak men om et Menneske ogsaa i denne Forstand nok saa godt kjendte sig selv, om han end nok saa godt vidste, at faae sig selv saa fordeelagtigt som muligt udsat paa Rente i Livet; mon han derfor kjendte sig selv? Men kjendte han ikke sig selv, saa var jo dog hans Liv i dybere Forstand en Skuffelse. Skulde det nu ogsaa være sjeldent, at et Menneske i disse kloge Tider forskyldte en saadan Skuffelse? Hvilken var nemlig hiin kløgtige Selvkundskab anden end denne, at han kjendte sig selv i Forhold til noget Andet, men ikke kjendte sig selv i Forhold til sig selv, det vil sige: hans hele Selvkundskab blev trods den tilsyneladende Tilforladelighed aldeles svævende, da den kun angik Forholdet mellem et tvivlsomt Selv og et tvivlsomt Andet; thi dette Andet kunde forandres, saa en Anden blev den Stærkeste, den Skjønneste, den Rigeste, og dette Selv kunde forandres, saa han selv blev fattig, blev styg, blev afmægtig; og denne Forandring kunde indtræde hvert Øieblik. Blot nu dette Andet tages bort, saa er han jo skuffet; og dersom dette Andet er noget Saadant, at det kan tages bort, er han jo skuffet, selv om det ikke toges bort, fordi hans hele Livs Betydning var grundet paa noget Andet. Det er nemlig ingen Skuffelse af det, som kan skuffe, at det skuffer: men snarere en Skuffelse, naar det lader det være.

En saadan Selvkundskab er da saare ufuldkommen, og langtfra at betragte Mennesket efter hans Fuldkommenhed; thi det var jo en besynderlig Fuldkommenhed, om hvilken der, efterat man maaskee havde priset den i de stærkeste Udtryk, slutteligen maatte siges: tillige er den en Skuffelse? Ad denne Vei kommer man ikke til at betragte Mennesket efter hans Fuldkommenhed, og for at begynde derpaa maa man begynde med at rive sig løs fra enhver saadan Betragtning, hvilket er vanskeligt nok, som det at rive sig ud af en Drøm, at man ikke tager feil og fortsætter Drømmen: drømmende, at man er vaagnet; vidtløftigt nok i en vis Forstand, fordi et Menneskes egentlige Selv synes ham at ligge saa fjernt, at hele Verden ligger ham langt nærmere; forfærdeligt nok, fordi den dybere Selvkundskab begynder med, hvad Den, der ikke vil forstaae det, maatte kalde det ængstende Bedrag: istedetfor hele Verden at faae sig selv, istedetfor Herre at blive Trængende, istedetfor at formaae Alt at formaae slet Intet. Ak, hvor vanskeligt, at man ikke her igjen falder hen i Drømme, og drømmer, at man gjør det ved sin egen Kraft.

Naar da Mennesket vender sig mod sig selv for at forstaae sig selv, træder han ligesom hiint første Selv iveien, standser det, der jo var vendt ud ad, i at hige og tragte efter den Omverden, der er dets Gjenstand, kalder det tilbage fra det Udvortes. For at bevæge det første Selv til denne Tilbagekaldelse lader det dybere Selv Omverdenen blive, hvad den er, blive tvivlsom. Saaledes er det jo ogsaa, at Verdenen om os er Ustadighed og ethvert Øieblik kan forandres til det Modsatte; og der er ikke fundet det Menneske, der ved sin Magt eller ved sit Ønskes Besværgelse kan tvinge denne Vexel. Det dybere Selv danner nu hiin Omverdenens svigefulde Bøielighed saaledes, at den ikke mere bliver hiint første Selv attraaværdig. Enten maa da det første Selv see at dræbe det dybere Selv, at faae det glemt, hvorved da Alt er opgivet; eller det maa indrømme, at det dybere Selv har Ret; thi at ville udsige Bestandighed om Det, som bestandig vexler, er jo en Modsigelse; og saasnart man bekjender, at det vexler, kan det jo vexle i samme Øieblik. Hvor Meget end hiint første Selv krymper sig derved, der er aldrig fundet nogen saa snild Ordgyder eller saa snedig Tankefalskner, at han kan afkræfte det dybere Selvs evige Paastand; der gives kun een Udvei, det er at bringe det dybere Selv til Taushed ved at lade Ustadighedens Brusen overdøve det. – Hvad er nu skeet? Det første Selv er standset, kan aldeles ikke røre sig. Ak, Omverdenen kan i Virkeligheden være saa begunstigende, saa haandgribelig trofast, saa tilsyneladende usvigelig, at Enhver vil borge for en lykkelig Fremgang, naar man blot begynder; det hjælper ikke. Det Menneske, der er Vidne til hiin Strid i hans Indre, maa give det dybere Selv Ret: i dette Minut kan Alt være forandret, og Den, der ikke opdager det, han løber bestandigt paa det Uvisse. Der har i Verden aldrig været saa hurtig en Tunge, at den kunde bedaare det dybere Selv, saasnart det blot faaer Lov at komme til Orde. Ak, det er en smertefuld Tilstand: det første Selv sidder og seer efter alle de vinkende Frugter; og det er jo saa klart, at blot man griber til, saa lykkes Alt, det vil ethvert Menneske indrømme – men det dybere Selv sidder alvorligt, betænkeligt som Lægen ved den Syges Leie, om dog ogsaa med sin forklarede Mildhed, fordi det veed, at denne Sygdom er ikke til Døden men til Livet. Det første Selv har nu en bestemt Higen; det veed med sig selv, at det er i Besiddelse af Betingelserne; Omverdenen, saaledes som det forstaaer den, er saa begunstigende som muligt; de vente ligesom blot paa hinanden: det lykkelige Selv og Lykkens Begunstigelser – ak, hvilket lysteligt Liv! men det dybere Selv viger ikke, det tinger ikke, giver ikke sit Minde, slutter ikke Forlig, det siger blot: endnu i dette Øieblik kan jo Alt være forandret. Dog Menneskene hjælpe hiint første Selv med Forklaringen, de kalde paa ham, de forklare, at saaledes gaaer det til i Livet, at der ere nogle Mennesker, som ere lykkelige, og de skulle nyde Livet, og han er een af dem. Da banker Hjertet, han vil afsted... at et Barn, der har en streng Fader, maa blive hjemme, det maa man finde sig i, thi Faderen er jo den Stærkeste; men han er jo intet Barn, og hiint dybere Selv er jo ham selv, og dog synes det strengere end den strengeste Fader, man kan slet ikke indsmigre sig hos det, det vil enten tale uforbeholdent, eller slet ikke. Da er der Fare paafærde, det mærke de Begge, baade det første og det dybere Selv, og da sidder dette bekymret som den erfarne Lods, medens der hemmeligen holdes Raad, om det ikke var bedst, at kaste Lodsen overbord, da han skaffer Modvind. Dog skeer det ikke, men hvad er Følgen? Det første Selv kan ikke røre sig af Stedet, og dog, dog er det tydeligt, at Glædens Øieblik iler, at Lykken allerede flygter; thi det sige jo Menneskene, at naar man ikke benytter Øieblikket strax, da er det snart bag efter. Og hvem er saa Skyld deri? Hvem Anden end hiint dybere Selv? Men selv dette Skrig, det hjælper ikke. – Hvad er dog dette for en unaturlig Tilstand; hvad skal det Hele betyde? Naar noget Saadant foregaaer i et Menneskes Sjel, betyder dette da ikke, at hans Sind begynder at blive svagt? Ak nei, det betyder noget ganske Andet, det betyder, at Barnet skal til at vænnes fra. Thi man kan jo være tredive Aar og derover, fyrretive Aar og dog kun Barn, ja man kan døe som et gammelt Barn. Men det at være Barn er saa deiligt! Saa ligger man ved Timelighedens Bryst i Endelighedens Vugge, og Sandsynligheden sidder ved Vuggen og synger for Barnet. Gaaer Ønsket ikke i Opfyldelse og Barnet bliver uroligt, saa dysser Sandsynligheden paa det og siger: lig nu blot stille og sov, saa gaaer jeg ud og kjøber Noget til Dig; og næste Gang saa kommer Touren til Dig. Saa sover Barnet ind igjen og saa er Smerten glemt, og Barnet rødmer atter i nye Ønskers Drømme, uagtet det troede, at det var en Umulighed at glemme Smerten; nu, det forstaaer sig, hvis han ikke havde været et Barn, havde han vel ikke glemt Smerten saa let, og det vilde have viist sig, at det ikke var Sandsynligheden, der havde siddet ved Vuggen; men det dybere Selv der havde siddet hos ham ved Dødsleiet i Selvfornægtelsens Dødsstund, da det selv opstod til en Evighed.

Naar da det første Selv falder det dybere Selv tilfode, da ere de forligte og følges ad. Da siger vel det dybere Selv: »det er sandt, nær havde jeg glemt det over vor megen Strid, hvad var det nu Du saa inderligen ønskede; i dette Øieblik troer jeg ikke, at der er Noget til Hinder for Ønskets Opfyldelse, naar Du blot ikke glemmer den lille Hemmelighed, vi To have med hinanden. Seer Du nu, nu kan Du da vel være fornøiet«. Maaskee svarer det første Selv: »ja nu bryder jeg mig ikke saaledes om det; nei, jeg bliver aldrig glad som før, o, som dengang, da min Sjel higede derefter; og egentligen forstaaer Du mig ikke«. »Det troer jeg heller ei; og det var vel heller ei ønskeligt, at jeg forstod Dig saaledes, at jeg higede ligesaa meget som Du. Dog skulde Du have tabt Noget ved at Du ikke saaledes bryder Dig derom? betænk paa den anden Side, hvis nu Omverdenen havde bedraget Dig, og det veed Du jo den kunde, mere sagde jeg ikke, jeg sagde blot det er muligt, hvorved jeg jo rigtignok tillige sagde, at det Du ansaae for Vished egentligen ogsaa kun var en Mulighed – hvad saa, saa havde Du fortvivlet, og Du havde ikke havt mig at stole paa; thi Du erindrer nok, at Skibsraadet næsten var betænkt paa, at kaste mig overbord. Skulde Du nu ikke være bedre tjent med at have tabt Noget af hiin brændende Begjer, og vundet, at Livet ikke kan bedrage Dig; at tabe paa den Maade, er det ikke at vinde?«

Hiin lille Hemmelighed, som vi To have med hinanden, saaledes sagde det dybere Selv. Hvilken er vel denne Hemmelighed, min Tilhører? Hvilken anden, end at et Menneske i Forhold til det Udvortes slet Intet formaaer. Vil han umiddelbart gribe efter det Udvortes, da kan det i samme Nu forandres og han være bedragen; derimod kan han tage det med den Bevidsthed, at det ogsaa kunde forandres, og han er ikke bedragen, selv om det forandredes; thi han har det dybere Selvs Samtykke. Vil han handle umiddelbart i det Udvortes, udføre Noget, da kan Alt i samme Øieblik blive til Intet; derimod kan han med denne Bevidsthed handle, og selv om det blev Intet, han er ikke bedragen, thi han har det dybere Selvs Samtykke.

Dog om end saaledes det første og det dybere Selv ere forligte, og det fælleds Sind vendt bort fra det Udvortes, saa er det jo kun Betingelsen for at kunne komme til at lære sig selv at kjende. Men for at han virkeligen skal kjende sig selv, er der ny Strid og nye Farer. Den Stridende lade sig blot ikke forfærde og afskrække ved den Tanke som var det at trænge en Ufuldkommenhed, naar Talen er om at trænge til Gud, som var det at trænge en beskæmmende Hemmelighed man helst skjulte, naar Talen dog er om at trænge til Gud; som var det at trænge en sørgelig Nødvendighed, man søgte at afvinde en blidere Side ved selv at udsige det, naar Talen dog var om at trænge til Gud. Ved den dybere Selverkjendelse lærer man netop, at man trænger til Gud; men det ved første Øiekast Nedtrykkende heri vilde afskrække En fra at begynde, hvis man ikke i Tide var opmærksom derpaa og begeistret ved Tanken, at dette netop er Fuldkommenheden, da det er langt ufuldkomnere og kun en Misforstaaelse, at man ikke skulde trænge til Gud. Thi selv om et Menneske havde fuldkommet de herligste Bedrifter, dersom han dog meente, at det var Alt ved hans egen Kraft, dersom han ved at overvinde sit Sind var bleven større end Den, der indtog en Stad, naar han dog meente, at det var skeet ved hans egen Kraft, da var hans Fuldkommenhed væsentligen kun en Misforstaaelse; men en saadan Fuldkommenhed var jo kun lidet priselig. Den derimod, som indsaae, at han ikke formaaede det Mindste, end ikke at være glad over den glædeligste Begivenhed, uden Gud, han er Fuldkommenheden nærmere; og Den, som forstod dette, og aldeles intet Smerteligt fandt deri, men kun Salighedens Overflødighed, som intet lønligt Ønske gjemte, der dog hellere vilde være glad paa egen Haand, ingen Beskæmmelse følte ved, at Menneskene mærkede det paa ham, at han saaledes selv slet Intet formaaede, ingen Betingelse gjorde med Gud, end ikke den, at hans Afmagt skulde bevares skjult for Andre, men i hvis Hjerte Glæden bestandig seirede ved, saa at sige, jublende at kaste sig i Guds Arme, i navnløs Forundring over Gud, som jo formaaer Alt – ja han var den Fuldkomne, som Apostelen Paulus bedre og kortere beskriver ham: han »roser sig af sin Svaghed,« og har end ikke gjort saa mange og tvetydige Erfaringer, at han veed at udtrykke sig vidtløftigere. – Ikke at kjende sig selv, sige jo Menneskene, er en Skuffelse og Ufuldkommenhed; men ofte ville de ikke forstaae, at Den, der virkeligen kjender sig selv, netop fatter, at han Intet formaaer.

I det Udvortes formaaede han Intet; men indvortes, formaaer han der ei heller Noget? Skal en Formaaen virkelig være en saadan, da maa den jo have en Modstand, thi har den ingen Modstand, da er den enten almægtig eller en Indbildning. Men skal han have en Modstand, hvorfra skal da denne komme? I det Indvortes kan Modstanden jo kun komme fra ham selv. Saa strider han da med sig selv i det Indvortes, ikke som før, hvor det dybere Selv stred med det første Selv, for at hindre det i, at bekymre sig om det Udvortes. Opdager et Menneske ikke denne Strid, da er han i en Misforstaaelse, og hans Liv altsaa ufuldkomment; men opdager han den, da vil han atter forstaae, at han selv slet Intet formaaer.

Det synes underligt, at det er dette, et Menneske skal lære af sig selv; hvortil da prise Selverkjendelsen? Og dog er det saaledes; og af hele Verden kan et Menneske ikke lære, at han aldeles Intet formaaer. Selv om hele Verden forenede sig om at knuse og tilintetgjøre den Svageste, han kunde dog endnu bestandig beholde en ganske svag Forestilling om, at han selv formaaede Noget, under andre Omstændigheder, naar Overmagten ikke var saa stor. At han slet Intet formaaer, kan han kun opdage ved sig selv; og hvad enten han seirer over hele Verden, eller han snubler over et Halmstraa; det bliver dog tilbage, at han veed dette med sig selv eller kan vide det, at han selv slet Intet formaaer. Vil Nogen forklare det anderledes, da har han jo ikke med de Andre at gjøre, men kun med sig selv, og da er enhver Udflugt gjennemskuet. Det er saa svært, mene Menneskene, at kjende sig selv, især naar man er meget begavet, og har mangfoldige Anlæg og Evner, og nu skal vide Beskeed om alt dette. O, den Selverkjendelse, om hvilken vi tale, er jo ikke vidtløftig; og hver Gang man ret fatter denne korte og fyndige Sandhed, at man selv slet Intet formaaer, da har man erkjendt sig selv.

[…]

\workheading{Frygt og Bæven (af Johannes de silentio)}{1843}

\worknote{Antologiens ''Fear and Trembling''. Hong-udvalget er stærkt sammensat: bogens Forord, et uddrag af ''Foreløbig Expectoration'' (Troens Ridder — »ligner en Rodemester«) samt et uddrag af Problema I. De store spring følger Hongs margental og er markeret med […].}

\sectionheading{Forord  ·  SV¹ III, 57–59}

Ikke blot i Handelens, men ogsaa i Ideernes Verden foranstalter vor Tid ein wirklicher Ausverkauf. Alt faaes for en saadan Spot-Priis, at det bliver et Spørgsmaal, om der tilsidst er Nogen, der vil byde. Enhver speculativ Marqueur, der samvittighedsfuldt pointerer den nyere Philosophies betydningsfulde Gang, enhver Privatdocent, Repetent, Student, enhver Udflytter og Indsidder i Philosophien bliver ikke staaende ved at tvivle om Alt, men gaaer videre. Maaskee vilde det være utidigt og ubetimeligt at spørge dem, hvor de dog egentlig komme hen, men høfligt og beskedent er det vel at ansee det for afgjort, at de have tvivlet om Alt, da det jo ellers var en besynderlig Tale, at de gik videre. Denne foreløbige Bevægelse have de da Alle gjort og formodentlig saa let, at de ikke finde det fornødent at lade et Ord falde om, hvorledes; thi end ikke den, der ængstelig og bekymret søgte en lille Oplysning, fandt en saadan, et veiledende Vink, en lille diætetisk Forskrift om, hvorledes man forholder sig under denne uhyre Opgave. »Men Cartesius har jo gjort det?« Cartesius, en ærværdig, ydmyg, redelig Tænker, hvis Skrifter vist Ingen kan læse uden den dybeste Rørelse, han har gjort, hvad han har sagt, og sagt, hvad han har gjort. Ak! Ak! Ak! det er en stor Sjeldenhed i vor Tid! Cartesius har, som han selv ofte nok gjentager, ikke tvivlet i Forhold til Troen. (»Memores tamen, ut jam dictum est, huic lumini naturali tamdiu tantum esse credendum, quamdiu nihil contrarium a Deo ipso revelatur.... Præter cætera autem, memoriæ nostræ pro summa regula est infigendum, ea quæ nobis a Deo revelata sunt, ut omnium certissima esse credenda; et quamvis forte lumen rationis, quam maxime clarum et evidens, aliud quid nobis suggerere videretur, soli tamen auctoritati divinæ potius quam proprio nostro judicio fidem esse adhibendam. Confer Principia philosophiæ, pars prima § 28 og § 76). Han har ikke raabt Brand og gjort det til Pligt for Alle at tvivle, thi Cartesius var en stille eensom Tænker, ikke en brølende Gadevægter; han har beskedent bekjendt, at hans Methode kun havde Betydning for ham selv og havde tildeels sin Grund i hans tidligere forqvaklede Viden. (Ne quis igitur putet, me hîc traditurum aliquam methodum, quam unusquisque sequi debeat ad recte regendam rationem; illam enim tantum, quam ipsemet secutus sum, exponere decrevi... Sed simul ac illud studiorum curriculum absolvi (sc. juventutis), quo decurso mos est in eruditorum numerum cooptari, plane aliud coepi cogitare. Tot enim me dubiis totque erroribus implicatum esse animadverti, ut omnes discendi conatus nihil aliud mihi profuisse judicarem, quam quod ignorantiam meam magis magisque detexissem. Confer Dissertatio de methodo pagina 2 og 3). – Hvad hine gamle Grækere, der dog ogsaa forstode sig lidt paa Philosophien, antog at være en Opgave for hele Livet, fordi den tvivlende Færdighed ikke erhverves i Dage og Uger, hvad den gamle udtjente Strider opnaaede, der havde bevaret Tvivlens Ligevægt gjennem alle Besnærelser, uforfærdet negtet Sandsens Vished og Tankens Vished, ubestikkelig trodset Selvkjærlighedens Angst og Medfølelsens Insinuationer – dermed begynder i vor Tid Enhver.

I vor Tid bliver Enhver ikke staaende ved Troen, men gaaer videre. Et Spørgsmaal om, hvor de komme hen, vilde maaskee være en Dumdristighed, derimod er det vel et Tegn paa Belevenhed og Dannelse, at jeg antager, at Enhver har Troen, da det ellers bliver en besynderlig Tale: at gaae videre. I hine gamle Dage var det anderledes, da var Troen en Opgave for hele Livet, fordi man antog, at den troende Færdighed ikke erhverves hverken i Dage eller Uger. Naar da den prøvede Olding nærmede sig sit Endeligt, havde stridt den gode Strid og bevaret Troen, da var hans Hjerte ungt nok til ikke at have glemt hiin Angst og Bævelse, der tugtede Ynglingen, som Manden vel beherskede, men som intet Menneske ganske voxer fra – uden forsaavidt det skulde lykkes ved saa tidlig som mulig at gaae videre. Hvor da hine ærværdige Skikkelser naaede hen, der begynder i vor Tid Enhver for at gaae videre.

\sectionheading{Foreløbig Expectoration (uddrag)  ·  SV¹ III, 87–101}

Man mener i Almindelighed, at det, Troen frembringer, ikke er Kunstværk, at det er grovt og plumpt Arbeide, kun for de mere klodsede Naturer; dog er det langt anderledes. Troens Dialektik er det Fineste og det Mærkværdigste af Alt, den har en Elevation, som jeg vel kan gjøre mig en Forestilling om, men heller ikke mere. Jeg kan gjøre det store Tramplin-Spring, hvorved jeg gaaer over i Uendeligheden, min Ryg er som en Liniedandsers, vreden i min Barndom, derfor falder det mig let, jeg kan een, to, tre gaae paa Hovedet i Tilværelsen, men det næste kan jeg ikke; thi det Vidunderlige kan jeg ikke gjøre; men kun forbauses ved. Ja dersom Abraham i det Øieblik, han svingede sit Been over Æselets Ryg, har sagt til sig selv: nu er Isaak tabt, jeg kunde ligesaa gjerne offre ham her hjemme som reise den lange Vei til Morija – saa behøver jeg ikke Abraham, medens jeg nu bøier mig syv Gange for hans Navn og halvfjersindstive Gange for hans Gjerning. Dette har han nemlig ikke gjort, hvad jeg kan bevise deraf, at han blev glad ved at modtage Isaak, ret inderlig glad, at han ingen Forberedelse behøvede, ingen Tid til at samle sig paa Endeligheden og dens Glæde. Hvis det ikke stod sig saaledes med Abraham, da havde han maaskee elsket Gud, men ikke troet; thi den, der elsker Gud uden Tro, han reflekterer paa sig selv, den, der elsker Gud troende, han reflekterer paa Gud.

Paa denne Spidse staaer Abraham. Det sidste Stadium, han taber af Sigte, er den uendelige Resignation. Han gaaer virkelig videre og kommer til Troen; thi alle disse Vrængbilleder af Troen, den jammerlige lunkne Dorskhed, der tænker: det har vel ingen Nød, det er ikke værd at sørge før Tiden; det usle Haab, der siger: man kan ikke vide, hvad der vil skee, det var dog muligt – disse Vrængbilleder høre hjemme i Livets Elendighed, og dem har allerede den uendelige Resignation uendelig foragtet.

Abraham kan jeg ikke forstaae, jeg kan i en vis Forstand Intet lære af ham uden at forbauses. Naar man indbilder sig, at man ved at betænke Udfaldet af hiin Historie, skulde lade sig bevæge til at troe, saa bedrager man sig selv, og vil bedrage Gud for Troens første Bevægelse; man vil suge Leve-Viisdom ud af Paradoxet. Maaskee lykkes det En og Anden; thi vor Tid bliver ikke staaende ved Troen, ikke ved dens Mirakel, at gjøre Vand til Viin, den gaaer videre, den gjør Viin til Vand.

Var det dog ikke bedst, at blive staaende ved Troen, og er det ikke oprørende, at Enhver vil gaae videre? Naar man i vor Tid, og det forkyndes jo paa forskjellig Maade, ikke vil blive staaende ved Kjærligheden, hvor kommer man da hen? Til jordisk Kløgt, smaalig Beregning, til Usselhed og Elendighed, til alt Det, der kan gjøre Menneskets guddommelige Herkomst tvivlsom. Var det ikke bedst, at man blev staaende ved Troen, og at den, der staaer, saae til, at han ikke faldt; thi Troens Bevægelse maa bestandig gjøres i Kraft af det Absurde, dog vel at mærke saaledes, at man ikke taber Endeligheden, men vinder den heel og holden. Jeg for mit Vedkommende kan vel beskrive Troens Bevægelser, men jeg kan ikke gjøre dem. Naar man vil lære at gjøre Svømmebevægelserne, da kan man lade sig hænge i Seler under Loftet, man beskriver vel Bevægelserne, men man svømmer ikke; saaledes kan jeg beskrive Troens Bevægelser, men naar jeg bliver kastet ud i Vandet, da svømmer jeg vel (thi jeg hører ikke til Vaderne), men jeg gjør andre Bevægelser, jeg gjør Uendelighedens Bevægelser, medens Troen gjør det Modsatte, gjør, efter at have gjort Uendelighedens Bevægelser, Endelighedens. Held den, der kan gjøre disse Bevægelser, han gjør det Vidunderlige, og jeg skal aldrig blive træt af at beundre ham, om det er Abraham eller det er Trællen i Abrahams Huus, om det er en Professor i Philosophi eller en stakkels Tjenestepige, er mig absolut ligegyldigt, jeg seer blot paa Bevægelserne. Men dem seer jeg ogsaa paa, og lader mig ikke narre, hverken af mig selv eller af noget Menneske. Ridderne af den uendelige Resignation kjender man let, deres Gang er svævende, dristig. De derimod, der bære Troens Klenodie, skuffe let, fordi deres Ydre har en paafaldende Lighed med det, som saavel den uendelige Resignation som Troen dybt foragter – med Spidsborgerlighed.

Jeg tilstaaer oprigtigt, jeg har i min Praxis ikke fundet noget paalideligt Exemplar, uden at jeg derfor vil negte, at maaskee hvert andet Menneske er et saadant Exemplar. Imidlertid har jeg dog i flere Aar eftersporet det forgjeves. Man reiser i Almindelighed Verden rundt for at see Floder og Bjerge, nye Stjerner, spraglede Fugle, vanskabte Fiske, latterlige Menneskeracer, man hengiver sig til den dyriske Stupor, der gloer paa Tilværelsen, og man mener at have seet Noget. Dette beskæftiger mig ikke. Vidste jeg derimod, hvor der levede en saadan Troens Ridder, da vilde jeg paa min Fod vandre til ham; thi dette Vidunder beskæftiger mig absolut. Jeg vilde intet Øieblik slippe ham, hvert Minut passe paa, hvorledes han bar sig ad med Bevægelserne; jeg vilde betragte mig selv som forsørget i Livet, og dele min Tid mellem at see paa ham og selv gjøre Øvelser, og saaledes anvende al min Tid paa at beundre ham. Som sagt jeg har ikke fundet nogen Saadan, imidlertid kan jeg vel tænke ham. Her er han. Bekjendtskabet stiftes, jeg bliver forestillet for ham. I det Moment, jeg først tager ham paa mine Øine, kaster jeg ham i samme Nu fra mig, gjør selv et Spring tilbage, slaaer Hænderne sammen og siger halv høit: »Herre Gud! er det Mennesket, er det virkeligt ham, han seer jo ud som en Rodemester.« Imidlertid er det dog ham. Jeg slutter mig lidt nærmere til ham, passer paa den mindste Bevægelse, om der ikke skulde vise sig en lille uensartet Brøks-Telegraphering fra det Uendelige, et Blik, en Mine, en Gestus, et Vemod, et Smil, der forraadte det Uendelige i sin Uensartethed med det Endelige. Nei! Jeg examinerer hans Skikkelse fra Top til Taa, om der ikke skulde være en Revne, igjennem hvilken det Uendelige tittede ud. Nei! Han er heelt igjennem solid. Hans Fodfæste? er kraftigt, tilhører ganske Endeligheden, ingen pyntet Borgermand, der Søndag-Eftermiddag gaaer ud paa Fresberg, træder grundigere paa Jorden, han tilhører ganske Verden, ingen Spidsborger kan tilhøre den mere. Intet er at opdage af dette fremmede og fornemme Væsen, hvorpaa man kjender Uendelighedens Ridder. Han glæder sig ved Alt, deeltager i Alt, og hver Gang man seer ham deeltage i det Enkelte, skeer det med en Vedholdenhed, som betegner det jordiske Menneske, hvis Sjæl hænger fast ved Sligt. Han passer sin Gjerning. Naar man da seer ham, skulde man troe, han var en Skriverkarl, der havde fortabt sin Sjæl i det italienske Bogholderi, saa punktlig er han. Han ferierer om Søndagen. Han gaaer i Kirke. Intet himmelsk Blik eller noget Tegn paa det Incommensurable forraader ham; hvis man ikke kjendte ham, var det umuligt, at udsondre ham af den øvrige Mængde; thi hans sunde kraftige Psalmesang beviser i det Høieste, at han har et godt Bryst. Om Eftermiddagen gaaer han til Skoven. Han fryder sig over Alt, hvad han seer, over Menneske-Vrimlen, de nye Omnibusser, Sundet, – naar man møder ham paa Strandveien skulde man troe, han var en Kræmmersjæl, der slog sig løs, netop saaledes glæder han sig; thi han er ikke Digter, og jeg har forgjeves søgt at aflure ham den poetiske Incommensurabilitet. Henad Aften gaaer han hjem, hans Gang er ufortrøden som et Postbuds. Underveis tænker han paa, at hans Kone vist har en apparte lille Ret varm Mad til ham, naar han kommer hjem, for Exempel et stegt Lammehoved med Grønt til. Hvis han mødte en Ligesindet, da kunde han vedblive lige til Østerport at samtale med ham om denne Ret med en Lidenskab, der vilde passe for en Restaurateur. Tilfældigvis eier han ikke 4 Skilling, og dog troer han fuldt og fast, at hans Kone har hiin lækkre Ret til ham. Har hun den, da skal det være et misundelsesværdigt Syn for fornemme Folk, begeistrende for Menigmand at see ham spise; thi hans Appetit er stærkere end Esaus. Hans Kone har den ikke – besynderligt nok – han er aldeles den Samme. Paa Veien kommer han forbi en Byggeplads, han træffer en anden Mand. De tale et Øieblik sammen, han opfører i et Nu en Bygning, han disponerer over alle Kræfter dertil. Den Fremmede forlader ham med den Tanke, det var vist en Capitalist, medens min beundrede Ridder tænker: jo kom det derpaa an, kunde jeg sagtens faae det. Han ligger i et aabent Vindue og betragter Pladsen, paa hvilken han boer, Alt hvad der foregaaer, at en Rotte smutter ned under et Rendesteensbræt, at Børnene lege, alt beskæftiger ham med en Ro i Tilværelsen, som var han en Pige paa 16 Aar. Og dog er han ikke Genie; thi Geniets Incommensurabilitet har jeg forgjeves søgt at udspionere hos ham. Han ryger sin Pibe i Aftenstunden; naar man seer ham, skulde man sværge paa, det var Spekhøkeren ligeoverfor, der vegeterer i Skumringen. Han lader 5 være lige med en Sorgløshed, som var han en letsindig Døgenigt, og dog kjøber han hvert Øieblik, han lever, den beleilige Tid til den dyreste Priis; thi han gjør end ikke det Mindste uden i Kraft af det Absurde. Og dog, dog, ja jeg kunde blive rasende derover, om ikke af anden Grund saa af Misundelse, dog har dette Menneske gjort og gjør hvert Øieblik Uendelighedens Bevægelse. Han tømmer i den uendelige Resignation Tilværelsens dybe Vemod, han kjender Uendelighedens Salighed, han har fornummet Smerten af at forsage Alt, det Kjæreste, man har i Verden, og dog smager Endeligheden ham fuldt saa godt som den, der aldrig kjendte noget Høiere, thi hans Forbliven i Endeligheden havde intet Spoer af en forknyt beængstet Dressur, og dog har han denne Tryghed til at fryde sig ved den, som var den det Visseste af Alt. Og dog, dog er hele den jordiske Skikkelse, han frembringer, en ny Skabning i Kraft af det Absurde. Han resignerede uendeligt paa Alt, og da greb han Alt igjen i Kraft af det Absurde. Han gjør bestandig Uendelighedens Bevægelse, men han gjør det med en saadan Correcthed og Sikkerhed, at han bestandig faaer Endeligheden ud, og der er intet Secund, hvor man ahner noget Andet. Det skal være den vanskeligste Opgave for en Dandser, at springe sig ind i en bestemt Stilling, saaledes, at der intet Secund er, da han griber efter Stillingen, men i Springet selv staaer i Stillingen. Maaskee kan ingen Dandser gjøre det – dette gjør hiin Ridder. Mængden af Mennesker lever fortabt i verdslig Sorg og Glæde, disse ere Oversidderne, som ikke komme med i Dandsen. Uendelighedens Riddere ere Dandsere og have Elevation. De gjøre Bevægelsen op efter og falde ned igjen, og ogsaa dette er en ikke usalig Tidsfordriv og ikke uskjønt at see paa. Men hver Gang de falde ned, kunne de ikke strax antage Stillingen, de vakle et Øieblik, og denne Vaklen viser, at de dog ere Fremmede i Verden. Den er mere eller mindre paafaldende, eftersom de have Kunst, men selv den kunstigste af disse Riddere kan dog ikke skjule denne Vaklen. Man behøver ikke at se dem i Luften, man behøver blot at see dem i det Øieblik, de berøre og have berørt Jorden – og man kjender dem. Men at kunne falde saaledes ned, at det i samme Secund seer ud som stod og gik man, at forvandle Springet i Livet til Gang, absolut at udtrykke det Sublime i det Pedestre – det kan kun hiin Ridder, – og dette er det eneste Vidunder.

Dog dette Vidunder kan saa let skuffe, jeg vil derfor beskrive Bevægelserne i et bestemt Tilfælde, der kan belyse disses Forhold til Virkeligheden; thi derom er det at Alt dreier sig. En Ungersvend forelsker sig i en Prindsesse, og hele hans Livs Indhold ligger i denne Kjærlighed, og dog er Forholdet et saadant, at den umulig lader sig realisere, umuligt lader sig oversætte fra Idealiteten til Realiteten. Elendighedens Trælle, Frøerne i Livets Sump, skrige naturligvis: en saadan Kjærlighed er en Daarskab, og den rige Bryggerenke er et fuldkomment saa godt og solidt Parti. Lad dem uforstyrret qvække i Sumpen. Saaledes gjør Ridderen af den uendelige Resignation ikke, han opgiver ikke Kjærligheden, ikke for al Verdens Herlighed. Han er ingen Gjæk. Han forvisser sig først om, at den virkelig er ham Livets Indhold, og hans Sjæl er for sund og for stolt til at ødsle det Mindste paa en Ruus. Han er ikke feig, han frygter ikke for at lade den snige sig ind i hans lønligste, hans mest afsides Tanker, at lade den snoe sig i utallige Slyngninger om ethvert Ligament i hans Bevidsthed – bliver Kjærligheden ulykkelig, da vil han aldrig kunne rive sig ud af den. Han føler en salig Vellyst ved at lade Kjærligheden gjennemgyse hver hans Nerve, og dog er hans Sjæl høitidelig som Dens, der har tømt Giftbægeret og føler, hvorledes Saften gjennemtrænger hver en Blodsdraabe, – thi dette Øieblik er Liv og Død. Naar han da saaledes har indsuget hele Kjærligheden i sig og fordybet sig i den, da mangler han ikke Mod til at forsøge og vove Alt. Han overskuer Livets Forhold, han sammenkalder de hurtige Tanker, der som slagvante Duer lyde hvert hans Vink, han svinger Staven over dem, og de styrte sig i alle Retninger. Men naar de nu alle vende tilbage, alle som Sorgens Bud og forklare ham, at det er en Umulighed, da bliver han stille, han takker dem af, han bliver ene og da foretager han sig Bevægelsen. Dersom hvad, jeg her siger, skal have nogen Betydning, da gjælder det om, at Bevægelsen skeer normalt. Ridderen vil da for det første have Kraft til at concentrere hele Livets Indhold og Virkelighedens Betydning i eet eneste Ønske. Mangler et Menneske denne Concentration, denne Sluttethed, er hans Sjæl fra Begyndelsen adsplittet i det Mangfoldige, da kommer han aldrig til at gjøre Bevægelsen, han vil i Livet handle klogt som de Pengemænd, der anbringe deres Capital i alle forskjellige Papirer, for at vinde paa det ene, naar de tabe paa det andet – kort, saa er han ikke Ridder. Dernæst vil Ridderen have Kraft til at concentrere hele Tanke-Operationens Resultat i een Bevidstheds-Akt. Mangler han denne Sluttethed, er hans Sjæl fra Begyndelsen adsplittet i det Mangfoldige, da vil han aldrig faae Tid til at gjøre Bevægelsen, han vil bestandig løbe Ærinder i Livet, aldrig gaae ind i Evigheden; thi selv i det Øieblik, da han er nærmest derved, vil han pludselig opdage, at han har glemt Noget, hvorfor han maa tilbage. I næste Øieblik, vil han tænke, er det muligt, og det er ogsaa ganske sandt; men ved slige Betragtninger kommer man aldrig til at gjøre Bevægelsen, men ved Hjelp af dem synker man dybere og dybere i Mudret.

Ridderen gjør da Bevægelsen, men hvilken? Vil han glemme det Hele; thi ogsaa deri ligger jo en Slags Concentration? Nei! thi Ridderen modsiger ikke sig selv, og det er en Modsigelse at forglemme hele sit Livs Indhold og dog blive den Samme. At blive en Anden, føler han ingen Drift til, og anseer det ingenlunde for det Store. Kun de lavere Naturer glemme sig selv og blive noget Nyt. Saaledes har Sommerfuglen aldeles glemt, at den var Kaalorm, maaskee kan den igjen glemme, at den var Sommerfugl saa aldeles, at den kan blive en Fisk. De dybere Naturer glemme aldrig sig selv og blive aldrig til Andet end hvad de vare. Ridderen vil da erindre Alt; men denne Erindren er netop Smerten, og dog er han i den uendelige Resignation forsonet med Tilværelsen. Kjærligheden til hiin Prindsesse blev for ham Udtrykket for en evig Kjærlighed, antog en religiøs Charakteer, forklarede sig i en Kjærlighed til det evige Væsen, der vel negtede Opfyldelsen, men dog atter forsonede ham i den evige Bevidsthed om dens Gyldighed i en Evigheds-Form, som ingen Virkelighed kan fratage ham. Daarer og unge Mennesker snakke om, at Alt er muligt for et Menneske. Det er imidlertid en stor Vildfarelse. Aandelig talt er Alt muligt, men i Endelighedens Verden er der Meget, der ikke er muligt. Dette Umulige gjør imidlertid Ridderen muligt derved, at han udtrykker det aandeligt, men aandeligt udtrykker han det derved, at han giver Afkald derpaa. Ønsket, der vilde føre ham ud i Virkeligheden men strandede paa Umuligheden, bøies nu indefter, men er derfor ikke tabt, heller ikke glemt. Snart er det Ønskets dunkle Rørelser i ham, der vække Erindringen, snart vækker han den selv; thi han er for stolt til at ville, at det, der var hele hans Livs Indhold, skulde have været et flygtigt Moments Sag. Han holder denne Kjærlighed ung, og den tiltager med ham i Aar og i Skjønhed. Derimod behøver han ingen Endelighedens Anledning til dens Fremvæxt. Fra det Øieblik, han har gjort Bevægelsen, er Prindsessen tabt. Han behøver ikke disse erotiske Nervezittringer ved at see den Elskede og saa videre, han behøver heller ikke i endelig Forstand bestandig at tage Afsked med hende, fordi han i evig Forstand erindrer hende, og han veed meget godt, at de Elskende, der ere saa forhippede paa endnu engang til Afsked at see hinanden for sidste Gang, have Ret i at være forhippede, Ret i at mene, at det er sidste Gang; thi de glemme snarest hinanden. Han har fattet den dybe Hemmelighed, at ogsaa i at elske et andet Menneske bør man være sig selv nok. Han tager intet endeligt Hensyn mere til hvad Prindsessen gjør, og netop dette beviser, at han har gjort Bevægelsen uendeligt. Her kan man faae Leilighed til at see, om Bevægelsen hos den Enkelte er sand eller tilløiet. Der var den, der ogsaa troede, at han havde gjort Bevægelsen, men see Tiden gik, Prindsessen gjorde noget Andet, hun giftede sig for Exempel med en Prinds, da tabte hans Sjæl Resignationens Elasticitet. Han viste derved, at han ikke havde gjort Bevægelsen rigtigt; thi den, der har resigneret uendeligt, han er sig selv nok. Ridderen hæver ikke sin Resignation, han bevarer sin Kjærlighed ligesaa ung, som den var i det første Øieblik, han slipper den aldrig fra sig, netop fordi han har gjort Bevægelsen uendeligt. Hvad Prindsessen gjør kan ikke forstyrre ham, det er kun de lavere Naturer, der have Loven for deres Handlinger i et andet Menneske, Præmisserne til deres Handlinger udenfor sig selv. Er Prindsessen derimod en Ligesindet, da vil det Skjønne fremkomme. Hun vil da introducere sig selv i den Ridderorden, i hvilken man ikke optages ved Ballotation, men af hvilken Enhver er Medlem, der har Mod til at indføre sig selv, den Ridderorden, der derved beviser sin Udødelighed, at den ingen Forskjel gjør paa Mand og Qvinde. Ogsaa hun vil bevare sin Kjærlighed ung og sund, ogsaa hun vil have forvundet sin Qvide, om hun end ikke, som der staaer i Visen: hver Nat ligger ved sin Herres Side. Disse Tvende vil da i al Evighed passe for hinanden, med en saadan taktfast harmonia præstabilita, at hvis nogensinde Øieblikket kom, et Øieblik, der dog ikke endeligt beskæftiger dem, thi da ældes de, hvis nogensinde Øieblikket kom, der tillod at give Kjærligheden sit Udtryk i Tiden, da vilde de være istand til at begynde netop der, hvor de vilde have begyndt, hvis de oprindeligen havde været forenede. Den, der forstaaer dette, hvad enten han er en Mand eller en Qvinde, han kan aldrig blive bedragen, thi det er kun de lavere Naturer, der indbilde sig, at de ere bedragne. Ingen Pige, der ikke er saa stolt, forstaaer egentlig at elske, men er hun saa stolt, da kan Alverdens List og Kløgt ikke bedrage hende.

I den uendelige Resignation er der Fred og Hvile; ethvert Menneske, der vil det, der ikke har fornedret sig selv ved, hvad der er endnu forfærdeligere end at være for stolt – at lade haant om sig selv, kan optugte sig selv til at gjøre denne Bevægelse, der i sin Smerte forsoner med Tilværelsen. Den uendelige Resignation er hiin Skjorte, om hvilken der tales i et gammelt Folkesagn. Traaden er spunden under Taarer, bleget ved Taarer, Skjorten syet i Taarer, men da beskytter den ogsaa bedre end Jern og Staal. Det Ufuldkomne i Folkesagnet er, at en Tredie kan forarbeide dette Linned. Hemmeligheden i Livet er, at Enhver maa sye sig den selv, og Mærkeligheden er, at en Mand kan sye den fuldt saa vel som en Qvinde. I den uendelige Resignation er der Fred og Hvile og Trøst i Smerten, det vil sige naar Bevægelsen er gjort normalt. Det var mig imidlertid ikke vanskeligt at skrive en heel Bog, naar jeg vilde gjennemgaae de forskjellige Misforstaaelser, de bagvendte Stillinger, de sjuskede Bevægelser, jeg har stødt paa blot i min lille Praxis. Man troer meget lidt paa Aand, og dog gjelder det netop om Aand for at gjøre denne Bevægelse, det gjelder om, at den ikke er et eensidigt Resultat af en dira necessitas, og jo mere denne er tilstede, desto tvivlsommere bliver det altid, om Bevægelsen er normal. Naar man saaledes vil mene, at den kolde, ufrugtbare Nødvendighed nødvendigt maa være tilstede, saa siger man dermed, at Ingen kan opleve Døden, før han virkelig døer, hvilket forekommer mig en cras Materialisme. Dog i vor Tid bekymrer man sig mindre om at gjøre rene Bevægelser. Dersom Een, der skulde lære at dandse, vilde sige: nu har i Aarhundreder den ene Slægt efter den anden lært Positioner, det er paa den høie Tid, at jeg drager mig dette til Fordeel og uden videre begynder paa Françaiser, saa vilde man vel lee lidt ad ham; men i Aandens Verden finder man det yderst plausibelt. Hvad er da Dannelse? Jeg troede, det var det Cursus, den Enkelte gjennemløb for at indhente sig selv; og den, der ikke vil gjennemgaae dette Cursus, ham hjælper det saare lidet, om han blev født i den mest oplyste Tidsalder.

Den uendelige Resignation er det sidste Stadium, der gaaer forud for Troen, saaledes, at Enhver, der ikke har gjort denne Bevægelse, ikke har Troen; thi først i den uendelige Resignation bliver jeg mig selv klar i min evige Gyldighed, og først da kan der være Tale om i Kraft af Troen at gribe Tilværelsen.

Vi ville nu lade Troens Ridder give Møde i det omtalte Tilfælde. Han gjør aldeles det samme, som den anden Ridder, han giver uendeligt Afkald paa den Kjærlighed, der er hans Livs Indhold, han er forsonet i Smerten; men da skeer Vidunderet, han gjør endnu en Bevægelse, forunderligere end Alt, thi han siger: jeg troer dog, at jeg faaer hende, i Kraft nemlig af det Absurde, i Kraft af, at for Gud er Alting muligt. Det Absurde hører ikke til de Differentser, der ligge indenfor Forstandens eget Omfang. Det er ikke identisk med det Usandsynlige, det Uventede, det Uformodede. I det Øieblik, Ridderen resignerede, da forvissede han sig om Umuligheden, menneskelig talt, dette var Forstandens Resultat, og han havde Energi nok til at tænke det. I uendelig Forstand var det derimod muligt, det er ved at resignere derpaa, men denne Besidden er jo tillige en Opgiven, men dog er denne Besidden for Forstanden ingen Absurditet; thi Forstanden vedblev at beholde Ret i, at i den Elendighedens Verden, hvor den hersker, var og blev det en Umulighed. Denne Bevidsthed har Troens Ridder ligesaa klar; det Eneste, der altsaa kan frelse ham, er det Absurde, og dette griber han ved Troen. Han erkjender altsaa Umuligheden og i samme Øieblik troer han det Absurde; thi vil han uden med al sin Sjæls Lidenskab og af sit ganske Hjerte at erkjende Umuligheden, indbilde sig at have Troen, da bedrager han sig selv, og hans Vidnesbyrd har intetsteds hjemme, da han end ikke er kommen til den uendelige Resignation.

Troen er derfor ingen æsthetisk Rørelse, men noget langt Høiere, netop fordi den har Resignationen forud for sig, den er ikke Hjertets umiddelbare Drift, men Tilværelsens Paradox. Naar saaledes en ung Pige tiltrods for alle Vanskeligheder dog holder sig forvisset om, at hendes Ønske vel bliver opfyldt, saa er denne Forvisning slet ikke Troens, og det uagtet hun er opdragen af christelige Forældre, og maaskee et heelt Aar har gaaet til Præsten. Hun er forvisset i al sin barnlige Naivetet og Uskyldighed, ogsaa denne Forvisning adler hendes Væsen, og giver hende en overnaturlig Størrelse, saa at hun som en Thaumaturg kan besværge Tilværelsens endelige Kræfter, og bringe Stene selv til at græde, medens hun paa den anden Side i sin Befippethed ligesaa godt kan løbe til Herodes som til Pilatus og røre al Verden med sine Bønner. Hendes Forvisning er saare elskelig, og man kan lære Meget af hende, men een Ting lærer man ikke af hende, man lærer ikke at gjøre Bevægelser; thi hendes Forvisning tør ikke i Resignationens Smerte see Umuligheden under Øine.

[…]

Til at resignere hører der ikke Tro, thi det, jeg i Resignation vinder, er min evige Bevidsthed, og dette er en reen philosophisk Bevægelse, som jeg trøster mig til at gjøre, naar forlanges, og som jeg kan tugte mig selv til at gjøre, thi hver Gang nogen Endelighed vil voxe mig overhovedet, da hungrer jeg mig selv ud, indtil jeg gjør Bevægelsen; thi min evige Bevidsthed er min Kjærlighed til Gud, og den er mig høiere end Alt. Til at resignere hører der ikke Tro, men til at faae det allermindste mere end min evige Bevidsthed hører der Tro, thi dette er det Paradoxe. Man forvexler ofte Bevægelserne. Man siger, at man behøver Troen for at give Afkald paa Alt, ja man hører det endnu Besynderligere, at et Menneske klager over, at han har tabt Troen, og naar man seer efter paa Scalaen, for at see, hvor han er, seer man besynderlig nok, at han kun er kommen til det Punkt, at han skal gjøre Resignationens uendelige Bevægelse. Ved Resignationen giver jeg Afkald paa Alt, denne Bevægelse gjør jeg ved mig selv, og naar jeg ikke gjør det, da er det fordi jeg er feig og blødagtig og uden Begeistring og ikke føler Betydningen af den høie Værdighed, ethvert Menneske er anviist, at være sin egen Censor, hvilket er meget fornemmere end at være General-Censor for hele den romerske Republik. Denne Bevægelse gjør jeg ved mig selv, og det, jeg derfor vinder, er mig selv i min evige Bevidsthed, i salig Forstaaelse med min Kjærlighed til det evige Væsen. Ved Troen giver jeg ikke Afkald paa Noget, tvertimod ved Troen faaer jeg Alt, netop i den Forstand, i hvilken det hedder, at den, som har Tro som et Sennepskorn, kan flytte Bjerge. Der hører et reent menneskeligt Mod til at give Afkald paa hele Timeligheden, for at vinde Evigheden, men denne vinder jeg, og kan i al Evighed ikke give Afkald paa, det er en Selvmodsigelse; men der hører et paradox og ydmygt Mod til nu at gribe hele Timeligheden i Kraft af det Absurde, og dette Mod er Troens. Ved Troen gav Abraham ikke Afkald paa Isaak, men ved Troen fik Abraham Isaak. I Kraft af Resignationen skulde hiin rige Yngling have bortgivet Alt, men naar han da havde gjort det, da skulde Troens Ridder have sagt til ham: I Kraft af det Absurde skal Du faae hver en Hvid igjen, kan Du troe det. Og denne Tale skal ingenlunde være den forhen rige Yngling ligegyldig; thi dersom han bortgav sit Gods, fordi han er kjed deraf, da var det daarligt bevendt med hans Resignation.

Timeligheden, Endeligheden er det, hvorom Alt dreier sig. Jeg kan ved egen Kraft resignere paa Alt, og da finde Fred og Hvile i Smerten, jeg kan finde mig i Alt, selv om hiin rædsomme Dæmon, forfærdeligere end Knokkelmanden, der forskrækker Menneskene, selv om Vanvidet holdt Narredragten op for mit Øie, og jeg forstod dens Mine, at det var mig, der skulde iføre mig den, jeg kan endnu frelse min Sjæl, dersom det ellers er mig mere om at gjøre, at min Kjærlighed til Gud seirer i mig, end min jordiske Lykke. Et Menneske kan endnu i dette sidste Øieblik samle hele sin Sjæl i eet eneste Blik mod den Himmel, fra hvilken al god Gave kommer, og dette Blik skal være forstaaeligt for ham og for den, hvem det søger, at han dog blev sin Kjærlighed tro. Da skal han rolig iføre sig Dragten. Den, hvis Sjæl ikke har denne Romantik, har solgt sin Sjæl, hvad enten han nu fik et Kongerige for den, eller en ussel Sølvpenge. Men ved egen Kraft kan jeg ikke faae det Mindste af hvad der tilhører Endeligheden; thi jeg bruger bestandig min Kraft til at resignere paa Alt. Ved egen Kraft kan jeg opgive Prindsessen, og jeg skal ikke blive nogen Suurmuler, men finde Glæde og Fred og Hvile i min Smerte, men ved egen Kraft kan jeg ikke faae hende igjen, thi jeg bruger netop min Kraft til at resignere. Men ved Troen, siger hiin vidunderlige Ridder, ved Troen skal Du faae hende i Kraft af det Absurde.

[…]

Den sidste Bevægelse, Troens paradoxe Bevægelse, kan jeg ikke gjøre, det være nu Pligt eller hvad det er, uagtet jeg hellere end gjerne vilde det. Om et Menneske har Lov til at sige dette, det maa overlades til ham; det bliver en Sag mellem ham og det evige Væsen, der er Troens Gjenstand, om han i denne Henseende kan træffe en mindelig Overeenskomst. Hvad ethvert Menneske kan, det er, han kan gjøre den uendelige Resignations Bevægelse, og jeg for mit Vedkommende vilde ikke tage i Betænkning, at erklære Enhver den for feig, der vil indbilde sig, at han ikke kan det. Med Troen er det en anden Sag. Men hvad ethvert Menneske ikke har Lov til, det er, at indbilde Andre, at Troen er noget Ringe eller at den er en let Sag, medens den er det Største og det Sværeste.

Man opfatter Fortællingen om Abraham paa en anden Maade. Man lovpriser Guds Naade, at han skjænkede ham Isaak igjen, det Hele var kun en Prøvelse. […]

\sectionheading{Problema I: Gives der en teleologisk Suspension af det Ethiske? (uddrag)  ·  SV¹ III, 104–115}

Det Ethiske er som saadant det Almene, og som det Almene Det, der er gjeldende for Enhver, hvilket fra en anden Side lader sig udtrykke saaledes, at det er gjeldende i ethvert Øieblik. Det hviler immanent i sig selv, har Intet uden for sig, der er dets τελος, men er selv τελος for Alt, hvad det har udenfor sig, og naar det Ethiske har optaget dette i sig, da kommer det ikke videre. Umiddelbar sandselig og sjælelig bestemmet er den Enkelte den Enkelte, der i det Almene har sit τελος, og dette er hans ethiske Opgave, bestandig at udtrykke sig selv i dette, at ophæve sin Enkelthed for at blive det Almene. Saasnart den Enkelte vil gjøre sig gjeldende i sin Enkelthed ligeoverfor det Almene, da synder han, og kan kun ved at anerkjende dette atter forsone sig med det Almene. Hver Gang den Enkelte, efterat have været traadt ind i det Almene, føler en Tilskyndelse til at gjøre sig gjeldende som den Enkelte, da er han i Anfægtelse, af hvilken han kun arbeider sig ud, ved angrende at opgive sig selv som den Enkelte i det Almene. Er dette det Høieste, der lader sig sige om Mennesket og om hans Tilværelse, saa har det Ethiske samme Beskaffenhed som et Menneskes evige Salighed, hvilken i al Evighed og i hvert Øieblik er hans τελος, da det vilde være en Modsigelse, at den skulde kunne opgives (det vil sige teleologisk suspenderes) da den, saasnart den suspenderes, forskjærtses, medens hvad der suspenderes ikke er forskjærtset, men netop bevaret i det Høiere, der er dets τελος.

[…]

Troen er nemlig dette Paradox, at den Enkelte er høiere end det Almene, dog vel at mærke saaledes, at Bevægelsen gjentager sig, at han altsaa, efterat have været i det Almene, nu som den Enkelte isolerer sig som høiere end det Almene. Dersom dette ikke er Troen, saa er Abraham tabt, saa har Troen aldrig været til i Verden, netop fordi den altid har været til. Thi hvis det Ethiske det vil sige det Sædelige er det Høieste, og der intet Incommensurabelt bliver tilbage i Mennesket paa anden Maade, end at dette Incommensurable er det Onde det vil sige det Enkelte, der skal udtrykkes i det Almene, saa behøver man ikke andre Kategorier end hvad den græske Philosophi havde, eller hvad der ved en consequent Tænkning lader sig uddrage af disse. Dette burde Hegel ikke have lagt Skjul paa; thi han har dog havt græske Studier.

[…]

Forskjellen mellem den tragiske Helt og Abraham er let iøinefaldende. Den tragiske Helt bliver endnu indenfor det Ethiske. Han lader et Udtryk af det Ethiske have sit τελος i et høiere Udtryk af det Ethiske, han nedsætter det ethiske Forhold mellem Fader og Søn eller Datter og Fader til en Følelse, der har sin Dialektik i sit Forhold til Sædelighedens Idee. Der kan da ikke her være Tale om en teleologisk Suspension af selve det Ethiske.

Med Abraham forholder det sig anderledes. Han overskred ved sin Gjerning hele det Ethiske, og havde et høiere τελος udenfor, i Forhold til hvilket han suspenderede dette. Thi jeg gad dog nok vide, hvorledes man vil bringe Abrahams Gjerning i Forhold til det Almene, om der lader sig opdage nogensomhelst anden Berøring mellem det, Abraham gjorde, og det Almene end den, at Abraham overtraadte det. Det er ikke for at frelse et Folk, ikke for at hævde Statens Idee, at Abraham gjør det, ikke for at forsone vrede Guder. Kunde der være Tale om, at Guddommen var vred, da var han jo kun vred paa Abraham, og Abrahams hele Gjerning staaer i intet Forhold til det Almene, er et reent privat Foretagende. Medens derfor den tragiske Helt er stor ved sin sædelige Dyd, er Abraham stor ved en reen personlig Dyd. Der er intet høiere Udtryk for det Ethiske i Abrahams Liv, end dette, at Faderen skal elske Sønnen. Det Ethiske i Betydning af det Sædelige kan der slet ikke være Tale om. Forsaavidt det Almene var tilstede, var det jo krypt i Isaak, skjult saa at sige i Isaaks Lænd, og maatte da raabe med Isaaks Mund: Gjør det ikke, Du tilintetgjør Alt.

Hvorfor gjør Abraham det da? For Guds Skyld og aldeles identisk hermed for sin egen Skyld. For Guds Skyld gjør han det, fordi Gud fordrer dette Bevis paa hans Tro, for sin egen Skyld gjør han det, at han kan føre Beviset. Eenheden heraf er aldeles rigtigt udtrykt i det Ord, hvormed man altid har betegnet dette Forhold: det er en Prøvelse, en Fristelse. En Fristelse; men hvad vil dette sige? Det, der ellers frister et Menneske, er jo det, der vil holde ham tilbage fra at gjøre sin Pligt, men her er Fristelsen selve det Ethiske, der vil holde ham tilbage fra at gjøre Guds Villie. Men hvad er da Pligten? Pligten er jo netop Udtrykket for Guds Villie.

Her viser Nødvendigheden sig af en ny Kategori for at forstaae Abraham. Et saadant Forhold til Guddommen kjender Hedenskabet ikke. Den tragiske Helt træder ikke i noget privat Forhold til Guddommen, men det Ethiske er det Guddommelige, og derfor lader det Paradoxe deri sig mediere i det Almene.

Abraham kan ikke medieres, hvilket ogsaa kan udtrykkes saaledes: han kan ikke tale. Saasnart jeg taler, udtrykker jeg det Almene, og naar jeg ikke gjør det, saa kan Ingen forstaae mig. Saasnart da Abraham vil udtrykke sig i det Almene, saa maa han sige, at hans Situation er en Anfægtelse, thi han har intet høiere Udtryk af det Almene, der staaer over det Almene, han overtræder.

[…]

Det er stort, naar Digteren, idet han stiller sin tragiske Helt frem for Menneskenes Beundring, at han da tør sige: Græder over ham, thi han fortjener det; thi det er Stort at fortjene deres Taarer, der fortjene at udgyde Taarer; det er Stort, at Digteren tør holde Mængden i Ave, tør tugte Menneskene, at Enhver prøver sig selv, om han er værdig til at græde over Helten, thi Flæbernes Spildevand er en Nedværdigelse af det Hellige. – Dog større end Alt dette er det, at Troens Ridder tør sige endog til det ædle Menneske, der vil græde over ham: græd ikke over mig, men græd over Dig selv.

Man bliver rørt, man søger tilbage i hine skjønne Tider, søde kjelne Længsler føre Een til Ønskets Maal, at see Christus vandre om i det forjættede Land. Man glemmer Angesten, Nøden, Paradoxet. Var det saa let en Sag ikke at tage feil? Var det ikke Forfærdeligt, at dette Menneske, der gik mellem de Andre, at han var Gud, var det ikke forfærdeligt at sidde tilbords med ham? Var det saa let en Sag at blive en Apostel? Men Udfaldet, 18 Aarhundreder, det hjælper, det hjælper til dette lumpne Bedrag, hvormed man bedrager sig selv og Andre. Jeg føler ikke Mod til at ønske at være Samtidig med saadanne Begivenheder, men derfor dømmer jeg ikke strængt om dem, der toge feil, ikke ringe om dem, der saae det Rette.

Dog jeg vender tilbage til Abraham. I den Tid før Udfaldet var Abraham enten i hvert Minut en Morder, eller vi staae ved det Paradox, der er høiere end alle Mediationer.

Abrahams Historie indeholder da en teleologisk Suspension af det Ethiske. Han er som den Enkelte bleven høiere end det Almene. Dette er Paradoxet, som ikke lader sig mediere. Det er ligesaa uforklarligt, hvorledes han kom ind i det, som det er uforklarligt, hvorledes han bliver i det. Forholder det sig ikke saaledes med Abraham, da er han end ikke en tragisk Helt, men en Morder. At ville vedblive at kalde ham Troens Fader, at tale derom til Mennesker, som ikke bekymre sig om Andet end om Ord, er tankeløst. En tragisk Helt kan et Menneske blive ved egne Kræfter, men Troens Ridder ikke. Naar et Menneske tiltræder den tragiske Helts i en vis Forstand tunge Vei, da skal Mange kunne raade ham; den, der gaaer Troens trange Vei, ham kan Ingen raade, Ingen forstaae. Troen er et Vidunder, og dog er intet Menneske udelukket derfra; thi det, hvori alt Menneskeliv enes, er i Lidenskab, og Troen er en Lidenskab.

\workheading{Gjentagelsen (af Constantin Constantius)}{1843}

\worknote{Antologiens ''Repetition''. Bogens mest sammensatte uddrag: den indledende bestemmelse af gjentagelse og erindring; Constantins beretning om Berlin-reisen (farce-afsnittet om Königstädter-Theatret er udeladt); hans betragtning over den unge Mand og Job; samt fire af den unge Mands breve (11. okt., 15. nov., 17. febr., 31. maj). Alle spring er markeret med […].}

\sectionheading{[Gjentagelse og Erindring]  ·  SV¹ III, 173–174}

Da Eleaterne negtede Bevægelsen, optraadte, hvad Enhver veed, Diogenes som Opponent; han optraadte virkelig; thi han sagde ikke et Ord; men gik blot nogle Gange frem og tilbage, hvorved han meente tilstrækkeligen at have modbeviist hine. Da jeg i længere Tid havde beskæftiget mig, leilighedsviis idetmindste, med det Problem, om en Gjentagelse er mulig og hvilken Betydning den har, om en Ting vinder eller taber ved at gjentages, faldt det mig pludselig ind: Du kan jo reise til Berlin, der har Du engang før været, og nu overbevise Dig om en Gjentagelse er mulig og hvad den har at betyde. I mit Hjem var jeg næsten gaaet istaae over dette Problem. Man sige hvad man vil derom, det vil komme til at spille en saare vigtig Rolle i den nyere Philosophi; thi Gjentagelse er et afgjørende Udtryk for hvad »Erindring« var hos Grækerne. Som da disse lærte, at al Erkjenden er en Erindren, saaledes vil den nye Philosophi lære, at hele Livet er en Gjentagelse. Den eneste nyere Philosoph, der har havt en Ahnelse herom, er Leibnitz. Gjentagelse og Erindring er den samme Bevægelse, kun i modsat Retning; thi hvad der erindres, har været, gjentages baglænds; hvorimod den egentlige Gjentagelse erindres forlænds. Derfor gjør Gjentagelsen, hvis den er mulig, et Menneske lykkeligt, medens Erindringen gjør ham ulykkeligt, under den Forudsætning nemlig, at han giver sig Tid til at leve, og ikke strax i sin Fødselsstund seer at finde paa et Paaskud til at liste sig ud igjen af Livet, for Exempel at han har glemt Noget.

\sectionheading{Constantin Constantius' beretning: reisen til Berlin (uddrag)  ·  SV¹ III, 189–211}

Hvad Betydningen angaaer, som Gjentagelsen har for en Ting, saa lader der sig sige Meget uden at gjøre sig skyldig i en Gjentagelse. Da Professor Ussing i sin Tid holdt en Tale i det 28de Mai-Selskab, og en Yttring i samme ikke behagede, hvad gjorde da Professoren, som dengang altid var resolut og gewaltig, han slog i Bordet og sagde: jeg gjentager. Han ventede altsaa dengang, at hvad han sagde, vandt ved Gjentagelsen. For nogle Aar siden hørte jeg en Præst ved to festlige Leiligheder holde aldeles den samme Tale. Havde han været af Professorens Mening, saa havde han den anden Gang, idet han besteg Prækestolen, slaaet i Pulpeten og sagt: jeg gjentager, hvad jeg allerede forrige Søndag sagde. Dette gjorde han ikke, og lod sig slet ikke mærke med Noget. Han var ikke af Professor Ussings Mening, og hvo veed, maaskee er Hr. Professoren selv ikke mere af den Mening, at det var godt for hans Tale, om den atter gjentoges. Da ved en Hoffest Dronningen havde fortalt en Historie, og alle Hofmændene leet ad den, en døv Minister inclusive, da reiste denne sig, og udbad sig den Naade ogsaa at maatte fortælle en Historie, og fortalte den samme. Spørgsmaal: hvilken Anskuelse havde han af Gjentagelsens Betydning. Naar en Lærer i en Skole siger: det er nu anden Gang, jeg gjentager, at Jespersen skal sidde roligt; og samme Jespersen faaer en Anmærkning for gjentagen Urolighed, saa er Gjentagelsens Betydning aldeles modsat.

Dog ved saadanne Exempler vil jeg ikke videre opholde mig, men gaae over til at tale lidt om den Opdagelses-Reise, jeg foretog mig for at forsøge Gjentagelsens Mulighed og Betydning. Uden noget Menneskes Vidende (for at nemlig al mulig Snakken ikke skulde gjøre mig udygtig til Experimentet og paa anden Maade kjed af Gjentagelsen) gik jeg med Dampskibet til Stralsund og tog en Plads i Schnellposten til Berlin. Der er blandt de Lærde forskjellige Meninger om, hvilken Plads der er den beqvemmeste i en Diligence; min Ansicht er følgende: det er Elendighed tilhobe. Forrige Gang havde jeg en Yderplads forlænds inden i Vognen (denne ansees af Nogle for den store Gevinst) og blev nu i 36 Timer saaledes sammenrystet med mine Nær-Paarørende, at jeg, da jeg ankom til Hamburg, ikke blot havde tabt Forstanden, men Benene med. Vi 6 Personer, som sadde inde i Vognen, arbeidedes i de 36 Timer saaledes sammen til eet Legeme, at jeg fik en Forestilling om, hvad der hændte Molboerne, der efter at have siddet længe sammen, ikke kunde kjende deres egne Been. For om muligt idetmindste at blive Lem paa et mindre Legeme, valgte jeg en Plads i Coupeen. Det var en Forandring. Imidlertid gjentog Alt sig. Postillonen blæste, jeg lukkede mine Øine, hengav mig i Fortvivlelsen, og tænkte som jeg ved saadan Leilighed pleier: Gud veed om Du holder det ud, om Du virkelig kommer til Berlin og i saa Fald om Du nogensinde bliver Menneske mere, istand til at frigjøre Dig i Isolationens Enkelthed, eller Du beholder et Minde om, at Du er et Lem paa et større Legeme.

Til Berlin ankom jeg da. Strax ilede jeg til mit gamle Logis for at forvisse mig om, hvorvidt en Gjentagelse var mulig. Jeg tør forsikkre enhver deeltagende Læser, at det var lykkedes mig den forrige Gang at faae en af de behageligste Leiligheder i Berlin, det tør jeg nu end bestemtere forsikkre, da jeg har seet flere. Gensd'arme-Pladsen er vel den smukkeste i Berlin, das Schauspielhaus, de to Kirker tage sig ypperligt ud især ved Maaneskin, sete fra et Vindue. Erindringen herom bidrog meget til, at jeg kom afsted. Man stiger op i første Etage i et med Gas oplyst Huus, man aabner en lille Dør, man staaer i Entreen. Til Venstre har man en Glasdør, der fører ind til et Cabinet. Man gaaer ligefrem, man er i et Forværelse. Indenfor dette er tvende Værelser, aldeles eensformede, aldeles eens møblerede, saaledes, at man i Speilet seer Værelset dobbelt. Det inderste Værelse er smagfuldt oplyst. En Armstage staaer paa et Arbeidsbord, en let tegnet Lænestol, betrukken med rødt Fløiel, staaer ved dette. Det forreste Værelse er ikke oplyst. Her blander Maanens blege Lys sig med den stærkere Belysning fra det indre Værelse. Man sætter sig hen paa en Stol ved Vinduet, man betragter den store Plads, man seer de Forbigaaendes Skygger ile hen over Murene, Alt forvandler sig til en scenisk Decoration. En drømmende Virkelighed dæmrer i Sjælens Baggrund. Man føler en Lyst til at kaste Kappen om sig, at snige sig sagte langs Muren med speidende Blikke, opmærksom paa enhver Lyd. Man gjør det ikke, man seer blot sig selv forynget gjøre det. Man har røget sin Cigar; man trækker sig tilbage i det indre Værelse, man begynder at arbeide. Midnat er forbi. Man slukker Lysene, man tænder en lille Natstage. Maanebelysningen seirer ublandet. En enkelt Skygge viser sig end mørkere, et enkelt Fodtrin behøver lang Tid for at forsvinde. Himlens skyfri Hvælving seer saa vemodig og tankedrømmende ud, som var Verdens Undergang forbi, og Himlen uforstyrret beskæftiget med sig selv. Man gaaer atter ud i Forværelset, i Entreen, ind i hiint lille Cabinet, man sover – hvis man hører til de Lykkelige, der kunne sove. – Ak! men her var ingen Gjentagelse mulig. Min Vert, Materialisten, er hatte sich verändert, i den prægnante Forstand, hvori Tydsken tager dette Ord, og, saavidt jeg veed, bruges »at forandre sig« paa en lignende Maade i enkelte af Kjøbenhavns Gader – det vil sige han havde giftet sig. Jeg vilde ønske ham til Lykke, men da jeg ikke er det tydske Sprog saa mægtig, at jeg veed at vende mig i en snæver Vending, ei heller havde de ved saadan Leilighed brugelige Talemaader promte, saa indskrænkede jeg mig til en pantomimisk Bevægelse. Jeg lagde Haanden paa Hjertet og saae paa ham, medens øm Deeltagelse stod læselig i mit Ansigt. Han trykkede min Haand. Efterat vi saaledes havde forstaaet hinanden, gik han over til at bevise Ægteskabets æsthetiske Gyldighed. Dette lykkedes ham overordentlig, netop ligesaa godt som den forrige Gang, at bevise Pebersvendenes Fuldkommenhed. Naar jeg taler tydsk, er jeg det føieligste Menneske af Verden.

Min forrige Vert vilde dog gjerne tjene mig, og jeg vilde gjerne boe hos ham; altsaa, jeg tog mod eet Værelse og Entreen. Da jeg kom hjem den første Aften, da jeg havde tændt Lysene, tænkte jeg: Ak! Ak! Ak! er det Gjentagelsen. Jeg blev ganske forstemt, eller om man saa vil, stemt netop som Dagen fordrede det; thi Skjebnen havde maget det saa besynderligt, at jeg ankom til Berlin paa den allgemeine Buß- und Bettag. Berlin var sønderknust. Vel kastede man ikke hinanden Støv i Øinene med de Ord: memento o homo! quod cinis es et in cinerem revertaris; men desuagtet stod den hele By i eet Støv. Jeg troede først, at det var en Regjerings-Foranstaltning; men senere overbeviste jeg mig om, at det var Vinden, der paatog sig denne Uleilighed, og uden Persons Anseelse fulgte sit Lune eller sin onde Vane; thi i Berlin er det idetmindste hver anden Dag Aske-Onsdag. Dog dette vedkommer mindre mit Forehavende. Denne Opdagelse betraf ikke »Gjentagelsen«; thi den forrige Gang jeg var i Berlin, havde jeg ikke bemærket dette Phænomen, formodentlig fordi det var Vinter.

Naar man har faaet sig bequemt og hyggeligt indrettet i sin Bolig, naar man saaledes har et fast Punkt, hvorfra man kan styrte frem, et betrygget Skjulested, hvor man kan trække sig tilbage, for i Eensomhed at fortære sit Bytte, Noget jeg især sætter Priis paa, da jeg som visse Rovdyr ikke kan spise, naar Nogen seer derpaa – saa gjør man sig bekjendt med hvad der kan være af Mærkeligheder i Byen. Er man en Reisende ex professo, et Iilbud, der reiser for at lugte til Alt, hvad Andre have lugtet til, eller for at skrive Navnene paa Mærkværdighederne i sin Dagbog, og til Gjengjeld sit i de Reisendes store Stambog, saa antager man en Lohndiener og kjøber das ganze Berlin for 4 Groschen. Ved denne Methode bliver man en upartisk Iagttager, hvis Udsagn bør staae til Troende ved enhver Politi-Protocol. Er man derimod paa sin Reise ikke i nødvendigt Ærinde, saa lader man staae til; faaer stundom Noget at see, som Andre ikke see; overseer det Vigtigste; faaer et tilfældigt Indtryk, der kun har Betydning for En selv. En saadan sorgløs Løsgænger har da i Almindelighed ikke stort at meddele Andre og gjør han det, løber han let Fare for at svække de gode Tanker, som gode Mennesker kunde have om hans Moralitet og Sædelighed. Hvis et Menneske havde reist i længere Tid i Udlandet, men aldrig været auf der Eisenbahn, skulde han da ikke være udstødt af ethvert bedre Selskab! Hvad om et Menneske havde været i London og aldrig kjørt i Tunnel! Hvad om et Menneske kom til Rom, forelskede sig i et lille Parti af Byen, der var ham et uudtømmeligt Stof til Glæde, forlod Rom uden at have seet en eneste Mærkværdighed! Berlin har 3 Theatrer. Hvad der i Operaen præsteres for Opera og Ballet, skal være großartig; hvad der i Skuespilhuset præsteres, skal være belærende, dannende, ikke blot til Lyst. Jeg veed det ikke. Derimod veed jeg, at der i Berlin er et Theater, som hedder Königstädter Theatret. Dette Theater besøge de officielle Reisende sjeldnere, skjøndt dog lidt oftere end, hvad der ogsaa har sin Betydning, de gemytlige Forlystelsessteder, der ligge mere afsides, og hvor en Dansk kan faae Leilighed til at opfriske Mindet om Lars Mathiesen og Kehlet. Da jeg kom til Stralsund og læste i Avisen, at »der Talismann« skulde opføres paa hiint Theater, blev jeg strax veltilmode. Erindringen om det vaagnede i min Sjæl, og den første Gang, jeg var der, var det som fremkaldte dette første Indtryk kun en Erindring i min Sjæl, der viste langt tilbage i Tiden.

[…]

Jeg ilede ud i Theatret. Der var ingen Loge at faae alene, ikke engang Pladsen i Nummer 5 og 6 links. Jeg maatte gaae til Høire. Der traf jeg et Selskab, som ikke med Bestemthed vidste, om det skulde more eller ennuere sig, et saadant Compagni kan man med Bestemthed ansee for kjedsommeligt. Der var næsten ikke en eneste tom Loge. Den unge Pige var ikke til at opdage, eller om hun har været der, har jeg ikke kunnet kjende hende, fordi hun var i Selskab. Beckmann kunde ikke bringe mig til at lee. I en halv Time holdt jeg ud, da forlod jeg Theatret og tænkte: der er slet ingen Gjentagelse til. Dette gjorde et dybt Indtryk paa mig. Jeg er ikke saa ganske ung, ikke aldeles ubekjendt med Livet, havde allerede længe førend jeg forrige Gang kom til Berlin, afvænt mig med at calculere paa det Uvisse. Imidlertid troede jeg dog, at den Nydelse, jeg havde havt i hiint Theater, skulde være af en durablere Art, netop fordi man maatte have lært at lade sig paa mange Maader reducere af Tilværelsen og dog hjælpe sig, før man egentlig fik Sands for den; men desto sikkrere maatte den vel ogsaa være. Skulde Tilværelsen være endnu svigefuldere end en Fallent! Han giver dog 50 Procent, eller 30, Noget giver han dog. Det Comiske er dog det Mindste, man kan forlange, skulde det heller ikke lade sig gjentage!

Med disse Tanker gik jeg hjem. Mit Arbeidsbord var sat frem. Den Fløielslænestol existerede endnu. Men da jeg saae den, blev jeg saa forbittret, at jeg nær havde slaaet den i Stykker, saa meget mere som Alle i Huset vare gaaede i Seng, og Ingen kunde tage den bort. Hvad kan en Fløiels Lænestol hjælpe, naar den øvrige Omgivelse ikke svarer dertil, det er jo som om en Mand vilde gaae nøgen med Trekantethat paa. Da jeg var gaaet i Seng, uden at have havt en eneste fornuftig Tanke, var det saa lyst i Værelset, at jeg i eet væk saae den Fløielslænestol, deels vaagen, deels i Drømme, indtil jeg om Morgenen stod op og gjorde Alvor af hvad jeg havde besluttet, at lade den henkaste i en Afkrog.

Mit Hjem var blevet mig uhyggeligt, netop fordi det var en forkeert Gjentagelse, min Tanke var ufrugtbar, min bekymrede Phantasi fremtryllede bestandig til tantalisk Forlystelse Erindringer om, hvorledes Tankerne forrige Gang tilbøde sig, og denne Erindringens Klinte qvalte enhver Tanke i Fødselen. Jeg gik ud, hen til den Conditor, hos hvem jeg den forrige Gang kom hver Dag, for at nyde den Drik, der, naar den efter Digterens Forskrift, er »reen og varm og stærk og umisbrugt«, altid kan staae ved Siden af det, hvormed Digteren sammenligner den – Venskabet; jeg holder idetmindste paa Kaffeen. Maaskee var Kaffeen ligesaa god som forrige Gang, man skulde næsten troe det, mig smagte den ikke. Solen brændte gloende paa Boutiquens Vinduer, Localet var qvalmt, omtrent som Luften i en Casserolle, ganske til at henkoges i; en Trækvind, der som en smal Passat gjennemborede Alt, forbød mig at tænke paa nogen Gjentagelse, selv om der ellers havde tilbudet sig Leilighed.

Om Aftenen gik jeg i den Restauration, hvor jeg forrige Gang pleiede at komme, og formodentlig ved Vanens Hjælp endog havde befundet mig vel. Da jeg kom der hver Aften, kjendte jeg Alt paa det Nøieste, jeg vidste, naar de tidige Gjester gik derfra, hvorledes de hilste paa Broderskabet, som de forlod, om de satte Hatten paa i det inderste Værelse, eller i det yderste, eller først idet de aabnede Døren, eller først udenfor Døren. Ingen undgik min Opmærksomhed, jeg rykkede som Proserpina et Haar af hvert Hoved, endog af de skaldede. – Det var aldeles det Samme, de samme Vittigheder, de samme Høfligheder, den samme Deeltagelse, Localet aldeles det samme – kort det Samme i det Samme. Salomon siger, at en Qvindes Trætte er som Tagdryp, hvad mon han vilde sige om dette Stillleben. Forfærdelige Tanke, her var en Gjentagelse mulig.

Den næste Aften var jeg i Königstädter-Theatret. Det Eneste, der gjentog sig, var Umuligheden af en Gjentagelse. Unter den Linden var Støvet utaaleligt, ethvert Forsøg paa at blande sig ind mellem Menneskene, og saaledes tage et Menneskebad, virkede høist afskrækkende. Hvorledes jeg end vendte og dreiede mig, var Alt forgjeves. Den lille Dandserinde, som den forrige Gang havde fortryllet mig ved en Ynde, der saa at sige stod paa Springet, hun havde gjort Springet; den blinde Mand udenfor Brandenburger-Thor, min Harpespiller – thi jeg var nok den eneste, der bekymrede sig om ham – havde faaet en graameleret Frakke istedenfor den lysegrønne, der var min Vemods Længsel, i hvilken han saae ud som en Taarepiil, han var tabt for mig og vunden for det Almeen-Menneskelige; Pedellens beundrede Næse var bleven bleg; Professor A. A. havde faaet et Par nye Buxer, der næsten sad militairisk. – –

Da dette havde gjentaget sig nogle Dage, blev jeg saa forbittret, saa kjed af Gjentagelsen, at jeg besluttede at reise hjem igjen. Min Opdagelse var ikke betydelig og dog var den besynderlig; thi jeg havde opdaget, at Gjentagelsen slet ikke var til, og det havde jeg forvisset mig om, ved paa alle mulige Maader at faae det gjentaget.

Mit Haab stod til mit Hjem. Justinus Kerner fortæller etsteds om en Mand, der blev kjed af sit Hjem, at han lod sin Hest sadle for at ride ud i den vide Verden. Da han havde redet et lille Stykke, slog Hesten ham af. Denne Vending blev afgjørende for ham; thi idet han vendte sig om for at bestige sin Hest, faldt hans Øie endnu engang paa det Hjem, han vilde forlade, og han saae og see! det var saa skjønt, at han strax vendte tilbage. I mit Hjem turde jeg regne temmelig sikker paa at finde Alt færdigt til Gjentagelsen. Jeg har altid havt stor Mistillid til alle Omvæltninger, ja det gaaer saavidt, at jeg af den Grund endog hader al mulig Reengjøren, og fremfor Alt den huslige Sæbenaf. Jeg havde da efterladt de strængeste Instruxer, for at see mine conservative Principer opretholdte ogsaa i min Fraværelse. Dog hvad skeer. Min tro Tjener var af en anden Mening. Han regnede paa, at naar han ganske kort efter min Afreise begyndte denne Rystelse, skulde den vel være ophørt til min Hjemkomst, og han vel nok være Manden for at bringe Alt paa det punktligste i Orden. Jeg kommer, jeg ringer paa min Dør, min Tjener aabner. Det var et betydningsfuldt Øieblik. Min Tjener blev som et Liig, gjennem den halvaabnede Dør til Værelserne saae jeg det Forfærdelige: der var vendt op og ned paa Alt. Jeg forstenede. I sin Befippelse vidste han ikke, hvad han skulde gjøre, hans onde Samvittighed slog ham og – han slog Døren i igjen for mig. Dette var for meget, min Nød havde naaet sit Høieste, mine Principer segnede, jeg maatte befrygte det værste, at blive ligesom Commerceraad Grønmeyer behandlet som et Gjenfærd. Jeg indsaae, at der ingen Gjentagelse er til, og min tidligere Betragtning af Livet havde seiret.

[…]

Jo ældre man bliver, jo mere Forstand man faaer paa Livet og Smag for det Behagelige og Evne til at goutere, kort, jo competentere man bliver, desto mindre tilfreds. Tilfreds, aldeles, absolut og i alle Maader tilfreds bliver man aldrig, og at være nogenlunde tilfreds, er ikke Umagen værd, saa er det bedre at være aldeles utilfreds. Enhver, der grundigen har overveiet den Sag, vil vist give mig Ret i, at det aldrig forundes et Menneske, end ikke saa meget som en halv Time i hele hans Liv, at være absolut tilfreds i alle tænkelige Maader. At dertil nemlig fordres noget Mere end at have Føde og Klæder, behøver jeg vel ikke at sige. Jeg har engang været nær efter det. Jeg stod op en Morgen og befandt mig ualmindelig vel; dette Velbefindende tiltog mod al Analogi op ad Formiddagen, præcise Klokken 1 var jeg paa det Høieste og ahnede det svimlende Maximum, der ikke findes anført paa nogen Velværens Gradestok, ikke engang paa et poetisk Thermometer. Legemet havde tabt sin jordiske Tyngde, det var som havde jeg intet Legeme, netop fordi enhver Funktion nød hele sin Tilfredsstillelse, enhver Nerve frydede sig paa egne og det Heles Vegne, medens ethvert Pulsslag som Organismens Uro kun mindede om og angav Øieblikkets Vellyst. Min Gang var svævende, ikke som Fuglens Flugt, der gjennemskærer Luften og forlader Jorden, men som Vindens Bølgen over Sæden, som Havets længselssalige Vuggen, som Skyens drømmende Hengliden. Mit Væsen var Gjennemsigtighed som Havets dybe Grunden, som Nattens selvtilfredse Taushed, som Middagens monologiske Stilhed. Enhver Stemning hvilede i min Sjæl med melodisk Resonants. Enhver Tanke tilbød sig, og enhver Tanke tilbød sig med Salighedens Festlighed, det taabeligste Indfald ikke mindre end den rigeste Idee. Ethvert Indtryk var ahnet før det kom, og vaagnede derfor i mig selv. Hele Tilværelsen var som forelsket i mig, og Alt bævede i skjæbnesvangert Rapport til min Væren, Alt var omineust i mig, og Alt gaadefuldt forklaret i min mikrokosmiske Salighed, der forklarede Alt i sig, selv det Ubehagelige, den kjedsommeligste Bemærkning, Synet af det Modbydelige, det fataleste Sammenstød. Som sagt præcise Klokken 1 var jeg paa det Høieste, hvor jeg ahnede det Allerhøieste, da begynder pludselig noget at gnave i mit ene Øie, om det var et Øienhaar, et Fnug, et Støvgran, jeg veed det ikke, men dette veed jeg, at jeg i samme Øieblik styrtede ned næsten i Fortvivlelsens Afgrund, Noget, Enhver let vil forstaae, der har været saa høit oppe som jeg, og idet han var paa dette Punkt tillige har været beskæftiget med det Principspørgsmaal, hvorvidt overhovedet den absolute Tilfredsstillelse er at opnaae. Siden den Tid opgav jeg ethvert Haab om nogensinde at befinde mig tilfreds absolut og i alle Maader, opgav det Haab, jeg engang havde næret, vel ikke til alle Tider at være absolut tilfreds, men dog i enkelte Øieblikke, om saa disse Eenheder af Øieblikke ikke bleve flere end at, som Shakespeare siger, »en Øltappers Regnekunst var tilstrækkelig for at opsummere dem.«

Saavidt var jeg allerede kommen før jeg lærte hiint unge Menneske at kjende. […]

\sectionheading{Constantin om den unge Mand og Gjentagelsen (uddrag)  ·  SV¹ III, 214–222}

Der gik nogen Tid hen; min Tjener havde som en huslig Eva gjort godt igjen, hvad han tidligere forskyldte. En monoton og eensformig Orden var tilveiebragt i min hele Oekonomi. Alt hvad der ikke kunde gaae, – det stod paa sit bestemte Sted, og hvad der kunde gaae, det gik sin beregnede Gang: mit Stueuhr, min Tjener og jeg selv, der med afmaalte Skridt gik op og ned ad Gulvet. Skjønt jeg nemlig havde forvisset mig om, at der ingen Gjentagelse er til, bliver det dog altid vist og sandt, at man ved Urokkelighed og tillige ved at sløve sin Iagttagelses-Evne kan opnaae en Eensformighed, der har en langt mere bedøvende Magt end de lunefuldeste Adspredelser, og som tillige i Tidens Løb bliver stærkere og stærkere ligesom en Besværgelsesformular. Ved Udgravningen af Herculanum og Pompeii fandt man Alt paa sin Plads, saaledes som de respektive Eiere forlode det; dersom jeg havde levet paa hiin Tid, da vilde Oldgranskerne maaskee med Forbauselse have truffet et Menneske, der gik med afmaalte Skridt op og ned ad Gulvet. For at vedligeholde denne bestaaende og bestandige Orden, anvendte jeg ethvert Middel, gik endog til visse Tider som Keiser Domitian om i Værelset bevæbnet med en Fluesmække, forfølgende enhver revolutionair Flue. Derimod blev der fredet om 3 Fluer, som til bestemte Tider fløi surrende gjennem Stuen. Saaledes levede jeg, som jeg troede, glemmende Verden og glemt, da jeg en Dag modtager et Brev fra min unge Ven. Det efterfulgtes af flere, altid med omtrent en Maaneds Mellemrum, uden at jeg dog deraf tør drage nogen Slutning med Hensyn til Afstanden af hans Opholdssted. Selv har han Intet ville oplyse, og det kunde vel være en Mystification, der da forsaavidt var forsigtigt gjennemført ved at lade Tiden fluctuere mellem næsten 5 Uger og kun 1 Dag over 3 Uger. Han ønsker ikke at besvære mig med en Correspondance, og selv om jeg var villig til at besvare eller dog svare paa hans Breve, ønsker han ikke at modtage noget saadant – han vil blot udgyde sig.

[…]

Naar han imidlertid troer, at jeg aldeles havde glemt ham, saa gjør han mig igjen Uret. Ved hans pludselige Forsvinden frygtede jeg virkelig for, at han i Fortvivlelse havde gjort en Ulykke paa sig selv. En saadan Hændelse pleier sjeldent at blive længe skjult, da jeg derfor Intet hørte eller læste, sluttede jeg, at han vel maatte være i Live, hvor han saa end stak. Pigen, som han lod i Stikken, vidste aldeles Intet. En Dag udeblev han og lod aldeles Intet høre fra sig. Hendes Overgang til Smerten var ikke pludselig; thi først successivt vaagnede den bange Ahnelse, og først successivt besindede Smerten sig paa sig selv, hvilket bevirkede, at hun blidelig henslumrede i en drømmende Uklarhed over hvad der var foregaaet og hvad det kunde have at betyde. Pigen var mig et nyt Stof for Iagttagelsen. Min Ven hørte ikke til dem, der veed at pine Alt ud af den Elskede og saa kaster hende bort; tvertimod befandtes hun ved hans Forsvinden i den ønskeligste Tilstand, sund, svulmende, beriget med hele hans poetiske Udbytte, kraftig næret ved den digteriske Illusions kostelige Hjertestyrkning. Det er et sjeldnere Tilfælde at træffe en forladt Pige i denne Tilstand. Da jeg saae hende nogle Dage efter, var hun endnu væver som en frisk fanget Fisk; ellers er en saadan Pige gjerne udhungret som en Fisk, der har levet i et Hyttefad. I min Samvittighed var jeg da overbeviist om, at han maatte være i Live, og var nu ret glad ved, at han ikke havde grebet til det fortvivlede Middel at udgive sig for død. Det er utroligt, saa megen Forvirring der kan komme ind i det Erotiske, naar det behager den ene Part at ville være død af Sorg eller at ville være død for at slippe fra det Hele. En Pige vilde ifølge sin egen høitidelige Erklæring døe af Sorg over, at hendes Elsker var en Bedrager. See! han var ingen Bedrager, og meente det maaskee meget bedre, end hun fattede. Men hvad han ellers maaskee vilde have gjort i Tidens Fylde, det kunde han nu ikke beslutte sig til, blot fordi hun eengang havde tilladt sig at ængste ham med hiin Forsikkring, fordi hun, som han sagde, havde anvendt et oratorisk Kneb mod ham, eller i ethvert Fald sagt, hvad en Pige aldrig bør sige, hvad enten hun nu troer, at han virkelig er en Bedrager; thi da bør hun være for stolt, eller hun har endnu Tro paa ham; thi da bør hun indsee, at hun kan gjøre ham himmelraabende Uret. At ville være død for at slippe fra det Hele er det ussleste Middel, der lader sig tænke, og indholder den mest krænkende Fornærmelse mod en Pige. Hun troer, han er død, hun anlægger Sorg, hun græder og begræder den Afdøde ærligt og oprigtigt. Hun maae jo næsten faae Modbydelighed for sine egne Følelser, naar hun engang senere opdager, at han lever og mindst af Alt har været betænkt paa Døden. Eller hvis hun først i et andet Liv fatter en Mistanke, ikke om han virkelig er død; thi det bliver jo uomtvisteligt, men om han er død dengang han sagde og hun sørgede. En saadan Situation vilde være en Opgave for en Apokalyptiker, der har forstaaet sin Aristophanes (jeg mener den græske, ikke de enkelte Mænd, der udnævnes ligesom doctores cerei i Middelalderen) og sin Lucian. Man kunde en tidlang vedligeholde Forvexlingen; thi død var han jo og død blev han. Den sørgende Pige vilde da vaagne til at begynde hvor de slap, indtil hun opdagede, at der var en lille Mellemsætning.

[…]

Det Problem, ved hvilket han er standset, er hverken mere eller mindre end Gjentagelsen. At han ikke søger Oplysning derom i den græske Philosophi, heller ikke i den nyere, deri gjør han Ret; thi Grækerne gjøre den modsatte Bevægelse, og en Græker vilde her vælge at erindre, uden at hans Samvittighed vilde ængste ham; den nyere Philosophi gjør ingen Bevægelse, den gjør i Almindelighed kun Ophævelse, og forsaavidt den gjør en Bevægelse, ligger denne altid i Immanentsen, Gjentagelsen derimod er og bliver en Transcendents. En Lykke er det, at han ikke søger nogen Forklaring hos mig; thi jeg har opgivet min Theori, jeg driver. Gjentagelsen er ogsaa mig for transcendent. Jeg kan omseile mig selv; men jeg kan ikke komme ud over mig selv, dette archimediske Punkt kan jeg ikke opdage. Min Ven søger da heldigviis ikke Oplysning hos nogen verdensberømt Philosoph eller hos nogen professor publicus ordinarius; han tyer til en privatiserende Tænker, der engang eiede Verdens Herlighed, men senere trak sig tilbage fra Livet – med andre Ord han tager sin Tilflugt til Job, der ikke figurerer paa et Katheder og ved betryggende Gesticulationer indestaaer for sine Sætningers Sandhed, men som sidder paa Arnestedet og skraber sig med Potteskaar, og uden at afbryde denne Haandgjerning henkaster flygtige Vink og Bemærkninger. Her mener han at have fundet hvad han søgte, i denne lille Kreds af Job og Kone samt 3 Venner lyder, efter hans Mening, Sandheden herligere og glædeligere og sandere end i et græsk Symposium.

Selv om han endnu vilde søge min Veiledning, det er forgjeves. En religiøs Bevægelse kan jeg ikke gjøre, det er min Natur imod. Derfor negter jeg dog ikke Realiteten deraf, eller at man kan lære saare meget af et ungt Menneske. Lykkes den for ham, da er han fri for alt det Pirrelige i sit Forhold til mig. Kun kan jeg ikke negte, at jeg, jo mere jeg betragter Sagen, paa ny fatter en Mistanke til Pigen, at hun paa een eller anden Maade har tilladt sig at ville fange ham i hans Melancholi. I saa Fald vilde jeg ikke være i hendes Sted. Det gaaer galt. Tilværelsen hævner altid en saadan Adfærd paa det strængeste.

\sectionheading{Den unge Mands breve (udvalg)  ·  SV¹ III, 234–255}

Den 11. October.

Min tause Medvider!

Mit Liv er bragt til det Yderste; jeg væmmes ved Tilværelsen, den er smagløs uden Salt og Mening. Om jeg var hungrigere end Pierrot, gad jeg dog ikke æde den Forklaring, som Menneskene byde. Man stikker Fingeren i Jorden, for at lugte, hvad Land man er i, jeg stikker Fingeren i Tilværelsen – den lugter af ingen Ting. Hvor er jeg? Hvad vil det sige: Verden? Hvad betyder dette Ord? Hvo har narret mig ind i det Hele, og lader mig nu staae der? Hvem er jeg? Hvorledes kom jeg ind i Verden; hvorfor blev jeg ikke adspurgt, hvorfor ikke gjort bekjendt med Skikke og Vedtægter, men stukket ind i Geleddet, som var jeg kjøbt af en Seelenverkoper? Hvorledes blev jeg Interessent i den store Entreprise, som man kalder Virkeligheden? Hvorfor skal jeg være Interessent? Er det ikke en fri Sag? Og skal jeg nødsages til at være det, hvor er da Dirigenten, jeg har en Bemærkning at gjøre? Er der ingen Dirigent? Hvor skal jeg henvende mig med min Klage? Tilværelsen er jo en Debat, maa jeg bede, at min Betragtning kommer med under Overveielse? Skal man tage Tilværelsen som den er, var det saa ikke bedst, at man fik at vide, hvorledes den er? Hvad vil det sige: en Bedrager? Siger ikke Cicero, at man udfinder en saadan ved at spørge: cui bono? Jeg lader Enhver spørge, og spørger Enhver, om jeg har havt nogen Gavn af, at gjøre mig selv og en Pige ulykkelig. Skyld – hvad vil det sige? Er det Hexeri? Vides det ikke med Bestemthed, hvorledes det gaaer til, at et Menneske bliver skyldig? Vil der Ingen svare? Er det da ikke af yderste Vigtighed for alle DHerrer Deeltagere?

[…]

Den 15. November.

Min tause Medvider!

Dersom jeg ikke havde Job! Det er umuligt at beskrive og nuancere, hvilken Betydning og hvor mangfoldig en Betydning han har for mig. Jeg læser ham ikke som man læser en anden Bog med Øiet, men jeg lægger Bogen ligesom paa mit Hjerte, og med Hjertets Øie læser jeg den, forstaaer i en clairvoyance det Enkelte paa den forskjelligste Maade. Som Barnet lægger Lærebogen under sit Hoved, for at være sikker paa, at han ikke har glemt sin Lectie, naar han vaagner om Morgenen, saaledes tager jeg Bogen med mig om Natten i Sengen. Hvert Ord af ham er Føde og Klæde og Lægedom for min elendige Sjæl. Snart vækker et Ord af ham mig af min Lethargi, saa jeg vaagner til ny Uro, snart stiller det den ufrugtbare Rasen i mit Indre, ender det Rædsomme i Lidenskabens stumme Qvalmen. De har dog læst Job? Læs den, læs den atter og atter. Jeg nænner ikke engang at skrive et eneste Udbrud af i et Brev til Dem, skjøndt jeg har min Glæde af at tage ny og ny Afskrift af Alt, hvad han har sagt, snart med danske, snart med latinske Bogstaver, snart i een Format snart i en anden. Enhver saadan Afskrift bliver lagt som et Gudshaandsplaster paa mit syge Hjerte. Og paa hvem laae vel Guds Haand som paa Job! Men citere ham – det kan jeg ikke. Det var at ville give min Besyv med, det var at ville gjøre hans Ord til mine i en Andens Nærværelse. Naar jeg er alene, da gjør jeg det, tilegner mig Alt; men saasnart Nogen er tilstede, da veed jeg vel, hvad et ungt Menneske har at gjøre, naar gamle Folk tale.

[…]

Den 17. Februar.

Min tause Medvider!

Her sidder jeg. Paa Uskyldigheden? som det hedder i Tyvesproget; eller paa Kongens Naade? Jeg veed det ikke; kun det veed jeg, at jeg sidder, og at jeg ikke rører mig af Stedet. Her staaer jeg, paa Spidsen eller ved Foden? det veed jeg ikke; kun det veed jeg, at jeg staaer og staaer suspenso gradu nu i en heel Maaned, uden at trække Foden til mig eller gjøre en eneste Bevægelse.

Jeg venter paa et Tordenveir – og paa Gjentagelsen. Og dog, naar blot Tordenveiret kommer, allerede dermed er jeg glad og ubeskrivelig salig, selv om min Dom blev, at ingen Gjentagelse var mulig.

Hvad skal dette Tordenveir bevirke? Det skal gjøre mig duelig til at være en Ægtemand. Det vil knuse hele min Personlighed, jeg er færdig, det vil gjøre mig næsten ukjendelig for mig selv, jeg vakler ikke, skjønt jeg staaer paa eet Been. Min Ære er reddet, min Stolthed indløst; og hvorledes det end vil omskabe mig, jeg haaber dog, at Erindringen herom som en uudtømmelig Trøst vil blive hos mig, blive, naar det er skeet, jeg i en vis Forstand frygter mere end et Selvmord, fordi det efter en ganske anden Maalestok vil forstyrre mig. Kommer Tordenveiret ikke, saa bliver jeg lumsk; saa døer jeg slet ikke, men lader engang som var jeg død, for at Slægt og Venner kunne begrave mig. Naar man da lægger mig i Kisten, da putter jeg i al Stilhed min Forventning til mig. Det faaer Ingen at vide, thi ellers vilde man vel vogte sig for at begrave et Menneske, der endnu var Liv i.

Forøvrigt gjør jeg Alt hvad der staaer i min Magt for at danne mig til en Ægtemand. Jeg sidder og beklipper mig selv, tager alt det Incommensurable bort, for at blive commensurabel. […]

[…]

Den 31. Mai.

Min tause Medvider.

Hun er gift; med hvem veed jeg ikke; thi da jeg læste det i Bladet, blev jeg som rørt af et Slag og tabte Avisen, og har siden ikke havt Taalmodighed til et nærmere Eftersyn. Jeg er atter mig selv; her har jeg Gjentagelsen; jeg forstaaer Alt, og Tilværelsen forekommer mig skjønnere end nogensinde. Det kom jo ogsaa som et Tordenveir, om jeg end skylder hendes Høimodighed, at det skete. Hvo det end er, hun har valgt, jeg vil end ikke sige: foretrukket; thi i Qvalitet af Ægtemand er Enhver at foretrække for mig, hun har dog viist Høimodighed imod mig. Om han var den skjønneste Mand i Verden, et Indbegreb af al Elskværdighed, istand til at fortrylle enhver Pige, om hun kunde bringe hele Kjønnet til Fortvivlelse ved at give ham sit Ja, hun har dog handlet høimodigt, om ikke ved andet saa ved aldeles at have glemt mig. Hvad er dog skjønt som qvindelig Høimodighed. Lad den jordiske Skjønhed visne, lad hendes Øies Glands udslukkes, lad hendes ranke Skud bøie sig i Aarene, lad Lokkerne tabe den besnærende Magt, naar de skjules af den ydmyge Kappe, lad hendes kongelige Blik, der beherskede Verden, med moderlig Kjærlighed kun omslutte og vogte paa den Kreds, om hvilken hun freder – en Pige, der saaledes har været høimodig, ældes aldrig. Lad Tilværelsen lønne hende som den har gjort det, lad den give hende hvad hun elskede mere, den gav ogsaa mig hvad jeg elskede mere – mig selv, og gav mig dette ved hendes Høimod.

Jeg er atter mig selv. Dette »Selv,« som en Anden ikke vilde tage op paa Landeveien, eier jeg igjen. Den Splid, der var i mit Væsen, er hævet; jeg slutter mig atter sammen. Sympathiens Angester, der fandt Medhold og Næring i min Stolthed, trænge sig ikke mere ind til at adsplitte og separere.

Er der da ikke en Gjentagelse? Fik jeg ikke Alt dobbelt? Fik jeg ikke mig selv igjen, netop saaledes, at jeg dobbelt maatte føle Betydningen deraf? Og hvad er en Gjentagelse af jordisk Gods, der er ligegyldigt mod Aandens Bestemmelse, i Sammenligning med en saadan Gjentagelse? Kun Børnene fik Job ikke dobbelt, fordi et Menneskeliv ikke saaledes lader sig fordoble. Her er kun Aandens Gjentagelse mulig, om den end i Timeligheden aldrig bliver saa fuldkommen som i Evigheden, der er den sande Gjentagelse.

Jeg er atter mig selv; Maskineriet er sat i Bevægelse. Sønderhugne ere de Besnærelser, hvori jeg var hildet; brudt er den Trylleformel, der havde hexet sig ind i mig, saa jeg ikke kunde komme tilbage til mig selv. Der er Ingen mere, der løfter sine Hænder mod mig, min Frigjørelse er sikker, jeg er født til mig selv; thi saalænge Ilithyia folder sine Hænder, kan den Fødende ikke føde.

Det er forbi, min Jolle er flot, i næste Minut er jeg atter der, hvor min Sjæls Higen var, der, hvor Ideerne bruse med elementarisk Rasen, hvor Tankerne staae larmende op som Nationerne i Folkevandringen, der, hvor der til en anden Tid er en Stilhed, som Sydhavets dybe Taushed, en Stilhed, saa man hører sig selv tale, om Bevægelsen kun foregaaer i Eens Indre; der hvor man hvert Øieblik sætter Livet ind, hvert Øieblik mister det og atter vinder det.

Ideen tilhører jeg. Naar den vinker mig, da følger jeg, naar den giver Stævne, da venter jeg Dage og Nætter, der kalder Ingen paa mig til Middag, der venter Ingen med Aftensmaden. Naar Ideen kalder, da forlader jeg Alt, eller rettere jeg har Intet at forlade, jeg sviger Ingen, jeg bedrøver Ingen ved at være den tro, min Aand bedrøves ikke ved at jeg maa bedrøve en Anden. Naar jeg vender hjem, da læser Ingen i min Mine, Ingen udfritter min Skikkelse. Ingen affrister mit Væsen en Forklaring, som jeg end ikke selv kan give nogen Anden, om jeg er salig i Glæde, eller nedsjunken i Nød, om jeg har vundet Livet eller tabt det.

Beruselsens Bæger rækkes mig atter, jeg indaander allerede dets Duft, jeg fornemmer allerede dets skummende Musik – dog først en Libation for hende, der frelste en Sjæl, som sad i Fortvivlelsens Eensomhed: priset være qvindeligt Høimod! – Leve Tankens Flugt, leve Livsfaren i Ideens Tjeneste, leve Kampens Nød, leve Seirens festlige Jubel, leve Dandsen i det Uendeliges Hvirvel, leve Bølgeslaget, der skjuler mig i Afgrunden, leve Bølgeslaget, der slynger mig op over Stjernerne.

\workheading{Philosophiske Smuler (af Johannes Climacus)}{1844}

\worknote{Antologiens ''Philosophical Fragments''. Hong-udvalget bringer titelbladets Propositio, hele første kapitel ''Tanke-Projekt'' (til og med ''… at blive til fra ikke at være?'') samt den afsluttende ''Moralen''. Kapitel II–V og Mellemspillet er ikke medtaget.}

\sectionheading{Propositio  ·  SV¹ IV, 173. 179}

Kan der gives et historisk Udgangspunkt for en evig Bevidsthed; hvorledes kan et saadant interessere mere end historisk; kan man bygge en evig Salighed paa en historisk Viden?

Spørgsmaalet gjøres af den Uvidende, der end ikke veed, hvad der har givet Anledning til at han spørger saaledes

\sectionheading{Capitel I: Tanke-Projekt  ·  SV¹ IV, 179–190}

A

Hvorvidt kan Sandheden læres? Med dette Spørgsmaal ville vi begynde. Det var et socratisk Spørgsmaal, eller blev det ved det socratiske Spørgsmaal, hvorvidt kan Dyden læres? thi Dyden bestemtes igjen som Indsigt (confer Protagoras, Gorgias, Meno, Euthydem). Forsaavidt Sandheden skal læres, maa den jo forudsættes ikke at være, altsaa idet den skal læres, søges den. Her møder nu den Vanskelighed, som Socrates i Meno (§ 80 Slutning) gjør opmærksom paa som paa en »stridslysten Sætning«, at et Menneske umuligt kan søge hvad han veed, og lige saa umuligt søge hvad han ikke veed; thi hvad han veed, kan han ikke søge, da han veed det, og hvad han ikke veed, kan han ikke søge, thi han veed jo end ikke, hvad han skal søge. Vanskeligheden gjennemtænker Socrates derved, at al Læren og Søgen kun er Erindren, saaledes, at den Uvidende blot behøver at mindes, for ved sig selv at besinde sig paa, hvad han veed. Sandheden bringes saaledes ikke ind i ham, men var i ham. Denne Tanke udfører Socrates videre, og i den concentrerer egentligen den græske Pathos sig, da den vorder et Beviis for Sjelens Udødelighed, vel at mærke retrogradt, eller et Beviis for Sjelens Præexistents.

I Betragtning heraf viser det sig, med hvilken vidunderlig Conseqvents Socrates blev sig selv tro og kunstnerisk realiserede, hvad han havde forstaaet. Han var og blev Gjordemoder; ikke fordi han »ikke havde det Positive«, men fordi han indsaae, at hiint Forhold er det høieste, et Menneske kan indtage mod det andet. Og deri beholder han jo i al Evighed Ret; thi selv om der nogensinde er givet et guddommeligt Udgangspunkt, mellem Menneske og Menneske bliver dette det sande Forhold, naar man reflekterer paa det Absolute og ikke gantes med det Tilfældige, men af Hjertens Grund renoncerer paa at forstaae den Halvhed, der synes Menneskenes Lyst og Systemets Hemmelighed. Derimod var Socrates en af Guden selv examineret Gjordemoder; den Gjerning han fuldkommede var et guddommeligt Ærende (confer Platos Apologi); om han end forekom Menneskene en Særling (ατοπωτατος. Theaet. § 149); og det var guddommeligt meent, hvad Socrates ogsaa forstod, at Guden forbød ham at føde (μαιευεσϑαι με ο ϑεος αναγϰαζει, γενναν δε απεϰωλυσεν. Theaet. § 150); thi mellem Menneske og Menneske er det μαιευεσϑαι det Høieste, det at føde tilhører jo Guden.

Socratisk seet, er ethvert Udgangspunkt i Tiden eo ipso et Tilfældigt, et Forsvindende, en Anledning; Læreren er ei heller mere, og giver han sig og sin Lærdom hen paa nogen anden Maade, da giver han ikke, men fratager, da er han end ikke den Andens Ven, mindre hans Lærer. Dette er den socratiske Tænknings Dybsind, dette hans ædle gjennemførte Humanitet, der ikke gjør slet og forfængeligt Compagni med gode Hoveder, men ogsaa føler sig ligesaa beslægtet med en Feldbereder, hvorfor han jo snarligen »forvissede sig om, at Physiken ikke er Menneskets Sag, og derfor begyndte at philosophere om det Ethiske i Værkstederne og paa Torvet« (Diogenes af L. II, 5, 21), men philosopherede lige absolut, hvem han end talede med. Med halve Tanker, med Prutten og Tingen, med Hævden og Opgiven, som skyldte den Enkelte til en vis Grad et andet Menneske Noget, men da atter igjen til en vis Grad ikke; med løse Ord, der forklare Alt, undtagen: hvilken denne visse Grad er; med alt Sligt gaaer man ikke ud over Socrates, naaer ei heller Aabenbaringsbegrebet, men bliver i Passiaren. For den socratiske Betragtning er ethvert Menneske sig selv det Centrale, og hele Verden centraliserer kun til ham, fordi hans Selverkjendelse er en Guds-Erkjendelse. Saaledes forstod Socrates sig selv, saaledes maatte efter hans Anskuelse ethvert Menneske forstaae sig selv, og i Kraft deraf maatte han forstaae sit Forhold til den Enkelte, altid lige ydmygt og lige stolt. Dertil havde Socrates Mod og Besindighed at være sig selv nok, men ogsaa i Forhold til Andre kun at være Anledning endog for det dummeste Menneske. O, sjeldne Høimod, sjelden i vor Tid, hvor Præsten er lidt mere end Degnen, hvor hvert andet Menneske er Autoritet, medens alle disse Forskjelligheder og al den megen Autoritet medieres i en fælleds Galskab, og i et commune naufragium; thi medens aldrig noget Menneske har været i Sandhed Autoritet, eller gavnet nogen Anden ved at være det, eller har formaaet i Sandhed at tage Clienten med sig, saa lykkes det bedre paa en anden Maade; thi det slaaer aldrig feil, at een Nar, idet han selv gaaer, tager flere andre med sig.

Forholder det sig saaledes med at lære Sandheden, da kan det, at jeg har lært af Socrates eller af Prodikos eller af en Tjenestepige, kun beskæftige mig historisk, eller, forsaavidt jeg er en Plato i Sværmeri, digterisk. Men dette Sværmeri, om det end er skjønt, om jeg end ønsker mig selv og ønsker Enhver denne ευϰαταφορια εις παϑος, som kun Stoikeren kunde advare imod, om jeg end ikke har socratisk Høimod og socratisk Selvfornægtelse til at tænke dens Intethed – dette Sværmeri er dog kun en Illusion, det vilde Socrates sige, ja en Uklarhed, i hvilken den jordiske Forskjellighed gjærer næsten vellystigt. Det kan heller ei interessere mig anderledes end historisk, at Socrates' eller Prodikos' Lære var den og den, thi Sandheden, i hvilken jeg hviler, var i mig selv og kom frem ved mig selv, og end ikke Socrates formaaede at give mig den, saa lidet som Kudsken formaaer at trække Hestens Byrde, om han end ved Svøben kan hjælpe den dertil. Mit Forhold til Socrates og Prodikos kan ikke beskæftige mig med Hensyn til min evige Salighed, thi denne er givet retrogradt i Besiddelsen af den Sandhed, hvilken jeg fra Begyndelsen havde, uden at vide det. Vilde jeg tænke mig i et andet Liv at træffe sammen med Socrates, Prodikos, eller Tjenestepigen, da blev atter her ingen af dem mere end Anledning, hvad Socrates uforfærdet udtrykker derved, at han selv i Underverdenen kun vilde spørge; thi den finale Tanke af al Spørgen er, at den Spurgte dog selv maa have Sandheden, og faae den ved sig selv. Det timelige Udgangspunkt er et Intet; thi i samme Øieblik som jeg opdager, at jeg fra Evighed har vidst Sandheden, uden at vide det, i samme Nu er hiint Øieblik skjult i det Evige, indoptaget deri, saaledes, at jeg, saa at sige, end ikke kan finde det, selv om jeg søgte det, fordi der intet Her og Der er, men kun et ubique et nusquam.

B

Skal det nu forholde sig anderledes, da maa Øieblikket i Tiden have afgjørende Betydning, saaledes at jeg intet Øieblik hverken i Tid eller Evighed vil kunne glemme det, fordi det Evige, som før ikke var, blev til i dette Øieblik. Lader os nu under denne Forudsætning betragte Forholdene med Hensyn til Spørgsmaalet om, hvorvidt Sandheden kan læres.

a) Den forudgaaende Tilstand

Vi begynde med den sokratiske Vanskelighed, hvorledes man kan søge Sandheden, da det jo er lige umuligt, enten man har den, eller man ikke har den. Den socratiske Tænkning ophævede egentligen Disjunctionen, idet det viste sig, at i Grunden har ethvert Menneske Sandheden. Dette var hans Forklaring; vi have seet, hvad deraf fulgte med Hensyn til Øieblikket. Skal dette nu faae afgjørende Betydning, da maa den Søgende lige indtil Øieblikket ikke have havt Sandheden, end ikke i Uvidenhedens Form, thi da bliver Øieblikket kun Anledningens; ja han maa end ikke være den Søgende; thi saaledes maae vi udtrykke Vanskeligheden naar vi ikke ville forklare den socratisk. Han maa da altsaa være bestemmet som udenfor Sandheden (ikke kommende til den, som Proselyt, men gaaende fra den), eller som Usandhed. Han er da Usandheden. Men hvorledes skal man nu minde ham, eller hvad skal det hjælpe ham, at minde ham om, hvad han ikke har vidst, og altsaa heller ikke kan besinde sig paa.

b) Læreren

Skal Læreren være Anledningen, der minder den Lærende, da kan han jo ikke bidrage til at han erindrer, at han egentligen veed Sandheden, thi den Lærende er jo Usandheden. Det Læreren da kan vorde ham Anledning til at han erindrer, er, at han er Usandheden. Men ved denne Besindelse er den Lærende jo netop udelukket fra Sandheden, mere end da han var uvidende om, at han var Usandheden. Paa denne Maade støder altsaa Læreren, netop ved at minde ham, den Lærende bort fra sig, kun at den Lærende, ved saaledes at vendes ind i sig selv, ikke opdager, at han forud vidste Sandheden, men opdager sin Usandhed, i Henseende til hvilken Bevidstheds-Akt det Socratiske gjælder, at Læreren kun er Anledning, i hvo han saa end er, selv om han er en Gud; thi min egen Usandhed kan jeg kun opdage ved mig selv, thi først naar jeg opdager det, er det opdaget, tidligere ikke, om saa hele Verden vidste det. (Under den antagne Forudsætning om Øieblikket bliver dette den eneste Analogi til det Socratiske.)

Skal nu den Lærende erholde Sandheden, da maa Læreren bringe ham den, og ikke blot dette, men maa give ham Betingelsen med for at forstaae den; thi dersom den Lærende selv var sig selv Betingelsen for at forstaae Sandheden, da behøver han jo blot at erindre; thi det er med Betingelsen for at forstaae Sandheden som med Det at kunne spørge om den, Betingelsen og Spørgsmaalet indeholder det Betingede og Svaret. (Forholder dette sig ikke saaledes, da er Øieblikket kun at forstaae socratisk).

Men Den, der giver den Lærende ikke blot Sandheden, men Betingelsen med, han er ikke Lærer. Al Underviisning beroer paa, at Betingelsen dog til syvende og sidst er tilstede; mangler denne, saa formaaer en Lærer Intet; thi i andet Fald maa han jo ikke omdanne, men omskabe den Lærende, førend han begynder at lære ham. Men dette formaaer intet Menneske, skal det da skee, maa det være ved Guden selv.

Forsaavidt nu den Lærende er til, er han jo skabt, og forsaavidt maa Gud have givet ham Betingelsen for at forstaae Sandheden (thi ellers var han jo tidligere kun Dyr, og hiin Lærer, der med Betingelsen gav ham Sandheden, gjorde ham først til Menneske); men forsaavidt Øieblikket skal have afgjørende Betydning (og hvis dette ikke antages, staae vi jo ved det Socratiske), maa han være uden Betingelsen, altsaa være den berøvet. Dette kan ikke være skeet ved Guden (thi det er en Modsigelse), ei heller ved et Tilfælde (thi det er en Modsigelse, at hvad der er lavere skulde kunne overvinde det Høiere), det maa da være skeet ved ham selv. Kunde han have tabt Betingelsen saaledes, at det ikke var skeet ved ham, og være i Tabets Tilstand uden at det skeer ved ham selv, da har han kun tilfældigt besiddet Betingelsen, hvilket er en Modsigelse, da Betingelsen for Sandheden er en væsentlig Betingelse. Usandheden er da ikke blot udenfor Sandheden, men er polemisk mod Sandheden, hvilket udtrykkes derved, at han selv har forspildt og forspilder Betingelsen.

Læreren er da Guden selv, der virkende som Anledning foranlediger, at den Lærende mindes om, at han er Usandheden, og er det ved egen Skyld. Men denne Tilstand, at være Usandheden og være det ved egen Skyld, hvad kunne vi kalde den? lader os kalde den Synd.

Læreren er da Guden, der giver Betingelsen og giver Sandheden. Hvad skulle vi nu kalde en saadan Lærer; thi derom ere vi jo enige, at vi allerede langt have overskredet Bestemmelsen af en Lærer. Forsaavidt den Lærende er i Usandheden, men er det ved sig selv (og paa anden Maade kan han jo ifølge det Foregaaende ikke være det), kunde det synes at han var fri; thi at være hos sig selv, det er jo Frihed. Og dog er han jo ufri og bunden og udelukket; thi at være fri for Sandheden, det er jo at være udelukket, og at være udelukket ved sig selv, det er jo at være bunden. Men da han er bunden ved sig selv, kan han da ikke løsne sig selv, eller befrie sig selv; thi hvad der binder mig, det samme maa jo ogsaa kunne frigive mig, naar det vil, og da dette er ham selv, saa maa han jo kunne det. Først maatte han da vel ville det. Men sæt nu, han saa dybt mindedes, hvad hiin Lærer blev Anledning til (og dette tør jo aldrig glemmes) at han erindrede; sæt nu, han vilde det. I saa Fald (dersom han ved at ville det, kunde det ved sig selv) blev da det at han havde været bunden en forbigangen Tilstand, der i Frigjørelsens Øieblik var sporløst forsvunden, og Øieblikket fik ikke afgjørende Betydning; han havde været uvidende om, at han havde bundet sig selv, og nu frigjorde han sig selv. Saaledes tænkt faaer Øieblikket da ingen afgjørende Betydning, og dog var det jo dette, vi vilde antage som Hypothesis. Altsaa efter Hypothesen vil han ikke kunne frigjøre sig selv. (Og saaledes er det ogsaa i Sandhed; thi han bruger Frihedens Kraft i Ufrihedens Tjeneste, da han jo er frit i den, og saaledes voxer Ufrihedens forenede Kraft og gjør ham til Syndens Træl). – Hvad skulle vi nu kalde en saadan Lærer, der giver ham Betingelsen igjen og med den Sandheden? lader os kalde ham en Frelser, thi han frelser jo den Lærende ud af Ufriheden, frelser ham fra sig selv; en Forløser, thi han løser jo den, der havde fanget sig selv, og Ingen er jo saa forfærdelig fangen, og intet Fangenskab saa umuligt at bryde ud fra, som det hvori Individet holder sig selv! Og dog er dette jo endnu ikke nok sagt; thi ved Ufriheden havde han jo forskyldt Noget, og giver da hiin Lærer ham Betingelsen og Sandheden, da er han jo en Forsoner, der borttager den Vrede, som var over det Forskyldte.

En saadan Lærer vil da den Lærende aldrig kunne glemme; thi i samme Øieblik synker han ned til sig selv igjen, ligesom Den, der eengang i Besiddelse af Betingelsen, ved at glemme, at Gud er til, sank i Ufriheden. Hvis de traf sammen i et andet Liv, da vilde hiin Lærer atter kunne give Den Betingelsen, som ikke havde modtaget den; men Den, der engang havde modtaget den, for ham vilde han være en Anden. Betingelsen var jo et Betroet, hvorfor Modtageren altid blev Regnskab skyldig. Men en saadan Lærer, hvad skulle vi kalde ham? En Lærer kan jo bedømme den Lærende, om han gjør Fremskridt eller ei, men dømme ham, kan han ikke; thi han maa jo være socratisk nok til at indsee, at det Væsentlige kan han ikke give den Lærende. Hiin Lærer er da egentlig ikke Lærer, men han er Dommer. Selv naar den Lærende meest har iført sig Betingelsen og ved den fordybet sig i Sandheden, han kan dog aldrig glemme hiin Lærer, eller lade ham socratisk forsvinde, hvilket dog er langt dybsindigere end al utidig Smaalighed og skuffet Sværmeri, ja er det Høieste, hvis hiint Andet ikke er Sandhed.

Og nu Øieblikket. Et saadan Øieblik er af en egen Natur. Det er vel kort og timeligt som Øieblikket er det, forbigaaende som Øieblikket er det, forbigangent, som Øieblikket er det, i det næste Øieblik, og dog er det Afgjørende, og dog er det fyldt af det Evige. Et saadant Øieblik maa dog have et særligt Navn, lad os kalde det: Tidens Fylde.

c) Discipelen

Naar Discipelen er Usandheden (og ellers gaae vi jo tilbage til det Socratiske), men dog er Menneske, og han nu faaer Betingelsen og Sandheden, da vorder han jo ikke først Menneske, thi det var han; men han vorder et andet Menneske, ikke i den spøgende Forstand, som blev han en Anden af samme Qvalitet som tidligere, men han bliver et Menneske af en anden Qvalitet, eller, som vi og kunne kalde det, et nyt Menneske.

Forsaavidt han var Usandheden, var han jo bestandigt ifærd med at gaae bort fra Sandheden; ved i Øieblikket at modtage Betingelsen, antog hans Gang den modsatte Retning, eller han vendtes om. Lader os kalde denne Forandring Omvendelse, om dette end er et hidtil ikke brugt Ord; men derfor vælge vi det just, for ikke at forstyrres; thi det er jo som skabt for den Forandring, om hvilken vi tale.

Forsaavidt han var i Usandheden ved egen Skyld, kan denne Omvendelse ikke foregaae uden at den optages i hans Bevidsthed, eller uden at han vorder sig bevidst, at det var ved hans egen Skyld; og med denne Bevidsthed tager han Afsked fra det Tidligere. Men hvorledes tager man Afsked, uden med Sorg i Sindet? Dog er denne Sorg jo her over, at han saa længe havde været i den tidligere Tilstand. Lader os kalde en saadan Sorg Anger; thi hvad Andet er vel Anger, der vel seer sig tilbage, men dog saaledes, at den netop derved fremskynder Gangen til, hvad der ligger foran!

Forsaavidt han var i Usandheden og nu med Betingelsen modtager Sandheden, da foregaaer der jo en Forandring med ham, som fra ikke at være til at være. Men denne Overgang fra ikke at være til at være, det er jo Fødselens. Dog Den, der er til, kan jo ikke fødes, og dog fødes han. Lader os kalde denne Overgang Gjenfødelsen, hvorved han anden Gang kommer til Verden ganske som ved Fødselen, et enkelt Menneske, der endnu Intet kjender til den Verden, i hvilken han fødes ind, om den er beboet, om der er andre Mennesker i den; thi man kan vel døbes en masse, men aldrig gjenfødes en masse. Som Den, der ved socratisk Fødselshjælp fødte sig selv, derover glemte alt Andet i Verden, og i dybere Forstand ikke skyldte noget Menneske Noget, saaledes skylder jo den Igjenfødte ikke noget Menneske Noget, men hiin guddommelige Lærer Alt, og maa, som hiin over sig selv glemte den hele Verden, atter over denne Lærer glemme sig selv.

Dersom da Øieblikket skal have afgjørende Betydning, og uden dette tale vi jo kun socratisk, i hvad vi end sige, om vi end bruge mange og besynderlige Ord, om vi end ved ikke at forstaae os selv skulde mene at være komne langt videre end hiin eenfoldige Vise, der ubestikkeligen dømte mellem Guden, Menneskene og sig selv, ubestikkeligere end Minos, Æachus og Rhadamantus; – saa er Bruddet skeet, og Mennesket kan ikke komme tilbage, og skal ikke lystes ved at erindre, hvad Mindelsen vil bringe ham i Erindring, og endnu mindre skal han formaae ved egen Kraft paany at drage Guden over paa sin Side.

Men lader det her Udviklede sig tænke? Vi ville ikke haste med Svaret, og ikke blot Den blev jo Svar skyldig, der formedelst Overveielsens Vidtløftighed aldrig kom til at svare, men ogsaa Den, der vel udviste en vidunderlig Hurtighed i at svare, men ikke den ønskelige Langsomhed i at betænke Vanskeligheden, førend han forklarede den. Førend vi svare, ville vi da spørge, hvo det er, som skal besvare Spørgsmaalet. Det at være født, lader det sig tænke? Ja hvorfor ikke; men hvo er det som skal tænke det; Den, som er født, eller Den, som ikke er født? Det Sidste er jo en Urimelighed, som da heller ei kan være falden Nogen ind; thi den Fødte kan jo ikke faae dette Indfald. Naar da den Fødte tænker sig født, da tænker han jo denne Overgang fra ikke at være til at være. Saaledes maa det vel ogsaa forholde sig med Igjenfødelsen. Eller gjør Dette Sagen vanskeligere, at den forud for Igjenfødelsen gaaende Ikke-Væren indeholder mere Væren end den Ikke-Væren, der gaaer forud for Fødselen? Men hvo er det da der skal tænke det? Det maa jo være den Igjenfødte, thi at den Ikke-Gjenfødte skulde gjøre det, var jo en Urimelighed, og var det ikke latterligt nok, om den Ikke-Igjenfødte vilde faae dette Indfald?

Dersom et Menneske oprindeligen er i Besiddelse af Betingelsen for at forstaae Sandheden, da tænker han, at Gud er til, derved at han selv er til. Dersom han er i Usandheden, da maa han jo tænke dette om sig selv, og Erindringen skal ikke kunne hjælpe ham uden til at tænke dette. Om han skal komme videre, maa Øieblikket afgjøre (om dette end allerede var virksomt i at lade ham indsee, at han er Usandheden). Forstaaer han ikke dette, da er han at henvise til Socrates, om han end ved sin Formening at være gaaet langt videre vil gjøre denne Vise mangt et Bryderi, som de gjorde det, der bleve saa opbragte paa ham, naar han fratog dem en eller anden Dumhed (επειδαν τινα ληρον αυτων αφαιρωμαι), at de ordentligen vilde bide ham (confer Theaetet § 151). – I Øieblikket bliver Mennesket sig bevidst, at han var født; thi hans Foregaaende, til hvilket han da ikke skal henholde sig, var jo, ikke at være; i Øieblikket bliver han sig Gjenfødelsen bevidst; thi hans foregaaende Tilstand var jo, ikke at være. Hvis hans foregaaende Tilstand havde været at være, da havde i intet af Tilfældene Øieblikket faaet afgjørende Betydning for ham, som ovenfor udviklet. Medens da den græske Pathos concentrerer sig paa Erindringen, concentrerer vort Projekts Pathos sig paa Øieblikket, og hvad Under, eller er det ikke en høist pathetisk Sag, at blive til fra ikke at være?

\sectionheading{Moralen  ·  SV¹ IV, 272}

Moralen

Dette Projekt gaaer udisputeerligt videre end det Socratiske, hvilket viser sig paa ethvert Punkt. Om det derfor er sandere end det Socratiske, er et ganske andet Spørgsmaal, hvilket ikke lader sig afgjøre i samme Aandedrag, da her jo blev antaget et nyt Organ: Troen, og en ny Forudsætning: Syndsbevidstheden, en ny Afgjørelse: Øieblikket, og en ny Lærer: Guden i Tiden, uden hvilke jeg sandeligen ikke skulde have vovet til Visitation at fremstille mig for hiin gjennem Aartusinder beundrede Ironiker, hvem jeg trods Nogen nærmer mig med Begeistringens Hjertebanken. Men at gaae videre end Socrates, naar man dog væsentligen siger det Samme som han, kun ikke slet saa godt, det er idetmindste ikke socratisk.

\workheading{Johannes Climacus eller De omnibus dubitandum est}{1842–43}

\worknote{Efterladt, ufuldført fortælling (Pap. IV B 1). Antologiens ''Johannes Climacus…''. Hong-udvalget bringer den fortællende Indledning, begyndelsen af Pars prima (de tre sætninger om Tvivl) og af Pars secunda første kapitels § 1 om bevidsthedens modsigelse, hvor manuskriptet afbrydes.}

\sectionheading{Indledning  ·  Pap. IV B 1, 104–112}

I Byen H..... levede for nogle Aar siden en ung Student, ved Navn Johannes Climacus, hvis Attraae ingenlunde var at blive bemærket i Verden, da det tvertimod var hans Glæde at henleve skjult og i Stilhed. De, der kjendte ham lidt nærmere, søgte at forklare hans indesluttede Væsen, der flygtede enhver videre Berøring med Menneskene, deraf, at han enten var melancholsk eller forelsket. De, der antoge det Sidste, havde i en vis Forstand ikke Uret, om de end feilede, hvis de antoge, at det var en Pige, der var Gjenstand for hans Drømmen. Saadanne Følelser vare aldeles fremmede for hans Hjerte, og som hans Ydre var fiint og ætherisk, næsten gjennemsigtigt, saaledes var hans Sjæl i samme Grad altfor aandelig bestemmet til at fængsles af en Qvindes Skjønhed. Forelsket var han, sværmerisk forelsket – i Tanken eller rettere i Tænkningen. Ingen elskende Yngling kan inderligere bevæges over den ubegribelige Overgang, hvormed Elskoven vaagner i hans Bryst, over det Lynets Slag, hvorved Gjenkjærligheden bryder frem i den Elskedes, end han over den begribelige Overgang, i hvilken den ene Tanke føier sig ind i den anden, en Overgang, der for ham var det lykkelige Øieblik, i hvilket det skete, hvad han havde ahnet og forventet i Sjælens Stilhed. Naar da hans Hoved tankefuldt bøiede sig som et modent Ax, da var det ikke fordi han hørte den Elskedes Stemme, men fordi han lyttede til Tankernes hemmelige Hvidsken; naar hans Blik var drømmende, da var det ikke fordi han ahnede hendes Billede, men fordi Tankens Bevægelse blev synlig for ham. Det var hans Lyst at begynde med en enkelt Tanke, fra den ad Consequentsens Vei at stige Trin for Trin op til en høiere; thi Consequentsen var ham en scala paradisi, og hans Salighed forekom ham endnu herligere end Englenes. Naar han nemlig var kommen op til den høiere Tanke, da var det ham en ubeskrivelig Glæde, en Lidenskabens Vellyst, hovedkuls at styrte sig ned i de samme Consequentser, indtil han naaede hen til det Punkt, fra hvilket han var gaaet ud. Dog dette lykkedes ham ikke altid efter Ønske. Naar han ikke fik netop saa mange Stød, som der var Leed i Consequentsen, da blev han bedrøvet, thi da var Bevægelsen ufuldkommen. Han begyndte nu forfra. Lykkedes det ham, da gysede hans Sjæl i Vellyst; han kunde af Glæde ikke sove om Natten, men vedblev hele Timer at gjøre den samme Bevægelse; thi dette Tankens Opogned, Nedogop var en Fryd uden Lige. I de lykkelige Tider var hans Gang let, næsten svævende; til andre Tider var den ængstelig og usikker. Saa længe han nemlig arbeidede for at stige op, saa længe Consequentsen endnu ikke havde formaaet at bane sig Veien, da var han betynget; thi de mange Consequentser, han havde færdige, uden at de dog endnu vare blevne ham ret klare og nødvendige, frygtede han for at tabe. Naar man seer et Menneske bære en Mængde skrøbelige Ting, stablede ovenpaa hverandre, da forundres man ikke over, at han gaaer usikkert og hvert Øieblik griber efter Ligevægten; naar man ikke seer Stabelen, da smiler man, saaledes som ogsaa Mange smilede ad Johannes Climacus uden at ahne, at hans Sjæl bar en langt høiere Stabel end hvad der ellers forbauser, at hans Sjæl var ængstelig for, at en eneste Conseqvents skulde falde ud; thi da gik det Hele istykker. At Menneskene smilede ad ham, bemærkede han ikke, ligesaa lidet som, at til andre Tider en Enkelt glad vendte sig om for at see efter ham, naar han let som i en Dands ilede gjennem Gaden. Han agtede ikke paa Menneskene, tænkte heller ei paa, at de kunde agte paa ham, han var og blev fremmed i Verden.

Som nu Climacuss Væsen for den Enkelte, der ikke kjendte ham nøiere, maatte have noget Paafaldende, saaledes var det ingenlunde uforklarligt for den, der vidste lidt Beskeed om hans tidligere Liv; thi som han nu var i sit 21d Aar, saaledes havde han til en vis Grad altid været. Den Disposition, hans Sjæl havde, var i hans Barndom ikke bleven forstyrret, men udviklet ved gunstige Omstændigheder. Hans Hjem tilbød ikke mange Adspredelser, og da han saa godt som aldrig kom ud, blev han tidlig vant til at sysselsætte sig med sig selv og med sine egne Tanker. Hans Fader var en meget stræng Mand, tilsyneladende tør og prosaisk, medens han under denne Vadmelsfrakke skjulte en glødende Phantasie, som end ikke hans høie Alderdom formaaede at sløve. Naar Johannes stundom bad om Tilladelse til at gaae ud, da fik han oftest et Afslag; hvorimod Faderen en enkelt Gang til Vederlag foreslog ham ved hans Haand at spadsere op og ned af Gulvet. Dette var ved første Øiekast en fattig Erstatning, og dog gik det med den ligesom med den Vadmelsfrakke, den skjulte noget ganske Andet i sig. Forslaget blev antaget, og det blev ganske overladt Johannes at bestemme, hvor de skulde gaae hen. De gik da ud af Porten, til et nærliggende Lystslot, eller ud til Strandkanten, eller om i Gaderne, alt eftersom Johannes vilde det; thi Faderen formaaede Alt. Medens de da gik op og ned af Gulvet, fortalte Faderen Alt hvad de saae; de hilste paa de Forbigaaende; Vognene larmede dem forbi og overdøvede Faderens Stemme; Kagekonens Frugter vare mere indbydende end nogensinde. Han fortalte Alt saa nøiagtigt, saa levende, saa præsentisk indtil den ubetydeligste Detail, der var Johannes bekjendt, saa udførligt og anskueligt hvad der var ham ubekjendt, at han, efter en halv Time at have spadseret med Faderen var saa overvældet og træt som om han en heel Dag havde været ude. Faderens Trolddomskonst lærte Johannes ham snart af. Hvad der før foregik episk, det skete nu dramatisk; de samtalede paa Touren. Gik de ad bekjendte Veie, da passede de gjensidigen paa hinanden, at Intet blev overseet; var Veien Johannes fremmed, da combinerede han, medens Faderens almægtige Phantasie var istand til at danne Alt, at bruge ethvert barnagtigt Ønske som en Ingredients i det Drama, der foregik. For Johannes var det som blev Verden til under Samtalen, som var Faderen vor Herre, og han selv hans Yndling, der fik Lov at blande sine taabelige Indfald deri ligesaa lystigt som han vilde; thi han blev aldrig afviist, Faderen aldrig forstyrret, det kom Alt med og altid til Johanness Tilfredshed.

Medens saaledes Livet i det fædrene Huus bidrog til at udvikle hans Phantasie, lærte ham at finde Smag i Ambrosia, saa var den Dannelse, der blev ham til Deel i Skolen, i Harmonie hermed. Den latinske Grammatiks ophøiede Myndighed, Reglens guddommelige Værdighed udviklede et nyt Sværmerie. Især tiltalte dog den græske Grammatik ham. Over den glemte han at læse Homer høit for sig selv, som han ellers pleiede for at glæde sig ved Digtets Rythmer. Den græske Lærer foredrog Grammatiken paa en mere philosophisk Maade. Naar det da blev forklaret ham, at Accusativ for Exempel er Udstrækningen i Tid og Rum, at Præpositionen ikke styrer Casus, men at Forholdet gjør det, da udvidede Alt sig for ham. Præpositionen forsvandt, Udstrækningen i Tid og Rum blev som et uhyre indholdsløst Billede for Intuitionen. Hans Phantasie var atter sysselsat men paa en anden Maade end før. Hvad der paa Spadseretourene forlystede ham, var det fyldte Rum, han kunde ikke faae det tæt nok om sig. Hans Phantasie var saa produktiv, at den kunde hjælpe sig med lidet. Udenfor det ene Vindue i Dagligstuen voxede der omtrent ti Græs Straa. Her opdagede han stundom et lille Dyr der løb imellem Stenglerne. Disse Straae bleve til en uhyre Skov, der dog havde hele den Tæthed og Dunkelhed som Græsset havde det. Istedenfor det fyldte Rum fik han nu det tomme Rum, han stirrede atter, men saae Intet uden den uhyre Udstrækning.

Medens saaledes en næsten vegetativ Henslumren i Phantasie – deels mere æsthetisk, deels mere intellectuelt – udvikledes hos ham, blev ogsaa en anden Side af Sjælen stærk dannet, hans Sands nemlig for det Pludselige, det Overraskende. Dette skete dog ikke ved de Tryllemidler, der ellers maa tjene til at spænde Børns Opmærksomhed, men ved noget langt Høiere. Med en almægtig Phantasie forbandt Faderen en uimodstaaelig Dialektik. Naar da ved en eller anden Leilighed Faderen indlod sig paa en Disput med en Anden, da var Johannes lutter Øre, og det saa meget mere som Alt foregik med en næsten festlig Orden. Faderen lod altid Modparten tale heelt ud, spurgte ham for en Forsigtigheds Skyld, om han havde mere at sige, førend han begyndte sit Tilsvar. Johannes havde fulgt Modpartens Foredrag med spændt Opmærksomhed, var paa sin Maade medinteresseret i Udfaldet. Pausen indtraadte, Faderens Replik fulgte, og see! i en Haandevending var Alt anderledens. Hvorledes det gik til, blev Johannes en Gaade, men hans Sjæl forlystede sig ved dette Skuespil. Modparten talte igjen, Johannes var endnu mere opmærksom for ret at holde Alt fast; Modparten perorerede, Johannes kunde næsten høre sit Hjerte banke, saa utaalmodigen ventede han, hvad der vel vilde skee. – Det skete; i et Nu var Alt vendt om, det Forklarlige gjort uforklarligt, det Visse tvivlsomt, Modsætningen indlysende. Naar en Hai vil gribe sit Bytte, da maa den kaste sig om paa Ryggen, da dens Mund sidder paa dens Bug; den er mørk paa Ryggen, sølvhvid under Bugen. Det skal være et herligt Syn, at see denne Afvexling i Farve; den skal stundom blinke saa stærkt, at det næsten gjør ondt i Øiet, og dog forlyster det at see derpaa. En lignende Afvexling blev Johannes Vidne til, naar han hørte Faderen dispütere. Han glemte igjen det Sagte, baade hvad Faderen og hvad Modparten havde sagt, men denne Gysen i Sjælen glemte han ikke. Analogier hertil lod Skolelivet ham ikke savne; han saae, hvorledes et Ord kunde forandre en heel Sætning, hvorledes et conjunctivum midt i en indikativisk Sætning kunde kaste en forandrende Belysning over det Hele. Jo ældre han blev, jo mere Faderen indlod sig med ham, jo mere opmærksom blev han paa hiint Ubegribelige, det var som om Faderen stod i en hemmelig Forstaaelse med det Johannes vilde sige, og derfor ved et eneste Ord kunde forvirre ham Alt. Naar da Faderen ikke blot opponerede, men selv foredrog Noget, da saae han, hvorledes han bar sig ad, hvorledes han successivt kom hen til det han vilde. Han ahnede nu, at Grunden hvorfor Faderen med et Ord kunde vende Alt om, maatte ligge i, at han selv i den Succession, i hvilken han havde Tanken, maatte have glemt Noget.

Hvad andre Børn have i Poesiens Tryllerie og Eventyrets Overraskelse, det havde han i Intuitionens Hvile og Dialektikens Changements. Det glædede Barnet, det blev Drengens Leeg, det blev Ynglingens Lyst. Saaledes havde hans Liv en sjelden Continuitet; det kjendte ikke de forskjellige Overgange, der ellers pleie at betegne de enkelte Tidsafsnit. Idet Johannes blev ældre havde han intet Legetøi at lægge til Siden; thi han havde lært at lege med det, der skulde være hans Livs alvorlige Beskjæftigelse, og dog havde det derved ingenlunde tabt det Tillokkende. Et Pigebarn leger saa længe med Dukken, at den tilsidst forvandles til den Elskede; thi Qvindens hele Liv er Elskov. En lignende Continuitet havde hans Liv; thi hans hele Liv var Tænkning.

Climacus blev Student, tog anden Examen, blev 20 Aar gammel, og dog foregik der ingen Forandring med ham, han var og blev fremmed for Verden. Imidlertid flyede han ikke Mennesker, tvertimod søgte han at træffe Ligesindede. Dog udtalte han sig ikke, lod sig aldrig mærke med, hvad der foregik i hans Indre, dertil var hans Erotik for dyb; det var ham som maatte han rødme, hvis han talte derom, han frygtede, at faae for Meget at vide eller for Lidet at vide. Derimod var han altid opmærksom naar Andre talte. Som en dybt forelsket ung Pige ikke gjerne taler om sin Kjærlighed, men i næsten piinlig Spænding lytter til, naar andre unge Piger tale om deres, for da i Stilhed at prøve, om hun er ligesaa lykkelig eller endnu lykkeligere, for at opsnappe ethvert ledende Vink, saaledes agtede Johannes taus paa Alt. Naar han da kom hjem, da overveiede han, hvad de Philosopherende havde sagt; thi det var naturligviis disses Selskab han søgte.

At ville være Philosoph, at ville hellige sig udelukkende til Speculationen, var ikke faldet ham ind; dertil var han endnu for letsindig. Vel reves hans Sjæl ikke snart til Eet, snart til et Andet, Tænkningen var og blev hans Lidenskab; men endnu manglede den Besindelse for at faae et dybere Sammenhæng. Det Ubetydeligste og det Betydeligste fristede ham lige meget til derfra at begynde sine Operationer; Resultatet var ham ikke magtpaaliggende, kun Bevægelserne interesserede ham. Undertiden blev han vel opmærksom paa, hvorledes han fra ganske forskjellige Punkter naaede hen til eet og samme; men dette tildrog sig ikke i dybere Forstand hans Opmærksomhed. Hans Lyst var bestandig blot at trænge igjennem. Over Alt, hvor han ahnede en Labyrinth, der maatte han finde Veien. Begyndte han derpaa, da kunde Intet bringe ham til at høre op. Faldt det ham vanskeligt, blev han træt deraf før Tiden, da pleiede han at bruge et saare simpelt Middel. Han lukkede sig inde paa sit Værelse, gjorde Alt saa festligt som muligt og sagde nu høit og lydeligt: jeg vil det. Af Faderen havde han lært, at man kan, hvad man vil; og Faderens Liv havde ikke ladet Theorien i Stikken. Denne Erfaring havde givet Johanness Sjæl en ubeskrivelig Stolthed. At der skulde være Noget, man ikke kunde, uagtet man vilde det, var ham utaaleligt. Men hans Stolthed var heller ei en afmægtig Villen; thi naar han havde sagt disse energiske Ord, da var han villig til Alt, han havde da et endnu høiere Maal: at gjennemtrænge Vanskelighedens Slyngninger med sin Villie. Dette var igjen et Eventyr, der begeistrede ham. Hans Liv var saaledes altid eventyrligt, om han end ikke til sine Eventyr behøvede Skove og Reiser, men blot, hvad han havde, et lille Værelse paa eet Fag.

Skjøndt nu hans Sjæl tidlig var ført hen til Idealiteten, saa var dog derved hans Tro og Tillid til Virkeligheden ingenlunde svækket. Den Idealitet, ved hvilken han næredes, laae ham saa nær, Alt gik saa naturligt til, at denne blev hans Virkelighed, og han igjen i Virkeligheden uden om sig maatte vente at finde Idealiteten. Hertil bidrog Faderens Tungsind. At Faderen var en overordentlig Mand, var det Johannes sidst fik at vide. At han forbausede ham, at intet andet Menneske gjorde det i den Grad, det vidste han; dog han kjendte jo saa Faae, at han ingen Maalestok kunde anlægge. At Faderen menneskelig talt var noget Ualmindeligt, det lærte han ikke i det fædrene Huus. Naar engang imellem en ældre prøvet Ven besøgte Familien og indlod sig i en fortroligere Samtale med Faderen, da hørte Johannes ham ofte sige: jeg duer til Intet, kan Intet gjøre, mit eneste Ønske var at finde en Plads i en mild Stiftelse«. Det var ikke Spøg, der var ikke Spoer af Ironi i Faderens Ord, tvertimod der var en mørk Alvor deri, der ængstede Johannes. Det var heller ingen løst henkastet Bemærkning; thi Faderen kunde nu bevise, at det allerubetydeligste Menneske var et Genie imod ham. Intet Modbeviis udrettede Noget, da hans uimodstaaelige Dialektik var istand til at bringe Een til at glemme, hvad der laae allernærmest, nødsage Een til at stirre paa den af ham opstillede Betragtning, som om der intet Andet var til i Verden. Johannes, hvis hele Anskuelse af Livet, var ligesom skjult i Faderen, da han selv kun fik saare Lidet at see, indvikledes herved i en Modsigelse, da det en lang Tid undgik ham, at Faderen om ikke ved Andet saa ved den Virtuositet, med hvilken han kunde beseire enhver Opponent og bringe ham til at forstumme, modbeviiste sig selv. Johanness Tillid til Virkeligheden var da ikke svækket, han havde ikke indsuget Idealiteten af Skrifter, der ikke lade den, de opfostre, uvidende om, at den Herlighed, de beskrive, dog ikke findes i Verden; han var ikke dannet ved en Mand, der forstod at gjøre sin Viden kostbar, men vel at gjøre den saa ubetydelig og værdiløs som mulig.

\sectionheading{Pars prima (indledning): de tre Sætninger om Tvivl  ·  Pap. IV B 1, 113–114}

Ved at agte paa Andres Tale, mærkede han nok, at han heller ikke var truffen paa Skrifter af de store Tænkere blandt de nyere Philosopher. Disses Navne hørte han atter og atter nævne med Begeistring, næsten med Tilbedelse. Dette glædede ham ubeskriveligt, om han end ikke vovede sig til at læse dem, da han hørte, at de vare meget vanskelige, saa Studiet af dem krævede Aar og Dag. Det var ikke Feighed eller Magelighed, der standsede ham, men en smertelig Følelse, der havde ledsaget ham fra hans tidlige Barndom, at han ikke var saaledes som andre Mennesker. Han var langt fra at føle nogen Glæde i denne Forskjellighed, tvertimod han følte den som et Tryk, hvorunder han vel vilde komme til at lide hele sit Liv. Han forekom sig selv som et Barn, der med megen Smerte var født til Verden, og som ikke kunde glemme denne Smerte, om Moderen end havde glemt den af Glæde over hans Fødsel.

Hvad Læsning angik, var Johannes nu kommen i en besynderlig Modsigelse. De Skrifter, han kjendte, vilde ikke tilfredsstille ham, uden at han dog vovede, at skyde Skylden derfor paa Skrifterne; de udmærkede Skrifter vovede han sig ikke til at læse. Han læste da mindre og mindre, fulgte sit Hang til i Stilhed at sysle med Tanken, blev mere og mere undselig af Frygt for, at de fornemme Tænkere skulde smile ad ham, naar de hørte, at han ogsaa vilde tænke, ligesom fornemme Damer smile ad en ringe Pige, naar hun formaster sig til ogsaa at ville kjende Elskovens Salighed. Han taug, men hørte desto opmærksommere.

Naar han da hørte de Andre tale, bemærkede han, at en enkelt Hovedtanke kom igjen atter og atter. En saadan greb han da og gjorde til Gjenstand for sin Tænken. Saaledes kom Skjebnen ham til Hjælp med at forskaffe ham Stof netop paa den Maade, han behøvede det. Jo renere saa at sige jo jomfrueligere Opgaven var, desto kjærere for ham, jo mindre Andre havde hjulpet paa Tanken, desto lykkeligere var han, desto bedre gik Alt for ham. Han følte det vel som en Ufuldkommenhed, at han bedst kunde tænke en Tanke, naar den kom til ham som nyfalden Snee, uden at være gaaet igjennem Andres Hænder, han ansaae det for det Store, at kunne som de Andre tumle sig i de mangfoldige Tænkeres mangfoldige Tanker; dog snart glemte han denne Smerte i Tænkningens Glæde.

Ved at lytte til de Andres Tale, blev han især opmærksom paa een Sætning, der atter og atter forekom, der gik fra Mund til Mund, altid lovpriist, altid venereret. Mangfoldige Gange hørte han det gjentage: de omnibus dubitandum est. Nu stod han da ved den Sætning, der kom til at spille en afgjørende Rolle i hans Liv. Denne Sætning blev for hans Liv, hvad ellers et Navn ofte er i et Menneskes Historie, man kan i Korthed sige Alt ved at nævne dette Navn.

Denne Sætning blev en Opgave for hans Tænken. Om det skulde vare længe eller kort inden han fik den gjennemtænkt, det vidste han ikke, men det vidste han, at han før det Øieblik ikke vilde slippe den, om det saa skulde koste ham Livet.

Hvad der end mere begeistrede ham, var det Forhold, i hvilket denne Sætning i Almindelighed blev sat til det at blive Philosoph. Om han da skulde kunne blive det, vidste han ikke, men han vilde gjøre Alt. Med stille Høitidelighed blev der da decreteret, at han skulde begynde. Han opbyggede sig ved at mindes Dions Begeistring, der, da han med nogle Faae gik ombord for at begynde Krigen med Dionys, sagde: mig er det nok at have deeltaget deri. Om jeg skulde døe i det Øieblik, jeg steeg i Land, uden at have udrettet Noget, jeg vil dog ansee denne Død for lykkelig og hæderlig.

Forholdet mellem hiin Sætning og Philosophien søgte han nu at blive sig bevidst. Beskjæftigelsen hermed vilde være ham et opmuntrende Forspil, jo mere klar han blev derover, jo mere begeistret vilde han gaae til Sagen selv. Saaledes lukkede han sig nu inde i sig selv med hiin philosophiske Sætning, medens han dog tillige agtede nøie paa ethvert Vink, han om den kunde opsnappe. Mærkede han, at hans egen Tankegang var forskjellig fra de Andres, da indprægede han sig deres i Hukommelsen, gik hjem og begyndte forfra. At deres Tankegang i Almindelighed var meget kort, var ham vel paafaldende, men han saae kun deri et nyt Fortrin hos dem.

Han begyndte nu sine Operationer, og sammenstillede strax de 3 Hoved-Udsagn, han havde hørt angaaende denne Sætnings Forhold til Philosophien.

Disse 3 Sætninger løde saaledes: 1. Philosophien begynder med Tvivl; 2 man maa have tvivlet for at komme til at philosophere; 3. den nyere Philosophie begynder med Tvivl.

\sectionheading{Pars secunda, Cap. I: Hvad det er at tvivle? — § 1 (uddrag)  ·  Pap. IV B 1, 141–150}

Idet Johannes begyndte paa denne Overveielse, indsaae han vel, at hvis han forlangte et empirisk Svar paa dette Spørgsmaal, da vilde Livet tilbyde en Mangfoldighed, der indenfor Extremets hele Omfang vilde skjule en kun forvirrende Vidtløftighed. Ikke blot kunde nemlig det, der hos den Enkelte fremkaldte Tvivlen være saare forskjellig, men det kunde være det Modsatte; thi hvis Een for at vække Tvivlen hos en Anden vilde foredrage Tvivlen, saa kunde han netop derved fremkalde Troen, ligesom omvendt Troen kunde fremkalde Tvivlen. Paa Grund af denne paradoxe Dialektik, der, hvad han alt tidligere havde været opmærksom paa, ingen Analogie havde i nogen Videns Sphære, da al Viden staaer i et ligefrem og immanent Forhold til sin Gjenstand og den Vidende, ikke i et omvendt og transcendent Forhold til en Tredie, indsaae han lettelig, at enhver empirisk Betragtning her vilde føre til Intet. Naar han da søgte et Svar paa hiint Spørgsmaal maatte han gaae en anden Vei. Han maatte see at udfinde Tvivlens ideelle Mulighed i Bevidstheden. Denne maatte jo blive den samme, hvor forskjellig end det foranledigende Phænomen var, da den, uden selv at forklares af Phænomenet, forklarede Phænomenets Virkning. Hvad der i den Enkelte frembragte Tvivlen, kunde være saa forskjelligt som det være vilde, hvis der i den Enkelte ikke var denne Mulighed, saa var Intet istand til at fremkalde den. Da fremdeles det foranledigende Phænomens Forskjellighed kunde være contrair, saa maatte Muligheden være en total, væsentlig for den menneskelige Bevidsthed.

Han søgte da at orientere sig i Bevidstheden saaledes som denne er i sig selv som den, der forklarer enhver enkelt Bevidsthed, uden dog selv at være en enkelt. Han spurgte, hvorledes Bevidstheden var beskaffen, naar den havde Tvivlen udenfor sig. I Barnet er Bevidstheden, men denne har Tvivlen udenfor sig. Hvorledes er da Barnets Bevidsthed bestemmet? Den er egentlig slet ikke bestemmet, hvilket ogsaa kan udtrykkes saaledes, den er umiddelbar. Umiddelbarheden er netop Ubestemmetheden. I Umiddelbarheden er intet Forhold, thi saasnart Forholdet er der, er Umiddelbarheden hævet. Umiddelbart er derfor Alt sandt, men denne Sandhed er i næste Øieblik Usandhed; thi umiddelbart er Alt usandt. Kan Bevidstheden forblive i Umiddelbarheden, saa er Spørgsmaalet om Sandhed hævet.

Hvorledes fremkommer Spørgsmaalet om Sandheden? Ved Usandheden; thi i det Øieblik jeg spørger om Sandheden, har jeg allerede spurgt om Usandheden. I Spørgsmaalet om Sandheden er Bevidstheden bragt i Forhold til noget Andet, og det der muliggjør dette Forhold er Usandheden.

Hvad er først Umiddelbarheden eller Middelbarheden? Det var et captieust Spørgsmaal. Han kom derved til at tænke paa det Svar, som Thales skal have givet Een, der spurgte om Natten eller Dagen var bleven først til: Natten er een Dag før. Η νυξ, εφη, μια ημερα προτερον. (confer Diog: af L. Liber 1. § 36.). Kan da Bevidstheden ikke forblive i Umiddelbarheden? Dette var et taabeligt Spørgsmaal; thi hvis den kunde det, var der slet ingen Bevidsthed til. Men hvorledes hæves da Umiddelbarheden? Ved Middelbarheden, der hæver Umiddelbarheden, ved at forudsætte den. Hvad er da Umiddelbarheden? Det er Realiteten. Hvad er Middelbarheden? Det er Ordet. Hvorledes hæver denne hiin? Ved at udtale den; thi Det, der udtales, er altid forudsat.

Umiddelbarheden er Realiteten, Sproget er Idealiteten, Bevidstheden er Modsigelsen. I det Øieblik jeg udsiger Realiteten, er Modsigelsen der; thi det jeg siger er Idealiteten.

Tvivlens Mulighed ligger da i Bevidstheden, hvis Væsen er en Modsigelse, der er frembragt ved og selv frembringer en Dupplicitet.

En saadan Dupplicitet har nødvendigviis 2 Udtryk. Duppliciteten er Realiteten og Idealiteten, Bevidstheden er Forholdet. Jeg kan enten bringe Realiteten i Forhold til Idealiteten, eller Idealiteten i Forhold til Realiteten. I Realiteten alene er der ingen Tvivlens Mulighed; idet jeg udtrykker den i Sproget, er Modsigelsen der, da jeg slet ikke udtrykker den, men frembringer noget Andet. Forsaavidt det Sagte skal være et Udtryk for Realiteten har jeg bragt denne i Forhold til Idealiteten, forsaavidt det Sagte er et af mig Frembragt, har jeg bragt Idealiteten i Forhold til Realiteten. Saalænge denne Udvexling foregaaer uden gjensidig Berøring, er Bevidstheden kun til efter sin Mulighed. I Idealiteten er ligesaa fuldt som i Realiteten Alt sandt. Som jeg derfor kan sige, at umiddelbart er Alt sandt, saa kan jeg ogsaa sige, at umiddelbart er Alt virkeligt; thi først i det Øieblik Idealiteten bringes i Forhold til Realiteten fremkommer Muligheden. I Umiddelbarheden er det Falskeste og det Sandeste lige sandt, i Umiddelbarheden er det Muligste og det Umuligste lige virkeligt. Saalænge denne Udvexling foregaaer uden Sammenstød er Bevidstheden egentlig ikke til, og dette uhyre Falsum foraarsager ingen Ophævelser. Realiteten er ikke Bevidstheden, Idealiteten ligesaa lidet, og dog er Bevidstheden ikke til uden begge, og denne Modsigelse er Bevidsthedens Tilbliven og dens Væsen.

Førend han gik videre, betænkte han, om ikke det han her kaldte Bevidsthed var hvad man ellers vilde kalde Reflexion. Han fastsatte Bestemmelsen i denne Henseende saaledes: Reflexionen er Muligheden af Forholdet, Bevidstheden er Forholdet, hvis første Form er Modsigelsen. Han bemærkede tillige, at deraf kom det, at Reflexionens Bestemmelser altid vare dichotomiske. Saaledes er: Idealitet og Realitet, Sjæl og Legeme, at erkjende – det Sande, ville – det Gode, elske – det Skjønne, Gud og Verden og saa videre Reflexions Bestemmelser. I Reflexionen berøre de hinanden saaledes, at et Forhold bliver muligt. Bevidsthedens Bestemmelser derimod ere trichotomiske, hvilket ogsaa Sproget beviser; thi naar jeg siger: jeg bliver mig dette Sandseindtryk bevidst, saa siger jeg en Trehed. Bevidstheden er Aand, og dette er det mærkelige, at naar i Aandens Verden Een deles bliver den 3 aldrig 2. Bevidstheden forudsætter derfor Reflexionen. Dersom dette ikke forholdt sig saaledes, saa blev det umuligt at forklare Tvivlen. Vel syntes Sproget at stride hermed; thi i de fleste Sprog, saavidt det var ham bekjendt, sættes det Ord at tvivle i etymologisk Henseende i Forhold til: To. Dog meente han, at derved blot var antydet Tvivlens Forudsætning, saa meget mere som det var ham klart, at saasnart jeg som Aand bliver til To er jeg eo ipso Tre. Var der ikke Andet end Dichotomier til, saa var der ingen Tvivl til; thi Tvivlens Mulighed ligger netop i det Tredie, der sætter de Tvende i Forhold til hinanden. Man kunde derfor ikke sige, at Reflexionen frembringer Tvivlen, med mindre man vilde udtrykke sig bagvendt, man maatte sige, Tvivlen forudsætter Reflexionen, uden at dog dette prius er temporairt. Tvivlen opstaaer derved, at der bliver et Forhold mellem To, men for at dette skal skee, maa de To være, medens dog Tvivlen, som et høiere Udtryk, gaaer forud ikke følger efter.

Reflexionen er Muligheden af Forholdet. Dette kan ogsaa udtrykkes saaledes: Reflexionen er interesseløs. Bevidstheden derimod er Forholdet og derved Interessen, hvilken Dobbelthed fuldelig og med prægnant Tvetydighed er udtrykt i Ordet Interesse (interesse). Al Viden derfor der er interesseløs (mathematisk; æsthetisk; metaphysisk) er kun Tvivlens Forudsætning. Saasnart Interessen er hævet er Tvivlen ikke overvundet, men neutraliseret, og enhver saadan Viden kun en Tilbagegang. Forsaavidt derfor Nogen meente ved en saa kaldet objektiv Tænkning at overvinde Tvivlen, var dette en Misforstaaelse; thi Tvivlen er en høiere Form end al objektiv Tænkning; thi den forudsætter denne, men har et Mere, et Tredie, hvilket er Interessen eller Bevidstheden. I denne Henseende forekom de græske Skeptikeres Adfærd ham langt consequentere end den moderne Overvindelse af Tvivlen. Disse indsaae meget vel, at Tvivlen ligger i Interessen, og meente derfor ganske consequent at hæve Tvivlen ved at forvandle Interessen til Apathie. Denne Fremgangsmaade var dog en Consequents, hvorimod det var en Inconsequents, der syntes at have sin Grund i Uvidenhed om hvad Tvivl er, der bevægede den nyere Philosophie til systematisk at ville betvinge Tvivlen. Selv om Systemet var absolut fuldendt, selv om Virkeligheden overgik Bebudelserne, Tvivlen var dog ikke overvundet, den begynder først; thi Tvivlen ligger i Interessen, og enhver systematisk Erkjenden er interesseløs. Man seer deraf, at Tvivlen er Begyndelsen til den høieste Form af Tilværelse, fordi den kan have alt Andet til sin Forudsætning. De græske Skeptikere indsaae ypperligt, at det at tale om Tvivlen, naar Interessen er hævet, er urimeligt, men de vilde formodentlig tillige have indseet, at det er et Ordspil at tale om en objektiv Tvivl; thi lad Idealiteten og Realiteten i al Evighed stride mod hinanden, saa længe der ingen Bevidsthed, ingen Interesse er, ingen Bevidsthed er, der har Interesse af denne Strid, saa længe er der ingen Tvivl; lad dem være forsonede, Tvivlen kan lige fuldt vedblive.

Bevidstheden er da Forholdet, et Forhold, hvis Form er Modsigelsen. Men hvorledes opdager Bevidstheden Modsigelsen? Dersom hiint omtalte Falsum kunde vedblive, at Idealiteten og Realiteten i al Troskyldighed communicerede med hinanden, saa vilde Bevidstheden aldrig fremkomme; thi Bevidstheden fremkommer netop ved Sammenstødet, ligesom den forudsætter Sammenstødet. Umiddelbart er der intet Sammenstød, men middelbart er det der. Saasnart der bliver Spørgsmaal om en Gjentagelse, saa er Sammenstødet der; thi Gjentagelse er kun tænkelig af hvad der har været før.

I Realiteten som saadan er der ingen Gjentagelse. Dette kommer ikke deraf, at Alt er forskjelligt, ingenlunde. Om Alt i Verden var absolut eens, i Realiteten er der ingen Gjentagelse, fordi den blot er i Momentet. Om Verden istedenfor at være Skjønheden var lutter ligestore, eensformige Kampestene, der var dog ingen Gjentagelse. Jeg vilde i al Evighed see i ethvert Moment en Kampesteen, men om det var den samme, som jeg før havde seet, derom var intet Spørgsmaal. I Idealiteten alene er ingen Gjentagelse; thi Ideen er og bliver den samme, og kan som saadan ikke gjentages. Naar Idealiteten og Realiteten berøre hinanden, da fremkommer Gjentagelsen. Idet jeg da i Momentet fE seer Noget, træder Idealiteten til og vil forklare, at det er en Gjentagelse. Her er Modsigelsen; thi det, der er, er tillige paa en anden Maade. At det Udvortes er, det seer jeg, men i samme Øieblik bringer jeg det i Forhold til Noget, der ogsaa er, Noget der er det samme, og som tillige vil forklare, at det Andet er det samme. Her er en Fordobling her er Spørgsmaal om en Gjentagelse. Idealiteten og Realiteten støde altsaa sammen; i hvilket Medium? I Tiden? det er jo en Umulighed. I Evigheden? Det er jo en Umulighed. Hvori da? I Bevidstheden, der er Modsigelsen. Spørgsmaalet er ikke interesseløs, som hvis man spurgte, om ikke hele Tilværelsen var et Afbillede af Ideen, og om forsaavidt ikke, i en vis forflygtiget Forstand, den synlige Tilværelse var en Gjentagelse. Spørgsmaalet er her nærmere om en Gjentagelse i Bevidstheden, altsaa om Erindringen. Erindringen har den samme Modsigelse. Erindringen er ikke Idealiteten, den er Idealiteten der har været, den er ikke Realiteten, den er Realiteten, der har været, hvilket igjen er en dobbelt Modsigelse; thi Idealiteten kan ifølge sit Begreb ikke have været, Realiteten ifølge sit Begreb ligesaa lidet.

[…]

Anmærkning. De græske Sophisters Sætning, at Alt er sandt. Platos Bestræbelser for at modbevise dem, især ved at vise, at det Negative er til. (confer Sophisten.) – Schleiermachers Lære med Hensyn til Følelsen, at Alt er sandt. (confer hans Dogmatik i Begyndelsen; nogle Modbemærkninger af Erdmann i Bruno Baurs Tidsskrift 3 Bind 1 Hæfte pagina 11.). Heraklits Sætning Alt er og Alt er ikke, hvilken Aristoteles forklarer saaledes: Alt er sandt. confer Tennemann Geschichte der Philosophie 1ste Bind pagina 237. note

Anmærkning. Hvad Johannes her gjør sig Rede for er ikke uden Betydning. Den nyere Philosophies Terminologie er ofte forvirrende. Den taler saaledes om sinnliches Bewußtsein, wahrnehmendes Bewusstsein, Verstand og saa videre, uagtet den langt hellere maatte kalde det Sandsning, Erfaring, thi i Bevidsthed ligger mere. Det var overhovedet ret interessant at see, hvorledes Hegel vilde danne Overgangen fra Bevidsthed til Selvbevidsthed fra Selvbevidsthed til Fornuft. Naar Overgangen blot bestaaer i en Overskrift, saa er det let nok.

\workheading{Begrebet Angest (af Vigilius Haufniensis)}{1844}

\worknote{Antologiens ''The Concept of Anxiety''. Hong-udvalget bringer af Caput I §§ 5–6 (uskyldighed og angest; angesten som arvesyndens forudsætning) samt af Caput III's indledning om tiden og øieblikket. Caput II og de lange lærde fodnoter (bl.a. om Parmenides) er ikke medtaget; spring er markeret med […].}

\sectionheading{Caput I, §§ 5–6: Begrebet Angest (uddrag)  ·  SV¹ IV, 313–322}

Uskyldigheden er Uvidenhed. I Uskyldigheden er Mennesket ikke bestemmet som Aand, men sjelelig bestemmet i umiddelbar Eenhed med sin Naturlighed. Aanden er drømmende i Mennesket. Denne Opfattelse er ganske i Overeensstemmelse med Bibelens, der ved at negte Mennesket i Uskyldigheden Kjendskab til Forskjellen mellem Godt og Ondt bryder Staven over alle katholsk-fortjenstlige Phantasterier.

I denne Tilstand er der Fred og Hvile; men der er paa samme Tid noget Andet, hvilket ikke er Ufred og Strid; thi der er jo Intet at stride med. Hvad er det da? Intet. Men hvilken Virkning har Intet? Det føder Angest. Dette er Uskyldighedens dybe Hemmelighed, at den paa samme Tid er Angest. Drømmende projekterer Aanden sin egen Virkelighed, men denne Virkelighed er Intet, men dette Intet seer Uskyldigheden bestandig udenfor sig.

Angest er en Bestemmelse af den drømmende Aand, og hører som saadan hjemme i Psychologien. Vaagen er Forskjellen mellem mig selv og mit Andet sat, sovende er den suspenderet, drømmende er den et antydet Intet. Aandens Virkelighed viser sig bestandig som en Skikkelse, der frister dens Mulighed, men er borte saa snart den griber efter den, og er et Intet, der kun kan ængste. Mere kan den ikke, saa længe den kun viser sig. Begrebet Angest seer man næsten aldrig behandlet i Psychologien, jeg maa derfor gjøre opmærksom paa, at det er aldeles forskjelligt fra Frygt og lignende Begreber, der referere sig til noget bestemt, medens Angest er Frihedens Virkelighed som Mulighed for Muligheden. Man vil derfor ikke finde Angest hos Dyret, netop fordi det i sin Naturlighed ikke er bestemmet som Aand.

Naar vi ville betragte de dialektiske Bestemmelser i Angest, da viser det sig, at disse netop have den psychologiske Tvetydighed. Angest er en sympathetisk Antipathie og en antipathetisk Sympathie. Man seer, tænker jeg, letteligen, at det er i en ganske anden Forstand en psychologisk Bestemmelse end hiin concupiscentia. Sprogbrug bekræfter det fuldkommen, man siger: den søde Angest, den søde Beængstelse, man siger: en underlig Angest, en sky Angest og saa videre

Den Angest, der er sat i Uskyldigheden, er da for det første ingen Skyld, for det andet er den ingen besværende Byrde, ingen Lidelse, der ikke lod sig bringe i Samklang med Uskyldighedens Salighed. Naar man vil iagttage Børn, vil man finde denne Angest bestemtere antydet som en Søgen efter det Eventyrlige, det Uhyre, det Gaadefulde. At der gives Børn, hos hvilke den ikke findes, beviser Intet; thi Dyret har den heller ikke, og jo mindre Aand jo mindre Angest. Denne Angest hører Barnet saa væsentligen til, at han ikke vil undvære den; om den end ængster ham, fængsler den ham dog i sin søde Beængstelse. Hos alle de Nationer, hos hvilke det Barnlige er bevaret som Aandens Drømmen er denne Angest; og jo dybere den er, jo dybere er Nationen. Det er kun en prosaisk Dumhed, der mener, at det er en Desorganisation. Angest har her samme Betydning som Tungsind paa et langt senere Punkt, hvor Friheden, efter at have gjennemløbet de ufuldkomne Former af sin Historie, i dybeste Forstand skal komme til sig selv.

Ligesom da Angestens Forhold til sin Gjenstand, til Noget, hvilket er Intet (Sprogbrug siger ogsaa prægnant: at ængstes for Intet) er aldeles tvetydigt, saaledes vil den Overgang, der her kan gjøres fra Uskyldighed til Skyld, netop være saa dialektisk, at den viser, at Forklaringen er, hvad den skal være, psychologisk. Det qualitative Spring er udenfor al Tvetydighed, men den, der gjennem Angest bliver skyldig, han er jo uskyldig; thi det var ikke ham selv, men Angesten, en fremmed Magt, der greb ham, en Magt, han ikke elskede, men ængstedes for; – og dog er han jo skyldig, thi han sank i Angesten, som han dog elskede idet han frygtede den. Der gives i Verden intet Tvetydigere end dette, og derfor er denne den eneste psychologiske, medens den, for atter at gjentage det, aldrig falder paa at ville være Forklaringen, der forklarer det qualitative Spring. Enhver Forestilling om, at Forbudet fristede ham, eller at Forføreren bedrog ham, har kun for en overfladisk Iagttagelse den tilstrækkelige Tvetydighed, forvansker Ethiken, frembringer en quantitativ Bestemmen, og vil ved Psychologiens Hjælp gjøre Mennesket en Compliment paa det Ethiskes Bekostning, hvilken Compliment Enhver, der er ethisk udviklet, maa frabede sig som en ny end profundere Forførelse.

At Angesten kommer tilsyne er det, hvorom Alt dreier sig. Mennesket er en Synthese af det Sjelelige og det Legemlige. Men en Synthese er utænkelig, naar de Tvende ikke enes i et Tredie. Dette Tredie er Aanden. I Uskyldigheden er Mennesket ikke blot Dyr, som han da overhovedet, hvis han noget Øieblik i sit Liv var blot Dyr, aldrig vilde blive Menneske. Aanden er altsaa tilstede, men som umiddelbar, som drømmende. Forsaavidt den nu er tilstede, er den paa en Maade en fjendlig Magt; thi den forstyrrer bestandig det Forhold mellem Sjel og Legeme, der vel har Bestaaen, men dog ikke har Bestaaen, forsaavidt det først faaer det ved Aanden. Paa den anden Side er den en venlig Magt, der jo netop vil constituere Forholdet. Hvilket er da Menneskets Forhold til denne tvetydige Magt, hvorledes forholder Aanden sig til sig selv og til sin Betingelse? Den forholder sig som Angest. Blive af med sig selv kan Aanden ikke; gribe sig selv, kan den heller ikke, saa længe den har sig selv udenfor sig selv; synke ned i det Vegetative kan Mennesket heller ikke, thi han er jo bestemmet som Aand; flye Angesten kan han ikke, thi han elsker den; egentlig elske den kan han ikke, thi han flyer den. Nu er Uskyldigheden paa sin Spidse. Den er Uvidenhed, men ikke en dyrisk Brutalitet, men en Uvidenhed, der er bestemmet af Aand, men som netop er Angest, fordi dens Uvidenhed er om Intet. Her er ingen Viden om Godt og Ondt og saa videre; men Videns hele Virkelighed projekterer sig i Angesten som Uvidenhedens uhyre Intet.

Endnu er Uskyldigheden, men der behøver blot at lyde et Ord, saa er Uvidenheden concentreret. Dette Ord kan Uskyldigheden naturligviis ikke forstaae, men Angesten har ligesom faaet sit første Bytte, istedenfor Intet har den faaet et gaadefuldt Ord. Naar det saaledes hedder i Genesis, at Gud sagde til Adam: »blot af Kundskabens Træ paa Godt og Ondt maa Du ikke spise«, saa følger det jo af sig selv, at Adam egentlig ikke forstod dette Ord; thi hvor skulde han forstaae Forskjellen paa Godt og Ondt, da denne Adskillelse jo først fulgte med Nydelsen.

Naar man nu antager, at Forbudet vækker Lysten, da faaer man en Viden, istedenfor Uvidenhed, thi Adam maa da have havt en Viden om Friheden, siden Lysten var den at bruge den. Denne Forklaring er derfor bag efter. Forbudet ængster ham, fordi Forbudet vækker Frihedens Mulighed i ham. Hvad der gik Uskyldigheden forbi som Angestens Intet, det er nu kommet ind i ham selv og er atter her et Intet, den ængstende Mulighed af at kunne. Hvad det er, han kan, derom har han ingen Forestilling; thi ellers forudsætter man jo, hvad i Almindelighed skeer, det Senere, Forskjellen mellem Godt og Ondt. Kun Muligheden af at kunne er der som en høiere Form af Uvidenhed, som et høiere Udtryk af Angest, fordi det i en høiere Forstand er og er ikke, fordi han i en høiere Forstand elsker og flyer det.

Efter Forbudets Ord følge Dommens Ord: da skal Du visseligen døe. Hvad det vil sige, at døe, fatter naturligviis Adam slet ikke, hvorimod der jo Intet er til Hinder for, hvis man antager det sagt til ham, at han har faaet Forestillingen om det Forfærdelige. Selv Dyret kan jo i denne Henseende forstaae det mimiske Udtryk og Bevægelsen i den Talendes Stemme uden at forstaae Ordet. Dersom man lader Forbudet vække Lysten, saa maa man ogsaa lade Straffens Ord vække en afskrækkende Forestilling. Dette forvirrer imidlertid. Forfærdelsen her bliver kun Angest; thi det Udsagte har Adam ikke forstaaet, og har altsaa kun igjen Angestens Tvetydighed. Den uendelige Mulighed af at kunne, som Forbudet vakte, rykkes nu nærmere ved, at denne Mulighed udviser en Mulighed som sin Følge.

Saaledes er Uskyldigheden bragt til sit Yderste. Den er i Angesten i Forhold til det Forbudne og Straffen. Den er ikke skyldig, og dog er der en Angest, som var den tabt.

Længere kan Psychologien ikke komme, men dette kan den naae, og fremfor Alt dette kan den i sin Iagttagelse af Menneskelivet vise atter og atter.

Jeg knyttede mig her i Slutningen til den bibelske Fortælling. Jeg lod Forbudet og Straffens Røst komme udenfra. Dette har naturligviis piint mangen Tænker. Den Vanskelighed er dog kun at smile ad. Uskyldigheden kan jo godt tale; forsaavidt eier den i Sproget Udtrykket for alt Aandeligt. Forsaavidt behøver man blot at antage, at Adam har talet med sig selv. Den Ufuldkommenhed i Fortællingen, at en Anden taler til Adam om hvad han ikke forstaaer, falder da bort. Fordi Adam har kunnet tale, deraf følger jo ikke i dybere Forstand, at han har kunnet forstaae det Udsagte. Fremfor Alt gjælder dette om Forskjellen mellem Godt og Ondt, hvilken vel er i Sproget, men kun er for Friheden. Uskyldigheden kan godt sige denne Forskjel, men Forskjellen er ikke for den, og har for den kun den Betydning, vi i det Foregaaende have viist.

§ 6 Angest som Arvesyndens Forudsætning og som forklarende Arvesynden retrogradt i Retning af dens Oprindelse

§ 6

Angest som Arvesyndens Forudsætning og som forklarende Arvesynden retrogradt i Retning af dens Oprindelse

Lad os saa nøiere gjennemgaae Fortællingen i Genesis, idet vi gjøre et Forsøg paa at opgive den fixe Idee, at det er en Mythe, og idet vi minde os om, at ingen Tid har været saa flink i at frembringe Forstands-Myther som vor, der selv frembringer Myther, medens den vil udrydde alle Myther.

Adam var da skabt, havde givet Dyrene Navne (her er altsaa Sproget, om end paa en lignende ufuldkommen Maade, som Børnene lære det ved i Fibelbrættet at kjende et Dyr), men havde ikke fundet Selskab for sig. Eva var skabt, dannet af hans Ribbeen. Hun stod i saa inderligt Forhold til ham som muligt, men dog var det endnu et udvortes Forhold. Adam og Eva er blot en numerisk Gjentagelse. Om der i den Forstand var tusinde Adam'er betyder ikke mere end der er een. Dette med Hensyn til Slægtens Nedstammen fra eet Par. Naturen ynder ikke en betydningsløs Overflødighed. Antager man derfor, at Slægten nedstammer fra flere Par, saa har der været et Øieblik, da Naturen havde en intetsigende Overflødighed. Saasnart Generations-Forholdet er sat, er intet Menneske en Overflødighed; thi hvert Individ er sig selv og Slægten.

Nu følger Forbudet og Dommen. Men Slangen var trædskere end alle Markens Dyr, den forlokkede Qvinden. Om man nu end vil kalde dette en Mythe, da maa man dog ikke glemme, at den slet ikke forstyrrer Tanken eller forvirrer Begrebet saaledes som en Forstands-Mythe gjør det. Mythen lader det foregaae udvortes, der er indvortes.

Hvad her først er at paaagte er, at Qvinden bliver først forført, at hun derpaa forfører Manden. Jeg skal senere søge at udvikle i et andet Capitel, i hvilken Forstand Qvinden er, som man siger, det svagere Kjøn, samt at Angest er hende mere tilhørende end Manden.

Flere Gange i det Foregaaende er der mindet om, at den Opfattelse, der fremsættes i dette Skrift, ikke negter Syndighedens Forplantelse i Generationen, eller med andre Ord at Syndigheden i Generationen har sin Historie, kun siges der, at denne bevæger sig i quantitative Bestemmelser, medens bestandig Synden kommer ind ved Individets qualitative Spring. En Betydning af Generationens Quantiteren kan man allerede see her. Eva er det Deriverede. Vel er hun skabt ligesom Adam, men hun er skabt ud af en foregaaende Skabning. Vel er hun uskyldig ligesom Adam, men der er ligesom en Ahnelse om en Disposition, der vel ikke er, men dog kan synes som et Vink om den ved Forplantelsen satte Syndighed, der er det Deriverede, hvilket prædisponerer den Enkelte uden dog at gjøre ham skyldig.

Hvad der i § 5 er sagt om Forbudets og Dommens Ord maa her mindes. Den Ufuldkommenhed i Fortællingen, hvorledes Nogen skulde falde paa at sige til Adam, hvad han væsentlig ikke kan forstaae, falder bort, naar vi betænke, at den Talende er Sproget, og det altsaa er Adam selv, der taler.

Nu staaer Slangen tilbage. Jeg er ingen Ven af Aandrighed, og skal volente deo modstaae Slangens Fristelser, der, som den i Tidens Begyndelse fristede Adam og Eva, i Tidens Løb fristede Skribenterne – til at være aandrige. Jeg tilstaaer hellere reent ud, at jeg ingen bestemt Tanke kan knytte til den. Vanskeligheden med Slangen er desuden en ganske anden, den nemlig at lade Fristelsen komme udvortes fra. Dette strider ligefrem mod Bibelens Lære, mod det bekjendte classiske Sted hos Jakob, at Gud frister Ingen og fristes og ikke af Nogen, men at Enhver fristes af sig selv. Naar man nemlig troer, at have reddet Gud ved at lade Mennesket fristes af Slangen, og forsaavidt mener at være i Overeensstemmelse med Jakob, »at Gud frister Ingen«, saa støder man an mod det Næste, at Gud ikke fristes af Nogen; thi Slangens Attentat mod Mennesket var tillige en indirecte Fristelse mod Gud ved at blande sig ind i Forholdet mellem Gud og Mennesket; og man støder an mod det Tredie, at ethvert Menneske fristes af sig selv.

Nu følger Syndefaldet. Dette kan Psychologien ikke forklare; thi det er det qualitative Spring. Men lad os et Øieblik betragte Følgen, som den angives i hiin Fortælling, for endnu engang at fæste Blikket paa Angesten som Arvesyndens Forudsætning.

Følgen var en dobbelt, at Synden kom ind i Verden, og at det Sexuelle blev sat, og det Ene skal være uadskilleligt fra det Andet. Dette er af yderste Vigtighed for at vise Menneskets oprindelige Tilstand. Var han nemlig ikke en Synthese, der hvilede i et Tredie, saa kunde Eet ikke have to Følger. Var han ikke en Synthese af Sjel og Legeme, der bæres af Aand, kunde det Sexuelle aldrig komme ind med Syndigheden.

Projektmagerier ville vi lade udenfor, og ganske simpelt antage den sexuelle Forskjelligheds Tilstedeværen før Syndefaldet, kun at den dog ikke var, fordi den ikke er i Uvidenhed. I denne Henseende have vi Medhold i Skriften.

I Uskyldigheden var Adam som Aand drømmende Aand. Synthesen er altsaa ikke virkelig; thi det Forbindende er netop Aanden, og denne er endnu ikke sat som Aand. Hos Dyret kan den sexuelle Forskjellighed være udviklet instinktagtigt, men saaledes kan et Menneske ikke have den, netop fordi han er Synthese. I det Øieblik da Aanden sætter sig selv, sætter den Synthesen, men for at sætte Synthesen, maa den først adskillende gjennemtrænge den, og det yderste af det Sandselige er netop det Sexuelle. Denne Yderlighed kan Mennesket først naae, i det Øieblik Aanden bliver virkelig. Før den Tid er han ikke Dyr, men heller ikke egentlig Menneske, først i det Øieblik han bliver Menneske, bliver han det ogsaa derved, at han tillige er Dyr.

Syndigheden er da ikke Sandseligheden, ingenlunde, men uden Synden ingen Sexualitet og uden Sexualitet ingen Historie. En fuldkommen Aand har hverken det Ene eller det Andet, hvorfor jo ogsaa den sexuelle Differents er hævet i Opstandelsen, og hvorfor heller ingen Engel har Historie. Selv om Michael havde optegnet alle de Ærinder, i hvilke han var bleven udsendt, og som han havde fuldkommet, er dette jo dog ikke hans Historie. Først i det Sexuelle er Synthesen sat som Modsigelse, men tillige, som enhver Modsigelse, som Opgave, hvis Historie i samme Øieblik begynder. Denne er Virkeligheden, forud for hvilken Frihedens Mulighed gaaer. Men Frihedens Mulighed er ikke at kunne vælge det Gode eller det Onde. En saadan Tankeløshed er ligesaa lidet i Skriftens som i Tænkningens Medfør. Muligheden er at kunne. I et logisk System er det nemt nok at sige, at Muligheden gaaer over til Virkelighed. I Virkeligheden er det ikke saa let, og der behøves en Mellembestemmelse. Denne Mellembestemmelse er Angesten, hvilken ligesaa lidet forklarer det qualitative Spring, som den ethisk retfærdiggjør det. Angest er ikke en Bestemmelse af Nødvendigheden, men heller ei af Friheden, den er en hildet Frihed, hvor Friheden ikke er fri i sig selv, men hildet, ikke i Nødvendigheden, men i sig selv. Er Synden kommen nødvendigt ind i Verden (hvilket er en Modsigelse), saa er der ingen Angest. Er Synden kommen ind ved en Akt af et abstrakt liberum arbitrium (der da ligesaa lidet som senere i Verden har existeret i Begyndelsen, da det er en Tanke-Uting), saa er der heller ikke Angest. Logisk at ville forklare Syndens Indkommen i Verden, er en Taabelighed, der kun kan falde Folk ind, der ere latterligen bekymrede for at faae en Forklaring.

Dersom jeg her havde Lov at gjøre et Ønske, saa vilde jeg ønske, at ingen Læser var dybsindig nok til at spørge: hvis Adam nu ikke havde syndet? I det Øieblik Virkeligheden er sat, gaaer Muligheden ved Siden som et Intet, der frister alle tankeløse Mennesker. At Videnskaben dog ikke kan beslutte sig til at holde Menneskene i Ave og sig selv i Tømme! Naar Een gjør et dumt Spørgsmaal, da see man vel til, at man ikke svarer ham; thi ellers er man nok saa dum som Spørgeren. Det Taabelige i hiint Spørgsmaal ligger ikke saa meget i Spørgsmaalet, som i at man dermed henvender sig til Videnskaben. Naar man bliver hjemme med det, som den kloge Else med sine Projekter, sammenkalder ligesindede Venner, saa har man dog nogenlunde forstaaet sin Dumhed. Videnskaben kan derimod ikke forklare Sligt. Enhver Videnskab ligger enten i en logisk Immanents, eller i Immanentsen indenfor en Transcendents, hvilken den ikke kan forklare. Synden er nu netop hiin Transcendents, hiin discrimen rerum, i hvilken Synden kommer ind i den Enkelte som den Enkelte. Anderledes kommer Synden ikke ind i Verden, og er aldrig kommen anderledes ind. Naar da den Enkelte er taabelig nok til at spørge om Synden som om noget ham uvedkommende, saa spørger han som en Nar; thi enten veed han slet ikke, hvad der er Tale om, og kan umuligen faae det at vide, eller han veed det og forstaaer det, samt at ingen Videnskab kan forklare ham det. Imidlertid har Videnskaben dog stundom været føielig nok til at imødekomme sentimentale Ønsker med grublende Hypotheser, om hvilke den da tilsidst selv indrømmede, at de ikke forklarede tilstrækkeligt. Dette er nu ogsaa ganske sandt; men Forvirringen er, at Videnskaben ikke energisk afviste taabelige Spørgsmaal, og derimod bestyrkede overtroiske Mennesker i, at der engang skulde komme en videnskabelig Projektmager, der var Mand for at hitte paa det Rette. Man taler om, at det er 6000 Aar siden, at Synden kom ind i Verden, aldeles paa samme Maade som at det er 4000 Aar siden, at Nebucadnezar blev en Oxe. Naar man opfatter Sagen saaledes, er det intet Under, at Forklaringen bliver derefter. Hvad der i een Henseende er det Simpleste af Verden, det gjør man til det Vanskeligste. Hvad det eenfoldigste Menneske forstaaer paa sin Viis og ganske rigtigt, fordi han forstaaer, at det ikke just er 6000 Aar siden, at Synden kom ind i Verden, det har Videnskaben ved Projektmageres Kunst faaet udsat som et Priisspørgsmaal, der endnu aldrig er blevet fyldestgjørende besvaret. Hvorledes Synden kom ind i Verden, forstaaer ethvert Menneske ene og alene ved sig selv; vil han lære det af en Anden, da vil han eo ipso misforstaae det. Den eneste Videnskab, der kan gjøre lidt, er Psychologien, der dog selv indrømmer, at den ikke forklarer, og ikke kan og ikke vil forklare Mere. Dersom nogen Videnskab kunde forklare det, saa er Alt forvirret. At Videnskabsmanden skal glemme sig selv, er ganske sandt; men derfor er det ogsaa saa heldigt, at Synden ikke er noget videnskabeligt Problem, og derfor ingen Videnskabsmand, saalidet som nogen Projektmager, forpligtet til at glemme, hvorledes Synden kom ind i Verden. Vil han det, vil han høimodigen glemme sig selv, da bliver han i sin Iver for at forklare hele Menneskeheden ligesaa comisk som hiin Hofraad, der i den Grad offrede sig for at aflægge sit Visitkort hos Creti og Pleti, at han derudover tilsidst glemte, hvad han selv heed. Eller hans philosophiske Begeistring gjør ham saa selvforglemmende, at han behøver en skikkelig nøktern Ægteviv, hvem han kan spørge, som Soldin spurgte Rebecca, da ogsaa han i begeistret Selvforglemmelse tabte sig i Passiarens Objektivitet: Rebekka, er det mig, som taler?

At min høistærede Samtids beundrede Videnskabs-Mænd, der i deres, for den hele Menighed vitterlige Bekymring for og Søgen efter Systemet, vel ogsaa ere bekymrede for at finde en Plads i det for Synden, ville finde dette høist uvidenskabeligt, er aldeles i sin Orden. Lad kun Menigheden søge med, eller dog indeslutte hine dybe Søgere i sine fromme Forbønner, de finde Pladsen ligesaa vist som Den, der søger efter den brændende Tamp, finder den, naar han ikke mærker, at den brænder i Haanden paa ham.

\sectionheading{Caput III: Tiden og Øieblikket (uddrag)  ·  SV¹ IV, 350–357}

Det blev i de tvende foregaaende Capitler bestandig fastholdt, at Mennesket er en Synthese af Sjel og Legeme, der constitueres og bæres af Aand. Angesten var, at jeg skal bruge et nyt Udtryk, der siger det Samme, som der sagdes i det Foregaaende, og som tillige peger hen til det Følgende, Angesten var Øieblikket i det individuelle Liv.

Der er en Kategorie, som bestandig bruges i den nyere Philosophie i logiske ikke mindre end i historisk-philosophiske Undersøgelser, det er: Overgangen. Nærmere Forklaring faaer man dog aldrig. Man benytter den frisk væk, og medens Hegel og den hegelske Skole fik Verden til at studse over den store Tanke, Philosophiens forudsætningsløse Begynden, eller at forud for Philosophien maa der ikke gaae Andet end den fuldkomne Forudsætningsløshed af Alt, generer man sig ingenlunde ved at benytte Overgangen, Negationen, Mediationen det vil sige Bevægelses-Principerne i den hegelske Tænken saaledes, at disse ikke tillige finde deres Plads i den systematiske Fremadskriden. Naar dette ikke er en Forudsætning, veed jeg ikke hvad Forudsætning er; thi at benytte Noget, som man intetsteds forklarer, er jo at forudsætte det. Systemet skulde have den vidunderlige Gjennemsigtighed og Indadskuen, at det omphalo-psychitisk urokkeligt skuede paa det centrale Intet saa længe, til Alt forklarede sig, og dets hele Indhold blev til ved sig selv. Denne indadvendte Offentlighed var jo Systemets. Imidlertid befindes dette ikke at være saa, og den systematiske Tanke synes at hylde Hemmelighedsfuldhed med Hensyn til sine inderste Rørelser. Negationen, Overgangen, Mediationen ere tre formummede, mistænkelige, hemmelige Agenter (agentia), der afstedkomme alle Bevægelser. Urolige Hoveder vilde Hegel vel aldrig kalde dem, da det er med hans allerhøieste Tilladelse at de drive deres Spil, saa ugeneret, at man selv i Logiken bruger Udtryk og Vendinger hentede fra Overgangens Timelighed: derpaa, naar, som værende er det det, som vordende er det det og saa videre

Dog dette være nu som det være vil, lad Logiken see til at hjælpe sig selv. Ordet Overgang er og bliver en Aandrighed i Logiken. I den historiske Friheds Sphære har den hjemme, thi Overgangen er en Tilstand, og er virkelig. Vanskeligheden af at anbringe Overgangen i det reent Metaphysiske har Plato meget godt indseet, og derfor har den Kategorie Øieblikket kostet ham saa megen Anstrængelse. At ignorere Vanskeligheden er vistnok ikke at »gaae videre« end Plato, at ignorere den, fromt bedragende Tænkningen, for at faae Speculationen flot og Bevægelsen i Logiken i Gang, er at behandle Speculationen som et temmelig endeligt Anliggende. Dog erindrer jeg engang at have hørt en Speculerende sige, at man ikke maatte tænke for meget paa Vanskelighederne iforveien; thi da kom man aldrig til at speculere. Naar det saaledes blot gjælder om at komme til at speculere, ikke om at Eens Speculation bliver virkelig Speculation, saa er det jo ganske resolut sagt, at man blot skal see at komme til at speculere, ligesom det er priseligt, om en Mand, der ikke havde Evne til at kjøre til Dyrehaven i egen Vogn, vilde sige: Sligt bør man ikke bryde sig om, man kan godt kjøre paa en Kaffemølle. Og saaledes er det jo ogsaa: Begge de Kjørende komme forhaabentligen til Dyrehaven. Derimod kommer den neppeligen til at speculere, der var resolveret nok til ikke at bryde sig om Befordringen, naar han barestens kunde komme til at speculere.

I den historiske Friheds Sphære er Overgangen en Tilstand. Imidlertid maa man, for retteligen at forstaae dette, ikke glemme, at det Ny kommer ved Springet. Hvis dette nemlig ikke fastholdes, da faaer Overgangen en quantiterende Overvægt over Springets Elasticitet.

Mennesket var altsaa en Synthese af Sjel og Legeme, men er tillige en Synthese af det Timelige og det Evige. At dette ofte nok er blevet sagt, har jeg Intet at indvende mod; thi det er ikke mit Ønske, at opdage Nyheder, men vel min Glæde og min forelskte Syssel at tænke over, hvad der synes ganske simpelt.

Hvad den sidste Synthese angaaer, da er det strax paafaldende, at den er dannet anderledes end den første. I den første var Sjel og Legeme Synthesens tvende Momenter, og Aanden det Tredie, dog saaledes, at der først egentlig var Tale om Synthesen idet Aanden sattes. Den anden Synthese har kun to Momenter: det Timelige og det Evige. Hvor er her det Tredie? Og er der intet Tredie, da er der egentligen ingen Synthese, thi en Synthese, som er en Modsigelse, kan ikke fuldkommes som Synthese uden i et Tredie; thi det, at Synthesen er en Modsigelse, udsiger jo netop, at den ikke er. Hvad er da det Timelige?

Naar man rigtigen bestemmer Tiden som den uendelige Succession, saa ligger det tilsyneladende nær ogsaa at bestemme den som nærværende, forbigangen og tilkommende. Imidlertid er denne Distinktion urigtig, hvis man mener, at den ligger i Tiden selv; thi den fremkommer først ved Tidens Forhold til Evigheden og ved Evighedens Reflexion i denne. Dersom man nemlig i Tidens uendelige Succession kunde finde et Fodfæste det vil sige et Nærværende, som var det Delende, saa var Inddelingen ganske rigtig. Men netop fordi ethvert Moment, som Summen af Momenter er det, er Proces (en Gaaen-forbi), saa er intet Moment et nærværende, og forsaavidt er der i Tiden hverken et Nærværende eller et Forbigangent eller et Tilkommende. […]

\workheading{Forord (af Nicolaus Notabene)}{1844}

\worknote{Antologiens ''Prefaces''. Hong-udvalget bringer bogens eget forord — den selvironiske fortælling om forfatteren, hvis kone forbyder ham at skrive, og som derfor kun skriver forord. De otte nummererede forord (I–VIII) er ikke medtaget.}

\sectionheading{Forord (uddrag)  ·  SV¹ V, 6–14}

I den nyere Videnskab har Forordet faaet sit Banesaar. Fra dens Standpunkt seet kommer den ældre Forfatter let til at gjøre en bedrøvelig Skikkelse, som man ikke veed om man skal lee eller græde over, fordi hans Ubehjælpelighed i at komme til Sagen gjør ham comisk, og hans Naivetet, som var der Nogen, der brød sig om ham, gjør ham rørende. Nuomstunder kan en saadan Situation ikke gjentage sig; thi naar man begynder Bogen med Sagen og Systemet med Intet, saa skjønnes der ikke at blive noget tilovers at sige i en Fortale. Denne Tingenes Orden har givet Anledning til, at jeg er bleven opmærksom paa, at Forord er en ganske egen Art af literær Frembringelse, og da den puffes ud er det paa den høie Tid at den emanciperer sig, ligesom alt Andet. Paa den Maade kan det endnu blive godt. Det Incommensurable, som man i ældre Tid nedlagde i Forordet til en Bog, kan nu finde sin Plads i et Forord, der ikke er Forord til nogen Bog. Saaledes troer jeg, at Striden er bilagt til gjensidig Contentement og Fornøielse; kan Forordet og Bogen ikke spændes med hinanden, saa lad den Ene give den Anden Skilsmissebrev.

Den nyeste videnskabelige Methode har gjort mig opmærksom paa, at det maatte komme til et Brud; min Fortjeneste bliver det, at gjøre Alvor af det med Brudet; her nu blot et Phænomen, der tyder hen paa den dybere Grund. Enhver æsthetisk udviklet Forfatter har vist havt Øieblikke, hvori det ikke vilde lyste ham at skrive en Bog, men hvor han ret gad skrive et Forord til en Bog, ligegyldig om den var af ham selv eller af en Anden. Dette tyder paa, at Forordet er væsentlig forskjelligt fra Bogen, og at det at skrive et Forord er noget ganske Andet end at skrive en Bog; thi i modsat Fald vilde Trangen enten kun yttre sig, naar man havde skrevet en Bog, eller naar man tænkte sig at ville skrive den, saaledes som man overfladeligen tænker sig det, og derfor opkaster det Spørgsmaal, om man skal skrive Forordet først eller sidst. Saasnart man imidlertid er i et af disse Tilfælde, saa har man enten havt en Sag, eller tænker sig at have den. Men naar man nu ogsaa uden dette kan være i Lysten efter at skrive et Forord, saa indsees let, at dette ikke maa handle om en Sag, thi i saa Fald bliver Forordet selv en Bog, og Spørgsmaalet om Forordet og Bogen skubbet tilbage. Forordet som saadant, det emanciperede Forord, maa da ingen Sag have at afhandle, men handle om Intet, og forsaavidt det synes at behandle Noget og handle om Noget, maa dette dog være en Tilsyneladelse og en fingeret Bevægelse.

Hermed er Forordet bestemmet reent lyrisk og bestemmet efter sit Begreb, medens Forordet i vulgair og traditionel Forstand er en Ceremonie efter Tid og Sæd. Et Forord er Stemning. At skrive et Forord er ligesom at hvæsse Leen, ligesom at stemme Guitaren, ligesom at snakke med et Barn, ligesom at spytte ud af Vinduet. Man veed ikke, hvorledes det gaaer til, Lysten kommer paa Een, Lysten efter eventyrligen at zittre i Produktivitetens Stemning, Lysten efter at skrive et Forord, Lysten efter disse leves sub noctem susurri. At skrive et Forord er ligesom at ringe paa en Mands Dør for at gjække ham; ligesom at gaae en ung Piges Vindue forbi og see paa Brostenene, det er ligesom at slaae med sin Stok i Luften efter Vinden, ligesom at svinge med Hatten, uagtet man Ingen hilser. At skrive et Forord er ligesom at have gjort Noget, der berettiger til at fordre en vis Opmærksomhed; ligesom at have Noget paa Samvittigheden, der frister Fortroligheden; ligesom at inclinere i Dandsen, skjøndt man ikke bevæger sig, ligesom at lægge den venstre Schenkel flad an, stramme Tømmen til Høire, høre Gangeren sige: Pst, og selv blæse hele Verden et Stykke, det er ligesom at være med uden at have den mindste Gene af, at man er med, ligesom at staae paa Valdbybakke og see efter Vildgæssene. At skrive et Forord er ligesom at være ankommen med Dagvognen til den første Station, holde i det mørke Skuur, anende, hvad der skal vise sig, see Porten og dermed Himlen aabnet, skue for sig Landeveien, der bestandigen har mere foran sig, øine Skovens forventende Hemmelighed, Fodstiens forføreriske Forsvinden; høre Posthornets Lyd og Echos vinkende Indbydelse, høre Kudskens vældige Smeld, og Skovens forvirrede Gjentagelse, og de Reisendes muntre Samtaler. At skrive et Forord er ligesom at være ankommen, at staae i den hyggelige Stue, hilse paa Længselens forønskede Skikkelse, sidde i Lænestolen, stoppe Piben, tænde den – og have saa uendeligt Meget at tale med hinanden om. At skrive et Forord er ligesom at mærke paa sig selv, at man er ifærd med at blive forelsket, Sjelen sødt uroliget, Gaaden opgiven, enhver Hændelse et Vink om Forklaringen. At skrive et Forord er ligesom at bøie Grenen tilside i Jasminhytten og see hende, der sidder i Løn: min Kjerlighed. Saaledes o! saaledes er det at skrive Forord; og hvorledes er Den, der skriver det? Han gaaer ud og ind mellem Menneskene som en Gjæk i Vinteren og en Nar om Sommeren, han er god Dag og Farvel i een Person, altid glad og sorgløs, fornøiet med sig, ret en letsindig Døgenicht, ja en umoralsk Person, thi han gaaer ikke paa Børsen for at skrabe Penge sammen, men gaaer kun derigjennem; han taler ikke paa Generalforsamlinger, fordi Luften er for indespærret; han udbringer ikke Skaaler i noget Samfund, fordi man skal gjøre Anmeldelse flere Dage iforveien; han render ikke Ærender for Systemet; han afbetaler ikke Statsgjelden, ja han bliver end ikke alvorlig derover; han gaaer gjennem Livet som en Skomagerdreng gaaer fløitende gjennem Gaden, selv om Den, der skal bruge Støvlerne, staaer og venter, saa maa han vente, saa længe der er en eneste Glidebane tilbage, eller det mindste Seeværdige at opdage. Saaledes, ja saaledes er Den, der skriver Forord.

See alt Dette kan nu Enhver tænke over, som han vil, ligesom det falder ham ind og naar det falder ham ind. Med mig er det en anden Sag, thi mig binder et Løfte og en Forpligtelse til ene og alene at beskæftige mig med denne Art af Produktivitet. Hvorledes det hænger sammen hermed, skal jeg strax fortælle Læseren; thi det er netop her paa rette Sted, og ligesom Bagvadskelse i et Kaffeselskab, Noget, der hører ret egentligen hjemme i et Forord.

Skjøndt nemlig lykkelig gift som Faa, og maaskee ogsaa taknemmelig for min Lykke som Faa, er jeg dog i mit Ægteskab stødt paa Vanskeligheder, hvis Opdagelse skyldes min Kone, thi jeg anede Intet. Der var hengaaet nogle Maaneder efter Brylluppet, jeg var efterhaanden bleven nogenlunde indøvet i det ægteskabelige Samlivs Methode, da vaagnede lidt efter lidt igjen en Lyst hos mig, som jeg altid har næret, og som jeg i al Troskyldighed meente, at turde hengive mig til: Beskæftigelsen med et eller andet literairt Arbeide. Gjenstanden var valgt, hvad jeg selv eiede af Bøger i denne Retning fremlagt, enkelte Værker laante fra det kongelige Bibliothek, mine Excerpter ordnede for Overskuelsen, min Pen saa at sige dyppet. Neppe har imidlertid min Kone fattet en Mistanke om, at noget Saadant var i Gjære, før hun omhyggeligen iagttog mine Bevægelser. Leilighedsviis lod hun falde et forblommet Ord, anbragte en dunkel Hentydning til, at min Travlhed paa Studerekamret, min længere Forbliven der, mit literaire Hovedbrud ikke var ganske efter hendes Hoved. Jeg holdt imidlertid Ørene stive, lod som jeg ikke forstod hende, hvad jeg i Begyndelsen virkeligen heller ikke gjorde; da overrumpler hun mig en Dag og afnøder mig den officielle Tilstaaelse, at jeg var ifærd med at ville være Forfatter. Havde hendes Adfærd hidindtil mere været en Recognosceren, saa concentrerede hun sig nu bestemtere og bestemtere, indtil hun endeligen erklærede aabenbar Krig, et quidem saa aabenbar, at hun havde i Sinde at confiskere Alt, hvad jeg skrev, for at benytte det paa en bedre Maade som Underlag under hendes Broderie, til Krøller og saa videre En Forfatters Stilling kan neppe være mere fortvivlet end min; thi selv Den, der staaer under speciel Censur, kan dog haabe, at faae sit Arbeide bragt dertil, at det »maa trykkes;« men min Frembringelse qvæles bestandig i Fødselen. Ogsaa paa en anden Maade blev det mig tydeligere og tydeligere, hvor fortvivlet min Stilling var; thi neppe havde jeg opdaget, at jeg var bleven Gjenstand for en Presseforfølgelse, førend, som naturligt er, det blev mig klart, hvad tidligere aldeles ikke var faldet mig ind, at det vilde være til ubodelig Skade for Menneskeheden, om mine Skrifter ikke kom for Dagen. Hvad er nu herved at gjøre? Jeg har ikke som en censureret Forfatter Regres til Cancelliet, Stænderforsamlingen, det høistærede Publikum, eller til en Efterverdens Ihukommelse, jeg lever og døer, staaer og falder med min Kone. Nu er jeg vel af mine Samtidige anseet for en god og meget øvet Disputator, der nok kan plaidere min Sag, men denne Dygtighed kommer mig kun her lidet til Baade, thi om jeg end kan disputere med Fanden selv, med min Kone kan jeg ikke disputere. Hun har nemlig kun een Syllogisme, eller rettere ikke engang een. Hvad lærde Folk kalde Sophisterie, det kalder hun, der ikke befatter sig med at være lærd, Drillerie. Fremgangsmaaden er nu ganske simpel, at sige for Den, der veed at gaae rigtigt frem. Hver Gang jeg siger Noget, som ikke behager hende, hvad enten det er i Syllogismens Form eller ikke, en lang Tale eller en kort Bemærkning, Formen er ligegyldig, men naar det Udsagte ikke behager hende, da seer hun paa mig med en Mine, der paa engang er elskværdig, indtagende, godmodig, fortryllende, men ogsaa triumpherende, tilintetgjørende og siger: det er bare Drillerie. Følgen heraf er, at al min Færdighed i at disputere bliver en Luxus-Artikel, der slet ikke spørges i mit huslige Liv. Kan jeg, den øvede Dialektiker, nogenlunde repræsentere Rettens Gang, der, efter Digterens Ord, er saare lang, saa er min Kone ligesom det kongelige danske Cancellie kurz und bündig, kun er hun deri forskjellig fra hiint høie Collegium, at hun er meget elskværdig; thi det er netop denne Elskværdighed, der hjemler hende en Myndighed, som hun hvert Øieblik veed at hævde paa en fortryllende Maade.

Saaledes staaer Sagen. Jeg er aldrig kommen længere end til en Indlednings-Paragraph. Da denne var af almindelig Natur og i mine Tanker saa lykkelig affattet, at den vilde more hende, hvis jeg ikke var Forfatteren, faldt det mig ind, om jeg ikke skulde kunne vinde hende for Sagen ved at læse den for hende. Jeg var forberedt paa, at hun vilde afvise mit Andragende, og benytte sig af den Fordeel, »at det vel nu endog skulde komme saa vidt, at jeg ikke blot gav mig af med at skrive, men hun skulde forpligtes til at høre Forelæsninger.« Ingenlunde. Hun tog saa venligt mod mit Forslag som muligt, hun hørte, hun loe, hun beundrede. Jeg troede Alt vundet. Hun traadte hen til Bordet, ved hvilket jeg sad, lagde fortroligen sin Arm om min Hals, bad mig, endnu engang at forelæse en Passus. Jeg begynder at læse, holder Manuscriptet saa høit op, at hun med Øiet kunde følge mig. Ypperligt. Jeg er ude af mig selv, men endnu ikke ganske ude af hiin Passus, da Manuscriptet pludselig staaer i lys Lue. Uden at jeg bemærkede det, havde hun skudt det ene Lys ind under Manuscriptet. Ilden havde Overmagten, der var Intet at redde, min Indlednings-Paragraph gik i Luerne – under almindelig Jubel, thi min Kone jublede for os Begge; som et overgivent Barn klappede hun i Hænderne, og derpaa kastede hun sig om min Hals med en Lidenskab, som havde jeg været adskilt fra hende, ja tabt for hende. Jeg kunde ikke faae et Ord indført. Hun bad mig om Tilgivelse, at hun saaledes havde kæmpet for sin Kjerlighed, bad med en Bevægethed, der næsten bragte mig selv til at troe, at jeg havde været ifærd med at blive den forlorne Ægtemand. Hun forklarede, at hun ikke kunde udholde, at jeg saaledes forandredes. »Din Tanke, sagde hun, tilhører mig, den maa tilhøre mig; Din Opmærksomhed er mit daglige Brød, Dit Bifald, Dit Smiil, Din Spøg er mit Liv, min Begeistring; tilstaae mig den, o! negt mig ikke, hvad der tilkommer mig med Rette; for min Skyld, for min Glædes Skyld, at jeg med Glæde maa kunne gjøre, hvad der er min eneste Glæde, at tænke paa Dig, og finde al min Tilfredsstillelse i at kunne Dag ud og Dag ind vedblive at beile til Dig, som Du engang beilede til mig.«

Hvad berettiger nu en Kone til saadan Adfærd, en Kone, der ikke blot er elskværdig i alle Deres Øine, som kjende hende, men fremfor Alt elskelig i mine Øine, liflig saa lang Dagen er. Hendes Betragtning er in contento følgende: en Ægtemand, der er Forfatter, er ikke stort bedre end en Ægtemand, der hver Aften gaaer i Klubben, ja endogsaa værre; thi Den, der gaaer i Klubben, han maa dog selv tilstaae, at det er et Brud, men det at være Forfatter er en fornem Utroskab, der ikke kan vække Fortrydelse, om Følgerne end ere værre. Den, der gaaer i Klubben, han er dog kun fraværende saa længe han er fraværende, men en Forfatter – »ja Du veed det formodentlig ikke selv; men der er foregaaet en heel Forandring med Dig. Du spinder Dig ind i Tankefuldhed fra Morgen og til Aften, og især mærker man det ved Middagsbordet. Der sidder Du og seer hen for Dig som en Geist eller som Keiser Nebucadnezar, der læser den usynlige Skrift. Naar jeg da selv har tilberedet Kaffen til Dig, har sat den paa Bakken, kommer glad hen til Dig, staaer foran Dig, neier for Dig – saa, saa er jeg nær ved at tabe Bakken af Forskrækkelse, og fremfor Alt, saa har jeg tabt min Munterhed og min Glæde og kan ikke neie for Dig.«

Som nu min Kone ved enhver Leilighed veed at faae anbragt sit catoniske præterea censeo, om hun end ikke gjør det saa kjedeligt som Cato, saa maa ogsaa Alt tjene hende til Argument. Hendes Argumentation er ligesom en Naturbesværgelse. Hvis det hændte mig ved en Disputats, at Modparten fremkom med lignende Argumenter, da vilde jeg formodentligen vende ham Ryggen, og sige om ham, hvad Magisteren siger hos Holberg: en Ignorant, der ikke veed at distinguere mellem ubi prædicamentale og ubi transcendentale. Med min Kone er det en anden Sag, hendes Argumentation gaaer frisk væk fra Haanden – og til Hjertet, hvorfra den da egentligen kommer. Hun har i denne Henseende lært mig at forstaae, hvorledes en Katholik kan opbygges ved en latinsk Gudstjeneste; thi hendes Argumentation er, betragtet som saadan, hvad Latin er for Den, der ikke forstaaer det, og dog opbygger hun altid, bevæger og rører mig.

»At være Forfatter, naar man er Ægtemand, siger hun, er aabenbar Utroskab, lige tvertimod hvad Præsten sagde; thi dette er Ægteskabets Gyldighed, at Manden holder hart ved sin Hustrue og saa ikke videre.« Svarer jeg dertil, at hun nok ikke har hørt rigtigt efter, hvad Præsten sagde, at man næsten skulde troe, hun var saa forsømt, at hun maatte gaae om igjen til Præsten; at Ægteskabet er en særlig Pligt, ja en »besynderlig« Pligt, samt at alle Pligter kunne inddeles i de almindelige og de besynderlige, og ere Pligter mod Gud, os selv og mod Næsten, saa kommer hun slet ikke i Forlegenhed. Det Hele bliver erklæret for Drillerie, og »forøvrigt har hun ikke glemt hvad der staaer i Lærebogen om Ægteskabet, at det er Mandens Pligt i Særdeleshed.« Forgjeves søger jeg at oplyse, at hun er i en sproglig Vildfarelse, at hun construerer hine Ord ulogisk, ugrammatikalsk, mod alle Fortolknings-Principer, da der paa hiint Sted kun tales om Mandens Pligter i Særdeleshed med Hensyn til Ægteskabet, ligesom i den næste Paragraph om Konens i Særdeleshed. Forgjeves. Hun henholder sig til sit Foregaaende, »at det at være Forfatter, naar man er Ægtemand, er den værste Art af Utroskab.« Nu er det endog bleven den »værste« Utroskab. Vil jeg da bringe hende i Erindring, at Manden ifølge alle guddommelige og menneskelige Love er den Herskende, at min Ansættelse i Livet i andet Fald bliver saare ringe, da jeg kun bliver et encliticon til hende, hvilket dog er at fordre for meget, saa bebreider hun mig min Ubillighed, »da jeg godt veed, at hun Intet fordrer, at hendes eneste Begjering er i Forhold til mig at være slet Intet.« Protesterer jeg herimod, fordi, naar jeg endelig kun skal være et encliticon, det bliver mig af Vigtighed, at hun bliver saa meget som muligt, at jeg ikke skal blive endnu mindre ved at være encliticon til Intet, saa seer hun paa mig og siger: bare Drillerie.

Consequent er min Kone, fix i sin Idee. Jeg har søgt at smigre hende, at det dog var fornøieligt, at see mit, vort Navn berømmet; at hun er Musen, der begeistrer mig. Hun vil Intet høre. Det Første anseer hun for den største Ulykke og min fuldkomne Fortabelse, da hun af sit ganske Hjerte vilde ønske, at en eftertrykkelig Critik viste mig hjem igjen. Det Sidste troer hun ikke, ønsker det endnu mindre, og beder af sin ganske Sjel Gud at forbyde, at hun saaledes selv skulde forskylde Tabet af sin ægteskabelige Lykke. Hun er utilgængelig, og summa summarum, »naar Alt er hørt,« bliver denne: »enten, siger hun, en ret Ægtemand – eller ogsaa..... ja saa er Resten ligegyldig.«

Medens nu Læseren vel ligesom jeg vil finde hendes Argumentation noget svag, og at hun aldeles overseer alle de egentligen omqvæstionerede Problemer, Grændsestridighederne nemlig mellem det Ægteskabelige og det Individuelle, hvilke kunne give et ligesaa dybsindigt som skarpsindigt Hoved nok at bestille, saa har hun endnu et Argument in subsidio, hvilket Læseren maaskee vil give mere Vægt. Efterat vi en Dag havde gjennemkæmpet vort Mellemværende, og Striden, som sædvanligen, forklaret sig i en redintegratio amoris, tog hun mig tilsidst fortroligen under Armen, saae saa indsmigrende paa mig som muligt og sagde: »Kjære! jeg har ikke villet sige Dig det ganske uforbeholdent, fordi jeg haabede paa anden Maade at faae Dig til at afstaae fra Dit Forehavende, og haabede at kunne spare Dig for en Krænkelse, men da det ikke skal lykkes, saa vil jeg sige Dig det med al den Oprigtighed, Du kan fordre af Din Kone: jeg troer ikke, Du duer til at være Forfatter – men derimod, ja lee nu kun lidt ad mig, men derimod har Du Genie og Talent og overordentlig Begavelse til at blive min Mand saaledes, at jeg uafbrudt vilde beundre Dig, medens jeg selv glad følte min egen Ringhed og kun lod min Kjerlighed fremkomme for Dig med Taksigelse.« Argumentet indlod hun sig imidlertid ikke paa nærmere at udføre. Saasnart jeg her vilde indlade mig paa et om, hvorvidt, og hvorledes, da havde hun en anden Forklaring, »at det engang vilde fortryde mig, at have været hende utro ved at blive Forfatter, og da vilde jeg ikke due til at slaae Fortrydelsen hen, men lide dens Bitterhed.«

Og hvad blev saa Enden paa denne Strid; hvo seirede, min hostis domesticus eller Forfatteren? Det er vel ikke vanskeligt at gjette, om det end bliver Læseren et Øieblik vanskeligt, da han jo læser dette, og altsaa seer, at jeg blev Forfatter. Enden blev, at jeg lovede ikke at ville være Forfatter. Men som man ved lærde Disputatser, naar Autor har afvæbnet alle Eens Indvendinger, tilsidst fremkommer med en eller anden linguistisk Ubetydelighed for dog at faae Ret i Noget, og Autor høfligen giver Een Ret, for dog at give Een Ret i Noget, saaledes forbeholdt jeg mig Tilladelse til at turde skrive »Forord.« Jeg beraabte mig i denne Henseende paa Analogier, at Mænd, der havde lovet deres Koner aldrig at snuse Tobak mere, havde som Vederlag faaet Lov til at have saa mange Tobaksdaaser, som de vilde. Hun tog mod Forslaget, maaskee i den Mening, at man ikke kunde skrive et Forord uden at skrive en Bog, hvilket jeg jo ikke tør; med mindre man er en berømt Forfatter, der efter Anmodning skriver et saadant, hvilket jo umuligen kunde blive mit Tilfælde.

Dette angaaende mit Løfte og min Forpligtelse. Det Lidet eller de Smaating, som jeg herved udgiver, har jeg kunnet skrive salva conscientia. Dog har jeg gjort det min Kone uafvidende, idet jeg har benyttet et Ophold paa Landet dertil. Min Bøn til Kritiken er, at den vil fare lempelig frem med mig; thi sæt den fandt, at det var som min Kone sagde, at jeg ikke duede til at være Forfatter; sæt den ubarmhjertigen gjennemheglede mig, sæt min Kone fik det at vide, saa vilde jeg nok forgjeves søge Opmuntring og Trøst hos mit Livs Ledsagerinde. Formodentlig vilde hun juble af Glæde over, at faae sin Krig frem, og over, at jeg saaledes var bleven tagen i Skole, finde sin Tro paa den retfærdige Styrelse bekræftet, sin Idee bestyrket, at det at være Forfatter, naar man er Ægtemand, er den værste Utroskab.

\workheading{Tre Taler ved tænkte Leiligheder}{1845}

\worknote{Antologiens ''Three Discourses on Imagined Occasions''. Hong-udvalget bringer et uddrag af tredje tale, ''Ved en Grav'' — den tænkte begravelse af en stille borger, og dernæst betragtningen over alvoren og tanken om døden (''endnu idag''). Store spring markeret med […].}

\sectionheading{Ved en Grav (uddrag)  ·  SV¹ V, 226–253}

Saa er det da forbi! – Og naar nu Den, som her først traadte hen til Graven, fordi han er den Nærmeste, naar han efter Talens korte Øieblik er bleven den Sidste ved Graven, ak, fordi han er den Nærmeste: saa er det forbi. Om han vilde blive derude, han erfarer dog ikke hvad den Døde foretager sig, thi den Døde er en stille Mand; om han i sin Bekymring vilde kalde paa hans Navn, om han i sin Sorg vilde sidde lyttende, han erfarer dog Intet, thi i Graven er der Stilhed, og den Døde er en taus Mand; og om han erindrende gik hver Dag til hans Grav, den Døde erindrer ikke ham –.

Thi i Graven er der ingen Erindring, end ikke af Gud. See, dette vidste den Mand, om hvem der nu maa siges, at han ikke mere erindrer Noget, til hvem det nu var for silde at sige det. Men som han vidste det, saa gjorde han derefter, og derfor erindrede han Gud medens han levede. Hans Liv gik hen i hæderlig Ubemærkethed, ikke Mange vare vidende om hans Tilværelse, kun Enkelte blandt de Faa kjendte ham. Han var Borger her af Staden, stræbsom i sin beskedne Gjerning forstyrrede han Ingen ved at svigte Samfundets Forpligtelser, forstyrrede Ingen ved utidig Bekymring om det Hele. Saaledes gik Aar efter Aar, eensformigt, dog ikke tomt; han blev Mand, han blev gammel, han blev bedaget: Gjerningen var og blev den samme, een og samme Syssel i de forskjellige Levealdre. Han efterlader sig en Hustru, glad fordum ved at forenes med ham, nu en gammel Kone, der begræder den Tabte, en ret Enke, der, forladt, har sit Haab til Gud. Han efterlader sig en Søn, der lærte at elske ham og at finde Tilfredshed i Vilkaaret og Faderens Gjerning; engang glad som Barn i Faderens Huus, fandt han det som Yngling aldrig for trangt, nu er det ham et Sørgehuus. – En saadan ubemærket Mands Død spørges ikke vidt, og naar Nogen om føie Tid gaaer forbi det Huus, hvor han boede i Ringhed, og læser hans Navn over Døren, fordi den borgerlige Bedrift fortsættes under hans Navn, da er det jo, som var han ikke død. Som han sov mildt og fredeligt hen, saaledes er i den omgivende Verden hans Død en Bortgang i Stilhed. Brav som Borger, redelig i sin Næringsvei, sparsom i sit Huus, godgjørende efter Evne, deeltagende i Oprigtighed, sin Hustru tro, sin Søn en Fader: see, alt Dette, og al den Sandhed hvormed det her kan siges, spænder ikke Forventningen paa en betydningsfuld Udgang, her er det et Livs Foretagende, paa hvilket en rolig Død blev den skjønne Udgang. – Dog havde han een Gjerning endnu, den blev udført med samme Troskab i Hjertets Eenfoldighed: han erindrede Gud. Han var Mand, gammel, han blev bedaget, saa døde han, men Erindringen om Gud blev den samme, en Veiledning i alt hans Foretagende, en stille Glæde i den fromme Betragtning. Ja, var der slet Ingen, som savnede ham i Døden, ja, var han nu ikke hos Gud, da vilde Gud savne ham i Livet, og vide hans Bopæl og søge ham der, thi den Afdøde vandrede for hans Aasyn og var bedre kjendt af ham end af nogen Anden. Han erindrede Gud, og han blev duelig til sit Arbeide, han erindrede Gud og han blev glad ved sit Arbeide og glad ved Livet, han erindrede Gud, og han blev lykkelig i sit beskedne Hjem med sine Kjære, han forstyrrede Ingen ved Ligegyldighed mod en offentlig Gudsdyrkelse, han forstyrrede Ingen ved ubetimelig Iver, men Guds Huus var ham et andet Hjem – nu er han vandret hjem.

Men i Graven er der ingen Erindring – derfor bliver den tilbage, den bliver hos de Tvende, der vare ham kjære i Livet: de ville erindre ham. Og naar nu Den, som traadte først hen til Graven, fordi han er den Nærmeste, efter Talens korte Øieblik er bleven den Sidste ved Graven, fordi han jo er den Nærmeste, naar han gaaer erindrende herfra, da gaaer han hjem til den sørgende Enke: og Navnet over Døren bliver en Erindring. Saa vil der i nogen Tid nu og da komme en Kjøber, der tilfældigt eller mere deeltagende spørger efter Manden; og naar han hører om hans Død vil Kjøberen sige: naa, er han død. Og naar saa alle de gamle Kjøbere har gjort det eengang, saa har Omgivelsens Liv intet Middel mere for at bevare Erindringen om ham. Men den gamle Enke vil ingen Paamindelse behøve for at erindre, og den stræbsomme Søn vil ikke finde det Sinkende at erindre. Og naar saa Ingen spørger mere om ham, da vil dog Navnet over Døren, naar Huset ikke synligt er et Sørgehuus, naar ogsaa i Huset Sorgen er formildet og det daglige Savn med Trøsten har indøvet Erindringen: da vil Navnet over Døren for de Tvende betyde, at ogsaa de have een Gjerning mere: at erindre den Afdøde.

Nu er Talen forbi. Kun en Handling staaer tilbage: med de tre Spader Jord at indvie den Afdøde, som Alt hvad der er kommen af Jord, til Jord igjen – og saa er det forbi.

[…]

Er det vist, at Døden er til, som det er; er det vist, at det med dens Afgjørelse er forbi; er det vist at Døden selv aldrig indlader sig paa at give nogen Forklaring: nu vel, da gjelder det at forstaae sig selv, og Alvorens Forstaaelse er, at er Døden Natten, saa er Livet Dagen, kan der ikke arbeides om Natten, saa kan der arbeides om Dagen; og Alvorens korte men tilskyndende Raab, ligesom Dødens korte, er: endnu idag. Thi Døden i Alvor giver Livskraft som intet Andet, den gjør aarvaagen som intet Andet. Det sandselige Menneske virker Døden paa, til at sige: lader os æde og drikke, thi imorgen skulle vi døe; men dette er Sandselighedens feige Livslyst, hiin foragtelige Tingenes Orden, hvor man lever for at spise og drikke, ikke spiser og drikker for at leve. Det dybere Menneske virker Dødens Forestilling maaskee paa til Afmagt, saa han segner slappet i Stemning; men den Alvorlige giver Dødens Tanke den rette Fart i Livet og det rette Maal, hvorefter han retter Farten. Og ingen Buestræng kan strammes saaledes og formaaer at give Pilen en saadan Fart, som Dødens Tanke formaaer at fremskynde den Levende, naar Alvoren spænder den. Da griber Alvoren det Nærværende endnu idag, forsmaaer ingen Opgave som for ringe, vrager ingen Tid som for kort, arbeider af yderste Evne, om den dog er villig til at smile over sig selv, hvis denne Anstrængelse skulde være Fortjeneste for Gud, og villig til i Afmagten at forstaae, at et Menneske er slet Intet, og at Den, der arbeider af yderste Evne, kun faaer ret Leilighed til at forundre sig over Gud.

Saa lad da Døden beholde sin Magt, »at det er forbi«, men Livet ogsaa beholde sin Ret til at arbeide medens det er Dag; og lad den Alvorlige søge Dødens Tanke som Bistand dertil. Den Vankelmodige er blot et Vidne til den bestandige Grændsestrid mellem Liv og Død, hans Liv kun Tvivlens Angivelse af Forholdet, hans Livs Udgang en Skuffelse, men den Alvorlige har sluttet Venskab med de Stridende, og hans Liv har i Dødens alvorlige Tanke den meest trofaste Forenede. Er der da end den Lighed for alle Døde, at det nu er forbi, der er dog een Forskjel, min Tilhører, den raaber til Himlen denne Forskjel, Forskjellen paa, hvilket Liv det var, som nu med Døden er forbi. Saaledes er det altsaa ikke forbi, og trods al Dødens Forfærdelse, nei, understøttet af Dødens alvorlige Tanke siger den Alvorlige: det er ikke forbi. Men frister denne lyse Udsigt, skimter han atter blot efter den i Betragtningens Dæmring, fjerner den ham fra Opgaven, bliver Tiden ikke Dyrtid, er Besiddelsen ham tryg; da er han atter ikke alvorlig. Siger Døden: maaskee endnu idag; da siger Alvoren: det være nu maaskee idag eller ikke, jeg siger: endnu idag.

[…]

Den Alvorlige betragter sig selv; er han ung da lærer Tanken om Døden ham, at det er et ungt Menneske, som her bliver dens Bytte, hvis den kommer idag, men han fjanter ikke i almindelig Tale om Ungdom som Dødens Bytte. Den Alvorlige betragter sig selv, han veed altsaa, hvorledes Den er beskaffen, som her vilde blive Dødens Bytte, hvis den kom idag; han betænker sin egen Gjerning og veed altsaa, hvilken Gjerning det var, som her afbrødes, hvis Døden kom idag. Saa hører Legen op, saa er Gaaden gjettet. Den almindelige Betragtning af Døden forvirrer kun Tanken, ligesom det at ville erfare i Almindelighed. Dødens Vished er Alvoren, dens Uvished er Underviisningen, Indøvelsen af Alvoren; den Alvorlige er Den, som ved Uvisheden opdrages til Alvor i Kraft af Visheden. Hvorledes lærer vel et Menneske Alvor? Mon derved, at en Alvorlig foresiger ham Noget, at han kunde lære dette? Ingenlunde. Har Du ikke selv saaledes lært af en alvorlig Mand, saa tænk Dig, hvorledes det gaaer til. See, den Lærende bekymrer sig (thi uden Bekymring ingen Lærende) om en eller anden Gjenstand med hele sin Sjel; og saaledes er jo Dødens Vished en Bekymringens Gjenstand. Nu henvender den Bekymrede sig til Alvorens Lærer; og saaledes er jo Døden, ikke en Skræmsel uden for Indbildningen. Den Lærende vil nu Dette eller Hiint, han vil gjøre det saaledes under disse Forudsætninger: »og ikke sandt, saa lykkes det?« Men den Alvorlige svarer slet Intet, og endeligen siger han, dog uden at spotte, med Alvorens Rolighed: »ja, det er muligt!« Den Lærende bliver allerede lidt utaalmodig; han udkaster en ny Plan, forandrer Forudsætningerne, og slutter sin Tale paa en endnu mere indtrængende Maade. Men den Alvorlige tier, seer roligt paa ham og siger endeligen: »ja, det er muligt!« Nu bliver den Lærende lidenskabelig, han griber til Bønner, eller hvis han er saaledes udrustet til snilde Tankevendinger, ja han fornærmer maaskee endog den Alvorlige, og bliver selv ganske forvirret, og Alt synes Forvirring om ham; men da han med disse Vaaben og i denne Tilstand stormer ind paa den Alvorlige, maa han udholde hans uforandrede rolige Blik og finde sig i hans Taushed, thi den Alvorlige seer blot paa ham og siger endeligen: »ja det er muligt.« Saaledes med Døden. Visheden er den Uforanderlige, og Uvisheden er det korte Ord: det er muligt; og enhver Betingelse, der vil gjøre Dødens Vished til en betinget Vished for den Ønskende, enhver Overeenskomst, der vil gjøre Dødens Vished til en betinget Vished for den Besluttende, enhver Aftale, der vil betinge Dødens Vished ved Tid og Time for den Handlende, enhver Betingelse, enhver Overeenskomst, enhver Aftale strander paa dette Ord; og al Lidenskabelighed og al Snildhed og al Trods gjøres afmægtig ved dette Ord, indtil den Lærende gaaer i sig selv. Men netop deri ligger Alvoren, og netop dertil var det, at Visheden og Uvisheden vilde hjælpe den Lærende. Faaer Visheden Lov at staae hen for hvad den nu kan være, som en almindelig Overskrift over Livet, ikke, som det skeer ved Uvishedens Hjælp, som Anvendelsens Paategning paa det Enkelte og Daglige, da læres Alvoren ikke. Uvisheden træder til og peger stadigen som Læreren paa Lærens Gjenstand, og siger til den Lærende: læg vel Mærke til Visheden: da bliver Alvoren til. Og ingen Lærer formaaer saaledes at lære Discipelen at lægge Mærke til hvad der siges, som Dødens Uvished, naar den peger paa Dødens Vished; og ingen Lærer formaaer saaledes at holde Discipelens Tanker samlede paa Underviisningens ene Gjenstand, som Tanken om Dødens Uvished, naar den indøver Tanken om Dødens Vished.

[…]

Den som her har talt, han er ung, endnu i den Lærendes Alder; han fatter kun Underviisningens Vanskelighed og Strenghed, o, at det maatte lykkes ham at gjøre det saaledes, at han netop derved blev værdig til engang at turde glæde sig ved Lærerens Venskab! Den som her har talet, han er jo ikke Din Lærer, min Tilhører, han lader Dig jo blot, ligesom han selv er det, være Vidne til, hvorledes et Menneske søger at lære Noget af Tanken om Døden, hiin Alvorens Læremester, der ved Fødselen blev Enhver beskikket til Lærer for hele Livet, og som i Uvisheden altid er rede til at begynde Underviisningen, naar den forlanges. Thi Døden kommer ikke, fordi Nogen kalder paa den (det var kun Spøg, at saaledes den Svagere bød over den Stærkere), men saasnart Nogen aabner Uvisheden Indgang, saa er Læreren der. Læreren, som engang kommer for at holde Prøve og overhøre Discipelen, hvad enten han har villet nytte hans Underviisning eller ikke. Og denne Dødens Prøve, eller for med et mere brugt fremmed Ord at betegne det Samme, denne Livets sidste Examen er lige svær for Alle. Det er ikke som ellers, at den lykkeligt Begavede har let ved at bestaae, den ringe Begavede svært, nei, Døden afpasser Prøven i Forhold til Evnen, o, saa nøiagtigt, og Prøven bliver lige svær, fordi det er Alvorens Prøve.

\workheading{Stadier paa Livets Vei}{1845}

\worknote{Antologiens ''Stages on Life's Way'' (Studier af Forskjellige, udgivne af Hilarius Bogbinder). Hong-udvalget bringer store, adskilte uddrag: fra ''In vino veritas'' (William Afham) motto og Forerindringens betragtning over Erindring og Hukommelse; Gjestebudets iscenesættelse; Modehandlerens tale om Kvinden og Moden; og fra Frater Taciturnus' afsluttende ''Til Læseren'' bestemmelsen af de tre Existents-Sphærer samt Slutningsordet (sophisterne, Modersmaalet, Danmark). Alle spring markeret med […].}

\sectionheading{»In vino veritas«: Forerindring (William Afham)  ·  SV¹ VI, 14–25}

Solche Werke sind Spiegel; wenn ein Affe hinein guckt, kann kein Apostel heraus sehen.

— Lichtenberg

Hvilken skjøn Syssel at tilberede sig en Hemmelighed, hvor forførerisk at nyde den, og dog hvor betænkeligt stundom at have nydt den, hvor let, at den ikke bekommer Een vel! Dersom nemlig Nogen troer, at en Hemmelighed er uden videre transportabel, at den kan tilhøre en Ihændehaver, da feiler han, thi det gjælder her: fra den Spisende kommer Spisen; men dersom Nogen troer, at man ved Nydelsen kun paadrager sig den Vanskelighed ikke at forraade den, da feiler han ogsaa, thi man forskylder jo tillige det Ansvar ikke at glemme den. Dog er det end modbydeligere at erindre halvt og forvandle sin Sjel til et Transit-Oplag for havareret Gods. I Forhold til Andre være da Glemsel det Silkegardin, der drages for, Erindringen den vestalske Jomfru, som gaaer ind bag Gardinet; bag Gardinet er Glemselen atter, hvis det ikke er en ret Erindring, thi da er Glemselen udelukket.

Erindringen maa ikke blot være nøiagtig, den maa ogsaa være lykkelig; Erindringens Afpropning maa have gjemt det Oplevedes Duft inden den forsegles. Som Druen ikke lader sig perse til enhver Tid, som Veierliget i Persningens Tid har stor Indflydelse paa Vinen, saaledes lader det Oplevede sig heller ei erindre eller inderindres til enhver Tid, eller under enhver Omgivelse.

At erindre er ingenlunde identisk med at huske. Man kan saaledes godt huske en Begivenhed til Punkt og Prikke uden derfor at erindre den. Hukommelsen er kun en forsvindende Betingelse. Ved Hukommelsen stiller det Oplevede sig frem for at modtage Erindringens Indvielse. Forskjellen skjønnes allerede i Menneskealdernes Forskjellighed. Oldingen taber Hukommelsen, der overhovedet er den Evne, som tabes først. Dog har Oldingen noget Digterisk, han er i Folkeforestillingen prophetisk, gudbeaandet. Men Erindringen er jo ogsaa hans bedste Kraft, hans Trøst, der trøster ham med det poetiske Fjernsyn. Barndommen har omvendt i høi Grad Hukommelse og Nemme, slet ikke Erindring. Istedenfor at sige: Alderdommen glemmer ikke hvad Ungdommen nemmer, kunde man maaskee sige: hvad Barnet husker, erindrer Oldingen. Oldingens Brille slibes til at see nærved med. Naar Ungdommen bruger Briller er Glasset til at see paa Afstand, thi den mangler Erindringens Kraft, hvilken er: at fjerne, at bringe paa Afstand. Imidlertid er Alderdommens lykkelige Erindring ligesom Barnets lykkelige Nemme Naturens Naadegave, der med Forkjærlighed omfatter de tvende mest hjælpeløse og dog i en vis Forstand lykkeligste Afsnit af Livet. Men derfor er ogsaa Erindringen saavelsom Hukommelsen stundom kun Ihændehaver af Tilfældigheder.

Skjøndt Forskjellen mellem Hukommelse og Erindring er stor, forvexles de ofte. I Menneskelivet giver denne Forvexling Leilighed til at studere Individets Dybde. Erindringen er nemlig Idealiteten, men som saadan ganske anderledes anstrængende og ansvarlig end den ligegyldige Hukommelse. Erindringen vil hævde et Menneske den evige Continuitet i Livet og sikkre ham, at hans jordiske Tilværelse bliver uno tenore, eet Aandedrag, og udsigeligt i Eet. Derfor frabeder den sig, at Tungen skal nødsages til atter og atter at slaae Sladder for at eftergjøre Livs-Indholdets Sladder. Dette er Betingelsen for Menneskets Udødelighed, at Livet er uno tenore. Besynderligt nok er, saavidt jeg veed, Jacobi den Eneste, hos hvem man finder Yttringer om det Forfærdelige i at tænke sig selv udødelig. Det var ham stundom, som skulde Udødelighedens Tanke, hvis han fastholdt den lidt længer i det enkelte Øieblik, forvirre ham Forstanden. Er Grunden den, at Jacobi var nervesvag? En stærk Mand, der har faaet haard Hud i Haanden alene ved at slaae i Prædikestolen eller Kathedret, hver Gang han beviste Udødeligheden, føler ingen saadan Forfærdelse, og dog forstaaer han sig jo paa Udødeligheden, thi at have haard Hud betyder paa Latin at forstaae sig paa Noget tilgavns. Saasnart man imidlertid forvexler Hukommelse og Erindring, bliver Tanken ikke saa forfærdelig. For det Første fordi man er modig, mandhaftig og robust, og for det Andet fordi man slet ikke tænker Tanken. Der er saa vist mangen Mand, som har skrevet Erindringer af sit Liv, hvori der ikke var Spor af Erindring, og dog vare jo Erindringerne hans Udbytte for Evigheden. I Erindringen trækker Mennesket paa det Evige. Det Evige er humant nok til at honorere enhver Fordring og ansee Enhver for vederhæftig. Men det Evige kan ikke for, at et Menneske gjør sig selv til Nar – og husker istedenfor at erindre og som en Følge deraf glemmer istedenfor at erindre, thi hvad der huskes det glemmes ogsaa. Men Hukommelsen gjør igjen Livet ugeneret. Ugeneret gjennemgaaer man de latterligste Metamorphoser; selv i fremrykkede Aar leger man endnu Blindebuk, spiller endnu i Livets Lotteri og kan endnu blive til hvad det skal være, uagtet man har været utrolig mange Ting. Derpaa døer man – og saa bliver man udødelig. Og skulde man dog ikke ogsaa netop ved at have levet saaledes have sikkret sig rigelig Nok at erindre en heel Evighed? Ja, hvis Erindringens Hovedbog ikke var andet end en Kladde, hvori man smører det Første det Bedste ind. Men Erindringens Bogholderi er besynderligt. Man kunde sætte nogle saadanne Stykker frem til Opgave – dog ikke i Selskabsregning. En Mand taler Dag ud og Dag ind paa Generalforsamlinger og bestandigt om hvad Tiden fordrer, dog ikke catonisk-kjedsommelig ved Gjentagelsen, men bestandig interessant og pikant følger han med Øieblikket og siger aldrig det Samme, item i Selskab kræver han sig og udmaaler sit Veltalenheds Forraad snart med Stryg- snart med Top-Maal, bestandigt hilset med Haandklap, eengang om Ugen idetmindste er der Noget at læse om ham i Bladet, endog om Natten gavner han, sin Kone nemlig, ved selv i Søvne at tale om Tidens Fordring som om han var paa Generalforsamling; en anden Mand tier før han taler, og bringer det saa vidt, at han slet ikke kommer til at tale: de leve lige længe, her spørges om Facit: hvo har mest at erindre? En Mand forfølger een Tanke, een eneste, alene beskæftiget med den; en anden er Forfatter i syv Videnskaber og »afbrydes i sin betydningsfulde Virksomhed (det er en Journalist, som taler) just som han vilde omskabe Veterinair-Videnskaben«: de leve lige længe, her spørges om Facit: hvo har mest at erindre?

Egentlig kan man kun erindre det Væsentlige, thi Oldingens Erindren er, som sagt, underlagt Tilfældighed; om Analogierne til hans Erindren gjelder det Samme. Det Væsentlige betinges ikke blot ved sig selv, men ogsaa ved sit Forhold til den Paagjeldende. Den, der har brudt med Ideen, kan ikke handle væsentlig, kan intet Væsentligt foretage sig; det skulde da være at angre, der er den eneste ny Idealitet. Alt, hvad han ellers gjør, er uvæsentligt, trods de ydre Kjendetegn. At tage sig en Kone er jo noget Væsentligt; men Den, som engang har jadsket i Elskov, han kan gjerne slaae sig for Panden og for Hjertet og paa R– af lutter Alvor og Høitidelighed; det bliver dog Fjas. Om hans Ægteskab angik et heelt Folk, og der blev ringet med Klokkerne, og Paven holdt Vielsen, det er dog ikke for ham noget Væsentligt, men væsentligen Fjas. Den ydre Larmen gjør slet Intet fra eller til, saa lidet som Fanfaren, og at der præsenteres Gevær gjør Tallotteriets Trækning til en væsentlig Handling for den Dreng, som trækker Lodderne. For at der skal handles væsentligen kommer det ikke væsentligen an paa, at der slaaes paa Tromme. – Men det Erindrede kan man heller ei glemme. Det Erindrede er ikke, som det Huskte er ligegyldigt mod Hukommelsen, ligegyldigt mod Erindringen. Det Erindrede kan man kaste bort, det vender tilbage ligesom Thors Hammer, og ikke blot saaledes, men det har en Længsel efter Erindringen som en Due, ja som den Due, der hvor ofte den end sælges til Andre aldrig kan blive en Andens Eiendom, fordi den bestandigt flyver hjem. Men saa har ogsaa Erindringen selv udruget det Erindrede, og denne Rugen er skjult og forborgen, lønlig og derfor ukrænket af nogen profan Viden: saaledes vil Fuglen ikke ruge over sit Æg, naar nogen Fremmed har berørt det.

Hukommelsen er umiddelbar og kommes umiddelbart tilhjælp, Erindringen kun reflecteret. Derfor er det en Kunst at erindre. I Modsætning til at huske ønsker jeg med Themistocles at kunne glemme; men at erindre og at glemme ere ikke Modsætninger. Erindringens Kunst er ikke let, fordi den i Tilberedelsens Øieblik kan blive forskjellig, medens Hukommelsen kun har Fluctuationen mellem at huske rigtigt og at huske galt. Hvad er saaledes Hjemvee? Det er noget Husket, som erindres. Simpelt frembringes Hjemvee derved, at man er borte. Kunsten vilde være, skjøndt man er hjemme, at kunne føle Hjemvee. Dertil fordres Færdighed i Illusion. At leve hen i en Illusion, hvori det bestandigt dæmrer, aldrig dages, at reflectere sig ud af al Illusion er ikke saa vanskeligt, som at reflectere sig til Illusion, samt at kunne lade den virke paa sig med al Illusionens Magt, uagtet man er Vidende. At trylle det Forbigangne til sig er ikke saa vanskeligt, som at trylle det Nærmeste fra sig for Erindringen. Dette er egentlig Erindringens Kunst og Reflectionen i anden Potens.

For at tilveiebringe sig en Erindring hører der Kjendskab til Stemningers, Situationers, Omgivelsers Modsætninger. En erotisk Situation, hvori Pointen var Landlivets hyggelige Afsideshed, lader sig stundom bedst erindre og inderindre i et Theater, hvor Omgivelsen og Larmen tvinger Modsætningen frem. Altid er dog den ligefremme Modsætning ikke den lykkelige. Hvis det ikke var uskjønt at bruge et Menneske som Middel, var det maaskee stundom den lykkelige Modsætning for at erindre et erotisk Forhold at skaffe sig en ny Kjærlighedshistorie blot for at erindre. – Modsætningen kan være yderst reflekteret. Yderspidsen af Reflexions-Forholdet mellem Hukommelse og Erindring er at bruge Hukommelsen mod Erindringen. Tvende Mennesker kunne af den modsatte Grund ikke ville see et Sted igjen, der minder om en Begivenhed. Den Ene aner slet ikke, at der er Noget til, som hedder Erindring, men frygter blot Hukommelsen. Ude af Øie ude af Sind tænker han, naar han blot ikke seer, saa har han glemt. Den Anden vil netop erindre, derfor vil han ikke see. Kun mod de ubehagelige Erindringer bruger han Hukommelsen. Den, der forstaaer sig paa Erindringen, men ikke forstaaer dette, har vel Idealiteten, men mangler Forfarenhed i at bruge consilia evangelica adversus casus conscientiæ. Han vil vel endog ansee Raadet for et Paradox, og være sky for at udholde den første Smerte, der dog altid er at foretrække ligesom det første Tab. Naar Hukommelsen atter og atter opfriskes beriger den Sjelen med en Masse Details, som adspreder Erindringen. Angeren er saaledes en Erindren af Skyld. Reent psychologisk seet troer jeg virkelig, at Politiet er Forbryderen behjælpelig i ikke at komme til at angre. Ved den idelige Optegnen og Repeteren af hans Levnetsløb faaer Forbryderen en saadan Hukommelses-Færdighed i at opramse sit Liv, at Erindringens Idealitet fordrives. Der hører stor Idealitet til virkeligt at angre og især til at angre strax; thi Naturen kan ogsaa hjælpe et Menneske, og den sildige Anger, der i Henseende til at huske er ubetydelig, er ofte den tungeste og dybeste. – Betingelsen for al Productivitet er det at kunne erindre. Vil man ikke mere være produktiv, behøver man blot at huske det Samme, man erindrende vilde producere, og Productiviteten er gjort umulig, eller den vil blive En saa ækel, at man opgiver den jo før jo hellere.

Erindringens Fælledsskab existerer egentligen ikke. En Art quasi-Fælledsskab er en Modsætnings Form, som den Erindrende bruger for sit eget Vedkommende. Det frister stundom Erindringen bedst frem, at man lader som betroede man sig til en Anden, blot for bagved denne Fortrolighed at gjemme en ny Reflexion, hvori Erindringen bliver til for Een selv. Med Hensyn til Hukommelsen kan man godt slaae sig sammen til gjensidig Assistance. I denne Henseende ere Festmaaltider og Geburtsdags-Glæder, Elskovspanter og dyrebare Amindelser hensigtsmæssige ligesom at lægge Krølle i en Bog for at huske, hvor man slap, og ved Hjælp af Krøllerne at kunne være sikker paa at have læst hele Bogen igjennem. Erindringens Perse derimod maa Enhver træde alene. I og for sig ligger deri langtfra nogen Forbandelse. Forsaavidt man altid er ene om en Erindring, er enhver Erindring en Hemmelighed. Selv om Flere ere interesserede i hvad der er den Erindrende Erindringens Gjenstand, han er dog ene om sin Erindring, den tilsyneladende Offentlighed er blot illusorisk.

Det her Fremsatte er mig selv til Erindring om Tanker og Tankesysler, der mange Gange og paa mange Maader have beskæftiget min Sjel. Anledningen til at de henkastedes er, at jeg nu føler mig stemt til at ville indfrie en oplevet Begivenhed for Erindringen, til at ville optegne, hvad der allerede i nogen Tid har ligget færdigt husket og partielt ogsaa erindret. Det jeg har at huske er i Omfang lidet, Hukommelsens Arbeide forsaavidt let; derimod har jeg havt Vanskelighed ved ret at faae det ud for Erindringen, netop fordi det for mig er blevet noget ganske Andet end for De herrer Deeltagere, der formodentlig vilde smile ved at see en saadan Ubetydelighed tillagt nogetsomhelst Værd, et kaadt Indfald, en fortvivlet Idee, som de selv vilde kalde det. Ja hvor lidet Hukommelsen her har for mig at betyde, seer jeg deraf, at det stundom er mig, som havde jeg slet ikke oplevet det, men selv digtet det.

Vel veed jeg, at jeg ikke saa snart skal glemme hiint Gjestebud, hvori jeg blev deeltagende uden at være Deeltager; men desuagtet kan jeg nu ikke beslutte mig til at slippe det, uden at have sikkret mig et omhyggeligt απομνημονευμα af hvad der for mig virkelig var memorabile.

[…]

\sectionheading{Gjestebudet  ·  SV¹ VI, 26–34}

Det var en af de sidste Dage i Juli Maaned om Aftenen henved Klokken 10, at Deeltagerne samlede sig til hiint Gjestebud. Datum og Aarstal har jeg glemt; Sligt interesserer da ogsaa kun Hukommelsen, ikke Erindringen. Stemning og hvad der sorterer under Stemning er alene Erindringens Gjenstand; og som den ædle Viin vinder ved at passere Linien, fordi Vandpartiklerne fordampe, saaledes vinder Erindringen ogsaa ved at tabe Hukommelsens Vandpartikler: dog bliver Erindringen ligesaa lidet derved en Indbildning som den ædle Viin bliver det. – Deeltagerne vare fem: Johannes med Tilnavn Forføreren, Victor Eremita, Constantin Constantius og endnu Tvende, hvis Navn jeg just ikke har glemt, hvad der var temmeligt ligegyldigt, men hvis Navn jeg ikke fik at vide. Det var som havde disse Tvende intet proprium; thi de nævnedes bestandigt kun ved et Epitheton. Den Ene blev kaldet: det unge Menneske. Han var vel heller ei mere end nogle og tyve Aar, slank og fiint bygget, meget mørkladen. Udtrykket i hans Ansigt var tænksomt, men mere end dette behagede hans Mine, der elskelig og indtagende udprægede en Sjele-Reenhed, der harmonerede ganske med hans hele Skikkelses næsten qvindelige vegetative Blødhed og Gjennemsigtighed. Men denne ydre Skjønhed glemte man dog igjen over det næste Indtryk, eller man beholdt det kun in mente medens man betragtede en Yngling, der kun dannet, eller for at bruge et endnu zartere Udtryk, opdægget ved Tanken, næret ved sin Sjels eget Indhold, Intet havde havt med Verden at skaffe, hverken var bleven vakt og opflammet eller uroliget og forstyrret. Liig en Søvngænger havde han Loven for sin Optræden i sig selv, og hans elskelige velvillige Mine angik Ingen, men afspeilede kun Sjelens Grundstemning. Den Anden kaldte de Modehandleren, hvilket var hans borgerlige Stilling. Af ham var det en Umuelighed at faae et heelt Indtryk. Han var klædt efter den allernyeste Mode, krøllet og parfumeret, duftende af eau de Cologne. Hans Optræden var i eet Øieblik ikke uden Aplomb, men i næste Øieblik antog hans Gang en vis dandsende Festlighed, en vis Svæven, som hans Førlighed dog indtil videre satte en Grændse for. Selv naar han var mest ondskabsfuld i sin Tale, havde hans Stemme bestandigt noget af den Boutique-Behagelighed, den Galanteri-Sødlighed, som vist var ham selv yderst modbydelig, og kun tilfredsstillede hans Trods. Naar jeg nu tænker paa ham, forstaaer jeg ham vel bedre, end da jeg saae ham træde ud af Vognen og uvilkaarligt kom til at lee. Imidlertid er der dog Modsigelse tilbage. Han har forgjort eller forhexet sig selv, ved sin Villies Trolddom tryllet sig ind i en næsten fjantet Skikkelse, men har ikke ganske tilfredsstillet sig selv deri, hvorfor Reflexionen af og til titter frem.

Det var en af de sidste Dage i Juli Maaned om Aftenen henved Klokken 10, at Deeltagerne samlede sig til hiint Gjestebud. Datum og Aarstal har jeg glemt; Sligt interesserer da ogsaa kun Hukommelsen, ikke Erindringen. Stemning og hvad der sorterer under Stemning er alene Erindringens Gjenstand; og som den ædle Viin vinder ved at passere Linien, fordi Vandpartiklerne fordampe, saaledes vinder Erindringen ogsaa ved at tabe Hukommelsens Vandpartikler: dog bliver Erindringen ligesaa lidet derved en Indbildning som den ædle Viin bliver det. – Deeltagerne vare fem: Johannes med Tilnavn Forføreren, Victor Eremita, Constantin Constantius og endnu Tvende, hvis Navn jeg just ikke har glemt, hvad der var temmeligt ligegyldigt, men hvis Navn jeg ikke fik at vide. Det var som havde disse Tvende intet proprium; thi de nævnedes bestandigt kun ved et Epitheton. Den Ene blev kaldet: det unge Menneske. Han var vel heller ei mere end nogle og tyve Aar, slank og fiint bygget, meget mørkladen. Udtrykket i hans Ansigt var tænksomt, men mere end dette behagede hans Mine, der elskelig og indtagende udprægede en Sjele-Reenhed, der harmonerede ganske med hans hele Skikkelses næsten qvindelige vegetative Blødhed og Gjennemsigtighed. Men denne ydre Skjønhed glemte man dog igjen over det næste Indtryk, eller man beholdt det kun in mente medens man betragtede en Yngling, der kun dannet, eller for at bruge et endnu zartere Udtryk, opdægget ved Tanken, næret ved sin Sjels eget Indhold, Intet havde havt med Verden at skaffe, hverken var bleven vakt og opflammet eller uroliget og forstyrret. Liig en Søvngænger havde han Loven for sin Optræden i sig selv, og hans elskelige velvillige Mine angik Ingen, men afspeilede kun Sjelens Grundstemning. Den Anden kaldte de Modehandleren, hvilket var hans borgerlige Stilling. Af ham var det en Umuelighed at faae et heelt Indtryk. Han var klædt efter den allernyeste Mode, krøllet og parfumeret, duftende af eau de Cologne. Hans Optræden var i eet Øieblik ikke uden Aplomb, men i næste Øieblik antog hans Gang en vis dandsende Festlighed, en vis Svæven, som hans Førlighed dog indtil videre satte en Grændse for. Selv naar han var mest ondskabsfuld i sin Tale, havde hans Stemme bestandigt noget af den Boutique-Behagelighed, den Galanteri-Sødlighed, som vist var ham selv yderst modbydelig, og kun tilfredsstillede hans Trods. Naar jeg nu tænker paa ham, forstaaer jeg ham vel bedre, end da jeg saae ham træde ud af Vognen og uvilkaarligt kom til at lee. Imidlertid er der dog Modsigelse tilbage. Han har forgjort eller forhexet sig selv, ved sin Villies Trolddom tryllet sig ind i en næsten fjantet Skikkelse, men har ikke ganske tilfredsstillet sig selv deri, hvorfor Reflexionen af og til titter frem.

[…]

Stedet var valgt i en Skovegn, et Par Mile fra Kjøbenhavn. Salen, hvori der skulde spises, var ny decoreret, og paa enhver Maade gjort ukjendelig; et mindre Værelse, der ved en Corridor skiltes fra Salen, var indrettet for et Orchester. Skodder og Forhæng vare anbragte for alle Vinduer, og bag disse stode Vinduerne aabne. Intimationen vilde Constantin skulde ligge i, at de kom kjørende i Aftenstunden. Selv om man veed, at man kjører til et Gjestebud, og Phantasien derfor et Øieblik forsøger sig i det Yppige, saa er dog Natur-Omgivelsens Indtryk for mægtigt til, at det ikke skulde seire. At dette ikke skulde skee, var det Eneste, Constantin frygtede, thi som der ingen Magt er, der saaledes veed at forskjønne Alt, som Phantasien, saa er der heller ingen Magt, der saaledes kan forstyrre Alt, naar det slaaer feil for Een, idet den skal berøre Virkeligheden. Men at kjøre i en Sommeraften giver ikke Phantasien Retningen imod det Yppige, men lige den modsatte. Om man end ikke seer og hører det, danner Phantasien dog uvilkaarlig et Billede af Aftenens hjemlige Længsel, saa man seer Piger og Karle vandre hjem fra Markarbeidet, hører Høstvognens iilsomme Raslen, forklarer selv det fjerne Brøl fra Engen som en Længsel. Saaledes frister Sommeraftenen det Idylliske frem, qvæger selv et higende Sind med sin Beroligelse, bevæger selv den flyvende Phantasie til med autochtonisk Hjemvee at dvæle ved Jorden, som det Sted man er kommen fra, lærer det umættelige Sind at nøies med Lidet, stiller Een tilfreds, thi i Aftenstunden staaer Tiden stille og Evigheden dvæler. – Saaledes ankom de i Aftenstunden: de Indbudne; thi Constantin var tagen noget tidligere ud. Victor E., der laae paa Landet i Nærheden, kom ridende, de Andre kjørende, og just som deres Vogn kjørte til Siden, svingede en holsteensk Vogn ind ad Porten: et lystigt Compagnie af fire Haandværkssvende, hvilke bleve beværtede for dernæst i det afgjørende Øieblik at haves i Beredskab som et Nedrivningscorps: saaledes ere Brandfolkene af modsat Grund tilstede i Theatret for strax at slukke.

Saa længe man er Barn har man Phantasie nok til, om det var i en Time i det mørke Værelse at kunne holde sin Sjel paa Spidsen, paa Forventningens Spidse; naar man er Ældre virker Phantasien let til at gjøre Een kjed af Juletræet, inden man faaer det at see.

Fløidørrene aabnedes; Virkningen af den straalende Belysning, Køligheden, der strømmede dem imøde, Vellugtens krydrede Bedaarelse, Anretningens Smagfuldhed overvældede et Øieblik de Indtrædende, og da der i det samme fra Orchestret lød Toner af Ballet i Don Juan, forklaredes de Indtrædendes Skikkelser, og som af Ærbødighed for en usynlig Aand, der omgav dem, standsede de et Øieblik som Den, hvem Beundring har vakt og som er opstanden til Beundring.

[…]

Hvad Talernes Indhold angik, da foreslog Constantin, at der skulde tales om Elskov eller om Forholdet mellem Mand og Qvinde, dog maatte der ikke fortælles Kjerlighedshistorier, men vel maatte der gjerne ligge Historier til Grund for Opfattelsen.

Betingelserne bleve antagne. – Alle en Verts retfærdige og billige Fordringer til Gjester bleve opfyldte: de spiste, drak og drak og bleve drukne, som det hedder i det Hebraiske, det vil sige de drak tappert.

Desserten blev fremsat. Havde Victor hidtil ikke faaet sin Fordring opfyldt, at høre et Vandsprings Pladsken, hvad han da heldigviis for ham havde glemt igjen siden hiin Samtale, saa brusede nu Champagne i Overflod. Klokken slog 12; da paabød Constantin Taushed, hilsede det unge Menneske med et Bæger og disse Ord: quod felix sit faustumque, og bød ham at tale først.

[…]

\sectionheading{Modehandlerens Tale  ·  SV¹ VI, 65–78}

Neppe havde Victor endt, førend Modehandleren sprang op, rev en Viinflaske omkuld, der stod for ham, og nu begyndte han saaledes.

Vel talt kjære Drikkebrødre, vel talt, jo mere jeg hører Eder tale, desto mere forvisses jeg om, at I ere Medsammensvorne, jeg hilser Eder som saadanne, jeg forstaaer Eder som saadanne, thi Sammensvorne forstaaer man langtfra. Og dog, hvad veed I, hvad er Eders Smule Theori, som I give Udseende af Erfaring, Eders Smule Erfaring, som I syer om til en Theori, og endeligen i hver anden Stund troe I dog et Øieblik og hildes et Øieblik. Nei jeg kjender Qvinden – fra hendes svage Side, det vil sige jeg kjender hende. Jeg skyer ingen Forfærdelse i mit Studium og skyer intet Middel for at forvisse mig om hvad jeg har forstaaet; thi jeg er en Rasende, og Rasende maa man være for at forstaae hende, og var man det ikke før, naar man har forstaaet hende, bliver man det. Som Røveren har sit Smuthul ved den larmende Landevei, og Myreslugeren sin Tragt ved det løse Sand, og Caperen sit Skjulested ved det brusende Hav, saa har jeg min Modeboutique midt i Menneske-Vrimlen forførerisk, uimodstaaelig for en Qvinde som Venusbjerget for Manden. Her i en Modeboutique lærer man hende at kjende, praktisk og fra Grunden af uden al theoretisk Ophævelse. Ja betydede Moden ikke Andet end at en Qvinde i Begjerligheds Brynde kastede Alt af sig, saa var det dog Noget. Men saaledes er det ikke, Moden er ikke aabenbar Vellyst, ikke tolereret Udsvævelse, men en Uanstændighedens Snighandel, der er autoriseret som Velanstændighed. Og som i det hedenske Preussen den giftefærdige Pige bar en Klokke, hvis Ringen var til Signal for Mændene, saaledes er en Qvindes Existents paa Moden et evindelig Klokkespil ikke for Udsvævende men for slikvurne Vellystninge. Man troer Lykken er en Qvinde o! ja den er jo ustadig; men den er dog ustadig i Noget, thi den kan give Meget, forsaavidt er den ikke en Qvinde. Nei, Moden er en Qvinde, thi Moden er Ustadighed i Nonsens, der kun kjender een Conseqvents, at den altid bliver mere og mere pinegal. Een Time i min Boutique er mere værd end Aar og Dage udenfor, hvis man ønsker at lære Qvinden at kjende; i min Modeboutique, thi den er den eneste i Residentsen, der er ikke Tanke om Concurrence; hvo skulde vove en Væddestrid med Den, der ganske har offret sig og offrer sig som Ypperstepræst i denne Afgudsdyrkelse. Nei der er ikke et distingueret Selskab uden at mit Navn er det første og det sidste, og der er ikke et borgerligt Selskab, hvor jo mit Navn, naar det nævnes, indgyder hellig Ærefrygt som Kongens, og der er ingen Dragt saa vanvittig, at den jo, naar den er fra min Boutique ledsages af en Hvidsken, idet den gaaer gjennem Salen; og der er ikke en Dame af Byrd, der vover at gaae min Boutique forbi, og ikke en borgerlig Pige uden at hun gaaer den sukkende forbi og tænker: dersom jeg blot havde Raad. Nu hun var ei heller bedragen. Jeg bedrager Ingen; det Fineste, det Kostbareste leverer jeg, til de billigste Priser, ja under Prisen sælger jeg, saaledes ønsker jeg ikke at vinde, nei jeg sætter aarligt store Summer til. Og dog vil jeg vinde, jeg vil det, jeg giver min sidste Hviid hen for at bestikke, for at underkjøbe Modens Organer, at mit Spil kan vindes. Det er mig en Vellyst uden Lige at tage de kosteligste Tøier frem, at skære til, at klippe ægte Brüsler-Kniplinger ud, for da at danne en Narredragt, jeg udsælger til de billigste Priser, ægte Stoffer og paa Moden. I troer maaskee, at det kun er i enkelte Øieblikke hun ønsker at være paa Moden. Langtfra, hun vil det altid, og det er hendes eneste Tanke. Thi Qvinden har Aand, men den er anbragt ligesaa godt som den forlorne Søns Midler, og Qvinden har Reflexion i ubegribelig høi Grad, thi der er intet saa Helligt, at hun jo i samme Nu finder det commensurabelt for Pynten, og Pyntens fornemste Udtryk er Moden; hvad Under at hun finder det commensurabelt, thi Moden er jo det Hellige; og der er intet saa ubetydeligt, at hun jo atter veed at bringe det i Forhold til Pynten, og Pyntens ideløseste Udtryk er Moden; og der er Intet, Intet i hele hendes Paaklædning, ikke det mindste Baand, uden hun har en Forestilling om dets Forhold til Moden, og uden hun øieblikkeligt opdager, om den forbigaaende Dame har bemærket det; thi for hvem pynter hun sig, uden for andre Damer! Selv i min Boutique, hvor hun jo kommer for at vorde udstyret paa Moden, selv der er hun paa Moden. Ligesom der er en særskilt Badedragt og Ridedragt, saaledes er der ogsaa en egen Art af Paaklædning, der er moderne til at gaae i Boutiquen med. Denne Dragt er ikke saaledes skjødesløs som den Negligee, i hvilken en Dame behager at lade sig overraske tidligere om Formiddagen. Pointen er da hendes Qvindelighed og Coquetteriet i at lade sig overraske. Modedragten derimod er beregnet paa at være skjødesløs, lidt letfærdig uden at være generet derved, fordi en Modehandler staaer i et andet Forhold til hende end en Cavaleer. Coquetteriet ligger i at vise sig saaledes for en Mand, der paa Grund af sin Stilling ikke tør fordre Damens qvindelige Anerkjendelse, men maa nøies med de uvisse Indtægter, der rigelig falde af, men uden at hun tænker derpaa, eller uden at det falder hende ind ligeover for en Modehandler at ville være Dame. Pointen er derfor den, at Qvindeligheden paa en Maade er udeladt, og Coquetteriet gjort ugyldigt i den distinguerede Dames fornemme Overlegenhed, der vilde smile hvis Nogen vilde alludere til et saadant Forhold. I sin Negligee ved en Visit skjuler hun sig og forraader sig gjennem Skjulet, i Boutiquen blotter hun sig med den yderste Nonchalence, thi det er kun en Modehandler – og hun er en Qvinde. Nu falder Schawlet lidt ned og giver en lille Bløße, dersom jeg ikke veed, hvad det betyder, og hvad hun vil, saa er mit Renomee fortabt, nu snærper hun apriorisk, nu gesticulerer hun aposteriorisk, nu svinker hun i Hofterne, nu speiler hun sig, og seer mit beundrende Fjæs i Speilet, nu læsper hun, nu tripper hun, nu svæver hun, nu slæber hun Foden løsagtigt, nu synker hun blødagtigt i en Lænestol, medens jeg i ydmyg Stilling præsenterer hende en Flaçon og med min Tilbedelse svaler over hendes Hede, nu slaaer hun skjelmsk efter mig med Haanden, nu taber hun sit Lommetørklæde; uden en eneste Bevægelse lader hun endog Armen forblive i sin løst nedhængende Stilling, medens jeg bøier mig dybt og tager det op, og byder det og faaer et lille protegerende Nik. Saaledes bærer en Dame sig ad, der er paa Moden, naar hun er i Boutiquen. Om Diogenes bevægede den Qvinde, der laae bedende i en noget uanstændig Stilling ved sin Bemærkning, om hun ikke troede, at Guderne kunde see hende bagfra, veed jeg ikke; men det veed jeg, at hvis jeg vilde sige til hendes knælende Naade, deres Kjoles Folder falde ikke paa Moden, da vilde hun frygte dette mere end at krænke Guderne. Vee! den Udskudte, den Askepot, der ikke forstaaer dette. Pro dii immortales hvad er en Qvinde ogsaa, naar hun ikke er paa Moden, per deos obsecro hvad er hun naar hun er paa Moden!

Om det er sandt? Ja forsøg det: Lad Elskeren, idet den Elskede synker salig til hans Bryst og hvidsker uforstaaeligt: Din for evig, skjulende sit Hoved ved hans Barm, lad ham sige til hende: søde Catinka din Frisure er slet ikke paa Moden. Maaskee tænke Mændene ikke derpaa, men Den, der veed det og nyder Anseelse af at vide det, er den farligste Mand i Kongeriget. Hvilke salige Timer, Elskeren tilbringer med den Elskede før Brylluppet, veed jeg ikke, men de salige Timer, hun tilbringer i min Boutique, gaae hans Næse forbi. Uden mit Kongebrev og min Sanction er dog et Bryllup en ugyldig Akt eller ogsaa et meget plebeiisk Foretagende. Lad Øieblikket allerede være kommet da de skal møde for Alteret, lad hende træde frem med den bedste Samvittighed i Verden, at Alt er kjøbt hos mig og for mig gjort Prøve paa Alt, hvis jeg vilde styrte til og sige: men min Gud min naadige Frøken den Myrthekrands sidder reent forkeert, – saa blev maaskee Ceremonien udsat. Men alt Sligt veed Mændene ikke, man maa være Modehandler for at vide det. Der hører en saa uhyre Reflexion til at controllere en Qvindes Reflexion, at kun en Mand, der offrer sig dertil, formaaer det, og formaaer det, hvis han oprindelig er begavet. Lykkelig er derfor den Mand, der ikke indlader sig med nogen Qvinde, hun tilhører ham dog ikke, om hun end ikke tilhører nogen anden Mand, thi hun tilhører hiint Phantom, dannet af qvindelig Reflexions unaturlige Omgang med qvindelig Reflexion: Moden. See derfor skulde en Qvinde altid sværge ved Moden, saa var der Fynd i hendes Eed; thi Moden er dog det Eneste, hun altid tænker paa, det Eneste, hun kan tænke sammen med og ind i Alt. Der er fra min Boutique til den fornemme Verden udgaaet det glade Budskab for alle distinguerede Damer, at Moden byder Brugen af en egen Art Hovedtøi, naar man tager i Kirke, og at dette Hovedtøi igjen maa være noget forskjelligt til Høimesse og til Aftensang. Naar da Klokkerne ringe, da holder Eqvipagen for min Dør. Hendes Naade stiger ud (thi ogsaa dette er forkyndt, at Ingen kan sætte det Hovedtøi rigtigt uden jeg Modehandleren); jeg styrter hende imøde i dybe Buk, fører hende ind i mit Cabinet; medens hun blødagtigt vegeterer, bringer jeg Alt i Orden. Hun er færdig, hun har speilet sig; hurtigt som et Gudernes Sendebud iler jeg forud, har aabnet Cabinetsdøren og bukker, iler til Boutiquesdøren, lægger min Arm paa Brystet som en østerlandsk Slave, men opmuntret af et naadigt Knix vover jeg endog at tilkaste hende et tilbedende og beundrende Haandkys – hun sidder i Vognen, see! hun har glemt Psalmebogen, jeg iler ud og rækker den ind ad Vinduet til hende, tillader mig endnu engang at minde hende om, at holde Hovedet en lille Kjende til Høire, og selv at rette lidt ved det, hvis hun ved at stige ud skulde bringe Hovedtøiet en Smule af Lave. Hun kjører hen og opbygges.

I troe maaskee, at det kun er fornemme Damer, der hylde Moden, langtfra. See! til mine Sye-Jomfruer, paa hvis Toilet jeg Intet sparer, at Modens Dogmer maae forkyndes med Eftertryk fra min Boutique. De danne et Chor af Halvgale, og jeg selv som Ypperstepræst foregaaer med et lysende Exempel, og bortødsler Alt, blot for ved Modens Hjelp at gjøre enhver Qvinde latterlig. Thi naar en Forfører brouter af, at enhver Qvindes Dyd er venibel for den rette Kjøber, da troer jeg ham ikke, men jeg troer, at enhver Qvinde i føie Tid skal være fanatiseret af Modens vanvittige og besmittende Selv-Reflexion, der fordærver hende ganske anderledes, end om hun var forført. Jeg har forsøgt det mere end eengang. Formaaer jeg det ikke selv, saa hidser jeg et Par af Modens Trælqvinder, som høre til hendes Stand, paa hende; thi som man afretter Rotter til at bide Rotter, saa er den fanatiserede Qvindes Bid ligesom Tarantellens. Og fremfor Alt er det farligt, naar en Mand træder understøttende til. Om jeg tjener Djævelen eller jeg tjener Guden, veed jeg ikke, men jeg har Ret, jeg vil have Ret, jeg vil det, saa længe jeg eier en eneste Hviid, jeg vil det saa længe til Blodet sprøiter ud af mine Fingre. Physiologen tegner en Qvindes Skabning af for at vise Snørelivets forfærdelige Følger, ved Siden tegner han den normale Skikkelse. Det er rigtigt, men kun den ene har Virkelighedens Gyldighed: de gaae alle med Snøreliv. Skildrer saaledes den Modesyges elendige forkrøblede Forskruethed, beskriv denne snigende Reflexion, der fortærer hende, og skildrer den qvindelige Blufærdighed, der mindst af Alt veed noget om sig selv, gjør det godt og Du har tillige dømt Qvinden og i Virkeligheden dømt hende forfærdeligt. Opdager jeg engang en saadan Pige, der nøisom og ydmyg ikke er fordærvet ved uanstændig Omgang med Qvinder, hun skal dog falde. Jeg bringer hende i mine Garn, nu staaer hun paa Offerpladsen det vil sige i min Boutique. Med det haanligste Blik, som fornem Nonchalance kan væbne sig med, maaler jeg hende, hun forgaaer af Skræk, en Latter fra Naboværelset, hvor mine afrettede Haandlangere sidde, tilintetgjør hende. Naar jeg da har faaet hende udmaiet paa Moden, naar hun seer vanvittigere ud end et Daarekistelem, saa vanvittig som Een, der end ikke kunde blive optagen i en Daarekiste, da gaaer hun salig fra mig, intet Menneske end ikke en Gud kunde forfærde hende, thi hun er jo paa Moden.

Forstaaer I mig nu, forstaaer I, hvorfor jeg kalder Eder Medsammensvorne, om end kun langt fra. Forstaaer I nu min Opfattelse af Qvinden. Alt i Livet er en Modesag, Gudsfrygt er en Modesag, og Kjærlighed og Fiskebeensskjørter og en Ring i Næsen. Saa vil jeg af yderste Evne komme den ophøiede Genius til Hjælp, der begjærer at lee ad det latterligste af alle Dyr. Har Qvinden reduceret Alt paa Moden, saa vil jeg ved Modens Hjælp prostituere hende, som hun har fortjent; jeg raster ikke, jeg Modehandleren, min Sjel raser naar jeg tænker paa min Opgave, hun skal endnu komme til at gaae med en Ring i Næsen. Søger derfor ingen Elsket, opgiver Elskov som det farligste Naboskab, thi ogsaa Eders Elskede skulle komme til at gaae med en Ring i Næsen.

[…]

Taflet var hævet. Der behøvedes kun et Vink af Constantin; med militairisk Taktfasthed forstode Deeltagerne hinanden, naar det gjaldt at gjøre et Sving og et Omsving. […]

\sectionheading{Til Læseren: Existents-Sphærerne (Frater Taciturnus)  ·  SV¹ VI, 443}

Der gives 3 Existents-Sphærer: den æsthetiske, den ethiske, den religieuse. Det Metaphysiske er Abstraktionen, og der er intet Menneske, som existerer metaphysisk. Det Metaphysiske, det Ontologiske er, men det er ikke til, thi naar det er til er det i det Æsthetiske, i det Ethiske, i det Religieuse, og naar det er, er det Abstraktionen af eller et Prius for det Æsthetiske, det Ethiske, det Religieuse. Den ethiske Sphære er kun Gjennemgangssphære, og derfor er dens høieste Udtryk Angeren som en negativ Handlen. Den æsthetiske Sphære er Umiddelbarhedens, den ethiske er Fordringens (og denne Fordring er saa uendelig, at Individet altid gaaer fallit), den religieuse er Opfyldelsens, men vel at mærke, ikke en saadan Opfyldelse, som naar man fylder Guld i en Stok eller i en Pose, thi Angeren har netop gjort uendelig Plads, og deraf den religieuse Modsigelse: paa eengang at ligge paa 70,000 Favne Vand og dog være glad.

Som den ethiske Sphære er en Gjennemgang, hvilken man dog ikke eengang for alle gjennemgaaer, som Angeren er dens Udtryk, saa er Angeren det mest dialektiske. Intet Under derfor, at man frygter den, thi giver man den een Finger, tager den hele Haanden. Som Jehova i det gamle Testament hjemsøger Forældrenes Brøde paa Børnene i sildigste Slægter, saaledes gaaer Angeren tilbage bestandigt forudsættende Gjenstanden for sin Undersøgelse. I Angeren er Bevægelsens Ryk, og derfor Alt vendt om. Dette Ryk betyder netop Forskjellen mellem det Æsthetiske og det Religieuse, som den mellem det Udvortes og Indvortes.

[…]

\sectionheading{Slutningsord  ·  SV¹ VI, 450–456}

Min kjære Læser! – dog til hvem taler jeg? Der er maaskee slet ingen bleven tilbage. Det er vel gaaet mig omvendt, som det gik hiin ædle Konge, hvem Sorgens Budskab lærte at haste, hvis ilsomme Ridt til den Elskede i Dødsfaren den uforglemmelige Vise har gjort uforglemmeligt, naar den synger om de 100 Svende, der fulgte ham fra Skanderborg, de 15, der rede med ham over Randbøl-Hede, men der han kom over Ribe Bro, da var den Herre alene – saaledes er det vel omvendt og af modsat Grund gaaet med mig, der fængslet til een Tanke ikke er kommen af Stedet: Alle ere redne fra mig. I Begyndelsen holdt vel den velvillige Læser sin rappe Ganger an og troede, at det var en Pasgænger, jeg reed, men da jeg slet ikke kom af Stedet, blev Hesten (Læserens nemlig) eller om man saa vil, Rytteren utaalmodig, og jeg blev alene tilbage: en Urider eller Søndagsrytter, hvem Alle ride fra.

Forsaavidt er der da slet Intet at haste efter, jeg har Tiden og Dagen for mig selv, og kan uforstyrret uden at incommodere Nogen tale med mig selv om mig selv. Den Religieuse er efter min Anskuelse: den Vise. Men Den, som indbilder sig at være det uden at være det, er en Daare; men Den, som seer een Side af det Religieuse, er en Sophist. Af disse Sophister er jeg een, og var jeg end duelig til at æde de Andre, vilde jeg dog ikke blive federe, hvilket ikke er uforklarligt som i Ægypten med de magre Køer, thi Sophisterne ere i Henseende til det Religieuse ikke fede Køer, men magre Sild. Jeg seer det Religieuse fra alle Sider, og forsaavidt har jeg bestandigt een Side mere end den Sophist, som kun seer een Side, men det, som gjør mig til Sophist, er, at jeg ikke bliver den Religieuse. Den Mindste i Religieusitetens Sphære er uendeligt større end den største Sophist. Min Smerte derover have Guderne lindret ved at forunde mig mangen skjøn Betragtning og væbne mig med en vis Grad af Vittighed, hvilken vil fratages mig, dersom jeg bruger den mod det Religieuse. – Sophisterne ere at henføre til 3 Classer. 1) De, der fra det Æsthetiske faae et umiddelbart Forhold til det Religieuse. Religionen bliver her Poesi, Historie; Sophisten selv er begeistret for det Religieuse, men digterisk begeistret, i denne Begeistring er han villig til at bringe ethvert Offer, sætter endog vel Livet til, men bliver derfor ikke den Religieuse. I Maximum af Anseelse forvexler han sig selv og lader sig forvexle med en Prophet og en Apostel. 2) De, som fra det umiddelbar Ethiske træde i et umiddelbart Forhold til det Religieuse. For dem bliver Religionen en positiv Pligtlære, istedenfor at Angeren er det Ethiskes høieste Gjerning og netop negativ. Sophisten bliver uforsøgt i den uendelige Reflexion, et Dydsmynster i et positivt Indbegreb. Her ligger hans Begeistring, og uden Svig gjør han sig en Glæde af at begeistre Andre til det Samme. 3) De, som sætte det Metaphysiske i et umiddelbart Forhold til det Religieuse. Religionen bliver her Historien, som er færdig, Sophisten er færdig med Religionen og bliver i sit Maximum Opfinder af Systemet. – Det, hvorfor Sophisterne blive beundrede af Menneskenes Mængde, er, at de høimodigen ikke bekymre sig om sig selv i Sammenligning med den poetiske Intuition, i hvilken den ene fortaber sig, i Sammenligning med den positive Stræben til et Maal udenfor sig, som vinker den anden, i Sammenligning med det uhyre Resultat, som den Tredie, ved at sætte det Færdige sammen, erhverver. Men det Religieuse ligger netop i religieust uendeligt at bekymre sig om sig selv, og ikke om Syner; uendeligt at bekymre sig om sig selv, og ikke om et positivt Formaal, hvilket er negativt og endeligt, fordi det uendeligt negative er den eneste adæquate Form for det Uendelige; uendeligt at bekymre sig om sig selv, og altsaa ikke mene sig selv færdig, hvilket er negativt og Fortabelse. – Dette veed jeg, men jeg veed det i Aandens Ligevægt og er derfor en Sophist som de Andre, thi denne Ligevægt er en Forsyndelse mod det Religieuses hellige Lidenskab. Men denne Ligevægt i Eenheden af det Comiske og Tragiske, der er den uendelige Bekymring om sig selv i græsk Forstand (ikke den uendelige religieuse Bekymring om sig selv) er ikke uden Betydning for Belysningen af det Religieuse. I en vis Forstand er jeg saaledes langt fjernere fra det Religieuse end de 3 Classer af Sophister, der alle have begyndt paa det, men i en anden Forstand er jeg det nærmere, fordi jeg tydeligere seer, hvor det Religieuse er, og altsaa ikke griber feil i at gribe noget Enkelt, men feiler i ikke at gribe det.

Saaledes forstaaer jeg mig selv. Nøiet med det Mindre – haabende at det Større muligen engang skal forundes mig, beskæftiget med Aandens Syssel, hvori ethvert Menneske, efter min Mening, maa have rigeligt nok for det længste Liv, selv om dette var dannet af lutter længste Dage, – er jeg glad ved Tilværelsen, glad ved den lille Verden, der er min Omgivelse.

[…]

Nogle af mine Landsmænd mene, at Modersmaalet ikke skulde være dygtigt til at udtrykke vanskelige Tanker. Dette synes mig en besynderlig og utaknemlig Mening, som det ogsaa synes mig besynderligt og overdrevent at ville ivre for det, saa man næsten glemmer at glæde sig ved det, at forfægte en Uafhængighed saa ivrigt, at Iveren næsten synes at tyde paa, at man allerede føler sig afhængig, og at det stridige Ord tilsidst bliver det Spændende, ikke Sprogets Fryd det Vederqvægende. Jeg føler mig lykkelig ved at være bunden til mit Modersmaal, bunden som maaskee kun Faa er det, bunden som Adam var til Eva, fordi der ingen anden Qvinde var, bunden fordi det har været mig en Umulighed at lære noget andet Sprog og derved en Umulighed at fristes til at lade stolt og fornemt om det medfødte, men ogsaa glad ved at være bunden til et Modersmaal, der er riigt i indre Oprindelighed, naar det udvider Sjelen, og lyder vellystigt i Øret med sin søde Klang; et Modersmaal, der ikke stønner forfangent i den vanskelige Tanke, og derfor er det maaskee Nogen troer, at det ikke kan udtrykke den, fordi det gjør Vanskeligheden let ved at udtale den; et Modersmaal, der ikke puster og lyder anstrænget, naar det staaer for det Uudsigelige, men sysler dermed i Spøg og i Alvor indtil det er udsagt; et Sprog, der ikke finder langt borte, hvad der ligger nær, eller søger dybt nede, hvad der er lige ved Haanden, fordi det i lykkeligt Forhold med Gjenstanden gaaer ud og ind som en Alf, og bringer den for Dagen som et Barn den lykkelige Bemærkning, uden ret at vide af det; et Sprog, der er hæftigt og bevæget, hver Gang den rette Elsker veed mandligt at hidse Sprogets qvindelige Lidenskab, selv bevidst og seierrigt i Tankestriden, hver Gang den rette Hersker veed at føre det an, smidigt som en Bryder, hver Gang den rette Tænker ikke slipper det og ikke slipper Tanken; et Sprog, der om det end paa et enkelt Sted synes fattigt, dog ikke er det, men forsmaaet som en beskeden Elskerinde, der jo har den høieste Værd og fremfor Alt ikke er forjadsket; et Sprog, der ikke uden Udtryk for det Store, det Afgjørende, det Fremtrædende, har en yndig, en tækkelig, en livsalig Forkjærlighed for Mellemtanken og Bibegrebet og Tillægsordet, og Stemningens Smaasnakken, og Overgangens Nynnen, og Bøiningens Inderlighed og den dulgte Velværens forborgne Frodighed; et Sprog, der forstaaer Spøg nok saa godt som Alvor: et Modersmaal, der fængsler sine Børn med en Lænke, som »er let at bære – ja! men tung at bryde.«

Nogle af mine Landsmænd mene, at det er en gammel Erindring, paa hvilken Danmark tærer. Dette synes mig en besynderlig og utaknemlig Mening, som Ingen kan billige, der hellere vil være venlig og glad end tvær og modvillig, thi kun dette tærer. Andre mene, at der forestaaer Danmark en mageløs Fremtid; Nogle, som troe sig miskjendte og upaaskjønnede, trøste sig ogsaa med en bedre Efterslægt. Men Den, der er lykkelig ved det Nærværende og er flink i Opfindsomhed, naar det gjelder at være tilfredsstillet ved dette, har jo ikke mange Øieblikke for mageløse Forventninger, og saa lidet som han griber efter dem, saa lidet lader han sig forstyrre af dem. Og Den, som ikke mener sig paaskjønnet af sin Samtid, han fører dog en underlig Tale, ved at forjætte en bedre Efterslægt. Thi selv om det var saa, at han ikke var paaskjønnet, og selv om det var saa, at han blev bekjendt i en Eftertid, der priste ham, det er dog det Uretfærdige og det Hildede, at sige om denne Eftertid, at den derfor er bedre end Samtiden, bedre nemlig, fordi den synes bedre om ham. Saa stor Forskjel er der ikke mellem den ene Slægt og den anden, netop den Slægt, han dadler, er jo i det Tilfælde, at den priser, hvad en tidligere Slægt af Samtidige miskjendte. – Nogle af mine Landsmænd mene, at det at være Forfatter i Danmark er et fattigt Levebrød og en kummerlig Ansættelse. De mene ikke blot, at dette er Tilfældet med en saadan tvetydig Forfatter som mig, der ikke har en eneste Læser og kun nogle faa til midt i Bogen, hvem de derfor end ikke tænke paa i deres Dom, men de mene, at det ogsaa er Tilfældet med de Udmærkede. Nu Landet er jo kun lille. Men var det saa ringe en Ansættelse i Grækenland at være Øvrighed, uagtet det kostede Penge at være det! Sæt det var saaledes, sæt det blev saaledes, at det tilsidst blev en Forfatters Lod i Danmark at betale noget Vist om Aaret for Arbeidet med at være Forfatter: nu vel, hvad om det ogsaa var saa, at den Fremmede maatte sige: i Danmark er det en bekostelig Sag at være Forfatter, der er derfor heller ikke Forfattere i Hobetal, men da har man igjen heller ikke, hvad vi Udlændinge kalde: Stüberfängere, hvad den danske Literatur er saa ukjendt med, at Sproget end ikke har et Udtryk derfor. […]

\workheading{Afsluttende uvidenskabelig Efterskrift (af Johannes Climacus)}{1846}

\worknote{Antologiens ''Concluding Unscientific Postscript'' — det største og mest sammensatte uddrag i hele antologien. Medtaget er: den indledende selvbiografiske betragtning (Frederiksberg Have, cigaren, ''at gjøre Vanskeligheder overalt'') stillet foran Piecens Problem, som Hong gør; afsnittet om Lessing; den store afhandling om at Sandheden er Subjektiviteten (med den objektive uvished, ''70,000 Favne Vand'', det Socratiske og det Absurde, bønnen til afguden); afsnittene om Mulighed og Virkelighed, den subjektive Tænker, og den existentielle Pathos (Lidelsen, Humor som Religieusitetens Incognito, Religieusiteten A); samt de afsluttende ''Forstaaelsen med Læseren'' og ''En første og sidste Forklaring'', hvor pseudonymerne vedkendes. Climacus' fodnoter er udeladt (kun hovedteksten gengives); de mange spring følger Hongs udvalg og er markeret med […].}

\sectionheading{Indledning: Piecens Problem  ·  SV¹ VII, 155. 6–8}

Det er vel nu en fire Aar siden, at jeg fik det Indfald at ville forsøge mig som Forfatter. Jeg husker det ganske tydeligt, det var en Søndag, ganske rigtigt, ja det var en Søndag-Eftermiddag, jeg sad som sædvanlig ude hos Conditoren i Frederiksberg Have, hiin vidunderlige Have, der for Barnet var Trylleriets Land, hvor Kongen boede med Dronningen, hiin deilige Have, der for Ynglingen var den lykkelige Adspredelse i Folkets glade Lystighed; hiin venlige Have, hvor der nu for den Ældre er saa hjemligt i veemodig Opløftelse over Verden og hvad Verdens er, hvor selv Kongeværdighedens misundte Herlighed er hvad den jo er derude, en Dronnings Erindring af sin afdøde Herre; der sad jeg, som sædvanlig og røg min Cigar. Destoværre er dette den eneste Lighed jeg har kunnet udfinde mellem min Smule philosophiske Stræbens Begynden, og hiin poetiske Helts mirakuløse Begyndelse: at det var paa et offentligt Sted. Ellers er Ligheden slet ingen, og skjøndt Forfatter af »Smulerne« er jeg saa ubetydelig, at jeg staaer udenfor Literaturen; end ikke Subscriptions-Plans Litteraturen har jeg forøget eller kan med Sandhed siges i den at indtage en betydningsfuld Plads.

Jeg havde en halv Snees Aar været Student; skjøndt aldrig doven, var dog al min Virksomhed kun som en glimrende Uvirksomhed, en Art Beskjæftigelse, for hvilken jeg endnu har en stor Forkjerlighed, og i Forhold til hvilken jeg maaskee endog har lidt Genialitet. Jeg læste Meget, tilbragte det Øvrige af Dagen med at drive og tænke, eller med at tænke og drive, men derved blev det ogsaa; den produktive Spire i mig gik med til daglig Brug, fortæredes i sin første Grønnen. En uforklarlig Overtalelsens Magt holdt mig bestandigt lige stærkt og lige snildt tilbage, fængslet ved dens Overtalelse. Denne Magt var min Indolents; den er ikke som Elskovens heftige Higen eller som Begeistringens stærke Tilskyndelse, den er snarere som en Hustru, der holder En tilbage, og hos hvem man har det meget godt, saa godt, at det aldrig falder En ind at ville gifte sig; og saa meget er ogsaa vist, skjøndt jeg ellers ikke er ukjendt med Livets Beqvemmeligheder, af alle Beqvemmeligheder er Indolents den allerbeqvemmeste.

Saa sad jeg da der og røg min Cigar, indtil jeg henfaldt i Tanker. Blandt andre erindrer jeg disse: Du gaaer nu, sagde jeg til mig selv, og bliver et gammelt Menneske, uden at være Noget og uden egentligen at foretage Dig Noget. Overalt derimod hvor Du seer Dig om i Literaturen eller i Livet, seer Du de Feiredes Navne og Skikkelser, de dyrebare og med Acclamation hilsede Mennesker fremtrædende eller omtalte, de mange Tidens Velgjørere, der vide at gavne Menneskeheden ved at gjøre Livet lettere og lettere, Nogle ved Jernbaner, Andre ved Omnibusser og Dampskibe, Andre ved Telegrapheringer, Andre ved letfattelige Oversigter og korte Meddelelser af alt Videværdigt, og endeligen de sande Tidens Velgjørere, ved at gjøre Aands-Existentsen i Kraft af Tanke systematisk lettere og lettere og dog betydningsfuldere og betydningsfuldere: hvad gjør Du? Her afbrødes min Selvbetragtning, thi Cigaren var udrøget, og der maatte tændes en ny. Saa røg jeg igjen, og da farer pludseligen denne Tanke gjennem min Sjæl: Du maa gjøre Noget, men da det for Dine indskrænkede Evner vil være umuligt at gjøre Noget lettere end det er blevet, saa maa Du med samme menneskekjærlige Begeistring som de Andre paatage Dig at gjøre Noget sværere. Dette Indfald behagede mig overordentligt, det smigrede mig tillige med at jeg trods Nogen vilde ved min Stræben blive elsket og agtet af hele Menigheden. Naar nemlig Alle forene sig om paa alle Maader at gjøre Alt lettere, saa bliver der kun een Fare mulig, den nemlig, at Letheden blev saa stor, at den blev altfor let; saa bliver der kun eet Savn tilbage, om end endnu ikke følt, naar man vil savne Vanskeligheden. Af Kjerlighed til Menneskeheden, af Fortvivlelse over min forlegne Stilling, ved Intet at have udrettet og ved Intet at kunne gjøre lettere end det er gjort, af sand Interesse for dem, som gjøre Alt let, fattede jeg da det som min Opgave: overalt at gjøre Vanskeligheder. Det faldt mig tillige besynderligt paa, om jeg dog ikke egentligen havde min Indolents at takke for, at denne Opgave blev min. Thi langt fra som en Alladdin ved et Lykketræf at have fundet den, maa jeg snarere antage, at min Indolents, ved at forhindre mig i betimeligen at gribe til med at gjøre let, har paanødet mig det Eneste, der var blevet tilbage.

[…]

Problemet som blev fremsat i hiin Piece, uden at denne dog gav sig ud for at have løst det, da den jo blot vilde fremsætte det, lød saaledes: kan der gives et historisk Udgangspunkt for en evig Bevidsthed; hvorledes kan et saadant interessere mere end historisk; kan man bygge en evig Salighed paa en historisk Viden? (confer Titelbladet). I Piecen (pagina 162) blev der sagt: »som bekjendt er nemlig Christendommen det eneste historiske Phænomen, der uagtet det Historiske, ja netop ved det Historiske har villet være den Enkelte Udgangspunktet for hans evige Bevidsthed, har villet interessere ham anderledes end blot historisk, har villet begrunde ham hans Salighed paa hans Forhold til noget Historisk.« Det i Problemet Omspurgte er saaledes i historisk Costume Christendommen. Problemet er saaledes i Forhold til Christendommen. Mindre problematisk i en Afhandlings Form vilde Problemet lyde saaledes: om de apologetiske Forudsætninger for Troen, Approximations Overgange og Tilløb til Troen, den qvantiterende Indledning til Troens Afgjørelse. Det, der saa vilde være at afhandle, blev en Mængde Overveielser, som behandles eller behandledes af Theologerne i Indlednings-Videnskaben, i Indledningen til Dogmatiken, i Apologetiken.

For dog ikke at afstedkomme Forvirring, maa strax erindres om, at Problemet ikke er om Christendommens Sandhed, men om Individets Forhold til Christendommen, altsaa ikke om det ligegyldige Individs systematiske Iver for at arrangere Christendommens Sandheder i §§, men om det uendeligt interesserede Individs Bekymring betræffende sit Forhold til en saadan Lære. Saa simpelt som muligt (at jeg experimenterende skal bruge mig selv): »jeg Johannes Climacus, barnefødt her af Byen, nu tredive Aar gammel, et slet og ret Menneske ligesom Folk er fleest, antager, at der for mig lige saa vel som for en Tjenestepige og en Professor er et høieste Gode ivente, som kaldes en evig Salighed; jeg har hørt, at Christendommen betinger En dette Gode: nu spørger jeg, hvorledes kommer jeg i Forhold til denne Lære.« »Mageløse Fripostighed,« hører jeg en Tænker sige, »rædsomme Forfængelighed i det verdenshistorisk-bekymrede, det theocentriske, det speculativt-betydningsfulde nittende Aarhundrede, at turde lægge en saadan Vægt paa sit eget lille Jeg.« Det gyser i mig; hvis jeg ikke havde hærdet mig mod adskillige Forfærdelser, vilde jeg vel stikke Halen mellem Benene. Men jeg veed mig fri for enhver Skyld i denne Henseende; thi det er ikke mig der ved mig selv er bleven saa fripostig, det er netop Christendommen der nødsager mig dertil. Den lægger en ganske anderledes Vægt paa mit lille Jeg, og paa ethvert nok saa lille Jeg, da den vil gjøre ham evig salig, hvis han er saa heldig at komme ind i den. Uden nemlig at have fattet Christendommen, da jeg jo fremsætter Spørgsmaalet, har jeg dog forstaaet saa Meget, at den vil gjøre den Enkelte evig salig, og at den netop hos den Enkelte forudsætter denne uendelige Interesse for hans Salighed som conditio sine qua non, en Interesse, hvormed han hader Fader og Moder, og saaledes da vel ogsaa blæser af Systemer og verdenshistoriske Oversigter. Saa meget har jeg forstaaet, skjøndt staaende udenfor, at den eneste utilgivelige Majestæts-Forbrydelse mod Christendommen er, naar det enkelte Individ uden videre antager sit Forhold til den for givet. Hvor beskedent det end kunde synes saaledes at gaae med i Kjøbet, saa anseer Christendommen netop dette for Frækhed. Jeg maa derfor paa det ærbødigste frabede mig alle theocentriske Hjælpere og Hjælperes Hjælperes Bistand for paa den Maade at hjælpe mig ind i Christendommen. Saa vil jeg hellere blive hvor jeg er, ved min uendelige Interesse, ved Problemet, ved Muligheden. Det er nemlig ikke umuligt, at Den som er uendeligt interesseret for sin evige Salighed kan engang vorde evig salig, derimod er det vel umuligt, at Den som har tabt Sandsen derfor, (hvilken jo dog ikke kan være nogen anden end en uendelig Bekymring), at han kan blive evig salig. Ja naar man engang har tabt den, er det maaskee umuligt at faae den igjen. De fem daarlige Brudejomfruer, de havde tabt Forventningens uendelige Lidenskab. Saa gik Lampen ud. Da blev der gjort Anskrig, at Brudgommen kom. Nu løb de hen hos Kræmmeren og kjøbte ny Olie og vilde begynde paa en frisk og lade Alt være glemt. Det forstaaer sig, Alt var ogsaa glemt. Døren blev lukket, og de udelukkede, og Brudgommen sagde til dem, da de bankede paa: jeg kjender Eder ikke. Dette var ingen Spydighed af Brudgommen, men en Sandhed, thi i aandelig Forstand vare de blevne ukjendelige ved at have tabt den uendelige Lidenskab.

Det objektive Problem vilde da være: om Christendommens Sandhed. Det subjektive er: om Individets Forhold til Christendommen. Ganske simpelt: hvorledes jeg Johannes Climacus kan blive deelagtig i den Salighed som Christendommen forjætter. Problemet angaaer ganske alene mig: deels fordi det, hvis det er rigtigt fremsat, vil paa samme Maade angaae Enhver, deels fordi alle de Andre jo har Troen endogsaa som noget Givet, som en Ubetydelighed, de end ikke agte stort værd, eller som en Ubetydelighed, der først ved nogle Beviser bliver maiet ud til at være Noget. Saa er da Problemets Fremsættelse ikke ubeskedent af mig, men kun en Art Galskab.

For at nu mit Problem kan blive tydeligt, skal jeg først fremsætte det objektive Problem, og vise hvorledes dette behandles. Herved vil det Historiske skee sin Ret. Derpaa skal jeg fremsætte det subjektive. Det er i Grunden mere end den lovede Fortsættelse, som Iførelse af det historiske Costume, thi dette er givet blot ved at nævne Ordet: Christendom. Den første Deel er den lovede Fortsættelse, den anden Deel et fornyet Forsøg i samme Retning som Piecen, et nyt Tilløb til Smulernes Problem.

\sectionheading{Noget om Lessing (uddrag)  ·  SV¹ VII, 55–98}

Uden da at turde beraabe mig paa Lessing, uden med Bestemthed at turde angive ham som Hjemmelsmand, uden at forpligte Nogen til paa Grund af Navnkundigheden pligtskyldigst at ville eller forsikkrende at ville have forstaaet hvad der bringer den Forstaaende i et betænkeligt Forhold til min Navnløshed, som vel har noget ligesaa Afskrækkende som Lessings Navnkundighed har Tiltrækkende: agter jeg nu at fremsætte hvad jeg Fanden gale mig henfører til Lessing uden at være sikker paa, at han vedstaaer det; hvad jeg snart kunde fristes til drillende i Overgivenhed at ville paadutte ham at han har sagt det, om ikke ligefrem; hvad jeg paa en anden Maade sværmersk kunde ønske beundrende at turde takke ham for; hvad jeg atter med stolt Paaholdenhed og Selvfølelse kun i Genereusitet henfører til ham; og hvad jeg atter frygter for at fornærme eller besvære ham med ved at bringe hans Navn i Forhold dertil. Ja, man finder dog sjeldent en Forfatter, der er saa behagelig i Omgang som Lessing. Og hvoraf kommer det? Deraf, tænker jeg, at han er saa sikker paa sig selv. Al denne trivielle og magelige Omgang mellem en Udmærket og en mindre Udmærket, at den ene er Genie, Mester, den anden Lærling, Bud, Leietjener og saa videre, er her forhindret. Om jeg med Djævels Vold og Magt vilde være Lessings Discipel, jeg kan det ikke, han har forhindret det. Som han selv er fri, saa vil han frigjøre Enhver, det tænker jeg, i sit Forhold til ham, frabedende sig Lærlingens Uddunstninger og Labansstreger, befrygtende at blive til Latter ved Repetenterne: et eftersnakkende Ecchos sædvanlige Gjengivelse af det Sagte.

1. Den subjektive existerende Tænker er opmærksom paa Meddelelsens Dialektik.

Medens den objektive Tænkning er ligegyldig ved det tænkende Subjekt og dets Existents, er den subjektive Tænker som existerende væsentligen interesseret i sin egen Tænkning, er existerende i den. Derfor har hans Tænkning en anden Art af Reflexion, nemlig Inderlighedens, Besiddelsens, hvorved den tilhører Subjektet og ingen Anden. Medens den objektive Tænkning sætter Alt i Resultat, og forhjælper den hele Menneskehed til at snyde ved at afskrive og opramse Resultatet og Facitet, sætter den subjektive Tænkning alt i Vorden, og udelader Resultatet, deels fordi dette netop tilhører ham, da han har Veien, deels fordi han som existerende bestandigen er i Vorden, hvilket da ethvert Menneske er, der ikke har ladet sig narre til at blive objektiv, til overmenneskeligt at blive Speculationen.

Inderlighedens Reflexion er den subjektive Tænkers Dobbelt-Reflexion. Tænkende tænker han det Almene, men som existerende i denne Tænkning, som erhvervende denne i sin Inderlighed bliver han mere og mere subjektiv isoleret.

Forskjelligheden mellem den subjektive og den objektive Tænkning maa ogsaa yttre sig i Meddelelsens Form, det vil sige, den subjektive Tænker maa strax blive opmærksom paa, at Formen maa kunstnerisk have lige saa megen Reflexion som han selv existerende i sin Tænkning har det. Kunstnerisk vel at mærke, thi Hemmeligheden bestaaer ikke i, at han ligefrem udsiger Dobbelt-Reflexionen, da et saadant Udsagn netop er en Modsigelse.

Den objektive Tænkning er aldeles ligegyldig mod Subjektiviteten og derved mod Inderligheden og Tilegnelsen; dens Meddelelse er derfor ligefrem. Det følger af sig selv, at den ingenlunde derfor behøver at være let, men den er ligefrem, den har ikke Dobbelt-Reflexionens Svig og Kunst, den har ikke hiin subjektive Tænknings gudfrygtige og humane Omsorg i at meddele sig, den lader sig ligefrem forstaae, den lader sig ramse. Den objektive Tænkning er derfor blot opmærksom paa sig selv, og er derfor ingen Meddelelse, idetmindste ingen kunstnerisk Meddelelse, forsaavidt der dog altid vilde fordres at tænke den Modtagende og at agte paa Meddelelsens Form i Forhold til Modtagerens Misforstaaelse. Den objektive Tænkning er ligesom de fleste Mennesker saa inderlig god og meddeelsom; den meddeler sig uden videre og griber i det Høieste til Forsikkringer om sin Sandhed, Anbefalinger, Forjættelser om at alle Mennesker engang ville antage denne Sandhed – saa sikker er den. Eller maaskee snarere saa usikker, thi Forsikkringen og Anbefalingen og Forjættelsen, som jo rigtignok er for de Andres Skyld, som skulle antage den, turde maaskee ogsaa være for Lærerens Skyld, der trænger til Stemmefleerhedens Sikkerhed og Tilforladelighed. Nægter Samtiden ham den, saa trækker han paa Eftertiden – saa sikker er han. Denne Sikkerhed har noget tilfælleds med den Uafhængighed, der, uafhængig af Verden, behøver Verden som Vidne til sin Uafhængighed, for at være vis paa, at man er uafhængig.

[…]

3. Hvad der under dette og det følgende Punkt bliver Gjenstand for Omtalen, lader sig bestemtere henføre til Lessing, forsaavidt Udsagnet ligefrem lader sig citere, men dog atter ikke med nogen ligefrem Bestemthed, da Lessing ikke er docerende, men subjektivt eviterende, uden at ville forpligtige Nogen til for hans Skyld at antage det, uden at ville forhjælpe Nogen til en ligefrem Continuitet med Ophavsmanden. Maaskee har Lessing selv forstaaet, at Sligt ikke lader sig docere ligefrem; idetmindste lader Lessings Adfærd sig forklare saaledes, og maaskee er Forklaringen rigtig, maaskee.

Lessing har sagt (S.W. 5te Bind pagina 80), at tilfældige historiske Sandheder aldrig kunne blive Beviis for evige Fornuft-Sandheder; samt (pagina 83) at den Overgang, hvorved man paa en historisk Efterretning vil bygge en evig Sandhed, er et Spring.

Nu til Lessing. Stedet findes i en lille Opsats: Über den Beweis des Geistes und der Kraft; an den Herrn Director Schumann. L. opponerer mod hvad jeg vilde kalde at qvantitere sig ind i en qvalitativ Afgjørelse; han bekæmper den ligefremme Overgang fra historisk Tilforladelighed til Afgjørelse af en evig Salighed. Han nægter ikke (thi han veed strax at gjøre Concessioner for at Categorierne kunne blive tydelige), at hvad der i Skriften fortælles om Undere og Spaadomme, er ligesaa tilforladeligt som andre historiske Efterretninger, ja saa tilforladeligt som overhovedet historiske Efterretninger kunne være: aber nun, wenn sie nur eben so zuverlässig sind, warum macht man sie bei dem Gebrauche auf einmal unendlich zuverlässiger? (pagina 79), nemlig fordi man paa dem vil bygge Antagelsen af en Lære, der betinger en evig Salighed, altsaa paa dem vil bygge en evig Salighed. L. er villig til, ligesom alle Andre, at troe, at en Alexander har levet, der har undertvunget hele Asien, »aber wer wollte auf diesen Glauben hin irgend etwas von großem und dauerhaftem Belange, dessen Verlust nicht zu ersetzen wäre, wagen?« (pagina 81).

Det er bestandigt Overgangen, den ligefremme Overgang fra det historisk Tilforladelige til den evige Afgjørelse, Lessing bekæmper. Derfor stiller han sig saaledes, at han gjør en Forskjel mellem Efterretninger om Undere og Spaadomme – og Samtidigheden med disse. (Denne Forskjel har Smulerne været opmærksom paa ved digterisk experimenterende at tilveiebringe Samtidigheden, og saaledes fjerne hvad der blev kaldet det Efterhistoriske.) Af Efterretningerne, det vil sige af deres indrømmede Tilforladelighed, følger Intet, siger Lessing, men, tilføier han, havde han været samtidig med Miraklerne og Spaadommene, saa vilde det have hjulpet ham. Vel underrettet, som Lessing altid er, protesterer han derfor imod et halv svigefuldt Citat af Origenes, der er blevet anført for at give hiint Beviis for Christendommens Sandhed Relief. Han protesterer ved at føie Slutningen af Origenes's Ord til, hvoraf sees, at Orthodoxismus antager, at der endnu i hans Tid skeete Mirakler, og at han tillægger disse Mirakler, med hvilke han altsaa er samtidig, Beviiskraft ligesom de, om hvilke han læser.

Da Lessing har stillet sig saaledes, ligeoverfor en given Fremstilling, faaer han ikke Leilighed til at udhæve det dialektiske Problem, om Samtidigheden vilde hjælpe Noget, om den kunde være mere end Anledning, hvad den historiske Efterretning ogsaa kan være. Lessing synes at antage det Modsatte; maaskee er dog dette Skin tilveiebragt for at give sin Fægtning e concessis mere dialektisk Tydelighed ligeoverfor en bestemt enkelt Mand. Smulerne søgte derimod at vise, at Samtidigheden hjælper slet Intet, fordi der i al Evighed ingen ligefrem Overgang er, hvilket jo ogsaa vilde have været en grændseløs Uretfærdighed mod alle Senere, en Uretfærdighed og Forskjel, der var meget værre end den mellem Jøder og Græker, Omskaarne og Uomskaarne, hvilken Christendommen har hævet.

Lessing har selv samlet sit Problem i følgende Ord, hvilke han har spatieret: zufällige Geschichtswahrheiten können der Beweis von nothwendigen Vernunftswahrheiten nie werden. Hvad her støder, er det Prædikat: zufällige. Dette er vildledende, det kunde synes at lede til den absolute Distinction mellem væsentlige og tilfældige historiske Sandheder, en Distinction, der dog kun er en Subdivision. Distingveres der absolut tiltrods for det høiere Prædikats (»historiske«) Identitet, saa kunde deraf synes at følge, at i Forhold til væsentlige historiske Sandheder vilde den ligefremme Overgang være at danne. Jeg kunde nu blive hidsig og sige: det er umuligt, at Lessing kan være saa inconseqvent, ergo – og min Hidsighed vist overbevise Mange. Imidlertid indskrænker jeg mig til et høfligt Maaskee, der antager, at Lessing i Prædikatet: tilfældige har skjult Alt, men kun sagt Noget, saa at »tilfældige« ikke er et relativt-distingverende Prædikat eller inddelende, men et Genus-Prædikat: historiske Sandheder, der som saadanne ere tilfældige. Hvis ikke, ligger her den hele Misforstaaelse, der atter og atter gaaer igjen i den nyeste Philosophie: at lade det Evige uden videre blive historisk, og at kunne begribe det Historiskes Nødvendighed. Alt hvad der bliver historisk er tilfældigt; thi netop derved, at det bliver til, bliver historisk, har det sit Moment af Tilfældighed, thi Tilfældighed er just den ene Faktor i al Tilblivelse. – Deri ligger igjen Incommensurabiliteten mellem en historisk Sandhed og en evig Afgjørelse.

Saaledes forstaaet er Overgangen, hvorved et Historisk og Forholdet til dette bliver afgjørende for en evig Salighed, en μεταβασις εις αλλο γενος (Lessing siger endog, hvis dette ikke er det, saa veed jeg ikke hvad Aristoteles har forstaaet derved pagina 82), et Spring baade for den Samtidige og for den Senere. Den er et Spring, og det er dette Ord som Lessing har brugt indenfor den tilfældige Indskrænkning, der betegnes ved en illusorisk Forskjel imellem Samtidighed og Ikke-Samtidighed. Ordene lyde saaledes: Das, das ist der garstige breite Graben, über den ich nicht kommen kann so oft und ernstlich ich auch den Sprung versucht habe (pagina 83). Maaskee er hiint Ord: Spring kun en Vending i Stilen, maaskee udføres derfor Metaphoren for Phantasien ved at tilføie Prædikatet: breit, som havde ikke end det mindste Spring den Egenskab, at det gjør Graven uendelig bred, som var det ikke lige vanskeligt for den, der slet ikke kan springe, enten Graven er bred eller smal, som var det ikke den dialektisk lidenskabelige Afsky for et Spring, der gjør Graven saa uendelig bred, ligesom Lady Machbeths Lidenskab gjør Blodpletten saa uhyre stor, at ikke Oceanet kan aftvætte den. Maaskee er det ogsaa en List af Lessing at anbringe det Ord: ernstlich, thi i Forhold til det at springe, især naar Metaphoren udføres for Phantasien, er Alvoren pudseerlig nok, fordi den staaer i intet eller i et comisk Forhold til Springet, da det jo ikke, udvortes forstaaet, er Gravens Brede der forhindrer det, men indvortes den dialektiske Lidenskab, der gjør Graven uendelig bred. At have været ganske nær ved Noget, har allerede sin comiske Side, men at have været ganske nær ved at gjøre Springet, er slet Intet, netop fordi Springet er Afgjørelsens Categorie. Og nu med yderste Alvor at have villet gjøre Springet – ja det er en Skjelm den Lessing, thi det er nok snarere med yderste Alvor han har gjort Graven bred: er dette ikke ligesom at gjøre Nar af Folk! Dog kan man, som bekjendt, Springet betræffende, ogsaa paa en anden mere yndet Maade gjøre Nar af Folk: man lukker Øinene, tager sig selv i Nakken, à la Münchhausen, og saa – saa staaer man paa den anden Side, paa hiin Side af den sunde Menneskeforstand i det systematisk forjættede Land.

[…]

Altsaa: a) et logisk System kan der gives; b) men der kan ikke gives noget Tilværelsens System.

a.

α. Skal der imidlertid construeres et logisk System, saa maa der fornemligen agtes paa, at Intet optages, som er underkastet Tilværelsens Dialektik, som altsaa kun er ved at være til, eller ved at have været til, ikke er ved at være. Heraf følger ganske simpelt, at hiin mageløse og mageløst beundrede Opfindelse af Hegel, at bringe Bevægelse ind i Logiken (ogsaa fraseet, at man paa hvert andet Sted savner endog blot hans eget Forsøg paa at bilde En ind at den er der), netop er at forvirre Logiken. Det er jo ogsaa besynderligt, at lægge Bevægelsen til Grund i en Sphære, hvor Bevægelse er utænkelig, eller at lade Bevægelse forklare Logiken, medens Logiken ikke kan forklare Bevægelsen. Betræffende dette Punkt er jeg imidlertid saa heldig at kunne henvise til en Mand, der tænker sundt og er lykkeligt dannet ved Grækerne (sjeldne Egenskaber i vor Tid!), en Mand, der har vidst at frigjøre sig og sin Tænkning fra ethvert slæbende og krybende Forhold til Hegel, hvis Navnkundighed Enhver ellers vil profitere af om ikke paa anden Maade, saa ved at gaae videre, det vil sige ved at have optaget Hegel i sig; en Mand, der hellere har villet nøies med Aristoteles og med sig selv – jeg mener Trendlenburg (Logische Untersuchungen). Hans Fortjeneste er blandt Andet at have fattet Bevægelsen som den uforklarlige Forudsætning, som det Fælleds, hvori Væren og Tænken enes og som den fortsatte Gjensidighed. Jeg kan her ikke gjøre noget Forsøg paa at eftervise hans Opfattelses Forhold til Grækerne, til det Aristoteliske, eller til, hvad der besynderligt nok, i en vis Forstand, skjøndt kun populairt, har megen Lighed med hans Fremstilling: et lille Partie i Plutarchs Skrift om Isis og Osiris. Det er ingenlunde min Mening, at den hegelske Philosophie ikke har havt gavnlig Indflydelse paa Trendlenburg, men det Heldige er, at han har indseet, at det gaaer ikke an at ville forbedre paa Hegels Bygning, gaae videre og saa videre (en løgnagtig Maade paa hvilken mangen Stymper i vor Tid anmasser sig Hegels Celebritet og lazaronagtigt fraterniserer ham); og paa den anden Side, at han, ædruelig som en græsk Tænker, uden at love Alt, uden at ville lyksaliggjøre Menneskeheden, præsterer Meget og lyksaliggjør Den, der behøvede hans Veiledning i Kjendskab til Grækerne.

I et logisk System maa Intet optages, som har et Forhold til Tilværelse, som ikke er ligegyldigt mod Existents. Den uendelige Overvægt, som det Logiske ved at være det Objektive har over al Tænkning, limiteres igjen derved, at det subjektivt seet er en Hypothese, netop fordi det i Virkelighedens Forstand er ligegyldigt mod Tilværelse. Denne Dobbelthed er det Logiskes Forskjel fra det Mathematiske, der slet intet Forhold har til eller fra Tilværelse, men kun har Objektiviteten – ikke Objektiviteten og det Hypothetiske som Eenheden og Modsigelsen, hvori det negativt forholder sig til Existents.

Det logiske System maa ikke være en Mystifikation, et Bugtalerie, hvori Tilværelsens Indhold kommer underfundigt og subrept frem, hvor den logiske Tanke studser og finder, hvad Hr. Professoren eller Licentiaten her har havt bag Øret. Strængere lader sig her dømme imellem ved at besvare det Spørgsmaal, i hvilken Forstand Categorien er Tilværelsens Abbreviatur, om den logiske Tænkning er abstrakt efter Tilværelse eller abstrakt uden noget Forhold til Tilværelse. Dette Spørgsmaal kunde jeg ønske at behandle lidt udførligere paa et andet Sted; og selv om det ikke blev fyldestgjørende besvaret, er det dog altid Noget, at der bliver spurgt saaledes.

[…]

b.

Et Tilværelsens System kan ikke gives. Altsaa er et saadant ikke til? Ingenlunde. Dette ligger ei heller i det Sagte. Tilværelsen selv er et System – for Gud, men kan ikke være det for nogen existerende Aand. System og Afsluttethed svare til hinanden, men Tilværelse er netop det Modsatte. Abstract seet lader System og Tilværelse sig ikke tænke sammen, fordi den systematiske Tanke for at tænke Tilværelse maa tænke den som ophævet, altsaa ikke som tilværende. Tilværelse er det Spatierende, der falder ud fra hinanden: det Systematiske er Afsluttetheden, der slutter sammen.

I Virkeligheden indtræder nu en Skuffelse, et Sandsebedrag, som ogsaa Smulerne søgte at paavise, og hvortil jeg maa henvise, navnligen til det Spørgsmaal, om det Forbigangne er mere nødvendigt end det Tilkommende. Idet nemlig en Tilværelse er tilbagelagt, er den jo færdig, er den jo afsluttet, og forsaavidt er den da hjemfalden til den systematiske Opfattelse. Ganske rigtigt – men for hvem? Den der selv er existerende kan jo ikke vinde hiin Afsluttethed udenfor Tilværelsen, som er tilsvarende til den Evighed, til hvilken det Forbigangne er indgaaet. Om en Tænker godhedsfuld vil være distrait nok til at glemme, at han er selv existerende, saa er Spekulation og Distraction dog ikke ganske det Samme. Tvertimod, det at han selv er existerende antyder Existentsens Krav paa ham, at hans, ja hvis han er stor, at hans Tilværelse i Samtiden kan igjen som en Forbigangenhed have en Afsluttethedens Gyldighed for den systematiske Tænker. Men hvo er da denne systematiske Tænker? Ja det er Den, som selv er udenfor Tilværelse og dog i Tilværelsen, som i sin Evighed er for evigt afsluttet, og dog indeslutter Tilværelsen i sig – det er Gud. Hvortil Bedraget! Fordi Verden nu har bestaaet i 6000 Aar, har derfor Tilværelse ikke aldeles det samme Krav paa den Existerende, som den altid har havt, hvilket ikke er, at han i en Indbildning skal være en contemplerende Aand, men i Virkeligheden en existerende Aand. Al Forstaaelse kommer bag efter. Medens den nu Existerende unægtelig kommer bag efter i Forhold til de 6000 Aar, der gik forud, saa vilde, hvis vi antog, at han kom til systematisk at forstaae disse, det besynderlige Ironiske fremkomme, at han ikke kom til at forstaae sig selv existerende, fordi han selv ingen Existents fik, fordi han selv Intet havde, der skulde forstaaes bag efter. Heraf vilde følge, at en saadan Tænker enten maatte være vor Herre eller et phantastisk Qvodlibet. Enhver indseer vel det Usædelige heri, og saa indseer vist tillige Enhver, at det er i sin Orden, hvad en anden Forfatter har bemærket om det hegelske System: at man ved ham fik et System, det absolute System færdigt – uden at have en Ethik. Lad os kun smile af Middelalderens ethisk-religieuse Phantasterier i Ascese og andet Saadant, men lad os for Alting ikke glemme, at den speculative lav-comiske Udsvævelse med at blive Jeg-Jeg, og saa dog qva Menneske ofte en saadan Philister, at ingen Begeistret gad have ført et saadant Liv – er ligesaa latterlig.

\sectionheading{Den subjektive Sandhed, Inderligheden; Sandheden er Subjektiviteten (uddrag)  ·  SV¹ VII, 157–208}

Dersom der ved Væren i de tvende angivne Bestemmelser forstaaes den empiriske Væren, saa er Sandheden selv forvandlet til et Desideratur og Alt sat i Vorden, fordi den empiriske Gjenstand ikke er færdig, og den existerende erkjendende Aand jo selv er i Vorden, og saaledes Sandheden en Approximeren, hvis Begyndelse ikke kan absolut sættes, netop fordi der ingen Slutning er, hvilket har tilbagevirkende Kraft; hvorimod enhver Begyndelse (hvis den ikke, ved ikke at være sig dette bevidst, er en Vilkaarlighed), naar den gjøres, ikke skeer i Kraft af den immanente Tænkning, men gjøres i Kraft af en Beslutning, væsentligen i Kraft af Tro. At den erkjendende Aand er en existerende, og at ethvert Menneske er et saadant for sig Existerende, kan jeg ikke ofte nok gjentage; thi at man phantastisk har overseet dette, er Skyld i megen Forvirring. Ingen misforstaae mig. Jeg er nu saaledes en stakkels existerende Aand ligesom alle andre Mennesker, men hvis jeg paa en lovlig og redelig Maade kunde hjælpes til at blive noget Overordentligt, det rene Jeg-Jeg, saa er jeg altid villig til at takke for Skjenken og Velgjerningen. Kan det derimod kun skee paa den omtalte Maade ved at sige ein, zwei, drei kokolorum, eller ved at binde et Baand om den lille Finger og naar det er Fuldmaane kaste det hen paa et afsides Sted: saa bliver jeg hellere hvad jeg er, et stakkels existerende enkelt Menneske.

Væren maa da i hine Definitioner fattes meget mere abstrakt, som den abstrakte Gjengivelse af, eller det abstrakte Forbillede paa hvad Væren in concreto er som empirisk Væren. Saaledes forstaaet er der Intet til Hinder for, at Sandheden abstrakt bestemmes abstrakt som et Færdigt; thi Overeensstemmelsen mellem Tænken og Væren er abstrakt seet altid færdig, da Vordens Begyndelse netop ligger i Concretionen, fra hvilken Abstraktionen abstrakt seer bort.

Men forstaaes Væren saaledes, saa er Formelen en Tautologie, det vil sige, Tænken og Væren betyder Eet og det Samme, og den Overeensstemmelse, om hvilken der tales, er blot den abstrakte Identitet med sig selv. Ingen af Formlerne siger derfor mere end at Sandheden er, naar dette forstaaes saaledes at Copulaet accentueres, Sandheden er, det vil sige Sandheden er en Fordoblelse, Sandheden er det Første, men Sandhedens Andet, at den er, er det Samme som det Første, denne dens Væren er Sandhedens abstrakte Form. Paa denne Maade er der udtrykt, at Sandheden ikke er noget Enkelt, men aldeles abstrakt en Fordoblelse, der dog i samme Øieblik er hævet.

Abstraktionen kan nu blive ved at omskrive dette saa meget den vil, den kommer aldrig videre. Saasnart Sandhedens Væren bliver empirisk concret, er Sandheden selv i Vorden, er vel igjen, anet, Overeensstemmelsen mellem Tænken og Væren, og er det vel saaledes virkelig for Gud, men er det ikke saaledes for nogen existerende Aand, da denne selv existerende er i Vorden.

For den existerende Aand qva existerende Aand bliver Spørgsmaalet igjen om Sandheden; thi den abstrakte Besvarelse er kun for det Abstraktum, hvilket en existerende Aand bliver ved at abstrahere fra sig selv qva existerende, hvad den kun momentviis kan gjøre, medens den dog selv i saadanne Øieblikke betaler Tilværelsen sin Gjeld ved alligevel at existere. Altsaa det er en existerende Aand, der spørger om Sandheden, formodentlig for at ville existere i den, men i ethvert Tilfælde er Spørgeren sig bevidst at være et existerende enkelt Menneske. Saaledes troer jeg at kunne gjøre mig forstaaelig for enhver Græker, og for ethvert fornuftigt Menneske. Om en tydsk Philosoph følger sin Lyst til at skabe sig, og først skaber sig om til et overfornuftigt Noget, ligesom Guldmagere og Troldmænd udpynte sig phantastisk, for da paa en yderst tilfredsstillende Maade at besvare Spørgsmaalet om Sandheden: vedkommer ikke mig, saa lidet som hans tilfredsstillende Svar, der vistnok er yderst tilfredsstillende – naar man phantastisk er bleven udklædt. Om derimod en tydsk Philosoph gjør det eller ikke, vil Enhver let forvisse sig om, der begeistret samler sin Sjel paa at ville lade sig veilede af en saadan Viis, uden Critik blot følgsomt benyttende hans Veiledning ved at ville danne sin Existents efter den; netop naar man saaledes begeistret forholder sig som en Lærende til en saadan tydsk Professor, saa realiserer man det fortræffeligste Epigram over ham, thi en saadan Speculant er Intet mindre tjent med end en Lærendes redelige og begeistrede Iver for at udtrykke og realisere, for existerende at tilegne sig hans Viisdom, da den er et Noget, som Hr. Professoren har indbildt sig selv, skrevet Bøger om, men aldrig selv forsøgt, ja det er end ikke faldet ham ind, at det skulde gjøres. Der gives Speculanter, der, ligesom hiin Portskriver skrev hvad han ikke selv kunde læse, i den Formening at det blot var hans Sag at skrive, blot skrive og skrive hvad der, naar det, om jeg saa tør sige, skal læses ved Hjelp af Gjerning, viser sig at være Nonsens, med mindre det maaskee skulde være bestemt for phantastiske Væsener.

Idet der for den existerende Aand qva existerende bliver Spørgsmaal om Sandheden, kommer hiin Sandhedens abstrakte Reduplikation igjen, men Existentsen selv, Existentsen selv i Spørgeren, der jo existerer, holder de tvende Momenter ud fra hinanden, og Reflexionen udviser tvende Forhold. For den objektive Reflexion bliver Sandheden et Objektivt, en Gjenstand, og det gjælder om at see bort fra Subjektet; for den subjektive Reflexion bliver Sandheden Tilegnelsen, Inderligheden, Subjektiviteten, og det gjelder netop om existerende at fordybe sig i Subjektiviteten.

Men hvad saa? Skal vi saa blive ved denne Disjunktion, eller tilbyder her ikke Mediationen sin godhedsfulde Bistand, saa Sandheden bliver Subjekt-Objektet? Hvorfor ikke? Men kan da ogsaa Mediationen hjælpe den Existerende, saa længe han existerer, til selv at blive Mediationen, der jo er sub specie æterni, medens den stakkels Existerende er existerende? Det kan dog ikke hjælpe at gjøre Nar af et Menneske, at lokke ham med Subjekt-Objektet, naar han selv er forhindret i at komme i den Tilstand, hvor han kan forholde sig dertil, forhindret derved at han selv er i Vorden ved at existere. Hvad kan det hjælpe at forklare hvorledes evigt den evige Sandhed er at forstaae, naar Den, der skal benytte Forklaringen, er forhindret i at forstaae den saaledes, derved at han er existerende, og kun en Phantast, naar han indbilder sig at være sub specie æterni, naar han altsaa skal bruge netop den Forklaring: hvorledes den evige Sandhed er at forstaae i Tidens Bestemmelse af Den, der, ved at existere, selv er i Tiden, hvad jo den velbaarne Professor selv indrømmer, om ikke ellers saa naar han hvert Qvartal hæver sin Gage. Ved Mediationens Subjekt-Objekt ere vi blot komne tilbage til Abstraktionen, thi Sandhedens Bestemmelse som Subjekt-Objekt er aldeles det Samme som Sandheden er, det vil sige Sandheden er en Fordoblelse. Den høie Viisdom har altsaa igjen blot været distrait nok til at glemme, at det var en existerende Aand der spurgte om Sandheden. Eller er maaskee den existerende Aand selv Subjekt-Objektet? I saa Fald maatte jeg da spørge, hvor et saadant existerende Menneske er, der tillige er Subjekt-Objektet? Eller skal vi maaskee her igjen først forvandle den existerende Aand til et noget Overhovedet, og saa forklare Alt, undtagen det der spurgtes om, hvorledes et existerende Subjekt in concreto forholder sig til Sandheden, eller det der saa maatte spørges om, hvorledes da det enkelte existerende Subjekt forholder sig til dette Noget, der synes at have ikke lidet tilfælleds med en Papirs-Drage, eller med det Stykke Sukker, som Hollænderne hang op under Loftet, og som alle slikkede af.

Altsaa vende vi igjen tilbage til Reflexionens tvende Veie, og have ikke glemt at det er en existerende Aand, der spørger, et ganske enkelt Menneske, og kunne nu heller ikke glemme, at det at han er existerende netop vil forhindre ham i at gaae begge Veie paa eengang, og hans bekymrede Spørgsmaal vil forhindre ham i letsindigt-phantastisk at blive Subjekt-Objekt. Hvilken af Veiene er nu Sandhedens Vei for den existerende Aand, thi kun det phantastiske Jeg-Jeg er paa eengang færdigt med begge Veiene, eller gaaer methodice frem af begge Veiene paa eengang, hvilken Maade at gaae paa for et existerende Menneske er saa umenneskelig, at jeg ikke tør anbefale den.

Da den Spørgende netop udhæver det at han er Existerende, saa vil den Vei naturligviis især være at anprise, der især accentuerer det at existere.

Den objektive Reflexions Vei gjør Subjektet til det Tilfældige og derved Existents til et Ligegyldigt, Forsvindende. Bort fra Subjektet gaaer Veien til den objektive Sandhed, og medens Subjektet og Subjektiviteten bliver ligegyldig, bliver Sandheden det ogsaa, og netop dette er dens objektive Gyldighed, thi Interessen er, ligesom Afgjørelsen, Subjektiviteten. Den objektive Reflexions Vei fører nu til abstrakt Tænkning, til Mathematik, til historisk Viden af forskjellig Art, bestandig fører den bort fra Subjektet, hvis Tilværelse eller Ikke-Tilværelse bliver, objektivt ganske rigtigt, uendelig ligegyldig, ganske rigtigt, thi Tilværelse og Ikke-Tilværelse har som Hamlet siger kun subjektiv Betydning. I sit Maximum vil denne Vei føre til en Modsigelse, og forsaavidt Subjektet ikke ganske bliver sig selv ligegyldigt, er dette jo kun et Tegn paa, at hans objektive Stræben ikke er objektiv nok, i sit Maximum vil den føre til den Modsigelse, at kun Objektiviteten er blevet til, og Subjektiviteten gaaet ud, det vil sige den existerende Subjektivitet, der har gjort et Forsøg paa at blive hvad man i abstrakt Forstand kalder Subjektiviteten, den abstrakte Objektivitets abstrakte Form. Og dog er den Objektivitet, der er bleven til, subjektivt seet, i sit Maximum enten en Hypothese eller en Approximation, fordi al evig Afgjørelse netop ligger i Subjektiviteten.

Den objektive Vei mener imidlertid at have en Sikkerhed, som den subjektive Vei ikke har (og det forstaaer sig, Existents, det at existere og objektiv Sikkerhed kan ikke tænkes sammen), den mener at undgaae en Fare, som venter den subjektive Vei, og denne Fare er i sit Maximum Afsindighed. Ved den blot subjektive Bestemmelse af Sandheden bliver Galskab og Sandhed tilsidst ikke til at adskille, fordi de begge kunde have Inderligheden. Men ved at blive objektiv bliver man ikke gal. Maaskee maatte jeg dog her tillade mig en lille Bemærkning, som ikke synes overflødig i en objektiv Tidsalder. Inderlighedens Udebliven er ogsaa Galskab. Den objektive Sandhed som saadan afgjør ingenlunde, at Den, som udtaler den, er forstandig, tvertimod kan den endogsaa forraade at Manden er gal, uagtet det er aldeles sandt og især objektiv sandt hvad han siger. Jeg skal tillade mig at fortælle en Begivenhed, der, uden nogensomhelst Tillæmpning ved mig, er direkte hentet fra en Daarekiste. En Patient i en saadan Anstalt vil løbe bort og udfører ogsaa virkelig sit Forehavende ved at springe ud af Vinduet. Han befinder sig nu i Anstaltens Have, og vil tiltræde Befrielsens Vandring, da falder det ham ind (skal jeg nu sige, at han var klog nok eller gal nok til at faae dette Indfald?): naar Du nu kommer til Byen, bliver Du kjendt, og man vil formodentlig strax transportere Dig ud igjen, det gjælder da om, at Du ved Din Tales objektive Sandhed fuldelig overbeviser ethvert Menneske om, at hvad Din Forstand angaaer, da er Alt vel. Som han gaaer og tænker over dette, seer han en Keglekugle liggende paa Jorden, den tager han op, putter i Kjoleskjødets Lomme. For hvert Skridt han gjør, slaaer denne Kugle ham med Respekt at sige paa R–, og for hver Gang den saaledes rammer ham, siger han: bum, Jorden er rund. Han kommer til Hovedstaden, besøger strax en af sine Venner, han vil overbevise ham, at han ikke er gal, gaaer derfor op og ned af Gulvet og siger bestandigt: bum, Jorden er rund. Men er da Jorden ikke rund, fordrer Daarekisten endnu et Offer for denne Antagelses Skyld som i hine Tider, da Alle antoge den for at være saa flad som en Pandekage? Eller er den Mand gal, der ved at sige en almindelig antaget og almindelig agtet objektiv Sandhed haaber at bevise, at han ikke er gal? Og dog blev det Lægen netop derved tydeligt, at Patienten endnu ikke var helbredet, medens Helbredelsen dog vel ikke kunde dreie sig om, at faae ham til at antage, at Jorden var flad. Men det er ikke Alle, der ere Læger, og Tidens Fordring har betydelig Indflydelse paa det Spørgsmaal om Galskab, ja man skulde næsten stundom fristes til at antage, at den moderne Tid, der har moderniseret Christendommen, ogsaa har moderniseret Pilati Spørgsmaal, og at Tidens Trang til at finde Noget at hvile i, forkynder sig i det Spørgsmaal: hvad er Galskab? Naar en Privat-Docent for hver Gang Kjoleskjødet paaminder ham om at sige Noget, siger de omnibus dubitandum est, og skriver frisk væk paa et System, hvor man i hvert andet Punktum finder tilstrækkeligt indre Beviis for, at den Mand aldrig har tvivlet om Noget: saa ansees han ikke for gal. – Don Quixote er Forbilledet paa den subjektive Galskab, hvor Inderlighedens Lidenskab omfatter en fix enkelt endelig Forestilling. Naar Inderligheden derimod udebliver, saa indtræder den ramsende Galskab, hvilken er nok saa comisk, og som det var at ønske, at en experimenterende Psycholog vilde fremstille ved at tage en Haandfuld saadanne Philosopher og sætte sammen. Naar Afsindigheden er Inderlighedens Vildelse, er det Tragiske og det Comiske, at det Noget, som uendeligt angaaer den Ulykkelige, er en fixeret Enkelthed, som ikke angaaer noget Menneske. Naar Afsindigheden derimod er Inderlighedens Udebliven, er det Comiske, at det Noget, den Lyksalige veed, er det Sande, det Sande der angaaer hele Menneskeheden, men som slet ikke angaaer den meget ærede Ramser. Denne Art af Afsindighed er mere umenneskelig end den anden; man gyser ved at see hiin Første ind i Øiet, for ikke at opdage Vildhedens Dyb, men man tør slet ikke see paa den Anden af Frygt for at opdage, at han ikke har rigtige Øine, men Glasøine, og Haar af en Gulvmaatte, kort at han er et Kunstprodukt. Træffer man tilfældig sammen med en saadan Sindssyg, hvis Sygdom netop er at han intet Sind har, da hører man med en kold Gru paa ham; man veed ikke om man tør troe, at det er et Menneske, man taler med, eller maaskee en »Spadserestok«, en kunstig Opfindelse af Døbler, som skjuler et Positiv i sig. Og at have drukket Duus med Skarpretteren, kan jo altid være ubehageligt nok for en ærekjær Mand, men at være geraadet i en fornuftig og speculativ Samtale med en Spadserestok – det er næsten til at blive gal over.

Den subjektive Reflexion vender sig indefter mod Subjektiviteten og vil i denne Inderliggjørelse være Sandhedens, og saaledes at ligesom i det Foregaaende, naar Objektiviteten var bragt frem, Subjektiviteten var forsvundet, her selve Subjektiviteten bliver det Sidste, og det Objektive det Forsvindende. Her glemmes nu intet Øieblik, at Subjektet er existerende, og at det at existere er en Vorden, og at derfor hiin Sandhedens Identitet af Tænken og Væren er en Abstraktionens Chimaire, og i Sandhed kun en Skabningens Forlængsel, ikke fordi Sandheden ikke er dette, men fordi den Erkjendende er en Existerende, og saaledes Sandheden ikke kan være dette for ham, saa længe han existerer. Fastholdes ikke dette, saa gaae vi strax ved Hjælp af Speculationen ind i det phantastiske Jeg-Jeg, hvilket vel den nyere Speculation har brugt, men ikke forklaret om, hvorledes det enkelte Individ forholder sig dertil, og Herre Gud, mere end et enkelt Individ er nu engang intet Menneske.

Naar den Existerende virkeligen kunde være udenfor sig selv, saa vilde Sandheden være noget afsluttet for ham, men hvor er dette Punkt? Jeg-Jeget er et mathematisk Punkt, der slet ikke er til, forsaavidt kan Enhver gjerne indtage dette Standpunkt, den Ene staaer den Anden ikke i Veien. Kun momentviis kan det enkelte Individ existerende være i en Eenhed af Uendelighed og Endelighed, der er udover det at existere. Dette Moment er Lidenskabens Øieblik. Den moderne Speculation har opdrevet Alt for at Individet objektivt kan komme ud over sig selv; men dette lader sig slet ikke gjøre; Existentsen holder igjen, og dersom Philosopher nutildags ikke vare blevne Skriverkarle i Tjeneste hos en phantastisk Tænknings Stundesløshed, saa vilde den have indseet, at Selvmord var den eneste praktiske nogenlunde Fortolkning af dens Forsøg. Men den skrivende moderne Speculation lader ringe om Lidenskab; og dog er Lidenskab for en Existerende netop Existentsens Høieste – og Existerende ere vi nu engang. I Lidenskaben er det existerende Subjekt uendeliggjort i Phantasiens Evighed, og dog tillige netop allerbestemtest sig selv. Det phantastiske Jeg-Jeg er ikke Uendelighed og Endelighed i Identitet, thi hverken det Ene eller det Andet er virkeligt, det er en phantastisk Overeenskomst i Skyen, en ufrugtbar Omfavnelse, og det enkelte Jegs Forhold til dette Luftsyn aldrig angivet.

Al væsentlig Erkjenden angaaer Existents, eller kun den Erkjenden, hvis Forhold til Existents er væsentlig, er væsentlig Erkjenden. Den Erkjenden, som ikke ind efter i Inderlighedens Reflexion angaaer Existents, er væsentlig seet tilfældig Erkjenden, dens Grad og Omfang væsentlig seet ligegyldig. At den væsentlige Erkjenden væsentlig forholder sig til Existents, vil dog ikke betyde hiin ovenfor anførte Abstraktionens Identitet mellem Tænken og Væren, ei heller objektivt, at Erkjendelsen forholder sig til noget Tilværende som sin Gjenstand; men vil sige, at Erkjendelsen forholder sig til den Erkjendende, der væsentligen er en Existerende, og at derfor al væsentlig Erkjenden væsentlig forholder sig til Existents og til det at existere. Kun den ethiske og den ethisk-religieuse Erkjenden er derfor væsentlig Erkjenden. Men al ethisk og al ethisk-religieus Erkjenden er væsentlig sig forholdende til, at den Erkjendende existerer.

Mediationen er et Luftsyn ligesom Jeg-Jeg. Abstrakt seet er Alt, og der vorder Intet. I Abstraktionen kan da umulig Mediationen finde sin Plads, da den har Bevægelse som sin Forudsætning. Den objektive Viden kan vel have det Tilværende til Gjenstand, men da det vidende Subjekt er existerende, og selv i Vorden ved at existere, saa maa Speculationen først forklare, hvorledes et enkelt existerende Subjekt forholder sig til Erkjendelsen af Mediationen, hvad han i det Øieblik er, om han for Exempel i det Øieblik ikke noget nær er distrait; hvor han er, om han ikke er i Maanen? Man taler bestandigt om Mediationen og Mediationen; er da Mediationen et Menneske ligesom Peer Degn antager at Imprimatur er det? Hvorledes bærer et Menneske sig ad med at blive noget Saadant; læser man sig til denne Værdighed, dette store Philosophicum; eller bortgiver Magistraten den ligesom Klokker- og Gravertjenester? Man forsøge blot paa at indlade sig med disse og andre saadanne skikkelige Spørgsmaal af et skikkeligt Menneske, der jo ogsaa gjerne vilde være Mediationen, naar han kunde blive det paa en lovlig og ærlig Maade, og ikke enten ved at sige ein zwei drei kokolorum, eller ved at glemme, at han selv er et existerende Menneske, for hvem altsaa det at existere er noget Væsentligt, og det ethisk-religieust at existere et passende qvantum satis. En Speculant vil det maaskee forekomme abgeschmackt at spørge saaledes, men det gjelder især om, at man ikke polemiserer paa urette Sted, og altsaa ikke selv phantastisk-objektiv begynder et pro og contra, om der er Mediation eller ikke, men at man holder fast ved hvad det er at være Menneske.

Jeg skal nu, for at tydeliggjøre Vei-Forskjellen mellem den objektive og den subjektive Reflexion, vise den subjektive Reflexions Søgen tilbage ind efter i Inderlighed. Inderlighedens Høieste i et existerende Subjekt er Lidenskab, til Lidenskab svarer Sandheden som et Paradox, og det at Sandheden bliver Paradoxet, er netop begrundet i dens Forhold til et existerende Subjekt. Saaledes svarer det Ene til det Andet. Ved at glemme, at man er existerende Subjekt, gaaer Lidenskaben ud, og Sandheden bliver til Gjengjeld ikke noget Paradox, men det erkjendende Subjekt bliver fra at være et Menneske et phantastisk Noget, og Sandheden en phantastisk Gjenstand for dettes Erkjenden

Naar der objektivt spørges om Sandheden, reflekteres der objektivt paa Sandheden som en Gjenstand, til hvilken den Erkjendende forholder sig. Der reflekteres ikke paa Forholdet, men paa at det er Sandheden, det Sande han forholder sig til. Naar dette, han forholder sig til, blot er Sandheden, det Sande, saa er Subjektet i Sandheden. Naar der subjektivt spørges om Sandheden, reflekteres der subjektivt paa Individets Forhold; naar blot dette Forholds Hvorledes er i Sandhed, saa er Individet i Sandhed, selv om det saaledes forholdt sig til Usandheden. Lad os tage som Exempel Guds Erkjendelse. Objektivt reflekteres der paa, at det er den sande Gud; subjektivt paa, at Individet forholder sig til et Noget saaledes, at hans Forhold i Sandhed er et Guds-Forhold. Paa hvilken af Siderne er nu Sandheden? Ak, maae vi ikke her tye til Mediationen, og sige: den er paa ingen af Siderne, den er i Mediationen? Fortræffeligt sagt, dersom blot saa Nogen kunde sige, hvorledes en Existerende bærer sig ad med at være i Mediationen; thi at være i Mediationen er at være færdig, at existere er at vorde. En Existerende kan heller ikke være to Steder paa eengang, være Subjekt-Objekt. Naar han er nærmest efter at være to Steder paa eengang, er han i Lidenskab, men Lidenskab er kun momentviis, og Lidenskab er netop Subjektivitetens Høieste. – Den Existerende, der vælger den objektive Vei, gaaer nu ind i hele den approximerende Overveielse, der objektivt vil bringe Gud frem, hvilket i al Evighed ikke naaes, fordi Gud er Subjekt, og derfor kun for Subjektiviteten i Inderlighed. Den Existerende, der vælger den subjektive Vei, fatter i samme Øieblik den hele dialektiske Vanskelighed ved at han skal bruge nogen Tid, maaskee lang Tid, for at finde Gud objektivt; han fatter denne dialektiske Vanskelighed i hele dens Smerte, fordi han skal bruge Gud i samme Øieblik, fordi ethvert Øieblik er spildt, hvori han ikke har Gud. I samme Øieblik har han Gud ikke i Kraft af nogen objektiv Overveielse, men i Kraft af Inderlighedens uendelige Lidenskab. Den Objektive generes ikke af saadanne dialektiske Vanskeligheder som den: hvad det vil sige, at udvise en heel Forskningens Tid til at finde Gud – da det jo var muligt, at Forskeren døde imorgen, og hvis han levede, han jo dog ikke kunde betragte Gud som Noget, man efter Leilighed tager med, da Gud netop er Noget man a tout prix tager med, hvilket i Lidenskabens Forstaaelse netop er Inderlighedens sande Forhold til Gud.

Det er paa dette dialektisk saa vanskelige Punkt, at Veien svinger af for Den, der veed hvad det er at dialektisere og existerende at dialektisere, hvilket er noget Andet end at sidde som et phantastisk Væsen ved Skriverbordet og skrive hvad man selv aldrig har gjort, hvilket er noget Andet end at skrive de omnibus dubitandum, og selv existerende være ligesaa lettroende som det sandseligste Menneske – her er det Veien svinger af, og Forandringen er denne, at medens den objektive Viden gaaer i Mag frem af Approximationens lange Vei, selv ikke tilskyndet af Lidenskab, er for den subjektive Viden ethvert Ophold livsfarligt, og Afgjørelsen saa uendelig vigtig, at det strax er saa nødende, som var Leiligheden allerede gaaet ubenyttet hen.

Naar nu Regnestykket er dette: hvor er der meest Sandhed, enten paa Dens Side (og at være paa eengang ligeligt paa begge Sider er, som sagt, ikke forundt en Existerende, men kun en lyksaliggjørende Indbildning for et indbildt Jeg-Jeg), der ene objektivt søger den sande Gud og Guds-Forestillingens approximerende Sandhed, eller paa Dens Side, der er uendelig bekymret om, at han i Sandhed forholder sig til Gud med Trangens uendelige Lidenskab: saa kan Svaret ikke være tvivlsomt for Nogen, der ikke er aldeles forkludret ved Hjælp af Videnskab. Dersom En, der lever midt i Christendommen, gaaer op i Guds Huus, i den sande Guds Huus, med den sande Forestilling om Gud i Viden, og nu beder, men beder i Usandhed; og naar En lever i et afgudisk Land, men beder med Uendelighedens hele Lidenskab, skjøndt hans Øie hviler paa en Afguds Billede: hvor er saa meest Sandhed? Den Ene beder i Sandhed til Gud, skjøndt han tilbeder en Afgud; den Anden beder i Usandhed til den sande Gud, og tilbeder derfor i Sandhed en Afgud.

Naar En forsker objektivt efter Udødeligheden, en Anden sætter Uendelighedens Lidenskab ind paa Uvisheden: hvor er der saa meest Sandhed, og hvo har meest Vished? Den Ene er eengang for alle gaaet ind i en Approximeren, som aldrig ender, thi Udødelighedens Vished ligger jo netop i Subjektiviteten; den Anden er udødelig og kæmper netop derfor ved at stride mod Uvisheden. Lad os betragte Socrates. Nuomstunder fusker Enhver jo paa nogle Beviser, Een har flere, en Anden færre. Men Socrates! Han henstiller Spørgsmaalet objektivt problematisk: dersom der er en Udødelighed. Han var da altsaa en Tvivler i Sammenligning med en af de moderne Tænkere af tre Beviser? Ingenlunde. Paa dette »dersom« sætter han hele sit Liv ind, han vover at døe, og han har indrettet hele sit Liv med Uendelighedens Lidenskab saaledes, at det maatte findes antageligt – dersom der er en Udødelighed. Gives der noget bedre Beviis for Sjelens Udødelighed? Men De, der have tre Beviser, indrette slet ikke deres Liv derefter; hvis der er en Udødelighed, maa den væmmes ved deres Maade at leve paa: gives der noget bedre Modbeviis mod de tre Beviser? Uvishedens »Smule« hjalp Socrates, fordi han selv hjalp til med Uendelighedens Lidenskab; de tre Beviser gavne de Andre slet Intet, fordi de dog ere og blive Dødbidere, og have ved deres tre Beviser i Mangel af at bevise Andet netop beviist det. Saaledes har maaskee ogsaa en Pige ved et svagt Haab om at være elsket af den Elskede eiet al Forelskelsens Sødme, fordi hun selv satte Alt ind paa dette svage Haab: mangen Ægtemadame derimod, som mere end eengang har ligget under for Elskovens stærkeste Udtryk, har vel havt Beviser, og dog besynderligt nok, ikke eiet det quod erat demonstrandum. Den socratiske Uvidenhed var saaledes det med Inderlighedens hele Lidenskab fastholdte Udtryk for, at den evige Sandhed forholder sig til en Existerende, og derfor maa blive ham et Paradox, saalænge han existerer, og dog turde det være muligt, at der i den socratiske Uvidenhed i Socrates var mere Sandhed end i hele Systemets objektive Sandhed, der coqueterer med Tidens Fordringer og lemper sig efter Privat-Docenter.

Objektivt accentueres: hvad der siges; subjektivt: hvorledes det siges. Allerede æsthetisk gjelder denne Distinction, og udtrykkes bestemt saaledes, at hvad der er Sandhed, kan i Den og Dens Mund blive Usandhed. Denne Distinction er i disse Tider særligen at agte paa, thi skulde man i en eneste Sætning udtrykke Forskjellen mellem Oldtiden og vor Tid, saa maatte man vel sige: at i Oldtiden var der kun Enkelte der vidste det Sande, nu veed Alle den, men Inderligheden staaer i et omvendt Forhold dertil. Æsthetisk opfattes Modsigelsen, som fremkommer ved at Sandheden bliver Usandhed i Den og Dens Mund, bedst comisk. Ethisk-religieus accentueres igjen: hvorledes; dog forstaaes det ikke om Anstand, Modulation, Foredrag og saa videre, men det forstaaes om den Existerendes Forhold til det Udsagte i selve sin Existents. Objektivt spørges der blot om Tankebestemmelserne, subjektivt om Inderligheden. I sit Maximum er dette Hvorledes Uendelighedens Lidenskab, og Uendelighedens Lidenskab er selve Sandheden. Men Uendelighedens Lidenskab er netop Subjektiviteten, og saaledes er Subjektiviteten Sandheden. Objektivt seet er der ingen uendelig Afgjørelse, og saaledes er det objektivt rigtigt, at Forskjellen mellem Godt og Ondt er hævet med Modsigelsens Grundsætning, og derved ogsaa den uendelige Forskjel mellem Sandhed og Løgn. Kun i Subjektiviteten er Afgjørelse, hvorimod det at ville blive objektiv er Usandheden. Uendelighedens Lidenskab er det Afgjørende, ikke dens Indhold, thi dens Indhold er netop den selv. Saaledes er det subjektive Hvorledes og Subjektiviteten Sandheden.

Men det Hvorledes, der subjektivt accentueres, er tillige, netop fordi Subjektet er existerende, dialektisk i Retning af Tid. I Lidenskabens Afgjørelses Øieblik, hvor Veien svinger af fra den objektive Viden, seer det ud som var dermed den uendelige Afgjørelse færdig. Men i samme Øieblik er den Existerende i Timeligheden, og det subjektive Hvorledes er forvandlet til en Stræben, som er impulseret og gjentagent forfrisket af Uendelighedens afgjørende Lidenskab, men som dog er en Stræben.

Naar Subjektiviteten er Sandheden, maa Sandhedens Bestemmelse tillige indeholde i sig et Udtryk for Modsætningen til Objektiviteten, en Erindring fra hiint Veiskille, og dette Udtryk angiver da tillige Inderlighedens Spændstighed. Her er en saadan Definition paa Sandhed: den objektive Uvished, fastholdt i den meest lidenskabelige Inderligheds Tilegnelse, er Sandheden, den høieste Sandhed der er for en Existerende. Der hvor Veien svinger af (og hvor det er, lader sig ikke objektivt sige, da det netop er Subjektiviteten), bliver den objektive Viden sat i Bero. Objektivt har han da kun Uvisheden, men netop dette strammer Inderlighedens uendelige Lidenskab, og Sandheden er netop dette Vovestykke, med Uendelighedens Lidenskab at vælge det objektivt Uvisse. Jeg betragter Naturen for at finde Gud, jeg seer jo ogsaa Almagt og Viisdom, men jeg seer tillige meget Andet som ængster og forstyrrer. Summa summarum heraf bliver den objektive Uvished, men netop derfor er Inderligheden saa stor, fordi Inderligheden omfatter den objektive Uvished med Uendelighedens hele Lidenskab. I Forhold til en mathematisk Sætning for Exempel er Objektiviteten given, men derfor er dens Sandhed ogsaa en ligegyldig Sandhed.

Men den givne Bestemmelse af Sandhed er en Omskrivning af Tro. Uden Risico ingen Tro. Tro er netop Modsigelsen mellem Inderlighedens uendelige Lidenskab og den objektive Uvished. Kan jeg objektivt gribe Gud, saa troer jeg ikke, men netop fordi jeg ikke kan det, derfor maa jeg troe; og vil jeg bevare mig i Troen, maa jeg bestandig passe paa, at jeg fastholder den objektive Uvished, at jeg, i den objektive Uvished er »paa de 70,000 Favne Vand«, og dog troer.

I den Sætning, at Subjektiviteten, Inderligheden er Sandheden, er indeholdt den socratiske Viisdom, hvis udødelige Fortjeneste er netop at have agtet den væsentlige Betydning af at existere, af at den Erkjendende er existerende, hvorfor han i sin Uvidenhed i høieste Forstand indenfor Hedenskabet var i Sandheden. At fatte dette, at Speculationens Ulykke netop er at den atter og atter glemmer, at den Erkjendende er en Existerende, kan allerede være vanskeligt nok i vor objektive Tid. »Men at gaae videre end Socrates, naar man end ikke har fattet det Socratiske, det er idetmindste ikke Socratisk.« Confer Smulernes Moral.

Lad os nu ligesom i Smulerne fra dette Punkt forsøge en Tankebestemmelse, der virkelig gaaer videre. Om den er sand eller den er usand, har jeg her Intet med at gjøre, da jeg blot er experimenterende, men saa Meget maa der fordres, at det er tydeligt, at det Socratiske er forstaaet deri, saa jeg da idetmindste ikke kommer bag ved Socrates igjen.

Naar Subjektiviteten, Inderligheden er Sandheden, saa er Sandheden objektivt bestemmet Paradoxet; og det at objektivt Sandheden er Paradoxet, viser netop, at Subjektiviteten er Sandheden, da jo Objektiviteten støder fra, og Objektivitetens Frastød eller Udtrykket for Objektivitetens Frastød er Inderlighedens Spændstighed og Kraftmaaler. Paradoxet er den objektive Uvished, der er Udtrykket for Inderlighedens Lidenskab, hvilken netop er Sandheden. Saaledes det Socratiske. Den evige, væsentlige Sandhed, det vil sige den, der forholder sig væsentlig til en Existerende ved væsentlig at angaae det at existere (al anden Viden er socratisk seet tilfældig, dens Grad og Omfang ligegyldig), er Paradoxet. Dog er den evige væsentlige Sandhed selv ingenlunde Paradoxet, men er det ved at forholde sig til en Existerende. Den socratiske Uvidenhed er Udtrykket for den objektive Uvished, den Existerendes Inderlighed er Sandheden. For her strax at anticipere, bemærkes Følgende: den socratiske Uvidenhed er et Analogon til det Absurdes Bestemmelse, kun at der i det Absurdes Støden fra, er endnu mindre objektiv Vished, da der kun er Vished om, at det er absurdt, og netop derfor uendelig større Spændstighed i Inderligheden; den socratiske Inderlighed i at existere er et Analogon til Troen, kun at dennes Inderlighed, som svarende ikke til Uvidenhedens Frastød, men til det Absurdes, er uendelig dybere.

Socratisk er den evige væsentlige Sandhed ingenlunde paradox i sig selv, men kun ved at forholde sig til en Existerende. Dette er udtrykt i en anden socratisk Sætning: at al Erkjenden er en Erindren. Denne Sætning er en Antydning af Speculationens Begynden, men Socrates forfulgte den derfor heller ikke, væsentligen blev den platonisk. Her er det, hvor Veien svinger af, og Socrates væsentligen accentuerer det at existere, medens Plato, glemmende dette, fortaber sig i Speculation. Socrates's uendelige Fortjeneste er netop den at være en existerende Tænker, ikke en Speculant, der glemmer hvad det er at existere. For Socrates betyder derfor det, at al Erkjenden er en Erindren, i Afskedens Øieblik og som en bestandig ophævet Mulighed af at speculere, for ham betyder det det Dobbelte: 1) at den Erkjendende væsentligen er integer, og at der for ham ingen anden Mislighed er med Hensyn til Erkjendelsen af den evige Sandhed end den, at han existerer, hvilken Mislighed er ham saa væsentlig og afgjørende, at den betyder, at det at existere, Inderliggjørelsen i og ved at existere, er Sandheden; 2) at Existentsen i Timeligheden ikke har nogen afgjørende Betydning, fordi der bestandigt er Muligheden af erindrende at tage sig selv tilbage i Evigheden, om denne Mulighed end bestandigt ophæves derved, at Inderliggjørelse i at existere udfylder Tiden.

Det uendelig Fortjenstlige ved det Socratiske var netop at accentuere, at den Erkjendende er Existerende, og at det at existere det Væsentlige. At gaae videre ved ikke at forstaae dette, er kun en maadelig Fortjeneste. Dette maae vi da altsaa have in mente, og saa see om Formelen dog ikke lod sig forandre saaledes, at der virkelig gaaes videre end det Socratiske.

Altsaa Subjektiviteten, Inderligheden er Sandheden; gives der nu et inderligere Udtryk derfor? Ja, hvis den Tale: Subjektiviteten, Inderligheden er Sandheden, begynder saaledes: Subjektiviteten er Usandheden. Man forhaste sig ikke. Speculationen siger ogsaa, at Subjektiviteten er Usandheden, men siger det lige i den modsatte Retning, nemlig hen til at Objektiviteten er Sandheden. Speculationen bestemmer Subjektiviteten negativt hen til Objektiviteten. Den anden Bestemmelse forhindrer sig selv derimod, idet den vil begynde, hvilket netop gjør Inderligheden langt inderligere. Socratisk er Subjektiviteten Usandheden, hvis den ikke vil fatte, at Subjektiviteten er Sandheden, men for Exempel vil være objektiv. Her derimod er Subjektiviteten, idet den vil begynde paa at blive Sandheden, ved at blive subjektiv, i den Vanskelighed, at den er Usandheden. Saa gaaer Arbeidet tilbage, tilbage nemlig i Inderlighed. Det er saa langtfra, at Veien saaledes skulde være hen til det Objektive, at Begyndelsen blot ligger endnu dybere i Subjektiviteten.

Men Usandheden kan Subjektet ikke være evigt, eller evigt forudsættes at have været, den maa være bleven det i Tiden, eller blive det i Tiden. Det socratiske Paradox laa i, at den evige Sandhed forholdt sig til en Existerende, men nu har Existentsen anden Gang mærket den Existerende; der er foregaaet en saa væsentlig Forandring med ham, at han ingenlunde kan socratisk erindrende tage sig selv tilbage i Evigheden. At gjøre dette er at speculere, at kunne gjøre dette, men at ophæve Muligheden deraf ved at fatte Inderliggjørelsen i Existents, er det Socratiske; men nu er Vanskeligheden den, at det der fulgte Socrates som en ophævet Mulighed, er bleven en Umulighed. Var det at speculere allerede i Forhold til det Socratiske en mislig Fortjeneste, nu er det kun Confusion.

Paradoxet fremkommer naar den evige Sandhed og det at existere sættes sammen, men for hver Gang det at existere mærkes, bliver Paradoxet tydeligere og tydeligere. Socratisk seet var den Erkjendende en Existerende, men nu er den Existerende saaledes mærket, at Existentsen har foretaget en væsentlig Forandring ved ham.

Lad os nu kalde Individets Usandhed Synd. Evigt seet kan han ikke være Synd eller forudsættes evigt at have været i den. Altsaa ved at blive til (thi Begyndelsen var jo, at Subjektiviteten er Usandheden) bliver han Synder. Han fødes ikke som Synder i den Forstand, at han forudsættes som Synder førend han fødtes, men han fødes i Synd og som Synder. Dette kunne vi jo kalde Arvesynden. Men har Existentsen saaledes faaet Magt over ham, da er han forhindret i ved Erindring at tage sig selv tilbage i Evigheden. Var det allerede paradox, at den evige Sandhed forholdt sig til en Existerende, nu er det absolut paradox, at den forholder sig til en saadan Existerende. Men jo vanskeligere det er gjort ham, erindrende at tage sig ud af Existents, desto inderligere kan hans Existeren blive i Existentsen; og naar det er gjort ham umuligt, naar han stikker saaledes i Existents, at Erindringens Bagdør for evigt er lukket, saa bliver Inderligheden den dybeste. Men lader os aldrig glemme, at det socratisk Fortjenstfulde netop var at udhæve, at den Erkjendende er existerende, thi jo vanskeligere Sagen bliver, desto mere fristes man at ile ad Speculationens lette Vei fra Forfærdelser og Afgjørelser til Navnkundighed, Ære, gode Dage og saa videre Fattede allerede Socrates Misligheden ved at tage sig selv speculerende ud af Existents tilbage i Evigheden, da der dog ingen Mislighed var ved den Existerende uden den, at han existerede, og saa den, at det at existere var det Væsentlige: nu er det umuligt. Fremad maa han, tilbage er det umuligt.

Subjektiviteten er Sandheden. Ved at den evige væsentlige Sandhed forholdt sig til den Existerende, blev Paradoxet til. Lad os nu gaae videre, lad os antage, at den evige væsentlige Sandhed selv er Paradoxet. Hvorledes fremkommer Paradoxet? Ved at den evige væsentlige Sandhed og det at existere sættes sammen. Naar vi da altsaa sætte det sammen i Sandheden selv, saa bliver Sandheden et Paradox. Den evige Sandhed er bleven til i Tiden. Dette er Paradoxet. Blev Subjektet i det nærmest Foregaaende forhindret i at tage sig selv tilbage i Evigheden ved Synden, nu skal det ikke bekymre sig derover, thi nu er den evige væsentlige Sandhed ikke bag ved, men kommen foran det ved selv at existere eller have existeret, saa hvis Individet ikke existerende, i Existentsen, faaer fat paa Sandheden, faaer det den aldrig.

Skarpere kan Existents aldrig accentueres end den nu er bleven det. Speculationens Svig med at ville erindre sig ud af Existentsen er gjort umulig. Her kan kun være Tale om at fatte dette, enhver Speculation, der vil være Speculation, viser eo ipso, at den ikke har fattet dette. Individet kan støde alt dette fra sig og tye til Speculationen, men antage det og saa ville hæve det ved Speculationen er umuligt, fordi det er lige beregnet paa at forhindre Speculationen.

Naar den evige Sandhed forholder sig til en Existerende, bliver den Paradoxet. Paradoxet støder i den objektive Uvished og Uvidenheden fra i den Existerendes Inderlighed. Men da Paradoxet ikke i sig selv er Paradoxet, støder det ikke inderlig nok fra; thi uden Risico ingen Tro, jo mere Risico, jo mere Tro; jo mere objektiv Tilforladelighed, jo mindre Inderlighed (thi Inderligheden er netop Subjektiviteten); jo mindre objektiv Tilforladelighed, jo dybere er den mulige Inderlighed. Naar Paradoxet selv er Paradoxet, støder det fra i Kraft af det Absurde, og Inderlighedens Lidenskab, der svarer dertil, er Troen. – Men Subjektiviteten, Inderligheden er Sandheden, thi ellers have vi glemt det socratisk Fortjenstlige; men der gives intet stærkere Udtryk for Inderligheden, end naar Retiraden ud af Existentsen erindrende ind i Evigheden er gjort umulig: da med Sandheden imod sig som Paradoxet, i Syndens Angest og med denne Smerte, med Objektivitetens uhyre Risico – at troe. Men uden Risico ingen Tro, end ikke den socratiske, end mindre den, vi her tale om.

Naar Socrates troede, at Gud er til, da fastholdt han den objektive Uvished med Inderlighedens hele Lidenskab, og i denne Modsigelse, i denne Risico er netop Troen. Nu er det anderledes, istedetfor den objektive Uvished, er her Visheden om at det objektivt seet er det Absurde, og dette Absurde fastholdt i Inderlighedens Lidenskab er Troen. Den socratiske Uvidenhed er som en vittig Spøg i Sammenligning med det Absurdes Alvor, og den socratiske existerende Inderlighed som en græsk Sorgløshed i Sammenligning med Troens Anstrængelse.

Hvilket er nu det Absurde? Det Absurde er, at den evige Sandhed er bleven til i Tiden, at Gud er blevet til, er født, har voxet og saa videre, er blevet til aldeles som det enkelte Menneske, ikke til at skjelne fra et andet Menneske, thi al umiddelbar Kjendelighed er for-socratisk Hedenskab og jødisk seet Afgudsdyrkelse; og enhver Bestemmelse af Det, som virkelig gaaer videre end det Socratiske, maa væsentligen have et Mærke af, at det staaer i Forhold til dette, at Guden er blevet til, fordi Tro, sensu strictissimo, som udviklet blev i Smulerne, refererer sig til Tilblivelse. Naar Socrates troede, at Gud er til, da indsaae han vel, at der hvor Veien svinger af, er der en objektiv Approximerens Vei, ved for Exempel at betragte Naturen, Verdenshistorien og saa videre Hans Fortjeneste var netop at flye denne Vei, hvor den qvantiterende Sirenesang fortryller og narrer den Existerende. I Forhold til det Absurde er den objektive Approximeren, der er som den Comedie Misforstaaelse paa Misforstaaelse, som i Almindelighed opføres af Privat-Docenter og Speculanter.

Det Absurde er netop ved det objektive Frastød Troens Kraftmaaler i Inderlighed. Altsaa det er en Mand, der vil have Troen; lad nu Comedien begynde. Han vil have Troen, men ved Hjælp af den objektive Overveielse og Approximeren vil han sikkre sig. Hvad skeer? ved Hjælp af Approximeringen bliver det Absurde noget Andet, det bliver sandsynligt, det bliver mere sandsynligt, det bliver maaskee overmaade og særdeles sandsynligt. Nu er det da, han skal til at troe det, og han tør sige om sig selv, at han troer ikke som Skomagere og Skrædere og Eenfoldige, men først efter lang Overveielse. Nu skal han da til at troe det; men see, nu er det netop blevet umuligt at troe det. Det næstendeels Sandsynlige, det Sandsynlige, det overmaade og særdeles Sandsynlige, det kan han næstendeels og saa godt som vide, overmaade og særdeles næstendeels vide – men troe det, lader sig ikke gjøre; thi det Absurde er netop Troens Gjenstand og det Eneste, der lader sig troe. – Eller det er en Mand, der siger, han har Troen, men han vil nu gjøre sig sin Tro tydelig, han vil forstaae sig i sin Tro. Nu begynder Comedien igjen. Troens Gjenstand bliver næstendeels sandsynlig, den bliver saa godt som sandsynlig, den bliver sandsynlig, den bliver overmaade og særdeles sandsynlig. Han er færdig, han tør sige om sig selv, at han ikke troer som Skomagere og Skrædere eller andre Eenfoldige, men han har tillige forstaaet sig i sin Troen. Forunderlige Forstaaelse; han har tvertimod faaet noget Andet at vide om Troen end han troede, og faaet at vide, at han ikke mere troer, da han næstendeels veed, saa godt som veed, overmaade og særdeles næstendeels veed.

Forsaavidt det Absurde har Tilblivelsens Moment i sig, vil en Approximations-Vei ogsaa være den, der forvexler Tilblivelsens absurde Faktum, der er Troens Gjenstand, med et simpelt historisk Faktum, og altsaa søger historisk Vished for det, der netop er det Absurde, fordi det indeholder den Modsigelse, at Det der kun lige stik mod al menneskelig Forstand kan blive det Historiske er blevet det. Denne Modsigelse er netop det Absurde, der kun kan troes; faaes der en historisk Vished, saa faaer man blot Vished for, at det, der er vist, er ikke det Omspurgte. Et Vidne kan bevidne, at han har troet det, og altsaa at det saa langtfra at være en historisk Vished er lige mod hans Forstand, men et saadant Vidne støder jo fra i samme Forstand som det Absurde, og et Vidne, der ikke saaledes støder fra, er eo ipso en Bedrager, eller en Mand, der taler om noget ganske Andet, og et saadant Vidne kan da heller ikke hjælpe uden til at faae Vished om noget ganske Andet. 100,000 enkelte Vidner, der netop ved deres Vidnesbyrds særlige Beskaffenhed (at de have troet det Absurde) blive enkelte Vidner, blive ei heller en masse til noget Andet, saa det Absurde blev mindre absurd – hvorfor? fordi 100,000 Mennesker hver for sig havde troet, at det var absurd? Tvertimod, hine 100,000 Vidner støde igjen fra aldeles ligesom det Absurde. – Dog dette behøver jeg ikke her videre at udføre, det har jeg i Smulerne (især der hvor Forskjelligheden ophæves mellem Discipelen paa første og anden Haand), og i første Deel af denne Bog, udført omhyggeligt nok, at al Approximation gavner til Intet, da det tvertimod gjelder om, at skaffe Indledningsbetragtninger, Tilforladeligheder, Beviser af Virkninger, det hele Opløb af Assistentshuusforvaltere og Vederhæftige bort, for at faae det Absurde tydeligt – at man da kan troe, om man vil det – jeg siger kun, at det maa være den yderste Anstrængelse.

[…]

Heri ligger alt Hedenskab, at Gud forholder sig ligefremt til Mennesket som det Paafaldende til den Forundrede. Men Aands-Forholdet i Sandhed til Gud det vil sige Inderligheden er netop først betinget ved Inderliggjørelsens Gjennembrud, der svarer til den guddommelige Underfundighed, at Gud intet, intet Paafaldende har, ja i den Grad er langtfra at være paafaldende, at han er usynlig, saa man slet ikke falder paa, at han er til, medens hans Usynlighed igjen er hans Allestedsnærværelse. Men en Allestedsnærværende er jo En man seer overalt, som for Exempel en Politiebetjent: hvor svigefuldt da, at en Allestedsnærværende netop derpaa er kjendelig, at han er usynlig, ene og alene derpaa, thi hans Synlighed er netop at ophæve Allestedsnærværelsen. Dette Forhold mellem Allestedsnærværelse og Usynlighed er ligesom Forholdet mellem Hemmelighed og Aabenbaring, at Hemmeligheden er Udtrykket for, at Aabenbaringen er Aabenbaring i strængere Forstand, at Hemmeligheden netop er det Eneste, den kan kjendes paa, da ellers en Aabenbaring bliver saadan Noget som en Politiebetjents Allestedsnærværelse. – Vil Gud aabenbare sig i menneskelig Skikkelse og give et ligefremt Forhold ved for Exempel at iføre sig Skikkelse af en Mand der er sex Alen høi, saa vil hiin experimenterede Selskabsmand og Fuglekonge nok blive opmærksom. Men Aands-Forholdet i Sandhed, naar Gud ikke vil bedrage, fordrer netop, at Skikkelsen slet intet Paafaldende har, saa Selskabsmanden maa sige: der er ikke det allermindste at see. Naar Guden slet intet Paafaldende har, saa bedrages maaskee Selskabsmanden derved, at han slet ikke bliver opmærksom. Men deri er Guden uden Skyld, og det Bedrags Virkelighed er tillige bestandigt Muligheden af Sandheden. Men har Guden det Paafaldende, saa bedrager han derved, at Mennesket bliver opmærksom paa Usandheden, og denne Opmærksomhed er tillige Sandhedens Umulighed. – I Hedenskabet er det ligefremme Forhold Afguderiet, i Christendommen veed jo Enhver, at Gud ikke saaledes kan vise sig. Men denne Viden er ingenlunde Inderligheden, og i Christendommen kan det vel hænde en Udenadvidende, at han bliver aldeles »uden Gud i Verden«, som det ikke var Tilfældet i Hedenskabet, thi der var dog Afgudsdyrkelsens usande Forhold. Og vel er Afguderie et sørgeligt Surrogat, men at Artiklen Gud aldeles gaaer ud, er dog endnu galere.

Altsaa end ikke Gud forholder sig ligefremt til den deriverede Aand (og dette er Skabelsens Vidunderlighed, ikke at frembringe Noget som er Intet ligeoverfor Skaberen, men at frembringe Noget som er Noget, og som i den sande Gudsdyrkelse kan benytte dette Noget til ved sig selv at blive Intet for Gud), mindre da, at det ene Menneske kan forholde sig saaledes til det Andet i Sandhed. Naturen, Skabningens Totalitet, er Guds Gjerning, og dog er Gud ikke der, men inden i det enkelte Menneske er der en Mulighed (han er efter sin Mulighed Aand), som i Inderligheden vækkes til et Guds-Forhold, og saa er det muligt at see Gud overalt. Den sandselige Distinction mellem det Store, det Forbausende, en sydlig Nations meest himmelraabende Superlativ er en Tilbagegang til Afguderie i Sammenligning med Inderlighedens Aands-Forhold. Er dette ikke ligesom hvis en Forfatter skrev 166 Folio-Bind, og Læseren læste og læste, ligesom man seer og seer paa Naturen, men opdagede ikke, at hiint uhyre Værks Mening laae i Læseren selv; thi Forbauselse over de mange Bind og over at der stod 500 Linier paa Siden, hvilken er ligesom Forbauselsen over hvor stor Naturen er og hvor utallige Dyreslægter der er, er jo ikke Forstaaelsen.

Mellem Aand og Aand er et ligefremt Forhold utænkeligt i Forhold til den væsentlige Sandhed; antages Forholdet, betyder det egentligen, at den ene Part er ophørt at være Aand, hvilket mangt et Genie ikke betænker, der baade hjælper Folk en masse ind i Sandheden, og er godmodig nok til at mene, at Acclamation, Villighed til at høre, Navns-Underskrift og saa videre er at have antaget Sandheden. Akkurat ligesaa vigtig som Sandheden, og Eet af To endnu vigtigere, er Maaden, paa hvilken Sandheden antages, og det hjælper kun lidet om En fik Millioner til at antage Sandheden, naar de netop ved Antagelsens Maade bleve satte over i Usandheden. Og derfor er al Godmodighed, al Overtalelse, al Prutten, al ligefrem Tiltrækning ved Hjælp af sin egen Person, i Betragtning af at man lider saa meget for den Sag, at man græder over Menneskene, at man er saa begeistret og saa videre, alt Sligt er Misforstaaelse, er et Falsum i Forhold til Sandheden, hvorved man, i Forhold til som man har Evne, hjælper en Slump Mennesker til at faae Sandhedens Skin.

See, Socrates var en Lærer i det Ethiske, men han var opmærksom paa, at der intet ligefremt Forhold er mellem Læreren og den Lærende, fordi Inderligheden er Sandheden, og Inderligheden i de Tvende netop Veien bort fra hinanden. Fordi han indsaae dette, derfor var det formodentlig, at han var saa glad over sit fordeelagtige Ydre. Hvilket var dette? Ja gjæt engang! I vor Tid sige vi jo om en Præst, at han har ret et fordeelagtigt Ydre, glæde os derved og forstaae, at han er en smuk Mand, at Præstekjolen klæder ham godt, at han har et klangfuldt Organ og en Skabning, som enhver Skrædder, eller hvad det var jeg vilde sige, enhver Tilhører maa have Glæde af. Ak, ja, naar man er saaledes udrustet af Naturen og saaledes paaklædt af Skrædderen, saa kan man sagtens være Religionslærer, og være det med Held; thi Religionslærernes Vilkaar er høist forskjelligt, ja mere end man tænker paa, naar der høres Besværing over, at nogle Præstekald ere saadanne store Levebrød, andre meget smaa; Forskjellen er endnu større, at nogle Religionslærere blive korsfæstede – men Reglionen er ganske den samme. Og den indeholdte Læres reduplicerede Gjentagelse i Forestillingen om, hvorledes Læreren skal være, bryder man sig ikke stort om. Man foredrager Orthodoxie og pynter Læreren ud i hedensk-æsthetiske Bestemmelser. Man fremstiller Christus i Bibelens Udtryk; at han bar al Verdens Synd vil ikke ret bevæge Menigheden, dog forkynder Taleren det, og for ret at gjøre Modsætningen stærk, skildrer han Christi Skjønhed (thi Modsætningen mellem Uskyld og Synd er ikke stærk nok), og den troende Menighed røres ved denne aldeles hedenske Bestemmelse af Guden i menneskelig Skikkelse: Skjønhed. – Dog tilbage til Socrates. Han havde ikke et saadant fordeelagtigt Ydre, som det skildrede; han var meget styg, havde klodderagtige Fødder og fremfor Alt en Mængde Organer paa Panden og andre Steder, der maatte forvisse ethvert Menneske om, at han var et demoraliseret Subjekt. See, dette var hvad Socrates forstod ved sit fordeelagtige Ydre, som han var saa kisteglad over, at han vilde have anseet det for en Chicane af Guden for at forhindre ham i at være Sædelærer om han havde givet ham et behageligt Ydre som en kjelen Citharspiller, et smægtende Blik som en Schäfer, saa smaae Fødder som en Bal-Direkteur i det Venskabelige, og in toto et saa fordeelagtigt Ydre som nogen Tjenestesøgende i Adresseavisen eller en theologisk Candidat, der har sat sit Haab til et privat Kald, kan ønske sig det. Hvorfor mon nu hiin gamle Lærer var saa glad over sit fordeelagtige Ydre, uden fordi han indsaae, at det maatte hjælpe til at fjerne den Lærende, saa denne ikke blev hængende i et ligefremt Forhold til Læreren, maaskee beundrede ham, maaskee lod sine Klæder sye paa samme Maade, men maatte fatte ved Modsætningens Frastød, hvilket igjen i en høiere Sphære hans Ironie var, at den Lærende har væsentlig med sig selv at gjøre, og at Sandhedens Inderlighed ikke er den kammeratlige Inderlighed, hvormed to Busenfreunde gaae med hinanden under Armen, men den Adskillelse, hvori hver for sig selv er existerende i det Sande.

\sectionheading{Mulighed høiere end Virkelighed; Virkelighed høiere end Mulighed (uddrag)  ·  SV¹ VII, 273–297}

Aristoteles bemærker i sin Poetik, at Poesien er høiere end Historien, fordi Historien kun fremstiller hvad der er skeet, Poesien hvad der kunde og burde være skeet, det vil sige Poesien raader over Muligheden. I Forhold til Virkelighed er, poetisk og intellectuelt, Muligheden høiere, det Æsthetiske og det Intellectuelle interesseløst. Men der er kun een Interesse, den at existere; Interesseløsheden er Udtrykket for Ligegyldigheden mod Virkelighed. Den Ligegyldighed er glemt i det cartesianske cogito – ergo sum, hvilket uroliger det Intellectuelles Interesseløshed, og fornærmer Speculationen, som skulde der af den følge noget Andet. Jeg tænker, ergo tænker jeg, om jeg er eller det er (i Virkelighedens Forstand, hvor jeg betyder et enkelt existerende Menneske, og det betyder et bestemt enkelt Noget), er uendelig ligegyldigt. At det jeg tænker er i Tænkningens Forstand, behøver jo intet Beviis, eller at bevises ved nogen Slutning, da det jo er beviist. Saasnart jeg begynder at ville gjøre min Tænken teleologisk i Forhold til noget Andet, er Interessen med i Spillet. Saasnart den er der, saa er det Ethiske med tilstede og fritager mig for yderligere Uleilighed med at bevise min Tilværelse, forhindrer mig i ethisk svigefuldt og metaphysisk uklart at gjøre Slutningens Sving, idet det forpligter mig til at existere.

Medens det Ethiske i vor Tid mere og mere ignoreres, har denne Ignoreren tillige havt den skadelige Følge, at det har forvirret Poesien og Speculationen, der have sluppet Mulighedens interesseløse Ophøiethed for at gribe efter Virkeligheden: istedetfor at give hver Sit har man forvirret paa dobbelt Maade. Poesien gjør det ene Forsøg efter det andet at virke som Virkelighed, hvilket er aldeles upoetisk; Speculationen vil atter og atter indenfor sit Omfang naae Virkeligheden, forsikkrer at det Tænkte er det Virkelige, at Tænkningen ikke blot er istand til at tænke, men til at give Virkelighed, hvilket er lige omvendt; og samtidigen glemmes det mere og mere hvad det er at existere. Tiden og Menneskene blive mere og mere uvirkelige, deraf disse Surrogater, der skal erstatte det Tabte. Det Ethiske opgiver man mere og mere, den Enkeltes Liv bliver ikke blot poetisk, men verdenshistorisk uroliget og derved hindret i at existere ethisk; saa maa Virkelighed paa andre Maader skaffes tilveie. Men denne misforstaaede Virkelighed er ligesom, hvis en Generation eller Individerne i den vare blevne for tidligt gamle og der nu paa kunstig Maade maatte skaffes Ungdommelighed tilveie. Istedetfor at det at existere ethisk er Virkeligheden, er Tiden bleven saa overveiende betragtende, at ikke blot Alle ere det, men at tilsidst dette bliver forfalsket som var det Virkelighed. Man smiler ad Kloster-Livet, og dog levede ingen Eremit saa uvirkeligt som man nuomstunder lever, thi en Eremit abstraherede vel fra den hele Verden, men han abstraherede ikke fra sig selv; man veed at beskrive et Klosters phantastiske Beliggenhed i det Afsides, i Skovens Eensomhed, i Horizontens fjerne Blaanen, men den rene Tænkens phantastiske Beliggenhed tænker man ikke paa. Og dog er Eneboerens pathetiske Uvirkelighed langt at foretrække for den rene Tænkers comiske Uvirkelighed; og dog er Eneboerens lidenskabelige Glemsel, der tager ham hele Verden bort, langt at foretrække for den verdenshistoriske Tænkers comiske Distraction, der glemmer sig selv.

Ethisk seet er Virkelighed høiere end Mulighed. Det Ethiske vil netop tilintetgjøre Mulighedens Interesseløshed ved at gjøre det at existere til den uendelige Interesse. Det Ethiske vil derfor forhindre ethvert Confusions-Forsøg, som for Exempel ethisk at ville betragte Verden og Menneskene. Ethisk at betragte lader sig nemlig ikke gjøre, thi der er kun een ethisk Betragtning, det er Selvbetragtning. Det Ethiske slutter sig øieblikkeligen om den Enkelte med Fordring til ham, at han skal existere ethisk; det sqvadronnerer ikke om Millioner og Generationer, det tager ikke Menneskeheden paa Slump, ligesaa lidet som Politiet arresterer den rene Menneskehed. Det Ethiske har med de enkelte Mennesker at gjøre, og vel at mærke med hver Enkelt. Veed Gud hvor mange Haar der er paa et Menneskes Hoved, saa veed det Ethiske hvor mange Mennesker der er til, og den ethiske Folketælling er ikke i Interesse af en Totalsum, men i Interesse af hver Enkelt. Det Ethiske fordrer sig selv af ethvert Menneske, og naar det dømmer, da dømmer det igjen hver Enkelt, kun en Tyran og et afmægtigt Menneske nøies med at decimere. Det Ethiske griber den Enkelte og fordrer af ham, at han skal entholde sig fra al Betragten, især af Verden og Menneskene; thi det Ethiske som det Indvortes lader sig slet ikke betragte af Nogen der staaer udenfor, det lader sig kun realisere af det enkelte Subjekt, der da kan være vidende om hvad der boer i ham, den eneste Virkelighed der ikke bliver til en Mulighed ved at vides og ikke kun kan vides ved at tænkes, da det er hans egen Virkelighed, hvilken han som tænkt Virkelighed det vil sige som Mulighed vidste førend den blev Virkelighed, medens han i Forhold til en Andens Virkelighed Intet vidste om den førend han ved at faae den at vide tænkte den, det vil sige forvandlede den til Mulighed.

I Forhold til enhver Virkelighed uden for mig gjelder det at jeg kun tænkende kan faae fat paa den. Skulde jeg virkeligen faae fat paa den, maatte jeg kunne gjøre mig til den Anden, til den Handlende, gjøre den mig fremmede Virkelighed til min egen Virkelighed, hvilket er en Umulighed. Gjør jeg nemlig den mig fremmede Virkelighed til min egen, saa betyder det ikke, at jeg ved at være vidende om den bliver han, men det betyder en ny Virkelighed, der tilhører mig som forskjellig fra ham.

[…]

Sandhedens Hvorledes er netop Sandheden. Det er derfor Usandhed at besvare et Spørgsmaal i et Medium, hvor Spørgsmaalet ikke kan komme frem. Saaledes at forklare Virkeligheden indenfor Muligheden, indenfor Muligheden at distingvere mellem Mulighed og Virkelighed. Ved at man da ikke spørger æsthetisk og intellectuelt, men kun ethisk, og ethisk igjen i Retning af sin egen Virkelighed, om Virkelighed, er ethisk ethvert Individ udsondret for sig. Ironie og Hyklerie som Modsætningerne, men som begge udtrykkende den Modsigelse, at det Udvortes ikke er det Indvortes (Hyklerie ved at synes god, Ironie ved at synes slet), indskærpe betræffende det betragtende Spørgsmaal om det ethiske Indvortes, at Virkelighed og Bedrag ere lige mulige, at Bedraget kan naae lige saa vidt som Virkeligheden. Kun Individet selv kan vide, hvilket der er hvilket. At spørge om dette ethiske Indvortes i et andet Individ er allerede uethisk, forsaavidt det er en Adspredelse. Men spørges der alligevel, saa er Vanskeligheden der, at jeg kun ved at tænke den Andens Virkelighed faaer fat i den, altsaa ved at oversætte den i Muligheden, hvor da Muligheden af Bedraget lige saa godt er at tænke. – Dette er et gavnligt Forstudium for at existere ethisk: at lære, at det enkelte Menneske staaer alene.

Æsthetisk og intellectuelt at spørge om Virkelighed, er en Misforstaaelse; ethisk at spørge om et andet Menneskes Virkelighed er en Misforstaaelse, da der kun skal spørges om den egne. Troens (sensu strictissimo, der refererer sig til noget Historisk) Forskjellighed fra det Æsthetiske, Intellectuelle, det Ethiske, viser sig her. Uendeligt interesseret at spørge om en Virkelighed, der ikke er Ens egen, er at ville troe, og udtrykker det paradoxe Forhold til det Paradoxe. Æsthetisk lader der sig ikke spørge saaledes uden i Tankeløshed, da æsthetisk Mulighed er høiere end Virkelighed; intellectuelt ikke, da intellectuelt Mulighed er høiere end Virkelighed, ethisk heller ikke, fordi Individet ethisk ene og alene er uendeligt interesseret i sin egen Virkelighed. – Troens Analogie med det Ethiske er den uendelige Interesserethed, hvorved den Troende absolut er forskjellig fra en Æsthetiker og en Tænker, men atter forskjellig fra en Ethiker derved, at han er uendeligt interesseret i en Andens Virkelighed (for Exempel at Guden virkelig har været til).

Æsthetisk og intellectuelt gjelder det, at kun da er en Virkelighed forstaaet og tænkt, naar dens esse er opløst i dens posse. Ethisk gjelder det, at kun da er Muligheden forstaaet, naar hvert et posse virkeligen er et esse. Naar det Æsthetiske og Intellectuelle seer efter, protesterer det ethvert esse, som ikke er et posse; naar det Ethiske seer efter, fordømmer det ethvert posse, som ikke er et esse, et posse nemlig i Individet selv, da det ethisk ikke har med andre Individer at gjøre. – I vor Tid blandes Alt sammen, man besvarer det Æsthetiske ethisk, Troen intellectuelt og saa videre Man er færdig med Alt, og dog er man langtfra agtsom paa i hvilken Sphære ethvert Spørgsmaal finder sit Svar. I Aandens Verden frembringer dette en endnu større Confusion end naar i den borgerlige Verden for Exempel et geistligt Anliggende blev besvaret for Exempel af Brolægningscommissionen.

Er Virkeligheden da Udvortesheden? Ingenlunde. Æsthetisk og intellectuelt indskærpes det ganske rigtigt, at det Udvortes kun er Bedraget for Den, som ikke fatter Idealiteten. Frater Taciturnus siger (loco citato pagina 341): »Viden af det Historiske hjælper En blot ind i et Sandsebedrag, der bedaares af det Stofartige. Hvad er det jeg veed historisk? Det Stofartige. Idealiteten veed jeg ved mig selv, og veed jeg den ikke ved mig selv, veed jeg den slet ikke, al historisk Viden hjælper ikke. Idealitet er ikke Løsøre, der kan transporteres fra Een til en Anden, eller Noget som gaaer med i Kjøbet, naar man tager store Partier. Veed jeg at Cæsar var stor, saa veed jeg hvad det Store er, og dette seer jeg paa, ellers veed jeg ikke, at Cæsar var stor. Historiens Fortælling, at vederhæftige Mænd forsikkre det, at der ingen Risico er forbunden med at antage denne Mening, da det skal være vist, at han var en stor Mand, at Udfaldet beviser det, hjælper slet ikke. At troe Idealiteten paa en Andens Ord er ligesom at lee af en Vittighed, ikke fordi man har forstaaet den, men fordi en Anden har sagt, at det var vittigt. Naar saa er, saa kan i Grunden Vittigheden for den i Kraft af Troen og Agtelsen Leende lige saa godt være usagt, han kan lee med samme Emphasis.« – Hvad er da Virkeligheden? Den er Idealiteten. Men æsthetisk og intellectuelt er Idealiteten Muligheden (Tilbageførelsen ab esse ad posse). Ethisk er Idealiteten Virkeligheden i Individet selv. Virkeligheden er Indvortesheden uendeligt interesseret i at existere, hvilket det ethiske Individ er for sig selv.

Naar jeg forstaaer en Tænker, saa er netop i samme Grad som jeg forstaaer ham, hans Virkelighed (at han existerer som et enkelt Menneske; at han selv virkelig har forstaaet det saaledes og saa videre; eller at han selv virkelig har realiseret det og saa videre) aldeles ligegyldig. Deri har Philosophien og Æsthetiken Ret, og det gjelder netop om, at dette ret fastholdes. Men deri ligger endnu intet Forsvar for den rene Tænken som Meddelelses-Medium. Fordi nemlig hans Virkelighed er mig den Lærende ligegyldig, ligesom omvendt min ham, deraf følger ingenlunde, at han selv tør være ligegyldig ved sin egen Virkelighed. Dette maa hans Meddelelse bære Præg af, vel ikke ligefremt, thi det lader sig mellem Mand og Mand ikke meddele ligefrem (da et saadant Forhold er den Troendes paradoxe Forhold til Troens Gjenstand), og lader sig ikke forstaae ligefremt, men maa indirecte være til at forstaae indirecte.

Naar de enkelte Sphærer ikke holdes afgjørende ude fra hinanden, forvirres Alt. Naar man saaledes er nysgjerrig i Forhold til en Tænkers Virkelighed, finder det interessant, at man veed Noget derom og saa videre, saa er man intellectuelt at dadle, fordi det i Intellectualitetens Sphære netop er Maximum, at Tænkerens Virkelighed er aldeles ligegyldig. Men ved saaledes at være vrøvlevorn i Intellectualitetens Sphære faaer man en forvirrende Liighed med en Troende. En Troende er netop uendeligt interesseret i en Andens Virkelighed. Dette er for Troen det Afgjørende, og denne Interesserethed ikke saadan lidt Nysgjerrighed, men den absolute Afhængighed at Troens Gjenstand.

Troens Gjenstand er en Andens Virkelighed; dens Forhold en uendelig Interesserethed. Troens Gjenstand er ikke en Lære, thi saa er Forholdet intellectuelt, og det gjelder om ikke at fuske, men at naae det intellectuelle Forholds Maximum. Troens Gjenstand er ikke en Lærer, der har en Lære, thi naar en Lærer har en Lære, er eo ipso Læren vigtigere end Læreren, og Forholdet intellectuelt, hvor det gjelder om ikke at fuske, men at naae det intellectuelle Forholds Maximum. Men Troens Gjenstand er Lærerens Virkelighed, det at Læreren virkelig er til. Troens Svar er derfor absolut enten Ja eller Nei. Thi Troens Svar er ikke i Forhold til en Lære, om den er sand eller ikke, ikke i Forhold til en Lærer, om hans Lære er sand eller ikke, men er Svaret paa Spørgsmaalet om et Faktum: antager Du, at han virkeligen har været til? Og vel at mærke, Svaret er med uendelig Lidenskab. I Forhold nemlig til et Menneske er det tankeløst at lægge saa uendelig megen Vægt paa, om han har været til eller ikke. Dersom derfor Troens Gjenstand er et Menneske, saa er det hele en Narrestreg af et taabeligt Menneske, der end ikke har fattet det Æsthetiske og Intellectuelle. Troens Gjenstand er derfor Gudens Virkelighed i Betydning af Existents. Men at existere betyder først og fremmest at være en Enkelt, og derfor er det, at Tænkningen maa see bort fra Existents, fordi det Enkelte ikke lader sig tænke, men kun det Almene. Troens Gjenstand er da Gudens Virkelighed i Existents det vil sige som en Enkelt det vil sige at Guden har været til som et enkelt Menneske.

Christendommen er ingen Lære om Eenheden af det Guddommelige og Menneskelige, om Subjekt-Objektet, for ikke at nævne de øvrige logiske Omskrivninger af Christendommen. Dersom nemlig Christendommen var en Lære, saa var Forholdet til den ikke Troens, thi til en Lære gives der kun et intellectuelt Forhold. Christendommen er derfor ingen Lære, men det Faktum, at Guden har været til.

Tro er saaledes ikke en Sinke-Lectie i Intellectualitetens Sphære, et Asyl for de daarlige Hoveder. Men Tro er en Sphære for sig selv, og enhver Misforstaaelse af Christendommen strax kjendelig paa, at den forvandler den til en Lære, og drager den ind i Intellectualitetens Omfang. Hvad der i Intellectualitetens Sphære gjelder som Maximum, at blive aldeles ligegyldig mod Lærerens Virkelighed, det gjelder omvendt i Troens Sphære, dens Maximum er den quam maxime uendelige Interesserethed i Lærerens Virkelighed.

[…]

Gud tænker ikke, han skaber; Gud existerer ikke, han er evig. Mennesket tænker og existerer, og Existents adskiller Tænken og Væren, holder dem ude fra hinanden i Succession.

Hvad er abstrakt Tænkning? Det er den Tænken, hvor der ingen Tænkende er. Den seer bort fra Alt Andet end Tanken, og kun Tanken er i sit eget Medium. Existentsen er ikke tankeløs, men i Existents er Tanken i et fremmed Medium. Hvad vil saa det sige, i den abstrakte Tænknings Sprog at spørge om Virkelighed i Betydning af Existents, da Abstraktionen netop seer bort derfra? – Hvad er concret Tænkning? Det er den Tænken, hvor der er en Tænkende, og et bestemt Noget (i Betydning af Enkelt) der tænkes, hvor Existentsen giver den existerende Tænker Tanken, Tid og Rum.

[…]

Hvad vil det sige at Væren er høiere end Tænken? Er dette Udsagn Noget, der skal tænkes, saa er jo eo ipso Tænken igjen høiere end Væren. Lader det sig tænke, saa er Tænkningen høiere; lader det sig ikke tænke, saa er intet Tilværelsens System muligt. Det hjælper slet ikke hverken at være høflig eller grov mod Væren, hverken at lade den være noget Høiere, der dog følger af Tænkningen og syllogistisk naaes, eller noget saa Ringe, at det uden videre følger med Tænkningen. Naar man saaledes har sagt: Gud maa have alle Fuldkommenheder, eller det høieste Væsen maa have alle Fuldkommenheder; at være er ogsaa en Fuldkommenhed, ergo maa det høieste Væsen være, eller Gud maa være: saa er den hele Bevægelse svigefuld. Dersom nemlig i den første Deel af denne Tale Gud virkelig ikke tænkes værende, saa kan Talen slet ikke komme istand. Den vil da lyde saaledes: et høieste Væsen, som vel at mærke ikke er til, maa være i Besiddelse af alle Fuldkommenheder, deriblandt ogsaa af den at være til, ergo er et høieste Væsen, som ikke er til, det er til. Dette vilde være en besynderlig Slutning. Det høieste Væsen maa enten ikke være til i Begyndelsen af Talen for at blive til i Slutningen, og saa kan det ikke blive til, eller det var til, og saa kan det jo ikke blive til, saa er Slutningen en svigefuld Form for Prædikats-Udvikling, en svigefuld Omskrivelse af en Forudsætning. I andet Fald maa Slutningen holdes reent hypothetisk: dersom et høieste Væsen antages at være, maa det ogsaa antages at være i Besiddelse af alle Fuldkommenheder; at være er en Fuldkommenhed, ergo maa det være – dersom det nemlig antages at være. Ved at slutte indenfor en Hypothese kan man dog vel aldrig komme til at slutte ud af Hypothesen. Som for Exempel dersom det og det Menneske er en Hykler, vil han bære sig ad som en Hykler; en Hykler vil gjøre Det og Det, ergo har det og det Menneske gjort Det og Det. Saaledes ogsaa med Slutningen om Gud. Naar Slutningen er færdig, er Guds Væren netop ligesaa hypothetisk som den var, men indenfor er der avanceret et Slutnings-Forhold mellem et høieste Væsen og at være som Fuldkommenhed, ligesom i det andet Tilfælde mellem det at være Hykler og en enkelt Yttring deraf. Confusionen er den samme som at forklare Virkelighed i den rene Tænken; §'en er overskreven Virkelighed, man har forklaret Virkelighed, men glemt at det Hele er indenfor den rene Tænkens Mulighed. Dersom en Mand begyndte en Paranthes, men denne blev saa lang at han selv glemte det, det hjælper dog ikke: saasnart man læser det op, bliver det meningsløst, uden videre at lade Indskuds-Sætningen forvandle sig til Hovedsætningen.

Naar Tænkningen vender sig mod sig selv for at tænke over sig selv, fremkommer der en Skepsis, som bekjendt. Hvorledes standses denne Skepsis, der har sin Grund i, at Tænkningen istedenfor at være tjenende ved at tænke Noget, selvisk vil tænke sig selv? Naar en Hest løber løbsk og løber durch, saa lod det sig, fraseet hvad Skade der midlertidig kunde skee, vel høre om En vilde sige: Lad den kun løbe, den bliver nok træt. Med Hensyn til Tænkningens Selv-Reflexion lader dette sig ikke sige, thi den kan blive ved saalænge det skal være, og løber rundt. Schelling standsede Selv-Reflexionen, forstod den intellectuelle Anskuelse ikke som en Opdagelse indenfor Selv-Reflexionen, som naaedes ved at fare fort, men som et nyt Udgangspunkt. Hegel anseer dette for en Feil og taler absprechend nok om den intellectuelle Anskuelse – saa kom da Methoden. Selv-Reflexionen bliver saa længe ved til den hæver sig selv, Tænkningen trænger seierrig igjennem og faaer atter Realitet, Tænken og Værens Identitet er vundet i den rene Tænken. Hvad vil det sige, at Selv-Reflexionen bliver ved saa længe indtil den hæver sig selv? For at man skal opdage Selv-Reflexionens Mislighed, behøver den ikke at blive længe ved, men paa den anden Side, saa længe den bliver ved, er det aldeles den samme Mislighed. Hvad skal saa det sige: saa længe indtil? Det er ikke Andet end en bestikkende Tale, der vil ved Qvantiteren bestikke Læserens Forestilling, som var det bedre at forstaae, at Selv-Reflexionen hævede sig selv, naar det varede længe inden det skeete. Denne Qvantiteren er et Sidestykke til Astronomernes uendelig smaae Vinkler, der tilsidst blive saa smaae (Vinkler), at man kan kalde dem parallelle Linier. Fortællingen om, at Selv-Reflexionen bliver »saa længe ved indtil« bortleder Opmærksomheden fra hvad der dialektisk er Hovedsagen: hvorledes Selv-Reflexionen hæves. Naar man siger om En, han vedblev saa længe at sige en Usandhed for Spøg, indtil han selv troede det var Sandhed: saa ligger den ethiske Accent paa Overgangen, men det Formildende, det Adspredende er dette saa længe; man glemmer næsten Overgangens Afgjørelse, fordi det varer saa længe. I Fortællingen, i det Beskrivende, i det rhetoriske Foredrag frembringer det abstrakte »saa længe indtil« en stor illusorisk Virkning, denne være nu som et optisk Bedrag (for Exempel Judiths Bog c. 10, v. 11: »Og Judith gik ud, hun og hendes Pige med hende, men de Mænd af Staden saae efter hende, indtil hun kom ned af Bjerget, indtil hun kom igjennem Dalen, og de ikke kunde see hende mere;« Pigen sad ved Stranden og saae efter den Elskede – indtil hun ikke saae ham mere), eller som Tidens phantastiske Forsvinden, fordi der ingen Maalestok er og Intet at maale med i det Abstrakte »saa længe indtil.« (Da seirede Lysten og han foer vild fra Sandhedens Vei – indtil Angerens Bitterhed standsede ham; – der skal Mesterskab i psychologisk Tegning til for at frembringe ved Concretion en saa stor Virkning som dette abstrakte Indtil, der lokker Phantasien.) Men dialektisk er denne phantastiske Længde aldeles af ingen Betydning. Da en græsk Philosoph blev adspurgt, hvad Religion var, bad han om Udsætningstid; da Terminen kom, forlangte han igjen denne udsat og saa fremdeles; han vilde dermed antyde, at Spørgsmaalet ikke lod sig besvare. Dette var græsk og skjønt og sindrigt. Dersom han derimod i Betragtning af, at det havde varet saa længe, vilde have meent, derved i fjerneste Maade at være kommen Besvarelsen nærmere: saa var jo dette en Misforstaaelse, ligesom naar en Debitor forbliver saa længe i Gjelden indtil den er betalt – derved at det har varet saa længe, at den ikke er blevet betalt. Det Abstrakte »saa længe – indtil« har noget besynderligt Bestikkende ved sig. Vilde En sige: Selv-Reflexionen hæver sig selv, og nu søge at vise hvorledes, saa vilde neppe Nogen forstaae det, men naar man siger: Selv-Reflexionen bliver ved saa længe indtil den hæver sig selv, saa tænker man maaskee: ja det er en anden Sag, det er der Noget i; man bliver angest og bange for denne Længde, man taber Taalmodigheden, man tænker: lad gaae – og saa begynder den rene Tænken. – Forsaavidt kan den rene Tænken have Ret i, at den ikke begynder bittweise som de ældre maadelige Philosopher; thi Læseren takker Gud til, at den begynder, af Frygt for den forfærdelige Længde – indtil.

[…]

Dersom det Tænkte var Virkelighed, saa vilde altsaa det saa fuldkomment som muligt Udtænkte, naar jeg endnu ikke havde handlet, det vilde være Handlingen. Paa den Maade blev der slet ingen Handling, men det Intellectuelle sluger det Ethiske. At jeg nu skulde mene, at det Udvortes er det der gjør Handlingen til Handling, er en Taabelighed; og paa den anden Side at ville vise, hvor ethisk Intellectualiteten er, at den endog gjør Tanken til Handling, er en Sophisme, der forskylder en Dobbelthed i Brugen af det Ord: at tænke. Skal der overhovedet være en Forskjel mellem Tænken og Handlen, da kan denne kun fastholdes ved at anvise Tænkningen Muligheden, Interesseløsheden, Objektiviteten – Handlingen Subjektiviteten. Men nu viser der sig let et Confinium. Naar jeg saaledes tænker at jeg vil gjøre det og det, saa er denne Tænken vel endnu ikke en Handling, og er i al Evighed qvalitativt forskjellig derfra, men dog er den en Mulighed, i hvilken Virkelighedens og Handlingens Interesse allerede reflekterer sig. Derfor er ogsaa Interesseløsheden, Objektiviteten, ifærd med at forstyrres, fordi Virkeligheden og Ansvaret vil have fat derpaa. (Der gives saaledes en Synd i Tanke). – Virkeligheden er ikke den udvortes Handling, men en Indvorteshed, i hvilken Individet ophæver Muligheden og identificerer sig med det Tænkte for at existere deri. Dette er Handling. Intellectualiteten synes saa rigoristisk ved at gjøre Tanken selv til Handling, men denne Rigorisme er blind Allarm, fordi det, at Intellectualiteten faaer Lov at hæve Handlingen overhovedet, er en Slappelse. Saa gjelder det som i tidligere paaviste Analogier: indenfor en total Slappelse at være regoristisk er kun Illusion og væsentligen kun en Slappelse. Dersom En for Exempel vilde kalde Synd Uvidenhed, og nu indenfor denne Bestemmelse med Rigorisme opfatte de enkelte Synder, saa er dette aldeles illusorisk, thi sagt indenfor den totale Bestemmelse, at Synd er Uvidenhed, bliver enhver Bestemmelse væsentlig letsindig, fordi den totale Bestemmelse er Letsindighed. – Med Hensyn til det Onde skuffer Forvexlingen af Tænken og Handlen lettere; men seer man nøiere, viser det sig, at Grunden er den, det Godes Nidkjærhed paa sig selv, der i den Grad fordrer sig af Individet, at den bestemmer en Tænken af det Onde som Synd. Men lad os tage det Gode. At have tænkt noget Godt, man vil gjøre, er det at have gjort det? Ingenlunde, men det er heller ikke det Udvortes, der gjør Udslaget; thi En der ikke eier en Hvid kan være lige saa barmhjertig som Den, der bortgiver et Kongerige. Da Leviten reiste den Ulykkelige forbi, der var bleven overfalden af Røvere paa Veien fra Jericho til Jerusalem, faldt det ham maaskee ind, da han endnu var lidt borte fra den Ulykkelige, at det dog var skjønt at hjælpe en Lidende; han tænkte maaskee endog allerede, hvilken Løn dog en saadan god Gjerning har i sig, han reed maaskee langsommere fordi han faldt i Tanker; men alt som han kom nærmere og nærmere, da viste Vanskelighederne sig, og han reed forbi. Nu reed han vel rask til, for hurtigt at komme bort, bort fra Tanken om Veiens Usikkerhed, bort fra Tanken om Røvernes mulige Nærhed, og bort fra Tanken om hvor let den Ulykkelige kunde komme til at forvexle ham med Røverne, der lod ham ligge. Han handlede altsaa ikke. Men sæt Fortrydelsen hentede ham ind paa Veien, sæt han iilsomt vendte om, uden at frygte hverken Røvere eller andre Vanskeligheder, blot frygtende for at komme for sildig; sæt han kom for sildig, da den barmhjertige Samaritan allerede havde faaet den Lidende bragt i Herberge: havde han da ikke handlet? Visseligen, og dog kom han ikke til at handle i det Udvortes. – Lad os tage en religieus Handling. At troe Gud, er det at tænke over hvor herligt det maa være at troe, at tænke over hvilken Fred og Tryghed Troen kan skjenke? Ingenlunde. Selv at ønske, hvor dog Interessen, Subjektets Interesse er langt tydeligere, er ikke at troe, ikke at handle. Individets Forhold til den tænkte Handling er dog endnu bestandigt kun en Mulighed, han kan gaae fra. – At der gives Tilfælde i Forhold til det Onde, hvor Overgangen næsten ikke er til at mærke, nægtes ikke, men disse Tilfælde maa forklares paa en egen Maade. Det ligger i, at Individet er i en Vanes Magt, saa at han ved ofte at have dannet Overgangen fra Tænken til Handlen tilsidst har tabt Magten dertil i en Vanes Trældom, der for hans Regning gjør den hurtigere og hurtigere.

Mellem den tænkte Handling og den virkelige, mellem Mulighed og Virkelighed er der maaskee i Indhold aldeles ingen Forskjel, i Form er den altid væsentlig. Virkeligheden er Interesseretheden ved at existere deri.

At Handlingens Virkelighed saa ofte forvexles med allehaande Forestillinger, Forsætter, Tilløb til Beslutninger, Stemnings Forspil og saa videre, at der overhovedet meget sjeldent virkeligt handles, nægtes ikke, tvertimod antages, at dette meget har bidraget til Confusionen. Men man tage en Handling sensu eminenti, da viser Alt sig tydeligt. Det Udvortes ved Luthers Handling er at træde frem paa Rigsdagen i Worms, men fra det Øieblik han med hele sin Subjektivitets lidenskabelige Afgjørelse existerede i at ville, da ethvert Muligheds Forhold til denne Handling af ham maatte betragtes som Anfægtelse: da havde han handlet. Da Dion gik ombord for at styrte Tyrannen Dionysius, skal han have sagt, at selv om han døde underveis havde han dog udført en herlig Bedrift – han havde altsaa handlet. At Afgjørelsen i det Udvortes skulde være høiere end den i det Indvortes, er svage og feige og listige Menneskers foragtelige Tale om det Høieste. At antage, at Afgjørelsen i det Udvortes kan for evigt afgjøre Noget, saa det aldrig kan omgjøres, men ikke Afgjørelsen i det Indvortes, er Foragt for det Hellige.

At give Tænkningen Suprematiet over alt Andet er Gnosticisme; at gjøre Subjektets ethiske Virkelighed til den eneste Virkelighed kunde synes Akosmisme. At det saa vil forekomme en travl Tænker, der skal forklare Alt, et gesvindt Hoved, der overfarer hele Verden, beviser kun at han har meget ringe Forestilling om, hvad det Ethiske betyder for Subjektet. Naar Ethiken fratog en saadan travl Tænker den hele Verden og lod ham beholde sit eget Selv, saa vilde han formodentligen tænke: »er det Noget, en saadan Ubetydelighed er ikke værd at holde paa, lad saa den gaae med som Alt det Andet«; – og saa, saa er det Akosmisme. Men hvorfor vil dog en saadan travl Tænker tale og tænke saa despekteerligt om sig selv? Ja var Meningen den, at han skulde opgive hele Verden, og nøies med et andet Menneskes ethiske Virkelighed: ja, da havde han Ret i at lade haant om Byttet. Men den egne ethiske Virkelighed bør ethisk betyde Individet selv mere end Himmel og Jord med Alt hvad som derudi befindes, mere end Verdenshistoriens 6000 Aar og end Astrologien, Veterinair-Videnskaben samt Alt, hvad Tiden fordrer, hvilket æsthetisk og intellectuelt er en uhyre Bornerethed. Er det ikke saa, er det værst for Individet selv; thi saa har han slet Intet, slet ingen Virkelighed; thi til alt Andet har han kun et Muligheds-Forhold netop som Maximum.

Overgangen fra Mulighed til Virkelighed er, som Aristoteles rigtigt lærer, ϰινησις, en Bevægelse. Dette lader sig slet ikke sige i Abstraktionens Sprog eller forstaae deri, da dette netop hverken kan give Bevægelsen Tid eller Rum, hvilke forudsætte den, eller hvilke den forudsætter. Der er en Standsning, et Spring. Vil Nogen sige, at dette kommer deraf, at jeg tænker paa noget Bestemt, og ikke abstraherer, da jeg i saa Fald vilde indsee, at der intet Brud er: saa er mit gjentagne Svar: ganske rigtigt, abstrakt tænkt er der intet Brud, men heller ingen Overgang; thi abstrakt seet er Alt. Naar Existentsen derimod giver Bevægelsen Tid og jeg eftergjør dette, saa viser Springet sig som da et Spring kan vise sig: at det maa komme, eller at det har været. Lad os tage et Exempel af det Ethiske. Det er ofte nok sagt, at det Gode har sin Løn i sig selv, at det forsaavidt ikke blot er det Rigtigste, men ogsaa det Klogeste, at ville det Gode. En klog Eudaimonist kan meget godt indsee dette, han kan tænkende i Mulighedens Form komme det Gode saa nær som det er muligt, fordi i Muligheden som i Abstraktionen er Overgang kun et Skin. Men naar Overgangen skal blive virkelig, da udaander al Klogskab i Anfægtelsen. Den virkelige Tid adskiller ham det Gode og Lønnen saa længe saa evigt, at Klogskaben ikke kan faae det sammen igjen, og Eudaimonisten betakker sig. Ja vist er det at ville det Gode det Klogeste, men ikke i Klogskabens Forstand, men i det Godes Forstand. Overgangen er tydelig nok som et Brud, ja som en Lidelse. – I Prædikeforedraget forekommer ofte det Sandsebedrag, der eudaimonistisk forvandler Overgangen til at blive en Christen til et Skin, hvorved Tilhøreren bedrages og Overgangen forhindres.

Subjektiviteten er Sandheden; Subjektiviteten er Virkeligheden.

\sectionheading{Den subjektive Tænker; hans Opgave; hans Form (uddrag)  ·  SV¹ VII, 304–311}

Den subjektive Tænker er Dialektiker i Retning af det Existentielle; han har Tanke-Lidenskab til at fastholde den qvalitative Disjunktion. Men paa den anden Side skal den qvalitative Disjunktion bruges blank og bar, anvendes den aldeles abstrakt paa det enkelte Menneske, saa kan man udsætte sig for den latterlige Fare at sige noget uendeligt Afgjørende, og have Ret i hvad man siger og dog ikke sige det Mindste. Det er derfor i psychologisk Henseende ret mærkeligt, at see den absolute Disjunktion svigefuldt brugt netop til Udflugt. Naar der sættes Dødsstraf paa enhver Forbrydelse, saa ender det med, at slet ingen Forbrydelser straffes. Saaledes ogsaa med den absolute Disjunktion, blank og bart anvendt, den bliver ligesom et stumt Bogstav; den lader sig ikke udtale, eller den lader sig udtale, men siger Intet. Den absolute Disjunktion, som tilhørende Existentsen, har derfor den subjektive Tænker med Tanke-Lidenskab, men han har den som den sidste Afgjørelse, der forhindrer at ikke Alt lægger sig ud i en Qvantiteren. Han har den saaledes vel ved Haanden, men ikke saaledes, at han, ved abstrakt at recurrere til den, netop forhindrer Existents. Den subjektive Tænker har derfor tillige æsthetisk Lidenskab og ethisk Lidenskab, derved faaes Concretionen. Alle Existents-Problemer ere lidenskabelige, thi Existents, naar man bliver sig den bevidst, giver Lidenskab. At tænke over dem, saa man udelader Lidenskaben, er slet ikke at tænke over dem, er at glemme Pointet, at man jo selv er en Existerende. Dog er den subjektive Tænker ikke Digter, om han end tillige er Digter, ikke Ethiker, om han end tillige er Ethiker, men tillige Dialektiker, og væsentligen selv existerende, hvorimod Digterens Existents er uvæsentlig i Forhold til Digtet, og ligesaa Ethikerens i Forhold til Læren, og Dialektikerens i Forhold til Tanken. Den subjektive Tænker er ikke Videnskabsmand, han er Kunstner. At existere er en Kunst. Den subjektive Tænker er æsthetisk nok til at hans Liv faaer æsthetisk Indhold, ethisk nok til at regulere det, dialektisk nok til tænkende at beherske det.

Den subjektive Tænkers Opgave er at forstaae sig selv i Existents. Den abstrakte Tænkning taler jo rigtignok om Modsigelse, og om Modsigelsens immanente Fremstøden, uagtet den, ved at see bort fra Existents og det at existere, hæver Vanskeligheden og Modsigelsen. Men den subjektive Tænker er en Existerende, og dog er han Tænkende; han abstraherer ikke fra Existents og fra Modsigelsen, men han er i den, og dog skal han tænke. I al sin Tænken har han da at tænke det med, at han selv er en Existerende. Men saa vil han igjen ogsaa altid have nok at tænke paa. Den rene Menneskehed er man snart færdig med, og Verdenshistorien ogsaa, thi selv saadanne uhyre Portioner, som China, Persien og saa videre, sluger det hungrige Uhyre, den verdenshistoriske Proces, som ingen Ting. Det at troe er man, abstrakt seet, snart færdig med, men den subjektive Tænker, der, idet han tænker, tillige er hos sig selv i Existents, vil finde det uudtømmeligt, naar hans Tro skal declineres i Livets mangfoldige casibus. En Spas er det heller ikke, thi Existents er det Vanskeligste for en Tænker, naar han skal forblive i den, da Øieblikket er commensurabelt for de høieste Afgjørelser og dog igjen et lille forsvindende Minut i de mulige 70 Aar. Poul Møller har rigtigt bemærket, at en Hofnar bruger mere Vittighed i eet Aar end mangen vittig Forfatter i hele sit Liv. Og hvoraf kommer det, undtagen deraf, at den Første er en Existerende, der hvert Øieblik paa Dagen maa have Vittigheden til sin Disposition, den Anden En der momentviis er vittig.

[…]

I en vis Forstand taler den subjektive Tænker ligesaa abstrakt som den abstrakte Tænker, thi denne taler om den rene Menneskehed, den rene Subjektivitet, den Anden om det ene Menneske (unum noris, omnes). Men dette ene Menneske er et existerende Menneske, og Vanskeligheden ikke udeladt.

At forstaae sig selv i Existents er ogsaa det christelige Princip, kun at dette »selv« har faaet langt rigere og langt dybere Bestemmelser, hvilke ere endnu vanskeligere at forstaae sammen med det at existere. Den Troende er en subjektiv Tænker, og Forskjellen kun, som ovenfor blev viist, mellem den Eenfoldige og den eenfoldige Vise. Dette »sig selv« er atter ikke her den rene Menneskehed, den rene Subjektivitet og andet Saadant, hvorved Alt bliver let, da Vanskeligheden tages bort, og hele Sagen overflyttes i Abstraktionens Schattenspiel. Vanskeligheden er større end for Grækeren, fordi der er endnu større Modsætninger sat sammen, fordi Existents er accentueret paradox som Synd, og Evigheden paradox som Guden i Tiden. Vanskeligheden er at existere deri, ikke abstrakt at tænke sig ud deraf og abstrakt at tænke over for Exempel en evig Gudvorden og andet Saadant, der fremkommer, naar man tager Vanskeligheden bort. Den Troendes Existents er derfor endnu mere lidenskabelig end den græske Philosophs (der selv i Forhold til sin Ataraxie i høi Grad behøvede Lidenskab), thi Existents giver Lidenskab, men Existents paradox giver Lidenskabens Høieste.

[…]

Ethvert Menneske maa væsentligen antages at være i Besiddelse af hvad der væsentligen hører til at være Menneske. Den subjektive Tænkers Opgave er at forvandle sig selv til et Instrument, der tydeligt og bestemt udtrykker det Menneskelige i Existents. At trøste sig i denne Henseende til Differents er en Misforstaaelse, thi det at være lidt bedre Hoved og andet Saadant er kun Ubetydelighed. At vor Tid har taget sin Tilflugt til Generationen og forladt Individerne, har ganske rigtigt sin Grund i en æsthetisk Fortvivlelse, der ikke har naaet det Ethiske. Man har indseet, at det at være et nok saa udmærket enkelt Menneske ikke batter Noget, fordi ingen Differents batter. Saa har man valgt en ny Differents: at være født i det nittende Aarhundrede. Enhver gjør da saa hurtigt som muligt et Forsøg paa at bestemme sin Smule Existents i Forhold til Generationen, og trøster sig. Men det nytter ikke, og er kun et høiere og mere glimrende Bedrag. Og som der vel i Oldtiden og ellers i hver Generation har været Daarer, der i forfængelig Indbildning have forvexlet sig selv med en eller anden stor og udmærket Mand, villet være Den og Den: saa er vor Tids Forskjellighed, at Daarerne end ikke nøies med at forvexle sig med en stor Mand, men forvexle sig med Tiden, Aarhundredet, Generationen, Menneskeheden. – At ville være et enkelt Menneske (hvad man unægteligt er) ved Hjælp og i Kraft af sin Differents, er Blødagtighed; men at ville være et enkelt existerende Menneske (hvad man unægteligt er) i samme Forstand som enhver Anden kan være det: er den ethiske Seier over Livet og over alt Blendværk, den Seier, der maaskee af alle er den vanskeligste i det theocentriske nittende Aarhundrede.

Den subjektive Tænkers Form, hans Meddelelses Form er hans Stiil. Hans Form maa være ligesaa mangfoldig som de Modsætninger ere, han holder sammen. Det systematiske ein, zwei, drei, er en abstrakt Form, som derfor ogsaa maa komme i Forlegenhed hver Gang den skal anvendes paa det Concrete. I samme Grad som den subjektive Tænker er concret, i samme Grad maa ogsaa hans Form være concret dialektisk. Men som han selv ikke er Digter, ikke Ethiker, ikke Dialektiker, saaledes er hans Form heller ingen af disses ligefrem. Hans Form maa først og sidst forholde sig til Existents, og i denne Henseende maa han raade over det Digteriske, det Ethiske, det Dialektiske, det Religieuse. Sammenlignet med en Digter vil hans Form være forkortet, sammenlignet med en abstrakt Dialektiker vil hans Form være bred. Concretionen i det Existentielle er nemlig abstrakt seet Brede. Det Humoristiske for Exempel er i Forhold til abstrakt Tænkning Brede, men i Forhold til concret Existents-Meddelelse ingenlunde Brede, hvis det ikke i sig selv er bredt. En abstrakt Tænkers Person er ligegyldig i Forhold til Tanken, men existentielt maa en Tænkende fremstilles væsentlig som Tænkende, men saaledes at han i det han foredrager sin Tanke tillige skildrer sig selv. Spøg er i Forhold til abstrakt Tænkning Brede, men ikke i Forhold til concret Existents-Meddelelse, hvis Spøgen ikke selv er bred. Men digterisk Ro til at skabe i Phantasiens Medium og æsthetisk udføre interesseløst har den subjektive Tænker ikke, fordi han selv væsentligen er en Existerende i Existentsen og ikke har Phantasie-Mediet til den æsthetiske Frembringelses Illusion. Den digteriske Ro er Brede i Forhold til den subjektive Tænkers Existents-Meddelelse. Bipersoner, Scenerie og saa videre hvad der hører med til den æsthetiske Frembringelses Afrundethed i sig selv er Brede; thi den subjektive Tænker har blot een Scene, Existentsen, og Intet med Egne og andet Saadant at gjøre. Scenen er ikke i Phantasiens Trylleland, hvor Poesien elsker Fuldendelsen frem; Scenen er heller ikke i Engeland, og Omhu at anvende paa den historiske Nøiagtighed; Scenen er Inderligheden i at existere som Menneske, Concretionen er Existents-Categoriernes Forhold til hinanden. Historisk Nøiagtighed og historisk Virkelighed er Brede.

Men Existents-Virkelighed lader sig ikke meddele; og den subjektive Tænker har i sin egen ethiske Existents sin egen Virkelighed. Naar Virkelighed skal forstaaes af Trediemand, maa den forstaaes som Mulighed, og en Meddeler, der er sig dette bevidst, vil derfor agte paa, at hans Existents-Meddelelse, netop for at være i Retning af Existents, maa være i Mulighedens Form. En Fremstilling i Mulighedens Form lægger Modtageren det saa nær som det er muligt mellem Menneske og Menneske at existere deri. Lad mig endnu engang belyse dette. Man skulde troe, at man ved at fortælle, at Den og Den virkeligt har gjort Det og Det (det Store og Udmærkede), lagde en Læser det nærmere at ville gjøre det Samme, at ville existere i det Samme, end naar man blot fremstiller det som muligt. Fraseet, hvad der paa sit Sted er paaviist, at Læseren dog kun kan forstaae Meddelelsen ved at opløse Virkelighedens esse i posse, da han ellers blot indbilder sig at forstaae, fraseet dette, saa kan det at Den og Den virkeligt har gjort Det og Det ligesaa godt virke sinkende som tilskyndende. Læseren forvandler blot Den, om hvem der tales (ved Hjælp af at han er en virkelig Person) til den sjeldne Undtagelse; han beundrer ham og siger: men jeg er for ringe til at gjøre noget Saadant. Beundring kan nu være meget berettiget i Forhold til Differentserne, men er aldeles Misforstand i Forhold til det Almene. At En kan svømme over Canalen og at en Anden kan 24 Sprog, og en Tredie gaae paa Hovedet og saa videre: det kan man beundre si placet; men dersom Den der fremstilles skal være stor i Retning af det Almene, ved sin Dyd, sin Tro, sit Høimod, sin Trofasthed, sin Udholdenhed og saa videre, da er Beundring et svigefuldt Forhold, eller kan let blive det. Hvad der er stort i Retning af det Almene maa derfor ikke fremstilles som Gjenstand for Beundring men som Fordring. I Mulighedens Form bliver Fremstillingen en Fordring. Istedenfor, som almindeligen gjøres, at fremstille det Gode i Virkelighedens Form, at Den og Den virkelig har levet og virkelig har gjort det, og saa forvandle Læseren til en Betragter, en Beundrer, en Vurderingsmand, skal det fremstilles i Mulighedens Form; saa er det lagt Læseren saa nær som muligt, om han vil existere deri. Muligheden opererer med det ideale Menneske (ikke i Retning af Differents, men i Retning af det Almene), der forholder sig til ethvert Menneske som Fordring. I samme Grad som man urgerer, at det var dette bestemte Menneske, gjør man de Andre Exceptionen lettere. Man behøver just ikke at være nogen Psycholog for at vide, at der gives en Svigefuldhed, der vil excipere mod det ethiske Indtryk netop ved Hjælp af Beundring. Istedenfor at det ethiske og det religieuse Forbillede skal vende Betragterens Blik ind i ham, skal støde fra, hvilket netop skeer ved at sætte Muligheden som den fælleds imellem dem, drager Fremstillingen i Virkelighedens Form æsthetisk en Mængdes Øine paa sig; og der bliver drøftet og prøvet, og vendt og dreiet om nu virkeligt og saa videre, og beundret og flæbet, at nu virkelig og saa videre Det at Job troede, for at tage dette Exempel, skal fremstilles saaledes, at det for mig kommer til at betyde, om jeg nu ogsaa vil erhverve Troen; men ingenlunde skal det betyde, at jeg er paa Comedie eller Medlem af et høistæret Publikum, der skal undersøge, om nu virkeligt, og applaudere, at nu virkeligt. Det er saaledes en lavkomisk Bekymring som en følende Menighed og sammes enkelte Lemmer stundom har betræffende den beskikkede Sjelesørger: om han nu virkelig; og en lavkomisk Glæde og Beundring over at have en Sjælesørger, om hvem det er vist, at han virkelig og saa videre Det er i al Evighed usandt, at Nogen er bleven hjulpen til at gjøre det Gode derved, at en anden virkelig har gjort det; thi kom han selv til virkelig at gjøre det: da var det ved at forstaae den Andens Virkelighed som Mulighed. Da Themistocles blev søvnløs ved Forestillingen om Miltiades's Triumpher, da var det Det, at han forstod Virkeligheden som Mulighed, der gjorde ham søvnløs; havde han havt travlt med, om nu virkelig Miltiades, og været nøiet med, at nu virkelig Miltiades havde gjort det: saa var han neppe bleven søvnløs, men vel en søvnig Beundrer eller höchstens en lydende Beundrer, men ingen Miltiades Nummer 2. Og ethisk forstaaet er der Intet man sover saa godt paa, som paa Beundring over en Virkelighed. Og ethisk forstaaet, dersom Noget kan purre et Menneske op, saa er det Muligheden, naar den idealt fordrer sig selv af et Menneske.

\sectionheading{Den existentielle Pathos: Lidelsen; Humor; Religieusiteten A (uddrag)  ·  SV¹ VII, 434–499}

Den religieuse Lidelses Betydning er Afdøen fra Umiddelbarheden; dens Virkelighed er dens væsentlige Vedvaren; men den tilhører Inderligheden og maa ikke udtrykke sig i det Udvortes (Klosterbevægelsen). Naar vi nu tage en Religieus, den skjulte Inderligheds Ridder, og anbringe ham i Existents-Mediet, saa vil der, idet han forholder sig til Omverdenen, fremkomme en Modsigelse, og denne maa han blive sig bevidst. Modsigelsen ligger ikke i, at han er forskjellig fra alle Andre (thi denne Selvmodsigelse er netop Loven for det Comiskes Nemesis over Klosterbevægelsen), men Modsigelsen er, at han med hele denne Inderlighed skjult i sig, med dette Lidelsens og Velsignelsens Svangerskab i det Indre, seer ganske ud som andre Mennesker – og derved skjules jo netop Inderligheden, at han seer ganske ud som Andre. Noget Comisk er her tilstede, thi her er en Modsigelse, og hvor der er en Modsigelse er det Comiske ogsaa tilstede. Dette Comiske er imidlertid ikke for Andre, som Intet vide af det, men er for den Religieuse selv, naar Humor er hans Incognito, som Frater Taciturnus siger (confer Stadier paa Livets Vei). Og dette er vel værd nærmere at forstaae; thi næst efter den Forvirring i den nyere Spekulation, at Tro er Umiddelbarhed, er maaskee den meest forstyrrende, at Humor er det Høieste, thi Humor er endnu ikke Religieusitet, men dennes Confinium; hvorom man allerede i det Foregaaende vil finde nogle Bemærkninger, som jeg maa bede Læseren at erindre.

Dog er Humor den Religieuses Incognito? Er hans Incognito ikke dette, at der slet Intet er at mærke, slet Intet der kunde vække Mistanke om den skjulte Inderlighed, end ikke saa meget som det Humoristiske? I sit høieste Maximum, hvis det lader sig naae i Existents, vilde dette vel være saa, men saa længe Striden og Lidelsen i Inderligheden vedvedbliver, saa længe vil det ikke lykkes ham ganske at skjule Inderligheden, men han vil ikke udtrykke den ligefrem, og negativt forhindre det ved Hjælp af det Humoristiske. En Iagttager, der gik ud mellem Menneskene for at opdage den Religieuse, vilde derfor følge det Princip, at Enhver, hos hvem han opdagede det Humoristiske, vilde blive gjort til Gjenstand for hans Opmærksomhed. Men dersom han har gjort sig Inderlighedens Forhold tydelige, vil han ogsaa vide, at han kan blive narret, thi den Religieuse er ikke Humoristen, men i sit Ydre er han Humoristen. Saaledes vilde en Iagttager, der søger den Religieuse og vil kjende ham paa det Humoristiske, blive narret, hvis han traf mig; han vilde finde det Humoristiske, men være narret, hvis han gjorde nogen Slutning deraf, thi jeg er ikke den Religieuse, men ene og alene Humorist. Maaskee mener Nogen, at det er en skrækkelig Anmasselse at tillægge mig selv Navn af en Humorist, og mener tillige, at hvis jeg virkeligen var en Humorist, saa skulde han nok vise mig Hæder og Ære: jeg skal hverken opholde mig derover eller derved, thi den, der gjør Indvendingen, er aabenbart antagende Humor for det Høieste. Jeg derimod siger jo, at uendeligt høiere end Humoristen er den Religieuse stricte sic dictus, og qvalitativ forskjellig fra Humoristen. Og hvad det forresten angaaer, at han ikke vil ansee mig for Humorist: nu vel, saa skal jeg være villig til at overføre Iagttagerens Situation fra mig til den, der gjør Indvendingen, lade Iagttageren være bleven opmærksom paa ham: Resultatet bliver det samme – Iagttageren er narret.

Der er tre Existents-Sphærer: den æsthetiske, den ethiske, den religieuse. Til disse svare to Confinier: Ironie er Confiniet imellem det Æsthetiske og Ethiske; Humor Confiniet mellem det Ethiske og det Religieuse.

Lad os tage Ironie. Saasnart en Iagttager opdager en Ironiker, vil han blive opmærksom, thi det er muligt, at Ironikeren er en Ethiker; men han kan ogsaa narres, thi det er ikke sagt, at Ironikeren er en Ethiker. Den Umiddelbare er strax kjendelig, og saasnart han er kjendt, er det givet, at han ikke er nogen Ethiker; thi han har ikke gjort Uendelighedens Bevægelse. Den ironiske Replik, naar denne er correct (og Iagttageren antages at være en forsøgt Mand, der veed Besked om at narre og forstyrre den Talende for at see, om det er noget udenad lært, eller der er rigeligt ironisk Valuta, som en existerende Ironiker altid vil have det), forraader, at den Talende har gjort Uendelighedens Bevægelse, men heller ikke mere. Ironien fremkommer ved ideligt at sætte Endelighedens Enkeltheder sammen med den ethiske uendelige Fordring og lade Modsigelsen blive til. Den der med Færdighed kan gjøre det, saa han ikke lader sig fange i nogen Relativitet, hvor hans Færdighed bliver sky, han maa have gjort en Uendelighedens Bevægelse, forsaavidt maa det være muligt at han kan være Ethiker. Iagttageren vil derfor end ikke kunne fange ham paa, at han ikke formaaer at opfatte sig selv ironisk, thi ogsaa dette formaaer han, at tale om sig selv som om en Trediemand, at sætte sig selv som en forsvindende Enkelthed sammen med den absolute Fordring, ja at sætte dem sammen; hvor forunderligt, at det Udtryk som betegner Existentsens sidste Vanskelighed, der netop er at sætte det absolute Forskjellige sammen (som Forestillingen om Gud med det at tage i Dyrehaven), at det samme Udtryk i Sproget ogsaa betegner Drilleriet! Men uagtet dette er givet, er det endnu ikke givet, at han er Ethikeren. Ethiker er han kun ved i sig selv at forholde sig til den absolute Fordring. En saadan Ethiker bruger Ironien som sit Incognito. Socrates var i denne Betydning en Ethiker, dog vel at mærke, hen til Grændsen af det Religieuse, hvorfor der jo ogsaa i det Foregaaende (andet Afsnit Capitel 2) blev paaviist det til Troen Analoge i hans Existents. Hvad er saa Ironie, hvis man vil kalde Socrates en Ironiker, og ikke som Magister Kierkegaard bevidst eller ubevidst kun vil drage den ene Side frem? Ironie er Eenheden af ethisk Lidenskab, der i Inderlighed uendeligt accentuerer det egne Jeg i Forhold til den ethiske Fordring – og af Dannelse, der i Udvorteshed uendeligt abstraherer fra det egne Jeg, som en Endelighed med blandt alle de andre Endeligheder og Enkeltheder. Denne Abstraktion gjør, at Ingen mærker det Første, og deri ligger netop Kunsten, og dermed betinges det Førstes sande Uendeliggjørelse. Mængden af Mennesker lever omvendt; de have travlt med at være Noget, naar Nogen seer derpaa; de ere om muligt Noget i deres egne Øine, saasnart Andre see paa dem; men inderst inde, hvor den absolute Fordring seer paa dem, der have de ikke Smag for at accentuere det egne Jeg.

Ironie er en Existents-Bestemmelse, og Intet er da latterligere end naar man troer det er en Tale-Form, eller naar en Forfatter priser sig lykkelig ved en Gang imellem at udtrykke sig ironisk. Den der har væsentligen Ironie har den saa lang Dagen er, ikke bunden i nogen Form, fordi den er Uendeligheden i ham.

Ironie er Aandens Dannelse, og følger derfor nærmest efter Umiddelbarheden; saa kommer Ethikeren, saa Humoristen, saa den Religieuse.

Men hvorfor bruger nu Ethikeren Ironie som sit Incognito? Fordi han fatter den Modsigelse, der er mellem den Maade, paa hvilken han existerer i sit Inderste, og at han i sit Ydre ikke udtrykker det; thi vel bliver Ethikeren aabenbar, forsaavidt han udtømmer sig i den faktiske Virkeligheds Opgaver, men det gjør den Umiddelbare jo ogsaa, og det hvorved han er Ethikeren, er den Bevægelse, ved hvilken han sætter sit Liv udefter ind efter sammen med det Ethiskes uendelige Fordring, og dette sees ikke ligefrem. For ikke at lade sig forstyrre af Endeligheden, af alt det Relative i Verden, sætter Ethikeren det Comiske mellem sig og Verden, og sikkrer sig derved, at han ikke selv bliver comisk ved naiv Misforstand af sin ethiske Lidenskab. En umiddelbart Begeistret vræler ud i Verden aarle og silde, altid cothurnisk paa Støvlerne plager han Folk med sin Begeistring, han mærker slet ikke, at det ikke begeistrer dem, med mindre de da slaae ham; han er vel nok underrettet og Ordren lyder paa en fuldkommen Omskabning – af hele Verden, ja her er det han har hørt feil, thi Ordren lyder paa en fuldkommen Omskabelse af sig selv. Er en saadan Begeistret samtidig med en Ironiker, tager denne sig naturligviis ham comisk til Indtægt. Ethikeren derimod er ironisk nok til meget godt at see, at det, der beskæftiger ham absolut, ikke beskæftiger de Andre absolut; han opfatter selv dette Misforhold, og sætter det Comiske mellem, for desto inderligere at kunne fastholde det Ethiske i sig selv. Nu begynder Comedien; thi Menneskenes Dom over en Saadan vil altid blive: for ham er Intet vigtigt. Og hvorfor ikke? Fordi det Ethiske for ham er absolut vigtigt, thi deri er han forskjellig fra Menneskene i Almindelighed, for hvem saa mange Ting ere vigtige, ja næsten Alt vigtigt – men Intet absolut vigtigt. – Dog, som sagt, en Iagttager kan blive narret, hvis han antager en Ironiker for en Ethiker, thi Ironien er kun Muligheden.

Saaledes igjen med Humoristen og den Religieuse, da jo ifølge det Foregaaende det Religieuses egen Dialektik forbyder det ligefremme Udtryk, forbyder den kjendelige Forskjellighed, protesterer Udvorteshedens Commensurabilitet, og dog ærer, naar galt skal være, Klosterbevægelsen langt over Mediationen. Humoristen sætter bestandigt (ikke i Betydning af Præstens »altid«, men til enhver Tid paa Dagen, hvor han er og hvad han tænker eller foretager sig) Guds-Forestillingen sammen med Andet og bringer Modsigelsen frem – men han forholder sig ikke selv i religieus Lidenskab (stricte sic dictus) til Gud, han forvandler sig selv til et spøgende og dog dybsindigt Gjennemgangssted for alle disse Omsætninger, men han forholder sig ikke selv til Gud. Den Religieuse gjør det Samme, han sætter Guds-Forestillingen sammen med Alt og seer Modsigelsen, men i sit Inderste forholder han sig til Gud; medens den umiddelbare Religieusitet hviler i den fromme Overtro, ligefrem at see Gud i Alt; og den Opvakte næsviist engagerer Gud til at være med hvor han er, saa man, bare man seer den Opvakte, kan være sikker paa at Gud er med, da den Opvakte har ham i Lommen. Religieusitet med Humor som Incognito er derfor: Eenhed af absolut (dialektisk inderliggjort) religieus Lidenskab og Aands-Modenhed, der kalder Religieusiteten tilbage fra al Udvorteshed i Inderlighed, og atter deri er jo igjen den absolute religieuse Lidenskab. Den Religieuse opdager, at hvad der beskæftiger ham absolut kun lidet synes at beskæftige Andre, men slutter Intet, da han deels ikke har Tid, og deels ikke med Bestemthed kan vide, om alle disse Mennesker dog ikke ere Riddere af den skjulte Inderlighed; han lader sig af Omverdenen nødsage til at gjøre, hvad den dialektiske Inderliggjørelse fordrer af ham, at sætte Skjulet mellem Menneskene og sig, for at værge og betrygge Lidelsens og Guds-Forholdets Inderlighed. Deraf følger dog ikke, at en saadan Religieus bliver uvirksom, tvertimod, han gaaer jo netop ikke ud af Verden, men bliver i den, thi dette er jo netop hans Incognito. Men sin Virken ud efter inderliggjør han ind efter for Gud ved at vedgaae at han formaaer Intet, ved at afbryde ethvert teleologisk Forhold til den i Retningen ud efter, enhver Provenu af den i Endeligheden, om han dog arbeider af yderste Evne; og dette er netop Begeistring. En Opvakt øger altid Guds Navn til ud efter, hans Troens Vished er sikker nok. Men Troens Vished er jo kjendelig paa Uvisheden, og som dens Vished er den høieste af alle, saa er denne samme Vished den meest ironiske af alle, ellers er den ikke Troens Vished. Det er vist, at Alt hvad der behager Gud, lykkes den Fromme, det er vist, o saa vist, ja Intet er saaledes vist som dette. Men nu det Næste, og behag at mærke, at Undersøgelsen ikke er paa Papiret, men i Existentsen, og den Troende en enkelt Existerende i Existentsens Concretion. Altsaa dette er det evig Visse, at hvad der behager Gud lykkes den Fromme. Men nu det Næste; hvad er det, der behager Gud? Er det Dette eller Hiint, er det denne Livsstilling han skal vælge, denne Pige han skal ægte, dette Arbeide han skal begynde, dette Foretagende han skal opgive? Ja maaskee, og maaskee ikke. Er dette ikke ironisk nok? Og dog er det evigt vist, og der er Intet saa vist, som dette, at hvad der behager Gud det lykkes den Fromme. Ja, men derfor skal den Religieuse ikke saa meget bekymre sig om det Udvortes, men eftertragte de høieste Goder, Sjelens Fred, sin Sjels Frelse: dette behager altid Gud. Og det er vist, o saa vist som at Gud lever, at hvad der behager Gud, det lykkes den Fromme. Det behager altsaa Gud, at han skal gjøre det; men naar skal det lykkes ham? Er det strax, eller om et Aar, eller maaskee først ved det jordiske Livs Ende, kan Striden og Prøvelsen ikke vare saa længe? Maaskee, maaskee ikke. Er dette ikke ironisk nok? Og dog er det vist, o saa vist, at hvad der behager Gud lykkes den Fromme; brister denne Vished, saa brister Troen, men ophører Uvisheden, der er dens Kjende og dens Form, saa ere vi ikke komne videre i Religieusitet, men tilbage til barnagtige Former. Saasnart Uvisheden ikke er Vishedens Form, saasnart Uvisheden ikke bestandigt holder den Religieuse svævende, for bestandigt at gribe Visheden, saasnart Visheden ligesom plomberer den Religieuse, ja saa er han naturligviis ifærd med at blive Masse. – Men af den skjulte Inderlighed med Humor som Incognito synes det at følge, at den Religieuse er sikkret mod at blive Martyr, hvad den Opvakte hellere end gjerne vil. Ja vist er den skjulte Inderligheds Ridder sikkret, han er et Dæggebarn i Sammenligning med den Opvakte, der gaaer freidigt Martyriet imøde – med mindre Martyriet ligger i den Tilintetgjørelsens Lidelse, der er Afdøen fra Umiddelbarheden, selve det Guddommeliges Modstand mod en Existerende, der forhindres i at forholde sig absolut, og endeligen Livet i Verden med denne Inderlighed uden at have et Udtryk derfor. Psychologisk gjelder ganske simpelt den Lov, at den samme Kraft, der, naar den vendes ud efter, formaaer det og det, den samme Kraft behøves der en endnu større Kraft til for at forhindre i at gribe ud efter. Thi Kraften vendt ud efter, Modstanden udenfra: saa er Modstanden kun halvt at anslaae for Modstand, halvt er den Understøttelse. Den skjulte Inderlighed har sit Martyrium i sig selv. Men saa er det jo muligt, at hvert andet Menneske er en saadan Ridder af den skjulte Inderlighed? Ja, hvorfor ikke. Hvem kan det skade. Maaskee En og Anden, der dog har nogen Religieusitet og finder det uforsvarligt, at dette ikke skal paaskjønnes; altsaa En, der ikke vilde formaae at udholde det Syn, at den meest lidenskabelige Inderlighed i det Udvortes skuffende ligner sin Modsætning. Men hvorfor vælger en saadan Religieus ikke Klosteret, der er der endog Avancement og Befordring, en Rangforordning for de Religieuse. Den skjulte Inderligheds sande Ridder kan det ikke forstyrre, ham beskæftiger det alene at være det, mindre at synes (forsaavidt han maa bruge nogen Anstrengelse for at forhindre det), slet Intet, om alle andre Mennesker ansees for at være det.

[…]

Nu staae vi ved Grændsen. Den Religieusitet der er den skjulte Inderlighed er eo ipso utilgængelig for den comiske Opfattelse. Det Comiske kan den ikke have uden for sig, netop fordi den er skjult Inderlighed, og altsaa ikke kan komme i Modsigelse med Noget. Modsigelsen som Humor dominerer, det Comiskes høieste Omfang, har den selv bragt sig til Bevidsthed, og har hos sig som det Lavere. Saaledes er den absolut væbnet mod det Comiske, eller ved det Comiske sikkret mod det Comiske.

Naar Religieusiteten i Kirke og Stat stundom har villet have Lovgivning og Politie til Hjælp for at sikkre sig mod det Comiske, saa kan dette være meget velmeent, men et Spørgsmaal er, hvorvidt det Bestemmende i sidste Instants er noget Religieust; og en Uretfærdighed er det mod det Comiske at ansee dette for en Fjende af det Religieuse. Saa lidet som Dialektik, saa lidet er det Comiske en Fjende af det Religieuse, hvilket tvertimod Alt tjener og lystrer. Men den Religieusitet, der væsentligen gjør Prætension paa Udvorteshed, væsentligen gjør Udvortesheden commensurabel, skal see sig vel for, og mere frygte sig selv (at den ikke blive Æsthetik) end frygte det Comiske, der kunde berettiget hjælpe den til at faae Øinene op. Meget i Catholicismen kan tjene til Exempel herpaa. Og hvad Individet angaaer, da gjelder det: den Religieuse, der vil at Alle skulle være alvorlige, vel endog netop saaledes alvorlige som han, fordi han er dum-alvorlig: han er i en Modsigelse; og den Religieuse, som ikke kunde taale, hvis saa skulde være, at Alle loe af hvad der beskjæftiger ham absolut, han mangler Inderlighed, og vil derfor trøstes ved Sandsebedrag, ved at mange ere af samme Mening, ja af samme Ansigt som han, og opbygges ved at øge det Verdenshistoriske til sin Smule Virkelighed, »da der jo nu overalt begynder at røre sig et nyt Liv, den bebudede Nyaarstid med Syn og Hjerte for Sagen.«

Den skjulte Inderlighed er utilgængelig for det Comiske, hvilket ogsaa sees deraf, at hvis en saadan Religieus kunde hidses til pludseligen at gjøre sin Religieusitet gjeldende i det Udvortes, hvis han for Exempel glemte sig selv og kom i Strid med en comparativ Religieus, og atter glemte sig selv og Inderlighedens absolute Fordring ved comparativt at ville være mere religieus end den Anden: saa er han comisk, og Modsigelsen er: paa eengang at ville være synlig og usynlig. Mod anmassende Former af det Religieuse bruger Humor berettiget det Comiske, netop fordi en Religieus nok selv maa vide Udveien, naar han bare vil. Tør dette ikke forudsættes, saa bliver en saadan Opfattelse mislig i samme Forstand som Opfattelsen af den Stundesløse comisk vilde være det, hvis det var givet, at han virkelig var sindssvag.

Loven for det Comiske er ganske simpel: det er overalt, hvor der er Modsigelse, og hvor Modsigelsen er smerteløs derved, at den sees hævet, thi vel hæver det Comiske ikke Modsigelsen (gjør den tvertimod aabenbar), men det berettigede Comiske kan det, ellers er det ikke berettiget. Talentet ligger i in concreto at kunne fremstille det. Prøven paa det Comiske er at eftersee, hvilket Forhold mellem Sphærerne det comiske Udsagn indeholder; er Forholdet ikke rigtigt, er det Comiske uberettiget; og en Comik, der slet intetsteds har hjemme, eo ipso uberettiget. Det Sophistiske i Forhold til det Comiske har derfor sit Tilhold i Intet, i den rene Abstraktion, og er udtrykt af Gorgias ved den Abstraktion: at tilintetgjøre Alvor ved Comik og Comik ved Alvor (confer Aristoteles Rhetorik 3, 18). Det Qvit som Alt her ender med er Skidt, og Misligheden let opdaget, at en Existerende har forvandlet sig selv til et phantastisk x; thi det maa jo dog være en Existerende, der vil bruge denne Fremgangsmaade, der kun gjør ham selv latterlig, naar man mod ham bruger Besværgelses-Formularen mod Spekulanter fra det Foregaaende: maa jeg have den Ære at spørge hvem jeg har den Ære at tale med, om det er et Menneske og saa videre Gorgias lander nemlig med sin Opdagelse i den rene Værens phantastiske Overdrev; thi naar han ved det Ene tilintetgjør det Andet, saa bliver der Intet tilbage. Dog har vel Gorgias nærmest villet betegne en Lommeprokurators Snildhed, der seirer ved at forandre Vaaben i Forhold til Modstanderens Vaaben; men en Lommeprokurator er ingen berettiget Instants i Forhold til det Comiske, han maa skyde en hvid Pind efter Berettigelsen – og nøies med Profiten, hvilket ogsaa som bekjendt har været alle Sophisters Yndlingsresultat – Penge, Penge, Penge, eller hvad der staaer paa lige Trin dermed.

I den religieuse Sphære, naar denne holdes reen i Inderlighed, er det Comiske tjenende. Man kunde sige, at Anger for Exempel er en Modsigelse, ergo er der noget Comisk, vel ikke for det Æsthetiske, eller for den endelige Forstandighed, som er lavere, eller for det Ethiske, som har sin Kraft i denne Lidenskab, eller for Abstraktionen, som er phantastisk og derved lavere (det var fra dette Standpunkt at ville opfatte den comisk, der blev afviist som Nonsens i det Foregaaende), men for det Religieuse selv, som veed et Middel derimod, en Udvei. Men dette er ikke saa, det Religieuse veed intet Middel mod Angeren, som seer bort fra Angeren, det Religieuse bruger tvertimod bestandigt det Negative som den væsentlige Form, Syndsbevidstheden er saaledes et bestemt Medhenhørende til Bevidstheden om Syndsforladelsen. Det Negative er ikke eengang for alle og saa det Positive, men det Positive er bestandigt i det Negative, og det Negative er Kjendet, saa det regulerende Princip: ne quid nimis her ikke kan finde Anvendelse. Naar det Religieuse opfattes æsthetisk, naar der i Middelalderen prækes Aflad for 4 Skilling., og man antager, at dermed er den Sag qvit, hvis man vil fastholde denne Fiction: saa bliver Angeren comisk at opfatte, saa er den i Angeren Sønderrevne comisk ligesom den Stundesløse, forudsat, at han da har de 4 Skilling., thi Udveien er jo saa let, og i Fictionen er det jo antaget, at det er Udveien. Men alt dette Galimathias er Følge af at det Religieuse er blevet Farce. Men i samme Grad som man afskaffer det Negative i den religieuse Sphære, eller lader det have været eengang for alle og saa dermed nok, i samme Grad vil det Comiske gjøre sig gjeldende mod det Religieuse og med Rette – fordi det Religieuse er bleven Æsthetik og dog vil være det Religieuse.

[…]

Humor sætter Skyldens evige Erindren sammen med Alt, men forholder sig ikke selv i denne Erindren til en evig Salighed. Nu staaer vi ved den skjulte Inderlighed. Skyldens evige Erindren kan ikke udtrykkes i det Udvortes, der er incommensurabelt derfor, da ethvert Udtryk i det Udvortes endeliggjør Skylden. Men den skjulte Inderligheds evige Erindren af Skylden er dog heller ikke Fortvivlelse; thi Fortvivlelse er bestandigt det Uendelige, det Evige, det Totale i Utaalmodighedens Øieblik, og al Fortvivlelse er en Art Arrigskab. Nei, den evige Erindren er Kjendet paa Forholdet til en evig Salighed, saa langt som muligt fra at være det ligefremme Kjende, men dog altid nok til at forhindre Fortvivlelsens Springen fra.

Humor opdager det Comiske ved at sætte den totale Skyld sammen med al den Relativitet mellem Mand og Mand. Det Comiske ligger i at den totale Skyld er det Tilgrundliggende, der bærer hele denne Comedie. Dersom nemlig den væsentlige Uskyld eller Godhed laae til Grund for det Relative, er det ikke comisk, thi det er ikke comisk at man indenfor den positive Bestemmelse bestemmer Mere og Mindre. Men naar Relativiteten hviler paa den totale Skyld, hviler altsaa Mere og Mindre paa hvad der er mindre end Intet, og dette er Modsigelsen, som det Comiske opdager. Forsaavidt Penge er et Noget, er Relativiteten mellem Rigere og Fattigere ikke comisk, men dersom det er Regne-Penge er det comisk at der er en Relativitet. Naar der til Grund for Menneskenes travle Løben omkring ligger en Mulighed af at undgaae Faren, er Travlheden ikke comisk, men dersom det for Exempel er paa et Skib som synker, er der noget Comisk i al denne Renden omkring, thi Modsigelsen er, at trods al Bevægelse bevæger man sig ikke bort fra det Sted, hvor Undergangen er.

Den skjulte Inderlighed maa ogsaa opdage det Comiske, ikke derved at den Religieuse er forskjellig fra Andre, men derved, at han, skjøndt tungest besværet ved at bære Skyldens evige Erindren, er ligesom alle Andre. Han opdager det Comiske, men da han i den evige Erindren bestandigt forholder sig til en evig Salighed, er det Comiske et bestandigt Forsvindende.

Det fremsatte Problem (confer andet Afsnit, Capitel IV.) var et existentielt, og som saadant pathetisk-dialektisk. Den første Deel (A) er behandlet, den pathetiske Deel: Forholdet til en evig Salighed. Nu skal der gaaes over til det Dialektiske (B), der er det Afgjørende for Problemet. Thi den Religieusitet, som hidtil er behandlet, og som fremtidigen for Kortheds Skyld vil forekomme under Benævnelsen: Religieusiteten A, er ikke den specifik christelige. Paa den anden Side er det Dialektiske kun forsaavidt det Afgjørende, som det sættes sammen med det Pathetiske til ny Pathos.

I Almindelighed er man ikke samtidigen opmærksom paa begge Dele. Det religieuse Foredrag vil representere det Pathetiske og slaae en Streg over det Dialektiske, og er derfor, hvor velmenende det end er, stundom en confus tumultuarisk Pathos af Allehaande, Æsthetik, Ethik, Religieusitet A og Christendom, er derfor stundom selvmodsigende, »men der er deilige Steder deri,« især deilige for Den, som skal handle og existere derefter. Det Dialektiske hævner sig ved skjult og ironisk at spotte Gestus og de store Ord, og fremfor Alt ved den ironiske Dom over et religieust Foredrag, at det godt lader sig høre, men ikke lader sig gjøre. Videnskaben vil tage sig af det Dialektiske, og fører det til den Ende over i Abstraktionens Medium, hvorved Problemet atter forfeiles, da det er et Existents-Problem, og den egentlige dialektiske Vanskelighed forsvinder ved at forklares i Abstraktionens Medium, der seer bort fra Existents. Er det tumultuariske religieuse Foredrag for følsomme Mennesker, der ere hurtige til at svede og til at svede ud, saa er den speculative Opfattelse for rene Tænkere; men ingen af Delene er for handlende og, i Kraft af at handle, existerende Mennesker.

Distinktionen mellem det Pathetiske og det Dialektiske maa dog nærmere bestemmes, thi Religieusiteten A er ingenlunde udialektisk, men den er ikke paradox dialektisk. Religieusiteten A er Inderliggjørelsens Dialektik; den er Forholdet til en evig Salighed ikke betinget ved et Noget, men er Forholdets dialektiske Inderliggjørelse, altsaa kun betinget ved Inderliggjørelsen, der er dialektisk. Religieusiteten B derimod, som den fremtidigen vil blive kaldet, eller den paradoxe Religieusitet, som den er bleven kaldet, eller den Religieusitet, der har det Dialektiske paa andet Sted, gjør Betingelser, saaledes at Betingelserne ikke ere Inderliggjørelsens dialektiske Fordybelser, men et bestemt Noget, der nærmere bestemmer den evige Salighed (medens i A Inderliggjørelsens nærmere Bestemmelse er den eneste nærmere Bestemmelse), ikke idet den nærmere bestemmer Individets Tilegnelse af den, men nærmere bestemmer den evige Salighed dog ikke som Opgave for Tænkning, men netop paradox som frastødende til ny Pathos.

Religieusiteten A maa først være tilstede i Individet førend der kan være Tale om at blive opmærksom paa det dialektiske Bind Naar Individet i den existentielle Pathos' meest afgjørende Udtryk forholder sig til en evig Salighed: saa kan der være Tale om at blive opmærksom paa, hvorledes det Dialektiske paa andet Sted (secundo loco) støder ham ned i det Absurdes Pathos. Man vil derfor see, hvor taabeligt det er, naar en Mand uden Pathos vil forholde sig til det Christelige; thi førend der overhovedet kan være Tale om at være blot i den Situation, at man kan blive opmærksom paa det, maa man først existere i Religieusiteten A. Det Forkeerte er imidlertid ofte nok skeet: man har uden videre taget sig Christus og Christendommen og det Paradoxe og det Absurde, kort alt det Christelige, til Indtægt i æsthetisk Galimathias, ret som om Christendommen var et gefundenes Fressen for Dumrianer, fordi den ikke kan tænkes, og ret som ikke netop den Bestemmelse, at den ikke kan tænkes, er den vanskeligste af alle at fastholde, naar man skal existere deri, den vanskeligste at fastholde – især for gode Hoveder.

Religieusiteten A kan være i Hedenskabet, og i Christendommen være Enhvers, der ikke afgjørende er Christen, han være nu døbt eller ikke. Det forstaaer sig, at blive en wohlfeil Udgave af en Christen i al Magelighed er langt lettere, og tillige jo saa godt som det Høieste, han er jo døbt, han har jo faaet et Exemplar af Bibelen og Psalmebogen forærende: er han saa ikke en Christen, en evangelisk luthersk Christen? Dog, det bliver da Vedkommendes Sag; min Mening er, at Religieusiteten A (i hvis Confinium jeg har min Existents) er saa anstrengende for et Menneske, at der i den altid er nok af Opgave; min Hensigt er det, at gjøre det vanskeligt at blive en Christen, dog ikke vanskeligere end det er, vanskeligt heller ikke for dumme Mennesker og let for gode Hoveder, men qvalitativt vanskeligt og væsentligen lige vanskeligt for ethvert Menneske, thi det er, væsentligen seet, lige vanskeligt for ethvert Menneske at opgive sin Forstand og sin Tænkning og at holde sin Sjel paa det Absurde, og comparativt vanskeligst for Den, der har megen Forstand, naar man erindrer, at ikke Enhver, der ikke tabte sin Forstand over Christendommen, derved beviser, at han har den. Min Hensigt er det, det vil da sige, saadan som en Experimenterende, der gjør Alt for sin egen Skyld, kan have en Hensigt. Ethvert Menneske, det viseste og det eenfoldigste, kan qvalitativt lige væsentligen (det Comparative giver Misforstaaelsen, som naar et godt Hoved sammenligner sig med en Eenfoldig, istedenfor at forstaae, at den samme Opgave er for hver især og ikke for de Tvende i Sammenligning) gjøre Distinktionen mellem hvad han forstaaer og hvad han ikke forstaaer (det forstaaer sig, den vil være Frugten af hans høieste Anstrengelse, denne anstrengede Sluttethed, og der ligger 2000 Aar mellem Socrates og Hamann: denne Distinktions Hævdere), kan opdage, at der er Noget, som er trods mod hans Forstand og Tænkning. Naar han paa dette Absurde sætter hele sit Liv ind, saa gjør han Bevægelsen i Kraft af det Absurde, og han er væsentligen bedragen, dersom det Absurde, han har valgt, lader sig bevise ikke at være det Absurde. Dersom dette Absurde er Christendommen, saa er han troende Christen; men forstaaer han det, at det ikke er det Absurde, saa er han eo ipso ikke mere troende Christen (han være nu døbt, confirmeret, Ihændehaver af Bibelen og Psalmebogen, om det saa var den nye forventede Psalmebog), indtil han atter tilintetgjør Forstaaelsen som Sandsebedrag og Misforstaaelse, og forholder sig til det christelige Absurde. Dersom nemlig Religieusiteten A ikke gaaer imellem som terminus a qvo for den paradoxe Religieusitet: saa er Religieusiteten A høiere end B, thi saa er Paradoxet, det Absurde og saa videre, ikke at forstaae sensu eminenti (at det absolut ikke kan forstaaes hverken af Kloge eller af Dumme), men brugt æsthetisk om det Forunderlige blandt meget Andet, det Forunderlige, som vel er forunderligt, men som man dog kan begribe. Speculationen maa (forsaavidt den da ikke vil afskaffe al Religieusitet for at introducere os en masse i den rene Værens forjættede Land) conseqvent være af den Mening, at Religieusiteten A er høiere end B, da den er Immanentsens: men hvorfor saa kalde den christelig? Christendommen vil ikke nøies med at være en Evolution indenfor Menneskenaturens Totalbestemmelse, et saadant Engagement er for lidt at byde Guden; den vil heller ikke eengang være den Troende det Paradoxe og saa underhaanden, lidt efter lidt, skaffe ham Forstaaelsen, thi Troens Martyrium (at korsfæste sin Forstand) er ikke et Øiebliks, men netop Vedvarenhedens Martyrium.

[…]

Anmærkning. Forsaavidt det Opbyggelige er et væsentligt Prædikat for al Religieusitet, vil Religieusiteten A ogsaa have sit Opbyggelige. Overalt hvor Guds-Forholdet findes af den Existerende i Subjektivitetens Inderlighed er det Opbyggelige, der tilhører Subjektiviteten, medens man ved at blive objektiv giver Afkald derpaa, der, skjøndt tilhørende Subjektiviteten, dog ikke er det Vilkaarlige saa lidet som Elskov og det at elske, hvilket man dog ogsaa giver Afkald paa ved at blive objektiv. Skyld-Bevidsthedens Totalitet er det meest Opbyggelige i Religieusiteten A. Det Opbyggelige i Religieusiteten A's Sphære er Immanentsens, er den Tilintetgjørelse, i hvilken Individet skaffer sig selv tilside for at finde Gud, da det nemlig er Individet selv der er til Hinder. Det Opbyggelige er altsaa her ganske rigtigt kjendeligt paa det Negative, paa Selvtilintetgjørelsen, der i sig finder Guds-Forholdet, der gjennemlidende synker i Guds-Forholdet, grunder i det, fordi Gud er i Grunden, naar blot alt det der er i Veien er ryddet bort, enhver Endelighed og først og fremmest Individet selv i dets Endelighed, i dets Retshaverie mod Gud. Æsthetisk er Opbyggelsens hellige Hvilested udenfor Individet, det søger Stedet; i den ethisk-religieuse Sphære er Individet selv Stedet, naar Individet har tilintetgjort sig selv.

Dette er det Opbyggelige i Religieusiteten A's Sphære. Passer man ikke paa det, og paa at have denne Bestemmelse af det Opbyggelige imellem: saa forvirres Alt igjen, idet man bestemmer det Paradox-Opbyggelige, som da forvexles med et æsthetisk Forhold ud efter. I Religieusiteten B er det Opbyggelige et Noget uden for Individet, Individet finder ikke Opbyggelsen ved at finde Guds-Forholdet i sig, men forholder sig til Noget uden for sig selv, for at finde Opbyggelsen. Det Paradoxe ligger i, at dette tilsyneladende æsthetiske Forhold, at Individet forholder sig til Noget uden for sig, dog skal være det absolute Guds-Forhold; thi i Immanentsen er Gud hverken et Noget, men Alt, og er uendeligt Alt, eller udenfor Individet, thi Opbyggelsen ligger netop i, at han er i Individet. Det Paradox-Opbyggelige svarer derfor til Bestemmelsen af Gud i Tiden som enkelt Menneske, thi naar saa er, forholder Individet sig til et Noget uden for sig. At dette ikke lader sig tænke, er jo netop det Paradoxe. Et Andet er, om den Enkelte ikke stødes tilbage derfra: det bliver hans Sag. Men fastholdes Paradoxet ikke saaledes, saa er Religieusiteten A høiere, og hele Christendommen at skyde tilbage i æsthetiske Bestemmelser, trods Christendommens Paastand, at det Paradoxe, den taler om, ikke lader sig tænke, forskjelligt altsaa fra et relativt Paradox, der höchstens vanskeligt lader sig tænke. Speculationen maa det indrømmes at holde paa Immanentsen, om det end maa forstaaes anderledes end Hegels rene Tænken, men Speculationen maa ikke kalde sig christelig. Religieusiteten A er derfor af mig aldrig bleven kaldet christelig eller Christendom.

[…]

I Forhold til Bestemmelsen af Individets dialektiske Inderliggjørelse rangere alle Existents-Opfattelser. Jeg skal nu, forudsættende hvad derom er udviklet i dette Skrift, blot recapitulere, mindende om, at Speculationen naturligviis spiller udenfor, da den som objektiv og abstrakt er ligegyldig mod Bestemmelsen af det existerende Subjekt, og i det Høieste har med den rene Menneskehed at gjøre; medens Existents-Meddelelser forstaae noget Andet ved unum, naar det hedder unum noris omnes, forstaaer noget Andet ved »Dig selv« naar det hedder »kjend Dig selv«, forstaaer derved et virkeligt Menneske, og antyder derved at den ikke beskjæftiger sig med de anekdotmæssige Differentser mellem Peer og Povel. – Er Individet udialektisk i sig, og har sin Dialektik udenfor sig: saa har vi de æsthetiske Opfattelser. Er Individet dialektisk ind efter i sig selv, i Selvhævdelse, saaledes at altsaa den sidste Grund ikke bliver dialektisk i sig, da det til Grund liggende Selv bruges til at overvinde og hævde sig selv: saa har vi den ethiske Opfattelse. Er Individet bestemmet dialektisk ind efter i Selvtilintetgjørelse for Gud: saa har vi Religieusiteten A. Er Individet paradox-dialektisk, enhver Rest af oprindelig Immanents tilintetgjort, og al Sammenhæng afskaaret, Individet bestedt i Existentsens Yderste: saa har vi det Paradox-Religieuse. Denne paradoxe Inderlighed er den størst mulige, thi selv den meest dialektiske Bestemmelse, naar den dog er indenfor Immanentsen, har ligesom en Mulighed af Udflugt, af en Springen fra, af en Tilbagetagen ind i det Evige bag ved; det er, som var dog Alt ikke sat ind. Men Bruddet gjør Inderligheden til den størst mulige.

I Forhold til Opfattelsen af det: at existere rangere igjen de forskjellige Existents-Meddelelser. (Speculationen som abstrakt og objektiv seer aldeles bort fra det at existere og Inderlighed; og er, da Christendommen endog paradox accentuerer det at existere, den størst mulige Misforstaaelse af Christendommen). Umiddelbarheden, det Æsthetiske finder ingen Modsigelse i det at existere; at existere er Eet, Modsigelsen noget Andet, som kommer udenfra. Det Ethiske finder Modsigelsen, men indenfor Selvhævdelsen. Det Religieuse A fatter Modsigelsen som Lidelse i Selvtilintetgjørelsen, dog indenfor Immanentsen, men ethisk accentuerende det at existere forhindrer det den Existerende i abstrakt at blive i Immanentsen eller i at blive abstrakt ved at ville blive i Immanentsen. Det Paradox-Religieuse bryder med Immanentsen, og gjør det at existere til den absolute Modsigelse ikke indenfor Immanentsen men mod Immanentsen. Der er intet immanent til Grund liggende Slægtskab mellem det Timelige og det Evige, fordi det Evige selv er kommet i Tiden og der vil constituere Slægtskabet.

\sectionheading{Tillæg: Forstaaelsen med Læseren  ·  SV¹ VII, 536–537}

Undertegnede, Johannes Climacus, der har skrevet denne Bog, udgiver sig ikke for at være en Christen; han er jo fuldt op beskjæftiget med, hvor vanskeligt det maa være at blive det; men endnu mindre er han En, der efter at have været Christen ved at gaae videre er ophørt at være det. Han er en Humorist; nøiet med Øieblikkets Vilkaar, haabende at noget Høiere vil forundes ham, føler han sig særdeles lykkelig ved, naar galt skal være, netop at være født i det spekulative, theocentriske Aarhundrede. Ja, vor Tid det er en Tid for Spekulanter og store Mænd med mageløse Opdagelser; og dog troer jeg, at Ingen af disse Ærede kan have det saa godt som en privatiserende Humorist har det i al Stilhed, hvad enten han afsides slaaer sig for sit Bryst eller leer ret hjerteligen. Han kan derfor meget godt være Forfatter, naar han blot passer paa, at det bliver for hans egen Fornøielse, at han bliver i Afsidesheden, at han ikke kommer ind med i Trængselen, omkommer i Tidens Vigtighed, som en Nysgjerrig ved en Ildebrand, bliver ansat til at pumpe, eller blot generes af den Tanke, at han kunde staae i Veien for nogen af de forskjellige Udmærkede, der har og skal have og maa have og vil have Betydning.

Hele Skriftet dreier sig i Experimentets Afsideshed om mig selv, ene og alene om mig selv. »Jeg, Joh. Cl., nu tredive Aar gammel, barnefødt i Kjøbenhavn, et slet og ret Menneske som Folk er fleest, har hørt tale om, at der er et høieste Gode i Vente, som kaldes en evig Salighed, og at Christendommen vil betinge En dette ved Forholdet til den: jeg spørger nu, hvorledes bliver jeg en Christen« (confer Indl.).

\sectionheading{En første og sidste Forklaring}

For en Form og for Ordens Skyld vedgaaer jeg herved, hvad der realiter neppe kan have Interesse for Nogen at vide, at jeg er, som man siger, Forfatter af: Enten–Eller(Victor Eremita),Kjøbenhavn i Februar 1843; Frygt og Bæven(Johannes de silentio), 1843; Gjentagelsen(Constantin Constantius) 1843; Om Begrebet Angest(Vigilius Hafniensis) 1844; Forord(Nicolaus Notabene) 1844; Philosophiske Smuler(Johannes Climacus) 1844; Stadier paa Livets Vei(Hilarius Bogbinder: William Afham, Assessoren, Frater Taciturnus) 1845; Afsluttende Efterskrift til de philosophiske Smuler (Johannes Climacus) 1846; en Artikel i »Fædrelandet« Nummer 1168, 1843 (Victor Eremita); to Artikler i »Fædrelandet«, Januar 1846 (Frater Taciturnus).

Min Pseudonymitet eller Polyonymitet har ikke havt en tilfældig Grund i min Person (visseligen da ikke i Frygt for Lovenes Straf, i hvilken Henseende jeg ikke er mig bevidst at have forbrudt Noget, og har Bogtrykkeren, samt Censor qua Embedsmand, samtidigen med Skriftets Udgivelse altid været officielt underrettet om, hvo Forfatteren var), men en væsentlig i selve Frembringelsen, der for Replikkens, for den psychologisk varierede Individualitets-Forskjelligheds Skyld digterisk krævede den Hensynsløshed i Godt og Ondt, i Sønderknuselse og Overgivenhed, i Fortvivlelse og Overmod, i Lidelse og Jubel og saa videre, der kun er ideelt begrændset af den psychologiske Conseqvents, hvilken ingen faktisk virkelig Person i Virkelighedens sædelige Begrændsning tør tillade sig eller kan ville tillade sig. Det Skrevne er da vel Mit, men kun forsaavidt jeg har lagt den producerende digterisk-virkelige Individualitet ham hans Livs-Anskuelse i Munden ved Replikkens Hørlighed. Thi mit Forhold er end yderligere end en Digters, der digter Personer og selv dog i Forordet er Forfatteren. Jeg er nemlig upersonligt eller personligt i tredie Person en Souffleur, der digterisk har frembragt Forfattere, hvis Forord atter er deres Frembringelse, ja hvis Navne ere det. Der er saaledes i de pseudonyme Bøger ikke et eneste Ord af mig selv; jeg har ingen Mening om dem uden som Trediemand, ingen Viden om deres Betydning uden som Læser, ikke det fjerneste private Forhold til dem, som dette da er umuligt at have til en dobbelt-reflekteret Meddelelse. Et eneste Ord af mig personligt i mit eget Navn vilde være den anmassende Selvforglemmelse, der forskyldte med dette ene Ord, dialektisk seet, væsentligen at have tilintetgjort Pseudonymerne. Saalidet som jeg i Enten–Eller er Forføreren eller Assessoren, saalidet er jeg Udgiveren Victor Eremita, netop lige saa lidet; han er en digterisk-virkelig subjektiv Tænker, som man jo gjenfinder ham i »in vino veritas«. Jeg er i »Frygt og Bæven« ligesaa lidet Johannes de silentio, som den Troens Ridder, han fremstiller, netop ligesaa lidet, og atter netop ligesaa lidet Forfatter af Forordet til Bogen, hvilket er en digterisk-virkelig subjektiv Tænkers Individualitets-Replik. Jeg er i Lidelses-Historien (Skyldig? – Ikke-Skyldig?) ligesaa lidet Experimentets Quidam som Experimentatoren, netop ligesaa lidet, da Experimentatoren er en digterisk-virkelig subjektiv Tænker og den Experimenterede hans Frembringelse i psychologisk Conseqvents. Jeg er saaledes det Ligegyldige, det vil sige det er ligegyldigt hvad og hvorledes jeg er, netop fordi igjen det Spørgsmaal om det nu ogsaa i mit Inderste saaledes er mig selv ligegyldigt, hvad og hvorledes jeg er, er et denne Frembringelse absolut Uvedkommende. Hvad der derfor ellers i Forhold til mangt et ikke dialektisk-redupliceret Foretagende kan have sin lykkelige Betydning i skjøn Overeensstemmelse med den Udmærkedes Foretagende, det vilde her, i Forhold til en maaskee ikke umærkelig Frembringelses aldeles ligegyldige Pleiefader, ikkun virke forstyrrende. Mit Facsimile, mit Portrait og saa videre vilde ligesom det Spørgsmaal om jeg gik med Hat eller Kaskjet, kun kunne blive Gjenstand for Deres Opmærksomhed, for hvem det Ligegyldige var blevet vigtigt – maaskee til Vederlag for at det Vigtige var blevet dem ligegyldigt. Juridisk og literairt er Ansvaret mit, men dialektisk-let forstaaet er det mig, der har foranlediget Frembringelsens Hørlighed i Virkelighedens Verden, hvilken naturligviis ikke kan indlade sig med digterisk-virkelige Forfattere, og derfor ganske conseqvent med absolut Ret juridisk og literairt holder sig til mig. Juridisk og literairt, thi al digterisk Frembringelse vilde eo ipso været gjort umulig eller meningsløs og utaalelig, hvis Replikken skulde være den Producerendes (ligefremt forstaaet) egne Ord. Mit Ønske, min Bøn er det derfor, at man, hvis det skulde falde Nogen ind at ville citere en enkelt Yttring af Bøgerne, vil gjøre mig den Tjeneste, at citere den respektive pseudonyme Forfatters Navn, ikke mit, det vil sige dele saaledes imellem os, at Yttringen qvindeligt tilhører den Pseudonyme, Ansvaret borgerligt mig. Jeg har fra Begyndelsen meget godt indseet og indseer at min personlige Virkelighed er et Generende, som Pseudonymerne pathetisk-selvraadigen maatte ønske bort jo før jo hellere, eller gjort saa ubetydeligt som muligt, og dog igjen ironisk-opmærksomme maatte ønske at have med som den frastødende Modstand. Thi mit Forhold er Eenheden af at være Secretairen og ironisk nok, dialektisk reduppliceret Forfatter af Forfatteren eller Forfatterne. Medens derfor vistnok Enhver, der overhovedet har bekymret sig om Sligt, hidtil uden videre har anseet mig for Forfatteren af de pseudonyme Bøger førend Forklaringen kom, saa vil Forklaringen maaskee i første Øieblik foranledige den besynderlige Virkning, at jeg, der dog selv bedst maa vide det, er den eneste, der kun meget tvivlsomt og tvetydigt anseer mig selv for Forfatteren, fordi jeg er den uegentlige Forfatter, hvorimod jeg ganske egentligen og ligefremt er Forfatter for Exempel af de opbyggelige Taler og af hvert Ord i disse. Den digtede Forfatter har sin bestemte Livs-Anskuelse, og den Replik, der saaledes forstaaet muligen kunde være betydningsfuld, vittig, vækkende, vilde maaskee i et bestemt faktisk enkelt Menneskes Mund lyde besynderlig, latterlig, modbydelig. Om Nogen paa den Maade, ukjendt med en fjernende Idealitets dannede Omgang, ved misforstaaet Paatrængenhed mod min faktiske Personlighed har forvandsket sig Indtrykket af de pseudonyme Bøger, har narret sig selv, virkeligen har narret sig selv ved at faae min personlige Virkelighed at trækkes med istedetfor en digterisk-virkelig Forfatters dobbelt-reflekterede lette Idealitet at dandse med, paralogistisk-nærgaaende har bedraget sig selv ved meningsløst at faae min private Enkelthed ud af qvalitative Modsætningers eviterende dialektiske Dobbelthed: er sandeligen ikke min Skyld, der netop sømmeligen og i Interesse af Forholdets Reenhed fra min Side har gjort, saa godt jeg kunde, Alt for at forhindre, hvad en nysgjerrig Portion af Læseverdenen, Gud veed i Interesse for hvem, lige fra Begyndelsen har gjort Alt for at opnaae.

Leiligheden synes at indbyde dertil, ja næsten at fordre det endog af den Modstræbende: nu, saa vil jeg benytte den til en aaben og ligefrem Yttring ikke som Forfatter, thi det er jeg jo ikke saaledes, men som Den, der har havt Arbeidet med, at Pseudonymerne kunde blive det. Først vil jeg takke den Styrelse, der paa saa mangfoldig Maade har begunstiget min Stræben, begunstiget den uden maaskee een eneste Dags Afbrydelse fra Anstrængelsen gjennem 4¼ Aar, og forundt mig meget Mere end hvad jeg nogensinde havde ventet, om jeg end tør give mig selv i Sandhed det Vidnesbyrd, af yderste Evne at have sat Livet ind; mere end hvad jeg idetmindste havde ventet, selv om det Præsterede forekommer Andre en vidtløftig Ubetydelighed. Saaledes med inderlig Tak til Styrelsen finder jeg det ikke Forstyrrende, at jeg just ikke kan siges at have udrettet Noget, eller hvad der er ligegyldigere, opnaaet Noget i den udvortes Verden; jeg finder det ironisk i sin Orden, i Frembringelsens og mit tvetydige Forfatterskabs Medfør, at Honoraret idetmindste har været temmelig socratisk. – Dernæst vil jeg, efter sømmeligen forud at have bedet om Undskyldning og Tilgivelse, hvis det skulde synes Nogen upassende, at jeg taler saaledes, hvad han dog maaskee atter selv vilde finde upassende om jeg undlod: jeg vil med en erindrende Taknemlighed mindes min afdøde Fader, det Menneske jeg skylder meest af Alle, ogsaa i Henseende til min Arbeiden. – Fra Pseudonymerne skilles jeg derpaa med Tvivlsomhedens gode Ønske for deres fremtidige Skjebne, at denne, hvis den skal blive dem gunstig, maa blive det, netop som de kunne ønske; jeg kjender dem jo fra den fortrolige Omgang, jeg veed det, mange Læsere kunne de ikke vente eller ønske: gid de da lykkeligen maatte finde de enkelte ønskelige. – Min Læser, hvis jeg tør tale om en saadan, vilde jeg udbede mig et glemsomt Minde hos i Forbigaaende, et Kjende paa, at det er mig han mindes, fordi han mindes mig som Bøgerne uvedkommende, hvad Forholdet kræver, ligesom Paaskjønnelsen derfor bydes oprigtigen her i Afskedens Øieblik, da jeg forøvrigt forbindtligst takker Enhver, der har tiet, med dyb Ærefrygt Firmaet Kts – at det har talet.

Forsaavidt de Pseudonyme paa nogensomhelst Maade skulde have fornærmet nogen agtværdig, eller vel endog nogen af mig selv beundret Mand; forsaavidt de Pseudonyme paa nogensomhelst Maade skulde have forstyrret eller tvetydiggjort noget virkeligt Godt i det Bestaaende: da er Ingen villigere end jeg, der jo bærer Ansvaret for Brugen af den paaholdne Pen, til at gjøre Afbigt. Hvad jeg saadan kjender til Pseudonymerne berettiger mig naturligviis ikke til noget Udsagn, men da heller ikke til nogen Tvivl om deres Samtykke, da deres Betydning (denne vorde nu virkeligen en hvilken) ubetinget ikke ligger i at gjøre noget nyt Forslag, nogen uhørt Opdagelse, eller stifte et nyt Partie og ville gaae videre, men netop i Modsætningen, i ingen Betydning at ville have, i paa en Afstand som er Dobbelt-Reflexionens Fjernhed at ville solo læse de individuelle, humane Existents-Forholds Urskrift, det Gamle, Bekjendte og fra Fædrene Overleverede, igjennem endnu engang om muligt paa en inderligere Maade.

Og gid saa ingen Halvbefaren vil lægge dialektisk Haand paa dette Arbeide, men lade det staae som det nu staaer.

Kjøbenhavn i Februar 1846.

S. Kierkegaard.

\workheading{En omreisende Æsthetikers Virksomhed}{1845}

\worknote{Antologiens ''The Activity of a Traveling Esthetician…''. Frater Taciturnus' polemiske artikel i Fædrelandet (27. december 1845), rettet mod P. L. Møller og hans behandling af ''Stadier paa Livets Vei'' i årbogen Gæa — det stykke der udløste Corsar-striden. Hele artiklen gengives; Frater Taciturnus' to fodnoter (bl.a. den om Experimentet og »Vanviddets Grændse«) er udeladt.}

\sectionheading{En omreisende Æsthetikers Virksomhed, og hvorledes han dog kom til at betale Gjæstebudet (Fædrelandet, 27. dec. 1845)  ·  SV¹ XIII, 423–431}

Skjønt de egenlige Nytaarsgratulanter mere og mere udvide Tiden for deres Virksomhed, som egenligen var indskrænket til Nytaarsdag, er den dog væsenligen begrændset af en 8 Dage. Anderledes med vor driftige og foretagende Literat Hr. P. L. Møller i Qvalitet af Nytaarsgratulant. Han begynder længe iforveien at gaae omkring og gratulere og indsamle milde Gaver til sin brillante Nytaarsgave (Gæa); ja, han reiser endog omkring paa Landet. Faaer han ikke Noget, eller kun Lidet, eller de Navnkundiges fyldige og fuldvægtige Bidrags ringe Omfang udviser, at han mangler noget Fyld i Nytaarsgaven, – Fylde faaer det nu at være med, – saa rykker han Samtalerne ind, som han har ført derude paa Landet. I Grunden er det en ret besparende Maade at reise paa, som aldrig er faldet mig ind, der altid har anseet det at reise for meget bekosteligt, og som maaskee ikke vilde falde mange Andre end Hr. P. L. Møller ind, thi, Herre Gud, man kan da ogsaa drive Oekonomien for vidt. Man reiser, som nu denne Gang Hr. P. L. M. har gjort det (ifølge Gæa 1846) ned til Sorø; man besøger ProfessorHauch, bliver af den udmærkede Digter modtagen med dansk Gjæstfrihed; man tager for sig af Retterne, og medens ellers meget nærige Folk stikke lidt af Fødevarerne til sig, et Stykke af Stegen i Lommen, noget Kage i Hatten, er Hr. P. L. M. graadig nok til at tage hele Samtalen med sig og lade den trykke – saa er det betalt – ja mere end betalt, thi man har jo ingen Udgifter havt med Anretningen, det er bare Indtægt. Havde Trop vidst Besked om denne Maade at reise paa, da skulde han ikke have lidt saa meget af Mangel, thi selv om han var bleven affærdiget i den berømte Mands Forværelse, han kunde jo dog altid tjent lidt ved at lade den berømte Mands Ord trykke. Ingen snyltende Probenreuter kan gjøre saa fordelagtig en Reise, thi han kan da blot tage Bestillingerne med; ja ingen slugen Tiendetager kan gjøre en fordelagtigere – Hr. P. L. M. har den Fordel, at intet Ord bliver spildt, det kommer i Nytaarsgaven.

I Samtalen, som vor omreisende Æsthetiker har havt dernede, er nu ogsaa min Frembringelse bleven Gjenstand for Omtale. Paa den Maade er da jeg ogsaa kommen til at contribuere min Skærv til Nytaarsgaven ved at foranledige ham til nogle Udgydelser efter Bordet. Lad ham have det; i Grunden er mit Bidrag meget uegenligt, thi da Alt, hvad han siger, ikke blot er Confusion, (en Omskrivning af Opgavens Vanskelighed, hvilken Bogen selv langt stærkere har fremhævet, til en Indvending mod Opgavens Løsning), men endog vrimler af factiske Usandheder i Henseende til det mest Afgjørende, kan jeg egenligen ikke sige, at det er min Bog, han taler om, uden forsaavidt han nævner dens Navn, og ved Opfyldelsen erindrer mig om dens prophetiske Motto: »solche Wercke sind Spiegel, wenn ein Affe hineinguckt, kann kein Apostel heraussehen« (Lichtenberg). Er der noget Ubehageligt ved Hr. P. L. Møllers Driftighed, er det mere det Fornærmelige mod en Digter som Professor Hauch og hans private Liv. At Scenen er i Hr. Professor Hauchs Huus og han deltagende i Samtalen, giver naturligvis denne Interessen. Men det synes dog noget nærgaaende, at gjøre sig betalt paa den Maade for – ja hvorfor? for at være bleven modtagen med Gjæstfrihed af en berømt Mand. Heldigvis falder der imidlertid ikke en eneste Yttring af Professor Hauch i det, som angaaer min Frembringelse, hvad der glæder mig lige saa meget for Professorens som for min Skyld. En Yttring af ham, den være nu gunstig eller ugunstig, har altid sin Betydning som enhver legitim Magthavers, den skal og bør ikke svækkes og misliggjøres ved Tvetydighed, at man ikke kan vide, hvilket der er hvilket, fordi Hr. Professor vel maaskee havde sagt det, men ikke har sagt det i »Gæa«, og P. L. Møller vel havde sagt det, men dog ikke sagt det, da han i »Gæa« kun havde sagt, at Hr. Professor Hauch havde sagt det i sin Stue. Hvilket vrængende Misforhold mellem en Beføiets Dom og denne Utilforladelighed!

Nu derimod er Alt i sin Orden; jeg har saa vist Intet at indvende, hverken mod at Hr. P. L. M. virkeligen har sagt det, thi det maa han jo selv bedst vide; eller mod at Yttringen er hans virkelige Mening, hvilken jeg i det mindste ikke er nysgjerrig efter. Naar Talen er om en Præstation i Retning af Religieusitetens Dobbelt-Dialektik i en Gjennemgangskrises Yderste er det, at det er Hr. P. L. M.'s virkelige Mening, det Forsvar, jeg altid skal nøies med. En Tilbagekaldelse af hans Mening vilde ikke kunne betyde mig saa meget som en høitidelig Forsikkring om, at han virkelig har en særdeles ugunstig Mening.

Samtalen, der angaaer min Frembringelse, føres mellem Hr. P. L. M. og, besynderligt nok, »en fornuftig Mand«. Egenligen kan man dog ikke kalde det en Samtale, men snarere tvende Foredrag. Man skulde næsten troe, at Hr. P. L. M. havde forberedt sig paa den Tale, han agtede at holde, thi i en Samtale er det noget stivt at tale 3 Pagina i Træk. Antages dette, bliver hans Reise at betragte som mindre profitabel. Maaskee betragter han det selv som en Opoffrelse, at han er reist ned til Sorø og i Hr. Professor Hauchs Hus »efter Bordet« æsthetisk har concurreret med »en fornuftig Mand«. Maaskee betragter han dog Talen som en Godtgjørelse, hvilket paa en Maade kunde være blevet tydeligere, dersom den var kommen lidt tidligere: at nemlig Hr. P. M. til Taksigelse havde paataget sig at læse fra Bordet. Profit bliver der dog altid derved. Lad os antage, at den ubetydeligste Person havde den Ære at tale med Keiseren af Rusland. Gud veed, det var noget Galimathias han sagde, men dog faaer det en vis Interesse ved at være sagt i Keiserens Palads ved en Audiens. Saaledes ogsaa her med Hr. P. L. M.'s Tale. Den skylder Professor Hauchs Hus Interessen ene og alene. Havde Hr. P. L. M. holdt den samme Tale paa Dagvognen til Postillonen, vilde den slet ingen Interesse have. Interessen ligger i Stedet, det Interessante i Modsigelsen mellem Stedets Betydningsfuldhed og Talens Ubetydelighed; holdt til Postillonen, altsaa paa rette Sted, er Interessen forsvunden og det Interessante. Som Hr. P. L. M.s virkelige Mening er den hans virkelige Mening – ikke videre. »Den fornuftige Mand« batter kun lidt; og besynderligt nok ender ogsaa han sin Tale med de Ord: »saaledes forekommer det i det mindste mig«, hvilket forekommer mig meget besynderligt i »en fornuftig Mands« Mund. Det er nemlig ikke besynderligt, naar for Exempel Hr. Petersen er en fornuftig Mand, at Hr. Petersen siger noget Saadant, men naar man slet ikke er Andet end rent abstract »en fornuftig Mand«, en fornuftig Mand under Gaaseøine, saa svarer dertil vanskeligt: det forekommer mig, da denne Formel antages at gjøre den individuelle Tilfældighed gjeldende.

Sagen er saaledes aldeles ubetydelig, kun ubehagelig for Professor Hauchs Skyld; fordi en udmærket Mand paa en usømmelig Maade bringes i en uværdig Indfatning; men saa har den tillige en mig næsten fordelagtig Side, som jeg skal fremhæve, fordi jeg troer det ikke er mere Dialektisk, end at det med Lethed lader sig læse i et Blad. Naar en Experimentator, der ellers Gudskelov befinder sig meget vel lever i en Tid, da Alle have tvivlet om Alt, overvundet Tvivlen, holdt Reflexionen ud, fundet Mediationen, ladet Religieusiteten bag sig som overvunden formodenligen ved at udelade dens Forfærdelser, saa bliver han opmærksom; han tænker her maa være Ugler i Mosen. Og hvad gjør han saa? Ja selv bevarer han sig sund og frisk, spiser, drikker, rider, kjører, en Ven af Spøg og Munterhed; thi ellers duer han ikke til at experimentere. Ved Siden deraf skriver han et saadant Experiment, som det, jeg nu har skrevet, hvor han strænger Reflexionen lidt an. Han frabeder sig Kritik, og forstaaer dette ganske ligefrem, som han netop har sagt det i Bogen, thi ligefrem maatte det jo betyde, at Kritiken skulde hjælpe ham til at forstaae sit eget Experiment. Derimod kan han psychologisk gjerne ønske Yttring, ikke just for Yttringens Skyld ligefrem, men for den Leilighed, han derved faaer til at gjøre et lille Indblik, om han har gjort Menneskene Uret med ikke at troe, at de virkeligen have holdt Reflexionen ud, thi have de det, da vil en saadan Anstrængelse i Reflexion være dem et Smørrebrød, eller det er, som han formoder. Hr. P. L. Møller tilbyder sig nu med sin Yttring. Men en Experimentator er en rolig og uforfærdet Mand, selv hvor glad han bliver over et muligt forønsket Udfald, seer han først efter, i hvilken Tilstand den befinder sig, der ved sine Yttringer skal forraade Noget, at ikke et rent tilfældigt Ildebefindende svækker Experimentets Betydning. Altsaa. Hr. P. L. Møller; han befindes at være i legemlig sund Tilstand, han er reist med Dagvognen ned til Sorø, forfrisket og styrket af Landluften. Videre. Han skulde dog ikke befinde sig i en forsulten og udhungret Tilstand, der kunde lade en Besvimelse befrygte? Nei, han har faaet godt at spise hos Professor Hauch. Ja, saa er Alt i sin Orden. Dog, han skulde vel ikke have spiist for meget? Nei, han skal jo passe paa Samtalen for at lade den trykke, altsaa maa det denne Gang være i hans Interesse som Gjæst ikke at tage for Meget til sig, eller overlæsse sin Mave. Ja, saa er Alt i sin Orden. Hvad skeer? Blot ved Tanken om Bogen, (han har den ikke ved Haanden, det er nede i Sorø), blot ved Tanken om Bogen, eller maaskee endog ved Tanken om, at have havt Bogen i Haanden, bliver han benauet, han faaer ondt, i en exalteret Tilstand siger han Noget, som han næsten selv befrygter »er for meget sagt.« Ingenlunde, Hr. P. L. Møller. Alt dette er ganske normale Phænomener, det eneste Abnorme er, at Patienten ikke taler i hele Menneskehedens Navn. Videre: her er jo endog Omstændigheder, der maatte hjælpe den Lidende til at holde igjen; han er i Selskab, det er efter Bordet, hvor saa let enhver Anden, end en besindig Experimentator, der med Bestemthed veed og veed Grunden, hvorfor Hr. P. L. M. ikke kan antages at have spist for meget, kunde falde paa at forklare Opstødet paa en anden Maade. Men nei, han faaer ondt. I Experimentatorens Rapport vilde dette lyde saaledes: denne Mand har ikke holdt Reflexionen ud, skulde han ellers give sig Mine af det og squadronere om Mediationen, saa veed jeg Beskeed.

Dersom blot Hr. P. L. Møller havde været en lidt betydeligere Person, og det virkelig havde været sikkert, at han havde læst Bogen: saa vilde jeg ligefrem have taget mig Phænomenet til Indtægt, og argumenteret derfra.

En Forfatter, der er sig sin Opgaves dialektiske Vanskelighed og Anstrængelse bevidst, venter naturligvis kun faa Læsere og ønsker det forsaavidt ogsaa, ikke skrømtet og coquet, i et Forord, men vedstaaer det Samme i sit Indre, og bruger derfor endog sit eget Jeg, ikke just à la Andersen, men snarere lidt socratisk, til drillende at støde fra med. Han nøies glad med faa Læsere, med een, han nøies med mindre, thi han nøies med at være Forfatter, fortryllet ved den Uendelighedens Modsigelse: at nøies med Tænkningens guddommelige Nydelse. Den existentielle Dialektik især i Dobbelt-Reflexionens Form lader sig ikke meddele ligefrem. Om det behagede mig, at erklære hele min Frembringelse for Galimathias, maatte Den, der skal være Læser, formaae aldeles ikke at lade sig forstyrre deraf, men selv see efter ved at reproducere de dialektiske Bevægelser; han maa heller ikke forstyrres, om det behagede hele Verden at erklære det for Galimathias, saa lidet som jeg vilde gjøre det, han maa ved sig selv afgjøre Sagen. En saadan Forfatter kan da forsaavidt ikke have Noget imod, at en snaksom Bordgjæst ved at lade sin Tale trykke bidrager om muligt til at støde Læsere bort. Langtfra ikke, det er tvertimod godt. Derfor har jeg heller ikke repliceret mod ham. Den, hvis Hoved holder den dialektiske Anstrængelse ud, kan sagtens forblive uforstyrret. Betræffende et Skrift, som ved sin dialektiske Charakter væsenligen kun kan have faa Læsere, har en snaksom Fyr altid en Fordel, naar Sagen skal føres for det store Publicum, hvis man i Sandhed kan kalde det en Fordel, der kun er en Fordel i Usandhed. Man lærer meget af Holberg; man lærer saaledes, at Feilen hos Erasmus Montanus egenlig ikke er Mangel paa Evne, men Mangel paa Selvbeherskelse til at forhindre, at han kommer i Samtale med Per Degn, at han selv foranlediger det Skin, at det for Omgivelsen seer ud, som talte disse Tvende sammen om de lærde og curieuse Materier, som de nu kan beskæftige en Erasmus Montanus. See, et saadant Skin skal jeg ikke foranledige, som talte jeg med Hr. P. L. Møller om Experimentet og dettes Dialektik. Nei, det vi To tale om, er om ganske andre Ting, det er om Reisen til Sorø, Postillonen paa Dagvognen, Mad og Drikke, Pakæsler og andre saadanne populaire Æmner, som ikke overskride Hr. P. L. Møllers Fatteevne. Hvad jeg her har skrevet, vil Experimentets virkelige Læser let opdage er af en anden Art, og noget Saadant, at det uden videre kan læses af Enhver, saa jeg ikke fornærmer nogen Læser af et Blad ved at føre ham ind i Undersøgelser, der ikke kan interessere ham, og som ikke kan afhandles i et Blad. En Opfattelse af Hr. P. L. Møllers Reise til Sorø er ikke dialektisk vanskelig; man skal heller ikke grunde dybt; thi det er netop Veien til Misforstaaelsen. Men saa meget som jeg finder Uendelighedens Glæde i Tankens Syssel, og veed, at jeg har fundet den, fordi jeg nøies med den, selv om intet Menneske vilde dele den med mig, saa lidet har jeg opgivet et psychologisk Kjendskab til de virkelige Mennesker. Et saadant virkeligt Menneske er Hr. P. L. Møller. Men paatrængende, som han er, og Mange bekjendt, har jeg troet, at en saadan lille Opfattelse kunde læses ikke uden Interesse i en Avis. Jeg har tillige virkeligen troet at gjøre Een og Anden en Tjeneste dermed, en Tjeneste, som jeg dog ikke forlanger at paaskjønnes, allermindst derved, at Nogen af den Grund vilde læse lidt i mit Skrift eller kjøbe et Exemplar. Thi, som jeg sagde det i min Skrivelse til Læseren (Pagina 309): »for en rund Sum kjøber man ikke Entreen til disse Forestillinger«, saa gjentager jeg det her uden Vrede; thi hvorfor skulde Den være vred, der har fundet sin Uendelighedens Lyst tilfredsstillet, fundet hvad der vil beskæftige ham Dag ud og Dag ind, selv om det behagede Guden at forøge hver Dags Længde med endnu 12 Timer! Den, der er fængslet af hvad der fængsler ham evigt, om han end har saare Meget tilbage at vinde, forstyrres ikke, og uforandret gjentager jeg Afskedens Ord (Pagina 377): »nøiet med det Mindre – haabende, at det Større muligen engang skal forundes mig, er jeg glad ved Tilværelsen, glad ved den lille Verden, der er min Omgivelse.« Da jeg skrev det, vidste jeg jo meget godt, at Hr. P. L. Møller og Corsaren ere til, og dog vidste jeg meget godt, hvad jeg skrev. Saadanne Personer høre ikke til min Omgivelse, og hvor paatrængende og nærgaaende de end ere; det hjælper ikke, det forstyrrer ikke min Glæde over den lille Verden, der er min Omgivelse; tvertimod Paatrængenheden hjælper mig kun til inderligere at glæde mig ved min Omgivelse.

Gid jeg nu blot snart maatte komme i »Corsaren«. Det er virkelig haardt for en stakkels Forfatter at staae saaledes udpeget i dansk Literatur, at han (antaget, at vi Pseudonymer ere Een) er den eneste, som ikke udskjeldes der. Min Foresatte, Hilarius Bogbinder, er bleven smigret i »Corsaren«, dersom jeg ikke husker feil; Victor Eremita har endog maattet opleve den Tort, at blive udødeliggjort – i »Corsaren«! Og dog, jeg har jo allerede været der; thi ubi spiritus, ibi ecclesia: ubiP. L. Møller, ibi »Corsaren«. Vor Landstryger ender derfor sit »Besøg i Sorø« ganske rigtigt med et af disse væmmelige Corsar-Angreb paa fredelige, respectable Mænd, der i hæderlig Ubemærkethed røgte hver sit Kald i Statens Tjeneste; paa udmærkede Mænd, der ved Meget have gjort sig fortjente, ved Intet latterlige; thi Forfattere maae som offenlige Personer finde sig i Meget, og deriblandt ogsaa i denne Paadutten af Slægtskab med Folk, der, da de lade Noget trykke, jo ogsaa ere Forfattere.

Frater Taciturnus,

Høvidsmand for 3die Afdeling af »Stadier paa Livets Vei«.

\workheading{En literair Anmeldelse}{1846}

\worknote{Antologiens ''Two Ages: The Age of Revolution and The Present Age''. Kierkegaards anmeldelse af Thomasine Gyllembourgs roman ''To Tidsaldre''. Hong-udvalget bringer det berømte kapitel ''Nutiden'' — kulturkritikken af den reflekterende, lidenskabsløse, nivellerende tidsalder, Publikum, Snakken og de Ukjendelige — fra kapitlets begyndelse til overgangen tilbage til romanen (''Dog jeg afbryder''). Gengivet efter SKS' tekst i sammenhæng.}

\sectionheading{Nutiden  ·  SV¹ VIII, 64–100}

Opgaven er atter her, efter hvad jeg har forstaaet, i en almindeligere Betragtning, der er i Novellens critiske Tjeneste, at avancere netop de Momenter, som Forfatteren med novellistisk Kunst har gjengivet.

Nutiden er væsentligen den forstandige, den reflecterende, den lidenskabsløse, den flygtigt i Begeistring opblussende og kløgtigt i Indolents udhvilende.

Dersom man, ligesom man har det i Forhold til Forbrugen af Brændeviin og saa videre, havde Tabeller over Forbrugen af Forstand fra Generation til Generation: saa vilde man forbauses ved at see, hvilken uhyre Qvantitet der nuomstunder forbruges, hvilket Qvantum af Betænkninger og Overveielser og Hensyn selv en privatiserende lille Familie bruger, om den dog har sit rigelige Udkomme, hvilket Qvantum endog Børn og Ungdommen bruger, thi som Børne-Korstoget ligner Middelalderen, saa ligner Børne-Klogskaben Nutiden. Mon der er et Menneske mere, der blot eengang gjør en gevaltig Dummestreg? Neppe en Selvmorder nuomstunder gjør det af med sig selv i Desperation, men han overveier dette Skridt saa længe og saa forstandigt, at han qvæles af Forstandighed, saa det endog kunde gjøres tvivlsomt, om han virkelig maatte kaldes en Selvmorder, forsaavidt det dog især var Overveielsen, der tog Livet af ham. En Selvmorder med Overlæg var han ikke, snarere en Selvmorder ved Overlæg. Det vilde derfor blive den vanskeligste af alle Opgaver, at skulle være Aktor mod en saadan Tid, thi hele Generationen er prokurator-dygtig, og dens Konst, dens Forstandighed, dens Virtuositet bestaaer i aldrig handlende at ville lade Sagen gaae til Doms og til Afgjørelse.

Maa man sige om Revolutions-Tiden, at den farer vild, saa maa man sige om Nutiden, at den farer ilde. Individet og Generationen kommer bestandigt sig selv og hinanden paa tvært; og derfor vilde det være umuligt for en Aktor at faae noget Faktum constateret, fordi der intet er. Af Indicier i Overflod maatte man slutte, at der enten var skeet eller snarest muligt vilde skee noget Overordentligt. Og dog vilde man netop slutte feil, thi Indicierne er just Nutidens eneste Forsøg i Kraft-Præstation, og dens Opfindsomhed og Virtuositet i at entreprenere fortryllende Blendværk, den opblussende Begeistrings Overilelse ved Hjælp af en projekteret Form-Forandrings skuffende Gjenvei er i Retning af Klogskab og negativ Anvendelse af Kraft lige saa høit at anslaae som Revolutionens Præstation i Retning af energisk og omskabende Lidenskab. Trættet af den chimairiske Anstrengelse udhviler saa Tidsalderen sig momentviis i fuldkommen Indolents. Dens Tilstand er som den henimod Morgenen Indslumredes: store Drømme, saa Dvaskhed, saa et vittigt eller kløgtigt Indfald til Undskyldning for at man bliver liggende.

Det enkelte Individ (selv hvor velmenende Mange af disse kunne være, selv hvor megen Kraft de muligen kunne have, hvis de kunde komme til at bruge den) har ikke sluttet Lidenskab nok i sig selv til at rive sig ud af Reflexionens Garn og Reflexionens forføreriske Uvisse; og Omgivelsen, Samtiden har ikke Begivenhed eller forenet Lidenskab, men danner i negativ Forening en Reflexions-Modstand, der først et Øieblik gøgler med skuffende Udsigt, og saa bedrager ved at styrke med glimrende Udflugt: at man dog har gjort det Klogeste ved at lade være at handle. Vis inertiæ er det Tilgrundliggende for Tidens Tergiversation, og enhver Lidenskabsløs gratulerer sig selv som første Opfinder – og bliver endnu kløgtigere. Blev der i Revolutions-Tiden frit udleveret Vaaben, blev der i Korstogenes Tid offentlig uddeelt Bedriftens Insignie: saa beværtes der i Nutiden overalt frit med Klogskabsregler, Hensyns-Beregninger og saa videre Dersom det kunde antages, at en heel Generation havde den diplomatiske Opgave at holde Tiden hen, saa det bestandigt blev forhindret at der skete Noget, og dog bestandigt saae ud som skete der Noget: saa kan man ikke frakjende vor Tid at den præsterer noget lige saa Vidunderligt som Revolutions-Tiden. Hvis En vilde anstille det Experiment med sig selv, at glemme Alt hvad han vidste om Tiden og om den ved Vanen opskruede faktiske Relativitet, og nu komme ligesom fra en ganske anden Verden, hvis han saa vilde læse en eller anden Bog, en Artikel i et Blad, eller vel endog blot tale med en Forbigaaende: han vilde faae det Indtryk: Død og Pine endnu i Aften maa der skee Noget – eller der maa være skeet Noget i forgaars Aftes.

I Modsætning til Revolutions-Tiden som handlende er Nutiden Avertissements-Tiden, de blandede Bekjendtgjørelsers Tid: der skeer ikke Noget, derimod skeer der strax Bekjendtgjørelse. Et Oprør vilde i Nutiden være det Utænkeligste af Alt; en saadan Kraft-Yttring vilde forekomme Tidens beregnende Forstandighed latterlig. Derimod vilde maaskee en politisk Virtuos være istand til en ganske anderledes forbausende Konst-Præstation. Han vilde være istand til at skrive en Indbydelse, som foreslog en General-Forsamlings Afholdelse for at beslutte en Revolution, saa forsigtigt, at selv Censor maatte lade den passere; og derpaa vilde han om Aftenen være istand til at frembringe paa Forsamlingen et saa skuffende Indtryk, at det var den, som havde de gjort Oprør; hvorpaa de ganske roligen vilde skilles ad – efterat have passeret en høist behagelig Aften. En enorm grundig Lærdoms Erhvervelse vilde næsten være utænkelig blandt Yngre i vor Tid, man vilde finde det latterligt. Derimod vilde en videnskabelig Virtuos være istand til at præstere et ganske anderledes Konststykke. Han vilde i en Subscriptions-Plan være istand til at henkaste nogle Lineamenter til et altomfattende System, og gjøre det saaledes, at det paa Læseren (af Subscriptions-Planen) frembragte det Indtryk, at han allerede havde læst Systemet. Thi Encyclopædisternes Tid er forbi, deres nemlig som med Jern-Flid skrev Folianter, nu er Touren kommen til de let bevæbnede Encyclopædister, der en passant gjør det af med hele Tilværelsen og alle Videnskaber. En dyb religieus Forsagelse af Verden og hvad Verdens er, fastholdt i daglig Fornegtelse, vilde være utænkelig blandt Yngre i vor Tid: derimod vilde hver anden theologisk Candidat have Virtuositet nok til at gjøre noget langt Vidunderligere. Han vilde være istand til at projektere et Selskabs Stiftelse, hvis Hensigt ikke var mindre end at frelse alle Fortabte. De store og gode Handlingers Tid er forbi, Nutiden er Anticipationernes. Ingen vil nøies med at gjøre noget Bestemt, Enhver vil smigres af Reflexionen med den Indbildning, at han idetmindste maa opdage en ny Verdens-Deel. Vor Tid er Anticipationens, selv Anerkjendelsen optages forud. Som et ungt Menneske, der beslutter sig til fra 1ste September af ret for Alvor at læse til sin Embedsexamen, for at styrke sig hertil beslutter midlertidigt at feriere i August Maaned: saaledes synes, hvad der dog er betydeligt vanskeligere at forstaae, den nuværende Generation at have fattet den alvorlige Beslutning, at den næste Generation skal for Alvor tage fat paa Arbeidet, og for at der slet ikke skal være noget Forstyrrende eller Sinkende for den, saa tager den nærværende sig – af Festmaaltiderne. Kun er Forskjellen den, at det unge Menneske ungdommeligt forstaaer sig selv i Letsindigheden, Nutiden bliver endog alvorlig – ved Festmaaltiderne.

Handling og Afgjørelse er der i Nutiden lige saa lidt, som der er Svømningens farefulde Lyst for dem, der svømme paa Lavden. Men som den Voxne, der selv forlystes ved at tumle sig i Bølgerne, raaber til den Yngre: kom herud, spring kun rask til – saaledes ligger Afgjørelsen ligesom i Tilværelsen (skjøndt den naturligviis er i Individet) og tilraaber den Yngre, der endnu ikke er afmattet ved Reflexionens Formeget og overlæsset med Reflexionens Indbildninger: kom herud, spring freidigt til; selv om det endog var et letsindigt Spring, naar det blot er afgjørende – dersom Du er dygtig til at være Mand, saa skal Faren og Tilværelsens strenge Dom over Din Letsindighed hjælpe Dig til at blive det.

Dersom Klenodiet, som er Alles Lyst, laae langt ude paa den ganske tynde Iisskorpe, saa Livsfaren altsaa holdt Vagt ved det, vaagende over, at det blev et livsfarligt Vovestykke at gaae saa langt ud, thi (lad os antage dette Besynderlige, som dog kun er besynderligt i Billedet) noget nærmere inde var Isen ganske sikker og bundfrossen: saa vilde i en lidenskabelig Tid Mængden tiljuble den Modige Bifald, naar han vovede sig derud; den vilde gyse for ham og med ham i Afgjørelsens Livsfare; den vilde sørge over ham i Undergangen; den vilde forgude ham hvis han vandt Klenodiet. I en lidenskabsløs reflekteret Tid vilde Sagen blive en ganske anden. Man vilde forstandigen i gjensidig Anerkjendelse af gjensidig Klogskab Alle blive enige om, at det nok ikke var Umagen værd at vove sig saa langt ud, ja at det vilde være uforstandigt og latterligt; og derpaa vilde man forvandle Begeistringens Vovestykke til en Konst-Præstation – for dog at gjøre Noget, »thi Noget maa der gjøres.« Man vilde da tage derud, man vilde paa det Trygge med Kjendermine vurdere de kunstfærdige Skøiteløbere, der kunde løbe næstendeels lige til den yderste Grændse (det vil sige saa langt at Isen endnu var sikker, og Faren endnu ikke begyndt) og saa svinge af. Blandt Skøiteløberne vilde der nu være en og anden særdeles Perfectioneret; han vilde endog kunne gjøre det Konststykke, lige ved Grændsens Yderste at tage endnu et skuffende farevarslende Tilløb, saa Tilskuerne raabte: »I Guder, han er gal, han sætter Livet til.« Men see, han var saa udmærket perfectioneret, at han ganske rigtigt kunde svinge af ved Grændsens Aller-Yderste det vil sige hvor Isen endnu var ganske sikker og Livsfaren endnu ikke begyndt. Aldeles som i et Theater vilde Mængden raabe Bravo og hilse med Acclamation, drage hjem med den store Helte-Konstner i deres Midte, og hædre ham med et velsmagende Festmaaltid. Forstandigheden havde i den Grad taget Overhaand, at den havde forvandlet selve Opgaven til en uvirkelig Konst-Opgave, og Virkeligheden til et Theater. Ved Festmaaltidet om Aftenen vilde Beundringen lyde høit. Og medens ellers Beundringens sande Forhold ere disse, at den Beundrende opløftes ved Tanken om at være Menneske, ligesom den Udmærkede, ydmyges ved Forestillingen om ikke selv at have kunnet udføre det Store, ethisk opmuntres ved Forbilledet til efter Evne at følge den Udmærkedes Exempel: saa vilde Forstandigheden atter have forvandlet Beundringens Forhold. Den ved Festmaaltidet beundrende Klinkende vilde endog i Fanfarens og de 9 Hurras overstadige Øieblik have en Forstandighedens kløgtige Forestilling om, at det dog nok egentligen ikke var synderligt bevendt med den Beundredes Bedrift, at det i Grunden var tilfældigt, at Selskabet gjordes for ham, da enhver af Deeltagerne ved nogen Øvelse i skuffende Vendinger omtrent vilde kunne præstere det Samme. Kort, istedenfor at hente Styrkelse i Skjønsomhed og Opmuntring til det Gode fra Beundringens Festlighed, vilde den Klinkende snarere gaae hjem med en yderligere Disposition til den farligste af alle Sygdomme, men ogsaa den fornemste, den paa Commerce ligesom at beundre hvad man selv anseer for ubetydeligt, fordi det Hele var bleven en dramatisk Spøg, og Beundringens stimulerede Kling-Klang i hemmelig Forstaaelse med, at man næstendeels lige saa godt kunde beundre sig selv.

Eller dersom en Mand endeligen stillede sig i Spidsen for et begeistret Foretagende, og der da, hvad let nok opnaaes, (thi opblussende Begeistring og kløgtig Apathie svare altid til hinanden) løb en Deel Mennesker sammen om ham, dersom han nu i Spidsen for denne Skare drog ud – under Seiers-Sange, indtil han nærmede sig Afgjørelsen og Faren – naar han da vendte sig om for at tale et begeistrende Ord til Skaren: da vilde hele Scenen være forandret. Deeltagerne havde kløgtigt forvandlet sig selv til en Flok Tilskuere, der nu med Kløgtighedens store Selvtilfredshed gav sig Mine af, at det var dem, der listigen og ironisk havde faaet ham narret ud til at være begeistret, og nu vare gaaet med for at see paa ham og lee ad ham. Og denne ualmindelige Kløgtighed vilde i gjensidig Anerkjendelse tilfredsstille dem uendeligt i Sammenligning med nogen Bedrift; det var i Kløgtighedens Forstand brillant derfra. Der vilde ikke høres et eneste Ord om Ustadighed, Feighed, nei man vilde bryste sig ved Kløgtighedens glimrende Indbildning, og man vilde atter have gjort sig sin Helbredelse vanskeligere. Saa vilde Anføreren maaskee ogsaa tabe Modet, og det Hele blive saa fordærveligt derfra som muligt, fordi det blev en fingeret Bevægelse og et Incitament for feig Indbildskhed.

At en Mand staaer eller falder paa sin Gjerning, gaaer af Brug, derimod sidde Alle og slaae sig brillant igjennem ved Hjælp af nogen Reflexion, og saa ved Hjælp af: at de meget godt veed, hvad der skal gjøres. Men see, hvad To og To i Samtale, hvad de Enkelte som Læsere, som Demonstrerende paa en Generalforsamling ypperligt forstaae i Reflexionens og Betragtningens Form: det vilde de aldeles ikke kunne forstaae i Handlingens Form. Hvis En gik omkring og hørte hvad det var man sagde der skulde gjøres, og nu, i Ironiens Interesse, mir nichts und Dir nichts gjorde Noget deraf: saa vilde Alle studse, de vilde finde det overilet; og saasnart de betragtende samtalede, vilde de forstaae, at det var det, der skulde gjøres.

Nutiden, i dens opblussende Begeistring, og saa igjen i den apathiske Indolents, der i det Høieste gider spadse, ligger det Comiske meget nær; men Den, der forstaaer det Comiske, seer let, at det Comiske ligger paa et ganske andet Sted, end Nutiden indbilder sig, og at Satiren i vor Tid, hvis det skal være muligt at den kan gjøre lidt Gavn og den ikke skal afstedkomme ubodelig Skade, maa have til Borgen en consequent og vel begrundet ethisk Livs-Anskuelse, en opoffrende Uegennyttighed, en høibaaren Fornemhed, der forsager Øieblikket; ellers bliver Medicinen uden Sammenligning uendeligt langt farligere end Sygdommen. Det Comiske ligger netop i, at en saadan Tid endnu vil være vittig og slaae stort paa i det Comiske; thi dette er ganske rigtigt den sidste og meest gøglende Udflugt. Hvad har vel en udreflekteret Tid at trodse paa i Henseende til det Comiske? Som lidenskabsløs har den intet Følelsens Valuta i det Erotiske, intet Begeistringens og Inderlighedens Valuta i det Politiske og Religieuse, intet Huslighedens, Pietetens, Beundringens Valuta i det Daglige og Omgangs-Livet. Men den Vittighed, der intet Valuta eier, den spotter Tilværelsen over, selv om Hoben leer skingrende. At ville være vittig, naar man ikke eier Inderlighedens Rigdom, er at ville ødsle paa Luxus og undvære Livs-Fornødenhederne, er, som Ordsproget siger, at sælge sine Buxer og kjøbe en Paryk. Men en lidenskabsløs Tid eier intet Valuta; Alt bliver en Omsætning i Repræsentativer. Det er saadan visse Talemaader og Betragtninger, tildeels sande og forstandige, men afsjelede, der cursere mellem Folk; men ingen Helt, ingen Elsker, ingen Tænker, ingen Troens Ridder, ingen Høimodig, ingen Fortvivlet indestaaer ved primitivt at have oplevet Det for Fuldgyldigheden. Og som man fra Seddel-Pengenes Hvisken i Omsætningen mellem Mand og Mand kan længes efter at høre Myntens Klangfuldhed: saaledes kan man i Nutiden længes efter lidt Primitivitet. Men hvad er mere primitivt end Vittighed, mere primitivt, idetmindste mere overraskende selv end Foraarets første Knop, og det første Grønts skøre Stengel. Ja selv om Foraaret kom efter Aftale, det vilde dog være Foraaret, men en Vittighed efter Aftale vilde være en Modbydelighed. Sæt saa, det til Afløsning fra en opblussende Begeistrings Febricitet kom saa vidt, at Vittighed, hiin guddommelige Accidentse – naar den kommer, hiin Tilgift paa Gudens Vink fra det Uforklarliges gaadefulde Oprindelse, saa end ikke den Vittigste, der har levet, tør sige: imorgen, men tilbedende siger: hvis det behager Guden – sæt at Vittigheden blev forvandlet til sin meest luvslidte Modsætning, en triviel Livs-Fornødenhed, saa det blev en profitabel Marchandiser-Næringsvei, at fabrikere og lave og omsye og opkjøbe gamle og nye Vittigheder: hvilket frygteligt Epigram over den vittige Tid!

Saa bliver Penge tilsidst det Attraaede, det er jo ogsaa et Repræsentativ og en Abstraktion. Selv en Yngling i Nutiden vilde neppe misunde En hans Evner, eller hans Konst, eller den skjønne Piges Elskov, eller hans Berømthed, nei, men hans Penge vilde han misunde ham. Giv mig dem, vil Ynglingen sige, saa er jeg hjulpen. Og denne Yngling, han vilde ikke fare vild i Letsindighed, han vil ikke forskylde, hvad Angeren gjengjelder, han vil ikke have Noget at bebreide sig, men han vil døe i den Indbildning, at havde han havt Penge, da skulde han have levet, da skulde han maaskee ogsaa have drevet det til noget Stort. Lad os tænke paa Novellen. En forelsket ung Mand, Ferdinand Bergland, er forelsket, men hans Forstand og hans Reflexion kommer ham paa tvært – og Afgjørelsen er negativ. End ikke Elskovens Umiddelbarhed i Nutiden er sorgløs som Markens Lilie, og i de Elskendes Øine herligere end Salomons Pragt. En erotisk Critik og en mismodig Forstandighed bøier Sorgløsheden og forfalsker Værdien – og Religieusitetens himmelske Herlighed kan ikke komme til at hjælpe – til det Høieste. En Elsker forlader den Forlovede af Frygt for Næringssorger. Hvor anderledes i Revolutions-Tiden, hvor Lusard letsindigen næsten sorgløs forlader den vanærede Pige, overladende hende alle Sorgerne; thi at Claudine realiserer Opgaven, lever saa godt som af Intet, ikke tænker paa Næringssorgerne, men kun paa Lusard: det bliver jo hendes Fortjeneste. Men Afgjørelse kom det dog til; og Afgjørelsens Trang er det netop Reflexionen fordriver eller vil fordrive, og under Følgerne deraf lider Individet i den dyspeptiske abnorme Forstandighed. Forgjeves fanger Afgjørelsen i Livet efter Individet, forgjeves venter Velsignelsen paa Afgjørelsens Øieblik: man veed klogt og dog bedragen at unddrage sig; og varer dette længe ved, og de da endeligen fanges, da ere de ligesom Piger, der for længe have været forlovede, hvilke sjeldent ere ret skikkede til Ægteskabet.

Efter disse paapegende Bemærkninger er det vel i sin Orden, som det er i Novellens interesseløse Tjeneste, fra Sammenligningen med Revolutions-Tiden at føre Nutiden tilbage paa dialektiske Kategorie-Bestemmelser og paa disses Consequentser, disse være nu netop i det givne Øieblik ganske faktiske eller ikke. Forøvrigt er der, saa lidet som i den digterisk frembringende, og at jeg saa skal sige foreskrivende Novelle, ligesaa lidet i den efterskrivende og som Tjenende sig underordnende Critik Spørgsmaal eller Tale om, hvilken Tidsalder der er den bedste, den betydningsfuldeste. Der er kun Spørgsmaal om Tidsalderens Hvorledes, og dette Hvorledes vindes ved en universellere Anskuelse, hvis Slutnings-Consequentser naaes ved en Slutning ab posse ad esse, og verificeres ved en iagttagende Erfarings ab esse ad posse. Hvad Betydningsfuldheden angaaer, da er det jo muligt, at den Nutiden foresatte Reflexions-Opgave vil forklare sig i en høiere Existents-Form; og hvad Godheden angaaer, da er det jo vist, at den i Reflexionen Hildede kan være lige saa velmenende, som den lidenskabeligt Snarraadige, ligesom omvendt, at der kan være lige saa Meget, der taler til Undskyldning for Den, der i Lidenskaben farer vild, som for Den, der selv har en underfundig Viden om, at han lader sig bedrage af sin Reflexion, medens hans Feil aldrig bliver vitterlig. Ogsaa dette er Reflexionens farefulde Mislighed, at man ikke kan see, om det er en gjennem Overveielse vunden Beslutning, der frelser fra det Onde, eller det er Overveielsens Udmattelse, der svækkende forhindrer fra det Onde. Men vist er det, at som al forøget Kundskab forøger Græmmelsen, saa forøger Reflexionen den ogsaa; og fremfor Alt er det vist, at som for det enkelte Individ saa for en heel Generation, er ingen Opgave og Anstrængelse vanskeligere, end at arbeide sig ud af Reflexionens Fristelser, netop fordi disse ere saa dialektiske, fordi en eneste kløgtig Opfindelse er istand til pludseligen at give Sagen en ny Vending, fordi Reflexionen er istand til hvert Øieblik at omforklare og lade En slippe nogenlunde, fordi det selv i den sidste Reflexions-Afgjørelses Øieblik er muligt at gjøre Alt om – efter at man altsaa har udholdt langt flere Anstrængelser end nogen Snarraadig behøver til at være midt i det. Men alt dette er dog kun igjen Reflexions-Undskyldning, og Stillingen i Reflexionen uforandret, fordi den kun er forandret i Reflexion. Selv det, at Nutiden ved at sammenstilles med en afsluttet Tidsalder paa en Maade skeer Uret, forsaavidt Nutiden jo netop ligger i Vordens Vanskelighed, er dog kun en Reflexions-Bestemmelse, thi derfor har den jo ogsaa Haabets Uvisse.

En lidenskabelig tumultuarisk Tid vil kaste Alt overende, omstøde Alt; en revolutionair men lidenskabsløs og reflekterende forvandler Kraftyttringen til et dialektisk Konststykke: at lade Alt bestaae, men underfundigt fraliste det Betydning; istedenfor i et Oprør culminerer den i at afmatte Forholdenes inderlige Virkelighed i en Reflexions Spænding, der dog lader Alt bestaae og har forvandlet hele Tilværelsen til en Tvetydighed, der i sin Facticitet er, medens dialektisk Sviig privatissime underskyder en hemmelig Læsemaade – at den ikke er.

Sædelighed er Charakteer, Charakteer er det Indgravede (χαρασσω), men Havet har ingen Charakteer og Sandet heller ikke og abstrakt Forstandighed heller ikke, thi Charakteren er netop Inderligheden. Usædelighed er som Energie ogsaa Charakteer. Tvetydighed er det derimod, naar man hverken er det Ene eller det Andet; og Tvetydighed i Tilværelsen er det, naar Qvaliteternes qvalitative Disjunktion svækkes ved en gnavende Reflexion. Lidenskabens Oprør er elementarisk, Tvetydighedens Opløsning er en stilfærdig men Dag og Nat travl Sorites. Adskillelsen mellem Godt og Ondt enerveres ved et letfærdigt, fornemt, theoretisk Kjendskab til det Onde, ved en hoffærdig Klogskab, der veed, at det Gode ikke paaskjønnes eller lønner sig i Verden – saa det næsten bliver Dumhed. Ingen henrives af det Gode i stor Bedrift, Ingen overiles af det Onde i himmelraabende Synd, den Ene skal forsaavidt ikke lade den Anden høre Noget, og dog bliver der af den Grund netop desto mere at snakke om, thi Tvetydigheden er et pirrende Incitament, og ganske anderledes ordriig end Glæden over det Gode og Afskyen for det Onde.

Livs-Forholdenes Springfjædre, der kun i den qvalitativt adskillende Lidenskab ere hvad de ere, taber Elasticiteten; det Forskjelliges Fjernhed fra sit Forskjellige i Qvalitets-Udtrykket er ikke Loven for Inderlighedens Forhold til hinanden i Forholdet. Inderligheden mangler, og Forholdet er forsaavidt ikke til, eller Forholdet er en dvask Cohæsion. Den negative Lov er nemlig: ikke at kunne undvære hinanden og ikke at kunne holde sammen, den positive: at kunne undvære hinanden og at kunne holde sammen, eller dog positivt: ikke at kunne undvære hinanden paa Grund af Sammenholdet. Istedenfor Inderlighedens Forhold indtræder et andet Forhold: det Forskjellige forholder sig ikke til sit Forskjellige, men de staae ligesom og passe paa hinanden med Øinene, og denne Spænding er egentligen Forholdets Ophøren. Det er ikke Beundringen, der glad og freidig, prompt i Paaskjønnelsens Udtryk tager Hatten af for Udmærkelsen, og nu oprøres ved dennes Stolthed og Anmasselse, ei heller er Forholdet det omvendte, ingenlunde, Beundring og Udmærkelse blive næstendeels et Par galante Ligemænd, der passe paa hinanden med Øinene. Det er ikke Borgeren, der med undersaatlig Hylding frimodigt ærer Kongen og nu forbittres over hans Herskesyge, ingenlunde, det at være Borger bliver noget Andet, bliver det at være Trediemand; Borgeren forholder sig ikke i Forholdet, men han er Tilskuer ved Udregningen af den Opgave: Forholdet mellem en Konge og Undersaat; thi det gaaer en Tidlang med at nedsætte Comiteer og Comiteer, saa længe der dog endnu bestandigt er en Mængde, der med fuld Lidenskab vil, hver især, være dette Bestemte, han skal være; men tilsidst ender det med, at hele Tidsalderen bliver en Comitee. Det er ikke Faderen, der i Harme samler sin faderlige Myndighed i en eneste Forbandelse, ikke Sønnen, der trodser, en Adskillelse der dog maaskee kunde ende i Forsoningens Inderlighed, nei Forholdet er saadan upaaklageligt, thi det er snarere ifærd med at ophøre, fordi de ikke væsentligen forholde sig til hinanden i Forholdet, men Forholdet er blevet et Problem, hvor Parterne ligesom i et Spil passe paa hinanden, istedenfor at forholde sig til hinanden, tæller, som man siger, hinanden Forholds-Yttringen i Munden, istedenfor Forholdets resolute Hengivelse; thi det gaaer en Tidlang med at Flere og Flere maae renoncere paa det stillere Livs beskedne, men dog saa fyldige og gudvelbehagelige Opgaver, for at realisere noget Høiere, for i et høiere Forhold at tænke over Forholdene, men tilsidst bliver hele Generationen en Repræsentation – der repræsenterer... ja det er ikke godt at sige hvem; der tænker over Forholdene... ja det er ikke godt at sige for hvis Skyld. Det er ikke en opsætsig Ungdom, der dog skjelver og bæver for Skolemesteren, nei snarere er Forholdet en vis Ligelighed i gjensidig Udvexling mellem Lærer og Discipel om, hvorledes en god Skole bør indrettes. Det at gaae i Skole betyder ikke at skjælve og bæve, betyder ei heller ene og alene at lære, men betyder tillige næstendeels at være interesseret i det Problem om Skole-Underviisningen. Adskillelsens Forhold mellem Mand og Qvinde brydes ikke i formastelig Tøilesløshed, ingenlunde, Sømmeligheden iagttages saaledes, at der bestandigt om den leflende »uskyldige« Grændsestrids enkelte Yttring maa siges, at den er en Ubetydelighed. Hvad maatte man nu kalde et saadant Forhold? Jeg tænker en Spænding, men vel at mærke ikke en Spænding, der strammer Kraften til en Katastrophe, men en Spænding, hvori Tilværelsen afmattes; Fyrigheden, Begeistringen, Inderligheden tabes, der gjør Afhængighedens Lænke og Herredømmets Krone let, der gjør Barnets Lydighed og Faderens Myndighed glad, der gjør Beundringens Underkastelse og Udmærkelsens Ophøiethed frimodig, der giver Læreren enegyldig Betydning og Discipelen saaledes Leilighed til at lære, der ener Qvindens Svaghed og Mandens Kraft i Hengivenhedens lige Styrke. Forholdet bestaaer vel, men det mangler Spændkraft til at samle sig i Inderlighed for at forenes i Samdrægtighed. Forholdene yttre sig vel som tilværende, og dog som fraværende, ikke fyldigt, snarere i en vis slæbende halvvaagen Uafbrudthed. Man tillade mig med et ganske simpelt Billede at oplyse hvad jeg mener. Jeg kom engang i en Familie, der havde et Stueuhr, hvis Værk paa en eller anden Maade var kommen i Uorden. Men Uordenen yttrede sig ikke ved, at Spiralen paa eengang snurrede ud, og Kjeden sprang, ei heller ved at det ophørte at slaae; tvertimod det vedblev at slaae, men paa en besynderlig abstrakt-normal og dog forvirrende Maade. Det slog ikke tolv Slag naar Klokken var 12, og saa eet Slag, naar Klokken var 1, men det slog eet Slag ad Gangen med et bestemt Mellemrum. Paa den Maade slog det hele Dagen igjennem, og angav dog aldrig Klokkeslettet. Og saaledes er det i en afmattende Spænding: Forholdene bestaae; med en abstrakt Uafbrudthed, som forhindrer Bruddet, yttrer der sig Noget, som man maa kalde Forholdets Yttringer, og dog angives Forholdene ikke blot ikke præcise men næsten meningsløse. Det Dyssende er Forholdenes Bestaaen, deres Facticitet; det Farlige er, at netop dette begunstiger Reflexionens underfundige Gnaven. Thi mod Oprør kan man bruge Magt, for vitterligt Falskmaal er Straffen i Vente, men dialektisk Hemmelighedsfuldhed er vanskelig at udrydde; der skal allerede et forholdsviis finere Øre til for at spore Reflexionens lydløse Listen ad Tvetydighedens Snigveie.

Det Bestaaende bestaaer, men Reflexionens Lidenskabsløshed finder sin Beroligelse i, at det er en Tvetydighed. Man vil ikke have Konge-Magten afskaffet, ingenlunde, men hvis man lidt efter lidt kunde faae den forvandlet til en Indbildning, saa vilde man med Fornøielse raabe Hurra for Kongen. Man vil ikke have Udmærkelsen styrtet, ingenlunde, men hvis man paa samme Tid kunde avancere en Viden om, at den er en Indbildning – saa vilde man beundre. Man vil lade hele den christelige Terminologie bestaae, men underhaanden være vidende om, at der ikke skal tænkes noget Afgjørende derved. Og man vil være angerløs, thi man nedriver jo Intet. Man vilde ligesaa lidet have en vældig Konge som en energisk Friheds-Helt, eller en religieus Befuldmægtiget, nei man vil aldeles uskyldigt lade det Bestaaende bestaae, men i en reflekteret Viden være mere eller mindre vidende om dets Ikke-Bestaaen. Og saa vilde man være stolt i den Indbildning, at dette er Ironie, ret som var den sande Ironiker ikke netop den skjult Begeistrede i en negativ Tid (ligesom Helten er den aabenbart Begeistrede i en positiv Tid), ret som var den sande Ironiker ikke Offrende, da jo dog hin Stormesteren endte med at blive straffet paa Livet.

Reflexions-Spændingen constituerer sig tilsidst til Princip, og som i en lidenskabelig Tid Begeistring er det enende Princip, saaledes bliver i en lidenskabsløs og meget reflekterende Tid Misundelse det negativt-enende Princip. Dette maa dog ikke strax forstaaes i ethisk Betydning som en Sigtelse, nei, Reflexionens Idee, hvis man kan tale saa, er Misundelse, og Misundelsen er derfor en dobbelt, er den selviske i Individet og da atter Omgivelsens mod ham. Reflexionens Misundelse i Individet forhindrer den pathetiske Afgjørelse i ham; og er det da næsten som vilde det lykkes ham, saa standser Omgivelsens Reflexions-Modstand ham. Reflexionens Misundelse holder Villien og Kraften ligesom i Fangenskab. Først maa da Individet gjennembryde Fængslet, hvori den egne Reflexion holder ham, og naar saa det er lykkedes, staaer han endnu ikke i det Frie, men i den store Fængsels-Bygning, som Omgivelsens Reflexion danner, og har ved Reflexions-Forholdet i sig atter et Forhold til denne, som kun religieus Inderlighed kan frigjøre ham fra, selv hvormeget han end gjennemskuer Forholdets Usandhed. Men at det er et Fængsel, hvori Reflexionen holder Individet og Tiden, at det er Reflexionen, der gjør det, og ikke Tyranner og hemmeligt Politie, ikke Præster og Aristokrater, denne Forstaaelse forhindrer Reflexionen af al Magt og vedligeholder den smigrende Indbildning, at Reflexions-Muligheden er noget ganske anderledes Stort end den fattige Afgjørelse. Den selviske Misundelse fordrer i Ønskets Form for meget af Individet selv, og forhindrer ham derved, den forkjeler ham ligesom en svag Moders Forkjærlighed, thi Misundelsen paa sig selv forhindrer Individet i at give sig hen. Omgivelsens Misundelse, hvori jo Individet atter selv deeltager mod Andre, er i negativ-critisk Forstand misundelig.

Men jo længere dette gaaer hen, desto mere vil Reflexionens Misundelse constituere sig til ethisk Misundelse. Indespærret Luft udvikler altid Gift, og saaledes udvikler Reflexions-Indespærretheden, naar ingen Handling, ingen Begivenhed lufter ud, den fordømmelige Misundelse. Medens de bedre Kræfter holde hinanden Stangen i en Reflexions-Spænding, kommer Usselheden op, dens Frækhed imponerer som en Slags Yttring af Kraft, og dens Foragtelighed vil blive dens beskyttende Privilegium, netop fordi den derved unddrager sig Misundelsens Opmærksomhed.

Forøvrigt ligger det dybt grundet i den menneskelige Natur, at den ikke i eet væk kan holde sig paa Høiden og vedblive at beundre, den fordrer en Afvexling. Selv den meest begeistrede Tid fordrer derfor en Misundelsens Spøg med det Udmærkede. Dette er aldeles i sin Orden, og kan skee i sin Orden, naar den, der efter Fastelavnsløierne med det Udmærkede, idet han atter fæster Beundringens Blik derpaa, er istand til at finde det uforandret; thi ellers har han ved Spøgen tabt meget Mere, end Spøgen var værd. Saaledes kan Misundelsen finde sit Raaderum selv i en begeistret Tid. Ja om end Tiden er mindre begeistret, naar den dog endnu har Kraft til at give Misundelsen Charakteer, og er enig med sig selv om, hvad dens Udtryk betyder, kan den have sin, om end farefulde Betydning. Saaledes var for Exempel Ostracismen i Grækenland et Udtryk for Misundelsen, et Slags Ligevægtens Nødværge mod det Udmærkede. Man udøvede den da, men man var enig med sig selv om, hvad der da dialektisk ligger i Forholdet, at Ostracismen var en Udmærkelse. Det kunde derfor i Forhold til Gjengivelsen af en noget tidligere Tid i Grækenland være ironisk i Aristophanes's Aand, at lade en aldeles ubetydelig Mand blive viist i Landflygtighed ved Ostracismen. Dette Ironiske vilde være en endnu høiere Comik end den, der for Exempel, lader en saadan Ubetydelig ironisk blive Hersker, netop fordi Landflygtighed ved Ostracisme allerede er det negative Udtryk for Udmærkelsen; hvorfor det igjen vilde blive en høiere ironisk Comik, at lade det ende med, at Folket kaldte den Forviste tilbage, fordi det ikke kunde undvære ham, der da maatte blive en reen Gaade for de Mennesker, mellem hvilke han levede i sit Exil, da de naturligviis slet intet Udmærket kunde opdage hos ham. I »Ridderne« gjengiver Aristophanes den fuldkomne Forraadnelses-Tilstand, hvor Pøbelagtigheden ender med, ligesom man dyrker Dalei Lamas Excrementer, at tilbede eller tilbedende at anskue sig selv i den første den bedste af Udskuddet, et Forhold, der i Opløsthed svarer indenfor Folkeregjeringens Bestemmelse til det, at sætte Keiserværdigheden til Auction. Men naar Misundelsen endnu har Charakteer, er Ostracismen en negativ Udmærkelse. Den Mand, der sagde til Aristides, at han stemmede for hans Landflygtighed, »fordi han ikke kunde lide, at Aristides kaldtes den eneste Retfærdige,« han negtede egentligen ikke Aristides Udmærkelsen, men han tilstod Noget om sig selv, at han istedenfor i Beundringens lykkelige, i Misundelsens ulykkelige Forelskelse forholdt sig til det Udmærkede, men han forkleinede det ikke.

Jo mere derimod Reflexionen faaer Overhaand til at udvikle Indolents, desto farligere bliver Misundelsen, fordi den ikke har Charakteer til at blive sig selv sin egen Betydning bevidst, men charakteerløst i Ustadighedens og Feighedens luskende Forhold til Omstændighederne omforklarer den samme Yttring paa den forskjelligste Maade, vil at den skal være Spøg, og hvis den seer det mislykket, at den skal være Fornærmelse, og hvis det mislykkes, at den skal være Ingenting, vil at den skal være Vittighed, og hvis det ikke slaaer til, forklarer, at den heller ikke var saaledes meent, at den var ethisk Satire, som man jo skal bryde sig om, og hvis dette mislykkes, forklarer, at den var Ingenting, som Ingen skal bryde sig om. Misundelsen constituerer sig som Charakteerløshedens Princip, der fra Usselhed vil snige sig op til at være Noget, bestandigt dækkende sig ved den Concession, at den er Ingenting. Charakteerløshedens Misundelse forstaaer ikke, at det Udmærkede er det Udmærkede, den forstaaer sig ikke i, selv negativt at anerkjende det, men den vil have det ned, have det Udmærkede forkleinet, saa det virkeligen ikke mere er det Udmærkede; og Misundelsen retter sig mod det Udmærkede, som er, og mod det, som vil komme.

Den sig etablerende Misundelse er Nivelleringen, og medens en lidenskabelig Tid fremskynder, hæver og styrter, ophøier og nedtrykker, gjør en reflekteret lidenskabsløs Tid det Modsatte, den qvæler og forhindrer, den nivellerer. At nivellere er en stille mathematisk abstrakt Beskjæftigelse, der undgaaer al Ophævelse. Medens den opblussende flygtige Begeistring mismodigt kunde ønske om det saa var en Ulykke, blot for at fornemme Tilværelsens Kræfter, er dens Afløsning Apathien ikke tjent med nogen Forstyrrelse, saa lidet som den nivellerende Ingenieur er det. Er et Oprør i sit Maximum som en Vulkans Udbrud, saa man ikke kan høre Ørenlyd: saa er Nivellementet i sit Maximum som en Døds-Stilhed, hvor man kan høre sit eget Aandedrag, en Døds-Stilhed, hvorover Intet kan hæve sig, men Alt synker afmægtigt ned i den.

I Spidsen for et Oprør kan der staae en enkelt Mand, men i Spidsen for Nivelleringen kan ingen enkelt Mand staae, thi saa vilde han jo blive Herskeren, og være undgaaet Nivelleringen. Det enkelte Individ kan i sin lille Kreds være medvirksomt i Nivelleringen, men denne er en abstrakt Magt, og Nivelleringen er Abstraktionens Seier over Individerne. Nivelleringen er i den moderne Tid Reflexionens Tilsvarende til Skjebnen i Oldtiden. Oldtiden er dialektisk i Retning af Fremragenhed (den enkelte Store – og saa Mængden, een Fri – og saa Trælle); Christendommen er indtil videre dialektisk i Retning af Repræsentation (Fleertallet anskuer sig selv i den Repræsenterende, frigjøres ved Bevidstheden om at det er dem, han repræsenterer, i en Art Selvbevidsthed); Nutiden er dialektisk i Retning af Ligeligheden, og dennes, i Forfeilethed, meest consequente Gjennemførelse er Nivellementet, som den negative Eenhed af Individernes negative Gjensidighed.

Enhver vil let see, at Nivelleringen har sin dybe Betydning i Generationskategoriens Overmagt over Individualitetskategorien. Medens i Oldtiden Individernes Mængde ligesom var til for at bestemme Prisen paa, hvor meget det udmærkede Individ var værd, saa er Myntfoden nu saaledes forandret, at der ligeligt gaaer omtrent saa og saa mange Mennesker paa eet Individ, saa det blot gjelder om at sikkre sig det behørige Antal – saa har man Betydning. Den Enkelte i Mængden betydede i Oldtiden slet Intet, den Udmærkede betydede alle disse; Nutiden tenderer til den mathematiske Ligelighed, at der saa omtrent ligeligt gjennem alle Stænder gaaer saa og saa Mange paa eet Individ. Den Udmærkede turde tillade sig Alt, de Enkelte i Mængden slet Intet; nu forstaaer man, at der gaaer saa og saa mange paa eet Individ, og ganske consequent tæller man sig sammen (man kalder det jo rigtignok at forene sig, men det er et Galanteri) i Forhold til det Ubetydeligste. Blot om at realisere et Indfald tæller man sig nogle Stykker sammen, og saa gjør man det det vil sige saa tør man gjøre det. Deraf kommer det tilsidst, at selv en fortrinligere Begavet dog ikke kan frigiøre sig fra Reflexionen, fordi han snart i det Ubetydeligste bliver sig bevidst som en Brøks-Deel, og gaaer Glip af Religieusitetens uendelige Frigjørelse. Selv om en Forening af Flere havde Mod til at gaae Døden imøde: det vilde i vor Tid derfor ikke være sagt, at hver Enkelt havde Mod dertil, thi Det, den Enkelte frygtede mere end Døden, var Reflexionens Dom over ham, Reflexionens Indsigelse mod ham, at han som Enkelt vilde vove Noget. Den Enkelte tilhører ikke Gud, ikke sig selv, ikke den Elskede, ikke sin Kunst, ikke sin Videnskab, nei, som en Vorned tilhører et Gods, saaledes bliver den Enkelte sig bevidst i eet og alt at tilhøre en Abstraktion, som Reflexionen ordner ham ind under. Om en Forening af Flere i vor Tid kunde beslutte sig til, hver især, at give hele sin Formue til et eller andet velgjørende Øiemed: deraf fulgte ikke, at den Enkelte kunde beslutte sig dertil, og atter ikke fordi han var tvivlraadig med Hensyn til Afkaldet paa sin Formue, men fordi han frygtede Reflexions-Dommen meget mere end Armoden. Om Ti kunde blive enige om at vedstaae Elskovs fulde ubeskaarne Gyldighed, eller Begeistringens ved ingen lammende Hensyn uindskrænkede Berettigelse, deraf fulgte ikke, at hver af de Ti var istand dertil, thi høiere endnu end Elskovens Salighed og end Begeistringens Vidnen med deres Aand vilde de dog tvetydigt elske Reflexions-Dommen – derfor maatte de være Ti om det, som det er en Modsigelse at være mere end Een om. Socialitetens forgudede positive Princip i vor Tid er netop det Fortærende, det Demoraliserende, der i en Reflexions Trældom gjør selv Dyderne til vitia splendida. Og hvoraf kan dette komme, uden deraf, at den religieuse Individualitets Udsondring for Gud i Evighedens Ansvar forbigaaes. Naar Forfærdelsen her begynder, søger man Trøst i Compagnie, og saa fanger Reflexionen Individet for hele Livet. Og De, der end ikke fornam Begyndelsen af denne Crise, de gaae uden videre ind under Reflexions-Forholdet.

Nivellementet er ikke en Enkelts Handling, men et Reflexions-Spil i en abstrakt Magts Haand. Ligesom man beregner Diagonalen i Kræfternes Parallellogram, saaledes kan man beregne Loven for Nivellementet. Thi den Enkelte, der nivellerer Nogle, tages selv igien med, og saa fremdeles. Medens derfor den Enkelte egoistisk synes at vide, hvad han gjør, maa man sige om dem Alle, de vide ikke hvad de gjøre, thi ligesom der i Begeistringens Samdrægtighed kommer et Mere ud, som ikke er de Enkeltes, saa kommer ogsaa her et Mere ud. Man maner en Dæmon op, som ingen Enkelt kan magte; og medens den Enkelte i Nivellerings-Lystens korte Øieblik selvisk nyder Abstraktionen, saa underskriver han tillige sin egen Undergang. Den Begeistredes Fremtrængen kan ende med Undergang, men den Nivellerendes Seier er eo ipso hans Undergang. Nivelleringens Skepsis kan ingen Tid standse, Tiden, Nutiden altsaa heller ikke, thi i det Øieblik, den vil standse den, vil den atter udvikle Loven. Den kan kun standses ved, at Individet i individuel Udsondring vinder Religieusitetens Uforfærdethed. Jeg saae engang et Slagsmaal, hvor tre Mennesker paa en skammelig Maade mishandlede en fjerde. Hoben stod og betragtede det med Harme; Uvilliens mumlende Udtryk begynder at give Fart til Handling: da samler et Par Stykker af Mængden sig, tager fat paa den ene af Angriberne, og kaster ham til Jorden og saa videre Altsaa Hævnerne udviklede den samme Lov som Angriberne. Hvis det tør tillades mig at anbringe min ringe Person med, saa skal jeg fortælle Historien ud. Jeg traadte til og søgte dialektisk at forklare en af Hævnerne det Inconsequente i deres Adfærd, men det lod til at være ham aldeles umuligt at indlade sig paa Sligt, han sagde blot gjentagende: »Det har han ærligt fortjent, at en saadan Slyngel skal være Tre om Een.« Det Comiske ligger nær nok især for Den, der ikke havde seet Begyndelsen og altsaa hørte en Mand sige om en anden: at han (den eenlige) var Tre om Een, og hørte det netop i det Øieblik da det Modsatte var Tilfældet: at de vare Tre om ham. Det Første vilde være comisk ved Modsigelsen i samme Forstand som »da Vægteren sagde til en eenlig Person: vil De være saa god at gaae fra hinanden;« det Andet vilde være comisk ved Selvmodsigelsen. Hvad jeg derimod forstod var, at det nok var bedst at opgive Haabet om at ende denne Skepsis, at den ikke skulde fortsættes mod mig.

Ingen enkelt Mand (den Udmærkede i Retning af Fremragenhed og Skjebnens Dialektik) vil kunne standse Nivelleringens Abstraktion, thi den er et negativt Høiere, og Heltenes Tid er forbi. Ingen Congregation vilde være istand til at standse Nivelleringens Abstraktion, fordi Congregationen selv er gjennem Reflexionens Sammenhæng i Nivelleringens Tjeneste. End ikke Nationaliteternes Individualitet vil kunne standse den, thi Nivelleringens Abstraktion reflekterer paa en høiere Negativitet: den rene Menneskehed. Nivelleringens Abstraktion, denne Menneske-Slægtens Selvantændelse, foranlediget ved den Friktion, der opstaaer, naar den individuelle Inderligheds Udsondring i Religieusitet udebliver, vil blive staaende som man siger det om en Passat, som fortærer Alt, men ved hvilken Individerne, hver især, igjen kan opdrages religieust, i høieste Forstand kan hjælpes til, i Nivelleringens examen rigorosum at vinde Religieusitetens Væsentlighed i sig selv. For den Yngre, der, hvor fast han end for sit Vedkommende hænger ved hvad han beundrer som Udmærket, fra Begyndelsen fatter, at Nivelleringen er hvad den selviske Enkelte og den selviske Slægt tænkte til det Onde, men ogsaa hvad for den Enkelte, hver især, hvis han saa vil det i Oprigtighed med Gud, kan blive Udgangspunktet for det høieste Liv – for ham vil det være i Sandhed dannende at leve i Nivelleringens Tid. Samtidigheden vil for ham være i høieste Forstand religieust udviklende og tillige æsthetisk og intellectuelt uddannende, idet det Comiske vil gjøre sig absolut gjeldende. Thi det høieste Comiske er netop, at det enkelte Individ skal uden nogen Mellembestemmelse henføres under den uendelige Abstraktion af den rene Menneskehed, idet alle Organisationens Individualitets-Concretioner, der ved Relativiteten temperere det Comiske og styrke ved relativ Pathos, ere fortærede. Men dette er atter Udtrykket for, at Frelsen kun er ved Religieusitetens Væsentlighed i det enkelte Individ. Og begeistrende for ham vil det være, at fatte, at netop Vildfarelsen aabner den Enkelte, hver især, hvis han høimodigt vil det, Adgang til det Høieste. Men Nivelleringen maa blive staaende, den maa til, ligesom Forargelsen maa indkomme i Verden, men vee den, ved hvem den kommer ind.

Det er ofte nok sagt, at en Reformation skal begynde med, at Enhver reformerer sig selv; men det er ikke saaledes skeet, thi Reformationens Idee har affødt en Helt, der maaskee dyrt nok af Guden har kjøbt sin Bestalling som Helt. Ved da ligefremt at slutte sig til ham faae Individerne for bedre, ja for godt Kjøb det Dyrekjøbte, men de faae heller ikke det Høieste. Nivelleringens Abstraktion derimod er et Princip, ligesom Østens skarpe Vind, der ikke indlader sig med det enkelte Individ i noget intimere Forhold, men kun i Abstraktions-Forholdet, der er ligeligt for Alle. Ingen Helt lider da for Andre og hjælper Andre, Nivelleringen bliver selv den skarpe Tugtemester, der tager sig af Opdragelsen. Og Den, der lærer Maximum af Opdragelsen og bliver Maximum, han bliver ikke den Udmærkede, Helt, den Fremragende, dette forhindrer Nivelleringen, der er consequent indtil sit Yderste, og han forhindrer det selv, fordi han har fattet Nivelleringens Betydning, nei han bliver kun et væsentligt Menneske i den fyldige Ligeligheds Forstand. Dette er Religieusitetens Idee. Men Opdragelsen er stræng, og Udbyttet tilsyneladende meget lille; tilsyneladende, thi dersom Individet ikke vil lære i Religieusitetens Væsentlighed for Gud at nøies med sig selv, nøies med, istedetfor at herske over Verden, at herske over sig selv, nøies med som Præst at være sin egen Tilhører, som Forfatter sin egen Læser og saa videre, dersom han ikke vil lære at begeistres ved dette som det Høieste, fordi det udtrykker Ligheden for Gud og Ligheden med Alle: saa slipper han ikke ud af Reflexionen, saa oplever han maaskee i Forhold til sin Begavethed, et skuffende Øieblik, hvor han troer det er ham, der nivellerer, indtil han selv segner under Nivelleringen. Det hjælper ikke at bebude og varsle en Holger Danske eller en Morten Luther, deres Tid er forbi, og det er dog i Grunden Individernes Magelighed, der ønsker en saadan, Endelighedens Utaalmodighed, der vil have billigt paa anden Haand istedenfor det Høieste, der kjøbes dyrt paa første Haand. Det hjælper ikke at stifte Selskab paa Selskab, thi negativt er der sat noget Høiere ind, om den kortsynede Selskabsmand end ikke kan see det. Individualitets-Principet i sin umiddelbare og skjønne Formation præfigurerer Generationen ved den Udmærkede, den Fremragende, og lader de underordnede Individualiteter slutte sig grupperende om Repræsentanten. Individualitets-Principet i sin evige Sandhed bruger Generationens Abstraktion og Ligelighed som Nivellerende og udvikler derved Individet i dets egen Medvirken religieust til et væsentligt Menneske. Thi saa afmægtig som Nivelleringen er mod det Evige, saa overmægtig er den mod enhver Midlertidighed. Reflexionen er en Slynge, hvori man fanges, men ved Religieusitetens begeistrede Spring bliver Forholdet et andet, saa bliver den Slyngen, som kaster En i det Eviges Favn. Og Reflexionen er og bliver den meest haardnakkede Creditor i Tilværelsen; listigen har den hidtil kjøbt alle mulige Livs-Anskuelser op, men den væsentlige Religieusitets evige Livs-Anskuelse kan den ikke kjøbe, derimod kan den friste med glimrende Blendværk fra alt det Andet, mistrøste med Reminiscentser fra alt det Andet. Men ved Springet paa Dybet lærer man at hjælpe sig selv, lærer at elske alle Andre lige saa høit som sig selv, hvad enten man saa anklages for Anmasselse og Stolthed – at man ikke vil lade sig hjælpe, eller for Selviskhed – at man ikke underfundigen vil bedrage Andre ved at hjælpe dem, det vil sige ved at hjælpe dem til at gaae Glip af det Høieste. – Vil Nogen sige, at hvad jeg her har fremsat, veed Enhver og kan Enhver sige, da er dette mit Svar: saa meget desto bedre, jeg ønsker ingen Fremragenhed, jeg har Intet mod, at Enhver veed det, med mindre det, at Enhver veed det og Enhver kan sige det, skulde betyde, at det skal fratages mig og deponeres i det negative Fællesskab. Faaer jeg blot Lov at beholde det, saa taber det for mig ikke i Værd derved, at Enhver veed det.

I Grunden har den moderne Tid længe tenderet til Nivelleringen gjennem mange Omvæltninger, hvilke dog alle ikke vare Nivelleringen, fordi de alle ikke vare abstrakte nok, men havde en Virkelighedens Concretion. Tilnærmelsesviis kan der nivelleres ved, at det Fremragende styrer mod det Fremragende, saa begge svækkes; tilnærmelsesviis kan der nivelleres derved, at et Fremragende holdes neutraliseret af et andet Fremragende; ved, at Foreningen af det i sig Svagere bliver stærkere end det enkelte Fremragende; tilnærmelsesviis kan der nivelleres ved en enkelt Stand, for Exempel Geistlighed, Borgerstand, Bondestand, ved Folket selv: men alt dette er dog kun Abstraktionens Rørelser indenfor Individualitets Concretioner.

For at Nivelleringen egentligen skal komme istand, maa der først bringes et Phantom tilveie, dens Aand, en uhyre Abstraktion, et altomfattende Noget som er Intet, et Luftsyn – dette Phantom er Publikum. Kun i en lidenskabsløs, men reflekteret Tid kan dette Phantom udvikle sig ved Hjælp af Pressen, naar denne selv bliver en Abstraktion. I begeistrede Tider, i lidenskabelige tumultuariske Tider, selv naar et Folk vil realisere den ufrugtbare Ørkens Idee at ødelægge og sløife Alt: der er dog intet Publikum. Der er Partier og der er Concretion. Pressen vil i saadanne Tider antage Concretionens Charakteer i Forhold til Adsplittelsen. Men som stillesiddende Professionister især ere udsatte for at udvikle phantastiske Sandsebedrag, saaledes vil en lidenskabsløs, stillesiddende, reflekteret Tid, naar Pressen skal være det eneste, der selv svagt skal holde et Slags Liv i denne Døsighed, udvikle dette Phantom. Publikum er den egentlige Nivellerings-Mester, thi naar der nivelleres tilnærmelsesviis, nivelleres der ved Noget, men Publikum er et uhyre Intet.

Publikum er et Begreb, som slet ikke kan forekomme i Oldtiden, fordi Folket selv en masse in corpore maatte træde op i Handlingens Situation, maatte bære Ansvaret for hvad den Enkelte af deres Midte afstedkom, medens igjen den Enkelte personligen som denne Bestemte maatte være tilstede, maatte underkaste sig Øieblikkets Standret i Bifald eller Misbilligelse. Først naar intet kraftigt Samliv giver Concretionen Fylde, vil Pressen danne dette Abstraktum Publikum, der bestaaer af uvirkelige Enkelte, som aldrig forenes eller kunne forenes i nogen Situationens eller Organisationens Samtidighed, og som dog fastholdes som et Hele. Publikum er et Corps, talrigere end alle Folk tilsammen, men dette Corps kan aldrig blive mønstret, ja det kan end ikke saa meget som have en eneste Repræsentant, fordi det selv er en Abstraktion. Og desuagtet bliver Publikum, naar Tiden er lidenskabsløs og reflekteret og udviskende Alt det Concrete, det Hele, som skal omfatte Alle. Men dette Forhold er atter netop Udtrykket for, at den Enkelte er anviist sig selv.

I det virkelige Øiebliks og den virkelige Situations Samtidighed med de Virkelige, der hver ere Noget, er der for den Enkelte det Understøttende. Men Publikums Tilværelse danner ingen Situation og ingen Forsamling. Den Enkelte, der læser, er jo ikke Publikum, og saa læser lidt efter lidt mange Enkelte, maaskee alle Enkelte, men der er ingen Samtidighed. Publikum kan bruge Aar og Dag for ligesom at blive samlet, og naar det saa er samlet, er det dog ikke til. Abstraktionen, som Individerne paralogistisk danne, støder ganske rigtigt Individerne fra sig istedenfor at hjælpe dem. Den, der i det virkelige Øiebliks og den virkelige Situations Samtidighed med de Virkelige, selv ingen Mening har, antager samme Mening som Majoriteten, eller er han mere stridbar, som Minoriteten. Men Majoriteten og Minoriteten er vel at mærke virkelige Mennesker, og deri ligger det Understøttende i Tilholdet til disse. Publikum derimod er en Abstraktion. At antage samme Mening som disse og disse bestemte Mennesker betyder, at man veed, at disse ville være underkastede samme Farer som En selv, at de ville fare vild med En, hvis Meningen er vildfarende og saa videre Men at antage samme Mening som Publikum er en svigefuld Trøst, thi Publikum er kun til in abstracto. Medens derfor ingen Majoritet nogensinde har været saa sikker paa at beholde Retten og Seiren som Publikum er det, saa er dette kun lidet trøsteligt for den Enkelte, thi Publikum er et Phantom, som ikke tillader nogen personlig Tilnærmelse. Dersom En idag antager Publikums Mening og imorgen bliver udpeben, saa bliver han udpeben af Publikum. En Generation, et Folk, en Folkeforsamling, et Samfund, en Mand har dog et Ansvar ved at være Noget, kan skamme sig ved Ustadighed og Troløshed, men Publikum bliver Publikum. Et Folk, en Forsamling, et Menneske kan forandre sig saaledes, at man maa sige: han er ikke mere den Samme; men Publikum kan blive lige det Modsatte og er dog det Samme – er Publikum. Men netop ved denne Abstraktion og denne abstrakte Optugtelse dannes Individet (forsaavidt det ikke allerede ved sin egen Inderlighed er dannet) dersom det ikke gaaer under, til i Religieusitetens høieste Forstand at nøies med sig selv og sit Guds-Forhold, til istedenfor Enigheden med Publikum, der fortærer alle de relative Individualitets-Concretioner, at substituere det at være enig med sig selv, til istedenfor at tælle og tælle, at finde Hvilen i sig selv for Gud. Og Dette vil være det Modernes absolute Forskjel fra Oldtiden: at det Totale ikke er Concretionen, der understøtter, der danner den Enkelte, uden dog at udvikle ham absolut, men er en Abstraktion, der i sin abstrakte Ligelighed frastødende hjælper ham til at dannes absolut – hvis han ikke omkommer. Det Trøstesløse i Oldtiden var, at den Udmærkede var, hvad de Andre ikke kunde være, det Begeistrende vil blive, at Den, der religieust vandt sig selv, kun er hvad Alle kunne være.

Publikum er ikke et Folk, ikke en Generation, ikke en Samtid, ikke en Menighed, ikke et Selskab, ikke disse bestemte Mennesker, thi alt Sligt er kun ved Concretionen det som det er; ja ikke en eneste af dem, der høre til Publikum, har noget væsentligt Engagement; i nogle Timer af Dagen hører han maaskee med til Publikum, nemlig i de Timer, i hvilke han Ingenting er, thi i de Timer, i hvilke han er det Bestemte han er, hører han ikke til Publikum. Dannet af saadanne Enere, af de Enkelte i de Øieblikke, hvor de Ingenting ere, er Publikum noget uhyre Noget, det abstrakte Øde og Tomme, som er Alle og Ingen. Men af den samme Grund kan Enhver anmasse sig at have et Publikum, og ligesom den romerske Kirke chimairisk udvidede sig ved at udnævne Biskopper in partibus infidelium: saaledes er Publikum Noget, som Enhver, selv en fuld Matros, der foreviser »en Perspektivkasse«, kan tilegne sig, og den fulde Matros har dialektisk consequent absolut samme Ret dertil som den meest Udmærkede, absolut Ret til at sætte alle disse mange, mange Nuller foran sit Eettal. Publikum er Alt og Intet, er den farligste af alle Magter og den meest intetsigende; man kan tale til en heel Nation i Publikums Navn, og dog er Publikum mindre end et eneste nok saa ringe virkeligt Menneske. Bestemmelsen Publikum er det Reflexionens Blendværk, der gøglende har gjort Individerne indbildske, fordi Enhver kan anmasse sig dette Uhyre, i Sammenligning med hvilket Virkelighedens Concretioner synes fattige, Publikum er Forstands-Tidens Eventyr, der gjør de Enkelte phantastisk til mere end at være Konge over et Folk; men Publikum er atter den grusomme Abstraktion, ved hvilken Individerne skal religieust opdrages – eller gaae under.

Pressens Abstraktion (thi et Blad, en Avis er ingen statsborgerlig Concretion og kun i abstrakt Forstand et Individuum) i Forening med Tidens Lidenskabsløshed og Reflekterethed afføder Abstraktionens Phantom: Publikum, der er den egentlig Nivellerende. Ogsaa dette kan, fraseet dets negative Betydning for Religieusiteten, have sin Betydning. Men jo mindre Idee der er i en Tid, jo mere den afmattet ved en opblussende Begeistring udhviler sig i Indolents, hvis vi endog vilde tænke os, at Pressen blev svagere og svagere, fordi ingen Begivenhed, ingen Idee griber Tiden: desto lettere vil Nivelleringen blive en fordærvelig Lyst, en Sandsepirring, som et Øieblik kildrer og kun gjør det Onde værre, og Frelsens Vilkaar tungere og Undergangens Sandsynlighed større. Og har man ofte skildret Eneherredømmets Demoralisation, revolutionaire Tiders Forfald, saa er en lidenskabsløs Tidsalders Forfald noget ligesaa fordærveligt, om end formedelst Tvetydigheden mindre paafaldende. Og at tænke det, kan derfor vel have sin Interesse og sin Betydning. Flere og flere Enkelte ville da i Indolentsens Blødagtighed aspirere til at blive Ingenting – for at blive Publikum, dette abstrakte Hele, der er dannet paa den latterlige Maade, at Participanten er Trediemand. Denne dvaske Mængde, som Intet selv forstaaer og Intet selv vil gjøre, dette Gallerie-Publikum søger nu Tidsfordriv, og hengiver sig da til den Indbildning, at Alt hvad Nogen gjør, skeer, for at det kan faae noget at snakke om. Dvaskheden sidder fornemt med Benene overkors, og Enhver, der vil arbeide, Kongen og Embedsmanden og Folkets Lærer og den dygtigere Journalist og Digteren og Kunstneren, Alle blive ligesom spændte for, for at slæbe denne Dvaskhed frem, der fornemt troer, at de andre ere Hestene. Vilde jeg tænke mig dette Publikum som en Person (thi om end og enkelte Bedre momentviis høre til Publikum, de have dog i sig selv en organiserende Concretion, der holder dem fast, selv om de ikke vinde Religieusitetens Høieste) saa vilde jeg nærmest tænke paa en eller anden romersk Keiser, en stor velnæret Figur, der lider af Keedsomhed, og derfor blot ønsker Latterens Sandsepirring, thi Vittighedens guddommelige Gave er ikke jordisk nok. Saa slendrer da til Afvexling denne Person omkring, mere dvask end ond, men negativ herskesyg. Enhver, der har læst Oldtidens Forfattere, veed, hvor mange Ting en Keiser kunde hitte paa, for at forkorte Tiden. Saa holder da Publikum sig en Hund for Morskabs Skyld. Denne Hund er den literaire Foragtelighed. Viser der sig nu en bedre, maaskee endog en Udmærket, saa hidses Hunden paa, og Morskaben begynder. Den bidske Hund river ham i Kjoleskjøderne, tillader sig alle uvorne Kaadheder – indtil Publikum bliver keed deraf og siger: nu kan det være nok. Da har Publikum nivelleret. Den Bedre, den Stærkere er mishandlet, – og Hunden, ja den bliver ved at være en Hund, som Publikum selv foragter. Der nivelleredes altsaa ved et Tredie, Intethedens Publikum nivellerede ved et Tredie, som ved sin Foragtelighed allerede var mere end nivelleret og mindre end Intet. Og Publikum er angerløst, thi det er jo ikke Publikum – det var jo Hunden; ligesom man siger til Børn: det var Katten. Og Publikum er angerløst, thi det var jo ingen væsentlig Forkleinelse – det var blot lidt Moerskab. Dersom nemlig Nivelleringens Redskab havde været en udmærket Dygtighed, saa havde det indolente Publikum været narret, thi saa var jo Redskabet igjen blevet et Forstyrrende, men naar man holder det Bedre nede ved Foragteligheden og Foragteligheden nede ved den selv: saa er dette Intethedens Quit. Og Publikum vil være angerløst, thi det er jo egentligen ikke dem, som holde Hunden, de subskribere blot; de hidse den heller ikke directe paa, de fløite heller ikke ligefrem af den; i Tilfælde af Søgsmaal vilde Publikum sige: Hunden er ikke min, det er en herreløs Hund; og i Tilfælde af, at Hunden blev optagen og sendt ud paa Veterinairskolen for at slaaes ihjel, kunde Publikum endnu sige: det var virkeligen godt, at den slemme Hund blev dræbt, det har vi alle ønsket – endogsaa Subskribenterne.

Maaskee er Nogen, som tænkende sætter sig ind i et saadant Forhold, tilbøielig til at fæste sin Opmærksomhed paa den Bedre, der leed Mishandlingen, og mene, at der vederfores ham en stor Ulykke. Denne Betragtning kan jeg aldeles ikke billige, thi den, der vil hjælpes til det Høieste, han har netop godt af at gjennemgaae Sligt og maatte snarere ønske det, om man end paa hans Vegne kan oprøres. Nei det Forfærdelige er noget Andet, det er Tanken om de mange Menneskeliv, der gaae spildte, eller let gjøre det. Jeg vil end ikke tale om de Fortabte, eller dog indtil Fortabelse Vildledte, der for Penge spille Hundens Rolle, men de mange Ubefæstede, Letsindige, Sandselige, der i fornem Indolents ikke faae noget dybere Indtryk af Tilværelsen end dette meningsløst grinende, alle de Daarlige, der føres i ny Fristelse, idet de i deres Indskrænkethed endog blive sig selv vigtige ved at have Medlidenhed med de Angrebne, uden at fatte, at ved et saadant Forhold ere de Angrebne altid de Stærkeste, uden at fatte, at her gjelder det med et forfærdeligt og dog ironisk Eftertryk: græder ikke over ham, men græder over Eder selv.

Dette Forhold er Nivelleringens laveste, thi den svarer altid til, hvad Divisor er, for hvem alle gjøres lige: saaledes er det evige Liv ogsaa en Art Nivellering, og dog er det ikke saa, fordi Divisor er det: at være et i religieus Forstand væsentligt Menneske.

Fra disse dialektiske Kategorie-Bestemmelser og sammes Consequentser, disse være nu i det givne Øieblik faktiske eller ikke, fra det dialektiske Sigte paa Nutiden gaaer jeg nu over til dialektisk at naae de concretere Prædikater, i hvilke den novellistiske Gjengivelse af Nutiden viser sig i sin Reflex i det Huslige og Omgangs-Livet. Det er Skyggesiden, her kommer til at vise sig, og om end dennes Facticitet ikke lader sig negte, saa er det dog ogsaa vist, at saalidet som Reflexionen selv er det Onde, saa vist maa ogsaa en meget reflekteret Tid have sin Lys-Side, netop fordi megen Reflekterethed betinger en høiere Betydningsfuldhed end umiddelbar Lidenskab; betinger den – naar Begeistringen træder til og fører Reflexions-Kræfterne ind i Afgjørelsen; og fordi megen Reflekterethed giver en større Gjennemsnits-Dygtighed i Forudsætningerne til at handle – naar Religieusiteten træder til i Individet og overtager Forudsætningerne. Reflexionen er ikke det Onde, men Tilstanden i Reflexionen og Stilstanden i Reflexionen er Misligheden og Fordærvelsen, der ved at forvandle Forudsætningerne til Udflugter foranlediger Tilbagegang.

Nutiden er væsentligen den forstandige, den lidenskabsløse, derfor har den hævet Modsigelsens-Grundsætning. Af denne Betragtning lader den Mangfoldighed af Træk sig deducere, som Novellens Forfatter med sin fine Kunst, med sin ophøiede Ligevægt saa interesseløst har gjengivet. Thi Forfatterens Mening erfarer man naturligviis intetsteds, han gjengiver blot Reflexen. Overhovedet kan man om en lidenskabsløs men reflekteret Tid i Sammenligning med en lidenskabelig sige: den vinder i Extensitet, hvad den taber i Intensitet. Men denne Extensitet kan atter blive Betingelsen for en høiere Form, naar en tilsvarende Intensitet igjen indtager, hvad der extensivt er disponeret over.

At hæve Modsigelsens Grundsætning, er i Existents Udtrykket for at være i Modsigelse med sig selv. Det skaberiske Almægtige, der ligger i den absolute Disjunktions-Lidenskab, som bringer Individet i den sluttede Enighed med sig selv, forvandles til Forstands-Reflexionens Extensitet: ved at vide og være alt Muligt, at være i Modsigelse med sig selv det vil sige slet Intet at være. Modsigelsens Grundsætning styrker Individet i Troskab med sig selv, saa han, ligesom hiint standhaftige Tretal, Socrates saa skjønt taler om, hellere vil udholde Alt end blive til et Fiirtal eller vel endog til et meget stort rundt Tal, saa han hellere vil være Lidet i Troskab med sig selv end Allehaande i Modsigelse med sig selv.

Hvad er det at snakke? Det er Ophævelsen af den lidenskabelige Disjunktion mellem at tie og tale. Kun Den, der væsentligen kan tie, kan væsentligen tale, kun Den, der væsentligen kan tie, kan væsentligen handle. Tausheden er Inderligheden. Snakken anticiperer den væsentlige Talen, og Reflexionens Yttring svækker ved Forprang Handlingen. Men Den, der væsentligen kan tale, fordi han kan tie, han vil ikke have det Mangfoldige at tale om, men det Ene, og han vil finde Tid til at tale og til at tie. Snaksomheden vinder extensivt: den faaer alt Muligt at snakke om, og bliver ved i eet væk. Naar i en Tid Individerne ikke i stille Nøisomhed, i gemytlig Tilfredshed, i religieus Inderlighed ere vendte ind ad, men i Reflexions-Forholdet søge ud efter og søge hinanden; naar ingen stor Begivenhed knytter Traad-Enderne i en Katastrophes Samdrægtighed: saa indtræder Snakken. Den store Begivenhed giver den lidenskabelige Tid (thi det Ene svarer da til det Andet) Noget at tale om; alle ville tale om det samme Ene, Digterne kun synge derom, Samtalerne kun gjenlyde deraf, de Forbigaaendes Hilsener indeholde Hentydninger til dette Ene. Det er Eet og det Samme. Snakken derimod har ganske anderledes Meget at snakke om. Og naar saa den store Begivenhed var forbi, naar Tausheden indtraadte, saa var der dog Noget at erindre, Noget at tænke paa, medens man tier, medens en ny Slægt taler om ganske andre Ting. Men Snakken gruer for Taushedens Øieblik, der vilde gjøre Tomheden aabenbar.

Hvad der viser sig som Loven i Forhold til digterisk Frembringelse, det Samme er ogsaa, i det mindre Forhold, Loven for ethvert Menneskes Liv i Omgang og Dannelse. Enhver, der primitivt oplever Noget, oplever tillige ved Idealiteten Mulighederne af det Samme og Muligheden af det Modsatte. Disse Muligheder er hans digteriske lovlige Eiendom. Hans egen private personlige Virkelighed derimod ikke. Hans Tale, hans Produceren er saaledes netop baaren af Tausheden. Hans Tales, hans Producerens ideale Fuldendthed vil correspondere med Tausheden, og Taushedens absolute Udtryk vil være, at Idealiteten indeholder den qualitative modsatte Mulighed. Saasnart den Producerende maa priisgive sin egen Virkelighed, dens Facticitet, saa er han ikke væsentligen productiv; hans Begyndelse vil være hans Ende, og det første Ord allerede være en Forsyndelse mod Idealitetens hellige Undseelse. En saadan digterisk Frembringelse er derfor ogsaa, æsthetisk vurderet, en Art privat Snaksomhed, og er let kjendelig paa, at den ikke har sin Modsætning hos sig i Ligelighed. Thi Idealiteten er Ligelighed mod det Modsatte. Den, der for Exempel ved at blive ulykkelig blev productiv, hvis han virkelig indviedes ved Idealiteten, han vil med samme Forkjerlighed producere det Lykkelige som det Ulykkelige. Men Tausheden, hvormed han slutter af mod sin egen personlige Virkelighed, er netop Betingelsen for at vinde Idealiteten, ellers vil han, trods alle Forsigtigheds-Midler med at forlægge Scenen til Africa og saa videre, dog være privat kjendelig i sin eensidige Forkjerlighed. Thi en Forfatter maa jo have sin private Personlighed som enhver Anden, men denne skal være hans αδυτον; og som man spærrer Indgangen i et Huus ved at stille to Soldater med korslagte Geværer: saaledes danner Idealitetens Ligelighed ved de qvalitative Modsætningers dialektiske Kors den Spærring, der gjør enhver Adgang umulig.

Men hvad der saaledes gjelder i de store Forhold, og der vel viser sig tydeligst, hvorfor det blev fremdraget: det Samme gjelder efter den mindre Maalestok i det Mindre, Taushed er atter her Betingelsen for den dannede Omgang i Samtalen mellem Mand og Mand. Jo mere et Menneske i Tausheden har Idealitet og Idee, desto mere vil han selv i daglig Omgang være istand til at gjenføde det daglige Liv og de daglige Menneskers Liv saaledes, at det er, som talede han kun paa Afstand selv om dette og dette bestemte Faktum. Jo mindre Idealitet, jo mere Udvorteshed, desto mere vil Samtalen blive en ubetydelig Opramsen og Refereren af Navne paa Personer, af »ganske tilforladelige« Privat-Efterretninger om hvad Den og Den med Navns Nævnelse har sagt og saa videre, en snaksom Fortrolighed betræffende hvad man selv vil og ikke vil, hvad man nu agter at præstere, hvad man vilde have sagt ved hiin Leilighed, hvilken Pige man gjør Cour til, hvorfor man dog ikke vil gifte sig og saa videre Taushedens Indadvendthed er Betingelsen for den dannede Tale i Omgang, Inderlighedens vrængende Udadvendthed er at snakke, er Udannethed. I Novellen vil man finde ypperlige Exempler paa en saadan Snakken; det er trivielle Ubetydeligheder, men det er altid om bestemte navngivne Personer, hvis trivielle Livsforhold interessere det Lige paa Grund af Navnet. Ligesom Klokker Link troer, han har talet med Sophie, skjøndt hun ikke har sagt Andet end Nei: saaledes høres der mangen lang Samtale, der seer ud som blev der sagt Noget, ja Meget, fordi der nævnes Navne. Den, der snakker, snakker vel om Noget, som jo Ønsket er at faae Noget at snakke om, men dette Noget er ikke i Idealitetens Forstand, thi saa taler man, det er i Trivialitetens faktiske Forstand et Noget, at Hr. Madsen er bleven forlovet og har foræret sin Kjæreste et persisk Schavl, at Digteren Petersen vil skrive en ny Digtsamling, at Skuespiller Marcussen kom til at udtale et Ord forkeert igaaraftes. Sæt Loven kunde overholdes, vi kan jo antage det, sæt altsaa, der kom en Lov ud, som ikke forbød Menneskene at tale, men blot befalede, at Alt hvad der taltes om, skulde der tales om, som var det skeet for 50 Aar siden: saa var det ude med alle Snakkerne, de vilde fortvivle; og derimod vilde det ikke væsentligen kunne forstyrre dem, der væsentligen kunne tale. At en Skuespiller har udtalt et Ord feil, kan væsentligen kun interessere, hvis der i Fortalelsen selv er noget mærkeligt, og saa er det ligegyldigt om det er 50 Aar siden – men Frøken Gusta for Exempel vilde fortvivle, hun som netop selv den Aften var i Theateret, i Loge med Kommerceraadinde Waller, og det var netop hende, der laae Mærke til det, og til at endog en af Chor-Personalet kom til at lee deraf og saa videre Det vilde jo rigtignok være Synd og en Grusomhed mod alle disse snaksomme Mennesker, der jo dog ogsaa skal leve, men derfor er Loven jo kun et Posito.

Ved denne Snakken hæves nu Distinktionen mellem det Private og Offentlige i en privat-offenlig Snaksomhed, der omtrent svarer til hvad Publikum er. Thi Publikum er det Offentlige, der interesserer sig for det allermeest Private. Hvad Ingen vilde vove at foredrage en Forsamling, hvad Ingen kunde tale om, hvad neppe de Snakkende ret ville vedgaae, at de have snakket om, det kan godt skrives for Publikum og vides i Qvalitet af Publikum.

Hvad er Formløshed? Det er den ophævede lidenskabelige Distinktion mellem Form og Indhold, der derfor vel, til Forskjel fra Galskab og Dumhed, kan have Sandhed til Indhold, men det Sande, den indeholder, kan aldrig være væsentligt sandt. Extensivt vil den kunne udbrede sig altomfattende eller altantastende, i Modsætning til det væsentlige Indhold, der intensivt i Selvfordybelsen har sin Bestemmethedens, om man saa vil, fattige Begrændsning.

Formløshedens Almindelighed i en lidenskabsløs men reflekteret Tid udtrykker sig forøvrigt, foruden ved den leflende Omgang mellem det Forskjelligste med det Forskjelligste, ved sin Modsætning: en overveiende Hang og Lyst til at handle »for Principets Skyld.« Principium er, som Ordet siger, det Første det vil sige det Substantielle, Ideen i Følelsens, Begeistringens uaabnede Form, der ved sin indre Drift fremskynder Individet. Dette mangler den Lidenskabsløse; for ham bliver Principet et Udvortes, for hvis Skyld han gjør baade det Ene og det Andet og det Modsatte. Den Lidenskabsløses Liv er ikke et sig aabenbarende og udfoldende Princip, tvertimod er hans indre Liv et gesvindt Noget, der ideligen er paa Farten og jager efter at gjøre Noget »for Principets Skyld«. Principet i denne Forstand bliver et uhyre Noget, et Abstraktum ligesom Publikum. Og medens Publikum er saadan noget uhyre Noget, at selv alle Nationer opstillede paa eengang og selv alle Sjele i Evigheden ikke ere saa talrige som Publikum, og dog enhver selv den fulde Matros har Publikum, saa er det ogsaa saaledes med Principet. Det er noget uhyre Noget, som selv det ubetydeligste Menneske øger til den ubetydeligste Handling, og bliver sig selv uendelig vigtig. Et skikkeligt, ubetydeligt Menneske bliver pludseligen en Heros for »Principets Skyld«, og Forholdet frembringer i Grunden en lige saa latterlig Virkning, som hvis en Mand, eller hvis det blev Mode, Enhver gik med en Kaskjet, hvis Skygge var 30 Alen lang. Dersom en Mand »for Principets Skyld« vilde lade sye en lille Knap for Brystlommen i sin Frakke: saa vilde denne ubetydelige og ret hensigtsmæssige Forsigtighedsregel pludseligen faae en uhyre Betydning – det var ikke usandsynligt, at der blev stiftet et Selskab i den Anledning.

Netop Dette: »for Principets Skyld« hæver igjen et Decorums lidenskabelige Disjunktion; thi Decorum ligger, som viist blev, i Umiddelbarheden (den første eller den vundne), i Følelsen, i Begeistringens indre Drift og indre Consequents med sig selv. For Principets Skyld kan man gjøre Alt og det bliver i Grunden ligegyldigt, ligesom Ens Existents bliver ubetydelig, selv om man for Principets Skyld understøtter Alt, hvad der kaldes Tidens Fordring, selv om man derved, ved at være Statist, bliver lige saa bekjendt i denne Qvalitet af Träger der öffentlichen Meinung, som hine Personer i Lirekassen, der komme frem og bukke med Tallerkenen. For Principets Skyld kan man gjøre Alt, deeltage i Alt og selv være et umenneskeligt Ubestemt. En Mand kan for Principets Skyld interessere sig i, at et Bordel anlægges (der er jo anstillet mange statsborgerlige Betragtninger af Sundheds-Politiet desangaaende); og den samme Mand kan for Principets Skyld interessere sig for en ny Psalmebog, der skal være Tidens Fordring. Og saa uretfærdigt som det vilde være af det Første at slutte, at Manden var udsvævende, saa overilet vilde det maaskee være af det Sidste at slutte, at han vilde læse eller synge i Psalmebogen. Paa den Maade bliver Alt tilladeligt for Principets Skyld, og ligesom Politiet »paa Embeds-Vegne« kommer mange Steder, hvor ingen Anden kommer, og ligesom man med Hensyn til Politi-Embedsmændenes Person Intet er berettiget til at slutte af deres Tilstedeværelse, saaledes kan man gjøre Alt for Principets Skyld, og undgaae ethvert personligt Ansvar. Man nedriver det, man selv beundrer, »for Principets Skyld«, hvilket er Nonsens; thi det Fødende, det Skabende er altid latent polemisk, fordi det skal have Plads, men det Nedrivende er jo Intet, og et Nedrivnings-Princip er Tomhed, hvad skal det saa med Plads? Undseelsen eller Fortrydelsen eller Ansvaret kan imidlertid vanskeligere faae fat paa en saadan Adfærd; thi det var jo for Principets Skyld.

Hvad er Overfladiskhed, og dennes Lyst: »Repræsentations-Lysten«? Overfladiskhed er den ophævede lidenskabelige Distinktion mellem Skjulthed og Aabenbarelse, den er en Aabenbarelse af Tomhed, der dog extensivt vinder Gøgleriets skuffende Fordeel over den væsentlige Aabenbarelse, der har Fordybelsens eensartede Væsentlighed, medens Overfladiskhed har det mangehaande eller allehaande Skin. Og Repræsentations-Lysten er Reflexions-Indbildningens Selvforgabelse. Inderlighedens Skjulthed faaer ikke Tid til at sætte et Væsentligt af, der kan vorde en Aabenbarelse, men plumres længe før Tiden, og til Vederlag trækker den selviske Reflexion Alles Øine om muligt hen paa dette mangfoldige Skin, som Frue Commerceraadinde Waller gjør det.

Hvad er Leflerie? Det er den ophævede Lidenskabelige Distinction mellem væsentligen at elske, og væsentligen at være udsvævende. Hverken den væsentlige Elsker eller den væsentlige Udsvævende forskylder Lefleriet, der gantes med Muligheden. Lefleriet er derfor en Indulgents i at turde tangere det Onde, og et Aflad fra at realisere det Gode. Det at handle for Principets Skyld er saaledes ogsaa Leflerie, fordi det forvansker den sædelige Handlen til en Abstraktion. Men extensivt har Lefleriet Fordelen, thi der kan lefles med alt Muligt, og væsentligen kan kun een Pige elskes, og erotisk rigtigt forstaaet (selv om Lysten i Forvildelsens Tid forblinder den Flygtige), er al Adderen en Subtraheren, jo Flere En føier til jo mere trækker han fra.

Hvad er det at raisonere? Det er den ophævede lidenskabelige Disjunktion mellem Subjektivitet og Objektivitet. Som abstrakt Tænkning er Raisonnementet ikke dialektisk dybt nok, som Mening og Overbeviisning er det uden Individualitets Fuldblodighed. Men extensivt gaaer den Raisonnerende af med Skin-Fordelen; thi en Tænker kan omfatte sin Videnskab, en Mand kan have en Mening om hvad der hører til et bestemt Fag, kan have en Overbeviisning i Kraft af en bestemt Livs-Anskuelse, men den Raisonnerende raisonnerer om alt Muligt.

Anonymitet har i vor Tid en langt prægnantere Betydning end man maaskee tænker paa, den har næsten epigrammatisk Betydning. Man ikke blot skriver anonymt, men man skriver anonymt med Navn under, ja man taler anonymt. Som en Forfatter concentrerer hele sin Sjel i Stilen, saaledes concentrerer et Menneske væsentligen sin Personlighed i Talen, dog maa dette forstaaes med den indskrænkende Undtagelse, som Claudius i et lignende Forhold minder om, idet han siger, at naar man besværger en Bog, saa kommer Aanden ud – med mindre der ingen Aand er. Nuomstunder kan man virkeligen tale med Mennesker, og man maa sige, at deres Yttringer ere overmaade forstandige, medens Samtalen dog gjør det Indtryk, som talte man med en Anonym. Det samme Menneske kan sige det meest Modsatte, kan roligt yttre, hvad der i hans Mund er den bittreste Satire over hans egen Existents, Yttringen selv er meget forstandig, lod sig godt høre paa en Generalforsamling, for at komme med i en Discussion, ved hvilken der fabrikeres Noget, ligesom der paa Fabrikken bliver Papiir af Klude. Men alle disse mange Yttringer udgjøre dog ikke tilsammen en personlig-menneskelig Tale, saaledes som denne kan føres selv af den Eenfoldigste, der kun kan tale om meget Lidt, men dog taler. Frue Commerceraadinde Wallers Yttring om det Dæmoniske er meget sand, og dog gjør hun i Yttringens Øieblik netop det Indtryk paa En, som var hun en Anonym, et Positiv, en Spilledaase. Yttringerne blive saa objektive, Yttringernes Omfang saa Alt omfattende, at det tilsidst bliver aldeles tilfældigt, hvem der siger det, et Forhold, der i Henseende til det at tale menneskeligt aldeles svarer til at handle for Principets Skyld. Og ligesom Publikum er en reen Abstraktion, saa bliver tilsidst den menneskelige Tale det ogsaa, der bliver Ingen mere som taler, men en objektiv Reflexion afsætter efterhaanden et atmosphærisk Noget, en Abstraktions-Lyden, der vil gjøre den menneskelige Tale overflødig, ligesom Maskinerne gjøre Arbeiderne overflødige. I Tydskland har man endog Haandbøger for Elskende, saa ender det vel med, at et elskende Par kan sidde og tale anonymt med hinanden. I Alt har man Haandbøger, og Dannelsen i Almindelighed bestaaer snart i at være perfectioneret i et større eller mindre Indbegreb af saadanne Haandbøgers Betragtninger, og man excellerer i Forhold til sin Færdighed i at tage det Enkelte frem, ligesom Sætteren tager Lettere frem.

Saaledes er Nutiden væsentlig forstandig, den er maaskee i Gjennemsnit vidende som ingen Generation før har været det, men den er lidenskabsløs. Enhver veed Meget, vi vide Alle, hvilken Vei der skal gaaes og de mange Veie der kan gaaes, men Ingen vil gaae. Overvandt endeligen En Reflexionen i sig selv og kom til at handle, saa vilde i samme Øieblik tusinde Reflexioner udenfra danne ham en Modstand, thi kun Forslag til yderligere Overveielse modtages med opblussende Begeistring, Handling med Indolents. Nogle vilde i fornem Selvtilfredshed finde de Handlendes Begeistring latterlig, Andre vilde blive misundelige, fordi det var ham, der gjorde Begyndelsen, da de jo dog lige saa godt som han vidste, hvad der skulde gjøres – men dog ikke havde gjort det. Andre vilde benytte den Omstændighed, at der dog var En, der havde handlet, til at faae afsat en Mængde critiske Betragtninger, afhændet et Lager af Raisonnements om, hvor langt forstandigere der kunde være handlet; Andre vilde faae travlt med at gjette Udfaldet, og med om muligt at influere lidt paa Foretagendet i Retning af deres Hypothese. Man fortæller om to engelske Lords, der kom ridende, at de mødte en forulykket Rytter, som var ifærd med at falde af Hesten, der styrtede afsted i Spring, medens Rytteren skreeg om Hjælp. Den ene Lord kastede et Blik til den anden og sagde: »Hundrede Guineer, han falder af;« »antaget,« svarede den anden. Derpaa satte de deres Væddeløbere i Carriere og ilede forud for at faae Bommene oplukkede og alle Hindringer ryddede af Veien. Saaledes vilde Nutidens Forstandighed kun med mindre Millionair-Spleen-Heroisme være en Personifikation, der nysgjerrig, critisk og verdensklog i Maximum vilde have Lidenskab til at indgaae et Væddemaal. Livets existentielle Opgaver have tabt Virkelighedens Interesse, ingen Illusion omfreder Inderlighedens guddommelige Væxt, der modner til Afgjørelser; den Ene er nysgjerrig efter den Anden, alle vente ubesluttede og forfarne i Udflugter paa, at der skal komme En, der vil Noget – for saa at parere paa hans Haand.

Og da der nu i vor Tid, hvor der gjøres saa lidt, gjøres saa overordentligt Meget i Prophetier, Apokalypser, Fingerpeeg, Indblik i Fremtiden: saa gaaer det vel ikke an Andet end at være med, om jeg dog har den Ugenerethedens Fordeel fremfor de Manges byrdefulde Ansvar med at prophetere og varsle, at jeg kan være sikker paa, at Ingen faaer det Indfald at troe det. Min Fordring er det derfor heller ikke, at Nogen skulde endog blot sætte et Kryds i Almanaken, eller uleilige sig med at lægge Mærke til, om det gaaer i Opfyldelse; thi gaaer det i Opfyldelse, saa vil han faae Andet at tænke paa end min Tilfældighed, og skeer det ikke, nu ja, saa bliver jeg dog en Prophet i moderne Forstand, thi en Prophet i moderne Forstand er En, der forudsiger Noget, videre ikke. Og det forstaaer sig, mere kan en Prophet i en vis Forstand heller ikke gjøre. Styrelsen var det jo, der føiede Opfyldelsen til de ældre Propheters Udsagn, vi moderne Propheter kunde maaskee i Mangel af Styrelsens Tilføiende med Thales tilføie: hvad vi forudsige, det vil enten skee eller det vil ikke skee, thi ogsaa os har Guden forundt Prophetiens Gave.

Det er da saa langt som muligt fra, at Socialitetens, Menighedens Idee vil blive Tidens Frelse, at den tvertimod er den Skepsis, der maa til, for at Individualitets-Udviklingen ret kan gaae for sig, idet hvert Individ enten fortabes eller, optugtet ved Abstraktionen, religieust vinder sig selv. Associations-Principet (der i det Høieste kan have sin Gyldighed i Forhold til de materielle Interesser) er i vor Tid ikke affirmativt men negativt, er en Udflugt, en Adspredelse, et Sandsebedrag, hvis Dialektik er: idet det styrker Individerne enerverer det dem, det styrker ved det Numeriske ved Sammenholdet, men dette er, ethisk, en Svækkelse. Først naar det enkelte Individ i sig selv har vundet ethisk Holdning trods hele Verden, først da kan der være Tale om i Sandhed at forene sig, ellers bliver Foreningen af de, hver i sig, Svage noget lige saa uskjønt og fordærveligt, som at Børn gifte sig. Fordum havde Herskeren, den Udmærkede, de Fremragende hver en Mening, de Andre havde den Sluttethed og Afgjorthed, at de ikke turde eller ikke kunde have en Mening; nu kan Enhver have en Mening, men de maae slaae sig numerisk sammen om at have den. 25 Underskrifter paa det Taabeligste er en Mening; det eminenteste Hoveds grundigst gjennemtænkte Mening er et Paradox. Den offentlige Mening er et uorganisk Noget, en Abstraktion. Men naar Contexten er bleven meningsløs, saa hjælper det ikke med store Overskuelser, saa gjør man bedst i at tage Talens enkelte Dele frem; naar Munden slaaer Sladder, saa hjælper det ikke at ville holde et sammenhængende Foredrag, men man gjør bedst i at tage hvert Ord for sig – og saaledes ogsaa med Individernes Forhold.

Omforandringen vil da ogsaa være denne. Medens i de ældre Formationer (af Forholdet mellem Generation og Individ) Underofficiererne, Officiererne, Compagnicheferne, Generalerne, Helten (det vil sige de Udmærkede, de Fremragende efter deres forskjellige Grad, Lederne) vare kjendelige, og hver (i Forhold til sin Myndighed) med sit lille Detachement malerisk og organisk ordnede sig ind under det Hele, selv understøttet af og understøttende det Hele: saa vil nu de Udmærkede, Lederne (efter deres respective Grad) være uden Myndighed, netop fordi de guddommeligen have forstaaet Nivelleringens diaboliske Princip, de ville være ukjendelige, som Politiet er civilt, og bære deres respektive Distinktion skjult, og kun understøtte negativt det vil sige frastødende, medens Abstraktionens uendelige Ligelighed dømmer hvert Individ, examinerer det i dets Isolation. Denne Formation er lige den dialektiske Modsætning til Propheter og Dommere; og som disses Fare var, ikke at blive respekteret i Forhold til deres respektive Myndighed, saa er de Ukjendeliges Fare, at blive kjendte, at blive forlokkede til at faae Anseelse og Betydning som Auctoritet, hvorved de vilde standse den høieste Udvikling. De ere nemlig ukjendelige, eller ligesom geheime Agenter ikke i Følge af en privat Instrux fra Guddommen, thi dette er jo netop Propheternes, Dommernes Forhold, men de ere ukjendelige (uden Myndighed) i Følge af selv at have fattet det Almene i Ligeligheden for Gud, og i Følge af at de hvert Øieblik fatte dette i Ansvaret, hvilket forhindrer dem fra i Distraction at forskylde Form-Realisationens Inconsequents i Forhold til Anskuelsens Consequents. Denne Formation er den dialektiske Modsætning til den organiserende, der gjør Generationen, præformeret i de Udmærkede, til Understøttende for Individerne, da den nu som en Abstraktion, understøttet negativt af de Ukjendelige, vender sig polemisk mod Individerne – for at frelse Hver især religieust.

Og naar saa Generationen, der jo selv har villet nivellere, har villet emanciperes, og revoltere, har villet sløife Myndigheden og derved selv, i Associationens Skepsis, har afstedkommet Abstraktionens trøstesløse Skovbrand, naar Generationen ved Nivelleringen gjennem Associationens Skepsis har faaet Individualiteterne bort og alle de organiske Concretioner, faaet Menneskeheden isteden og Talforholdenes Ligelighed mellem Mand og Mand; naar Generationen saa et Øieblik har forlystet sig ved den abstrakte Uendeligheds vide Udsigt, som ingen Fremragenhed, end ikke den mindste, begrændsende forstyrrer, men der kun er »idel Luft og Hav«: saa begynder Arbeidet, da Individerne skulle hjælpe sig selv, hver især. Thi det skal ikke være som fordum, at Individerne, naar det begyndte at svimle lidt for deres Øine, saae hen til den nærmest staaende Udmærkede, for at orientere sig. Det er nu forbi; de skulle enten fortabes i den abstrakte Uendeligheds Svimmel, eller frelses uendeligt i Religieusitetens Væsentlighed. Mange Mange ville maaskee skrige i Fortvivlelsen, men det skal ikke hjælpe, nu er det for sildig. Er Myndigheden og Magten fordum bleven misbrugt i Verden, og har den bragt en Revolutions-Nemesis over sig, saa var det jo Afmagten og Svagheden, der selv vilde staae paa egne Been, og derfor nu bragte denne Nemesis over sig. Og ingen af de Ukjendelige tør understaae sig at hjælpe ligefrem, at yttre sig ligefrem, at lære ligefrem, i Spidsen for Mængden at gaae i Afgjørelsen (istedenfor negativt understøttende at hjælpe Individet ud i Afgjørelsen, hvor han selv er); dette vilde være hans Afsked, fordi han vilde fuske i menneskelig Medlidenheds kortsynede Opfindsomhed istedenfor at lystre Guddommens Ordre, den vrede, men dog saa naaderige Guddom, thi Udviklingen er dog et Fremskridt, fordi alle de Individer, som frelses, vinde Religieusitetens specifike Vægtfylde, vinde dens Væsentlighed paa første Haand fra Gud. Da vil det hedde: »see Alt er beredt, see Abstraktionens Grusomhed gjør Endeligheden aabenbar i sin Skuffelse som saadan, see det Uendeliges Afgrund aabner sig, see Nivelleringens skarpe Lee lader Alle, hver især, springe over Klingen – see Guden venter! Saa spring da til i Guddommens Favn.« Men om det var den meest Betroede af de Ukjendelige og om det var den Qvinde, der havde baaret ham under sit Hjerte, og om det var den Pige, han glad offrede Livet for, han skal ikke turde og ikke kunne hjælpe dem, de maae selv gjøre Springet, og Guddommens uendelige Kjerlighed skal ikke blive dem et Forhold paa anden Haand. Og dog skulle de Ukjendelige (i Forhold til deres respektive Grad) have dobbelt Arbeide i Sammenligning med de Udmærkede (af samme Grad) i en ældre Formation; thi de Ukjendelige maae bestandigt arbeide – og tillige arbeide for at skjule det.

Men Nivelleringens trøstesløse Abstraktion vil bestandigt fortsættes ved sine Tjenere, for at forhindre, at det dog ikke skal ende med, at en ældre Formation kommer frem igjen. Disse Nivelleringens Tjenere ere den onde Magts Tjenere, thi Nivelleringen selv er ikke af Guddommen, og ethvert godt Menneske vil have Øieblikke, hvor han kunde græde over dens Trøstesløshed, men Guddommen tillader det, og vil med Individerne det vil sige med Hver især bringe det Høieste ud. Nivelleringens Tjenere ere kjendte af de Ukjendelige, men Magt eller Myndighed mod dem tør de Ukjendelige ikke bruge, thi da gaaer Udviklingen tilbage, fordi det i samme Øieblik vilde være aabenbart for en Trediemand, at den Ukjendelige var en Auctoritet, og saa var jo denne Trediemand forhindret i det Høieste. Kun ved en lidende Handling vil den Ukjendelige turde fremhjælpe Nivelleringen, og ved den samme lidende Handling dømme Redskabet. Han tør ikke overvinde Nivelleringen ligefrem, dette vilde være hans Afskeed, da det var at handle i Retning af Auctoritet, men han vil overvinde den lidende, og derved atter udtrykke Loven for sin Existents, som ikke er at herske, styre, lede, men lidende at tjene, at hjælpe indirecte. De, der ikke have gjort Springet, ville forstaae den Ukjendeliges lidende Handling som hans Nederlag, og de, der have gjort Springet, ville have en Forestilling om, at det var hans Seier, men de ville ingen Vished have, thi Visheden kunde de kun faae ved ham, og giver han et eneste Menneske den ligefrem: saa betyder det hans Afskeed, fordi han blev Guddommen utro og vilde spille Auctoritet, fordi han ikke lystrende Guddommen lærte at elske Menneskene uendeligt ved at tvinge sig selv, ikke svigefuldt ved herskende at tvinge dem, selv om de bade derom.

\workheading{Opbyggelige Taler i forskjellig Aand}{1847}

\worknote{Antologiens ''Upbuilding Discourses in Various Spirits''. Hong-udvalget tager begyndelsen af Første Deel, ''Ved Anledningen af et Skriftemaal'' (kendt som ''Purity of Heart'' / Hjertets Reenhed): tilegnelsen til hiin Enkelte, forordet, bønnen, temaet (efter Jakobs Brev 4,8) og det første afsnit (''Skal det være muligt … maa han ville det Gode''). Selve skriftemaals-betragtningen mellem bøn og tema samt resten af talen er udeladt; spring markeret med […].}

\sectionheading{Ved Anledningen af et Skriftemaal: Hjertets Reenhed er at ville Eet (uddrag)  ·  SV¹ VIII, 116–143}

»Hiin Enkelte« helliges dette lille Skrift.

Forord.

Skjøndt denne lille Bog (Leilighedstale kan den kaldes, uden dog at have Leiligheden, som gjør Taleren og gjør ham myndig, eller Leiligheden, der gjør Læseren og gjør ham lærende) i Virkelighedens Forhold er som en Indbildning og liig en Drøm om Dagen: er den dog ikke uden Tillid og ikke uden Haab om Opfyldelsen. Den søger hiin Enkelte, til hvem den ganske giver sig hen, af hvem den ønsker at modtages som var den opkommen i hans eget Hjerte; hiin Enkelte, hvem jeg med Glæde og Taknemlighed kalder min Læser; hiin Enkelte, der i Villighed læser langsomt, læser gjentagent, og som læser høit – for sin egen Skyld. Finder den ham, da er Forstaaelsen fuldkommen i Adskillelsens Fjernhed, naar han beholder den og Forstaaelsen for sig selv i Tilegnelsens Inderlighed.

Naar en Qvinde forarbeider et Klæde til helligt Brug, da virker hun hver en Blomst, om muligt, skjøn som Markens Deilige, hver en Stjerne, om muligt, funklende som Nattens Blink; da sparer hun Intet, men anbringer det Kostbareste, hun eier; da sælger hun hver en anden Fordring til Livet for at kunne kjøbe den uafbrudte Dagens og Nattens beleilige Tid til sit eneste, sit elskede Arbeide. Men naar da Klædet er færdigt og anbragt efter sin hellige Bestemmelse: da bedrøves hun inderligen, om Nogen saae feil og saae paa hendes Kunst istedenfor paa Klædets Betydning, eller saae feil og saae en Mangel istedenfor at see paa Klædets Betydning. Thi den hellige Betydning kunde hun ikke virke ind med i Klædet, hun kunde og ikke sye den paa som en Prydelse mere. Denne Betydning ligger nemlig i Betragteren og i Betragterens Forstaaelse, naar han, i Adskillelsens uendelige Fjernhed, over sig selv og over sit eget Selv har uendeligt glemt Syersken og hvad hendes er. Det var tilladeligt, det var sømmeligt, det var Pligt, det var en kjær Pligt, det var Glædens Høieste for Syersken at gjøre Alt for at gjøre Sit; men det var en Brøde mod Gud, en krænkende Misforstaaelse mod den stakkels Syerske, om Nogen saae feil og saae hvad der kun kan være med for at oversees, hvad der kun kan være med – ikke for at tildrage sig Opmærksomheden, men tvertimod for at Savnet deraf ikke skal forstyrrende tildrage sig Opmærksomheden.

S. K.

Ved Anledningen af et Skriftemaal.

Fader i Himlene! Hvad er dog et Menneske uden Dig! Hvor er Alt hvad han veed, var det end det Mangfoldiges Mængde, kun et afbrudt Stykke, hvis han ikke kjender Dig; hvor er al hans Stræben, var den end omfattende en Verden, kun et halvgjort Arbeide, hvis han ikke kjender Dig, Du den Ene, som er Eet og som er Alt! Saa give Du i Forstanden Viisdom til at fatte det Ene, i Hjertet Oprigtighed til at annamme Forstaaelsen, i Villien Reenhed ved kun at ville Eet; saa give Du i gode Dage Vedholdenhed til at ville Eet, i Adspredelser Samling til at ville Eet, i Lidelser Taalmod til at ville Eet. O, Du som giver baade at begynde og at fuldkomme, Du give aarle, naar Dagen gryer, den Unge Beslutningen at ville Eet; naar Dagen helder give Du Oldingen en fornyet Ihukommelse af sin første Beslutning, at det Sidste maatte være som det Første, det Første som det Sidste, være Dens Liv, der kun har villet Eet. Ak, men det er jo ikke saaledes; der kom jo Noget derimellem, der ligger jo en Syndens Skilsmisse derimellem: der lægger sig jo daglig og Dag fra Dag Noget derimellem: Forsinkelsen, Standsningen, Afbrydelsen, Vildfarelsen, Fortabelsen. Saa give Du i Angeren Frimodighed til atter at ville Eet. Vel er det en Afbrydelse fra den sædvanlige Gjerning, vel er det en Standsning fra Arbeidet, som var det en Hviledag, naar den Angrende og kun i Angeren besværet Arbeidende er stille i Syndens Bekjendelse, er ene for Dig i Selvanklagen, o, men det er jo en Afbrydelse, som søger tilbage til sin Begyndelse, at den maatte paa ny binde mellem det Adskilte, at den i sin Sorg maatte indhente det Forsømte, at den i sin Bekymring maatte fuldkomme hvad der ligger foran. O, Du som giver baade at begynde og at fuldkomme, Du give Seier paa Nødens Dag, at hvad der ikke lykkedes den i Ønsket Brændende, ikke den i Forsættet Besluttede, at det dog maatte lykkes den i Angeren Bedrøvede: kun at ville Eet.

[…]

Saa lad os da ved Anledningen af et Skriftemaal tale over denne Sætning:

Hjertets Reenhed er at ville Eet,

idet vi til Grund for Betragtningen lægge Apostelen Jakobs Ord, i hans Brev 4de Capitel, Vers 8:

Holder Eder nær til Gud, saa skal han holde sig nær til Eder! Renser Hænderne, I Syndere, og luttrer Hjerterne I Tvesindede;

thi kun de Rene af Hjerte kunne see Gud og altsaa holde sig nær til ham, og bevare Reenheden ved at han holder sig nær til dem; og Den, der i Sandhed kun vil Eet, han kan kun ville det Gode, og Den der kun vil Eet, naar han vil det Gode, han kan kun ville det Gode i Sandhed.

Lad os tale derom, men lad os først glemme Anledningen for at enes i Forstaaelsen af Sætningen, og i hvad det apostoliske formanende Ord (»luttrer Hjerterne I Tvesindede«) modsætter sig: Tvesindetheden; da ville vi til Slutning nærmere benytte Anledningen.

I. Skal det være muligt, at et Menneske kan ville Eet, da maa han ville det Gode.

Kun at ville Eet – men maa dette nu ikke blive en vidtløftig Tale; hvis Nogen ret vil betænke denne Sag, maa han da ikke først Stykke for Stykke gjennemgaae ethvert Maal, som et Menneske kan sætte sig i Livet, nævne alt det Meget enkeltviis, som et Menneske kan ville? Og ikke nok hermed, maa han ikke dernæst, da dog enhver saadan Betragtning let bliver af for almindelig Art, maa han ikke forsøge sig i at ville det Ene efter det Andet, for at udfinde, hvilket det Ene er, han kan ville, naar det gjelder kun at ville Eet? Ja, hvis Nogen saaledes vilde begynde, da skulde han vel aldrig blive færdig, eller rettere, hvor var det muligt, at han kunde blive færdig, da han netop begyndte med at gaae feil af Veien, og dog blev ved at gaae frem og gaae frem ad Vildfarelsens Vei, der jo kun paa en sørgelig Maade fører til det Gode – naar Vandreren nemlig vender sig om og gaaer tilbage; thi som det Gode kun er Eet, saa føre og alle Veie til det Gode, selv Vildfarelsens – naar den Omvendte gaaer samme Vei tilbage.

[…]

At ville Eet kan da ikke betyde at ville Det, som kun tilsyneladende er Eet. Det Verdslige er nemlig i sit Væsen ikke Eet, da det er det Uvæsentlige; dets saakaldede Eenhed er ingen væsentlig men en Tomhed, over hvilken det Mangfoldige skjuler. Det Verdslige er derfor i Skuffelsens stakkede Øieblik det Mangfoldige og saaledes ikke Eet; derpaa forvandler det sig til sit Modsatte – saa langt er det fra at være og blive Eet. Eller hvad er Lysten i sit grændseløse Yderste vel Andet end Væmmelse? Hvad er jordisk Ære i sit svimlende Høieste vel Andet end Foragt for Tilværelsen? Hvad er Riigdommens høieste Overmaal Andet end Armod, thi fylder vel al Verdens Guld gjemt i Pengebegjerlighed saa meget, eller ikke uendelig meget mindre end den ringeste Skjærv gjemt i den Fattiges Nøisomhed! Hvad er verdslig Almagt Andet end Afhængighed; hvilken Slave i Lænker var saa ufri som en Tyran! Nei, det Verdslige er ikke Eet; mangfoldigt som det er, forvandles det i Livet til sit Modsatte, i Døden til Intet, i Evigheden til Forbandelse for Den, som har villet dette Ene. Kun det Gode er Eet i sit Væsen og det Samme i hver sin Yttring. Lad Kjærlighed oplyse det. Den, der i Sandhed elsker, han elsker ikke engang for alle, ei heller bruger han en Deel af sin Kjærlighed og saa igjen en Deel; thi at bytte den er at forbytte den. Nei han elsker med hele sin Kjærlighed, den er heel tilstede i enhver Yttring, han giver den bestandigt heel ud, og dog beholder han den bestandig heel i sit Hjerte. Forunderlige Riigdom! Naar den Gjerrige har samlet al Verdens Guld – i Smudsighed, da er han bleven fattig: naar den Elskende giver sin hele Kjærlighed ud beholder han den heel – i Hjertets Reenhed. – Skal et Menneske i Sandhed ville Eet, da maa jo det Ene, han vil, være et Saadant, at det bliver uforandret i alle Forandringer, saa han ved at ville det, kan vinde Uforanderlighed; forandres det ideligen, da vorder han selv foranderlig, tvesindet og ustadig. Men denne idelige Foranderlighed er jo netop Ureenheden.

Men det at ville Eet er ei heller hiin den formastelige, ugudelige Begeistrings kraftige Vildfarelse: at ville det Store, ligegyldigt enten det er godt eller ondt. Ogsaa en saadan Villende, om han end er nok saa fortvivlet, er dog tvesindet; eller er Fortvivlelse ikke netop Tvesindethed; eller hvad er det at fortvivle Andet end at have to Villier! Thi hvad enten han, den Svage, fortvivler over, at han ikke kan løsrive sig fra det Onde; eller han, den Formastelige, fortvivler over, at han ikke ganske kan løsrive sig fra det Gode: de ere begge Tvesindede, de have begge to Villier, Ingen af dem vil i Sandhed Eet, hvor fortvivlet de end synes at ville det. Hvad enten det var en Qvinde, hvem Ønsket bragte til Fortvivlelse, eller det var en Mand, der fortvivlede i Trods; hvad enten et Menneske fortvivlede – fordi han fik sin Villie, eller han fortvivlede fordi han ikke fik sin Villie: enhver Fortvivlet har to Villier, een som han forgjeves vil følge ganske, og een som han forgjeves ganske vil være qvit. Saaledes har Gud, bedre end nogen Konge, sikkret sig mod ethvert Oprør. Thi det er vel skeet, at en Konge ved et Oprør er bleven stødt fra Thronen: men enhver Oprører mod Gud bringer det i det Høieste til selv at fortvivle. Fortvivlelsen er Grændsen – hertil og ikke videre! Fortvivlelsen er Grændsen, her mødes den feige frygtagtige Selvkjærligheds Arrigskab, og det stolte trodsige Sinds Formastelighed, her mødes de i lige Afmagt.

Kun altfor snart lærer den egne Erfaring og Erfaringen paa Andre, hvor langt dog de fleste Menneskers Liv er fra at være hvad et Menneskeliv burde være. Alle have de store Øieblikke, de see sig selv i Mulighedens Tryllespeil, som Haabet holder for dem medens Ønsket smigrer; men hurtigt glemme de Synet i det Dagligdagse. Eller de tale maaskee begeistrede Ord, »thi Tungen er et lidet Lem, dog den pukker storligen« – men Talen tager jo netop Begeistringen forfængelig ved høirøstet at forkynde hvad der i Taushed skal udøves, og hurtigt glemmes de begeistrede Ord i Livets Ubetydelighed; det glemmes, at det var om dette Menneske Sligt sagdes, det glemmes, at det var ham selv der sagde det. Saa vaagner maaskee engang Erindringen med Forfærdelse, og Fortrydelsen synes at ville give ny Kraft: ak, dette blev ogsaa kun et stort Øieblik. De have Alle Forsætter, Planer, Beslutninger for Livet, ja for Evigheden; men Forsættet taber hurtigt Ungdommens Kraft og bliver affældigt; og Beslutningen staaer ikke fast og staaer ikke imod, den vakler og forandres med Omstændighederne; og Hukommelsen slaaer feil – indtil de ved Vane og Omgang lære, som man siger, at trøste sig med hinanden, indtil de endog, istedetfor at finde det forrædersk, finde det opbyggeligt, om Nogen forkynder Undskyldningernes matte Trøstegrunde, der opmuntrende styrke til Dorskhed. Thi der gives Mennesker, som finde det opbyggeligt, at Fordringen gjøres gjeldende i al sin Ophøiethed, i al sin Strenghed, saa den fordrende trænger ind i Sjelens Inderste; Andre finde det opbyggeligt, at der træffes en Overeenskomst i Jammerlighed med Gud og Fordringen – og Sproget. Der gives Mennesker, som finde det opbyggeligt, om dog Nogen vil kalde paa dem; men der ere ogsaa søvnige Sjele, som kalde det, ikke behageligt, men endog opbyggeligt, at der lulles for dem.

Vel er dette nu sørgeligt; men da er der en Viisdom, som ikke er herovenfra, den er jordisk og kjødelig og djævelsk. Den har opdaget denne almindelige menneskelige Svaghed og Dorskhed; den vil hjælpe. Den seer, at det er Villien, det kommer an paa, og forkynder nu høit: »uden at ville Eet bliver dog et Menneskes Liv jammerlig Middelmaadighed, ussel Elendighed; Eet maa han ville, ligegyldigt enten det er godt eller ondt; Eet maa han ville, deri ligger et Menneskes Storhed.« Dog er det ikke vanskeligt at gjennemskue denne kraftige Vildfarelse. Den hellige Skrift lærer til Saliggjørelse, at Synden er Menneskets Fordærvelse, og Frelsen derfor kun i Reenhed ved at ville det Gode. Hiin jordiske og djævelske Kløgt forvender dette til fristende Fortabelse: Svagheden er et Menneskes Ulykke, Kraften det eneste Frelsende. »Naar den urene Aand farer ud af et Menneske, vandrer den gjennem øde og tomme Steder, men finder ikke Hvile; da vender den igjen tilbage, og har i Selskab med sig« hiin urene Kløgt, Ørkenens og de tomme Steders Viisdom, hiin urene Kløgt, der nu udjager Dorskhedens og Middelmaadighedens Aand – »saa det Sidste bliver værre end det Første«. Hvorledes skal man beskrive et saadant Menneskes Væsen? Der siges om Sangeren, at han kan sprenge sin Stemme ved at overskrige sig selv: saaledes er et saadant Menneskes Væsen sprengt ved at overskrige sig selv og Samvittighedens Stemme. Der siges om Den, som svimlende staaer paa et høit Sted, at Alt løber ud i Eet for hans Øie: saaledes har et saadant Menneske svimlet sig ud i det Uendelige, hvor Det løber i Eet, hvad der for evig er adskilt, saa kun det Store bliver, det er det Øde og Tomme, der jo altid føder Svimmelhed. Men tvesindet er dog et saadant Menneske, hvor fortvivlet han end synes at ville Eet. Naar han, den Selvraadige, maatte raade, da skulde der kun være Eet, han være den Eneste, der ikke var tvesindet, han den Eneste, der havde afkastet enhver Lænke, han den eneste Fri. Men Fri er – Syndens Træl dog ikke; afkastet har han heller ikke Lænken, »fordi han spotter den«; i Tvang er han og derfor tvesindet; og raade maa han nu engang ikke. Der er en Magt, som tvinger ham; løsrive sig fra den kan han ikke, ja han kan end ikke ganske ville det, ogsaa denne Magt er ham negtet. Dersom Du, min Tilhører, saae et saadant Menneske, det er vel sjeldnere, som jo Svaghed og Middelmaadighed unegteligt er det Almindeligere; dersom Du traf ham i hvad han selv vilde kalde et svagt Øieblik, ak, hvad Du maatte kalde et bedre Øieblik – naar han ingen Hvile havde fundet i det Øde, naar Svimlen et Øieblik gik over, og han fornam en qvalfuld Længsel efter det Gode; naar han rystet i sit Inderste ikke uden Veemod tænkte paa den Eenfoldige, der om end i Skrøbelighed dog vil det Gode: da opdagede Du, at han havde to Villier, og hans qvalfulde Tvesindethed. Fortvivlet, som han var, meente han: fortabt er fortabt, men han kunde ikke lade være end eengang at vende sig om i Længsel efter det Gode, hvor rædsomt han end selv havde forbittret sig denne Længsel, en Længsel, der beviser, at, som et Menneske dog i al sin Trods ikke har Magt til ganske at løsrive sig fra det Gode, fordi det er det Stærkere, saa har han end ikke Magt til ganske at ville det. Maaskee hørte Du endog hiin Fortvivlede sige: »der gaaer dog noget Godt under i mig«. Naar et Menneske finder sin Død i Bølgerne, da synker han, uden dog endnu at være død, han stiger atter op igjen, tilsidst gaaer der en Boble ud af hans Mund: naar dette er skeet, da synker han i Døden. Hiin Boble var det sidste Aandedrag, det sidste Forraad af Luft, der kunde gjøre ham lettere end Havet: saaledes med hiint Ord; i hiint Ord udaandede det sidste Haab om Frelse, i hiint Ord opgav han sig selv; fortiet, var der endnu i denne Tanke et Haab om Frelse, gjemt i Sjelen var der endnu i denne Tanke et muligt Forhold til Frelsen; naar Ordet er udtalt, i Fortrolighed til et andet Menneske, (o, rædsomme Misbrug af Fortroligheden, om dog den Fortvivlede kun misbruger den mod sig selv!) naar dette Ord er hørt, da synker han for stedse. Ak, det er forfærdeligt at see et Menneske storme frem mod sin egen Undergang, det er forfærdeligt at see ham dandse paa Afgrundens Rand uden at ane det; men denne Klarhed over sig selv og over sin egen Undergang er dog endnu forfærdeligere. Det er forfærdeligt at see et Menneske søge Trøst ved at styrte sig i Fortvivlelsens Hvirvel, men denne Besindighed er dog endnu forfærdeligere; at et Menneske ikke i Skrig om Hjælp i Dødsens Angest raaber: jeg gaaer under, frels mig, men at han roligt vil være Vidne til sin egen Fortabelse. O, yderste Forfængelighed, ikke at ville drage Menneskenes Øine paa sig ved sin Skjønhed, ved Riigdom, ved Evner, ved Magt, ved Æren, men at ville tiltrygle sig deres Opmærksomhed ved sin Fortabelse, ved at ville sige om sig selv, hvad i det Høieste Medlidenheden veemodig turde sige om en Saadan ved hans Grav: der gik dog noget Godt under i ham! O rædsomme Tvesindethed, i Fortabelsen at ville drage en Slags Fordeel af, at det Gode er til, det Eneste, man ikke har villet! Thi nu viste den sig jo, den anden Villie, om den end var saa svag, at den blev et Lefleri i Fortabelsen, et Forsøg paa at blive mærkværdig – ved Fortabelse.

At ville Eet kan da ikke betyde at ville, hvad der i sit Væsen ikke er Eet, men kun ved en rædsom Usandhed synes det, hvad der ved Løgnen er Eet, som Den, der kun vil dette ene, er en Løgner, som Den, der fremkogler dette Eet, er Løgnens Fader. Det Øde og Tomme er ikke Eet i Sandhed, men er i Sandhed Intet, og er Fortabelsen i Den, som kun vil dette Ene. Skal derimod et Menneske i Sandhed kun ville Eet, da maa dette Ene i sit inderste Væsens Sandhed være Eet; det maa ved en evig Adskillelse fjerne det Ueensartede fra sig, at det i Sandhed kan vedblive at være Eet og være det Samme, og derved danne Den til Lighed med sig, som kun vil dette Ene.

I Sandhed at ville Eet kan saaledes kun betyde at ville det Gode,fordi ethvert andet Eet ikke er et Eet, og den Villende, der kun vil det, derfor maa blive tvesindet; thi, som det Attraaede er, saaledes bliver den Attraaende. Eller var det muligt, at et Menneske ved at ville det Onde, kunde ville Eet, selv om det var muligt, at et Menneske kunde forhærde sig til kun at ville det Onde? Er ikke det Onde, ligesom de Onde, uenigt med sig selv, splidagtigt i sig selv? Tag et saadant Menneske, fjern ham fra det borgerlige Samfund, indespær ham i eensomt Aflukke: er han ikke der splidagtig med sig selv, som det slette Sammenhold mellem saadanne Ligesindede er en splidagtig Forening. Men selv om den Gode levede i en Afkrog af Verden og aldrig saae noget Menneske, han er dog enig med sig selv og enig med Alle, fordi han vil Eet, og fordi det Gode er Eet. Til at ville det Gode maa da Enhver den føres, der i Sandhed skal ville Eet, om det end stundom kan være saa, at et Menneske begynder med at ville Eet, der dog ikke i dybeste Forstand er det Gode, men vel noget Uskyldigt, og nu lidt efter lidt omdannes til ret i Sandhed at ville Eet ved at ville det Gode. Saaledes har vel stundom Elskov hjulpet et Menneske paa den rette Vei. Trofast vilde han kun Eet, sin Forelskelse; for den vilde han leve og døe, for den offre Alt, i den ene have sin Salighed. Dog en Forelskelse er endnu ikke i dybeste Forstand det Gode, men muligt blev den ham en dannende Opdrager, som tilsidst førte ham, ved Besiddelsen af den Elskede eller maaskee ved Tabet, til i Sandhed at ville Eet og ville det Gode. Saaledes opdrages et Menneske paa mange Maader; ogsaa oprigtig Elskov er en Opdragelse til det Gode. – Maaskee var der Den, hvem Begeistringen greb for en bestemt Stræben. Begeistret vilde han kun Eet, han vilde leve og døe for denne Stræben, han vilde offre Alt for den, i den ene vilde han have sin Salighed; thi Elskov og Begeistring tage ikke til Takke med et deelt Hjerte. Dog hans Stræben var maaskee endnu ikke i dybeste Forstand det Gode: saa blev Begeistringen ham Læremesteren, hvilken han vel voxede fra, men hvem han dog ogsaa skyldte meget. Thi, som sagt, alle Veie føre til det Gode, naar et Menneske i Sandhed kun vil Eet; og naar der dog er nogen Sandhed i, at han vil Eet, saa tjener dette ham ogsaa til Gode. Men Faren er der, at den Forelskede og den Begeistrede svinger urigtigt af, saa han svinger hen til det Store, istedenfor at føres til det Gode. Vel er nemlig det Gode i Sandhed det Store, men det Store er ikke altid det Gode; og man kan beile til en Qvindes Gunst ved at ville Noget, naar det blot er stort; det kan smigre Pigens Stolthed og hun kan lønne med sin Tilbedelse: men Gud i Himlene er ikke som en ung Piges Daarlighed, han lønner ikke det Store med Beundring, men den Godes Løn er at turde tilbede i Sandhed.

\workheading{Kjerlighedens Gjerninger}{1847}

\worknote{Antologiens ''Works of Love''. Hong-udvalget bringer forordet og bønnen samt uddrag af tre overveielser: ''Du skal elske'' (om at det at elske er en Pligt), ''»Du« skal elske Næsten'' (forkjerlighed kontra næstekjerlighed, Vandmærket) sammenstillet med slutningen af ''Vor Pligt at elske de Mennesker, vi see'', og ''Kjerlighed opbygger''. Spring markeret med […].}

\sectionheading{Forord; Bøn  ·  SV¹ IX, 7–12}

Disse christelige Overveielser, som ere Frugten af megen Overveielse, ville forstaaes langsomt men da ogsaa let, medens de vistnok ville blive meget vanskelige, om Nogen ved flygtig og nysgjerrig Læsning gjør sig dem meget vanskelige. »Hiin Enkelte«, som først overveier med sig selv, om han vil læse eller han ikke vil læse, hvis han saa vælger at læse, han overveie kjerligt, om dog ikke Vanskeligheden og Letheden, naar de bringes betænksomt sammen paa Vægtskaalen, forholde sig rigtigt til hinanden, saa det Christelige ikke udgives efter falsk Vægt, ved at gjøre Vanskeligheden eller ved at gjøre Letheden for stor.

Det er »christelige Overveielser« derfor ikke om – »Kjerlighed«, men om – »Kjerlighedens Gjerninger«.

Det er »Kjerlighedens Gjerninger«, ikke som var nu hermed alle dens Gjerninger optalte og beskrevne, o langtfra; ikke som var selv den enkelte beskrevne nu eengang for alle beskreven, Gud være lovet langtfra! Hvad der i sin hele Rigdom er væsentligen uudtømmeligt, det er ogsaa i sin den mindste Gjerning væsentligen ubeskriveligt, just fordi det væsentligen er heelt tilstede overalt, og væsentligen ikke til at beskrive.

S. K.

Bøn.

Hvor skulle der kunne tales retteligen om Kjerlighed, hvis Du var glemt, Du Kjerlighedens Gud, af hvem al Kjerlighed er i Himmelen og paa Jorden; Du, der Intet sparede, men gav Alt hen i Kjerlighed; Du, der er Kjerlighed, saa den Kjerlige kun er hvad han er ved at være i Dig! Hvor skulde der kunne tales retteligen om Kjerlighed, hvis Du var glemt, Du, der gjorde det aabenbar, hvad Kjerlighed er, Du, vor Frelser og Forsoner, som gav Dig selv hen for at frelse Alle! Hvor skulde der kunne tales retteligen om Kjerlighed, hvis Du var glemt, Du Kjerlighedens Aand, Du, som Intet tager af Dit Eget, men minder om hiint Kjerlighedens Offer, minder den Troende om at elske, som han er elsket, og sin Næste som sig selv! O, evige Kjerlighed, Du, som overalt er tilstede, og aldrig uden Vidnesbyrd, hvor Du paakaldes, Du lade Dig heller ei uden Vidnesbyrd i hvad her skal tales om Kjerlighed, eller om Kjerlighedens Gjerninger. Thi vel er der kun nogle Gjerninger, som det menneskelige Sprog særligt og smaaligt kalder Kjerlighedsgjerninger; men i Himlen er det jo saaledes, at ingen Gjerning kan tækkes der, uden den er en Kjerlighedens Gjerning: oprigtig i Selvfornegtelse, en Kjerlighedens Trang og just derfor uden Fortjenstlighedens Fordring!

\sectionheading{Du skal elske (uddrag)  ·  SV¹ IX, 21–46}

Du skal elske.

Mth. XXII, 39. Men det andet Bud er ligesom dette: Du skal elske Din Næste som Dig selv.

Enhver Tale, især da et Brudstykke af en Tale, forudsætter gjerne Noget, hvorfra der gaaes ud; Den, som vil tage Talen eller Udsagnet til Overveielse, gjør derfor vel i først at finde denne Forudsætning, for da at begynde med den. Saaledes er der nu ogsaa i vor forelæste Text indeholdt en Forudsætning, som rigtignok kommer sidst, men dog er Begyndelsen. Naar der nemlig siges »Du skal elske Din Næste som Dig selv«, da er deri indeholdt hvad der er forudsat, at ethvert Menneske elsker sig selv. Dette forudsætter Christendommen altsaa, der ingenlunde, som hine høitflyvende Tænkere, begynder uden Forudsætning, og heller ei med en smigrende Forudsætning. Og turde vi vel negte, at det er saa som Christendommen forudsætter? Men paa den anden Side, skulde Nogen kunne misforstaae Christendommen, som var det dens Mening at lære, hvad verdslig Kløgt, eenstemmigt – ak og dog just splidagtigt, lærer, »at Enhver er sig selv nærmest«; skulde Nogen kunne misforstaae dette, som var det Christendommens Mening at lyse Selvkjerlighed i Hævd? Det er jo tvertimod dens Mening at fravriste os Mennesker Selvkjerligheden. Denne ligger nemlig i at elske sig selv, men skal man elske Næsten »som sig selv«, saa vrister jo Budet, som med en Dirk, Selvkjerlighedens Lukke op og fravrister dermed Mennesket den. Hvis Budet om at elske Næsten var udtrykt paa en anden Maade end med dette lille Ord »som Dig selv«, der paa eengang er saa let at haandtere og dog har Evighedens Spændkraft, da vilde Budet ikke saaledes kunne magte Selvkjerligheden. Dette »som Dig selv« vakler ikke i Sigtet og trænger saa med Evighedens Urokkelighed dømmende ind til det inderste Skjulested, hvor et Menneske elsker sig selv; det lader ikke Selvkjerligheden den mindste Undskyldning tilovers, den mindste Udflugt aaben. Hvor underligt! Der kunde jo holdes lange og skarpsindige Taler om, hvorledes et Menneske skal elske sin Næste; og naar saa Talerne vare hørte, skulde Selvkjerligheden dog kunne hitte paa Undskyldning og finde Udflugt, fordi Sagen ikke ganske var udtømt, alle Tilfælde ikke opregnede, fordi der bestandigt var Noget glemt, eller Noget ikke nøiagtigt og bindende nok udtrykt og beskrevet. Men dette »som Dig selv« – ja, ingen Bryder kan saaledes omslutte Den, med hvem han skal brydes, som dette Bud omslutter Selvkjerligheden, der ikke kan røre sig af Stedet. Sandeligen, naar Selvkjerligheden har stridt med dette Ord, der dog er saa let at forstaae, at Ingen skal behøve at bryde sit Hoved dermed, da skal den formærke, at den har stridt med den Stærkere. Som Jacob haltede efter at have stridt med Gud, saa skal Selvkjerligheden være brudt, hvis den har stridt med dette Ord, der dog ikke vil lære et Menneske, at han ikke skal elske sig selv, men tvertimod netop lære ham den rette Selvkjerlighed. Hvor underligt! Hvilken Strid er langvarig, forfærdelig, indviklet som Selvkjerlighedens Kamp for at værge om sig, og dog afgjør Christendommen Alt med eet eneste Slag. Det Hele er hurtigt som en Haandevending, Alt er afgjort, som Opstandelsens evige Afgjørelse, »i en Hast, i et Øieblik« (1 Cor. 15, 52): Christendommen forudsætter, at Mennesket elsker sig selv, og føier derpaa blot Ordet til om Næsten »som Dig selv«. Og dog ligger der en Evighedens Forandring mellem det Første og det Sidste.

Men skulde dette ogsaa være det Høieste, skulde det ikke være muligt at elske et Menneske høiere end sig selv? Denne Digter-Begeistringens Tale høres jo i Verden; skulde det da maaskee være saa, at det var Christendommen, der ikke formaaede at svinge sig saa høit, saa den, formodentligen ogsaa i Betragtning af, at den henvender sig til eenfoldige dagligdags Mennesker, kummerligt blev staaende ved Fordringen om at elske Næsten »som sig selv«, ligesom den jo istedetfor hiin høitflyvende Kjærligheds besjungne Gjenstand »en Elsket, en Ven«, sætter det tilsyneladende meget lidt Digteriske »Næsten« – thi Kjerlighed til Næsten har vist ingen Digter besjunget, saa lidet som det at elske »som sig selv«: skulde dette maaskee være saaledes? Eller skulde vi, idet vi dog gjøre den besjungne Kjerlighed en Indrømmelse i Sammenligning med den befalede, fattigt prise Christendommens Besindighed og Forstand paa Livet, at den mere ædrueligt og mere fastende holder sig ved Jorden, maaskee i samme Mening som Ordsproget »elsk mig lidt og elsk mig længe«? Dette være langt fra os. Christendommen veed dog nok bedre Besked, om hvad Kjerlighed er og om det at elske, end nogen Digter; netop derfor veed den ogsaa, hvad der maaskee undgaaer Digterne, at den Kjerlighed, de besynge, skjult er Selvkjerlighed, og at just deraf dens berusede Udtryk lader sig forklare: at elske et andet Menneske høiere end sig selv. Elskov er endnu ikke det Evige, den er Uendelighedens skjønne Svimmel, dens høieste Udtryk er Gaadefuldhedens Dumdristighed; deraf kommer det, at den endog forsøger sig i et endnu mere svimlende Udtryk »at elske et Menneske høiere end Gud«. Og denne Dumdristighed behager Digteren over al Maade, den er ham en Liflighed i hans Øre, den begeistrer ham til Sange. Ak, Christendommen lærer, at dette er Gudsbespottelse. – Og om Venskab gjælder igjen hvad der gjælder om Elskov, forsaavidt ogsaa dette ligger i Forkjerligheden: at elske dette eneste Menneske fremfor alle Andre, at elske ham i Modsætning til alle Andre. Baade Elskovens og Venskabets Gjenstand har derfor ogsaa Forkjerlighedens Navn »den Elskede«, »Vennen«, hvilke elskes i Modsætning til hele Verden. Derimod er det den christelige Lære at elske Næsten, at elske hele Slægten, alle Mennesker, endog Fjenden, og ikke gjøre Undtagelse, hverken Forkjerlighedens eller Afskyens.

Der er kun Een, et Menneske med Evighedens Sandhed kan elske høiere end sig selv, det er Gud. Derfor hedder det heller ikke »Du skal elske Gud som Dig selv«, men det hedder »Du skal elske Herren din Gud i dit ganske Hjerte og i din ganske Sjel og i dit ganske Sind«. Gud skal et Menneske elske ubetinget lydende og elske ham tilbedende. Det er Ugudelighed, om noget Menneske vovede at elske sig selv saaledes, eller vovede at elske et andet Menneske saaledes, eller vovede at tillade et andet Menneske at elske ham saaledes. Dersom den Elskede eller Vennen bad Dig om Noget, som Du, netop fordi Du elskede redeligt, bekymret havde overveiet, at det var til hans Skade: da skal Du bære et Ansvar, hvis Du elsker ved at adlyde istedenfor at elske ved at negte Ønskets Opfyldelse. Men Gud skal Du elske i ubetinget Lydighed, selv om det han fordrer af Dig, maatte synes Dig Din egen Skade, ja skadeligt for hans Sag; thi Guds Viisdom er uden Sammenligningens Forhold til Din, og Guds Styrelse er uden Ansvarets Forpligtelse i Forhold til Din Klogskab; Du har blot elskende at lyde. Et Menneske derimod skal Du kun, – dog nei, det er jo det Høieste – altsaa et Menneske skal Du elske som Dig selv; kan Du bedre indsee hans Bedste end han, da skal Du ingen Undskyldning have i, at det Skadelige var hans eget Ønske, var hvad han selv bad om. Dersom dette ikke var Tilfældet, da kunde der ganske rigtigt være Tale om at elske et andet Menneske høiere end sig selv; thi det vilde være: tiltrods for sin Indsigt, at det var ham skadeligt, adlydende at gjøre det, fordi han forlangte det, eller tilbedende, fordi han ønskede det. Men dette har Du netop ikke Lov til; Du har Ansvaret, hvis Du gjør det, ligesom og den Anden har Ansvaret, hvis han paa saadan Maade vil misbruge sit Forhold til Dig.

Altsaa – »som Dig selv«. Dersom den listigste Bedrager, der nogensinde har levet (eller vi kunne jo digte ham endnu listigere end han nogensinde har levet), for om muligt at faae Loven til at gjøre mange Ord og til at blive vidtløftig, thi da har Bedrageren snart seiret, vilde vedblive Aar ud og Aar ind »fristende« at spørge den »kongelige Lov« »hvorledes skal jeg elske min Næste?« da skal det ordknappe Bud uforandret vedblive at gjentage det korte Ord »som Dig selv«. Og dersom nogen Bedrager – bedrog sig selv hele Livet igjennem ved allehaande Vidtløftigheder denne Sag betræffende, da skal Evigheden blot foreholde ham det korte Lovens Ord »som Dig selv«. Sandeligen Ingen skal kunne undflye Budet; kommer dets »som Dig selv« Selvkjerligheden saa nær paa Livet som muligt, saa er »Næsten« atter en Bestemmelse, der i Nærgaaenhed er Selvkjerligheden saa livsfarlig som muligt. At det er en Umulighed at snoe sig derfra, indseer Selvkjerligheden selv. Den eneste Udflugt er, hvad jo ogsaa Pharisæeren i sin Tid, for at gjøre sig retfærdig, forsøgte: at lade det være tvivlsomt, hvo Ens Næste er – for at faae ham fra Livet.

Hvo er da Ens Næste?Ordet er aabenbart dannet af Nærmeste, altsaa er Næsten Den, der er Dig nærmere end alle Andre, dog ikke i Forkjerlighedens Forstand, thi at elske Den, der i Forkjerlighedens Forstand er En nærmere end alle Andre, er Selvkjerlighed – »gjøre ikke ogsaa Hedningerne det Samme?« Næsten er da Dig nærmere end alle Andre. Men er han Dig ogsaa nærmere end Du er Dig selv? Nei, det er han ikke, men han er just, eller han skal just være Dig lige saa nær. Begrebet »Næste« er egentligen Fordoblelsen af Dit eget Selv; »Næsten« er hvad Tænkerne vilde kalde det Andet, Det, hvorpaa det Selviske i Selvkjerligheden skal prøves. Forsaavidt behøver, for Tankens Skyld, Næsten end ikke at være til. Om et Menneske levede paa en øde Øe, dersom han dannede sit Sind efter Budet, saa vilde han, ved at forsage Selvkjerligheden, kunne siges at elske Næsten. Vel er »Næsten« i sig en Mangfoldighed, thi »Næsten« betyder »alle Mennesker«, og dog er i en anden Forstand eet Menneske nok for at Du kan indøve Loven. Det er nemlig en Umulighed i selvisk Forstand, bevidst, at være To om at være Selv; det maa Selvkjerligheden være ene om. Der behøves heller ikke Tre, thi er der To, det vil sige er der eet andet Menneske, som Du i christelig Forstand elsker »som Dig selv«, eller i hvem Du elsker »Næsten«, saa elsker Du alle Mennesker. Men hvad det Selviske ubetinget ikke kan taale er: Fordoblelse, og Budets Ord »som Dig selv« er just Fordoblelsen. Den i Elskov Glødende kan ingenlunde paa Grund eller i Kraft af denne Gløden taale Fordoblelsen, der her vilde være: at opgive Elskoven, hvis den Elskede fordrede dette. Den Elskende elsker altsaa ikke den Elskede »som sig selv«, thi han er fordrende, men dette »som sig selv« indeholder just Fordringen til ham – ak, og dog mener den Elskende endog at elske det andet Menneske høiere end sig selv.

»Næsten« er saaledes Selvkjerligheden saa nær paa Livet som muligt; er der kun to Mennesker, er det andet Menneske Næsten; er der Millioner, er Enhver af disse Næsten, der igjen er En nærmere end »Vennen« og »den Elskede«, forsaavidt disse, som Forkjerlighedens Gjenstand, saa smaat holde sammen med Selvkjerligheden i En. At Næsten saaledes er til og er En saa nær, er man i Almindelighed ogsaa vidende om, naar man mener at have Rettigheder i Forhold til ham, at kunne fordre Noget af ham. Dersom Nogen i denne Forstand spørger, hvo er min Næste, da vil hiint Christi Svar til Pharisæeren kun paa en egen Maade indeholde Besvarelsen, thi i Svaret gjøres egentligt først Spørgsmaalet til det modsatte, hvorved betydes, hvorledes et Menneske har at spørge. Christus siger nemlig, efterat have fortalt Parabelen om den barmhjertige Samaritan, til Pharisæeren (Luc. 10, 36.) »hvilken af disse Tre tykkes Dig nu at have været hans Næste, der var falden blandt Røvere«, og Pharisæeren svarer »rigtigt« »den, som gjorde Barmhjertighed mod ham«; det er, ved at anerkjende Din Pligt opdager Du let, hvo Din Næste er. Pharisæerens Svar er indeholdt i Christi Spørgsmaal, der ved sin Form nødsagede Pharisæeren til at svare saaledes. Den jeg har Pligten mod er min Næste, og naar jeg fuldkommer min Pligt viser jeg, at jeg er Næsten. Christus taler nemlig ikke om det at kjende Næsten, men om det selv at blive Næsten, at bevise sig at være Næsten, som Samaritanen beviste det ved sin Barmhjertighed; thi ved den beviste han jo ikke, at den Overfaldne var hans Næste; men at han var den Overfaldnes Næste. Leviten og Præsten vare i nærmere Forstand den Overfaldnes Næste, men disse vilde Intet vide derom; Samaritanen derimod, der ved Fordomme var foranlediget til Misforstaaelse, forstod dog rigtigt, at han var den Overfaldnes Næste. At vælge en Elsket, at finde en Ven, ja det er et vidtløftigt Arbeide, men Næsten er let at kjende, let at finde, hvis man blot selv vil – anerkjende sin Pligt.

Budet lød saaledes »Du skal elske Din Næste som Dig selv«, men naar Budet forstaaes retteligen, siger det tillige det Omvendte: Du skal elske Dig selv paa den rette Maade. Dersom Nogen derfor ikke af Christendommen vil lære at elske sig selv paa den rette Maade, kan han heller ikke elske Næsten; han kan saa maaskee som det hedder »i Liv og Død« – holde sammen med et andet eller flere andre Mennesker, men dette er ingenlunde at elske Næsten. At elske sig selv paa den rette Maade og at elske Næsten svare ganske til hinanden, er i Grunden Eet og det Samme. Naar Lovens »som Dig selv« har fravristet Dig den Selvkjerlighed, hvilken Christendommen, sørgeligt nok, maa forudsætte at være i ethvert Menneske, saa har Du netop lært at elske Dig selv. Loven er derfor: Du skal elske Dig selv saaledes som Du elsker Næsten, naar Du elsker ham som Dig selv. Hvo der har noget Kjendskab til Menneskene vil vist tilstaae, at som han ofte har ønsket at kunne bevæge dem til at opgive Selvkjerligheden, at han saaledes ogsaa ofte har maattet ønske, om det var muligt at lære dem at elske dem selv. Naar den Travle ødsler sin Tid og sin Kraft i forfængelige, ubetydelige Idrætters Tjeneste, er dette da ikke, fordi han ikke har lært ret at elske sig selv? Naar den Letsindige bortgiver sig selv næsten som et Intet med i Øieblikkets Narrespil, er dette da ikke, fordi han ikke har Forstand paa ret at elske sig selv? Naar den Tungsindige ønsker at blive af med Livet, ja med sig selv, er dette da ikke, fordi han ikke vil lære strengt og alvorligt at elske sig selv? Naar et Menneske, fordi Verden eller et andet Menneske troløs lod ham forraadt, overgiver sig til Fortvivlelse, hvilken er da hans Skyld (thi hans uskyldige Lidelse tale vi jo her ikke om) uden den ikke at elske sig selv paa den rette Maade? Naar et Menneske selvplagersk mener at gjøre Gud en Tjeneste ved at martre sig selv, hvilken er da hans Synd uden denne, ikke ret at ville elske sig selv? Ak, og naar et Menneske formasteligt lægger Haand paa sig selv, er dette da ikke just hans Synd, at han ikke retteligen elsker sig selv i den Forstand, i hvilken et Menneske skal elske sig selv? O, der tales i Verden saa meget om Forræderie og Troløshed, og Gud bedre det, at dette destoværre kun er altfor sandt, men lad os dog derover aldrig glemme, at den farligste Forræder blandt alle er den, ethvert Menneske har i sig selv. Dette Forræderie, hvad enten det bestaaer i, at han selvisk elsker sig, eller det bestaaer i, at han selvisk ikke vil elske sig selv paa den rette Maade, dette Forræderie er jo rigtignok en Hemmelighed, der gjøres intet Anskrig i den Anledning som ellers i Anledning af Forræderie og Troløshed; men er det ikke just derfor desto vigtigere, at der atter og atter mindes om Christendommens Lære: at et Menneske skal elske sin Næste som sig selv, det er, som han skal elske sig selv.

Budet om Kjerlighed til Næsten taler da i eet og samme Ord, »som Dig selv«, om denne Kjerlighed og om Kjerligheden til sig selv – og nu er Talens Indledning standset ved det, som den ønsker at gjøre til Gjenstand for Betragtningen. Det nemlig, hvorved Budet om Kjerlighed til Næsten og om Kjerlighed til sig selv bliver eenstydigt, er ikke blot dette »som Dig selv«, men endnu mere det Ord »Du skal«. Herom er det vi ville tale:

Du skal elske,

thi dette er just Kjendet paa den christelige Kjerlighed og er dens Eiendommelighed, at den indeholder denne tilsyneladende Modsigelse: at det at elske er Pligt.

[…]

»Du skal elske«. Kun naar det er Pligt at elske, kun da er Kjerligheden for evig betrygget mod enhver Forandring; evig frigjort i salig Uafhængighed; evig lykkelig sikkret mod Fortvivlelse.

Hvor glad, hvor lykkelig, hvor ubeskrivelig tillidsfuld end Driftens og Tilbøielighedens Kjerlighed, den umiddelbare Kjerlighed som saadan kan være, den føler dog netop i sit skjønneste Øieblik en Trang til, om muligt at binde sig end fastere. Derfor sværge de tvende, de tilsværge hinanden Tro eller Venskab; og naar vi tale festligst, sige vi ikke om Tvende »de elske hinanden«, vi sige »de tilsvore hinanden Tro« eller »de tilsvore hinanden Venskab«. Men hvorved sværger da denne Kjerlighed? Vi ville nu ikke forstyrre Opmærksomheden og adsprede ved at minde om det meget Forskjellige, som denne Kjerligheds Talsmænd »Digterne« ved Indvielsen vide Besked om – thi i Forhold til denne Kjerlighed er det Digteren, som tager Løftet af de Tvende, Digteren, som forener de Tvende, Digteren, som foresiger de Tvende Eden og lader dem sværge, kort det er Digteren, der er Præsten. Sværger nu denne Kjerlighed ved Noget, som er høiere end den selv? Nei, det gjør den ikke. Just dette er den skjønne, den rørende, den gaadefulde, den digteriske Misforstaaelse, at de Tvende ikke selv opdage det; og netop derfor er Digteren deres eneste, deres elskede Fortrolige, fordi han heller ikke opdager det. Naar denne Kjerlighed sværger, da giver den egentligen selv Det Betydning, hvorved den sværger; det er Kjerligheden selv, der kaster Glandsen over Det, hvorved den sværger, saa den altsaa ikke blot ikke sværger ved noget Høiere, men egentligen sværger ved Noget, der er ringere end den selv. Saa ubeskrivelig riig er denne Kjerlighed i sin elskelige Misforstaaelse; thi netop fordi den er sig selv en uendelig Rigdom, en ubegrændset Tilforladelighed, kommer den, idet den vil sværge, til at sværge ved noget Ringere, men opdager det ikke selv. Deraf kommer det igjen, at denne Sværgen, der jo skulde være og ogsaa oprigtigt mener sig selv at være den høieste Alvor, dog er den yndigste Spøg. Og den gaadefulde Ven, Digteren, hvis fuldkomne Fortrolighed er denne Kjerligheds høieste Forstaaelse, han forstaaer det heller ikke. Dog er det jo let at forstaae, at skal man i Sandhed sværge, da maa man sværge ved noget Høiere, saa kun Gud i Himlene er den Eneste, der i Sandhed er i det Tilfælde, at han ene kan sværge ved sig selv. Dette kan imidlertid Digteren ikke forstaae, det er, den Enkelte, som er Digter, kan vel forstaae det, men han kan ikke forstaae det forsaavidt han er Digter, og »Digteren« kan ikke forstaae det; thi Digteren kan forstaae Alt – i Gaader, og vidunderligt forklare Alt – i Gaader, men han kan ikke forstaae sig selv, eller forstaae, at han selv er en Gaade. Skulde han tvinges til at forstaae det, da vilde han, hvis han ikke blev opbragt og forbittret, veemodigt sige: gid man ikke havde paanødt mig denne Forstaaelse, der forstyrrer mig mit Skjønneste, forstyrrer mig mit Liv, medens jeg ingen Brug kan gjøre af den. Og deri har Digteren forsaavidt Ret, thi den sande Forstaaelse er Livsspørgsmaalets Afgjørelse for hans Tilvær. Der er paa den Maade to Gaader, den første er de Tvendes Kjerlighed, den anden er Digterens Forklaring af den, eller at Digterens Forklaring ogsaa er en Gaade.

Saaledes sværger denne Kjerlighed, og da tilføie de Tvende til Eden, at de ville elske hinanden »for evig«. Naar dette ikke føies til, saa forener Digteren ikke de Tvende, han vender sig ligegyldig bort fra en saadan timelig Kjerlighed, eller han vender sig spottende imod den, medens han for evig tilhører hiin evig Kjerlighed. Der er da egentligen to Foreninger, først de Tvende, som for evig ville elske hinanden, og saa Digteren, som for evig vil tilhøre de Tvende. Og deri har jo Digteren Ret, at naar tvende Mennesker ikke ville elske hinanden for evig, saa er deres Kjerlighed ikke værd at tale om, end mindre værd at besynge. Derimod mærker Digteren ikke Misforstaaelsen, at de Tvende ved deres Kjerlighed sværge at ville elske hinanden for evig, istedenfor ved Evigheden at tilsværge hinanden Kjerlighed. Evigheden er det Høiere; skal man sværge, da maa man sværge ved det Høiere, men skal man sværge ved Evigheden, da sværger man ved Pligten »at skulle elske«. Ak, men hiin de Elskendes Yndling, Digteren, han, der selv er det end Sjeldnere end de tvende ret Elskende, hvilke hans Længsel søger, han, der selv er Elskelighedens Vidunder, han er ogsaa som det kjælne Barn, han kan ikke taale dette »skal«, saasnart det siges, bliver han enten utaalmodig eller han giver sig til at græde.

Altsaa denne umiddelbare Kjerlighed, den har vel i den skjønne Indbildnings Forstand det Evige i sig, men den er ikke bevidst begrundet paa det Evige, og derfor kan den forandres. Selv om den ikke forandredes, den kan dog forandres, thi den er jo det Lykkelige; men om det Lykkelige gjælder hvad der gjælder om Lykken, hvad der, naar man tænker det Evige, ikke kan tænkes uden Veemod, ligesom det er sagt med Gysen: »Lykken er, naar den har været«. Dette vil sige, saa længe den bestod eller var til, var en Forandring mulig; først naar den er bleven forbigangen, kan man sige, den bestod. »Man prise intet Menneske lykkelig, saa længe han lever«; saa længe han lever, kan nemlig Lykken forandre sig, først naar han er død, og Lykken ikke har forladt ham, medens han levede, først da viser det sig, at han – har været lykkelig. Hvad der blot er til, hvad der ingen Forandring har undergaaet, det har bestandigt Forandringen uden for sig; den kan bestandigt indtræde, endnu i det sidste Øieblik kan den komme, og først naar Livet er tilendebragt, kan man sige: Forandringen kom ikke – eller maaskee kom den. Hvad der ingen Forandring har undergaaet har vel Bestaaen, men har ikke Bestandighed; forsaavidt det har Bestaaen, er det til, men forsaavidt det ikke i Forandringen har vundet Bestandighed, kan det ikke blive samtidigt med sig selv, og er da enten lykkeligt uvidende om dette Misforhold eller stemt i Veemod. Thi det Evige er det Eneste, som kan være og blive og forblive samtidigt med enhver Tid; Timeligheden derimod adskiller i sig selv, og det Nærværende kan ikke blive samtidigt med det Tilkommende, eller det Tilkommende med det Forbigangne, eller det Forbigangne med det Nærværende. Om hvad der derfor ved at undergaae Forandringen vandt Bestandighed, om det kan man ikke blot, naar det har bestaaet, sige »det bestod«, men man kan sige »det har bestaaet, medens det bestod.« Just dette er Betryggelsen, og et ganske andet Forhold end Lykkens. Naar Kjerligheden har undergaaet Evighedens Forandring ved at være bleven Pligt, da har den vundet Bestandighed, og det følger da af sig selv, at den bestaaer. Det følger nemlig ikke af sig selv, at hvad der bestaaer i dette Øieblik ogsaa bestaaer i næste Øieblik, men det følger af sig selv, at det Bestandige bestaaer. Vi tale jo om, at Noget bestaaer sin Prøve, og vi prise det, naar det har bestaaet sin Prøve; men denne Tale er dog om det Ufuldkomne, thi det Bestandiges Bestandighed skal og kan ikke vise sig i at bestaae en Prøve – det er jo det Bestandige, og kun det Forgængelige kan give sig – Skin af Bestandighed ved at bestaae i en Prøve. Om Prøve-Sølv vilde derfor ingen falde paa at sige, at det skal bestaae sin Prøve i Aarene, thi det er jo Prøve-Sølv. Saaledes ogsaa med Kjerlighed. Den Kjerlighed, der blot har Bestaaen, hvor lykkelig, hvor livsalig, hvor tillidsfuld, hvor digterisk den end er, den maa dog bestaae sin Prøve i Aarene; men den Kjerlighed, der undergik Evighedens Forandring ved at blive Pligt, den vandt Bestandighed, den er Prøve-Sølv. Er den derfor maaskee mindre anvendelig, mindre brugbar i Livet? Er da Prøve-Sølv det mindre Brugbare? Dog vel ikke; men Sproget, uvilkaarligt, og Tanken, bevidst, ærer paa en egen Maade Prøve-Sølvet, thi om dette siges blot »man bruger det«, der tales slet ikke om Prøve, man fornærmer det ikke ved at ville prøve det, det veed man jo forud, at Prøve-Sølv bestaaer. Naar man derfor bruger en mindre tilforladelig Sammensætning, saa nødsages man til at være mere nærgaaende og til at tale mindre eenfoldigt, man nødsages til næsten tvetydigt at sige det Dobbelte »man bruger det, og medens man bruger det, prøver man det tillige«, thi det er dog bestandigt muligt, at det kunde forandre sig.

Altsaa, kun naar det er Pligt at elske, kun da er Kjerligheden evig betrygget. Denne Evighedens Betryggethed jager al Angest ud og gjør Kjerlighed fuldkommen, fuldkommen betrygget. Thi i hiin Kjerlighed, som kun har Bestaaen, hvor tillidsfuld den end er, er der dog endnu en Angest, en Angest for Forandringens Mulighed. Den selv forstaaer det ikke, lige saa lidet som Digteren, at det er Angest, thi Angesten er skjult, og kun den brændende Higen Yttringen, hvorpaa dog just kjendes, at Angesten er skjult i Grunden. Hvoraf kommer det vel ellers, at den umiddelbare Kjerlighed er saa tilbøielig til, ja saa forelsket i at gjøre Prøve paa Kjerligheden? Dette er jo just, fordi Kjerligheden ikke, ved at blive Pligt, i dybeste Forstand har undergaaet »Prøven«. Derfra denne, hvad Digteren vilde kalde, søde Uro, der dumdristigere og dumdristigere vil gjøre Prøve. Den Elskende vil prøve den Elskede, Vennen vil prøve Vennen; Prøvelsen har vistnok sin Grund i Kjerligheden, men denne heftigt flammende Lyst til at prøve, dette Ønskets Higen efter at maatte blive sat paa Prøve forklarer dog, at Kjerligheden sig selv ubevidst er usikker. Atter her er en gaadefuld Misforstaaelse i den umiddelbare Kjerlighed og i Digterens Forklaringer. De Elskende og Digteren mene, at denne Lyst til at ville prøve Kjerligheden netop er et Udtryk for, hvor sikker den er. Men er dette vel saa? Det er ganske rigtigt, at det Ligegyldige ønsker man ikke at prøve; men deraf følger jo ikke, at det at ville prøve det Elskede er et Udtryk for Sikkerheden. De Tvende elske hinanden, de elske hinanden for evig, de ere saa visse, at de – gjøre Prøve derpaa. Er denne Vished den høieste? Er Forholdet ikke her, netop ligesom naar Kjerligheden sværger og dog sværger ved hvad der er lavere end Kjerligheden? Saaledes er jo her de Elskendes høieste Udtryk for deres Kjerligheds Bestandighed et Udtryk for, at den kun har Bestaaen, thi hvad der kun har Bestaaen, det prøver man, det sætter man paa Prøve. Men naar det er Pligt at elske, da behøves der ingen Prøve og fornærmelig Dumdristighed med at ville prøve, da er Kjerligheden høiere end enhver Prøve, den har allerede mere end bestaaet i Prøven, i samme Forstand, som Troen »mere end seirer«. Det at prøve forholder sig altid til Mulighed, det er dog altid muligt, at det, der prøves, ikke bestod i Prøven. Dersom derfor En vilde prøve om han har Troen, eller prøve paa at faae Troen, da vil dette egentligen sige, han vil forhindre sig selv i at faae Troen, han vil bringe sig selv i en Higens Uro, hvor Troen aldrig vindes, thi »Du skal troe«. Dersom en Troende vilde bede Gud at sætte hans Tro paa en Prøve, da er dette ikke et Udtryk for, at denne Troende har Troen i overordentlig høi Grad (at mene det er en Digter-Misforstaaelse, som det ogsaa er en Misforstaaelse at have Troen i en »overordentlig« Grad, da den ordentlige Grad netop er det Høieste), men det er et Udtryk for, at han ikke ganske har Troen, thi »Du skal troe«. Aldrig fandtes der nogen høiere Betryggelse, og aldrig fandtes der Evighedens Hvile i noget Andet end i dette »skal«. Det er dog, hvor livsalig den end er, en urolig Tanke »at prøve«, og det er Uroen, der vil indbilde Dig, at dette er en høiere Forvisning; thi det at prøve er i sig selv opfindsomt og ikke til at udtømme, saa lidet som Klogskaben nogensinde har kunnet beregne alle Tilfælde, medens derimod, som den Alvorlige saa fortræffeligt siger, »Troen har beregnet alle Tilfælde«. Og naar man skal, saa er det for evig afgjort; og naar Du vil forstaae, at Du skal elske, saa er Din Kjerlighed for evig betrygget.

Og Kjerligheden er tillige ved dette »skal« evig betrygget mod enhver Forandring. Thi den Kjerlighed, der kun har Bestaaen, den kan forandres, den kan forandres i sig selv, og den kan forandres fra sig selv.

Den umiddelbare Kjerlighed kan forandres i sig selv, den kan forandres til sit Modsatte, til Had. Had er en Kjerlighed, som er bleven til sit Modsatte, en Kjerlighed, som er gaaet til Grunde. Nede i Grunden brænder bestandigt Kjerligheden, men Flammen er Hadets; først, naar Kjerligheden er udbrændt, først da er ogsaa Hadets Flamme slukket. Som der er sagt om Tungen, at »det er den samme Tunge, med hvilken vi velsigne og forbande«, saaledes maa man ogsaa sige, at det er den samme Kjerlighed, der elsker og hader; men netop fordi det er den samme Kjerlighed, netop derfor er den ikke i Evighedens Forstand den sande, der bliver uforandret den samme, medens hiin umiddelbare, naar den er forandret, dog i Grunden er den samme. Den sande Kjerlighed, der undergik Evighedens Forandring ved at blive Pligt, forandres aldrig, den er eenfoldig, den elsker – og hader aldrig, hader aldrig – den Elskede. Det kunde synes som var hiin umiddelbare Kjerlighed den stærkere, fordi den kan gjøre det Dobbelte, fordi den baade kan elske og hade; det kunde synes som havde den en ganske anden Magt over sin Gjenstand, naar den siger »vil Du ikke elske mig, saa vil jeg hade Dig«: dog er dette kun et Sandsebedrag. Thi er vel Forandrethed en stærkere Magt end Uforanderlighed; og hvo er den Stærkere, den som siger »vil Du ikke elske mig saa vil jeg hade Dig«, eller den, som siger »vil Du hade mig, saa vil jeg dog vedblive at elske Dig?« Ja vist er det forfærdende og forfærdeligt, at Kjerlighed forandres til Had, men hvem er det egentligen forfærdeligt for, mon ikke for den Paagjældende selv, hvem det hændte, at hans Kjerlighed forandrede sig til Had.

Den umiddelbare Kjerlighed kan forandres i sig selv, den kan ved Selvantændelse blive til Iversyge, blive fra den største Lykke til den største Qval. Saa farefuld er den umiddelbare Kjerligheds Hede, hvor stor dens Lyst end er, saa farefuld, at denne Hede let kan blive en Sygdom. Det Umiddelbare er ligesom det Gjærende, der jo just derfor kaldes saaledes, fordi det endnu ingen Forandring har undergaaet, og derfor heller ikke fra sig har udsondret den Gift, som dog just er det Glødende i det Gjærende. Antænder Kjerligheden sig selv ved denne Gift, istedenfor at udsondre den, da fremkommer Iversygen, ak, Ordet siger det jo selv, den er en Iver for at blive syg, en Sygdom af Iver. Den Iversyge hader ikke Kjerlighedens Gjenstand, langtfra, men han martrer sig selv ved den Gjenkjerlighedens Ild, der luttrende skulde rense hans Kjerlighed. Den Iversyge opfanger, næsten tryglende, hver en Straale af Kjerligheden i den Elskede, men alle disse Straaler samler han gjennem Iversygens Brændglas paa sin Kjerlighed, og han forbrænder langsomt. Den Kjerlighed derimod, der undergik Evighedens Forandring ved at blive Pligt, den kjender ikke Iversyge; den elsker ikke blot, som den bliver elsket, men den elsker. Iversygen elsker som den bliver elsket; ængstelig piint ved Forestillingen, om den bliver elsket, er den lige saa iversyg paa den egne Kjerlighed, om den dog ikke er uforholdsmæssig i Forhold til den Andens Ligegyldighed, som den er iversyg paa Yttringen af den Andens Kjerlighed; ængstelig piint ved Selvbeskæftigelsen tør den hverken ganske troe den Elskede, ei heller ganske give sig hen, for ikke at give for meget, og brænder sig derfor bestandigt, som man brænder sig paa Det, der ikke er brændende – uden for den ængstelige Berøring. Sammenligningen er Selvantændelsen. Det kunde synes, som var der en ganske anden Ild i den umiddelbare Kjerlighed, siden den kan blive til Iversyge; ak, men denne Ild er jo just det Forfærdelige. Det kunde synes, som holdt Iversygen sin Gjenstand ganske anderledes fast, naar den med hundrede Øine vogter paa den, medens den eenfoldige Kjerlighed ligesom kun har eet Øie for sin Kjerlighed. Men mon dog Adsplittelse er stærkere end Eenhed, mon et sønderrevet Hjerte er stærkere end et fuldt og udeelt, mon en ideligt ængstelig Griben holder sin Gjenstand fastere end Eenfoldighedens enige Kræfter! Og hvorledes er nu hiin eenfoldige Kjerlighed sikkret mod Iversygen? Mon ikke derved, at den ikke elsker sammenlignelsesviis? Den begynder ikke med umiddelbart at elske fortrinsviis, den elsker; derfor kan den heller aldrig komme til sygeligt at elske sammenlignelsesviis, den elsker.

Den umiddelbare Kjerlighed kan forandres fra sig selv, den kan forandres i Aarene, hvad jo ofte nok sees. Saa taber Kjerligheden sin Fyrighed, sin Glæde, sin Lyst, sin Oprindelighed, sit friske Liv; som den Flod, der sprang ud af Klippen, naar den længere henne afmattes i det stillestaaende Vands Dvaskhed, saaledes afmattes Kjerligheden i en Vanes Lunkenhed og Ligegyldighed. Ak, af alle Fjender er maaskee Vanen den lumskeste, og fremfor Alt er den lumsk nok til aldrig at lade sig blive seet, thi Den, der saae Vanen, han er frelst fra Vanen; Vanen er ikke som andre Fjender, hvilke man seer og mod hvilke man stridende værger sig, Striden er egentligen den med sig selv om at faae den at see. Der er et for sin Lumskhed bekjendt Rovdyr, som listende overfalder de Sovende; medens det da suger Blodet ud af den Sovende, vifter det Kølighed over ham og gjør ham Søvnen endnu lifligere. Saaledes er Vanen – eller den er endnu værre; thi hiint Dyr søger sit Bytte blandt de Sovende, men det har intet Middel til at dysse de Vaagne i Søvn. Det har derimod Vanen; den lister sig søvndyssende paa et Menneske, og naar dette er skeet, da blodsuger den en Sovende, medens den vifter Kølighed over ham og gjør ham Søvnen endnu lifligere. – Saaledes kan den umiddelbare Kjerlighed forandres fra sig selv, og blive ukjendelig – thi Had og Iversyge kjendes dog paa Kjerligheden. Saa mærker stundom Mennesket selv, som naar en Drøm farer forbi og er glemt, at Vanen har forandret ham; han vil da gjøre det godt igjen, men han veed ikke, hvor han skal gaae hen og kjøbe ny Olie til at opflamme Kjerligheden. Saa bliver han mismodig, fortræden, kjed af sig selv, kjed af sin Kjerlighed, kjed af, at den er saa ussel som den er, kjed af, at han ikke kan faae det forandret, ak, thi Evighedens Forandring havde han ikke itide agtet paa, og nu har han endog tabt Evne til at udholde Helbredelsen. O, man seer stundom med Sorg et Menneske, der engang levede i Velmagt, nu forarmet, og dog, hvor langt mere sørgeligt end denne Forandring er det, naar man seer en Kjerlighed forandret til dette næsten Modbydelige. – Dersom derimod Kjerligheden undergik Evighedens Forandring ved at blive Pligt, da kjender den ikke Vane, da kan Vane aldrig faae Magt over den. Som der siges om det evige Liv, at der er ikke Suk og ikke Taarer, saa kunde man tilføie: der er og ikke Vane; sandeligen dermed sige vi ikke noget mindre herligt. Vil Du frelse Din Sjel eller Din Kjerlighed fra Vanens Lumskhed – ja Menneskene troe, at der er mange Midler til at holde sig vaagen og sikkret, men sandeligen der er kun eet: Evighedens »skal«. Lad hundrede Kanoners Torden tre Gange dagligt minde Dig om at modstaae Vanens Magt; hold som hiin mægtige Keiser i Østen en Slave, der minder Dig dagligt, hold hundrede; hav en Ven, der minder Dig, hver Gang han seer Dig, hav en Hustru, der minder Dig tidligt og sildigt i Kjerlighed: men tag Dig iagt, at ikke ogsaa dette bliver en Vane! Thi Du kan blive vant til at høre hundrede Kanoners Torden, saa Du kan sidde tilbords og høre den mindste Ubetydelighed langt tydeligere end de hundrede Kanoners Torden, som Du blev – vant til at høre. Og Du kan blive vant til, at hundrede Slaver minde Dig hver Dag, saa Du hører det ikke mere, fordi Du ved Vanen har faaet det Øre, hvormed man hører og dog ikke hører. Nei, kun Evighedens »Du skal« – og det hørende Øre, som vil høre dette »skal«, kan frelse Dig fra Vanen. Vane er den sørgeligste Forandring, og paa den anden Side enhver Forandring kan man vænne sig til; kun det Evige, og altsaa hvad der undergik Evighedens Forandring ved at blive Pligt, er det Uforanderlige, men det Uforanderlige kan netop ikke blive Vane. Hvor fast end en Vane sætter sig, den bliver aldrig det Uforanderlige, selv om Mennesket bliver uforbederligt; thi Vanen er bestandigt Det, som skulde forandres, det Uforanderlige derimod Det, som hverken kan eller skal forandres. Men det Evige bliver aldrig gammelt og aldrig Vane.

Kun naar det er Pligt at elske, kun da er Kjerligheden evig frigjort i salig Uafhængighed. Men er da hiin umiddelbare Kjerlighed ikke fri, har den Elskende ikke netop sin Frihed i Kjerligheden? Og paa den anden Side, skulde det være Talens Mening at anprise Selvkjerlighedens trøstesløse Uafhængighed, der blev uafhængig, fordi den ikke havde Mod til at binde sig, altsaa fordi den blev afhængig af sin Feighed; den trøstesløse Uafhængighed, der svæver, fordi den intet Tilhold fandt, og er som »Den, der vanker hid og did, en bevæbnet Røver, der tager ind hvor det aftnes;« den trøsteløse Uafhængighed, der uafhængigt ikke bærer Lænker – idetmindste ikke synligt! O, langtfra, vi have jo tvertimod i den foregaaende Tale mindet om, at Udtrykket for den høieste Rigdom er at have en Trang; og saaledes er ogsaa Frihedens sande Udtryk, at den er en Trang i den Frie. Den, i hvem Kjerligheden er en Trang, han føler sig vistnok fri i sin Kjerlighed, og just Den, der føler sig ganske afhængig, saa han vilde tabe Alt ved at tabe den Elskede, just han er uafhængig. Dog paa een Betingelse, at han ikke forvexler Kjerligheden med Besiddelsen af den Elskede. Dersom En vilde sige »enten elske eller døe«, og derved betegne, at et Liv uden at elske ikke var værd at leve, da vilde vi ganske give ham Ret. Men hvis han ved det Første forstod at besidde den Elskede, og altsaa meente enten besidde den Elskede eller døe, enten vinde denne Ven eller døe, saa maatte vi sige, at en saadan Kjerlighed i usand Forstand er afhængig. Saasnart Kjerligheden, i sit Forhold til sin Gjenstand, i Forholdet ikke forholder sig lige saa meget til sig selv, medens den dog er ganske afhængig, saa er den i usand Forstand afhængig, saa har den Loven for sin Tilværelse uden for sig selv og er altsaa i forkrænkelig, i jordisk, i timelig Forstand afhængig. Men den Kjerlighed, der undergik Evighedens Forandring ved at blive Pligt, og elsker, fordi den skal elske, den er uafhængig, den har Loven for sin Tilværelse i selve Kjerlighedens Forhold til det Evige. Denne Kjerlighed kan aldrig blive i usand Forstand afhængig, thi det Eneste, den er afhængig af, er Pligten, og Pligten er det eneste Frigjørende. Den umiddelbare Kjerlighed gjør et Menneske fri og i næste Øieblik afhængig. Det er ligesom med et Menneskes Tilblivelse; ved at blive til, ved at blive et »Selv«, bliver han fri, men i næste Øieblik er han afhængig af dette Selv. Pligten derimod gjør et Menneske afhængig og i samme Øieblik evig uafhængig. »Ikkun Loven kan give Friheden.« Ak, man mener saa ofte, at Friheden er til, og at det er Loven, som binder Friheden. Dog er det just omvendt; uden Loven er Friheden slet ikke til, og det er Loven, som giver Friheden. Man mener ogsaa, at det er Loven, der gjør Forskjel, fordi der, naar Loven ikke er, slet ingen Forskjel er. Dog er det omvendt; naar det er Loven, som gjør Forskjel, saa er det just Loven, som gjør Alle lige for Loven.

Saaledes frigjør dette »skal« Kjerligheden i salig Uafhængighed; en saadan Kjerlighed staaer og falder ikke med sin Gjenstands Tilfældighed, den staaer og falder med Evighedens Lov – men saa falder den jo aldrig; en saadan Kjerlighed afhænger ikke af Dette eller Hiint, den afhænger kun af – det eneste Frigjørende, altsaa er den evigt uafhængig. Med denne Uafhængighed kan ingen sammenlignes. Stundom priser Verden den stolte Uafhængighed, som formener ingen Trang at føle til at blive elsket, om den dog tillige mener »at behøve andre Mennesker – ikke for at elskes af dem, men for at elske dem, for dog at have Nogen at elske.« O, hvor usand er ikke denne Uafhængighed! Den føler ingen Trang til at blive elsket, og dog behøver den Nogen at elske; den trænger altsaa til et andet Menneske – for at kunne tilfredsstille sin stolte Selvfølelse. Er dette ikke, som naar Forfængeligheden mener at kunne undvære Verden og dog trænger til Verden, det er, til at Verden bliver vidende om, at Forfængeligheden ikke trænger til Verden. Men den Kjerlighed, der undergik Evighedens Forandring ved at blive Pligt, føler vistnok en Trang til at blive elsket, og denne Trang med dette »skal« er derfor en evigt samstemmende Overeensstemmelse; men den kan undvære, hvis saa skal være, medens den dog vedbliver at elske: er dette ikke Uafhængighed? Denne Uafhængighed er kun afhængig af Kjerligheden selv ved Evighedens »skal«, den er ikke afhængig af noget Andet, og derfor heller ikke afhængig af Kjerlighedens Gjenstand, saa snart denne viser sig at være et Andet. Dog betyder dette ikke, at den uafhængige Kjerlighed saa ophørte, forvandlede sig til stolt Selvtilfredshed; dette er Afhængighed. Nei, Kjerligheden bliver, det er Uafhængighed. Uforandrethed er den sande Uafhængighed; enhver Forandring, det være Svaghedens Besvimen eller Stolthedens Kneisen, den sukkende eller den selvtilfredse, er Afhængighed. Dersom En, naar den Anden siger »jeg kan ikke elske Dig længere« stolt svarer »da kan jeg ogsaa lade være at elske Dig«: er dette Uafhængighed? Ak, det er jo Afhængighed, thi det, om han skal vedblive at elske eller ikke afhænger af, om den Anden vil elske. Men Den, som svarer »da skal jeg dog vedblive at elske Dig,« hans Kjerlighed er evig frigjort i salig Uafhængighed. Han siger det ikke stolt – afhængig af sin Stolthed, nei, han siger det ydmygt, ydmygende sig under Evighedens »skal«, netop derfor er han uafhængig.

Kun naar det er Pligt at elske, kun da er Kjerligheden evig lykkelig sikkret mod Fortvivlelse. Den umiddelbare Kjerlighed kan blive ulykkelig, kan komme til Fortvivlelse. Det kunde atter synes et Udtryk for denne Kjerligheds Styrke, at den har Fortvivlelsens Kraft, men dette er dog kun Tilsyneladelse; thi Fortvivlelsens Kraft, hvor meget den end anprises, er dog Afmagt, dens Høieste er netop dens Undergang. Dog det, at den umiddelbare Kjerlighed kan komme til at fortvivle, viser, at den er fortvivlet, at den, selv naar den er lykkelig, elsker med Fortvivlelsens Kræfter – elsker et andet Menneske »høiere end sig selv, høiere end Gud«. Om Fortvivlelse maa siges: kun Den kan fortvivle, der er fortvivlet. Naar den umiddelbare Kjerlighed fortvivler over Ulykken, da bliver det kun aabenbart, at den – var fortvivlet, at den i sin Lykke ogsaa havde været fortvivlet. Fortvivlelsen ligger i, med uendelig Lidenskab at forholde sig til et Enkelt; thi med uendelig Lidenskab kan man kun, hvis man ikke er fortvivlet, forholde sig til det Evige. Den umiddelbare Kjerlighed er saaledes fortvivlet; men naar den bliver lykkelig, som det hedder, skjules det for den, at den er fortvivlet, naar den bliver ulykkelig, bliver det aabenbart, at den – var fortvivlet. Den Kjerlighed derimod, som undergik Evighedens Forandring ved at blive Pligt, kan aldrig fortvivle, netop fordi den ikke er fortvivlet. Fortvivlelse er nemlig ikke Noget, som kan hænde et Menneske, en Begivenhed ligesom Lykke og Ulykke. Fortvivlelse er et Misforhold i hans Væsens Inderste – saa langt, saa dybt kan ingen Skjebne eller Begivenhed gribe ind, den kan kun gjøre det aabenbart, at Misforholdet – var der. Derfor er der kun een Sikkerhed mod Fortvivlelse: at undergaae Evighedens Forandring ved Pligtens »skal«; Enhver, som ikke undergik denne Forandring, han er fortvivlet; Lykke og Medgang kan skjule over det, Ulykke og Modgang derimod, ikke som han mener gjøre ham fortvivlet, men gjøre det aabenbart, at han – var fortvivlet. Forsaavidt man taler anderledes, er det, fordi man paa en letfærdig Maade forvexler de høieste Begreber. Det, der nemlig gjør et Menneske fortvivlet, er ikke Ulykken, men er, at han mangler det Evige; Fortvivlelse er at mangle det Evige; Fortvivlelse er, ikke at have undergaaet Evighedens Forandring ved Pligtens »skal«. Fortvivlelse er altsaa ikke Tabet af den Elskede, det er Ulykken, Smerten, Lidelsen; men Fortvivlelsen er Manglen af det Evige.

Hvorledes sikkres da Budets Kjerlighed mod Fortvivlelse? Ganske simpelt, ved Budet, ved dette »Du skal elske«. Deri ligger nemlig først og fremmest, Du maa end ikke elske paa den Maade, at Tabet af den Elskede vilde gjøre det aabenbart, at Du var fortvivlet, det er, Du maa slet ikke elske fortvivlet. Er dermed det at elske forbudt? Ingenlunde, det var dog vel og en underlig Tale, at det Bud, der siger »Du skal elske«, ved sin Befaling skulde forbyde at elske. Altsaa Budet forbyder kun at elske paa den Maade, som ikke er befalet; væsentligen er Budet ikke forbydende, men befalende, at Du skal elske. Altsaa Kjerligheds Budet sikkrer ikke mod Fortvivlelse ved Hjælp af matte, lunkne Trøstegrunde, at man ikke maa tage sig Noget for nær og saa videre Er og vel en saadan kummerlig Kløgt, der »har ophørt at sørge«, er den vel mindre Fortvivlelse end den Elskendes Fortvivlelse, er den ikke snarere en endnu værre Art Fortvivlelse! Nei, Kjerligheds Budet forbyder Fortvivlelsen – ved at befale at elske. Hvo skulde have dette Mod uden Evigheden, hvo er berettiget til at sige dette »skal« uden Evigheden, der netop i det Øieblik, da Kjerligheden vil fortvivle over sin Ulykke, befaler at elske! Hvor kan denne Befaling have hjemme uden i Evigheden? Thi naar det i Timeligheden er gjort umuligt at besidde den Elskede, da siger Evigheden »Du skal elske«, det er, da frelser Evigheden Kjerligheden fra Fortvivlelse netop ved at gjøre den evig. Lad det være Døden, der skiller de Tvende ad – naar da den Efterladte vil synke i Fortvivlelse: hvad skal saa hjælpe? Timelig Hjælp er en endnu sørgeligere Art Fortvivlelse; men saa hjælper Evigheden. Naar den siger »Du skal elske«, saa siger den dermed »Din Kjerlighed har en evig Gyldighed«; men den siger det ikke trøstende, thi dette vilde ikke hjælpe, den siger det befalende, just fordi der er Fare paafærde. Og naar Evigheden siger »Du skal elske«, saa bliver det jo dens Sag at indestaae for, at det lader sig gjøre. O, hvad er al anden Trøst mod Evighedens, hvad er al anden Sjelesorg mod Evighedens! Vilde den tale mildere og sige »trøst Dig«, da skulde vel den Sørgende have Indvendinger tilrede; men – ja det er ikke fordi Evigheden stolt ingen Indvending vil taale – af Omsorg for den Sørgende befaler den »Du skal elske«. Vidunderlige Trøstetale, vidunderlige Medlidenhed; thi menneskelig talt er det jo dog det Besynderligste, næsten som en Spot, at sige til den Fortvivlende, at han skal gjøre det, som var hans eneste Ønske, men hvis Umulighed bringer ham til Fortvivlelse. Behøves der noget andet Beviis for, at Kjerligheds-Budet er af guddommelig Oprindelse! Har Du forsøgt det, eller forsøg det, træd hen til en saadan Sørgende i det Øieblik, da Tabet af den Elskede vil overvælde ham, og see saa, hvad Du kan finde paa at sige; tilstaae det, Du vil trøste; det Eneste, Du ikke vil falde paa, er, at sige »Du skal elske«. Og paa den anden Side, forsøg det, om det ikke i det første Øieblik, naar det siges, næsten forbittrer den Sørgende, fordi det synes det meest Upassende at sige ved denne Leilighed. O, men Du, som gjorde den alvorlige Erfaring, Du, som i det tunge Øieblik fandt Tomheden og Modbydeligheden i de menneskelige Trøstegrunde – men ingen Trøst; Du, som med Forfærdelse opdagede, at end ikke Evighedens Formaning kunde holde Dig fra at synke: Du lærte at elske dette »skal«, som frelser fra Fortvivlelse! Hvad Du maaskee ofte i mindre Forhold havde sandet, at den sande Opbyggelse er, at der tales strengt, det lærte Du her i dybeste Forstand: at kun dette »skal« evig lykkeligt frelser fra Fortvivlelse. Evig lykkeligt – ja, thi kun Den er frelst fra Fortvivlelse, som er evig frelst fra Fortvivlelse. Den Kjerlighed, som undergik Evighedens Forandring ved at blive Pligt, er ikke fritagen fra Ulykke, men den er frelst fra Fortvivlelse, i Lykke og Ulykke lige meget frelst fra Fortvivlelse.

See, Lidenskaben ophidser, jordisk Kløgt køler, men hverken denne Hede, ei heller denne Kulde, ei heller Blandingen af denne Hede og denne Kulde er Evighedens rene Luft. Der er i denne Hede noget Hidsigt, og der er i denne Kulde noget Skarpt, og der er i Blandingen noget Ubestemt, eller en ubevidst Lumskhed som i Foraarets farlige Tid. Men dette »Du skal elske« tager alt det Usunde bort og bevarer det Sunde for Evigheden. Saaledes er det overalt, dette Evighedens »skal« er det Frelsende, det Luttrende, det Forædlende. Sid hos en dybt Sørgende; det kan et Øieblik lindre, hvis Du har denne Evne at kunne give Lidenskaben Fortvivlelsens Udtryk, som end ikke den Sørgende selv formaaer; men det er dog det Usande. Det kan et Øieblik kølende friste, hvis Du har Klogskab og Erfaring til at skaffe midlertidig Udsigt, hvor den Sørgende ingen seer; men det er dog det Usande. Derimod dette »Du skal sørge«, det er baade det Sande og det Skjønne. Jeg skal ikke have Lov til at forhærde mig mod Livets Smerte, thi jeg skal sørge; men jeg skal heller ikke have Lov til at fortvivle, thi jeg skal sørge; og dog skal jeg heller ikke have Lov til at ophøre at sørge, thi jeg skal sørge. Saaledes ogsaa med Kjerlighed. Du skal ikke have Lov at forhærde Dig mod denne Følelse, thi Du skal elske; men Du skal heller ikke have Lov at elske fortvivlet, thi Du skal elske; og lige saa lidet skal Du have Lov at forqvakle denne Følelse i Dig, thi Du skal elske. Du skal bevare Kjerligheden, og Du skal bevare Dig selv, ved og i at bevare Dig selv bevare Kjerligheden. Der, hvor det blot Menneskelige vil storme frem, der holder Budet igjen; der, hvor det blot Menneskelige vil tabe Modet, der styrker Budet; der, hvor det blot Menneskelige vil blive træt og kløgtigt, der opflammer Budet og giver Viisdom. Budet fortærer og udbrænder det Usunde i Din Kjerlighed, men ved Budet skal Du igjen kunne opflamme den, naar den menneskelig talt vilde aflade. Der hvor Du mener sagtens selv at kunne raade, der tage Du Budet med paa Raad; der hvor Du fortvivlet vil raade Dig selv, der skal Du tage Budet med paa Raad; men der hvor Du ikke veed at raade, der skal saa Budet raade for, at Alt dog bliver godt.

\sectionheading{»Du« skal elske Næsten; Vor Pligt at elske de Mennesker vi see (uddrag)  ·  SV¹ IX, 63–160}

»Du« skal elske Næsten.

Saa gaae da hen og gjør det, tag Forskjelligheden og dens Lighed bort, at Du kan elske »Næsten«. Tag Forkjerlighedens Forskjel bort, at Du kan elske Næsten. Du skal derfor ikke ophøre at elske den Elskede, o langtfra. Saa var jo ogsaa det Ord »Næsten« det største Bedrag, som er opfundet, hvis Du, for at elske Næsten, skulde gjøre Begyndelsen med at opgive at elske dem, for hvem Du har Forkjerligheden. Desuden vilde det jo ogsaa være en Modsigelse, thi, da Næsten er alle Mennesker kan dog vel ingen være udelukket – skulle vi nu sige: mindst den Elskede? Nei, thi dette er Forkjerlighedens Sprog. Altsaa er det blot Forkjerligheden, der skulde tages bort – og dog vel ikke igjen anbringes i Forholdet til Næsten, saa Du med en forskruet Forkjerlighed vilde elske Næsten i Modsætning til den Elskede. Nei, som man siger til den Eenlige: agt paa Dig selv, at Du ikke føres i Selvkjerlighedens Snare, saaledes maa der siges til de tvende Elskende: agter paa, at I ikke just ved Elskov føres i Selvkjerlighedens Snare. Thi jo mere afgjørende og udelukkende Forkjerligheden slutter sig om eet eneste Menneske, jo længere er den fra at elske Næsten. Du, Mand, før ikke Din Hustrue i den Fristelse over Dig at glemme at elske Næsten; Du, Hustrue, før ikke Din Mand i denne Fristelse! De Elskende mene vel i Elskov at have det Høieste, o, men det er ikke saaledes, thi deri have de endnu ikke det Evige betrygget ved det Evige. Vel vil »Digteren« forjætte de Elskende Udødeligheden, hvis de ere ret Elskende; men hvo er da »Digteren«, hvad hjælper det med hans Borgen, han, der ikke kan borge for sig selv. Den »kongelige Lov« derimod, Kjerligheds-Budet forjætter Livet, det evige Liv, og dette Bud siger just »Du skal elske Din Næste.« Og som dette Bud vil lære ethvert Menneske, hvorledes han skal elske sig selv, saaledes vil det ogsaa lære Elskoven og Venskabet den rette Kjerlighed: bevar i at elske Dig selv Kjerlighed til Næsten, bevar i Elskov og Venskab Kjerlighed til Næsten. Det vil maaskee støde Dig – nu, Du veed jo, at der ved det Christelige altid staaer Forargelsens Mærker. Men troe det dog; troe ei, at den Lærer, der ikke udslukkede nogen rygende Tande, at han vilde udslukke nogen ædel Ild i et Menneske; troe, at han, der var Kjerlighed, netop vil lære ethvert Menneske at elske; troe, at hvis alle Digtere forenede sig i een Sang til Elskovens og Venskabets Priis, at det var Intet, hvad de havde at sige i Sammenligning med Budet »Du skal elske, Du skal elske Din Næste som Dig selv!« Lad ikke være at troe, fordi Budet næsten forarger Dig, fordi Talen ikke lyder indsmigrende som »Digterens«, der med sine Sange indynder sig i Din Lykke, men frastødende og forfærdende, som vilde den skrække Dig ud af Forkjerlighedens elskede Tilholdssted – lad derfor ikke være at troe det, betænk, at just fordi Budet er saaledes og Talen saaledes, just derfor kan Gjenstanden være Troens Gjenstand! Hengiv Dig ikke til den Indbildning, at Du kunde tinge, at Du ved at elske nogle Mennesker, Slægt og Venner, elskede Næsten – thi dette er at opgive Digteren uden at gribe det Christelige, og for at forhindre denne Tingen var det, at Talen søgte at stille Dig mellem »Digterens« Stolthed, der foragter al Tingen, og mellem den kongelige Lovs guddommelige Majestæt, der gjør al Tingen til Skyld. Nei, elsk den Elskede trofast og ømt, men lad Kjerlighed til Næsten være det helliggjørende i Eders Forenings Pagt med Gud; elsk Din Ven oprigtigt og hengivent, men lad Kjerlighed til Næsten være hvad I lære af hinanden i Venskabets Fortrolighed med Gud! See, Døden afskaffer alle Forskjelligheder, men Forkjerlighed forholder sig altid til Forskjel, dog gaaer Veien til Livet og til det Evige gjennem Døden og gjennem Forskjellighedernes Afskaffelse: derfor fører ene Kjerlighed til Næsten i Sandhed til Livet. Som Christendommens glade Budskab er indeholdt i Læren om Menneskets Slægtskab med Gud, saa er dens Opgave Menneskets Lighed med Gud. Men Gud er Kjerlighed, derfor kunne vi kun ligne Gud i at elske, som vi ogsaa kun, efter en Apostels Ord, kunne være »Guds Medarbeidere i – Kjerlighed«. Forsaavidt Du elsker den Elskede, ligner Du ikke Gud, thi for Gud er ingen Forkjerlighed, hvad Du vel mangen Gang til Din Ydmygelse, men ogsaa mangen Gang til Din Opreisning har betænkt. Forsaavidt Du elsker Din Ven, ligner Du ikke Gud, thi for Gud er ingen Forskjel. Men naar Du elsker »Næsten«, da ligner Du Gud.

Saa gaae da hen og gjør derefter, lad Forskjellighederne fare, at Du kan elske Næsten. Ak, maaskee behøves dette end ikke at siges Dig, maaskee fandt Du i Verden ingen Elsket, ingen Ven paa Din Vei, saa Du gaaer eensomt; eller maaskee tog Gud af Din Side og gav Dig den Elskede, men Døden tog og tog hende fra Din Side, den tog igjen og tog Din Ven, men gav Dig ingen igjen, saa du nu gaaer eensomt, saa Du ingen Elsket har til at dække Din svage Side og ingen Ven ved Din høire; eller maaskee skilte Livet Eder ad, om I end bleve uforandrede – i Adskillelsens Eensomhed; ak, maaskee skilte Forandringen Eder ad, saa Du sørgende gaaer eensomt, fordi Du fandt, men fandt igjen, hvad Du fandt, forandret! Hvor trøstesløst! Ja, spørg kun »Digteren«, hvor trøstesløst det er, at leve eensomt, at have levet eensomt, uden at være elsket og uden at have nogen Elsket; spørg kun Digteren om han veed noget Andet, end at det er trøstesløst, naar Døden gaaer de Elskende imellem, eller, naar Livet skiller Ven fra Ven, eller naar Forandringen skiller dem som Fjender fra hinanden; thi vistnok elsker Digteren Eensomhed, han elsker den – for i Eensomheden at opdage Elskovens og Venskabets savnede Lykke, ligesom Den søger et mørkt Sted, der undrende vil betragte Stjernerne. Og dog, dersom det var uden egen Skyld, at et Menneske ingen Elsket fandt, og dersom han søgte, men, uden egen Skyld, forgjeves at finde en Ven, og dersom Tabet, Adskillelsen, Forandringen var uden egen Skyld: veed i saa Fald »Digteren« noget Andet, end at det er trøstesløst? Men saa er Digteren jo selv underlagt Forandring, naar han, Glædens Forkynder, paa Nødens Dag ikke veed Andet end Trøstesløshedens Klageskrig. Eller vil Du ikke kalde det Forandring, vil Du kalde det Troskab af Digteren, at han – trøstesløst sørger med den trøstesløst Sørgende: nu vel, derom ville vi ikke stride. Men hvis Du vil sammenligne denne menneskelige Trofasthed med Himlens og Evighedens, da vil Du vel selv indrømme, at den er en Forandring. Thi Himlen glæder sig ikke blot, fremfor nogen Digter, med den Glade, men Himlen sørger ikke blot med den Sørgende, nei Himlen har ny, har saligere Glæde i Beredskab for den Sørgende. Saaledes har Christendommen altid Trøst, og dens Trøst er deri forskjellig fra al menneskelig Trøst, at denne er sig bevidst kun at være en Erstatning for Tabet af Glæden: den christelige Trøst er Glæden. Menneskeligt talt er Trøst en sildigere Opfindelse: først kom Lidelsen og Smerten og Glædens Tab, og saa, bagefter, ak langt om længe kom Mennesket paa Spor efter Trøsten. Og om den Enkeltes Liv gjælder det Samme: først kommer Lidelsen og Smerten og Glædens Tab, og saa, bagefter, ak stundom langt om længe, kommer Trøsten. Men den christelige Trøst kan aldrig siges at komme bagefter, thi da den er Evighedens Trøst er den ældre end al timelig Glæde; saasnart denne Trøst kommer, kommer den med Evighedens Forspring, og opsluger ligesom Smerten, thi Smerten og Glædens Tab er det Øiebliklige – selv om dette Øieblik var Aar – er det Øieblikkelige der druknes i det Evige. Og den christelige Trøst er heller ei nogen Erstatning for Glædens Tab, da den er Glæden; al anden Glæde er dog i sidste Grund kun Trøstesløshed i Sammenligning med Christendommens Trøst. Ak, saa fuldkommen var og er Menneskets Liv ikke paa Jorden, at Evighedens Glæde kunde forkyndes ham som Glæde, han havde og han har selv forspildt det; deraf kommer det, at Evighedens Glæde kun kan forkyndes ham som Trøst. Som det menneskelige Øie ikke taaler at see Solens lys uden gjennem et dunklende Glas: saaledes kan Mennesket heller ei taale Evighedens Glæde uden gjennem dette Dunklende, at den forkyndes som Trøst. – Altsaa, hvilken Din Skjebne end var i Elskov og Venskab, hvilket end Dit Savn var, hvilket Dit Tab, hvilken Dit Livs Trøstesløshed i Fortrolighed med »Digteren«: det Høieste staaer endnu tilbage – elsk Næsten! Ham kan Du let finde, som viist blev; ham kan Du ubetinget altid finde, som viist blev; ham kan Du aldrig tabe. Thi den Elskede kan handle saaledes mod Dig, at han er tabt, og Du kan tabe en Ven; men hvad end Næsten gjør mod Dig, Du kan aldrig tabe ham. Sandt nok, Du kan ogsaa vedblive at elske den Elskede og Vennen, hvorledes De end handle mod Dig, men Du kan ikke i Sandhed vedblive at kalde dem den Elskede og Vennen, naar de, desto værre, i Sandhed have forandret sig. »Næsten« derimod kan ingen Forandring tage fra Dig, thi det er ikke »Næsten«, der holder Dig fast, men det er Din Kjerlighed, der holder »Næsten« fast; dersom Din Kjerlighed til Næsten bliver uforandret, saa bliver ogsaa Næsten uforandret ved at være til. Og Døden kan ikke berøve Dig »Næsten«, thi tager den een, saa giver Livet Dig strax een igjen. Døden kan berøve Dig en Ven, fordi Du i at elske en Ven egentlig holder sammen med Vennen; men i at elske Næsten holder Du sammen med Gud, derfor kan Døden ikke berøve Dig Næsten. – Hav derfor tabt Alt i Elskov og Venskab, hav aldrig eiet Noget af denne Lykke: Du beholder dog det Bedste tilbage i at elske Næsten.

Kjerlighed til Næsten har nemlig Evighedens Fuldkommenheder. Er det og vel en Fuldkommenhed ved Kjerligheden, at dens Gjenstand er det Fortrinlige, det Udmærkede, det Eneste? Det troede jeg var en Fuldkommenhed ved Gjenstanden, og denne Gjenstandens Fuldkommenhed som en spidsfindig Mistanke mod Kjerlighedens Fuldkommenhed. Er det en fortrinlig Egenskab ved Din Kjerlighed, om den kun kan elske det Overordentlige, det Sjeldne? Jeg troede, det var et Fortrin ved det Overordentlige og det Sjeldne, at det er det Overordentlige og Sjeldne, men ikke ved Kjerligheden. Og er Du ikke ogsaa af samme Mening? Thi har Du aldrig tænkt over Guds Kjerlighed; dersom da det var Kjerlighedens Fortrin, at elske det Overordentlige, saa var Gud, hvis jeg saa tør sige, i Forlegenhed, thi for ham er det Overordentlige slet ikke til. Det Fortrin, kun at kunne elske det Overordentlige, er altsaa snarere som en Anklage, ikke mod det Overordentlige, heller ei mod Kjerligheden, men mod den Kjerlighed, der kun kan elske det Overordentlige. Eller er det et Fortrin ved et Menneskes kjælne Befindende, at han kun kan føle sig vel paa eet eneste Sted i Verden, omgiven af enhver Begunstigelse? Naar Du seer et Menneske, der saaledes har indrettet sig i Livet, hvad er det saa Du priser? Dog vel Indretningens Beqvemmelighed. Men har Du ikke lagt Mærke til, at det virkeligt er saa, at hvert Ord i Din Lovtale over denne Herlighed, egentlig lyder som en Spot over den Stakkel, der kun kan leve i denne herlige Omgivelse? Altsaa Gjenstandens Fuldkommenhed er ikke Kjerlighedens Fuldkommenhed. Og netop fordi »Næsten« ingen af de Fuldkommenheder har, som den Elskede, Vennen, den Beundrede, den Dannede, den Sjeldne, den Overordentlige i saa høi Grad har, netop derfor har Kjerlighed til Næsten alle de Fuldkommenheder, som Kjerligheden til den Elskede, Vennen, den Dannede, den Beundrede, den Sjeldne, den Overordentlige ikke har. Lad da Verden stride saa meget den vil om, hvilken Kjerlighedens Gjenstand der er den fuldkomneste: derom kan der aldrig være Strid, at Kjerlighed til Næsten er den fuldkomneste Kjerlighed. Al anden Kjerlighed har derfor og den Ufuldkommenhed, at der er to Spørgsmaal og forsaavidt ogsaa nogen Tvetungethed: der er først Spørgsmaal om Gjenstanden og saa om Kjerligheden, eller der er dog baade Spørgsmaal om Gjenstanden og Kjerligheden. Men betræffende Kjerlighed til Næsten er der kun eet Spørgsmaal, det om Kjerligheden, og der er kun eet Evighedens Svar: denne er Kjerligheden; thi denne Kjerlighed til Næsten forholder sig ikke som en Art til andre Arter af Kjerlighed. Elskov bestemmes ved Gjenstanden, Venskab bestemmes ved Gjenstanden, kun Kjerlighed til Næsten bestemmes ved Kjerligheden. Da »Næsten« nemlig er ethvert Menneske, ubetinget ethvert Menneske, saa er jo alle Forskjelligheder borttagne fra Gjenstanden, og denne Kjerlighed altsaa just kjendelig paa, at dens Gjenstand er uden nogen Forskjellighedens nærmere Bestemmelse, hvilket vil sige, at denne Kjerlighed kun kjendes paa Kjerligheden. Er dette ikke den høieste Fuldkommenhed! Thi forsaavidt Kjerligheden kan kjendes paa og maa kjendes paa noget Andet, saa er dette Andet, i selve Forholdet, ligesom en Mistanke til Kjerligheden, at den ikke er omfattende nok og forsaavidt heller ei i evig Forstand uendelig; dette Andet er et Kjerligheden selv ubevidst Anlæg til Sygelighed. I denne Mistanke boer derfor skjult den Angest, som gjør Elskov og Venskab afhængig af sin Gjenstand, den Angest, som kan antænde Iversygen, den Angest, som kan bringe til Fortvivlelse. Men Kjerlighed til Næsten er uden Forholdets Mistanke, og kan derfor heller ikke blive Mistænksomhed i den Elskende. Dog er denne Kjerlighed ikke stolt uafhængig af sin Gjenstand, dens Ligelighed fremkommer ikke derved, at Kjerligheden stolt vender tilbage i sig selv ved Ligegyldighed mod Gjenstanden, nei, Ligeligheden fremkommer derved, at Kjerligheden ydmygt vender sig ud efter, omfattende Alle, og dog elskende Hver især, men Ingen særligen.

[…]

Betragt engang Verden, der ligger for Dig i al dens brogede Mangfoldighed; det er, som naar Du seer et Skuespil, kun at Mangfoldigheden er langt, langt større. Hver Enkelt af disse Utallige er ved sin Forskjellighed noget Bestemt, forestiller noget Bestemt, men væsentligen er han noget Andet; dog dette faaer Du ikke at see her i Livet, her seer Du ikkun, hvad den Enkelte forestiller og hvorledes han gjør det. Det er som i Skuespillet. Men naar paa Skuepladsen Teppet saa falder, saa er Den, der spillede Kongen, og Den, der spillede Tiggeren, og saaledes fremdeles hver især, de ere alle lige meget, alle Eet og det Samme: Skuespillere. Og naar i Døden Teppet er faldet for Virkelighedens Skueplads (thi dette er en forvirrende Sprogbrug, naar man taler om, at i Dødens Øieblik Teppet drages op for Evighedens Skueplads; Evigheden er nemlig ingen Skueplads, den er Sandheden), saa ere de ogsaa alle Eet, de ere Mennesker, og ere alle det, de væsentligen vare, hvad Du ikke saae paa Grund af Forskjelligheden, Du saae, de ere Mennesker. Kunstens Skueplads er som en fortryllet Verden, men tænk Dig, at det en Aften ved en almindelig Aandsfraværelse løb surr for alle de Spillende, saa de meente, at de virkeligen vare hvad de forestillede: var dette ikke, hvad man, i Modsætning til Kunstens Fortryllelse, maatte kalde en ond Aands Fortryllelse, en Forgjørelse. Og saaledes ogsaa, hvis i Virkelighedens Fortryllelse (thi vi ere jo alle fortryllede, ved at være tryllede ind hver i sin Forskjellighed) Grundtanken forvirrede sig for os, saa vi meente, at det vi forestille, det var hvad vi væsentligen ere. Ak, men er dette dog ikke just Tilfældet. At Jordlivets Forskjellighed kun er som Skuespillerens Dragt, eller kun som en Reisedragt, at Hver især burde sørge for og vogte paa, at have Baandene, hvormed dette Overtøi er befæstet, løst bundne og fremfor Alt ikke i Haardknude, paa det, at Dragten i Forvandlingens Øieblik med Lethed kan afkastes: det synes glemt. Og dog have vi jo alle Kunstforstandighed nok til at stødes, dersom paa Skuepladsen Skuespilleren i Forvandlingens Øieblik, naar han skal afkaste Forklædningen, maatte løbe ud for at faae Baandene løste. Ak, men i Virkelighedens Liv, der snører man Forskjellighedens Overtøi saa fast, at det ganske skjuler, at denne Forskjellighed er Overtøi, fordi Ligelighedens indre Herlighed aldrig eller saa saare sjeldent skinner igjennem, hvad den dog bestandigt skulde og burde. Thi Skuespillerens Kunst er den bedragende, Kunsten Bedraget, det at kunne bedrage det Store, det at lade sig bedrage det lige saa Store, derfor maa man netop ikke kunne og ikke ville gjennem Dragten see Skuespilleren; derfor er det Kunstens Høieste, naar Skuespilleren bliver Eet med det han forestiller, fordi dette er Bedragets Høieste. Men Livets Virkelighed, om den end ikke, som Evigheden, er Sandheden, burde dog være af Sandheden, og derfor burde det dog bestandigt glimte igjennem Forklædningen det Andet, som Enhver væsentligen er. Ak, men i Virkelighedens Liv, der voxer den Enkelte i Timelighedens Væxt ganske sammen med Forskjelligheden, og dette er det Modsatte af Evighedens Væxt, der voxer fra Forskjelligheden, den Enkelte forvoxer, enhver Saadan er, i Evighedens Forstand, en Vanskabning. Ak, i Virkeligheden groer den Enkelte saaledes sammen med sin Forskjellighed, at Døden tilsidst maa bruge Magt, for at rive den af ham. – Men skal man i Sandhed elske Næsten, maa det ethvert Øieblik være erindret, at Forskjelligheden er en Forklædning. Thi, som sagt, Christendommen har ikke villet storme frem for at afskaffe Forskjelligheden, hverken Fornemhedens eller Ringhedens, den har heller ei villet verdsligt træffe en verdslig Overeenskomst mellem Forskjellighederne; men den vil, at Forskjelligheden skal hænge løst om den Enkelte, løst som den Kappe Majestæten afkaster, for at vise, hvo han er; løst som den pjaltede Dragt, i hvilken et overnaturligt Væsen har skjult sig. Naar nemlig Forskjelligheden hænger saaledes løst, da skimtes bestandigt i hver Enkelt hiint væsentlige Andet, det for Alle Fælles, det evigt Lignende, Ligningen. Dersom dette var saaledes, dersom hver Enkelt levede saaledes, da havde Timeligheden naaet sit Høieste. Som Evigheden kan den ikke være; men denne forventningsfulde Høitidelighed, der uden at standse Livets Gang, hver Dag forynger sig ved det Evige og ved Evighedens Ligelighed, hver Dag frelser Sjelen ud af Forskjelligheden, i hvilken den dog bliver: dette vilde være Evighedens Afglands. Da skulde Du i Livets Virkelighed vel see Herskeren, glad og ærbødig bringe ham Din Hylding; men Du skulde dog i Herskeren see den indre Herlighed, den Herlighedens Lighed, over hvilken hans Pragt kun skjuler. Da skulde Du vel, maaskee i Sorg over ham, lidende mere end han, see Tiggeren, men Du skulde dog i ham see den indre Herlighed, Herlighedens Lighed, over hvilken hans ringe Overtøi skjuler. Ja, da skulde Du, hvor Du vendte Dit Øie hen, see »Næsten«. Thi der er og har fra Verdens Begyndelse intet Menneske været til, der er Næsten i den Forstand som Kongen er Kongen, den Lærde den Lærde, Din Slægtning Din Slægtning, det er i Særlighedens eller hvad der er det Samme i Forskjellighedens Forstand; nei, ethvert Menneske er »Næsten«. I at være Konge, Tigger, Lærd, Riig, Fattig, Mand, Qvinde og saa videre ligne vi ikke hinanden, deri ere vi jo netop forskjellige; men i at være »Næsten« ligne vi ubetinget alle hinanden. Forskjelligheden er Timelighedens Forvirrende, der mærker ethvert Menneske forskjelligt, men »Næsten« er Evighedens Mærke – paa ethvert Menneske. Tag mange Ark Papir, skriv det Forskjellige paa hvert Enkelt, saa ligner det ene ikke det andet; men tag saa igjen hvert enkelt Ark, lad Dig ikke forstyrre af Forskjellighedens Paaskrift, hold det op for Lyset, saa seer Du et fælles Mærke paa dem alle. Og saaledes er Næsten det fælles Mærke, men Du seer det kun ved Hjælp af Evighedens Lys, naar det gjennemlyser Forskjelligheden.

[…]

Naar det er Pligt i at elske at elske de Mennesker, man seer, da gjælder det om, at man i at elske det enkelte virkelige Menneske ikke underskyder en indbildt Forestilling om, hvorledes man mener eller kunde ønske, at dette Menneske skulde være. Den, som nemlig gjør det, han elsker dog ikke det Menneske, han seer, men igjen noget Usynligt, sin egen Forestilling eller andet Saadant.

Der er i Forhold til det at elske en Adfærd, som til Kjerligheden har en betænkelig Tilsætning af Tvetydighed og Kræsenhed. Eet er det jo at vrage og vrage og aldrig finde nogen Gjenstand for sin Kjerlighed, et Andet i at elske, hvad man selv kalder Gjenstanden for sin Kjerlighed, nøiagtigt og oprigtigt at fuldkomme denne Pligt at elske, hvad man seer. Sandeligen, det er jo altid et Ønske værd og atter et Ønske værd, at Den, vi skulle elske, maa være i Besiddelse af de elskelige Fuldkommenheder; vi ønske det ikke blot for vor egen Skyld, men ogsaa for den Andens. Fremfor Alt er det værd at ønske og bede, at Den, vi elske, altid maatte handle og være saaledes, at vi ganske kunne billige og samtykke det. Men i Guds Navn, lad os ikke glemme, at dette ikke er vor Fortjeneste, om han er saaledes, endnu mindre vor Fortjeneste, at fordre det af ham – skulde der være Tale om nogen vor Fortjeneste, hvad dog er usømmeligt og usømmelig Tale i Forhold til Kjerlighed, saa vilde den jo netop være, at elske lige trofast og ømt.

Men der er en Kræsenhed, der bestandigt ligesom arbeider Kjerligheden imod og vil forhindre den i at elske hvad den seer, idet Kræsenheden, usikker i Blikket og dog i en anden Forstand saa nøieregnende, forflygtiger den virkelige Skikkelse, eller støder an mod den og nu underfundigt fordrer at see noget Andet. Der gives Mennesker, om hvilke man maa sige, at de ikke have vundet Skikkelse, at deres Virkelighed ikke har fastnet sig, fordi de ere i deres Inderste uenige med sig om, hvad de ere og hvad de ville være. Men man kan ogsaa ved den Maade, paa hvilken man seer, gjøre et andet Menneskes Skikkelse vaklende, eller uvirkelig, fordi Kjerligheden, som skulde elske det Menneske den seer, ikke ret kan beslutte sig, men snart vil have en Feil bort ved Gjenstanden og snart en Fuldkommenhed til, som var, om jeg saa tør sige, Kjøbet ikke endnu ganske afsluttet. Dog Den, som i at elske saaledes er tilbøielig til at være kræsen, han elsker ikke det Menneske han seer, og gjør sig let selv sin Kjerlighed modbydelig som besværlig for den Elskede.

Den Elskede, Vennen, er jo ogsaa i almindeligere Forstand et Menneske, og er som saadan til for os Andre, men for Dig skulde han væsentligen kun være til som den Elskede, hvis Du skal fuldkomme den Pligt at elske det Menneske, Du seer. Dersom der da er en Dobbelthed i Dit Forhold, saa han deels blot er Dig i almindeligere Forstand dette enkelte Menneske, deels særligen den Elskede, da elsker Du ikke det Menneske, Du seer. Det er snarere, som havde Du i den Forstand to Øren, at Du ikke som ellers hører Eet med begge, men hører Eet med det ene, et Andet med det andet. Du hører med det ene Øre, hvad han siger, og om det nu er viist og rigtigt og skarpsindigt og aandfuldt og saa videre, ak, og kun med det andet Øre hører Du, at det er den Elskedes Stemme. Du betragter ham med det ene Øie prøvende, forskende, mønstrende, ak, og kun med det andet Øie seer Du, at han er den Elskede. O, men saaledes at dele, det er ikke at elske det Menneske, man seer. Er det ikke, som var der bestandigt en Tredie tilstede, selv naar de Tvende ere ene, en Tredie, der koldt prøver og vrager, en Tredie, der forstyrrer Inderligheden, en Tredie, der stundom dog vel maa gjøre den Paagjeldende modbydelig ved sig selv og sin Kjerlighed, at den er saaledes kræsen, en Tredie, der vilde ængste den Elskede, hvis han vidste, at denne Tredie er tilstede. Hvad betyder det ogsaa, at denne Tredie er tilstede? Betyder det, at hvis... hvis nu Dette eller Hiint ikke var efter Ønske, saa kunde Du ikke elske; betyder altsaa den Tredie Skilsmissen, Adskillelsen, saa altsaa Adskillelsens Tanke er med – i Fortroligheden, ak, ligesom naar i Hedenskabet, vanvittigt, det ødelæggende Væsen er taget med ind i Guddommens Enighed. Betyder denne Tredie, at Kjerlighedens-Forhold dog i en vis Forstand intet Forhold er, at Du staaer over Forholdet og prøver den Elskede? Betænker Du, at der i saa Fald prøves noget Andet, om Du virkeligen har Kjerlighed, eller rettere, at der er noget Andet afgjort, at Du egentligen ikke har Kjerlighed? Thi Livet har jo Prøver nok, og disse Prøver skulde just finde de Elskende, finde Ven og Ven forenede for at bestaae i Prøven. Men skal Prøven drages ind med i Forholdet, saa er der begaaet et Forræderie. Sandeligen, denne hemmelighedsfulde Indesluttethed er den farligste Art Troløshed; et saadant Menneske bryder ikke sin Tro, men han gjør det bestandigt svævende, om han er bunden ved sin Tro. Er det ikke Troløshed, naar Din Ven rækker Dig Haanden, og der da i Dit Haandtryk er et Ubestemt, som var det ham, der trykkede Din Haand, men det var tvivlsomt, om hvorvidt han i dette Øieblik svarede saaledes til Din Forestilling, at Du gjengjeldte det paa samme Maade. Er det at være i et Forhold, ligesom hvert Øieblik at begynde forfra paa at komme i Forholdet, er dette at elske det Menneske, Du seer, hvert Øieblik at see prøvende paa ham, som var det første Gang Du saae ham? Det er modbydeligt at see den Kræsne, som vrager al Mad, men det er ogsaa modbydeligt at see Den, som vel spiser den Mad, der velvilligt bydes ham, og dog i en vis Forstand ikke spiser, men ligesom bestandigt, om han dog bliver mæt, kun smager paa Maden, eller gjør sig Umage at smage en lækkrere Ret i sin Mund, medens han mættes ved den tarveligere.

Nei, skal et Menneske fuldkomme den Pligt i at elske at elske de Mennesker, han seer, da maa han ikke blot blandt de virkelige Mennesker finde dem, han elsker, men han maa udrydde al Tvetydighed og Kræsenhed i at elske dem, at han i Alvor og Sandhed elsker dem som de ere, i Alvor og Sandhed fatter Opgaven: at finde den nu engang givne eller valgte Gjenstand elskelig. Vi mene ikke hermed at anprise en barnagtig Forgabethed i den Elskedes Tilfældigheder, endnu mindre en kjælen Eftergivenhed paa urette Sted; langtfra, det Alvorlige ligger just i, at Forholdet selv vil med forenet Kraft kæmpe mod det Ufuldkomne, overvinde det Mangelfulde, fjerne det Ueensartede. Dette er Alvoren, det Kræsne er at gjøre Forholdet selv tvetydigt. Den Ene bliver ikke ved sin Svaghed eller ved sin Feil gjort fremmed for den Anden, men Foreningen betragter det Svagere som det Fremmede, hvilket er begge lige vigtigt at overvinde og fjerne. Det er ikke Dig, der paa Grund af den Elskedes Svaghed skal ligesom fjerne Dig fra ham eller gjøre Dit Forhold fjernere, men tvertimod skal de Tvende holde desto fastere og inderligere sammen, for at fjerne Svagheden. Saasnart Forholdet gjøres tvetydigt, saa elsker Du ikke det Menneske Du seer, saa er det jo som fordrede Du noget Andet, for at kunne elske; men naar Feilen eller Svagheden gjør Forholdet inderligere, ikke som skulde nu Feilen fastholdes, men just for at overvinde den, saa elsker Du det Menneske, Du seer. Du seer Feilen, men det, at Dit Forhold saa bliver inderligere, det viser just, at Du elsker det Menneske, hos hvem Du dog seer Feilen eller Svagheden eller Ufuldkommenheden.

Som der gives hykkelske Taarer, en hykkelsk Sukken og Klagen over Verden, saaledes gives der ogsaa en hykkelsk Sorg over den Elskedes Svagheder og Ufuldkommenheder. Det er saa nemt og saa blødagtigt at ønske den Elskede i Besiddelse af alle mulige Fuldkommenheder, og naar der da mangler Noget, saa er det igjen saa nemt og saa blødagtigt at sukke og sørge og blive sig selv vigtig ved sin formentlige saa rene og saa dybe Sorg. Det er overhovedet en maaskee almindeligere Form af Vellystighed, selvisk at ville stadse med den Elskede eller Vennen, og saa at ville fortvivle over enhver Ubetydelighed. Men skulde det være at elske de Mennesker man seer! O, nei de Mennesker man seer, og altsaa saaledes ogsaa med os, naar Andre see os, de ere ikke saaledes fuldkomne, og dog er det saa ofte Tilfældet, at et Menneske i sig selv udvikler denne kjælne Skrøbelighed, der kun er beregnet paa at elske Fuldkommenhedernes fuldstændige Indbegreb, og dog, uagtet vi Mennesker alle ere ufuldkomne, seer man saa sjeldent den sunde, stærke, dygtige Kjerlighed, der er beregnet paa at elske de ufuldkomnere, det er, de Mennesker, vi see.

Naar det er Pligt i at elske at elske de Mennesker vi see, da er der ingen Grændse for Kjerligheden; skal Pligten fuldkommes, da maa Kjerligheden være grændseløs, det er uforandret, hvorledes end Gjenstanden forandres.

\sectionheading{Kjerlighed opbygger (uddrag)  ·  SV¹ IX, 201–215}

Kjerlighed opbygger.

1 Cor. VIII, 1. Men Kjerligheden opbygger.

Al menneskelig Tale, selv den hellige Skrifts guddommelige Tale om det Aandelige er væsentligen overført Tale, og dette er ganske i sin, eller i Tingenes og Tilværelsens Orden, da Mennesket, om han end fra Fødselens Øieblik er Aand, dog først senere bliver sig bevidst som Aand, og saaledes sandseligt-sjeleligt har udlevet et vist Afsnit forinden. Men dette første Afsnit skal saa ikke bortkastes, naar Aanden vaagner, lige saa lidet som Aandens Vaagnen forkynder sig paa en sandselig eller sandselig-sjelelig Maade i Modsætning til det Sandselige og Sandselig-Sjelelige. Det første Afsnit bliver netop at overtage af Aanden, og saaledes benyttet, saaledes lagt til Grund, bliver det det Overførte. Det aandelige og det sandselige-sjelelige Menneske sige derfor i en vis Forstand det Samme, og dog er der en uendelig Forskjel, da den Sidste ikke aner det overførte Ords Hemmelighed, medens han dog bruger det samme Ord, men ikke overført. Der er en Verdens Forskjel mellem de Tvende, den ene har gjort Overgangen eller ladet sig føre over paa hiin Side, medens den Anden bliver paa denne Side, og dog er der det Forbindende mellem dem, at begge bruge det samme Ord. Den, i hvem Aanden er vaagnet, han forlader jo derfor ikke den synlige Verden, han er endnu bestandigt, skjøndt sig bevidst som Aand, i Synlighedens Verden og selv sandseligt synlig: saaledes bliver han ogsaa i Sproget, kun at hans Sprog er det overførte; men det overførte Ord er jo ikke et splinternyt Ord, det er tvertimod det allerede givne Ord. Som Aanden er usynlig, saa er ogsaa dens Sprog en Hemmelighed, og Hemmeligheden stikker just deri, at den bruger de samme Ord som Barnet og den Eenfoldige, men bruger dem overført, hvorved Aanden negter sig at være det Sandselige eller Sandselig-Sjelelige, men ikke negter det paa en sandselig eller sandselig-sjelelig Maade. Forskjellen er ingenlunde en Paafaldenhedens Forskjel. Vi ansee det netop derfor med Rette som et Tegn paa falsk Aandighed, at prunke med Paafaldenhedens Forskjel – hvilket just er Sandselighed, hvorimod Aandens Væsen er det Overførtes stille hviskende Hemmelighed – for Den, som har Øren at høre med.

Et af de overførte Udtryk, som den hellige Skrift hyppigst bruger, eller et af de Ord, som den hellige Skrift hyppigst bruger overført er: at opbygge. Og det er allerede – ja det er saa opbyggeligt at see, hvorledes den hellige Skrift ikke trættes ved dette simple Ord, hvorledes den ikke aandrigt eftertragter Afvexling og nye Vendinger, men tvertimod, hvad der er Aandens sande Væsen, fornyer Tanken i det samme Ord. Og det er – ja det er saa opbyggeligt at see, hvorledes Skriften med det simple Ord faaer betegnet det Høieste og paa den inderligste Maade. Det er fast som hiint Bespisningens Mirakel med det lidet Forraad, der dog ved Velsignelsen strakte rigeligen til og efterlod Overflod. Og det er – ja det er opbyggeligt, om det lykkes Nogen, istedenfor at have travlt med at gjøre nye Opdagelser, som travlt skulle fortrænge det Gamle, ved ydmygt at nøies med Skriftordet, ved taknemligt og inderligt at tilegne sig det fra Fædrene Overleverede, at stifte et nyt Bekjendtskab med det – gamle Bekjendte. Som Børn have vi vel alle ofte leget Fremmed: sandeligen, dette er just Alvor, aandeligt forstaaet at kunne vedblive denne i Alvor opbyggelige Spøg, at lege Fremmed med det gamle Bekjendte.

At opbygge er et overført Udtryk, dog ville vi nu, med denne Aandens Hemmelighed i Tanken, see, hvad dette Ord betegner i ligefrem Tale. At opbygge er dannet af »at bygge« og Tillægsordet »op«, paa hvilket altsaa Eftertrykket maa ligge. Enhver, der opbygger, han bygger, men ikke Enhver, der bygger, opbygger. Naar saaledes en Mand bygger en Fløi til sit Huus, saa siger man ikke, han opbygger en Fløi, men han bygger til. Dette »op« synes altsaa at angive Retningen i Høiden, Retningen op ad. Dog er dette heller ei Tilfældet. Dersom saaledes en Mand fører en Bygning, der er tredive Alen høi, endnu ti Alen høiere, saa sige vi dog ikke, han opbyggede Huset ti Alen høiere, vi sige han byggede til. Her begynder allerede Ordets Betydning at blive mærkelig; thi det sees, at det dog heller ikke er Høiden det kommer an paa. Dersom derimod en Mand opførte et, om end nok saa lavt og lille Huus, men fra Grunden af, saa sige vi, han opbyggede et Huus. At opbygge er altsaa fra Grunden af at opføre Noget i Høiden. Dette »op« angiver vistnok Retningen som Høide; men kun naar Høiden tillige omvendt er Dybde, sige vi at opbygge. Deraf kommer det, at hvis en Mand opbygger i Høiden og fra Grunden, men Dybden ikke rigtigen svarer til Høiden, saa sige vi vel, at han opbyggede, men tillige at han opbyggede slet, hvorimod vi ved at »bygge slet« forstaae noget Andet. Eftertrykket kommer saaledes i Forhold til det at opbygge især til at ligge paa: at bygge fra Grunden af. Vel kalde vi ikke det at bygge i Grunden at opbygge, vi sige ikke at opbygge en Brønd; men skal der være Tale om at opbygge, saa maa, ligegyldigt, hvor høi eller hvor lav Bygningen end bliver, Arbeidet være fra Grunden af. Vi kunne derfor sige om En: han begyndte at opbygge et Huus, men han blev ikke færdig. Derimod kunne vi aldrig sige om En, der føiede nok saa meget til Bygningen i Høiden, naar det ikke skete fra Grunden af, at han opbyggede. Hvor underligt! Dette »op« i Ordet at opbygge angiver Høide, men angiver Høide omvendt som Dybde; thi at opbygge er at bygge fra Grunden af. Derfor siger Skriften ogsaa om den uforstandige Mand, at han »byggede uden Grundvold«; men om den Mand, som hører Ordet til sand Opbyggelse, eller efter Skriftens Ord, Den, som hører Ordet og gjør derefter, om ham siges der, at han er liig et Menneske, der byggede et Huus »og grov dybt« (Luc. 6, 48.). Da derfor et Vandløb kom og Stormen stødte an paa dette forsvarligt opbygte Huus, da glædede vi os Alle ved det opbyggelige Syn, at Stormen ikke kunde rokke det. Ak, thi som sagt, med Hensyn til det at opbygge kommer det især an paa at bygge i Grunden. Det er priseligt, at en Mand, førend han begynder, overveier, »hvor høit han kan opføre Taarnet«, men hvis han skal opbygge, da lad ham endeligen passe paa at grave dybt; thi om Taarnet end, hvis det var muligt, hævede sig til Skyerne, dersom det var uden Grundvold, var det egentligen ikke opbygget. At opbygge aldeles uden Grundvold er umuligt, thi det er at bygge i Luften. Derfor siger man sprogligt rigtigt at bygge Luftcasteller, man siger ikke at opbygge Luftcasteller, hvilket vilde være en skjødesløs og forkeert Sprogbrug. Thi selv i Udtrykket for det Intetsigende maa der være en Overeensstemmelse mellem de enkelte Ord, som der ikke er mellem »i Luften« og »at opbygge«, da det første tager Grundvolden bort, og det sidste henholder sig til dette »fra Grunden af«; Sammensætningen vilde derfor være en usand Overdrivelse.

Saaledes med det Udtryk at opbygge i ligefrem Tale; nu erindre vi, at det er et overført Udtryk, og gaae da over til denne Overveielses Gjenstand:

Kjerlighed opbygger.

Men er det at opbygge, aandeligt forstaaet, et saaledes eiendommeligt Tillægsord for Kjerlighed, at det ene og alene tilkommer den? Det kan jo ellers gjælde om et Tillægsord, at der gives flere Gjenstande, som ligeligt, eller om end i forskjellig Grad, dog alle have Adkomst til det ene og samme Tillægsord. Dersom dette er Tilfældet med det at opbygge, saa vilde det være urigtigen saa særligen, som denne Overveielse gjør det, at fremhæve det i Forhold til Kjerlighed, det vilde jo være et Misforstaaelsens Forsøg paa at paadutte Kjerligheden en Anmasselse, som vilde den være ene om eller tilrive sig det, som den deelte med andre – og Kjerlighed er jo just villig til at dele med Andre, da den »aldrig søger sit Eget« (1 Cor. 13, 5.). Dog er det i Sandhed saa, at det at opbygge er udelukkende eiendommeligt for Kjerlighed; men paa den anden Side, denne Egenskab opbyggelig har igjen den Eiendommelighed, at den kan hengive sig i Alt, være med i Alt – just lige som Kjerlighed. Saaledes seer man, at Kjerligheden i denne sin eiendommelige Egenskab ikke afsondrer sig, ei heller trodser paa nogen Selvstændighed og Forsigværen i Række med Andet, men ganske giver sig hen; det Eiendommelige er just, at den udelukkende har den Egenskab ganske at give sig hen. Der er Intet, Intet, som jo kan gjøres eller siges saaledes, at det bliver opbyggeligt, men hvad det saa end er, naar det er opbyggeligt, saa er Kjerlighed med. Formaningen lyder derfor, netop hvor den selv indrømmer Vanskeligheden af at give bestemt Regel, »gjører Alt til Opbyggelse«. Den kunde lige saa gjerne lyde »gjører Alt i Kjerlighed«, og den havde udtrykt ganske det Samme. Det ene Menneske kan gjøre lige det Modsatte af det andet Menneske, men gjøre de hver det Modsatte – i Kjerlighed, saa bliver det Modsatte opbyggeligt. Der er i Sproget intet Ord, som i og for sig er opbyggeligt, og der er intet Ord i Sproget, som jo kan siges opbyggeligt og blive opbyggeligt, naar Kjerlighed er med. Det er derfor saa saare langtfra (ak det er just ukjerlig og splidagtig Vildfarelse), at det at opbygge skulde være Noget, der er enkelte Begavedes Fortrin, som Kundskab og Digtergave og Skjønhed og andet Saadant, at lige tvertimod ethvert Menneske ved sit Liv og Vandel, ved sin Adfærd i det Daglige, ved sin Omgang med de Jævnlige, ved sit Ord, sin Yttring burde og kunde opbygge, og vilde gjøre det, hvis Kjerligheden ret var i ham.

Dette bemærke vi ogsaa selv, thi vi bruge det Ord »opbyggelig« i det videste Omfang; men hvad vi maaskee ikke gjøre os selv Rede for, er, at vi dog kun bruge det, overalt hvor Kjerlighed er med. Men denne er den rette Sprogbrug: at være nøieregnende i ikke at bruge dette Ord uden hvor Kjerlighed er med, og ved denne Begrændsning igjen at gjøre dets Omfang ubegrændseligt, da Alt kan være opbyggeligt i samme Forstand som Kjerlighed kan være med overalt. – Naar vi saaledes see et eenligt Menneske ved priisværdig Tarvelighed sparsomt at komme ud af det med Lidet, saa ære og prise vi ham, vi glædes, vi befæstes i det Gode ved dette Syn, men vi sige egentligen ikke, at det er et opbyggeligt Syn. Naar vi derimod see, hvorledes en Husmoder, som har Mange at sørge for, ved Tarvelighed og viis Sparsommelighed kjerligt veed at lægge Velsignelse i det Lidet, saa der dog bliver nok til Alle: da sige vi, at det er et opbyggeligt Syn. Det Opbyggelige ligger i, at vi paa samme Tid, som vi see Tarveligheden og Sparsommeligheden, hvilke vi ære, see Husmoderens kjerlige Omsorg. Derimod sige vi, at det er et lidet opbyggeligt, et uhyggeligt Syn, at see Den, som paa en Maade hungrer i Overflod, og som dog slet Intet har tilovers for Andre. Vi sige det er et rystende Syn, vi væmmes ved hans Yppighed, vi gyse ved Forestillingen om Nydelsessygens rædsomme Hevn, at hungre i Overflod; men det, at vi forgjeves søge den mindste Yttring af Kjerlighed, bestemmer os, idet vi sige, det er lidet opbyggeligt. – Naar vi see en talrig Familie indkneben i en lille Leilighed, og vi dog see den beboe en hyggelig, venlig, – rummelig Leilighed, saa sige vi, det er et opbyggeligt Syn, fordi vi see den Kjerlighed, der maa være i de Enkelte og i hver Enkelt, da jo een Ukjerlig allerede var nok til at tage hele Pladsen, vi sige det, fordi vi see, at der virkelig bliver Rum, hvor der er Hjerterum. Og derimod er det saa lidet opbyggeligt at see en urolig Aand beboe Palladset uden at finde Hvile i een eneste af de mange Sale, og uden dog at kunne afsee eller undvære det mindste Aflukke. – Ja, hvad kan ikke saaledes være opbyggeligt! At see et Menneske sove skulde man ikke troe kunde være opbyggeligt. Og dog, dersom Du seer Barnet sove ved Moderens Bryst – og Du seer Moderens Kjerlighed, seer, at hun ligesom har ventet paa og nu benytter Øieblikket, medens Barnet sover, til ret at glæde sig ved det, fordi hun neppe tør lade Barnet mærke, hvor usigeligt hun elsker det: saa er det et opbyggeligt Syn. Er Moderens Kjerlighed ikke synlig, søger Du forgjeves i hendes Aasyn og Mine at opdage det mindste Udtryk af Moderkjerlighedens Glæde over eller Omsorg for Barnet, seer Du kun Dorskhed, Ligegyldighed, som er glad ved saa længe at være Barnet qvit: saa er Synet heller ikke opbyggeligt. At see Barnet alene sove er et venligt, et velgjørende, et beroligende Syn, men det er ikke opbyggeligt. Vil Du alligevel kalde det opbyggeligt, saa er det fordi Du dog alligevel seer Kjerligheden tilstede, saa er det fordi Du seer Guds Kjerlighed omsvæve Barnet. – At see den store Kunstner fuldende sit Mesterværk, det er et herligt, et opløftende Syn, men det er ikke opbyggeligt. Sæt, dette Mesterstykke var et Underværk – hvis nu Kunstneren af Kjerlighed til et Menneske slog det i Stykker: saa vilde dette Syn være opbyggeligt.

Overalt, hvor det Opbyggelige er, er Kjerlighed, og overalt, hvor Kjerlighed er, er det Opbyggelige. Derfor siger Paulus, at et Menneske uden Kjerlighed, selv om han talte med Menneskers og Engles Tungemaal, dog er som et lydende Malm og en klingende Bjelde. Hvad er vel ogsaa mindre opbyggeligt end en klingende Bjelde! Det Verdslige, hvor herligt og hvor høirøstet det end er, er uden Kjerlighed, og derfor ikke opbyggeligt; det ubetydeligste Ord, den mindste Gjerning med Kjerlighed eller i Kjerlighed er opbyggelig. Derfor opblæser Kundskab. Og dog kan jo Kundskab og Kundskabs Meddelelse ogsaa være opbyggelig; men er den det, er det fordi Kjerlighed er med. At rose sig selv synes lidet opbyggeligt, og dog kan dette jo ogsaa være opbyggeligt; gjør ikke Paulus det stundom, men han gjør det i Kjerlighed, og derfor som han selv siger: »til Opbyggelse«. Det vilde da derfor være den uudtømmeligste Tale af alle Taler, den, om hvad der kan være opbyggeligt, da Alt kan være det, det vilde være den uudtømmeligste Tale, ak, ligesom det er den sørgeligste Klage, der kan føres over Verden: at man seer og hører saa lidet Opbyggeligt. Om det nemlig er sjeldent at see Rigdom, gjør hverken fra eller til, vi ønske jo og helst at see almindelig Velstand. Om det er sjeldent at see et Mesterværk, nu, det gjør da i en vis Forstand hverken fra eller til, og i denne Henseende kunne Menneskene for største Delen selv hverken gjøre fra eller til. Anderledes med det Opbyggelige. Der lever jo i ethvert Øieblik den talløse Mængde af Mennesker, det er muligt, at Alt, hvad ethvert Menneske foretager sig, Alt, hvad ethvert Menneske siger, kan være opbyggeligt: ak, og dog er det saa sjeldent at see eller høre noget Opbyggeligt!

Kjerlighed opbygger. Lad os nu betænke hvad der blev udviklet i Indledningen, hvorved vi strax have sikkret os mod, at Talen ikke forvilder sig ved at vælge en uoverkommelig Opgave, forsaavidt Alt kan være opbyggeligt. At opbygge er at opføre Noget fra Grunden af. I den ligefremme Tale om et Huus, en Bygning, veed Enhver, hvad der forstaaes ved Grunden og Grundvolden. Men hvad er nu, aandeligt forstaaet, Aands-Livets Grund og Grundvold, som skal bære Bygningen? Det er just Kjerlighed; Kjerlighed er Alts Ophav, og aandeligt forstaaet er Kjerligheden Aands-Livets dybeste Grund. I ethvert Menneske, i hvem Kjerlighed er, er, aandeligt forstaaet, Grundvolden lagt. Og Bygningen, som, aandeligt forstaaet, skal opføres, er atter Kjerlighed; og det er Kjerlighed, som opbygger. Kjerlighed opbygger og dette er: den bygger Kjerlighed op. Saaledes er Opgaven begrændset; Talen spreder sig ikke i det Enkelte og det Mangfoldige, begynder ikke forvirret paa Noget, som den ganske vilkaarligt etsteds maatte afbryde for dog at faae en Ende, nei den samler sig og Opmærksomheden paa det Væsentlige, paa det Ene og Samme i alt det Mangfoldige, Talen bliver først og sidst om Kjerlighed, just fordi det at opbygge er Kjerlighedens eiendommeligste Bestemmelse. Kjerlighed er Grunden, Kjerlighed er Bygningen, Kjerlighed opbygger. Det at opbygge er at bygge Kjerlighed op, og det er Kjerlighed, der opbygger. Vi tale rigtignok stundom i en almindeligere Forstand om at opbygge, vi tale, i Modsætning til den Fordærvethed, som kun vil nedrive, eller i Modsætning til den Forvirrethed, som kun kan nedrive og adsplitte, om, at den Dygtige opbygger, Den, der veed at styre og lede, Den, som veed tilgavns at undervise i sit Fag, Den, der er Mester i sin Kunst. Enhver saadan opbygger i Modsætning til at nedrive. Men al denne Opbyggelse, i Kundskab, i Indsigt, i Kunstfærdighed, i Retskaffenhed og saa videre, er dog, forsaavidt den ikke bygger Kjerligheden op, ikke i dybeste Forstand Opbyggelse. Thi, aandeligt, er Kjerlighed Grunden, og at opbygge er jo at opføre fra Grunden af.

Naar altsaa Talen er om den Kjerlighedens Gjerning at opbygge, saa maa dette enten betyde, at den Kjerlige nedlægger Kjerligheden i et andet Menneskes Hjerte; eller det maa betyde, at den Kjerlige forudsætter, at Kjerligheden er i det andet Menneskes Hjerte, og just ved denne Forudsætning bygger Kjerligheden i ham op – fra Grunden af, forsaavidt han jo kjerligt forudsætter den i Grunden. Een af Delene maa være det at opbygge. Men mon nu det ene Menneske kan nedlægge Kjerlighed i det andet Menneskes Hjerte? Nei, dette er et overmenneskeligt Forhold, et utænkeligt Forhold mellem Menneske og Menneske, i den Forstand kan menneskelig Kjerlighed ikke opbygge. Det er Gud, Skaberen, som maa nedlægge Kjerlighed i ethvert Menneske, han, der selv er Kjerlighed. Det er jo derfor netop ukjerligt og ingenlunde opbyggeligt, om Nogen anmassende bilder sig ind at ville og at kunne skabe Kjerligheden i det andet Menneske; al travl og vigtig Iver i denne Henseende bygger hverken Kjerligheden op eller er selv opbyggelig. Det første Forhold af at opbygge var da utænkeligt, altsaa maae vi tænke paa det andet Forhold. Saaledes have vi vundet den Forklaring af hvad det er, at Kjerlighed opbygger, ved hvilken vi ville dvæle: den Kjerlige forudsætter, at Kjerligheden er i det andet Menneskes Hjerte, og just ved denne Forudsætning bygger han Kjerligheden i ham op – fra Grunden af, forsaavidt han jo kjerligt forudsætter den i Grunden.

Talen kan da ikke være om, hvad den Kjerlige, der vil opbygge, nu skal gjøre for at omskabe det andet Menneske, eller for at tvinge Kjerligheden frem i ham, men den er om, hvorledes den Kjerlige opbyggeligt tvinger sig selv. See, allerede dette er opbyggeligt at betænke, at den Kjerlige opbygger ved at tvinge sig selv. Kun den Ukjerlige indbilder sig at skulle opbygge ved at tvinge den Anden; den Kjerlige forudsætter bestandigt, at Kjerligheden er tilstede, just derved opbygger han. En Bygmester tænker ringe om de Steen og det Gruus, han skal bruge til Bygningen, og en Lærer forudsætter, at Discipelen er uvidende, en Tugtemester forudsætter, at det andet Menneske er fordærvet; men den Kjerlige, som opbygger, han har kun een Fremgangsmaade, at forudsætte Kjerligheden, hvad der yderligere er at gjøre kan bestandigt kun være bestandigt at tvinge sig til at forudsætte Kjerligheden. Saaledes lokker han det Gode frem, han opelsker Kjerligheden, han opbygger. Thi Kjerlighed kan og vil kun behandles paa een Maade, ved at elskes frem; at elske den frem er at opbygge. Men at elske den frem er jo netop at forudsætte, at den er tilstede i Grunden. Det kan derfor friste Mennesket at være Bygmesteren, at være Læreren, at være Tugtemesteren, fordi det synes at være at herske over Andre; men det at opbygge, som Kjerlighed gjør det, kan ikke friste, thi det er netop at være den Tjenende; derfor har kun Kjerlighed Lyst til at opbygge, fordi den er villig til at tjene. – Bygmesteren kan vise paa sit Arbeide, og sige »det er mit Værk«, Læreren paa sin Discipel; men Kjerligheden, der opbygger, har Intet at vise paa, thi dens Arbeiden bestaaer jo just blot i at forudsætte. Atter dette er saa opbyggeligt at betænke. Sæt, det lykkedes den Kjerlige at bygge Kjerligheden op i et andet Menneske, naar da Bygningen staaer der: da staaer den Kjerlige afsides for sig selv, beskæmmet siger han »det har jeg jo bestandigt forudsat«. Ak, den Kjerlige har slet ingen Fortjeneste. Bygningen bliver ikke som et Mindesmærke om Bygmesterens Kunst, eller som Discipelen en Erindring om Lærerens Underviisning; den Kjerlige har jo Intet gjort, han har kun forudsat, at Kjerligheden var i Grunden. Den Kjerlige arbeider saa stille og saa høitideligt, og dog er Evighedens Kræfter i Bevægelse; ydmygt gjør Kjerligheden sig ubemærket netop naar den arbeider meest, ja dens Arbeiden er jo som gjorde den slet Intet. Ak, Travlhed og Verdslighed er dette den største Daarskab: at det i en vis Forstand slet Intet at gjøre skulde være det vanskeligste Arbeide. Og dog er det saaledes. Thi vanskeligere er det at beherske sit Sind end at indtage en Stad, og vanskeligere at opbygge, som Kjerlighed gjør det, end at udføre det meest forbausende Værk. Er det vanskeligt i Forhold til sig selv at beherske sit Sind, hvor vanskeligt da i Forhold til et andet Menneske at tilintetgjøre sig selv ganske, og dog gjøre Alt og lide Alt. Skal det ellers være vanskeligt at begynde uden Forudsætning, sandelig, allervanskeligst er at begynde at opbygge med den Forudsætning, at Kjerligheden er tilstede, og at ende med den samme Forudsætning: saa er jo Ens hele Arbeiden forud gjort som til Intet, idet nemlig Forudsætningen først og sidst er Selvfornegtelse, eller at Bygmesteren er skjult og som Intet. Vi kunne derfor kun sammenligne denne Kjerlighedens Opbyggen med Naturens Arbeiden i det Skjulte. Medens Mennesket sover, sove Naturens Kræfter hverken Nat eller Dag; Ingen tænker paa, hvorledes de holde ud – medens Alle forlyste sig ved Engens Deilighed og Markens Frugtbarhed. Saaledes bærer Kjerligheden sig ad; den forudsætter, at Kjerligheden er tilstede, som Spiren i Kornet, og lykkes det den da at bringe det til Væxt, saa har Kjerligheden skjult sig, som den var skjult, medens den arbeidede aarle og silde. Dog dette er jo just i Naturen det Opbyggelige: Du seer al denne Herlighed, og da griber det Dig opbyggende, naar Du kommer til at tænke paa det Underlige, at Du slet ikke seer Den, som frembringer det. Dersom Du kunde see Gud med sandselige Øine, dersom han, om jeg saa tør sige, stod ved Siden af og sagde »det er mig, som har frembragt alt Dette«: saa er det Opbyggelige forsvundet.

Kjerlighed opbygger ved at forudsætte, at Kjerligheden er tilstede. Saaledes opbygger den ene Kjerlige den anden, og her er det da let nok at forudsætte den, hvor den vitterligt er tilstede. Ak, men imidlertid er Kjerlighed dog aldrig fuldkommen tilstede i noget Menneske, forsaavidt er det jo muligt at gjøre noget Andet end at forudsætte den, at opdage een eller anden Feil og Skrøbelighed ved den. Og naar saa Nogen, ukjerligt, har opdaget dette, saa vil han maaskee, som det hedder, tage det bort, tage Splinten bort, for ret at bygge Kjerligheden op. Men Kjerlighed opbygger. Den, der elsker meget, ham tilgives meget; men jo fuldkomnere den Kjerlige forudsætter Kjerligheden at være, des fuldkomnere en Kjerlighed elsker han frem. Der er i al Verdens Forhold intet Forhold, hvor der saaledes er Lige for Lige, hvor det, der kommer ud, saa nøiagtigt svarer til hvad der blev forudsat. Man gjøre ingen Indsigelse, man beraabe sig ikke paa Erfaring, thi dette er jo just ukjerligt, vilkaarligt at fastsætte en Dag, hvor det nu skal vise sig, hvad der kom ud. Sligt forstaaer Kjerligheden sig ikke paa, den er evigt forvisset om Forudsætningens Opfyldelse; er dette ikke Tilfældet, saa er Kjerligheden ifærd med at afmattes.

Kjerlighed opbygger ved at forudsætte, at Kjerligheden er tilstede i Grunden, derfor opbygger Kjerlighed ogsaa der, hvor, menneskeligt talt, Kjerligheden synes at mangle, og hvor det, menneskeligt forstaaet, synes først og fremmest fornødent at nedrive, sandeligen ikke for Lystens Skyld, men for Frelsens Skyld. Det Modsatte af at opbygge er at nedrive. Denne Modsætning viser sig aldrig tydeligere end naar Talen er om, at Kjerlighed opbygger; thi i hvilkensomhelst anden Forbindelse der tales om at opbygge har det dog den Lighed med det at nedrive, at det er at gjøre Noget ved en Anden. Men naar den Kjerlige opbygger, da er det lige det Modsatte af at nedrive, thi den Kjerlige gjør Noget ved sig selv: han forudsætter, at Kjerligheden er tilstede i det andet Menneske – hvilket dog vel er lige det Modsatte af at gjøre Noget ved det andet Menneske. At nedrive tilfredsstiller kun altfor let det sandselige Menneske; at opbygge i den Forstand, at man gjør Noget ved det andet Menneske, kan ogsaa tilfredsstille den Sandselige; men at opbygge ved at overvinde sig selv, det tilfredsstiller kun Kjerlighed. Og dog er dette den eneste Maade at opbygge paa. Men i den velmeente Ivren for at nedrive og opbygge, glemmer man, at tilsidst dog intet Menneske formaaer at lægge Kjerlighedens Grund i det andet Menneske.

See, her viser det sig just, hvor vanskelig den Bygnings-Kunst er, som Kjerligheden øver, og som findes beskreven i hiint priste Sted hos Apostelen Paulus (1 Cor. 13); thi hvad der siges om Kjerlighed er just nærmere Bestemmelser af, hvorledes den bærer sig ad med at opbygge. »Kjerlighed er langmodig«, derved opbygger den; thi Langmodighed er jo just Udholdenhed i at forudsætte, at Kjerligheden dog er i Grunden. Den, som dømmer, selv om dette skete langsomt, Den, som dømmer det andet Menneske at mangle Kjerlighed, han tager Grundvolden bort – han kan ikke opbygge; men Kjerligheden opbygger ved Langmodighed. Derfor »bærer den ikke Nid« eiheller »Nag«, thi Nid og Nag negte Kjerligheden i det andet Menneske og fortære derved, hvis det var muligt, Grundvolden. Kjerligheden, som opbygger, bærer derimod det andet Menneskes Misforstaaelse, hans Utaknemlighed, hans Vrede – det er allerede nok at bære paa, hvor skulde saa Kjerligheden tillige kunne bære Nid og Nag! Saaledes fordeles det i Verden: Den, som bærer Nid og Nag, han bærer igjen ikke det andet Menneskes Byrder, men den Kjerlige, som ikke bærer Nid og Nag, han bærer Byrderne. Hver bærer sin Byrde, den Nidske og den Kjerlige, de blive i en vis Forstand begge Martyrer, thi som en from Mand har sagt: ogsaa den Nidske er en Martyr – men Djævelens. »Kjerlighed søger ikke sit Eget«, derfor opbygger den. Thi Den, som søger sit Eget, han maa jo skaffe alt Andet tilside, han maa rive ned for at faae Plads for sit Eget, som han vil opbygge. Men Kjerlighed forudsætter, at Kjerligheden er tilstede i Grunden, derfor opbygger den. »Den glæder sig ikke over Uretfærdighed«, men Den, som vil nedrive, eller dog vil blive sig selv vigtig ved den Mening, at det er nødvendigt at rive ned, han maa jo siges at glæde sig ved Uretfærdighed – ellers var der jo ikke Noget at nedrive. Kjerlighed derimod glæder sig ved at forudsætte, at Kjerligheden er tilstede i Grunden, derfor opbygger den. »Kjerlighed fordrager Alt«, thi hvad er det at fordrage Alt, det er i Alt dog tilsidst at finde Kjerligheden, som forudsættes i Grunden. Naar vi sige om et Menneske, der har en meget stærk Helbred, at han, i Forhold til Spise og Drikke, kan fordrage Alt, saa mene vi dermed, at hans Sundhed faaer det Nærende ud selv af det Usunde (ligesom den Sygelige har Skade selv af den sunde Spise), vi mene, at hans Sundhed faaer Næring ud af hvad der mindst synes nærende. Saaledes fordrager Kjerlighed Alt, bestandigt forudsættende, at Kjerligheden dog er tilstede i Grunden – og derved opbygger den. »Kjerlighed troer Alt«, thi at troe Alt er jo just, skjøndt det ikke sees, ja skjøndt det Modsatte sees, at forudsætte, at Kjerligheden dog er tilstede i Grunden, selv i den Vildledte, selv i den Fordærvede, selv i den Hadefuldeste. Mistroiskhed tager just Grundvolden bort ved at forudsætte, at Kjerlighed ikke er tilstede, derfor kan Mistroiskhed ikke opbygge. »Kjerlighed haaber Alt«, men at haabe Alt, det er jo, skjøndt det ikke sees, ja skjøndt det Modsatte sees, at forudsætte, at Kjerligheden dog er tilstede i Grunden, og at den nok vil vise sig, i den Vildfarende, i den Vildledte, selv i den Fortabte. See, den forlorne Søns Fader var maaskee den Eneste, som ikke vidste, at han havde en forloren Søn, thi Faderens Kjerlighed haabede Alt. Broderen vidste strax, at han var haabløst forloren. Men Kjerlighed opbygger; og Faderen vandt igjen den forlorne Søn, just fordi han, som haabede Alt, forudsatte, at Kjerligheden var tilstede i Grunden. Der blev trods Sønnens Forvildelse intet Brud fra Faderens Side (og et Brud er jo lige det Modsatte af at opbygge) han haabede Alt, derfor opbyggede han i Sandhed ved sin faderlige Tilgivelse, just fordi Sønnen ret levende fornam, at den faderlige Kjerlighed havde holdt ud med ham, saa der intet Brud havde været. »Kjerlighed taaler Alt.« Thi at taale Alt er just at forudsætte, at Kjerligheden er tilstede i Grunden. Naar vi sige, at Moderen taaler alle Barnets Uartigheder, sige vi saa dermed, at hun som Qvinde betragtet taalmodigt lider Ondt. Nei, vi sige noget Andet, at hun som Moder bestandigt bliver ved at huske paa, at det er Barnet, og altsaa ved at forudsætte, at Barnet dog nok elsker hende, og at det nok vil vise sig. Ellers talte vi jo og om, hvorledes Taalmod taaler Alt, ikke hvorledes Kjerlighed taaler Alt. Thi Taalmod taaler Alt og tier, og dersom Moderen saaledes taalte Barnets Uartigheder, saa sagde vi egentligen dermed, at Moderen og Barnet dog vare blevne hinanden fremmede. Men Kjerlighed taaler Alt, tier taalmodigt – men forudsætter i al Stilhed, at Kjerligheden dog nok er tilstede i det andet Menneske.

Saaledes opbygger Kjerlighed. »Den opblæses ikke, den bruger ikke Fremfusenhed, den forbittres ikke.« Den er ikke opblæst af den Mening, at den skulde skabe Kjerligheden i det andet Menneske, den er ikke forbittret og fremfusende, utaalmodigt, næsten haabløs travl med hvad den først maa nedrive for saa at bygge op igjen. Nei, den forudsætter bestandigt, at Kjerligheden er tilstede i Grunden. Derfor er det ubetinget det opbyggeligste Syn at see Kjerlighed opbygge, et Syn som selv Engle opbygges ved, og derfor er det ubetinget det Opbyggeligste, hvis det lykkedes et Menneske retteligen at tale om: hvorledes Kjerlighed opbygger. Der er mangt et venligt, mangt et velgjørende, mangt et fortryllende, mangt et gribende, mangt et opløftende, mangt et fængslende, mangt et overtalende Syn og saaledes fremdeles; men der er kun eet opbyggeligt Syn, at see Kjerlighed opbygge. Hvad Du end derfor har seet af Rædsel eller Afskyelighed i Verden, som Du ønskede at kunne glemme, fordi det vil nedbryde Dit Mod, Din Tillid, give Dig Afsmag for Livet og Væmmelse ved at leve: betænk blot, hvorledes Kjerlighed opbygger, og Du er opbygget til at leve. Der er saare mange forskjellige Gjenstande at tale om, men der er kun een opbyggelig: hvorledes Kjerlighed opbygger. Hvad der derfor end kan være hændet Dig saa forbittrende, at Du kunde ønske aldrig at være født og jo før jo hellere at forstumme i Døden: betænk blot, hvorledes Kjerlighed opbygger, og Du er atter opbygget til at tale. Der er kun eet opbyggeligt Syn og kun een opbyggelig Gjenstand; dog kan Alt siges og gjøres opbyggeligt, thi overalt, hvor det Opbyggelige er, er Kjerlighed, og overalt, hvor Kjerlighed er, er det Opbyggelige, og saasnart Kjerlighed er tilstede, opbygger den.

Kjerlighed opbygger ved at forudsætte, at Kjerligheden er tilstede. Har Du ikke selv erfaret dette min Tilhører? Dersom nogensinde noget Menneske har talt saaledes til Dig eller handlet saaledes mod Dig, at Du ret følte Dig opbygget derved, da var det fordi, Du ret levende fornam, hvorledes han forudsatte Kjerligheden at være i Dig. Eller hvorledes tænker Du Dig vel ogsaa det Menneske skulde være, der i Sandhed kunde opbygge Dig? Ikke sandt, Du vilde ønske ham Indsigt og Kundskab og Gave og Erfaring, men dette meente Du dog ikke, at det afgjørende kom an paa, derimod paa, at han var et tilforladelig kjerligt Menneske, det er, i Sandhed et kjerligt Menneske. Du mener altsaa, at det for at opbygge afgjørende og væsentligt kommer an paa at være kjerlig eller at have Kjerlighed, i saadan Grad, at man kan forlade sig derpaa. Men hvad er nu Kjerlighed? Kjerlighed er at forudsætte Kjerlighed; at have Kjerlighed er at forudsætte Kjerlighed hos Andre, at være kjerlig er at forudsætte, at Andre ere kjerlige. Lad os forstaae hinanden. De Egenskaber, et Menneske kan have, maae enten være Egenskaber han har for sig, selv om han gjør Brug af dem mod Andre, eller Egenskaber for Andre. Viisdom er en for sig værende Egenskab, Magt og Gave og Kundskab og saa videre ligeledes for sig værende Egenskaber. At være viis vil ikke sige, at forudsætte, at Andre ere Vise, tvertimod kan det jo være meget viist og sandt, om den i Sandhed Vise antager, at langtfra alle Mennesker ere vise. Ja, fordi »viis« er en forsigværende Egenskab, er der i Tanken ikke Noget til Hinder for at antage, at der kunde leve eller have levet en Viis, der turde sige, at han antog alle Andre for uvise. I Tanken (at være viis – og at antage, at alle Andre ere uvise) er ingen Modsigelse. I Livets Virkelighed vilde en saadan Yttring være Hovmod, men i Tanken blot som saadan er der ingen Modsigelse. Dersom derimod En vilde mene, at han var kjerlig, men tillige, at alle Andre ikke vare kjerlige, saa vilde vi sige: nei stop, her er en Modsigelse i selve Tanken; thi at være kjerlig er jo just at antage, at forudsætte, at andre Mennesker ere kjerlige. Kjerlighed er ikke en forsigværende Egenskab, men en Egenskab, ved hvilken eller i hvilken Du er for Andre. I daglig Tale sige vi jo rigtignok, idet vi opregne et Menneskes Egenskaber, at han er viis, forstandig, kjerlig – og vi mærke ikke, hvilken Forskjel der er mellem den sidste Egenskab og de første. Hans Viisdom, hans Erfarenhed, hans Forstandighed har han for sig, selv om han gjør Andre Gavn med dem; men hvis han er i Sandhed kjerlig, da er det ikke ham, der har Kjerlighed, i samme Forstand som han har Viisdom, men det er just hans Kjerlighed, at forudsætte, at vi Andre have Kjerlighed. Du priser ham som den Kjerlige, Du mener, at det er en Egenskab han har, som det ogsaa er, Du føler Dig opbygget ved ham, just fordi han er kjerlig, men Du mærker ikke, at Forklaringen er, at hans Kjerlighed betyder, at han forudsætter Kjerlighed hos Dig, og at just derved opbygges Du, at just derved Kjerligheden i Dig bygges op. Dersom det virkeligt var saa, at et Menneske kunde være kjerlig uden at dette betydede at forudsætte Kjerlighed hos Andre, saa vilde Du heller ikke i dybeste Forstand føle Dig opbygget, hvor tilforladeligt det end var, at han var kjerlig, Du vilde ikke i dybeste Forstand føle Dig opbygget, saa lidet som Du i dybeste Forstand opbygges, hvor tilforladeligt det end er, at han er viis, forstandig, erfaren, lærd. Dersom det var muligt, at han i Sandhed kunde være kjerlig uden at dette betydede at forudsætte Kjerlighed hos Andre, saa kunde Du heller ikke ganske forlade Dig paa ham, thi det Tilforladelige ved den Kjerlige er netop, at selv naar Du mistvivler over Dig selv, om der er Kjerlighed i Dig, er han kjerlig nok til at forudsætte, eller rettere er han den Kjerlige, der forudsætter den. – Men Du fordrede jo, at et Menneske for i Sandhed at opbygge, skulde være i Sandhed kjerlig. Og at være kjerlig har nu viist sig betydende: at forudsætte Kjerlighed hos Andre. Saa siger Du jo ganske det Samme, som Talen har udviklet.

Saaledes vender Overveielsen tilbage til sin Begyndelse. At opbygge er at forudsætte Kjerlighed; at være kjerlig er at forudsætte Kjerlighed; kun Kjerlighed opbygger. Thi at opbygge er at opføre Noget fra Grunden af, men, aandeligt, er Kjerlighed Alts Grund. Lægge Kjerlighedens Grund i et andet Menneskes Hjerte kan intet Menneske, dog er Kjerligheden Grunden, og opbygge kan man kun fra Grunden af, altsaa kan man kun opbygge ved at forudsætte Kjerlighed. Tag Kjerlighed bort, saa er der Ingen, som opbygger, og Ingen, som opbygges.

\workheading{Christelige Taler}{1848}

\worknote{Antologiens ''Christian Discourses''. Hong-udvalget bringer tre hele taler: ''Ringhedens Bekymring'' (af Første Deel, Hedningenes Bekymringer), ''Det Glædelige i: at Modgang er Medgang'' (af Anden Deel, Stemninger i Lidelsers Strid) og fredags-nadvertalen over 1 Joh. 3,20 (af Fjerde Deel).}

\sectionheading{Ringhedens Bekymring  ·  SV¹ X, 42–51}

Ringhedens Bekymring.

Bekymrer Eder ikke for, hvad I skulle iføres – efter alt saadant søge Hedningene.

Denne Bekymring har Fuglen ikke.Spurvene inddeles i Graa-Spurve og Guld-Spurve, men denne Forskjel, denne Inddeling: ringe – fornem, er ikke til for dem eller for nogen eneste af dem. Den Fugl, der i Flokken flyver foran, eller til Høire, den følge vel de andre, der er Forskjel paa Forrest og Bagest, paa Høire og Venstre, men den Forskjel: ringe – fornem er ikke til; i Flugtens dristige Sving, naar Flokken yndigt svinger sig frit i de luftige Skikkelser, vexler det ogsaa med Først og Sidst, med Høire og Venstre. Og naar de tusinde Stemmer synge i Chor, saa er der vel een, som slaaer Tonen an, denne Forskjel er der, men: ringe – fornem, denne Forskjel er ikke, og frit lever Glæden i Stemmernes Vexlen. Det fornøier »den Enkelte« saa ubeskriveligt at synge i Chor med de Andre, men dog synger den ikke for at fornøie de Andre, den fornøier sig selv med sin Sang og med de Andres Sang, bryder derfor ganske vilkaarligt af, holder et Øieblik inde, indtil det atter lyster den at stemme i med – og at høre sig selv.

Fuglen har da ikke denne Bekymring. Hvoraf kommer nu det? Det kommer deraf, at Fuglen er hvad den er, er sig selv, nøiet med at være sig selv, fornøiet med sig selv. Fuglen veed neppe tydeligt eller er ganske paa det Rene med sig selv om, hvad den er, end mindre, at den skulde vide Noget om Andre. Men den er fornøiet med sig selv, og med hvad den er – hvad dette saa end forresten er; thi det har den ikke Tid til at tænke over, eller til blot at begynde paa, saa fornøiet er den ved at være det den er. For at være, for at have Glæde af at være, skal den ikke gaae ad den lange Vei ved først at faae Noget at vide om de Andre, for derved at faae at vide, hvad den selv er. Nei den har Sit fra første Haand, den skyder den fornøieligste Gjenvei; den er hvad den er; for den er der intet Spørgsmaal om at være eller ikke at være; ved Hjælp af Gjenveien slipper den forbi alle hiin Forskjelligheds Bekymringer. Om den er Fugl ligesom alle andre Fugle, om den er »lige saa god Fugl« som de andre af samme Art, ja blot om den er ligesom dens Mage: alt Sligt tænker den slet ikke paa, saa utaalmodig er den i sin Glæde over at være. Ingen ung Pige, der staaer paa Springet og skal til Dands, kan være saa utaalmodig efter at komme afsted, som Fuglen er efter at være i Gang med at være hvad den er. Thi den har intet, ikke det korteste Øieblik at give bort, naar dette skulde sinke den fra at være; det korteste Øieblik vilde for den være en dræbende lang Tid, naar den i det ikke maatte være hvad den er; den vilde døe for Utaalmodighed, hvis der gjordes den den mindste mindste Indsigelse mod at faae Lov til uden videre at være. Den er, hvad den er, men den er; den lader fem være lige, og saa er den. Dette er dog vel saaledes. Saae Du end ikke den kongelige Fugls stolte Flugt, naar Du seer den lille Fugl, der sidder og gynger sig paa et Ax, og morer sig selv med at synge: er der vel mindste Spor af Ringhedens Bekymring! Thi Du vil dog vel ikke indvende, hvad jo netop er det Lærerige: at den er høit paa Straa. Og vil Du det, saa tag Straaet paa hvilket den sidder. Mere livlig end Lilien i Glæde over at være er Fuglen dog ganske som Lilien i uskyldig Selvtilfredshed. Saae Du end ikke den pragtfulde Lilie, der kneiser ydmygt i al sin Deilighed, naar Du seer den uanseelige, der staaer i en Grøft og narres med Vinden som var de To Ligemænd; naar Du seer den, efterat et Uveir har gjort Alt for at lade den føle dens Ubetydelighed – dersom Du betragter den, idet den atter stikker Hovedet frem for at see, om det dog ikke snart skulde til at ville være godt Veir igjen: synes Dig, at der er mindste Spor af Ringhedens Bekymring! Eller naar den staaer ved det mægtige Træes Fod og seer med Forundring op ad det, synes Dig, at der er mindste Spor af Ringhedens Bekymring hos denne, den forundrede Lilie; eller troer Du den vilde føle sig selv mindre, hvis Træet var end een Gang saa stort! Er det ikke snarere, som havde den i al Uskyldighed den Indbildning, at Alt var til for dens Skyld!

Saa let har Fuglen og Lilien ved at være, saa let tage de sig i det at leve, saa naturligt falder Begyndelsen dem, eller det at komme til at begynde. Thi dette er Liliens og Fuglens lykkelige Begunstigelse, at det er gjort dem saa let med at begynde paa at være til, at de, strax de ere blevne til, strax ere begyndte, strax ere i fuld Fart med at være, saa der slet ikke behøves, at Noget gaaer forud for Begyndelsen, saa de slet ikke forsøges i hiin blandt Menneskene meget omtalte og som saa farefuld skildrede Vanskelighed, Begyndelsens Vanskelighed.

Men hvorledes er nu Fuglen Læremesteren, hvor er Underviisningens Tilknytningspunkt? Mon ikke deri: at gjøre Omveien efter Begyndelsen, det er efter at finde Begyndelsen, denne Omvei, der kan blive saa saare lang, at gjøre den saa kort som muligt, for saa hurtigt som muligt at komme til sig selv, at være sig selv.

Denne Bekymring har den ringe Christen ikke. Dog er han deri forskjellig fra Fuglen, at han maa forsøges i denne Begyndelsens Vanskelighed; thi han er vidende om den Forskjel: ringe – fornem. Han veed, og han veed, at de Andre vide det Samme om ham, at han er et ringe Menneske, og han veed, hvad dette vil sige. Han veed ogsaa hvad der forstaaes ved Jordlivets Fortrin, hvilke disse saa saare forskjellige ere, ak og at de alle ere ham nægtede, at de, medens de ellers ere til for at udtrykke, hvad de Andre ere i disse Fortrin, i Forhold til ham ere ligesom for at betegne, hvor ringe han er. Thi med hvert et Fortrin, den Fornemme føier til, desto fornemmere bliver han, og med hvert et Fortrin den Ringe maa tilstaae at være ham nægtet, desto ringere bliver han jo paa en Maade. Hvad der er til for at betegne hvor meget den Fornemme er, det synes fra den anden Side at være til for at betegne, hvor lidet den Ringe er. – O, vanskelige Begyndelse paa at være til eller for at komme til at være til: at være til, derpaa at blive til, for saa først at være til; o, listigt skjulte Snare, som ikke spændes efter nogen Fugl! Thi det seer jo ud, som skulde et Menneske, for at begynde paa at være sig selv, først blive færdig med, hvad de Andre ere, og derved faae at vide, hvad han da selv er – for saa at være det. Dog, gaaer han i denne Øienforblindelsens Snare, saa kommer han netop aldrig til at blive sig selv. Han gaaer da fremad og fremad som Den, der gaaer ad en Vei, om hvilken de Forbigaaende sige ham, at den ganske rigtigt fører til Byen, men glemme at sige ham, at hvis han vil til Byen maa han vende sig om; thi han gaaer ad den Vei, som fører til Byen, ad den Vei gaaer han – bort fra Byen.

Dog den ringe Christen gaaer ikke i denne Øienforblindelsens Snare, han seer med Troens Øine, han er med Troens Hurtighed, der søger Gud, ved Begyndelsen, er for Gud sig selv, nøiet med at være sig selv. Han har faaet at vide af Verden eller af de Andre, at han er et ringe Menneske, men han har ikke givet sig hen i denne Viden, han taber sig ikke verdsligt i den, gaaer ikke ganske op i den; ved med det Eviges Tilbageholdenhed at holde sig til Gud er han bleven sig selv. Han er som Den, der har to Navne, eet for alle Andre, og et andet for sine Nærmeste. I Verden, i Omgang med de Andre er han det ringe Menneske, Andet giver han sig ikke ud for, og Andet tages han heller ikke for; men for Gud er han sig selv. Det seer i Omgangens Berøring bestandigt ud som skulde han hvert Øieblik vente paa af de Andre at faae at vide, hvad han nu er i dette Øieblik. Men han venter ikke; han haster med at være for Gud, fornøiet med for Gud at være sig selv. Han er et ringe Menneske med i Menneskenes Mængde, og hvad han saaledes er, afhænger af Forholdet; men han er i at være sig selv ikke afhængig af Mængden, han er for Gud sig selv. Thi af »de Andre« faaer et Menneske naturligviis dog egentligen kun at vide, hvad de Andre ere – det er paa den Maade Verden vil bedrage et Menneske fra at blive sig selv. »De Andre« veed igjen heller ikke, hvad de selv ere, men bestandigt kun hvad »de Andre« ere. Der er kun Een, som ganske kjender sig selv, som i og for sig selv veed, hvad han selv er, det er Gud; og han veed ogsaa, hvad ethvert Menneske i sig selv er, thi det er han netop ene ved at være for Gud. Det Menneske, der ikke er for Gud, er heller ikke sig selv, hvad man kun kan være ved at være i Den der er i og for sig selv. Man kan, naar man er sig selv ved at være i Den, som er i og for sig selv, være i Andre eller for Andre, men man kan ikke ved at være blot for Andre være sig selv.

Den ringe Christen er for Gud sig selv. Saaledes er Fuglen ikke sig selv; thi Fuglen er hvad den er. Ved Hjælp af denne Væren er den i ethvert Øieblik sluppen forbi Begyndelsens Vanskelighed. Men saa naaede den heller ikke til den vanskelige Begyndelses herlige Ende: i Fordoblelse at være sig selv. Fuglen er som en Ener, Mennesket, der er sig selv, er mere end en Tier. Fuglen slipper heldigt forbi Begyndelsens Vanskelighed, faaer derfor ingen Forestilling om, hvor ringe den er. Men saa er den jo netop uden Sammenligning ringere end Den, som veed, hvor ringe han er. Ringhedens Forestilling er ikke til for Fuglen; men den ringe Christen er ikke væsentligen til for denne Forestilling, han vil ikke væsentligen være til for den, thi væsentligen er han og vil han for Gud være sig selv. Fuglen er altsaa egentligen den Ringe. Den ringe Christen er i Modsætning til sin Ringhed sig selv, uden dog derfor taabeligt at ville ophøre at være det ringe Menneske, han er i Forhold til Andre; han er i Ringheden sig selv. Og saaledes er den ringe Christen i Ringhed uden Ringhedens Bekymring. Thi hvori ligger Ringhed? I Forholdet til »de Andre«. Men hvori grunder dens Bekymring? I ikkun at være til for de Andre, i ikke at vide om Andet end om Forholdet til de Andre. Fuglen er slet ikke vidende om Forholdet til de Andre, forsaavidt ikke Ringe, og igjen forsaavidt uden Ringhedens Bekymring; men den er naturligviis saa heller ikke vidende om, at den har et høiere Forhold.

Og hvad er saa den ringe Christen, som for Gud er sig selv? Han er Menneske. Forsaavidt han er Menneske er han i en vis Forstand ligesom Fuglen, der er hvad den er. Dog dette ville vi her ei videre dvæle ved.

Men han er tillige Christen, hvad der jo ligger i Spørgsmaalet om, hvad den ringe Christen er. Forsaavidt er han ikke som Fuglen. Thi Fuglen er hvad den er. Men saaledes kan man ikke være Christen. Er man Christen, saa maa man være bleven det. Den ringe Christen er altsaa bleven til Noget i Verden. Ak, Fuglen, den kan ikke blive til Noget, den er, hvad den er. Den ringe Christen var, som Fuglen, Menneske, men saa blev han Christen. Han blev til Noget i Verden. Og han kan bestandigt blive Mere og Mere; thi han kan bestandigt blive mere og mere Christen. Som Menneske var han skabt i Guds Billede; men som Christen har han Gud til Forbillede. Denne uroligende Tanke, der i eet væk kalder paa En, et Forbillede, den kjender Fuglen ikke; den er hvad den er, Intet, Intet forstyrrer den denne dens Væren. Ja, det er sandt, Intet forstyrrer den – end ikke den salige Tanke at have Gud til sit Forbillede. Et Forbillede er rigtignok en Opfordring, men hvilken Salighed! Vi tale jo allerede om Lykke, naar vi tale om, at der inden i Digteren er Noget, som opfordrer ham til Sang. Men Forbilledet fordrer endnu stærkere, tilskynder hver Den end stærkere, hvem det kom for Øie, for hvem det er til. Forbilledet er en Forjættelse, ingen anden er saa sikker, thi Forbilledet er jo Opfyldelsen. – For Fuglen er der intet Forbillede, men for den ringe Christen er Forbilledet til, og han er til for sit Forbillede – han kan bestandigt komme til at ligne det mere og mere.

Den ringe Christen, der for Gud er sig selv, er som Christen til for sit Forbillede. Han troer, at Gud har levet paa Jorden, at Han har ladet sig føde i ringe og fattige Kaar, ja i Vanære, og derpaa som Barn levet sammen med den simple Mand, der kaldtes Hans Fader, og den foragtede Jomfrue, der var Hans Moder. At Han derpaa vandrede omkring i en Tjeners ringe Skikkelse, ikke til at adskille fra andre ringe Mennesker end ikke ved sin paafaldende Ringhed, indtil Han endte i den yderste Elendighed, korsfæstet som en Forbryder – og saa rigtignok efterlod sig et Navn; men den ringe Christens Begjæring er jo blot at turde i Liv og Død tilegne sig Hans Navn, eller Navn efter Ham. Den ringe Christen troer, som det er fortalt, at Han til sine Disciple valgte ringe Mennesker af den simpleste Stand, og at Han til sin Omgang søgte dem, som Verden forskjød og foragtede; at Han gjennem sit Livs forskjellige Omvexlinger, da Menneskene vilde høit ophøie Ham, og da de vilde fornedre Ham om muligt end mere, end Han selv havde fornedret sig, blev de ringe Mennesker tro, til hvem Han var knyttet ved nærmere Forbindelser, de ringe Mennesker tro, hvem Han havde knyttet til sig, de foragtede Mennesker tro, hvem man just udstødte af Synagogen, fordi Han havde hjulpet dem. Den ringe Christen troer, at dette ringe Menneske eller at dette Hans Liv i Ringhed har viist, hvad et ringe Menneske har at betyde, ak, og hvad et, menneskelig talt, fornemt Menneske dog egentligen har at betyde, hvor uendelig Meget det kan betyde at være et ringe Menneske, og hvor uendeligt Lidt det at være et fornemt Menneske, hvis man ikke er Andet. Den ringe Christen troer, at dette Forbillede just er til for ham, han, som jo er et ringe Menneske, kjempende maaskee med Armod og trange Kaar, eller det endnu Ringere, foragtet og forstødt. Han tilstaaer vel, at han jo ikke er i det Tilfælde selv at have valgt denne overseete eller foragtede Ringhed, og forsaavidt ikke ligner Forbilledet; men dog fortrøster han sig til, at Forbilledet er til for ham, Forbilledet, der ved Hjælp af Ringheden barmhjertigt ligesom paanøder sig ham, som vilde det sige »ringe Menneske, kan Du ikke see, at dette Forbillede er for Dig«. Vel har han ikke med egne Øine selv seet Forbilledet, men han troer, at Han har været til. Der havde jo i en vis Forstand heller ikke været Noget at see – uden Ringheden (thi Herligheden, den maa troes); og Ringheden kan han da nok gjøre sig en Forestilling om. Han har ikke med egne Øine seet Forbilledet, han gjør heller ikke noget Forsøg paa at lade Sandsen danne sig et saadant Billede. Dog seer han ofte Forbilledet. Thi hver Gang han i Troens Glæde over dette Forbilledes Herlighed ganske forglemmer sin Armod, sin Ringhed, sin Foragtethed, da seer han Forbilledet – og da seer han selv tilnærmelsesviis ud som Forbilledet. Dersom nemlig i et saadant saligt Øieblik, naar han er fortabt i sit Forbillede, et andet Menneske seer paa ham, saa seer dette andet Menneske kun et ringe Menneske for sig: saaledes var det just ogsaa med Forbilledet, man saae kun det ringe Menneske. Han troer og haaber stedse mere og mere at skulle nærme sig til Lighed med dette Forbillede, der først hisset skal vise sig i sin Herlighed; thi her paa Jorden kan det kun være i Ringhed og kun sees i Ringheden. Han troer, at dette Forbillede, hvis han stadigt kæmper for at ligne det, bringer ham anden Gang og i endnu nøiere Slægtskab med Gud, at han ikke blot har Gud til Skaber, som alle Skabninger have, men har Gud til Broder.

Men saa er jo denne ringe Christen dog noget meget Ophøiet? Ja vist er han det, noget saa Ophøiet, at man ganske taber Fuglen af Sigte. Han er, som Fuglen, ringe uden Ringhedens Bekymring, trykket i en vis Forstand, som Fuglen ikke er det, ved Bevidstheden om sin Ringhed – dog er han høit ophøiet. Men ringheden taler han ikke om, og aldrig da bedrøvet, den minder ham jo blot om Forbilledet, medens han tænker paa Forbilledets Ophøiethed – og naar han gjør det, saa minder han selv tilnærmelsesviis om Forbilledet.

Men den ringe Hedning, han har denne Bekymring.

Den ringe Hedning, han er uden Gud i Verden, og altsaa aldrig væsentlig sig selv, (hvad man kun er ved at være for Gud) og altsaa heller aldrig nøiet med at være sig selv, hvad man dog vel ikke er, naar man ikke er sig selv. Han er ikke sig selv, ikke nøiet med at være sig selv, ei heller, som Fuglen, det han er: han er misfornøiet med det han er; sig selv en Plage sukker han over, og anklager han sin Skjebne.

Hvad er han da? Han er den Ringe, slet ikke Andet, det er, han er hvad »de Andre« gjøre ham til, og hvad han gjøre sig selv til ved kun at være for Andre. Hans Bekymring er: Intet at være, ja slet ikke at være. Saa langtfra er det, at han er som Fuglen, der er hvad den er. Og derfor er hans Bekymring igjen: at blive Noget i Verden. At være til for Gud, det er ikke Noget, mener han – det tager sig heller ikke ud i Verden ved Modsætning til eller i Sammenligning med Andre. At være Menneske, det er ikke Noget at være, mener han – det er jo at være Ingenting; thi deri er da slet ingen Forskjel fra eller Fortrin for alle andre Mennesker. At være Christen, det er ikke Noget at være, mener han – det er vi jo Alle. Men at blive Justitsraad – det var Noget at være. Og Noget maa han for Alt blive til i Verden, det er til at fortvivle over, slet intet at være.

»Det er til at fortvivle over«, han taler som var han ikke allerede fortvivlet. Den ringe Hedning er fortvivlet, og Fortvivlelsen er hans Bekymring. Man antager jo, at i enhver Stat den Ringe ordentligviis er fritagen fra at bære de Byrder, som de mere Begunstigede maae bære. Men den fortvivlede Ringe, Hedningen, han vil, selv om han er det, ikke være fritagen, han bærer den tungeste af alle Byrder. Man taler om, at Kongen bærer Kronens Vægt, den Høitstaaende Vægten af Styrelsens Ansvar, Den, hvem Mange ere betroede, Forsorgens Vægt; men saa er dog Kongen ogsaa Kongen, den Høitstaaende den Høitstaaende, den Betroede den Betroede. Men den fortvivlede Ringe, Hedningen, han slæber sig ihjel paa Vægten af hvad han ikke er, han, ja det er Afsindighed, han forløfter sig paa hvad han ikke bærer. Om det er saa, at det er Kongen, der, som Grunden, bærer alle de Andre, eller det er alle de Andre, der bære Kongen som den Øverste, ville vi her ikke undersøge: men den fortvivlede Ringe, Hedningen, han bærer alle de Andre. Denne uhyre Tyngde »alle de Andre«, den tynger paa ham, og med Fortvivlelsens fordoblede Vægt; thi det er ikke gjennem Forestillingen om, at han er Noget, at den tynger paa ham, nei den tynger paa ham gjennem den Forestilling, at han er Intet. Sandeligen, saa umenneskeligt har aldrig nogen Stat eller noget Samfund handlet mod noget Menneske, saa den vilde, at En paa Vilkaaret af at være Intet skulde bære Byrden af Alle. Saa umenneskeligt handler kun den fortvivlede Ringe, Hedningen, mod sig selv. Dybere og dybere synker han i fortvivlet Bekymring, men han finder intet Fodfæste for at bære – han er jo Intet, hvad han til sin egen Qval bliver sig bevidst gjennem Forestillingen om hvad de Andre ere. Latterligere og latterligere, o nei, ynkeligere og ynkeligere, eller rettere ugudeligere og ugudeligere, mere og mere Umenneske bliver han i sin daarlige Stræben for dog at blive til Noget, Noget, om det var nok saa Lidt, men Noget, som det efter hans Begreb var værd at være.

Saaledes synker den fortvivlede Ringe, Hedningen, under Sammenligningens uhyre Vægt, som han selv paalægger sig. Dette, at være et ringe Menneske, hvad der for den ringe Christen hører ham med til det at være Christen som den neppe hørlige, svage Beaandelse foran Bogstavet hører med til Bogstavet, hvilket egentligen høres (og saaledes taler den ringe Christen om sin jordiske Ringhed, han taler blot om den, idet han udtaler, at han er Christen) – dette er for Hedningen hans Bekymring Dag og Nat, med dette beskjeftiger al hans Digten og Tragten sig. Uden Evighedens Udsigt, aldrig styrket af Himlens Haab, aldrig sig selv, forladt af Gud lever han fortvivlet, som var han til Straf dømt at leve disse 70 Aar, martret ved Tanken om at være Intet, martret ved det Frugtesløse i hans Stræben efter at blive til Noget. For ham har Fuglen intet Trøsteligt, Himlen ingen Trøst – og det forstaaer sig, Jordlivet har da for ham heller ingen Trøst. Om ham kan man ikke sige, at han trælbunden bliver ved Jorden, overtalt af Jordlivets Fortryllelse, der bragte ham til at glemme Himmelen. Nei, det er jo snarere som gjorde Timeligheden Alt for at støde ham fra sig, ved at gjøre ham til Intet. Og dog vil han tilhøre Timeligheden paa det usleste Vilkaar; han vil ikke slippe den, han klamrer sig fast ved det Intet at være, fastere og fastere, da han forgjeves søger verdsligt at blive Noget, han klamrer sig med Fortvivlelse fastere og fastere om det – han indtil Fortvivlelse ikke vil være. Og saaledes lever han, ikke ved Jorden, men som var han nedstødt i Underverdenen. See, hiin Konge, hvem Guderne straffede, han leed den rædsomme Straf, at hver Gang han hungrede, da viste de lystelige Frugter sig, men naar han greb efter dem, da forsvandt de. Den fortvivlede Ringe, Hedningen, lider end qvalfuldere i Selvmodsigelsen. Thi medens han, martret ved Intet at være, forgjeves søger at blive Noget, saa er han egentligen ikke blot Noget men Meget; det er ikke Frugterne, der unddrage sig ham, det er ham selv, der unddrager sig endog for at være hvad han er. Thi er han ikke Menneske – og kan han ikke blive Christen!

Lad os saa til Slutning tænke paa Fuglen; den er med i Evangeliet, og skal være med i Talen. Den ringe Fugl er uden Ringhedens Bekymring; den ringe Christen er i Ringhed uden Ringhedens Bekymring, og saa – høit ophøiet over al jordisk Høihed; den ringe Hedning er i Bekymringen langt under sig selv, var han end den Ringeste. Fuglen seer ikke nøie paa hvad den er; den ringe Christen seer nøie paa hvad han er som Christen; den ringe Hedning stirrer indtil Fortvivlelse paa, at han er ringe. »Hvad.. ringe..«, siger Fuglen, »lad os aldrig tænke paa Sligt, det flyver man fra.« »Hvad.. ringe«, siger den Christne, »jeg er Christen!« »Ak, ringe,« siger Hedningen. »Jeg er hvad jeg er,« siger Fuglen; »hvad jeg skal vorde, er endnu ikke aabenbart,« siger den ringe Christen; »jeg er Intet, og jeg vorder ikke til Noget,« siger den ringe Hedning. »Jeg er til,« siger Fuglen; »i Døden begynder Livet,« siger den ringe Christen; »jeg er Intet, og jeg bliver i Døden til Intet,« siger den ringe Hedning. I Sammenligning med den ringe Christen er Fuglen et Barn; i Sammenligning med den ringe Hedning er den et lykkeligt Barn. Som den frie Fugl, naar den svinger sig høiest i Glæde over at være til, saaledes hæver den ringe Christen sig endnu høiere; som den fangne Fugl, naar den modløs og angst spræller sig ihjel i Nettet, saaledes afsjeler den ringe Hedning sig selv, endnu ynkeligere, i Intethedens Fangenskab. Der er efter Christendommens Lære kun een Ophævelse: den at være Christen; alt Andet er Ringe: Ringhed og Høihed. Der er, naar man er Ringe, kun een Vei til Høihed: den at blive Christen. Den Vei kjender Fuglen ikke, den bliver hvad den er. Men saa er der ogsaa en anden Vei, den kjender Fuglen heller ikke, ad den gaaer Hedningen. Fuglens Vei i at være er gaadefuld, den fandtes aldrig; den Christnes Vei er funden af Ham, som er Veien, og den er salig at finde; Hedningens Vei ender i Mørket, Ingen fandt tilbage ad den. Fuglen slipper hiin Omvei og lykkeligt al Fare forbi; den ringe Christen gaaer ikke ad Omveien og frelses salig til Herligheden; den ringe Hedning vælger Omveien og »gaaer saa sin egen Vei« til Fortabelsen.

\sectionheading{Det Glædelige i: at Modgang er Medgang  ·  SV¹ X, 152–160}

Det Glædelige i: at Modgang er Medgang.

Modgang er Medgang. Men, hører jeg Nogen sige, dette er jo kun en Spøg, let nok at forstaae; thi naar man blot betragter Alt omvendt, saa er det i sin Rigtighed: ligefrem forstaaet er Modgang Modgang, omvendt er Modgang Medgang. Sligt er kun Spøg, ligesom det at gjætte Gaader, eller som naar en Tusindkonstner siger »Intet er lettere at gjøre end dette, blot man er vant til at gaae paa Hovedet istedetfor paa Benene.« Nu ja, men er det derfor ogsaa saa let at gjøre det? Og fordi det falder den i Livets Virkelighed uforsøgte Tanke, som ikke kjender til noget Tryk, saa let at svinge sig op og ned og ned og op, at gjøre omkring baade til Høire og Venstre: er det derfor ogsaa saa let, naar Modgang tynger paa den Tanke, som skulde gjøre Svinget, er det saa let, naar Tanken skal mægte at dreie Den om, der i Lidelse og Modgang bestandigt vil indtage den modsatte Stilling? Det vil sige, for Tanken, den ørkesløse og herreløse Tanke, Tanken saadan i Almindelighed, Tanken, der intetsteds har hjemme og Ingens Tanke er, Tanken, der fægter i Luften med ubenævnte Benævnelser og intetbestemmende Bestemmelser »her – der«, »Høire – Venstre«, »ligefrem – omvendt«: for Tanken som Løsgjenger er det let nok at gjøre Konststykket. Men naar det er den benævnte Tanke, naar det er min Tanke, eller naar det er Din Tanke, og – naar Du er en Lidende, saa det altsaa skal blive Alvor med at Tanken, som let nok kan vende sig, faaer denne Magt over Dig til at vende Dig om trods alt det Meget, der mangfoldigt holder forhindrende igjen: er det saa saa let?

Og fordi det er en Spøg med at kunne gaae paa Hovedet istedetfor paa Benene, er det derfor ogsaa en Spøg, at betragte Alt omvendt? O, langtfra, eller rettere, o, lige omvendt, det er just Alvoren, Evighedens Alvor. Hvad der er Spøg, en intetsigende Spøg, saa længe det skal være Tanken saadan i Almindelighed – naar det bliver Alvor derved, at det er Din Tanke, der skal vende Dig om: det er saa just Evighedens Alvor. Evigheden, der dog vel er Alvorens Ophav og Borgen, siger »dette er Opgaven, thi det er just min, Evighedens, Betragtning af Livet at see Alt omvendt. Du skal vænne Dig til at betragte Alt omvendt, og Du Lidende, vil Du for Alvor trøstes, trøstes saa endog Glæden seirer, da maa Du lade mig, Evigheden, hjælpe Dig – men saa maa Du ogsaa betragte Alt omvendt.« Dette er Evighedens Alvor, dette er Evighedens Trøst for den Lidende, den Lov, som Evigheden foreskriver, den Betingelse, som Evigheden gjør, til hvilken alle Forjættelser ere knyttede. Thi Evigheden kjender kun een Fremgangsmaade: at betragte Alt omvendt. Saa lad os da betragte Forholdet omvendt, og derved finde

det Glædelige i: at Modgang er Medgang.

Men lad os gaae saaledes frem, at vi først stræbe at faae den Lidende rigtigt stillet, at han maatte faae Synet for Omvendtheden, at han maatte ville give sig hen i denne Betragtning og give den Magt over sig: saa følger det Glædelige vel af sig selv.

Hvad er Medgang? Medgang er hvad der er mig behjælpeligt i at naae mit Maal, hvad der fører mig til Maalet; og Modgang er, hvad der vil forhindre mig i at naae mit Maal.

Men hvilket er nu Maalet? Vi have gjort den ene Tanke antagelsesviis fast ved at bestemme hvad Medgang og Modgang er; men idet vi skulle til at bestemme den anden Tanke (Maalets), viser det sig jo let, at dersom Maalet er det Forskjellige, er det Modsatte, saa maa ogsaa Medgang og Modgang forandres i Forhold hertil.

Vi staae ved Begyndelsen. Men i en anden Forstand staae vi ikke ved Begyndelsen. Talen henvender sig til en Lidende. Men en Lidende skal ikke først nu til at begynde sit Liv, tværtimod, han er midt i det, ak, ikke blot midt i Livet, men midt i Livets Lidelse. Naar saa er, maa han jo sagtens vide, hvad Modgang er, han den noksom Forsøgte. Maaskee. Men vi bleve jo enige om, at det, hvorvidt han veed hvad Modgang er, afhænger af, om han veed, hvilket Maalet er. Kun Den, der har den sande Forestilling om, hvilket det Maal er, som er sat Mennesket, kun han veed ogsaa, hvad der er Modgang og hvad der er Medgang. Den, der har den usande Forestilling om Maalet, han har ogsaa en usand Forestilling om Medgang og Modgang, han kalder Det Medgang det som fører ham til – det usande Maal, og altsaa forhindrer ham i at naae Maalet (det sande Maal). Men det, som forhindrer En i at naae Maalet, det er jo Modgang.

Der kan nu af Mennesker stræbes efter meget Forskjelligt, men væsentligen er der kun tvende Maal: eet Maal, som et Menneske ønsker, begjerer at naae, og et andet, som han skal naae. Det ene Maal er Timelighedens, det andet Evighedens; de ere hinanden modsatte, men altsaa maa ogsaa Medgang og Modgang i Forhold hertil blive det Omvendte. Hvis denne Tale henvendte sig til en Yngling, skulde den stræbe at gjøre ham dette om de tvende Maal ret indlysende, at han maatte begynde sit Liv med at vælge det rette Maal, begynde med at være rigtigt stillet. Maaskee skulde det alligevel ikke lykkes Talen; thi Ynglingens Sjel vil vel være i en betænkelig Forstaaelse med Timelighedens Maal, og altsaa med den usande Forestilling om Medgang og Modgang. Og nu en Lidende, som altsaa ikke staaer ved Begyndelsen, tværtimod er langt henne i det, han veed, ak, kun altfor godt, hvad Modgang er; men Spørgsmaalet er, som sagt, om han ogsaa ret veed, hvilket Maalet er. Jo heftigere han taler om sin Lidelse, og om hvorledes Alt gaaer ham imod, desto tydeligere bliver det kun, at han har den usande Forestilling om Maalet. Og har han den usande Forestilling om Maalet, kan han jo heller ikke tale sandt om Medgang og Modgang.

Der maa altsaa, hvis han skal hjælpes, fordres af ham, at han endnu engang ret dybt besinder sig paa, hvilket Maal der er sat Mennesket, at han ikke, skuffet af den Indbildning, godt at vide, hvilket Maalet er, farer fort i at klage. Thi vistnok lider Du Modgang, Du kan ikke naae det Maal, som Du saa gjerne ønskede at naae, men hvis nu Maalet er det usande!

Hvad fordres der altsaa? Der fordres af den Lidende, at han standser sin vildfarende Tanke, at han dernæst besinder sig paa, hvilket Maalet er, det er, der fordres, at han vender sig om. Thi i Forhold til Synden fordres der Omvendelse; i Forhold til Evighedens Trøst fordres det Samme, dog i en mildere Skikkelse, at man vender sig om: til Synderen siger Lovens Strenghed forfærdende »omvend Dig«; til den Lidende siger Evigheden mildt deeltagende »o, vend Dig dog blot om.« Altsaa der fordres, at han vender sig om. Allerede her viser Evigheden sig som det Omvendte af Timeligheden. Evigheden forudsætter nemlig, at det naturlige Menneske ingenlunde veed, hvilket Maalet er, at det tværtimod har den usande Forestilling. Timeligheden forudsætter, at ethvert Menneske sagtens veed, hvilket Maalet er, saa Forskjellen mellem Menneske og Menneske blot bliver, om det lykkes dem at naae det eller ikke. Derimod antager Evigheden, at Forskjellen mellem Menneske og Menneske er, at den Ene veed, hvilket Maalet er, og styrer derefter, den Anden veed det ikke, – og styrer der efter, det er, styrer feil. O, Du Lidende, hvo Du end er, det falder Dig vistnok kun altfor let at gjøre Dig forstaaelig for Menneskene i Almindelighed, naar Du klager over Din Lidelse – selv om de ingen Trøst have for Dig, de forstaae Dig dog: men Evigheden vil ikke forstaae Dig saaledes; og dog er det ved den Du skal hjælpes.

Saa vend Dig da om! O, lad Dig dog sige, Herre Gud det er jo saa indlysende, at skal et Menneske komme til Maalet, maa han vide, hvilket Maalet er, og være rigtigt stillet; det er jo saa indlysende, skal et Menneske frydes ved den herlige Udsigt, saa maa han vende sig til den Side, hvor den er at see, ikke til den modsatte Side. Bliv ikke utaalmodig, siig ikke »jeg veed saamænd nok hvad Modgang er«, gjør intet Forsøg paa ved Skildring af Din Lidelse ogsaa at forfærde os, saa vi ogsaa vendte os forkeert og tabte Synet. Thi er Din Lidelse saa forfærdelig, nu hvorfor vil Du saa stirre paa den; og er dette just det Forfærdelige at Du ikke kan lade være at stirre paa den, umuligt er det dog ikke. Siig ikke »naar man lider som jeg lider, saa veed man hvad Modgang er, og kun Den, der lider som jeg lider, veed hvad Modgang er«. Nei, siig ikke dette, o, men hør efter; for ikke at saare Dig, tale vi jo anderledes, vi negte jo ikke, at Du veed, hvad Modgang er, vi tale om, at Du dog nok ikke veed, hvilket Maalet er.

Og naar Du saa har vendt Dig om, og faaet Øie paa Maalet (Evighedens), saa lad Maalet, hvad det er og skal være, blive Dig saa vigtigt, at der intet Spørgsmaal er om, hvorledes Veien er, men ene om at naae Maalet, saa Du faaer Mod til at forstaae, at hvordan saa Veien end er, den sletteste af alle, den piinagtigste af alle – dersom den fører Dig til Maalet, at den saa er Medgang. Ikke sandt, naar der er et Sted, hvor det er Dig saa vigtigt at komme hen, fordi Du saa ubeskriveligt gjerne vil derhen, saa siger Du »enten jeg skal kjøre baglænds eller forlænds, enten jeg skal ride eller gaae eller krybe – det er det Ligegyldige, naar jeg blot kommer derhen.« Det er dette Evigheden først og fremmest vil, den vil gjøre Dig Maalet saa vigtigt, at det faaer ganske Magt over Dig og Du derved Magt over Dig selv til at faae Din Tanke, Dit Sind, Dit Øie bort fra Besværligheden, Vanskeligheden, fra hvorledes Du kommer derhen, fordi det er Dig det ene Vigtige at komme derhen.

Altsaa, af Ærbødighed for Maalet er det nu blevet Dig ligegyldigt, om hvad der ellers kaldes Medgang eller hvad der ellers kaldes Modgang skal føre Dig til Maalet: hvad der fører Dig til Maalet er Medgang. Hvor forandret! Troer Du, at det sandselige Menneske kunde være ligegyldigt herved? Hvad vilde det trøste ham, at Modgang førte ham til »Maalet«, naar han kun er bekymret om det Maal, hvortil Medgang fører!

Dog maaskee Du endnu ikke kan lade være at see Dig om efter den Forskjel: hvad man ellers kalder Modgang og Medgang. Du har vundet den rette Stilling, men endnu ikke Ro i den. Vel, Evigheden vil hjælpe Dig videre. Thi naar nu just hvad man ellers kalder Modgang ene eller dog fortrinsviis fører til »Maalet«, er der saa nogen Grund til at see Dig om? Naar det var saa, lad os antage det, at Du ene eller dog bedst ved at kjøre baglænds kunde komme til det Sted, hvor Du saa gjerne vilde hen: var det saa ogsaa rigtigt at sige »enten jeg skal kjøre forlænds eller jeg skal kjøre baglænds, det er det Ligegyldige?« Det var dog vel rigtigere sagt »saa er det jo godt, at jeg er kommen til at kjøre baglænds.« Og saaledes ogsaa, dersom det var muligt, at hvad man ellers kalder Medgang lettere kunde føre Dig til »Maalet«, saa blev der jo Plads for et Ønske. Men nu skal Intet friste Dig – thi Modgang fører just til »Maalet«. Og ikke sandt, Du vil jo staae ved Dit Ord, at hvad der fører til Maalet, det er Medgang. Altsaa er Modgang Medgang.

Lad os nu gjøre os dette ret tydeligt, at ikke baade hvad man kalder Medgang og Modgang lige godt føre til »Maalet«, men at ene eller dog fortrinsviis hvad man kalder Modgang fører til »Maalet«. Hvad kan forhindre et Menneske i at naae »Maalet«? Dog vel netop det Timelige, og hvorledes meest? Naar hvad man ellers kalder Medgang fører ham til at naae Timelighedens Maal. Idet han nemlig ved Hjælp af Medgang naaer Timelighedens Maal, er han længst fra at naae »Maalet«. Mennesket skulde stræbe efter Evighedens Maal, men det Timelige har ved Hjælp af Medgang sinket ham. Det, at Timeligheden gaaer ham med, fører jo ikke til det Evige, altsaa ikke til »Maalet«. Dersom Noget gjør det, saa maa det være lige det Omvendte, at Timeligheden gaaer ham imod. Men det, at Timeligheden gaaer ham imod, er jo, hvad man saadan kalder Modgang. – Naar det hedder »tragter først efter Guds Rige«, saa er hermed Evighedens Maal sat Mennesket som det, hvorefter han skal tragte. Skal nu dette gjøres, og nøie efter Ordene (o, og Evigheden lader sig ikke spotte, ei heller bedrage!), saa gjælder det jo for Alt om, at Mennesket ikke kommer til først at tragte efter noget Andet. Hvad er nu det Andet, han kan tragte efter? Det er det Timelige. Hvis han da skal tragte først efter Guds Rige, maa han altsaa frit forsage alle Timelighedens Maal. O, vanskelige Opgave; naar Beleiligheden tilbydes maaskee i Overflod, naar Alt er vinkende, naar hvad man kalder Medgang er flux tilrede for, blot han ønsker det, at føre ham til Besiddelsen af alle Timelighedens lystelige Goder: da at forsage alt dette! Men den Lidende har jo Modgang, derfor kaldes han en Lidende. Og hvad man kalder Modgang, det forhindrer jo den Lidende i at naae disse Timelighedens Maal. Modgang gjør ham det vanskeligt, maaskee umuligt. O, hvor tungt saaledes at see Vanskeligheder optaarne sig for Ønsket, hvor tungt, at Ønskets Opfyldelse blev umulig! Ikke sandt? Ja, det behøver jeg vel ikke at spørge Dig om, men ikke sandt (og Gud give det var saa), det er nok snarere Dig, der nu vil spørge mig, om jeg nu selv har glemt, hvorom Talen er. Saa siig Du det da, det var blot det jeg ønskede, siig Du os blot, hvorom Talen er, medens jeg med Glæde hører efter og hører Dig sige: naar hvad man kalder Medgang er det Sinkende, som forhindrer fra at naae »Maalet«, saa er det jo godt, at hvad man kalder Modgang gjør En det vanskeligt eller umuligt – at sinkes, det er, saa fører Modgang En jo netop til »Maalet.«

O, Du Lidende, hvo Du end er, løsriv Dig dog blot et Øieblik fra Din Lidelse, og fra de Tanker, den vil paanøde Dig, forsøg at tænke ganske uhildet over Livet. Tænk Dig da et Menneske i Besiddelse af alle Lykkens Goder, begunstiget paa enhver Maade – men tænk Dig, at dette Menneske dog derhos er alvorlig nok til at have rettet sit Sind paa Evighedens Maal. Han forstaaer altsaa, at alt Dette, som er ham givet, skal han forsage. Han er ogsaa dertil villig, men see, da vaagner der en Mismodets Bekymring i hans Sjel, en ængstelig Selvbekymring, om han dog maaskee ikke bedrager sig selv, og det kun er en Indbildning med hans Forsagelse, da han jo bliver i Besiddelsen af alle Goderne. Han tør ikke kaste fra sig hvad der er givet ham; thi han forstaaer, at dette kunde være en formastelig Overdrivelse, der let kunde blive ham til Fordærvelse, istedetfor til Gavn: og sørgmodigt har han fattet en bekymret Mistillid til sig selv, om han ikke muligen bedrager Gud, og al hans Forsagelse er et Spilfægteri. Da kunde han vel ønske, at det maatte fratages ham Altsammen, at det maatte blive Alvor for ham med at opgive det Timelige for at gribe det Evige. Og skeer det ikke, da udvikles der maaskee i hans Inderste en Sindets Sygdom, et uheldbredeligt Tungsind, som har sin Grund i, at han i dybere Forstand blev raadvild over sig selv.

Har Du aldrig tænkt Dig dette? Det vilde dog vistnok just for Dig være et rigtigt Synspunkt, da det fjerner Dig det længst mulige fra dit Eget. Thi betragt fra dette Synspunkt Dit Vilkaar. Du har jo havt og har Modgang nok: Du har altsaa kun den Opgave at forsage hvad der er Dig negtet; istedetfor, at han har den Opgave at forsage hvad ham er givet. Dernæst er Du fritagen for den Bekymring, om Du virkelig, det er her i udvortes Forstand, har opgivet det; thi da Du ikke besidder det, saa er i denne Henseende Sagen let nok. Hvor langt mere er Du altsaa ikke hjulpen? Det er Dig negtet, det som vil forhindre Dig fra at naae »Maalet«, Du har ikke selv kastet det bort, og derved paadraget Dig et Ansvar, som maaskee i et afgjørende Øieblik skulde gjøre Dig Dit Liv saa saare vanskeligt, fordi Du befandt Dig kraftesløs ligeoverfor den Opgave, Du vilkaarligen havde sat Dig. Nei, i Forhold til Dig har Styrelsen paataget sig alt Ansvaret, det er Styrelsen, der har negtet Dig det. Du har da ene at komme Styrelsen til Hjælp, Styrelsen, som har hjulpet Dig. Thi Modgang er Medgang, og Du har jo Modgang.

Altsaa Modgang er Medgang. Det er evig vist, ikke al Satans Kløgt skal formaae at gjøre det tvivlsomt. Du kan ogsaa godt forstaae det. Derimod troer Du maaskee ikke ret paa, at det er saa. Men troer Du da (for at byde Dig lidt svagere Føde, dersom Skriftens Tale om først at tragte efter Guds Rige skulde være Dig for stærk), troer Du da, at den Digter, hvis Sange fryde Menneskeheden, troer Du, at han kunde have digtet disse Sange, hvis ikke Modgang og tunge Lidelser havde været med til at stemme Sjelen! Thi det er just i Modgang, »naar Hjertet sidder meest beklemt, da bliver Frydens Harpe stemt.« Eller troer Du, at Den, der i Sandhed vidste at trøste Andre, troer Du, han havde formaaet dette, hvis ikke Modgang for ham havde været den fornødne Medgang, der havde hjulpet ham til Færdighed i denne skjønne Konst! Maaskee fandt han det ogsaa selv i Begyndelsen tungt nok, næsten grusomt, at hans Sjel skulde martres, for da at blive opfindsom i at hitte paa Trøst for Andre. Men tilsidst lærte han vel at forstaae, at uden Modgang kunde han ikke være bleven og ikke være Den han var; han lærte at troe paa, at Modgang er Medgang.

Saa tro Du da ogsaa paa, at Modgang er Medgang. At forstaae det er let nok – men at troe paa det er vanskeligt; lad Dig ikke bedrage af den forfængelige Viisdom, der vil indbilde Dig, at det er let at troe, vanskeligt at forstaae. Men tro paa det. Saa længe Du ikke troer derpaa, saa længe er og bliver Modgang Modgang. Det hjælper Dig ikke, at det er evig vist, at Modgang er Medgang, saa længe Du ikke troer derpaa, er det ikke saa for Dig. See, den Ældre veed i Forhold til Barnet, det Raad mod Nelder: tag blot rask til, saa brænder Du Dig ikke. Det maa jo synes Barnet det urimeligste af Alt; thi, tænker Barnet, brænder man sig blot man rører ved den, hvor meget mere, naar man tager haardt. Saa siger man det da til Barnet. Men idet Barnet skal gribe til, saa har det dog ikke ret Mod – det tager alligevel ikke rask nok til, og det brænder sig. Saaledes med dette, at Modgang er Medgang: er Du ikke i Troen besluttet, saa faaer Du dog kun Modgang ud af det.

Tro altsaa, at Modgang er Medgang. Det er vist, der ventes blot paa, at Du skal troe det. Lad Dig ikke i din Tro forstyrre af Andre, »hav Troen hos Dig selv for Gud« (Rom. 14, 22). Naar den Søfarende er forvisset om, at den Vind, som nu blæser, fører ham til Maalet – om saa end alle Andre vilde kalde den Modvind, hvad bryder han sig derom, han kalder den Medbør. Thi Medbør er den Vind, som fører En til Maalet; og Medgang er Alt, hvad der fører En til Maalet; og Modgang fører En til »Maalet«: altsaa er Modgang Medgang.

At det er glædeligt behøver ikke at udvikles. Den, som troer paa, at Modgang er Medgang, ham behøver Talen sandeligen ikke at forklare, at det er glædeligt. Og Den, som ikke ret troer derpaa, ham er det vigtigere intet Øieblik at spilde, men at gribe Troen. Altsaa derom behøve vi ikke at tale, eller dog kun med et Ord. Tænk Dig da, som ved en Jagt over hele Verden, alt hvad der ellers kaldes Trøstegrund opdrevet og samlet, alle disse Trøstegrunde, som, ja jeg troer det er saa, de Lykkelige have opfundet for at blive de Ulykkelige qvit; og tænk Dig saa til Sammenligning Evighedens Trøst, denne korte Trøst, som Bekymringen har opfundet, ligesom den jo ogsaa har opfundet, at det er en Bekymret, en Lidende, der skal trøste Andre, ikke en Lykkelig – denne korte Trøst: Modgang er Medgang. Ikke sandt, Du vil finde det ganske i sin Orden, og i en vis Forstand vel betænkt, at de menneskelige Trøstegrunde dog ikke give sig Skin af at kunne gjøre den Sørgende glad, men blot paatage sig saadan at trøste ham, hvad de da gjøre slet nok. Evigheden derimod, naar den trøster, saa gjør den glad, dens Trøst er egentlig Glæden, er den egentlige Glæde. Det er med de menneskelige Trøstegrunde, som naar den Syge, der allerede har havt mange Læger, faaer en ny, der maaskee finder paa noget Nyt, som et Øieblik frembringer en Smule Forandring, medens saa dog strax Alt igjen er det Gamle. Nei, naar Evigheden bliver hentet til den Syge, da helbreder den ham ikke blot ganske, men gjør ham sundere end alle de Sunde. Det er med de menneskelige Trøstegrunde, som naar Lægen hitter paa en ny, maaskee lidt beqvemmere Art Krykke for Den, der gaaer paa Krykker: give ham sunde Fødder at gaae paa, og Kraft i Knæerne, det kan Lægen ikke. Men naar Evigheden bliver hentet, saa kastes Krykkerne bort, han kan saa ikke blot gaae, o, nei, i en anden Forstand maa man sige, han gaaer ikke mere – saa let gaaer han. Thi Evigheden, den giver Fødder at gaae paa. Naar det i Modgang synes umuligt at komme af Stedet, naar det i Lidelsens Afmagt er som kunde man ikke røre en Fod: saa gjør Evigheden Modgang til Medgang.

Der er i al Modgang kun een Fare, naar den Lidende ikke vil troe, at Modgang er Medgang. Dette er Fortabelse; kun Synden er Menneskets Fordærvelse.

\sectionheading{1 Joh. 3,20 (Tale ved Altergangen om Fredagen)  ·  SV¹ X, 301–309}

1 Joh. III, 20.

Bøn.

Stor er Du, o Gud; skjøndt vi kun kjende Dig som i en mørk Tale og som i et Speil, vi tilbede dog undrende Din Storhed – hvor meget mere skulle vi engang prise den, idet vi lære den fuldeligere at kjende! Naar jeg under Himmelens Hvælving staaer omgiven af Skabningens Under, da priser jeg rørt og tilbedende Din Storhed, Du, som let bærer Stjernerne i det Uendelige, og faderligt bekymrer Dig om Spurven. Men naar vi ere forsamlede her i Dit hellige Huus, da ere vi jo ogsaa overalt omgivne af hvad der i end dybere Forstand minder om Din Storhed. Thi stor er Du, Verdens Skaber og Opholder; men da Du, o Gud, tilgav Verdens Synd og forligte Dig med den faldne Slægt, ak, da var Du jo dog end større i Din ubegribelige Forbarmelse! Hvorledes skulle vi da ikke troende prise og takke og tilbede Dig her i Dit hellige Huus, hvor Alt minder os derom, særligen dem, som idag ere forsamlede for at annamme Syndernes Forladelse, og for paa ny at tilegne sig Forligelsen med Dig i Christo!

1 Joh. III, 20... om end vort Hjerte fordømmer os, da er Gud større end vort Hjerte.

Om end vort Hjerte fordømmer os. Da Pharisæerne og de Skriftkloge havde bragt en Qvinde, som var greben i aabenbar Synd, til Christum i Templet for at anklage hende, og da de derpaa, beskæmmede ved Hans Svar, vare alle gangne bort, sagde Christus til hende »fordømte Ingen Dig,« men hun sagde »Herre, Ingen«. Der var altsaa Ingen, der fordømte hende. Saaledes ogsaa her i denne Helligdom, der er Ingen, der fordømmer Dig; om Dit Hjerte fordømmer Dig, maa Du ene selv vide. Ingen Anden tør vide det; thi ogsaa denne Anden er jo idag beskjeftiget med sit eget Hjerte, om det fordømmer ham. Hvorvidt Dit Hjerte fordømmer Dig, det vedkommer ingen Anden; thi ogsaa denne Anden har kun med sit eget Hjerte at skaffe, dets anklagende eller dets frikjendende Tanker. Hvorledes Du blev tilmode, idet disse Ord forelæstes »om end vort Hjerte fordømmer os«, vedkommer ingen Anden; thi ogsaa denne Anden henfører jo andægtigt Alt til sig selv, tænker kun paa, hvorledes han blev tilmode, om Ordet overraskede ham som en pludselig Tanke, eller han hørte, ak hvad han havde sagt sig selv, eller han hørte hvad han dog meente ikke at passe paa ham. Vel kan nemlig et Hjerte anklage sig selv, men deraf følger endnu ikke, at det maa fordømme sig selv; og vi lære jo ikke tungsindig Overdrivelse, saa lidet som vi lære letsindigt Aflad. Men naar der skulde tales over de forelæste Ord, hvor skulde man finde bedre Tilhørere end paa en saadan Dag som denne, og end saadanne, som de, der her idag ere komne hid, ikke fra Verdens Adspredelser, men fra Skriftestolens Samling, hvor de jo hver især have gjort Gud Regnskab, hver især have ladet deres Hjerte være Anklageren, hvad det ogsaa bedst kan være, da det er Medvideren, og hvad det jo ogsaa helst maa blive itide, at det ikke engang forfærdeligt skal blive det mod Menneskets egen Villie. Dog er der jo Forskjel paa Skyld og Skyld; der er Forskjel mellem at være fem hundrede Penninge skyldig og kun halvtredssindstive; det ene Menneske kan have langt, langt mere at bebreide sig end det andet; der kan ogsaa være Den, som maa sige sig selv, at hans Hjerte fordømmer ham; maaskee kan der ogsaa være en Saadan her tilstede, eller maaskee er der ingen Saadan her tilstede: men Trøst trænge vi jo dog alle til. Og dette kan dog vel ikke være mistrøstende for Nogen, at Trøstens Ord er saa riigt paa Forbarmelse, at det tager Enhver med, dette kan jo ikke være mistrøstende for Nogen, selv om hans Hjerte ikke fordømmer ham. Det er dog væsentligen den samme Trøst vi alle behøve, vi, hvis Hjerte ikke frifinder os: Guds Storhed, at han er større end vort Hjerte.

Guds Storhed er i at tilgive, i at forbarme sig, og i denne sin Storhed er han større end det Hjerte, som fordømmer sig selv. See, det er denne Guds Storhed, om hvilken der især skal tales paa de hellige Steder; thi herinde kjende vi jo Gud anderledes, nærmere, om man saa tør sige fra en anden Side end derudenfor, hvor han vistnok er aabenbar, at kjende i sine Gjerninger, medens han her er at kjende som han har aabenbaret sig selv, som han vil kjendes af den Christne. Tegnene, paa hvilke Guds Storhed i Naturen kjendes, kan Enhver undrende see, eller rettere, der er egentligen intet Tegn, thi Gjerningerne selv ere Tegnene; saaledes kan jo Enhver see Regnbuen, og maa undres, naar han seer den. Men Tegnet paa Guds Storhed i at forbarme sig er kun for Troen; dette Tegn er jo Sacramentet. Guds Storhed i Naturen er aabenbar; men Guds Storhed i at forbarme sig er en Hemmelighed, der maa troes. Netop fordi den ikke ligefrem er aabenbar for Enhver, netop derfor kaldes den og er den aabenbarede. Guds Storhed i Naturen vækker strax Studsen og saa Tilbedelse; Guds Storhed i at forbarme sig er først til Forargelse og saa for Troen. Da Gud havde skabt Alt, saae han og see »det var Alt saare godt«; og ved hvert hans Værk staaer ligesom tilføiet: priis, lov, tilbed Skaberen. Men ved hans Storhed i at forbarme sig staaer tilføiet: salig Den, som ikke forarges.

Al vor Tale om Gud er, som naturligt, menneskelig Tale. Hvor meget vi end stræbe at forebygge Misforstaaelsen ved atter at tilbagekalde hvad vi udsige – hvis vi ikke ganske ville tie, vi maae dog bruge Menneskets Maalestok, idet vi, Mennesker, tale om Gud. Hvilken er nu den sande menneskelige Storhed? Dog vel Hjertets Storhed. Vi sige egentligen ikke, at Den er stor, der har megen Magt og Herredom, ja om der levede eller havde levet en Konge, hvis Herskermagt var over den hele Jord – hvor hurtig end Forbauselsen er til strax at kalde ham stor: det dybere Menneske lader sig ikke forstyrre af Udvortesheden. Og derimod, om det var det ringeste Menneske, der nogensinde har levet, – naar Du bliver Vidne til hans Handling i Afgjørelsens Øieblik, at han i Sandhed handler ædelt, høimodigt af ganske Hjerte tilgiver sin Fjende, i Selvfornegtelse bringer det yderste Offer, eller naar Du bliver Vidne til den inderlige Langmodighed, med hvilken han fra Aar til Aar kjerligt fordrager Ondt: da siger Du »han er dog stor, i Sandhed han er stor.« Altsaa Hjertets Storhed er den sande menneskelige Storhed; men Hjertets Storhed er netop at overvinde sig selv i Kjerlighed.

Naar vi da nu, Mennesker som vi ere, ville gjøre os en Forestilling om Guds Storhed, saa maae vi tænke paa den sande menneskelige Storhed, altsaa paa Kjerligheden, og paa Kjerligheden som tilgiver og forbarmer sig. Men hvad vil nu dette sige, skulde Meningen være, at vi dog ville sammenligne Gud med et Menneske, var end dette Menneske det ædleste, det reneste, det forsonligste, det kjerligste, der nogensinde havde levet? Langtfra. Saaledes taler Apostelen heller ikke. Han siger ikke, at Gud er større end det kjerligste Menneske, men at han er større end det Hjerte, som fordømmer sig selv. Saaledes ligner altsaa Gud og Mennesket hinanden kun omvendt. Ikke ad den ligefremme Ligheds Trin stor, større, størst, naaer Du Sammenligningens Mulighed, den er kun mulig omvendt; det er jo heller ikke ved mere og mere at opløfte sit Hoved, at et Menneske nærmer sig mere og mere Gud, men omvendt ved dybere og dybere at kaste sig ned i Tilbedelse. Det sønderknuste Hjerte, som fordømmer sig selv, kan ikke faae, søger forgjeves at finde et Udtryk, som er det stærkt nok for dermed at betegne dets Skyld, dets Elendighed, dets Besmittethed: endnu større er Gud i at forbarme sig! Forunderlige Sammenligning! Al menneskelig Reenhed, al menneskelig Forbarmelse duer ikke til Sammenligning; men et angrende Hjerte, der fordømmer sig selv, med dette sammenlignes Guds Storhed i at forbarme sig, kun at den er end større: saa dybt som dette Hjerte kan nedsætte sig selv og dog aldrig sig selv dybt nok, saa uendelig ophøiet, eller uendeligt mere ophøiet er Guds Storhed i at forbarme sig! See, Sproget lige som sprænges og brister for at betegne Guds forbarmende Storhed; Tanken søgte forgjeves en Sammenligning, da fandt den endeligen, hvad der dog menneskelig talt ingen Sammenligning er, et angerfuldt Hjertes Sønderknuselse – endnu større er Guds Forbarmelse. Et angergivent Hjerte, naar det i Sønderknuselse fordømmer sig selv; ja, dette Hjerte undte sig ikke Ro, ikke i noget eneste Øieblik, det fandt intet Skjulested, hvor det kunde flygte hen fra sig selv, det fandt ingen Undskyldning mulig, fandt det en ny, den forfærdeligste Skyld, at søge en Undskyldning, det fandt ingen, ingen Lindring, selv det meest forbarmende Ord, hvilket al Barmhjertighedens Inderlighed formaaede at udtænke, det lød for dette Hjerte, der ikke turde og ikke vilde lade sig trøste, som en ny Fordømmelse over det: saa uendelig er Guds Storhed i at forbarme sig, eller den er endnu større. Den halter, denne Sammenligning – det gjør Mennesket altid efter at have stridt med Gud. Den er søgt, denne Sammenligning, vist nok, thi den blev funden ved gudfrygtigt at vrage al menneskelig Lighed; tør Mennesket ikke gjøre sig noget Billede af Gud, da visseligen heller ikke indbilde sig, at det Menneskelige ligefrem kunde være Sammenligningen. Ingen forhaste sig med at søge, Ingen forraske sig med at ville have fundet en Sammenligning for Guds Storhed i at forbarme sig; enhver Mund skal tilstoppes, Enhver slaae sig for sit Bryst – thi der er kun een Sammenligning, som dog er nogenlunde, et ængstet Hjerte, som fordømmer sig selv.

Men Gud er større end dette Hjerte! Saa lad Dig dog trøste. Maaskee lærte Du tidligere af Erfaring, hvor tungt det er med et saadant Hjerte, at stedes for Pharisæers og Skriftkloges Dom, eller at komme til den Misforstaaelse, der kun veed at sønderrive det end mere, eller til den Smaalighed, der ængster Hjertet end mere sammen, Du, der i saa høi Grad trængte til En, der var stor. Gud i Himlene han er større; han er ikke større end Pharisæere og Skriftkloge, ei heller større end Misforstaaelse og Smaalighed, ei heller større end det Menneske, der dog vidste at tale et lindrende Ord til Dig, hos hvem Du fandt nogen Lise, fordi han ikke var smaalig, ikke vilde nedtrykke Dig end mere, men reise Dig op – Gud er ikke større end han (trøstesløse Sammenligning!), nei Gud er større end Dit eget Hjerte! O, hvad enten det var en Sjelens Sygdom, der saaledes natligt formørkede Dit Sind, saa Du tilsidst i Dødsens Angest, bragt næsten til Afsindighed ved Forestillingen om Guds Hellighed, meente at maatte fordømme Dig selv; hvad heller det var det Forfærdelige, der tyngede saaledes paa Din Samvittighed, at Dit Hjerte fordømte Dig selv: Gud er større! Vil Du ikke troe, tør Du ikke troe uden at see Tegn, nu det tilbydes jo. Han, der kom til Verden og døde, Han døde ogsaa for Dig, ogsaa for Dig. Han døde ikke for Menneskene saadan i Almindelighed, o, lige det Modsatte, døde Han særligen for Nogen, da var det jo for den Ene, ikke for de Ni og halvfemsindstive – ak, og Du er jo for Elendig til at være paa Slump med i det runde Tal, paa Dig faldt jo Elendighedens og Skyldens Eftertryk saa forfærdeligt, at Du tælles udenfor. Og Han, der døde for Dig da Du var Ham en Fremmed, skulde Han forlade sine Egne! Elskede Gud Verden saaledes, at han gav sin eenbaarne Søn for at Ingen skulde fortabes, hvorledes skulde han da ikke bevare dem, som bleve Dyrekjøbte! O, martre Dig ikke selv; hvis det er Tungsindets Angester, der hilde Dig, saa veed Gud Alt – og han er stor! Og er det den tunge Skylds Centnervægt, der hviler paa Dig: han, der, hvad ikke opkom i noget Menneskes Hjerte, af sig selv forbarmede sig over Verden, han er stor! Martre Dig ikke selv, husk paa hiin Qvinde, at der Ingen var som fordømte hende, og betænk, at dette Samme ogsaa kan udtrykkes paa en anden Maade: Christus var tilstede. Just fordi Han var tilstede, derfor var der Ingen, der fordømte hende. Han reddede hende ud af Pharisæernes og de Skriftkloges Fordømmelse; de gik beskæmmede bort, thi Christus var tilstede: der var Ingen, der fordømte hende. Saa blev Christus ene tilbage med hende – men der var Ingen, der fordømte hende. Just det, at Han ene blev tilbage med hende, betyder i en langt dybere Forstand, at der Ingen er, der fordømmer hende. Det vilde kun lidet have hjulpet hende, at Pharisæerne og de Skriftkloge gik bort, de kunde jo være komne igjen med Fordømmelsen. Men Frelseren blev tilbage ene med hende: derfor var der Ingen, der fordømte hende. Ak, der er kun een Skyld, som Gud ikke kan tilgive, det er den ikke at ville troe paa hans Storhed!

Thi han er større end det Hjerte, som fordømmer sig selv. Men derimod staaer der Intet om, at han er større end det verdslige, letsindige, daarlige Hjerte, der taabeligt regner paa Guds indbildte Storhed i at tilgive. Nei, Gud er og kan være lige saa nøieregnende, som han er stor og kan være stor i at forbarme sig. Saaledes ener altid Guds Væsen det Modsatte, ligesom i hiint Mirakel med de fem smaa Brød. Folket havde Intet at spise – ved et Mirakel skaffes der Overflod; men see, derpaa byder Christus, at man omhyggeligt opsamler hver en Levning. Hvor guddommeligt! Thi eet Menneske kan ødsle, et andet være sparsommeligt; men dersom der var et Menneske, som ved et Mirakel hvert Øieblik kunde – guddommeligt skaffe Overflod, troer Du ikke, at han – menneskeligt havde ladet ringe om Smulerne, troer Du, at han – guddommeligt havde opsamlet Smulerne! Saaledes ogsaa med Guds Storhed i at forbarme sig, neppe har et Menneske blot Forestilling om, hvor nøieregnende Gud kan være. Lad os ikke bedrage os selv, ikke lyve for os selv, og, hvilket er det Samme, forringe Guds Storhed, ved at ville gjøre os selv bedre end vi ere, mindre skyldige, eller ved at benævne vor Skyld med letsindigere Navne; dermed forringe vi nemlig Guds Storhed, som er i at tilgive. Men lader os heller ikke afsindigt ville synde endnu mere for at gjøre Tilgivelsen end større; thi Gud er lige saa stor i at være nøieregnende.

Og saa lader os da her i Dit hellige Huus prise Din Storhed, o Gud, Du som ubegribeligt forbarmede Dig og forligte Verden med Dig. See, udenfor der forkynde Stjernerne Din Majestæt, og Alts Fuldkommenhed Din Storhed; men herinde er det de Ufuldkomme, er det Syndere, der prise Din endnu større Storhed! – Ihukommelsens Maaltid er atter beredet, saa være Du da forinden erindret og takket for Din Storhed i at forbarme Dig.

\workheading{Lilien paa Marken og Fuglen under Himlen (Tre gudelige Taler)}{1849}

\worknote{Antologiens ''The Lily in the Field and the Bird of the Air''. Hong-udvalget bringer begyndelsen af første tale, ''Taushed, eller lære at tie'' — om at lære tavsheden af lilien og fuglen og at blive taus for Gud (efter Matth. 6,33). Talens fortsættelse og de øvrige to taler (Lydighed, Glæde) er ikke medtaget.}

\sectionheading{Taushed, eller lære at tie (uddrag)  ·  SV¹ XI, 14–20}

Taushed, eller lære at tie.

Thi vistnok er det Talen, der udmærker Mennesket fremfor Dyret, og da, hvis Nogen saa vil, langt fremfor Lilien. Men fordi det er et Fortrin at kunne tale, deraf følger ikke, at det ikke skulde være en Kunst at kunne tie, eller at det skulde være en ringe Kunst; tvertimod, just fordi Mennesket kan tale, just derfor er det en Kunst at kunne tie, og just fordi dette hans Fortrin saa let frister ham, derfor er det en stor Kunst at kunne tie. Men det kan han lære af de tause Læremestere: Lilien og Fuglen.

»Søger først Guds Rige og hans Retfærdighed.«

Men hvad vil dette sige, hvad har jeg at gjøre, eller hvilken Stræben er den, om hvilken kan siges, at den søger, at den tragter efter Guds Rige? Skal jeg see at faae et Embede, svarende til mine Evner og Kræfter, for at virke deri? Nei, Du skal først søge Guds Rige. Skal jeg da give al min Formue til de Fattige? Nei, først skal Du søge Guds Rige. Skal jeg da gaae ud og forkynde denne Lære i Verden? Nei, Du skal først søge Guds Rige. Men saa er det jo i en vis Forstand Intet jeg skal gjøre? Ja, ganske vist, det er i en vis Forstand Intet; Du skal i dybeste Forstand gjøre Dig selv til Intet, blive til Intet for Gud, lære at tie; i denne Taushed er Begyndelsen, som er, først at søge Guds Rige.

Saaledes kommer man, gudeligt, i en vis Forstand baglænds til Begyndelsen. Begyndelsen er ikke det, man begynder med, men det, man kommer til; og man kommer baglænds til den. Begyndelsen er denne Kunst at blive taus; thi at være taus som Naturen er det, er ingen Kunst. Og dette saaledes i dybeste Forstand at blive taus, taus lige overfor Gud, det er Gudsfrygts Begyndelse; thi som Gudsfrygt er Viisdoms Begyndelse, saa er Taushed Gudsfrygts Begyndelse. Og som Gudsfrygt er mere end Viisdoms Begyndelse, er »Viisdom«, saaledes er Taushed mere end Gudsfrygts Begyndelse, er »Gudsfrygt«. I denne Taushed forstummer gudfrygtigt Ønskets og Begjeringens mange Tanker; i denne Taushed forstummer gudfrygtigt Taksigelsens Ordrigdom.

Det er Menneskets Fortrin for Dyret at kunne tale; men i Forhold til Gud kan det let blive til Fordærvelse for Mennesket, som kan tale, at ville tale. Gud er i Himmelen, Mennesket paa Jorden: derfor kunne de ikke godt tale sammen. Gud er Alviisdom, Det Mennesket veed er løs Snak: derfor kunne de ikke godt tale sammen. Gud er Kjerlighed, Mennesket er, som man siger det til et Barn, endog sit eget Vel betræffende en lille Tosse: derfor kunne de ikke godt tale sammen. Kun i megen Frygt og Bæven kan Mennesket tale med Gud; i megen Frygt og Bæven. Men at tale i megen Frygt og Bæven er af en anden Grund vanskeligt; thi som Angest gjør, at Stemmen legemligt svigter, saaledes gjør vel ogsaa megen Frygt og Bæven, at Talen forstummer i Taushed. Det veed den rette Bedende; og Den, der ikke var den rette Bedende, han lærte maaskee just dette i Bønnen. Der var Noget, der laae ham saa meget paa Sinde, en Sag, som var ham saa vigtig, det var ham saa magtpaaliggende ret at gjøre sig forstaaelig for Gud, han var bange for, at han skulde i Bønnen have glemt Noget, ak, og hvis han havde glemt det, saa var han bange for, at Gud ikke af sig selv skulde huske paa det: derfor vilde han samle sit Sind til at bede ret inderligt. Og hvad skeete ham saa, hvis han ellers bad inderligt? Det Forunderlige skeete ham; efterhaanden som han blev inderligere og inderligere i Bønnen, havde han Mindre og Mindre at sige, og tilsidst blev han ganske taus. Han blev taus, ja hvad der om muligt endnu mere end Taushed er modsat det at tale, han blev en Hørende. Han meente at det at bede er at tale; han lærte at det at bede ikke blot er at tie, men er at høre. Og saaledes er det; at bede er ikke at høre sig selv tale, men er at komme til at tie, og at blive ved at tie, at bie, til den Bedende hører Gud.

Derfor er det, at Evangeliets Ord, søger først Guds Rige, opdragende ligesom binder Munden paa Mennesket, ved paa ethvert enkelt Spørgsmaal han gjør, om det er Det han skal gjøre, at svare: nei, Du skal først søge Guds Rige. Og derfor er det man kan omskrive dette Evangeliets Ord saaledes: Du skal begynde med at bede, ikke som om – hvad vi jo have viist – Bønnen altid begyndte med Taushed, men fordi naar Bønnen ret er blevet Bøn, saa er den blevet Taushed. Søger først Guds Rige det er: beder! Dersom Du spørger, ja om Du i Spørgsmaal gjennemgik alt Enkelt, spørgende, er det dette, jeg skal gjøre, og naar jeg gjør det, er det saa at søge Guds Rige: da maa dertil svares: nei, Du skal først søge Guds Rige. Men at bede, det vil sige, retteligen at bede, det er at blive taus, og det er at søge først Guds Rige.

Denne Taushed kan Du lære hos Lilien og Fuglen. Det vil sige, deres Taushed er ingen Kunst, men naar Du bliver taus som Lilien og Fuglen, saa er Du ved den Begyndelse, hvilken er, først at søge Guds Rige.

Hvor høitideligt er der ikke derude under Guds Himmel hos Lilien og Fuglen, og hvorfor? Spørg »Digteren«; han svarer: fordi der er Taushed. Og ud til denne høitidelige Taushed længes han, bort fra Verdsligheden i Menneske-Verdenen, hvor der er saa megen Talen, bort fra hele det verdslige Menneske-Liv, som kun paa en sørgelig Maade beviser, at Mennesket ved Talen udmærker sig fremfor Dyrene. »Thi«, vil Digteren sige, »er det vel at udmærke sig, nei saa foretrækker jeg langt, langt Tausheden derude; jeg foretrækker den, nei, der er ingen Sammenligning, den udmærker sig uendeligt fremfor Menneskene, som kunne tale«. I Naturens Taushed mener Digteren nemlig at fornemme Guddommens Stemme; i Menneskenes travle Talen mener han ikke blot ikke at fornemme Guddommens Stemme, men ikke engang at fornemme, at Mennesket er i Slægtskab med Guddommen. Digteren siger: Talen er Menneskets Fortrin for Dyret, ja ganske vist – dersom han kan tie.

Men at kunne tie, det kan Du lære derude hos Lilien og Fuglen, hvor der er Taushed, og ogsaa noget Guddommeligt i denne Taushed. Derude er Taushed; og ikke blot naar Alt tier i den tause Nat, men ogsaa naar Dagen igjennem de tusinde Strænge ere i Bevægelse, og Alt er som et Hav af Lyd, er der dog Taushed derude: hver især gjør det saa godt, at Ingen af dem, og ikke Alle tilsammen gjøre noget Brud paa den høitidelige Taushed. Derude er Taushed. Skoven er taus; selv naar den hvisker, er den dog taus. Thi Træerne, endog hvor de staae tættest i Mængde, holde hinanden, hvad Menneskene saa sjeldent, trods givet Løfte, holde hinanden: det bliver mellem os. Havet er taust; selv naar det raser larmende, er det dog taust. I første Øieblik hører Du maaskee feil, og Du hører det larme. Dersom Du haster og gaaer med den Besked, gjør Du Havet Uret. Hvis Du derimod giver Tid og hører nøiere efter, hører Du – forunderligt! – Du hører Tausheden; thi Eensformighed er dog ogsaa Taushed. Naar Taushed i Aftenen hviler over Landskabet, og Du fra Engen hører den fjerne Brølen, eller Du fjernt fra Bondemandens Huus hører Hundens hjemlige Stemme: saa kan man ikke sige, at dette Brøl eller denne Stemme forstyrrer Tausheden, nei Dette hører med til Tausheden, er i hemmelighedsfuld, forsaavidt igjen i taus Overeenskomst med Tausheden, det forøger den.

Og lad os nu nærmere betragte Lilien og Fuglen, af hvilke vi skulle lære. Fuglen tier og bier: den veed, eller rettere, den troer fuldt og fast, at Alt skeer i sin Tid, derfor bier Fuglen; men den veed, at den tilkommer det ikke at vide Tid eller Dag, derfor tier den. Det vil nok skee til den beleilige Tid, siger Fuglen, dog nei, det siger Fuglen ikke, den tier; men dens Taushed er talende, og denne dens Taushed siger, at den troer det, og fordi den troer det, derfor tier den og bier. Naar saa Øieblikket kommer, saa forstaaer den tause Fugl, at det er Øieblikket; den benytter det, og den blev aldrig beskjæmmet. Saaledes ogsaa med Lilien, den tier og bier. Den spørger ikke utaalmodigt »naar kommer Foraaret?« thi den veed, det kommer i den beleilige Tid, den veed, at det var den selv mindst tjenligt, om den fik Lov at bestemme Aarets Tider; den siger ikke »naar faae vi dog Regn?« eller »naar faae vi dog Solskin?«, eller »nu fik vi for megen Regn«, eller »nu var Heden for stærk;« den spørger ikke forud, hvorledes Sommeren vil blive iaar, hvor lang eller hvor kort: nei, den tier og bier – saa eenfoldig er den, men bedraget blev den dog aldrig, hvad jo ogsaa kun kan hænde Klogskaben, ikke Eenfoldigheden, der ikke bedrager og ikke bedrages. Saa kommer da Øieblikket, og naar saa Øieblikket kommer, saa forstaaer den tause Lilie, at nu er Øieblikket, og den benytter det. O, I Eenfoldighedens dybsindige Læremestere, mon det ogsaa skulde være muligt, naar man taler, at træffe »Øieblikket«? Nei, kun ved at tie træffer man Øieblikket; idet man taler, blot man siger eet Ord, gaaer man Glip af Øieblikket; kun i Tausheden er Øieblikket. Og derfor hændte det vel saa sjeldent et Menneske, at han ret kom til at forstaae, naar Øieblikket er, og til ret at benytte Øieblikket, derfor, fordi han ikke kan tie. Han kan ikke tie og bie, deraf lader det sig maaskee forklare, om Øieblikket slet ikke kommer for ham; han kan ikke tie, deraf lader det sig maaskee forklare, at han ikke mærkede Øieblikket, da det kom for ham. Thi Øieblikket, skjøndt svangert med sin rige Betydning, lader dog ikke noget Bud gaae forud, for at melde dets Ankomst, dertil kommer det for hurtigt, naar det kommer, der er jo ikke et Øiebliks Tid forinden; ei heller kommer Øieblikket, hvor betydningsfuldt det end er i sig selv, med Larm eller med Anskrig, nei det kommer sagte, med lettere Fjed end nogen Skabnings letteste Gang, thi det kommer med det Pludseliges lette Fjed, listende kommer det: derfor maa man være ganske taus, hvis man skal fornemme »nu er det der«; og i næste Øieblik er det forbi, derfor maa man have været ganske taus, hvis det skal lykkes En at benytte det. Men dog afhænger Alt af »Øieblikket«. Og dette er vistnok Ulykken i langt, langt de fleste Menneskers Liv, at de aldrig fornam »Øieblikket«, at i deres Liv det Evige og det Timelige kun skiltes ad, og hvorfor? fordi de ikke kunde tie.

Fuglen tier og lider. Hvor megen Hjertesorg den end har, den tier. Selv Ørkenens eller Eensomhedens tungsindige Klagesanger tier. Den sukker tre Gange, derpaa tier den, sukker atter tre Gange; men væsentligen tier den. Thi hvad det er, siger den ikke, den klager ikke, den anklager Ingen, den sukker for igjen at komme til at tie. Det er nemlig som vilde Tausheden sprænge den, derfor maa den sukke, for at kunne tie. Fritagen for Lidelse er Fuglen ikke; men den tause Fugl fritager sig for hvad der gjør Lidelsen tungere, Andres misforstaaede Deeltagelse, for hvad der gjør Lidelsen langvarigere, den megen Talen om Lidelsen, for hvad der gjør Lidelsen til det som er værre end Lidelse, til Utaalmodighedens og Sørgmodighedens Synd. Thi tro ikke, at det dog kun er lidt Falskhed af Fuglen, at den tier, naar den lider, at den dog i sit Inderste, hvor taus den end er mod Andre, ikke tier, at den der klager over sin Skjebne, anklager Gud og Mennesker, og lader »Hjertet synde i Sorgen«. Nei, Fuglen tier og lider. Ak, det gjør Mennesket ikke. Men hvoraf kommer det vel, at menneskelig Lidelse, sammenlignet med Fuglens, synes saa frygtelig? mon ikke deraf, at Mennesket kan tale? Nei, ikke deraf, thi det er jo et Fortrin, men deraf, at Mennesket ikke kan tie. Det er nemlig ikke som den Utaalmodige, og, endnu heftigere, den Fortvivlede mener at forstaae det, naar han – og dette er allerede en Misbrug af Talen og Stemmen – naar han siger eller skriger: »gid jeg havde en Stemme som Stormens, at jeg kunde udsige al min Lidelse, som jeg fornemmer den!« O, det var kun et daarligt Middel, i samme Grad vil han kun stærkere fornemme den. Nei, men hvis Du kunde tie, hvis Du havde Fuglens Taushed, saa skulde Lidelsen vel blive mindre.

Og som Fuglen, saaledes Lilien, den tier. Om den end staaer og lider idet den visner, den tier; forstille sig kan det uskyldige Barn ikke – det forlanges da heller ikke, og det er en Lykke for den, at den ikke kan det, thi sandeligen den Kunst at kunne forstille sig kjøbes dyrt – forstille sig kan den ikke, den kan ikke for, at den skifter Farve, og derved forraader, hvad man jo kjender paa denne Bleghedens Farveskiften, at den lider; men tie gjør den. Den vilde gjerne holde sig opreist for at skjule hvad den lider, dog dertil har den ikke Kræfter, ikke dette Herredømme over sig selv, dens Hoved synker mat og bøiet, den Forbigaaende – hvis ellers nogen Forbigaaende har saa megen Deeltagelse, at han agter paa den! – den Forbigaaende forstaaer, hvad dette betyder, det er talende nok; men Lilien tier. Saaledes Lilien. Men hvoraf kommer det vel, at menneskelig Lidelse sammenlignet med Liliens synes saa frygtelig, mon ikke deraf, at den ikke kan tale? Hvis Lilien kunde tale, og hvis den da, ak som Mennesket, ikke havde lært den Kunst at tie: mon da ikke ogsaa dens Lidelse vilde blive frygtelig? Men Lilien tier. For Lilien er det at lide det at lide, hverken mere eller mindre. Dog just naar det at lide hverken er mere eller mindre end det at lide, saa er Lidelsen saa meget som det er muligt enkelt- og eenfoldiggjort, og gjort saa liden som det er muligt. Mindre kan Lidelsen ikke blive, da den jo dog er, og altsaa er hvad den er. Men derimod kan Lidelsen blive uendelig meget større, naar den ikke nøiagtigt bliver hverken mere eller mindre end det den er. Naar Lidelsen er hverken mere eller mindre, altsaa naar Lidelsen kun er det Bestemte, den er: er den, selv om den var den største Lidelse, det Mindste den kan være. Men naar det bliver ubestemt, hvor stor Lidelsen egentligen er, saa bliver Lidelsen større; denne Ubestemthed forøger Lidelsen uendelig. Og denne Ubestemthed fremkommer just ved dette Menneskets tvetydige Fortrin at kunne tale. Lidelsens Bestemthed derimod, at den er hverken mere eller mindre end hvad den er, naaes kun igjen ved at kunne tie; og denne Taushed kan Du lære af Fuglen og Lilien.

Derude hos Lilien og Fuglen er der Taushed. Men hvad udtrykker denne Taushed? Den udtrykker Ærbødighed for Gud, at det er ham, der raader, og ham ene, hvem Viisdom og Forstand tilkommer. Og just fordi denne Taushed er Ærefrygt for Gud, er, som den kan være det i Naturen, Tilbedelse, derfor er denne Taushed saa høitidelig. Og fordi denne Taushed er saaledes høitidelig, derfor er det, man fornemmer Gud i Naturen – hvad Under vel, naar Alt tier af Ærbødighed for ham! selv om han ikke taler, det at Alt tier af Ærbødighed for ham, virker jo paa En, som om han talte.

Hvad Du derimod ved ingen »Digters« Hjælp kan lære af Tausheden derude hos Lilien og Fuglen, hvad kun Evangeliet kan lære Dig, er, at det er Alvor, at det skal være Alvor, at Fuglen og Lilien skal være Læremesteren, at Du skal tage efter dem, lære af dem, ganske alvorligt, at Du skal blive taus som Lilien og Fuglen.

Og allerede dette er jo Alvoren – naar det forstaaes retteligen, ikke som den drømmende Digter, eller som Digteren, der lader Naturen drømme om sig – dette: at Du ude hos Lilien og Fuglen fornemmer, at Du er for Gud, hvad oftest glemmes saa ganske i Talen og Samtalen med andre Mennesker. Thi naar vi blot tale To sammen, end mere naar vi ere Ti eller Flere, glemmes det saa let, at Du og Jeg, vi To, eller at vi Ti ere for Gud. Men Lilien, der er Læremesteren, er dybsindig. Den indlader sig slet ikke med Dig, den tier, og ved at tie vil den betyde Dig, at Du er for Gud, at Du husker paa, at Du er for Gud – at Du ogsaa i Alvor og Sandhed maatte blive taus for Gud.

\workheading{Tvende ethisk-religieuse Smaa-Afhandlinger (af H. H.)}{1849}

\worknote{Antologiens ''Two Ethical-Religious Essays''. Hong-udvalget bringer den anden afhandling, ''Om Forskjellen mellem et Genie og en Apostel'', fra det berømte sted om Paulus (som geni kontra som apostel) og de tre punkter om Myndighed, Immanens og Teleologi. Den akademiske indledning og den første afhandling (om at lade sig ihjelslaae for Sandheden) er ikke medtaget.}

\sectionheading{Om Forskjellen mellem et Genie og en Apostel (uddrag)  ·  SV¹ XI, 96–109}

Om Forskjellen mellem et Genie og en Apostel.

Som Genie kan Paulus hverken udholde Sammenligning med Plato eller Shakspeare; som Forfatter af skjønne Lignelser kommer han temmelig langt ned; som Stilist er han et aldeles obscurt Navn – og som Tapetmager: ja, det maa jeg sige, jeg veed ikke hvor høit han i denne Henseende kan komme. See, dum Alvorlighed gjør man altid rigtigst i at forvandle til Spøg, og saa kommer Alvoren, Alvoren: at Paulus er Apostel; og som Apostel har han igjen intet, slet intet Svogerskab hverken med Plato eller Shakespeare eller Stilister eller Tapetmagere, de ere alle (Plato lige saa fuldt som Shakespeare og som Tapetmager Hansen) uden nogen Sammenligning med ham.

Et Genie og en Apostel er det qvalitative Forskjellige, ere Bestemmelser, som høre hjemme hver i sin qvalitative Sphære: Immanentsens og Transcendentsens: 1) Geniet kan derfor vel have noget Nyt at bringe, men dette forsvinder igjen i Slægtens almindelige Assimilation, ligesom Differentsen »Genie« forsvinder, saasnart man tænker Evigheden; Apostelen har paradoxt noget Nyt at bringe, hvis Nyhed, netop fordi den er væsentlig paradox og ikke en Anticipation i Forhold til Slægtens Udvikling, bliver bestandigt, ligesom en Apostel bliver i al Evighed en Apostel, og ingen Evighedens Immanents sætter ham væsentligen paa lige Linie med alle Mennesker, da han er væsentligen paradox forskjellig. 2) Geniet er hvad han er ved sig selv det vil sige ved det han er i sig selv; en Apostel er hvad han er ved sin guddommelige Myndighed. 3) Geniet har kun immanent Teleologie; Apostelen er absolut paradoks teleologisk stillet.

1. Al Tænken aander i Immanentsen, hvorimod det Paradoxe og Troen danner en qvalitativ Sphære for sig. Immanent er, i Forholdet mellem Menneske og Menneske qua Menneske, enhver Differents et for den væsentlige og evige Tænkning Forsvindende, et Moment, der momentant vel har sin Gyldighed, men væsentligen forsvinder i Evighedens væsentlige Lighed. Genie er, som Ordet selv siger (ingenium, det Medfødte, Primitivitet (primus), Originalitet (origo), Oprindelighed og saa videre) Umiddelbarheden, Natur-Bestemmelsen, Geniet fødes. Allerede længe førend der kan være Tale om, hvorvidt Geniet nu vil henføre sin sjeldne Begavethed til Gud eller ikke, er det Genie, og er Genie, selv om det ikke gjør det. Med Geniet kan der foregaae den Forandring, at det udvikler sig til at være hvad det ϰατα δυναμιν er, at det kommer i bevidst Besiddelse af sig selv. Forsaavidt man for at betegne det Nye, et Genie kan have at bringe, bruger Udtrykket »Paradox«, bruges dette dog kun i uvæsentlig Forstand om det transitoriske Paradox, om Anticipationen, der fortætter sig til noget Paradoxt, hvilket dog igjen forsvinder. Et Genie kan i sin første Meddelen være paradox, men jo mere han kommer til sig selv, jo mere forsvinder det Paradoxe. Et Genie kan maaskee være et Aarhundrede forud for sin Tid og derfor staae som et Paradox, men tilsidst vil Slægten dog saaledes assimilere sig det eengang Paradoxe, saa det ikke mere er paradoxt. –

Anderledes med en Apostel. Ordet selv tyder paa Forskjellen. En Apostel fødes ikke; en Apostel er en Mand, der af Gud kaldes og beskikkes, af ham sendes i Ærende. En Apostel udvikles ikke saaledes, at han successivt bliver hvad han ϰατα δυναμιν er. Thi forud for det at blive en Apostel gaaer ingen potentiel Mulighed; ethvert Menneske er væsentligen lige nær til at blive det. En Apostel kan aldrig saaledes komme til sig selv, at han bliver sig sit Apostel-Kald bevidst som Moment i sin egen Livs-Udvikling. Apostel-Kaldet er et paradox Faktum, som i første og sidste Øieblik af hans Liv staaer paradox udenfor hans personlige Identitet med sig selv som den Bestemte han er. En Mand er maaskee forlængst kommen til Skjels-Aar og Alder, da bliver han kaldet til Apostel. Han bliver ved dette Kald ikke et bedre Hoved, han faaer ikke mere Phantasie, ikke større Skarpsindighed og saa videre, ingenlunde, han bliver sig selv, men ved det paradoxe Faktum af Gud sendt i et bestemt Ærende. Ved dette paradoxe Faktum er Apostelen i al Evighed gjort paradox forskjellig fra alle andre Mennesker. Det Nye, han kan have at forkynde, er det væsentlige Paradoxe. Hvor længe det end bliver forkyndt i Verden, det bliver væsentligen lige Nyt, lige Paradoxt, ingen Immanents kan assimilere sig det. Apostelen forholdt sig jo ikke som det ved Naturbegavelse udmærkede Menneske, der var forud for sin Samtid, han var maaskee hvad vi kalde et eenfoldigt Menneske, men ved et paradox Faktum blev han kaldet til at forkynde dette Nye. Selv om Tænkningen vilde mene at kunne assimilere sig Læren: Maaden, paa hvilken Læren kom ind i Verden, kan den ikke assimilere sig; thi det væsentlige Paradox er netop Protesten mod Immanentsen. Men Maaden, paa hvilken en saadan Lære kom ind i Verden, er netop det qvalitative Afgjørende, hvad kun ved Svig eller Tankeløshed kan oversees.

2. Et Genie vurderes reent æsthetisk efter hvad dets Indhold, dets Vægtfylde befindes at være; en Apostel er hvad han er ved at have guddommelig Myndighed. Den guddommelige Myndighed er det qvalitative Afgjørende. Det er ikke ved æsthetisk eller philosophisk at vurdere Lærens Indhold, at jeg skal eller kan komme til det Resultat: ergo er Den, der har foredraget denne Lære, kaldet ved en Aabenbaring, ergo er han en Apostel. Forholdet er lige omvendt: den ved en Aabenbaring Kaldede, hvem en Lære betroes, argumenterer fra, at det er en Aabenbaring, fra, at han har Myndighed. Jeg skal ikke høre paa Paulus, fordi han er aandrig eller mageløs aandrig, men jeg skal bøie mig under Paulus, fordi han har guddommelig Myndighed; og i ethvert Tilfælde maa det blive Paulus's Ansvar, at han sørger for, at han frembringer dette Indtryk, hvad enten nu Nogen vil bøie sig under hans Myndighed eller ikke. Paulus skal ikke beraabe sig paa sin Aandrighed, thi saa er han en Nar; han skal ikke indlade sig paa en reen æsthetisk eller philosophisk Discussion af Lærens Indhold, thi saa er han Distrait. Nei han skal beraabe sig paa sin guddommelige Myndighed og netop ved den, medens han villigt offrer sit Liv og Alt, forhindre alle næsvise æsthetiske og philosophiske Ligefremheder mod Lærens Indhold og Form. Paulus skal ikke anbefale sig og sin Lære ved Hjælp af de skjønne Billeder, omvendt skulde han vel sige til den Enkelte: »hvad enten nu Lignelsen er skjøn eller den er luslidt og afdanket, det er ligemeget, Du skal betænke, at Det jeg siger er mig betroet ved en Aabenbaring, saa det er Gud selv eller den Herre Jesus Christus der taler, og Du skal ikke formastelig indlade Dig paa at critisere Formen. Jeg kan ikke, jeg tør ikke tvinge Dig til at lyde, men jeg gjør Dig, gjennem Din Samvittigheds Forhold til Gud, evigt ansvarlig for Dit Forhold til denne Lære derved, at jeg har forkyndt den som mig aabenbaret og altsaa forkyndt den med guddommelig Myndighed.«

Myndigheden er det qvalitativt Afgjørende. Eller er der ikke, endog blot indenfor det menneskelige Livs Relativitet, om den end immanent forsvinder, en Forskjel mellem et Kongebud og en Digters eller en Tænkers Ord? Og hvilken er Forskjellen, uden den, at Kongebudet har Myndighed, og derfor forbyder al æsthetisk og critisk Næsvished med Hensyn til Form og Indhold. Digteren, Tænkeren derimod har, end ikke indenfor denne Relativitet, nogen Myndighed, hans Udsagn vurderes reent æsthetisk eller philosophisk ved at vurdere Indhold og Form. Men hvad er det vel der fra Grunden har forvirret det Christelige, uden det, at man først i Tvivlen saa noget nær er bleven uvis, om der er en Gud til, og saa, at man i Oprørskhed mod alle Autoriteter har glemt hvad Myndighed er, og dens Dialektik. En Konge er saaledes sandselig til, at man sandseligt kan forvisse sig derom, og gjøres det fornødent, kan dog maaskee Kongen ret sandseligt forvisse En om, at han er til. Men saaledes er Gud ikke til. Dette har Tvivlen benyttet til at bringe Gud paa lige Linie med alle dem, som ingen Myndighed har, paa lige Linie med Genier, Digtere og Tænkere, hvis Udsagn slet og ret blot vurderes æsthetisk eller philosophisk; og dersom det saa er godt sagt, saa er Manden Genie – og dersom det saa er ualmindeligt godt sagt, saa, saa er det Gud, der har sagt det!!!

Paa den Maade practiseres Gud egentligen udenfor. Hvad skal han gjøre? Standser Gud et Menneske paa hans Vei, kalder ham ved en Aabenbaring, og sender ham udrustet med guddommelig Myndighed til de øvrige Mennesker: da sige disse til ham: fra hvem er Du? Han svarer: fra Gud. Men see, Gud kan nu ikke paa en saadan sandselig Maade hjælpe sin Udsending, som en Konge kan det, der giver ham Soldater eller Politie-Betjente med, eller hans Ring, eller hans Haandskrift som Alle kjende, kort Gud kan ikke være Menneskene til Tjeneste med at skaffe dem en sandselig Vished for, at en Apostel er en Apostel – dette vilde ogsaa være Nonsens. Selv Miraklet, hvis Apostelen har denne Gave, giver ingen sandselig Vished; thi Miraklet er Troens Gjenstand. Og desuden er det da Nonsens at faae sandselig Vished om, at en Apostel er en Apostel (den paradoxe Bestemmelse af et Aands Forhold), ligesom det er Nonsens at faae sandselig Vished om, at Gud er til, da jo Gud er Aand. Altsaa Apostelen siger, at han er fra Gud. De Andre svare: ja saa, saa lad os see, om Din Læres Indhold er guddommelig, thi saa ville vi antage den, samt at den er Dig aabenbaret. Paa den Maade er baade Gud og Apostelen narret. Den Kaldedes guddommelige Myndighed skulde netop være det sikkre Værn, der betryggede Læren og holdt den paa det Guddommeliges majestætiske Afstand fra Næsviisheder, istedetderfor maa Lærens Indhold og Form lade sig recensere og oversnøfle – for at man ad den Vei kan komme til det Resultat, om det nu var en Aabenbaring eller ikke; og imidlertid maa formodentligen Apostelen og Gud vente i Porten eller hos Portneren, indtil Sagen bliver afgjort af de Vise paa 1ste Sahl. Den Kaldede skulde efter Guds Bestemmelse bruge sin guddommelige Myndighed til at jage alle Næsvise bort, der ikke vil lystre men raisonnere; og istedet derfor faae Menneskene i en Fart Apostelen forvandlet til en Examinand, der saadan kommer tiltorvs med en ny Lære.

Hvad er da Myndighed? Er Myndighed Lærens Dybsind, dens Fortrinlighed, dens Aandrighed? Ingenlunde. Dersom Myndighed saaledes blot skulde betegne, i en anden Potens eller redupliceret, at Læren er dybsindig: saa er der netop ingen Myndighed; thi hvis da en Lærende heelt og fuldt ved Forstaaelsen tilegnede sig denne Lære, saa blev der jo ingen Forskjel tilbage mellem Læreren og den Lærende. Myndighed er derimod Noget, som bliver uforandret, som man ikke kan erhverve ved paa det fuldeste at have forstaaet Læren. Myndighed er en specifik Qvalitet, der træder til andetstedsfra og netop qvalitativt gjør sig gjældende, naar Udsagnets eller Gjerningens Indhold æsthetisk er sat i Indifferents. Lad os tage et Exempel, saa simpelt som muligt, hvor dog Forholdet viser sig. Naar Den, der har Myndigheden til at sige det, siger til et Menneske: gaa! og naar Den, der ikke har Myndigheden, siger: gaa! saa er jo Udsagnet (gaa!) og dets Indhold identisk, æsthetisk vurderet er det om man saa vil lige godt sagt, men Myndigheden gjør Forskjellen. Dersom Myndigheden ikke er det Andet (το ετερον), dersom den paa nogen Maade blot skulde betegne en Potenseren indenfor Identiteten, saa er der netop ingen Myndighed. Dersom saaledes en Lærer er sig begeistret bevidst, at han selv existerende udtrykker og har udtrykt, med Opoffrelse af Alt, den Lære han forkynder: saa kan denne Bevidsthed vel give ham en vis og en fast Aand, men den giver ham ikke Myndighed. Hans Liv som Beviis for Lærens Rigtighed er ikke det Andet (το ετερον) men er en simpel Fordoblelse. Det at han lever efter Læren beviser ikke, at den er rigtig; men fordi han er selv overbeviist om Lærens Rigtighed, derfor lever han efter den. Derimod, hvad enten en Politie-Officiant for Exempel er en Slyngel, eller han er en retskaffen Mand, saasnart han er i Function, har han Myndighed.

For nærmere at belyse dette for den paradox-religieuse Sphære saa vigtige Begreb: Myndighed, skal jeg forfølge Myndighedens Dialektik. I Immanentsens Sphære lader Myndigheden sig slet ikke tænke, eller den lader sig kun tænke som forsvindende. Forsaavidt der da i de politiske, borgerlige, sociale, huslige, disciplinaire Forhold er Tale om Myndighed, eller Myndighed udøves, er Myndigheden dog kun et transitorisk Moment, et Forsvindende, der enten allerede i Timeligheden senere forsvinder, eller forsvinder, forsaavidt Timeligheden og Jordlivet selv er et transitorisk Moment, der forsvinder med alle dens Differentser. Til Grund for ethvert Forhold mellem Menneske og Menneske qua Menneske kan kun tænkes at ligge Forskjelligheden indenfor Immanentsens Identitet, det er, den væsentlige Lighed. Det ene Menneske kan ikke tænkes (saa hører al Tænkning op, som den ganske conseqvent gjør det i det Paradox-religieuses og Troens Sphære) ved en specifik Qvalitet at være forskjellig fra alle andre. Alle menneskelige Differentser mellem Menneske og Menneske qua Menneske forsvinde for Tænkningen som Momenter i Identitetens Totale og Qvalitet. I Momentet skal jeg være saa god at respektere og lystre Differentsen; men det er mig tilladt, religieust at opbygges ved den Vished, at i Evigheden forsvinde Differentserne, den der udmærker mig og den der nedtrykker mig. Som Undersaat skal jeg ære og lyde Kongen med udeelt Sjel, men det er mig tilladt religieust at opbygges ved den Tanke, at væsentligen er jeg Borger i Himlen, og at jeg, hvis jeg engang der træffer sammen med den afdøde Majestæt, da skal jeg ikke være forpligtet i undersaatlig Lydighed mod ham.

Saaledes altsaa med Forholdet mellem Menneske og Menneske qua Menneske. Men mellem Gud og Menneske er der en evig væsentlig qvalitativ Forskjel, hvilken ingen uden en formastelig Tænkning kan lade forsvinde i den Blasphemie, at Gud og Menneske vel i Endelighedens transitoriske Moment differentiere sig, saa det bør sig Mennesket her i dette Liv at lyde og tilbede Gud, men i Evigheden skulde Forskjellen forsvinde i den væsentlige Lighed, saa Gud og Mennesket blev Ligemænd i Evigheden, ligesom Kongen og Kammertjeneren.

Mellem Gud og Menneske er og bliver der altsaa en evig, væsentlig qvalitativ Forskjel. Det paradox-religieuse Forhold (hvilket ganske rigtigt ikke lader sig tænke, men kun troe), fremkommer da, naar Gud beskikker et enkelt Menneske til at have guddommelig Myndighed, vel at mærke i Forhold til det ham af Gud Betroede. Den saaledes Kaldede forholder sig ikke i Forholdet: Menneske til Menneske qua Menneske; han forholder sig ikke i en qvantitativ Differents (som Genie, udmærket Begavet og saa videre) til andre Mennesker. Nei, han forholder sig paradox ved at have en specifik Qvalitet, som ingen Immanents kan tilbagekalde i Evighedens Lighed; thi den er væsentlig paradox og efter Tænkningen (ikke før, foran Tænkningen) mod Tænkningen. Har en saadan Kaldet en Lære at bringe efter guddommelig Ordre, og et andet Menneske, lad os tænke det, af sig selv og ved sig selv udfandt det Samme: saa blive disse Tvende dog i al Evighed ikke lige; thi den første er ved sin paradox specifike Qvalitet (den guddommelige Myndighed) forskjellig fra ethvert andet Menneske, og fra den immanent til Grund for alle andre menneskelige Differentser liggende Bestemmelse af den væsentlige Lighed. Bestemmelsen »en Apostel« hører hjemme i Transcendentsens Sphære, den paradox-religieuse Sphære, som, ganske conseqvent, ogsaa har et qvalitativt forskjelligt Udtryk for andre Menneskers Forhold til en Apostel: de forholde sig nemlig troende til ham, medens al Tænkning ligger og er og aander i Immanentsen. Men Troen er ikke en transitorisk Bestemmelse, saalidet som Apostelens paradoxe Qvalification var et Transitorisk.

I Forholdet mellem Menneske og Menneske qua Menneske var der altsaa ingen bestaaende eller bestandig Myndighedens Differents tænkelig, den var et Forsvindende. Lader os imidlertid et Øieblik dvæle ved nogle Exempler paa saadanne saakaldte og i Timelighedens Vilkaar sande Myndigheds-Forhold mellem Menneske og Menneske qua Menneske, for derved at blive opmærksom paa den væsentlige Betragtning af Myndigheden. En Konge antages jo at have Myndighed. Hvoraf kommer det nu, at man endog støder sig over, at en Konge er aandrig, Kunstner og saa videre? Det kommer dog vel deraf, at man hos ham væsentligen accentuerer den kongelige Myndighed, og i Sammenligning med denne finder almindeligere Bestemmelser af menneskelig Differents at være et Forsvindende, et Uvæsentligt, et forstyrrende Tilfældigt. Et Regjerings-Collegium antages at have Myndighed i sin bestemte Kreds. Hvoraf kommer det nu, at man vilde stødes, hvis et saadant Collegium i sine Decreter for Exempel var virkelig aandrig, vittig, dybsindig? Fordi man ganske rigtigt qvalitativt accentuerer Myndigheden. At spørge om en Konge er Genie – for i saa Fald at ville lyde ham, er i Grunden Majestæts-Forbrydelse; thi i Spørgsmaalet er indeholdt en Tvivl i Retning af Underkastelse under Myndighed. At ville lyde et Collegium, dersom det kan sige Vittigheder, er i Grunden at gjøre Nar af Collegiet. At ære sin Fader, fordi han er et udmærket Hoved, er Impietet.

Dog, som sagt, i Forholdet mellem Menneske og Menneske qua Menneske er Myndigheden, om den end er, et Forsvindende, og Evigheden afskaffer al jordisk Myndighed. Men nu i Transcendentsens Sphære. Lad os tage Exemplet saa simpelt, men netop derfor ogsaa saa iøinefaldende som muligt. Naar Christus siger: »der er et evigt Liv;« og naar theologisk Candidat Petersen siger »der er et evigt Liv«: saa sige de Begge det Samme, der er ikke i det første Udsagn indeholdt mere Deduction, Udvikling, Dybsind, Tankefylde end i det sidste; begge Udsagnene ere, æsthetisk vurderede, lige gode. Og dog er der vel en evig qvalitativ Forskjel! Christus er, som Gud-Mennesket, i Besiddelse af den specifike Myndigheds Qvalitet, hvilken ingen Evighed kan mediere, eller lægge Christus paa lige Linie med den væsentlige menneskelige Ligelighed. Christus lærte derfor med Myndighed. At spørge om Christus er dybsindig, er Blasphemie, og er Forsøg paa underfundigt (det være nu bevidst eller ubevidst) at tilintetgjøre Ham; thi i Spørgsmaalet er indeholdt en Tvivl i Retning af Hans Myndighed, og gjort et Forsøg paa i næsviis Ligefremhed at ville vurdere og censurere Ham, som var Han til Examen og skulde overhøres, istedetfor at Han er den, hvem al Magt er given i Himlen og paa Jorden.

Dog sjeldent, saare sjeldent hører eller læser man nuomstunder et religieust Foredrag, der er ganske correct. De bedre af dem fuske dog gjerne en Smule i, hvad man kunde kalde det ubevidste eller det velmeente Oprør, idet de forsvare og hævde det Christelige af yderste Kraft – i feile Kategorier. Lad mig tage et Exempel, det første det bedste. Jeg tager det helst af en Tydsker, saa veed jeg da, at Ingen, ikke den Dummeste og ikke den Arrigste, kan falde paa, at jeg skriver dette om en Sag, der i mine Tanker er uendelig vigtig, – for at sigte til een eller anden Geistlig. Biskop Sailer prædiker i en Homilie paa femte Søndag i Fasten over Texten Joh. VIII. 47–51. Han vælger disse to Vers: »Wer von Gott ist, der höret Gottes Wort, og, Wer mein Wort hält, der siehet den Tod nicht; og siger derpaa: es sind in diesen Worten des Herrn drei große Räthsel gelöset, mit denen sich die Menschen von jeher den Kopf so oder anders zerbrochen haben.« Der har vi det. Det Ord: Räthsel, og især drei große Räthsel, og saa i næste Sætning, mit denen die Menschen den Kopf sich zerbrochen haben, fører strax Tanken hen til det Dybsindige i intellectuel Forstand, det Grundende, det Grublende, det Speculerende. Men hvorledes kan dog vel et simpelt apodictisk Udsagn være dybsindigt? et apodictisk Udsagn, der kun er hvad det er derved, at Den og Den har sagt det; et Udsagn, der slet ikke vil forstaaes eller udgrundes, men blot troes. Hvorledes kan et Menneske falde paa, at en Gaade skulde være løst, i Retning af det grublende og grundende Dybsindige, ved et ligefremt Udsagn, ved et Assertum? Spørgsmaalet er jo: er der et evigt Liv? Svaret er: der er et evigt Liv. Hvor i al Verden stikker nu det Dybsindige? Dersom Christus ikke er Den, som har sagt det, og dersom Christus ikke er Den, han har sagt sig at være, saa maa jo, hvis Udsagnet i sig selv er dybsindigt, det Dybsindige alligevel være at finde. Lad os tage Hr. theologisk Candidat Petersen, der jo ogsaa siger: der er et evigt Liv. Hvo i al Verden vilde det dog vel falde paa at beskylde ham for Dybsindighed paa Grund af et ligefremt Udsagn? Det Afgjørende ligger saaledes ikke i Udsagnet, men i at det er Christus der har sagt det; men det Confunderende er, at man ligesom for at lokke Menneskene til at troe, taler noget om det Dybsindige og det Dybsindige. En christelig Præst maa, hvis han vil tale correct, ganske simpelt sige »vi har Christi Ord for, at der er et evigt Liv: dermed er den Ting afgjort. Her er ikke Tale hverken om Hovedbrud eller Speculation, men om, at det er Christus, der, ikke i Egenskab af Dybsindig, men med sin guddommelige Myndighed har sagt det.« Lad os gaae videre, lad os antage, at En troer, at der er et evigt Liv, fordi Christus har sagt det, saa gaaer han jo netop troende uden om al det Dybsindige og det Grublende og Grundende, »hvormed man bryder sit Hoved«. Derimod lad os tage En, der dybsindigt vil bryde sit Hoved med det Spørgsmaal om Udødeligheden: mon han ikke vil have Ret i at benegte, at det ligefremme Udsagn er et dybsindigt Svar paa Spørgsmaalet? Det Plato siger om Udødeligheden er virkelig dybsindigt, vundet ved en dyb Grunden; men stakkels Plato har heller ingen Myndighed.

Sagen er imidlertid denne. Tvivlen og Vantroen, der gjør Troen forfængelig, har blandt Andet ogsaa gjort Menneskene generede ved at lyde, ved at bøie sig under Myndighed. Denne Oprørskhed sniger sig ind selv i de Bedres Tankegang, dem maaskee ubevidst, og saa begynder alt dette Forskruede, som i Grunden er Forræderie, om det Dybe og det Dybe, og det Vidunderlig-deilige, som man kan skimte og saa videre Skulde man derfor med eet bestemt Prædikat betegne det christelig-religieuse Foredrag, som det nu høres og læses, da maatte man sige, det er affecteret. Naar man ellers taler om en Præsts Affectation, tænker man maaskee paa, at han pynter og sminker sig, eller at han sødligt smægter med Stemmen, eller at han norsk snurrer paa R og rynker Brynet, eller at han anstrænger sig i Kraftstillinger og i Opvækkelsens Spring og saa videre Dog alt Sligt er mindre vigtigt, om det end altid var ønskeligt, at det ikke var. Men det Fordærvelige er, naar Prædike-Foredragets Tankegang er affecteret, naar dets Rettroenhed naaes ved at lægge Accenten paa et reent forkeert Sted, naar det i Grunden opfordrer til at troe paa Christus, prædiker Troen paa Ham paa Grund af Det, som slet ikke kan blive Troens Gjenstand. Dersom en Søn vilde sige »jeg lyder min Fader, ikke fordi han er min Fader, men fordi han er Genie, eller fordi hans Befalinger altid ere dybsindige og aandrige«: saa er denne sønlige Lydighed affecteret. Sønnen accentuerer noget reent Forkeert, accentuerer det Aandrige, det Dybsindige ved en Befaling, medens en Befaling netop er ligegyldig mod denne Bestemmelse. Sønnen vil lyde i Kraft af Faderens Dybsindighed og Aandrighed; og i Kraft deraf kan han netop ikke lyde ham, thi hans critiske Forhold i Retning af, om nu Befalingen er dybsindig og aandrig, underminerer Lydigheden. Og saaledes er det ogsaa Affectation, naar der er saa megen Tale om at tilegne sig Christendommen og troe Christus, formedelst Lærens Dybsindighed og Dybsindighed. Man tilegner sig Rettroenheden ved at accentuere noget aldeles Forkeert. Den hele moderne Spekulation er derfor affecteret ved at have afskaffet Lydigheden paa den ene Side og Myndigheden paa den anden Side, og ved saa desuagtet at ville være rettroende. En Præst, der er aldeles correct i sit Foredrag, maa tale saaledes, idet han anfører et Ord af Christus: »Dette Ord er af Den, hvem, ifølge hans eget Udsagn, al Magt er given i Himlen og paa Jorden. Du maa nu, min Tilhører, overveie med Dig selv, om Du vil bøie Dig under denne Myndighed eller ikke, antage og troe Ordet eller ikke. Men vil Du det ikke, da gaae for Gud i Himlens Skyld ikke hen og antag Ordet, fordi det er aandrigt eller dybsindigt, eller vidunderlig-deiligt, thi dette er Gudsbespottelse, det er at ville criticere Gud.« Saasnart nemlig Myndighedens, den specifik-paradoxe Myndigheds Dominant er sat i, saa er alle Forhold qvalitativt forandrede, saa er den Art Tilegnelse, som ellers er tilladelig og ønskelig, en Brøde og Formastelighed.

Men hvorledes kan nu Apostelen bevise, at han har Myndigheden? Kunde han sandseligt bevise det, saa var han netop ingen Apostel. Han har intet andet Beviis end sit eget Udsagn. Og saaledes maa det netop være; thi ellers kom jo den Troende i et ligefremt Forhold til ham, ikke i et paradoxt. I de transitoriske Forhold af Myndighed mellem Menneske og Menneske qua Menneske vil Myndigheden i Regelen være sandselig kjendelig paa Magten. En Apostel har intet andet Beviis end sit eget Udsagn, og i det Høieste sin Villighed til med Glæde at ville lide Alt for dette Udsagns Skyld. Hans Tale vil i den Henseende være kort »jeg er kaldet af Gud; gjører I nu ved mig, hvad I vil, hudfletter mig, forfølger mig, men mit sidste Ord bliver mit første: jeg er kaldet af Gud; og jeg gjør Eder evigt ansvarlige for hvad I gjøre mod mig.« Dersom det i Virkeligheden var saa, lad os tænke det, at en Apostel havde i verdslig Forstand Magten, havde stor Indflydelse og mægtige Forbindelser, ved hvilke Kræfter man seirer over Menneskenes Meninger og Domme, – dersom han da brugte den, havde han eo ipso forspildt sin Sag. Ved at bruge Magten bestemte han nemlig sin Stræben i væsentlig Identitet med de andre Menneskers, og dog er en Apostel kun hvad han er ved sin paradoxe Heterogenitet, ved at have guddommelig Myndighed, hvilken han absolut uforandret kan have, selv om han, som Paulus siger, af Menneskene ikke agtes mere værd end det Snavs, hvorpaa de træde.

3. Geniet har kun immanent Teleologie; Apostelen er paradox teleologisk stillet.

Dersom ellers noget Menneske kan siges at være absolut teleologisk stillet, saa er det en Apostel. Den ham meddeelte Lære er ikke en Opgave, der er givet ham at grunde over, den er ikke givet ham for hans egen Skyld, han er tvertimod i Ærende og har at forkynde Læren og at bruge Myndigheden. Saa lidet som Den, der bliver sendt i Byen med et Brev, har med Brevets Indhold at gjøre, men kun med at overbringe det; saa lidet som Ministeren, der sendes til et fremmed Hof, har noget Ansvar for Budskabets Indhold, men kun at besørge det rigtigt: saaledes har en Apostel hovedsageligen ene og alene at være tro i Tjenesten, hvilken er at rygte sit Ærende. Deri ligger væsentligen en Apostels opoffrende Liv, selv om han aldrig blev forfulgt, deri, »at han selv fattig kun har at gjøre Andre rige«, at han aldrig tør give sig Tid eller Ro eller Sorgløshed til i gode Dage, i otium at beriges ved Det, ved hvis Forkyndelse han beriger Andre. Han er, aandelig forstaaet, som den travle Husmoder, der neppe faaer Tid til selv at spise, for at lave Maden til de mange Munde. Og om han, idet han begyndte, turde haabe et langt Liv, hans Liv vil dog indtil det Sidste blive uforandret, thi der vil altid være Nye og Nye at forkynde Læren. Skjøndt en Aabenbaring er det paradoxe Faktum, som overgaaer Menneskets Forstand, kan man dog nok forstaae saa meget, hvad ogsaa overalt har viist sig: at et Menneske kaldes ved en Aabenbaring til at gaae ud i Verden, til at forkynde Ordet, til at handle og lide, til det uafbrudte virksomme Liv som Herrens Sendebud. At derimod et Menneske ved en Aabenbaring skulde kaldes til at sidde hen i uskiftet Bo, i et virksomt literairt far niente, til at være momentant aandrig og derpaa Samler og Udgiver af sin Aandrigheds Uvisse: er noget nær en blasphemisk Tanke.

Anderledes med et Genie; det har kun immanent Teleologie, det udvikler sig selv, og idet det udvikler sig selv, projekterer denne dets Selv-Udvikling sig som dets Virken. Det faaer da nok Betydning, maaskee endog stor Betydning, men det er ikke sig selv teleologisk stillet i Forhold til Verden og Andre. Et Genie lever i sig selv; og det kan humoristisk leve i en tilbagetrukken Selvtilfredshed uden derfor at tage sin Begavethed forfængeligt, naar det blot, uden Hensyn til om Andre have Gavn deraf eller ikke, med Alvor og Flid udvikler sig selv, følgende sin egen Genius. Geniet er derfor ingenlunde uvirksomt, det arbeider maaskee i sig selv mere end ti Forretningsmænd, det præsterer maaskee saare Meget, men hver dets Præstation har intet τελος udenfor. Dette er paa engang Geniets Humanitet og dets Stolthed: Humaniteten ligger i, at det ikke bestemmer sig teleologisk i Forhold til noget andet Menneske, som var der Nogen, der trængte til det; Stoltheden ligger i, at det immanent forholder sig til sig selv. Det er beskedent af Nattergalen, at den ikke fordrer, at Nogen skal høre paa den; men det er ogsaa stolt af Nattergalen, at den slet ikke vil være vidende om, hvorvidt Nogen hører paa den eller ikke. Geniets Dialektik vil især være til Anstød i vor Tid, hvor Mængden, Massen, Publikum og andre saadanne Abstraktioner intendere at vende op og ned paa Alt. Det høistærede Publikum, den herskesyge Mængde vil, at Geniet skal udtrykke, at han er til for dens eller for deres Skyld; det høistærede Publikum, den herskesyge Mængde seer blot den ene Side af Geniets Dialektik, støder an paa Stoltheden, og mærker ikke, at dette Samme ogsaa er Ydmygheden og Beskedenheden. Det høistærede Publikum, den herskesyge Mængde vilde derfor ogsaa tage en Apostels Existents forfængeligt. Thi vel er det sandt, at han absolut er til for Andres Skyld, udsendes for Andres Skyld; men det er ikke Mængden og ikke Menneskene og ikke det høistærede Publikum, og ikke engang det høistærede dannede Publikum, der er hans Herre eller hans Herrer – det er Gud; og Apostelen er Den, der har guddommelig Myndighed til at befale baade Mængden og Publikum.

Geniets humoristiske Selvtilfredshed er Eenheden af beskeden Resignation i Verden og stolt Opløftelse over Verden, er Eenheden af at være en unyttig Overflødighed og en kostelig Prydelse. Er Geniet en Kunstner, da frembringer han sit Kunstværk, men hverken han eller hans Kunstværk har noget τελος udenfor. Eller det er en Forfatter, der tilintetgjør ethvert teleologisk Forhold til en Omverden og bestemmer sig humoristisk som Lyriker. Det Lyriske har ganske rigtigt intet τελος uden for sig; om En skriver en Pagina Lyrik eller Folianter Lyrik, det gjør hverken fra eller til med Hensyn til at bestemme hans Virksomheds Retning. Den lyriske Forfatter bryder sig kun om Frembringelsen, nyder Frembringelsens Glæde, maaskee ofte gjennem Smerte og Anstrengelse; men han har intet med Andre at skaffe, han skriver ikke for at, for at oplyse Menneskene, for at hjælpe dem paa rette Vei, for at sætte Noget igjennem, kort han skriver ikke: for at. Og saaledes med ethvert Genie. Intet Genie har et »for at«; Apostelen har absolut paradox et »for at«.

\workheading{Sygdommen til Døden (af Anti-Climacus)}{1849}

\worknote{Antologiens ''The Sickness unto Death''. Hong-udvalget bringer den berømte begyndelse — bestemmelsen af Selvet (''Mennesket er Aand … Selvet er et Forhold, der forholder sig til sig selv'') og af fortvivlelsen — samt fra Andet Afsnit ''Fortvivlelse er Synden'' (Synd er for Gud fortvivlet ikke at ville være sig selv; det theologiske Selv; den socratiske Definition af Synd og dens fuldstændiggørelse), forargelsen over syndsforladelsen, og slutningsordet om at troens formel er modsætningen til synd. Spring markeret med […].}

\sectionheading{Sygdommen til Døden er Fortvivlelse; Fortvivlelse er Synden (uddrag)  ·  SV¹ XI, 127–241}

Sygdommen til Døden er Fortvivlelse.

A. Fortvivlelse er en Sygdom i Aanden, i Selvet.

Mennesket er Aand. Men hvad er Aand? Aand er Selvet. Men hvad er Selvet? Selvet er et Forhold, der forholder sig til sig selv, eller er det i Forholdet, at Forholdet forholder sig til sig selv; Selvet er ikke Forholdet, men at Forholdet forholder sig til sig selv. Mennesket er en Synthese af Uendelighed og Endelighed, af det Timelige og det Evige, af Frihed og Nødvendighed, kort en Synthese. En Synthese er et Forhold mellem To. Saaledes betragtet er Mennesket endnu intet Selv.

I Forholdet mellem To er Forholdet det Tredie som negativ Eenhed, og de To forholde sig til Forholdet, og i Forholdet til Forholdet; saaledes er under Bestemmelsen Sjel Forholdet mellem Sjel og Legeme et Forhold. Forholder derimod Forholdet sig til sig selv, saa er dette Forhold det positive Tredie, og dette er Selvet.

Et saadant Forhold, der forholder sig til sig selv, et Selv, maa enten have sat sig selv, eller være sat ved et Andet.

Er Forholdet, der forholder sig til sig selv, sat ved et Andet, saa er Forholdet vistnok det Tredie, men dette Forhold, det Tredie, er saa dog igjen et Forhold, forholder sig til hvad der har sat hele Forholdet.

Et saadant deriveret, sat Forhold er Menneskets Selv, et Forhold, der forholder sig til sig selv, og i at forholde sig til sig selv forholder sig til et Andet. Deraf kommer det, at der kan blive to Former for egentlig Fortvivlelse. Havde Menneskets Selv sat sig selv, saa kunde der kun være Tale om een Form, den ikke at ville være sig selv, at ville af med sig selv, men der kunde ikke være Tale om, fortvivlet at ville være sig selv. Denne Formel er nemlig Udtrykket for hele Forholdets (Selvets) Afhængighed, Udtrykket for, at Selvet ikke ved sig selv kan komme til eller være i Ligevægt og Ro, men kun ved i at forholde sig til sig selv at forholde sig til Det, som har sat hele Forholdet. Ja, det er saa langt fra, at denne anden Form af Fortvivlelse (fortvivlet at ville være sig selv) blot betegner en egen Art af Fortvivlelse, at tværtimod al Fortvivlelse tilsidst kan opløses i og tilbageføres paa den. Dersom en Fortvivlende er, som han mener det, opmærksom paa sin Fortvivlelse, ikke taler meningsløst om den, som om Noget, der hænder ham (omtrent som naar Den, der lider af Svimmel, ved et nerveust Bedrag taler om en Tyngde paa Hovedet, eller om at det er som faldt der Noget ned paa ham og saa videre, hvilken Tyngde, og hvilket Tryk dog ikke er noget Udvortes, men en omvendt Reflexion af det Indvortes) – og nu af al Magt ved sig selv og ene ved sig selv vil hæve Fortvivlelsen: saa er han endnu i Fortvivlelsen, arbeider sig med al sin formeentlige Arbeiden kun desto dybere i en dybere Fortvivlelse. Fortvivlelsens Misforhold er ikke et simpelt Misforhold, men et Misforhold i et Forhold, der forholder sig til sig selv, og er sat af et Andet, saa Misforholdet i hiint for sig værende Forhold tillige reflekterer sig uendeligt i Forholdet til den Magt, som satte det.

Dette er nemlig Formelen, som beskriver Selvets Tilstand, naar Fortvivlelsen ganske er udryddet: i at forholde sig til sig selv, og i at ville være sig selv grunder Selvet gjennemsigtigt i den Magt, som satte det.

B

Fortvivlelses Mulighed og Virkelighed

B

Fortvivlelses Mulighed og Virkelighed

Er Fortvivlelse et Fortrin eller en Mangel? Reent dialektisk er den begge Dele. Hvis man vilde fastholde den abstrakte Tanke Fortvivlelse, uden at tænke nogen Fortvivlet, saa maatte man sige: den er et uhyre Fortrin. Muligheden af denne Sygdom er Menneskets Fortrin for Dyret, og dette Fortrin udmærker ham ganske anderledes end den opreiste Gang, thi det tyder paa den uendelige Opreisthed eller Ophøiethed, at han er Aand. Muligheden af denne Sygdom er Menneskets Fortrin for Dyret; at være opmærksom paa denne Sygdom er den Christnes Fortrin for det naturlige Menneske; at være helbredet fra denne Sygdom den Christnes Salighed.

Altsaa det er et uendeligt Fortrin at kunne fortvivle; og dog er det ikke blot den største Ulykke og Elendighed at være det, nei det er Fortabelse. Saaledes er Forholdet mellem Mulighed og Virkelighed ellers ikke; er det et Fortrin at kunne være Det og Det, saa er det et endnu større Fortrin, at være det, det vil sige, det at være forholder sig som en Stigen til det at kunne være. Hvad Fortvivlelse derimod angaaer, saa forholder det at være sig som et Fald til det at kunne være; saa uendeligt som Mulighedens Fortrin er, saa dybt er Faldet. Det der altsaa i Forhold til Fortvivlelse er det Stigende er ikke at være det. Dog er atter denne Bestemmelse tvetydig. Det er ikke med det ikke at være fortvivlet som med det ikke at være halt, blind og deslige. Dersom det ikke at være fortvivlet hverken betyder mere eller mindre end det ikke at være det, saa er det netop at være det. Det ikke at være fortvivlet maa betyde den tilintetgjorte Mulighed af at kunne være det; hvis det skal være sandt, at et Menneske ikke er fortvivlet, maa han i ethvert Øieblik tilintetgjøre Muligheden. Saaledes er Forholdet mellem Mulighed og Virkelighed ellers ikke. Thi vel sige Tænkerne, at Virkelighed er den tilintetgjorte Mulighed, men det er ikke ganske sandt, den er den udfyldte, den virksomme Mulighed. Her derimod er Virkeligheden (ikke at være fortvivlet), som derfor ogsaa er en Negtelse, den afmægtige, tilintetgjorte Mulighed; ellers er Virkelighed i Forhold til Mulighed en Bekræftelse, her en Benegtelse.

Fortvivlelse er Misforholdet i en Syntheses Forhold, som forholder sig til sig selv. Men Synthesen er ikke Misforholdet, den er blot Muligheden, eller, i Synthesen ligger Muligheden af Misforholdet. Var Synthesen Misforholdet, saa var Fortvivlelse slet ikke til, saa vilde Fortvivlelse være Noget, der laae i Menneskenaturen som saadan, det er, saa var det ikke Fortvivlelse; den vilde være Noget, der hændte Mennesket, Noget han leed, som en Sygdom, i hvilken Mennesket falder, eller som Døden, der er Alles Lod. Nei, det at fortvivle ligger i Mennesket selv; men var han ikke Synthese, kunde han slet ikke fortvivle, og var Synthesen ikke oprindeligt fra Guds Haand i det rette Forhold, kunde han heller ikke fortvivle.

Hvorfra kommer saa Fortvivlelsen? Fra Forholdet, hvori Synthesen forholder sig til sig selv, idet Gud, der gjorde Mennesket til Forholdet, ligesom slipper det ud af sin Haand, det er, idet Forholdet forholder sig til sig selv. Og deri, at Forholdet er Aand, er Selvet, deri ligger Ansvaret, under hvilket al Fortvivlelse er og er hvert Øieblik den er, hvor meget, og hvor sindrigt end skuffende sig selv og Andre, den Fortvivlede taler om sin Fortvivlelse som en Ulykke, ved en Forvexling som i hiint anførte Tilfælde af Svimmelhed, med hvilken Fortvivlelse, om end qvalitativ forskjellig, har meget tilfælles, da den under Bestemmelsen Sjel er hvad Fortvivlelse er under Bestemmelsen Aand, og svanger paa Analogier til Fortvivlelse.

Naar saa Misforholdet, Fortvivlelsen, er indtraadt, følger det saa af sig selv, at det vedbliver? Nei, det følger ikke af sig selv; hvis Misforholdet vedbliver, følger det ikke af Misforholdet, men af Forholdet, som forholder sig til sig selv. Dette vil sige, for hver Gang Misforholdet yttrer sig, og i ethvert Øieblik det er til, maa der gaaes tilbage til Forholdet. See, man taler om, at et Menneske paadrager sig en Sygdom, for Exempel ved Uforsigtighed. Saa indtræder da Sygdommen, og fra det Øieblik gjør den sig gjældende og er nu en Virkelighed, hvis Oprindelse bliver mere og mere forbigangen. Det vilde være baade grusomt og umenneskeligt, om man i eet væk vilde vedblive at sige »i dette Øieblik paadrager Du, den Syge, Dig denne Sygdom«, det er, om man i ethvert Øieblik vilde opløse Sygdommens Virkelighed i dens Mulighed. Det er sandt, at han paadrog sig Sygdommen, men det gjorde han kun een Gang, Sygdommens Vedbliven er en simpel Affølge af, at han een Gang paadrog sig den, dens Fremgang ikke i ethvert Øieblik at henføre til ham som Aarsag; han paadrog sig den, men man kan ikke sige, han paadrager sig den. Anderledes med det at fortvivle; ethvert Fortvivlelsens virkelige Øieblik er at tilbageføre til Mulighed, hvert Øieblik han er fortvivlet, paadrager han sig det; det er bestandigt den nærværende Tid, der vorder intet i Forhold til Virkeligheden tilbagelagt Forbigangent; i ethvert Fortvivlelsens virkelige Øieblik bærer den Fortvivlede alt det Foregaaende i Muligheden som et Nærværende. Dette kommer deraf, at at fortvivle er en Bestemmelse af Aand, forholder sig til det Evige i Mennesket. Men det Evige kan han ikke blive af med, nei, i al Evighed ikke; han kan ikke een Gang for alle kaste det fra sig, Intet er umuligere; han maa i ethvert Øieblik, hvor han ikke har det, have kastet eller kaste det fra sig – men det kommer igjen, det er, hvert Øieblik han er fortvivlet, paadrager han sig det at fortvivle. Thi Fortvivlelsen følger ikke af Misforholdet, men af Forholdet, som forholder sig til sig selv. Og Forholdet til sig selv kan et Menneske ikke blive af med, saa lidet som med sit Selv, hvad da forøvrigt er Eet og det Samme, da jo Selvet er Forholdet til sig selv.

C

Fortvivlelse er: »Sygdommen til Døden«

C

Fortvivlelse er: »Sygdommen til Døden«

Dette Begreb Sygdommen til Døden maa dog tages paa en egen Maade. Ligefrem betyder det en Sygdom, hvis Ende, hvis Udgang er Døden. Saaledes taler man om en dødelig Sygdom eenstydigt med en Sygdom til Døden. I den Forstand kan Fortvivlelse ikke kaldes Sygdommen til Døden. Men christeligt forstaaet er Døden selv en Gjennemgang til Livet. Forsaavidt er, christelig, ingen jordisk, legemlig Sygdom til Døden. Thi vistnok er Døden det sidste af Sygdommen, men Døden er ikke det Sidste. Skal der i strængeste Forstand være Tale om en Sygdom til Døden, maa det være en, hvor det Sidste er Døden, og hvor Døden er det Sidste. Og dette er just Fortvivlelse.

Dog er Fortvivlelse i en anden Forstand endnu bestemtere Sygdommen til Døden. Det er nemlig saa langt som muligt fra, at man, ligefrem forstaaet, døer af denne Sygdom, eller at denne Sygdom ender med den legemlige Død. Tværtimod, Fortvivlelsens Qval er netop ikke at kunne døe. Den har saaledes mere tilfælles med den Dødssyges Tilstand, naar han ligger og trækkes med Døden og ikke kan døe. Saaledes er det at være syg til Døden det ikke at kunne døe, dog ikke som var der Haab om Livet, nei Haabløsheden er, at selv det sidste Haab, Døden, ikke er. Naar Døden er den største Fare, haaber man paa Livet; men naar man lærer den end forfærdeligere Fare at kjende, haaber man paa Døden. Naar saa Faren er saa stor, at Døden er blevet Haabet, saa er Fortvivlelse den Haabløshed end ikke at kunne døe.

I denne sidste Betydning er da Fortvivlelse Sygdommen til Døden, denne qvalfulde Modsigelse, denne Sygdom i Selvet, evig at døe, at døe og dog ikke at døe, at døe Døden. Thi at døe betyder at det er forbi, men at døe Døden betyder at opleve det at døe; og lader dette sig eet eneste Øieblik opleve, saa er det dermed for evig at opleve det. Skulde et Menneske døe af Fortvivlelse, som man døer af en Sygdom, saa maatte det Evige i ham, Selvet, kunne døe i samme Forstand som Legemet døer af Sygdommen. Men dette er en Umulighed; Fortvivlelsens Døen omsætter sig bestandigt i en Leven. Den Fortvivlede kan ikke døe; »saa lidet som Dolken kan dræbe Tanker«, saalidet kan Fortvivlelsen fortære det Evige, Selvet, der ligger til Grund for Fortvivlelsen, hvis Orm ikke døer, og hvis Ild ikke udslukkes. Dog er Fortvivlelse just en Selvfortærelse, men en afmægtig Selvfortærelse, der ikke formaaer hvad den selv vil. Men hvad den selv vil er at fortære sig selv, hvilket den ikke formaaer, og denne Afmagt er en ny Form af Selvfortærelse, i hvilken dog Fortvivlelsen atter ikke formaaer hvad den vil, at fortære sig selv, det er en Potentsation eller Loven for Potentsationen. Dette er det Hidsende, eller det er den kolde Brand i Fortvivlelsen, dette Nagende, hvis Bevægelse er bestandigt ind efter, dybere og dybere i afmægtig Selvfortærelse. Det er saa langt fra at det er nogen Trøst for den Fortvivlede, at Fortvivlelsen ikke fortærer ham, det er lige det Modsatte, denne Trøst er just Qvalen, er just hvad der holder Naget ilive og Liv i Naget; thi just derover – ikke fortvivlede – men fortvivler han: at han ikke kan fortære sig selv, ikke kan blive af med sig selv, ikke kan blive til Intet. Dette er den potentserede Formel for Fortvivlelsen, Feberens Stigen i denne Selvets Sygdom.

En Fortvivlende fortvivler over Noget. Saaledes seer det et Øieblik ud, men det er kun et Øieblik; i samme Øieblik viser den sande Fortvivlelse sig eller Fortvivlelsen sig i sin Sandhed. Idet han fortvivlede over Noget, fortvivlede han egentlig over sig selv, og vil nu af med sig selv. Naar saaledes den Herskesyge, hvis Løsen er »enten Cæsar eller slet Intet«, ikke bliver Cæsar, saa fortvivler han derover. Men dette betyder noget Andet: at han, just fordi han ikke blev Cæsar, nu ikke kan udholde at være sig selv. Han fortvivler altsaa egentligen ikke over, at han ikke blev Cæsar, men over sig selv at han ikke blev Cæsar. Dette Selv, der, hvis det var blevet Cæsar, havde været ham al hans Lyst, i en anden Forstand forøvrigt lige saa fortvivlet, dette Selv er ham nu det Utaaleligste af Alt. Det er i dybere Forstand ikke det, der er ham det Utaalelige, at han ikke blev Cæsar, men dette Selv, som ikke blev Cæsar, er ham det Utaalelige, eller endnu rigtigere, det der er ham det Utaalelige, er, at han ikke kan blive af med sig selv. Naar han var blevet Cæsar, saa var han fortvivlet blevet af med sig selv; men nu blev han ikke Cæsar, og kan fortvivlet ikke blive af med sig selv. Væsentligen er han lige fortvivlet, thi han har ikke sit Selv, han er ikke sig selv. Ved at være blevet Cæsar var han dog ikke blevet sig selv, men af med sig selv; og ved ikke at blive Cæsar fortvivler han over ikke at kunne blive af med sig selv. Det er derfor en overfladisk Betragtning (der da formodentligen aldrig har seet en Fortvivlet, end ikke sig selv), naar man om en Fortvivlet siger, som var dette hans Straf: han fortærer sig selv. Thi netop det er det han fortvivlet, og netop det er det, han til sin Qval ikke kan, da der ved Fortvivlelsen er gaaet Ild i Noget, som ikke kan brænde, eller ikke forbrænde, i Selvet.

At fortvivle over Noget er altsaa endnu ikke egentlig Fortvivlelse. Det er Begyndelsen, eller det er, som naar Lægen siger om en Sygdom, den har endnu ikke erklæret sig. Det Næste er den erklærede Fortvivlelse, at fortvivle over sig selv. En ung Pige fortvivler af Elskov, altsaa hun fortvivler over Tabet af den Elskede, at han døde eller at han blev hende utro. Dette er ingen erklæret Fortvivlelse, nei, hun fortvivler over sig selv. Dette hendes Selv, som hun, hvis det var blevet »hans« Elskede, paa den livsaligste Maade var blevet af med, eller havde tabt, dette Selv er hende nu en Plage, naar det skal være et Selv uden »ham«; dette Selv, der var blevet hende, forøvrigt i en anden Forstand ligesaa fortvivlet, hendes Rigdom, er nu blevet hende en modbydelig Tomhed, da »han« er død, eller det er blevet hende en Afsky, da det minder hende om, at hun er bedragen. Forsøg det nu, siig til en saadan Pige: Du fortærer dig selv, og du skal høre hun svarer: »o, nei Qvalen er just, at jeg ikke kan det.«

At fortvivle over sig, fortvivlet at ville af med sig selv, er Formelen for al Fortvivlelse, saa derfor den anden Form for Fortvivlelse, fortvivlet at ville være sig selv, kan føres tilbage til den første, fortvivlet ikke at ville være sig selv, ligesom vi i det Foregaaende opløste den Form fortvivlet ikke at ville være sig selv i den fortvivlet at ville være sig selv (confer A.). En Fortvivlende vil fortvivlet være sig selv. Men naar han fortvivlet vil være sig selv, saa vil han jo ikke af med sig selv. Ja, saa synes det; men naar man seer nærmere til, seer man dog, at Modsigelsen er den samme. Det Selv, han fortvivlet vil være, er et Selv han ikke er (thi at ville være det Selv, han i Sandhed er, er jo lige det Modsatte af Fortvivlelse), han vil nemlig løsrive sit Selv fra den Magt, som satte det. Men dette formaaer han trods al Fortvivlen ikke; trods al Fortvivlelsens Anstrængelse er hiin Magt den stærkere, og tvinger ham til at være det Selv, han ikke vil være. Men saaledes vil han jo dog af med sig selv, af med det Selv, han er, for at være det Selv, han selv har hittet paa. At være Selv, som han vil det, det vilde, om end i en anden Forstand lige saa fortvivlet, være ham al hans Lyst; men at tvinges til at være Selv, som han ikke vil være det, det er hans Qval, den er, at han ikke kan blive af med sig selv.

Socrates beviste Sjelens Udødelighed af, at Sjelens Sygdom (Synden) ikke fortærer denne, som Legemets Sygdom fortærer Legemet. Saaledes kan man ogsaa bevise det Evige i et Menneske deraf, at Fortvivlelsen ikke kan fortære hans Selv, at dette just er Modsigelsens Qval i Fortvivlelsen. Var der intet Evigt i et Menneske, saa kunde han slet ikke fortvivle, men kunde Fortvivlelsen fortære hans Selv, saa var der alligevel ingen Fortvivlelse til.

Saaledes er Fortvivlelse, denne Sygdom i Selvet, Sygdommen til Døden. Den Fortvivlede er dødssyg. Det er i en ganske anden Forstand end det gjælder om nogen Sygdom, de ædleste Dele, Sygdommen har angrebet; og dog kan han ikke døe. Døden er ikke Sygdommens Sidste, men Døden er i eet væk det Sidste. At frelses fra denne Sygdom ved Døden er en Umulighed, thi Sygdommen og dens Qval – og Døden er just ikke at kunne døe.

Denne er Tilstanden i Fortvivlelse. Om det end nok saa meget undgaaer den Fortvivlede, om det end nok saa meget (hvilket da især maa gjælde om den Art Fortvivlelse, der er Uvidenhed om at være Fortvivlelse) lykkes den Fortvivlede ganske at have tabt sit Selv, og tabt det saaledes, at det ikke mærkes i mindste Maade: Evigheden vil saa dog gjøre det aabenbart, at hans Tilstand var Fortvivlelse, og fornagle ham hans Selv, saa Qvalen dog bliver, at han ikke kan blive af med sit Selv, og det bliver aabenbart, at det var en Indbildning med, at det var lykkedes ham. Og saaledes maa Evigheden gjøre, fordi det at have et Selv, at være et Selv, er den største, den uendelige Indrømmelse, der er gjort Mennesket, men tillige Evighedens Fordring paa ham.

B

Denne Sygdoms (Fortvivlelsens) Almindelighed

B

Denne Sygdoms (Fortvivlelsens) Almindelighed

Som vel Lægen maa sige, at der maaskee ikke lever eet eneste Menneske der er ganske sund, saaledes maatte man, hvis man ret kjendte Mennesket, sige, at der ikke lever eet eneste Menneske, uden at han jo dog er lidt fortvivlet, uden at der dog inderst inde boer en Uro, en Ufred, en Disharmoni, en Angest for et ubekjendt Noget, eller for et Noget han end ikke tør stifte Bekjendtskab med, en Angest for en Tilværelsens Mulighed eller en Angest for sig selv, saa han dog, som Lægen taler om at gaae med en Sygdom i Kroppen, gaaer med en Sygdom, gaaer og bærer paa en Aandens Sygdom, der en enkelt Gang glimtviis, ved og med en ham selv uforklarlig Angest, lader sig mærke med, at den er derinde. Og i ethvert Tilfælde har der intet Menneske levet og der lever intet Menneske udenfor Christenheden, uden at han er fortvivlet, og i Christenheden Ingen, forsaavidt han ikke er sand Christen; og forsaavidt han ikke ganske er dette, er han dog noget fortvivlet.

Denne Betragtning vil vistnok synes mange et Paradox, en Overdrivelse, derhos en mørk og forstemmende Anskuelse. Dog er den Intet af alt Dette. Den er ikke mørk, den søger tværtimod at skaffe Lys i hvad man almindeligviis lader staae hen i en vis Dunkelhed; den er ikke forstemmende, tværtimod opløftende, da den betragter ethvert Menneske under Bestemmelsen af den høieste Fordring til ham, at være Aand; den er heller ikke et Paradox, tværtimod en conseqvent gjennemført Grund-Anskuelse, og forsaavidt da heller ingen Overdrivelse.

Den almindelige Betragtning af Fortvivlelse bliver derimod staaende ved Tilsyneladelsen, og er saaledes en overfladisk Betragtning, det er, ingen Betragtning. Den antager, at ethvert Menneske jo bedst maa vide med sig selv, om han er fortvivlet eller ikke. Den der da siger sig selv at være det, han ansees for fortvivlet, men Den der mener sig selv ikke at være det, han ansees heller ikke derfor. Som en Følge heraf bliver Fortvivlelse et sjeldnere Phænomen, istedetfor at den er det ganske Almindelige. Det er ikke det Sjeldne, at En er fortvivlet; nei, det er det Sjeldne, det saare Sjeldne, at En i Sandhed ikke er det.

Men den vulgaire Betragtning forstaaer sig meget daarligt paa Fortvivlelse. Den overseer saaledes blandt Andet ganske (for blot at nævne dette, der dog rigtigt forstaaet bringer Tusinder og Tusinder og Millioner ind under Bestemmelsen Fortvivlelse), den overseer ganske, at det just er en Form af Fortvivlelse det ikke at være det, det ikke at være sig bevidst, at man er det. Det gaaer i en langt dybere Forstand den vulgaire Betragtning i Forhold til at opfatte Fortvivlelse, som det stundom gaaer den i Forhold til at bestemme, om et Menneske er syg eller ikke – i en langt dybere Forstand; thi den vulgaire Betragtning forstaaer sig endnu langt mindre paa, hvad Aand er (og uden det kan man heller ikke forstaae sig paa Fortvivlelse) end paa Sygdom og Sundhed. I Almindelighed antager man, at et Menneske, naar han ikke selv siger, at han er syg, er rask, end sige, naar han selv siger, at han er rask. Lægen derimod betragter Sygdommen anderledes. Og hvorfor? Fordi Lægen har en bestemt og udviklet Forestilling om, hvad det er at være sund, og efter denne prøver han et Menneskes Tilstand. Lægen veed, at som der er en Sygdom, der kun er Indbildning, saaledes ogsaa en Sundhed; han anvender derfor i sidste Tilfælde først Midler, for at faae Sygdommen til at blive aabenbar. Overhovedet har Lægen, just fordi han er Læge (den Indsigtsfulde) ikke ubetinget Tiltro til Menneskets eget Udsagn om sit Befindende. Hvis det var Tilfældet, at hvad ethvert Menneske sagde om sit Befindende, om han er sund eller syg, om hvor han lider o. s. v., var ubetinget til at stole paa, saa var det at være Læge en Indbildning. Thi en Læge har ikke blot at foreskrive Lægemidlerne, men først og fremmest at kjende Sygdommen, og altsaa igjen først og fremmest at kiende, om den formeentligen Syge virkelig er syg, eller om den formeentligen Sunde maaskee er virkelig syg. Saaledes ogsaa med den Sjelekyndiges Forhold til Fortvivlelse. Han veed, hvad Fortvivlelse er, han kjender den og nøies derfor ikke hverken med det Udsagn af et Menneske, at han ikke er fortvivlet, eller, at han er det. Det maa nemlig bemærkes, at i en vis Forstand ere end ikke altid De fortvivlede, der sige sig at være det. Man kan jo affectere Fortvivlelse, og man kan tage feil og forvexle Fortvivlelse, som er en Aandens Bestemmelse, med allehaande transitorisk Forstemthed, Sønderrevethed, der atter gaaer over uden at bringe det til Fortvivlelse. Imidlertid anseer rigtignok den Sjelekyndige ogsaa dette for Former af Fortvivlelse; han seer meget godt at det er Affectation – men just denne Affectation er Fortvivlelse; han seer meget godt, at denne Forstemthed og saa videre ikke har stort at betyde – men just Dette, at den ikke har og ikke faaer stort at betyde, er Fortvivlelse.

Den vulgaire Betragtning overseer fremdeles, at Fortvivlelse, sammenlignet med en Sygdom, er anderledes dialectisk end hvad man ellers kalder Sygdom, fordi det er en Aandens Sygdom. Og dette Dialectiske, rigtigt forstaaet, bringer atter Tusinder ind under Bestemmelsen Fortvivlelse. Dersom nemlig en Læge i et givet Øieblik har forvisset sig om, at Den og Den er sund – og han saa i et senere Øieblik bliver syg: saa kan Lægen have Ret i, at dette Menneske var dengang rask, nu derimod er han syg. Anderledes med Fortvivlelse. Saasnart Fortvivlelse viser sig, saa viser det sig, at Mennesket var fortvivlet. Forsaavidt kan man i intet Øieblik afgjøre Noget om et Menneske, der ikke er frelst ved at have været fortvivlet. Thi naar, dersom Det indtræder, som bringer ham til Fortvivlelse, saa bliver det i samme Øieblik aabenbart, at han hele sit foregaaende Liv igjennem har været fortvivlet. Ingenlunde kan man derimod sige, naar En faaer Feber, at det nu bliver aabenbart, at han har havt Feber hele sit Liv igjennem. Men Fortvivlelse er en Aandens Bestemmelse, forholder sig til det Evige, og har derfor noget af det Evige i sin Dialectik.

Fortvivlelse er ikke blot anderledes dialectisk end en Sygdom, men i Forhold til Fortvivlelse ere alle Kjendetegn dialectiske, og derfor skuffes den overfladiske Betragtning saa let i at bestemme, om Fortvivlelse er tilstede eller ikke. Det ikke at være fortvivlet kan nemlig betyde just at være fortvivlet, og det kan betyde at være frelst fra at være fortvivlet. Tryghed og Beroligelse kan betyde at være fortvivlet, just denne Tryghed, denne Beroligelse kan være Fortvivlelsen; og det kan betyde at have overvundet Fortvivlelse og vundet Fred. Det er ikke med det ikke at være fortvivlet, som med det ikke at være syg; thi det ikke at være syg kan dog ikke være det at være syg, men det ikke at være fortvivlet kan just være det at være fortvivlet. Det er ikke med Fortvivlelse som med en Sygdom, at Ildebefindendet er Sygdommen. Ingenlunde. Ildebefindendet er atter dialectisk. Det aldrig at have fornummet dette Ildebefindende, er just at være fortvivlet.

Dette vil sige, og har sin Grund i, at som Aand betragtet (og skal man tale om Fortvivlelse maa man betragte Mennesket under Bestemmelsen Aand) er Menneskets Tilstand altid critisk. Man taler i Forhold til Sygdom om en Crisis, men ikke i Forhold til Sundhed. Og hvorfor ikke? Fordi legemlig Sundhed er en umiddelbar Bestemmelse, som først bliver dialectisk i Sygdommens Tilstand, hvor der saa bliver Tale om Crisen. Men aandelig, eller naar Mennesket betragtes som Aand, er baade Sundhed og Sygdom critisk; der gives ingen umiddelbar Aandens Sundhed.

Saasnart man ikke betragter Mennesket under Bestemmelsen Aand (og hvis ikke, kan man heller ikke tale om Fortvivlelse), men kun som sjelelig-legemlig Synthese, saa er Sundhed en umiddelbar Bestemmelse, og Sjelens eller Legemets Sygdom først den dialectiske Bestemmelse. Men Fortvivlelse er just, at Mennesket ikke er sig bevidst at være bestemmet som Aand. Selv hvad der menneskelig talt er det Skjønneste og Elskeligste af Alt, en qvindelig Ungdommelighed, der er idel Fred og Harmoni og Glæde: det er dog Fortvivlelse. Dette er nemlig Lykke, men Lykke er ingen Bestemmelse af Aand, og langt, langt inde, inderst inde i Lykkens skjulteste Forborgenhed, der boer ogsaa Angesten, som er Fortvivlelsen; den vil saa gjerne have Lov til at blive derinde, thi det er Fortvivlelsen dens kjereste, dens meest udsøgte Sted at boe: inderst inde i Lykken. Al Umiddelbarhed er trods sin illusoriske Tryghed og Ro Angest, og derfor ganske conseqvent meest angest for Intet; man gjør ikke Umiddelbarheden saa angest ved den meest gysende Beskrivelse af det forfærdeligste Noget, som ved et underfundigt, næsten skjødesløst, men dog med Reflectionens sikkert beregnende Sigte henkastet halvt Ord om et Ubestemt, ja man gjør Umiddelbarheden allermeest angest ved paa en listig Maade at paadutte den, at den selv nok veed, hvad det er man taler om. Thi vistnok veed Umiddelbarheden det ikke; men aldrig fanger Reflexionen saa sikkert som naar den danner sin Snare af Intet, og aldrig er Reflexionen saaledes sig selv, som naar den er – Intet. Der skal en eminent Reflexion, eller rettere der skal en stor Tro til for at kunne udholde Intets, det er, den uendelige Reflexion. Altsaa selv det Skjønneste og Elskeligste af Alt, en qvindelig Ungdommelighed er dog Fortvivlelse, er Lykke. Det lykkes vel derfor heller ikke at slippe Livet igjennem paa denne Umiddelbarhed. Og lykkes det denne Lykke at slippe igjennem, ja, det hjælper kun lidet, thi det er Fortvivlelse. Fortvivlelse er nemlig, just fordi den er heelt dialectisk, den Sygdom, om hvilken det gjælder, at det er den største Ulykke aldrig at have havt den – en sand Guds Lykke at faae den, om det end er den allerfarligste Sygdom, naar man ikke vil helbredes fra den. Ellers kan der dog kun være Tale om, at det er en Lykke at blive helbredet fra en Sygdom, Sygdommen selv er Ulykken.

Det er derfor saa langt som muligt fra, at den vulgaire Betragtning har Ret, som antager, at Fortvivlelse er det Sjeldne, den er tværtimod det ganske Almindelige. Det er saa langt som muligt fra, at den vulgaire Betragtning har Ret, som antager, at Enhver, der ikke mener eller føler sig at være fortvivlet, heller ikke er det, og at kun Den er det, der siger sig at være det. Tværtimod Den, der uden Affectation siger sig at være det, er dog lidt, er et Dialektisk nærmere ved at helbredes end alle De, som ikke ansees for og ikke ansee sig selv for fortvivlede. Men just dette er, hvad vistnok den Sjelekyndige vil give mig Ret i, det Almindelige, at de fleste Mennesker leve uden ret at blive sig bevidste at være bestemte som Aand – og deraf al den saakaldte Tryghed, Tilfredshed med Livet og saa videre, og saa videre, hvad just er Fortvivlelse. De, der derimod sige sig at være fortvivlede, ere ordentligviis enten Dem, der have en saa meget dybere Natur, at de maae blive sig bevidste som Aand, eller Dem, hvem tunge Begivenheder og forfærdelige Afgjørelser have hjulpet hen til at blive sig bevidste som Aand – det er enten Dem eller Dem; thi saare sjelden er vistnok Den, der i Sandhed ikke er fortvivlet.

O, der tales saa meget om menneskelig Nød og Elendighed – jeg søger at forstaae det, har ogsaa kjendt Adskilligt deraf nærved; der tales saa meget om at forspilde et Liv: men kun det Menneskes Liv var forspildt, der levede saaledes hen, bedragen af Livets Glæder eller af dets Sorger, at han aldrig evigt afgjørende blev sig bevidst som Aand, som Selv, eller hvad der er det Samme, aldrig blev opmærksom paa, og i dybeste Forstand fik Indtrykket af, at der er en Gud til, og »han«, han Selv, hans Selv til for denne Gud, hvilken Uendelighedens Vinding aldrig naaes uden gjennem Fortvivlelse. Ak, og denne Elendighed, at der leve saa Mange saaledes hen, bedragne for den saligste af alle Tanker, denne Elendighed, at man beskjæftiger sig, eller i Forhold til Menneskenes Mængde beskjæftiger dem med alt Andet, bruger dem til at afgive Kræfterne i Livets Skuespil, men aldrig minder dem om denne Salighed, at man sammenhober dem – og bedrager dem, istedetfor at splitte dem ad, at hver Enkelt maa vinde det Høieste, det Eneste, der er værd at leve for, og nok for en Evighed at leve i: mig synes, jeg kunde græde en Evighed over, at denne Elendighed er til! O, og dette er i mine Tanker eet Rædselens Udtryk mere for denne af alle forfærdeligste Sygdom og Elendighed: dens Skjulthed, ikke blot, at Den, der lider af den, kan ønske at skjule den og kan formaae det, at den kan boe i et Menneske saaledes, at Ingen, Ingen opdager den, nei, at den kan være saaledes skjult i et Menneske, at han ikke selv veed af det! O, og naar saa engang Timeglasset er udrundet, Timelighedens Timeglas; naar Verdslighedens Larm er forstummet, og den rastløse eller uvirksomme Travlhed fik en Ende; naar Alt er stille om Dig, som det er i Evigheden – hvad enten Du var Mand eller Qvinde, Riig eller Fattig, Afhængig eller Uafhængig, Lykkelig eller Ulykkelig; hvad enten Du bar Kronens Glands i Høihed, eller i ringe Ubemærkethed kun Dagens Møie og Hede; hvad enten Dit Navn skal mindes saa længe Verden staaer, og altsaa er mindet saalænge den stod, eller Du, uden Navn, som Navnløs løb med i den talløse Mængde; hvad enten den Herlighed, som omgav Dig overgik al menneskelig Beskrivelse, eller Dig overgik den strengeste og meest vanærende menneskelige Dom: Evigheden spørger Dig, og hver Enkelt i disse Millioner Millioner, kun om Eet, om Du har levet fortvivlet eller ikke, om saaledes fortvivlet, at Du ikke vidste af at Du var fortvivlet, eller saaledes, at Du skjult bar denne Sygdom i Dit Inderste, som din nagende Hemmelighed, som en syndig Kjerligheds Frugt under dit Hjerte, eller saaledes, at Du, en Rædsel for Andre, rasede i Fortvivlelse. Og hvis saa, hvis Du har levet fortvivlet, hvad Du end forresten vandt eller tabte, saa er for Dig Alt tabt, Evigheden vedkjender sig Dig ikke, den kjendte Dig aldrig, eller end forfærdeligere, den kjender Dig som Du er kjendt, den sætter Dig fast ved Dit Selv i Fortvivlelsen!

[…]

Fortvivlelse er Synden.

Synd er: for Gud, eller med Forestillingen om Gud fortvivlet ikke at ville være sig selv, eller fortvivlet at ville være sig selv. Synd er saaledes den potentserede Svaghed eller den potentserede Trods: Synd er Potentsationen af Fortvivlelse. Det hvorpaa Eftertrykket ligger, er: for Gud, eller at Guds-Forestillingen er med; det, der dialectisk, ethisk, religieust gjør Synd til hvad Juristerne kalde »qvalificeret« Fortvivlelse, er Guds-Forestillingen.

Skjøndt der i dette Afsnit, mindst da i A, ikke er Plads til eller Sted for psychologisk Skildring, maa dog her som det meest dialektiske Confinium mellem Fortvivlelse og Synd anføres, hvad man kunde kalde en Digter-Existents i Retning af det Religieuse, en Existents, der har noget tilfælles med Resignationens Fortvivlelse, kun at Guds-Forestillingen er med. En saadan Digter-Existents vil, hvad der sees af Kategoriernes Conjunktion og Stilling være den eminenteste Digter-Existents. Christelig betragtet er (trods al Æsthetik) enhver Digter-Existents Synd, den Synd: at digte istedetfor at være, at forholde sig til det Gode og Sande gjennem Phantasi istedetfor at være det, det er, existentielt stræbe at være det. Den Digter-Existents, om hvilken vi her tale, er deri forskjellig fra Fortvivlelse, at den har Guds-Forestillingen med sig, eller er for Gud; men den er uhyre dialektisk, og som et uigjennemtrængeligt dialectisk Virvar over, om hvorvidt den dunkelt er sig bevidst at være Synd. En saadan Digter kan have meget dyb religieus Trang, og Forestillingen om Gud er medoptaget i hans Fortvivlelse. Han elsker Gud over Alt, Gud som er ham hans eneste Trøst i hans hemmelige Qval, og dog elsker han Qvalen, den vil han ikke slippe. Han vil saa gjerne være sig selv for Gud, dog ikke i Henseende til det faste Punkt, hvori Selvet lider, der vil han fortvivlet ikke være sig selv; han haaber paa at Evigheden vil tage det bort, og her i Timeligheden kan han, hvormeget han end lider derunder, ikke beslutte sig til at overtage det, ikke troende ydmyge sig derunder. Og dog vedbliver han at forholde sig til Gud, og dette er hans eneste Salighed; det vilde være ham den største Rædsel at skulle undvære Gud, »det vilde være til at fortvivle over;« og dog tillader han sig egentligen, men maaskee ubevidst, at digte Gud en lille Smule anderledes end Gud er, lidt mere a la en kjerlig Fader, som altformeget føier Barnets – eneste Ønske. Som En der blev ulykkelig i Elskov, og derved blev Digter, livsaligt priser Elskovens Lykke: saaledes bliver han Religieusitetens Digter. Han blev ulykkelig i Religieusitet, han forstaaer dunkelt, at det der fordres af ham er at slippe denne Qval, det vil sige troende at ydmyge sig under den og overtage den som hørende med til Selvet – thi han vil holde den ude fra sig og just derved holder han den fast, skjøndt han rigtignok (hvad der, som hvert Ord af en Fortvivlet, er bagvendt rigtigt og altsaa at forstaae omvendt) mener, at dette skal betyde, at skille sig af med den, saa godt som muligt, at slippe den, saa godt som det er et Menneske muligt. Men troende at overtage den, det kan han ikke, det er, det vil han i sidste Forstand ikke, eller, her ender hans Selv i Dunkelhed. Men som hiin Digters Beskrivelse af Elskov, saaledes har denne Digters Beskrivelse af det Religieuse en Fortryllelse, et lyrisk Sving, som ingen Ægtemands og ingen Velærværdigheds. Det er heller ikke usandt hvad han siger, ingenlunde, hans Fremstilling er just hans lykkeligere, hans bedre Jeg. Han er i Forhold til det Religieuse en ulykkelig Elsker, det vil sige han er ikke i streng Forstand en Troende; han har kun det Første af Troen: Fortvivlelsen, og i den en brændende Længsel efter det Religieuse. Hans Kollision er egentligen denne: er han den Kaldede, er Pælen i Kjødet Udtrykket for, at han skal bruges til det Overordentlige, er det for Gud ganske i sin Orden med det Overordentlige, han er blevet? eller er Pælen i Kjødet Det, han skal ydmyge sig under, for at naae det almene Menneskelige? – Dog nok herom, jeg kan med Sandhedens Eftertryk sige: til hvem taler jeg? Saadanne psychologiske Undersøgelser i nte Potents, hvo bryder sig om Sligt; de Nürnberger-Billeder, som Præsten maler, kan man bedre forstaae, de ligne Alle og Enhver, skuffende, Folk som de ere fleest, og, aandeligt forstaaet, Ingenting.

Første Kapitel

Gradationerne i Bevidstheden om Selvet

(Bestemmelsen: for Gud)

Første Kapitel

Gradationerne i Bevidstheden om Selvet

(Bestemmelsen: for Gud)

I det foregaaende Afsnit er der stadigt efterviist en Gradation i Bevidsthed om Selvet; først kom Uvidenhed om at have et evigt Selv (C. B. a.), saa en Viden om at have et Selv, hvori der dog er noget Evigt (C. B. b.), og herindenfor (α. 1. 2. β) blev igjen efterviist Gradationer. Denne hele Betragtning maa nu dialektisk vendes paa en ny Maade. Sagen er denne. Den Gradation i Bevidstheden om Selvet, vi hidtil have beskjæftiget os med, er indenfor Bestemmelsen: det menneskelige Selv, eller det Selv, hvis Maalestok er Mennesket. Men en ny Qvalitet og Qvalification faaer dette Selv derved, at det er Selvet lige over for Gud. Dette Selv er ikke mere det blot menneskelige Selv, men er, hvad jeg, haabende ikke at blive misforstaaet, vilde kalde det theologiske Selv, Selvet lige over for Gud. Og hvilken uendelig Realitet faaer Selvet ikke ved at være sig bevidst at være til for Gud, ved at blive et menneskeligt Selv, hvis Maalestok er Gud! En Røgter, der (hvis dette var muligt) er Selv lige over for Kjøer, er et meget lavt Selv; en Hersker, der er Selv lige over for Trælle, ligeledes, og egentligen intet Selv – thi i begge Tilfælde mangler der Maalestok. Barnet, der hidtil blot har havt Forældrenes Maalestok, bliver Selv ved som Mand at faae Staten til Maalestok; men hvilken uendelig Accent falder der paa Selvet ved at faae Gud til Maalestok! Maalestokken for Selvet er altid: hvad Det er, lige over for hvilket det er Selv, men Dette er jo igjen Definitionen paa hvad »Maalestok« er. Som man kun kan addere eensartede Størrelser, saaledes er enhver Ting qvalitativt Det, hvormed den maales; og Det, der qvalitativt er dens Maalestok, er ethisk dens Maal; og Maalestokken og Maalet er qvalitativt Det, Noget er, med Undtagelse da af Forholdet i Frihedens Verden, hvor imidlertid En ved ikke qvalitativt at være Det, som er hans Maal og hans Maalestok, selv maa have forskyldt denne Disqvalification, saa Maalet og Maalestokken dog – dømmende bliver den samme, gjør det aabenbart, hvad det er han ikke er, Det nemlig, som er hans Maal og Maalestok.

Det var en meget rigtig Tanke, til hvilken saa ofte en ældre Dogmatik recurrerede, medens en senere Dogmatik saa ofte opholdt sig over den, fordi den manglede Forstand derpaa og Sands derfor – det var en meget rigtig Tanke, om der end stundom gjordes forkeert Anvendelse af den: at Det, der gjorde Synden til det Frygtelige var, at den var for Gud. Deraf beviste man saa Helvedesstraffenes Evighed. Senere blev man kløgtig, og sagde: Synd er Synd; Synden er ikke større, fordi den er mod Gud eller for Gud. Besynderligt! selv Juristerne tale om qvalificerede Forbrydelser, selv Juristerne gjøre Forskjel, om en Forbrydelse for Exempel er udøvet mod en Embedsmand, eller mod en Privatmand, gjøre Forskjel i Straf for et Fadermord og et almindeligt Mord.

Nei, deri havde den ældre Dogmatik Ret, at det at Synden var mod Gud, uendelig potentserede den. Feilen laae i, at man betragtede Gud som noget Udvortes, og i, at det var, som om man antog, at der kun stundom syndedes mod Gud. Men Gud er ikke noget Udvortes i den Forstand som en Politibetjent. Det der maa sees paa, er, at Selvet har Forestillingen om Gud, og at det saa dog ikke vil som han vil, saa dog er ulydigt. Ei heller er det kun stundom der syndes for Gud; thi enhver Synd er for Gud, eller rettere, Det, der egentligen gjør den menneskelige Skyld til Synd, er, at den Skyldige havde Bevidstheden om at være til for Gud.

Fortvivlelsen potentseres i Forhold til Bevidsthed om Selvet; men Selvet potentseres i Forhold til Maalestokken for Selvet, og uendeligt naar Gud er Maalestokken. Jo mere Guds-Forestilling, jo mere Selv; jo mere Selv, jo mere Guds-Forestilling. Først naar et Selv, som dette bestemte enkelte, er sig bevidst at være til for Gud, først da er det det uendelige Selv; og dette Selv synder saa for Gud. Hedenskabets Selviskhed er derfor, trods Alt hvad der kan siges om den, dog ikke nær saa qvalificeret som Christenhedens, forsaavidt ogsaa her er Selviskhed; thi Hedningen havde ikke sit Selv lige over for Gud. Hedningen og det naturlige Menneske har det blot menneskelige Selv som Maalestok. Man kan derfor vel have Ret i fra et høiere Synspunkt at see Hedenskabet liggende i Synd, men Hedenskabets Synd var egentlig den fortvivlede Uvidenhed om Gud, om at være til for Gud; den er: »at være uden Gud i Verden.« Fra en anden Side er det derfor sandt, at Hedningen ikke i strengeste Forstand syndede, thi han syndede ikke for Gud; og al Synd er for Gud. Det er fremdeles ogsaa i een Forstand ganske vist, at mange Gange er en Hedning hjulpet til at slippe saadan upaaklageligt gjennem Verden, just fordi hans pelagiansk-letsindige Forestilling frelste ham; men saa er hans Synd en anden: denne pelagiansk-letsindige Opfattelse. Derimod er det paa den anden Side ogsaa ganske vist, at mangen Gang er et Menneske, just ved at være strengt opdraget i Christendom, i en vis Forstand styrtet i Synd, fordi den hele christelige Anskuelse var ham for alvorlig, især i en tidligere Tid af hans Liv; men saa er dog i en anden Forstand dette igjen noget til Hjelp for ham, denne dybere Forestilling om, hvad Synd er.

Synd er: for Gud fortvivlet ikke at ville være sig selv, eller for Gud fortvivlet at ville være sig selv. Men er denne Definition, om den end i andre Henseender dog maaskee indrømmes at have sine Fortrin (og deriblandt det vigtigste af alle at være den eneste skriftmæssige; thi Skriften definerer altid Synd som Ulydighed), er den ikke for aandelig? Dertil maa først og fremmest svares: en Definition af Synd kan aldrig være for aandelig (naar den da ikke bliver saa aandelig, at den afskaffer Synd); thi Synd er just en Bestemmelse af Aand. Og dernæst: hvorfor skulde da den være for aandelig? Fordi den ikke taler om Mord, Tyveri, Utugt o. D.? Men taler den da ikke derom? Er det ikke ogsaa en Selvraadighed mod Gud, en Ulydighed, der trodser hans Bud? Men derimod, naar man i at tale om Synd kun taler om saadanne Synder, glemmes saa let, at alt Sligt saadan til en vis Grad, menneskelig talt, kan være i sin Orden, og dog det hele Liv Synd, den bekjendte Art Synd: de glimrende Laster, Selvraadighed, der enten aandløst eller frækt forbliver eller vil være uvidende om, i hvilken uendelig langt dybere Forstand et menneskeligt Selv er Gud forpligtet i Lydighed betræffende hver dets lønligste Ønske og Tanke, betræffende Lydhørethed til at opfatte og Villighed til at følge hvert det mindste Guds Vink om, hvad hans Villie er med dette Selv. Kjødets Synder er det lavere Selvs Selvraadighed; men hvor ofte uddrives ikke een Djævel ved Djævelens Hjælp, og det Sidste bliver værre end det Første. Thi saaledes gaaer det jo just til i Verden: først synder et Menneske af Skrøbelighed og Svaghed; og saa – ja saa maaskee lærer han at flye til Gud, og hjælpes til Troen, der frelser fra al Synd; men derom tale vi nu ikke her – saa fortvivler han over sin Svaghed, og bliver enten en Pharisæer, som fortvivlet driver det til en vis legal Retfærdighed, eller han fortvivlet styrter sig i Synden igjen.

Definitionen omfatter derfor vistnok enhver tænkelig og enhver virkelig Form af Synd; den udhæver dog vistnok rigtigt det Afgjørende, at Synd er Fortvivlelse (thi Synd er ikke Kjød og Blods Vildhed, men er Aandens Samtykke deri), og er: for Gud. Som Definition er den Bogstavregning; det vilde i dette lille Skrift være paa urette Sted og desuden et Forsøg, som maatte mislykkes, hvis jeg vilde begynde paa at beskrive de enkelte Synder. Hovedsagen er her blot, at Definitionen som et Net omfatter alle Former. Og det gjør den, hvilket ogsaa kan sees, naar man gjør Prøve paa den ved at stille Modsætningen, Definitionen paa: Tro, efter hvilken jeg i hele dette Skrift styrer som efter det sikkre Sømærke. Tro er: at Selvet i at være sig selv og i at ville være sig selv gjennemsigtigt grunder i Gud.

Men dette er ofte nok blevet overseet, at Modsætningen til Synd ingenlunde er Dyd. Dette er en tildeels hedensk Betragtning, som nøies med en blot menneskelig Maalestok, som netop ikke veed, hvad Synd er, at al Synd er for Gud. Nei, Modsætningen til Synd er Tro, som det derfor hedder Rom. 14, 23: Alt hvad der ikke er af Tro er Synd. Og dette er en af de for hele Christendommen meest afgjørende Bestemmelser, at Modsætningen til Synd ikke er Dyd, men Tro.

Tillæg

At Definitionen paa Synd har Forargelsens Mulighed i sig; en almindelig Bemærkning om Forargelse.

Tillæg

At Definitionen paa Synd har Forargelsens Mulighed i sig; en almindelig Bemærkning om Forargelse.

Modsætningen: Synd – Tro er den christelige, der, christelig, omdanner alle ethiske Begrebsbestemmelser, giver dem eet Udtræk mere. Til Grund for Modsætningen ligger det afgjørende Christelige: for Gud, hvilken Bestemmelse igjen har det Christeliges afgjørende Criterium: det Absurde, Paradoxet, Forargelsens Mulighed. Og at dette paavises ved enhver Bestemmelse af det Christelige, er af yderste Vigtighed, da Forargelsen er det Christeliges Værge mod al Speculation. Hvor er da her Forargelsens Mulighed? Den er, at et Menneske skulde have den Realitet: som enkelt Menneske at være til lige over for Gud, og altsaa igjen, hvad deraf følger, at Menneskets Synd skulde beskjæftige Gud. Dette om det enkelte Menneske – for Gud faaer Speculationen aldrig i sit Hoved; den universaliserer blot de enkelte Mennesker phantastisk i Slægten. Just derfor var det ogsaa, at en vantro Christendom udfandt, at Synd er Synd, om den er lige over for Gud eller ikke, gjør jo hverken fra eller til. Det vil sige, man vilde have Bestemmelsen: for Gud bort, og dertil opfandt man en høiere Viisdom, der dog, besynderligt nok, hverken var mere eller mindre end, hvad den høiere Viisdom vistnok oftest er, det gamle Hedenskab.

Der er nu saa ofte talet om, at man forargedes paa Christendommen, fordi den er saa mørk og dunkel, forargedes paa den, fordi den er saa streng og saa videre; det bliver dog nok rigtigst eengang at oplyse, at det hvorfor Mennesket egentligen forarges paa Christendommen er, fordi den er for høi, fordi dens Maal ikke er Menneskets Maal, fordi den vil gjøre Mennesket til noget saa Overordentligt, at han ikke kan faae det i sit Hoved. Dette vil ogsaa en ganske simpel psychologisk Udvikling af, hvad Forargelse er, opklare, og tillige vise, hvor uendelig taabelig man har baaret sig ad med at forsvare Christendommen, saa man tog Forargelsen bort; hvor dumt eller frækt man har ignoreret Christi egen Anviisning, der ofte og saa bekymret advarer mod Forargelsen, det er, selv paapeger, at dens Mulighed er der, og skal være der; thi skal den ikke være der, er den ikke et evig væsentlig Medhenhørende til det Christelige, saa er det jo menneskelig Nonsens af Christus: istedetfor at tage den bort, at gaae bekymret og advare mod den.

Dersom jeg tænkte mig en fattig Dagleier og den mægtigste Keiser, der nogensinde har levet, og denne den mægtigste Keiser pludselig fik det Indfald at sende Bud efter Dagleieren, der da aldrig havde drømt om, »og i hvis Hjerte det aldrig var opkommet,« at Keiseren vidste, han var til, der da vilde prise sig ubeskrivelig lykkelig blot ved engang at faae Lov til at see Keiseren, hvad han i saa Fald vilde fortælle Børn og Børnebørn som den vigtigste Begivenhed i sit Liv – dersom Keiseren sendte Bud efter ham og lod ham vide, at han vilde have ham til Svigersøn: hvad saa? Saa vilde Dagleieren paa menneskelig Viis blive noget eller meget forlegen, flau og generet derved, det vilde (og dette er det Menneskelige), menneskeligt, forekomme ham som noget høist besynderligt, afsindigt Noget, hvorom han da mindst af Alt turde tale til noget andet Menneske, da han selv allerede i sit stille Sind ikke var langt fra den Forklaring, som Nabo og Gjenbo snarest muligt vilde faae yderst travlt med: at Keiseren vilde gjøre Nar af ham, saa Dagleieren vilde blive til Latter for hele Byen, han aftegnet i Bladet, Historien om hans Formæling med Keiserens Datter solgt af Visekjerlingerne. Dog det med at blive Keiserens Svigersøn, det maatte da vel snart vorde en udvortes Virkelighed, saa Dagleieren paa sandselig Maade kunde forvisse sig om, hvorvidt det var Keiserens Alvor dermed, eller han blot vilde have det stakkels Menneske til Bedste, gjøre ham ulykkelig for hele hans Levetid, og hjælpe ham til at ende i en Daarekiste; thi det quid nimis er tilstede, som saa uendeligt let kan slaae om til sit Modsatte. En lille Gunstbeviisning, det vilde Dagleieren kunde faae i sit Hoved; det vilde blive forstaaet i Kjøbstaden, af det høistærede dannede Publicum, af alle Visekjerlinger, kort af de 5 Gange 100,000 Mennesker, som boede i hiin Kjøbstad, der rigtignok i Henseende til Folkemængde var endog en meget stor Stad, derimod i Henseende til at have Forstand paa og Sands for det Overordentlige, en meget lille Kjøbstad – men dette, med at blive Svigersøn, ja, det var meget for meget. Og sæt nu, at det ikke var en udvortes Virkelighed, der var Tale om, men om en indvortes, saa altsaa Facticiteten ikke kunde hjælpe Dagleieren til Vished, men Troen selv var den eneste Facticitet, og Alt altsaa overladt til Troen, om han havde ydmygt Mod nok til at turde troe det (thi frækt Mod kan ikke hjælpe til at troe): hvor mange Dagleiere var der saa vel, som havde dette Mod? Men Den, der ikke havde dette Mod, han vilde forarges; det Overordentlige vilde for ham lyde næsten som en Spot over ham. Han vilde saa maaskee ærligt og ligefrem tilstaae: Sligt er mig for høit, jeg kan ikke faae det i mit Hoved, det er, at jeg skal sige det lige ud, mig en Daarskab.

Og nu Christendommen! Christendommen lærer, at dette enkelte Menneske, og saaledes ethvert enkelt Menneske, hvad han saa forresten er, Mand, Kone, Tjenestepige, Minister, Kjøbmand, Barbeer, Studiosus og saa videre, dette enkelte Menneske er til for Gud – dette enkelte Menneske, som maaskee vilde være stolt af een Gang i sit Liv at have talt med Kongen, dette Menneske, som ikke bilder sig lidt ind af, at leve paa en fortrolig Fod med Den og Den, dette Menneske er til for Gud, kan tale med Gud, hvad Øieblik han vil, sikker paa at blive hørt af ham, kort, dette Menneske tilbydes det at leve paa den fortroligste Fod med Gud! Videre, for dette Menneskes, ogsaa for dette Menneskes Skyld kommer Gud til Verden, lader sig føde, lider, døer; og denne lidende Gud, han næsten beder og bønfalder dette Menneske om dog at tage mod Hjælpen, som tilbydes ham! Sandeligen er Noget til at tabe Forstanden over, saa dog vel dette! Hver Den, der ikke har ydmygt Mod til at turde troe det, han forarges. Men hvorfor forarges han? fordi det er ham for høit, fordi han ikke kan faae det i sit Hoved, fordi han ikke kan faae sin Frimodighed lige over for det, og derfor maa have det bort, gjort til Intet, til Galskab og Nonsens, thi det er som skulde det qvæle ham.

Thi hvad er Forargelse? Forargelse er ulykkelig Beundring. Den er derfor i Slægt med Misundelse, men den er en Misundelse, som vender sig mod En selv, i endnu strengere Forstand, værst mod sig selv. Det naturlige Menneskes Snæverhjertethed kan ikke unde sig selv det Overordentlige, som Gud har tiltænkt ham; saa forarges han.

Forargelsens Grad beroer nu paa, hvad Lidenskab i Forhold til Beundring et Menneske har. Mere prosaiske Mennesker uden Phantasi og Lidenskab, som altsaa heller ikke ere ret skikkede til at beundre, de forarges vel ogsaa, men de indskrænke sig til at sige: Sligt kan jeg ikke faae i mit Hoved, jeg lader det staae. Det er Skeptikerne. Men jo mere Lidenskab og Phantasi et Menneske har, jo nærmere han altsaa i en vis Forstand, i Muligheden nemlig, er ved at kunne blive troende, NB. ved tilbedende at ydmyge sig under det Overordentlige, desto mere lidenskabelig Forargelse, der tilsidst ikke kan nøies med Mindre end at faae dette udryddet, tilintetgjort, traadt i Snavset.

Vil man lære at forstaae sig paa Forargelse, saa studere man den menneskelige Misundelse, et Studium, som jeg giver op uden om, og indbilder mig at have gjort grundigt. Misundelse er skjult Beundring. En Beundrende, der føler, at han ikke kan blive lykkelig ved at give sig hen, han vælger at blive misundelig paa Det, han beundrer. Saa taler han et andet Sprog; i hans Sprog hedder det nu, at Det, som han egentlig beundrer, er Ingenting, noget dumt og flaut og sært og overspændt Noget. Beundring er lykkelig Selvfortabelse, Misundelse ulykkelig Selvhævdelse.

Saaledes ogsaa med Forargelse; thi hvad der i Forholdet mellem Menneske og Menneske er: Beundring – Misundelse, det er i Forholdet mellem Gud og Menneske: Tilbedelse – Forargelse. Summa summarum af al menneskelig Viisdom er dette »gyldne«, eller maaskee rigtigere dette pletterede: ne quid nimis, for lidt og for meget fordærver Alt. Dette tages og gives mellem Mand og Mand som Viisdom, honoreres med Beundring; dets Cours fluctuerer aldrig, hele Menneskeheden garanterer dets Værd. Saa lever der engang imellem et Geni, som gaaer en lille Smule ud derover, han erklæres gal – af de Kloge. Men Christendommen gjør et uhyre Kjæmpeskridt ud over dette: ne quid nimis, ind i det Absurde; der begynder Christendommen – og Forargelsen.

Man seer nu, hvor (for at der dog kunde blive noget Overordentligt tilbage) hvor overordentlig dumt det er at forsvare Christendommen, hvor liden Menneskekundskab det forraader, hvorledes det, om end ubevidst, spiller under Dække med Forargelsen, ved at gjøre det Christelige til noget saadant kummerligt Noget, der am Ende skulde reddes ved et Forsvar. Derfor er det vist og sandt, at den, der først opfandt, i Christenheden at forsvare Christendommen, de facto er en Judas No. 2; ogsaa han forraader med et Kys, kun at hans Forræderi er Dumhedens. At forsvare Noget er altid at disrecommandere det. Lad En have et Pakhuus fuldt af Guld, lad ham være villig til at udgive hver en Dukat til de Fattige – men lad ham derhos være dum nok til at begynde dette sit velgjørende Foretagende med et Forsvar, hvori han af tre Grunde beviser det Forsvarlige deri: og det skal ikke være langt fra, at Menneskene næsten finde det tvivlsomt, om han gjør noget Godt. Men nu det Christelige! Ja Den, der forsvarer det, han har aldrig troet paa det. Troer han, saa er Troens Begeistring – ikke et Forsvar, nei, den er Angrebet, og Seieren; en Troende er en Seierherre.

Saaledes med det Christelige og Forargelsen. Dens Mulighed er da ganske rigtigt tilstede ved den christelige Definition paa Synden. Den er dette: for Gud. En Hedning, det naturlige Menneske er meget villig til at indrømme, at der er Synd til, men dette »for Gud«, som dog egentlig gjør Synden til Synd, det er ham meget for meget. Det er (om end paa en anden Maade end den her paaviste) ham at gjøre Meget for Meget af det at være Menneske; lidt mindre, saa er han villig til at gaae ind derpaa – »men for Meget er for Meget«.

Andet Capitel

Den socratiske Definition af Synd

Andet Capitel

Den socratiske Definition af Synd

Synd er Uvidenhed. Dette er, som bekjendt, den socratiske Definition, der som alt Socratisk altid er en Instants, der er værd at agte paa. Imidlertid er det gaaet i Forhold til dette Socratiske som i Forhold til saa meget andet Socratisk, man har lært at føle en Trang til at gaae videre. Hvor utallig Mange have ikke følt Trang til at gaae videre end den socratiske Uvidenhed – formodentligen fordi de følte, det var dem en Umulighed at blive staaende ved den; thi hvor mange ere vel i hver Generation De, som vare istand til endog blot een Maaned at udholde, existentielt at udtrykke Uvidenhed om Alt.

Den socratiske Definition vil jeg derfor ingenlunde affærdige med, at man ikke kan blive staaende ved den; men med det Christelige in mente vil jeg benytte den til at faae dette frem i sin Skarphed – just fordi den socratiske Definition er saa ægte græsk; saa her som altid enhver anden Definition, der ikke i strengeste Forstand er streng christelig, det er, enhver Mellem-Definition viser sig i sin Tomhed.

Misligheden ved den socratiske Definition er nu, at den lader det ubestemt, om hvorledes Uvidenheden selv nærmere er at forstaae, dens Oprindelse og saa videre Dette vil sige, selv om Synd er Uvidenhed (eller hvad Christendommen maaskee snarere vilde kalde Dumhed), hvilket i een Forstand slet ikke lader sig negte: er dette en oprindelig Uvidenhed, er Tilstanden altsaa Dens, der ikke har vidst, og hidtil ikke har kunnet vide Noget om Sandheden; eller er det en frembragt, en senere Uvidenhed? Er den det Sidste, saa maa jo Synden egentligen stikke i noget Andet end i Uvidenheden, den maa stikke i den Virksomhed i Mennesket, ved hvilken han har arbeidet paa at fordunkle sin Erkjendelse. Men ogsaa dette antaget, kommer denne haardnakkede og meget seiglivede Mislighed igjen, idet Spørgsmaalet bliver, om Mennesket da i det Øieblik, han begyndte paa at fordunkle sin Erkjendelse, er sig dette tydelig bevidst. Er han sig det ikke tydeligt bevidst, saa er jo Erkjendelsen allerede noget fordunklet, førend han begyndte derpaa; og Spørgsmaalet kommer blot atter igjen. Antages derimod, at han, da han begyndte at fordunkle Erkjendelsen, er sig dette tydeligt bevidst, saa ligger jo Synden (om den end er Uvidenhed, forsaavidt denne er Resultatet) ikke i Erkjendelsen men i Villien, og det Spørgsmaal, der maatte fremkomme, er om Erkjendelsens og Villiens Forhold til hinanden. Alt Sligt (og her kunde man da vedblive at spørge i mange Dage) indlader den socratiske Definition sig egentlig ikke paa. Socrates var vistnok en Ethiker (hvad Oldtiden jo ubetinget vindicerer ham, Ethikens Opfinder), den første, som han er og bliver den Første i sit Slags; men han begynder med Uvidenheden. Intellectuelt er det Uvidenheden han tenderer til, det Intet at vide. Ethisk forstaaer han ved Uvidenhed noget ganske Andet, og begynder saa med den. Men derimod er, som naturligt, Socrates ingen væsentlig religieus Ethiker, endnu mindre, hvad der er det Christelige, en Dogmatiker. Derfor kommer han egentligen slet ikke ind i den hele Undersøgelse, med hvilken Christendommen begynder, ikke ind i det Prius, i hvilket Synden forudsætter sig selv, og som christeligt forklares i Dogmet om Arvesynden, til hvilket Dogme vi i denne Undersøgelse dog blot komme som til Grændsen.

Socrates kommer derfor egentligen ikke til Bestemmelsen Synd, hvilket vistnok er en Mislighed ved en Definition paa Synd. Hvorledes dette? Er nemlig Synd Uvidenhed, saa er jo Synd egentlig ikke til; thi Synd er jo netop Bevidsthed. Er Synden det at være uvidende om det Rigtige, saa man derfor gjør det Urigtige, saa er Synd ikke til. Er dette Synd, saa antages jo, hvad Socrates ogsaa antog, at det Tilfælde ikke forekommer, at En, vidende det Rette, gjør det Urette, eller vidende, at det er det Urette, gjør dette Urette. Altsaa, er den socratiske Definition rigtig, saa er Synd slet ikke til. Men see, dette, just dette er christelig ganske i sin Orden, i en dybere Forstand ganske rigtigt, i christelig Interesse quod erat demonstrandum. Netop det Begreb, ved hvilket Christendommen meest afgjørende qvalitativt differentierer sig fra Hedenskabet, er: Synden, Læren om Synden; og derfor antager Christendommen ogsaa ganske conseqvent, at hverken Hedenskabet eller det naturlige Menneske veed hvad Synd er, ja, den antager, at der maa en Aabenbaring fra Gud til for at gjøre det aabenbart, hvad Synd er. Ikke er det nemlig saa, som en overfladisk Betragtning antager, at Læren om Forsoningen er den qvalitative Differents mellem Hedenskab og Christendom. Nei, Begyndelsen maa gjøres langt dybere, med Synden, med Læren om Synden, som Christendommen ogsaa gjør. Hvilken farlig Indvending derfor mod Christendommen, hvis Hedenskabet havde en Definition paa Synd, som Christendommen maatte erkjende for rigtig.

Hvad er det da for en Bestemmelse, Socrates i at bestemme Synden mangler? Det er: Villien, Trods. Den græske Intellectualitet var for lykkelig, for naiv, for æsthetisk, for ironisk, for vittig – for syndig til at kunne faae det i sit Hoved, at En med sit Vidende kunde lade være at gjøre det Gode, eller med sit Vidende, med Viden om det Rette, gjøre det Urette. Græciteten statuerer et intellectuelt kategorisk Imperativ.

Det Sande heri bør nu ingenlunde oversees, og gjøres vel fornødent at indskjærpes i Tider som disse, der ere løbne vild i saare megen tomt opblæsende og ufrugtbar Viden, saa det vistnok nu, aldeles som paa Socrates's Tid, kun endnu mere, gjordes Behov, at Menneskene udhungredes en lille Smule socratisk. Det er baade til at lee og til at græde over, saavel alle disse Forsikkringer om at have forstaaet og begrebet det Høieste, som ogsaa den Virtuositet, med hvilken Mange in abstracto vide at fremstille det, i en vis Forstand ganske rigtigt – det er baade til at lee og til at græde over, naar man saa seer, at al denne Viden og Forstaaen slet ingen Magt udøver over Menneskenes Liv, at dette ikke i fjerneste Maade udtrykker hvad de have forstaaet, men snarere lige det Modsatte. Man udbryder uvilkaarligt ved Synet af dette lige saa sørgelige som latterlige Misforhold: men hvor i al Verden er det dog muligt at de kunne have forstaaet det, er det ogsaa sandt at de have forstaaet det? Her svarer hiin gamle Ironiker og Ethiker: o, Kjere, tro aldrig det; de have ikke forstaaet det, thi havde de i Sandhed forstaaet det, saa udtrykte deres Liv det ogsaa, saa gjorde de hvad de havde forstaaet.

At forstaae og forstaae er altsaa to Ting? Ganske vist; og den som har forstaaet dette – dog vel at mærke ikke i Betydning af den første Art Forstaaen, han er eo ipso indviet i alle Ironiens Hemmeligheder. Det er egentligen denne Modsigelse, der beskjæftiger Ironien. At opfatte det comisk, at et Menneske virkelig er uvidende om Noget, er en meget lav Art Comik, og under Ironiens Værdighed. Noget dybere Comisk er der egentlig ikke i, at der har levet Mennesker, som antoge at Jorden staaer stille – naar de ikke vidste bedre. Det vil formodentlig paa en lignende Maade igjen gaae vor Tid i Forhold til en mere physikalsk vidende Tid. Modsigelsen er mellem to forskjellige Tider, der fattes et dybere Coincidentspunkt, en saadan Modsigelse er ikke væsentlig og altsaa heller ei væsentlig comisk. Nei, men at et Menneske staaer og siger det Rigtige – og altsaa har forstaaet det; og naar han saa skal handle, gjør det Urigtige – og altsaa viser, at han ikke har forstaaet det: ja det er uendelig comisk. Det er uendelig comisk, at et Menneske rørt indtil Taarer, saa altsaa ikke blot Sveden men Taarerne hagle ned af ham, kan sidde og læse eller høre Fremstillingen af Selvfornegtelse, af det Ædle i at offre sit Liv for Sandheden – og saa i næste Øieblik, ein, zwei, drei, vupti, næsten endnu med Taarerne i Øinene, er i fuld Gang med, i sit Ansigts Sveed, efter fattig Leilighed, at hjælpe Usandheden til at seire. Det er uendelig comisk, at en Taler med Sandhed i Stemme og Mimik, dybt grebet og dybt gribende, rystende kan fremstille det Sande, træde alt det Onde, alle Helvedes-Magter under Øinene, med en Aplomb i Skikkelsen, med en Freidighed i Blikket, en Rigtighed i Pas'ene som er beundringsværdig – det er uendelig comisk, at han, næsten i samme Øieblik, næsten endnu med »Adriennen paa«, kan feigt og frygtagtigt løbe af Veien for den mindste Ulempe. Det er uendelig comisk, at En kan forstaae hele Sandheden, om hvor ussel og smaalig Verden er og saa videre – at han kan forstaae det, og saa ikke kan kjende Det igjen, han har forstaaet; thi næsten i samme Øieblik gaaer han selv hen og er med i den samme Smaalighed og Usselhed, tager Ære af den og æres af den, det er, anerkjender den. O, naar man seer En, der forsikkrer aldeles at have forstaaet, hvorledes Christus gik omkring i en ringe Tjeners Skikkelse, fattig, foragtet, bespottet, som Skriften siger: bespyttet – naar jeg da seer den Samme saa omhyggeligt at flye hen, hvor der verdsligt er godt at være, indrette sig der paa det tryggeste, naar jeg seer ham saa ængstelig, som var det Livet om at gjøre, at undflye ethvert en ugunstig Vinds Pust fra Høire eller Venstre, saa lyksalig, saa høistlyksalig, saa kisteglad – ja, for at gjøre det complet, saa kisteglad, at han endog rørt takker Gud derfor – ved ubetinget at være æret og anseet af Alle, Alle: da har jeg ofte sagt til mig selv og ved mig selv »Socrates, Socrates, Socrates, skulde det være muligt, at dette Menneske har forstaaet hvad han siger sig at have forstaaet?«. Saaledes har jeg sagt, ja jeg har tillige ønsket, at Socrates havde Ret. Thi det er mig dog, som var Christendommen for streng, jeg kan ei heller bringe det i Samklang med min Erfaring, at gjøre en Saadan til en Hykler. Nei, Sokrates, Dig kan jeg forstaae; Du gjør ham til en Spasmager, en Slags lystig Broder, Du gjør ham til en Prise for Latteren, Du har Intet mod, det har endog Dit Bifald, at jeg serverer og anretter ham comisk – det vil da sige, naar jeg gjør det godt.

Socrates, Socrates, Socrates! Ja, man maa nok nævne Dit Navn tre Gange, det var ikke for meget at nævne det ti Gange, dersom det kunde hjælpe Noget. Man mener, Verden behøver en Republik, og man mener at behøve en ny Samfunds-Orden, og en ny Religion: men Ingen tænker paa, at en Socrates er det dog nok denne, just ved megen Viden forvirrede, Verden behøver. Dog det forstaaer sig, tænkte Nogen derpaa, end sige, hvis Mange tænkte derpaa, da behøvedes han mindre. Det en Vildfarelse meest behøver, er altid Det, den mindst tænker paa – naturligt, thi ellers var det jo ikke en Vildfarelse.

Altsaa en saadan ironisk-ethisk Correction kunde vor Tid saare godt behøve, er maaskee egentligen det Eneste, den har Behov – thi det er aabenbart Det den mindst tænker paa; det gjordes høiligt fornødent, at vi, istedetfor at gaae videre end Socrates, blot kom tilbage til dette Socratiske: at forstaae og forstaae er to Ting, – ikke som et Resultat, der tilsidst hjælper Menneskene i den dybeste Elendighed, da det just hæver Forskjellen mellem at forstaae og forstaae, men som den ethiske Opfattelse af Livets Daglighed.

Den socratiske Definition hjælper sig da saaledes. Naar En ikke gjør det Rette, saa har han heller ikke forstaaet det; hans Forstaaen er en Indbildning; hans Forsikkring om at have forstaaet en feil Direction; hans gjentagne Forsikkring om Fanden gale ham at have forstaaet, en uhyre, uhyre Fjernhed ad den størst mulige Omvei. Men saa er jo Definitionen rigtig. Gjør En det Rigtige, saa synder han dog vel ikke; og gjør han ikke det Rigtige, saa har han heller ikke forstaaet det; havde han i Sandhed forstaaet det, vilde det snart bevæge ham til at gjøre det, snart gjøre ham til en Klangfigur for hans Forstaaen: ergo er Synd Uvidenhed.

Men hvor stikker da Misligheden? Den stikker i, hvad det Socratiske, men kun til en vis Grad, selv er opmærksomt paa, og afhjælper, at der mangler en dialektisk Bestemmelse betræffende Overgangen fra det at have forstaaet Noget til det at gjøre det. I denne Overgang begynder det Christelige; ved saa at gaae ad den Vei, kommer det til at vise Synden liggende i Villien, til Begrebet Trods; og for saa at gjøre Enden rigtig fast, føies Dogmet om Arvesynden til – ak thi Spekulationens Hemmelighed med det at begribe er just Det at sye uden at fæste Ende og uden at slaae Knude paa Traaden, og derfor kan den, vidunderligt, blive ved at sye og sye, det vil sige at trække Traaden igjennem. Christendommen derimod fæster Ende ved Hjælp af Paradoxet.

I den rene Idealitet, hvor der ikke er Tale om det enkelte virkelige Menneske, der er Overgangen nødvendig (i Systemet gaaer jo ogsaa Alt for sig med Nødvendighed), eller der er slet ingen Vanskelighed forbunden med Overgangen fra det at forstaae til det at gjøre. Dette er Græciteten (dog ikke det Socratiske, thi dertil er Socrates for meget Ethiker). Og ganske dette Samme er egentlig hele den nyere Philosophies Hemmelighed; thi det er dette: cogito ergo sum, at tænke er at være; (christeligt derimod hedder det: Dig skeer, som Du troer, eller, som Du troer, saa er Du, at troe er at være). Man vil saaledes see, at den nyere Philosophi hverken er mere eller mindre end Hedenskab. Dette er nu endda ikke det Værste; at være i Slægtskab med Socrates ikke det Ringeste. Men det aldeles Usocratiske i den nyere Philosophi er, at den vil bilde sig og os ind, at dette er Christendom.

I Virkelighedens Verden derimod, hvor der er Tale om det enkelte Menneske, der er denne lille bitte Overgang fra at have forstaaet til at gjøre, den er ikke altid cito citissime, ikke, at jeg, i Mangel af at tale philosophisk, skal tale tydsk, geschwind wie der Wind. Tværtimod, her begynder en meget vidtløftig Historie.

I Aands-Livet er der ingen Stilstand (egentlig heller ingen Tilstand, Alt er Actualitet); dersom altsaa et Menneske ikke i samme Secund, som han har erkjendt det Rette, gjør det – ja saa gaaer nu for det første Erkjendelsen af Kog. Og dernæst saa bliver Spørgsmaalet, hvad synes Villien om det Erkjendte. Villien er et Dialektisk, og har igjen under sig hele den lavere Natur i Mennesket. Synes denne nu ikke om det Erkjendte, saa følger deraf vel ikke, at Villien gaaer hen og gjør det Modsatte af hvad Erkjendelsen forstod, saa stærke Modsætninger forekomme vistnok sjeldnere; men saa lader Villien nogen Tid gaae hen, det bliver et Interim, det hedder: vi vil dog see det an til imorgen. Under alt Dette bliver Erkjendelsen dunklere og dunklere, og det Lavere seirer mere og mere; ak, thi det Gode maa gjøres strax, strax idet det er erkjendt (og derfor gaaer det i den rene Idealitet saa let med Overgangen fra at tænke til at være, thi der er Alt strax), men det Lavere har sin Styrke i at trække ud. Saa smaat har Villien ikke noget mod, at dette skeer, den seer næsten igjennem Fingre dermed. Og naar saa Erkjendelsen er blevet behørig dunkel, saa kan Erkjendelsen og Villien bedre forstaae hinanden; tilsidst samtykke de ganske, thi nu er Erkjendelsen gaaet over paa Villiens Side, og erkjender, at det er ganske rigtigt som den vil det. Og saaledes lever maaskee en stor Mængde Mennesker; de arbeide saa smaat paa at fordunkle deres ethiske og ethisk-religieuse Erkjenden, der vil føre dem ud i Afgjørelser og Conseqventser, som det Lavere i dem ikke elsker; derimod udvide de deres æsthetiske og metaphysiske Erkjenden, hvilket, ethisk, er Adspredelse.

Dog med Alt dette ere vi endnu ikke komne videre end til det Socratiske; thi, vilde Socrates sige, skeer dette, saa viser det jo, at et saadant Menneske dog ikke har forstaaet det Rette. Dette vil sige, til at udsige, at En med sit Vidende gjør det Urette, med Viden om det Rette gjør det Urette, har Græciteten ikke Mod, saa hjælper den og siger: naar En gjør det Urette, har han ikke forstaaet det Rette.

Ganske rigtigt, og videre kan heller intet Menneske komme; intet Menneske kan ved sig selv og af sig selv udsige, hvad Synd er, just fordi han er i Synden; al hans Talen om Synden er i Grunden Besmykkelse for Synden, en Undskyldning, en syndig Formildelse. Derfor begynder Christendommen ogsaa paa en anden Maade, med, at der maa en Aabenbaring fra Gud til, for at oplyse Mennesket om, hvad Synd er, at Synden dog ikke ligger i, at Mennesket ikke har forstaaet det Rette, men i at han ikke vil forstaae det, og i at han ikke vil det.

Allerede med Hensyn til den Distinktion: ikke at kunne forstaae og ikke at ville forstaae, oplyser Socrates egentligen Intet, medens han derimod er Stormesteren for alle Ironikere i at operere ved Hjælp af den Distinktion mellem at forstaae og forstaae. Socrates forklarer, at Den, som ikke gjør det Rette, heller ikke har forstaaet det; men Christendommen gaaer lidt længere tilbage og siger, det er fordi han ikke vil forstaae det, og dette igjen, fordi han ikke vil det Rette. Og dernæst lærer den, at et Menneske gjør det Urette (den egentlige Trods), uagtet han forstaaer det Rette, eller lader være at gjøre det Rette, uagtet han forstaaer det; kort, den christelige Lære om Synden er lutter Nærgaaenhed mod Mennesket, Beskyldning paa Beskyldning, den er den Paastand, hvilken det Guddommelige som Aktor tillader sig at nedlægge mod Mennesket.

Men kan noget Menneske begribe dette Christelige? Ingenlunde, det er jo ogsaa det Christelige, altsaa til Forargelse. Det maa troes. At begribe er Menneskets Omfang i Forhold til det Menneskelige; men at troe er Menneskets Forhold til det Guddommelige. Hvorledes forklarer da Christendommen dette Ubegribelige? Ganske conseqvent, paa en ligesaa ubegribelig Maade, ved Hjælp af, at det er aabenbaret.

Christelig forstaaet ligger Synden altsaa i Villien, ikke i Erkjendelsen; og denne Villiens Fordærvethed gaaer ud over den Enkeltes Bevidsthed. Dette er det ganske Conseqvente; thi ellers maatte jo betræffende hver Enkelt det Spørgsmaal opkomme, hvorledes Synden begyndte.

Her er da igjen Forargelsens Kjendetegn. Forargelsens Mulighed ligger i: at der skal en Aabenbaring fra Gud til for at oplyse Mennesket om, hvad Synd er, og hvor dybt den stikker. Det naturlige Menneske, Hedningen, tænker som saa: »lad gaae, jeg indrømmer, at jeg ikke har forstaaet alle Ting i Himlen og paa Jorden, skal der en Aabenbaring til, saa lad den oplyse os om det Himmelske; men at der skulde en Aabenbaring til for at oplyse hvad Synd er, er da det Urimeligste af Alt. Jeg udgiver mig ikke for at være noget fuldkomment Menneske, langt fra, men det veed jeg dog, og jeg er jo villig til at indrømme det, hvor langt jeg er fra Fuldkommenhed: skulde jeg saa ikke vide, hvad Synd er?« Men Christendommen svarer: nei, det er hvad Du allermindst veed, hvor langt Du er fra Fuldkommenhed, og hvad Synd er. – See, i den Forstand er rigtignok, christelig, Synd Uvidenhed, den er Uvidenhed om hvad Synd er.

Definitionen paa Synd, som blev givet i forrige Capitel, maa derfor endnu saaledes fuldstændiggjøres: Synd er, efter ved en Aabenbaring fra Gud at være oplyst om hvad Synd er, for Gud fortvivlet ikke at ville være sig selv, eller fortvivlet at ville være sig selv.

[…]

Altsaa Fortvivlelse om Syndernes Forladelse er Forargelse. Og Forargelse er Syndens Potentsation. Dette tænker man i Almindelighed slet ikke paa; man regner i Almindelighed vel neppe Forargelse til Synd, om hvilken man da ikke taler, men om Synder, blandt hvilke Forargelse ikke finder Plads. Endnu mindre opfatter man Forargelse som Syndens Potentsation. Dette kommer af, at man ikke, christeligt, danner Modsætningen mellem: Synd – Tro, men mellem Synd – Dyd.

[…]

Derimod er denne Modsætning gjort gjældende i hele dette Skrift, der strax i første Afsnit A, A opstillede Formelen for den Tilstand, hvori der slet ingen Fortvivlelse er: i at forholde sig til sig selv og i at ville være sig selv grunder Selvet gjennemsigtigt i den Magt, som satte det. Hvilken Formel igjen, hvorom oftere er mindet, er Definitionen paa Tro.

\workheading{Indøvelse i Christendom}{1850}

\worknote{Af Anti-Climacus. Antologiens uddrag (s. 373–384) er sat sammen af fire adskilte steder, som Hongs SV¹-margental viser.}

\sectionheading{Udgiverens Forord  ·  SV¹ XII, 15}

I dette Skrift, hidrørende fra Aaret 1848, er Fordringen til det at være Christen af Pseudonymen tvunget op til et Idealitetens Høieste.

Dog siges, fremstilles, høres bør jo Fordringen; der bør, christeligt, ikke slaaes af paa Fordringen, ei heller den forties – istedetfor at gjøre Indrømmelse og Tilstaaelse sig selv betræffende.

Høres bør Fordringen; og jeg forstaaer det Sagte som sagt alene til mig – at jeg maatte lære ikke blot at henflye til »Naaden«, men at henflye til den i Forhold til Benyttelsen af »Naaden«.

S. K.

\sectionheading{Denne Fremstillings Indhold — Nr. II  ·  SV¹ XII, 91–93}

Som Begrebet »Tro« er en ganske eiendommelig christelig Bestemmelse, saaledes er igjen »Forargelse« en ganske eiendommelig christelig Bestemmelse, der forholder sig til Tro. Forargelsens Mulighed er Skilleveien, eller er som det at staae paa Skilleveien. Man svinger fra Forargelsens Mulighed enten til Forargelsen eller til Troen; men man kommer aldrig til Troen uden fra Forargelsens Mulighed.

1. I Skrifter af nogle Pseudonymer er viist, at man i den nyere Philosophi forvirrende har talt om Tvivl, hvor man skulde tale om Fortvivlelse. Derfor har man saa ikke kunnet styre eller magte Tvivlen hverken i Videnskaben eller i Livet. »Fortvivlelse« viser derimod strax til Rette, ved at føre Forholdet ind under Personlighedens Bestemmelse (den Enkelte), og ind under det Ethiske. Men som man saaledes forvirrende har talt om »Tvivl« istedetfor at tale om »Fortvivlelse«, saaledes har man ogsaa havt for Skik at bruge Kategorien »Tvivl«, hvor man skulde tale om »Forargelse.« Forholdet, Personlighedens Forhold til Christendommen er ikke: at tvivle eller at troe, men: at forarges eller at troe. Hele den moderne Philosophi er, baade ethisk og christelig, baseret paa en Letfærdighed. Istedetfor at afskrække og kalde til Orden ved at tale om at fortvivle og forarges, har den vinket og indbudt Menneskene til at blive indbildske, ved at tvivle og at have tvivlet. Den moderne Philosophi er, som abstrakt, svævende i det Metaphysiskes Ubestemthed. Istedetfor nu at forklare dette om sig selv, og saa at henvise Menneskene (de enkelte Mennesker) til det Ethiske, det Religieuse, det Existentielle, har Philosophien givet det Skin af, at Menneskene kunne, som man ganske prosaisk siger, speculere sig ud af deres gode Skind og ind i det rene Skin.()

Forargelse forholder sig væsentlig til Sammensætningen af Gud og Menneske, eller til Gud-Mennesket. Speculationen har naturligviis meent at kunne »begribe« Gud-Mennesket – det kan man godt begribe; thi Speculationen tager Timelighedens, Samtidighedens, Virkelighedens Bestemmelse bort fra Gud-Mennesket. Overhovedet er det sørgeligt og forfærdeligt, at, hvad man ikke bruger for stærkt et Udtryk om, naar man siger, at det ikke er Andet end at gjøre Narrestreger og Nar af Folk, at det er blevet feteret som Dybsind. Nei, Situationen hører med til Gud-Mennesket, Situationen at et enkelt Menneske, som staaer ved Siden af Dig, er Gud-Mennesket. Gud-Mennesket er ikke Eenheden af Gud og Menneske; saadan Terminologi er dybsindig Øienforblindelse. Gud-Mennesket er Eenheden af Gud og et enkelt Menneske. At Menneskenes Slægt er eller skal være i Slægtskab med Gud, er gammelt Hedenskab; men at et enkelt Menneske er Gud, er Christendom, og dette enkelte Menneske er Gud-Mennesket. Der er hverken i Himlen eller paa Jorden eller i Afgrunden eller i den meest phantastiske Tænknings Forvildelser Mulighed af en, menneskelig talt, afsindigere Sammensætning. Saaledes viser det sig i Samtidighedens Situation; og til Gud-Mennesket er intet Forhold muligt uden begyndende med Samtidighedens Situation.

1. Herom maa jeg henvise til: »Kommer hid alle I, som arbeide og ere besværede«, Standsningen.()

Forargelse i strengeste Forstand, Forargelse ϰατ' εξοχην forholder sig altsaa til Gud-Mennesket, og har tvende Former. Den er enten i Retning af Høiheden, man forarges over, at et enkelt Menneske siger sig at være Gud, handler eller taler paa en Maade som forraader Gud (Dette afhandles under B). Eller Forargelse er i Retning af Ringheden, at Den, der er Gud, er dette ringe Menneske, lidende som et ringe Menneske (Dette afhandles under C). I første Form kommer Forargelsen saaledes, at jeg ingenlunde forarges over det ringe Menneske, men over at han vil, jeg skal troe han er Gud. Og har jeg troet det, saa kommer Forargelsen fra den anden Side, at han skulde være Gud, han dette ringe, afmægtige Menneske, som, da det kommer til Stykket, slet Intet formaaer. I det ene Tilfælde gaaes der ud fra Bestemmelsen Menneske, og Forargelsen er over Bestemmelsen Gud; i det andet gaaes der ud fra Bestemmelsen Gud, og Forargelsen er over Bestemmelsen Menneske.

Gud-Mennesket er Paradoxet, absolut Paradoxet; derfor er det ganske sikkert, at Forstanden skal komme til at staae stille derpaa. Mærker et Menneske ikke Forargelsen i Retning af Høiheden, saa vil han opdage den i Retning af Ringheden. Det var ikke utænkeligt, at et Menneske med overveiende Phantasi og Følelse, en Repræsentant for den barnlige eller barnagtige Christendom (thi for et Barn er Forargelsen ϰατ' εξοχην ikke til, og just derfor er Christendommen egentlig heller ikke til for Barnet) kunde gaae hen og mene, at han troede, at dette enkelte Menneske var Gud, uden at opdage Forargelsen. Det vil have sin Grund i, at en Saadan ikke har en udviklet Forestilling om Gud, men en barnlig eller barnagtig Phantasi om noget Overordentligt, noget uendelig høit Ophøiet, Helligt og Reent, en Forestilling om En, der saadan er større end alle Konger o. s. v., uden at deri just ligger Qvaliteten: Gud. Dette vil sige, en Saadan havde ingen Kategori, derfor lod det sig gjøre at mene, at han troede, et enkelt Menneske er Gud, uden at støde paa Forargelsen. Men dette samme Menneske vil saa støde an paa Ringheden.

Saaledes med Forargelsen; og saaledes er den ogsaa fremstillet i den hellige Skrift paa de Steder, hvor Christus selv advarer mod Forargelsen.

Men da er der i Skriften endnu Tale om en Forargelse paa Christus, hvis Mulighed er et historisk Forbigangent. Denne Forargelse forholder sig nemlig ikke til Christus som Christus, som Gud-Mennesket (dette er den væsentlige Forargelse, og dens tvende Former blive, saa længe Timeligheden bliver, saa længe til Troen afskaffes), men til ham som et slet og ret enkelt Menneske, der colliderer med et Bestaaende (Dette afhandles under A).

\sectionheading{Af Nr. III: Barnet og Billedet af den Korsfæstede  ·  SV¹ XII, 176–179}

Tænk Dig da et Barn, og glæd nu dette Barn ved at vise det nogle af hine i kunstnerisk Henseende ubetydelige men for Børn saa værdifulde Billeder, hvilke man kjøber i Kramboden. – Han her paa den fnysende Ganger, med den vaiende Fjær, med Hersker-Minen, i Spidsen for de Tusinder og Tusinder, som Du ikke seer, Haanden udrakt til Befaling, »fremad«, fremad over Toppen af Bjergene, som Du seer ligge for Dig, fremad til Seier: det er Keiseren, den eneste, Napoleon; og nu fortæller Du Barnet Lidt om Napoleon. – Han her er klædt som Jæger; han støtter sig ved sin Bue og seer hen for sig med et Blik, saa gjennemtrængende, saa sikkert, og dog saa bekymret. Det er Wilhelm Tell; Du fortæller nu Barnet Lidt om ham, og om dette mærkelige Blik, at han med det samme Blik paa een Gang har Øie for det elskede Barn, ikke at skyde det, og med det samme Blik Øie for Æblet paa Barnets Hoved, at han træffer. – Og saaledes viser Du nu Barnet flere Billeder til Barnets usigelige Glæde, da kommer Du til et, som med Flid var lagt imellem, det forestiller en Korsfæstet. Barnet vil ikke strax, end ikke ganske ligefrem, forstaae dette Billede, det vil spørge, hvad det betyder, hvorfor han hænger paa et saadant Træ. Saa forklarer Du Barnet, at det er et Kors, og det at hænge derpaa betyder at være korsfæstet, og at Korsfæstelse var i hiint Land den piinligste Dødsstraf, derhos en vanærende Dødsstraf, som kun anvendtes mod de groveste Forbrydere. Hvorledes vil nu dette virke paa Barnet? Barnet vil blive underligt tilmode, det vil vel egentligen undres over, hvor det kan falde Dig ind at lægge saadant et stygt Billede ind imellem alle de andre deilige, ind imellem alle disse Helte og Herlige, Billedet af en grov Forbryder. Thi som der Jøderne til Trods kom til at staae over hans Kors »Jødernes Konge«, saaledes er dette Billede, der bestandigt udkommer iaar, Slægten til Trods en Erindring, den aldrig kan og aldrig skal blive af med, han skal ikke fremstilles anderledes; og det skal være som var det denne Slægt, der korsfæstede ham, hver Gang denne Slægt første Gang viser Barnet af den nye Slægt dette Billede, første Gang forklarende hvorledes det gik til i Verden; og Barnet skal, første Gang det hører det, blive angst og bange for den Ældre og Verden og sig selv; og de andre Billeder, ja de skulle alle, som der staaer i Visen, vende sig omkring, saa forskjelligt er dette Billede. Imidlertid – og vi ere jo heller ikke endnu komne til det Afgjørende, Barnet har jo endnu ikke faaet at vide, hvo denne grove Forbryder var – vil sagtens Barnet, videbegjerligt som et Barn altid er, dog spørge, hvem er det, hvad havde han gjort, hvad? Fortæl saa Barnet, at denne Korsfæstede er Verdens Frelser. Dog dette vil Barnet ikke kunne knytte nogen bestemt Forestilling til, fortæl det derfor blot, at denne Korsfæstede var det kjerligste Menneske, der nogensinde har levet. O, i Handel og Vandel, hvor man jo kan det uden ad paa Remse med den Geschichte, i Handel og Vandel, hvor et halvt Ord, henkastet som en Antydning, er nok, saa veed Enhver det: der gaaer det saa gesvindt; men sandeligen det maatte være et underligt Menneske, eller rettere et Umenneske, der ikke uvilkaarligt slog Øiet ned og stod næsten som en arm Synder, i det Øieblik han skal fortælle et Barn det første Gang, et Barn, som altsaa aldrig har hørt et Ord derom, og selvfølgeligt altsaa heller aldrig anet noget Saadant. Men saa staaer jo den Ældre i det Øieblik som en Anklager, der anklager sig selv og hele Slægten! – Hvad Indtryk troer Du nu det vil gjøre paa Barnet, som naturligviis vil spørge, men hvorfor var man da saa slem mod ham, hvad?

See, nu er Øieblikket; hvis Du ikke allerede har gjort for stærkt et Indtryk paa Barnet, saa fortæl det nu om ham, den Ophøiede, der fra Høiheden vil drage Alle til sig. Fortæl det, at denne Ophøiede han er denne Korsfæstede. Fortæl Barnet, at han var Kjerlighed, at han kom til Verden af Kjerlighed, iførte sig en ringe Tjeners Skikkelse, levede kun for Eet, at elske og hjælpe Menneskene, især alle dem, som vare syge og sorrigfulde og lidende og ulykkelige. Fortæl saa Barnet, hvorledes det gik ham i Livet, hvorledes En af de Faa, der stode ham nær, forraadte ham, de faa Andre fornægtede ham, og alle Andre forhaanede og bespottede ham, indtil de tilsidst naglede ham til Korset – som det sees paa Billedet – begjerende, at hans Blod maatte komme over dem og deres Børn, medens han bad for dem, at dette ikke maatte skee, at den himmelske Fader vilde tilgive dem denne Skyld. Fortæl Barnet det ret levende, som havde Du aldrig selv hørt det før, eller før fortalt Nogen det; fortæl det, som havde Du selv digtet det Hele, men glem intet Træk, der er opbevaret, kun maa Du gjerne, idet Du fortæller det, glemme at det er opbevaret. Fortæl Barnet, at der samtidigen med denne Kjerlige levede en berygtet Røver, som var dømt fra Livet – ham fordrede Folket frigivet, ham raabte de Leve for, leve Barrabas; men over den Kjerlige raabte de »korsfæst, korsfæst«, saa altsaa denne Kjerlige ikke blot blev korsfæstet som en Forbryder, men som et saadant Uhyre af en Forbryder, at hiin berygtede Røver blev en Slags retskaffen Mand i Sammenligning med denne Kjerlige.

Hvad Virkning troer Du denne Fortælling vil gjøre paa Barnet? Dog, for ret at belyse det, hvorom Talen er, saa gjør en Prøve, fortsæt Fortællingen om denne Korsfæstede, at han derpaa tredie Dagen efter opstod fra de Døde, saa opfoer til Himmels, for at indgaae til Herligheden hos Faderen i Himlene – prøv det, Du skal see, dette vil Barnet for det første næsten overhøre; Fortællingen om hans Lidelse vil have gjort saa dybt et Indtryk paa Barnet, at det slet ikke føler sig oplagt til at høre om Herligheden, som fulgte. Thi man maa være betydelig forqvaklet, og gjennem mange Aar forvænt med at have faaet den hele Historie om hans Fornedrelse, Lidelse og Død letfærdig paa Remse, for, uden at føle nogen Standsning derved, strax at kunne gribe efter Høiheden.

Altsaa hvilken Virkning troer Du, at denne Fortælling vilde fremkalde paa Barnet? Først og fremmest vel den, at det ganske glemte de øvrige Billeder, Du havde viist det; thi nu havde det faaet ganske Andet at tænke paa. Og saa vilde Barnet vel falde i den dybeste Forundring over, at Gud i Himlene ikke havde gjort Alt for at forhindre, at dette skeete; eller at det skeete uden at Gud, om ikke før saa dog i det sidste Øieblik, for at forhindre hans Død, lod Ild regne ned fra Himlen; at det skeete, uden at Jorden aabnede sig for at opsluge de Ugudelige. Og saaledes maatte vi Ældre jo ogsaa forstaae det, hvis vi ikke forstode, at det var frivillig Lidelse, derfor desto tungere, at han, den Fornedrede, i ethvert Øieblik havde det i sin Magt at bede, saa havde Faderen sendt ham Englenes Legioner for at afværge det Forfærdelige. – Dette var vel det første Indtryk. Men efterhaanden som Barnet gik og tænkte paa denne Historie, vilde det vel blive mere og mere lidenskabeligt; det vilde kun tænke paa og tale om Vaaben og Krig – thi det havde Barnet fast besluttet, naar han blev stor, at ville hugge alle disse ugudelige Mennesker ihjel, som havde handlet saaledes med den Kjerlige; det havde Barnet besluttet, barnligt glemmende, at det var over 1800 Aar siden de levede.

\sectionheading{Af Nr. III: Beundreren og Efterfølgeren  ·  SV¹ XII, 227–230}

Herre Jesus Christus, ikke kom Du til Verden for at lade Dig tjene, og altsaa heller ikke for at lade Dig beundre eller i den Forstand tilbede. Selv var Du Veien og Livet – og Du har kun forlangt »Efterfølgere«. Saa opvække Du os da, hvis vi ere slumrede ind i denne Bedaarelse, frels os fra denne Vildfarelse, at ville beundre eller tilbedende beundre Dig, istedetfor at ville følge Dig efter og ligne Dig.

Joh. XII, 32: Og jeg, naar jeg bliver ophøiet fra Jorden, vil jeg drage Alle til mig.

I Christenheden høres der ofte nok Prædikener, Foredrag, Taler om, hvad der fordres af en Christi Efterfølger, hvad det vil sige at følge Christum efter o. D. Hvad der høres er almindeligviis ganske rigtigt og sandt, kun ved at høre nøiere efter opdager man en dybt skjult, uchristelig Grund-Forvirring og Mislighed. Den christelige Prædiken er nutildags hovedsageligen blevet »Betragtning«: lader os i denne Time betragte; jeg indbyder mine Tilhørere til Betragtninger over; Betragtningens Gjenstand er o. s. v. Men at betragte kan i een Forstand betyde at komme ganske nær til Noget, til Det nemlig man vil betragte, i en anden Forstand at holde sig meget fjernt, uendelig fjernt, personligt nemlig. Naar man viser En et Maleri og opfordrer ham til at betragte det, eller naar i Handel og Vandel En betragter f. Ex. et Stykke Klæde, da træder han ganske nær til Gjenstanden, i sidste Tilfælde endog tager og føler han paa den, kort, han kommer Gjenstanden saa nær som det er muligt; men i en anden Forstand gaaer han just ved denne Bevægelse ganske ud af sig selv, bort fra sig selv, glemmer sig selv, og Intet erindrer ham derom, da det jo er ham, der betragter Maleriet og Klædet, ikke Maleriet og Klædet, der betragter ham. Dette vil sige, ved at betragte gaaer jeg ind i Gjenstanden (jeg bliver objektiv), men jeg gaaer ud af eller bort fra mig selv (jeg ophører at være subjektiv). Saaledes har Prædikeforedraget ved sin Yndlings-Betragtning af det Christelige, hvilken just er »Betragtningen« og »Betragtningerne«, afskaffet hvad der christelig er det Afgjørende i Prædikeforedraget, det Personlige, dette »Du og Jeg«, den Talende og Den til hvem der tales, dette, at Den, der taler, selv er personligt i Bevægelse, en Stræbende, og ligesaa den Tiltalte, hvem han derfor opmuntrer, tilskynder, formaner, advarer, men Altsammen i Henseende til en Stræben, et Liv, dette, at Taleren har bestandigt, ikke at komme bort fra sig selv, men at komme tilbage til sig selv, og at være Tilhøreren behjælpelig, ikke i at komme bort fra sig selv, men i at komme tilbage til sig selv; Prædikeforedraget i vor Tid har først selv reent overseet og derpaa virket til at bringe det ganske i Glemme, at den christelige Sandhed egentligen ikke kan være Gjenstand for »Betragtning«. Thi den christelige Sandhed har, om jeg saa tør sige, selv Øine at see med, ja den er som lutter Øie; men det vilde jo være meget forstyrrende, ja det vilde jo saa være mig umuligt at betragte et Maleri eller et Stykke Klæde, naar jeg, idet jeg vilde til at betragte, opdagede, at det var Maleriet eller Klædet, der saae paa mig – og saaledes er det just Tilfældet med den christelige Sandhed, det er den der betragter mig, om jeg gjør, som den siger, jeg skal gjøre. See, derfor lader den christelige Sandhed sig ikke fremstille for Betragtningen, ei heller foredrage som Betragtning, den har nemlig selv, om jeg saa tør sige, Øren at høre med, ja, den er som lutter Øre, den hører efter, idet den Talende taler; man kan ikke tale om den som om en Fraværende eller et kun gjenstandsviis Nærværende, thi da den er fra Gud og Gud i den, saa er den i en ganske egen Forstand nærværende, idet der tales om den, og ikke som Gjenstand, snarere bliver Taleren Gjenstanden for den, han maner en Aand frem, der, idet han taler, overhører ham.

Derfor er det et Vovestykke at prædike; thi idet jeg bestiger det hellige Sted – hvad enten Kirken er overfyldt eller saa godt som tom, hvad enten jeg selv er opmærksom derpaa eller ikke: jeg har een Tilhører mere end hvad der er til at see, en usynlig Tilhører, Gud i Himlene, som jeg rigtignok ikke kan see, men som sandeligen kan see mig. Denne Tilhører, han hører nøie efter, om det er sandt, hvad jeg siger, om det er sandt i mig, altsaa han seer efter – og det kan han, just fordi han er usynlig, paa en Maade, saa det er umuligt at tage sig iagt for ham – han seer efter, om mit Liv udtrykker, hvad jeg siger. Og om jeg end ikke har Myndighed til at forpligte noget andet Menneske: hvert Ord, jeg har sagt fra Prædikestolen i Prædikenen, det har jeg forpligtet mig selv til – og Gud har hørt det. I Sandhed, et Vovestykke at prædike! De Fleste have vistnok en Forestilling om, at det er Noget, hvortil der fordres Mod, som Skuespiller at træde frem paa Skuepladsen, at vove sig ud i den Fare, at Alles Øine sigte paa En. Og dog er denne Fare i een Forstand, som Alt paa Skuepladsen, Blendværk; thi Skuespilleren er jo personligt udenfor, hans Opgave just at bedrage, at forstille sig, at forestille en Anden, og nøiagtigt at gjengive en Andens Ord. Den christelige Sandheds Forkynder derimod, han træder frem paa et Sted, hvor om ikke Alles Øine sigte paa ham, saa dog en Alvidendes Øie; han har den Opgave: at være sig selv, og i en Omgivelse, Guds Huus, der, lutter Øie og Øre, ene fordrer dette af ham, at han er sig selv, er sand. Er sand, det er, at han selv er hvad han forkynder, eller dog stræber efter at være det, eller dog er sanddru nok til at tilstaae om sig selv, at han ikke er det – ak, og hvor Mangen, der, for at forkynde Christendommen, bestiger det hellige Sted, er vel endog blot lydhør nok til at opdage det helligede Steds Uvillie og Spot over ham, at han, begeistret, rørt, grædende forkynder, hvad hans Liv udtrykker det Modsatte af!

Saa voveligt er det at være det Jeg, der prædiker, den Talende, et Jeg, der ved at prædike og idet han prædiker, forpligter sig selv ubetinget, lægger sit Liv frem, at man, om muligt, kunde see ham lige ind i Sjelen: at være dette Jeg, det var voveligt! Derfor fandt Præsten lidt efter lidt paa, at trække sit Øie ind i sig selv, for dermed at antyde, at Ingen havde at see paa ham. Det var jo, meente han, ikke sig selv, han talte om, det var om Sagen; og dette blev beundret som et overordentligt Fremskridt i Viisdom, at saaledes den Talende paa en Maade ophørte at være et Jeg, og blev, hvis det er muligt, Sagen. I ethvert Tilfælde blev det paa den Maade langt lettere at være Præst – den Talende prædikede ikke mere, han anvendte disse Øieblikke til at anstille nogle Betragtninger. Nogle Betragtninger! Det seer man paa den Talende; hans Blik dragersig ind i Øiet, han ligner ikke saa meget et Menneske som en af hine i Steen udhugne Skikkelser, der ikke have Øine. Herved befæsterhan et svælgende Dyb mellem Tilhøreren og sig, næsten lige saa svælgende som det, der er mellem Skuespiller og Tilskuer. Og Det, han foredrager, er »Betragtning«, derved befæster han, mellem sig og det han siger, igjen et svælgende Dyb som mellem Skuespilleren og Digteren; personligt er han saavidt muligt udenfor, medens han »anvender disse Øieblikke til at anstille Betragtninger.«

Saaledes gik det Jeg ud, som var den Talende; den Talende er ikke Jeg, han er Sagen, Betragtningen. Og idet dette Jeg gik ud, afskaffedes selvfølgeligt ogsaa dette »Du«, Du, den Hørende, at det er Dig, Du som der sidder, til hvem der tales. Ja, man er vel snart kommet saavidt, at man anseer det for »Personligheder« paa den Maade at tale til andre Mennesker. Ved Personligheder, at bruge Personligheder, at tillade sig Personligheder, derved forstaaer man jo en usømmelig, udannet Adfærd – og altsaa gaaer det jo ikke an at tale personligt (det talende Jeg) og til Personer (det hørende Du). Og gaaer det ikke an, saa er det at prædike afskaffet. Men det er det jo ogsaa – man anstiller kun Betragtninger. Og »Betragtningen« kommer hverken den Talende eller den Hørende for nær, Betragtningen sikkrer ganske tilforladeligt imod, at det ikke kommer til Personligheder: det er ikke mig, den Talende, der tales om, neppe er det mig, der taler, det er Betragtning, og det er ikke Dig, den Hørende, der tales til, det er Betragtning; om jeg gjør hvad jeg siger, det kommer ikke Dig ved, naar blot Betragtningen er rigtig, neppe kommer det mig selv ved, da jeg jo dog vel skylder mig selv samme Hensyn som enhver Anden, ikke at tillade mig Personligheder; om Du, den Hørende, gjør hvad der siges, kommer ikke mig ved, neppe Dig selv, det er Betragtning, og der er i det Høieste Spørgsmaal om, hvorvidt Betragtningen har tilfredsstillet Dig.

Denne Grund-Forandring i Prædikeforedraget, hvorved Christendommen blev afskaffet, er blandt Andet ogsaa Udtrykket for den Grund-Forandring, som indtraadte med den triumpherende Kirke og den bestaaende Christenhed, at Christus ordentligviis i det Høieste fik Beundrere, ikke Efterfølgere.

Ved at fremstille denne Forskjel, Forskjellen mellem: en Beundrer og en Efterfølger, vil denne Udvikling stræbe at belyse Christendommen, atter med stadigt Hensyn til de hellige Ord »fra Høiheden vil han drage Alle til sig«; thi det er atter her Høiheden og Ringheden, eller Forholdet til Høiheden og Forholdet til Ringheden, der gjør Udslaget. Er Christus kun til for os i Høiheden, er det glemt med hans Fornedrelse, eller havde han aldrig været til i Ringhed, saa kunde – i dette Tilfælde – han end ikke selv, for at være i Overeensstemmelse med sig selv, forlange Andet end Beundrere, tilbedende Beundrere; thi Høihed og Beundrer, guddommelig Høihed – tilbedende Beundrer, svare ganske til hinanden. Ja i Forhold til Høihed vilde det jo endog fra vor Side være Ubeskedenhed, Anmasselse, Forblindelse, tildeels Vanvid, at ville være Efterfølgere, istedetfor sømmeligen ikke at tragte efter hvad der maaskee ikke er givet os, fordi det er givet den Anden, sømmeligen at nøies med at beundre og tilbedende at beundre. Til Fornedrelse og Ringhed derimod svarer: Efterfølger.

Det er nu bekjendt nok, at Christus bestandigt bruger det Udtryk: Efterfølgere; han taler aldrig om, at han forlanger: Beundrere, tilbedende Beundrere, Tilhængere; og naar han bruger det Udtryk: Disciple, forklarer han det altid saaledes, at man seer, derved forstaaes Efterfølgere, at det ikke er Tilhængere af en Lære, men Efterfølgere af et Liv, der ikke, ved nogen tilfældig Høihed, gjør det at ville ligne det til Formastelse eller Vanvid. Det er ogsaa bekjendt nok, hvad andetsteds atter og atter er gjentaget, at det er den fornedrede Christus der taler, at hvert Ord vi have af Christus, er af ham, den Fornedrede. Nu maa man dog vel antage, at Christus selv fuldeligt har vidst, hvorfor han valgte netop dette Udtryk, der ene og ubetinget er i den inderligste og dybeste Overeensstemmelse med, hvad han bestandigt sagde om sig selv, eller sagde sig selv at være: Sandheden nemlig, og Veien og Livet, at han ikke var en Lærer i den Forstand, at han blot havde en Lære at foredrage, saa han kunde nøies med Tilhængere, der antoge Læren – i Livet derimod lode som Ingenting, eller lode fem være lige. Og man maa dog vel ogsaa antage, at han selv fuldeligen har vidst, hvorfor hans hele Liv paa Jorden, fra først til sidst, var beregnet paa kun at kunne faae »Efterfølgere«, og beregnet paa, at gjøre »Beundrere« umulige.

Christus kom til Verden med den Hensigt at frelse Verden, tillige med den Hensigt – dette ligger atter indeholdt i hans første Hensigt – at ville være »Forbilledet«, at ville efterlade et Fodspor for Den, som vilde slutte sig til ham, hvilken altsaa maatte blive en »Efterfølger«, det svarer jo til et »Fodspor«. Just derfor lod han sig føde i Ringhed, og levede derpaa fattig,forladt, foragtet, fornedret – ja, intet Menneske har levet saa fornedret som han; selv det ellers ringeste Menneske maatte ved at sammenligne sit Livs Vilkaar med hans komme til den Afgjørelse, at hans Liv dog, menneskelig talt, var at foretrække i Sammenligning med Christi Livs Vilkaar. Hvortil nu det, hvorfor denne Ringhed og Fornedrelse? Fordi Den, der i Sandhed skal være »Forbilledet« og kun forholde sig til Efterfølgere, han maa i een Forstand, fremad drivende, ligge bagved Menneskene, medens han i en anden Forstand vinkende staaer foran dem. Dette er Høihedens og Ringhedens Forhold i »Forbilledet«. Høiheden maa ikke være den ligefremme, hvilken er den verdslige, den jordiske, men den aandelige, og saaledes just Nægtelsen af den verdslige og jordiske Høihed. Ringheden maa være den ligefremme; thi den ligefremme Ringhed, naar man skal igjennem den, er just Veien, men tillige for det verdslige og jordiske Sind Omveien, som betrygger, at Høiheden ikke tages forfængeligt. »Forbilledet« ligger saaledes uendeligt nær, i Ringheden og Fornedrelsen, og dog uendeligt fjernt, i Høiheden, ja endnu fjernere end om det blot var til fjernt i Høiheden; thi at man for at naae det, for at bestemme sig i Lighed dermed, skal igjennem Ringheden og Fornedrelsen, at der slet ingen anden Vei er, er en endnu større, er egentligen den uendelige Fjernhed. Og saaledes er »Forbilledet« i een Forstand bag ved, i Ringhed og Fornedrelse dybere nedtrykt end nogensinde noget Menneske har været det, og i en anden Forstand for an, uendelig ophøiet. Men bagved maa Forbilledet være for at kunne fange og omfatte Alle; var der et eneste Menneske, som med Sandhed kunde byde under eller dukke under, ved at oplyse, at han i Ringhed og Fornedrelse var endnu lavere stillet, saa er Forbilledet ikke »Forbilledet«, saa er det kun et ufuldkomment Forbillede, det er, kun Forbillede for en stor Mængde Mennesker. Ubetinget bag ved maa Forbilledet være, bag ved Alle, og det maa være bag ved, for at drive Dem, der skulle dannes efter det, fremad.

Der ligger nemlig i Menneske-Slægten, og i hver Enkelt i Slægten, en dyb, bevidst eller ubevidst, Underfundighed i Forhold til hvad der skal være Forbilledet for dem, en Underfundighed, som er af det Onde. Er Den, som skal være Forbilledet, i Besiddelse af jordiske, verdslige, timelige Fortrin, hvad saa? Ja, saa er Forbilledet forkeert stillet, forkeert vendt, og saa benytter Slægten samt hver Enkelt i Slægten dette til igjen paa sin Side at gjøre forkeert omkring. Forbilledet bliver da puffet udenfor som en Opfordring til Digter-Beundring, istedetfor at Forbilledet skulde ligge bag ved, komme bag paa Menneskene som en Fordring til dem. Idet saa Forbilledet er blevet Gjenstand for Beundring, saa skulke Menneskene fra »Fordringen«, de sige: »han kan sagtens, han som er i Besiddelse af alle disse Fortrin og Begunstigelser; vare vi kun i hans Sted, da skulde vi være lige saa fuldkomne som han. Nu kunne vi ikke Andet end beundre ham, og det er vor Ære og Roes, at vi gjøre det, at vi nemlig ikke hengive os til Misundelse. Men Andet end beundre ham kunne vi ikke; thi han er i Besiddelse af Betingelser, som vi ikke ere, og som han ikke kan give os, hvor urimeligt da, at fordre det Samme af os, som han fordrer af sig selv«.

Christus er nu »Forbilledet«. Var han kommet til Verden i jordisk og timelig Høihed, saa var den størst mulige Usandhed afstedkommet. Istedetfor at blive »Fordringen« til hele Slægten og hver Enkelt i Slægten, var han blevet som en General-Undskyldning og Udflugt for hele Slægten og hver Enkelt i Slægten. Han var saa heller ikke blevet slagen ihjel – thi ogsaa dette bidrog til at ophidse de Samtidige imod ham, at de, om jeg saa tør sige, ikke kunde faae ham vendt som de vilde, at han »trodsigt og stedigt« vilde være den Fornedrede, og, hvad der forbittrer Menneskenes selvkjerlige Blødagtighed allermeest, blot vilde have »Efterfølgere« – nei, han var blevet Beundringens Gjenstand, og Forvirringen var blevet saa uhyre, som man neppe kan tænke sig den. Han sagde sig jo selv at være Sandheden, og da nu, efter Antagelsen, Menneskene beundrede ham, saa saae det jo ud, som elskede de ogsaa Sandheden, og det var næsten gjort umuligt at hitte Rede i dette. Forvirringen vilde nemlig i Samtidighedens Bestedelse være blevet lige saa stor som den er i den bestaaende Christenhed, hvor En i stærkeste Udtryk beundrer og tilbedende beundrer og beundrer og tilbeder Christus – medens hans Liv dog udtrykker lige det Modsatte af Christi Liv, som dette blev ført paa Jorden af ham, der just for at være »Forbilledet« fødtes og levede i Ringhed og Fornedrelse. Men den Beundrende har et ypperligt Skjul; »thi«, vil han sige, »mere kan man jo dog ikke fordre af mig end at jeg i de stærkeste Udtryk – og er der endnu stærkere Udtryk i Sproget, saa er jeg ogsaa villig til at bruge dem – erkjender og bekjender, at jeg beundrende tilbeder Christus som Sandheden. Mere kan der jo dog ikke fordres af mig; kan Du sige noget Høiere?«

See, derfor fødtes og levede Christus i Fornedrelse; der levede intet, ubetinget intet Menneske samtidigen med ham saa fornedret, der har aldrig levet noget Menneske saa fornedret, der var altsaa ubetinget for ethvert Menneske umuligt, at skulke sig fra »Fordringen« ved Undskyldning og Udflugt, foranlediget af, at »Forbilledet« jo var i Besiddelse af jordiske og verdslige Fortrin, som han ikke havde. I den Forstand at beundre Christus er en senere Tids usande Opfindelse ved Hjælp af »Høiheden«. Der var, saaledes forstaaet, ubetinget Intet at beundre, med mindre man da vil beundre Armod, Elendighed, Foragtethed o. D. End ikke det Yderste gik han fri for, at blive ynket, en ynkelig Gjenstand for Medlidenhed. Nei, der var sandeligen ikke det Mindste at beundre.

Og der var i Samtidighedens Bestedelse heller ingen Leilighed til at beundre; thi Christus havde kun det samme Vilkaar at byde Den, der sluttede sig til ham – og paa det Vilkaar har aldrig nogen Beundrer villet være med; det samme Vilkaar: at blive lige saa fattig, foragtet, forhaanet, bespottet, om muligt endnu lidt mere, i Betragtning af, at man oven i Kjøbet var Tilhænger af en saadan Foragtet, hvem da enhver Forstandig flyede.

Hvilken er da Forskjellen mellem »en Beundrer« og »en Efterfølger?« En Efterfølger er eller stræber efter at være Det, han beundrer; en Beundrer holder sig personligt udenfor, bevidst eller ubevidst opdager han ikke, at det Beundrede indeholder en Fordring til ham, at være eller dog at stræbe efter at være det Beundrede.

\workheading{To Taler ved Altergangen om Fredagen}{1851}

\worknote{Antologiens ''Two Discourses at the Communion on Fridays''. Hong-udvalget bringer forordet — det berømte hvilested for forfatterskabet ''ved Alterets Fod'', med tanken om menneskelig ligelighed (enhver er lige nær Gud: elsket af ham) — samt hele første tale over Luc. 7,47 (''hvem Lidet forlades, elsker lidet''). Anden tale (1 Peter 4,8) er ikke medtaget.}

\sectionheading{Forord  ·  SV¹ XII, 267}

Forord.

En trinviis fremadskridende Forfatter-Virksomhed, der tog sin Begyndelse med »Enten – Eller«, søger her sit afgjørende Hvilepunkt, ved Alterets Fod, hvor Forfatteren, personligt sig sin Ufuldkommenhed og Skyld selv bedst bevidst, ingenlunde kaldersig et Sandheds-Vidne, men kun en egen Art Digter og Tænker, der, »uden Myndighed«, intet Nyt har havt at bringe, men »villet læse de individuelle, humane Existents-Forholds Urskrift, det Gamle, Bekjendte og fra Fædrene Overleverede igjennem endnu engang, om muligt paa en inderligere Maade« (jvf. min Efterskrift til »Afsluttende Efterskrift«).

Saaledes vendt har jeg her intet videre at tilføie. Dog lad mig blot udtale Dette, hvad der paa en Maade er mit Liv, mig mit Livs Indhold, dets Fylde, dets Lykke, dets Fred og Tilfredshed – Dette eller denne Livs-Anskuelse, der er Menneskeligheds og Menneske-Liighedens Tanke: christeligt er ethvert Menneske (den Enkelte), ubetinget ethvert Menneske, endnu en Gang, ubetinget ethvert Menneske Gud lige nær; og hvorledes nær og lige nær? elsket af ham. Altsaa Liigheden er der, den uendelige mellem Menneske og Menneske. Er der nogen Forskjel, o, denne Forskjel, hvis den er, den er som Fredsommeligheden selv; uforstyrret forstyrrer den ikke i fjerneste Maade Liigheden, Forskjellen er: at Een betænker, at han er elsket, maaskee Dag ud og Dag ind, maaskee i 70 AarDag ud og Dag ind, maaskee kun med een Længsel, efter Evigheden, for ret at kunne tage fat og fare fort, sysler med denne salige Beskjeftigelse, at betænke, at han – ak, og ikke for sin Dyds Skyld! – er elsket; en Anden tænker maaskee ikke paa, at han er elsket, maaskee gaaer Aar efter Aar, Dag efter Dag, og han tænker ikke paa, at han er elsket, eller han er maaskee glad og taknemlig for at være elsket af sin Hustru, af sine Børn, af Venner, af Medlevende, men han tænker ikke paa, at han er elsket af Gud, eller han sukker maaskee over ikke at være elsket af Nogen, og han tænker ikke paa, at han er elsket af Gud. »Dog«, saaledes maatte den Første vel sige, »jeg er uskyldig, jeg kan jo ikke for, om en Anden overseer eller forsmaaer den Kjerlighed, der paa ham ødsles lige saa rigeligt som paa mig«. Uendelige guddommelige Kjerlighed, som ingen Forskjel gjør! Ak – menneskelige Utaknemlighed! – dersomdet saa var Liigheden mellem os Mennesker, i hvilken vi ganske ligne hinanden, at Ingen af os ret tænker paa, at han er elsket!

Idet jeg derpaa vender mig til den anden Side, ønskede jeg, og vilde jeg tillade mig (takkende for hvad Deeltagelse og Velvillie der kan være blevet mig beviist) ligesom at overrække og anbefale Skrifterne til det Folk, hvis Sprog jeg med sønlig Hengivenhed og med næsten qvindelig Forelskelse er stolt af at have den Ære at skrive, dog ogsaa fortrøstende mig til, at det heller ikke skal have Skam af, at jeg har skrevet det.

Kjøbenhavn, i Eftersommeren 1851.

S. K.

\sectionheading{Luc. 7,47: Hvem Lidet forlades, elsker lidet  ·  SV¹ XII, 271–279}

I. Luc. VII, 47.

Bøn.

Herre Jesus Christus! Du som vistnok ikke kom til Verden for at dømme, Du var dog, ved at være Kjerligheden der ikke blev elsket, Dommen over Verden. Vi kalde os Christne, vi sige, at vi vide Ingen at gaae til uden til Dig – ak til hvem skulle vi saa gaae hen, naar det, just ved Din Kjerlighed, ogsaa over os bliver Dommen, at vi elske lidet? til hvem, o Trøstesløshed, hvis dog ikke til Dig; til hvem, o Fortvivlelse, hvis Du virkelig ikke vilde forbarmende antage Dig os, tilgivende vor store Synd mod Dig og mod Kjerligheden, vi som syndede meget ved at elske lidet!

Luc. VII, 47: Men hvem Lidet forlades elsker lidet.

A. T.Ved Alteret lyder jo Indbydelsen »kommer hid alle I, som arbeide og ere besværede, jeg vil give Eder Hvile«. Den Enkelte følger saa Indbydelsen; han gaaer op til Alteret. Derpaa vender han tilbage, forladende Alteret – da er der et andet Ord, det kunde være Indskriften over Kirkens Dør indvendig, ikke til at læse af dem, som gaae ind i Kirken, men kun til at læse af dem, som gaae ud af Kirken, dette Ord: hvem Lidet forlades, elsker lidet. Det første Ord er Alterets Indbydelse, det andet er Alterets Retfærdiggjørelse, som blev der sagt: hvis Du ved Alteret ikke fornam Dine Synders, hver Din Synds Forladelse, da ligger det i Dig selv, Alteret er uden Skyld, Skylden Din, fordi Du kun elsker lidet. O, som det er en vanskelig Sag i at bede ret at kunne komme til Amen – thi Den, der aldrig har bedet, ham synes det let nok, let nok at blive snart færdig; men Den, der følte Trang til at bede og begyndte at bede, hændes det vel, at det bestandigt er ham, som havde han noget Mere paa sit Hjerte, som kunde han ikke hverken faae sagt Alt, eller faae det Alt sagt som han ønskede det sagt, og saa kommer han ikke til Amen – saaledes er det ogsaa en vanskelig Sag, ret at annamme Syndernes Forladelse ved Alteret. Der tilsiges Dig alle Dine Synders naadefulde Forladelse. Hør det ret, tag det ganske bogstaveligen, alle Dine Synders Forladelse; Du skal fra Alteret kunne gaae bort, saa, gudelig forstaaet, let om Hjertet som et nyfødt Barn, paa hvem Intet, Intet tynger, altsaa endnu lettere om Hjertet, forsaavidt meget har tynget paa Dit Hjerte; der er ved Alteret Ingen, som beholder Dig end ikke den mindste af Dine Synder, Ingen – hvis Du ikke selv gjør det. Saa kast dem da alle fra Dig; og Erindringen om dem – at Du ikke i den beholder dem; og Erindringen om, at Du har kastet dem fra Dig – at Du ikke i den beholder dem hos Dig: kast det Alt fra Dig, Du har slet Intet at gjøre uden troende at kaste af Dig, og at kaste af Dig hvad der jo tynger og besværer. Hvad kan være lettere! Ellers er det det Tunge, at skulle tage Byrder paa sig; men at turde, at skulle kaste af sig! Og dog hvor vanskeligt! Ja, endnu sjeldnere end En der paatog sig alle Byrder, endnu sjeldnere En der fuldkommede den tilsyneladende saa lette Opgave, efter at have modtaget Forsikkringen om sine Synders naadefulde Forladelse og Pantet derpaa, at føle sig ganske lettet for enhver – end den mindste Synd, eller for enhver – ogsaa den største Synd! Hvis Du kunde see i Hjerterne, Du vilde vistnok see, hvorledes Mangen, besværet, sukkende under den tunge Byrde, gaaer op til Alteret; og naar man saa gaaer derfra, hvis Du kunde see i Hjerterne, Du vilde muligt see, at der dog i Grunden ikke var en Eneste, der gik ganske lettet derfra, og Du vilde maaskee stundom see, at der var Den, der gik end mere besværet bort, besværet ved Tanken om, at han vel ikke havde været en værdig Gjest ved Alteret, siden han ingen Lindring fandt.

At dette er saa, ville vi ikke fordølge for hinanden, vi ville ikke tale saaledes, at Talen lader uvidende om, hvorledes det gaaer til i Virkeligheden, fremstiller Alt saa fuldkomment, at det slet ikke passer paa os virkelige Mennesker. O, nei, hvortil skulde saa Talen hjælpe! Men derimod, naar Talen gjør os saa ufuldkomne som vi ere, da hjælper den os til at bevares i en stadig Stræben, hverken berusede i Drømme at indbilde os, at Alt var afgjort med denne ene Gang, eller i stille Mismod at forsage, fordi det denne Gang ikke lykkedes os efter Ønske, ikke skeete, som vi havde bedet og begjeret.

Lader os da i de foreskrevne korte Øieblikke betænke dette Ord: men hvem Lidet forlades, elsker lidet, et Dommens Ord, men ogsaa et Trøstens Ord.

Og Du, m. T., lad Dig ikke forstyrre af, at jegtaler saaledes nu i dette Øieblik, inden Du gaaer op til Alteret, maaskee formenende og fordrende, at Den, der skal tale i dette Øieblik, han burde tale paa en anden Maade, anvende Alt for at berolige den Enkelte og gjøre ham tryg; hvis han da bagefter erfarede, at den hellige Handling dog ikke fuldkomment havde været til Glæde og Velsignelse for en Enkelt, da kunde han jo tale til ham paa en anden Maade. O, min Ven, deels er det sandeligen ikke saa, at det kun er en Enkelt hvem det Fuldkomne ikke lykkes, nei, det er kun en Enkelt hvem det Fuldkomne lykkes; deels er der en Bekymring, en inderlig Bekymring, som maaskee understøtter bedre til at det Høieste maa lykkes et Menneske, bedre end den for store Tillid og den for sorgløse Frimodighed. Der er en Længsel efter Gud, en Tillid til Gud, en Fortrøstning, en Haaben paa Gud, en Kjerlighed, en Frimodighed: men hvad der sikkrest finder ham er dog maaskee Sorgen efter Gud; Sorgen efter Gud – det er ingen flygtig Stemning, der strax forsvinder, idet man nærmer sig Gud, tværtimod den er maaskee dybest, just naar den kommer Gud nærmest, som den saaledes Sørgende er meest bange for sig selv, jo nærmere han kommer Gud.

Hvem Lidet forlades, elsker lidet. Det er Dommens Ord.

I Almindelighed fremstilles det vel saaledes: Retfærdigheden, det er den strenge Dom; Kjerligheden er den Milde, som ikke dømmer, og dømmer den, Kjerlighedens Dom er den milde Dom. Nei, nei, Kjerlighedens Dom er den strengeste Dom. Var ikke den strengeste Dom, der er gaaet over Verden, strengere end Syndfloden, strengere end Babels Forvirring, end Sodoma og Ghomorras Undergang, var den ikke Christi uskyldige Død, der dog var Kjerlighedens Offer? Og hvilken var Dommen? dog vel den: »Kjerligheden« blev ikke elsket. Saaledes ogsaa her. Dommens Ord lyder ikke: hvem Lidet forlades, han syndede meget, saa altsaa hans Synder vare for store og for mange, til at de kunne tilgives; nei, Dommen lyder: han elsker lidet. Det er altsaa ikke Retfærdigheden, der strengt negter Syndernes Tilgivelse og Forladelse; det er Kjerligheden, der mildt og forbarmende siger: jeg tilgiver Dig Alt – forlades Dig kun Lidet, saa er det fordi Du kun elsker lidet. Retfærdigheden sætter strengt Grændsen og siger: ikke videre, nu er Maalet fyldt, for Dig er ingen Forladelse; men derved bliver det ogsaa. Kjerligheden siger: Alt er Dig tilgivet – forlades Dig kun Lidet, da er det fordi Du kun elsker lidet; saa der altsaa kommer en ny Synd til, en ny Skyld, den at forskylde at kun Lidet tilgives, at forskylde det ikke ved de begangne Synder, men ved Mangelen paa Kjerlighed. Vil Du lære at frygte, da lær at frygte, ikke Retfærdighedens Strenghed, men Kjerlighedens Mildhed! – Retfærdigheden seer dømmende paa et Menneske, og Synderen kan ikke udholde dens Blik; men Kjerligheden, naar den seer paa ham, ja om han end unddrager sig dens Blik, slaaer Øiet ned, han fornemmer dog alligevel, at den seer paa ham; thi Kjerligheden trænger langt inderligere ind paa Livet, ind ved Livet, ind der, hvorfra Livet udgaaer, end Retfærdigheden, som frastødende befæster det svælgende Dyb mellem Synderen og sig, medens Kjerligheden jo staaer paa hans Side, anklager ikke, dømmer ikke, tilgiver og forlader. Retfærdighedens dømmende Stemme kan Synderen ikke udholde, han søger, om muligt, at lukke sit Øre for den; men om han end vilde, det er ham umuligt ikke at høre Kjerligheden, hvis Dom, o frygtelige Dom, er: Dine Synder ere Dig forladte! Frygtelige Dom, hvis Ord dog i og for sig ere Intet mindre end det Forfærdelige; og just derfor kan Synderen ikke lade være at høre, hvad der dog er Dommen. Hvor skal jeg henflye fra Retfærdigheden, tager jeg Morgenrødens Vinger og flyver til det yderste Hav, da er den der, og skjuler jeg mig i Afgrunden, da er den der, og saaledes paa ethvert Sted, dog nei, der er eet Sted, hvor jeg kan flye hen: til Kjerligheden; men naar Kjerligheden dømmer Dig, og Dommen er – o, Rædsel! – Dommen er: Dine Synder ere Dig forladte! Dine Synder ere Dig forladte – og dog er der Noget (ja dette Noget er i Dig, hvor i al Verden skulde det ellers finde sit Tilholdssted, naar Kjerligheden tilgiver Alt!), der er Noget i Dig, som gjør, at Du fornemmer, at de ere Dig ikke forladte. Hvad er saa den strengeste Doms Rædsel mod denne Rædsel! Hvad er saa Vredens strenge Dom, Forbandelsen, mod denne Dom: Dine Synder ere Dig forladte! Saa er jo næsten Retfærdigheden Mildhed, der siger, som Du siger: nei, de ere Dig ikke forladte! Hvad er »Brodermorderens« Lidelse, naar han flygtig og ustadig frygter, at Enhver skal kjende ham paa Retfærdighedens Mærke, der dømte ham – hvad er denne Lidelse mod den Ulyksaliges Qvaler, hvem det blev Dommen, ikke Frelsen: Dine Synder ere Dig forladte! Frygtelige Strenghed! At Kjerligheden, at det er Kjerligheden, den tilgivende Kjerlighed, der, ikke dømmende, nei, ak selv lidende derved, dog forvandles til Dommen; at Kjerligheden, den tilgivende Kjerlighed, der ikke vil, som Retfærdigheden, gjøre Skylden aabenbar, tværtimod vil skjule den ved at tilgive og forlade, at det dog er den, der, ak selv lidende derved,gjør Skylden frygteligere aabenbar end Retfærdigheden! – Tænk denne Tanke: det Selvforskyldte. Det er selvforskyldt, siger Retfærdigheden, at der for det Menneske ingen Tilgivelse er; den tænker derved paa hans mange Synder, thi Retfærdigheden kan Intet glemme. Kjerligheden siger: det er selvforskyldt – den tænker derved ikke paa hans mange Synder, o nei, den er villig til at glemme dem alle, den har glemt dem alle; og dog, det er selvforskyldt, siger Kjerligheden. Hvilket er det Forfærdeligste? dog vel det Sidste, der jo næsten lyder som afsindig Tale; thi Anklagen er ikke hans Synder, nei, Anklagen er: det er ham tilgivet, det er ham Alt tilgivet. Tænk en Synder, der synker i Afgrunden, hør dette Angestraab, naar han i sidste Suk lader Retfærdigheden, hvilken hans Liv spottede, vederfares Ret og siger: det er selvforskyldt. Forfærdeligt! Der er kun eet Forfærdeligere, hvis det ikke er Retfærdigheden han taler til, men Kjerligheden og siger: det er selvforskyldt.Lader Retfærdigheden sig ikke spotte, sandeligen endnu mindre Kjerligheden. Strengere end Retfærdighedens strengeste Dom over den største Synder er Kjerlighedens: ham forlades Lidet – fordi han kun elsker lidet.

Hvem Lidet forlades, elsker lidet. Det er Dommens Ord, men ogsaa et Trøstens Ord.

Ikke veed jeg, m. T., hvad Du forbrød, hvilken Din Skyld, Dine Synder ere; men den Skyld bleve vi dog nok alle, mere eller mindre, skyldige i: kun at elske lidet. Saa trøst Dig da ved Ordet, som jeg trøster mig ved det. Og hvorledes trøster jeg mig? Jeg trøster mig ved, at Ordet jo Intet udsiger om den guddommelige Kjerlighed, men kun Noget om min. Ordet siger ikke, at nu er den guddommelige Kjerlighed blevet træt af at være Kjerlighed, at den nu har forandret sig, træt af i ubeskrivelig Forbarmelse ligesom at ødsles paaMenneskenes utaknemlige Slægt eller paa mig Utaknemlige, at den nu er blevet noget Andet, en mindre Kjerlighed, dens Varme afkjølet, fordi Kjerligheden blev kold i Menneskenes utaknemlige Slægt eller i mig Utaknemlige. Nei, derom taler Ordet slet ikke. Og trøst Dig, som jeg trøster mig – hvorved? Derved, at Grunden, hvorfor Ordet ikke siger det, er, at det hellige Ord ikke lyver, saa det da ikke tilfældigviis eller grusomt er blevet fortiet i Ordet, medensdet dog i Virkeligheden er saa, at Guds Kjerlighed er blevet træt af at elske. Nei, naar Ordet ikke siger det, saa er det heller ikke saa; og om Ordet sagde det – nei Ordet kunde ikke sige det, thi Ordet kan ikke lyve. O, i dybeste Sorg saligste Trøst! Hvis Guds Kjerlighed i Sandhed havde forandret sig; hvis Du Intet havde hørt derom, men bekymret over Dig selv, at Du hidtil kun havde elsket lidet, med fromt Forsæt havde stræbt at bringe Kjerligheden i Dig til at flamme, og ved det Samme, ved hvilket Du bragte den i Flamme, næret Flammen – og nu, om end beskæmmet i Følelsen af, hvor ufuldkommen dog Din Kjerlighed var, nu vilde Du nærme Dig Gud, for, efter Skriftens Ord, at forlige Dig med ham: men han havde forandret sig! Tænk Dig en elskende Pige; bekymret tilstaaer hun sig selv, hvor lidet hun dog hidtil har elsket – nu saa vil jeg da ogsaa, siger hun til sig selv, blive lutter Kjerlighed. Og det lykkes hende; disse Bekymringens Taarer, hun fælder i Sorg over sig selv, disse Taarer – de slukke ikke Ilden, nei dertil ere de for brændende, o nei, just disse Taarer bringer Ilden til at flamme: men imidlertid havde den Elskede forandret sig, han var ikke mere den Kjerlige. O, een Bekymring for et Menneske! O, for et Menneske kan een Bekymring være nok, mere kan intet Menneske bære; skulde et Menneske, idet han i Selvbekymring maatte tilstaae sig selv, hvor lidet han hidtil havde elsket Gud, ængstes af den Tanke, om Gud nu ikke imidlertid kunde have forandret sig: saa, ja saa vil jeg fortvivle, og jeg vil fortvivle strax, thi saa er der Intet mere at bie paa, hverken i Tid eller i Evighed. Men derfor trøster jeg mig ved Ordet; og jeg spærrer mig enhver Udflugt, og jeg skaffer alle Undskyldninger til Side og alle Besmykkelser, og blotter Brystet, hvor jeg skal saares af Ordet, der dømmende trænger ind, dømmende »Du elskede kun lidet.« O, træng kun dybere ind, end dybere, Du helbredende Smerte »Du elskede slet ikke« – selv naar Dommen lyder saaledes, jeg fornemmer i een Forstand ingen Smerte, jeg fornemmer en ubeskrivelig Salighed; thi just min Dom, Dødsdommen over mig og min usle Kjerlighed, den indeholder tillige noget Andet: Gud er uforandret Kjerlighed.

Saaledes trøster jeg mig. Og jeg finder skjult en Trøst i Ordet, som vel ogsaa Du finder, m. T., just naar Du hører Ordet saaledes, at det saarer Dig. Thi der staaer ikke, hvem Lidet forlades, »elskede«lidet, nei der staaer: elsker lidet. O, naar Retfærdigheden dømmer, den gjør Regnskabet op, den slutter af, den bruger den forbigangne Tid, den siger: han elskede lidet, og siger dermed, at nu er den Sag for stedse afgjort, vi To ere skilte ad, har ikke mere med hinanden at skaffe. Ordet, Kjerlighedens Ord derimod lyder: hvem Lidet forlades, han elsker lidet. Han elsker lidet; han elsker, det vil sige: saaledes er det nu, nu i dette Øieblik – mere siger Kjerligheden ikke; uendelige Kjerlighed, at Du saaledes bliver Dig selv tro i end Din mindste Yttring! Han elsker lidet nu, i dette Nu. Men hvad er Nuet, og hvad er Øieblikket – hurtigt, hurtigt er det forbi; og nu, i det næste Øieblik, nu er Alt forandret, nu elsker han, om ikke meget, han stræber dog at elske meget; nu er Alt forandret, kun ikke »Kjerligheden«, den er uforandret, uforandret den Samme, der kjerligt har biet paa ham, kjerligt ikke nænt at lade det være afgjort med ham, ikke nænt at søge Skilsmissen fra ham, men er blevet hos ham, og nu er det ikke Retfærdigheden, der afsluttende siger: han elskede lidet, nu er det Kjerligheden, der glad i Himlensiger: han elskede lidet, det er, nu er det forandret, det var den Gang, nu, nu elsker han meget.

Men er det saa dog i Grunden ikke saa, at Syndens Forladelse fortjenes, vel ikke ved Gjerninger, men ved Kjerlighed; naar det siges, at hvem Lidet forlades, elsker lidet, ligger saa ikke deri, at det er Kjerligheden, der gjør Udslaget, om Ens Synder, og hvorvidt Ens Synder skulle tilgives – og saa fortjenesjo Syndernes Forladelse? O, nei. I det samme Evangelium noget tidligere (V. 42 Slutning) taler Christus om to Skyldnere, af hvilke den Ene var Meget, den Anden Lidet skyldig, og som begge fandt Tilgivelse – han spørger: hvilken af disse To bør elske meest? Og Svaret er: Den, hvem Meget blev tilgivet. Agt nu paa, hvorledes vi dog ikke komme ind i det Fortjenstliges usalige Egne, men hvorledes Alt bliver indenfor Kjerligheden! Naar Du elsker meget, forlades Dig Meget – og naar Dig Meget forlades, elsker Du meget. See her den salige, Frelsens Tilbagevenden i Kjerligheden! Først elsker Du meget, og Meget bliver Dig saa forladt – o, og see, saa tager Kjerligheden endnu stærkere paa, dette, at saa Meget blev Dig forladt, det fremelsker endnu engang Kjerligheden, og Du elsker meget, fordi Meget blev Digforladt! Det er med Kjerligheden her som med Troen. Tænk Dig En af hine Ulykkelige, hvem Christus helbredede ved et Mirakel.For at skulle helbredes maa han troe – nu troer han, og bliver helbredet. Nu er han helbredet – og nu bliver Troen end een Gang stærkere, nu da han er frelst. Det er ikke saaledes: han troede, og saa skeete Miraklet, og saa var det forbi; nei Opfyldelsen forøger ham end een Gang Troen, efter Opfyldelsen er hans Tro een Gang saa stærk, som da han troede, førend han blev frelst. Og saaledes med det at elske meget. Stærk er den Kjerlighed,guddommelig stærk i Svaghed, den Kjerlighed, som elsker meget, og hvilken saa Meget forlades; men endnu stærkere er Kjerlighedens anden Gang, naar den samme Kjerlighed anden Gang elsker, elsker, fordi Meget blev forladt.

M. T. Du erindrer vel Begyndelsen af denne Tale. Man kan i det høitidelige Øieblik forstyrre paa tvende Maader: ved at tale om det Uvedkommende, selv om dette forresten er baade vigtigt og Talen betydningsfuld; eller ved at tale forstyrrende om Det, som i et saadant Øieblik er En det nærmeste. »Hvem Lidet forlades, elsker lidet« – det kunde synes forstyrrende just i dette Øieblik, forinden Du nu gaaer op til Alteret, hvor Du jo annammer alle Dine Synders Forladelse. O, men som det Opbyggelige i sit Første altid er det Forfærdende, og som al sand Kjerlighed i sit Første altid er Uro, og som Kjerlighed til Gud i sit Første altid er Sorg: saaledes er hvad der synes det Forstyrrende ikke altid det Forstyrrende; det i Sandhed Beroligende er altid i sit Første det Foruroligende. Men er der nogen Sammenligning mellem disse tvende Farer: den, i bedragersk Tryghed at være beroliget, og den, at foruroliges ved at blive mindet om det Foruroligende – nu om hvilket Foruroligende? om det Foruroligende: at ogsaa det kan tilgives, hvis man hidtil kun elskede lidet. Det er selsomtmed det Foruroligende; Den, der ret dannedes ved dette, det er sandt, han seer ikke saa stærk ud som Den, der blev uvidende derom. Men i det sidste Øieblik da er han, just ved Afmagten, dog maaskee den Stærkeste, i det sidste Øieblik lykkes, just ved Afmagten, maaskee ham, hvad der ikke lykkes den Stærkeste.

Saa velsigne da Gud denne foruroligende Tale, at den kun til det Gode maatte have foruroliget Dig, at Du ved Alteret beroliget maatte fornemme, at Du annammer alle Dine Synders naadefulde Forladelse.

\workheading{Til Selvprøvelse, Samtiden anbefalet}{1851}

\worknote{Antologiens ''For Self-Examination''. Hong-udvalget tager hovedsagelig af Første Deel (over Jakobs Brev 1,22 — at være Ordets Gjørere og ikke blot dets Hørere; Guds Ord som Speilet; Nathans ''Du est Manden''), med forordet (''læs høit, om muligt''), Forbemærkningen om Socrates, og til slut billedet om rigmandens heste og den manglende kudsk. Første Deel gengives i sammenhæng til og med ''Du est Manden'' og overgangen til det Subjektive; derefter springes (markeret med […]) til slutbilledet om rigmandens heste. Anden og Tredie Deel er ellers ikke medtaget.}

\sectionheading{Forord; Forbemærkning; Første Deel (uddrag)  ·  SV¹ XII, 295–370}

Forord.

Min kjere Læser! læs, om muligt, høit! Gjør Du det, lad mig takke Dig derfor; gjør Du det ikke blot selv, bevæger Du ogsaa Andre dertil, lad mig takke dem hver især, og Dig atter og atter! Du vil ved at læse høit stærkest faae Indtrykket af, at Du har ene med Dig selv at gjøre, ikke med mig, der jo er »uden Myndighed«, ei heller med Andre, hvilket vilde være Adspredelse.

En Forbemærkning.

Der er et Ord, som oftere er faldet mig paa Sinde, et Ord af en Mand, hvem jeg vistnok ikke christeligt kan siges at skylde Noget, det var jo en Hedning, men hvem jeg dog mener personligt at skylde saare meget, han som ogsaa levede under Forhold, der, i mine Tanker, ganske svare til vore Tiders Vilkaar: jeg mener Oldtidens eenfoldige Vise. Om ham fortælles der, at da han var anklaget for Folket, kom en Taler til ham, som overrakte ham en omhyggelig udarbeidet Forsvars-Tale med Anmodning om at afbenytte den. Den eenfoldige Vise modtog den, læste den. Derpaa gav han Taleren den igjen tilbage, og sagde: det er en smuk og veludarbeidet Tale (altsaa det var ikke, fordi Talen var uheldig, var slet, at han gav den tilbage), men, vedblev han, jeg er nu blevet halvfjerdsindstyve Aar gammel, saa synes jeg, det ikke sømmer sig for mig at benytte en Talers Kunst. Hvad meente han? Han meente for det Første: mit Liv er for alvorligt til at kunne være tjent med at understøttes af en Talers Kunst; jeg har sat et Liv ind; selv om jeg ikke bliver dømt fra Livet, jeg har dog sat et Liv ind, og i Gudens Tjeneste har jeg røgtet mit Ærinde: lad mig da nu ikke i det sidste Øieblik tilintetgjøre Indtrykket af mig selv og mit Liv ved Hjælp af Kunst-Talere eller Tale-Kunster. Dernæst meente han. De Tanker, Forestillinger, Begreber jeg nu i tyve Aar (saa længe var det) kjendt af Enhver, latterliggjort af Eders comiske Digtere, anseet for en Særling, ideligt angreben af »Navnløse« (det er hans Ord) har udviklet i Samtale med den Første den Bedste paa Torvet – disse Tanker ere mit Liv, have beskjeftiget mig tidligt og sildigt, have de ingen Anden beskjeftiget, mig have de beskjeftiget uendelig, og naar jeg, hvad der har tildraget sig Eders særdeles Opmærksomhed, stundom har kunnet staae en heel Dag og stirre hen for mig, det var disse Tanker, der beskjeftigede mig: saa tænker jeg ogsaa, at jeg, uden Hjælp af Kunst-Talere og Tale-Kunster, hvis jeg paa Doms-Dagen overhovedet agter at sige Noget, kan sige et Par Ord; thi den Omstændighed, at jeg formodentligen vil blive dømt fra Livet, gjør da hverken fra eller til, hvad jeg siger vil naturligviis blive det Samme og om det Samme og paa den samme Maade som hidtil, og som jeg nu igaar talte med en Feldbereder paa Torvet: det Par Ord tænker jeg kan jeg nok sige uden Forberedelse eller Andres Bistand; det forstaaer sig, ganske uden Forberedelse er jeg heller ikke, da jeg har forberedt mig i tyve Aar, og ganske uden Bistand er jeg heller ikke, da jeg regner paa Gudens Bistand. Men, som sagt, det Par Ord... ja, det negter jeg ikke, det kan ogsaa blive vidtløftigere; hvis jeg fik tyve Aar igjen, vilde jeg kunne vedblive at tale om det Samme, jeg bestandigt har talet om; men Kunst-Talere og Tale-Kunster er ikke Noget for mig. – O, Du Alvorlige! Miskjendt maatte Du tømme Giftbægeret; Du blev ikke forstaaet. Saa døde Du. I over to tusinde Aar er Du blevet beundret »men om jeg er blevet forstaaet?« det er sandt!

Og nu det at prædike! burde det ikke ogsaa være saaledes alvorligt! Den, der skal prædike, han skal leve i de christelige Tanker og Forestillinger, det skal være hans daglige Liv – naar saa er, det er Christendommens Mening, saa vil Du ogsaa have Veltalenhed nok, og just den, der fordres, naar Du taler frisk væk uden særlig Forberedelse. Usand Veltalenhed er det derimod, hvis Nogen, uden selv forresten at beskjeftige sig med, uden at leve i disse Tanker, engang imellem først sidder og sanker møisommeligt saadanne Tanker, maaskee paa Literaturens Mark, derpaa arbeider dem sammen til en vel udarbeidet Tale, som saa læres ypperligt udenad, holdes fortræffeligt baade i Henseende til Stemme og Foredrag og i Henseende til Armbevægelser. Nei, som man i de velindrettede Huse ikke behøver at stige ned ad Trapperne for at hente Vand, men ved et Høitryk har det deroppe, blot man omdreier Hanen, saaledes er Den en ægte christelig Taler, der, fordi det Christelige er hans Liv, i ethvert Øieblik har Veltalenheden, just den sande Veltalenhed, nærværende, lige ved Haanden – dog følger det naturligviis af sig selv, at det da ikke er Meningen hermed at anvise Vrøvlehovederne Hæderspladsen, om det end nok saa meget er vist, at det er uden Forberedelse de vrøvle. Videre, Skriften siger: I skulle aldeles ikke sværge, Eders Tale være Ja og Nei, hvad derover er, er af det Onde: saaledes er der ogsaa en Veltalenhedens Kunst som er af det Onde, hvis den gjøres til det Høieste, da den er det Lavere. Thi Prædikenen skal ikke splidagtigt befæste Adskillelsen mellem de Begavede og de Ikke-Begavede, den skal, i Hellig-Aands Enighed, fæste Opmærksomheden ene og alene paa, at der skal gjøres efter det Sagte. Du Eenfoldige, og var Du end den af Alle meest Indskrænkede – dersom Dit Liv udtrykker det Lidet Du har forstaaet: Du taler mægtigere end alle Taleres Veltalenhed! Og Du, o Qvinde, om Du end ganske forstummer i yndig Taushed – dersom Dit Liv udtrykker, hvad Du hørte: Din Veltalenhed er mægtigere, sandere, mere overbevisende end alle Taleres Kunst!

Dette er saaledes. Men lader os see os for, at vi ikke gribe for høit; thi fordi det er sandt, deraf følger ikke, at vi formaae at gjøre det. Og Du, min Tilhører, Du vil betænke, at jo høiere det Religieuse tages, jo strengere bliver det; men deraf følger ikke, at Du kan bære det, Dig vilde det maaskee endog være til Forargelse og til Fordærvelse. Maaskee kan Du nok behøve denne lavere Form af det Religieuse, at der anvendes en vis Kunst paa Fremstillingen for at gjøre det mere tiltrækkende. Den strenge Religieuse, hans Liv er væsentligen Handling – og hans Fremstilling ganske anderledes nærgaaende og fastende end den mere omhyggeligt udarbeidede Tale. Min Tilhører, er Du af denne Mening, saa tag imod dette og læs det til Opbyggelse. Det er ikke for min Fuldkommenheds Skyld og ikke for Din Fuldkommenheds Skyld, at Talen er saa udarbeidet, tværtimod det er, gudelig forstaaet, en Ufuldkommenhed og Svaghed. Jeg tilstaaer, og endog Dig, min, ikke sandt, saa vil Du ogsaa tilstaae, ikke mig, nei det forlanges ikke, men Dig selv og Gud Din. Ak, vi som dog kalde os Christne, vi ere, christeligt forstaaet, saa forkjælede, saa langt fra at være, hvad Christendommen dog fordrer af dem, som ville kalde sig Christne, afdøde fra Verden, neppe have vi vel engang Forestilling om den Art Alvor; vi kunne endnu ikke undvære, forsage det Kunstneriske og dets Formildelse, ikke taale det sande Indtryk af Virkelighed: nu, saa lader os da idetmindste være oprigtige og tilstaae det. Forstaaer Nogen maaskee ikke strax, hvad jeg her siger og tilsigter: han være langsom til at dømme, han give Tid, vi skulle nok komme nærmere ind paa den Sag. O, men hvo Du end er, hav Tillid, giv Dig hen; Magt kan der jo ikke være Tale om at jeg skulde kunne bruge, jeg den af Alle Afmægtigste, men der skal ikke blive brugt den mindste Overtalelse eller Underfundighed eller List eller Lokken, for at føre Dig maaskee saa langt ud, at Du maatte (hvad Du dog vistnok alligevel egentligen ikke burde, og, hvis Din Tro var større, heller ikke vilde) fortryde, at Du havde givet Dig hen; tro mig (jeg siger det til min egen Beskæmmelse) ogsaa jeg er kun altfor forkjælet.

Bøn.

Fader i Himlene! Hvad er et Menneske, at Du kommer ham ihu, et Menneskes Barn, at Du bekymrer Dig om ham – og paa enhver Maade, i enhver Henseende! Sandeligen i Intet lod Du Dig uden Vidnesbyrd; og tilsidst gav Du ham Dit Ord. Mere kunde Du ikke gjøre; at tvinge ham til at benytte det, læse eller høre det, at tvinge ham til at gjøre derefter, det kunde Du ikke ville. O, og dog Du gjør mere. Thi ikke er Du som et Menneske; han gjør sjeldent Noget for Intet, men gjør han det for Intet, saa vil han da idetmindste ikke have Uleilighed deraf. Derimod Du, o Gud, Du giver Dit Ord som en Gave, det gjør Du, uendeligt Ophøiede – og vi Mennesker have Intet at give til Vederlag. Og finder Du da blot nogen Villighed hos den Enkelte, da er Du flux tilrede, og er først Den, der, med mere end menneskelig ja med guddommelig Taalmodighed, sidder og staver med den Enkelte, at han maa kunne forstaae Ordet retteligen; og saa er Du dernæst Den, der, atter med mere end menneskelig ja med guddommelig Taalmodighed, ligesom tager ham ved Haanden og hjælper ham, naar han stræber at gjøre derefter – Du vor Fader i Himlene!

Tiderne ere forskjellige; og om det end ofte gaaer med Tiderne som med et Menneske: han forandrer sig ganske – men det bliver lige galt, blot i et nyt Mønster: det er dog ligefuldt sandt, at Tiderne ere forskjellige; og forskjellige Tider kræve det Forskjellige.

Der var en Tid, da Evangeliet, »Naaden«, var blevet forvandlet til en ny Lov, strengere over Menneskene end den gamle. Alt var blevet noget forpiint, trælsomt og ulysteligt Noget, næsten som var der – trods Englenes Sang ved Christendommens Indkomst – ingen Glæde mere hverken i Himlen eller paa Jorden. I smaaligt Selvplageri havde man – saaledes hævner det sig! – gjort Gud lige saa smaalig. Man gik i Kloster, man blev der – ja, det er sandt; det var frivilligt, og dog var det Trældom, thi det var ikke i Sandhed frivilligt, man var ikke ganske enig med sig selv, ikke glad ved at være der, ikke fri, dog havde man heller ikke Frimodighed til at lade det være, eller til at forlade Klosteret igjen og blive fri. Alt var blevet Gjerning. Og som usunde Udvæxter paa Træerne, saaledes vare disse Gjerninger bedærvede af usunde Udvæxter, saa ofte kun Hykleri, Fortjenstligheds Indbildskhed, Ørkesløshed. Just deri var det Feilen stak, ikke saa meget i Gjerningerne. Thi lader os ikke overdrive, ikke benytte en forbigangen Tids Vildfarelse til ny Vildfarelse. Nei, tag denne Usundhed og Usandhed bort fra Gjerningerne, og lad os saa kun beholde Gjerningerne i Oprigtighed, i Ydmyghed, i gavnlig Virksomhed. Saaledes skulde det nemlig være med disse Gjerninger, som naar for Exempel en krigersk Yngling i Forhold til det farefulde Foretagende selv frivilligt kommer og beder den Øverste, sigende: o, maa jeg ikke faae Lov at komme med der! Hvis saaledes et Menneske vilde sige til Gud: »o, maa jeg ikke faae Lov at give alt mit Gods til de Fattige; at det skulde være noget Fortjenstligt, nei, nei, jeg erkjender dybt ydmyget, at hvis jeg engang bliver salig, jeg bliver det af Naade aldeles som Røveren paa Korset, men maa jeg ikke faae Lov til at gjøre det, at jeg ganske kan virke for Guds Riges Udbredelse blandt mine Medmennesker« – saa, ja saa er dette, for at jeg skal tale lutherisk, det er trods Satan, Bladene, det høistærede Publikum (thi Pavens Tid er nu forbi), trods alle kloge Mænds og Koners forstandige, geistlige eller verdslige, Indsigelser, det er Gud velbehageligt. Men saaledes var det ikke den Tid, vi tale om.

Da fremtraadte der en Mand, Morten Luther, fra Gud og med Troen; med Troen (thi sandeligen hertil behøvedes Tro) eller ved Tro indsatte han Troen i dens Rettigheder. Hans Liv udtrykte Gjerningerne, lad os aldrig glemme det, men han sagde: et Menneske frelses alene ved Troen. Faren var stor. Hvor stor den var i Luthers Øine, jeg veed intet stærkere Udtryk derfor, end at han besluttede: for at faae Orden i Sagen, maa Apostelen Jacob skubbes tilside. Tænk Dig en Luthers Ærbødighed for en Apostel – og saa at maatte vove dette for at faae Troen indsat i dens Ret!

Imidlertid hvad skeete? Der er bestandig en Verdslighed, som nok vil have Navn af at være Christen, men ønsker at blive det for saa billig Priis som muligt. Denne Verdslighed blev opmærksom paa Luther. Den hørte, af Forsigtighed hørte den endnu en Gang, at den da ikke skulde have hørt feil, derpaa sagde den »ypperligt, det er Noget for os; Luther siger: det er alene Troen det kommer an paa; at hans Liv udtrykker Gjerningerne, det siger han ikke selv, og da han nu er død, saa er det ikke mere en Virkelighed. Altsaa tage vi hans Ord, hans Lære – – og vi ere frie for alle Gjerninger, Leve Luther: wer nicht liebt Weiber, Wein, Gesang, Er wird ein Narr sein Leben lang. Dette er Betydningen af Luthers Liv, denne Guds Mand, der, tidssvarende, reformerede Christendommen«. Og om end ikke Alle slet saa verdsligt toge Luther forfængeligt – der er i ethvert Menneske en Tilbøielighed til, enten, naar Gjerningerne skal med, at ville være fortjenstlig, eller, naar Troen og Naaden skal gjøres gjældende, saa da ogsaa at ville være saavidt muligt ganske fri for Gjerningerne. »Mennesket«, denne fornuftige Guds Skabning, lader sig sandeligen ikke narre, er ikke en Bonde, der kommer tiltorvs, han seer sig for. »Nei, Eet af To«, siger Mennesket, »skal det være Gjerninger: vel, men maa jeg saa ogsaa bede om den mig lovligt tilkommende Profit af mine Gjerninger, at det er Fortjenstligt. Skal det være Naade: vel, men maa jeg saa ogsaa bede, at jeg bliver fri for Gjerningerne, ellers er det jo ikke en Naade. Skal det være Gjerninger og alligevel Naade, det er jo Galskab.« Ja vist er det Galskab, det var det sande Lutherske ogsaa, det var jo Christendom. Christendommens Fordring er: Dit Liv skulde saa anstrenget som muligt udtrykke Gjerningerne; saa fordres der Eet endnu, at Du ydmyger Dig og tilstaaer: men det er alligevel Naade, at jeg frelses. Man afskyede Middelalderens Vildfarelse: det Fortjenstlige. Naar man seer dybere i Sagen, vil man let see, at man havde en maaskee endnu større Forestilling om, at Gjerninger ere fortjenstlige, end Middelalderen havde; men man anbragte »Naaden« saaledes, at man fritog sig for Gjerningerne. Naar man saa har afskaffet Gjerningerne, kunde man dog ikke godt fristes til at ansee de Gjerninger, man ikke gjør, for noget Fortjenstligt. Luther vilde tage »Fortjenstligheden« bort fra Gjerningerne, og anbringe dem noget anderledes nemlig i Retning af at vidne for Sandhed; Verdsligheden, som forstod Luther fra Grunden, tog Fortjenstligheden ganske bort – samt Gjerningerne.

Og hvor ere vi nu? Jeg er »uden Myndighed«; langt fra mig at dømme nogen Eneste. Men da jeg dog ønsker denne Sag oplyst, saa vil jeg tage mig selv og et Øieblik prøve mit Liv efter blot een lutherisk Bestemmelse af Troen: Troen er en urolig Ting. Jeg antager da, at Luther er staaet op af sin Grav; allerede i flere Aar har han levet iblandt os men ukjendt, agtet paa det Liv, der føres af os, været opmærksom paa alle Andre saaledes ogsaa paa mig. Jeg antager, at han nu en Dag taler til mig og siger »er Du en Troende, har Du Troen?« Enhver, der kjender mig som Forfatter, vil see, at jeg endda maaskee var Den, der af Alle maatte komme bedst fra en saadan Examination; thi jeg har jo bestandigt sagt: »jeg har ikke Troen« – som en Fugls angstfulde Flugt foran et Uveir, saaledes har jeg udtrykt, her er Uraad, »jeg har ikke Troen«. Dette kunde jeg altsaa sige til Luther, sige til ham: »nei, kjere Luther, jeg har dog viist den Ærbødighed at sige: jeg har ikke Troen«. Dog dette vil jeg ikke gjøre gjældende; men som alle Andre kalde sig Christne og Troende, saa vil jeg ogsaa sige: »jeg er en Troende«, thi ellers faaer jeg jo ikke oplyst, hvad jeg ønsker oplyst. Altsaa jeg svarer: »jo, jeg er en Troende«. »Hvorledes« svarer Luther, »da har jeg ikke mærket Noget paa Dig, og jeg har dog agtet paa dit Liv; og Du veed, Troen er en urolig Ting. Hvortil har Troen, Du siger Du har, uroliget Dig, hvor har Du vidnet for Sandhed, og hvor mod Usandhed, hvilke Offre har Du bragt, hvad har Du lidt af Forfølgelse for Din Christendom; og hjemme i Dit huuslige Liv, hvor har Din Selvfornegtelse og Forsagelse været til at mærke?« »Ja, men, kjere Luther, jeg kan forsikkre Dig, jeg har Troen«. »Forsikkre, forsikkre, hvad er det for en Tale? I Forhold til at have Troen behøves ingen Forsikkring, hvis man har den (thi Troen er en urolig Ting, det mærkes strax); og ingen Forsikkring kan hjælpe, hvis man ikke har den«. »Ja, men tro mig dog blot, jeg kan forsikkre saa høitideligt som muligt....« »Aa, saa hold op med den Snak, hvad kan det hjælpe med Din Forsikkring!« »Ja, men dersom Du blot vilde læse et af mine Skrifter, skal Du see, hvor jeg kan fremstille Troen, saa veed jeg da, jeg maa have den.« »Jeg troer det Menneske er gal! Hvis det er saa, at Du kan fremstille Troen, det beviser blot, at Du er en Digter, og gjør Du det godt, at Du er en god Digter, men Intet mindre, end at Du er en Troende. Maaskee kan Du ogsaa græde, naar Du fremstiller Troen, det vilde saa bevise, at Du var en god Skuespiller; Du husker nok Historien om ham den Skuespiller i Oldtiden, der i den Grad formaaede at sætte sig ind i det Rørende, at han endog græd, naar han kom hjem fra Theatret, og græd flere Dage efter – det beviste kun, at han var en god Skuespiller. Nei, min Ven, Troen er en urolig Ting; den er en Sundhed, men stærkere og heftigere end den hidsigste Feber, og det hjælper ikke, at en Syg vil forsikkre, jeg har ikke Feber, naar Lægen føler det paa Pulsen, men heller ikke, at en Sund vil sige, jeg har Feber, naar Lægen ved at føle Pulsen føler det er ikke sandt – saaledes ogsaa, naar man ikke føler Troens Puls i Dit Liv, saa har Du heller ikke Troen. Naar man derimod fornemmer Troens Uro som Pulsen i Dit Liv, da kan Du siges at have Troen og at »vidne« om Troen. Og dette er saa igjen egentligen det at prædike; thi at prædike er hverken at fremstille Troen i Bøger, ei heller som Taler at fremstille den i »stille Timer«, der skulde jo, som jeg har sagt det i en Prædiken, egentligen »ikke prædikes i Kirker men paa Gaden«, og det skal heller ikke være en Taler, men et Vidne, det er: Troen, denne urolige Ting, skal være kjendelig i hans Liv.«

Ja, Troen er en urolig Ting. Lad mig, for dog at gjøre lidt opmærksom i denne Henseende, fremstille Troens Uro i en saadan Troes-Helt eller et Sandheds-Vidne. Saa er der da en given Virkelighed; det er der jo i ethvert Øieblik. Disse Tusinder og Tusinder og Millioner, de passe hver Sit, Embedsmanden passer Sit, og den Lærde passer Sit, og Kunstneren Sit, og den Næringsdrivende Sit, og Bagvaskeren Sit, og Lediggængeren, ikke mindre travlt, Sit, og saaledes fremdeles, Enhver passer Sit i dette Mangfoldighedens krydsende Spil, som er Virkeligheden. Imidlertid sidder, som nu Luther i en Klostercelle, eller paa et afsides Værelse, kort afsides sidder der et eenligt Menneske i Frygt og Bæven og megen Anfægtelse. Et eenligt Menneske! Ja, det er Sandhed. Det er Usandhed, hvad disse Tider have opfundet, at det er Tallet (det Numeriske), Mængden, eller det høistærede og høistærede dannede Publikum, fra hvem Reformationer udgaae – det vil da sige religieuse Reformationer, thi i Gadebelysningen, i Befordrings-Væsenet udgaae de dog maaskee bedst fra Publikum; men at en religieus Reformation skulde udgaae fra Publikum er Usandhed og, christelig forstaaet, en oprørsk Usandhed. Altsaa der sidder et eenligt Menneske i Anfægtelse. Maaskee nyder jeg dog nogen Anerkjendelse i min Samtid som Sjelekyndig (Psycholog); jeg kan bevidne, jeg har seet Mennesker, om hvem jeg tør sige: de have vistnok været meget udsatte for Fristelser; men aldrig har jeg seet Nogen, om hvem jeg turde sige: han er forsøgt i Anfægtelse. Og dog er et Aar at være udsat for Fristelse, det er Intet mod en Time i Anfægtelse. Saaledes sidder hiint eenlige Menneske; han sidder, eller, om Du vil, han gaaer, maaskee op og ned af Gulvet, som den fangne Løve i Buret; og dog, det han er fanget i er forunderligt, han er af Gud eller ved Gud fanget i sig selv. – Nu skal det da sættes ind i Virkelighed, det han har lidt for i Anfægtelsen. Troer Du han har Lyst dertil? Sandeligen, Enhver, der ad disse Veie kommer jublende: vær forvisset, han er ikke kaldet. Der er Ingen af de Kaldede, uden at han helst har villet være fri, Ingen uden at han, som et Barn kan tigge og bede for sig, har bedet for sig; men det hjalp ikke, han skal frem. – Saa veed han, at idet han nu gjør Skridtet, da reiser Forfærdelsen sig. Den der ikke er kaldet – idet Forfærdelsen reiser sig, bliver han saa angest, at han flyer tilbage. Men den Kaldede – o, min Ven, hellere end gjerne vilde han tilbage, gysende for Forfærdelsen; men idet han allerede har vendt sig om for at flye, seer han – han seer den endnu større Rædsel, der ligger bagved, Anfægtelsens Rædsel, han skal frem – saa gaaer han frem; han er nu ganske rolig, o, thi Anfægtelsens Rædsel er en frygtelig Tugtemester, der kan give Courage. – Forfærdelsen reiser sig. Alt hvad der ganske tilhører den givne Virkelighed væbner sig mod denne Anfægtelsernes Mand, hvem man dog ikke kan gjøre bange, og, forunderligt nok, fordi han er saa bange – for Gud; de angribe ham, hade ham, forbande ham. De Faa, der ere ham hengivne, de raabe til ham: »o, skaan Dig selv, Du gjør Dig selv og os Alle ulykkelige, o, hold nu op, reis ikke Forfærdelsen stærkere, stands det Ord, der svæver Dig paa Læben, tilbagekald heller det sidste.« O, min Tilhører, Troen er en urolig Ting.

Saa forkynder jeg da maaskee Tumult, Alts Omvæltning, Uorden? Sandeligen nei. Enhver, der ikke kjender min Forfatter-Virksomhed, maa nøies med denne Forsikkring. Enhver, der kjender min Forfatter-Virksomhed, maa vide, at jeg har arbeidet i modsat Retning.

Men christelig gives der to Arter Uorden. Den ene er Tumulten, Udvorteshedens Spektakel. Den anden Uorden er: Dødsstilheden, Uddøetheden; og denne er maaskee den farligste.

Mod den har jeg virket, og virket for at vække Uro i Retning af Inderliggjørelse. Lad mig nøiagtigt bestemme, hvor jeg saa at sige er. Der er blandt os en høiærværdig Olding, denne Kirkes øverste Geistlige; Det han, hans »Prædiken«, har villet, det Samme er det, jeg vil, kun en Tone stærkere, hvad der ligger i min Personligheds Forskjellighed, og hvad Tidens Forskjellighed kræver. Der er blandt os Nogle, som gjøre Fordring paa i strengeste Forstand at være Christne, at være det i Modsætning til os Andre: dem har jeg ikke kunnet slutte mig til. Deels mener jeg, at deres Liv ikke indfrier den Maalestok, som de selv foranledige eller nødsage En til at anlægge ved saa stærkt at fremhæve at de ere Christne – dog dette er mig det mindre Vigtige; deels er jeg for lidt Christen til at turde slutte mig til Nogen, som gjør en saadan Fordring. Er jeg end maaskee lidt – ja selv om saa var, hvis jeg maaskee var, ikke lidt fremmeligere end Adskillige af Gjennemsnittet blandt os: jeg er kun et Digterisk fremmeligere, det vil sige, jeg veed bedre, hvad Christendom er, veed bedre at fremstille det, o, men dette er (husk paa, hvad Luther sagde til mig!) en saare uvæsentlig Forskjel. Væsentligen tilhører jeg Gjennemsnittet. Og her er det, jeg har arbeidet for Uro i Retning af Inderliggjørelse.

Thi, christelig, er der to Arter sand Uro. Uroen i Troes-Heltene og Sandheds-Vidnerne, hvilken sigter til at reformere et Bestaaende. Saa langt har jeg aldrig vovet ud, det er ikke Noget for mig; og forsaavidt Nogen i Samtiden kunde synes at ville vove saa langt ud, var jeg ikke utilbøielig til at tage Polemik mod ham, for at bidrage til at det maatte blive aabenbart, om han er den Berettigede. Den anden Art Uro er i Henseende til Inderliggjørelse. Saaledes er jo ogsaa en sand Forelskelse en urolig Ting – men det falder ikke den Elskende ind at ville forandre et Bestaaende.

For denne Uro i Retning af Inderliggjørelse har jeg arbeidet. Men »uden Myndighed«. Istedetfor selv tomt opblæst at udgive mig for et Sandheds-Vidne, og foranledige Andre til dumdristigt at ville være det Samme, er jeg en umyndig Digter, der rører ved Hjælp af Idealerne. Som nu saaledes, at jeg strax skal give et Exempel, og tillige vise, hvorledes jeg blandt Andet ogsaa benytter Troes-Heltene. Du, min Tilhører, Du kalder Dig jo en Christen. Nu, og Du veed, at hvad der er det Bestemteste af Alt men tillige det Ubestemteste: Døden, at den ogsaa engang vil nærme sig Dig, og det bliver Din Død. Dog Du er en Christen; Du haaber altsaa og troer, at Du vil blive salig, lige saa salig som noget Sandheds-Vidne eller nogen af Troens Helte, der dog efter en ganske anden Maalestok have maattet kjøbe det at være Christen. Maaskee vilde derfor ogsaa En, der havde Myndighed, her tale anderledes til Dig, forfærdende sige, at Du var i en Indbildning med at være Christen, at Du farer til Helvede. Langt fra mig være at fremstille dette som en Overdrivelse af den Myndige, nei, jeg forstaaer det kun altfor vel, hvilken Anstrængethed der fordres for at turde vove at anbringe et saadant Enten–Eller paa en Anden. Men jeg, den Umyndige, jeg tør ikke tale saaledes; jeg troer, at Du vil blive lige saa salig som noget Sandheds-Vidne, nogen af Troens Helte. Men saa siger jeg til Dig. Tænk nu engang Dit og en Saadans Liv sammen. Betænk, hvad han har maattet offre, han som offrede Alt, det der falder tungest at offre i et første Øieblik, og det der i Længden falder tungest at have offret! Tænk, hvad han har lidt, hvor smerteligt og hvor langvarigt! O, dersom Du lever lykkelig i et elsket Hjem, dersom Din Hustru hænger fast ved Dig af sit ganske Hjerte og af al sin Magt, dersom Du har Glæde af Dine Børn: betænk hvad det dog vil sige, saaledes Dag efter Dag at leve hen i denne Fred og Hvile, velgjørende for et Menneskes Sjel, mere velgjørende end den svage Eftermiddagens sildigere Belysning er for Den, hvis Øie smerter; og betænk, at dette er Dit daglige Liv – og tænk saa paa Sandheds-Vidnet! Og dersom Du lever, ikke i Lediggang, langtfra, men saaledes, at Din Gjerning, der tager Din Tid, Din Flid, Din Kraft, dog ikke tager anderledes, end at der baade er Hvile nok fra Arbeidet, og Arbeidet selv mangen Gang forfriskende som Tidsfordriv; og dersom Du lever om ikke i Overflod, saa dog saa Du har Dit rigelige Udkomme; og dersom Du har Tid til saa mangen Nydelse, der vederqvægende udfylder Tiden og giver ny Livslyst; og dersom kort og godt Dit Liv er en stille daglig Nydelse – o, hans Liv var en piinagtig daglig Lidelse: saa døe I begge, og I blive lige salige! Og dersom Du i lykkelig Ubemærkethed kan glæde Dig ved Livet, uforstyrret og ubemærket faae Lov til at gaae og være Dig selv; og dersom Du, just fordi Du lever ubemærket, saa ofte faaer Leilighed til at lære Menneskene at kjende fra deres bedre, deres gode, deres elskelige Side; og dersom Du, naar Du færdes i Vrimlen, enten møder den Fremmede, der slet ikke kjender Dig, eller Du møder deres velvillige, deeltagende Blik, som kjende Dig; og dersom Du, naar Du finder Leilighed til at gjøre et andet Menneske en Tjeneste, en Velgjerning, lønnes med saa megen Glæde, at det er et Spørgsmaal, om Du ikke egentligen gjør Dig selv en Tjeneste, en Velgjerning; og dersom Du letforstaaende Dit Liv let er i Forstaaelse med de Andre, og let forstaaet: o, han maatte Dag ud og Dag ind (hvad der er uadskilleligt fra en saadan Virken) lade sig ligesom fortære og opæde af denne menneskelige Piatten, der altid hungrig altid begierer Noget at snakke om; han maatte Aar efter Aar daglig være foranlediget, nødsaget til at lære Menneskene at kjende fra den, mildest talt, bestialske, stundom den dybe Fordærvelses Side; han maatte idelig og idelig forvisse sig om, at han var kjendt af Enhver, og kjende dette paa, at han i Enhvers Blik mødte Uvillie, Modstand, Forbittrelse, Haan og Deslige; han gjorde vel mod en heel Samtid, og lønnedes med en heel Samtids Forbandelser; i Anfægtelsens Qvaler maatte han erhverve sit Livs Forstaaelse, og saa maatte han møisommeligen Dag efter Dag arbeide sig gjennem alle de Medlevendes Misforstaaelser, og alle disse Misforstaaelsens Qvaler – – saa døe I begge, og I blive lige salige! Betænk dette, og ikke sandt, saa vil Du sige til Dig selv, hvad jeg siger til mig: hvad enten jeg nu virkelig ikke skal vove saa langt ud, eller jeg kjæler for mig selv, at jeg slet ikke vover ud: Eet vil jeg gjøre, og om jeg end havde nok saa meget at bestille forresten, jeg skal skaffe Tid til at mindes disse Herlige hver eneste Dag. O, mig synes, at det dog er en himmelraabende Uret, at vi To blive lige salige! Men i ethvert Fald, mit Liv skal være en Erindring om dem! – Og see, her har Du strax et Exempel paa en Bevægelse, der er Uro i Retning af Inderliggjørelse.

Og denne Uro, det er det Mindste, den mildeste, den laveste Form af Gudelighed. Og dog, troer Du, at vi ere saa fuldkomne, at vi ikke behøve, at der arbeides i den Retning? Husk paa, hvorledes det gik mig med Luther! Om det vilde gaae de Andre ligesaa, hvis Luther kom til dem, jeg veed det ikke.

Men tænk Dig Luther i vor Tid, opmærksom paa vor Tilstand, troer Du ikke han vilde sige, som han siger i en Prædiken: »Verden er som den fulde Bonde, naar man hjælper ham op fra den ene Side paa Hesten, saa falder han ned paa den anden Side«. Troer Du ikke, han vilde sige: Apostelen Jacob maa drages lidt frem, ikke for Gjerningerne mod Troen, nei, nei, det var dog heller ikke Apostelens Mening, men for Troen, for, om muligt, at bevirke, at Trangen til »Naaden« føles dybt i sand ydmyg Inderlighed, og for, om muligt, at forhindre, at »Naaden«, at Troen og Naaden som det ene Frelsende og ene Saliggjørende ikke tages ganske forfængeligt, bliver et Skalkeskjul for en endog raffineret Verdslighed. Luther – denne Guds Mand, denne ærlige Sjel! – oversaae eller glemte dog maaskee et vist Noget, hvad en senere, hvad især vor Tid maaskee kun altfor stærkt indskærper. Han glemte – anden Gang, Du Ærlige! – han glemte, hvad han var for ærlig til selv at vide, hvilken ærlig Sjel han selv var, hvad jeg, og ikke for min Dyds Skyld men for Sandhedens Skyld, maa fremhæve. Det Lutherske er ypperligt, er Sandheden; jeg har i Forhold til dette ypperlige Lutherske kun een Betænkelighed. Den angaaer ikke det Lutherske, nei, den angaaer mig selv: at jeg har forvisset mig om, at jeg er ikke nogen ærlig Sjel, men en listig Krabat. Saa bliver det vistnok rigtigst, at der passes noget nøiere paa Undersætningen (Gjerningerne, Existensen, det at vidne og lide for Sandhed, Kjerlighedens Gjerninger og saa videre) Undersætningen i det Lutherske. Ikke som skulde nu Undersætningen gjøres til Oversætning, Troen og Naaden afskaffes eller nedsættes, Gud forbyde det, nei, just for Oversætningens Skyld, og saa fordi jeg er en Saadan som jeg er, bliver det vistnok rigtigst, at der passes noget nøiere paa Undersætningen i det Lutherske – thi i Forhold til »ærlige Sjele« behøves der Intet at gjøres.

Og Jacob siger: vorder ikke blot Ordets Hørere, men dets Gjørere.

Dog for at vorde dets Gjørere, maa man jo først være dets Hører eller Læser, hvilket Jacob ogsaa siger.

Og nu staae vi ved vor Text.

Saa ville vi tale om:

Hvad der fordres for til sand Velsignelse at betragte sig i Ordets Speil?

For det Første fordres der, at Du ikke maa see paa Speilet, betragte Speilet, men see Dig selv i Speilet.

Dette synes saa indlysende, at man skulde troe, det neppe behøvede at siges. Dog gjøres det vist fornødent; og hvad der bestyrker mig i den Mening er, at denne Bemærkning ikke er af mig selv, ei heller af hvad vi nutildags kalde en from Mand, en Mand der saadan har nogle fromme Stemninger, men af et Sandheds-Vidne, et Blod-Vidne; og saadanne Herlige de ere vel underrettede.

Han advarer mod det Feilsyn, at komme til at betragte Speilet, istedetfor at see sig selv i Speilet. Jeg benytter blot Bemærkningen, og spørger saa Dig, min Tilhører, er den ikke som myntet paa vore Tider og vore Forhold, og overhovedet paa Christenhedens senere Tider?

Thi »Guds Ord« er jo Speilet: men, men – – o, uoverskuelige Vidtløftighed: hvor meget hører der i strengere Forstand til »Guds Ord«, hvilke Bøger ere ægte, ere de ogsaa af Apostlene, og ere disse ogsaa troværdige, have de selv seet Alt, eller maaskee i Forhold til Adskilligt dog kun hørt det af Andre; og nu Læsemaaderne, 30,000 forskjellige Læsemaader; og saa denne Stimmel eller Trængsel af Lærde og Meninger, og lærde Meninger og ulærde Meninger om hvorledes det enkelte Sted er at forstaae...... ikke sandt, dette her seer noget vidtløftigt ud! Guds Ord er Speilet – jeg skal, ved at læse eller høre, see mig i Speilet: men see, det med Speilet forvirrer sig saaledes, at jeg nok aldrig kommer til at speile mig – idetmindste ikke, naar jeg slaaer ind paa den Vei. Næsten kunde man fristes til at antage, at her er en heel menneskelig Underfundighed med i Spillet (ak, og sandt er det, vi Mennesker ere i Forhold til Gud og det Guddommelige og den gudfrygtige Sandhed saa underfundige, det er slet ikke saaledes, som vi saadan sige til hinanden: at vi saa gjerne ville gjøre Guds Villie, naar vi blot kunde faae den at vide), næsten kunde man fristes til at antage, at dette er Underfundighed, at vi Mennesker ikke gjerne ville til at see os i hiint Speil, og at det er derfor vi have hittet paa alt Dette, som truer med at gjøre Speilet umuligt, alt Dette, som vi saa hædre med de prisende Navne af lærd og grundig og alvorlig Gransken og Grublen.

min Tilhører hvor høit holder Du Guds Ord i Priis? Siig nu ikke, at Du holder det saa høit i Priis, at intet Udtryk kan betegne det; thi man kan ogsaa tale saa høit, at man slet Intet siger. Lad os derfor, for at det kan blive til Noget, tage et simpelt menneskeligt Forhold; holder Du Guds Ord høiere i Priis, saa meget desto bedre.

Tænk Dig en Elskende, der har modtaget et Brev fra den Elskede – saa dyrebart som dette Brev er for den Elskende, saa dyrebart, antager jeg, at Guds Ord er for Dig; som den Elskende læser dette Brev, saaledes antager jeg, at Du læser, og at Du mener, at Du bør læse Guds Ord.

Dog maaskee siger Du: »ja, men den hellige Skrift er skrevet i et fremmed Sprog«. Det er dog vel egentligen nærmest de Lærde, der behøve at læse den hellige Skrift i Grundsproget; men vil Du det ikke anderledes, vil Du holde fast paa, at Du maa læse Skriften i Grundsproget: vel, vi kan godt blive i Billedet med Brevet fra den Elskede, kun føie vi en lille Bestemmelse til.

Jeg antager da, at dette Brev fra den Elskede er skrevet i et Sprog, som den Elskende ikke forstaaer; og der er Ingen dersteds, som kan oversætte ham det, og maaskee han ikke engang gjerne ønskede nogen saadan Hjælp, for ikke at indvie en Uvedkommende i sine Hemmeligheder. Hvad gjør han? Han tager en Ordbog, sætter sig til at stave Brevet igjennem, slaaer hvert Ord op, for derpaa at tilveiebringe en Oversættelse. Lad os antage, at som han sidder beskjeftiget med dette Arbeide, kommer en Bekjendt af ham ind til ham. Han veed, at dette Brev er kommet; idet han seer hen paa Bordet, seer han det ligge der, og siger »naa, Du sidder og læser det Brev, Du har faaet fra Din Elskede« – hvad mener Du den Anden vil sige? Han svarer »er Du fra Forstanden, troer Du det er at læse et Brev fra den Elskede! Nei, min Ven, jeg sidder her i Slid og Slæb for ved Hjælp af en Ordbog at faae det oversat; stundom er jeg færdig at briste af Utaalmodighed, Blodet styrter mig til Hovedet, saa jeg gjerne kunde smide Ordbogen hen ad Gulvet – og det kalder Du at læse, vil Du spotte mig! Nei, snart er jeg, Gud skee Lov, færdig med Oversættelsen, og saa, ja saa, saa skal jeg til at læse Brevet fra den Elskede, det er noget ganske Andet – dog til hvem taler jeg.... dumme Menneske, gaa bort fra mine Øine, jeg gider ikke see Dig, at Du kunde falde paa i den Grad at fornærme den Elskede og mig, at kalde det at læse et Brev fra hende! Og dog, bliv, bliv, Du veed nok, det er kun mit Spøg, ja jeg saae det endogsaa meget gjerne, at Du blev, men oprigtigt talt, jeg har ingen Tid, der er endnu Noget tilbage at oversætte, og jeg er saa utaalmodig efter at komme til at læse – derfor, bliv ikke vred, men gaa, at jeg kan blive færdig«.

Altsaa, den Elskende gjør, i Forhold til Brevet fra den Elskede, Forskjel mellem det at læse og det at læse, mellem det at læse med Ordbog, og det at læse Brevet fra den Elskede. Blodet styrter ham til Hovedet af Utaalmodighed, naar han sidder og hexer med at læse med Ordbog, han bliver som Rasende, da hans Ven vover at kalde denne lærde Læsen en Læsen af Brevet fra den Elskede. Nu er han færdig med Oversættelsen – nu læser han Brevet fra den Elskede. Han betragtede hele dette, om Du saa vil, lærde Forarbeide som et nødvendigt Onde, at han kunde komme til – at læse Brevet fra den Elskede.

Lad os ikke for tidligt slippe dette Billede. Lad os antage, at dette Brev fra den Elskede ikke blot, som vel saadanne Breve i Almindelighed gjøre, indeholdt Udtalelse af en Følelse, men at der var et Ønske indeholdt, Noget den Elskede ønskede at den Elskende skulde gjøre. Det var, lad os antage dette, meget der fordredes af ham, saare meget, der var, vilde enhver Trediemand sige, god Grund til at betænke sig: men den Elskende – i samme Secund flux afsted for at faae opfyldt den Elskedes Ønske. Lad os antage, at efter en Tids Forløb traf de Elskende sammen, og den Elskede sagde: »men, Kjere, det havde jeg jo slet ikke forlangt af Dig, Du maa have forstaaet det Ord forkeert, eller oversat det forkeert« – troer Du, at den Elskende nu fortrød, at han, istedetfor strax i samme Secund at ile med at efterkomme Ønsket, ikke først havde gjort sig nogle Betænkeligheder, og saa maaskee havde faaet et Par Ordbøger endnu til Hjælp, og saa faaet flere Betænkeligheder, og saa maaskee faaet Ordet oversat rigtigt, og altsaa været fri – troer Du, han fortryder denne sin Feiltagelse, troer Du, han behager den Elskede mindre? Tænk Dig et Barn, ret hvad man kalder en flink og dygtig Discipel; da Læreren en Dag har sat dem Lectien for til næste Dag, siger han: lad mig nu see, at I kan Eders Ting godt imorgen. Paa vor flinke Discipel gjør dette et dybt Indtryk. Han kommer hjem fra Skole, flux til Arbeidet. Men han har ikke hørt ganske nøie efter, hvor langt de har for, hvad gjør han? Denne Lærerens Formaning er det, der har gjort Indtryk paa ham, han læser vel dobbelt saa langt, som det viser sig, at han virkelig har for: troer Du, Læreren vil synes mindre godt om ham, fordi han ganske udmærket kan en dobbelt saa lang Lectie? Tænk Dig en anden Discipel; han hørte ogsaa Lærerens Formaning, han havde heller ikke hørt nøie, hvor langt de havde for. Da han saa kom hjem, sagde han: jeg maa først have at vide, hvor langt vi have for. Saa gik han hen til en af sine Kammerater, saa til en anden, han var heller ikke hjemme, derimod kom han her i Snak med en ældre Broder af ham – og saa langt om længe kom han hjem, og Tiden var gaaet, og han fik slet Ingenting læst!

Altsaa, den Elskende gjorde i Forhold til Brevet fra den Elskede Forskjel mellem at læse og at læse, fremdeles forstod han det at læse saaledes, at var der et Ønske indeholdt i Brevet, da skulde man begynde dets Opfyldelse strax, der var ikke eet Secund at spilde.

Tænk nu paa Guds Ord. Naar Du læser Guds Ord lærd – vi nedsætte ikke Lærdom, nei, langtfra – men husk vel paa: naar Du læser Guds Ord lærd, med Ordbog og saa videre, saa læser Du ikke Guds Ord – husk paa den Elskende, der sagde: »det er ikke at læse Brevet fra den Elskede.« Er Du da Lærd, saa pas dog endeligen paa, at Du ikke over al denne lærde Læsen (der ikke er at læse Guds-Ord) glemmer at læse Guds Ord. Er Du ikke lærd, o, misund ham ikke, glæd Dig, at Du strax kan komme til at læse Guds Ord! Og er der nu et Ønske, et Bud, en Befaling, saa – husk paa den Elskende! – saa flux afsted for at gjøre derefter. »Men«, siger Du maaskee, »der ere saa mange dunkle Steder i den hellige Skrift, hele Bøger, der næsten ere som Gaader«. Hertil vilde jeg svare: for at jeg skulde indlade mig med denne Indvending, maatte den gjøres af En, hvis Liv udtrykte, at han nøiagtigt havde efterkommet alle de Steder, som ere lette at forstaae; er dette Tilfældet med Dig? Dog saaledes vilde den Elskende bære sig ad med Brevet; hvis der vare dunkle Steder, men der ogsaa var tydeligt udtalte Ønsker, han vilde sige: »jeg maa flux efterkomme Ønsket, saa faaer jeg at see, hvad det bliver med de dunkle Steder, o, men hvorledes skulde jeg sætte mig hen og grunde over de dunkle Steder, og lade Ønsket uopfyldt, Ønsket, som jeg tydeligt forstod!« Dette vil sige; naar Du læser Guds Ord: Det der forpligter Dig, er ikke de dunkle Steder, men Det, Du forstaaer; og Det har Du øieblikkeligt at efterkomme. Var der kun eet eneste Sted, Du forstod i den hele Hellige Skrift: vel, saa har Du først at gjøre det; men ikke har Du først at sætte Dig hen, og grunde over de dunkle Steder. Guds Ord er givet, for at Du skal handle derefter, ikke for at Du skal øve Dig i at fortolke dunkle Steder. Læser Du ikke Guds Ord saaledes, at Du betænker, at den mindste Smule Du forstaaer øieblikkeligt forpligter Dig til at gjøre derefter: saa læser Du ikke Guds Ord. Saaledes meente den Elskende: »hvis jeg, istedetfor øieblikkeligt at ile for at opfylde Ønsket, som jeg forstaaer, vil sætte mig hen og grunde over, hvad jeg ikke forstaaer, saa læser jeg ikke Brevet fra den Elskede. Jeg kan træde med god Samvittighed frem for den Elskede og sige: der vare nogle dunkle Steder i Dit Brev, i Forhold til dem har jeg sagt, den Tid den Sorg; men der var et Ønske, som jeg forstod, det har jeg øieblikkeligt opfyldt. Derimod kan jeg ikke træde med god Samvittighed frem og sige: der vare nogle dunkle Steder i Dit Brev, som jeg ikke forstod, dem satte jeg mig hen at grunde over, og i Forhold til Dit Ønske, som jeg nok forstod, sagde jeg: den Tid den Sorg.« – Men maaskee frygter Du, at det skulde gaae Dig i Forhold til Guds Ord, som det gik den Elskende med Brevet, at Du (dog denne Frygt er vistnok ugrundet i Forhold til Guds Fordring) at Du skulde komme til at gjøre for meget, at Du ved at slaae op i endnu een Ordbog skulde see, at der dog ikke var fordret saa meget: o, min Ven, mishagede dette da den Elskede, at den Elskende var kommet til at gjøre for meget? Og hvad troer Du, at den Elskende vilde sige om det at nære en saadan Frygt? Han vilde sige: »Den, der nærer en saadan Frygt for at skulle komme til at gjøre for meget, han læser ikke Brevet fra den Elskede«; og jeg vilde sige: han læser heller ikke Guds Ord.

Lad os endnu ikke slippe dette Billed med Brevet fra den Elskede. Da han sad og var beskjeftiget med at oversætte det ved Hjælp af en Ordbog, blev han forstyrret, idet en Bekjendt kom ind til ham. Han blev utaalmodig, »men«, vilde han vistnok sige, »det var blot fordi jeg blev sinket, thi ellers kunde det da være det Samme, jeg læste jo ikke den Gang Brevet. Ja, var En kommet ind til mig, medens jeg sad og læste Brevet, det var noget ganske Andet, det havde været en Forstyrrelse. Dog det skal jeg nok sikkre mig mod; førend jeg begynder paa Sligt, lukker jeg min Dør i Laas og er ikke hjemme. Thi jeg vil være ene, uforstyrret ene med Brevet; er jeg ikke det, saa læser jeg heller ikke Brevet fra den Elskede.«

Han vil være ene, uforstyrret ene med Brevet – »ellers«, siger han, »læser jeg ikke Brevet fra den Elskede«.

Og saaledes med Guds Ord; Den, som ikke er ene med Guds Ord, læser ikke Guds Ord.

Ene med Guds Ord! Min Tilhører, lad mig her gjøre en Tilstaaelse om mig selv: jeg tør endnu ikke ret være ganske ene med Guds Ord, saa intet Sandsebedrag underskyder sig. Og tillad mig saa at sige Eet: jeg har aldrig seet Nogen, om hvem jeg turde troe, at han havde Oprigtighed og Mod til at være ene med Guds Ord, saa intet, intet Sandsebedrag underskyder sig.

Forunderligt! Naar der i Samtid træder en stærkere bevæget Mand frem, der bestemmer Prisen af det at være Christen blot en Femtedeel mod den Priis, som Evangeliet sætter: da raaber man »tager Eder iagt for det Menneske, læser ikke hvad han skriver, allermindst i Eensomhed, taler ikke med ham, allermindst eensomt, han er et farligt Menneske.« Men den hellige Skrift! Ja næsten ethvert Menneske eier den, man tager ikke i Betænkning at forære enhver Confirmand (altsaa i den farligste Alder) denne Bog. I Sandhed, der maa være mange Sandsebedrag med, man maa være forvænt ved at denne Bog nu engang er til, man maa læse den paa en ganske egen Maade – mindst saaledes, at man er ene med den.

At være ene med den hellige Skrift! Jeg tør det ikke! Naar jeg nu slog op i den: det første det bedste Sted – det fanger mig øieblikkeligt; det spørger mig (ja, det er som det var Gud selv, der spurgte mig): har Du gjort, hvad Du der læser? Og saa, saa.... ja saa er jeg fangen. Saa enten strax til Handling, eller øieblikkeligt ydmygende Indrømmelse.

O at være ene med den hellige Skrift – og hvis ikke, saa læser Du ikke den hellige Skrift.

At det at være ene med Guds Ord, at det er en farlig Sag, det er ogsaa stiltiende tilstaaet af just dygtigere Mennesker. Der var maaskee Den (et dygtigere, et alvorligere Menneske, om vi end ikke kunne prise hans Beslutning), som sagde til sig selv: »jeg duer ikke til at gjøre noget halvt – og denne Bog, Guds Ord, det er en yderst farlig Bog for mig, og det er en herskesyg Bog, giver man den en Finger, tager den hele Haanden, giver man den hele Haanden, tager den hele Manden og omkalfatrer maaskee pludselig hele mit Liv efter en uhyre Maalestok. Nei, uden (hvad jeg afskyer) uden at tillade mig et eneste spottende eller nedsættende Ord om denne Bog: jeg flytter den hen paa et afsides Sted, jeg vil ikke være ene med den.« Vi billige det ikke; men der er dog altid Noget heri, som vi billige, en vis Redelighed.

Men man kan ogsaa paa en ganske anden Maade værge sig mod Guds Ord, trodsende paa, at man godt vover at være ene med det, hvilket dog ikke er sandt. Thi tag den hellige Skrift, luk Din Dør – men tag saa ti Ordbøger, fem og tyve Fortolkninger: saa kan Du læse den lige saa rolig og ugeneret som Du læser Adresseavisen. Falder det Dig saa maaskee engang, forunderligt nok, ind, ligesom Du bedst sidder og læser et Sted: har jeg gjort dette, handler jeg herefter (det er naturligviis i Distraction, i et adspredt Øieblik hvor Du ikke er samlet i vanlig Alvor, Du kan falde paa Sligt), saa er Faren dog ikke saa stor. Thi see, maaskee er der flere Læsemaader, og maaskee bliver der just nu fundet et nyt Haandskrift: ih bevares – og Udsigt til nye Læsemaader, og maaskee ere fem Fortolkere af een Mening, og syv af en anden, og to af en mærkelig Mening, og tre vaklende eller har ingen Mening, og »jeg selv er ikke ganske enig med mig selv om Meningen af dette Sted, eller, for at sige min Mening, jeg er af samme Mening som de tre Vaklende, der ingen Mening har« og saa videre En Saadan kommer saa ikke i den Forlegenhed som jeg: enten strax at maatte gjøre efter Ordet, eller dog at maatte gjøre en ydmygende Tilstaaelse. Nei, han er rolig, han siger: »der er ikke Noget i Veien fra min Side, jeg skal nok gjøre derefter – dersom det blot først bliver bragt i Orden med Læsemaaden, og Fortolkerne blive nogenlunde enige.« Aha! Dermed har det nemlig vistnok lange Udsigter. Derimod opnaaede Manden, at det bliver dunkelt, om Feilen ikke stikker i ham, at det er ham, der ikke har Lyst til at fornegte Kjød og Blod og gjøre efter Guds Ord. O, sørgelige Misbrug af Lærdom, o, at det gjøres Menneskene saa let saaledes at bedrage sig selv!

Thi var der ikke saa mange Sandsebedrag og Selvbedrag, da vilde vistnok Enhver, som jeg gjør det, tilstaae: jeg tør neppe være ene med Guds Ord.

Ene med Guds Ord, det maa man være, som den Elskende vilde være ene med Brevet, thi ellers var det ikke at læse Brevet fra den Elskede – og ellers er det ikke at læse Guds Ord, eller at see sig i Speilet. Og det var jo det vi skulde og det vi først skulde, hvis vi til Velsignelse skulde betragte os i Ordets Speil, vi skulde ikke see paa Speilet men os i Speilet. Er Du lærd, husk paa, at hvis Du ikke læser Guds Ord anderledes, det vil hænde Dig, at Du, efter hele Dit Liv igjennem at have læst mange Timer hver evige Dag i Guds Ord, dog aldrig har læst – Guds Ord. Gjør da Forskjellen, saa Du kommer til (foruden den lærde Læsning) ogsaa at læse Guds Ord, eller tilstaa Dig idetmindste selv, at Du, trods den daglige lærde Læsen deri, ikke læser Guds Ord, at det vil Du overhovedet ikke have med at gjøre. Er Du ulærd: desto mindre er Du vel foranlediget til at see feil; altsaa flux til Sagen, intet Sinkeri med at betragte Speilet, men flux til at betragte Dig i Speilet.

Dog hvorledes læses Guds Ord vel i Christenheden? Skulde vi deles i to Classer – thi enkelte Undtagelser kan der ikke saaledes tages Hensyn til – saa maatte man vel sige: en større Deel læser aldrig Guds Ord, en mindre Deel læser det paa en eller anden Maade lærd, det er, læser dog ikke Guds Ord, men betragter Speilet. Eller for at sige det Samme paa en anden Maade: en større Deel betragter Guds Ord som et forældet Oldtids-Skrift, man sætter til en Side; en mindre Deel betragter Guds Ord som et yderst mærkeligt Oldtids-Skrift, paa hvilket man anvender en forbausende Flid og Skarpsindighed og saa videre – betragtende Speilet.

Tænk Dig et Land. Der udgaaer et Kongebud til alle Embedsmænd, Undersaatter, kort til hele Befolkningen. Hvad skeer? Der foregaaer med Alle en mærkelig Forandring: Alt forvandler sig til Fortolkere, Embedsmændene blive Forfattere, den ene Fortolkning lærdere, skarpsindigere, smagfuldere, dybsindigere, ingenieusere, vidunderligere, deiligere og vidunderlig deiligere end den anden, udkommer hver evige Dag; Critiken, som skal holde Oversigten, kan neppe holde Oversigten over denne uhyre Literatur, ja Critiken selv bliver en saa vidtløftig Literatur, at det ikke er muligt at holde Oversigten over Critiken: Alt er Fortolkning – men ingen læste Kongebudet saaledes, at han gjorde derefter. Og ikke blot dette, at Alt blev Fortolkning, nei man forrykkede tillige Synspunktet for hvad Alvor er, og gjorde det at sysle med Fortolkningen til den egentlige Alvor. Tænk Dig, at denne Konge ikke var en menneskelig Konge – thi en saadan vilde vel ogsaa nok forstaae, at man egentligen gjorde Nar af ham ved paa den Maade at vende Sagen; men da en menneskelig Konge er afhængig især af samtlige Embedsmænd og Undersaatter, saa vilde han vel blive nødsaget til at gjøre gode Miner til slet Spil, lade som var dette i sin Orden, saa den smagfuldeste Fortolker til Belønning blev ophøiet i Adelsstanden, og den dybsindigste blev hædret med en Orden og saa videre – tænk Dig, at denne Konge var en Almægtig, som altsaa ikke kommer i Forlegenhed, om saa alle Embedsmænd og Undersaatter spille falsk mod ham. Hvad troer Du nu denne almægtige Konge vil tænke om Sligt? Mon han ikke vilde sige: det, at de ikke efterkomme Budet, det kunde jeg endda tilgive; fremdeles, hvis de forenede indgik med et Bønskrift til mig, om at jeg vilde have Taalmodighed med dem, eller maaskee ganske forskaane dem for dette Bud, der faldt dem saa tungt: det kunde jeg tilgive dem. Men dette kan jeg ikke tilgive, at man endog forrykker Synspunktet for hvad Alvor er.

Og nu Guds Ord! »Mit Huus er et Bedehuus, men I have forvandlet det til en Røverkule.« Og Guds Ord, hvad er det efter sin Bestemmelse, og hvad have vi forvandlet det til? Al denne Fortolken og Fortolken og Videnskab og ny Videnskab, det bringes frem under den høitidelige, alvorsfulde Form: at det er for ret at forstaae Guds Ord – see nærmere til, og Du skal see, det er for at værge sig mod Guds Ord. Det er kun altfor let at forstaae Fordringen, som er indeholdt i Guds Ord (»giv alt dit Gods til de Fattige«, »naar En slaaer Dig paa den høire Kind, vend den venstre til«, »naar En tager Din Kappe, giv ham Kjortelen med«, »værer altid glade«, »agter det for idel Glæde, naar I falde udi adskillige Fristelser« og saa videre, Alt lige saa let at forstaae som den Bemærkning: »det er godt Veir idag«, en Bemærkning, som kun paa een Maade kunde blive vanskelig at forstaae, naar der blev en Literatur til for at fortolke den), ikke den meest indskrænkede sølle Stakkel kan med Sandhed negte at kunne forstaae Fordringen – men det kniber for Kjød og Blod at ville forstaae den og at skulle gjøre derefter. Og det er, i mine Tanker, menneskeligt, at et Menneske krymper sig ved, at lade Ordet ret faae Magt over ham – vil ingen Anden tilstaae det, jeg tilstaaer, det gjør jeg. Det er menneskeligt at bede Gud at have Taalmodighed, om man ikke strax kan, hvad man skal, men dog lovende at stræbe; det er menneskeligt, at bede Gud at have Medlidenhed, at Fordringen er En for høi: vil ingen Anden tilstaae det om sig, jeg tilstaaer, det gjør jeg. Men det er dog ikke menneskeligt at give Sagen en ganske anden Vending: at jeg listigt skyder, det ene Lag efter det andet, Fortolkning og Videnskab og atter Videnskab (omtrent som naar en Dreng anbringer en Serviet eller flere under Trøien, naar han skal have Prygl) at jeg skyder alt Dette mellem Ordet og mig, og saa giver denne Fortolken og Videnskabelighed Navn af Alvor og Sandheds-Iver, og saa lader denne Syslen svulme op til en saadan Vidtløftighed, at jeg aldrig kom til at faae Indtrykket af Guds Ord, aldrig kommer til at betragte mig i Speilet. Det seer ud, som drog al denne Forsken og Grublen og Gransken og Grunden Guds Ord ganske nær til mig: det Sande er, at jeg just saaledes, paa den listigste Maade, fjerner Guds Ord det længst mulige fra mig, uendelig langt længere end det er fra Den, der aldrig saae Guds Ord, uendelig langt længere end det er fra Den, der blev saa angest og bange for Guds Ord, at han kastede det det længst mulige bort.

Thi en endnu større Afstand fra det Fordrede (at see sig selv i Speilet), en endnu større Afstand derfra end det aldrig at see Speilet, en endnu større Afstand er det: Aar ud og Aar ind hver evige Dag at kunne sidde ganske rolig – og betragte Speilet.

For det Andet fordres der, at naar Du læser Guds Ord, for at see Dig i Speilet, saa maa Du (for at Du virkelig kan komme til at see Dig i Speilet) huske paa i eetvæk at sige til Dig selv: det er mig der tales til, det er mig der tales om.

Lad Dig ikke bedrage – eller, vær ikke selv underfundig. Thi vi Mennesker, vi ere i Forhold til Gud og Guds Ord o, saa listige, selv den dummeste iblandt os saa listig – ja, Kjød og Blod og Selvkjerlighed er meget listig.

Saa have vi da (at det er for at værge os mod Guds Ord, det sige vi ikke, vi ere da ikke gale, sagde vi det, saa havde vi jo ingen Profit af vor kloge Opfindelse), vi have opfundet, at det at tænke paa sig selv (hvad det vistnok i mange Tilfælde kan være, kun ikke naar det gjælder om at lade Guds Ord faae Magt over sig) at det er – tænk Dig, hvor listigt! – det er Forfængelighed, sygelig Forfængelighed! Pfui, skulde jeg være saa forfængelig! Thi det at tænke paa sig selv og sige: »det er mig«, det er, som vi Lærde sige, det er det Subjektive; og det Subjektive, det er Forfængelighed, denne Forfængelighed, ikke at kunne læse en Bog – Guds Ord! – uden at mene, at det er mig, den handler om. Skulde jeg ikke afskye at være forfængelig! Og skulde jeg da være saa dum ikke at gjøre det, naar jeg derved tillige sikkrer mig, at Guds Ord ikke kan komme til at faae fat i mig, fordi jeg ikke sætter mig i noget personligt (subjektivt) Forhold til Ordet, men derimod – o, Alvor, for hvilken jeg saa bliver høit priset af Menneskene! – forvandler Ordet til et upersonligt Noget (det Objektive, en objektiv Lære og deslige), til hvilket jeg – den baade Alvorlige og Dannede! – forholder mig objektivt, saa jeg da ikke er saa udannet og forfængelig, at jeg bringer min Personlighed med i Spillet, at jeg skulde troe det var mig, der taltes til, mig – og ieetvæk – mig, der taltes om. Ih langtfra mig være en saadan forfængelig Udannethed – og langtfra mig være ogsaa, hvad der jo ellers saa let kunde skee, at Ordet fik fat i mig, just i mig, fik Magt over mig, saa jeg ikke kunde værge mig mod det, saa det blev ved at forfølge mig, indtil jeg enten gjorde efter det, forsagende Verden, eller dog tilstod, at jeg ikke gjorde det – den retfærdige Straf over hver Den, der tillader sig paa en udannet Maade at omgaaes Guds Ord.

Nei, nei, nei! Dette, naar Du læser Guds Ord, da i Alt hvad Du læser, bestandigt at sige til Dig selv: det er mig der tales til, mig der tales om, det er Alvoren, just det er Alvoren. Ei heller har nogen Eneste af dem, hvem Christendommens Sag i høiere Forstand har været betroet, glemt, at indskærpe dette atter og atter som det meest Afgjørende, som ubetinget Betingelsen, hvis Du skal komme til at see Dig i Speilet. Altsaa, det er dette Du skal gjøre, Du skal ieetvæk under Læsningen sige til Dig selv: det er til mig, det er om mig der tales.

Hiin mægtige Keiser i Østen, hvis Vrede det lille navnkundige Folk havde tildraget sig, om ham fortælles, at han havde en Slave, der sagde til ham hver Dag: husk paa at tage Hevn. Det var ogsaa Noget at huske paa; mig synes det havde været bedre, at have en Slave, der hver Dag havde erindret ham om at glemme, hvilket dog heller ikke er godt; thi skal man hver Dag erindres om at glemme, saa bliver det jo ikke til Alvor med at glemme. Men i ethvert Tilfælde, denne Hersker forstod meget rigtigt, just fordi han var vred (og Vrede er en, om ikke priselig, Bestemmelse af Personlighed), hvorledes man bærer sig ad, naar der skal virkes personligt paa En.

Men endnu bedre end denne Hersker blev dog Kong David betjent – det forstaaer sig, det er af den Art Betjening, som man sjeldent selv frivilligt ønsker, man fristes snarere til at betragte den som en af Livets største Ubeqvemmeligheder.

Historien jeg sigter til er bekjendt. Kong David saae Bathseba. At see hende – og at see, at hendes Mand stod i Veien, det var Eet. Han maa altsaa bort. Og det skeete ogsaa; man veed ikke ret, hvorledes det gik til, der maa være en Styrelse, han faldt i Slaget, dog »det er saa Krigens Gang«, siger Kongen, han havde formodentlig i Dumdristighed selv valgt saa farlig en Post, at det var den visse Død – jeg siger blot, var der Nogen, der ønskede ham af Dage, han kunde, hvis han raadede over Sligt, aldrig have gjort det bedre end at anbringe ham paa den Post, der var den visse Død. Nu er han af Veien. Det gik meget nemt. Og nu er der altsaa heller ikke mere Noget i Veien for at komme i lovlig Besiddelse af hans Hustru. Noget i Veien, er Du sær, det er jo endogsaa en i høi Grad ædel, høimodig, ægte kongelig Handling, der vil begeistre hele Krigerstanden, at Kongen ægter den for Fædrelandet faldne Krigers Enke.

Saa kommer der en Dag en Prophet op til Kong David. Lad os gjøre Sagen ret nærværende og modernisere den lidt. Den ene er Kongen, den høiest staaende Mand i Folket, den anden Propheten, en anseet Mand i Folket, begge, naturligviis, dannede Mænd, og man kan da være forvisset om, at deres Omgang med hinanden, deres Samtale vil bære ubetinget Præget af Dannelse. Desuden de ere begge, især den ene af dem, berømte Forfattere, Kong David den navnkundigste Digter, og, hvad dermed følger, en Kjender, en smagfuld Smags-Dommer, der veed at vurdere Fremstillingen og Udtrykkenes Valg, og Anlæget i et Digt og Sprogformen og Tonearten og dets Gavnlighed eller Skadelighed for Sæderne og saa videre

Og det er ret heldigt, just den rette Mand at komme til; thi Propheten har digtet en Novelle, en Fortælling, som han vil have den Ære at foredrage for Hans Majestæt, den kronede Digter og Kjender af Digterværker.

»Der boede tvende Mænd i en Stad; den ene var saare riig, og havde stort Qvæg og smaat Qvæg i Mængde; den Fattige derimod havde ikkun et lidet Lam, som han havde kjøbt og opklækket, og som var blevet stort hos ham tilligemed hans Børn. Det aad af hans Haand og drak af hans Bæger, og det var som et Barn i Huset. Men der en Vandringsmand kom til den Rige, da sparede han sit smaae Qvæg og sit store Qvæg, og tog den Fattiges Faar; det slagtede han, og lavede det til Spise for den Fremmede, som var kommet til ham.«

Jeg tænker mig, David har hørt opmærksomt efter, har derpaa givet sin Mening tilkjende, naturligviis ikke indblandet sin Personlighed (Subjektivitet), men upersonligt (objektivt) vurderet dette nydelige lille Arbeide. Der har maaskee været en Enkelthed, han har meent kunde være anderledes, et heldigere valgt Udtryk har han maaskee foreslaaet, maaskee ogsaa viist en lille Feil i Anlæget, roest Prophetens mesterlige Foredrag af Fortællingen, hans Stemme, Minespil, kort udtalt sig saaledes, som i vore Tider vi Dannede pleie at bedømme en Prædiken for Dannede, det er, en Prædiken der ogsaa selv er objektiv.

Da siger Propheten til ham: Du est Manden.

See, Fortællingen, Propheten foredrog, det var en Historie; men dette: Du est Manden, det var en anden Historie – det var Overgangen til det Subjektive.

[…]

Min Tilhører, jeg har endnu et Ord, jeg vilde sige; men jeg vil indklæde det i en Fremstilling, som maaskee ved første Øiekast vil synes Dig mindre høitidelig. Dog gjør jeg det med Flid og velbetænkt, fordi jeg mener, at det maaskee just saaledes vil gjøre et sandere Indtryk paa Dig.

Der var engang en Rigmand; han lod i Udlandet kjøbe i dyre Domme et Par aldeles feilfrie og udmærkede Heste, som han vilde have til sin egen Fornøielse og Fornøielsen af selv at kjøre. Saa gik der vel omtrent et Aar eller to. Dersom Nogen, som tidligere havde kjendt disse Heste, nu saae ham kjøre dem, han vilde ikke kunne kjende dem igjen: Øiet var blevet mat og døsigt, deres Gang uden Holdning og Sluttethed, Intet kunde de taale, Intet udholde, neppe kunde han kjøre en Miil uden at han maatte tage ind underveis, stundom gik de istaa ligesom han allerbedst sad og kjørte, derhos havde de faaet allehaande Nykker og Vaner, og uagtet de naturligviis fik Foder i Overflod, skæmmede de sig Dag for Dag. Da lod han Kongens Kudsk kalde. Han kjørte dem i en Maaned: der var paa hele Egnen intet Par Heste, der bar Hovedet saa stolt, hvis Blik var saa fyrigt, hvis Holdning saa skjøn, intet Par Heste, der saaledes kunde holde ud at løbe om det var syv Miil i eet Træk, uden at der toges ind. Hvori stak det? Det er let at see: Eieren, der, uden at være Kudsk, gav sig af med at være Kudsk, han kjørte dem efter Hestenes Forstand paa hvad det er at kjøre; den kongelige Kudsk kjørte dem efter Kudskens Forstand paa hvad det er at kjøre.

Saaledes med os Mennesker. O, naar jeg tænker paa mig selv og paa de Utallige, jeg har lært at kjende, da har jeg ofte sagt med Veemod til mig selv: her er Evner og Kræfter og Forudsætninger nok – men Kudsken mangler. Gjennem længere Tid er fra Slægt til Slægt vi Mennesker blevne, om jeg saa tør sige, kjørte (for at blive i Billedet) efter Hestenes Forstand paa hvad det er at kjøre, vi ere styrede, dannede, opdragne efter Menneskets Begreb om, hvad det er at være Menneske. See, deraf hvad vi mangle: Opløftelse, og hvad deraf følger igjen, at vi kun kunne udholde saa lidet, utaalmodige bruge strax Øieblikkets Midler, og utaalmodigen øieblikkeligt ville see Løn for vort Arbeide, hvilket just derfor bliver derefter.

\workheading{Dømmer selv! (Til Selvprøvelse, anden Række)}{1851–52 (udg. 1876)}

\worknote{Antologiens ''Judge for Yourself!''. Hong-udvalget bringer af anden afhandling ''Christus som Forbilledet, eller Ingen kan tjene to Herrer'' stykket om lilien og fuglen (den lille Ludvig og moderen der egentlig trækker vognen — det er Gud der virker), om Efterfølgelsen og Forbilledet, og Luthers forhold til tro og gjerninger, endende med spørgsmaalet: ''mon Troen findes paa Jorden?''}

\sectionheading{Christus som Forbilledet (uddrag)  ·  SV¹ XII, 452–461}

Christus som Forbilledet, eller Ingen kan tjene to Herrer.

Agt da paa Fuglen! Den synger og qviddrer, og qviddrer saa – o, hør! – det ind med hvad den – agt vel derpaa! – siger til Sorgen, hvad der staaer i en gammel Psalme: »ja, ja imorgen.« Og saa glæder Fuglen sig »idag.« Saa tænker Sorgen: »bliv Du kun, jeg skal nok passe paa; imorgen før det gryer ad Dag, og før Du er kommet op af Reden, og før Fanden faaer Sko paa (thi jeg er endnu tidligere paafærde end han; jeg er en af hans Tjenere og Forløbere, der komme først for at forsøge at bane ham Indgang) saa kommer jeg.« Og imorgen – Fuglen er der ikke mere. Hvorledes den er der ikke mere? Nei, den er reist, den er borte. Hvorledes reist? Der var jo lagt Beslag paa dens Pas, og jeg veed da for Satan, at den ikke er reist uden Pas.« »Ja, saa maa man ikke tilstrækkelig have passet paa, thi den er reist den efterlod en Hilsen til Dem, det Sidste den sagde, var: Siig til Sorgen: Ja, ja imorgen.« Du er dog snild, en uforlignelig Professor i Levekunsten, Du bevingede Reisende! O, saaledes at kunne sige til Sorgen: ja, ja imorgen; og saaledes at kunne, efterat have sagt det, glæde sig idag, dobbelt næsten fordi man har sin Glæde af at have sagt det! Og saa ikke at narre Sorgen et Par Dage, hvad hjælper det, da lige saa gjerne først som sidst, men at sige det saa længe, til det bliver en forgjæves Gang, naar Sorgen endelig for Alvor kommer; at lade den rende forgjæves hver evige Dag med den Besked: imorgen, og naar det saa endelig maatte blive Alvor, at den saa for Alvor kommer til at gaae forgjæves. – Og Lilien! Den er tankefuld, den hælder lidt med Hovedet, den ryster paa Hovedet, det er til Sorgen: ja, ja imorgen. Og imorgen, saa har Lilien lovligt Forfald, den er ikke hjemme, den er borte, Keiseren har tabt sin Ret, hvis han havde nogen, og Sorgen kan lige saa gjerne strax rive Fordringen istykker, den gjælder ikke; og derved bliver det, om Sorgen saa blev rasende og siger: det gjælds ikke. O, saaledes at kunne sige til Sorgen: ja, ja imorgen; og saa at kunne blive saaledes rolig paa Stedet, yndig i sorgløs Glæde, om muligt gladere ved at have sin Spøg med Sorgen: imorgen! Ikke at narre den et Par Dage, en Uges-Tid; nei, at vedblive at sige til Sorgen, hver Gang den melder sig: det er for tidligt, Du kommer for tidligt, at sige det saa længe, indtil det saa naar den kommer, er – for silde! O, Mesterskabet i at leve! Man næsten gyser, medens man beundrer Mesteren; man næsten gyser, thi det er jo Liv og Død om at gjøre. Man næsten gyser, og dog nei, Mesterens Kunst er saa stor, at man – utaknemligt at være den store Kunstner! – end ikke mærker den mindste Gysen, men giver sig hen deri som i den yndigste, den lykkeligste Spøg.

Agt derfor paa Lilien og Fuglen! Visseligen, der er Aand i Naturen – især naar Evangeliet beaander den; thi da er Naturen idel Sindbillede og idel Lærdom for Mennesket, ogsaa den indblæst af Gud, og »nyttig til Lærdom, til Rettelse, til Optugtelse«. »Betragter Lilierne paa Marken; de sye ikke, de spinde ikke« – og dog, den bedst lærte Syerske, der syer om sig selv, og en Prindsesse, der af de kostbareste Tøier lader sye hos den bedst lærte Syerske, og Salomo i al sin Herlighed var ikke klædt som en af dem. Saa er der altsaa En, der spinder og syer for Lilierne? Det er der: Gud i Himlene. Men Mennesket han spinder og syer. »Ja, det lærer Nøden ham nok; Nød lærer nøgen Kone at spinde.« »Aah, pfui, at Du kan tænke saa ringe om Dit Arbeide, om det at være Menneske, om Gud, om Tilværelsen – som var det Hele en Straffeanstalt! Nei, betragt Lilierne paa Marken, lær af dem, lær at forstaae hvad Du veed: Du veed, at det er Mennesket, der spinder og syer, lær af dem at forstaae, at det dog egentligen, ogsaa da naar Mennesket spinder og syer, dog er Gud der spinder og syer. Troer Du, at Syersken, hvis hun forstaaer dette, derfor vil blive mindre flittig til eller ved sit Arbeide, at hun vil lægge Hænderne i Skjødet og tænke: naar det dog alligevel egentligen er Gud, der spinder og syer, saa er det bedst, at jeg bliver fri, fri for denne uegentlige Spinden og Syen? Hvis saa, da er denne Syerske en taabelig lille Jomfrue, for ikke at sige en næsvis Tøs, af hvem Gud ingen Glæde kan have, og som selv ingen Glæde kan have af Lilierne, og som ikke var bedre værd, end at vor Herre satte hende ud af Døren; og saa kan hun see, hvor hun blev af. Men hun, vor egen, kjere, vor barnlig fromme, vor elskelige Syerske, hun forstaaer, at blot naar hun selv syer, er det Gud vil sye for hende; og derfor bliver hun kun desto flittigere til sit Arbeide, at hun ved stadigt at sye, stadigt maa forstaae, at det – livsalige Spøg! – er Gud der syer, hvert Sting; at hun, ved stadigt at sye, stadigt maa forstaae, at det – o, Alvor! – er Gud der syer, hvert Sting. Og har hun forstaaet dette, underviist af Lilien og Fuglen, da har hun fattet Livets Betydning, og hendes Liv i høieste Forstand været betydningsfuldt; og naar hun engang er død, vil der med det størst mulige Eftertryk sandt kunne siges ved hendes Grav: hun har levet; om hun blev gift eller ikke, er ikke afgjørende.

»Seer til Fuglene under Himlen!« Hvorledes, Du er bekymret, Dit Sind nedbøiet, Dit Øie vendt mod Jorden! Hvad er dette? Saaledes skabte Gud ikke Mennesket; det maa Du jo vide af enhver Børnebog. Det der udmærker Mennesket frem for Dyret er den opreiste Gang. Altsaa, vær saa god, op med Hovedet! »O, lad mig blot med Fred!« »Nei, lad os gaae lempeligt frem. Det vilde maaskee ogsaa være for stærk en Bevægelse for Dit syge Sind, for brat en Overgang, om Du pludseligt fra Jorden saae op til Himlen. Saa lad os tage Fuglen til Hjælp. Den sidder ved Jorden, hvor Dit Blik er. Nu løfter den sig – den Smule kan Du vel ogsaa nok taale at hæve Dit Hoved, saa Dit Blik følger den. Den stiger – saa løft Dit Hoved lidt endnu; og lidt endnu. Nu er det godt: nu er Fuglen høit under Himlen – og Du rigtigt stillet: see til Fuglen under Himlen – o, og tilstaae Dig selv, at saa lidet som Himmelhvælvingen kan siges at trykke, saa lidet er Gud Den, der nedtrykker, nei, det Nedtrykkende kommer fra Jorden, eller fra det Jordiske i Dig; men som Himmelhvælvingen opløfter, saaledes er Gud Den der vil opløfte. »Seer til Fuglene under Himlen; de saae ikke, de høste ikke, de sanke ikke i Lade.« Dog lever Fuglen jo ikke af Luften, saa lidet som vi Mennesker. Der maa altsaa være En, der saaer og høster og sanker i Lade for dem? Det er der ogsaa, Gud, den store Forsørger, eller Forsyner, eller som vi ogsaa kalde ham: Forsynet. Han saaer og høster og sanker i Lade, og hele Verden er som hans store Forraadskammer. Kjedsommelige Mennesker have haft den kjedsommelige Tanke, for at kunne undvære Gud at ville gjøre hele Verden til een stor Ladebygning. Det er en daarlig Efterabelse. Nei, naar det er Gud der gjør det, saa er det fornøieligt – hvor fornøiet er ikke Fuglen under Himlen, som ikke saaer, ikke høster, ikke sanker i Lade! Men det gjør Mennesket, han pløier, han saaer, han høster og sanker i Lade. Lær saa blot af Fuglen under Himlen at forstaae hvad Du veed: Du veed, at det er Mennesket der saaer og høster – lær at forstaae, at det, naar Mennesket gjør det, dog egentligen er Gud som gjør det. »Hvilken Snak! Naar jeg i mit Ansigts Sveed gaaer ude paa Marken, og høster, saa Sveden hagler ned af mig: saa veed jeg dog nok med Vished, at det er mig, der høster, idetmindste er det mig, der sveder. Eller er det maaskee ogsaa egentligen Gud der sveder; eller hvis det er Gud der høster, hvorfor skal jeg saa svede? Din Tale er noget høitravende, upraktisk Vaas.« »Menneske, Menneske, forhærdede Menneske-Forstand, vil Du da aldrig af Fuglen lære at gaae fra Forstanden for at blive Menneske! Vil Du aldrig lære, i gudelig Opløftelse, som Fuglens, at forstaae: det at arbeide. Du vil vistnok endog blot saaledes komme Sandheden langt nærmere end Du er, ved engang at betragte Sagen omvendt, at Arbeide ikke saadan er Møie og Besvær, som man helst var fri for, at meget mere Gud har indrømmet Mennesket at kunne arbeide for at unde ham en Fornøielse, en Selvstændigheds-Følelse, der ikke kan kjøbes for dyrt med Ansigtets Sveed – thi om man sveder eller ikke, kan da ikke gjøre Udslaget, en Dandsende sveder jo ogsaa, man kalder ikke derfor det at dandse et Arbeide, Møie og Besvær. Dette er den ene gudelige Forstaaelse af det at arbeide – og saa er man saare langt fra at klage over Ansigtets Sveed. Tag et Barn og Forældrene i Forhold til det. Lille Ludvig, han bliver hver Dag kjørt en Tour i sin Barnevogn, en Fornøielse der gjerne varer en Time; og at det er en Fornøielse, forstaaer lille Ludvig godt. Og dog har Moderen hittet paa noget Nyt, som bestemt vil fornøie lille Ludvig endnu mere: om han ikke selv kunde trække Vognen? Og han kan! Hvorledes, han kan? Ja, see Tante, lille Ludvig kan selv trække Vognen! Lad os nu være Mennesker, ikke forstyrre Barnet! thi vi veed jo nok, at lille Ludvig ikke kan, at det egentligen er Moderen der trækker Vognen, og kun for ret at fornøie ham, er det hun leger den Leeg, at lille Ludvig kan selv. Og han, han puster og stønner. Sveder han maaskee ikke? Jo, min sandten sveder han, Sveden staaer paa hans Pande, i Ansigtets Sveed trækker han Vognen – men hans Aasyn er glædestraalende, glædedrukkent kunde man kalde det, og bliver det, om muligt endnu mere, for hver Gang Tanten siger: nei, see, lille Ludvig kan selv. Det var en mageløs Fornøielse. Det at svede? Nei, det at kunne selv. Saaledes med det at kunne arbeide; ret forstaaet, gudeligt forstaaet, er det idel Fornøielse, Noget Gud har hittet paa for at fornøie Mennesket, Noget hvorom Gud har sagt til sig selv: det vil bestemt fornøie ham meget mere end bestandigt at blive kjørt i Barnevognen. Det er Forestillingen der gjør Udslaget her som i Alt. Naar det er for Din Lyst for Fornøielses Skyld, saa klager Du ikke over at svede: vel, saa lad Dit Arbeide være Din Lyst, forstaae det saaledes, at det er Noget Gud har hittet paa for at fornøie Dig, o bedrøv ikke hans Kjerlighed, han troede, at det ret vil glæde Dig! Dette er en gudelig Forstaaelse af det at arbeide. – Men der er en endnu høiere gudelig Forstaaelse, hvilken vi lære af Fuglen: at det saa dog igjen er Gud, der arbeider, Gud der saaer og høster, naar Mennesket saaer og høster. Tænk paa lille Ludvig! Han er nu blevet Mand, og forstaaer altsaa nu godt Sammenhænget, at det var Moderen der trak Vognen; han har derfor nu en anden Glæde af denne Barndoms-Erindring: at tænke paa Moderens Kjerlighed, der saaledes kunde hitte paa Noget for at fornøie Barnet. Men nu er han Mand, nu kan han virkelig selv; han føres nu maaskee endog i Fristelse af dette: virkelig at kunne selv – indtil hiin Erindring fra Barndommen minder ham om, hvorvidt han ikke endnu, i en langt høiere Forstand, er i samme Tilfælde som Barnet, at det, naar Manden arbeider, dog egentligen er en Anden, er – Gud, der arbeider. Troer Du, han derfor vil blive uvirksom, lægge sig paa den lade Side og sige: naar det dog er Gud der arbeider, saa er det bedst, at jeg bliver fri? Hvis saa, saa er denne Mand en Nar, for ikke at sige en uforskammet Slyngel, af hvem Gud ingen Glæde kan have, og som ikke selv kan glæde sig ved Fuglen, og som ikke var bedre værd, end at vor Herre satte ham paa Porten; og saa kan han see, hvor han bliver af. Men den brave, retsindige, gudfrygtige Arbeider, han bliver kun desto stræbsommere, for desto stadigere at forstaae, at – livsalige Spøg! – at Gud er Medarbeider – o, Alvor! Skabt i Guds Billede som han er, med opløftet Hoved, seer han mod Himlen efter Fuglen – Spøgefuglen, af hvem han lærer Alvoren, at det er Gud, der saaer og høster og sanker i Lade. Men han synker ikke hen i Uvirksomhed, han seer strax til, han passer sit Arbeide – ellers faaer han jo ikke at see, at det er Gud der saaer og høster og sanker i Lade.

Du Lilie paa Marken, Du Fugl under Himlen! Hvad skylder et Menneske dog ikke Dig? Nogle af sine bedste og saligste Timer. Thi da Evangeliet indsatte Dig til Forbillede og Læremester, da blev Loven afskaffet, og Spøgen anviist sit Sted i Himmeriges Rige, saa vi ikke mere ere under Tugtemesteren men under Evangeliet: Betragter Lilierne paa Marken, seer til Fuglene under Himlen!

Men saa bliver maaskee det Hele om at følge Christum efter, om Efterfølgelsen, det bliver saa maaskee en Spøg? Han hjalp os selv ved ikke at sige: seer paa mig, men: betragter Lilierne seer til Fuglene; han viste bort fra sig, og vi – ja, det var da ikke at fortænke os Mennesker i – vi toge kun altfor gjerne mod Vinket, kløgtige som vi alle ere, naar det gjælder om at spare Kjød og Blod, kløgtige forstode vi kun altfor godt, hvad der blev os indrømmet ved at have saadanne Forbilleder; og vi bleve uudtømmelige i at pynte dette ud, kun med en vis hemmelig Gru tænkende paa Alvoren: Christi Efterfølgelse.

Nei, ikke ganske saaledes kan det tillades os, det vilde være at gjøre Evangeliet saa let, at det i Grunden – hvad just »Christi Efterfølgelse« er som beregnet paa at forhindre – blev Poesie.

Thi vistnok kan Lilien og Fuglen med Sandhed siges kun at tjene een Herre; men det er dog kun sindbilledligt, og Forpligtelsen for Mennesket til her »at følge efter« et digterisk Udtryk, ligesom ogsaa Lilien og Fuglen, som Læremestere betragtede, i strengeste Forstand ere uden Myndighed. Fremdeles, hvis et Menneske, med Lilien og Fuglen som Forbillede, levede som vi i det Foregaaende have fremstillet det, saa han i Alt tænkte Tanken om Gud ind med: saa er dette vistnok Fromhed, og en Fromhed der, ganske reen, vel aldrig er seet blandt Mennesker. Men i strengeste Forstand Christendom er dette endnu ikke; det er egentligen jødisk Fromhed. Det for Christendommen afgjørende er slet ikke udviist her: det at lide fordi man holder sig til Gud, eller, som det hedder, at lide for Læren – den egentlige Christi Efterfølgelse.

Ak ja, det er som aldeles glemt i Christenheden, hvad Christendom er. Naar En blot nogenledes vil fremstille den, er det ikke langtfra, at man indbilder sig, at det er en Grusomhed, et Menneske-Plagerie, han selv hitter paa – i den Grad følger nøiagtigt Lidelse for Ordet eller Læren med Christendommen, i den Grad, at naar En blot nogenledes sandt fremstiller den, vil han strax paadrage sig menneskelig Unaade. Og som sagt, trods de Millioner Exemplarer af det nye Testamente, der ere i Omløb, og trods at Enhver eier det nye Testamente, er døbt, confirmeret, kalder sig Christen, og trods at 1000 Præster prædike hver evige Søndag – det vil dog ikke være langtfra, at man paastaaer, at det er det Menneskes egen Opfindelse, naar han, ganske simpelt, fremdrager af det Nye Testamente Det, som tydeligt nok staaer der og med klare Ord, men som vi Mennesker fra Slægt til Slægt, høist ugenerede have udeladt, uden derfor at lade som Andet, end at det vi har beholdt under Navn af Christendom, er den pure, den sunde, uforfalskede Lære.

»Efterfølgelsen, Christi Efterfølgelse« er egentligen det Punkt, hvori Menneske-Slægten ømmer sig, her stikker hovedsagelig Vanskeligheden, her er det egentligen det bliver afgjort, om man vil antage Christendommen eller ikke. Trykkes der paa dette Punkt, trykkes der stærkt: i samme Grad færre Christne. Slaaes der af paa dette Punkt (saa Christendom bliver, intellectuelt, en Lære), gaaer Flere ind paa Christendommen; afskaffes det ganske (saa Christendom bliver, existentielt, let som Mythologie og Poesie, Efterfølgelsen en Overdrivelse, en latterlig Overdrivelse), saa breder Christendommen sig i den Grad, at Christenhed og Verden næsten falde sammen, eller Alle blive Christne, Christendommen har fuldkommen seiret, det vil sige, den er afskaffet.

O, at man itide havde passet paa dette, saa var Tilstanden i Christenheden en anden end den nu er. Men da den menneskelige Paastaaelighed, betræffende ikke at ville vide Noget af denne Snak om Efterfølgelse, blev mere og mere truende; da Leiesvende og Menneske-Trælle eller dog kun meget svagt Troende paatoge sig at være Ordets Forkyndere: saa blev Historien i Christenheden, fra Slægt til Slægt en stadig Slaaen af paa Prisen af det at være Christen. Til Sidst blev det en saadan Spotpriis, at vel snart den omvendte Virkning indtraadte, at Menneskene neppe gad have med Christendommen at gjøre, fordi den ved denne usande Mildhed var blevet saa vammel, at det æklede. At være Christen – ja, naar man da blot ikke bogstavelig stjæler, ikke bogstavelig gjør Tyverie til sin Næringsvei; thi at være Tyv i sin Næringsvei, det kan godt forenes med at være en alvorlig Christen, der gaaer til Alters en Gang om Aaret og et Par Gange om Aaret eller dog bestemt Nytaarsdag i Kirke. At være Christen – ja, naar man da i at bedrive Hoer ikke overdriver, eller, forladende den gyldne Middelvei, gaaer til Yderlighed; thi iagttagende – Decorum! – det vil sige skjult, med Smag og Dannelse, kan det endnu nok forenes med at være en alvorlig Christen, der idetmindste een Gang hører en Prædiken for hver 14 Gange han læser Komedier og Romaner. Og at der skulde være Noget i Veien for, at man, paa enhver Maade aldeles bestemmede sig i Lighed med Verdens Skikkelse, ved ethvert klogt Middel søgte at sikkre sig den størst mulige jordiske Fordeel, Nydelse og saa videre – at der skulde være Noget i Veien for, at det ypperligt lod sig forene med at være en alvorlig Christen: det vilde da være en latterlig Overdrivelse, en Uforskammethed, om Nogen vovede at byde os Sligt, en grændseløs Taabelighed af Den, der vovede det, da der dog ikke var eet eneste Menneske, der vilde reflectere paa – hvad der staaer eller at det staaer i det nye Testamente. Dette er en wohlfeil Udgave af det at være Christen, dog er det Virkelighedens Status; thi at Præster om Søndagen i en stille Time declamere om de høie Dyder og saa videre forandrer ikke Virkelighedens Status om Mandagen, da man forklarer sig en saadan Forkyndelse af, at det er Præstens Embede og Levebrød, og da mangen Præsts Liv vel ikke er anderledes end hiin Virkelighedens Status – men det er egentligen Existentsen der prædiker, det med Munden og Armene har ingen Art.

Imidlertid vare der jo ogsaa Dem, som holdt Christendommen høiere i Priis, men aldrig høiere end omtrent til hiin stille Fromhed, der under Naadens Mildhed oftere tænker paa Gud, venter alt Godt af hans faderlige Haand, søger Trøst hos ham i Livets Nød.

»At lide for Læren« – Christi Efterfølgelse: det er aldeles afskaffet, længst, længst bragt i Glemsel. Forsaavidt man i Prædikeforedraget ikke godt ganske kan undgaae at tale om Efterfølgelsen (om Nogle dog end har vidst, at indrette det saaledes, at det lod sig gjøre) saa gjør man det saaledes, at det egentlig Afgjørende forties, og der sættes noget Andet istedet: at man bør finde sig i Livets Gjenvordigheder med Taalmod og saa videre.

Men »Efterfølgelsen« er afskaffet. Den bestaaende Christenhed vilde vistnok, dersom den da ellers for Latter kunde komme til at høre, falde i den dybeste Forundring, hvis den fik at høre, at dette er det nye Testamentes og alle de sande Christnes Lære efter det nye Testamente, at til at være sand Christen hører at lide for Læren. At lide for Læren – i den Grad kun at tjene een Herre, saaledes at følge Forbilledet efter, at man lider for at være Christen! At lide for Læren – »nei nu troer jeg« vilde vistnok Christenheden sige »nu troer jeg at det Menneske reent er gaaet fra Forstanden, at fordre at man skal lide for Læren: i den Grad at forfalde til Christendom er da meget galere end at forfalde til Spil, Drik, Hoer. Lad det være som Præsterne forkynde det, at Christendommen er den milde Trøst, en Slags Assurance for Evigheden: det lader sig høre, det kan man være villig til at give sine Penge for – og maaskee er det dog dyrt nok med de høie Præstepenge, saa vi forsaavidt kunne siges at lide for Læren. Men at skulle betale for at faae det forkyndt, at man skal lide for Læren: det Menneske er pinegal.« Dog ligger Skylden ikke hos ham, det »Pinegale« er egentligen, at man i at forkynde Christendom har udeladt og fortiet hvad der ikke behager det verdslige og jordiske Sind, og saa foranlediget hele denne Verdslighed til at bilde sig ind at være Christen.

O, at der dog var blevet holdt igjen paa dette Punkt om Efterfølgelsen! At man, belært af tidligere Tiders Vildfarelser, havde sandt holdt igjen paa dette Punkt! Det er ikke skeet. Saa maa det skee. »Efterfølgelsen«, hvortil svarer: Christus som Forbillede, maa, hvis der skal blive Mening i Christenheden, igjen anbringes, men, som sagt, saaledes at man har lært Noget af tidligere Tiders Vildfarelse.

Uden at anbringe »Efterfølgelsen« er det umuligt at faae Magt med Tvivlene. Derfor er ogsaa Tilstanden i Christenheden: at have sat Tvivl istedetfor Tro. Og saa vil man standse Tvivlen med – Grunde; og endnu er man ikke standset i denne Retning, endnu har man ikke lært, at det er spildt Møie, ja at det er at nære Tvivlen, at give den Grund til at vedblive, endnu er man ikke opmærksom paa, at »Efterfølgelsen« er den eneste Magt, der, som en Politie-Magt, kan splitte alle Tvivlenes Opløb, og gjøre lyst, og tvinge til, hvis man ikke vil være »Efterfølgeren«, idetmindste at gaae hjem og holde sin Mund.

»Efterfølgelsen«, hvilket svarer til »Christus som Forbillede«, maa drages frem, anbringes, bringes i Erindring.

Lad os tage denne Sag fra Grunden, dog i al Korthed. Verdens Frelser vor Herre Jesus Christus er ikke kommet til Verden for at bringe en Lære; han har aldrig doceret. Da han ikke har bragt en Lære, har han da heller ikke søgt ved Grunde at bevæge Nogen til at gaae ind paa Læren, ikke ved Beviser at godtgjøre den. Hans Lære var væsentligen hans Liv, hans Tilværelse. Vilde Nogen være hans Discipel, da var, hvad jo ogsaa sees af Evangeliet, hans Fremgangsmaade en anden end den docerende. Han sagde til en Saadan omtrent: vov en afgjørende Handling, saa kan vi begynde. Hvad vil det sige? Det vil sige, Christen bliver man ikke ved at høre noget om Christendom, at læse Noget derom, at tænke Noget derover, eller ved, medens Christus levede, at see ham engang imellem eller at gaae og gabe paa ham hele Dagen: nei, der fordres en Bestedelse (Situation) – vov en afgjørende Handling; Beviset gaaer ikke forud men følger efter, er i og med Efterfølgelsen, der følger Christum efter. Naar Du nemlig har vovet den afgjørende Handling, saa Du bliver ueensartet med denne Verdens Liv, ikke kan have Dit Liv deri, støder sammen dermed: saa vil Du efterhaanden bringes i en saadan Spænding, at Du vil kunne blive opmærksom paa, hvad det er jeg taler om. Maaskee vil ogsaa Spændingen virke saaledes paa Dig, at Du forstaaer, at Du ikke kan holde den ud uden at tye til mig – og saa kan vi begynde. Kunde man vente det anderledes af »Sandheden«? Maatte den ikke udtrykke, at det er den Lærende der trænger til Læreren, »den Syge der trænger til Lægen«, ikke omvendt, som Christendommen i en senere Tid blev forkyndt, saa det var »Lægen der trænger til Patienterne«, Læreren der trænger til de Lærende, og derfor, selvfølgeligt, maa, som en anden Sælger, der jo ikke forlanger, at et høistæret Publikum skal kjøbe Katten i Sækken, være til Tjeneste med Grunde, Beviser, Recomendationerne fra Andre, som ere blevne hjulpne eller underviste og saa videre. Men den guddommelige Sandhed! Dog at den bærer sig anderledes ad, har ikke sin Grund i hvad man kun kalde guddommelig Fornemhed. Nei, nei, i den Henseende var Verdens Frelser vistnok villig til at fornedre sig som i Alt; men det kan ikke være anderledes.

Vi ville nu ikke dvæle ved, hvorledes Christendommen efterhaanden udbredte sig i Verden, vi haste hen til et bestemt Punkt, der er afgjørende for Tilstanden i den nuværende Christenhed.

Vi standse et Øieblik ved Middelalderen. Hvor store dens Vildfarelser end vare, dens Opfattelse af Christendommen har eet afgjørende Fortrin for vore Tiders. Middelalderen opfattede Christendommen i Retning af Handling, Liv, Existents-Omdannelse. Dette er det Værdifulde. Noget Andet er, at det var nogle besynderlige Handlinger, den hittede paa, at den kunde mene, at det at faste i og for sig var Christendom, det at gaae i Kloster, det at give Alt til de Fattige, for ikke at tale om, hvad vi nu neppe kunne omtale uden at smile, at det at hudflette sig selv, at gaae paa Knæerne, at staae paa eet Been og saa videre, at det skulde være den sande Efterfølgelse. Dette var Vildfarelse. Og som det gaaer, naar man er svinget ind paa en forkeert Vei og derpaa stræber fremad, at man fjerner sig længere og længere fra det Sande dybere og dybere ind i Vildfarelsen, at det bliver værre og værre: saaledes her. Hvad der var værre end den første Vildfarelse udeblev ikke: at man hittede paa det Fortjenstlige, meente at have ved sine gode Gjerninger Fortjeneste for Gud. Og det blev værre, man meente endog i den Grad at have Fortjeneste af sine gode Gjerninger, at det ikke blot kom En selv til Gode, at man som Kapitalist og Selvskyldner kunde lade det komme Andre til Gode. Og det blev værre, det blev en reen Handel; Mennesker, der aldrig en Gang selv havde tænkt paa dog at tilveiebringe nogle af disse saakaldte gode Gjerninger, fik dog nu fuldt op med gode Gjerninger at gjøre, forsaavidt de bleve satte i Virksomhed som Kræmmere, der for Penge solgte Andres gode Gjerninger til bestemte men billige Priser.

Da fremtræder Luther. Denne Tilstand, siger han, er Aandløshed. En forfærdelig Aandløshed, ellers maatte De, der ved gode Gjerninger mene at fortjene Himlens Salighed, see, at dette er den visse Vei enten til Formastelse, altsaa til at gaae Glip af Saligheden, eller til Fortvivlelse, altsaa til at gaae Glip af Saligheden. Thi at ville bygge paa gode Gjerninger – jo flere Du øver, jo strengere Du er mod Dig selv, desto mere udvikler Du blot i Dig Angest, og ny Angest. Ad den Vei – hvis et Menneske ikke er ganske aandløst, ad den Vei kommer han kun til lige det Modsatte af Fred for sin Sjel og Hvile, til Ufred og Uro. Nei, et Menneske retfærdiggjøres ene ved Troen. Og derfor i Guds Navn ad Helvede til med Paven og alle hans Hjælperes Hjælpere; og bort med Klosteret, samt al Eders Fasten, Pidsken, og alt det Abespil, som er brugt op under Navn af Efterfølgelse.

Men lad os ikke glemme, derfor afskaffede Luther ikke Efterfølgelsen, ei heller det Frivillige, hvad Kjælenskabet saa gjerne vil bilde os ind om Luther. Han anbragte Efterfølgelsen i Retning af at vidne for Sandhed, og (dog uden at indbilde sig, at det var fortjenstligt) udsatte sig der frivilligt for Farer nok. Det var jo ikke Paven, der angreb Luther, det var Luther, der angreb Paven; og Luthers Liv var, om han end ikke blev henrettet, dog et, menneskelig talt, offret Liv, et Liv offret til at vidne for Sandhed.

Den nuværende Christenhed, idetmindste den jeg taler om, slutter sig til Luther; noget Andet er, om Luther kunde vedkjende sig den, om ikke det Sving, som Luther gjorde, kun altfor let kan blive ind paa en forkeert Vei, saasnart der ikke er en Luther med, hvis Liv gjør det sande Sving til Sandhed. I ethvert Tilfælde bliver det vistnok rigtigst, dersom man vil see, om der ikke er Misligheder i den nuværende Tilstand, at søge tilbage til Luther og det Sving han gjorde.

Vildfarelsen som Luther svingede af fra var: Overdrivelse i Retning af Gjerninger. Og ganske rigtigt, han tog ikke Feil: et Menneske retfærdiggjøres ene og alene ved Troen. Saaledes talte og lærte – og troede han. Og at dette ikke var at tage Naaden forfængeligt, derom vidnede hans Liv. Ypperligt!

Men allerede den næste Slægt sagtnede, den svingede ikke med Rædsel af fra Overdrivelse i Retning af Gjerninger (hvad Luther havde levet i) ind i Troen. Nei, den gjorde det Lutherske til Doctrin; og saaledes satte ogsaa Troen af i Livs-Kraft. Saa sattes der af fra Slægt til Slægt; Gjerninger – ja, det veed Gud, det var der rigtignok ikke mere Tale om, det var Synd at beskylde denne senere Tid for Overdrivelse i Retning af Gjerninger, og saa fjantet var man da heller ikke, at man formastede sig til at ville have Fortjeneste af hvad man fritog sig for at gjøre. Men nu Troen – mon den findes paa Jorden?

\workheading{Bogen om Adler (af Petrus Minor)}{1846–55 (udg. posthumt)}

\worknote{Antologiens ''The Book on Adler'', om Myndighedens begreb med Magister Adler som anledning. Kierkegaard omarbeidede manuskriptet mere end noget andet af sine; Hong-oversættelsen (KW 24) følger den tredje integrale version (Pap. VII² B 235), og antologiens engelske uddrag er selv sammenstillet af to vidt adskilte steder i den version: fra Capitel II (''En Aabenbaring i Situationen i Nutiden'') og fra Tillæget ''Det dialektiske Forhold'' (Addendum I). Vor kilde er samme version-III-tekst (fra SKS' TEI-udgave af ''Bogen om Adler''), med samme ordlyd stykke for stykke; forskellen er kun redaktionel placering: det stof Hong sætter som selvstændigt Tillæg I (''Det dialektiske Forhold'') staaer her i Capitel I (''Den historiske Situation''). For sammenligning gengives begge Hongs kildesteder i fuld udstrækning, i hans rækkefølge: hele Capitel II-stykket (''Det Christelige har ingen Historie …''), derefter Capitel I-stykket om det dialektiske Forhold mellem det Almene, den Enkelte og den særlige Enkelte (den sande Extraordinaire), endende med Berserkegangen mellem Skjolde. Spring mellem de to steder markeret med […].}

\sectionheading{Det Christelige har ingen Historie (Cap. II); Den sande Extraordinaire (Cap. I / Det dialektiske Forhold)}

[Af Capitel II — »En Aabenbaring i Situationen i Nutiden«]

Det Christelige har ingen Historie, thi det Christelige er dette Paradox, at Gud eengang er blevet til i Tiden. Dette er Anstødet men ogsaa Udgangspunktet; og hvad enten det er 1800 Aar siden eller det er igaar: man kan lige godt være samtidig med det. Som Polar-Stjernen aldrig forandrer sin Stilling, og derfor ingen Historie har: saaledes staaer dette Paradox urokket og uforandret; og om Christendommen bestod i 10,000 Aar man er i afgjørende Forstand ikke kommen længere fra det end de Samtidige vare. Thi Afstanden er ikke til at maale med Tidens og Rummets Quantiteren, da den er den qvalitativt-afgjørende, at det er et Paradox. Christendommens Historie forholder sig heller ikke ligefremt til dette Christelige, som naar Træets Bestaaen ved dets Voxen forholder sig til Spiren. Det Christelige er et evigt Afsluttet som der ikke skal sættes Noget til eller tages Noget fra, og i hver Generation og i hvert Individ, hvis han i Sandhed er Christen, begyndes der for fra, fra hiint Paradox. Der begyndes ikke der, hvor den foregaaende Generation slap, men for fra.

Saasnart man derimod forvexler Christendommen og det Christelige, saasnart man begynder at tælle Aarene: saa begynder man at ville forandre det Usandsynlige til det Sandsynlige. Man siger: nu har Christendommen bestaaet (thi det Christelige er jo et Faktum for 1800 Aar siden) i 300 Aar, nu i 700, nu i 1800: ja saa maa den vist være det Sande. Ved denne Fremgangsmaade vinder man at forvirre Alt. Afgjørelsen (den at blive Christen) bliver for den Enkelte let en reen Ubetydelighed; falder det ham allerede let nok at antage Skik og Brug i den By, hvor han lever, fordi de Mange gjøre det: skulde det saa ikke falde ganske af sig selv at han blev Christen med

– naar Christendommen har bestaaet i 1800 Aar! Paa den anden Side bliver det Christelige svækket, gjort til en Ubetydelighed ved Hjælp af Afstanden, ved Hjælp af de 1800 Aar. Hvad der hvis det skete samtidigen vilde komme som Forfærdelse over Menneskene, vilde oprøre dem Tilværelsen fra Grunden, hvad de, hvis det skete samtidigen, enten vilde forarges paa vilde hade og forfølge og om muligt faae udryddet, eller troende tilegne det sig: det synes man nok man saadan kan antage og troe (det vil sige saadan lade staae hen) siden det er 1800 Aar siden. Thi Samtidigheden er Spændingen, der ikke tillader En at lade det staae hen, men tvinger En til enten at forarges eller tro; Afstanden derimod er det Aflad, der begunstiger Dorskhed i den Grad, at det troende at antage Noget saadan bliver identisk med at lade det staae hen. Hvoraf kommer det vel, at ingen Samtid kan komme ud af det med Sandhedens Vidner, og neppe er Manden død saa kan Alle saa ypperligt komme ud af det med ham? Det kommer der af, at saa længe han lever og de Samtidige med ham i Samtidighedens Situation: saa føle de Braaden af hans Existents, han tvinger dem til en mere anstrænget Afgjørelse. Men naar han er død, saa kan man godt være gode Venner med ham og beundre ham det vil sige saadan tankeløst og mageligt lade det Hele staae hen. Hvorfor mon Socrates sammenlignede sig selv med en Bremse, uden fordi han forstod, at hans Leven med de Samtidige var en Braad. Da han var død, forgudede de ham. Naar en Mand oplever en lille Begivenhed i sit Liv: saa lærer han Noget deraf, og hvorfor? Fordi Begivenheden ret kommer ham paa Livet. Derimod kan den samme Mand sidde i Theateret og see Tragediens store Optrin, han kan læse i Aviser om det Overordentlige, han kan høre Præsten: og det gjør Alt ikke ret noget Indtryk, og hvorfor? Fordi han ikke bliver samtidig med det, fordi han i de to første Tilfælde mangler Phantasie, i det sidste indre Erfaring for ret at blive samtidig med det Fremstillede; fordi han tænker som saa, det er nok mange Aar siden det skete.

Naar nu saaledes gjennem hele Aarrækker en desorienteret Orthodoxie, som ikke veed hvad den gjør, og en oprørsk Heterodoxie, der dæmonisk veed hvad den gjør og kun forsaavidt ikke veed hvad den gjør, have forenet sig om ved Hjælp af de 1800 Aar at forvirre Alt, at forskylde det ene Sandsebedrag galere end det andet, den ene Paralogisme værre end den anden, den ene μεταβασις εις αλλο

γενος mere confunderende end den anden: saa er Hovedsagen nu at kunne gjøre Terrainet ryddeligt, faae de 1800 bort, saa det Christelige foregaaer for os, som skete det idag. Det der har pustet Indvendingerne mod Christendommen og Forsvaret for den op til volumina er de 1800 Aar. Det der har bedøvet Forsvarerne og hjulpet Angriberne er de 16, de 17, de 1800 Aar. Det der har holdt Utalliges Liv hen i en Indbildning er de 1800 Aar. Ved Hjælp af de 1800 Aar har Forsvarerne bagvendt gjort Christendommen til en Hypothese, og Angriberne gjort den til Ingenting.

Hvad den Intet mindre end travle Johannes Climacus i denne Henseende har gjort for at opspore ethvert Sandsebedrag, attrapere enhver Paralogisme, anholde enhver svigefuld Vending, kan her ikke gjentages. Det er gjort saaledes af ham, at enhver videnskabeligere Dannet, naar han for Alvor vil anvende lidt Tid paa dagligt at indøve sig i det Dialektiske, let vil forstaae det. Anderledes er det rigtignok ikke gjort, og anderledes kunde det heller ikke gjøres. Sligt kan ikke fremsættes i en Avis og læses »medens man tager Skjæget af«, eller mens man er paa Lokumet. Saadanne travle Mennesker, der kun have Tid at læse det Øieblik de ere paa Lokumet, og altsaa i det Høieste kun en Smule otium naar de engang imellem have Diarrhoe maa det overlades Aviser at skrive for. Climacuss Fremstilling er anstrængende, som Sagen fører det med sig. Hans Fortjeneste er den: ved Hjælp af Dialektik med Phantasie at have trukket (som man siger det om en Kikkert) hiint urokkelige Christelige saa nær for Øiet, at Læseren forhindres i at komme til at see feil af de 1800 Aar. Hans Fortjeneste er ved Hjælp af Dialektik at have skaffet vüet, Perspektivet. At rette sit Øie mod en Stjerne er ikke saa vanskeligt, fordi Luften er som et tomt Rum, og der derfor saa godt som Intet er underveis der standser eller adspreder Blikket. Anderledes derimod, naar Retningen som Øiet skal tage er lige ud som ad en Vei, og der da tillige er Trængsel og Mylder og Opløb og Larm og Travlhed, gjennem hvilken Øiet skal trænge sig, for at faae vüet, medens ethvert Sideblik, ja enhver Blinken med Øiet absolut forstyrrer qvalitativt; og endnu vanskeligere bliver det, naar man tillige skal staae i en Omgivelse der pro virili arbeider paa at forhindre En i at faae vüet. – Og dog er Samtidigheden med det urokkede Christelige det Afgjørende. Denne Samtidighed er imidlertid at forstaae i samme Betydning som de Medlevende havde den da Christus levede.

Det der behøves er fremfor Alt at faae de uhyre Bibliotheker og Skriverier og de 1800 Aar til en Side for at faae vüet. Og dette er ingenlunde en dumdristig Fordring af en høitflyvende Dialektiker, det er den ganske tarvelige og ægte religieuse Fordring, som Enhver maa gjøre – ikke for Videnskabens og Publikums Skyld; men for sin egen Skyld, reent personligt for sin egen Skyld, hvis det er ham Alvor med at blive Christen; og det er hvad Christendommen selv maa fordre. Thi det Christelige vil netop staae urokket som Polar-Stjernen, og vil derfor ud af al det Vrøvlerie, der kun tager Livet af det.

Dog er den Samtidighed, om hvilken her er Tale, ikke en Apostels Samtidighed, forsaavidt han blev kaldet ved en Aabenbaring, men er blot den Samtidighed, som enhver Medlevende havde, den Mulighed i Samtidighedens Spænding at maatte forarges eller gribe Troen. Og til den Ende er det netop nødvendigt at der luftes ud saaledes, at det, som eengang, bliver muligt, at et Menneske for Alvor kan forarges, eller troende tilegne sig det Christelige, at det ikke bliver med det Christelige som med en Retssag, naar den har staaet hen i umindelige Tider: at man hverken veed ud eller ind paa Grund af den megen Viden. Samtidighedens Situation er det Spændende, som giver Kategorierne qvalitativ Elasticitet; og det maa da være en stor Dosmer, der ikke veed, hvilken uendelig Forskjel det gjør, naar man for sin egen Skyld i Samtidighedens Situation betænker Noget, og naar man saadan tænker over Noget i den Indbildning, at det er 1800 Aar siden; i den Indbildning, ja i den Indbildning, thi netop fordi det Christelige er det qvalitative Paradox, er det en Indbildning, at 1800 Aar er længere siden end igaar.

Naar nu Forholdet er saaledes i Christenheden, at det netop gjelder om, at faae Ende paa denne seiglivede Dorskhed med de 1800 Aar: saa kan man ikke nægte, at det, at der pludseligen traadte en Mand frem og beraabte sig paa en Aabenbaring, kunde afgive et ønskeligt Incitament, thi saa er der dannet en analog Samtidighedens Situation. Jo jeg takker, alle de dybsindige og spekulative og lærde og svedende Prækere, der godt kan forstaae, at man for 1800 Aar siden fik en Aabenbaring – de vilde geraade i Forlegenhed. Den som overhovedet kan forstaae, at et Menneske faaer en Aabenbaring, han maa dog vel lige saa godt forstaae det, om det er skeet for 6000 Aar siden eller det skeer om 6000 Aar, eller det skeer idag. Men Prækeren har maaskee levet christeligt af de 1800 Aar, har præket sig ind i, at han kunde forstaae det – fordi det var 1800 Aar siden. Dersom Sagen ikke var saa alvorlig, kan jeg ikke nægte, at jeg vilde ansee det for den kosteligste Comedie der overhovedet nogensinde kunde blive digtet i Verden: at lade hele den moderne Exegese og Dogmatik gjennemgaae sit Cursus – i Samtidighedens Situation. Alle disse svigefulde psychologiske Paafund, alt dette til en vis Grad og saa igjen til en vis Grad, Bravour-Dybsindet, og saa fremfor Alt den stadselige Mediation som forklarer – da det er 1800 Aar siden Det skete, som forklares: alt Dette vilde tage sig ypperligt ud i Samtidighed med det Omforklarede. Det er ganske vist, at langt bedre end al videnskabelig Bekæmpen vilde een eneste aristophanisk Comedie i den Stiil rydde op i den moderne christelige Videnskabs Confusion.

Da jeg derfor, uden endnu selv at have seet hans Prædikener og disses Forord, hørte at MagisterAdler var traadt op og havde beraabt sig paa en Aabenbaring, kan jeg ikke nægte jeg blev forbauset; om jeg lever længe, vil jeg maaskee ikke let igjen blive saa forbauset. Thi jeg er ikke uden Kjendskab til hvad der gjærer i Tiden, jeg følger med, om end i særskilt Kahyt. Da jeg hørte det, tænkte jeg: enten er dette Manden, vi behøve, den Udvalgte, der i guddommelig Oprindelighed har det ny Kildespring til at forfriske Christenhedens trevne Jordbund, eller det er en Forarget, men en udstuderet Gavtyv, der for at sløife Alt ogsaa en Apostels Værdighed, for at styrte Alt sammen, bringer en saadan Christenhed som den nærværende i den anstrængende Afgjørelse, at skulde gjennemgaae sin Dogmatik i Samtidighedens Situation.

Under Antagelse af det Sidste vilde det jo rigtignok have overrasket mig, om en Forarget virkeligen havde været saa klog. Skjøndt man nemlig ikke kan nægte de Forargede Talent og dæmonisk Inspiration, saa pleie de dog i Almindelighed derhos i total Forstand at være noget dumme det vil sige de veed egentligen ikke rigtigt, hvordan man skal gribe Sagen an for at gjøre Skade. De angribe Christendommen, men de stille sig selv uden for den, og netop derfor gjør de ingen Skade. Nei den Forargede maa see at komme Christendommen ganske anderledes paa Livet, see at skyde, ligesom en Muldvarp, op midt i Christenheden. Sæt Feuerbach istedenfor at angribe Christendommen, var gaaet underfundigere tilværks, sæt han i dæmonisk Taushed havde lagt sin Plan, og derpaa var traadt frem og forkyndt han havde haft en Aabenbaring, og sæt han nu, som en Forbryder formaaer at fastholde en Løgn, urokkeligt havde holdt dette fast, medens han tillige kløgtigt havde afluret Orthodoxien alle dens svage Sider, som han dog langtfra angreb men kun med en vis troskyldig Naivitet vidste at holde for Lyset, sæt han havde gjort det saa godt, at Ingen kunde blive klog paa hans List: han vilde have bragt Orthodoxien i den største Forlegenhed. Orthodoxien kæmper i det Bestaaendes Interesse for, at vedligeholde det Skin, at vi saadan alle ere Christne, at Landet er christeligt og Menighederne Christne. Naar da En stiller sig udenfor og angriber Christendommen, saa skulde jo, hvis han seirede, Menigheden uleiliges ud af sit behagelige Tummerumme med saadan at være Christen som Folk er flest, den skulde ud i den Afgjørelse at opgive Christendommen. Hvilken Uleilighed; nei saa er det bedre at blive ved det Gamle. See derfor udretter den Forargede Intet. Fremdeles, naar En angriber Christendommen og stiller sig udenfor, saa forsvarer Orthodoxien den ved Hjælp af de 1800 Aar; den taler i de høie Toner om Guds overordentlige Foranstaltninger i sin Tid det vil sige for 1800 Aar siden. Og det maa man nu sige om det Overordentlige og om Guds overordentlige Foranstaltninger, det gaaer lettere i Folk, jo længere det er siden. Altsaa den Forargede angriber Christendommen, Orthodoxien forsvarer den ved Hjælp af Afstanden, og Menigheden tænker som saa: naar det er 1800 Aar siden saa kan man nok forstaae, at der er skeet noget Overordentligt. Den Forargede udretter igjen Intet. Anderledes derimod naar han selv snildt traadte op med en Aabenbaring, naar han forbandet vel studeret i Orthodoxien vidste at skjule sin dæmoniske Kløgt i en underlig troskyldig Naivitet, hvormed han bestandigt bragte Orthodoxien i Vinden, medens han som en Borre hang sig fast ved Orthodoxien. Paa den ene Side turde Orthodoxien ikke godt være bekjendt at nægte, at det var det Orthodoxe han foredrog, paa den anden Side var det dog yderst fatalt, at faae det sagt paa en saa ligefrem Maade, der tvang Orthodoxien til i en Samtidighedens Situation at rykke bestemt ud med Sproget.

Det er oftere sagt, at dersom Christus nu traadte frem – i Christenheden, hvis han i endnu strengere Forstand end fordum »kom til sine Egne«: saa vilde han igjen blive korsfæstet, det skulde da være at Dødsstraffen var afskaffet, men i saa Fald vilde han komme til at lide den Straf, der var traadt isteden for Dødsstraffen; og Orthodoxien vilde især være ivrig for at faae ham arresteret og dømt. Og hvorfor vilde det vel igjen gaae saaledes? Fordi Samtidigheden giver det behørige qualitative Tryk; Distance derimod hjælper baade til at gjøre Noget til Ingenting, og til at gjøre noget til det Overordentlige omtrent i Forstand af Ingenting. Hvorfor forargedes vel næsten Alle paa Christus da han levede, uden fordi det Overordentlige skete lige for deres Øine, saa Den, der vilde tale derom, maatte sige: igaaraftes, imorges, i Eftermiddags skete det. Men naar Miraklet er skeet for 1800 Aar siden – ja saa kan man sagtens forstaae, at det er skeet og at det var et Mirakel. Blandt Præste-Foredragets mange kostelige og uskateerlige Syllogismer maa denne ansees for een af de kosteligste: at Det, som man ikke kan forstaae, hvis det skete idag, det kan man forstaae og troe naar det er skeet for 1800 Aar siden, naar det vel at mærke er det Vidunderlige, der til enhver Tid paa Dagen baade Klokken 4 og Klokken 5 overgaaer menneskelig Forstand. Vil man nemlig kun sige, at man kan forstaae, at hine Mænd for 1800 Aar siden troede at det var et Mirakel: saa kan man lige saa gjerne sige lige ud, at man selv ikke troer det. Dog foretrækker man at hjælpe sig ved svigefulde Vendinger, som nu denne, der seer saa troende ud, og dog netop nægter Miraklet, da den siger om hine Mænd, at de troede det det vil sige at de meente det det vil sige, at man selv ikke troer det.

At tro i eminent Forstand svarer ganske rigtigt til det Vidunderlige, det Absurde, det Usandsynlige, det der er Forstanden en Daarskab, og netop derfor er det aldeles ligegyldigt hvor længe det er siden eller om det er idag. Enhver der har fjerneste Begreb om Dialektik i sin Pande maa dog indsee, at Den, der troer det, naar det er skeet for 1800 Aar siden, han kan lige saa godt tro det, naar det skeer idag – med mindre han troer det, fordi det er 1800 Aar siden, hvilket slet ikke er at tro. Troer han det og det er skeet for 1800 Aar siden, saa er han netop i Troen paradox samtidig med det som skete det idag. Hvad forøvrigt Præste-Foredraget paa dette Punkt præsterer i Galimathias: ja lad os ikke tale derom; ei heller om, hvorledes de menige Christne ved dette Præsteforedrag beroliges i deres Saligheds Sag. Thi det gjøre de, og derimod er Intet at sige, det skulde da være, at det nok i vor Tid baade var vigtigere men ogsaa vanskeligere at gjøre Menigheden lidt urolig og bekymret for sin Saligheds Sag. Bare dog den Dialektik ikke var til, den gjør kun Bryderier. Hvilken skjønnere Lovtale kunde der vel tænkes over et Lands Geistlighed, end at den beroliger Menighederne i deres Saligheds Sag. Og det gjør Geistligheden i vor Tid. Stundom høres der dog Klage i et eller andet af vore fortræffelige Blade over, at en Vægter raaber for høit og forstyrrer Beboernes Rolighed og Nattesøvn; men Geistligheden høres der ingen Klage over, den beroliger Menigheden i dens Saligheds Sag! Hvis Christus nu kom til sine Egne, han vilde vel finde Menighederne sovende – ved Geistlighedens Hjælp beroligede i deres Saligheds Sag.

Som sagt et Dilemma havde jeg tænkt mig i Forhold til Opfattelsen af det Overordentlige, at en Mand beraaber sig paa et Aabenbarings-Faktum: at han enten var den Udvalgte eller en dæmonisk-kløgtig Forarget. Adlers Optræden har imidlertid forvisset mig om, at der maa være et Tredie, samt om at jeg ved at tage ham for hvad han udgav sig for har været til Nar, thi han er ingen af Delene. At han ikke er den Udvalgte, at det Hele med hans Aabenbaring er Misforstaaelse, skal jeg nok senere vise og bevise. Langtfra mig være det, ligefremt, i Kraft af nogen min Anskuelse eller Theorie, at ville benægte Muligheden af at Adler kan have haft en Aabenbaring, nei men det lader sig af Adlers senere Attituder og senere Skrifter tilstrækkeligt godtgjøre, at han i Grunden ikke selv troer det, uden at han dog har fundet sig foranlediget til angrende at tilbagekalde sit Første. En dæmonisk-kløgtig Forarget er han om muligt endnu mindre, dertil er der ikke det mindste Spor eller Symptom. Derfor er han imidlertid ingenlunde uden Betydning, og jeg veed Ingen i min Samtid der i strængere Forstand maa kaldes et Phænomen end Adler. Tilværelsens Kræfter har haft fat i ham, og han er som Phænomen en Anticipation af den Dialektik der giærer i Tiden. Men Phænomenet selv veed Intet at forklare det vil sige man maa selv være Lærer for at lære Noget af Adler. Sandsynligviis uden nogen væsentlig christelig Forudsætning er han gaaet hen og blevet Præst –

da kommer Crisen pludselig paa ham, dobbelt forfærdelig, fordi han paa en saa høitidelig Maade jo allerede havde udgivet sig for at være Christen, da han har paataget sig at være christelig Præst. Som ved et Slag gaaer pludseligen det Christelige op i ham; i den anstrengede Situation producerer han de forskjelligste Yttringer af forskjelligartet christelig Orthodoxie og af Kjætterie mellem hinanden. Hvad Andre ved Læsning og Erfaring veed Beskeed om, det piner ligesom ud af ham, bryder primitivt frem, men aldeles confust. Men primitivt er det vistnok i ham, thi han har upaatvivleligen den Gang ikke læst det Mindste i Retning af den strængere Orthodoxie, eller Pietismen. Adler er saaledes ret egentligen et Tegn, han er et høist alvorligt Beviis for at det Christelige er en Magt, som ikke er til at spøge med. Men paa den anden Side er han snarere end en Udvalgt en Henhvirvlet, som er slynget hen liig en advarende Forfærdelse. Istedenfor at kunne hjælpe os Andre er han snarere som den bange fortumlede Fugl, der med Angestens Vingeslag styrter foran Uveiret, som følger, medens man endnu kun hører en Hvislen; og hans mange Tanker ere som den forvirrede Fugleskare der i mellem hinanden flygter foran Uveiret. At man derfor skulde være berettiget til at opgive ham eller tænke ringe om hans Mulighed er langt fra min Mening. Trods de forfærdelige Overhalinger han har faaet, trods al den Skade han upaatvivleligt vilde have gjort, hvis han ikke heldigviis var bleven ignoreret, er det idetmindste min Mening, at der i ham er meget Haab tilbage. Upaatvivleligen har han som Candidat levet hen i den Indbildning, at det fattige theologiske Kjendskab som examinando fordres ved theologisk Embeds-Examen at det var Christendom. Da saa det Christelige kom over ham, er han geraadet i den besynderlige Situation at vide i en vis Forstand Beskeed derom, men ved Hjælp af en Nomenklatur. I Skyndingen griber han det stærkeste Udtryk for at betegne hvad der er vederfaret ham: og saa har vi hans Aabenbaring. Men saa meget er tillige vist, at han hvor længe han end lever og hvad han end præsterer først skylder at tilbagekalde sit Første, eller træde i dettes Charakteer og anullere det Mellemliggende.

[…]

[Af Capitel I — »Den historiske Situation«: det dialektiske Forhold, den sande Extraordinaire]

Han var da Præst – og nu først indtraadte den Begivenhed, der maatte stille ham som den særlige Enkelte extra ordinem. Collisionen er der; thi at Adler var saaledes situeret er let at see og kun for yderligere Tydeligheds Skyld skal jeg ganske kortelig angive de dialektiske Forhold mellem det a)Almene og b)den Enkelte og c)den særlige Enkelte det vil sige den Extraordinaire. Naar den Enkelte blot reproducerer det Bestaaende i sit Liv (naturligviis forskjelligt i Forhold til hvilke Kræfter og Evner, hvilken Dygtighed han har), saa forholder han sig som den normale Enkelte den ordentlige Enkelte til det Bestaaende; han udfolder det Bestaaendes Liv i sin Existeren; det Bestaaende er for ham det til Grund Liggende, der dannende gjennemtrænger og udvikler hans Evner i Lighed med sig; han forholder sig som en Enkelt, hvis Liv bøier sig efter det Bestaaendes Paradigma. Lad os imidlertid ikke glemme (thi Misfornøiede og Ildesindede ville gjerne udsprede falske Rygter): derfor er hans Liv ingenlunde aandløst. Han er ikke een med i Ramsen af Gloserne som gaae efter Paradigmet, nei han er fri og væsentligen selvstændig; og at være en saadan ordentlig Enkelt er i Ordenen i Reglen den høieste, men ogsaa den qvalitative betydningsfulde, Opgave der er sat noget Menneske. – Saasnart derimod den Enkelte lader sin Reflexion gribe saa dybt, at han vil reflektere over det Bestaaendes Grundforudsætning: saa er han ifærd med at intendere at ville være en særlig Enkelt, og saa længe han reflekterer saaledes er han vragende det Bestaaendes impressa vestigia, er extra ordinem paa eget An- og Tilsvar. Og naar den Enkelte vedbliver ad denne Vei og gaaer saa vidt, at han ikke som den ordentlige Enkelte reproduktivt fornyer det Bestaaendes Liv

i sig ved under evigt Ansvar at ville indordne sig deri, men vil fornye det Bestaaendes Liv ved at bringe et nyt Udgangspunkt for dette, et nyt Udgangspunkt i Forhold til det Bestaaendes Grund-Forudsætning, naar han ved umiddelbart at indordne sig under Gud maa forholde sig omskabende til det Bestaaende: saa er han den Extraordinaire, det vil sige saa bliver denne Plads ham at anvise, hvad enten han nu er berettiget eller ikke, han maa her seire og her finde sin Dom, det Almene skal netop lukke ham ude.

Som overalt saa gjelder det her fornemlig om at den qvalitative Dialektik respekteres med ethisk Alvor. I en charakteerløs Tid kan nemlig det Sophistiske fremkomme, at En, der intenderer at være en Extraordinair vil at denne Intention skal endog komme ham til Gode i det Almenes Tjeneste, saa han paa Grund der af endog bliver en ualmindelig blandt de Ordinaire. Usalige Forvirring, som har sin Grund i en tankeløs letfærdig Qvantiteren. Enten skal et Menneske ville tjene det Almene, det Bestaaende, udtrykke dette, og i saa Fald bliver hans Fortjeneste i Forhold til som han med Troskab og Nøiagtighed veed at indordne sig i det, veed at gjøre sit Liv til en skjøn og riig og tro Gjengivelse af det Bestaaende, udarbeide sig til Type for det Bestaaende; eller han skal for Alvor være en Extraordinair, og saa skal han extra ordinem, ud af Geleddet, ud af Rækkerne, hvor han ikke hører hjemme. Men i vor Tid forvirres Alt. En misfornøiet Embedsmand for Exempel vil endogsaa være noget Overordentligt, fordi han er Embedsmand – og tillige misfornøiet. Sørgelige, usædelige Confusion! Er han misfornøiet, fordi han har noget Nyt fra Gud at bringe os: saa ud af Geleddet, »Strikken om Halsen« og lad ham saa tale, saa er Situationen som en sand Extraordinair behøver den og maa fordre den, for at kunne gestikulere og faae Sangværket i Gang. Men har han ikke noget Nyt fra Gud at bringe os, saa skal det ingenlunde anslaaes ham til Fordeel, at han er misfornøiet – og tillige Embedsmand. Men Tidens Charakteerløshed og nyfigne Feighed gjør det tilsidst til en Slags vanærende Bornerethed, at være Noget tilgavns: enten en tro Embedsmand med Liv og Sjel, eller en Reformator med Sværdet over Hovedet, i Livsfare, i Selvfornægtelse. Εις ϰοιρανος

εστω, og saaledes ogsaa Een være den Extraordinaire. Vil en heel Slægt være Konge, og en heel Slægt fuske i at være Extraordinair, saa bliver det Jux. Og Følgen deraf bliver kun Sinkerie. Dersom Styrelsen havde tiltænkt Slægten en Extraordinair: saa vil Følgen blive, at der maa ventes, at der maaskee først maa sendes en Forløber, en Betjent, der har at gjøre Luft, har at cassere alle disse falske Propheter, har at bringe lidt Mening og Fynd ind igjen i de udmarvede og meningsløse Forhold. Naar nemlig hele Generationen er bleven reformatorisk, saa kan den sande Reformator slet ikke komme til at vise sig i sin Sandhed; det er, for at erindre om et tidligere Udtryk, ligesom naar ved en Ildebrand Alle commandere: saa kan Brandmajoren ikke commandere.

Det der saaledes gjør Forskjellen mellem den ordentlige Enkelte og den særlige Enkelte er Udgangspunktet; forresten kan det godt være, at en ordentlig Enkelt er, menneskeligt talt, større end en virkelig Extraordinair. Den sidste Maalestok, efter hvilken Menneskene rangere, er det Ethiske, i Forhold til hvilket Differentserne (selv den i særlig Forstand at være kaldet af Gud) er et Forsvindende; men det verdslige Sind lader omvendt Rangfølgen bestemmes efter Differentsen. Lad os tage et Exempel paa en saadan ordentlig Enkelt, og lad os ret glæde os ved, at vi har Exempler at pege paa; lad os nævne honoris causa, men ogsaa for at oplyse dette Forhold: Stats-Kirkens øverste Geistlige, visseligen ogsaa dens troeste TjenerSjellands beundrede Biskop

– og Enhver gjør vel i her at beundre; thi Den der her udtrykker ogsaa det Almene skal man netop gjøre sig en Glæde af at beundre, da man af ham kan lære. Biskop Mynster har ikke det Mindste af Det, som man i strængeste Forstand maatte kalde den særlige Enkeltes Signalement. Tværtimod, med ophøiet Rolighed, glad hvilende i sin Overbeviisning som rigeligt Indhold for et riigholdigt Liv, med formanende Eftertryk, med Alvorens ædruelige Besindighed, indtil Grændsen af en ædel lille ironisk Vending mod fortumlede Hoveder har denne Mand bestandigt vedkjendt sig, at det ikke var noget Nyt, han havde at bringe, at det tvertimod var det gamle Bekjendte; han har aldrig rokket ved det Bestaaendes Piller, tvertimod har han selv staaet urokket som en Grund-Pille. Og naar han reviderer første Oplag af sine tidligste Prædikener, saa »finder han Intet at forandre i det Væsentlige« (som havde han siden den Tid maaskee været saa heldig at overkomme een eller anden frisk arriveret systematisk Nyhed); og naar han engang paa sit Dødsleie, ikke for det ny Oplags Skyld, men til Vitterlighed reviderer alle Prædikenerne, saa finder han formodentligen igjen »Intet at forandre i det Væsentlige«. Nei, det var Altsammen det gamle Bekjendte – som dog i ham fandt et saa friskt og saa forfriskende Udspring, et saa ædelt, saa skjønt og saa riigt Udtryk, at han i et langt Liv bevæger Mange forunderligt nok ved det Gamle og at efter hans Død Mange, bevægede, ville længes efter denne Gamle og dette Gamle som man længes efter Ungdommens Liflighed, som man i Sommerens Varme kan længes efter Kildens Kølighed – forunderligt nok, at det skal være noget Gammelt og Bekjendt! Ja, sandeligen, dersom en Lære engang i sin første Begynden maa ønske en Apostel, der i strængeste Forstand som Extraordinair staaer udenfor Geleddet: o, da vil den samme Lære i en sildigere Tid ønske saadanne Husholdere, der intet Nyt have at bringe, der tvertimod kun i Alvor have deres Glæde af selv at udtrykke det Almene, og deres Glæde af at gaae med i Geleddet og lære os Andre at holde Takten – naar vi blot passe paa at see flinkt op til Høire. – En Viis har sagt om Piger: at de skulle giftes, naar de ere Piger i Alder men Koner i Forstand. En Mand har ikke et saadant Øieblik, om ham maa der siges at Opgaven er at bevare sig selv gjennem et heelt Liv. Men hvad er det at bevare sig selv? Det er, naar Blodet er varmt og Hjertet banker heftigt i Ungdommens Dage, da at kunne køle næsten med en Oldings Besindighed; og det er, naar Dagen hælder, naar det stunder mod Afskeden, da som Olding at kunne opflamme med Ungdommens Friskhed. Naar skal en Pige giftes? Oldtiden svarer: naar hun er Pige i Alder, men Kone i Forstand.« Naar skal en Mand være Lærer? Naar han er en Mand i Alder men i Besindighed en Olding. Og naar staaer han paa sit Høieste? naar han engang er en Olding i Alder og Besindighed, men hjertefrisk som en Yngling. Thi hvad er vel det at bevare sig selv, hvilket væsentligen forstaaet er Mandens høieste Opgave? Det er, naar Blodet er varmt og Hjertet banker heftigt i Ungdommens Dage, da at kunne svale næsten med en Oldings-Besindighed; og det er naar Dagen hælder, naar det stunder mod Afskeden da at kunne opflamme med Ungdommens

Ild.

– Men er dette ikke en Fornærmelse mod den Høiærværdige at sidde og skrive noget Saadant; dersom det Fremsatte er sandt, saa er Biskop Mynster jo ingen stor Mand, saa har han jo aldrig fulgt med Tiden, saa veed han jo ikke, hvad Tidens Fordring er, end mindre at han selv har været istand til at inventere den. Nei, inventeret Noget har han ikke; om han maaskee ikke har kunnet (han der dog nok som skarpsindig Psycholog kjender de menneskelige Daarskaber fra Grunden, og altsaa er i Besiddelse af Nøglen til det store Pakloft, hvor de diverse Tidens Fordringer ligger opdyngede), skal jeg ikke fordriste mig til at afgjøre; men vist er det, han har ikke inventeret Noget.

Det ny Udgangspunkt er Differentsen mellem den sande Ordinaire og den sande Extraordinaire; den væsentlige menneskelige Maalestok: det Ethiske have de begge tilfælles. Naar da den Enkelte virkelig er den sande Extraordinaire, og virkelig har et nyt Udgangspunkt, naar han fatter sit Livs betrængte Vanskelighed i hiin discrimen mellem det Almene og Enkeltheden extra ordinem: maa han være ubetinget kjendelig paa, at han er villig til at bringe Offer. Og dette maa han være villig til for sin egen Skyld og for det Almenes Skyld.

Netop fordi den Extraordinaire, hvis han er det i Sandhed, maa ved sit Guds-Forhold være sig bevidst, at han ϰατα

δυναμιν er stærkere end summa summarum af alt det Bestaaende, har han slet Intet med den Bekymring at gjøre: om han nu bare seirer. Nei denne Bekymring er han ganske fritagen for; men derimod har han den særlige Enkeltheds rædsomme Ansvar for hver Skridt han gjør, om han nu indtil det Mindste nøie følger sin Ordre, om han bestemt og ene og lyt hører Guds Stemme – det rædsomme Ansvar, hvis han hørte eller havde hørt feil. Netop af den Grund maa han ønske sig al mulig Modstand udenfra, ønske, at det Bestaaende har Kræfter for at kunne gjøre ham hans Liv til et tentamen rigorosum, thi denne Prøve og dens Smerte er dog Intet mod Ansvarets Rædsel: hvis han var eller havde været i en Vildfarelse! Dersom for Exempel en Søn skulde føle sig kaldet til at bringe en ny Anskuelse af det huslige Liv (og som en Søn er bundet i Pieteten, saa skal og bør enhver Enkelt være bunden i Pietet ligeoverfor det Almene): vilde han da ikke, hvis der var Sandhed i ham, netop ønske, at Faderen var den Stærke, som kunde træde ham imod med faderlig Myndigheds fulde Kraft? Sønnen vilde nemlig ikke saa meget frygte det at ligge under, hvis han havde Uret, altsaa at maatte ydmyget men frelst vende om til det Gamle, som han vilde gyse for den Rædsel: at seire – hvis han i Grunden havde Uret.

Saaledes med den sande Extraordinaire; han er det sorgløseste Menneske betræffende hiin de verdslige Heltes timelige Bekymring, om nu Det han har at forkynde skal seire i Verden; derimod er han ængstet som en arm Synder, sønderknuset hver Gang han betænker sit Ansvar, og om han nu paa nogen Maade skulde have taget feil; ja det er ham, som skulde Aandedrættet standse, saa tungt kan Ansvarets Vægt hvile paa ham. Netop derfor ønsker han Modstand: han – den Svage, han – den Stærke, der, skjøndt eenligt Menneske, ϰατα

δυναμιν i sin Svaghed er stærkere end hele det Bestaaendes forenede Magt, hvilken naturligviis har Magt til baade at hudflette og henrette ham som Ingenting. Naar Berserkegangen kom over vore nordiske Fædre, da lode de sig tvinge ind mellem Skjolde: saaledes ønsker ogsaa den sande Extraordinaire det Bestaaende Kraft til at danne behørig Modstand.

\workheading{Bladartikler og Øieblikket (Kirkekampen)}{1854–55}

\worknote{Antologiens ''Fædrelandet Articles and The Moment'' — skrifterne fra angrebet på den bestaaende Kirke i Kierkegaards sidste leveaar. Her gengives de ti stykker Hong-udvalget bringer, i rækkefølge: fra åbningsartiklen om Biskop Mynster til ''Min Opgave''.}

\sectionheading{Var Biskop Mynster et »Sandhedsvidne«? (Fædrelandet, dec. 1854)  ·  SV¹ XIV}

Var Biskop Mynster et »Sandhedsvidne«, et af »de rette Sandhedsvidner« – er dette Sandhed?

Iden Tale, Professor Martensen har »holdt 5te Søndag efter hellig tre Konger, Søndagen før Biskop Doctor Mynsters Jordefærd«, en Erindringstale, hvad den maaskee ogsaa forsaavidt kan kaldes, som den bringer ProfessorMartensen i Erindring til den ledige Bispestol – i denne Tale stilles Biskop Mynster frem som et Sandhedsvidne, som et af de rette Sandhedsvidner; der tales i saa stærke og afgjørende Udtryk som vel muligt. Med den afdøde Biskops Skikkelse, hans Liv og Levnet, Udgangen af hans Vandel for Øie formanes vi til at »efterligne de sande Veilederes, de rette Sandhedsvidners Tro« (pagina 5), deres Tro, thi den var, hvad der udtrykkeligt siges om Biskop Mynster »ikke blot i Ord og Bekjendelse, men i Gjerning og Sandhed« (pagina 9); den afdøde Biskop indføres af Professor Martensen i (pagina 6) »den hellige Kjæde af Sandhedsvidner, som strækker sig gjennem Tiderne fra Apostlernes Dage« osv.

Dette maa jeg gjøre Indsigelse mod – og nu da Biskop Mynster er død, kan jeg ville tale, her dog meget kort, og da slet ikke om, hvad der bestemte mig til at indtage det Forhold til ham, som jeg indtog.

Naar man ved »Forkyndelse« nærmest tænker paa hvad der siges, skrives, trykkes, Ordet, Prædikenen – at da (for blot at tage Eet) Biskop Mynsters Christendoms-Forkyndelse fortoner, tilslører, fortier, udelader noget af det mest afgjørende Christelige, det, der falder os Mennesker ubeleiligere, det, der vilde gjøre vort Liv anstrænget, forhindre os i at nyde Livet, det om at afdøe, om frivillig Forsagelse, om at hade sig selv, om at lide for Læren og saa videre: behøver man ikke at være synderlig skarpsynet for at kunne see, naar man lægger det nye Testamente ved Siden af Mynsters Prædiken.

Forsaavidt derimod ved »Forkyndelse« nærmest tænkes paa, hvorvidt nu Forkynderens Liv udtrykker det, han siger (og dette er vel at mærke christeligt det Afgjørende, hvorved Christendommen just har villet sikkre sig mod at faae charakterløse Docenter i Stedet for Vidner) – at da Biskop Mynsters Christendoms-Forkyndelse ikke var i Charakter, at han udenfor de stille Timer ikke var i Charakter, ikke engang af den hans Prædiken, der dog, som sagt, sammenholdt med det nye Testamente har slaaet betydeligt af paa det Christelige: dette behøver man atter ikke at være synderlig skarpsynet for at kunne see, hvis man ellers ved at høre og læse ham har behørigt Kjendskab til hans Prædiken; i 1848 og efter den Tid blev det til at see selv for blinde Beundrere, hvis de have behørigt Kjendskab til hans Prædiken, for at kunne vide, hvad denne, hvad de stille Timer foranledige En til at vente.

Saaledes var, naar man lægger det nye Testamente ved Siden, Biskop Mynsters Christendoms-Forkyndelse en, især for et Sandhedsvidne, mislig Christendoms Forkyndelse. Men da var der, mente jeg, det Sande hos ham, at han, hvad jeg holder mig fuldt forvisset om, var villig til at tilstaae for Gud og sig selv, at han ingenlunde, ingenlunde var et Sandhedsvidne – just denne Tilstaaelse var i mine Tanker det Sande.

Men skal Biskop Mynster fra Prædikestolen fremstilles, canoniseres som Sandhedsvidne, et af de rette Sandhedsvidner: saa maa der gjøres Indsigelse. Den Berlingske Tidende (den officielle Tidende, ligesom vel Professor Martensen er den officielle Prædikant) er, som jeg seer, af den Mening, at Professor Martensen (der med paafaldende Hasten kommer Begravelsen og saa ogsaa Monumentet i Forkjøbet) ved denne Tale har fra Prædikestolen reist den Afdøde et skjønt og værdigt Monument; jeg vilde foretrække at sige: et Professor Martensen selv værdigt Monument. Men i ethvert Tilfælde Monumenter kan man ikke ignorere, derfor maa der gjøres Indsigelse, hvad saa maaskee endog kunde bidrage til at gjøre Monumentet (for Professor Martensen) varigere.

Biskop Mynster et Sandhedsvidne! Du som læser dette, Du veed vel, hvad der christeligt forstaaes ved et Sandhedsvidne; men lad mig dog erindre Dig derom, at dertil fordres ubetinget det at lide for Læren. Og naar der skærpende siges: et af »de rette« Sandhedsvidner, saa maa altsaa Ordet tages i strengeste Forstand. Lad mig da med nogle faa Træk forsøge at antyde, hvad derved maa forstaaes, for at gjøre Dig det nærværende.

Et Sandhedsvidne, det er en Mand, hvis Liv fra først til sidst er ukjendt med Alt, hvad der hedder Nydelse – o, og hvad enten Meget eller Lidet forundtes Dig, Du veed, hvor godt det gjør, det man kalder Nydelse; men hans Liv var fra først til sidst ukjendt med Alt hvad der hedder Nydelse, derimod fra først til sidst indviet i Alt, hvad der hedder Lidelse – ak, og blev Du end fritagen for de langvarige, de qvalfuldere Lidelser, Du veed det dog af egen Erfaring, hvorledes et Menneske ømmer sig i det, man kalder Lidelse! Men deri var hans Liv fra først til sidst indviet, i hvad der vel endog sjeldnere tales om blandt Mennesker, fordi det sjeldnere forekommer, i Inderlighedens Kampe, i Frygt og Bæven, i Skælven, i Anfægtelser, i Sjels-Angester, i Aands-Qvaler; og saa dertil forsøgt i alle de Lidelser, som der almindeligere tales om i Verden. Et Sandhedsvidne, det er en Mand, der i Armod vidner for Sandhed, i Armod, i Ringhed og Fornedrelse, saa miskjendt, forhadt, afskyet, saa bespottet, forhaanet, udgriint – det daglige Brød har han maaskee ikke altid havt, saa fattig var han, men Forfølgelsens daglige Brød fik han rigeligt hver Dag; for ham var der aldrig Avancement og Befordring uden omvendt, ned efter Skridt for Skridt. Et Sandhedsvidne, et af de rette Sandhedsvidner, det er en Mand, der bliver hudflettet, mishandlet, slæbt fra det ene Fængsel i det andet, og saa tilsidst – sidste Avancement, hvorved han optages i første Classe af den christelige Rangforordning, blandt de rette Sandhedsvidner – saa tilsidst – thi det er jo et af de rette Sandhedsvidner, Professor Martensen taler om – saa tilsidst korsfæstet eller halshugget eller brændt eller stegt paa en Rist, hans afsjelede Legeme af Rakkeren henslængt paa et afsides Sted ubegravet – saaledes begraves et Sandhedsvidne! – eller brændt til Aske og henkastet for alle Vinde, at ethvert Spor af dette »Udskud«, som Apostelen siger han er bleven, maatte udslettes.

Dette er et Sandhedsvidne, hans Liv og Levnet, hans Død og Begravelse – og Biskop Mynster, siger Professor Martensen, var et af de rette Sandhedsvidner.

Er det Sandhed? Er det at tale saaledes, er det maaskee ogsaa at vidne for Sandhed, og Professor Martensen ved denne Tale selv traadt i Charakter af Sandhedsvidne, et af de rette Sandhedsvidner? Sandeligen, der er det, som er Christendommen og Christendommens Væsen mere imod end ethvert Kjetteri, ethvert Schisma, mere imod end alle Kjetterier og Schismaer tilsammen, og det er det: at lege Christendom. Men det er – aldeles, aldeles i samme Forstand, som Barnet leger Soldat – at lege Christendom: at tage Farerne bort (christeligt svarer »Vidne« og »Fare« til hinanden), istedetfor disse at anbringe Magten (til at være Fare for Andre), Goderne, Fordelene, en yppig Nydelse selv af de fineste Raffinements – og saa at lege den Leg, at Biskop Mynster var Sandhedsvidne, et af de rette Sandhedsvidner, lege den saa frygtelig alvorligt, at man da slet ikke kan standse Legen, men leger den fort ind i Himlen, leger Biskop Mynster med ind i den hellige Kjæde af Sandhedsvidner, som fra Apostlenes Dage strækker sig til vore Tider.

Efterskrift

Efterskrift

Denne Artikel har, som man vil see af dens Dato, ligget nogen Tid.

Saa længe der var Spørgsmaal om Besættelsen af Sjællands Bispestol, mente jeg at burde lade Professor Martensen offenligen uomtalt, da han jo, hvad enten han nu blev Biskop eller ikke, i ethvert Tilfælde var en Candidat til denne Plads, og formodenligt ønskede, at der, saalænge det stod paa, saavidt muligt Intet passerede, ham betræffende.

Med Professor Martensens Udnævnelse til Biskop faldt dette Hensyn bort. Men da nu engang Artiklen ikke kunde komme og derfor ikke var kommen strax, tænkte jeg som saa: der er jo ingen Grund til at haste. Desuden fremkaldte Biskop Martensens Udnævnelse Angreb paa ham fra andre Sider og af en ganske anden Art: dette Angreb vilde jeg høist unødigt træffe sammen med; jeg ventede da, jeg mente: der er, som sagt, slet ingen Grund til at haste og der tabes slet Intet ved at vente. Muligt vil endog Nogen finde, at der er vundet Noget, finde en dybere Betydning i, at Indsigelsen kommer med en saadan Langsomhed.

I Efteraaret 1854.

Men Indsigelse maa der gjøres mod dette, at Biskop Mynster skulde have været et Sandhedsvidne.

Biskop Mynster kan tilnærmelsesvis siges at have baaret en hel Slægt – det er derfor en Vanskelighed, som grændser til det Umulige, at bringe Klarhed ind i vore forvirrede religieuse Forhold og Begreber, saalænge der ikke falder en sandere Belysning over Sandheden af Biskop Mynsters Christendoms Forkyndelse, hvad i Grunden ogsaa skyldes mig; thi Biskop Mynster, just Biskop Mynster var, om man saa vil, mit Livs Ulykke, ikke endda det, at han ikke var et Sandhedsvidne (den Sag havde ikke været saa farlig) men det, at han foruden al anden Fordel, som han, efter største Maalestok, drog af at forkynde Christendom, ogsaa medtog den Nydelse, ved at declamere i stille Timer om Søndagen og ved verdslig-klogt at dække sig om Mandagen, at foranledige det Skin, at han var et Charaktermenneske, en Mand af Grundsætninger, en Mand, der staaer fast, naar Alt vakler, en Mand, der ikke svigter, naar Alle svigte osv. osv., medens Sandheden var, at han var i høi Grad verdslig klog, men svag, nydelsessyg og kun stor som Declamator – og mit Livs Ulykke, om man saa vil (hvad dog i en meget høi Forstand ved Styrelsens Kjærlighed blev til Gode for mig, blev min Lykke) min Ulykke var, at jeg, af en afdød Fader opdragen ved »Mynsters Prædiken«, ogsaa, af Pietet mod den afdøde Fader, honorerede denne falske Vexel i Stedet for at protestere den.

Han er nu død – Gud være lovet, at det kunde holdes hen, saalænge han levede! Det blev naaet, hvad jeg rigtignok mod Slutningen var nær ved at fortvivle om, men det blev dog naaet, hvad der var min Tanke, mit Ønske, hvad jeg ogsaa kan huske engang for mange Aar siden at have sagt til gamle Grundtvig: Biskop Mynster skal først leve ud, begraves med fuld Musik – det blev naaet, han blev jo, om jeg saa tør sige, begravet med fuld Musik.Til Monumentet for ham er der vel ogsaa nu omtrent indkommet hvad der vil indkomme.

Saa kan der ikke længer ties, Indsigelsen maa komme, ved sin Langsomhed kun desto alvorligere, Indsigelsen mod fra Prædikestolen, altsaa for Gud, at fremstille Biskop Mynster som Sandhedsvidne; thi usandt er det, men saaledes forkyndt bliver det en himmelraabende Usandhed.

I December 1854.

Dog maaskee er dette ved Biskop Mynsters mangeaarige Christendoms-Forkyndelse bragt i Glemsel. Thi ogsaa dette er, og en capital Mislighed ved hans Forkyndelse, ikke det, at han selv var Embedsmand, – christeligt trækker dette fra – Forkyndelsen, hans egen glimrende, paa Nydelser rige, Carriere, nei, ikke det, men det, at han vilde autorisere den Art Forkyndelse til den sande christelige, og derved, i Fortielse, gjøre den sande christelige (den af lidende Sandhedsvidner) til en Overdrivelse, i Stedet for omvendt, at gjøre Christendommen den Tilstaaelse, at den Forkyndelse, han repræsenterede, er Noget, som maa ved Aflad og Indulgens indrømmes os almindelige Mennesker, Noget, vi almindelige Mennesker hjælpe os med, fordi vi ere for selviske, for verdslige, for sandselige til at due til mere, Noget, som vi almindelige Mennesker hjælpe os med, og som saaledes forstaaet, saa ingenlunde – alle falske Reformatorer til Trods! – skal indbildsk og opblæst vrages, tvertimod respecteres.

\sectionheading{En Thesis — kun een eneste (Fædrelandet, jan. 1855)  ·  SV¹ XIV}

En Thesis – kun een eneste.

O, Luther, Du havde 95 Theses: forfærdeligt! Og dog i dybere Forstand jo flere Theses, jo mindre forfærdeligt. Sagen er langt forfærdeligere: der er kun een Thesis.

Det nye Testamentes Christendom er slet ikke til. Her er ikke Noget at reformere; det, det gjælder om, er at bringe Lys i en gjennem Aarhundreder fortsat, af Millioner (uskyldigere eller skyldigt) øvet christelig Criminal-Forbrydelse, hvorved man kløgtigt, under Navn af Christendommens Fuldkommengjørelse, har lidt efter lidt søgt at franarre Gud Christendommen, faaet Christendom til at være lige det Modsatte af, hvad den er i det nye Testamente.

For at den her i Landet almindelige, den officielle, Christendom skal kunne siges endog blot sandt at forholde sig til det nye Testamentes Christendom, maa den først, saa redeligt, uforbeholdent, høitideligt som muligt gjøre det kjendeligt, paa hvilken Afstand den er fra det nye Testamentes Christendom og uden med Sandhed at kunne kaldes en Stræben i Retning af at komme nærmere til det nye Testamentes Christendom.

Saa længe dette ikke skeer, saa længe man enten lader som Ingenting, som var Alt i sin Rigtighed og Det vi kalde Christendom det nye Testamentes Christendom, eller man gjør Konster for at skjule over Forskjellen, Konster for at opretholde Skinnet af at være det nye Testamentes Christendom: saa længe fortsættes den christelige Criminal-Forbrydelse; der kan ikke være Tale om at reformere, men om at skaffe Lys i denne christelige Criminal-Sag.

Og for at sige et Ord om mig selv: jeg er ikke, hvad Tiden maaskee kræver, en Reformator, paa ingen Maade, ei heller en spekulativ dybsindig Geist, en Seer, Prophet, nei, jeg er – med Tilladelse – jeg er et i sjelden Grad afgjort Politi-Talent. Forunderlige Samtræf, at jeg just skal blive samtidig med den Periode i Kirkens Historie, der, modernt, er »Sandhedsvidnernes« Periode, hvor Alle ere »hellige Sandhedsvidner«.

\sectionheading{Hvad jeg vil? (Fædrelandet, marts 1855)  ·  SV¹ XIV}

Hvad jeg vil?

Ganske simpelt: jeg vil Redelighed. Jeg er ikke, som man velmenende – thi Forbittrelsens og Raseriets og Afmagtens og Vrøvlets Opfattelse af mig kan jeg ikke tage Hensyn til – har villet fremstille mig, jeg er ikke en christelig Strenghed lige over for en given christelig Mildhed.

Paa ingen Maade, jeg er hverken Mildhed eller Strenghed – jeg er: en menneskelig Redelighed.

Den Formildelse, som er den almindelige Christendom her i Landet, vil jeg have holdt ved Siden af det nye Testamente for at faae at see, hvorledes disse to forholde sig til hinanden.

Viser det sig saa, kan jeg eller nogen Anden vise, at den lader sig holde lige over for det nye Testamentes Christendom: saa vil jeg med største Glæde gaae ind paa den.

Men Eet vil jeg ikke, for ingen, ingen Pris: jeg vil ikke ved Fortielser eller ved at gjøre Konster see at frembringe det Skin, at den almindelige Christendom i Landet og det nye Testamentes Christendom ligne hinanden.

See, det er dette jeg ikke vil; og hvorfor ikke? Nu, fordi jeg vil Redelighed, eller vil Du jeg skal tale anderledes, vel, fordi jeg troer, at er det muligt, kan end den alleryderligste Formildelse af det nye Testamentes Christendom holde Stik i Evighedens Dom: den kan umuligt holde Stik, naar man endog har anvendt Konster paa at skjule over Forskjellen mellem det nye Testamentes Christendom og denne Formildelse. Jeg mener som saa: er En den Naadige, nu vel saa lad mig vove at forlange af ham, at han skal tilgive mig al min Gjeld: men var end hans Naade guddommelig Naade, dette er for meget forlangt, at jeg ikke engang vil være sand, betræffende hvor stor Gjælden er.

Og det er denne Usandhed, jeg mener den officielle Christendom forskylder: den gjør ikke hensynsløst den christelige Fordring kjendelig, maaskee fordi den er bange for, at man vilde komme til at gyse ved at see, paa hvilken Afstand vi leve, og uden at vort Liv i fjerneste Maade kan kaldes en Stræben i Retning af at opfylde Fordringen. Eller for blot at tage et Exempel af hvad der forresten jo overalt er tilstede i det nye Testamentes Christendom, naar Christendommen fordrer til at frelse sit Liv evigt (og det er jo det, vi mene at opnaae som Christne): at hade sit eget Liv i denne Verden; er der nogen Eneste iblandt os, hvis Liv end i fjerneste Maade kunde kaldes end den svageste Stræben i denne Retning, medens der i Landet maaskee ere tusindvis »Christne«, som ikke engang ere vidende om denne Fordring. Altsaa, vi »Christne«, vi leve i ganske almindelig menneskelig Forstand elskende vort Liv. Hvis da Gud alligevel skal af »Naade« antage os for Christne: Eet maa dog fordres, at vi ved at være nøiagtigt vidende om Fordringen have en sand Forestilling om, hvor uendelig stor Naaden er, som vises os. Saa langt kan »Naaden« umuligt række, Eet maa den aldrig, den maa aldrig bruges til at fortie eller formindske Fordringen; i saa Fald endevender »Naaden« hele Christendommen. – Eller for at tage et Exempel af en anden Art. En Lærer i Christendom lønnes for Exempel med flere Tusinder. Naar vi nu fortie den christelige Maalestok og gaae ud fra det almindelige Menneskelige, at det jo er ganske naturligt, at en Mand skal have Løn for sit Arbeide, Løn, saa han kan leve med Familie, og som Embedsmand i anseet Stilling anseelig Løn: saa er flere Tusinder om Aaret slet ikke meget. Saasnart derimod den christelige Fordring: Armod gjøres gjældende, saa er Familie Luxus, flere Tusinder en saare høi Gage. Jeg siger ikke dette for at fradrage, hvis jeg formaaede det, en saadan Embedsmand en eneste Skilling, tvertimod, hvis han ønskede det, og jeg formaaede det, skulde han gjerne faae dobbelt saa mange Tusinder: men jeg siger, at Fortielsen af den christelige Fordring forandrer Synspunktet for hele hans Gage. Til Redelighed mod Christendommen vil fordres, at man selv bringer det i Erindring, at christeligt er Fordringen Armod, og at dette ikke er noget capricieust Indfald af Christendommen, men er, fordi Christendommen meget godt forstaaer, at kun i Armod kan den betjenes sandt, og at jo flere Tusinder en Christendoms Lærer har i Gage, jo mindre kan han tjene Christendommen. Derimod er det ikke redeligt at fortie Fordringen eller at gjøre Konster for at give det Udseende af, at den Art Levevei og Carriere er ganske det nye Testamentes Christendom. Nei, lad os tage Pengene; men for Gud i Himlens Skyld ikke det Næste, ikke ville skjule over den christelige Fordring, saa der, ved Fortielse eller ved Falsk, tilveiebringes et Slags Decorum, der er i aller-allerhøieste Grad demoraliserende og er at snigmyrde Christendommen.

Altsaa Redelighed vil jeg; men hidtil har det Bestaaende ikke af egen Drift villet gaae ind paa den Art Redelighed, og heller ikke villet lade sig paavirke af mig. Dog bliver jeg derfor hverken en Mildhed eller en Strenghed, nei jeg er og bliver: ganske simpelt en menneskelig Redelighed.

Lad mig vove det Yderste, for om muligt at blive forstaaet i, hvad det er jeg vil.

Jeg vil Redelighed. Er det da dette, Slægten eller Samtiden vil, vil den ærligt, redeligt, uforbeholdent, aabent, ligefrem gjøre Oprør mod Christendommen, sige til Gud »vi kan, vi vil ikke bøie os under denne Magt« – men vel at mærke, det gjøres ærligt, redeligt, uforbeholdent, aabent, ligefrem: nu vel, hvor besynderligt det kunde synes, jeg er med; thi Redelighed vil jeg. Og overalt, hvor der er Redelighed, kan jeg gaae med; et redeligt Oprør mod Christendommen kan kun gjøres, naar man redeligt vedgaaer, hvad Christendom er, og hvorledes man selv forholder sig dertil.

Er det dette man vil: ærligt, aabent, oprigtigt som det sømmer sig, naar man taler med sin Gud, hvad derfor Enhver er, der agter sig selv, og ikke saa dybt foragter sig selv, at han vil være uoprigtig lige over for Gud – altsaa vil man ærligt, oprigtigt, uforbeholdent, ganske gjøre Gud Tilstaaelse, betræffende hvorledes det egenligen hænger sammen med os Mennesker, at Slægten gjennem Tidernes Løb har tilladt sig at formilde og formilde Christendommen, indtil vi tilsidst har faaet den til at være lige det Modsatte af hvad den er i det nye Testamente – og at vi nu dog gjerne ville, hvis det lod sig gjøre, at dette maatte være Christendom: vil man det, saa er jeg med.

Men Eet vil jeg ikke, nei, for ingen, ingen, ingen Pris vil jeg det, Eet vil jeg ikke: jeg vil ikke, om det saa blot var med den sidste Fjerdedel af det sidste Led paa min lille Finger deltage i hvad der hedder: officiel Christendom, der ved Fortielse eller Konster gives Skin af at være det nye Testamentes Christendom, hvad jeg paa mine Knæ takker min Gud for at han forbarmende forhindrede mig at komme for langt ind i.

Vil saa den officielle Christendom i Landet af hvad her er sagt tage Anledning til at bruge Magt mod mig; jeg er rede; thi jeg vil Redelighed.

For denne Redelighed vil jeg vove. Derimod siger jeg ikke: at det er for Christendommen jeg vover. Antag det, antag, at jeg blev ganske bogstavelig et Offer: jeg blev dog ikke et Offer for Christendommen, men fordi jeg vilde Redelighed.

Men medens jeg ikke tør sige, at jeg vover for Christendommen, holder jeg mig dog fuldt og saligt forvisset om, at denne min Voven er Gud velbehagelig, har hans Samtykke. Ja, det veed jeg, det har hans Samtykke, at i en Verden af Christne, hvor Millioner og Millioner kalde sig Christne – at der et Menneske udtrykker: jeg tør ikke kalde mig en Christen; men Redelighed vil jeg, og til den Ende vil jeg vove.

\sectionheading{Et Genie — en Christen (Øieblikket Nr. 5)  ·  SV¹ XIV}

Et Genie – en Christen.

At Genie ikke er Noget ethvert Menneske er, er vel Noget ethvert Menneske vil indrømme. Men at en Christen er noget endnu Sjeldnere end et Genie – dette er gavtyveagtigt reent bragt i Glemme.

Forskjellen mellem et Genie og en Christen er, at Geniet er Naturens Overordentlige; det kan intet Menneske gjøre sig selv til. En Christen er Frihedens Overordentlige, eller, nøiagtigere, Frihedens Ordentlige, kun at det findes overordentlig sjeldent, er hvad Enhver af os skulde være. Derfor vil Gud, at Christendommen skal forkyndes ubetinget for Alle, derfor ere Apostlerne ganske simple Mennesker, Forbilledet i en ringe Tjeners Skikkelse, Alt for at betegne, at dette Overordentlige er det Ordentlige, er tilgængeligt for Alle – men dog er en Christen noget endnu Sjeldnere end et Genie.

Man lade sig kun ikke bedaare af den Omstændighed, at det er tilgængeligt for Alle, muligt for Alle, som fulgte deraf, at det var noget let Noget, og der var mange Christne. Nei, det maa være muligt for Alle, ellers var det ikke Frihedens Overordentlige; men en Christen bliver alligevel det end Sjeldnere end et Genie.

Antaget, at det er i sin Rigtighed med disse Batailloner og Millioner × Millioner Christne, fremkommer her en Indvending, som virkelig faaer Betydning: at Christendommens Forhold saa ere aldeles uden Analogie i den øvrige Tilværelse. Ellers see vi overalt de uhyre Proportioner, som ere Tilværelsens: Millioner Planters Mulighed som Blomsterstøv henveires, Millioner Muligheder af levende Væsener spildes og saa videre, paa eet Genie gaaer der vel Mennesker tusinde × tusinde og saa videre, altid den enorme Ødselhed. Kun i Forhold til Christendommen anderledes: i Forhold til hvad der er sjeldnere end et Genie, slaaer det til, Enhver der fødes er en Christen.

Ligeledes faaer ogsaa en anden Indvending stor Betydning, naar dette med Millioner Christne skal være Sandhed. Jorden er kun et lille Punkt i Universet – og dog skulde Christendommen være forbeholdt den, og til en saadan Spotpriis, at saadan Een og Hver der fødes er en Christen.

Anderledes stiller Sagen sig, naar det sees, at det at være Christen er en saadan Idealitet, at istedetfor den Sludder om Christenheden og Christendommens 1800-aarige Historie, og om at Christendommen er perfectibel, vel maa stilles den Sætning: Christendommen er egentligen ikke kommet ind i Verden, det blev ved Forbilledet og høiest Apostlene; men disse forkyndte den dog allerede saa stærkt i Retning af Udbredelse, at her allerede Misligheden begynder. Thi Eet er at virke saaledes for Udbredelse, at man uafbrudt tidligt og sildigt forkynder Læren for Alle, noget Andet er, at være for hurtig til at tillade Folk saadan i 100, i 1000-viis at antage Navn af Christne, udgive sig for Jesu Christi Disciple. Forbilledets Forkyndelse var en noget anden; thi saa ubetinget han forkyndte Læren for Alle, kun levende derfor, ligesaa ubetinget holdt han igjen paa, det at blive Discipel, at faae Lov til at kalde sig saa. Om en Folke-Forsamling havde ladet sig gribe af Christi Tale, derfor havde han vistnok ikke strax tilladt disse Tusinder at kalde sig Christi Disciple. Nei, han holdt stærkere igjen. Han vandt derfor ogsaa i 3½ Aar kun 11 – medens een Apostel paa een Dag vel i een Time vinder 3000 Christi Disciple. Enten er her Disciplen større end Mesteren, eller Sandheden er, at Apostelen er lidt for hurtig til at slaae til, lidt for hurtig i Retning af Udbredelse, saa Misligheden allerede begynder her.

Kun guddommelig Myndighed kunde imponere Menneske-Slægten saaledes, at det blev til ubetinget Alvor med ubetinget at ville det Evige. Kun Gud-Mennesket kan ene dette: at arbeide ubetinget for Udbredelse og at holde ubetinget lige saa stærkt igjen paa, hvad der skal forstaaes ved at være Discipel. Kun Gud-Mennesket vilde kunne udholde, hvis man vil tænke sig det, 1000 Aar igjennem og atter 1000 Aar ubetinget at virke for Lærens Udbredelse ved at forkynde den, selv om han ikke fik en eneste Discipel, hvis han kun kunde faae dem ved at forandre Vilkaaret. Apostelen har dog nogen selvisk Trang til den Lindring at faae Tilhængere, at blive Flere, hvad Gud-Mennesket ikke har, der ikke selvisk trænger til Tilhængere, derfor kun har Evighedens Priis, ingen Markeds-Priis.

Saaledes skete det, da Christus forkyndte Christendom: Slægten blev ubetinget imponeret.

Men naturam furca expelles, den kommer dog igjen. Det Mennesket tenderer til er at faae Forholdet vendt om. Som en Hund, der tvinges til at gaae paa To, i ethvert Øieblik har en Tendents til at komme til at gaae paa Fire igjen, og saa snart den seer sit Snit, strax gjør det, og blot venter paa at see sit Snit: saaledes er hele Christenheden Menneske-Slægtens Stræben efter at komme til igjen at gaae paa Fire, at blive Christendommen qvit, gavtyveagtigt under Navn af at det er Christendom, og med Paastand paa, at dette er Christendommens Fuldkommengjørelse.

Først vendte man »Forbilledets« anden Side ud, Forbilledet blev ikke mere Forbilledet, men Forsoneren; og istedetfor at see paa ham i Retning af Efterfølgelse, dvælede man ved hans Velgjerninger og ønskede at være i deres Sted mod hvem de vistes, hvilket er lige saa bagvendt, som, hvis Een fremstilledes som Forbillede i Gavmildhed, man da ikke vilde see paa ham i Betydning af at efterligne hans Gavmildhed, men i Betydning af at ønske at være i Deres Sted, mod hvem han viste Gavmildhed.

Altsaa Forbilledet gik ud. Saa afskaffede man ogsaa Apostelen som Forbillede. Derpaa ogsaa den første christne Tid som Forbillede. Og saaledes lykkedes det tilsidst at naae – at gaae paa Fire igjen, og at Det, just Det var den sande Christendom. Ved Hjælp af Dogmer sikkrede man sig mod Alt, hvad der blot med nogen Sandhed kunde kaldes christeligt Forbillede, og saa gik det for fulde Seil i Retning af – Perfectibiliteten.

\sectionheading{Kort og Spidst (Øieblikket Nr. 6)  ·  SV¹ XIV}

Kort og Spidst.

Christendommen lader sig fuldkommengjøre (er perfectibel); det gaaer fremad; nu er det naaet det Fuldkomne. Hvad der tilstræbtes som Idealet, men hvad dog selv den første Tid kun tilnærmelsesviis naaede, at de Christne er et Folk af Præster, det er nu fuldkommen naaet, især i Protestantismen, især i Danmark.

Dersom nemlig det vi kalde Præst er det at være Præst – ja saa ere vi Alle Præster!

I den pragtfulde Domkirke fremtræder den høivelbaarne høiærværdige Geheime-General-Ober-Hof-Prædikant, den fornemme Verdens udvalgte Yndling, han træder frem for en udvalgt Kreds af Udvalgte, og prædiker rørt over den af ham selv udvalgte Text »Gud har udvalgt det i Verden Ringe og Foragtede« – og der er Ingen, som leer.

Naar en Mand har Tandpine, siger Verden »stakkels Mand«; naar en Mands Kone bliver ham utro, siger Verden »stakkels Mand«; naar en Mand er i Pengeforlegenhed, siger Verden »stakkels Mand«. – Naar det behager Gud i en ringe Tjeners Skikkelse at ville lide i denne Verden, siger Verden »stakkels Menneske«; naar en Apostel i guddommeligt Ærinde har den Ære at lide for Sandhed, siger Verden »stakkels Menneske«: stakkels Verden!

»Havde Apostelen Paulus nogen Embedsstilling?« Nei, Paulus havde ingen Embedsstilling. »Tjente han da paa anden Maade mange Penge?« Nei, han tjente paa ingen Maade Penge. »Var han da idetmindste gift?« Nei, han var ikke gift. »Men saa er jo Paulus ingen alvorlig Mand!« Nei, Paulus er ingen alvorlig Mand.

Om en svensk Præst fortælles, at han, rystet ved Synet af den Virkning hans Tale frembragte paa Tilhørerne, som hensvømmede i Taarer, der fortælles at han beroligende sagde: græder ikke Børn, det turde være Løgn Altsammen.

Hvorfor siger Præsten det nu ikke mere? Behøves ikke, vi veed det – vi ere Alle Præster.

Men derfor kan vi jo gjerne græde, baade hans og vore Taarer være paa ingen Maade hykkelske, men meente, sande – som i Theatret.

Da Hedenskabet opløstes, levede der nogle Præster, Augurer kaldede. Om disse berettes, at den ene Augur ikke kunde see paa den anden uden at smile.

I »Christenhed« kan vel snart Ingen see en Præst, eller vel snart det ene Menneske ikke see paa det andet uden at smile – men vi ere jo ogsaa Alle Præster.

Er dette den samme Lære, naar Christus siger til den rige Yngling: sælg Alt, hvad Du haver – og giv Fattige det;

og naar Præsten siger: sælg Alt hvad Du haver og – giv mig det?

Genier ere som Tordenveir: de gaae mod Vinden, forfærde Menneskene, rense Luften.

Det Bestaaende har opfundet adskillige Tordenafledere.

Og det lykkes. Ja vist lykkes det; det lykkes at gjøre det næste Tordenveir desto alvorligere.

Man kan ikke leve af Ingenting. Det hører man saa ofte, især af Præster.

Og just Præsterne gjøre dette Konststykke: Christendommen er slet ikke til – dog leve de deraf.

\sectionheading{Frygt meest af Alt at være i en Vildfarelse! (Øieblikket Nr. 6)  ·  SV¹ XIV}

Frygt meest af Alt at være i en Vildfarelse!

Dette er som bekjendt Sokrates's Sætning; han frygtede af Alt meest at være i en Vildfarelse.

Christendommen, der visseligen i een Forstand ikke lærer Mennesket at frygte, ikke engang Dem, som kunne slaae En ihjel, lærer dog i en anden Forstand en endnu større Frygt end hiin sokratiske, lærer at frygte Den, som kan fordærve baade Sjel og Legeme i Helvede.

Dog først det Første, at blive opmærksom paa det nye Testamentes Christendom; og hertil vil da hjælpe Dig hiin sokratiske Frygt, at frygte meest af Alt at være i en Vildfarelse.

Har Du ikke denne Frygt, eller (for ikke at slaae saa høi en Tone an) er det ikke saaledes med Dig, er det ikke det Du vil, vil Du ikke stræbe at vinde Mod til at »frygte meest af Alt at være i en Indbildning«: saa indlad Dig aldrig med mig. Nei, bliv saa hos Præsterne, lad dem jo før jo hellere overbevise Dig ganske om, at hvad jeg siger er en Art Galskab (thi at det staaer i det nye Testamente er jo aldeles ligegyldigt, naar Præsten ved Ed er forpligtet paa det nye Testamente er Du jo fuldkommen sikkret, at Intet forties af hvad der staaer i det nye Testamente) bliv hos Præsterne, stræb af bedste Evne at sætte Dig det ret fast, at Biskop Mynster var et Sandhedsvidne, et af de rette, en af den hellige Kjæde, Biskop Martensen et Dito Dito, enhver Præst ligeledes, og den officielle Christendom den saliggjørende Sandhed; at det var derfor Christus i de rædsomste Qvaler, endog forladt af Gud, udaandede paa Korset, for at vi skulle have Lyst til at anvende vor Tid og Flid og Kraft paa klogt og smagfuldt at nyde dette Liv; at Hensigten med hans Komme til Verden egentligen var at oplive Børne-Avlingen, hvorfor det ogsaa er »upassende, at Nogen er Præst, som ikke er gift«; og at hans Livs uforglemmelige Betydning er, ved sin Død (Eens Død en Andens Brød!) at have (som en sand Velgjører!) muliggjort en ny Næringsvei, Præsternes, en Næringsvei, der maa ansees for en af de fordeelagtigste, ligesom den ogsaa tæller det største Antal Næringsdrivende, Spediteurer, Rhedere, hvis Geschäft er for en (i Forhold til Reisens Vigtighed, Veiens Længde, Ankomst-Stedets Herlighed, Opholdets Langvarighed) fast utroligt billig Godtgjørelse at udskibe Folk til Evighedens Salighed, en Geschäft der, eneste i sit Slags, har, sammenlignet med alle Udskibninger til America, Australien og saa videre, den uvurdeerlige Fordeel, som sikkrer Rhederiet mod endog Muligheden af at kunne komme i Miscredit, at man aldeles ingen Efterretning har fra de Udskibede.

Har Du derimod dog Mod til at ville have det Mod, der meest af Alt frygter det at være i en Vildfarelse, saa kan Du ogsaa faae Sandheden at vide betræffende det at blive Christen. Sandheden er: det at blive Christen er at blive, menneskeligt talt, ulykkelig for dette Liv; Forholdet er: jo mere Du indlader Dig med Gud, og jo mere han elsker Dig, jo mere vil Du blive, menneskeligt talt, ulykkelig for dette Liv, jo mere vil Du komme til at lide i dette Liv.

Og denne Tanke, der rigtignok lader en noget forstyrrende Belysning falde over (hvad der jo skal være det nye Testamentes Christendom!) hele det muntre, børne-avlende, carrieregjørende Præste-Laugs livslystige Trafik, og som et Lyn gjennemlysner dette phantastiske Blendværk, Maskeraden, Selskabs-Legen, Narrestregen med – alle Sandsebedragenes Tilhold! – »Christenhed« christne Stater, Lande, en christen Verden, denne Tanke er for et stakkels Menneske frygtelig, dræbende, fast overmenneskeligt anstrengende. Dette veed jeg paa to Maader af Erfaring. Deels deraf, at jeg egentligen ikke kan holde den Tanke ud, og derfor blot opdagende øiner denne sande christelige Bestemmelse af det at være Christen, medens jeg for mit Vedkommende hjælper mig til at udholde Lidelser med en langt lettere Tanke, en jødisk, ikke en i høieste Forstand christelig: at jeg lider for mine Synders Skyld; deels deraf, at jeg paa ganske særlig Maade ved mit eget Livs Forhold maatte føres til at blive opmærksom herpaa; hvis ikke, var jeg aldrig bleven opmærksom, og havde endnu mindre formaaet at bære Trykket af denne Tanke; men jeg blev, som sagt, hjulpet ved mit eget Livs Forhold.

Mit eget Livs Forhold vare mine Forkundskaber, ved Hjælp af disse blev jeg, alt som jeg udvikledes i Aarenes Løb, mere og mere opmærksom paa Christendommen, og paa Bestemmelsen: at blive Christen. Hvad er nemlig ifølge det nye Testamente det at blive Christen, hvortil den idelige og idelige Formaning mod ikke at forarges, og hvorfra de frygtelige Collisioner (at hade Fader, Moder, Hustru, Barn og saa videre), hvori det nye Testamente aander? mon ikke begge Dele, fordi Christendommen meget godt veed, at det at blive Christen er at blive menneskelig talt ulykkelig for dette Liv, dog saligt forventende en evig Salighed? Thi hvad er ifølge det nye Testamente det at blive elsket af Gud? det er at blive, menneskelig talt, ulykkelig for dette Liv, dog saligt forventende en evig Salighed – anderledes kan, ifølge det nye Testamente, Gud, som er Aand, ikke elske et Menneske; han gjør Dig ulykkelig, men han gjør det af Kjerlighed, salig Den, som ikke forarges! Og hvad er ifølge det nye Testamente det at elske Gud? det er at ville blive, menneskelig talt, ulykkelig for dette Liv, dog saligt forventende en evig Salighed – anderledes kan et Menneske ikke elske Gud, som er Aand. Og alene ved Hjælp heraf kan Du see, at det nye Testamentes Christendom slet ikke er til, at den Smule Religieusitet der findes i Landet høist er Jødedom.

\sectionheading{Hvad siger Brand-Majoren? (Øieblikket Nr. 6)  ·  SV¹ XIV}

Hvad siger Brand-Majoren?

At et Menneske, naar han paa nogen Maade har hvad man kalder Sag, Noget han for Alvor vil – og der da ere Andre, som sætte sig det til Opgave at modvirke, forhindre, skade, at han da maa tage sine Forholdsregler mod disse hans Fjender: er Enhver strax opmærksom paa. Men at der gives en, om man saa vil, skikkelig Velmenen, der maaskee er langt farligere og som beregnet paa at forhindre, at Sagen i Sandhed bliver Alvor: er ikke Enhver strax opmærksom paa.

Naar et Menneske pludseligt bliver syg, flux iler der Velmenende til Hjælp, den ene foreslaaer Det, den anden Det; fik de alle Samtlige Lov til at raade, var vel den Syges Død vis, den Enkeltes velmeente Raad kan maaskee allerede være betænkeligt nok. Og selv om Intet af dette skeer, hverken de samtlige Velmenendes eller den Enkeltes Raad følges: deres travle, befippede Tilstedeværelse er dog maaskee alligevel til Skade, forsaavidt de staae Lægen i Veien.

Saaledes ogsaa ved en Ildebrand. Neppe er Brandraabet hørt, før der stormer en Menneske-Masse hen til Stedet, rare, hjertelige, deeltagende, hjælpsomme Mennesker, den ene har en Strippe, den anden en Spølkom, den tredie en Pibesprøite, og saa videre, alle rare, hjertelige, deeltagende, hjælpsomme Mennesker, der saa gjerne ville hjælpe til med at slukke.

Men hvad siger Brand-Majoren? Brand-Majoren, han siger – ja, ellers er Brand-Majoren en meget behagelig og dannet Mand; men ved en Ildebrand er han, hvad man kalder, grov i sin Mund – han siger eller rettere han brøler det ud: »aah, reis ad hele Helvede til med samt Eders Stripper og Pibesprøiter.« Og naar saa disse Velmenende maaskee blive fornærmede, finder det høist usømmeligt at behandle dem saaledes, og forlanger idetmindste at behandles med Agtelse: hvad siger saa Brand-Majoren? Ja, ellers er Brand-Majoren en meget behagelig og dannet Mand, som veed at vise Enhver den Agtelse, ham tilkommer, men ved en Ildebrand er han noget anderledes – han siger: »hvor Satan er Politiet!« Og kommer der saa nogle Betjente siger han til dem: »skil mig af med disse forbandede Mennesker med deres Stripper og Pibesprøiter; og vil de ikke med det Gode, saa smør dem nogle over Rygstykkerne, at vi kan blive af med dem – og komme til.«

Altsaa ved en Ildebrand er hele Betragtningsmaaden en ganske anden end i det stille daglige Liv. Det, hvorved man i det stille daglige Liv opnaaer at være godt lidt: en godmodig, skikkelig Velmenen, Det honoreres ved en Ildebrand med Grovheder og tilsidst med nogle over Rygstykkerne.

Og dette er aldeles i sin Orden. Thi en Ildebrand det er Alvor; og overalt hvor det virkelig er Alvor, strækker det aldeles ikke til med denne skikkelige Velmenen. Nei Alvor indfører en ganske anden Lov: enten-eller; enten er Du Den, der her for Alvor kan gjøre Noget og for Alvor her har Noget at gjøre, eller, hvis Du ikke er i det Tilfælde, da er Alvoren just, at Du forføier Dig bort. Vil Du ikke ved Dig selv forstaae det, saa lader Brand-Majoren Dig det ved Politiet indbanke, hvad Du kan have særdeles godt af, og hvad dog maaskee kan bidrage til at gjøre Dig lidt alvorlig, svarende til den Alvor en Ildebrand er.

Men som ved en Ildebrand, saaledes ogsaa i Aandens Forhold. Overalt hvor der er en Sag der skal fremmes, et Foretagende der skal drives igjennem, en Idee der skal anbringes – man kan altid være forvisset om, at naar Den, som egentligen er Manden, den Rette, han som i høiere Forstand har og skal have Commandoen, han der har Alvoren og kan give Sagen den Alvor, den i Sandhed har – man kan altid være forvisset om, at naar han kommer, om jeg saa tør sige, til Stedet, at han da vil forefinde et gemytligt Vrøvle-Compagnie, der, under Navn af Alvor, ligger og fusker i at ville tjene denne Sag, fremme dette Foretagende, anbringe denne Idee; et Vrøvle-Compagnie, der naturligviis betragter det ikke at ville gjøre fælleds Sag med dem (hvilket just er: Alvoren) som et sikkert Beviis paa, at Vedkommende mangler Alvor. Jeg siger, naar den Rette kommer, vil han forefinde Dette, jeg kan ogsaa vende Sagen saaledes: Det om han er den Rette, afgjøres egentlig ved, hvorledes han forstaaer sig selv i Forhold til dette Vrøvle-Compagnie. Mener han, at det er Dem, der skal hjælpe, og at han skal styrke sig ved Forening med dem: er han eo ipso ikke den Rette. Den Rette seer strax, ligesom Brand-Majoren, at dette Vrøvle-Compagnie maa bort, at dets Tilstedeværelse og Virken er den farligste Bistand, Ildebranden kunde faae. Men i Aands-Forhold er det ikke som ved en Ildebrand, hvor Brand-Majoren blot behøver at sige til Politiet: skil mig af med disse Mennesker.

Saaledes i alle Aands-Forhold, og saaledes ogsaa paa det religieuse Gebeet. Man har ofte nok sammenlignet Historien med, hvad Chemikerne kalde en Proces. Billedet kan være ganske betegnende, vel at mærke, naar det forstaaes rigtigt. Man taler om en Affiltrerings-Proces; Vand affiltreres, afsætter i denne Proces de urenere Bestanddele. Ganske i modsat Forstand er Historien en Proces. Ideen anbringes – og gaaer nu ind i Historiens Proces. Men denne bestaaer uheldigviis ikke i – latterlige Antagelse! – at luttre Ideen, der aldrig er renere end i sit Første, nei, den bestaaer i, stadigt tiltagende, at forkluddre, forsluddre, forvrøvle Ideen, at forbruge Ideen, at – det Modsatte af at filtrere – anbringe de, oprindeligt manglende, urenere Bestanddele, indtil det tilsidst ved en Række Slægters begeistrede og af hinanden gjensidig anerkjendte successive Samvirken, er naaet, at Ideen er aldeles gaaet ud, det Modsatte af Ideen blevet Det, som nu kaldes Ideen, hvilket paastaaes at være naaet ved den historiske Proces, hvori Ideen luttres og forædles.

Naar saa endelig den Rette kommer, han som i høieste Forstand har Opgaven, maaskee tidligt udseet og langsomt opdraget til denne Forretning, som er at bringe Lys i Sagen, at faae sat Ild paa dette Vildnis, alt Vrøvlets, alle Sandsebedragenes, alle Gavtyvestregernes Tilhold – naar han kommer, han vil altid forefinde et Vrøvle-Compagnie, der i gemytlig Hjertelighed nok saadan mener, at det er forkeert, og at der maa gjøres Noget, eller som har indrettet sig paa at snakke om, at det er uhyre forkeert, at blive sig selv vigtig ved at snakke derom. Dersom han, den Rette, noget Secund seer feil, og mener, at det er dette Compagnie, der skal hjælpe: er han eo ipso ikke den Rette. Dersom han griber feil, og indlader sig med dette Compagnie: slipper Styrelsen ham øieblikkeligt som ubrugbar. Men den Rette, han seer med et halvt Øie, hvad Brand-Majoren seer, at det Compagnie, som velmenende vil hjælpe til at slukke en Ildebrand med en Strippe eller en Pibesprøite, at dette samme Compagnie, der her, hvor Talen ikke er om at slukke en Ildebrand, men just om at faae Ild sat paa, velmenende vil hjælpe til med en Svovlstik uden Svovl eller en vaad Fidibus – at dette Compagnie maa bort, at han ikke maa have det allermindste med dette Compagnie at gjøre, at han maa være saa grov i sin Mund mod dem, som vel muligt, han som maaskee ellers er Intet mindre. Men Alt gjælder om at blive det Compagnie qvit; thi dets Virken er, i Skikkelse af hjertelig Deeltagelse, at udmarve den egentlige Alvor af Sagen. Naturligviis vil saa Compagniet rase mod ham, mod denne frygtelige Hovmod og deslige Det maa ikke gjøre ham hverken fra eller til. Overalt hvor det i Sandhed skal være Alvor er Loven den: enten-eller; enten er jeg Den, som for Alvor har med Sagen at gjøre, kaldet dertil og ubetinget villig til at vove afgjørende, eller, hvis dette ikke er mit Tilfælde, saa er Alvoren den: aldeles ikke at befatte mig dermed. Intet er afskyeligere, nederdrægtigere, forraadende og bevirkende en dybere Demoralisation end dette: at ville være saadan lidt med i Forhold til hvad der skal være aut – aut, aut Cæsar aut nihil, at ville være saadan lidt med, saa hjerteligt smaat, pjatte derom, og saa ved denne Pjat tillyve sig at være bedre end De, der slet ikke befatte sig med det hele Anliggende – tillyve sig at være bedre og vanskeliggjøre Sagen for Den, der virkelig har Opgaven.

\sectionheading{»Først Guds Rige.« Et Slags Novelle (Øieblikket Nr. 7)  ·  SV¹ XIV}

»Først Guds Rige.« Et Slags Novelle.

Candidatus theologiaeLudvig From – han søger. Og naar man hører, at en »theologisk« Candidat søger, behøver man da ikke nogen levende Indbildningskraft for at forstaae, hvad det er han søger, naturligviis Guds Rige, som man jo skal søge først.

Nei, det er det dog ikke; det han søger er: et kongeligt Levebrød som Præst; og er der, hvad jeg med nogle faa Træk skal angive, først skeet saare Meget, førend han er naaet saavidt.

Først har han gaaet i den lærde Skole, fra hvilken han saa er demitteret. Derpaa har han først taget 2 Examiner, og, efter 4 Aars Læsning først taget Embedsexamen.

Han er da altsaa theologisk Candidat; og man skulde maaskee mene, at han, efter først at have tilbagelagt alt Dette, endeligen kan komme til at virke for Christendommen. Ja, paa det Lav. Nei, først maa han gaae ½ Aar i Seminariet; og naar det er gaaet, kan der da ikke være Tale om at kunne søge i de første 8 Aar, hvilke altsaa først maa være tilbagelagte.

Og nu staae vi ved Novellens Begyndelse: de 8 Aar ere forløbne, han søger.

Hans Liv, som hidtil ikke kan siges at have havt noget Forhold til det Ubetingede, antager pludseligt et saadant Forhold: han søger ubetinget Alt; skriver det ene Ark stemplet Papir fuldt efter det andet; render fra Herodes til Pilatus; anbefaler sig baade hos Ministeren og hos Portneren, kort han er ganske i et Ubetingets Tjeneste. Ja, en af hans Bekjendte, der i de sidste Par Aar ikke har seet ham, mener til sin Forbauselse at opdage, at han er blevet mindre, hvilket maaskee lader sig forklare af, at det er gaaet ham som Münchhausens Hund, der var en Mønde, men ved den megen Løben blev en Grævling.

Saaledes gaaer der 3 Aar. Vor theologiske Candidat trænger virkelig til Hvile, til, efter en saa uhyre anstrenget Virksomhed at blive sat ud af Virksomhed eller at komme til Ro i et Embede og pleies lidt af hans tilkommende Kone – thi han er imidlertid først blevet forlovet.

Endeligen – som Pernille siger til Magdelone – slaaer hans »Forløsnings« Time, saa han med Overbeviisningens hele Magt vil, af egen Erfaring, kunne »vidne« for Menigheden, at i Christendommen er der Frelse og Forløsning: han faaer et Embede.

Hvad skeer? Ved at indhente en endnu nøiere Efterretning end han hidtil har havt om Kaldets Indtægter opdager han, at disse ere circa 150 Rigsbankdaler mindre end han havde troet. Nu er 101 ude. Det ulykkelige Menneske næsten fortvivler. Han har allerede kjøbt stemplet Papir for at indgaae med en Ansøgning til Ministeren, at det maa tillades ham at betragtes som ikke kaldet – og for saa igjen at begynde forfra: da en af hans Bekjendte faaer ham til at opgive dette. Det bliver altsaa derved han beholder Kaldet.

Han er ordineret – og Søndagen kommet, da han skal fremstilles for Menigheden. Provsten, ved hvem dette skeer, er en mere end almindelig Mand, han har ikke blot (hvad de fleste Præster have, og oftest i samme Grad mere jo høiere de komme op i Rangen) et uhildet Blik for den jordiske Profit, men tillige et speculativt Blik paa Verdenshistorien, Noget han ikke beholder for sig selv, men lader Menigheden komme til Gode. Han har, genialt, valgt til Text de Ord af Apostelen Peter »see, vi have forladt alle Ting og fulgt Dig«, og forklarer nu Menigheden, at just i Tider som vore maae der saadanne Mænd til som Lærere, og anbefaler i Forbindelse hermed denne unge Mand, om hvem Provsten veed, hvor nær han var traadt tilbage for de 150 Rdlrs. Skyld.

Den unge Mand bestiger nu selv Prædikestolen – og Dagens Evangelium er (besynderligt nok!): tragter først efter Guds Rige.

Han holder sin Prædiken. »En meget god Prædiken« siger Biskoppen, som selv var tilstede, »en meget god Prædiken; og det frembragte ordentlig Virkning det hele Partie om »først« Guds Rige, den Maade paa hvilken han fremhævede dette Først«. »Men synes da Deres Høiærværdighed, at her var den ønskelige Overeensstemmelse mellem Talen og Livet; paa mig gjorde det næsten et satirisk Indtryk dette Først«. »Hvilken Urimelighed; han er jo kaldet til at forkynde Læren, den sunde, uforfalskede Lære om at søge først Guds Rige; og det gjorde han meget godt«.

Det er den Art Gudsdyrkelse man – under Eed! – vover at byde Gud, den rædsomste Forhaanelse:

Hvo Du end er, tænk blot dette Guds Ord »først Guds Rige«: og tænk saa paa denne Novelle, der er saa sand, saa sand, saa sand: og Du vil ikke behøve mere for at det bliver Dig tydeligt, at hele den officielle Christendom er et Dyb af Usandhed og Øienforblindelse, noget saa Vanhelligt, at det Eneste, der sandt er at sige derom er: Du har (hvis Du ellers deeltager i den offentlige Gudsdyrkelse) ved at lade være at deeltage i den, som den nu er, bestandigt een og en stor Skyld mindre, den ikke at deeltage i at holde Gud for Nar (confer »Dette skal siges, saa være det da sagt«).

Guds Ord lyder »først Guds Rige« og Fortolkningen, maaskee endog »Fuldkommengjørelsen« (thi man lader sig ikke lumpe) er: først alt Andet og sidst Guds Rige; langt om længe naaes først det Jordiske, og saa endelig kommer tilsidst en Prædiken om: først at søge Guds Rige. Saaledes bliver man Præst; og Præstens hele Praxis er saa en stadig Udøven af dette: først det Jordiske og saa – Guds Rige, først Hensynet til det Jordiske, om det behager Regjeringen, eller Majoriteten, eller man dog er nogle Stykker det vil sige først Hensynet til hvad Menneske-Frygt byder eller forbyder, og saa Guds Rige; først det Jordiske, først Pengene, og saa kan Du faae Dit Barn døbt, først Pengene, saa skal der blive kastet Jord paa og holdt Tale efter Taxten, først Pengene saa skal jeg komme til den Syge; først Pengene og saa: virtus post nummos, først Pengene og saa Dyden, saa Guds Rige, hvilket Sidste tilsidst kommer i den Grad tilsidst, at det slet ikke kommer, og det Hele bliver ved det Første: Pengene – i dette eneste Tilfælde føler man ikke Trangen til »at gaae videre«.

Saaledes forholder paa ethvert Punkt og i Alt officiel Christendom sig til det nye Testamentes Christendom. Og dette er saa ikke hvad man selv vedgaaer at være en Jammerlighed, nei, frækt trodses paa, at Christendommen er perfectibel, at man ikke kan blive staaende ved den første Christendom, at den er blot et Moment og saa videre

Derfor er Intet Gud saaledes imod som officiel Christendom og Deeltagelse deri med Paastand paa at det er at dyrke ham. Dersom Du troer, og det troer Du jo dog, at stjæle, røve, plyndre, hore, bagtale, fraadse og saa videre er Gud imod: den officielle Christendom og dens Gudsdyrkelse er ham uendeligt væmmeligere. At et Menneske kan være sjunken i en saadan dyrisk Dumhed og Aandløshed, at han vover at byde Gud slig Dyrkelse, hvor Alt er Tankeløshed, Aandløshed, Dorskhed; og at saa Mennesket frækt vover at ansee det for et Fremskridt i Christendom!

Dette er det min Pligt at sige, dette: »hvo Du end er, hvilket Dit Liv forresten – ved (hvis Du ellers deeltager) at lade være at deeltage i den offentlige Gudsdyrkelse, som den nu er, har Du bestandigt een og en stor Skyld mindre.« Selv bærer Du saa og bære Du saa Ansvaret for, hvorledes Du handler; men Du er gjort opmærksom!

\sectionheading{Man lever kun een Gang (Øieblikket Nr. 8)  ·  SV¹ XIV}

Man lever kun een Gang.

Dette Ord høres saa ofte i Verden. »Man lever kun een Gang; derfor kunde jeg ønske at see Paris, inden jeg døer, eller snarest muligt at samle en Formue, eller dog tilsidst at blive noget Stort i denne Verden – thi man lever kun een Gang.«

Sjeldnere forekommer, men det forekommer dog, et Menneske, som kun har eet Ønske, ganske bestemt kun eet Ønske. »Dette«, siger han, »dette kunde jeg ønske; o, at det maatte opfyldes mit Ønske, thi, ak, man lever kun een Gang!«

Tænk Dig nu et saadant Menneske paa hans Dødsleie. Ønsket blev ikke opfyldt, men hans Sjel hænger uforandret ved dette Ønske – og nu, nu er det ikke mere muligt. Da reiser han sig paa Leiet; med Fortvivlelsens Lidenskab udtaler han endnu engang sit Ønske »o, Fortvivlelse, det opfyldes ikke, Fortvivlelse, man lever kun een Gang!«

Dette synes forfærdeligt, og er det i Sandhed; men ikke som han mener det; thi det Forfærdelige er endda ikke, at Ønsket ikke blev opfyldt, det Forfærdelige er den Lidenskab, med hvilken han hænger ved det. Hans Liv er ikke spildt, fordi hans Ønske ikke blev opfyldt, paa ingen Maade; er hans Liv spildt, er det fordi han ikke vilde opgive sit Ønske, intet Høiere vilde lære af Livet end det om hans eneste Ønske, som afgjorde dets Opfyldelse eller Ikke-Opfyldelse Alt.

Det sande Forfærdelige er derfor et ganske Andet, hvis for Exempel et Menneske paa sit Dødsleie opdagede, eller hvis det dog blev ham paa Dødsleiet saa tydeligt, hvad han sit Liv igjennem dunklere havde forstaaet, men aldrig havde villet forstaae: at det i denne Verden at have lidt for Sandhed hører med til for at kunne blive evig salig – og man lever kun een Gang, den ene Gang, som nu er for ham forbi! Og man havde jo havt det i sin Magt; og Evigheden forandrer man ikke, Evigheden, hvilken man saa just døende gaaer imøde som sin Fremtid!

Vi Mennesker ere af Naturen tilbøielige til at betragte Livet paa denne Maade: vi ansee det at lide for et Onde, som vi paa enhver Maade stræbe at undgaae. Og hvis dette da lykkes os, mene vi engang paa vort Dødsleie med særdeles Grund og Føie at kunne takke Gud, at vi bleve forskaanede for at lide. Vi Mennesker mene, at det, det kommer an paa, er, at vi blot kan slippe lykkelig og vel igjennem denne Verden; og Christendommen mener, at alle Forfærdelserne egentligen komme fra den anden Verden, at denne Verdens Forfærdelser ere som en Barnagtighed i Sammenligning med Evighedens Forfærdelser, og at det just derfor ikke kommer an paa at slippe lykkelig og vel gjennem dette Liv, men paa ved at lide at forholde sig rigtigt til Evigheden.

Man lever kun een Gang; er, naar Døden kommer, Dit Liv vel brugt det vil sige brugt, saa det forholder sig rigtigt til Evigheden: Gud være evigt lovet; er det det ikke, det er evigt uopretteligt – man lever kun een Gang!

Man lever kun een Gang; saaledes er det her paa Jorden. Og medens Du nu lever denne ene Gang, hvis Udstrækning i Tid svinder ind med hver svindende Time, sidder Kjerlighedens Gud i Himlen, kjerligt elskende ogsaa Dig. Ja, elskende; derfor vilde han saa gjerne, at Du dog endelig vil, som han for Evighedens Skyld vil det med Dig: at Du dog kunde beslutte Dig til at ville lide, det er, at Du kunde beslutte Dig til at ville elske ham, thi ham kan Du kun elske ved at lide, eller hvis Du elsker ham, som han vil være elsket, vil Du komme til at lide. Husk paa, man lever kun een Gang; er det forsømt, kom Du ikke til at lide, unddrog Du Dig: det er evigt uopretteligt. Tvinge Dig, nei, det vil Kjerlighedens Gud for ingen Priis, han vilde saa faae noget ganske Andet end det han vil, hvor kunde det ogsaa falde Kjerlighed ind at ville tiltvinge sig at blive elsket! Men Kjerlighed er han, og af Kjerlighed er det, at han vil, Du skulde som han vil; og i Kjerlighed lider han, som kun uendelig og almægtig Kjerlighed kan det, hvad intet Menneske formaaer at fatte, saaledes lider han, naar Du ikke vil som han vil.

Gud er Kjerlighed; aldrig var det Menneske født, hvem denne Tanke, især naar den kommer ham saaledes nærmere, at det at Gud er Kjerlighed betyder, Du er elsket, hvem denne Tanke ikke overvælder i ubeskrivelig Salighed. I næste Øieblik naar Forstaaelsen kommer »dette betyder at komme til at lide«: frygteligt! »Ja, men det er af Kjerlighed Gud vil det, det er fordi han vil være elsket; og det at han vil være elsket af Dig, er hans Kjerlighed til Dig«: nu vel! – I næste Øieblik, saa snart det bliver Alvor med at lide: frygteligt! »Ja, men det er af Kjerlighed; Du aner ikke, hvorledes han lider, fordi han meget godt veed at det at lide smerter, men forandres kan han dog ikke, saa maatte han jo blive til noget Andet end Kjerlighed«: nu, vel! – I næste Øieblik, saasnart det bliver ret Alvor med at lide: frygteligt!

Dog pas paa, pas paa, at Tiden ikke, maaskee i unyttig Lidelse, gaaer unyttet hen, husk paa: man lever kun een Gang; kan det hjælpe Dig, saa see ogsaa Sagen saaledes: vær forvisset om, at Gud lider mere i Kjerlighed end Du lider, uden at han dog derfor kan forandres. Men frem for Alt husk paa: man lever kun een Gang; der er Tab som ere evigt uoprettelige, saa Evigheden – endnu frygteligere! – langt fra at udslette Erindringen om det Tabte, er en evig Erindren om det Tabte!

\sectionheading{Min Opgave (Øieblikket Nr. 10)  ·  SV¹ XIV}

Min Opgave.

»Jeg kalder mig ikke en Christen, siger ikke mig selv at være en Christen.« Det er dette jeg bestandigt maa gjentage, hvad Enhver, der vil forstaae min, ganske særlige Opgave, maa øve sig i at kunne fastholde.

Ja, jeg veed det nok, det lyder næsten som et Slags Galskab i denne christne Verden, hvor Alle og Enhver er Christen, hvor det at være Christen er Noget, som da naturligviis Enhver er: at der En siger om sig selv: jeg kalder mig ikke en Christen; og En, hvem Christendommen i den Grad beskjeftiger, som den beskjeftiger mig.

Men anderledes kan det ikke være; det Sandere maa altid tage sig ud som en Art Galskab – i Vrøvlets Verden; og at det er en Vrøvlets Verden, hvori jeg lever, at den er det blandt andet ogsaa just ved dette Vrøvl, at saadan udenvidere Enhver er Christen: det er vistnok.

Dog forandre mit Udsagn, hverken kan, vil eller tør jeg – ellers vilde der maaskee ogsaa indtræde en anden Forandring: at den Magt, en Almagt er det, som særligen bruger min Afmagt, slog Haanden af mig og lod mig seile i min egen Sø. Nei, forandre mit Udsagn hverken kan, vil eller tør jeg; jeg kan ikke tjene disse Legioner af næringsdrivende Gavtyve, jeg mener Præsterne, som ved at forfalske Bestemmelsen Christen have, for Næringsbrugets Skyld, faaet Millioner og Millioner Christne: jeg er ikke en Christen – og uheldigviis, jeg kan gjøre det aabenbart, at de Andre ere det heller ikke, ja, endnu mindre end jeg; thi de indbilde sig at være det, eller de tillyve sig at være det, eller de (som Præsterne) indbilde Andre, at de ere det, hvorved Præste-Næringsveien bliver florerende.

Det Synspunkt jeg har at udvise og udviser, er af en saadan Eiendommelighed, at jeg i Christenhedens 1800 Aar ganske bogstavelig intet Analogt, intet Tilsvarende har at holde mig til. Ogsaa saaledes – lige over for 1800 Aar – staaer jeg ganske bogstavelig ene.

Den eneste Analogie jeg har for mig er: Socrates; min Opgave er en socratisk Opgave, at revidere Bestemmelsen af det at være Christen: selv kalder jeg mig ikke en Christen (holdende Idealet frit), men jeg kan gjøre aabenbart, at de Andre ere det endnu mindre.

Du Oldtidens ædle Eenfoldige, Du det eneste Menneske, jeg beundrende anerkjender som Tænker: det er kun Lidt, der er opbevaret om Dig, blandt Mennesker den eneste sande Intellectualitetens Martyr, lige saa stor qua Charakteer som qua Tænker; men dette Lidet hvor uendelig Meget! Hvor jeg, bort fra disse Batailloner af Tænkere, som »Christenhed« stiller i Marken under Navn af christelige Tænkere – thi ellers har der dog i Aarhundredernes Løb levet i Christenhed nogle ganske enkelte betydelige Tænkere – kan længes efter blot en halv Time at kunne tale med Dig!

Det er i et Dyb af Sophistik »Christenhed« ligger, langt, langt værre end da Sophisterne florerede i Grækenland. Disse Legioner Præster og christelige Docenter, de ere alle Sophister, ernærende sig – her er jo ifølge Oldtiden Sophistens Charakteristiske – af at bilde Dem, som Intet forstaae, Noget ind og saa at gjøre dette Menneske-Tal til Instantsen for, hvad Sandhed, hvad Christendom er.

Men jeg kalder mig ikke en Christen. At dette er meget generende for Sophisterne, forstaaer jeg meget godt, forstaaer meget godt, at de langt hellere saae, at jeg med Pauker og Trompeter forkyndte mig selv at være den eneste sande Christen, forstaaer ogsaa meget godt, at man forsøger, usandt, at fremstille min Optræden saaledes. Men man narrer mig ikke! I en vis Forstand er jeg saare let at narre; jeg er næsten blevet narret i ethvert Forhold hvori jeg har været – men saa har det været, fordi jeg selv har villet det. Naar jeg ikke selv vil: er der i min Samtid Ingen, som narrer mig, et afgjort Politie-Talent, som jeg er.

Altsaa man narrer mig ikke: jeg kalder mig ikke en Christen. I en vis Forstand synes det saa let nok at blive mig qvit; thi de Andre ere jo saa alle ganske anderledes Karle, de ere sande Christne. Ja, ja, saaledes synes det. Men det er ikke saaledes; thi just fordi jeg ikke kalder mig en Christen, er det umuligt at blive mig qvit, mig, der saa har den forbistrede Egenskab, at jeg, just ogsaa ved Hjælp af at jeg kalder mig ikke en Christen, kan gjøre det aabenbart, at de Andre ere det endnu mindre.

O, Socrates! Havde Du med Pauker og Trompeter forkyndt Dig som den meest Vidende, vare Sophisterne snart blevne færdige med Dig. Nei, Du var den Uvidende; men Du havde tillige den forbistrede Egenskab, at Du(just ogsaa ved Hjælp af, at Du selv var den Uvidende) kunde gjøre aabenbart, at de Andre vare endnu mindre Vidende end Du, de vidste ikke engang, at de vare Uvidende.

Men som det gik Dig (efter hvad Du siger i Dit »Forsvar«, som Du ironisk nok har kaldet den grusomste Satire over en Samtid) at Du derved paadrog Dig mange Fjender, ved at gjøre det aabenbart, at de vare Uvidende; og som man paaduttede Dig, at Du selv maatte være, hvad Du kunde vise, at de Andre ikke vare, og derfor i Misundelse fattede Nag til Dig: saaledes er det ogsaa gaaet mig. Det har reist Forbittrelse mod mig, at jeg kan gjøre aabenbart, at de Andre ere endnu mindre Christne end jeg, som dog er saa meget forholdende mig til Christendom, at jeg sandt seer og vedgaaer, at jeg ikke er en Christen. Og man vil paanøde mig, at dette, at jeg ikke er en Christen, kun er en skjult Form for Hovmod, at jeg vel maa være, hvad jeg kan vise de Andre ikke ere. Men dette er en Misforstaaelse, det er aldeles sandt: jeg er ikke en Christen; og det er en overilet Slutning, at fordi jeg kan vise, at de Andre ikke ere Christne, at jeg derfor selv maa være det, lige saa overilet som af, at En er for Exempel ½ Qvarteer høiere end Andre strax at slutte ergo er han 6 Alen.

Min Opgave er at revidere Bestemmelsen: Christen. Der lever kun eet eneste Menneske, som har Forudsætninger til at kunne levere en virkelig Kritik over min Arbeiden: det er mig selv. Der var derfor noget Sandt i, hvad allerede for en Deel Aar siden nuværende Provst Kofoed-Hansen sagde til mig i Anledning af at han havde paatænkt at levere en Critik over Afsluttende Efterskrift, at han, ved at læse den i Afsluttende Efterskrift indeholdte Critik over den tidligere Forfatter-Virksomhed, opgav at ville criticere, hvor Forfatteren var den Eneste, der kunde præstere en virkelig Critik. Nei, der er i Samtiden ikke en Eneste, som har Forudsætninger til at kunne levere en Critik over min Arbeiden. Den Eneste, der leilighedsviis har sagt et sandere Ord om min Betydning, er Professor R. Nielsen; men dette Sande har han haft fra privat Samtale med mig.

Naar nu saa competente Instantser som for Exempel De herrer Israel Levin, Davidsen, Siesby, eller saa modnede Aander som Redakteur Bille, eller saa uomtaagede Tænkere som Grüne, eller saa aabne Charakterer som Anonymer og saa videre – naar alt Sligt skal, og lige over for en saa vel underrettet Domstol som Publikum, bedømme en saadan Eiendommelighed: følger det af sig selv, at det maa blive til – ja til hvad det er blevet til, hvad der kun smerter mig paa det lille Folks Vegne, der ved saadant Forhold virkelig gjøres latterligt qua Folk.

Men selv om een eller anden dog adskilligt bedre Underrettet paatager sig at ville sige Noget om mig og min Opgave: det bliver egentligen ikke til Andet, end at han, efter et flygtigt Blik paa Mit, nu i en Fart finder eet eller andet Tidligere, hvilket han erklærer for at være Tilsvarende.

Paa den Maade bliver det alligevel ikke til Noget. Hvad et Menneske med mit Otium, min Flid, mine Evner, min Dannelse (hvad jo Biskop Mynster har givet mig Attest for paa Prent) har anvendt ikke blot 14 Aar men i Grunden hele sit Liv paa, det Eneste han har levet og aandet for: at saa een eller anden Præst, høist en Professor, ikke skulde behøve mere end et flygtigt Blik for at kunne vurdere det, er dog en Taabelighed. Og at hvad der i den Grad har været særligt mærket, at det strax har haft Stemplet »den Enkelte – jeg er ikke en Christen«, hvad ganske vist ikke er forekommet i Christenhedens 1800 Aar, hvor Alt er stemplet »Menighed, Samfund – jeg er en sand Christen«: at saa een eller anden Præst, høist en Professor, strax skulde finde en Analogie hertil, er ogsaa en Taabelighed; ved et omhyggeligere Eftersyn vilde han netop opdage, at det er en Umulighed. Men det er det man ikke finder Umagen værd, man foretrækker et flygtigt Blik paa Mit og et lige saa flygtigt paa det Tidligere, og saa har man strax Analogier nok for Mit, hvad Publikum strax kan forstaae.

Dog er det som jeg siger: der er i Christenhedens 1800 Aar aldeles intet Tilsvarende, intet Analogt til min Opgave; det er i »Christenhed« første Gang.

Det veed jeg, veed ogsaa, hvad det har kostet, hvad jeg har lidt, hvilket dog kan udtrykkes ved eet eneste Ord: jeg var aldrig som Andre. O, i Ungdommens Dage af alle Qvaler rædsomste, intensiveste: ikke at være som Andre, aldrig at leve nogen Dag uden smerteligt at mindes om, at man ikke er som Andre, aldrig at kunne løbe med i Flokken, Ungdommens Lyst og Glæde, aldrig frit at kunne give sig hen, altid, saasnart man vilde vove det, smerteligt mindet om Lænken, Særskilthedens Afsondring, der, indtil Fortvivlelse smerteligt, skiller En fra Alt hvad der hedder menneskelig Liv og Munterhed og Glæde. Sandt, man kan ved de frygteligste Anstrengelser stræbe at skjule over, hvad man i den Alder forstaaer som sin Vanære, at man ikke er som de Andre; det kan maaskee til en vis Grad lykkes: men Qvalen er dog alligevel i Hjertet; og det er dog kun til en vis Grad, saa en eneste Uforsigtighed kan hævne sig frygteligt.

Med Aarene forvindes saa vistnok denne Smerte mere og mere; thi alt som man mere og mere bliver Aand, smerter det ikke, at man ikke er som Andre, Aand er just: ikke at være som Andre.

Og saa kommer da maaskee endelig et Øieblik, hvor den Magt, som saaledes engang har – ja, saa syntes det stundom – næsten mishandlet En, forklarer sig og siger: »har Du Noget at beklage Dig over, synes det Dig, at jeg, i Sammenligning med hvad der gjøres for andre Mennesker, har forfordelet Dig, om jeg end – af Kjerlighed – har maattet forbittre Dig Din Barndom, baade Din tidligere og sildigere Ungdom, synes det Dig, at jeg har narret Dig ved hvad Du fik i Stedet?« Og dertil kan saa kun Svaret blive »nei, nei, uendelige Kjerlighed«, medens dog vistnok Menneskenes Mængde i høieste Grad vilde betakke sig for Det, jeg paa en saa qvalfuld Maade er blevet.

I saadanne Qvaler som mine opdrages nemlig et Menneske til at kunne bære at være en Offret; og den uendelige Naade, der vistes og vises mig, er at være udseet til at være en Offret, at være udseet dertil, ja, og saa Eet endnu, under Almagtens og Kjerlighedens samvirkende Indflydelse at være udviklet til at kunne fastholde, at dette er den høieste Grad af Naade, som Kjerlighedens Gud kan vise mod Nogen, derfor kun mod de Elskede.

Min kjere Læser, Du seer, det gaaer ikke løs paa Profiten; det bliver først Tilfældet efter min Død, naar de eedfæstede Næringsdrivende ville tage ogsaa mit Liv til Indtægt for Saltmads-Tønden.

Christendommen ligger saa høit, at Det, den forstaaer ved Naade, er hvad alle Profane (procul, o procul este profani) ville af Alt meest betakke sig for. Løgne-Præsterne eller Præsterne faae Naaden forvandlet til Aflad; Naaden bestaaer i, at Mennesket har, ganske ligefrem, Profit af Gud, og Præsten Profit af Menneskene, hvem han indbilder dette, indbydende dem med Christi Ord »kommer hid Alle«, hvilke Ords sande Betydning er, at Indbydelsen unegteligt er til Alle, men at det, naar det kommer til Stykket, og det skal staae fast, hvad det er Christus indbyder dem til (i Efterfølgelse at blive en Offret) og dette ikke gjøres om til Noget, som behager Alle: saa vil det vise sig, som i Samtidigheden med Christus, at Alle af Alt meest betakke sig for dette, og at kun ganske undtagelsesviis en ganske sjelden Enkelt følger Indbydelsen, og af disse Enkelte atter kun en ganske Enkelt følger Indbydelsen saaledes, at han fastholder, at det er uendelig, en ubeskrivelig Naade, der vises ham: at offres. En ubeskrivelig Naade; thi det er den eneste Maade, paa hvilken Gud kan elske et Menneske og være elsket af et Menneske; men det er jo en uendelig Naade, at Gud vil det og vil tillade det. Saa blæse være med, at der, for en Forsigtigheds Skyld, for at fjerne alt Profant, er anbragt som Mellembestemmelse: at blive offret. Og det vilde ogsaa næsten være væmmeligt, lummert, qvalmt, benauende, at det at være elsket af Gud og at turde elske ham, skulde, aandløst, fæisk, belæmres med, at man havde Profit deraf!

Du menige Mand! Det nye Testamentes Christendom er noget uendelig Høit, men vel at mærke ikke saaledes det Høie, at det forholder sig til Forskjellen mellem Menneske og Menneske i Henseende til Begavethed og Deslige. Nei, det er for Alle. Enhver, ubetinget Enhver – hvis han ubetinget vil, vil ubetinget hade sig selv, vil ubetinget finde sig i Alt, lide Alt (og det kan jo Enhver, hvis han vil): saa er dette uendelig Høie tilgængeligt for ham.

Du menige Mand! Jeg har ikke afsondret mit Liv fra Dit; Du veed det, jeg har levet paa Gaden, er kjendt af Alle; videre, jeg er ikke blevet til Noget, tilhører ingen Stands Egoisme: saa hvis jeg tilhører Nogen, maatte jeg tilhøre Dig, Du menige Mand, Du som dog engang, forlokket af En, der, tjenende Dine Penge, affecterede at ville Dit Vel, har været villig nok til at finde mig og min Tilværelse latterlig, Du som dog mindst af Alle har Grund til at være utaalmodig over eller burde være uskjønsom ved, at jeg er med, hvad de Fornemmere snarere have haft Grund til, fordi jeg ikke afgjort har sluttet mig til dem, men kun holdt en løs Forbindelse med dem.

Du menige Mand! Jeg dølger det ikke for Dig, at efter mine Begreber er det at være Christen noget saa uendelig Høit, at der (hvad baade Christi Liv bekræfter, naar man agter paa hans Samtid, og hvad hans Forkyndelse, naar man tager den nøie, antyder) bestandigt kun ere Enkelte, som naaer det: dog er det muligt for Alle. Men Eet besværger jeg Dig for Gud i Himlens Skyld og ved Alt hvad Helligt er, flye Præsterne, flye dem, disse Afskyelige, hvis Næringsvei er at forhindre Dig i endog blot at blive opmærksom paa hvad sand Christendom er, og derved at forvandle Dig, omtaaget af Galimathias og Sandsebedrag, til hvad de forstaae ved en sand Christen, et betalende Medlem af Statskirke, Folkekirke og Deslige. Flye dem; kun at Du da passer paa at betale dem villigt og prompt hvad Penge de skal have. Med hvem man foragter, maa man da for ingen Priis have Penge-Mellemværende, at det maaskee skulde hedde, det var for at slippe for at betale, at man flyede dem. Nei, betal dem dobbelt, at Din Uenighed med dem ret kan blive tydelig: at Det, der beskjeftiger dem, slet ikke beskjeftiger Dig, Pengene; og at derimod Det, der ikke beskjeftiger dem, beskjeftiger Dig uendeligt, Christendommen.

\workheading{Om min Forfatter-Virksomhed; Synspunktet for min Forfatter-Virksomhed}{1851; 1859}

\worknote{Antologiens ''On My Work as an Author'' og ''The Point of View for My Work as an Author''. Her gengives åbningen af ''Om min Forfatter-Virksomhed'' (bevægelsen fra det æsthetiske til det religieuse; ''uden Myndighed''), samt fra ''Synspunktet'' tesen om at Christenheden er et uhyre Sandsebedrag og det berømte slutord (''hiin Enkelte''; digteren der døde af længsel efter evigheden for uafladeligt at takke Gud).}

\sectionheading{Om min Forfatter-Virksomhed (uddrag)  ·  SV¹ XIII, 493–501}

Om min Forfatter-Virksomhed

Naar et Land er lille, er selvfølgeligt i alle Forhold Proportionerne smaae i det lille Land. Saaledes literairt; Honoraret og Alt, hvad dertil hører vil kun være ubetydeligt; det at være Forfatter – naar man da ikke er Digter, og saa igjen dramatisk, eller man skriver Lærebøger eller paa anden Maade er Forfatter i Forhold til en Embedsstilling – omtrent den daarligst aflagte, den mindst betryggede, forsaavidt den utaknemmeligste Ansættelse. Lever der nu en Enkelt, som er i Besiddelse af Evne til at kunne være Forfatter, er han derhos saa lykkelig at have nogen Formue, saa bliver han da Forfatter, noget nær for egen Regning. Dette er imidlertid ganske i sin Orden, derom Intet videre at sige; saaledes skal den Enkelte i Gjerning elske sin Idee, det Folk han tilhører, den Sag han tjener, det Sprog han som Forfatter har den Ære at skrive. Det vil ogsaa skee saaledes, hvor der er Samdrægtighed mellem den Enkelte og Folket, hvilket saa igjen i det givne Tilfælde vil være en Smule skjønsomt mod denne Enkelte.

Om det Modsatte heraf paa nogen Maade skulde have været mit Tilfælde, om jeg af Nogen eller af Nogle skulde være blevet behandlet uskjønsomt: det vedkommer jo egentligen ikke mig, men ret egentligen dem og bliver deres Sag. Hvad der derimod vedkommer mig, og hvad der er mig saa kjert at det vedkommer mig, er at skulle og burde takke for, hvad Gunst og Velvillie og Imødekommen og Paaskjønnen der, i Almindelighed eller af de Enkelte særligen, er blevet mig beviist.

Den Bevægelse Forfatterskabet beskriver er: fra »Digteren« – fra det Æsthetiske, fra »Philosophen« – fra det Speculative til Antydningen af den meest inderlige Bestemmelse i det Christelige: fra den pseudonyme »Enten – Eller« gjennem »Afsluttende Efterskrift«, med mit Navn som Udgiver, til »Taler ved Altergangen om Fredagen«, af hvilke de to vare holdte i Frue Kirke. Denne Bevægelse er tilbagelagt eller beskrevet uno tenore, i eet Aandedrag, om jeg saa tør sige, saa Forfatterskabet, totalt betragtet, er religieust fra Først til Sidst, Noget, Enhver, der kan see, naar han vil see, ogsaa maa see. Og som den Naturkyndige paa Traadenes Krydsning i Spindet, strax kjender, hvilket kunstforstandigt lille Dyr det er, hvis Spind det er: saaledes vil den Indsigtsfulde ogsaa kjende, at til dette Forfatterskab svarer som Ophav En, der qva Forfatter »kun har villet Eet«. Den Indsigtsfulde vil tillige kjende, at dette Ene er det Religieuse, men det Religieuse heelt og holdent sat ind i Reflexion, dog saaledes, at det heelt og holdent tages ud af Reflexion tilbage i Eenfold, det er, han vil see, at den tilbagelagte Vei er: at naae, at komme til Eenfoldighed.

Og dette er ogsaa (i Reflexion, som det faktisk var det oprindeligt) den christelige Bevægelse. Christeligt er det ikke det Eenfoldige, fra hvilket man gaaer ud, for saa at blive interessant, vittig, dybsindig, Digter, Philosoph og saa videre Nei, lige omvendt; her begynder man, og bliver saa eenfoldigere og eenfoldigere, kommer til det Eenfoldige. Dette er, i »Christenhed«, christeligt Reflexions-Bevægelsen; man reflekterer sig ikke ind i Christendom, men man reflekterer sig ud af Andet og bliver, eenfoldigere og eenfoldigere, Christen. Havde Forfatteren været en riigt begavet Aand, eller, hvis han har været det, havde han været en dobbelt saa riigt begavet Aand, vilde han vel have behøvet en længere, eller en dobbelt saa lang Tid for, i Forfatter-Frembringelse, at beskrive denne Vei og naae dette Punkt.

Men som Det, der er blevet meddeelt (det Religieuses Tanke) er blevet sat heelt og holdent ind i Reflexion og atter taget tilbage ud af Reflexion: saaledes har ogsaa Meddelelsen afgjørende været mærket af Reflexion, eller der er brugt den Art Meddelelse, som er Reflexionens. »Ligefrem Meddelelse« er: ligefrem at meddele det Sande. »Meddelelse i Reflexion« er: at bedrage ind i det Sande; men da Bevægelsen er at komme til det Eenfoldige, maa dog atter Meddelelsen engang, tidligere eller sildigere, ende i den ligefremme Meddelelse. Der begyndtes, maieutisk, med æsthetisk Frembringelse, og den hele pseudonyme Productivitet er et saadant Maieutisk. Derfor var denne Productivitet jo ogsaa pseudonym, medens det ligefremme Religieuse – der fra Først af var tilstede i en Antydnings Glimten – bar mit Navn. Det ligefremme Religieuse var fra Først af tilstede; thi »to opbyggelige Taler 1843« ere jo samtidige med »Enten – Eller«. Og for at sikkre denne det ligefremme Religieuses Samtidighed, derfor fulgte saa samtidigt med enhver Pseudonym en lille Samling af »opbyggelige Taler«, indtil »Afsluttende Efterskrift« kom, der stiller Problemet, som ϰατ' εξοχην er »Problemet«, hele Forfatterskabets Problem: »det at blive Christen«. Fra det Øieblik ophører det ligefremme Religieuses Glimten; thi nu indtræder den blot religieuse Produktivitet: »opbyggelige Taler i forskjellig Aand«; »Kjerlighedens Gjerninger«; »Christelige Taler«. Men for omvendt at erindre om Begyndelsen kom (som svarende til hvad »to opbyggelige Taler« var ved Begyndelsen, da de volumineuse Skrifter vare æsthetiske) ved Slutningen (hvor i lang Tid Alt udelukkende og volumineust var blot religieus Produktivitet) en lille æsthetisk Artikel af »Inter et Inter« i Bladet »Fædrelandet« Juli 1848 Nummer 188–191: de to opbyggelige Talers Glimten ved Begyndelsen betydede, at det egentligen var Det, der skulde frem, Det, hvortil der skulde kommes; den lille æsthetiske Artikels Glimten ved Slutningen betydede ved Gjenskæret, at det bragtes til Bevidsthed, at det Æsthetiske fra Begyndelsen af var Det, hvorfra der skulde gaaes bort, Det, der skulde forlades. »Afsluttende Efterskrift« er Midtpunktet, og – hvad da naturligviis kun gjør Fordring paa at være en Curiositet – saa nøiagtigt, at selv Qvantum af Præstation foran og efter er noget nær lige stort, naar man da, som jo er i sin Orden, regner 18 opbyggelige Taler med til den blot religieuse Produktivitet; og selv Forfatterskabets Tid foran »Afsluttende Efterskrift« og efter den er noget nær lige lang.

Endeligen er atter denne Forfatterskabets Bevægelse afgjørende mærket af Reflexion, eller den er Reflexionens. Ligefrem begynder man med enkelte, nogle faa Læsere, og Opgaven eller Bevægelsen er saa at samle Mængde, at faae sig et Abstraktum: Publikum. Her begyndtes, maieutisk, med Sensation, og med hvad dertil hører, Publikum, som altid er med, hvor der er Noget paafærde; og Bevægelsen var, maieutisk, at ryste »Mængde« af for at faae fat i »den Enkelte« religieust forstaaet. Just i det samme Øieblik, da den Sensation, »Enten – Eller« vakte, var paa sit Høieste, just i det samme Øieblik udkom »to opbyggelige Taler« 1843, der anbragte Formlen, som siden blev stereotyp gjentaget »den søger hiin Enkelte, hvem jeg med Glæde og Taknemlighed kalder min Læser«. Og just i det critiske Øieblik, da »Afsluttende Efterskrift«, der, som sagt, stiller »Problemet«, var afleveret til Bogtrykkeren, for at Trykningen snarest muligt skulde paabegyndes, og Udgivelsen altsaa vel snart maatte paafølge: just i det Øieblik gjorde en Pseudonym, og paa det dertil hensigtsmæssigste Sted, i en Blad-Artikel, den størst mulige Anstrengelse for at støde fra i Retning af Publikum; og derpaa begyndte den afgjørende religieuse Produktivitet. Religieust holdt jeg anden Gang paa »hiin Enkelte«, til hvem saa det (efter Afsluttende Efterskrift) næste væsentlige Skrift »opbyggelige Taler i forskjellig Aand« eller sammes første Deel »Skriftetalen« var helliget. Da jeg første Gang anbragte Categorien »hiin Enkelte«, lagde man maaskee ikke Mærke til den; stort lagde man vel heller ikke Mærke til, at den stereotyp blev gjentaget i Forordet til hvert Hefte af opbyggelig Tale: da jeg anden Gang, eller i anden Potens, gjentog Budet og stod ved mit første Bud, var der gjort Alt, hvad jeg formaaede, for at Eftertrykkets hele Vægt maatte falde paa denne Categori. Bevægelsen er atter her: at komme til det Eenfoldige, Bevægelsen er: fra Publikum til »den Enkelte«. Religieust er der nemlig intet Publikum, men kun Enkelte; thi det Religieuse er Alvoren, og Alvoren er: den Enkelte, dog at ethvert Menneske, ubetinget ethvert Menneske, som han jo er det, kan være, ja skal være den Enkelte. Mig, den opbyggelige Forfatter, var og er det derfor en Glæde, at ogsaa just fra hiint Øieblik af De blive flere, som blive opmærksomme paa det om »den Enkelte«; det var og er mig en Glæde, thi vistnok har jeg Tro paa min Tankes Rigtighed trods hele Verden, men hvad jeg dog næst den sidst opgav, er min Tro paa de enkelte Mennesker. Og dette er min Tro, at saa meget Forvirret og Ondt og Afskyeligt der kan være i Menneskene, saasnart de blive det Ansvar- og Angerløse »Publikum«, »Mængde« og Deslige: lige saa meget Sandt og Godt og Elskværdigt er der i dem, naar man kan faae dem Enkelte. O, og i hvilken Grad skulde Menneskene ikke blive – Mennesker og Elskelige, naar de vilde blive Enkelte for Gud!

Saaledes forstaaer jeg det Hele nu; fra Begyndelsen har jeg ikke saaledes kunnet overskue, hvad der jo tillige har været min egen Udvikling. Til vidtløftig Fremstilling var da her mindst Stedet; her gjaldt det netop om ganske korteligen at kunne folde Det sammen i Eenfold, Det, som er udfoldet i de mange Bøger eller som udfoldet er de mange Bøger; og er denne korte Meddelelse nærmest foranlediget ved, at Forfatterskabets Første nu kommer anden Gang, det nye Oplag af »Enten – Eller«, hvilket jeg ikke tidligere har villet lade udgaae.

Personligt – saa vel naar jeg betænker mine indre Lidelser, som hvad jeg personligt kan have forskyldt – personligt beskjeftiger mig Eet ubetinget, er mig vigtigere end hele Forfatterskabet og ligger mig mere paa Hjerte, at udtrykke saa oprigtigt og saa stærkt som muligt, hvad jeg aldrig noksom kan takke for, og hvad jeg, naar jeg engang har glemt det hele Forfatterskab, uforandret evigt vil erindre: hvor uendelig meget Mere Styrelsen har gjort for mig, end jeg nogensinde havde ventet, kunde have ventet, eller turde have ventet.

»Uden Myndighed« at gjøre opmærksom paa det Religieuse, det Christelige, er Categorien for hele min Forfatter-Virksomhed totalt betragtet. At jeg var »uden Myndighed«, har jeg fra første Øieblik indskærpet og stereotyp gjentaget; jeg betragter mig selv helst som en Læser af Bøgerne, ikke som Forfatter.

»For Gud«, religieust, kalder jeg, naar jeg taler med mig selv, hele Forfatter-Virksomheden min egen Opdragelse og Udvikling, kun ikke i den Mening som var jeg da nu fuldkommen eller fuldkommen færdig i Henseende til at behøve Opdragelse og Udvikling.

\sectionheading{Synspunktet: Forklaringen og Slutord (uddrag)  ·  SV¹ XIII, 521–end}

§ 1. At »Christenheden« er et uhyre Sandsebedrag.

Enhver der med Alvor og derhos da med nogen Evne til at see, betragter hvad man saadan kalder Christenheden, eller Tilstanden i et saa kaldet christeligt Land, maa dog upaatvivleligt strax blive ganske betænkelig. Hvad vil det dog sige, at alle disse Tusinder og Tusinder udenvidere kalde sig Christne! Disse mange, mange Mennesker, hvis langt, langt overveiende Fleertal efter Alt hvad man kan skjønne have deres Liv i ganske andre Categorier, Noget man ved den simpleste Iagttagelse kan forvisse sig om! Mennesker, der maaskee aldrig engang gaae i Kirke, aldrig tænke paa Gud, aldrig nævne hans Navn uden naar de bande! Mennesker, for hvem det aldrig er gaaet op, at deres Liv skulde have nogen Forpligtelse mod Gud, Mennesker, der enten holde paa en vis borgerlig Ustraffelighed som Maximum eller endog ikke finde denne ganske fornøden! Dog alle disse Mennesker, endog de der paastaae, at der ingen Gud er til, de ere Alle Christne, kalde sig Christne, erkjendes som Christne af Staten, begraves som Christne af Kirken, dimitteres som Christne til Evigheden!

At her maa stikke en uhyre Confusion, et frygteligt Sandsebedrag, derom kan dog vist ingen Tvivl være. Men at røre ved det! Ja, jeg kjender godt Indvendingen. Thi der er nok baade Den og Den, der forstaaer hvad jeg mener, men som saa med en vis Godmodighed vilde klappe mig paa Skulderen og sige: »Kjære Ven, De er dog noget ung endnu; at ville begynde paa et saadant Foretagende, et Foretagende der, hvis det blot nogenlunde skulde lykkes, vilde kræve idetmindste en halv Snees vel disciplinerede Missionairer, et Foretagende, der egentligen løber ud paa hverken mere eller mindre end at ville indføre Christendommen igjen – i Christenheden. Nei, kjære Ven, lad os være Mennesker, et saadant Foretagende er baade over Deres og mine Kræfter. Dette Foretagende er lige saa afsindigt storartet som det at ville reformere »Mængden«, med hvilken ingen Forstandig indlader sig, men lader den gaae for hvad den er. At begynde paa Sligt er den visse Undergang«. Maaskee; men er eller var end Undergangen vis, vist er det ogsaa, at den Indvending har man ikke lært af Christendommen, thi da den kom ind i Verden, var det endnu mere afgjort den visse Undergang at begynde derpaa – dog begyndtes der; og vist er det ogsaa, at man heller ikke har lært den Indvending af Socrates, thi han indlod sig dog med »Mængden«, og vilde reformere den.

Saaledes omtrent forholder Sagen sig. Engang imellem gjør en Præst lidt Allarm paa Prædikestolen, at det ikke hænger rigtigt sammen med de mange Christne – men alle De, som høre ham og der ere tilstede, altsaa alle Dem, han taler til, de ere Christne; og dem, han taler om, dem taler han jo ikke til. Dette kaldes meest passende en fingeret Bevægelse. – Engang imellem træder der en opvakt Religieus frem; han stormer ind paa Christenheden, han støier og larmer, erklærer næsten Alle for ikke at være Christne – og han udretter Intet. Han tager sig ikke iagt for, at et Sandsebedrag ikke er saa let at hæve. Dersom det er saa, at det er en Indbildning, hvori de Fleste ere, naar de kalde sig Christne, hvad gjøre de saa mod en saadan Opvakt? Først og fremmest de bryde dem slet ikke om ham, de see ikke i hans Bog, men lægge den øieblikkeligt ad acta; eller virker han ved det levende Ord, saa gaae de om ad en anden Gade og høre ham slet ikke. Dernæst praktisere de ham ved Hjælp af en Begrebsbestemmelse udenfor, og indrette sig saa ganske trygt i Sandsebedraget. De gjøre ham til en Sværmer, hans Christendom til en Overdrivelse – am Ende bliver han den Eneste, eller een af de Faa, som ikke for Alvor (thi Overdrivelse er jo ogsaa Mangel paa Alvor) er Christen; de Andre ere alle alvorlige Christne.

Nei, et Sandsebedrag hæves aldrig direkte, og kun grundigt indirekte. Er det et Sandsebedrag med at Alle ere Christne – og skal der gjøres Noget, maa det gjøres indirekte, ikke af En, der høirøstet forkynder sig selv at være en overordentlig Christen, men af En, der bedre underrettet, erklærer sig selv end ikke at være en Christen. Det er, man maa komme bag paa Den, der er i Sandsebedraget. Istedetfor at ville have den Fordeel selv at være den sjeldne Christen, maa man lade den Hildede have Fordelen, at han er Christen, og selv have Resignation nok til at være Den, der er langt bag ved ham – ellers faaer man ham saa vist ikke ud af Sandsebedraget, det kan være vanskeligt nok endda.

Naar da nu, efter Antagelsen, de Fleste i Christenheden kun indbildt ere Christne, i hvilke Categorier leve de saa? De leve i æsthetiske, eller høist æsthetisk-ethiske.

Antaget saa, at en religieus Forfatter ret fra Grunden af er blevet opmærksom paa dette Sandsebedrag: Christenheden, og, saavidt hans Kræfter vel at mærke med Guds Bistand række, vil det tillivs: hvad har han saa at gjøre? Ja, først og fremmest ingen Utaalmodighed. Bliver han utaalmodig, saa stormer han ligefrem paa, og udretter – Intet. Ved direkte Angreb bestyrker man kun et Menneske i Sandsebedraget og forbittrer ham tillige. Der er overhovedet Intet, der kræver en saa lempelig Behandling som et Sandsebedrag, for at hæves. Foranlediger man paa nogen Maade den Hildede til at sætte Villie imod, saa er Alt tabt. Og det gjør man ved direkte Angreb, der desuden ogsaa indeholder det Anmassende, at fordre af et andet Menneske, at han skal tilstaae En eller ligeoverfor En gjøre den Tilstaaelse, som egentlig gavner mest, naar den Vedkommende gjør sig den selv i Stilhed. Dette naaes ved den indirekte Methode, der i Sandhedskjerlighedens Tjeneste lægger dialektisk Alt tilrette for den Hildede, og nu, som Kjerlighed altid er det, undseelig unddrager sig for at være Vidne til Tilstaaelsen, han saa ene for Gud gjør sig selv, at han dog har levet i en Indbildning.

Den religieuse Forfatter maa da altsaa først see at komme i Rapport til Menneskene. Det er, han maa begynde med æsthetisk Præstation. Det er Haandpengene. Jo brillantere Præstationen er, jo bedre for ham. Saa maa han dernæst være sikker paa sig selv, eller rettere (hvad der er det Sikkreste og det eneste Sikkre) i Frygt og Bæven forholde sig til Gud, at ikke det Omvendte skeer, saa han ikke bliver Den, der nu faaer Andre paa Gleed, men de Andre De, som faae Magt over ham, saa han ender med selv at blive stikkende i det Æsthetiske. Han maa altsaa have Alt i Beredskab, for, dog uden nogen Utaalmodighed, saa hurtigt som muligt, just naar han har faaet dem med sig, at bringe det Religieuse frem, saa de samme Mennesker med Hengivelsens Fart i det Æsthetiske, løbe lige Panden mod det Religieuse.

Det gjælder om, hverken for hurtigt eller for langsomt at anbringe det Religieuse. Gaaer der for lang Tid imellem, saa kommer strax Sandsebedraget op, at nu er den æsthetiske Forfatter blevet ældre og derfor religieus. Kommer det for hurtigt, bliver Virkningen ikke stærk nok.

[…]

[Slutning:]

»Dog fandt han ogsaa her i Verden hvad han søgte: var ingen Anden det, han var selv »hiin Enkelte«, og blev det mere og mere. Det var Christendommens Sag han tjente, hans Liv fra Barndommen af forunderligt dertil anlagt. Saa fuldkommede han Reflexionens Gjerning, at sætte Christendommen, det at blive Christen, heelt og holdent ind i Reflexion. Hans Hjertes Reenhed var: kun at ville Eet; hvad der i levende Live var de Samtidiges Anklage mod ham, at han ikke vilde slaae af, ikke give efter, netop det Samme er Eftertidens Lovtale over ham, at han ikke slog af, ikke gav efter. Men det storartede Foretagende bedaarede ham ikke; medens han dialektisk qua Forfatter holdt Overskuelsen over det Hele, forstod han christeligt, at det Hele for ham betydede, selv at opdrages i Christendom. Den dialektiske Bygning, han fuldførte, hvis enkelte Dele allerede ere Værker, kunde han ikke tilegne noget Menneske, endnu mindre vilde han tilegne sig den selv; skulde han have tilegnet Nogen den, var det blevet Styrelsen, hvem den dog alligevel Dag for Dag, Aar efter Aar var tilegnet af Forfatteren, der, historisk, døde af en dødelig Sygdom, men digterisk døde af Længsel efter Evigheden, for uafbrudt ikke at bestille Andet end takke Gud.«

\workheading{Guds Uforanderlighed. En Tale}{1855}

\worknote{Antologiens ''The Changelessness of God'' — det sidste selvstændige skrift Kierkegaard udgav. En tale over Jakobs Brev 1,17–21, tilegnet faderen Michael Pedersen Kierkegaard. Hele talen gengives: tilegnelsen, forordet, bønnen og betragtningen over Guds uforanderlighed som paa eengang forfærdende og beroligende.}

\sectionheading{Guds Uforanderlighed (Jakob 1,17–21)  ·  SV¹ XIV, 279–294}

Mindet om min afdøde Fader Michael Pedersen Kierkegaard, forhen Hosekræmmer her i Byen, helliget.

Forord.

Denne Tale er holdet i Citadellets Kirke den 18de Mai 1851. Texten er den første jeg har benyttet; senere er den oftere blevet fremdraget; nu vender jeg atter tilbage til den.

Bøn.

Du Uforanderlige, hvem Intet forandrer! Du i Kjerlighed Uforanderlige, som, just til vort Bedste, ikke lader Dig forandre: at vi ogsaa maatte ville vort eget Vel, ved Din Uforanderlighed lade os opdrage til, i ubetinget Lydighed, at finde Hvilen og at hvile i Din Uforanderlighed! Ikke er Du som et Menneske; skal han bevare blot nogen Uforanderlighed, maa han ikke have for Meget, der kan bevæge ham, og ikke lade sig for meget bevæge. Dig derimod bevæger, og i uendelig Kjerlighed, Alt; endog hvad vi Mennesker kalde en Ubetydelighed og gaae ubevægede forbi, Spurvens Trang, den bevæger Dig; hvad vi saa ofte neppe paaagte, et menneskeligt Suk, det bevæger Dig, uendelige Kjerlighed: men Intet forandrer Dig, Du Uforanderlige! O, Du, der i uendelig Kjerlighed lader Dig bevæge, Dig bevæge ogsaa denne vor Bøn, at Du velsigner den, saa Bønnen forandrer den Bedende i Overeensstemmelse med Din uforanderlige Villie, Du Uforanderlige!

Jacob 1, 17–21.

Al god Gave og al fuldkommen Gave er ovenfra, og kommer ned fra Lysenes Fader, hos hvilken er ikke Forandring eller Skygge af Omskiftelse. Efter sit Raad fødte han os formedelst Sandheds Ord, at vi skulde være en Førstegrøde af hans Skabninger. Derfor, mine elskelige Brødre, være hvert Menneske snar til at høre, langsom til at tale, langsom til Vrede; thi en Mands Vrede udretter ikke det, som er ret for Gud. Derfor aflægger al Skidenhed og al Levning af Ondskab, og med Sagtmodighed annammer Ordet, som er indplantet i Eder, og som er mægtigt til at gjøre Eders Sjele salige.

Min Tilhører, Du har hørt Texten forelæse. Hvor nær ligger det nu ikke at tænke paa Modsætningen: de timelige, de jordiske Tings Foranderlighed, og Menneskenes Foranderlighed! Nedtrykkende, Trættende, at Alt er Forkrænkelighed, at Menneskene ere Foranderlighed, Du, min Tilhører., og jeg! Sørgeligt, at Forandringen saa ofte er til det Værre! Fattige menneskelige Trøst, men dog en Trøst, at der da er endnu een Forandring ved det Foranderlige: at det faaer Ende!

Dog, hvis vi vilde tale saaledes, især i denne Forstemthedens Aand, altsaa ikke saaledes som naar der med Alvor tales om Forkrænkeligheden, om »den menneskelige Ustadighed«, da holdt vi os ikke blot ikke til Texten, nei, vi forlode den, ja, vi forandrede den. Thi i Texten tales der om det Modsatte, om Guds Uforanderlighed. Texten er idel Fryd og Glæde; som fra Bjergenes Tinde, hvor Taushed boer, saaledes er Apostelens Tale opløftet over al Jordlivets Foranderlighed; han taler om Guds Uforanderlighed, ikke om Andet. Om en »Lysenes Fader«, der boer histoppe, hvor ingen Omskiftelse naaer hen, end ikke Skyggen af den. Om »gode og fuldkomne Gaver«, der komme ned herovenfra, fra denne Fader, der, som »Lysenes« eller Lysets Fader, uendeligt veed at sikkre sig, at det sandeligen ogsaa er Godt og Fuldkomment, hvad der kommer fra ham, og som »Fader« Intet hellere vil, intet Andet tænker paa end uforandret at sende gode og fuldkomne Gaver. Og derfor, mine elskelige Brødre, være hvert Menneske »hurtig til at høre«, nemlig ikke efter Løst og Fast, men op efter, thi der oppe fra spørges bestandigt kun godt Nyt; »langsom til at tale«, thi den Snak vi Mennesker kan komme med især i Forhold hertil og især strax, den kan som oftest kun tjene til at gjøre de gode og fuldkomne Gaver mindre gode og fuldkomne; »langsom til Vrede«, at vi da ikke, naar Gaverne ikke synes os gode og fuldkomne, blive vrede, og bevirke, at det Gode og Fuldkomne, der var bestemt til vort Vel, ved vor egen Skyld bliver til Fordærvelse – det er Det, hvad en Mands Vrede kan udrette, og »en Mands Vrede udretter ikke det, som er ret for Gud«. »Derfor aflægger al Skidenhed og al Levning af Ondskab« – som man reengjør og smykker sit Huus og selv sidder pyntet, festligt ventende Besøget: at vi saaledes værdigen maatte modtage de gode og fuldkomne Gaver. »Og med Sagtmodighed annammer Ordet, som er indplantet i Eder, og som er mægtigt at gjøre Eders Sjele salige«. Med Sagtmodighed! I Sandhed, hvis det ikke var Apostelen, der talte, og hvis vi ikke strax efterkomme Befalingen »at være langsomme til at tale, langsomme til Vrede«, vi maatte vel sige: det var da en besynderlig Tale, ere vi da i den Grad Daarer, at vi behøve at formanes til Sagtmodighed i Forhold til Den, der kun vil vort Vel, det er jo som at spotte os, saaledes at anbringe det Ord »Sagtmodighed«. Thi see, dersom Een vilde slaae mig med Urette, og der da stod en Anden hos, der formanende sagde: find Dig med Sagtmodighed deri; dette er ligefrem Tale. Men tænk Dig det venligste Væsen, Kjerligheden selv; den har udsøgt en Gave, bestemt for mig, og Gaven er god og fuldkommen, ja som Kjerligheden selv; den kommer og vil skjenke mig denne Gave – da staaer der et andet Menneske hos, og han siger formanende til mig: lad mig nu see, at Du med Sagtmodighed finder Dig deri. Og dog saaledes er det med os Mennesker. En Hedning, og kun et Menneske, den eenfoldige Vise i Oldtiden, klager over, ofte at have erfaret, at naar han vilde fratage et Menneske een eller anden Taabelighed for at bibringe ham en bedre Viden, altsaa for at gjøre vel mod ham, at den Anden da kunde blive saa vred, at han, som den Eenfoldige spøgende siger i Alvor, gjerne vilde bide ham. Ak, og hvad har Gud ikke maattet erfare i de 6000 Aar, hvad erfarer han ikke hver Dag fra Morgen til Aften med hvert eneste af disse Millioner Mennesker; vi blive stundom meest vrede, naar han vil gjøre meest vel imod os. Ja, hvis vi Mennesker i Sandhed vidste vort eget Vel, og i dybeste Forstand i Sandhed vilde vort eget Vel, da behøvedes ingen Formaning til Sagtmodighed i dette Forhold. Men vi Mennesker (hvo har ikke i egen Erfaring sandet det!) vi ere dog i Forholdet til Gud som Børn. Og derfor behøves Formaningen om Sagtmodighed i Forhold til at modtage det Gode og det Fuldkomne – i den Grad er Apostelen forvisset om, at kun gode og fuldkomne Gaver komme ned fra ham, den evig Uforanderlige.

Forskjellige Synspunkter! Det blot Menneskelige (hvad jo Hedenskabet viser) taler mindre om Gud, og har overveiende Tilbøielighed til sørgmodigt kun at ville tale om de menneskelige Tings Foranderlighed; Apostelen vil ene og alene tale om Guds Uforanderlighed. Saaledes med Apostelen. For ham er Tanken om Guds Uforanderlighed ene og alene idel Trøst, Fred, Glæde, Salighed. Og dette er jo ogsaa evig sandt. Men lader os ikke glemme, at det, at det er saa for Apostelen, ligger i, at Apostelen er Apostelen, at han allerede længst havde i ubetinget Lydighed ganske givet sig hen i Guds Uforanderlighed, at han ikke stod ved Begyndelsen, men snarere ved Enden af Veien, den trange, men ogsaa den gode Vei, som han, forsagende Alt, havde valgt, uforandret fulgt uden at see sig tilbage, med stærkere og stærkere Skridt hastende mod Evigheden. Vi derimod som kun endnu ere Begyndere, under Opdragelsen, hos os maa Guds Uforanderlighed ogsaa have en anden Side at sees fra; og glemme vi dette, da løbe vi let Fare for at tage Apostelens Ophøiethed forfængeligt.

Saa ville vi da tale, om muligt baade i Forfærdelse og til Beroligelse om Dig, Du Uforanderlige, eller om Din Uforanderlighed.

Gud er uforanderlig. Almægtig skabte han denne synlige Verden – og gjorde sig usynlig; han iførte sig den synlige Verden som en Klædning, han forandrer den som man omskifter en Klædning – selv uforandret. Saaledes i den sandselige Verden. I Begivenhedernes Verden er han overalt tilstede i ethvert Øieblik; i en sandere Forstand end den aarvaagneste menneskelige Retfærdighed, der siges at være allesteds, er han, aldrig seet af nogen Dødelig, allestedsnærværende, overalt tilstede, ved det Mindste og ved det Største, ved hvad der kun uegentlig kan kaldes en Begivenhed og ved hvad der er den eneste Begivenhed, naar en Spurv døer, og naar Menneske-Slægtens Frelser fødes; han holder i ethvert Øieblik alt Virkeligt som Mulighed i sin almægtige Haand, har i ethvert Øieblik Alt i Beredskab, forandrer i et Nu Alt, Menneskenes Meninger, Domme, menneskelig Høihed og end Ringhed, han forandrer Alt – selv uforandret. Naar Alt er tilsyneladende Uforandrethed (thi det er kun tilsyneladende, at det Udvortes til en vis Tid er uforandret, det forandres altid), i Alts Omveltning forbliver han lige uforandret, ingen Omskiftelse berører ham, end ikke Skyggen af Omskiftelsen; i uforandret Klarhed er han, Lysenes Fader, evig uforandret. I uforandret Klarhed – ja just derfor er det han er uforandret, fordi han er idel Klarhed, en Klarhed som ingen Dunkelhed har i sig, og som ingen Dunkelhed kan komme nær. Med os Mennesker er det ikke saaledes, vi ere ikke saaledes Klarhed, og just derfor ere vi foranderlige: snart bliver Noget lysere i os, og snart fordunkles Noget, og vi forandres, nu vexler det udenom os, og Omskiftelsens Skygge glider forandrende over os, nu falder der fra Omverdenen en forandrende Belysning over os, medens vi selv under alt Dette igjen forandres i os selv. Men Gud er uforanderlig.

Denne Tanke er forfærdende, idel Frygt og Bæven. I Almindelighed fremhæves dette maaskee mindre; man klager over Menneskenes, over alt det Timeliges Foranderlighed, men Gud er uforanderlig, det er Trøsten, idel Trøst, siger endog Letsindigheden. Ja vist er Gud uforanderlig.

Men først og fremmest, er Du ogsaa i Forstaaelse med Gud, betænker Du ret alvorligt, stræber Du oprigtigt at forstaae – og dette er Guds evig uforanderlige Villie med Dig som med ethvert Menneske, at han skal stræbe derefter – stræber Du oprigtigt at forstaae, hvad der kan være Guds Villie med Dig? Eller lever Du saadan hen, og dette er ikke faldet Dig ind? Forfærdeligt, at saa han er den evig Uforanderlige! Thi med denne uforanderlige Villie maa Du dog engang, tidligt eller sildigt, støde sammen, denne uforanderlige Villie, der vilde, Du skulde betænke dette, fordi den vilde Dit Vel, denne uforanderlige Villie, som saa maa knuse Dig, hvis Du paa anden Maade støder sammen med den. Dernæst, Du som dog er i Forstaaelse med Gud, er Du ogsaa i god Forstaaelse med ham, er Din Villie, er den, og ubetinget, hans Villie, Dine Ønsker, hvert dit Ønske, hans Bud, Dine Tanker, den første og den sidste, hans Tanker: hvis ikke, forfærdeligt, at Gud er uforanderlig, evig, evig uforanderlig! Blot det at være ueens med et Menneske! Dog maaskee er Du den Stærkere og siger om den Anden: ja, ja han forandrer sig nok. Men dersom nu han er den Stærkere; dog maaskee mener Du, at kunne holde længere ud. Men dersom det nu er en heel Samtid; dog maaskee siger Du: 70 Aar er da ingen Evighed. Men den evig Uforanderlige – dersom Du var ueens med ham, det er jo en Evighed: forfærdeligt!

Tænk Dig en Vandringsmand; han er standset ved Foden af et uhyre, et uoverstigeligt Bjerg. Det er det han.... nei, han skal ikke, men det er det han vil over, thi hans Ønsker, hans Længsler, hans Begjeringer, hans Sjel – der har en lettere Art Befordring – er allerede ovre paa den anden Side, og det der mangler er blot at han følger efter. Tænk Dig han blev 70 Aar gammel; men Bjerget staaer uforandret, uoverstigeligt. Lad ham blive endnu en Gang 70 Aar; men Bjerget staaer uforandret ham i Veien, uforandret, uoverstigeligt. Saa forandres maaskee han under alt Dette, han døer fra sine Længsler, sine Ønsker, sine Begjeringer, han kjender neppe sig selv mere, saaledes træffer nu en senere Slægt ham siddende, forandret, ved Foden af Bjerget, der staaer uforandret, uoverstigeligt. Lad det være 1000 Aar siden; han den Forandrede, han er længst død, kun et Sagn fortælles om ham, det er det eneste, der er blevet tilbage – ja, og saa Bjerget, der staaer uforandret, uoverstigeligt. Og nu den evig Uforanderlige, for hvem 1000 Aar ere som en Dag, ak og selv dette er for meget sagt, de ere for ham som et Nu, ja egentligen ere de for ham som vare de ikke for ham – hvis Du end i fjerneste Maade vil ad en anden Vei end han vil Du skal: frygteligt!

Sandt nok, om Din, om min, om disse Tusinders og Tusinders Villie just ikke er saa ganske i Overeensstemmelse med Guds: det gaaer jo som det bedst kan derude i den saakaldte virkelige Verdens Travlhed, Gud lader sig jo egentligen ikke mærke med Noget; snarere er det vel saa, at hvis der var en Retfærdig – hvis der var en Saadan! – der betragtede denne Verden, en Verden, der, som Skriften siger, ligger i det Onde, han maatte blive mismodig over, at Gud ikke lader sig mærke med Noget. Men troer Du derfor, Gud har forandret sig, eller er dette, at han ikke lader sig mærke med Noget, det mindre Forfærdelige, naar det dog er vist, at han er evig uforanderlig? Mig synes det ikke saa. Betænk det dog, og siig da, hvilket er det Forfærdeligste, enten dette: den uendelig Stærkere, der, træt af at lade sig spotte, reiser sig i sin Magt og knuser de Gjenstridige – det er forfærdeligt, og saaledes fremstilles det ogsaa, naar der tales om, at Gud ikke lader sig spotte, og der da vises hen paa Tider, hvor hans Straffedomme tilintetgjørende gik over Slægten; men er det dog egentligen det Forfærdeligste? Er dette ikke det endnu Forfærdeligere: den uendelig Stærke, der – evig uforanderlig! – sidder ganske stille og seer, uden en Mines Forandring, næsten som var han ikke til, medens dog, saa maatte den Retfærdige vel klage, Usandhed har Fremgang, har Magt, Vold og Uret seirer og endog saaledes, at selv en Bedre kan fristes til at mene at maatte benytte Lidt af de samme Midler, hvis der skal være Haab om at udrette Noget for det Gode, saa det er som spottedes han, han den uendelig Stærke, den evig Uforanderlige, som dog hverken lader sig spotte eller forandre – er dette ikke det Forfærdeligste? Thi hvorfor troer Du vel, han er saa stille? Fordi han veed med sig selv, at han er evig uforanderlig. En, der ikke var evig sikker paa sig selv, at han er den Uforanderlige, han kunde ikke holde sig saaledes stille, han reiste sig i sin Magt; men kun den evig Uforanderlige kan sidde saaledes stille. Han giver Tid, det kan han ogsaa, han har Evigheden, og evigt er han uforanderlig; han giver Tid, det gjør han med velberaad Hu, saa kommer der et Evighedens Regnskab, hvor Intet er glemt, ikke eet eneste af de utilbørlige Ord, der bleve talte, og evigt er han uforanderlig. Dog det kan ogsaa være Barmhjertighed, at han saaledes giver Tid, Tid til Omvendelse og Forbedring; men frygteligt, hvis denne Tid ikke benyttes saaledes, thi da maatte Daarligheden og Letsindigheden i os hellere ønske, at han strax var ved Haanden med Straffen, end at han saaledes giver Tid, lader som Ingenting og dog er evig uforanderlig. Spørg en Opdrager – og vi ere dog Alle i Forholdet til Gud mere eller mindre Børn! – spørg Den, der har havt med vildfarende Mennesker at gjøre – og enhver af os har dog idetmindste een Gang været vildfarende, er det i længere eller kortere Tid, med større eller mindre Mellemrum – og Du skal høre, han vil sande, at det er en stor Hjælp for Letsindigheden eller rettere for at forhindre Letsindighed – og hvo turde sige sig ganske fri for Letsindighed! – at Straffens Lidelse om muligt følger i samme Nu som Overtrædelsen, saa den Letsindiges Hukommelse vænnes til at huske Straffen samtidigt med Skylden. Ja var det saa, at Forvildelse og Straf stode saaledes i Forhold til hinanden, at man, ligesom paa et dobbeltløbet Skydevaaben, trykkede paa een Fjeder, og i samme Nu man greb det forbudne Lystelige eller forskyldte Forvildelsen, i samme Nu fulgte Straffen: jeg troer Letsindigheden vilde vogte sig. Men jo længere Tid mellem Skylden og Straffen (hvilket i Sandhed forstaaet udtrykker Maalestokken for Sagens Alvor), jo mere fristende for Letsindighed, som kunde maaskee det Hele blive glemt, eller Retfærdigheden selv maaskee forandre sig og til den Tid faae ganske andre Begreber, eller som vilde det dog idetmindste være saa længe siden det skete, at det vil være umuligt at fremstille Sagen uforandret. Saa forandrer Letsindigheden sig, og ikke til det Bedre; saa bliver Letsindigheden tryg; og naar saa Letsindigheden er bleven tryg, saa drister den mere, og saa gaaer Aar efter Aar – Straffen udebliver, og Glemsel indtræder, og atter udebliver Straffen, men ny Forvildelse udebliver ikke, og den gamle Forvildelse er nu blevet mere ondartet; og saa er det forbi; saa slutter Døden af – og til alt Dette (det var kun Letsindighed!) var der en evig Uforanderlig Vidne: var det saa ogsaa Letsindighed? en evig Uforanderlig, og det er med ham Du vil have at gjøre Regnskabet op. I det Øieblik, da Timelighedens Viser, Minut-Viseren, viste 70 Aar, og det Menneske døde, i den Tid havde Evighedens Viser neppe flyttet sig en Kjende: i den Grad er Alt nærværende for Evigheden og for ham den Uforanderlige!

Og derfor hvo Du end er, husk paa, hvad jeg siger til mig selv, at for Gud er intet Betydeligt og intet Ubetydeligt, at i een Forstand er det Betydelige for ham ubetydeligt, i en anden Forstand selv den mindste Ubetydelighed noget uendelig Betydeligt. Er da Din Villie ikke i Overeensstemmelse med hans, betænk det, Du slipper aldrig fra ham, tak ham, om han ved Mildhed eller ved Strenghed lærer Dig at bringe Din Villie i Overeensstemmelse med hans – frygteligt, hvis han ikke lader sig mærke med Noget, frygteligt, hvis det kunde komme saa vidt med et Menneske, at han næsten trodser paa, at Gud enten ikke er til, eller at han har forandret sig, eller endog blot at han er for stor til at lægge Mærke til hvad vi kalde Ubetydeligheder; thi baade er han til, og han er evig uforanderlig, og hans uendelige Storhed er just den, at han seer end det Mindste, husker end det Mindste, ja, og hvis Du ikke vil som han, husker det uforandret en Evighed!

Der er altsaa for os letsindige og ustadige Mennesker idel Frygt og Bæven i denne Tanke om Guds Uforanderlighed. O, betænk det vel! hvad enten han lader sig mærke med Noget eller ikke, han er evig uforanderlig. Han er evig uforanderlig, betænk det vel, hvis Du har, som det siges, noget Udstaaende med ham, han er uforanderlig. Maaskee har Du lovet ham Noget, ved helligt Løfte forpligtet Dig..... men i Tidens Løb har Du forandret Dig, tænker nu sjeldnere paa Gud (har Du maaskee som Ældre faaet vigtigere Ting at tænke paa?) eller Du tænker maaskee anderledes om Gud, at han ikke bryder sig om Dit Livs Ubetydeligheder, at slig Tro er Barnagtighed: i ethvert Tilfælde Du har saadan glemt, hvad Du lovede ham, saa derpaa glemt, at Du lovede ham det, og saa tilsidst glemt, glemt – ja glemt, at han Intet glemmer, han den evig Uforanderlige, at det just er en Aarenes bagvendte Barnagtighed at mene, at Noget er ubetydeligt for Gud, og at Gud glemmer Noget, han den evig Uforanderlige!

I Forholdet mellem Menneske og Menneske klages der saa ofte over Foranderlighed, den Ene klager over den Anden, at han har forandret sig; men selv i Forholdet mellem Menneske og Menneske kan stundom den Enes Uforandrethed være som til Plage. Maaskee har En talt til et andet Menneske om sig selv. Maaskee var det lidt barnagtig, tilgivelig Tale, han førte. Men maaskee var Sagen ogsaa alvorligere: det fristede det daarlige, forfængelige Hjerte at tale i høie Toner om sin Begeistring, sin Følelses Varighed, sin Villen i denne Verden. Den Anden hørte roligt derpaa, han ikke engang smilede, eller forhindrede ham i at tale; han lod ham tale, han hørte, han taug, kun lovede han, hvad der blev forlangt, ikke at glemme det Sagte. Saa gik der Tid hen; og den Første havde længst glemt alt Dette; den Anden derimod havde ikke glemt det, ja, lad os tænke det endnu Besynderligere, han havde ladet sig bevæge af de Tanker, som den Første i et Stemningens Øieblik havde udtalt, ak, og ligesom givet fra sig, havde i redelig Stræben dannet sit Liv i Forhold dertil: hvilken Plage med hans Hukommelses Uforandrethed, han som kun altfor tydeligt viste, at han indtil det Mindste huskede, hvad der i hiint Øieblik blev sagt!

Og nu den evig Uforanderlige – og dette menneskelige Hjerte! O, dette menneskelige Hjerte, hvad gjemmer Du dog ikke i Dit hemmelighedsfulde Indelukke, ubekjendt for Andre – det var ikke det Værste – men stundom næsten ubekjendt for den Vedkommende selv! Fast er det jo, saasnart et Menneske blot er blevet noget tilaars, fast er det jo som et Gravsted, dette menneskelige Hjerte! Der ligge de begravne, begravne i Glemsel, Løfter, Forsætter, Beslutninger, hele Planer og Brudstykker af Planer og Gud veed hvad – ja saaledes tale vi Mennesker, thi vi Mennesker tænke sjeldent ved hvad vi sige; vi sige: der ligger Gud veed hvad. Og det sige vi saadan halvt letsindigt, halvt træt af Livet – og saa er det saa frygteligt sandt, at Gud veed hvad, han veed det, indtil det Mindste, det Du har glemt, han veed hvad der for Din Hukommelse har forandret sig, det veed han uforandret, han ikke engang erindrer det, som var det dog noget siden, nei han veed det, som var det idag, og han veed, om der i Forhold til noget af disse Ønsker, Forsætter, Beslutninger var saa at sige talt til ham derom – og han er evig uforandret og evig uforanderlig. O, kan et andet Menneskes Hukommelse falde et Menneske til Byrde – nu, den er dog vel aldrig saa ganske tilforladelig, og i ethvert Tilfælde det kan da ikke vare en Evighed, jeg bliver dog fri for dette andet Menneske og hans Hukommelse; men en Alvidende, og en evig uforanderlig Hukommelse, fra hvilken Du saa ikke slipper, mindst i Evigheden: frygteligt! Nei, evig uforanderligt er for ham Alt evigt nærværende, evigt lige nærværende, ingen hverken Morgenens eller Aftenens, hverken Ungdommens eller Alderdommens, hverken Glemsomhedens eller Undskyldningens omskiftende Skygge omskifter ham, nei for ham er ingen Skygge; ere vi, som det hedder, Skygger, han er evig Klarhed i sin evige Uforanderlighed; ere vi Skygger, som ile hen – min Sjel, see Dig dog vel for, thi enten Du vil eller ikke, iler Du hen til Evigheden, til ham, og han er evig Klarhed! Derfor ikke blot holder han Regnskab, men han er Regnskabet; det hedder, vi Mennesker skulle gjøre Regnskab, som var der maaskee en lang Tid dertil, og saa maaskee en uoverkommelig Masse af Vidtløftigheder, for at faae Regnskabet bragt istand: o, min Sjel, det er i ethvert Øieblik gjort; thi hans uforanderlige Klarhed er Regnskabet, indtil det Mindste fuldstændigt færdigt og bevaret af ham den evig Uforanderlige, der Intet har glemt af hvad jeg har glemt, ei heller, som jeg, husker Noget anderledes end det dog virkelig var.

Saaledes er der idel Frygt og Bæven i denne Tanke om Guds Uforanderlighed, næsten er det som var det langt, langt over et Menneskes Kræfter at have med en saadan Uforanderlighed at gjøre, ja som maatte denne Tanke styrte et Menneske i Angest og Uro indtil at fortvivle.

Men saa er det ogsaa saaledes, at der er Beroligelse og Salighed i denne Tanke; det er virkelig saa, at naar Du, trættet af al denne menneskelige, al denne timelige og jordiske Foranderlighed og Omvexling, trættet af Din egen Ustadighed, kunde ønske et Sted, hvor Du kunde hvile Dit trætte Hoved, Din trætte Tanke, Dit trætte Sind, for at hvile og hvile ud: o, i Guds Uforanderlighed der er Hvile! Naar Du derfor lader denne hans Uforanderlighed tjene Dig, som han vil, til Dit Bedste, Dit evige Bedste, naar Du lader Dig opdrage, saa Din Egenvillie (og det er fra den Foranderligheden egenligen kommer, endnu mere end udenfra) uddøer, jo før jo hellere, det hjælper Dig jo ikke, Du skal enten med det Gode eller med det Onde, tænk Dig det Forgjeves i at ville være ueens med en evig Uforanderlighed, vær som Barnet, naar det ret dybt fornemmer, at det lige over for sig har en Villie, hvor der kun hjælper Eet, at lystre – naar Du lader Dig ved hans Uforanderlighed opdrage, saa Du forsager Ustadighed og Foranderlighed og Lune og Egenraadighed: da hviler Du stedse tryggere og saligere og saligere i denne Guds Uforanderlighed. Thi at Tanken om Guds Uforanderlighed er salig, ja, hvo tvivler; men vogt blot paa, at Du bliver saaledes, at Du saligt kan hvile i denne Uforanderlighed! O, som Den, der har et lykkeligt Hjem, saaledes siger en Saadan: mit Hjem er evig betrygget, jeg hviler i Guds Uforanderlighed. Den Hvile kan Ingen forstyrre Dig uden Du selv; kunde Du blive ganske lydig i uforandret Lydighed, da skulde Du i ethvert Øieblik med samme Nødvendighed som et tungt Legeme synker til Jorden, eller med samme Nødvendighed som hvad der er let hæver sig mod Himlen, frit hvile i Gud.

Og lad det saa kun Alt forresten vexle, som det gjør. Skal Du finde Din Virksomhed paa en større Skueplads, vil Du efter en større Maalestok erfare Alts Foranderlighed; men paa en mindre Skueplads, paa den mindste, vil Du dog erfare det Samme, maaskee lige saa smerteligt. Du vil erfare, hvorledes Menneskene forandre sig, hvorledes Du selv forandrer Dig; stundom vil det ogsaa være, som om Gud forandrede sig, hvilket hører med til Opdragelsen. Herom, om Alts Foranderlighed vil en Ældre bedre kunne tale end jeg, medens maaskee hvad jeg kunde sige, kunde synes den ganske Unge noget Nyt. Dog dette ville vi ikke videre udføre, men overlade Livets Mangfoldighed at udfolde for Enhver, som det er ham bestemt, at han kan komme til at erfare, hvad før ham alle Andre have erfaret. Stundom vil Forandringen være saaledes, at Du mindes Ordet: Forandring behager – ja, usigeligt! Der vil ogsaa komme Tider, da Du selv opfinder et Ord, som Sproget har fortiet, og Du siger: Forandring behager ikke – hvor kunde jeg dog sige: Forandring behager. Naar saa er, vil Du særligen være foranlediget til (hvad Du dog vel heller ikke i det første Tilfælde vil glemme) at søge hen til ham, den Uforanderlige.

Min Tilhører! denne Time er nu snart forbi, og Talen. Hvis Du ikke selv vil det anderledes, snart vil ogsaa denne Time være glemt, og Talen. Og hvis Du ikke selv vil det anderledes, snart vil ogsaa Tanken om Guds Uforanderlighed være glemt i Foranderligheden. Dog deri er han dog vel ikke Skyld, han den Uforanderlige! Men forskylder Du ikke selv at glemme den, da vil Du i denne Tanke være forsørget for Dit Liv, for en Evighed.

Tænk Dig i Ørkenen en Eensom; forbrændt næsten af Solens Hede, forsmægtende finder han en Kilde. O, liflige Kølighed! Nu er jeg, Gud være lovet, siger han – og han fandt dog kun en Kilde, hvorledes maatte ikke Den tale, der fandt Gud! og dog maatte ogsaa han sige »Gud være lovet«, jeg fandt Gud! – nu er jeg, Gud være lovet, forsørget. Thi Din trofaste Kølighed, o elskede Kilde, er ikke underlagt Forandring. I Vinterens Kulde, hvis den naaede her hen, Du bliver ikke koldere, men bevarer nøiagtigt den samme Kølighed, Kildens Vand fryser ikke! I Sommersolens Middags-Brand, Du bevarer nøiagtigt Din uforandrede Kølighed, Kildens Vand lunknes ikke! Og der er intet Usandt i hvad han siger (han, der i mine Tanker ikke valgte nogen utaknemlig Gjenstand for en Lovtale, en Kilde, hvad Enhver bedre forstaaer, jo bedre han veed, hvad det vil sige: Ørken og Eensomhed), der er ingen usand Overdrivelse i hans Lovtale. Imidlertid, hans Liv tog en ganske anden Vending, end han havde tænkt. Han forvildede sig engang bort, blev saa revet ud i den vide Verden. Mange Aar efter vendte han tilbage. Hans første Tanke var Kilden – den var der ikke, den var udtørret. Et Øieblik stod han taus i Sorg; da fattede han sig og sagde: nei, jeg tilbagekalder dog ikke eet Ord af hvad jeg sagde til Din Lov, det var Sandhed Alt. Og prisede jeg Din liflige Kølighed, medens Du var, o elskede Kilde, saa lad mig ogsaa prise den, nu Du er forsvunden, at det maa være sandt, at der er Uforandrethed i et Menneskes Bryst. Ei heller kan jeg sige, at Du bedrog mig; nei, havde jeg fundet Dig, jeg er forvisset, Din Kølighed vilde være uforandret – og Mere havde Du ikke lovet.

Men Du, o Gud, Du Uforanderlige, Du er uforandret altid at finde, og lader Dig uforandret altid finde, Ingen reiser, hverken i Liv eller Død, saa langt bort, at Du ikke er at finde, at Du ikke er der, Du er jo overalt, – saaledes er der ikke Kilder paa Jorden, Kilderne ere kun paa enkelte Steder. Og desuden – overvældende Sikkerhed! – Du bliver jo ikke som Kilden paa Stedet, Du reiser med; ak, og Ingen forvilder sig saa langt bort, at han ikke kan finde tilbage til Dig, Du der ikke blot er som en Kilde, der lader sig finde, – fattige Beskrivelse af Dit Væsen! – Du, der er som en Kilde, der selv søger den Tørstende, den Forvildede, hvad man aldrig har hørt om nogen Kilde. Saaledes er Du uforandret altid og overalt at finde. O, og naarsomhelst et Menneske kommer til Dig, i hvilken Alder, til hvilken Tid paa Dagen, i hvilken Tilstand: dersom han kommer oprigtigt, han finder altid (som Kildens uforandrede Kølighed) Din Kjerlighed lige varm, Du Uforanderlige! Amen.
\end{document}
